\documentclass[10pt]{article}
\usepackage[verbose, a4paper, hmargin=2.5cm, vmargin=2.5cm]{geometry}

\usepackage{fontspec}
\usepackage{ctex}
\usepackage{paratype}

\usepackage{subfigure}

\usepackage{amsmath}
\usepackage{amsfonts}
\usepackage{amssymb}
\usepackage{esint}

\usepackage{graphicx}
\usepackage[export]{adjustbox}
\usepackage{mdframed}
\usepackage{booktabs,array,multirow}
\usepackage{adjustbox}
\usepackage{tabularx}
\usepackage{hyperref}
\hypersetup{colorlinks=true, linkcolor=blue, filecolor=magenta, urlcolor=cyan,}
\urlstyle{same}
\usepackage[most]{tcolorbox}
\definecolor{mygray}{RGB}{240,240,240}
\tcbset{
  colback=mygray,
  boxrule=0pt,
}
\graphicspath{ {./images/} }
\newcommand{\HRule}{\begin{center}\rule{0.9\linewidth}{0.2mm}\end{center}}
\newcommand{\customfootnote}[1]{
  \let\thefootnote\relax\footnotetext{#1}
}

\begin{document}
\section*{Game Theory Basics}

\begin{tcolorbox}
\section*{博弈论基础}
\end{tcolorbox}

Game theory is the science of interaction. This textbook, derived from courses taught by the author and developed over several years, is a comprehensive, straightforward introduction to the mathematics of non-cooperative games. It teaches what every game theorist should know: the important ideas and results on strategies, game trees, utility theory, imperfect information, and Nash equilibrium. The proofs of such results - in particular, existence of an equilibrium via fixed points and an elegant direct proof of the minimax theorem for zero-sum games - are presented in a self-contained, accessible way. Complementary to that are chapters on combinatorial games such as Go, and - as introductions to algorithmic game theory - traffic games and the geometry of two-player games. This detailed and lively text requires minimal mathematical background and includes many examples, exercises, and pictures. It is suitable for self-study or introductory courses in mathematics, computer science, or economics departments.

\begin{tcolorbox}
博弈论是研究交互行为的科学。这本教材源自作者多年授课的课程内容，是对非合作博弈数学理论的全面且直接的介绍。它涵盖了每个博弈论学者都应了解的内容:关于策略、博弈树(game tree)、效用理论、不完全信息和纳什均衡(Nash equilibrium)的重要概念与成果。这些成果的证明——特别是通过不动点证明均衡(equilibrium)的存在性，以及对零和博弈(zero-sum game)极小极大定理(minimax theorem)的简洁直接证明——都以自成体系且易于理解的方式呈现。与之互补(complementary)的是关于组合博弈(如围棋)的章节，以及作为算法博弈论入门的交通博弈和双人博弈几何的内容。这本详细且生动的教材所需的数学背景极少，包含大量示例、练习和插图。它适合自学，也适用于数学、计算机科学或经济学专业的入门课程。
\end{tcolorbox}

Bernhard von Stengel, educated in Germany and the US, is a mathematical game theorist at the London School of Economics and Political Science, and an authority on computational and geometric methods for solving games. He chaired the 2016 World Congress of the Game Theory Society, and is Co-Editor of the International Journal of Game Theory.

\begin{tcolorbox}
伯恩哈德·冯·施滕格尔(Bernhard von Stengel)在德国和美国接受教育，是伦敦政治经济学院的数学博弈论学者，也是解决博弈问题的计算和几何方法方面的权威。他担任了2016年博弈论学会世界大会的主席，还是《国际博弈论杂志》的联合主编。
\end{tcolorbox}

"This looks like a fine introduction to game theory, inter alia emphasizing methods for computing equilibria, and mathematical aspects in general. Especially worthy of note is the chapter devoted to correlated equilibria, a topic of central importance not normally covered in introductory texts."

\begin{tcolorbox}
“这似乎是一本优秀的博弈论入门教材，尤其强调了计算均衡的方法以及一般的数学方面。特别值得注意的是关于相关均衡的章节，这一重要主题通常不在入门教材中涉及。”
\end{tcolorbox}

Robert Aumann, The Hebrew University of Jerusalem,

\begin{tcolorbox}
罗伯特·奥曼(Robert Aumann)，耶路撒冷希伯来大学
\end{tcolorbox}

Nobel Memorial Prize in Economic Sciences 2005

\begin{tcolorbox}
2005年诺贝尔经济学奖
\end{tcolorbox}

"This book is a delightful adventure into the mathematics of game theory. Without any heavy apparatus, it lets us into the secrets of a whole range of exciting results that are usually thought too advanced for the common herd. It is not only undergraduate students who will benefit from reading this book. Professional game theorists will find it very useful too."

\begin{tcolorbox}
“这本书是一场对博弈论数学的精彩探索。无需复杂的工具，它就让我们了解到一系列令人兴奋的成果背后的奥秘，这些成果通常被认为对普通读者来说过于高深。不仅本科生会从阅读本书中受益，专业的博弈论学者也会发现它非常有用。”
\end{tcolorbox}

Ken Binmore, University College London

\begin{tcolorbox}
肯·宾默尔(Ken Binmore)，伦敦大学学院
\end{tcolorbox}

"This is a rather reader-friendly, engaging, and polished superior creation. It illustrates, explains, motivates every definition, theorem, proof. Interesting and unique choice of topics, such as a delightful introductory chapter on combinatorial games. Highly recommended."

\begin{tcolorbox}
“这是一部非常便于阅读、引人入胜且精心打磨的佳作。它对每个定义、定理和证明都进行了说明、解释和动机阐述。主题选择有趣且独特，比如关于组合博弈的精彩入门章节。强烈推荐。”
\end{tcolorbox}

Aviezri Fraenkel, Weizmann Institute of Science, Israel

\begin{tcolorbox}
阿维兹里·弗兰克尔(Aviezri Fraenkel)，以色列魏茨曼科学研究所
\end{tcolorbox}

"A masterful presentation of mathematical game theory in all its beauty and elegance, from basic notions to advanced techniques. It fills the gaps left by the many textbooks that cover concepts and applications, but devote only the bare minimum to the mathematical tools and insights, without which game theory would not have become the success it is today."

\begin{tcolorbox}
“对数学博弈论进行了精彩而优雅的呈现，从基本概念到高级技术一应俱全。它填补了许多只涵盖概念和应用，却对数学工具和见解仅作最低限度介绍的教材的空白，没有这些数学工具和见解，博弈论就不会取得如今的成功。”
\end{tcolorbox}

Sergiu Hart, The Hebrew University of Jerusalem

\begin{tcolorbox}
塞尔久·哈特(Sergiu Hart)，耶路撒冷希伯来大学
\end{tcolorbox}

"Game theory is the child of mathematicians, as this textbook demonstrates through self-contained, elegant proofs of all seminal theorems. The lively and rigorous exposition of carefully selected models, such as bargaining, combinatorial, and congestion games (the latter two rarely the stuff of textbooks), explains its success far beyond mathematics. To reach deep results on both sides of the theory, Bernhard von Sten-gel's marvellous learning tool uses uncompromising, yet accessible mathematics, and chooses examples to maximal effect."

\begin{tcolorbox}
“博弈论是数学家的结晶，正如这本教材通过对所有开创性定理自成体系且优雅的证明所展示的那样。对精心挑选的模型(如讨价还价博弈、组合博弈和拥塞博弈(congestion game)，后两者很少出现在教材中)进行生动而严谨的阐述，解释了博弈论在数学之外的广泛成功。为了深入探究该理论的各个方面，伯恩哈德·冯·施滕格尔的这部出色的学习工具采用了严谨但易于理解的数学方法，并精心选择示例以达到最佳效果。”
\end{tcolorbox}

Hervé Moulin, University of Glasgow

\begin{tcolorbox}
埃尔韦·莫林(Hervé Moulin)，格拉斯哥大学
\end{tcolorbox}

"This book is a gem. The presentation is clear and well structured, often with nice geometric illustrations. It moves step by step from basics to powerful concepts, methods, and results. It is ideal for students of mathematics, computer science, and economics who are curious about what game theory is and how it can be used." Jörgen Weibull, Stockholm School of Economics

\begin{tcolorbox}
“这本书是一颗宝石。内容呈现清晰、结构合理，常有精美的几何插图。它循序渐进地从基础内容过渡到强大的概念、方法和成果。对于对博弈论是什么以及如何应用感兴趣的数学、计算机科学和经济学专业的学生来说，它是理想之选。” 约根·魏布尔(Jörgen Weibull)，斯德哥尔摩经济学院
\end{tcolorbox}

"An exceptionally lucid introduction to the fundamentals of game theory, enlivened by examples that are sure to captivate students."

\begin{tcolorbox}
“对博弈论基础进行了异常清晰的介绍，大量示例定会吸引学生的兴趣。”
\end{tcolorbox}

Peyton Young, University of Oxford

\begin{tcolorbox}
佩顿·扬，牛津大学
\end{tcolorbox}

"This is a rigorous, yet accessible introduction to mathematical non-cooperative game theory. In addition to the coverage of the basic concepts and results, it includes special and advanced topics and applications usually not contained in game theory textbooks, such as combinatorial games, congestion games, and inspection games. The special emphasis on algorithmic and computational techniques makes this textbook, just like its author, a valuable bridge between game theory and computer sciences."

\begin{tcolorbox}
"这是一本严谨且易懂的数学非合作博弈论入门书籍。除了涵盖基本概念和结果外，它还包含了博弈论教材中通常没有的特殊和高级主题及应用，如组合博弈、拥塞博弈(congestion game)和监管博弈(Inspection game)。对算法和计算技术的特别强调，使这本教材就像其作者一样，成为博弈论和计算机科学之间的宝贵桥梁。"
\end{tcolorbox}

Shmuel Zamir, The Hebrew University of Jerusalem

\begin{tcolorbox}
什穆埃尔·扎米尔，耶路撒冷希伯来大学
\end{tcolorbox}

\section*{Game Theory Basics}

\begin{tcolorbox}
\section*{博弈论基础}
\end{tcolorbox}

Bernhard von Stengel

\begin{tcolorbox}
伯恩哈德·冯·施滕格尔
\end{tcolorbox}

London School of Economics and Political Science

\begin{tcolorbox}
伦敦政治经济学院
\end{tcolorbox}

\section*{CAMBRIDGE UNIVERSITY PRESS}

\begin{tcolorbox}
\section*{剑桥大学出版社}
\end{tcolorbox}

University Printing House, Cambridge CB2 8BS, United Kingdom One Liberty Plaza, 20th Floor, New York, NY 10006, USA 477 Williamstown Road, Port Melbourne, VIC 3207, Australia 314-321, 3rd Floor, Plot 3, Splendor Forum, Jasola District Centre, New Delhi - 110025, India 103 Penang Road, \#05-06/07, Visioncrest Commercial, Singapore 238467 Cambridge University Press is part of the University of Cambridge. It furthers the University's mission by disseminating knowledge in the pursuit of education, learning, and research at the highest international levels of excellence. www.cambridge.org Information on this title: www.cambridge.org/9781108843300 DOI: 10.1017/9781108910118 © Bernhard von Stengel 2022 This publication is in copyright. Subject to statutory exception and to the provisions of relevant collective licensing agreements, no reproduction of any part may take place without the written permission of Cambridge University Press. First published 2022 Printed in the United Kingdom by TJ Books Limited, Padstow, Cornwall A catalogue record for this publication is available from the British Library. Library of Congress Cataloging-in-Publication Data Names: Stengel, Bernhard von, author. Title: Game theory basics / Bernhard von Stengel, London School of Economics and Political Science, Department of Mathematics. Description: New York : Cambridge University Press, 2021. Includes bibliographical references and index.

\begin{tcolorbox}
英国剑桥市大学印刷所，邮编:CB2 8BS 美国纽约市自由广场一号20层，邮编:10006 澳大利亚墨尔本港威廉斯敦路477号，邮编:VIC 3207 印度新德里贾索拉区中心辉煌论坛3号楼3层314 - 321号，邮编:110025 新加坡槟榔路103号维景商业中心05 - 06/07室，邮编:238467 剑桥大学出版社是剑桥大学的一部分。它通过在追求教育、学习和研究的过程中传播知识，以达到国际最高卓越水平，从而推进大学的使命。网址:www.cambridge.org 本书相关信息:www.cambridge.org/9781108843300 数字对象标识符:10.1017/9781108910118 © 伯恩哈德·冯·施滕格尔 2022 本出版物受版权保护。除非有法定例外情况和相关集体许可协议的规定，否则未经剑桥大学出版社书面许可，不得复制任何部分。2022年首次出版 由英国康沃尔郡帕德斯托的TJ图书有限公司印刷 英国国家图书馆可提供本出版物的目录记录。美国国会图书馆编目数据 作者:施滕格尔，伯恩哈德·冯 书名:博弈论基础 / 伯恩哈德·冯·施滕格尔，伦敦政治经济学院数学系 描述:纽约:剑桥大学出版社，2021年。包括参考文献和索引。
\end{tcolorbox}

Identifiers: LCCN 2021024598 (print) | LCCN 2021024599 (ebook) | ISBN 9781108843300 (hardback) | ISBN 9781108843300 (ebook)

\begin{tcolorbox}
标识符:LCCN 2021024598(印刷版) | LCCN 2021024599(电子版) | ISBN 9781108843300(精装版) | ISBN 9781108843300(电子版)
\end{tcolorbox}

Subjects: LCSH: Game theory. | BISAC: MATHEMATICS / Discrete Mathematics Classification: LCC QA269 .S699 2021 (print) | LCC QA269 (ebook) | DDC 519.3-dc23

\begin{tcolorbox}
主题:LCSH:博弈论。 | BISAC:数学 / 离散数学 分类:LCC QA269 .S699 2021(印刷版) | LCC QA269(电子版) | DDC 519.3 - dc23
\end{tcolorbox}

LC record available at https://lccn.loc.gov/2021024598

\begin{tcolorbox}
美国国会图书馆记录可在https://lccn.loc.gov/2021024598查询
\end{tcolorbox}

LC ebook record available at https://lccn.loc.gov/2021024599

\begin{tcolorbox}
美国国会图书馆电子版记录可在https://lccn.loc.gov/2021024599查询
\end{tcolorbox}

ISBN 978-1-108-84330-0 Hardback

\begin{tcolorbox}
ISBN 978 - 1 - 108 - 84330 - 0 精装版
\end{tcolorbox}

ISBN 978-1-108-82423-1 Paperback

\begin{tcolorbox}
ISBN 978 - 1 - 108 - 82423 - 1 平装版
\end{tcolorbox}

Cambridge University Press has no responsibility for the persistence or accuracy of URLs for external or third-party internet websites referred to in this publication and does not guarantee that any content on such websites is, or will remain, accurate or appropriate.

\begin{tcolorbox}
剑桥大学出版社对本出版物中提及的外部或第三方互联网网站的网址的持久性或准确性不承担责任，也不保证此类网站上的任何内容现在或将来都是准确或合适的。
\end{tcolorbox}

\section*{Contents}

\begin{tcolorbox}
\section*{目录}
\end{tcolorbox}

Preface ix

\begin{tcolorbox}
前言  ix
\end{tcolorbox}

Aims and Contents ix

\begin{tcolorbox}
目标与内容  ix
\end{tcolorbox}

Mathematical and Scholarly Level xi

\begin{tcolorbox}
数学与学术水平  xi
\end{tcolorbox}

Acknowledgments xii

\begin{tcolorbox}
致谢  xii
\end{tcolorbox}

1 Nim and Combinatorial Games 1

\begin{tcolorbox}
1 尼姆游戏(Nim)与组合博弈  1
\end{tcolorbox}

1.1 Prerequisites and Learning Outcomes 1

\begin{tcolorbox}
1.1 先决条件与学习(learning)成果  1
\end{tcolorbox}

1.2 Nim 2

\begin{tcolorbox}
1.2 尼姆游戏(Nim)  2
\end{tcolorbox}

1.3 Top-down Induction 7

\begin{tcolorbox}
1.3 自上而下归纳(top - down induction)  7
\end{tcolorbox}

1.4 Game Sums and Equivalence of Games 10

\begin{tcolorbox}
1.4 博弈之和(game sum)与博弈等价性  10
\end{tcolorbox}

1.5 Nim, Poker Nim, and the Mex Rule 15

\begin{tcolorbox}
1.5 尼姆游戏(Nim)、扑克尼姆游戏(Poker Nim)与mex规则(mex rule)  15
\end{tcolorbox}

1.6 Sums of Nim Heaps 19

\begin{tcolorbox}
1.6 尼姆堆(Nim heap)之和  19
\end{tcolorbox}

1.7 Finding Nim Values 22

\begin{tcolorbox}
1.7 寻找尼姆值(Nim value)  22
\end{tcolorbox}

1.8 A Glimpse of Partizan Games 24

\begin{tcolorbox}
1.8 有偏博弈(partizan game)一瞥  24
\end{tcolorbox}

1.9 Further Reading 30

\begin{tcolorbox}
1.9 延伸阅读  30
\end{tcolorbox}

1.10 Exercises for Chapter 1 32

\begin{tcolorbox}
1.10 第1章练习  32
\end{tcolorbox}

2 Congestion Games 39

\begin{tcolorbox}
2 拥塞博弈(congestion game)  39
\end{tcolorbox}

2.1 Prerequisites and Learning Outcomes 40

\begin{tcolorbox}
2.1 先决条件与学习成果 40
\end{tcolorbox}

2.2 Introduction: The Pigou Network 40

\begin{tcolorbox}
2.2 引言:庇古网络(Pigou network) 40
\end{tcolorbox}

2.3 The Braess Paradox 43

\begin{tcolorbox}
2.3 布雷斯悖论(Braess paradox) 43
\end{tcolorbox}

2.4 Definition of Congestion Games 44

\begin{tcolorbox}
2.4 拥塞博弈(congestion game)的定义 44
\end{tcolorbox}

2.5 Existence of Equilibrium in a Congestion Game 46

\begin{tcolorbox}
2.5 拥塞博弈中的均衡(equilibrium in a congestion game)的存在性 46
\end{tcolorbox}

2.6 Atomic and Splittable Flow, Price of Anarchy 48

\begin{tcolorbox}
2.6 原子流与可拆分流(splittable flow)、无政府代价(Price of Anarchy) 48
\end{tcolorbox}

2.7 Further Reading 49

\begin{tcolorbox}
2.7 延伸阅读 49
\end{tcolorbox}

2.8 Exercises for Chapter 2 50

\begin{tcolorbox}
2.8 第2章习题 50
\end{tcolorbox}

3 Games in Strategic Form 53

\begin{tcolorbox}
3 策略形式(strategic form)的博弈 53
\end{tcolorbox}

3.1 Prerequisites and Learning Outcomes 54

\begin{tcolorbox}
3.1 先决条件与学习成果 54
\end{tcolorbox}

3.2 Games in Strategic Form 54

\begin{tcolorbox}
3.2 策略形式的博弈 54
\end{tcolorbox}

3.3 Best Responses and Equilibrium 57

\begin{tcolorbox}
3.3 最优反应(best response)与均衡 57
\end{tcolorbox}

3.4 Games with Multiple Equilibria 60

\begin{tcolorbox}
3.4 具有多重均衡的博弈 60
\end{tcolorbox}

3.5 Dominated Strategies 62

\begin{tcolorbox}
3.5 (被占优) 策略(dominated strategy) 62
\end{tcolorbox}

3.6 The Cournot Duopoly of Quantity Competition 66

\begin{tcolorbox}
3.6 古诺双寡头产量竞争博弈(Cournot duopoly of quantity competition) 66
\end{tcolorbox}

3.7 Games without a Pure-Strategy Equilibrium 72

\begin{tcolorbox}
3.7 不存在纯策略均衡的博弈 72
\end{tcolorbox}

3.8 Symmetric Games with Two Strategies per Player 72

\begin{tcolorbox}
3.8 每个参与者有两种策略的对称博弈(symmetric game) 72
\end{tcolorbox}

3.9 Further Reading 74

\begin{tcolorbox}
3.9 延伸阅读 74
\end{tcolorbox}

3.10 Exercises for Chapter 3 75

\begin{tcolorbox}
3.10 第三章练习 75
\end{tcolorbox}

4 Game Trees with Perfect Information 77

\begin{tcolorbox}
4 完美信息(perfect information)博弈树(game tree) 77
\end{tcolorbox}

4.1 Prerequisites and Learning Outcomes 78

\begin{tcolorbox}
4.1 先决条件与学习(learning)成果 78
\end{tcolorbox}

4.2 Definition of Game Trees 79

\begin{tcolorbox}
4.2 博弈树的定义 79
\end{tcolorbox}

4.3 Backward Induction 82

\begin{tcolorbox}
4.3 逆向归纳法(backward induction) 82
\end{tcolorbox}

4.4 Strategies in Game Trees 84

\begin{tcolorbox}
4.4 博弈树中的策略 84
\end{tcolorbox}

4.5 Reduced Strategies 86

\begin{tcolorbox}
4.5 精简策略(reduced strategy) 86
\end{tcolorbox}

4.6 Subgame-Perfect Equilibrium (SPE) 91

\begin{tcolorbox}
4.6 子博弈完美均衡(subgame-perfect equilibrium，SPE) 91
\end{tcolorbox}

4.7 Commitment Games 93

\begin{tcolorbox}
4.7 承诺博弈(commitment game) 93
\end{tcolorbox}

4.8 Further Reading 98

\begin{tcolorbox}
4.8 延伸阅读 98
\end{tcolorbox}

4.9 Exercises for Chapter 4 99

\begin{tcolorbox}
4.9 第四章练习 99
\end{tcolorbox}

5 Expected Utility 103

\begin{tcolorbox}
5 期望效用 103
\end{tcolorbox}

5.1 Prerequisites and Learning Outcomes 105

\begin{tcolorbox}
5.1 先决条件与学习成果 105
\end{tcolorbox}

5.2 Summary 106

\begin{tcolorbox}
5.2 总结 106
\end{tcolorbox}

5.3 Decisions Under Risk 109

\begin{tcolorbox}
5.3 风险下的决策 109
\end{tcolorbox}

5.4 Preferences for Lotteries 111

\begin{tcolorbox}
5.4 对彩票的偏好 111
\end{tcolorbox}

5.5 Ordinal Preferences for Decisions Under Certainty 113

\begin{tcolorbox}
5.5 确定性下决策的序数偏好 113
\end{tcolorbox}

5.6 Cardinal Utility Functions and Simple Lotteries 115

\begin{tcolorbox}
5.6 基数效用函数与简单彩票 115
\end{tcolorbox}

5.7 Consistency Axioms 120

\begin{tcolorbox}
5.7 一致性公理 120
\end{tcolorbox}

5.8 Existence of an Expected-Utility Function 124

\begin{tcolorbox}
5.8 期望效用函数的存在性 124
\end{tcolorbox}

5.9 Risk Aversion 129

\begin{tcolorbox}
5.9 风险厌恶 129
\end{tcolorbox}

5.10 Discussion and Further Reading 132

\begin{tcolorbox}
5.10 讨论与进一步阅读 132
\end{tcolorbox}

5.11 Exercises for Chapter 5 134

\begin{tcolorbox}
5.11 第5章练习 134
\end{tcolorbox}

6 Mixed Equilibrium 135

\begin{tcolorbox}
6 混合均衡 135
\end{tcolorbox}

6.1 Prerequisites and Learning Objectives 136

\begin{tcolorbox}
6.1 先决条件与学习目标 136
\end{tcolorbox}

6.2 Compliance Inspections 136

\begin{tcolorbox}
6.2 合规检查 136
\end{tcolorbox}

6.3 Bimatrix Games 141

\begin{tcolorbox}
6.3 双矩阵博弈 141
\end{tcolorbox}

6.4 The Best-Response Condition 145

\begin{tcolorbox}
6.4 最优反应条件 145
\end{tcolorbox}

6.5 Existence of Mixed Equilibria 149

\begin{tcolorbox}
6.5 混合均衡的存在性 149
\end{tcolorbox}

6.6 Finding Mixed Equilibria in Small Games 153

\begin{tcolorbox}
6.6 小博弈中混合均衡的求解 153
\end{tcolorbox}

6.7 The Upper-Envelope Method 158

\begin{tcolorbox}
6.7 上包络法 158
\end{tcolorbox}

6.8 Degenerate Games 161

\begin{tcolorbox}
6.8 退化博弈 161
\end{tcolorbox}

6.9 Further Reading 166

\begin{tcolorbox}
6.9 延伸阅读 166
\end{tcolorbox}

6.10 Exercises for Chapter 6 167

\begin{tcolorbox}
6.10 第6章习题 167
\end{tcolorbox}

7 Brouwer's Fixed-Point Theorem 169 7.1 Prerequisites and Learning Outcomes 171

\begin{tcolorbox}
7 布劳威尔不动点定理 169 7.1 先决条件与学习成果 171
\end{tcolorbox}

7.2 Labels 171

\begin{tcolorbox}
7.2 标签 171
\end{tcolorbox}

7.3 Simplices and Triangulations 173

\begin{tcolorbox}
7.3 单纯形与三角剖分 173
\end{tcolorbox}

7.4 Sperner's Lemma 180

\begin{tcolorbox}
7.4 斯佩纳引理 180
\end{tcolorbox}

7.5 The Knaster-Kuratowski-Mazurkiewicz Lemma 187

\begin{tcolorbox}
7.5 克纳斯特 - 库拉托夫斯基 - 马祖尔克维奇引理 187
\end{tcolorbox}

7.6 Brouwer's Fixed-Point Theorem on a General Compact Convex Set 190

\begin{tcolorbox}
7.6 一般紧致凸集(compact convex set)上的布劳威尔不动点定理 190
\end{tcolorbox}

7.7 The Freudenthal Triangulation 196

\begin{tcolorbox}
7.7 弗赖登塔尔三角剖分 196
\end{tcolorbox}

7.8 Further Reading 203

\begin{tcolorbox}
7.8 延伸阅读 203
\end{tcolorbox}

7.9 Exercises for Chapter 7 204

\begin{tcolorbox}
7.9 第7章习题 204
\end{tcolorbox}

8 Zero-Sum Games 205

\begin{tcolorbox}
8 零和博弈 205
\end{tcolorbox}

8.1 Prerequisites and Learning Outcomes 207

\begin{tcolorbox}
8.1 先决条件与学习成果 207
\end{tcolorbox}

8.2 Example: Soccer Penalty 207

\begin{tcolorbox}
8.2 示例:足球点球 207
\end{tcolorbox}

8.3 Max-Min and Min-Max Strategies 211

\begin{tcolorbox}
8.3 极大极小与极小极大策略 211
\end{tcolorbox}

8.4 A Short Proof of the Minimax Theorem 216

\begin{tcolorbox}
8.4 极小极大定理的简短证明 216
\end{tcolorbox}

8.5 Further Notes on Zero-Sum Games 218

\begin{tcolorbox}
8.5 关于零和博弈(zero - sum game)的进一步说明 218
\end{tcolorbox}

8.6 Further Reading 220

\begin{tcolorbox}
8.6 延伸阅读 220
\end{tcolorbox}

8.7 Exercises for Chapter 8 221

\begin{tcolorbox}
8.7 第8章习题 221
\end{tcolorbox}

9 Geometry of Equilibria in Bimatrix Games 223

\begin{tcolorbox}
9 双矩阵博弈(bimatrix game)中均衡的几何性质 223
\end{tcolorbox}

9.1 Prerequisites and Learning Outcomes 224

\begin{tcolorbox}
9.1 先决条件与学习(learning)目标 224
\end{tcolorbox}

9.2 Labeled Best-Response Regions 224

\begin{tcolorbox}
9.2 带标签(label)的最优反应区域 224
\end{tcolorbox}

9.3 The Lemke-Howson Algorithm 230

\begin{tcolorbox}
9.3 莱姆基 - 豪森算法(Lemke - Howson algorithm) 230
\end{tcolorbox}

9.4 Using Best-Response Diagrams 235

\begin{tcolorbox}
9.4 使用最优反应示意图(best - response diagram) 235
\end{tcolorbox}

9.5 Strategic Equivalence 239

\begin{tcolorbox}
9.5 策略等价(strategic equivalence) 239
\end{tcolorbox}

9.6 Best-Response Polyhedra and Polytopes 245

\begin{tcolorbox}
9.6 最优反应多面体和有界多面体(polytope) 245
\end{tcolorbox}

9.7 Complementary Pivoting 250

\begin{tcolorbox}
9.7 互补换基法(complementary pivoting) 250
\end{tcolorbox}

9.8 Degeneracy Resolution 256

\begin{tcolorbox}
9.8 退化问题的解决 256
\end{tcolorbox}

9.9 Further Reading 259

\begin{tcolorbox}
9.9 延伸阅读 259
\end{tcolorbox}

9.10 Exercises for Chapter 9 261

\begin{tcolorbox}
9.10 第9章习题 261
\end{tcolorbox}

10 Game Trees with Imperfect Information 263

\begin{tcolorbox}
10 不完美信息的博弈树(Game Trees with Imperfect Information) 263
\end{tcolorbox}

10.1 Prerequisites and Learning Outcomes 264

\begin{tcolorbox}
10.1 先决条件和学习(Learning)成果 264
\end{tcolorbox}

10.2 Information Sets 264

\begin{tcolorbox}
10.2 信息集(Information Sets) 264
\end{tcolorbox}

10.3 Extensive Games 270

\begin{tcolorbox}
10.3 扩展式博弈 270
\end{tcolorbox}

10.4 Strategies for Extensive Games and the Strategic Form 272

\begin{tcolorbox}
10.4 扩展式博弈的策略与策略形式(Strategic Form) 272
\end{tcolorbox}

10.5 Reduced Strategies 274

\begin{tcolorbox}
10.5 精简策略(Reduced Strategies) 274
\end{tcolorbox}

10.6 Perfect Recall 275

\begin{tcolorbox}
10.6 完美记忆(Perfect Recall) 275
\end{tcolorbox}

10.7 Behavior Strategies 279

\begin{tcolorbox}
10.7 行为策略(Behavior Strategies) 279
\end{tcolorbox}

10.8 Kuhn's Theorem: Behavior Strategies Suffice 283

\begin{tcolorbox}
10.8 库恩定理:行为策略(Behavior Strategies)足够 283
\end{tcolorbox}

10.9 Behavior Strategies in the Monty Hall Problem 287

\begin{tcolorbox}
10.9 蒙提霍尔问题中的行为策略(Behavior Strategies) 287
\end{tcolorbox}

10.10 Subgames and Subgame-Perfect Equilibria 291

\begin{tcolorbox}
10.10 子博弈(Subgames)与子博弈完美均衡 291
\end{tcolorbox}

10.11 Further Reading 295

\begin{tcolorbox}
10.11 延伸阅读 295
\end{tcolorbox}

10.12 Exercises for Chapter 10 296

\begin{tcolorbox}
10.12 第10章练习 296
\end{tcolorbox}

11 Bargaining 299

\begin{tcolorbox}
11 讨价还价 299
\end{tcolorbox}

11.1 Prerequisites and Learning Outcomes 300

\begin{tcolorbox}
11.1 先决条件和学习(Learning)成果 300
\end{tcolorbox}

11.2 Bargaining Sets 300

\begin{tcolorbox}
11.2 讨价还价集 300
\end{tcolorbox}

11.3 Bargaining Axioms 302

\begin{tcolorbox}
11.3 讨价还价公理 302
\end{tcolorbox}

11.4 The Nash Bargaining Solution 305

\begin{tcolorbox}
11.4 纳什讨价还价解 305
\end{tcolorbox}

11.5 Geometry of the Bargaining Solution 308

\begin{tcolorbox}
11.5 讨价还价解的几何表示 308
\end{tcolorbox}

11.6 Splitting a Unit Pie 311

\begin{tcolorbox}
11.6 分割单位蛋糕 311
\end{tcolorbox}

11.7 The Ultimatum Game 313

\begin{tcolorbox}
11.7 最后通牒博弈 313
\end{tcolorbox}

11.8 Alternating Offers Over Two Rounds 317

\begin{tcolorbox}
11.8 两轮交替出价 317
\end{tcolorbox}

11.9 Alternating Offers Over Several Rounds 320

\begin{tcolorbox}
11.9 多轮交替出价 320
\end{tcolorbox}

11.10 Stationary Strategies 324

\begin{tcolorbox}
11.10 平稳策略 324
\end{tcolorbox}

11.11 The Nash Bargaining Solution Via Alternating Offers 328

\begin{tcolorbox}
11.11 通过交替出价得到的纳什讨价还价解 328
\end{tcolorbox}

11.12 Further Reading 330

\begin{tcolorbox}
11.12 延伸阅读 330
\end{tcolorbox}

11.13 Exercises for Chapter 11 330

\begin{tcolorbox}
11.13 第11章练习 330
\end{tcolorbox}

12 Correlated Equilibrium 333

\begin{tcolorbox}
12 相关均衡 333
\end{tcolorbox}

12.1 Prerequisites and Learning Outcomes 333

\begin{tcolorbox}
12.1 先决条件和学习成果 333
\end{tcolorbox}

12.2 Examples of Correlated Equilibria 334

\begin{tcolorbox}
12.2 相关均衡的例子 334
\end{tcolorbox}

12.3 Incentive Constraints 337

\begin{tcolorbox}
12.3 激励约束 337
\end{tcolorbox}

12.4 Coarse Correlated Equilibrium 341

\begin{tcolorbox}
12.4 粗相关均衡 341
\end{tcolorbox}

12.5 Existence of a Correlated Equilibrium 344

\begin{tcolorbox}
12.5 相关均衡(correlated equilibrium)的存在性 344
\end{tcolorbox}

12.6 Further Reading 348

\begin{tcolorbox}
12.6 延伸阅读 348
\end{tcolorbox}

12.7 Exercises for Chapter 12 349

\begin{tcolorbox}
12.7 第12章习题 349
\end{tcolorbox}

References 351

\begin{tcolorbox}
参考文献 351
\end{tcolorbox}

Index 358

\begin{tcolorbox}
索引 358
\end{tcolorbox}

\section*{Preface}

\begin{tcolorbox}
\section*{前言}
\end{tcolorbox}

This book is an introduction to the mathematics of non-cooperative game theory. Each concept is explained in detail, starting from a main example, with a slow-paced proof of each theorem. The book has been designed and tested for self-study, and as an undergraduate course text with core chapters and optional chapters for different audiences. It has been developed over 15 years for a one-semester course on game theory at the London School of Economics and the distance learning program of the University of London, attended each year by about 200 third-year students in mathematics, economics, management, and other degrees. After studying this book, a student who started from first-year mathematics (the basics of linear algebra, analysis, and probability) will have a solid understanding of the most important concepts and theorems of non-cooperative game theory.

\begin{tcolorbox}
本书是对非合作博弈论数学知识的入门介绍。每个概念都从一个主要示例出发进行详细解释，并对每个定理进行循序渐进的证明。本书专为自学而设计并经过测试，也可作为本科课程教材，设有核心章节和供不同受众选择的可选章节。本书是作者在15年时间里，为伦敦政治经济学院的一学期博弈论课程以及伦敦大学远程学习项目编写的，每年约有200名数学、经济学、管理学及其他专业的三年级学生参加学习。学完本书后，具备一年级数学基础(线性代数、分析和概率论基础)的学生将扎实掌握非合作博弈论中最重要的概念和定理。
\end{tcolorbox}

The intended audience are primarily students in mathematics, computer science, and mathematical economics. For mathematicians, we provide complete self-contained proofs (in an economics course, these may be used as reliable background and reference material). For computer scientists, we introduce important ideas in algorithmic game theory, such as traffic equilibria, the "parity argument" for the existence and computation of Nash equilibria, and correlated equilibria. For economists, we use some important economic models as examples, such as Cournot's quantity competition, which is the oldest formal definition of a game. Commitment games are applied to Stackelberg leadership, and in Chapter 11 to the iterated-offers bargaining model. However, given the many good introductions to game theory for economists (for example, Gibbons, 1992, or Osborne, 2004), economic applications are not central to this book.

\begin{tcolorbox}
本书主要面向数学、计算机科学和数理经济学专业的学生。对于数学家，我们提供完整且自成体系的证明(在经济学课程中，这些证明可作为可靠的背景资料和参考材料)。对于计算机科学家，我们介绍算法博弈论中的重要思想，如交通均衡、用于证明纳什均衡(Nash equilibrium)存在性和计算的“奇偶性论证(parity argument)”以及相关均衡。对于经济学家，我们以一些重要的经济模型为例，如古诺产量竞争(Cournot's quantity competition)，这是最早对博弈的正式定义。承诺博弈(commitment game)被应用于斯塔克尔伯格领导模型(Stackelberg leadership)，在第11章中还应用于迭代报价议价模型。不过，鉴于已有许多面向经济学家的优秀博弈论入门书籍(如Gibbons，1992年或Osborne，2004年的著作)，经济应用并非本书的核心内容。
\end{tcolorbox}

\section*{Aims and Contents}

\begin{tcolorbox}
\section*{目标与内容}
\end{tcolorbox}

The first aim of this book is to let the student become fluent in game theory. The student will learn the modeling tools (such as game trees and the strategic form) and methods for analyzing games (finding their equilibria). This is provided in the core chapters on non-cooperative games: Chapter 3 on games in strategic form introduces important games such as the Prisoner's Dilemma or Matching Pennies, and the concept of Nash equilibrium (in the book nearly always just called "equilibrium" for brevity). Chapter 4 treats game trees with perfect information. Chapter 6 explains mixed strategies and mixed equilibrium. Game trees with imperfect information are covered in Chapter 10.

\begin{tcolorbox}
本书的首要目标是让学生熟练掌握博弈论。学生将学习建模工具(如博弈树(game tree)和策略形式(strategic form))以及分析博弈的方法(寻找博弈的均衡(equilibrium))。这些内容在关于非合作博弈的核心章节中有所呈现:第3章介绍策略形式的博弈，引入了诸如囚徒困境(Prisoner's Dilemma)或硬币匹配(Matching Pennies)等重要博弈，以及纳什均衡(Nash equilibrium)的概念(为简洁起见，本书中几乎总是简称为“均衡”)。第4章讨论具有完美信息(perfect information)的博弈树。第6章解释混合策略(mixed strategy)和混合均衡(mixed equilibrium)。第10章涵盖具有不完美信息的博弈树。
\end{tcolorbox}

A second aim is to provide the conceptual and mathematical foundations of the theory. Chapter 5 explains how the concept of expected utility represents a consistent preference for risky outcomes (and not "risk neutrality" for monetary payoffs, a common confusion). Chapter 7 proves Brouwer's fixed-point theorem, used by Nash (1951) to show the existence of an equilibrium point. Chapter 8 gives a self-contained two-page proof of von Neumann's minimax theorem for zero-sum games. These chapters form independent "modules" that can be used for reference and that can be omitted in a taught course that has less emphasis on mathematical proofs. If only one proof from these chapters is presented, it should be the short proof of the minimax theorem.

\begin{tcolorbox}
第二个目标是为该理论提供概念和数学基础。第5章解释了期望效用概念如何代表对风险结果的一致偏好(而非对货币收益的“风险中性”，这是常见的误解)。第7章证明了布劳威尔不动点定理，纳什(1951年)用该定理证明了均衡点的存在性。第8章用两页篇幅自成体系地证明了冯·诺伊曼针对零和博弈(zero-sum game)提出的极小极大定理(minimax theorem)。这些章节构成独立的“模块”，可用于参考，在不太注重数学证明的授课课程中可以省略。如果只展示这些章节中的一个证明，那应该是极小极大定理的简短证明。
\end{tcolorbox}

A third aim is to introduce ideas that every game theorist should know but which are not normally taught, and which make this book special. They are mathematical highlights, typically less known to economists:

\begin{tcolorbox}
第三个目标是介绍每个博弈论学者都应了解但通常不教授的思想，这也是本书的独特之处。这些都是数学亮点，经济学家通常不太了解:
\end{tcolorbox}

\begin{itemize}
\item An accessible introduction to combinatorial games in Chapter 1, pioneered by Berlekamp, Conway, and Guy (2001-2004), with a focus on impartial games and the central game of Nim.
\end{itemize}

\begin{tcolorbox}
\begin{itemize}
\item 第1章对组合博弈进行了通俗易懂的介绍，这一领域由Berlekamp、Conway和Guy(2001 - 2004年)开创，重点介绍无偏博弈(impartial game)和核心的尼姆游戏(Nim)。
\end{itemize}
\end{tcolorbox}

\begin{itemize}
\item Congestion games, with the famous Braess paradox where increasing network capacity can worsen congestion, and where an equilibrium is found with the help of a potential function (Chapter 2).
\end{itemize}

\begin{tcolorbox}
\begin{itemize}
\item 拥塞博弈(congestion game)，其中有著名的布雷斯悖论(Braess paradox)，即增加网络容量可能会加剧拥堵，并且可以借助势函数(potential function)找到均衡(第2章)。
\end{itemize}
\end{tcolorbox}

\begin{itemize}
\item An elegant constructive proof of Sperner's lemma, used to prove Brouwer's fixed-point theorem, and the Freudenthal simplicial subdivision that works naturally in any dimension (Sections 7.4 and 7.7 in Chapter 7).
\end{itemize}

\begin{tcolorbox}
\begin{itemize}
\item 对斯佩纳引理(Sperner's lemma)进行了优雅的构造性证明，该引理用于证明布劳威尔不动点定理，还介绍了在任何维度都能自然适用的弗罗伊登塔尔单纯剖分(Freudenthal simplicial subdivision)(第7章第7.4节和第7.7节)。
\end{itemize}
\end{tcolorbox}

\begin{itemize}
\item The geometry of two-player games (which is my own research specialty) and the algorithm by Lemke and Howson (1964) for finding a Nash equilibrium, which implies that generic games have an odd number of equilibria (Chapter 9).
\end{itemize}

\begin{tcolorbox}
\begin{itemize}
\item 两人博弈的几何性质(这是我自己的研究专长)以及莱姆克 - 豪森算法(Lemke and Howson，1964年)，用于寻找纳什均衡，该算法表明通有博弈(generic game)有奇数个均衡(第9章)。
\end{itemize}
\end{tcolorbox}

\begin{itemize}
\item An introduction to correlated equilibria in Chapter 12, with an elegant existence proof due to Hart and Schmeidler (1989) that does not require a fixed-point theorem but only von Neumann's minimax theorem.
\end{itemize}

\begin{tcolorbox}
\begin{itemize}
\item 第12章介绍了相关均衡(correlated equilibria)，并给出了一个优雅的存在性证明，该证明由哈特(Hart)和施迈德勒(Schmeidler)于1989年提出，不依赖不动点定理，仅使用了冯·诺伊曼的极小极大定理(minimax theorem)。
\end{itemize}
\end{tcolorbox}

These chapters are optional. I teach the course in variations but always include Chapter 1 on combinatorial games; they are much enjoyed by students, who would not encounter them in an introductory economics course on game theory. These chapters give also a short introduction to algorithmic game theory for computer science students.

\begin{tcolorbox}
这些章节为选学内容。我会以不同方式讲授这门课程，但始终会涵盖第1章关于组合博弈的内容；学生们很喜欢这些内容，因为在经济学专业的博弈论入门课程中他们不会接触到这些。这些章节还为计算机科学专业的学生简要介绍了算法博弈论。
\end{tcolorbox}

The general emphasis of the book is to teach methods, not philosophy. Game theorists tend to question and to justify the approaches they take, for example the concept of Nash equilibrium, and the assumed common knowledge of all players about the rules of the game. These questions are of course very important. In fact, these are probably the very issues that require a careful validation in a practical game-theoretic analysis. However, this problem is not remedied by a lengthy discussion of why one should play Nash equilibrium.

\begin{tcolorbox}
本书的重点是传授方法，而非探讨哲学。博弈论学者往往会质疑并论证他们所采用的方法，例如纳什均衡(Nash equilibrium)的概念，以及假设所有参与者对博弈规则具有共同知识(common knowledge)。这些问题当然非常重要。事实上，在实际的博弈论分析中，这些可能正是需要仔细验证的问题。然而，冗长地讨论为何要采用纳什均衡并不能解决这个问题。
\end{tcolorbox}

I think that a student of game theory should first learn the central models of the strategic and extensive form, and be fluent in analyzing them. That toolbox will then be useful when comparing different game-theoretic models and the solutions that they imply.

\begin{tcolorbox}
我认为，博弈论的学生首先应该学习战略式和扩展式的核心模型，并熟练掌握对它们的分析方法。当比较不同的博弈论模型及其所蕴含的解时，这些工具将非常有用。
\end{tcolorbox}

\section*{Mathematical and Scholarly Level}

\begin{tcolorbox}
\section*{数学与学术水平}
\end{tcolorbox}

General mathematical prerequisites are the basic notions of linear algebra (vectors and matrices), probability theory (independence, conditional probability), and analysis (continuity, closed sets). Rather than putting them in a rarely read appendix, important concepts are recalled where needed in text boxes labeled as Background material.

\begin{tcolorbox}
一般的数学先决条件包括线性代数(向量和矩阵)、概率论(独立性、条件概率(conditional probability))和分析学(连续性、闭集)的基本概念。重要的概念没有放在很少有人阅读的附录中，而是在需要时以“背景材料(background material)”的文本框形式回顾。
\end{tcolorbox}

For each chapter, the necessary mathematical background and the required previous chapters are listed in a first section on prerequisites and learning objectives.

\begin{tcolorbox}
每章的先决条件和学习目标部分会列出所需的数学背景和之前需要学习的章节。
\end{tcolorbox}

The mathematics of game theory is not technically difficult, but very conceptual, and requires therefore a certain mathematical maturity. For example, combinatorial games have a recursive structure, for which a generalization of the mathematical induction known for natural numbers is appropriate, called "top-down induction" and explained in Section 1.3.

\begin{tcolorbox}
博弈论的数学知识在技术上并不难，但概念性很强，因此需要一定的数学成熟度。例如，组合博弈具有递归结构(recursive structure)，对于这种结构，自然数的数学归纳法的一种推广形式——“自上而下归纳(top - down induction)”是适用的，这在第1.3节中会进行解释。
\end{tcolorbox}

Game-theoretic concepts have precise and often unfamiliar mathematical definitions. In this book, each main concept is typically explained by means of a detailed introductory example. We use the format of definitions, theorems and proofs in order to be precise and to keep some technical details hidden in proofs. The proofs are detailed and complete and can be studied line by line. The main idea of each proof is conveyed by introductory examples, and geometric arguments are supported by pictures wherever possible.

\begin{tcolorbox}
博弈论的概念具有精确且往往不常见的数学定义。在本书中，每个主要概念通常通过详细的引入性示例进行解释。我们采用定义、定理和证明的形式，以确保表述精确，并将一些技术细节隐藏在证明中。证明详细且完整，可以逐行研读。每个证明的主要思想通过引入性示例传达，并且在可能的情况下，用图形支持几何论证。
\end{tcolorbox}

Great attention is given to details that help avoid unnecessary confusions. For example, a random event with two real values \(x\) and \(y\) as outcomes will be described with a probability \(p\) assigned to the second outcome \(y\) , so that the interval \(\left\lbrack  {0,1}\right\rbrack\) for \(p\) corresponds naturally to the interval \(\left\lbrack  {x,y}\right\rbrack\) of the expectation \(\left( {1 - p}\right) x + {py}\) of that event.

\begin{tcolorbox}
我们非常注重避免不必要的混淆。例如，一个有两个实数值 \(x\) 和 \(y\) 作为结果的随机事件，会将概率 \(p\) 分配给第二个结果 \(y\) ，这样 \(p\) 的区间 \(\left\lbrack  {0,1}\right\rbrack\) 自然对应于该事件期望 \(\left( {1 - p}\right) x + {py}\) 的区间 \(\left\lbrack  {x,y}\right\rbrack\) 。
\end{tcolorbox}

This book emphasizes accessibility, not generality. I worked hard (and enjoyed) on presenting the most direct and overhead-free proof of each result. In studying these proofs, students may encounter for the first time important ideas from topology, convex analysis, and linear programming, and thus may become interested in the more general mathematical theory of these subjects.

\begin{tcolorbox}
本书强调易读性，而非一般性。我努力(并且乐在其中)为每个结果提供最直接、无冗余的证明。在学习这些证明的过程中，学生可能首次接触到拓扑学、凸分析和线性规划(linear programming)中的重要思想，从而可能对这些学科更一般的数学理论产生兴趣。
\end{tcolorbox}

Bibliographic references, historical discussions, and further extensions are deferred to a final section in each chapter, to avoid distractions in the main text. References to further reading may include popular science books that I found relevant and interesting. This book is primarily directed at undergraduate students and not at researchers, but I have tried to be historically accurate. For important works in French or German I have also cited English translations where I could find them.

\begin{tcolorbox}
参考文献、历史讨论和进一步的拓展内容放在每章的最后一节，以免干扰正文的阅读。推荐的进一步阅读材料可能包括我认为相关且有趣的科普书籍。本书主要面向本科生，而非研究人员，但我尽量保证历史信息的准确性。对于重要的法语或德语著作，只要能找到英文译本，我也会引用。
\end{tcolorbox}

References to statements inside this book are in upper case, like "Theorem 2.2", and to other works in lower case, like "theorem 13.6 of Roughgarden (2016)".

\begin{tcolorbox}
本书内的陈述引用使用大写，如“定理2.2”，对其他著作的引用使用小写，如“拉夫加登(Roughgarden，2016)的定理13.6”。
\end{tcolorbox}

I use both genders to refer to players in a game, in particular in two-player games to distinguish them more easily. Usage of the pronoun "she" has become natural, and (like "he") it may stand collectively for "he or she".

\begin{tcolorbox}
我在描述博弈中的参与者时会使用两种性别指代，特别是在两人博弈中，这样可以更轻松地区分参与者。使用代词“她”已经变得很自然，并且(和“他”一样)它可以泛指“他或她”。
\end{tcolorbox}

Each chapter concludes with a set of exercises. They test methods (such as how to find mixed equilibria), or the understanding of mathematical concepts used in proofs (in particular in Chapters 5 and 7). Complete solutions to the exercises are available from the publisher for registered lecturers.

\begin{tcolorbox}
每章结尾都有一组练习题。这些练习题用于检验方法(如如何找到混合均衡)，或对证明中使用的数学概念的理解(特别是第5章和第7章)。已注册的教师可以从出版商处获取练习题的完整解答。
\end{tcolorbox}

\section*{Acknowledgments}

\begin{tcolorbox}
\section*{致谢}
\end{tcolorbox}

Some of the material for this book has been adapted and developed from a course guide originally produced for the distance learning International Programme offered by the University of London (http://london.ac.uk/).Rahul Savani and George Zouros gave valuable comments on that first course guide. Many students, some of them anonymously, have pointed out errors and needs for clarification in the subsequent versions of lecture notes. Rahul Savani proof-read the final version with great care. Vanessa Wells drew the nice animal pictures on the cover and in Figures 8.1 and 9.6. I thank Robert Aumann, Paul Dütting, Aviezri Fraenkel, Sergiu Hart, Urban Larsson, Abraham Neyman, Richard Nowakowski, David Tranah, Peyton Young, and Shmuel Zamir for further suggestions.

\begin{tcolorbox}
本书的部分内容改编自伦敦大学(http://london.ac.uk/)远程学习国际项目最初编写的课程指南。拉胡尔·萨瓦尼(Rahul Savani)和乔治·佐罗斯(George Zouros)对第一版课程指南提出了宝贵意见。许多学生(有些匿名)指出了后续版本讲义中的错误，并提出了需要澄清的地方。拉胡尔·萨瓦尼仔细校对了最终版本。凡妮莎·韦尔斯(Vanessa Wells)绘制了封面以及图8.1和图9.6中可爱的动物图片。我感谢罗伯特·奥曼(Robert Aumann)、保罗·迪廷(Paul Dütting)、阿维兹里·弗伦克尔(Aviezri Fraenkel)、塞尔吉乌·哈特(Sergiu Hart)、乌尔班·拉尔森(Urban Larsson)、亚伯拉罕·内曼(Abraham Neyman)、理查德·诺瓦科夫斯基(Richard Nowakowski)、大卫·特拉纳(David Tranah)、佩顿·扬(Peyton Young)和什穆埃尔·扎米尔(Shmuel Zamir)提出的进一步建议。
\end{tcolorbox}

Martin Antonov is the current lead programmer of the Game Theory Explorer software (http://www.gametheoryexplorer.org), a joint project with Rahul Savani and (as part of the Gambit software) Theodore Turocy, for solving games in strategic and extensive form. From Game Theory Explorer, one can export games into the LaTeX formats used in this book for drawing payoff tables and game trees (available on request). 1

\begin{tcolorbox}
马丁·安东诺夫(Martin Antonov)是博弈论探索者软件(http://www.gametheoryexplorer.org)的现任首席程序员，该软件是与拉胡尔·萨瓦尼(Rahul Savani)以及(作为Gambit软件的一部分)西奥多·图罗西(Theodore Turocy)的联合项目，用于求解战略型和扩展型博弈。通过博弈论探索者软件，可以将博弈导出为本书用于绘制收益表和博弈树的LaTeX格式(可按需获取)。1
\end{tcolorbox}

\section*{Nim and Combinatorial Games}

\begin{tcolorbox}
\section*{尼姆游戏与组合博弈}
\end{tcolorbox}

Combinatorial game theory is about perfect-information two-player games, such as Checkers, Go, Chess, or Nim, which are analyzed using their rules. It tries to answer who will win in a game position (assuming optimal play on both sides), and to quantify who is ahead and by how much. The topic has a rich mathematical theory that relates to discrete mathematics, algebra, and (not touched here) computational complexity, and highly original ideas specific to these games.

\begin{tcolorbox}
组合博弈论研究的是完全信息的两人博弈，如跳棋、围棋、国际象棋或尼姆游戏(Nim)，通过其规则对这些博弈进行分析。它试图回答在一个博弈位置(假设双方都进行最优策略)谁会获胜，并量化哪一方领先以及领先多少。该主题有丰富的数学理论，与离散数学、代数以及(此处未涉及)计算复杂性相关，还有针对这些博弈的极具原创性的思想。
\end{tcolorbox}

Combinatorial games are not part of "classical" game theory as used in economics. However, they nicely demonstrate that game theory is about rigorous, and often unfamiliar, mathematical concepts rather than complex techniques.

\begin{tcolorbox}
组合博弈并非经济学中“经典”博弈论的一部分。然而，它们很好地证明了博弈论涉及的是严谨且往往不熟悉的数学概念，而非复杂的技术。
\end{tcolorbox}

This chapter is only an introduction to combinatorial games. It presents the theory of impartial games where in any game position both players have the same allowed moves. We show the powerful and surprisingly simple result (Theorem 1.14), independently found by Sprague (1935) and Grundy (1939), that every impartial game is equivalent to a "Nim heap" of suitable size.

\begin{tcolorbox}
本章仅为组合博弈的入门介绍。它介绍了无偏博弈(impartial game)理论，在任何博弈位置，双方玩家都有相同的可行移动。我们展示了一个强大且惊人简单的结果(定理1.14)，该结果由斯普拉格(Sprague，1935年)和格伦迪(Grundy，1939年)独立发现，即每个无偏博弈都等价于一个合适大小的“尼姆堆”(Nim heap)。
\end{tcolorbox}

In Section 1.8 we give a short glimpse into the more general theory of partizan games, where the allowed moves may depend on the player (e.g., one player can move the white pieces on the game board and the other player the black pieces).

\begin{tcolorbox}
在1.8节中，我们简要介绍更一般的有偏博弈(partizan game)理论，其中可行移动可能取决于玩家(例如，一名玩家可以移动棋盘上的白子，另一名玩家可以移动黑子)。
\end{tcolorbox}

For a deeper treatment, the final Section 1.9 of this chapter lists some excellent textbooks on combinatorial games. They treat impartial games as a special case of general combinatorial games. In contrast, we first treat the simpler impartial games in full.

\begin{tcolorbox}
为了更深入地探讨，本章的最后一节1.9列出了一些关于组合博弈的优秀教科书。它们将无偏博弈作为一般组合博弈的特殊情况进行处理。相比之下，我们首先全面处理较简单的无偏博弈。
\end{tcolorbox}

\section*{1.1 Prerequisites and Learning Outcomes}

\begin{tcolorbox}
\section*{1.1 先决条件与学习成果}
\end{tcolorbox}

The combinatorial games that we consider are finite, and dealing with them uses a type of mathematical induction that we explain further in Section 1.3. It is helpful to have seen induction about the natural numbers before. Although not essential, is is useful to know the algebraic concept of an abelian group, in particular addition modulo 2. We assume familiarity with the binary system where numbers are written in base 2, using only the two digits 0 and 1, rather than in the familiar base 10 .

\begin{tcolorbox}
我们考虑的组合博弈是有限的，处理这些博弈使用一种数学归纳法，我们将在1.3节进一步解释。之前了解过关于自然数的归纳法会有所帮助。虽然不是必需的，但了解阿贝尔群(abelian group)的代数概念，特别是模2加法，会很有用。我们假设熟悉二进制(binary system)，即数字以2为基数书写，仅使用0和1这两个数字，而不是熟悉的以10为基数。
\end{tcolorbox}

After studying this chapter, you should be able to:

\begin{tcolorbox}
学习本章后，你应该能够:
\end{tcolorbox}

\begin{itemize}
\item play Nim optimally;
\end{itemize}

\begin{tcolorbox}
\begin{itemize}
\item 最优地玩尼姆游戏；
\end{itemize}
\end{tcolorbox}

\begin{itemize}
\item explain the different concepts of options, game sums, equivalent games, Nim values, and the mex rule;
\end{itemize}

\begin{tcolorbox}
\begin{itemize}
\item 解释选项(option)、博弈之和(game sum)、等价博弈、尼姆值(Nim value)和mex规则等不同概念；
\end{itemize}
\end{tcolorbox}

\begin{itemize}
\item apply these concepts to play other impartial games like those described in the exercises;
\end{itemize}

\begin{tcolorbox}
\begin{itemize}
\item 将这些概念应用于玩其他无偏博弈，如练习中描述的博弈；
\end{itemize}
\end{tcolorbox}

\begin{itemize}
\item understand how some (not all) game positions in partizan games can be expressed as numbers that state how many moves the Left player is safely ahead, and why these numbers may be fractions with powers of two in the denominator.
\end{itemize}

\begin{tcolorbox}
\begin{itemize}
\item 理解有偏博弈中的一些(并非全部)博弈位置如何表示为数字，这些数字表示左玩家(Left player)安全领先的移动步数，以及为什么这些数字可能是分母为2的幂的分数。
\end{itemize}
\end{tcolorbox}

\section*{1.2 Nim}

\begin{tcolorbox}
\section*{1.2 尼姆游戏}
\end{tcolorbox}

Combinatorial games are two-player win-lose games of perfect information, that is, every player is perfectly informed about the state of play (unlike, for example, the card games Bridge or Poker that have hidden information). The games do not have chance moves like rolling dice or shuffling cards. When playing the game, the two players always alternate in making a move. Every play of the game ends with a win for one player and a loss for the other player (some games like Chess allow for a draw as an outcome, but not the games we consider here).

\begin{tcolorbox}
组合博弈是一种双人胜负博弈，具有完美信息，即每个玩家都完全了解游戏的状态(这与像桥牌或扑克这类具有隐藏信息的纸牌游戏不同)。这类游戏没有掷骰子或洗牌等随机动作。在游戏过程中，两名玩家总是轮流行动。每一局游戏的结果都是一名玩家获胜，另一名玩家失败(像国际象棋这类游戏允许平局，但我们这里讨论的游戏不包括这种情况)。
\end{tcolorbox}

The game has a (typically finite) number of positions, with well-defined rules that define the allowed moves to reach the next position. The rules are such that play will always come to an end because some player is unable to move. This is called the ending condition. We assume the normal play convention that a player unable to move loses. The alternative to normal play is misère play, where a player who is unable to move wins (so the previous player who has made the last move loses).

\begin{tcolorbox}
游戏有(通常是有限数量的)博弈位置，并有明确的规则来定义到达下一个位置所允许的行动。这些规则使得游戏总会结束，因为总会有玩家无法行动。这被称为终止条件。我们采用正常规则，即无法行动的玩家输掉游戏。与正常规则相对的是反向玩法，在反向玩法中，无法行动的玩家获胜(即上一个做出最后行动的玩家失败)。
\end{tcolorbox}

We study impartial games where the available moves in a game position do not depend on whose turn it is to move. If that is not the case, as in Chess where one player can only move the white pieces and the other player the black pieces, the game is called partizan.

\begin{tcolorbox}
我们研究无偏博弈，在无偏博弈中，博弈位置的可用行动不取决于轮到谁行动。如果不是这种情况，例如在国际象棋中，一名玩家只能移动白子，另一名玩家只能移动黑子，这种游戏就被称为有偏博弈。
\end{tcolorbox}

For impartial games, the game Nim plays a central role. A game position in Nim is given by some heaps of tokens, and a move is to remove some (at least one, possibly all) tokens from one of the heaps. The last player able to move wins the game, according to the normal play convention.

\begin{tcolorbox}
对于无偏博弈，尼姆游戏(Nim)起着核心作用。尼姆游戏的一个博弈位置由一些尼姆堆(Nim heap)的代币表示，一次行动是从其中一个尼姆堆中移除一些(至少一个，可能全部)代币。根据正常规则，最后能够行动的玩家赢得游戏。
\end{tcolorbox}

We analyze the Nim position1,2,3, which means three heaps with one, two, and three tokens, respectively. One possible move is to remove two tokens from the heap of size three, like here:

\begin{tcolorbox}
我们分析尼姆游戏位置1,2,3，这意味着有三个尼姆堆，分别有一、二和三个代币。一种可能的行动是从有三个代币的尼姆堆中移除两个代币，就像这里:
\end{tcolorbox}

\begin{center}
\includegraphics[max width=0.7\textwidth]{images/bo_d2e5cmbef24c73b03tn0_14_324_291_1004_78_0.jpg}
\end{center}
\hspace*{3em} 

which we write as a move from1,2,3to1,2,1. Because the move can be made in any one heap, the order of the heap sizes does not matter, so the position1,2,1 could also be written as1,1,2. The options of a game position are the positions that can be reached by a single legal move (according to the game rules) from the player to move. We draw them with moves shown as downward lines, like here,

\begin{tcolorbox}
我们将其写为从1,2,3到1,2,1的行动。因为行动可以在任何一个尼姆堆上进行，所以尼姆堆大小的顺序无关紧要，因此位置1,2,1也可以写成1,1,2。一个博弈位置的选项是当前玩家通过一次合法行动(根据游戏规则)能够到达的位置。我们用向下的线表示行动，就像这里，
\end{tcolorbox}

(1.1)

\begin{center}
\includegraphics[max width=0.7\textwidth]{images/bo_d2e5cmbef24c73b03tn0_14_346_678_901_196_0.jpg}
\end{center}
\hspace*{3em} 

where the first option 2,3 is obtained by removing from1,2,3the entire heap of size 1, the second option 1, 1, 3 by removing one token from the heap of size 2, and so on. The game tree is obtained by continuing to draw all possible moves in this way until play ends (game trees are studied in much more detail in Chapter 4). We may conflate options with obvious equal meaning, such as the positions1,1,2 and 1, 2, 1 that can be reached from 1, 2, 2. However, we do not draw moves to the same position from two different predecessors, such as1,1,2that can be reached from1,1,3and1,2,2. Instead, such a position like1,1,2will be repeated in the game tree, so that every position has a unique history of moves.

\begin{tcolorbox}
其中第一个选项2,3是通过从1,2,3中移除整个大小为1的尼姆堆得到的，第二个选项1,1,3是通过从大小为2的尼姆堆中移除一个代币得到的，依此类推。博弈树(game tree)是通过继续以这种方式绘制所有可能的行动直到游戏结束得到的(在第4章将更详细地研究博弈树)。我们可以合并具有明显相同含义的选项，例如可以从1,2,2到达的位置1,1,2和1,2,1。然而，我们不会从两个不同的前一个位置绘制到同一个位置的行动，例如可以从1,1,3和1,2,2到达的1,1,2。相反，像1,1,2这样的位置会在博弈树中重复出现，这样每个位置都有唯一的行动历史(history of moves)。
\end{tcolorbox}

In an impartial game, the available moves in a game position are by definition independent of the player to move. A game position belongs therefore to exactly one of two possible outcome classes, namely it is either a winning or a losing position. "Winning" or "losing" applies to the player whose turn it is to move, assuming optimal play. A winning position means that the player can force a win with a suitable first "winning" move (and subsequent winning moves at all later positions). A losing position means that every move from the current position leads to a winning position of the other player, who can then force a win, so that the current player will lose.

\begin{tcolorbox}
在无偏博弈中，根据定义，博弈位置的可用行动与当前行动的玩家无关。因此，一个博弈位置恰好属于两种可能的结果类(outcome class)之一，即它要么是必胜局势(winning position)，要么是失败位置(losing position)。“必胜”或“失败”适用于轮到行动的玩家，前提是进行最优玩法。必胜局势意味着玩家可以通过合适的第一步“必胜步”(winning move)(以及后续在所有后续位置的必胜步)迫使获胜。失败位置意味着从当前位置的每一步行动都会导致另一名玩家处于必胜局势，然后该玩家可以迫使获胜，因此当前玩家将失败。
\end{tcolorbox}

When drawn in full with all possible subsequent moves, the Nim position 1,2,3has already a rather large game tree. However, we can tell from its options in (1.1) that1,2,3is a losing position, using the following simple observations of how to play optimally in a Nim position with at most two heaps:

\begin{tcolorbox}
当完整绘制出所有可能的后续行动时，尼姆游戏位置1,2,3的博弈树已经相当大。然而，我们可以从(1.1)中的选项判断出1,2,3是一个失败位置，使用以下关于在最多有两个尼姆堆的尼姆游戏位置中如何进行最优玩法的简单观察:
\end{tcolorbox}

\begin{itemize}
\item If there are no Nim heaps left, this defines a losing position, by the normal play convention.
\end{itemize}

\begin{tcolorbox}
\begin{itemize}
\item 如果没有尼姆堆剩余，根据正常规则，这定义为一个失败位置。
\end{itemize}
\end{tcolorbox}

\begin{itemize}
\item A single Nim heap is a winning position; the winning move is to remove the entire heap.
\end{itemize}

\begin{tcolorbox}
\begin{itemize}
\item 一个尼姆堆是一个必胜局势；必胜步是移除整个尼姆堆。
\end{itemize}
\end{tcolorbox}

\begin{itemize}
\item Two Nim heaps are a winning position if and only if the heaps have different sizes; the winning move is to equalize their sizes by suitably reducing the larger heap, which creates a losing position of two equal heaps. If the two heaps are equal, then a Nim move means reducing one of the heaps, so that the two heaps become unequal, which is a winning position; hence, two equal heaps define a losing position, as claimed.
\end{itemize}

\begin{tcolorbox}
\begin{itemize}
\item 两个尼姆堆是必胜局势当且仅当这两个尼姆堆的大小不同；必胜步是通过适当地减少较大的尼姆堆使其大小相等，这会创建一个两个相等尼姆堆的失败位置。如果两个尼姆堆相等，那么一次尼姆游戏行动意味着减少其中一个尼姆堆，使得两个尼姆堆变得不相等，这是一个必胜局势；因此，两个相等的尼姆堆定义为一个失败位置，如所声称的那样。
\end{itemize}
\end{tcolorbox}

Every option of1,2,3shown in (1.1) is a winning position, because it allows moving to a two-heap losing position, as follows:

\begin{tcolorbox}
(1.1) 中所示的 1,2,3 的每个选项都是必胜局势，因为它允许移动到一个双堆失败位置，如下所示:
\end{tcolorbox}

(1.2)

\begin{center}
\includegraphics[max width=0.7\textwidth]{images/bo_d2e5cmbef24c73b03tn0_15_481_615_905_332_0.jpg}
\end{center}
\hspace*{3em} 

This picture shows sequences of two moves from position1,2,3, but not all of them, so this is not the top part of the full game tree. Only from the first position1,2,3, which is a losing position, all moves are shown. For the second move, only a single, winning move is shown, which leads again to a losing position. This suffices to prove that1,2,3is a losing position in Nim, by the following observation.

\begin{tcolorbox}
这张图展示了从位置 1,2,3 开始的两步走法序列，但不是所有走法，所以这不是完整博弈树的顶部部分。只有从第一个位置 1,2,3(这是一个失败位置)出发的所有走法都被展示出来了。对于第二步走法，只展示了一个必胜步，它又会导致一个失败位置。根据以下观察，这足以证明 1,2,3 在尼姆游戏中是一个失败位置。
\end{tcolorbox}

Lemma 1.1. In an impartial game, a game position is losing if and only if all its options are winning positions. A game position is winning if and only if at least one of its options is a losing position; moving there is a winning move.

\begin{tcolorbox}
引理 1.1:在无偏博弈中，一个博弈位置是失败位置当且仅当它的所有选项都是必胜局势。一个博弈位置是必胜局势当且仅当它的至少一个选项是失败位置；移动到那个位置就是一个必胜步。
\end{tcolorbox}

This lemma applies also to the game position that has no options at all, so trivially ("vacuously") all these options are winning (because there aren't any) and hence the position is losing. The lemma can also be taken as a definition of losing and winning positions, which is not circular but recursive because the options of a game position are simpler in the sense of being closer to the end of any play in the game. In Section 1.3 we will make this precise, and will give a formal proof of Lemma 1.1.

\begin{tcolorbox}
这个引理也适用于根本没有选项的博弈位置，所以显然(“空”意义上)所有这些选项都是必胜的(因为根本没有选项)，因此该位置是失败位置。这个引理也可以作为失败位置和必胜局势的定义，它不是循环定义而是递归定义，因为一个博弈位置的选项在更接近游戏任何一轮结束的意义上更简单。在 1.3 节中，我们将对此进行精确阐述，并给出引理 1.1 的正式证明。
\end{tcolorbox}

The concepts of winning and losing position assume optimal play. Binmore (2007, p. 49) writes: "A dull art movie called Last Year in Marienbad consists largely of the characters playing Nim very badly." In that movie (from 1961, directed by Alain Resnais), they play several times misère Nim, always from the starting position1,3,5,7. There is a scene where three successive positions are

\begin{tcolorbox}
必胜局势和失败位置的概念假设了最优玩法。宾莫尔(2007 年，第 49 页)写道:“一部沉闷的艺术电影《去年在马里昂巴德》很大程度上是角色们把尼姆游戏玩得很糟糕的情节。”在那部电影(1961 年上映，由阿伦·雷乃执导)中，他们多次玩反向尼姆游戏，总是从起始位置 1,3,5,7 开始。有一个场景中，三个连续的位置是
\end{tcolorbox}

\[
\text{ player I player II }
\]

\[
2,3,6 \rightarrow  2,3,5 \rightarrow  2,3,1
\]

and then player I, whose turn it is, eventually loses. We can conclude from this that player I indeed played badly. Namely, either2,3,1is a winning position, in which case player I failed to force a win. Or2,3,1is a losing position (which is in fact the case), in which case player I failed to move there from position 2,3,6.

\begin{tcolorbox}
然后轮到玩家 I 行动，最终玩家 I 输了。我们可以由此得出结论，玩家 I 确实玩得很糟糕。也就是说，要么 2,3,1 是一个必胜局势，在这种情况下玩家 I 未能确保获胜。要么 2,3,1 是一个失败位置(实际上就是这种情况)，在这种情况下玩家 I 未能从位置 2,3,6 移动到该位置。
\end{tcolorbox}

On the other hand, a game is more entertaining if it is not known how to play it optimally. In Chess, for example, which allows for a draw, one can show similarly to Lemma 1.1 that either White can force a win, or Black can force a win, or both players can force a draw. However, Chess is not yet "solved" in the sense that we know which outcome results from optimal play, which is why it is still an interesting game because players do not play perfectly.

\begin{tcolorbox}
另一方面，如果不知道如何最优地玩一个游戏，这个游戏会更有趣。例如在国际象棋中，存在和棋的情况，我们可以类似引理 1.1 那样证明，要么白方可以确保获胜，要么黑方可以确保获胜，要么双方都可以确保和棋。然而，国际象棋在我们知道最优玩法会导致何种结果的意义上还没有被“解决”，这就是为什么它仍然是一个有趣的游戏，因为玩家们不会完美地玩。
\end{tcolorbox}

Nim, however, has been solved by Bouton (1901). We demonstrate his winning strategy for the already formidable Nim position

\begin{tcolorbox}
然而，尼姆游戏已经被布顿(1901 年)解决了。我们将展示他针对已经很有挑战性的尼姆位置的必胜策略
\end{tcolorbox}

\begin{center}
\includegraphics[max width=0.4\textwidth]{images/bo_d2e5cmbef24c73b03tn0_16_547_726_554_226_0.jpg}
\end{center}
\hspace*{3em} 

which is4,5,9,14and a winning position. The trick is to write the heap sizes in the binary system. The following tables show that one winning move is to reduce the Nim heap of size 5 to size 3, which leads to the losing position4,3,9,14:

\begin{tcolorbox}
即 4,5,9,14，这是一个必胜局势。诀窍是将尼姆堆的大小用二进制表示。以下表格显示，一个必胜步是将大小为 5 的尼姆堆减少到大小为 3，这会导致失败位置 4,3,9,14:
\end{tcolorbox}

(1.3)

\begin{center}
\includegraphics[max width=0.7\textwidth]{images/bo_d2e5cmbef24c73b03tn0_16_376_1145_900_291_0.jpg}
\end{center}
\hspace*{3em} 

In (1.3), the top row represents the powers of 2 used to represent each heap size with binary digits, such as \(5 = 0 \cdot  8 + 1 \cdot  4 + 0 \cdot  2 + 1 \cdot  1\) . The bottom row is called the Nim sum of the heap sizes and is obtained by "addition modulo 2 without carry" for each column; that is, the Nim sum has a 0 if the number of 1 's in that column is even, and a 1 if that number is odd.

\begin{tcolorbox}
在 (1.3) 中，顶行代表用于用二进制数字表示每个尼姆堆大小的 2 的幂，例如 \(5 = 0 \cdot  8 + 1 \cdot  4 + 0 \cdot  2 + 1 \cdot  1\) 。底行称为尼姆堆大小的尼姆和，是通过对每一列进行“无进位模 2 加法”得到的；也就是说，如果该列中 1 的个数是偶数，尼姆和为 0，如果个数是奇数，尼姆和为 1。
\end{tcolorbox}

We claim that a Nim position is losing if and only if that Nim sum is zero for all columns, like on the right in (1.3); we call such a position a zero position. In order to show this, we have to show, in line with Lemma 1.1, that

\begin{tcolorbox}
我们声称，一个尼姆位置是失败位置当且仅当对于所有列，该尼姆和都为零，就像 (1.3) 右边那样；我们称这样的位置为零位置。为了证明这一点，我们必须根据引理 1.1 证明:
\end{tcolorbox}

(a) every move from a zero position leads to a position which is not zero and therefore winning, and

\begin{tcolorbox}
(a) 从一个零位置的每一步走法都会导致一个非零位置，因此是必胜局势，并且
\end{tcolorbox}

(b) from every winning position (with a Nim sum that is not zero) there is a move to a zero position.

\begin{tcolorbox}
(b) 从每个必胜局势(尼姆和不为零)都有一步走法可以到达一个零位置。
\end{tcolorbox}

It easy to see (a): A move changes exactly one Nim heap, which corresponds to one row in the table of binary numbers, and changes at least one binary digit of the binary representation of the heap size. Therefore, the resulting Nim sum is also changed in the corresponding digits and becomes nonzero.

\begin{tcolorbox}
很容易看出 (a):一步走法恰好改变一个尼姆堆，这对应于二进制数表中的一行，并且至少改变尼姆堆大小二进制表示中的一个二进制数字。因此，得到的尼姆和在相应的数字上也会改变，并且变为非零。
\end{tcolorbox}

Condition (b) amounts to finding at least one winning move for a position that is not zero, like on the left in (1.3). Such a winning move is found by these steps:

\begin{tcolorbox}
条件 (b) 相当于为一个非零位置(如 (1.3) 左边的位置)找到至少一个必胜步。可以通过以下步骤找到这样的必胜步:
\end{tcolorbox}

\begin{itemize}
\item Choose the leftmost " 1 " column in the Nim sum (which exists because not all columns are 0 ), which represents the highest power of two used in the binary representation of the Nim sum, here \(4 = {2}^{2}\) .
\end{itemize}

\begin{tcolorbox}
\begin{itemize}
\item 选择尼姆和(Nim sum)中最左边的“1”列(该列存在是因为并非所有列都是 0)，它代表尼姆和的二进制表示中使用的最高 2 的幂次，这里是 \(4 = {2}^{2}\) 。
\end{itemize}
\end{tcolorbox}

\begin{itemize}
\item By definition, there is an odd number of Nim heaps which use that power of two in their binary representation. In (1.3), these are the heaps of sizes 4, 5, and 14. Choose one of the corresponding rows; we choose the heap of size 5 .
\end{itemize}

\begin{tcolorbox}
\begin{itemize}
\item 根据定义，有奇数个尼姆堆(Nim heap)在其二进制表示中使用了该 2 的幂次。在 (1.3) 中，这些是大小为 4、5 和 14 的堆。选择其中一个对应的行；我们选择大小为 5 的堆。
\end{itemize}
\end{tcolorbox}

\begin{itemize}
\item In the chosen row, "flip" the binary digits for each 1 in the Nim-sum, here for the columns 4 and 2; the resulting changes are shown with boxes on the right in (1.3). The largest of the changed powers of two in that row has digit 1, which is changed to 0 . Even if all subsequent digits were changed from 0 to 1, the resulting binary number gives a new, smaller integer (possibly zero, which will mean removing the entire heap). The winning move is to reduce the Nim heap to that smaller size, here from 5 to 3 .
\end{itemize}

\begin{tcolorbox}
\begin{itemize}
\item 在所选行中，对尼姆和中每个为 1 的二进制位进行“翻转”，这里是第 4 列和第 2 列；(1.3) 中右侧的方框显示了所做的更改。该行中被更改的 2 的幂次中最大的那个位原来是 1，现在将其更改为 0。即使所有后续位都从 0 变为 1，得到的二进制数也会是一个新的、更小的整数(可能为零，这意味着移除整个堆)。必胜步(winning move)是将尼姆堆的大小减小到那个更小的值，这里是从 5 减小到 3。
\end{itemize}
\end{tcolorbox}

As a result of these steps, the Nim sum is changed from 1 to 0 exactly in the columns where there was a 1 before, with the digits 0 untouched, so the new Nim sum has a 0 in each column. The resulting Nim position is therefore a zero position, as required in (b).

\begin{tcolorbox}
通过这些步骤，尼姆和恰好会在之前为 1 的列中从 1 变为 0，而 0 的位保持不变，因此新的尼姆和每一列都是 0。所以得到的尼姆博弈位置(game position)是一个零位置，符合 (b) 的要求。
\end{tcolorbox}

This winning strategy in Nim can be done with pen and paper, but is probably difficult to perform in one's head in actual play. For small heap sizes, a helpful mental image is to group the tokens in each Nim heap (now shown as dots) into distinct powers of two:

\begin{tcolorbox}
尼姆游戏(Nim)中的这个必胜策略可以用笔和纸来完成，但在实际游戏中可能很难在脑海中完成。对于较小的堆大小，一个有用的心理图像是将每个尼姆堆中的代币(现在用点表示)分组为不同的 2 的幂次:
\end{tcolorbox}

\begin{center}
\includegraphics[max width=0.3\textwidth]{images/bo_d2e5cmbef24c73b03tn0_17_726_1477_477_343_0.jpg}
\end{center}
\hspace*{3em} 

(1.4)

This is an equivalent representation to the left table in (1.3), but without binary digits. It shows that the powers of two that occur an odd number of times are 4 and 2. A winning move has to change one heap so that afterwards they occur an even number of times. This can be done by changing the four dots in the heap of size 5 to two dots, as in \(:  :  \rightarrow   :\) . This is the winning move in (1.3) that removes two tokens from the heap of size 5 ; the same winning move can be made in the heap of size 4 . The third possible winning move is to remove six tokens (:: and :) from the heap of size 14. Caution: The resulting powers of two in a heap size must be distinct. The change \(:  :  \rightarrow   :\) is not allowed in the heap of size 14, because the intended result ::::: : : is not a binary representation because 2 appears twice; rather, the result of removing two tokens from the heap of size 14 is :::: :::, and then the overall created Nim position is not zero. Once the heap for the winning move is chosen, the number of tokens that have to be removed is unique.

\begin{tcolorbox}
这与 (1.3) 中左边的表格是等价的表示，但没有二进制位。它表明出现奇数次的 2 的幂次是 4 和 2。一个必胜步必须改变一个堆，使得之后这些幂次出现偶数次。这可以通过将大小为 5 的堆中的四个点变为两个点来实现，如 \(:  :  \rightarrow   :\) 所示。这就是 (1.3) 中从大小为 5 的堆中移除两个代币的必胜步；同样的必胜步也可以在大小为 4 的堆中进行。第三种可能的必胜步是从大小为 14 的堆中移除六个代币(:: 和 :)。注意:堆大小中得到的 2 的幂次必须是不同的。在大小为 14 的堆中不允许进行 \(:  :  \rightarrow   :\) 的更改，因为预期的结果 ::::: : : 不是一个二进制表示，因为 2 出现了两次；相反，从大小为 14 的堆中移除两个代币的结果是 :::: ::: ，然后整体得到的尼姆博弈位置不是零位置。一旦选择了用于必胜步的堆，必须移除的代币数量是唯一的。
\end{tcolorbox}

\section*{1.3 Top-down Induction}

\begin{tcolorbox}
\section*{1.3 自上而下归纳(top-down induction)}
\end{tcolorbox}

When talking about combinatorial games, we will often use for brevity the word game for "game position". Every game \(G\) has finitely many options \({G}_{1},\ldots ,{G}_{m}\) that are reached from \(G\) by one of the allowed moves in \(G\) , as in this picture:

\begin{tcolorbox}
在讨论组合博弈时，为简洁起见，我们经常用“博弈”这个词来表示“博弈位置”。每个博弈 \(G\) 都有有限个选项 \({G}_{1},\ldots ,{G}_{m}\) ，这些选项可以通过 \(G\) 中允许的一步移动从 \(G\) 到达，如图所示:
\end{tcolorbox}

(1.5)

\begin{center}
\includegraphics[max width=0.3\textwidth]{images/bo_d2e5cmbef24c73b03tn0_18_628_862_391_188_0.jpg}
\end{center}
\hspace*{3em} 

If \(m = 0\) then \(G\) has no options. We denote the game with no options by 0, which by the normal play convention is a losing game. Otherwise the options of \(G\) are themselves games, defined by their respective options according to the rules of the game. In that way, any game is completely defined by its options. In short, the starting position defines the game completely.

\begin{tcolorbox}
如果 \(m = 0\) ，那么 \(G\) 没有选项。我们用 0 表示没有选项的博弈，根据正常规则(normal play)约定，这是一个必败博弈。否则， \(G\) 的选项本身也是博弈，根据游戏规则由它们各自的选项来定义。这样，任何博弈都完全由其选项来定义。简而言之，起始位置完全定义了这个博弈。
\end{tcolorbox}

We introduce a certain type of mathematical induction for games, which is applied to a partial order (see the background material text box on the next page).

\begin{tcolorbox}
我们为博弈引入一种特定类型的数学归纳法，它应用于一个偏序(partial order)(见下一页的背景材料(background material)文本框)。
\end{tcolorbox}

Consider a set \(S\) of games, defined, for example, by a starting game and all the games that can reached from it via any sequence of moves of the players. For two games \(G\) and \(H\) in \(S\) , call \(H\) simpler than \(G\) if there is a sequence of moves that leads from \(G\) to \(H\) . We allow \(G = H\) where this sequence is empty. The relation of being "simpler than" defines a partial order which for the moment we denote by \(\leq\) . Note that \(\leq\) is antisymmetric because it is not possible to reach \(G\) from \(G\) by a nonempty sequence of moves because this would violate the ending condition. The ending condition for games implies the following property:

\begin{tcolorbox}
考虑一个博弈集合 \(S\) ，例如，它由一个起始博弈以及所有可以通过玩家的任何移动序列从该起始博弈到达的博弈定义。对于 \(S\) 中的两个博弈 \(G\) 和 \(H\) ，如果存在一个从 \(G\) 到 \(H\) 的移动序列，就称 \(H\) 比 \(G\) 更简单。我们允许 \(G = H\) 这种移动序列为空的情况。“比……更简单”的关系定义了一个偏序，目前我们用 \(\leq\) 表示。注意， \(\leq\) 是反对称的(antisymmetric)，因为不可能通过非空的移动序列从 \(G\) 回到 \(G\) ，因为这会违反终止条件(ending condition)。博弈的终止条件意味着以下性质:
\end{tcolorbox}

\[
\text{ Every nonempty subset of }S\text{ has a minimal element. } \tag{1.8}
\]

If there was a nonempty subset \(T\) of \(S\) without a minimal element, then we could produce an infinite play as follows: Start with some \(G\) in \(T\) . Because \(G\) is not minimal, there is some \(H\) in \(T\) with \(H < G\) , so there is some sequence of moves from \(G\) to \(H\) . Similarly, \(H\) is not minimal, so another game in \(T\) is reached from \(H\) . Continuing in this manner creates an infinite sequence of moves, which contradicts the ending condition.

\begin{tcolorbox}
如果 \(S\) 存在一个非空子集 \(T\) 没有极小元，那么我们可以按照如下方式构造一个无限对局:从 \(T\) 中的某个 \(G\) 开始。因为 \(G\) 不是极小元，所以 \(T\) 中存在某个 \(H\) 满足 \(H < G\) ，因此存在从 \(G\) 到 \(H\) 的一系列走法。类似地， \(H\) 也不是极小元，所以从 \(H\) 又能到达 \(T\) 中的另一个对局。按照这种方式继续下去会产生一个无限的走法序列，这与终止条件相矛盾。
\end{tcolorbox}

\section*{Background material: Partial orders}

\begin{tcolorbox}
\section*{背景材料:偏序}
\end{tcolorbox}

Definition 1.2 (Partial order, total order). A binary relation \(\leq\) on a set \(S\) is called a partial order if the following hold for all \(x,y,z\) in \(S\) :

\begin{tcolorbox}
定义 1.2(偏序、全序)。集合 \(S\) 上的二元关系 \(\leq\) 若对于 \(S\) 中的所有 \(x,y,z\) 都满足以下条件，则称其为偏序:
\end{tcolorbox}

\(x \leq  y\) and \(y \leq  z \Rightarrow  x \leq  z\;\left( { \leq  \text{ is transitive }}\right)\) ,

\begin{tcolorbox}
\(x \leq  y\) 和 \(y \leq  z \Rightarrow  x \leq  z\;\left( { \leq  \text{ is transitive }}\right)\) ，
\end{tcolorbox}

\(x \leq  x\) (≤ is reflexive), and

\begin{tcolorbox}
\(x \leq  x\) (≤ 是自反的)，并且
\end{tcolorbox}

\(x \leq  y\) and \(y \leq  x \Rightarrow  x = y\;\left( { \leq  \text{ is antisymmetric }}\right)\) .

\begin{tcolorbox}
\(x \leq  y\) 和 \(y \leq  x \Rightarrow  x = y\;\left( { \leq  \text{ is antisymmetric }}\right)\) 。
\end{tcolorbox}

If, in addition, for all \(x,y\) in \(S\)

\begin{tcolorbox}
此外，如果对于 \(S\) 中的所有 \(x,y\)
\end{tcolorbox}

\(x \leq  y\) or \(y \leq  x\;\left( { \leq  \text{ is total }}\right)\)

\begin{tcolorbox}
\(x \leq  y\) 或 \(y \leq  x\;\left( { \leq  \text{ is total }}\right)\)
\end{tcolorbox}

then \(\leq\) is called a total order .

\begin{tcolorbox}
那么 \(\leq\) 称为全序。
\end{tcolorbox}

For a given partial order \(\leq\) , we often say " \(x\) is less than or equal to \(y\) " if \(x \leq  y\) . We then define \(x < y\) (" \(x\) is less than \({y}^{\prime \prime }\) ) as follows:

\begin{tcolorbox}
对于给定的偏序 \(\leq\) ，如果 \(x \leq  y\) ，我们常说“ \(x\) 小于或等于 \(y\) ”。然后我们按如下方式定义 \(x < y\) (“ \(x\) 小于 \({y}^{\prime \prime }\) ”):
\end{tcolorbox}

\[
x < y\; \Leftrightarrow  \;x \leq  y\;\text{ and }\;x \neq  y. \tag{1.6}
\]

This relation \(<\) is also called the strict order that corresponds to \(\leq\) . Exercise 1.1 asks you to show how a partial order \(\leq\) can be defined from a relation \(<\) on \(S\) with suitable properties (transitivity and "irreflexivity") that define a strict order. Then \(\leq\) is obtained from \(<\) according to

\begin{tcolorbox}
这个关系 \(<\) 也称为与 \(\leq\) 对应的严格序。练习 1.1 要求你展示如何从集合 \(S\) 上具有合适性质(传递性和“非自反性”)的关系 \(<\) 定义一个偏序 \(\leq\) ，这些性质定义了一个严格序。然后根据以下方式从 \(<\) 得到 \(\leq\)
\end{tcolorbox}

\[
x \leq  y\; \Leftrightarrow  \;x < y\;\text{ or }\;x = y. \tag{1.7}
\]

An element \(x\) of \(S\) is called minimal if there is no \(y\) in \(S\) with \(y < x\) .

\begin{tcolorbox}
如果 \(S\) 中不存在 \(y\) 满足 \(y < x\) ，则 \(S\) 中的元素 \(x\) 称为极小元。
\end{tcolorbox}

Any set \(S\) of rational or real numbers has the familiar relation \(\leq\) , which is a total order. The most important partial order that is usually not total is the set inclusion \(\subseteq\) on a set \(S\) of sets. Then elements \(x,y\) of \(S\) as used in Definition 1.2 are usually written with capital letters like \(A,B\) . The partial order is written with the symbol \(\subseteq\) where \(A \subseteq  B\) holds if \(A\) is a subset of \(B\) (which allows for \(A = B\) ). The inclusion order is not total because two sets may be incomparable, that is, neither \(A \subseteq  B\) nor \(B \subseteq  A\) hold, for example if \(A = \{ 1,2\}\) and \(B = \{ 2,3,4\}\) . If \(S\) is the set of all subsets of, say, the finite set \(\{ 1,\ldots ,n\}\) for some positive integer \(n\) , then \(S\) , partially ordered by \(\subseteq\) , has the empty set \(\varnothing\) as a unique minimal element. If \(S\) is the set of nonempty subsets of \(\{ 1,\ldots ,n\}\) , then its minimal elements are the singletons \(\{ 1\} ,\{ 2\} ,\ldots ,\{ n\}\) .

\begin{tcolorbox}
任何有理数或实数集 \(S\) 都有我们熟悉的关系 \(\leq\) ，它是一个全序。通常不是全序的最重要的偏序是集合族 \(S\) 上的集合包含关系 \(\subseteq\) 。那么定义 1.2 中使用的 \(S\) 的元素 \(x,y\) 通常用大写字母如 \(A,B\) 表示。偏序用符号 \(\subseteq\) 表示，其中如果 \(A\) 是 \(B\) 的子集(允许 \(A = B\) )，则 \(A \subseteq  B\) 成立。包含序不是全序，因为两个集合可能不可比较，即 \(A \subseteq  B\) 和 \(B \subseteq  A\) 都不成立，例如当 \(A = \{ 1,2\}\) 且 \(B = \{ 2,3,4\}\) 时。如果 \(S\) 是某个正整数 \(n\) 对应的有限集 \(\{ 1,\ldots ,n\}\) 的所有子集构成的集合，那么由 \(\subseteq\) 给出偏序的 \(S\) 有唯一的极小元，即空集 \(\varnothing\) 。如果 \(S\) 是 \(\{ 1,\ldots ,n\}\) 的非空子集构成的集合，那么它的极小元是单元素集 \(\{ 1\} ,\{ 2\} ,\ldots ,\{ n\}\) 。
\end{tcolorbox}

For any partial order with the property (1.8), the following theorem applies.

\begin{tcolorbox}
对于任何具有性质 (1.8) 的偏序，以下定理成立。
\end{tcolorbox}

Theorem 1.3 (Top-down induction). Consider a set \(S\) with a partial order \(\leq\) that fulfills (1.8). Let \(P\left( x\right)\) be a statement that may be true or false for an element \(x\) of \(S\) , and assume that \(P\left( x\right)\) holds whenever \(P\left( y\right)\) holds for all \(y\) of \(S\) with \(y < x\) . Then \(P\left( x\right)\) is true for all \(x\) in \(S\) . In symbols, where " \(\forall x\) " means "for all elements \(x\) of \({S}^{\prime \prime }\) :

\begin{tcolorbox}
定理1.3(自上而下归纳)。考虑一个带有偏序 \(\leq\) 的集合 \(S\) ，该偏序满足 (1.8)。设 \(P\left( x\right)\) 是一个关于 \(S\) 中元素 \(x\) 的陈述，其可能为真也可能为假，并且假设只要对于 \(S\) 中所有满足 \(y < x\) 的 \(y\) ， \(P\left( y\right)\) 都成立时， \(P\left( x\right)\) 就成立。那么 \(P\left( x\right)\) 对于 \(S\) 中的所有 \(x\) 都为真。用符号表示，其中“ \(\forall x\) ”表示“对于 \({S}^{\prime \prime }\) 中的所有元素 \(x\) :
\end{tcolorbox}

\[
\left( {\forall x : \left( {\forall y < x : P\left( y\right) }\right)  \Rightarrow  P\left( x\right) }\right)  \Rightarrow  \left( {\forall x : P\left( x\right) }\right) . \tag{1.9}
\]

Proof. Suppose, as assumed, that condition \(P\left( x\right)\) holds whenever it holds for all elements \(y\) of \(S\) with \(y < x\) . Consider the set \(T = \{ z \in  S \mid  P\left( z\right)\) is false \(\}\) . Then \(P\left( x\right)\) is true for all \(x\) in \(S\) if \(T\) is the empty set, so assume \(T\) is not empty. Let \(x\) be a minimal element of \(T\) , which exists by (1.8). Then all \(y\) in \(S\) with \(y < x\) do not belong to \(T\) , that is, \(P\left( y\right)\) is true. But, by assumption (1.9), this implies that \(P\left( x\right)\) is also true, which contradicts \(x \in  T\) . So \(T\) is empty, as claimed.

\begin{tcolorbox}
证明。假设，如所假定的，只要对于 \(S\) 中所有满足 \(y < x\) 的元素 \(y\) ，条件 \(P\left( x\right)\) 都成立。考虑集合 \(T = \{ z \in  S \mid  P\left( z\right)\) 为假 \(\}\) 。如果 \(T\) 是空集，那么对于 \(S\) 中的所有 \(x\) ， \(P\left( x\right)\) 都为真，所以假设 \(T\) 不是空集。设 \(x\) 是 \(T\) 的一个极小元(根据 (1.8) 它存在)。那么 \(S\) 中所有满足 \(y < x\) 的 \(y\) 都不属于 \(T\) ，即 \(P\left( y\right)\) 为真。但是，根据假设 (1.9)，这意味着 \(P\left( x\right)\) 也为真，这与 \(x \in  T\) 矛盾。所以，如所声称的， \(T\) 是空集。
\end{tcolorbox}

Theorem 1.3 is called a method of induction because it assumes that a property \(P\left( x\right)\) can be inferred from the fact that \(P\left( y\right)\) is true for all \(y\) that are smaller than \(x\) . Then, assuming (1.8), the property holds for all \(x\) in \(S\) . An important application of this induction method is for proving the prime number decomposition theorem \(P\left( x\right)\) , which states that if \(x\) is an integer and \(x \geq  2\) , then \(x\) can be written as a product of prime numbers. In this induction proof, one first shows that \(x\) has a prime factor \(p\) , where either \(x = p\) (in which case we have shown \(P\left( x\right)\) and are done) or \(x = p \cdot  y\) for some integer \(y\) with \(2 \leq  y < x\) . In the latter case, \(P\left( y\right)\) holds by the "inductive hypothesis", that is, \(y = {p}_{1}{p}_{2}\cdots {p}_{k}\) for some prime numbers \({p}_{1},{p}_{2},\ldots ,{p}_{k}\) . Hence, \(x = p{p}_{1}{p}_{2}\cdots {p}_{k}\) , which also shows \(P\left( x\right)\) . By Theorem 1.3, \(P\left( x\right)\) holds throughout. The induction can be applied because any nonempty set of positive integers has a minimal element.

\begin{tcolorbox}
定理1.3被称为一种归纳方法，因为它假设一个性质 \(P\left( x\right)\) 可以从以下事实推出:对于所有小于 \(x\) 的 \(y\) ， \(P\left( y\right)\) 都为真。然后，假设 (1.8) 成立，该性质对于 \(S\) 中的所有 \(x\) 都成立。这种归纳方法的一个重要应用是证明质数分解定理 \(P\left( x\right)\) ，该定理指出，如果 \(x\) 是一个整数且 \(x \geq  2\) ，那么 \(x\) 可以写成质数的乘积。在这个归纳证明中，首先证明 \(x\) 有一个质因数 \(p\) ，其中要么 \(x = p\) (在这种情况下我们就证明了 \(P\left( x\right)\) ，证明完成)，要么对于某个满足 \(2 \leq  y < x\) 的整数 \(y\) 有 \(x = p \cdot  y\) 。在后一种情况下，根据“归纳假设”， \(P\left( y\right)\) 成立，即对于某些质数 \({p}_{1},{p}_{2},\ldots ,{p}_{k}\) 有 \(y = {p}_{1}{p}_{2}\cdots {p}_{k}\) 。因此， \(x = p{p}_{1}{p}_{2}\cdots {p}_{k}\) ，这也证明了 \(P\left( x\right)\) 。根据定理1.3， \(P\left( x\right)\) 始终成立。可以应用这种归纳法是因为任何非空的正整数集合都有一个极小元。
\end{tcolorbox}

The method is called "top-down" induction because we start from the "top" (here the element \(x\) ) and use the property \(P\left( y\right)\) in some way for elements \(y\) that are smaller than \(x\) , where we typically encounter \(y\) without knowing how much smaller \(y\) is compared to \(x\) , as in the prime number decomposition theorem. In contrast, ordinary induction for natural numbers \(n\) proceeds typically by proving \(P\left( 1\right)\) and then the implication \(P\left( n\right)  \Rightarrow  P\left( {n + 1}\right)\) , that is, one step at a time, which then implies that \(P\left( n\right)\) holds for all natural numbers \(n\) . Top-down induction is more powerful; it also applies to partial orders and not just to the totally ordered natural numbers (where it is often called "strong induction").

\begin{tcolorbox}
这种方法被称为“自上而下”归纳法，因为我们从“顶部”(这里是元素 \(x\) )开始，并以某种方式对小于 \(x\) 的元素 \(y\) 使用性质 \(P\left( y\right)\) ，通常我们遇到 \(y\) 时并不知道它比 \(x\) 小多少，就像在质数分解定理中一样。相比之下，自然数 \(n\) 的普通归纳法通常是先证明 \(P\left( 1\right)\) ，然后证明蕴含关系 \(P\left( n\right)  \Rightarrow  P\left( {n + 1}\right)\) ，即一次一步，这意味着 \(P\left( n\right)\) 对所有自然数 \(n\) 都成立。自上而下归纳法更强大；它不仅适用于全序的自然数(在那里它通常被称为“强归纳法”)，也适用于偏序。
\end{tcolorbox}

Where did the base case go, which is normally proved separately in an induction proof? The base case states that \(P\left( x\right)\) holds for any minimal element \(x\) (if \(\leq\) is total, then there is only one minimal element). The answer is that the base case is covered by the assumption in (1.9): The left-hand side of the implication (1.9) says that

\begin{tcolorbox}
通常在归纳证明中单独证明的基础情况去哪里了呢？基础情况表明 \(P\left( x\right)\) 对任何极小元 \(x\) 都成立(如果 \(\leq\) 是全序的，那么只有一个极小元)。答案是基础情况包含在 (1.9) 的假设中:蕴含关系 (1.9) 的左边表示
\end{tcolorbox}

\[
\left( {\forall y < x : P\left( y\right) }\right)  \Rightarrow  P\left( x\right)  \tag{1.10}
\]

holds for all \(x\) , in particular for all minimal \(x\) . But then \(\left( {\forall y < x : P\left( y\right) }\right)\) is vacuously true because there is no \(y\) with \(y < x\) . Therefore, for minimal \(x\) , condition (1.10) means that \(P\left( x\right)\) is true, which is the base case.

\begin{tcolorbox}
对所有 \(x\) 都成立，特别是对所有极小的 \(x\) 。但此时 \(\left( {\forall y < x : P\left( y\right) }\right)\) 显然为真，因为不存在满足 \(y < x\) 的 \(y\) 。因此，对于极小的 \(x\) ，条件 (1.10) 意味着 \(P\left( x\right)\) 为真，这就是基础情况。
\end{tcolorbox}

As a proof technique, we can often use Theorem 1.3 generally without any case distinctions about whether the considered elements are minimal or not. It also allows using a property about games recursively by referring to the same property applied to simpler games.

\begin{tcolorbox}
作为一种证明技巧，我们通常可以普遍使用定理 1.3 ，而无需区分所考虑的元素是否为极小元。它还允许通过引用应用于更简单博弈的相同性质来递归地使用关于博弈的性质。
\end{tcolorbox}

Using top-down induction, we can now prove Lemma 1.1, which states that any game \(G\) is either winning or losing. Assume that this holds for all games that are simpler than \(G\) , in particular the options of \(G\) . If all of these options are winning, then \(G\) is a losing position, because the other player can force a win no matter which move the player makes in \(G\) . If not all options of \(G\) are winning, then one of them, say \(H\) , is losing, and by moving to \(H\) the player forces a win, hence \(G\) is a winning position. This argument is not in any way deeper than what we said before. Theorem 1.3 merely asserts that it is rigorous (and that it uses the ending condition).

\begin{tcolorbox}
使用自上而下归纳法，我们现在可以证明引理 1.1 ，该引理指出任何博弈 \(G\) 要么是必胜局势，要么是失败位置。假设这对所有比 \(G\) 更简单的博弈都成立，特别是 \(G\) 的选项。如果所有这些选项都是必胜局势，那么 \(G\) 是一个失败位置，因为无论玩家在 \(G\) 中采取什么行动，另一个玩家都能迫使获胜。如果 \(G\) 的选项并非都是必胜局势，那么其中一个，比如 \(H\) ，是失败位置，通过移动到 \(H\) ，玩家可以迫使获胜，因此 \(G\) 是一个必胜局势。这个论证并不比我们之前所说的更深刻。定理 1.3 只是断言它是严谨的(并且它使用了终止条件)。
\end{tcolorbox}

\section*{1.4 Game Sums and Equivalence of Games}

\begin{tcolorbox}
\section*{1.4 博弈之和与博弈等价性}
\end{tcolorbox}

Combinatorial games often "decompose" into parts in which players can move independently, and the players then have to decide in which part to make their move. This is captured by the important concept of a sum of games.

\begin{tcolorbox}
组合博弈通常会“分解”成玩家可以独立行动的部分，然后玩家必须决定在哪个部分行动。这可以通过博弈之和这个重要概念来体现。
\end{tcolorbox}

Definition 1.4. Suppose that \(G\) and \(H\) are game positions with options (positions reached by one move) \({G}_{1},\ldots ,{G}_{k}\) and \({H}_{1},\ldots ,{H}_{m}\) , respectively. Then the options of the game sum \(G + H\) are

\begin{tcolorbox}
定义 1.4 。假设 \(G\) 和 \(H\) 是博弈位置，它们的选项(通过一步行动到达的位置)分别是 \({G}_{1},\ldots ,{G}_{k}\) 和 \({H}_{1},\ldots ,{H}_{m}\) 。那么博弈之和 \(G + H\) 的选项是
\end{tcolorbox}

\[
{G}_{1} + H,\ldots ,{G}_{k} + H,\;G + {H}_{1},\ldots ,G + {H}_{m}\text{ . } \tag{1.11}
\]

The first list of options \({G}_{1} + H,\ldots ,{G}_{k} + H\) in (1.11) simply means that the player makes his move in \(G\) , the second list \(G + {H}_{1},\ldots ,G + {H}_{m}\) that he makes his move in \(H\) ; the other part of the game sum remains untouched. As an example, a Nim position is simply the game sum of its individual Nim heaps, because the player moves in exactly one of the heaps.

\begin{tcolorbox}
(1.11) 中的第一组选项 \({G}_{1} + H,\ldots ,{G}_{k} + H\) 仅仅意味着玩家在 \(G\) 中行动，第二组选项 \(G + {H}_{1},\ldots ,G + {H}_{m}\) 意味着玩家在 \(H\) 中行动；博弈之和的另一部分保持不变。例如，一个尼姆游戏位置仅仅是其各个尼姆堆的博弈之和，因为玩家恰好可以在其中一个堆中行动。
\end{tcolorbox}

Definition 1.4 is a recursive definition, because the game sum is defined in terms of its options, which are themselves game sums (but they are simpler games).

\begin{tcolorbox}
定义 1.4 是一个递归定义，因为博弈之和是根据其选项来定义的，而这些选项本身也是博弈之和(但它们是更简单的博弈)。
\end{tcolorbox}

The sum of games turns out to define an abelian group on the (appropriately defined) set of games. It is a commutative and associative operation: for any games \(G,H,J,\)

\begin{tcolorbox}
事实证明，博弈之和在(适当定义的)博弈集合上定义了一个阿贝尔群。它是一个满足交换律和结合律的运算:对于任何博弈 \(G,H,J,\)
\end{tcolorbox}

\[
G + H = H + G\;\text{ and }\;\left( {G + H}\right)  + J = G + \left( {H + J}\right) . \tag{1.12}
\]

The first condition (commutativity) holds because the order of the options of a game, used in (1.11), does not matter. The second condition (associativity) holds because both \(\left( {G + H}\right)  + J\) and \(G + \left( {H + J}\right)\) mean in effect that the player decides to move in \(G\) , in \(H\) , or in \(J\) , leaving the other two parts of the game sum unchanged. We can therefore assume the equalities (1.12). More generally, in a sum of several games \({G}_{1},\ldots ,{G}_{n}\) the player moves in exactly one of these games, which does not depend on how these games are arranged, so that we can write this sum unambiguously without parentheses as \({G}_{1} + \cdots  + {G}_{n}\) .

\begin{tcolorbox}
第一个条件(交换律)成立，因为在(1.11)中使用的博弈选项的顺序无关紧要。第二个条件(结合律)成立，因为 \(\left( {G + H}\right)  + J\) 和 \(G + \left( {H + J}\right)\) 实际上都意味着玩家决定在 \(G\) 、 \(H\) 或 \(J\) 中移动，而博弈之和的另外两部分保持不变。因此，我们可以假设等式(1.12)成立。更一般地，在几个博弈 \({G}_{1},\ldots ,{G}_{n}\) 的和中，玩家恰好只在其中一个博弈中移动，这与这些博弈的排列方式无关，所以我们可以明确地将这个和写成 \({G}_{1} + \cdots  + {G}_{n}\) 而无需括号。
\end{tcolorbox}

The losing game 0 which has no options is a zero (neutral element) for game sums: It fulfills \(G + 0 = G\) for any game \(G\) , because the game 0 is "invisible" when added to \(G\) .

\begin{tcolorbox}
没有选项的必败博弈 0 是博弈之和的零元(中性元素):对于任何博弈 \(G\) ，它满足 \(G + 0 = G\) ，因为当 0 与 \(G\) 相加时，博弈 0 是“不可见的”。
\end{tcolorbox}

In order to obtain a group, every game \(G\) needs to have a "negative" game \(- G\) so that \(G + \left( {-G}\right)  = 0\) . However, this equality cannot hold as stated as soon as the game \(G\) has options, because then the game \(G + \left( {-G}\right)\) also has options but 0 has none. Instead, we need a more general condition

\begin{tcolorbox}
为了构成一个群，每个博弈 \(G\) 都需要有一个“负”博弈 \(- G\) ，使得 \(G + \left( {-G}\right)  = 0\) 成立。然而，一旦博弈 \(G\) 有选项，这个等式就不能按所述方式成立，因为此时博弈 \(G + \left( {-G}\right)\) 也有选项，而 0 没有选项。相反，我们需要一个更一般的条件
\end{tcolorbox}

\[
G + \left( {-G}\right)  \equiv  0, \tag{1.13}
\]

where \(G \equiv  H\) means that the two games \(G\) and \(H\) are equivalent, according to the following definition.

\begin{tcolorbox}
其中 \(G \equiv  H\) 表示根据以下定义，两个博弈 \(G\) 和 \(H\) 是等价的。
\end{tcolorbox}

Definition 1.5. Two games \(G,H\) are called equivalent, written \(G \equiv  H\) , if and only if for any other game \(J\) , the game sum \(G + J\) is losing if and only if \(H + J\) is losing.

\begin{tcolorbox}
定义 1.5:当且仅当对于任何其他博弈 \(J\) ，博弈之和 \(G + J\) 是必败的当且仅当 \(H + J\) 是必败的时，两个博弈 \(G,H\) 被称为等价的，记为 \(G \equiv  H\) 。
\end{tcolorbox}

Lemma 1.6. The relation \(\equiv\) in Definition 1.5 is indeed an equivalence relation between games, that is, it is reflexive, symmetric and transitive.

\begin{tcolorbox}
引理 1.6:定义 1.5 中的关系 \(\equiv\) 确实是博弈之间的一种等价关系(关系的等价性)，即它是自反的、对称的和传递的。
\end{tcolorbox}

Proof. Reflexivity means \(G \equiv  G\) for any game \(G\) , which holds trivially. Symmetry means that if \(G \equiv  H\) then \(H \equiv  G\) , which holds also by definition. Suppose that \(G \equiv  H\) and \(H \equiv  K\) for games \(G,H,K\) . Transitivity means that then \(G \equiv  K\) . Consider any other game \(J\) . If \(G + J\) is losing then \(H + J\) is losing because \(G \equiv  H\) , and then \(K + J\) is losing because \(H \equiv  K\) . Conversely, if \(K + J\) is losing then \(H + J\) is losing because \(H \equiv  K\) , and then \(G + J\) is losing because \(G \equiv  H\) . Hence, \(G + J\) is losing if and only if \(K + J\) is losing, that is, \(G \equiv  K\) . This shows that \(\equiv\) is transitive.

\begin{tcolorbox}
证明:自反性意味着对于任何博弈 \(G\) 有 \(G \equiv  G\) ，这显然成立。对称性意味着如果 \(G \equiv  H\) 则 \(H \equiv  G\) ，这也根据定义成立。假设对于博弈 \(G,H,K\) 有 \(G \equiv  H\) 和 \(H \equiv  K\) 。传递性意味着 \(G \equiv  K\) 。考虑任何其他博弈 \(J\) 。如果 \(G + J\) 是必败的，那么由于 \(G \equiv  H\) ， \(H + J\) 是必败的，又因为 \(H \equiv  K\) ， \(K + J\) 是必败的。反之，如果 \(K + J\) 是必败的，那么由于 \(H \equiv  K\) ， \(H + J\) 是必败的，再由于 \(G \equiv  H\) ， \(G + J\) 是必败的。因此， \(G + J\) 是必败的当且仅当 \(K + J\) 是必败的，即 \(G \equiv  K\) 。这表明 \(\equiv\) 是传递的。
\end{tcolorbox}

By taking \(J = 0\) in Definition 1.5, if the games \(G\) and \(H\) are equivalent then either both are losing or both are winning. Equivalence of \(G\) and \(H\) means that the games \(G\) and \(H\) always have the same outcome (winning or losing) and that this property is preserved in any game sum \(G + J\) and \(H + J\) for any other game \(J\) . This is much stronger than just belonging to the same outcome class. For example, if \(G\) and \(H\) are Nim heaps of different sizes, then both games are winning but they are not equivalent. We state this as a lemma.

\begin{tcolorbox}
在定义 1.5 中取 \(J = 0\) ，如果博弈 \(G\) 和 \(H\) 是等价的，那么它们要么都是必败的，要么都是必胜的。 \(G\) 和 \(H\) 等价意味着博弈 \(G\) 和 \(H\) 总是有相同的结果(必胜或必败)，并且对于任何其他博弈 \(J\) ，这个性质在任何博弈之和 \(G + J\) 和 \(H + J\) 中都保持不变。这比仅仅属于同一个结果类要强得多。例如，如果 \(G\) 和 \(H\) 是不同大小的尼姆堆(Nim heap)，那么这两个博弈都是必胜的，但它们并不等价。我们将此表述为一个引理。
\end{tcolorbox}

Lemma 1.7. Two Nim heaps are equivalent if and only if they have equal size.

\begin{tcolorbox}
引理 1.7:两个尼姆堆(Nim heap)等价当且仅当它们的大小相等。
\end{tcolorbox}

Proof. Let \(G\) and \(H\) be two Nim heaps of different sizes, and consider \(J = G\) in Definition 1.5. Then \(G + G\) is a losing game, but \(H + G\) is winning. This shows that \(G\) and \(H\) are not equivalent. If \(G\) and \(H\) are Nim heaps of equal size then they are the same game and therefore equivalent.

\begin{tcolorbox}
证明:设 \(G\) 和 \(H\) 是两个不同大小的尼姆堆(Nim heap)，考虑定义 1.5 中的 \(J = G\) 。那么 \(G + G\) 是一个必败博弈，但 \(H + G\) 是必胜的。这表明 \(G\) 和 \(H\) 不等价。如果 \(G\) 和 \(H\) 是大小相等的尼姆堆(Nim heap)，那么它们是同一个博弈，因此是等价的。
\end{tcolorbox}

However, any two losing games are equivalent, because they are all equivalent to the game 0 that has no options.

\begin{tcolorbox}
然而，任意两个必败游戏都是等价的，因为它们都等价于没有选项(option)的游戏 0。
\end{tcolorbox}

Lemma 1.8. \(G\) is a losing game if and only if \(G \equiv  0\) .

\begin{tcolorbox}
引理 1.8。 \(G\) 是必败游戏当且仅当 \(G \equiv  0\) 。
\end{tcolorbox}

Proof. If \(G \equiv  0\) then \(G\) is losing because 0 is losing, using \(J = 0\) in Definition 1.5. So it remains to prove the "only if" part.

\begin{tcolorbox}
证明。若 \(G \equiv  0\) ，则根据定义 1.5 中的 \(J = 0\) ，由于 0 是必败游戏，所以 \(G\) 是必败的。因此，只需证明“仅当”部分。
\end{tcolorbox}

Let \(G\) be a losing game. We want to show \(G \equiv  0\) , that is, for any other game \(J\) , the game \(G + J\) is losing if and only if \(0 + J\) is losing. Because \(0 + J = J\) , this amounts to the two implications

\begin{tcolorbox}
设 \(G\) 是必败游戏。我们要证明 \(G \equiv  0\) ，即对于任意其他游戏 \(J\) ，游戏 \(G + J\) 是必败的当且仅当 \(0 + J\) 是必败的。因为 \(0 + J = J\) ，这等价于两个蕴含关系
\end{tcolorbox}

\[
G + J\text{ is losing } \Rightarrow  J\text{ is losing } \tag{1.14}
\]

and (we add the assumption that \(G\) is losing for symmetry)

\begin{tcolorbox}
并且(为了对称，我们添加假设 \(G\) 是失败位置)
\end{tcolorbox}

\[
G\text{ and }J\text{ are losing } \Rightarrow  G + J\text{ is losing. } \tag{1.15}
\]

Condition (1.15) is rather obvious: If \(G\) and \(J\) are both losing games, there is no way to win the game sum \(G + J\) because the second player always has a counter-move. We prove (1.15) formally by top-down induction. Let \(G\) and \(J\) be losing. We want to show that \(G + J\) is losing, which means that every option of \(G + J\) is a winning game. Any such option is of the form \({G}^{\prime } + J\) where is \({G}^{\prime }\) is an option of \(G\) , or of the form \(G + {J}^{\prime }\) where \({J}^{\prime }\) is an option of \(J\) . Consider the first type of option \({G}^{\prime } + J\) . Because \(G\) is a losing game, its option \({G}^{\prime }\) is a winning position, so \({G}^{\prime }\) has some option \({G}^{\prime \prime }\) that is a losing position. Because \({G}^{\prime \prime }\) is a losing game that is simpler than \(G\) , we can assume that (1.15) holds, that is, \({G}^{\prime \prime }\) and \(J\) are losing \(\Rightarrow  {G}^{\prime \prime } + J\) is losing. So the option \({G}^{\prime } + J\) is winning because it has the losing option \({G}^{\prime \prime } + J\) , and this was an arbitrary option of \(G + J\) with a move to \({G}^{\prime } + J\) in the first game \(G\) of the sum \(G + J\) . The same applies for any option of \(G + J\) of the form \(G + {J}^{\prime }\) where \({J}^{\prime }\) is an option of \(J\) . So all options of \(G + J\) are winning games, and \(G + J\) is therefore a losing game, as claimed. This concludes the proof of (1.15).

\begin{tcolorbox}
条件 (1.15) 相当明显:如果 \(G\) 和 \(J\) 都是失败的博弈，那么就没有办法在博弈之和 \(G + J\) 中获胜，因为第二个玩家总是有应对的策略。我们通过自上而下归纳(top-down induction)正式证明 (1.15)。设 \(G\) 和 \(J\) 是失败的。我们要证明 \(G + J\) 是失败的，这意味着 \(G + J\) 的每个选项(option)都是必胜局势(winning position)。任何这样的选项都具有 \({G}^{\prime } + J\) 的形式，其中 \({G}^{\prime }\) 是 \(G\) 的一个选项，或者具有 \(G + {J}^{\prime }\) 的形式，其中 \({J}^{\prime }\) 是 \(J\) 的一个选项。考虑第一种类型的选项 \({G}^{\prime } + J\) 。因为 \(G\) 是一个失败的博弈，它的选项 \({G}^{\prime }\) 是一个必胜局势，所以 \({G}^{\prime }\) 有某个选项 \({G}^{\prime \prime }\) 是失败位置。因为 \({G}^{\prime \prime }\) 是一个比 \(G\) 更简单的失败博弈，我们可以假设 (1.15) 成立，即 \({G}^{\prime \prime }\) 和 \(J\) 是失败的， \(\Rightarrow  {G}^{\prime \prime } + J\) 是失败的。所以选项 \({G}^{\prime } + J\) 是必胜的，因为它有失败的选项 \({G}^{\prime \prime } + J\) ，并且这是博弈之和 \(G + J\) 的第一个博弈 \(G\) 中向 \({G}^{\prime } + J\) 移动的 \(G + J\) 的任意一个选项。同样的情况适用于 \(G + J\) 的任何形式为 \(G + {J}^{\prime }\) 的选项，其中 \({J}^{\prime }\) 是 \(J\) 的一个选项。所以 \(G + J\) 的所有选项都是必胜的博弈，因此 \(G + J\) 是一个失败的博弈，正如所声称的那样。这就完成了 (1.15) 的证明。
\end{tcolorbox}

We now show (1.14). Suppose that \(G + J\) is losing but that \(J\) is winning. Then there is a winning move to some option \({J}^{\prime }\) of \(J\) where \({J}^{\prime }\) is a losing game. But then by (1.15), which we just proved, this implies that \(G + {J}^{\prime }\) is a losing game, which is an option of \(G + J\) , and would represent a winning move from \(G + J\) , contradicting that \(G + J\) is losing. This shows that (1.14) holds, which completes the proof.

\begin{tcolorbox}
我们现在证明 (1.14)。假设 \(G + J\) 是失败的，但 \(J\) 是必胜的。那么存在一个必胜步(winning move)到 \(J\) 的某个选项 \({J}^{\prime }\) ，其中 \({J}^{\prime }\) 是一个失败的博弈。但根据我们刚刚证明的 (1.15)，这意味着 \(G + {J}^{\prime }\) 是一个失败的博弈，它是 \(G + J\) 的一个选项，并且代表了从 \(G + J\) 出发的一个必胜步，这与 \(G + J\) 是失败的相矛盾。这表明 (1.14) 成立，从而完成了证明。
\end{tcolorbox}

The next lemma shows that adding a game preserves equivalence between two games \(G\) and \(H\) .

\begin{tcolorbox}
下一个引理表明，添加一个博弈可以保持两个博弈 \(G\) 和 \(H\) 之间的等价性。
\end{tcolorbox}

Lemma 1.9. For all games \(G,H,K\) ,

\begin{tcolorbox}
引理 1.9. 对于所有的博弈 \(G,H,K\) ，
\end{tcolorbox}

\[
G \equiv  H \Rightarrow  G + K \equiv  H + K. \tag{1.16}
\]

Proof. Let \(G \equiv  H\) . We show that \(G + K \equiv  H + K\) , that is, \(\left( {G + K}\right)  + J\) is losing \(\Leftrightarrow\)  \(\left( {H + K}\right)  + J\) is losing for any other game \(J\) . By (1.12), this means \(G + \left( {K + J}\right)\) is losing \(\Leftrightarrow  H + \left( {K + J}\right)\) is losing, which holds because \(G\) and \(H\) are equivalent.

\begin{tcolorbox}
证明。设 \(G \equiv  H\) 。我们证明 \(G + K \equiv  H + K\) ，即对于任何其他博弈 \(J\) ， \(\left( {G + K}\right)  + J\) 是失败的 \(\Leftrightarrow\)  \(\left( {H + K}\right)  + J\) 是失败的。根据 (1.12)，这意味着 \(G + \left( {K + J}\right)\) 是失败的 \(\Leftrightarrow  H + \left( {K + J}\right)\) 是失败的，这是成立的，因为 \(G\) 和 \(H\) 是等价的。
\end{tcolorbox}

Using equivalence, any losing game \(Z\) can take the role of a zero (neutral) element for game sums, that is,

\begin{tcolorbox}
利用等价性，任何失败的博弈 \(Z\) 都可以充当博弈之和的零(中性)元素，即
\end{tcolorbox}

\[
Z\text{ is losing } \Rightarrow  G + Z \equiv  G \tag{1.17}
\]

for any game \(G\) . The reason is that any move in the component \(Z\) of the game sum \(G + Z\) has an immediate counter-move back to a losing position, which will eventually become the game 0 without options. Therefore, adding \(Z\) does not affect if \(G\) is winning or losing, and this is also preserved when adding any other game \(J\) . In particular, \(Z \equiv  0\) for any losing game \(Z\) , by taking \(G = 0\) in (1.17). We state (1.17) as a lemma, which is proved easily with the help of the previous lemmas.

\begin{tcolorbox}
对于任意博弈 \(G\) 。原因在于，博弈之和 \(G + Z\) 的分量 \(Z\) 中的任何一步行动都有一个直接的反制行动，可回到失败位置，最终该博弈将变为没有选项的博弈 0。因此，添加 \(Z\) 不会影响 \(G\) 是必胜还是必败，并且在添加任何其他博弈 \(J\) 时，这一性质也会保留。特别地，对于任何失败博弈 \(Z\) ，通过在 (1.17) 中取 \(G = 0\) ，有 \(Z \equiv  0\) 。我们将 (1.17) 表述为一个引理，借助之前的引理很容易证明该引理。
\end{tcolorbox}

Lemma 1.10. Let \(Z\) be a losing game. Then \(G + Z \equiv  G\) for any game \(G\) .

\begin{tcolorbox}
引理 1.10。设 \(Z\) 为失败博弈。那么对于任意博弈 \(G\) ，有 \(G + Z \equiv  G\) 。
\end{tcolorbox}

Proof. By Lemma 1.8, \(Z \equiv  0\) , and by (1.16) \(G + Z = Z + G \equiv  0 + G = G\) .

\begin{tcolorbox}
证明。根据引理 1.8，有 \(Z \equiv  0\) ，根据 (1.16) 有 \(G + Z = Z + G \equiv  0 + G = G\) 。
\end{tcolorbox}

The next lemma implies that any game \(G\) has a negative game \(- G\) so that (1.13) holds. For impartial games, this negative game is \(G\) itself, that is, \(G + G\) is always a losing position. This is the copycat principle for impartial games: the second player always has a move left, by copying the move of the first player. If \({G}^{\prime }\) is an option of \(G\) and the first player moves in \(G + G\) to \({G}^{\prime } + G\) or \(G + {G}^{\prime }\) , then the second player responds by moving to \({G}^{\prime } + {G}^{\prime }\) . This continues until the first player has run out of moves and loses.

\begin{tcolorbox}
下一个引理表明，任意博弈 \(G\) 都有一个负博弈 \(- G\) ，使得 (1.13) 成立。对于无偏博弈，这个负博弈就是 \(G\) 本身，即 \(G + G\) 总是一个失败位置。这就是无偏博弈的模仿原则:第二名玩家总是可以通过模仿第一名玩家的行动来保留一步行动。如果 \({G}^{\prime }\) 是 \(G\) 的一个选项，并且第一名玩家在 \(G + G\) 中行动到 \({G}^{\prime } + G\) 或 \(G + {G}^{\prime }\) ，那么第二名玩家通过行动到 \({G}^{\prime } + {G}^{\prime }\) 来回应。这一过程会一直持续，直到第一名玩家无步可走并失败。
\end{tcolorbox}

Lemma 1.11 (The copycat principle). \(G + G \equiv  0\) for any impartial game \(G\) .

\begin{tcolorbox}
引理 1.11(模仿原则)。对于任意无偏博弈 \(G\) ，有 \(G + G \equiv  0\) 。
\end{tcolorbox}

Proof. Let \(G\) be an impartial game. By Lemma 1.8, we have to show that \(G + G\) is a losing game. We prove this by top-down induction, by showing that every option of \(G + G\) is a winning game. Any such option of \(G + G\) is of the form \({G}^{\prime } + G\) or \(G + {G}^{\prime }\) for an option \({G}^{\prime }\) of \(G\) . Then the next player has the winning move to the game \({G}^{\prime } + {G}^{\prime }\) which is losing, by the inductive assumption because \({G}^{\prime } + {G}^{\prime }\) is simpler than \(G + G\) . So \(G + G\) is indeed a losing game.

\begin{tcolorbox}
证明。设 \(G\) 为无偏博弈。根据引理 1.8，我们需要证明 \(G + G\) 是一个失败博弈。我们通过自上而下归纳法来证明这一点，即证明 \(G + G\) 的每个选项都是必胜博弈。 \(G + G\) 的任何这样的选项都具有 \({G}^{\prime } + G\) 或 \(G + {G}^{\prime }\) 的形式，其中 \({G}^{\prime }\) 是 \(G\) 的一个选项。然后，下一名玩家有一个必胜步行动到博弈 \({G}^{\prime } + {G}^{\prime }\) ，根据归纳假设， \({G}^{\prime } + {G}^{\prime }\) 是失败的，因为 \({G}^{\prime } + {G}^{\prime }\) 比 \(G + G\) 更简单。所以 \(G + G\) 确实是一个失败博弈。
\end{tcolorbox}

Lemma 1.10 states that any losing game in a game sum is irrelevant for the outcome class of that sum. For example, by Lemma 1.11 a losing game of this sort may be the sum of two equal games in a larger game sum. This can simplify game sums substantially. For example, in the Nim position1,2,3,5,5,5, the game sum of two Nim heaps of size 5 is a losing game, and the position1,2,3is also losing, as we saw earlier, so the overall position is equivalent to a single Nim heap of size 5 .

\begin{tcolorbox}
引理 1.10 表明，博弈之和中的任何失败博弈对于该和的结果类都是无关紧要的。例如，根据引理 1.11，这种失败博弈可能是一个更大博弈之和中两个相等博弈的和。这可以大大简化博弈之和。例如，在尼姆游戏位置 1,2,3,5,5,5 中，两个大小为 5 的尼姆堆的博弈之和是一个失败博弈，并且如我们之前所见，位置 1,2,3 也是失败的，所以整个位置等价于一个大小为 5 的单个尼姆堆。
\end{tcolorbox}

The following condition is very useful to prove that two games are equivalent.

\begin{tcolorbox}
以下条件对于证明两个博弈等价非常有用。
\end{tcolorbox}

Lemma 1.12. Two impartial games \(G,H\) are equivalent if and only if \(G + H \equiv  0\) . Proof. If \(G \equiv  H\) , then by (1.16) and Lemma 1.11, \(G + H \equiv  H + H \equiv  0\) . Conversely, \(G + H \equiv  0\) implies \(G \equiv  G + H + H \equiv  0 + H = H\) .

\begin{tcolorbox}
引理 1.12。两个无偏博弈 \(G,H\) 等价当且仅当 \(G + H \equiv  0\) 。证明。如果 \(G \equiv  H\) ，那么根据 (1.16) 和引理 1.11，有 \(G + H \equiv  H + H \equiv  0\) 。反之， \(G + H \equiv  0\) 意味着 \(G \equiv  G + H + H \equiv  0 + H = H\) 。
\end{tcolorbox}

The following lemma states that two games \(G\) and \(H\) are equivalent if all their options are equivalent.

\begin{tcolorbox}
以下引理表明，如果两个博弈 \(G\) 和 \(H\) 的所有选项都等价，那么这两个博弈等价。
\end{tcolorbox}

Lemma 1.13. Consider two games \(G\) and \(H\) so that for every option of \(G\) there is an equivalent option of \(H\) and vice versa. Then \(G \equiv  H\) .

\begin{tcolorbox}
引理 1.13。考虑两个博弈 \(G\) 和 \(H\) ，使得对于 \(G\) 的每个选项，都有 \(H\) 的一个等价选项，反之亦然。那么 \(G \equiv  H\) 。
\end{tcolorbox}

Proof. We show that \(G + H\) is losing, by showing that every option of \(G + H\) is a winning game. Consider such an option \({G}^{\prime } + H\) where \({G}^{\prime }\) is an option of \(G\) . By assumption, there is an option \({H}^{\prime }\) of \(H\) with \({G}^{\prime } \equiv  {H}^{\prime }\) , that is, \({G}^{\prime } + {H}^{\prime } \equiv  0\) by Lemma 1.12, so \({G}^{\prime } + {H}^{\prime }\) is losing. Moving there defines a winning move in \({G}^{\prime } + H\) . By the same argument, any option \(G + {H}^{\prime }\) of \(G + H\) is winning. Hence, every option of \(G + H\) is winning, so \(G + H\) is a losing game. Therefore, \(G + H \equiv  0\) , and \(G \equiv  H\) by Lemma 1.12.

\begin{tcolorbox}
证明。我们通过证明 \(G + H\) 的每个选项都是必胜博弈，来证明 \(G + H\) 是失败博弈。考虑这样一个选项 \({G}^{\prime } + H\) ，其中 \({G}^{\prime }\) 是 \(G\) 的一个选项。根据假设，存在 \(H\) 的一个选项 \({H}^{\prime }\) 满足 \({G}^{\prime } \equiv  {H}^{\prime }\) ，即根据引理 1.12 有 \({G}^{\prime } + {H}^{\prime } \equiv  0\) ，所以 \({G}^{\prime } + {H}^{\prime }\) 是失败博弈。移动到该位置定义了 \({G}^{\prime } + H\) 中的一个必胜步。通过相同的论证， \(G + H\) 的任何选项 \(G + {H}^{\prime }\) 都是必胜的。因此， \(G + H\) 的每个选项都是必胜的，所以 \(G + H\) 是一个失败博弈。所以， \(G + H \equiv  0\) ，并且根据引理 1.12 有 \(G \equiv  H\) 。
\end{tcolorbox}

Note that Lemma 1.13 states only a sufficient condition for the equivalence of \(G\) and \(H\) . Games can be equivalent without that property. For example, \(G + G\) is equivalent to 0, but \(G + G\) has in general many options and 0 has none.

\begin{tcolorbox}
注意，引理 1.13 仅陈述了 \(G\) 和 \(H\) 等价的一个充分条件。博弈在不满足该性质的情况下也可能等价。例如， \(G + G\) 等价于 0，但一般来说 \(G + G\) 有很多选项，而 0 没有选项。
\end{tcolorbox}

We conclude this section with a comment on the abelian group of impartial games. The group operation is the sum + of games. In order for the group to work, we have to identify equivalent games, which are "compatible" with game sums by Lemma 1.9. Any losing game (which is equivalent to the game 0 with no options) is a zero element of the group operation. Every game is its own negative. Essentially (although we will not prove it that way), in our context the only groups with this property are given by the set of binary vectors with component-wise addition modulo 2. This may serve as an intuition why the optimal strategy of Nim as in (1.3) uses the binary representation based on powers of two (and not powers of three, say). Moreover, Nim plays a central game for all impartial games. We will prove this directly, without any reference to abelian groups, in the next section.

\begin{tcolorbox}
我们以对无偏博弈的阿贝尔群(abelian group)的评论来结束本节。群运算为博弈的和 +。为了使该群成立，我们必须识别等价博弈，根据引理 1.9，这些等价博弈与博弈之和是“相容的”。任何失败博弈(它等价于没有选项的博弈 0)是群运算的零元。每个博弈都是其自身的逆元。本质上(尽管我们不会这样证明)，在我们的上下文中，具有此性质的唯一群是由按分量模 2 加法的二进制向量集给出的。这可以解释为什么如 (1.3) 中尼姆游戏(Nim)的最优策略使用基于 2 的幂(而不是 3 的幂等)的二进制表示。此外，尼姆游戏在所有无偏博弈中起着核心作用。我们将在下一节直接证明这一点，而不涉及阿贝尔群。
\end{tcolorbox}

A second comment concerns misère games, where the last player to move loses. One losing position in misère Nim is a single heap with one token in it. However, adding such a heap \(Z\) , say, to a game \(G\) completely reverses the outcome class of \(G\) because the token can just be removed, so we certainly do not have \(G + Z \equiv  G\) as in Lemma 1.10, no matter how the equivalence \(\equiv\) is defined. Also, 1,1is a winning position but2,2is a losing position in misère Nim, so there is no easy copycat principle for misère games. Although misère Nim is well understood (see Exercise 1.3), general misère games are not. In contrast, game sums have a nice structure with the normal play convention.

\begin{tcolorbox}
第二个评论涉及反常博弈(misère games)，即最后行动的玩家失败的博弈。反向尼姆游戏(misère Nim)中的一个失败位置是只有一个筹码的单个尼姆堆(Nim heap)。然而，例如，将这样的一个堆 \(Z\) 添加到博弈 \(G\) 中会完全反转 \(G\) 的结果类，因为那个筹码可以被移除，所以无论等价关系 \(\equiv\) 如何定义，我们肯定不会有引理 1.10 中的 \(G + Z \equiv  G\) 。此外，在反向尼姆游戏中，(1,1) 是一个必胜局势，但 (2,2) 是一个失败位置，所以反常博弈没有简单的模仿原则(copycat principle)。尽管反向尼姆游戏已被很好地理解(见练习 1.3)，但一般的反常博弈并非如此。相比之下，在正常规则(normal play)约定下，博弈之和具有良好的结构。
\end{tcolorbox}

\section*{1.5 Nim, Poker Nim, and the Mex Rule}

\begin{tcolorbox}
\section*{1.5 尼姆游戏、扑克尼姆游戏(Poker Nim)和 mex 规则(mex rule)}
\end{tcolorbox}

In this section, we prove a sweeping statement (Theorem 1.14): Every impartial game is equivalent to some Nim heap. Note that this means a single Nim heap, even if the game itself is rather complicated, for example a general Nim position, which is a game sum of several Nim heaps. This can proved without any reference to playing Nim optimally. As we will see at the end of this section with the help of Figure 1.2, the so-called mex rule that underlies Theorem 1.14 can be used to discover the role of powers of two for playing sums of Nim heaps.

\begin{tcolorbox}
在本节中，我们证明一个全面的陈述(定理 1.14):每个无偏博弈都等价于某个尼姆堆。注意，这意味着即使博弈本身相当复杂，例如一般的尼姆游戏位置(它是几个尼姆堆的博弈之和)，也等价于单个尼姆堆。这可以在不涉及尼姆游戏最优策略的情况下证明。正如我们将在本节末尾借助图 1.2 看到的，作为定理 1.14 基础的所谓 mex 规则可用于发现 2 的幂在尼姆堆之和博弈中的作用。
\end{tcolorbox}

We use the following notation for Nim heaps. If \(G\) is a single Nim heap with \(n\) tokens, \(n \geq  0\) , then we denote this game by \(* n\) . This game is completely specified by its \(n\) options, and they are defined recursively as follows:

\begin{tcolorbox}
我们对尼姆堆使用以下符号表示。如果 \(G\) 是一个有 \(n\) 个筹码的单个尼姆堆， \(n \geq  0\) ，那么我们用 \(* n\) 表示这个博弈。这个博弈完全由它的 \(n\) 个选项确定，并且它们递归定义如下:
\end{tcolorbox}

\[
\text{ options of } * n : \; * 0, * 1, * 2,\ldots , * \left( {n - 1}\right) \text{ . } \tag{1.18}
\]

Note that \(* 0\) is the empty heap with no tokens, that is, \(* 0 = 0\) ; we will normally continue to just write 0 .

\begin{tcolorbox}
注意， \(* 0\) 是没有筹码的空堆，即 \(* 0 = 0\) ；我们通常仍只写 0。
\end{tcolorbox}

We can use (1.18) as the definition of \(* n\) . For example, the game \(* 4\) is defined by its options \(* 0, * 1, * 2, * 3\) . It is very important to include \(* 0\) in that list of options, because it means that *4 has a winning move. Condition (1.18) is a recursive definition of the game \(* n\) , because its options are also defined by reference to such games \(* k\) , for numbers \(k\) smaller than \(n\) . This game fulfills the ending condition because the heap gets successively smaller in any sequence of moves.

\begin{tcolorbox}
我们可以使用 (1.18) 作为 \(* n\) 的定义。例如，博弈 \(* 4\) 由其选项 \(* 0, * 1, * 2, * 3\) 定义。在这些选项列表中包含 \(* 0\) 非常重要，因为这意味着 *4 有一个必胜步。条件 (1.18) 是博弈 \(* n\) 的递归定义，因为它的选项也是通过参考此类博弈 \(* k\) 来定义的，其中数字 \(k\) 小于 \(n\) 。这个博弈满足终止条件，因为在任何移动序列中，堆的大小会依次减小。
\end{tcolorbox}

A general Nim position is a game sum of several Nim heaps. Earlier we had written such a position by just listing the sizes of the Nim heaps, such as1,2,3in (1.1). The fancy way to write this is now \(* 1 +  * 2 +  * 3\) , a sum of games.

\begin{tcolorbox}
一般的尼姆游戏局势是几个尼姆堆的博弈之和。早些时候，我们只是列出尼姆堆的大小来表示这样的局势，比如 (1.1) 中的 1, 2, 3。现在更正式的写法是 \(* 1 +  * 2 +  * 3\) ，即博弈的和。
\end{tcolorbox}

The game of Poker Nim is a variation of Nim. Suppose that each player is given, at the beginning of the game, some extra "reserve" tokens. Like Nim, the game is played with heaps of tokens. In a move, a player can choose, as in ordinary Nim, a heap and remove some tokens, which he can add to his reserve tokens. A second, new kind of move is to add some of the player's reserve tokens to some heap (or even to create an entire new heap with these tokens). These two kinds of moves are the only ones allowed.

\begin{tcolorbox}
扑克尼姆游戏是尼姆游戏的一种变体。假设在游戏开始时，每个玩家都会得到一些额外的“储备”代币。和尼姆游戏一样，这个游戏也是用代币堆来玩的。在一次移动中，玩家可以像普通尼姆游戏那样，选择一个堆并移除一些代币，然后将这些代币添加到自己的储备代币中。第二种新的移动方式是将玩家的一些储备代币添加到某个堆中(甚至可以用这些代币创建一个全新的堆)。这两种移动方式是唯一允许的操作。
\end{tcolorbox}

Suppose that there are three heaps, of sizes 1, 2, 5, and that the game has been going on for some time, so that both players have accumulated substantial reserves of tokens. It is player I's turn, who moves to1,2,3because that is a good move in ordinary Nim. But then player II adds 50 tokens to the heap of size 2, creating the position1,52,3, which seems complicated.

\begin{tcolorbox}
假设现在有三个堆，大小分别为 1、2、5，并且游戏已经进行了一段时间，两个玩家都积累了大量的储备代币。现在轮到玩家 I 行动，他移动到 1, 2, 3，因为这在普通尼姆游戏中是一步好棋。但随后玩家 II 向大小为 2 的堆中添加了 50 个代币，形成了 1, 52, 3 的局势，这看起来很复杂。
\end{tcolorbox}

What should player I do? After a moment's thought, he just removes the 50 tokens that player II has just added to the heap, reverting to the previous position. Player II may keep adding tokens, but will eventually run out of them, no matter how many she acquires in between, and then player I can proceed as in ordinary Nim.

\begin{tcolorbox}
玩家 I 应该怎么做呢？稍加思考后，他只需移除玩家 II 刚刚添加到堆中的 50 个代币，恢复到之前的局势。玩家 II 可能会继续添加代币，但无论她中途获得多少代币，最终都会用完，然后玩家 I 就可以像在普通尼姆游戏中那样继续行动。
\end{tcolorbox}

Hence, a player who can win in a position in ordinary Nim can still win in Poker Nim. He replies to the opponent's heap-reducing moves just as he would in ordinary \(\mathrm{{Nim}}\) , and reverses the effect of any heap-increasing move by using a heap-reducing move to restore the heap to the same size again. Strictly speaking, the ending condition is violated in Poker Nim because in theory the game could go on forever. However, a player in a winning position wants to end the game with his victory, and never has to put back any tokens; then the losing player will eventually run out of reserve tokens that she can add to a heap, so that the game terminates.

\begin{tcolorbox}
因此，在普通尼姆游戏中能在某个局势获胜的玩家，在扑克尼姆游戏中仍然可以获胜。他对对手减少堆大小的行动的回应方式与在普通 \(\mathrm{{Nim}}\) 游戏中一样，并且可以通过减少堆大小的行动来抵消对手增加堆大小的行动，使堆恢复到原来的大小。严格来说，扑克尼姆游戏违反了终止条件，因为从理论上讲，游戏可能会永远进行下去。然而，处于必胜局势的玩家希望以胜利结束游戏，并且永远不必放回任何代币；那么失败的玩家最终会用完可以添加到堆中的储备代币，从而使游戏终止。
\end{tcolorbox}

Consider now an impartial game where the options of player I are games that are equivalent to the Nim heaps \(* 0, * 1, * 2, * 5, * 9\) . This can be regarded as a rather peculiar Nim heap of size 3 which can be reduced to any of the sizes0,1,2, but which can also be increased to size 5 or 9 . The Poker Nim argument shows that this extra freedom is of no use to turn a losing game into a winning game.

\begin{tcolorbox}
现在考虑一个无偏博弈，玩家 I 的选项是与尼姆堆 \(* 0, * 1, * 2, * 5, * 9\) 等价的博弈。这可以被看作是一个相当特殊的大小为 3 的尼姆堆，它可以被减小到 0、1、2 中的任何一个大小，但也可以增加到 5 或 9。扑克尼姆游戏的论证表明，这种额外的自由度并不能将一个必败的游戏变成一个必胜的游戏。
\end{tcolorbox}

The mex rule says that if the options of a game \(G\) are equivalent to Nim heaps with sizes from a set \(S\) (like \(S = \{ 0,1,2,5,9\}\) above), then \(G\) is equivalent to a Nim heap of size \(m\) , where \(m\) is the smallest nonnegative integer not contained in \(S\) . This number \(m\) is written \(\operatorname{mex}\left( S\right)\) , where \(m\) ex stands for minimum excluded number. That is,

\begin{tcolorbox}
mex 规则指出，如果一个博弈 \(G\) 的选项等价于大小来自集合 \(S\) 的尼姆堆(就像上面的 \(S = \{ 0,1,2,5,9\}\) )，那么 \(G\) 等价于一个大小为 \(m\) 的尼姆堆，其中 \(m\) 是 \(S\) 中不包含的最小非负整数。这个数字 \(m\) 写作 \(\operatorname{mex}\left( S\right)\) ，其中 \(m\) ex 代表最小排除数。也就是说，
\end{tcolorbox}

\[
m = \operatorname{mex}\left( S\right)  = \min \{ k \geq  0 \mid  k \notin  S\} .
\]

For example, \(\operatorname{mex}\left( {\{ 0,1,2,3,5,6\} }\right)  = 4,\operatorname{mex}\left( {\{ 1,2,3,4,5\} }\right)  = 0\) , and \(\operatorname{mex}\left( \varnothing \right)  = 0\) .

\begin{tcolorbox}
例如， \(\operatorname{mex}\left( {\{ 0,1,2,3,5,6\} }\right)  = 4,\operatorname{mex}\left( {\{ 1,2,3,4,5\} }\right)  = 0\) ，以及 \(\operatorname{mex}\left( \varnothing \right)  = 0\) 。
\end{tcolorbox}

If \(G \equiv   * m\) for an impartial game \(G\) , then \(m\) is called the Nim value of \(G\) . The Nim value is unique because Nim heaps of different sizes are not equivalent by Lemma 1.7. Because the concept has been discovered independently by Sprague (1935) and Grundy (1939), other common names for the Nim value are also Sprague-Grundy value, sg-value, or Grundy value.

\begin{tcolorbox}
如果对于一个无偏博弈 \(G\) 有 \(G \equiv   * m\) ，那么 \(m\) 被称为 \(G\) 的尼姆值(Nim value)。尼姆值是唯一的，因为根据引理 1.7，不同大小的尼姆堆(Nim heap)并不等价。由于这一概念由斯普拉格(Sprague，1935 年)和格鲁迪(Grundy，1939 年)独立发现，尼姆值的其他常见名称还有斯普拉格 - 格鲁迪值(Sprague - Grundy value)、sg 值(sg - value)或格鲁迪值(Grundy value)。
\end{tcolorbox}

Theorem 1.14 (The mex rule). Consider an impartial game G. Then \(G\) has Nim value \(m\) (that is, \(G \equiv   * m\) ), where \(m\) is uniquely determined as follows: For each option \(H\) of \(G\) , let \(H\) have Nim value \({s}_{H}\) , and let \(S = \left\{  {{s}_{H} \mid  H\text{ is an option of }G}\right\}\) . Then \(m = \operatorname{mex}\left( S\right)\) , that is, \(G \equiv   * \left( {\operatorname{mex}\left( S\right) }\right)\) .

\begin{tcolorbox}
定理 1.14(mex 规则)。考虑一个无偏博弈 G。那么 \(G\) 具有尼姆值 \(m\) (即 \(G \equiv   * m\) )，其中 \(m\) 由以下方式唯一确定:对于 \(G\) 的每个选项 \(H\) ，设 \(H\) 具有尼姆值 \({s}_{H}\) ，并设 \(S = \left\{  {{s}_{H} \mid  H\text{ is an option of }G}\right\}\) 。那么 \(m = \operatorname{mex}\left( S\right)\) ，即 \(G \equiv   * \left( {\operatorname{mex}\left( S\right) }\right)\) 。
\end{tcolorbox}

Proof. The statement is true if \(S = \varnothing\) because then \(G = 0 =  * 0\) . Assume, per top-down induction, that each option \(H\) of \(G\) has Nim value \({s}_{H}\) , where \({s}_{H}\) is itself the mex of the Nim values of the options of \(H\) , and that \(S\) is the set of these Nim values \({s}_{H}\) as described. Let \(m = \operatorname{mex}\left( S\right)\) . We show \(G +  * m \equiv  0\) , which proves that \(G \equiv   * m\) by Lemma 1.12. The game \(G +  * m\) is losing if each of its options is winning.

\begin{tcolorbox}
证明。如果 \(S = \varnothing\) ，该陈述为真，因为此时 \(G = 0 =  * 0\) 。假设根据自上而下归纳(top - down induction)， \(G\) 的每个选项 \(H\) 都有尼姆值 \({s}_{H}\) ，其中 \({s}_{H}\) 本身是 \(H\) 的选项的尼姆值的 mex，并且 \(S\) 是如前所述的这些尼姆值 \({s}_{H}\) 的集合。设 \(m = \operatorname{mex}\left( S\right)\) 。我们证明 \(G +  * m \equiv  0\) ，根据引理 1.12，这证明了 \(G \equiv   * m\) 。如果一个博弈 \(G +  * m\) 的每个选项都是必胜局势，那么这个博弈就是失败位置。
\end{tcolorbox}

First, consider such an option of the form \(G +  * h\) for some \(h < m\) , that is, the player moves in the Nim heap \(* m\) of the game sum \(G +  * m\) . Then \(h \in  S\) because \(m = \operatorname{mex}\left( S\right)\) , and by construction of \(S\) there is an option \(H\) of \(G\) so that \(h = {s}_{H}\) , that is, \(H \equiv   * h\) . Hence, the other player has the counter-move from \(G +  * h\) to the losing position \(H +  * h\) , which shows that \(G +  * h\) is indeed winning.

\begin{tcolorbox}
首先，考虑对于某个 \(h < m\) 形式为 \(G +  * h\) 的这样一个选项，即玩家在博弈之和 \(G +  * m\) 的尼姆堆 \(* m\) 中移动。那么 \(h \in  S\) ，因为 \(m = \operatorname{mex}\left( S\right)\) ，并且根据 \(S\) 的构造， \(G\) 有一个选项 \(H\) 使得 \(h = {s}_{H}\) ，即 \(H \equiv   * h\) 。因此，另一个玩家有从 \(G +  * h\) 到失败位置 \(H +  * h\) 的应对策略，这表明 \(G +  * h\) 确实是必胜局势。
\end{tcolorbox}

Second, consider an option of \(G +  * m\) of the form \(H +  * m\) where \(H\) is an option of \(G\) , where \(H \equiv   * {s}_{H}\) and \({s}_{H} \in  S\) , and either \({s}_{H} < m\) or \({s}_{H} > m\) (the case \({s}_{H} = m\) is excluded because \(m \notin  S\) ). If \({s}_{H} < m\) , then the other player counters by moving from \(H +  * m\) to the losing position \(H +  * {s}_{H}\) . If \({s}_{H} > m\) , we use the inductive hypothesis that \({s}_{H}\) is the mex of the Nim values of the options of \(H\) , so one of these Nim values is \(m\) , with \({H}^{\prime } \equiv   * m\) for an option \({H}^{\prime }\) of \(H\) . Then the other player counters by moving from \(H +  * m\) to the losing position \({H}^{\prime } +  * m\) (this is the Poker Nim argument). This shows that \(H +  * m\) is also a winning position. Hence, \(G +  * m\) is a losing game and \(G +  * m \equiv  0\) , as required.

\begin{tcolorbox}
其次，考虑 \(G +  * m\) 的一个形式为 \(H +  * m\) 的选项，其中 \(H\) 是 \(G\) 的一个选项，其中 \(H \equiv   * {s}_{H}\) 且 \({s}_{H} \in  S\) ，并且要么 \({s}_{H} < m\) 要么 \({s}_{H} > m\) (排除 \({s}_{H} = m\) 的情况，因为 \(m \notin  S\) )。如果 \({s}_{H} < m\) ，那么另一个玩家通过从 \(H +  * m\) 移动到失败位置 \(H +  * {s}_{H}\) 来应对。如果 \({s}_{H} > m\) ，我们使用归纳假设，即 \({s}_{H}\) 是 \(H\) 的选项的尼姆值的 mex，所以这些尼姆值之一是 \(m\) ，对于 \(H\) 的一个选项 \({H}^{\prime }\) 有 \({H}^{\prime } \equiv   * m\) 。然后另一个玩家通过从 \(H +  * m\) 移动到失败位置 \({H}^{\prime } +  * m\) 来应对(这是扑克尼姆游戏(Poker Nim)的论证)。这表明 \(H +  * m\) 也是一个必胜局势。因此， \(G +  * m\) 是一个失败的博弈，并且如要求的那样 \(G +  * m \equiv  0\) 。
\end{tcolorbox}

A special case of the preceding theorem is that \(\operatorname{mex}\left( S\right)  = 0\) , which means that all options of \(G\) are equivalent to positive Nim heaps, so they are all winning positions, or that \(G\) has no options at all. Then \(G\) is a losing game, and indeed \(G \equiv   * 0\) .

\begin{tcolorbox}
前面定理的一个特殊情况是 \(\operatorname{mex}\left( S\right)  = 0\) ，这意味着 \(G\) 的所有选项都等价于正的尼姆堆，所以它们都是必胜局势，或者 \(G\) 根本没有选项。那么 \(G\) 是一个失败的博弈，实际上 \(G \equiv   * 0\) 。
\end{tcolorbox}

\(\Rightarrow\) The digraph game in Exercise 1.9 is excellent for understanding the mex rule.

\begin{tcolorbox}
\(\Rightarrow\) 练习 1.9 中的有向图博弈非常适合理解 mex 规则。
\end{tcolorbox}

Note that the mex rule has nothing to do with game sums (apart from the way Theorem 1.14 is proved, which is irrelevant). The mex rule evaluates the options of a game \(G\) , and they are not added when looking at them. If the options are themselves game sums (for example for the game Kayles in (1.22) below or the game in Exercise 1.7), then this is a separate matter that is used to find the Nim value of each option, and is specific to the rules of that game. The mex rule is general.

\begin{tcolorbox}
请注意，mex 规则与博弈之和无关(除了定理 1.14 的证明方式，这并不相关)。mex 规则评估博弈 \(G\) 的选项，并且在考虑这些选项时不会将它们相加。如果这些选项本身就是博弈之和(例如下面 (1.22) 中的凯尔斯博弈或练习 1.7 中的博弈)，那么这是一个单独的问题，用于找出每个选项的尼姆值，并且特定于该博弈的规则。mex 规则是通用的。
\end{tcolorbox}

By Theorem 1.14, any impartial game can be played like Nim, provided the Nim values of the positions of the game are known. The Nim values of the positions can be evaluated recursively by the mex rule. This is illustrated in the Rook-move game in Figure 1.1. Place a rook on a Chess board that extends arbitrarily to the right and bottom but ends at the left and at the top. In one move, the rook is moved either horizontally to the left or vertically upwards, for any number of squares (at least one) as long as it stays on the board. The first player who can no longer move loses, when the rook is on the top left square of the board. We number the rows and columns of the board by \(0,1,2,\ldots\) starting from the top left.

\begin{tcolorbox}
根据定理 1.14，任何无偏博弈都可以像尼姆游戏一样进行，前提是知道该博弈局势的尼姆值。局势的尼姆值可以通过 mex 规则递归地评估。这在图 1.1 的战车移动游戏中得到了说明。将一个战车放在一个国际象棋棋盘上，该棋盘向右和向下任意延伸，但在左边和顶部结束。在一步移动中，战车可以水平向左或垂直向上移动任意数量的方格(至少一个)，只要它留在棋盘上。当战车位于棋盘的左上角方格时，第一个无法再移动的玩家失败。我们从左上角开始用 \(0,1,2,\ldots\) 对棋盘的行和列进行编号。
\end{tcolorbox}

Figure 1.2 gives the Nim values for the positions of the rook on the Chess board. The top left square is equivalent to 0 because the rook can no longer move. The square below only allows reaching the square with 0 on it, so it is equivalent to *1 because \(\operatorname{mex}\left( {\{ 0\} }\right)  = 1\) . The square below gets *2 because its options are equivalent to 0 and *1. From any square in the leftmost column in Figure 1.2, the rook can only move upwards, so any such square in column 0 and row \(i\) corresponds obviously to a Nim heap \(* i\) . Similarly, the topmost row 0 has entry \(* j\) in column \(j\) .

\begin{tcolorbox}
图 1.2 给出了战车在国际象棋棋盘上局势的尼姆值。左上角的方格等价于 0，因为战车无法再移动。下面的方格只能到达值为 0 的方格，所以它等价于 *1，因为 \(\operatorname{mex}\left( {\{ 0\} }\right)  = 1\) 。再下面的方格得到 *2，因为它的选项等价于 0 和 *1。在图 1.2 中最左边的列中，战车只能向上移动，所以第 0 列和第 \(i\) 行的任何这样的方格显然对应于一个尼姆堆 \(* i\) 。类似地，最上面的第 0 行在第 \(j\) 列有条目 \(* j\) 。
\end{tcolorbox}

\begin{center}
\includegraphics[max width=0.5\textwidth]{images/bo_d2e5cmbef24c73b03tn0_29_332_203_730_571_0.jpg}
\end{center}
\hspace*{3em} 

Figure 1.1 Rook-move game, where the player may move the rook on the Chess board in the direction of the arrows.

\begin{tcolorbox}
图 1.1 战车移动游戏，玩家可以在国际象棋棋盘上沿箭头方向移动战车。
\end{tcolorbox}

In general, a position on the board is evaluated knowing all Nim values for the squares to the left and the top of it, which are the options of that position. As an example, consider the square in row 3 and column 2. To the left of that square, the entries \(* 3\) and \(* 2\) are found, and to the top the entries \(* 2, * 3,0\) . So the square itself is equivalent to \(* 1\) because \(\operatorname{mex}\left( {\{ 0,2,3\} }\right)  = 1\) . The square in row 3 and column 5, where the rook is placed in Figure 1.1, gets entry *6 .

\begin{tcolorbox}
一般来说，知道棋盘上某个方格左边和上边所有方格的尼姆值(这些是该局势的选项)，就可以评估该方格的局势。例如，考虑第 3 行和第 2 列的方格。在该方格的左边，找到条目 \(* 3\) 和 \(* 2\) ，在上面找到条目 \(* 2, * 3,0\) 。所以该方格本身等价于 \(* 1\) ，因为 \(\operatorname{mex}\left( {\{ 0,2,3\} }\right)  = 1\) 。在图 1.1 中放置战车的第 3 行和第 5 列的方格得到条目 *6。
\end{tcolorbox}

\begin{center}
\includegraphics[max width=0.6\textwidth]{images/bo_d2e5cmbef24c73b03tn0_29_583_1139_772_772_0.jpg}
\end{center}
\hspace*{3em} 

Figure 1.2 Equivalent Nim heaps \(* n\) for positions of the Rook-move game.

\begin{tcolorbox}
图 1.2 战车移动游戏局势的等价尼姆堆 \(* n\) 。
\end{tcolorbox}

Figure 1.2 shows that the only losing squares on the Chess board are on the diagonal, which are equivalent to the losing game 0 . It is also easy to see directly that the diagonal squares are exactly the losing positions of the Rook-move game: These include the top left square where the rook can no longer move. Whenever the rook is on the diagonal the next move has to move the rook away from the diagonal, which is a winning position because at the next move the rook can always be moved back to the diagonal.

\begin{tcolorbox}
图 1.2 显示，国际象棋棋盘上唯一的失败方格在对角线上，它们等价于失败的博弈 0。也很容易直接看出，对角线上的方格恰好是战车移动游戏的失败位置:这些包括战车无法再移动的左上角方格。每当战车在对角线上时，下一步移动必须将战车移离对角线，这是一个必胜局势，因为下一步总是可以将战车移回到对角线上。
\end{tcolorbox}

The astute reader will have noticed that the Rook-move game is just Nim with two heaps. A rook positioned in row \(i\) and column \(j\) can either reduce \(i\) by moving up, or \(j\) by moving left. So this position is the sum of Nim heaps \(* i +  * j\) . It is a losing position if and only if \(i = j\) , where the rook is on the diagonal.

\begin{tcolorbox}
敏锐的读者会注意到，战车移动游戏只是有两个堆的尼姆游戏。位于第 \(i\) 行和第 \(j\) 列的战车可以通过向上移动减少 \(i\) ，或者通过向左移动减少 \(j\) 。所以这个局势是尼姆堆 \(* i +  * j\) 的和。当且仅当 \(i = j\) (即战车在对角线上)时，这是一个失败位置。
\end{tcolorbox}

In addition to distinguishing winning and losing positions, the Nim values become very useful as soon as we place multiple rooks on the board, assuming that they are allowed to share any square, which corresponds to a sum of Rook-move games. For example, it is not obvious that two rooks, one placed in row 2 and column 0, and another in row 4 and column 6, define a losing position because both squares have Nim value 2 .

\begin{tcolorbox}
除了区分获胜位置和失败位置之外，一旦我们在棋盘上放置多个战车，尼姆值就变得非常有用。假设允许它们共用任何一个方格，这对应于多个战车移动游戏的博弈之和。例如，一个战车放在第2行第0列，另一个放在第4行第6列，这两个战车构成一个失败位置并不明显，因为这两个方格的尼姆值都是2。
\end{tcolorbox}

Theorem 1.14 states that any impartial game is equivalent to a single Nim heap of a certain size. In particular, the sum of two Nim heaps is equivalent to a single Nim heap. Assuming we don't know how to play Nim using the binary system as explained with the example (1.3), so far we only know that the game sum of two Nim heaps is winning if they have different sizes, and losing if they are equal. Figure 1.2 represents the computation of the Nim values of a pair of Nim heaps \(* i +  * j\) .

\begin{tcolorbox}
定理1.14指出，任何无偏博弈都等价于一个特定大小的单一尼姆堆。特别地，两个尼姆堆的博弈之和等价于一个单一尼姆堆。假设我们不知道如何像示例(1.3)中那样使用二进制来玩尼姆游戏，到目前为止，我们只知道两个尼姆堆的博弈之和在它们大小不同时是获胜的，在它们大小相等时是失败的。图1.2展示了一对尼姆堆 \(* i +  * j\) 的尼姆值的计算。
\end{tcolorbox}

There is a certain pattern to these Nim values: Clearly, they are symmetric when reflected at the diagonal because rows and columns can be interchanged. In a single row (or column), all the values are different because for any set \(S\) of integers, \(\operatorname{mex}\left( S\right)\) does not belong to \(S\) , so the value in a square is certainly different from the values to the left of that square. Now consider some power of two, say \({2}^{k}\) , and look at the first \({2}^{k}\) rows and columns, which are each numbered \(0,1,\ldots ,{2}^{k} - 1\) . This quadratic top left part of the table (for example, the top left \(4 \times  4\) squares if \(k = 2\) ) only has the entries \(0,1,\ldots ,{2}^{k} - 1\) in each row and column in some permutation. The role of powers of two for computing the Nim value of a sum of two Nim heaps is made precise in the next section, and will explain the use of the binary representation as in (1.3).

\begin{tcolorbox}
这些尼姆值有一定的规律:显然，当沿对角线反射时它们是对称的，因为行和列可以互换。在单行(或单列)中，所有的值都不同，因为对于任何整数集合 \(S\) ， \(\operatorname{mex}\left( S\right)\) 不属于 \(S\) ，所以一个方格中的值肯定与该方格左侧的值不同。现在考虑2的某个幂，比如 \({2}^{k}\) ，并查看前 \({2}^{k}\) 行和列，它们的编号都是 \(0,1,\ldots ,{2}^{k} - 1\) 。表格的这个左上角方形部分(例如，如果 \(k = 2\) ，则是左上角的 \(4 \times  4\) 个方格)在每行和每列中都只有 \(0,1,\ldots ,{2}^{k} - 1\) 以某种排列形式出现。在下一节中，将精确说明2的幂在计算两个尼姆堆之和的尼姆值时的作用，并解释如(1.3)中那样使用二进制表示的原因。
\end{tcolorbox}

\section*{1.6 Sums of Nim Heaps}

\begin{tcolorbox}
\section*{1.6 尼姆堆的和}
\end{tcolorbox}

In this section, we derive how to compute the Nim value for a general Nim position, which is a sum of different Nim heaps. This will be the Nim sum that we have defined using the binary representation, now cast in the language of game sums and equivalent games, and without assuming the binary representation.

\begin{tcolorbox}
在本节中，我们将推导如何计算一般尼姆局面(即不同尼姆堆的和)的尼姆值。这将是我们用二进制表示法定义的尼姆和，现在用博弈之和和等价博弈的语言来表述，并且不假设使用二进制表示法。
\end{tcolorbox}

For example, we know that \(* 1 +  * 2 +  * 3 \equiv  0\) , so by Lemma 1.12, \(* 1 +  * 2\) is equivalent to *3. In general, however, the sizes of the Nim heaps cannot simply be added to obtain the equivalent Nim heap, because \(* 2 +  * 3\) is also equivalent to \(* 1\) , and \(* 1 +  * 3\) is equivalent to \(* 2\) .

\begin{tcolorbox}
例如，我们知道 \(* 1 +  * 2 +  * 3 \equiv  0\) ，所以根据引理1.12， \(* 1 +  * 2\) 等价于 *3。然而，一般来说，不能简单地将尼姆堆的大小相加来得到等价的尼姆堆，因为 \(* 2 +  * 3\) 也等价于 \(* 1\) ，并且 \(* 1 +  * 3\) 等价于 \(* 2\) 。
\end{tcolorbox}

If \(* k \equiv   * n +  * m\) , then we call \(k\) the Nim sum of \(n\) and \(m\) , written \(k = n \oplus  m\) . The following theorem states that the Nim sum of distinct powers of two is their arithmetic sum. For example, \(1 = {2}^{0}\) and \(2 = {2}^{1}\) , so \(1 \oplus  2 = 1 + 2 = 3\) .

\begin{tcolorbox}
如果 \(* k \equiv   * n +  * m\) ，那么我们称 \(k\) 为 \(n\) 和 \(m\) 的尼姆和，记为 \(k = n \oplus  m\) 。以下定理表明，不同2的幂的尼姆和就是它们的算术和。例如， \(1 = {2}^{0}\) 且 \(2 = {2}^{1}\) ，所以 \(1 \oplus  2 = 1 + 2 = 3\) 。
\end{tcolorbox}

Theorem 1.15. Let \(n \geq  1\) , and \(n = {2}^{a} + {2}^{b} + {2}^{c} + \cdots\) , where \(a > b > c > \cdots  \geq  0\) . Then

\begin{tcolorbox}
定理1.15。设 \(n \geq  1\) ，且 \(n = {2}^{a} + {2}^{b} + {2}^{c} + \cdots\) ，其中 \(a > b > c > \cdots  \geq  0\) 。那么
\end{tcolorbox}

\[
* n \equiv   * \left( {2}^{a}\right)  +  * \left( {2}^{b}\right)  +  * \left( {2}^{c}\right)  + \cdots . \tag{1.19}
\]

We first discuss the implications of this theorem, and then prove it. The expression \(n = {2}^{a} + {2}^{b} + {2}^{c} + \cdots\) is an arithmetic sum of distinct powers of two. Any \(n\) is uniquely given as such a sum. It amounts to the binary representation of \(n\) , which, if \(n < {2}^{a + 1}\) , gives \(n\) as the sum of all powers \({2}^{a},{2}^{a - 1},{2}^{a - 2},\ldots ,{2}^{0}\) where each power of two is multiplied with 0 or 1, the binary digit for the respective position. For example,

\begin{tcolorbox}
我们首先讨论这个定理的含义，然后证明它。表达式 \(n = {2}^{a} + {2}^{b} + {2}^{c} + \cdots\) 是不同 2 的幂的算术和。任何 \(n\) 都唯一地表示为这样的和。这相当于 \(n\) 的二进制表示，如果 \(n < {2}^{a + 1}\) ，则将 \(n\) 表示为所有幂 \({2}^{a},{2}^{a - 1},{2}^{a - 2},\ldots ,{2}^{0}\) 的和，其中每个 2 的幂都乘以 0 或 1，即相应位置的二进制数字。例如，
\end{tcolorbox}

\[
9 = 8 + 1 = 1 \cdot  {2}^{3} + 0 \cdot  {2}^{2} + 0 \cdot  {2}^{1} + 1 \cdot  {2}^{0},
\]

so that 9 in decimal is written as 1001 in binary. Theorem 1.15 uses only the distinct powers of two \({2}^{a},{2}^{b},{2}^{c},\ldots\) that correspond to the digits 1 in the binary representation of \(n\) .

\begin{tcolorbox}
因此十进制的 9 在二进制中写成 1001。定理 1.15 仅使用与 \(n\) 的二进制表示中数字 1 对应的不同 2 的幂 \({2}^{a},{2}^{b},{2}^{c},\ldots\) 。
\end{tcolorbox}

The right-hand side of (1.19) is a game sum. Equation (1.19) states that the single Nim heap \(* n\) is equivalent to a game sum of Nim heaps whose sizes are distinct powers of two. If we visualize the tokens in a Nim heap as dots like in (1.4), then an example of this equation is

\begin{tcolorbox}
(1.19) 式的右边是一个博弈之和。等式 (1.19) 表明，单个尼姆堆 \(* n\) 等价于大小为不同 2 的幂的尼姆堆的博弈之和。如果我们像 (1.4) 那样将尼姆堆中的代币可视化为点，那么这个等式的一个例子是
\end{tcolorbox}

\[
\text{ *14 = ::::::::: ≡ ::::: + :: + : }
\]

for \(n = {14}\) . In addition to \(* n\) , consider a second Nim heap \(* m\) and represent it as its equivalent game sum of several Nim heaps, all of which have a size that is a power of two. Then \(* n +  * m\) is a game sum of many such heaps, where equal heaps cancel out in pairs because a game sum of two identical games is losing and can be omitted. The remaining heap sizes are all distinct powers of two, which can be added to give the size of the single Nim heap \(* k\) that is equivalent to \(* n +  * m\) . As an example, let \(n = 9 = 8 + 1\) and \(m = {14} = 8 + 4 + 2\) . Then \(* n +  * m \equiv   * 8 +  * 1 +  * 8 +  * 4 +  * 2 \equiv   * 4 +  * 2 +  * 1 \equiv   * 7\) , which we can also write as \(9 \oplus  {14} = 7\) . In particular, \(* 9 +  * {14} +  * 7\) is a losing game, which would be very laborious to show without the theorem.

\begin{tcolorbox}
对于 \(n = {14}\) 。除了 \(* n\) 之外，考虑第二个尼姆堆 \(* m\) ，并将其表示为几个尼姆堆的等价博弈之和，所有这些尼姆堆的大小都是 2 的幂。那么 \(* n +  * m\) 是许多这样的堆的博弈之和，其中相等的堆会成对抵消，因为两个相同博弈的博弈之和是必败的，可以省略。剩下的堆的大小都是不同的 2 的幂，将它们相加可以得到与 \(* n +  * m\) 等价的单个尼姆堆 \(* k\) 的大小。例如，设 \(n = 9 = 8 + 1\) 且 \(m = {14} = 8 + 4 + 2\) 。那么 \(* n +  * m \equiv   * 8 +  * 1 +  * 8 +  * 4 +  * 2 \equiv   * 4 +  * 2 +  * 1 \equiv   * 7\) ，我们也可以写成 \(9 \oplus  {14} = 7\) 。特别地， \(* 9 +  * {14} +  * 7\) 是一个必败的博弈，如果没有这个定理，证明起来会非常繁琐。
\end{tcolorbox}

One consequence of Theorem 1.15 is that the Nim sum of two integers never exceeds their arithmetic sum. Moreover, if both integers are less than some power of two, then so is their Nim sum.

\begin{tcolorbox}
定理 1.15 的一个推论是，两个整数的尼姆和永远不会超过它们的算术和。此外，如果两个整数都小于某个 2 的幂，那么它们的尼姆和也小于该 2 的幂。
\end{tcolorbox}

Lemma 1.16. Suppose that Theorem 1.15 holds for any \(n\) with \(n < {2}^{d}\) . Let \(0 \leq  p,q < {2}^{d}\) . Then \(* p +  * q \equiv   * r\) where \(0 \leq  r < {2}^{d}\) , that is, \(r = p \oplus  q < {2}^{d}\) .

\begin{tcolorbox}
引理 1.16。假设定理 1.15 对于任何满足 \(n < {2}^{d}\) 的 \(n\) 都成立。设 \(0 \leq  p,q < {2}^{d}\) 。那么 \(* p +  * q \equiv   * r\) 其中 \(0 \leq  r < {2}^{d}\) ，即 \(r = p \oplus  q < {2}^{d}\) 。
\end{tcolorbox}

Proof. Both \(p\) and \(q\) are sums of distinct powers of two, all smaller than \({2}^{d}\) . By Theorem 1.15, \(r\) is also a sum of such powers of two, where those powers that appear in both \(p\) and \(q\) , say \({2}^{k}\) , cancel out in the game sum because \(* \left( {2}^{k}\right)  +  * \left( {2}^{k}\right)  \equiv  0\) . Hence, the resulting sum \(r\) is also less than \({2}^{d}\) .

\begin{tcolorbox}
证明。 \(p\) 和 \(q\) 都是小于 \({2}^{d}\) 的不同 2 的幂的和。根据定理 1.15， \(r\) 也是这样的 2 的幂的和，其中那些同时出现在 \(p\) 和 \(q\) 中的幂，比如 \({2}^{k}\) ，在博弈之和中会抵消，因为 \(* \left( {2}^{k}\right)  +  * \left( {2}^{k}\right)  \equiv  0\) 。因此，得到的和 \(r\) 也小于 \({2}^{d}\) 。
\end{tcolorbox}

The following proof may be best understood by considering it along with an example, say \(n = 7\) .

\begin{tcolorbox}
通过结合一个例子，比如 \(n = 7\) ，来考虑下面的证明可能会更容易理解。
\end{tcolorbox}

Proof of Theorem 1.15. We proceed by induction. Consider some \(n\) , and assume that the theorem holds for all smaller \(n\) . Let \(n = {2}^{a} + q\) where \(q = {2}^{b} + {2}^{c} + \cdots\) . If \(q = 0\) , the claim holds trivially ( \(n\) is just a single power of two), so let \(q > 0\) . We have \(q < {2}^{a}\) . By inductive assumption, \(* q \equiv   * \left( {2}^{b}\right)  +  * \left( {2}^{c}\right)  + \cdots\) , so all we have to prove is that \(* n \equiv   * \left( {2}^{a}\right)  +  * q\) in order to show (1.19). We show this using Lemma 1.13, that is, by showing that each option of \(* n\) is equivalent to some option of the game \(* \left( {2}^{a}\right)  +  * q\) and vice versa. The options of \(* n\) are \(0, * 1, * 2,\ldots , * \left( {n - 1}\right)\) .

\begin{tcolorbox}
定理 1.15 的证明。我们采用归纳法进行证明。考虑某个 \(n\) ，并假设该定理对所有更小的 \(n\) 都成立。设 \(n = {2}^{a} + q\) ，其中 \(q = {2}^{b} + {2}^{c} + \cdots\) 。如果 \(q = 0\) ，该命题显然成立( \(n\) 只是 2 的单一幂次)，所以设 \(q > 0\) 。我们有 \(q < {2}^{a}\) 。根据归纳假设， \(* q \equiv   * \left( {2}^{b}\right)  +  * \left( {2}^{c}\right)  + \cdots\) ，所以为了证明 (1.19)，我们只需证明 \(* n \equiv   * \left( {2}^{a}\right)  +  * q\) 。我们使用引理 1.13 来证明这一点，即通过证明 \(* n\) 的每个选项都等同于博弈 \(* \left( {2}^{a}\right)  +  * q\) 的某个选项，反之亦然。 \(* n\) 的选项是 \(0, * 1, * 2,\ldots , * \left( {n - 1}\right)\) 。
\end{tcolorbox}

The options of \(* \left( {2}^{a}\right)  +  * q\) are of two kinds, depending on whether the player moves in the Nim heap \(* \left( {2}^{a}\right)\) or \(* q\) . The first kind of options are given by

\begin{tcolorbox}
\(* \left( {2}^{a}\right)  +  * q\) 的选项有两种类型，取决于玩家是在尼姆堆(Nim heap) \(* \left( {2}^{a}\right)\) 还是 \(* q\) 中移动。第一种类型的选项由以下给出
\end{tcolorbox}

\[
0 +  * q \equiv   * {r}_{0}
\]

\[
* 1 +  * q \equiv   * {r}_{1} \tag{1.20}
\]

\[
\vdots
\]

\[
* \left( {{2}^{a} - 1}\right)  +  * q \equiv   * {r}_{{2}^{a} - 1}
\]

where the equivalence of \(* i +  * q\) with some Nim heap \(* {r}_{i}\) , for \(0 \leq  i < {2}^{a}\) , holds by the inductive assumption even if we did not know the mex rule. Moreover, by Lemma 1.16 (with \(d = a\) , because for \(n < {2}^{a}\) Theorem 1.15 holds by the inductive assumption), both \(i\) and \(q\) are less than \({2}^{a}\) , and therefore also \({r}_{i} < {2}^{a}\) . On the right-hand side in (1.20), there are \({2}^{a}\) many Nim heaps \(* {r}_{i}\) for \(0 \leq  i < {2}^{a}\) . We claim they are all different, so that these options form exactly the set \(\{ 0, * 1, * 2,\ldots\) , \(\left. {*\left( {{2}^{a} - 1}\right) }\right\}\) . Namely, by adding the game \(* q\) to the heap \(* {r}_{i}\) ,(1.16) implies \(* {r}_{i} +  * q \equiv\)  \(* i +  * q +  * q \equiv   * i\) , so that \(* {r}_{i} \equiv   * {r}_{j}\) implies \(* {r}_{i} +  * q \equiv   * {r}_{j} +  * q\) , that is, \(* i \equiv   * j\) and hence \(i = j\) by Lemma 1.7, for \(0 \leq  i,j < {2}^{a}\) .

\begin{tcolorbox}
其中，对于 \(0 \leq  i < {2}^{a}\) ， \(* i +  * q\) 与某个尼姆堆(Nim heap) \(* {r}_{i}\) 的等价性，即使我们不知道 mex 规则，根据归纳假设也成立。此外，根据引理 1.16(对于 \(d = a\) ，因为对于 \(n < {2}^{a}\) ，定理 1.15 根据归纳假设成立)， \(i\) 和 \(q\) 都小于 \({2}^{a}\) ，因此 \({r}_{i} < {2}^{a}\) 也小于 \({2}^{a}\) 。在 (1.20) 的右侧，对于 \(0 \leq  i < {2}^{a}\) 有 \({2}^{a}\) 个尼姆堆(Nim heap) \(* {r}_{i}\) 。我们声称它们都是不同的，所以这些选项恰好构成集合 \(\{ 0, * 1, * 2,\ldots\) ， \(\left. {*\left( {{2}^{a} - 1}\right) }\right\}\) 。即，通过将博弈 \(* q\) 添加到堆 \(* {r}_{i}\) 中，(1.16) 意味着 \(* {r}_{i} +  * q \equiv\)  \(* i +  * q +  * q \equiv   * i\) ，所以 \(* {r}_{i} \equiv   * {r}_{j}\) 意味着 \(* {r}_{i} +  * q \equiv   * {r}_{j} +  * q\) ，即 \(* i \equiv   * j\) ，因此根据引理 1.7，对于 \(0 \leq  i,j < {2}^{a}\) 有 \(i = j\) 。
\end{tcolorbox}

The second kind of options of \(* \left( {2}^{a}\right)  +  * q\) are of the form

\begin{tcolorbox}
\(* \left( {2}^{a}\right)  +  * q\) 的第二种类型的选项形式为
\end{tcolorbox}

\[
* \left( {2}^{a}\right)  + 0 \equiv   * \left( {{2}^{a} + 0}\right)
\]

\[
* \left( {2}^{a}\right)  +  * 1 \equiv   * \left( {{2}^{a} + 1}\right)
\]

\[
\vdots
\]

\[
* \left( {2}^{a}\right)  +  * \left( {q - 1}\right)  \equiv   * \left( {{2}^{a} + q - 1}\right) ,
\]

where the heap sizes on the right-hand sides are given again by the inductive assumption. These heaps form the set \(\left\{  {*\left( {2}^{a}\right) , * \left( {{2}^{a} + 1}\right) ,\ldots , * \left( {n - 1}\right) }\right\}\) . Together with the first kind of options, they are exactly the options of \(* n\) . This shows that the options of \(* n\) and of the game sum \(* \left( {2}^{a}\right)  +  * q\) are indeed equivalent, which completes the proof.

\begin{tcolorbox}
其中右侧的堆大小再次由归纳假设给出。这些堆构成集合 \(\left\{  {*\left( {2}^{a}\right) , * \left( {{2}^{a} + 1}\right) ,\ldots , * \left( {n - 1}\right) }\right\}\) 。与第一种类型的选项一起，它们恰好是 \(* n\) 的选项。这表明 \(* n\) 的选项和博弈之和(game sum) \(* \left( {2}^{a}\right)  +  * q\) 的选项确实是等价的，从而完成了证明。
\end{tcolorbox}

In the proof of Theorem 1.15, if \(n = 7\) then \(n = {2}^{a} + q\) with \(a = 2\) and \(q = 3\) . Then the Nim heaps \(* {r}_{0}, * {r}_{1}, * {r}_{2}, * {r}_{3}\) in (1.20) are the entries in the first \({2}^{a}\) rows 0,1,2,3in column \(q = 3\) of the Rook-move table in Figure 1.2.

\begin{tcolorbox}
在定理1.15的证明中，如果 \(n = 7\) ，那么 \(n = {2}^{a} + q\) ，其中 \(a = 2\) 且 \(q = 3\) 。那么(1.20)中的尼姆堆 \(* {r}_{0}, * {r}_{1}, * {r}_{2}, * {r}_{3}\) 是图1.2中战车移动表第 \(q = 3\) 列前 \({2}^{a}\) 行0、1、2、3中的元素。
\end{tcolorbox}

This step in the proof of Theorem 1.15, which is based on Lemma 1.16, also explains to some extent why powers of two appear in the computation of Nim sums. The options of moving from \(* \left( {2}^{a}\right)  +  * q\) in (1.20) neatly produce exactly the numbers \(0,1,\ldots ,{2}^{a} - 1\) , which would not work when replacing \({2}^{a}\) with something else.

\begin{tcolorbox}
定理1.15证明中的这一步基于引理1.16，它在一定程度上也解释了为什么2的幂次会出现在尼姆和的计算中。在(1.20)中从 \(* \left( {2}^{a}\right)  +  * q\) 移动的选项恰好能产生数字 \(0,1,\ldots ,{2}^{a} - 1\) ，如果用其他东西代替 \({2}^{a}\) ，这就行不通了。
\end{tcolorbox}

The Nim sum of any set of numbers can be obtained by writing each number as a sum of distinct powers of two and then canceling repetitions in pairs. For the example in (1.3) and (1.4), we have

\begin{tcolorbox}
任何一组数字的尼姆和可以通过将每个数字写成不同2的幂次之和，然后成对消去重复项来得到。对于(1.3)和(1.4)中的例子，我们有
\end{tcolorbox}

\[
4 \oplus  5 \oplus  9 \oplus  {14} = 4 \oplus  \left( {4 + 1}\right)  \oplus  \left( {8 + 1}\right)  \oplus  \left( {8 + 4 + 2}\right)  = 4 + 2 = 6. \tag{1.21}
\]

Hence, the Nim position4,5,9,14is equivalent to a single Nim heap of size 6 . This equivalence holds because adding this Nim heap to the game sum of the four Nim heaps produces a losing game:

\begin{tcolorbox}
因此，尼姆局势4、5、9、14等价于一个大小为6的单个尼姆堆。这种等价关系成立是因为将这个尼姆堆加到四个尼姆堆的博弈之和中会产生一个失败的博弈:
\end{tcolorbox}

\[
\left( {*4 +  * 5 +  * 9 +  * {14}}\right)  +  * 6 \equiv  0.
\]

The same can be seen from the left table in (1.3). That is, Nim sums amount to writing the heap sizes in binary and adding the binary digits separately in each column (power of 2) modulo 2. The column entry of the Nim sum is 0 if the respective power of 2 appears an even number of times (all pairs cancel) and 1 if it appears an odd number of times.

\begin{tcolorbox}
从(1.3)中的左表也可以看出这一点。也就是说，尼姆和相当于将堆的大小写成二进制形式，并在每一列(2的幂次)中分别对二进制数字进行模2加法。如果相应的2的幂次出现偶数次(所有对都消去)，尼姆和的列元素为0；如果出现奇数次，则为1。
\end{tcolorbox}

\section*{1.7 Finding Nim Values}

\begin{tcolorbox}
\section*{1.7 寻找尼姆值}
\end{tcolorbox}

In this section, we analyze some impartial games using the mex rule in Theorem 1.14.

\begin{tcolorbox}
在本节中，我们使用定理1.14中的mex规则来分析一些无偏博弈。
\end{tcolorbox}

A game similar to the Rook-move game is the Queen-move game shown in Figure 1.3 where the rook is replaced by a Chess queen, which may move horizontally, vertically, and diagonally (left or up). The squares on the main diagonal are therefore no longer losing positions. This game can also be played with two heaps of tokens where in one move, the player may either remove tokens from one heap as in Nim, or reduce both heaps by the same number of tokens (so this is no longer a sum of two Nim heaps!). In order to illustrate that we are not just interested in the winning and losing squares, we add to this game a Nim heap of size 4 .

\begin{tcolorbox}
与战车移动游戏类似的一个博弈是图1.3所示的皇后移动博弈，其中战车被国际象棋中的皇后所取代，皇后可以水平、垂直和对角(向左或向上)移动。因此，主对角线上的方格不再是失败位置。这个博弈也可以用两堆代币来玩，在一次移动中，玩家可以像在尼姆游戏中那样从一堆中移除代币，或者从两堆中移除相同数量的代币(所以这不再是两个尼姆堆的和！)。为了说明我们不只是对胜负方格感兴趣，我们给这个博弈添加一个大小为4的尼姆堆。
\end{tcolorbox}

\begin{center}
\includegraphics[max width=0.8\textwidth]{images/bo_d2e5cmbef24c73b03tn0_34_264_189_1042_405_0.jpg}
\end{center}
\hspace*{3em} 

Figure 1.3 Game sum of a Queen-move game and a Nim heap of size four. The player may either move the queen in the direction of the arrows or take some of the four tokens from the heap.

\begin{tcolorbox}
图1.3 皇后移动博弈和一个大小为4的尼姆堆的博弈之和。玩家可以要么沿箭头方向移动皇后，要么从堆中取走四个代币中的一些。
\end{tcolorbox}

Figure 1.4 shows the equivalent Nim heaps for the positions of the Queen-move game, determined by the mex rule. The square in row 3 and column 4 occupied by the queen in Figure 1.3 has entry *2. So a winning move is to remove two tokens from the Nim heap to turn it into the heap *2, creating the losing position *2 + *2 . Because 2 is the mex of the Nim values of the options of the queen, these may include (as in Poker Nim) higher Nim values. Indeed, the queen can reach two positions equivalent to *4, in row 3 column 1, and row 0 column 4 . If the queen moves there this creates the game sum \(* 4 +  * 4\) which is losing, so these are two further winning moves in Figure 1.3.

\begin{tcolorbox}
图1.4显示了由mex规则确定的皇后移动博弈局势的等价尼姆堆。图1.3中皇后所在的第3行第4列的方格元素为*2。所以一个必胜步是从尼姆堆中取走两个代币，使其变成堆*2，形成失败局势*2 + *2。因为2是皇后选项的尼姆值的mex，这些选项可能包括(如在扑克尼姆游戏中)更高的尼姆值。实际上，皇后可以到达两个等价于*4的局势，即第3行第1列和第0行第4列。如果皇后移动到那里，就会形成博弈和 \(* 4 +  * 4\) ，这是一个失败的博弈，所以在图1.3中这是另外两个必胜步。
\end{tcolorbox}

\begin{center}
\includegraphics[max width=0.4\textwidth]{images/bo_d2e5cmbef24c73b03tn0_34_271_1234_617_396_0.jpg}
\end{center}
\hspace*{3em} 

Figure 1.4 Equivalent Nim heaps for positions of the Queen-move game.

\begin{tcolorbox}
图1.4 皇后移动博弈局势的等价尼姆堆。
\end{tcolorbox}

The impartial game Kayles is played as follows: Given a row of \(n\) bowling pins (numbered \(1,2,\ldots ,n\) ), a move knocks out one or two consecutive pins, as in this example of 5 pins where pins 2 and 3 are knocked out:

\begin{tcolorbox}
无偏博弈凯尔斯博弈的玩法如下:给定一排 \(n\) 个保龄球瓶(编号为 \(1,2,\ldots ,n\) )，一次移动可以击倒一个或两个连续的球瓶，就像这个有5个球瓶的例子，其中第2和第3个球瓶被击倒:
\end{tcolorbox}

(1.22)

\begin{center}
\includegraphics[max width=0.4\textwidth]{images/bo_d2e5cmbef24c73b03tn0_34_545_1824_563_140_0.jpg}
\end{center}
\hspace*{3em} 

If this game is called \({K}_{n}\) , then knocking out a single pin \(p\) creates the game sum \({K}_{p - 1} + {K}_{n - p}\) , and knocking out pins \(p\) and \(p + 1\) (in the picture, \(p = 2\) and \(n = 5\) ) creates the game sum \({K}_{p - 1} + {K}_{n - p - 1}\) (here \({K}_{1} + {K}_{2}\) ). The options of \({K}_{n}\) are these game sums for all possible \(p\) , where \(p\) only needs to be considered up to the middle pin due to the symmetry of the row of pins. Note that in this game, options happen to be sums of games. For the mex rule, only their Nim value is important. The following picture shows all options of Kayles with 1, 2, or 3 pins:

\begin{tcolorbox}
如果这个博弈被称为 \({K}_{n}\) ，那么击倒单个球瓶 \(p\) 会形成博弈之和 \({K}_{p - 1} + {K}_{n - p}\) ，而击倒球瓶 \(p\) 和 \(p + 1\) (在图中是 \(p = 2\) 和 \(n = 5\) )会形成博弈之和 \({K}_{p - 1} + {K}_{n - p - 1}\) (这里是 \({K}_{1} + {K}_{2}\) )。 \({K}_{n}\) 的选项是所有可能的 \(p\) 对应的这些博弈之和，由于球瓶行的对称性， \(p\) 只需要考虑到中间的球瓶即可。注意，在这个博弈中，选项恰好是博弈的和。对于 mex 规则，只有它们的尼姆值是重要的。下图展示了有 1、2 或 3 个球瓶的凯尔斯博弈的所有选项:
\end{tcolorbox}

\begin{center}
\includegraphics[max width=0.4\textwidth]{images/bo_d2e5cmbef24c73b03tn0_35_669_376_582_374_0.jpg}
\end{center}
\hspace*{3em} 

which shows \({K}_{1} \equiv   * \operatorname{mex}\left( 0\right)  =  * 1,{K}_{2} \equiv   * \operatorname{mex}\left( {1,0}\right)  =  * 2\) , and \({K}_{3} \equiv   * \operatorname{mex}\left( {2,1,1 \oplus  1}\right)  =\)  \(* 3\) because \(1 \oplus  1 = 0\) . The options of Kayles with 4 pins are

\begin{tcolorbox}
这表明 \({K}_{1} \equiv   * \operatorname{mex}\left( 0\right)  =  * 1,{K}_{2} \equiv   * \operatorname{mex}\left( {1,0}\right)  =  * 2\) ，并且 \({K}_{3} \equiv   * \operatorname{mex}\left( {2,1,1 \oplus  1}\right)  =\)  \(* 3\) ，因为 \(1 \oplus  1 = 0\) 。有 4 个球瓶的凯尔斯博弈的选项是
\end{tcolorbox}

\begin{center}
\includegraphics[max width=0.7\textwidth]{images/bo_d2e5cmbef24c73b03tn0_35_509_896_905_124_0.jpg}
\end{center}
\hspace*{3em} 

and therefore \({K}_{4} \equiv   * \operatorname{mex}\left( {3,2,1 \oplus  2,1 \oplus  1}\right)  =  * \operatorname{mex}\left( {3,2,3,0}\right)  =  * 1\) which leads to an unexpected winning move for 6 pins:

\begin{tcolorbox}
因此 \({K}_{4} \equiv   * \operatorname{mex}\left( {3,2,1 \oplus  2,1 \oplus  1}\right)  =  * \operatorname{mex}\left( {3,2,3,0}\right)  =  * 1\) ，这导致了有 6 个球瓶时一个意想不到的必胜步:
\end{tcolorbox}

\begin{center}
\includegraphics[max width=0.4\textwidth]{images/bo_d2e5cmbef24c73b03tn0_35_674_1167_580_123_0.jpg}
\end{center}
\hspace*{3em} 

To check this, it is illustrative to show directly that \({K}_{1} + {K}_{4}\) is a losing position, by showing that all its options are winning by finding for each one a counter-move to a losing position.

\begin{tcolorbox}
为了验证这一点，通过为 \({K}_{1} + {K}_{4}\) 的每个选项找到一个走向失败位置的应对策略，直接证明所有选项都是必胜的，从而说明 \({K}_{1} + {K}_{4}\) 是一个失败位置，这是很有启发性的。
\end{tcolorbox}

Further impartial games are described in the exercises.

\begin{tcolorbox}
练习中描述了更多的无偏博弈。
\end{tcolorbox}

\section*{1.8 A Glimpse of Partizan Games}

\begin{tcolorbox}
\section*{1.8 有偏博弈一瞥}
\end{tcolorbox}

A combinatorial game is called partizan if the available moves in a game position may be different for the two players. These games have a rich theory of which we sketch the first concepts here, in particular the four outcome classes, and numbers as strengths of positions in certain games, pioneered by Conway (2001).

\begin{tcolorbox}
如果在一个博弈位置中，两个玩家可用的行动可能不同，那么这个组合博弈就被称为有偏博弈。这些博弈有丰富的理论，我们在这里概述其最初的概念，特别是四种结果类，以及在某些博弈中作为位置强度的数字，这是由康威(2001 年)开创的。
\end{tcolorbox}

The two players are called Left and Right. Consider the game Domineering, given by a board of squares (the starting position is typically rectangular but does not have to be), where in a move Left may occupy two vertically adjacent free squares with a vertical domino, and Right two horizontally adjacent squares with a horizontal domino. (At least in lower case, the letter "l" is more vertical than the letter "r" to remember this.) Hence, starting from a 3-row, 2-column board, we have, up to symmetry,

\begin{tcolorbox}
两个玩家分别称为左玩家和右玩家。考虑纵横棋博弈，它由一个方格棋盘给出(起始位置通常是矩形，但不必如此)，在每一步中，左玩家可以用一个垂直的多米诺骨牌占据两个垂直相邻的空格，右玩家可以用一个水平的多米诺骨牌占据两个水平相邻的空格。(至少小写字母“l”比字母“r”更垂直，这样可以记住这一点。)因此，从一个 3 行 2 列的棋盘开始，在考虑对称性的情况下，我们有
\end{tcolorbox}

(1.23)This game is very different from the impartial game Cram (see Exercise 1.4) where each player may place a domino in either orientation. As (1.23) shows, in Domineering the options of Left and Right are usually different. As before, we assume the normal play convention that a player who can no longer move loses.

\begin{tcolorbox}
这个博弈与无偏博弈塞砖棋(见练习 1.4)非常不同，在塞砖棋中，每个玩家可以以任意方向放置多米诺骨牌。如(1.23)所示，在纵横棋中，左玩家和右玩家的选项通常是不同的。和之前一样，我们假设采用正常规则，即无法再行动的玩家输。
\end{tcolorbox}

\begin{center}
\includegraphics[max width=0.7\textwidth]{images/bo_d2e5cmbef24c73b03tn0_36_310_275_941_125_0.jpg}
\end{center}
\hspace*{3em} 

We always assume optimal play. In (1.23), this means Right would choose the first option and place the domino in the middle, after which Left can no longer move, whereas the second option that leaves a \(2 \times  2\) board would provide Left with two further moves.

\begin{tcolorbox}
我们总是假设玩家采取最优策略。在(1.23)中，这意味着右玩家会选择第一个选项，将多米诺骨牌放在中间，之后左玩家就无法再行动了，而第二个选项留下一个 \(2 \times  2\) 棋盘会给左玩家提供两步额外的行动。
\end{tcolorbox}

An impartial game can only be losing or winning as stated in Lemma 1.1. For partizan games, there are four possible outcome classes, which are denoted by the following calligraphic letters:

\begin{tcolorbox}
如引理 1.1 所述，无偏博弈只有失败或必胜两种情况。对于有偏博弈，有四种可能的结果类，用以下花体字母表示:
\end{tcolorbox}

\(\mathcal{L}\) : Left wins no matter who moves first.

\begin{tcolorbox}
\(\mathcal{L}\) :无论谁先行动，左玩家必胜。
\end{tcolorbox}

\(\mathcal{R}\) : Right wins no matter who moves first.

\begin{tcolorbox}
\(\mathcal{R}\) :无论谁先行动，右玩家必胜。
\end{tcolorbox}

\(\mathcal{P}\) : The first player to move loses, so the previous player wins.

\begin{tcolorbox}
\(\mathcal{P}\) :先行动的玩家失败，因此前一个玩家获胜。
\end{tcolorbox}

\(\mathcal{N}\) : The first player to move wins. (Sometimes called the "next player", although this is as ambiguous as "next Friday", because it is the current player who wins.)

\begin{tcolorbox}
\(\mathcal{N}\) :先行动的玩家获胜。(有时称为“下一玩家”，尽管这和“下周五”一样模棱两可，因为实际上是当前玩家获胜。)
\end{tcolorbox}

Every game belongs to exactly one of these outcome classes. The \(\mathcal{P}\) -positions are what we have called losing positions, and \(\mathcal{N}\) -positions are what we have called winning positions. For partizan games we have to consider the further outcome classes \(\mathcal{L}\) and \(\mathcal{R}\) . In Domineering, a single vertical strip of at least two squares, in its smallest form -, belongs to \(\mathcal{L}\) , whereas a horizontal strip such as belongs to \(\mathcal{R}\) . The starting \(3 \times  2\) board in (1.23) belongs to \(\mathcal{N}\) .

\begin{tcolorbox}
每个博弈都恰好属于这些结果类中的某一类。 \(\mathcal{P}\) 局势就是我们所说的失败位置，而 \(\mathcal{N}\) 局势就是我们所说的必胜局势。对于有偏博弈(partizan game)，我们还需要考虑另外的结果类 \(\mathcal{L}\) 和 \(\mathcal{R}\) 。在纵横棋(Domineering)中，一个至少有两个方格的单个垂直条，其最小形式为 - ，属于 \(\mathcal{L}\) 类，而像这样的水平条则属于 \(\mathcal{R}\) 类。(1.23) 中的起始 \(3 \times  2\) 棋盘属于 \(\mathcal{N}\) 类。
\end{tcolorbox}

The first move of Right in (1.23) creates two disjoint \(1 \times  2\) games which represent, as before, a game sum, here \(\rightarrow\) the game to the four outcome classes, games \(G\) and \(H\) are called equivalent, written \(G \equiv  H\) , if and only if \(G + J\) and \(H + J\) always belong to the same outcome class, for any other game \(J\) ; this extends Definition 1.5. Every losing game (any game in \(\mathcal{P}\) ) is equivalent to the zero game 0 that has no options, as in Lemma 1.8, by the same argument. As before, Lemma 1.10 holds, that is, adding a losing game does not change the outcome class of a game.

\begin{tcolorbox}
在 (1.23) 中，Right 的第一步行动创造了两个不相交的 \(1 \times  2\) 博弈，和之前一样，这代表了一个博弈之和(game sum)，这里 \(\rightarrow\) 将博弈划分为四个结果类，如果且仅当对于任何其他博弈 \(J\) ， \(G + J\) 和 \(H + J\) 总是属于同一结果类时，博弈 \(G\) 和 \(H\) 被称为等价的，记为 \(G \equiv  H\) ；这扩展了定义 1.5。每个失败博弈( \(\mathcal{P}\) 中的任何博弈)都等价于没有选项的零博弈 0，就像引理 1.8 中那样，依据相同的论证。和之前一样，引理 1.10 成立，即添加一个失败博弈不会改变一个博弈的结果类。
\end{tcolorbox}

For any partizan game \(G\) , there is a negative game \(- G\) so that \(G + \left( {-G}\right)  \equiv  0\) as stated in (1.13). This negative game is obtained by exchanging the options of Left and Right, and doing so recursively for their options. For Domineering, the negative of any game position \(G\) (as shown by the squares that have not yet been occupied by dominos) is obtained by turning the board by 90 degrees. For example,

\begin{tcolorbox}
对于任何有偏博弈 \(G\) ，存在一个负博弈 \(- G\) ，使得如 (1.13) 中所述的 \(G + \left( {-G}\right)  \equiv  0\) 成立。这个负博弈是通过交换 Left 和 Right 的选项，并对它们的选项进行递归交换得到的。对于纵横棋(Domineering)，任何博弈位置 \(G\) 的负博弈(由尚未被多米诺骨牌(domino)占据的方格表示)是通过将棋盘旋转 90 度得到的。例如，
\end{tcolorbox}

(1.24)

\begin{center}
\includegraphics[max width=0.4\textwidth]{images/bo_d2e5cmbef24c73b03tn0_37_712_349_510_153_0.jpg}
\end{center}
\hspace*{3em} 

because any move by a player in one component can be copied by the other player by the corresponding move in the other component. In that way, the second player always has a move left and has won when the position reached is \(0 + 0\) .

\begin{tcolorbox}
因为一个组件中玩家的任何行动都可以被另一个玩家在另一个组件中进行相应的行动所复制。通过这种方式，第二个玩家总是有行动可做，并且当达到 \(0 + 0\) 局势时获胜。
\end{tcolorbox}

Formally, a game \(G\) is defined by two sets of games \({\mathcal{G}}^{L}\) and \({\mathcal{G}}^{R}\) which define the options of Left and Right, respectively. Then \(G\) is written as \(\left\{  {{\mathcal{G}}^{L} \mid  {\mathcal{G}}^{R}}\right\}\) . In concrete examples, one lists directly the elements of the sets \({\mathcal{G}}^{L}\) and \({\mathcal{G}}^{R}\) rather than using set notation. For example, if \({\mathcal{G}}^{L} = {\mathcal{G}}^{R} = \varnothing\) , then we write this game as \(\{ \left| {\} \text{ rather than }\{ \varnothing }\right| \varnothing \}\) . This is the simplest possible game that has no options for either player, which we have denoted by 0 :

\begin{tcolorbox}
形式上，一个博弈 \(G\) 由两组博弈 \({\mathcal{G}}^{L}\) 和 \({\mathcal{G}}^{R}\) 定义，它们分别定义了 Left 和 Right 的选项。那么 \(G\) 可以写成 \(\left\{  {{\mathcal{G}}^{L} \mid  {\mathcal{G}}^{R}}\right\}\) 。在具体例子中，人们直接列出集合 \({\mathcal{G}}^{L}\) 和 \({\mathcal{G}}^{R}\) 的元素，而不是使用集合符号。例如，如果 \({\mathcal{G}}^{L} = {\mathcal{G}}^{R} = \varnothing\) ，那么我们将这个博弈写成 \(\{ \left| {\} \text{ rather than }\{ \varnothing }\right| \varnothing \}\) 。这是可能的最简单的博弈，任何玩家都没有选项，我们用 0 来表示:
\end{tcolorbox}

\[
0 = \{  \mid  \} \text{ . } \tag{1.25}
\]

The options in a game are defined recursively. We can write (1.23) as

\begin{tcolorbox}
博弈中的选项是递归定义的。我们可以将 (1.23) 写成
\end{tcolorbox}

(1.26)

\begin{center}
\includegraphics[max width=0.5\textwidth]{images/bo_d2e5cmbef24c73b03tn0_37_583_1122_755_129_0.jpg}
\end{center}
\hspace*{3em} 

and continue to define

\begin{tcolorbox}
并继续定义
\end{tcolorbox}

(1.27)

\begin{center}
\includegraphics[max width=0.4\textwidth]{images/bo_d2e5cmbef24c73b03tn0_37_684_1328_556_127_0.jpg}
\end{center}
\hspace*{3em} 

where the first option of Left is to put the domino in the top left which leaves Right with no further move, the second (inferior) option to put the domino at the bottom; Right has only one choice. Furthermore,

\begin{tcolorbox}
其中 Left 的第一个选项是将多米诺骨牌(domino)放在左上角，这使得 Right 没有进一步的行动，第二个(较差的)选项是将多米诺骨牌放在底部；Right 只有一种选择。此外，
\end{tcolorbox}

\[
\square  = \left\{  {0|}\right\}  ,\;\square \square  = \left\{  {0|0}\right\}  , \tag{1.28}
\]

which are the positions where Left respectively Right have exactly one move left and the other player none. The first of Right's options in (1.26) can also be defined by the recursive definition of game sums as

\begin{tcolorbox}
这些是 Left 或 Right 分别只剩下一步行动，而另一个玩家没有行动的局势。(1.26) 中 Right 的第一个选项也可以通过博弈之和的递归定义来定义为
\end{tcolorbox}

\[
\square \square  + \square \square  = \left\{  {\square \square }\right\}  ,
\]

and the second of Right's options in (1.26) is defined by

\begin{tcolorbox}
(1.26) 中 Right 的第二个选项定义为
\end{tcolorbox}

(1.29)The definitions (1.25) to (1.29) are recursive by referring to simpler games, with the shown board configurations as suitable placeholders. It only matters which player eventually loses in subsequent optimal play.

\begin{tcolorbox}
定义 (1.25) 到 (1.29) 通过引用更简单的博弈进行递归定义，所示的棋盘配置作为合适的占位符。重要的是在后续的最优玩法中哪个玩家最终失败。
\end{tcolorbox}

\[
\text{ 田 = \{B|□ \textbackslash }.}
\]

In (1.27), Left has the dominated option [1] that Left should never move to because it unnecessarily gives Right an extra move. Dominated options can be omitted, which simplifies the recursive description of a game in terms of its

\begin{tcolorbox}
在(1.27)中，先手玩家(左方)有一个被占优选项[1]，先手玩家永远不应该选择这个选项，因为这会不必要地让后手玩家(右方)多走一步。被占优选项可以被省略，这能简化博弈根据其选项进行的递归描述。
\end{tcolorbox}

options. Note: there are no dominated options in the board configuration 1 in

\begin{tcolorbox}
注意:在无偏博弈塞砖棋(Cram)的棋盘布局1中，没有被占优选项。
\end{tcolorbox}

the impartial game Cram. In Cram, this configuration is equivalent to a Nim heap \(* 2\) of size two with its two options \(* 1\) and 0, neither of which dominates the other. For example, in the game sum \(* 2 +  * 1\) , the move to \(* 1 +  * 1\) is a winning move which requires the option *1 of *2. Impartial games are exactly the games \(\left\{  {{\mathcal{G}}^{L} \mid  {\mathcal{G}}^{R}}\right\}\) where \({\mathcal{G}}^{L} = {\mathcal{G}}^{R}\) , and the option sets themselves contain only impartial games. The simplest impartial game is 0 by (1.25), the next simplest is \(\{ 0 \mid  0\}\) which is the Nim heap *1 of size one, which is also denoted by *(just a star).

\begin{tcolorbox}
在塞砖棋(Cram)中，这种布局等价于一个大小为2的尼姆堆(Nim heap) \(* 2\) ，它有两个选项 \(* 1\) 和0，且这两个选项彼此都不占优。例如，在博弈之和 \(* 2 +  * 1\) 中，移动到 \(* 1 +  * 1\) 是一个必胜步(winning move)，这需要用到*2的选项*1。无偏博弈(impartial game)恰好是那些满足 \(\left\{  {{\mathcal{G}}^{L} \mid  {\mathcal{G}}^{R}}\right\}\) (其中 \({\mathcal{G}}^{L} = {\mathcal{G}}^{R}\) )的博弈，并且选项集本身只包含无偏博弈。根据(1.25)，最简单的无偏博弈是0，次简单的是 \(\{ 0 \mid  0\}\) ，它是大小为1的尼姆堆*1，也可以用*(仅一个星号)表示。
\end{tcolorbox}

An important new concept for partizan games is that of a number, which applies to some (but not all) partizan games. If a game is identified with a number, then this number states how many moves player Left is safely ahead in the game. The simplest numbers are integers, but they can also be certain types of fractions. For example, the Domineering position - can be identified with the integer 1 . Furthermore, \(1 + 1 = 2\) , that is, the game sum of numbers corresponds to the arithmetic sum of these numbers. Similarly, the Domineering position identified with -1, in agreement with \(1 + \left( {-1}\right)  \equiv  0\) . For a positive integer \(n\) , if Left is \(n\) moves ahead and makes one of these moves, then Left is still \(n - 1\) moves ahead afterwards. By denoting the game itself by \(n\) , this gives rise to the recursive definition

\begin{tcolorbox}
对于有偏博弈(partizan game)，一个重要的新概念是数的概念，它适用于某些(但不是所有)有偏博弈。如果一个博弈可以用一个数来表示，那么这个数表示先手玩家(左方)在博弈中安全领先的步数。最简单的数是整数，但它们也可以是某些类型的分数。例如，纵横棋(Domineering)的局面 - 可以用整数1表示。此外， \(1 + 1 = 2\) ，即数的博弈之和对应于这些数的算术和。类似地，纵横棋(Domineering)中对应于 -1的局面，与 \(1 + \left( {-1}\right)  \equiv  0\) 一致。对于正整数 \(n\) ，如果先手玩家(左方)领先 \(n\) 步并走了其中一步，那么之后先手玩家(左方)仍然领先 \(n - 1\) 步。通过用 \(n\) 表示博弈本身，这就产生了递归定义
\end{tcolorbox}

\[
n = \{ n - 1 \mid  \} , \tag{1.30}
\]

which has its base case \(1 = \{ 0 \mid  \}\) as in (1.28). Games that are negative numbers are defined as negative games \(- n\) of positive numbers \(n\) . In the game \(- n\) , player Right is ahead by \(n\) moves.

\begin{tcolorbox}
它的基础情况是 \(1 = \{ 0 \mid  \}\) ，如(1.28)所示。负整数对应的博弈被定义为正整数对应博弈的负博弈(negative game) \(- n\) 。在博弈 \(- n\) 中，后手玩家(右方)领先 \(n\) 步。
\end{tcolorbox}

Caution: No impartial game is a number (except for the game 0 ). Nim heaps as defined in (1.18) are very different, and obey different rules such as \(* n +  * n \equiv  0\) . (Nim heaps are sometimes called nimbers to distinguish them from numbers.)

\begin{tcolorbox}
注意:除了博弈0之外，没有无偏博弈可以用数来表示。(1.18)中定义的尼姆堆(Nim heap)非常不同，并且遵循不同的规则，如 \(* n +  * n \equiv  0\) 。(尼姆堆有时被称为尼姆数(nimber)，以区别于数。)
\end{tcolorbox}

In order to define a number more generally, we need a partial order \(\geq\) on games (which is not the "simpler than" order), defined as follows:

\begin{tcolorbox}
为了更一般地定义数，我们需要在博弈上定义一个偏序(partial order) \(\geq\) (它不是“比……更简单”的顺序)，定义如下:
\end{tcolorbox}

\[
G \geq  H \Leftrightarrow  \left( {\forall J : H + J \in  \mathcal{L} \cup  \mathcal{N} \Rightarrow  G + J \in  \mathcal{L} \cup  \mathcal{N},}\right.  \tag{1.31}
\]

\[
\forall J : H + J \in  \mathcal{L} \cup  \mathcal{P} \Rightarrow  G + J \in  \mathcal{L} \cup  \mathcal{P})
\]

where a game belongs to \(\mathcal{L} \cup  \mathcal{N}\) if Left wins by moving first, and to \(\mathcal{L} \cup  \mathcal{P}\) if Left wins by moving second. With this understanding, (1.31) can be written as

\begin{tcolorbox}
其中，如果先手玩家(左方)先走能获胜，则一个博弈属于 \(\mathcal{L} \cup  \mathcal{N}\) ；如果先手玩家(左方)后走能获胜，则属于 \(\mathcal{L} \cup  \mathcal{P}\) 。有了这个理解，(1.31)可以写成
\end{tcolorbox}

\[
G \geq  H \Leftrightarrow  \left( {\forall J : \text{ Left wins in }H + J \Rightarrow  \text{ Leftwins in }G + J}\right) .
\]

That is, \(G \geq  H\) if replacing \(H\) by \(G\) never hurts Left, no matter what the context (always assuming optimal play). Then \(\geq\) is a partial order on games because it is obviously reflexive and transitive, and antisymmetric up to equivalence of games:

\begin{tcolorbox}
也就是说，如果无论在什么情况下(始终假设是最优玩法)，用 \(G\) 替换 \(H\) 都不会对先手玩家(左方)造成不利，那么就有 \(G \geq  H\) 。那么 \(\geq\) 是博弈上的一个偏序，因为它显然是自反的(reflexive)和传递的(transitive)，并且在博弈等价的意义下是反对称的(antisymmetric):
\end{tcolorbox}

Lemma 1.17. \(G \geq  H\) and \(H \geq  G \Leftrightarrow  G \equiv  H\) .

\begin{tcolorbox}
引理1.17. \(G \geq  H\) 且 \(H \geq  G \Leftrightarrow  G \equiv  H\) 。
\end{tcolorbox}

Proof. Recall that \(G \equiv  H\) means that \(G + J\) and \(H + J\) belong to the same outcome class \(\mathcal{L},\mathcal{R},\mathcal{P},\mathcal{N}\) for any game \(J\) , so " \(\Leftarrow\) " is obvious. To show " \(\Rightarrow\) ", assume \(G \geq  H\) and \(H \geq  G\) , and let \(J\) be any game. If \(H + J \in  \mathcal{L}\) then \(G + J \in  \left( {\mathcal{L} \cup  \mathcal{N}}\right)  \cap  \left( {\mathcal{L} \cup  \mathcal{P}}\right)  = \mathcal{L}\) by (1.31), because the outcome classes are disjoint. By symmetry, if \(G + J \in  \mathcal{L}\) then \(H + J \in  \mathcal{L}\) . Next, consider the outcome class \(\mathcal{P}\) : If \(H + J \in  \mathcal{P}\) then \(G + J \in  \mathcal{L} \cup  \mathcal{P}\) by (1.31). If we had \(G + J \in  \mathcal{L}\) then \(H + J \in  \mathcal{L}\) as just shown, but \(H + J \in  \mathcal{P}\) which is disjoint from \(\mathcal{L}\) , so we have \(G + J \in  \mathcal{P}\) . By symmetry, if \(G + J \in  \mathcal{P}\) then \(H + J \in  \mathcal{P}\) . Similar arguments apply for the outcome classes \(\mathcal{R}\) and \(\mathcal{N}\) .

\begin{tcolorbox}
证明。回想一下， \(G \equiv  H\) 意味着对于任何博弈 \(J\) ， \(G + J\) 和 \(H + J\) 属于同一个结果类 \(\mathcal{L},\mathcal{R},\mathcal{P},\mathcal{N}\) ，所以 “ \(\Leftarrow\) ” 是显然的。为了证明 “ \(\Rightarrow\) ”，假设 \(G \geq  H\) 且 \(H \geq  G\) ，并设 \(J\) 为任意博弈。如果 \(H + J \in  \mathcal{L}\) ，那么由 (1.31) 可得 \(G + J \in  \left( {\mathcal{L} \cup  \mathcal{N}}\right)  \cap  \left( {\mathcal{L} \cup  \mathcal{P}}\right)  = \mathcal{L}\) ，因为结果类是不相交的。由对称性可知，如果 \(G + J \in  \mathcal{L}\) ，那么 \(H + J \in  \mathcal{L}\) 。接下来，考虑结果类 \(\mathcal{P}\) :如果 \(H + J \in  \mathcal{P}\) ，那么由 (1.31) 可得 \(G + J \in  \mathcal{L} \cup  \mathcal{P}\) 。如果 \(G + J \in  \mathcal{L}\) ，那么正如刚才所证 \(H + J \in  \mathcal{L}\) ，但 \(H + J \in  \mathcal{P}\) 与 \(\mathcal{L}\) 不相交，所以 \(G + J \in  \mathcal{P}\) 。由对称性可知，如果 \(G + J \in  \mathcal{P}\) ，那么 \(H + J \in  \mathcal{P}\) 。类似的论证适用于结果类 \(\mathcal{R}\) 和 \(\mathcal{N}\) 。
\end{tcolorbox}

Define \(G < H\) as \(H \geq  G\) and not \(G \equiv  H\) . This allows the following recursive definition of a partizan game that is a number.

\begin{tcolorbox}
将 \(G < H\) 定义为 \(H \geq  G\) 且非 \(G \equiv  H\) 。这允许对作为数字的有偏博弈进行如下递归的定义。
\end{tcolorbox}

Definition 1.18. A game \(G = \left\{  {{\mathcal{G}}^{L} \mid  {\mathcal{G}}^{R}}\right\}\) is a number if all left options \({G}^{L} \in  {\mathcal{G}}^{L}\) of \(G\) and all right options \({G}^{R} \in  {\mathcal{G}}^{R}\) of \(G\) are numbers and fulfill \({G}^{L} < G < {G}^{R}\) .

\begin{tcolorbox}
定义 1.18。一个博弈 \(G = \left\{  {{\mathcal{G}}^{L} \mid  {\mathcal{G}}^{R}}\right\}\) 是一个数字，如果 \(G\) 的所有左选项 \({G}^{L} \in  {\mathcal{G}}^{L}\) 和 \(G\) 的所有右选项 \({G}^{R} \in  {\mathcal{G}}^{R}\) 都是数字，并且满足 \({G}^{L} < G < {G}^{R}\) 。
\end{tcolorbox}

The simplest number is \(0 = \{  \mid  \}\) because it has no left or right options, so the condition \({G}^{L} < G < {G}^{R}\) for \(G = 0\) holds vacuously. Similarly, the game defined in (1.30) for a positive integer \(n\) is a number because its left option fulfills \(n - 1 < n\) , and it has no right options.

\begin{tcolorbox}
最简单的数字是 \(0 = \{  \mid  \}\) ，因为它没有左选项或右选项，所以对于 \(G = 0\) 的条件 \({G}^{L} < G < {G}^{R}\) 显然成立。类似地，对于正整数 \(n\) 在 (1.30) 中定义的博弈是一个数字，因为它的左选项满足 \(n - 1 < n\) ，并且它没有右选项。
\end{tcolorbox}

The intuition of the inequalities \({G}^{L} < G < {G}^{R}\) in Definition 1.18 is that by moving in a number, a player gives up a safe advantage and thereby weakens his position, so moving to a left option \({G}^{L}\) weakens the position of Left ( \({G}^{L} < G\) ), and moving to a right option \({G}^{R}\) weakens the position of Right \(\left( {G < {G}^{R}}\right)\) .

\begin{tcolorbox}
定义 1.18 中不等式 \({G}^{L} < G < {G}^{R}\) 的直观理解是，在一个数字博弈中移动时，玩家放弃了安全优势，从而削弱了自己的局面，因此移动到左选项 \({G}^{L}\) 会削弱左方( \({G}^{L} < G\) )的局面，而移动到右选项 \({G}^{R}\) 会削弱右方 \(\left( {G < {G}^{R}}\right)\) 的局面。
\end{tcolorbox}

The advantage that is given up by moving in a number can be small, because games can also be numbers that are fractions. For example, the Domineering

\begin{tcolorbox}
在一个数字博弈中移动所放弃的优势可能很小，因为博弈也可以是分数形式的数字。例如，纵横棋
\end{tcolorbox}

position 1 in (1.27) is given by \(\{ 0, - 1 \mid  1\}\) (or, because -1 is dominated, \(\{ 0 \mid  1\}\) ),

\begin{tcolorbox}
(1.27) 中的局面 1 由 \(\{ 0, - 1 \mid  1\}\) 给出(或者，因为 -1 被占优，由 \(\{ 0 \mid  1\}\) 给出)。
\end{tcolorbox}

where \(0 = \{ \left| \right| \}\) and \(1 = \{ 0 \mid  \}\) according to (1.30), and is therefore a number by by Definition 1.18. That number is the fraction \(\frac{1}{2}\) , that is, Left is half a move ahead in this game, because \(\frac{1}{2} + \frac{1}{2} + \left( {-1}\right)  \equiv  0\) , that is,

\begin{tcolorbox}
其中 \(0 = \{ \left| \right| \}\) 和 \(1 = \{ 0 \mid  \}\) 根据 (1.30)，因此根据定义 1.18 它是一个数。那个数是分数 \(\frac{1}{2}\) ，也就是说，在这个博弈中先手玩家(左玩家)领先半手，因为 \(\frac{1}{2} + \frac{1}{2} + \left( {-1}\right)  \equiv  0\) ，即
\end{tcolorbox}

(1.32)

\begin{center}
\includegraphics[max width=0.3\textwidth]{images/bo_d2e5cmbef24c73b03tn0_39_738_1874_453_124_0.jpg}
\end{center}
\hspace*{3em} 

This can be inferred from the (essentially unique) sequence of optimal moves

\begin{tcolorbox}
这可以从(本质上唯一的)最优走法序列推断出来
\end{tcolorbox}

\begin{center}
\includegraphics[max width=0.7\textwidth]{images/bo_d2e5cmbef24c73b03tn0_39_519_2054_943_129_0.jpg}
\end{center}
\hspace*{3em} 

\begin{center}
\includegraphics[max width=0.7\textwidth]{images/bo_d2e5cmbef24c73b03tn0_40_377_199_944_127_0.jpg}
\end{center}
\hspace*{3em} 

which shows that Left as starting player loses, and similarly that Right as starting player loses.

\begin{tcolorbox}
这表明左玩家作为起始玩家会失败，同样，右玩家作为起始玩家也会失败。
\end{tcolorbox}

In terms of recursively defined games, starting from the definition of \(0 = \{  \mid  \}\) , (1.32) states

\begin{tcolorbox}
就递归定义的博弈而言，从 \(0 = \{  \mid  \}\) 的定义开始，(1.32) 表明
\end{tcolorbox}

\[
\{ 0 \mid  1\}  + \{ 0 \mid  1\}  + \left( {-1}\right)  \equiv  0. \tag{1.33}
\]

This holds for the same reason as for the Domineering moves in (1.32), just a bit more abstract. Recall that \(- 1 = \{  \mid  0\}\) (which is one safe move for Right). In the game sum \(\{ 0 \mid  1\}  + \{ 0 \mid  1\}  + \{  \mid  0\}\) in (1.33), Leftc an only move in the component \(\{ 0 \mid  1\}\) , so that the next game is \(0 + \{ 0 \mid  1\}  + \{  \mid  0\}  = \{ 0 \mid  1\}  + \{  \mid  0\}\) . Here Right has a choice, moving either in the first component to \(1 + \{  \mid  0\}\) , which is \(1 + \left( {-1}\right)  = 0\) and a loss for Left who moves next, or in the second component (using his safe move) to \(\{ 0 \mid  1\}  + 0 = \{ 0 \mid  1\}\) , which is a win for Left. Clearly, Right should choose the former. Hence, Left as starting player loses in \(\{ 0 \mid  1\}  + \{ 0 \mid  1\}  + \{ 0 \mid  1\}\) . By a very similar reasoning, Right as starting player loses, which proves (1.33).

\begin{tcolorbox}
这与 (1.32) 中纵横棋(Domineering)走法的原因相同，只是更抽象一些。回想一下 \(- 1 = \{  \mid  0\}\) (这对右玩家来说是一步安全走法)。在 (1.33) 的博弈之和 \(\{ 0 \mid  1\}  + \{ 0 \mid  1\}  + \{  \mid  0\}\) 中，左玩家只能在组件 \(\{ 0 \mid  1\}\) 中走棋，所以下一个博弈状态是 \(0 + \{ 0 \mid  1\}  + \{  \mid  0\}  = \{ 0 \mid  1\}  + \{  \mid  0\}\) 。在这里，右玩家有一个选择，要么在第一个组件中走到 \(1 + \{  \mid  0\}\) ，这是 \(1 + \left( {-1}\right)  = 0\) 且对接下来走棋的左玩家来说是失败局面，要么在第二个组件中(使用他的安全走法)走到 \(\{ 0 \mid  1\}  + 0 = \{ 0 \mid  1\}\) ，这对左玩家来说是胜利局面。显然，右玩家应该选择前者。因此，左玩家作为起始玩家在 \(\{ 0 \mid  1\}  + \{ 0 \mid  1\}  + \{ 0 \mid  1\}\) 中失败。通过非常类似的推理，右玩家作为起始玩家也会失败，这就证明了 (1.33)。
\end{tcolorbox}

As an extension of this argument, there is a game that can be identified with the fraction \(\frac{1}{{2}^{p}}\) for every positive integer \(p\) , defined recursively as follows:

\begin{tcolorbox}
作为这个论证的扩展，对于每个正整数 \(p\) ，都有一个可以与分数 \(\frac{1}{{2}^{p}}\) 对应的博弈，其递归定义如下:
\end{tcolorbox}

\[
\frac{1}{{2}^{p}} = \left\{  {0\left| {\;\frac{1}{{2}^{p - 1}}}\right. }\right\}  \;\text{ and hence }\;\frac{-1}{{2}^{p}} = \left\{  {\left. \frac{-1}{{2}^{p - 1}}\right| \;0}\right\}  . \tag{1.34}
\]

In order to show

\begin{tcolorbox}
为了证明
\end{tcolorbox}

\[
\frac{1}{{2}^{p}} + \frac{1}{{2}^{p}} + \frac{-1}{{2}^{p - 1}} \equiv  0, \tag{1.35}
\]

one can see that in this game sum it is better for both players to move in the component \(\frac{1}{{2}^{p}}\) rather than in \(\frac{-1}{{2}^{p - 1}}\) where the player gives up more, similar to the reasoning for (1.33). Hence, no matter who starts, after two moves the position (1.35) becomes the losing position \(\frac{1}{{2}^{p - 1}} + \frac{-1}{{2}^{p - 1}}\) , which shows (1.35).

\begin{tcolorbox}
可以发现，在这个博弈之和中，对于两个玩家来说，在组件 \(\frac{1}{{2}^{p}}\) 中走棋比在 \(\frac{-1}{{2}^{p - 1}}\) 中走棋更好，因为在 \(\frac{-1}{{2}^{p - 1}}\) 中玩家会放弃更多，这与 (1.33) 的推理类似。因此，无论谁先开始，两步之后局面 (1.35) 就会变成失败位置 \(\frac{1}{{2}^{p - 1}} + \frac{-1}{{2}^{p - 1}}\) ，这就证明了 (1.35)。
\end{tcolorbox}

The following is the more general, recursive definition of a game that is a fraction with a power of two in its denominator.

\begin{tcolorbox}
以下是一个更一般的、递归定义的博弈，其值是分母为 2 的幂的分数。
\end{tcolorbox}

Definition 1.19. Let \(p\) be a positive integer and \(m\) be an odd integer. Define the game \(\frac{m}{{2}^{p}}\) by

\begin{tcolorbox}
定义 1.19。设 \(p\) 是一个正整数， \(m\) 是一个奇数。通过以下方式定义博弈 \(\frac{m}{{2}^{p}}\)
\end{tcolorbox}

\[
\frac{m}{{2}^{p}} = \left\{  {\left. \frac{m - 1}{{2}^{p}}\right| \;\frac{m + 1}{{2}^{p}}}\right\}  . \tag{1.36}
\]

On the right-hand side of (1.36), \(m - 1\) and \(m + 1\) are even integers, and we assume that the resulting fractions are maximally canceled as fractions normally are, so they have a smaller denominator, or are integers as in (1.30). Hence, (1.36) is a recursive definition. An example of (1.36) is \(\frac{17}{8} = \left\{  {\frac{16}{8}\left| {\;\frac{18}{8}}\right. }\right\}   = \left\{  {2\left| {\;\frac{9}{4}}\right. }\right\}\) .

\begin{tcolorbox}
在 (1.36) 的右边， \(m - 1\) 和 \(m + 1\) 是偶数，并且我们假设得到的分数像通常那样进行了最大程度的约分，所以它们的分母更小，或者像 (1.30) 中那样是整数。因此，(1.36) 是一个递归定义。(1.36) 的一个例子是 \(\frac{17}{8} = \left\{  {\frac{16}{8}\left| {\;\frac{18}{8}}\right. }\right\}   = \left\{  {2\left| {\;\frac{9}{4}}\right. }\right\}\) 。
\end{tcolorbox}

Equation (1.33) justifies the definition \(\frac{1}{2} = \{ 0 \mid  1\}\) . In a similar way, one has to show that the game \(\frac{m}{{2}^{p}}\) in (1.36) is a number and that

\begin{tcolorbox}
方程 (1.33) 证明了定义 \(\frac{1}{2} = \{ 0 \mid  1\}\) 的合理性。类似地，必须证明 (1.36) 中的博弈 \(\frac{m}{{2}^{p}}\) 是一个数，并且
\end{tcolorbox}

\[
\frac{m}{{2}^{p}} + \frac{m}{{2}^{p}} + \frac{-m}{{2}^{p - 1}} \equiv  0 \tag{1.37}
\]

This is done by using the recursive definition of \(G = \frac{m}{{2}^{p}}\) in (1.36) as \(G = \left\{  {{G}^{L} \mid  {G}^{R}}\right\}\) with \({G}^{L} = \frac{m - 1}{{2}^{p - 1}}\) and \({G}^{R} = \frac{m + 1}{{2}^{p - 1}}\) . According to Definition 1.18, we also have to show \({G}^{L} < G\) and \(G < {G}^{R}\) . For details we refer to chapters 4 and 5 of Albert, Nowakowski, and Wolfe (2007). Exercise 1.14 asks to compute some numbers that represent positions of the game Black-White Hackenbush, as special cases of (1.37).

\begin{tcolorbox}
这可以通过使用 (1.36) 中 \(G = \frac{m}{{2}^{p}}\) 的递归定义来完成，即 \(G = \left\{  {{G}^{L} \mid  {G}^{R}}\right\}\) 且 \({G}^{L} = \frac{m - 1}{{2}^{p - 1}}\) 和 \({G}^{R} = \frac{m + 1}{{2}^{p - 1}}\) 。根据定义 1.18，我们还必须证明 \({G}^{L} < G\) 和 \(G < {G}^{R}\) 。具体细节请参考 Albert、Nowakowski 和 Wolfe(2007 年)的第 4 章和第 5 章。练习 1.14 要求计算一些代表黑白哈肯布什游戏(Black - White Hackenbush)位置的数，作为 (1.37) 的特殊情况。
\end{tcolorbox}

Using Definition 1.19, every rational number with a denominator \({2}^{p}\) for a positive integer \(p\) can be represented as a game. In that game, a player who makes a move gives up a safe advantage of \(\frac{1}{{2}^{p}}\) of a move. In a game sum of numbers, the player will move in the component with largest dominator where he gives up least. A game sum may have components that are not numbers, such as Nim heaps. The number avoidance theorem states that a player should move first in those components, and last in components that are numbers.

\begin{tcolorbox}
使用定义 1.19，分母为 \({2}^{p}\) (其中 \(p\) 为正整数)的每个有理数都可以表示为一个博弈。在该博弈中，进行一步行动的玩家会放弃 \(\frac{1}{{2}^{p}}\) 步的安全优势。在数的博弈之和中，玩家会在分母最大的分量中行动，因为在那里他放弃的最少。博弈之和可能包含非数的分量，例如尼姆堆(Nim heap)。避数定理指出，玩家应先在这些非数的分量中行动，最后在数的分量中行动。
\end{tcolorbox}

Not all partizan games are numbers. Indeed, there are games which are not numbers because moving there confers an extra advantage. For example, the \(2 \times  2\) Domineering board in (1.29) is equal to \(\{ 1 \mid   - 1\}\) . This is called a hot game because moving there gives the player an extra advantage in the form of additional moves. The \(3 \times  2\) Domineering board, which by (1.26) we found to be equal to \(\left\{  {\frac{1}{2} \mid   - 2}\right\}\) , is also a hot game. A number, in contrast, is a cold game.

\begin{tcolorbox}
并非所有有偏博弈(partizan game)都是数。实际上，有些博弈不是数，因为在这些博弈中行动会带来额外的优势。例如，(1.29) 中的 \(2 \times  2\) 纵横棋(Domineering)棋盘等于 \(\{ 1 \mid   - 1\}\) 。这被称为热情博弈(hot game)，因为在那里行动会以额外步数的形式给玩家带来额外的优势。由 (1.26) 可知等于 \(\left\{  {\frac{1}{2} \mid   - 2}\right\}\) 的 \(3 \times  2\) 纵横棋棋盘也是热情博弈。相反，数是冷静博弈(cold game)。
\end{tcolorbox}

If there are several hot components in a game sum, analyzing their temperature is of interest to decide where to move first. This has been applied to some situations of the board game Go with disconnected components. In Chess, different moves are rarely about disconnected parts of the game. However, Chess has the related concept of tempo of development which can outweigh a material advantage.

\begin{tcolorbox}
如果博弈之和中有几个热情分量，分析它们的热度对于决定先在何处行动很有意义。这已应用于一些具有不相连部分的围棋局面。在国际象棋中，不同的走法很少涉及游戏的不相连部分。然而，国际象棋有相关的发展节奏概念，这可能会超过物质优势。
\end{tcolorbox}

This concludes our peek into the theory of partizan games. Further comments on this topic, with literature references, are given in the next section.

\begin{tcolorbox}
至此，我们对有偏博弈理论的初步探讨结束。下一节将给出关于该主题的进一步评论以及文献参考。
\end{tcolorbox}

\section*{1.9 Further Reading}

\begin{tcolorbox}
\section*{1.9 进一步阅读}
\end{tcolorbox}

An undergraduate textbook on combinatorial games is Albert, Nowakowski, and Wolfe (2007), from which we have adopted the name "top-down induction". This book starts from many examples of games, considers in depth the game Dots and Boxes, and treats the theory of short games where (as we have assumed throughout) every position has only finitely many options and no game position is ever revisited (which in "loopy" games is allowed and could lead to a draw). An authoritative graduate textbook is Siegel (2013), which gives excellent overviews before treating each topic in depth, including newer research developments. It stresses the concept of equivalence of games that we have treated in Section 1.4.

\begin{tcolorbox}
Albert、Nowakowski 和 Wolfe(2007 年)编写的一本关于组合博弈的本科教材，我们采用了其中“自上而下归纳(top - down induction)”这一名称。这本书从许多博弈的例子入手，深入探讨了点格棋(Dots and Boxes)博弈，并处理了短博弈理论，在该理论中(正如我们始终假设的那样)每个位置只有有限个选项，并且不会重复访问任何博弈位置(在“循环”博弈中这是允许的，并且可能导致平局)。一本权威的研究生教材是 Siegel(2013 年)的著作，它在深入处理每个主题之前给出了精彩的概述，包括较新的研究进展。它强调了我们在 1.4 节中讨论的博弈等价概念。
\end{tcolorbox}

The classic text on combinatorial game theory is Winning Ways by Berlekamp, Conway, and Guy (2001-2004). Its wit, vivid colour drawings, and the wealth of games it considers make this a very attractive original source. However, the mathematics is hard, and it helps to know first what the topic is about; for example, equivalence of games is rather implicit and written as equality.

\begin{tcolorbox}
关于组合博弈理论的经典著作是 Berlekamp、Conway 和 Guy(2001 - 2004 年)的《Winning Ways》。其诙谐的语言、生动的彩色插图以及所涉及的大量博弈，使其成为极具吸引力的原始资料。然而，其数学内容较难，并且最好先了解该主题的大致内容；例如，博弈的等价性相当隐晦，被写成了相等。
\end{tcolorbox}

All these books consider directly partizan games. We have chosen to focus on the simpler impartial games, apart from our short introduction to partizan games in Section 1.8. For a more detailed study of games as numbers see Albert, Nowakowski, and Wolfe (2007, Section 5.1), where Definition 5.12 is our Definition 1.18.

\begin{tcolorbox}
所有这些书籍都直接讨论有偏博弈。除了在 1.8 节中对有偏博弈的简要介绍外，我们选择专注于较简单的无偏博弈(impartial game)。若要更详细地研究将博弈表示为数的内容，请参阅 Albert、Nowakowski 和 Wolfe(2007 年，第 5.1 节)，其中定义 5.12 就是我们的定义 1.18。
\end{tcolorbox}

The winning strategy for the game Nim based on the binary system was first described by Bouton (1901). The Queen move game is due to Wythoff (1907), described as an extension of Nim. The observation that every impartial game is equivalent to a Nim heap is independently due to Sprague (1935) and Grundy (1939). This is therefore also called the Sprague-Grundy theory of impartial games, and Nim values are also called "Sprague-Grundy values" or just "Grundy values".

\begin{tcolorbox}
基于二进制的尼姆游戏(Nim)必胜策略最早由布顿(Bouton，1901年)提出。皇后移动游戏由维托夫(Wythoff，1907年)提出，他将其描述为尼姆游戏的扩展。每个无偏博弈(impartial game)都等同于一个尼姆堆(Nim heap)这一发现分别归功于斯普拉格(Sprague，1935年)和格鲁迪(Grundy，1939年)。因此，这也被称为无偏博弈的斯普拉格 - 格鲁迪理论(Sprague - Grundy theory)，尼姆值(Nim value)也被称为“斯普拉格 - 格鲁迪值(Sprague - Grundy values)”或简称“格鲁迪值(Grundy values)”。
\end{tcolorbox}

The approach in Section 1.5 to the mex rule is inspired by Berlekamp, Conway, and Guy (2001-2004), chapter 4, which describes Poker Nim, the Rook and Queen move games, Northcott's game (Exercise 1.8), and chapter 5 with Kayles and Lasker's Nim (called Split-Nim in Exercise 1.12). Chomp (Exercises 1.5 and 1.10) was described by Gale (1974); for a generalization see Gale and Neyman (1982). The digraph game in Exercise 1.9 is another great example for understanding the mex rule, due to Fraenkel (1996).

\begin{tcolorbox}
1.5节中对mex规则的处理方法受到了伯利坎普(Berlekamp)、康威(Conway)和盖伊(Guy)(2001 - 2004年)所著书籍第4章的启发，该章介绍了扑克尼姆游戏(Poker Nim)、车和皇后移动游戏、诺斯科特游戏(练习1.8)，第5章介绍了凯尔斯博弈(Kayles)和拉斯克尼姆游戏(在练习1.12中称为分裂尼姆游戏，Split - Nim)。巧取游戏(Chomp，练习1.5和1.10)由盖尔(Gale，1974年)提出；关于其推广可参考盖尔和内曼(Gale and Neyman，1982年)的研究。练习1.9中的有向图博弈(digraph game)是理解mex规则的另一个很好的例子，它由弗伦克尔(Fraenkel，1996年)提出。
\end{tcolorbox}

The systematic study of sums of partizan games led John Conway, one of the authors of Winning Ways, to the insight that it is possible to define in an elegant way real numbers (in fact, a generalization called surreal numbers) as the strengths of positions in certain games. A number is a game \(G\) as in Definition 1.18, but where the sets \({\mathcal{G}}^{L}\) and \({\mathcal{G}}^{R}\) of left options \({G}^{L}\) and right options \({G}^{R}\) are allowed to be infinite. The condition \({G}^{L} < G < {G}^{R}\) generalizes "Dedekind cuts", a standard construction of the real numbers. The original ideas are described in Conway (2001), known as ONAG ("On Numbers and Games"), which is another classic. For an introduction see Conway (1977).

\begin{tcolorbox}
对有偏博弈(partizan game)之和的系统研究使《制胜之道》(Winning Ways)的作者之一约翰·康威(John Conway)认识到，可以用一种优雅的方式将实数(实际上是一种称为超实数的推广)定义为某些博弈中局势的强度。一个数是如定义1.18中的博弈 \(G\) ，但左选项 \({G}^{L}\) 的集合 \({\mathcal{G}}^{L}\) 和右选项 \({G}^{R}\) 的集合 \({\mathcal{G}}^{R}\) 可以是无限集。条件 \({G}^{L} < G < {G}^{R}\) 推广了“戴德金分割(Dedekind cut)”，这是实数的一种标准构造方法。原始思想在康威(Conway，2001年)所著的《论数字与博弈》(ONAG，“On Numbers and Games”)中有所描述，这也是一本经典著作。关于入门介绍可参考康威(Conway，1977年)的书。
\end{tcolorbox}

An informative and entertaining mathematical study of parlour games, with detailed historical notes, is Bewersdorff (2005). Combinatorial games are treated in part II of that book.

\begin{tcolorbox}
贝沃斯多夫(Bewersdorff，2005年)对客厅游戏进行了既有知识性又有趣味性的数学研究，并附有详细的历史注释。该书第二部分讨论了组合博弈。
\end{tcolorbox}

\section*{1.10 Exercises for Chapter 1}

\begin{tcolorbox}
\section*{1.10 第1章练习}
\end{tcolorbox}

Exercise 1.1. For this exercise, be careful to use only the definitions and not your intuitions that are familiar to you about the symbols \(\leq\) and \(<\) . Consider a set \(S\) with a binary relation \(<\) that is transitive and fulfills for all \(x\) in \(S\)

\begin{tcolorbox}
练习1.1。在做本题时，要注意仅使用定义，而不要依赖你对符号 \(\leq\) 和 \(<\) 的熟悉直觉。考虑一个具有二元关系 \(<\) 的集合 \(S\) ，该关系是传递的，并且对于 \(S\) 中的所有 \(x\) 都满足
\end{tcolorbox}

\[
\text{ not }x < x\;\text{ ( } < \text{ is irreflexive). } \tag{1.38}
\]

Define the relation \(\leq\) on \(S\) by (1.7), that is, \(x \leq  y \Leftrightarrow  x < y\) or \(x = y\) .

\begin{tcolorbox}
通过(1.7)式在 \(S\) 上定义关系 \(\leq\) ，即 \(x \leq  y \Leftrightarrow  x < y\) 或 \(x = y\) 。
\end{tcolorbox}

(a) Show that \(\leq\) is a partial order, and that \(<\) is obtained from \(\leq\) via (1.6), that is, \(x < y \Leftrightarrow  x \leq  y\) and \(x \neq  y\) .

\begin{tcolorbox}
(a) 证明 \(\leq\) 是一个偏序(partial order)，并且 \(<\) 可通过(1.6)式从 \(\leq\) 得到，即 \(x < y \Leftrightarrow  x \leq  y\) 且 \(x \neq  y\) 。
\end{tcolorbox}

(b) Show that if \(<\) is obtained from a partial order \(\leq\) via (1.6), then \(<\) is transitive and irreflexive, and (1.7) holds, that is, \(x \leq  y \Leftrightarrow  x < y\) or \(x = y\) .

\begin{tcolorbox}
(b) 证明如果 \(<\) 是通过(1.6)式从一个偏序 \(\leq\) 得到的，那么 \(<\) 是传递且非自反的，并且(1.7)式成立，即 \(x \leq  y \Leftrightarrow  x < y\) 或 \(x = y\) 。
\end{tcolorbox}

Exercise 1.2. Consider the game Nim with heaps of tokens. The players alternately remove some tokens from one of the heaps. The player to remove the last token wins. Try to prove part (a) without reference to the binary representation for playing Nim optimally.

\begin{tcolorbox}
练习1.2。考虑有若干堆代币的尼姆游戏。玩家轮流从其中一堆中取走一些代币。取走最后一个代币的玩家获胜。在证明(a)部分时，尽量不参考尼姆游戏的最优玩法的二进制表示。
\end{tcolorbox}

(a) For all positions with three heaps, where one of the heaps has only one token, describe exactly the losing positions. Justify your answer, for example with a proof by induction, or by theorems on Nim. [Hint: Start with the easy cases to find the pattern.]

\begin{tcolorbox}
(a) 对于所有三堆代币的局势，其中有一堆只有一个代币，准确描述失败位置(losing position)。说明你的答案的合理性，例如用归纳法证明，或者使用尼姆游戏的定理。[提示:从简单情况入手寻找规律。]
\end{tcolorbox}

(b) Consider a Nim game with three heaps whose sizes are consecutive integers, for example2,3,4. Show that the only case where this is a losing position is 1,2,3and that otherwise this is always a winning position. [Hint: Use (a).]

\begin{tcolorbox}
(b) 考虑一个三堆代币的尼姆游戏，其堆的大小是连续整数，例如2、3、4。证明这种局势是失败位置的唯一情况是1、2、3，否则总是必胜局势(winning position)。[提示:使用(a)的结论。]
\end{tcolorbox}

(c) Determine all initial winning moves for Nim with three heaps of sizes 8, 11, and 13, using the binary representation of the heap sizes.

\begin{tcolorbox}
(c) 使用堆大小的二进制表示，确定三堆大小分别为8、11和13的尼姆游戏的所有初始必胜步(winning move)。
\end{tcolorbox}

Exercise 1.3. Misère Nim is played just like Nim but where the last player to move loses. A losing position is therefore a single Nim heap with a single token in it (which the player then has to take). Another losing position is given by three heaps with a single token each.

\begin{tcolorbox}
练习1.3。反向尼姆游戏(misère Nim)的玩法与尼姆游戏(Nim)相同，但最后一个行动的玩家输。因此，一个包含单个棋子的单个尼姆堆(Nim heap)就是一个失败位置(losing position)(玩家必须取走该棋子)。另一个失败位置是三个各含一个棋子的尼姆堆。
\end{tcolorbox}

(a) Determine the winning and losing positions of misère Nim with one or two heaps, and the winning moves from a winning position.

\begin{tcolorbox}
(a) 确定有一个或两个尼姆堆的反向尼姆游戏的必胜局势(winning position)和失败位置，以及从必胜局势出发的必胜步(winning move)。
\end{tcolorbox}

(b) Are three heaps with sizes1,2,3winning or losing in misère Nim?

\begin{tcolorbox}
(b) 在反向尼姆游戏中，大小为1、2、3的三个尼姆堆是必胜局势还是失败位置？
\end{tcolorbox}

(c) Try to describe the losing positions in misère Nim for any number of heaps. [Hint: They are closely related to the losing positions in Nim with normal play.]

\begin{tcolorbox}
(c) 尝试描述任意数量尼姆堆的反向尼姆游戏的失败位置。[提示:它们与正常规则(normal play)的尼姆游戏的失败位置密切相关。]
\end{tcolorbox}

Exercise 1.4. The impartial game Cram is played on a board of \(m \times  n\) squares, where players alternately place a domino on the board which covers two adjacent squares that are free (not yet occupied by a domino), vertically or horizontally. The

\begin{tcolorbox}
练习1.4。无偏博弈(impartial game)塞砖棋(Cram)在一个 \(m \times  n\) 方格的棋盘上进行，玩家轮流在棋盘上放置一块多米诺骨牌(domino)，该骨牌覆盖两个相邻的空格(尚未被多米诺骨牌占据)，可以是垂直或水平放置。
\end{tcolorbox}

first player who cannot place a domino any more loses. Example play for a \(2 \times  3\) board:

\begin{tcolorbox}
第一个无法再放置多米诺骨牌的玩家输。 \(2 \times  3\) 棋盘的示例玩法:
\end{tcolorbox}

\begin{center}
\includegraphics[max width=0.7\textwidth]{images/bo_d2e5cmbef24c73b03tn0_44_357_317_897_124_0.jpg}
\end{center}
\hspace*{3em} 

(a) Who will win in \(3 \times  3\mathrm{{Cram}}\) ? [Hint: Use the symmetry of the game to investigate possible moves, and remember that it suffices to find one winning strategy.]

\begin{tcolorbox}
(a) 在 \(3 \times  3\mathrm{{Cram}}\) 的情况下谁会赢？[提示:利用游戏的对称性来研究可能的走法，并记住找到一种必胜策略就足够了。]
\end{tcolorbox}

(b) Who will win in \(m \times  n\) Cram when both \(m\) and \(n\) are even? (That player has a simple and precisely definable winning strategy which you have to find.)

\begin{tcolorbox}
(b) 当 \(m\) 和 \(n\) 都是偶数时，在 \(m \times  n\) 的塞砖棋中谁会赢？(该玩家有一个简单且可精确定义的必胜策略，你必须找到它。)
\end{tcolorbox}

(c) Who will win in \(m \times  n\) Cram when \(m\) is odd and \(n\) is even? (Easy with (b).)

\begin{tcolorbox}
(c) 当 \(m\) 是奇数且 \(n\) 是偶数时，在 \(m \times  n\) 的塞砖棋中谁会赢？(利用(b)的结果很容易得出。)
\end{tcolorbox}

Note (not a question): Because of the known answers from (b) and (c), this game is more interesting for "real play" on an \(m \times  n\) board where both \(m\) and \(n\) are odd. Play it with your friends on a \(5 \times  5\) board, for example. The situation often decomposes into independent parts, like contiguous fields of2,3,4,5,6squares, that have a known winner, which may help you analyse the situation.

\begin{tcolorbox}
注意(不是问题):由于从(b)和(c)中已知答案，这个游戏在 \(m \times  n\) 棋盘上“实际玩”时更有趣，其中 \(m\) 和 \(n\) 都是奇数。例如，和你的朋友在 \(5 \times  5\) 棋盘上玩这个游戏。局势通常会分解成独立的部分，比如2、3、4、5、6个方格的连续区域，这些区域的获胜者是已知的，这可能有助于你分析局势。
\end{tcolorbox}

Exercise 1.5. Consider the following game Chomp: A rectangular array of \(m \times  n\) dots is given, in \(m\) rows and \(n\) columns, like \(3 \times  4\) in the next picture on the left. A dot in row \(i\) (from the top) and column \(j\) (from left) is named(i, j). A move consists in picking a dot(i, j)and removing it and all other dots to the right and below it, which means removing all dots \(\left( {{i}^{\prime },{j}^{\prime }}\right)\) with \({i}^{\prime } \geq  i\) and \({j}^{\prime } \geq  j\) , as shown for \(\left( {i,j}\right)  = \left( {2,3}\right)\) in the middle picture, resulting in the picture on the right:

\begin{tcolorbox}
练习1.5。考虑以下巧取游戏(Chomp):给定一个由 \(m \times  n\) 个点组成的矩形阵列，有 \(m\) 行和 \(n\) 列，就像下一张图左边的 \(3 \times  4\) 那样。位于第 \(i\) 行(从顶部开始)和第 \(j\) 列(从左边开始)的点被命名为(i, j)。一步行动包括选择一个点(i, j)并移除它以及它右边和下方的所有其他点，这意味着移除所有满足 \({i}^{\prime } \geq  i\) 且 \({j}^{\prime } \geq  j\) 的点 \(\left( {{i}^{\prime },{j}^{\prime }}\right)\) ，如中间图中 \(\left( {i,j}\right)  = \left( {2,3}\right)\) 所示，结果如右图所示:
\end{tcolorbox}

\begin{center}
\includegraphics[max width=0.7\textwidth]{images/bo_d2e5cmbef24c73b03tn0_44_324_1314_985_174_0.jpg}
\end{center}
\hspace*{3em} 

Player I is the first player to move, players alternate, and the last player who removes a dot loses.

\begin{tcolorbox}
玩家I是第一个行动的玩家，玩家轮流行动，最后一个移除点的玩家输。
\end{tcolorbox}

One can think of these dots as (real) cookies: a move is to eat a cookie and all those to the right and below it, but the top left cookie is poisoned.

\begin{tcolorbox}
可以把这些点想象成(真正的)饼干:一步行动就是吃掉一块饼干以及它右边和下方的所有饼干，但左上角的饼干是有毒的。
\end{tcolorbox}

In the following questions, do not rush into any quick and false application of Nim values. Part (c) has an elegant but tricky answer.

\begin{tcolorbox}
在以下问题中，不要急于快速且错误地应用尼姆值(Nim value)。(c)部分有一个巧妙但棘手的答案。
\end{tcolorbox}

(a) Assuming optimal play, determine the winning player and a winning move for Chomp of size \(2 \times  2\) , size \(2 \times  3\) , size \(2 \times  n\) , and size \(m \times  m\) , where \(m \geq  3\) . Justify your answers.

\begin{tcolorbox}
(a) 假设进行最优玩法，确定大小为 \(2 \times  2\) 、大小为 \(2 \times  3\) 、大小为 \(2 \times  n\) 和大小为 \(m \times  m\) 的巧取游戏的获胜玩家和必胜步，其中 \(m \geq  3\) 。证明你的答案。
\end{tcolorbox}

(b) As it is described here, Chomp is a misère game where the last player to make a move loses. Suppose we want to play the same game but so that the normal play convention applies, where the last player to move wins. How should the array of dots be changed to achieve this? (Normal play with the array as given would be a boring game, with the obvious winning move of taking the top left dot(1,1).)

\begin{tcolorbox}
(b) 如这里所描述的，巧取游戏是一个反常规则的游戏，最后一个行动的玩家输。假设我们想玩同样的游戏，但采用正常规则，即最后一个行动的玩家赢。为了实现这一点，点的阵列应该如何改变？(对于给定的阵列采用正常规则会是一个无聊的游戏，明显的必胜步是取走左上角的点(1,1)。)
\end{tcolorbox}

(c) Show that when Chomp is played for a game of any size \(m \times  n\) , player I can always win (unless \(m = n = 1\) ). [Hint: You only have to show that a winning move exists, but you do not have to describe that winning move.]

\begin{tcolorbox}
(c) 证明当进行任意大小为 \(m \times  n\) 的巧取游戏时，玩家一总能获胜(除非 \(m = n = 1\) )。[提示:你只需证明存在必胜步，但无需描述该必胜步。]
\end{tcolorbox}

Exercise 1.6.

\begin{tcolorbox}
练习 1.6。
\end{tcolorbox}

(a) Complete the entries of equivalent Nim heaps for the Queen-move game in columns 5 and 6, rows 0 to 3, in the table in Figure 1.4 (assuming the queen can be anywhere on the board).

\begin{tcolorbox}
(a) 完成图 1.4 表格中第 5 列和第 6 列、第 0 行到第 3 行的皇后移动博弈等效尼姆堆的条目(假设皇后可以在棋盘上的任意位置)。
\end{tcolorbox}

(b) Describe all winning moves in the game sum of the Queen-move game and the Nim heap in Figure 1.3.

\begin{tcolorbox}
(b) 描述图 1.3 中皇后移动博弈和尼姆堆的博弈之和游戏中的所有必胜步。
\end{tcolorbox}

Exercise 1.7. Consider the game Cram from Exercise 1.4, played on a \(1 \times  n\) board for \(n \geq  2\) . Let \({D}_{n}\) be the Nim value of that game, so that the starting position of the \(1 \times  n\) board is equivalent to a Nim heap of size \({D}_{n}\) . For example, \({D}_{2} = 1\) because the \(1 \times  2\) board is equivalent to \(* 1\) .

\begin{tcolorbox}
练习 1.7。考虑练习 1.4 中的塞砖棋游戏，在 \(1 \times  n\) 的棋盘上进行，其中 \(n \geq  2\) 。设 \({D}_{n}\) 为该游戏的尼姆值，使得 \(1 \times  n\) 棋盘的起始位置等效于大小为 \({D}_{n}\) 的尼姆堆。例如， \({D}_{2} = 1\) ，因为 \(1 \times  2\) 棋盘等效于 \(* 1\) 。
\end{tcolorbox}

(a) How is \({D}_{n}\) computed from smaller values \({D}_{k}\) for \(k < n\) ? [Hint: Use sums of games and the mex rule. The notation for Nim sums is \(\oplus\) , where \(a \oplus  b = c\) if and only if \(* a +  * b \equiv   * c\) .]

\begin{tcolorbox}
(a) 如何从 \(k < n\) 时的较小值 \({D}_{k}\) 计算 \({D}_{n}\) ？[提示:使用博弈之和与 mex 规则。尼姆和的表示法为 \(\oplus\) ，其中当且仅当 \(* a +  * b \equiv   * c\) 时 \(a \oplus  b = c\) 。]
\end{tcolorbox}

(b) Give the values of \({D}_{n}\) up to \(n = {10}\) (or more, if you are ambitious - higher values come at \(n = {16}\) , and at some point they even repeat, but before you detect that you will probably have run out of patience). For which values of \(n\) in \(\{ 1,\ldots ,{10}\}\) is Cram on a \(1 \times  n\) board a losing game?

\begin{tcolorbox}
(b) 给出 \({D}_{n}\) 直到 \(n = {10}\) 的值(如果你有雄心壮志，也可以给出更多值——更高的值在 \(n = {16}\) 出现，并且在某个时候它们甚至会重复，但在你发现这一点之前，你可能已经失去耐心了)。对于 \(\{ 1,\ldots ,{10}\}\) 中的哪些 \(n\) 值， \(1 \times  n\) 棋盘上的塞砖棋是必败游戏？
\end{tcolorbox}

Exercise 1.8. Consider the following game on a rectangular board where a white and a black counter are placed in each row, like in this example:

\begin{tcolorbox}
练习 1.8。考虑在矩形棋盘上进行的以下游戏，每行放置一个白色和一个黑色棋子，例如:
\end{tcolorbox}

(*)

\begin{center}
\includegraphics[max width=0.3\textwidth]{images/bo_d2e5cmbef24c73b03tn0_45_750_1464_385_393_0.jpg}
\end{center}
\hspace*{3em} 

Player I is White and starts, and player II is Black. Players take turns. In a move, a player moves a counter of his color to any other square within its row, but may not jump over the other counter. For example, in (*) above, in row 8 White may move from e8 to any of the squares \(\mathrm{c}8,\mathrm{\;d}8,\mathrm{\;f}8,\mathrm{\;g}8\) , or h8. The player who can no longer move loses.

\begin{tcolorbox}
玩家一是白方且先行动，玩家二是黑方。玩家轮流行动。在一次行动中，玩家将自己颜色的棋子移动到其所在行的任何其他方格，但不能跳过另一个棋子。例如，在上面的 (*) 中，在第 8 行，白方可以从 e8 移动到 \(\mathrm{c}8,\mathrm{\;d}8,\mathrm{\;f}8,\mathrm{\;g}8\) 中的任何一个方格，或 h8。无法再移动的玩家输。
\end{tcolorbox}

(a) Who will win in the following position?

\begin{tcolorbox}
(a) 在以下局面中谁会获胜？
\end{tcolorbox}

\begin{center}
\includegraphics[max width=0.3\textwidth]{images/bo_d2e5cmbef24c73b03tn0_46_610_203_387_396_0.jpg}
\end{center}
\hspace*{3em} 

(b) Show that White can win in position (*) above. Give at least two winning moves from that position. Justify your answers. [Hint: Compare this with another game that is not impartial and that also violates the ending condition, but that nevertheless is close to Nim and ends in finite time when played well.]

\begin{tcolorbox}
(b) 证明白方在上面的 (*) 局面中可以获胜。从该局面给出至少两种必胜步。说明你的答案的合理性。[提示:将此游戏与另一个非无偏博弈且也违反终止条件的游戏进行比较，但该游戏仍然接近尼姆游戏，并且如果玩得好，会在有限时间内结束。]
\end{tcolorbox}

Exercise 1.9. Consider the following network (in technical terms, a directed graph or "digraph"). Each circle, here marked with one of the letters A to P, represents a node of the network. Some of these nodes (here A, F, G, H, and K) have counters on them. There can be any number of counters on a node, like currently the two counters on \(\mathrm{H}\) . In a move, one of the counters is moved to a neighbouring node in the direction of the arrow as indicated, for example from \(\mathrm{F}\) to \(\mathrm{I}\) (but not from \(\mathrm{F}\) to C, nor directly from F to D, say). Players alternate, and the first player no longer able to move loses.

\begin{tcolorbox}
练习 1.9。考虑以下网络(用技术术语来说，是一个有向图或“digraph”)。每个圆圈，这里用字母 A 到 P 之一标记，表示网络的一个节点。其中一些节点(这里是 A、F、G、H 和 K)上有棋子。一个节点上可以有任意数量的棋子，例如目前 \(\mathrm{H}\) 上有两个棋子。在一次行动中，其中一个棋子沿着箭头方向移动到相邻节点，例如从 \(\mathrm{F}\) 到 \(\mathrm{I}\) (但不能从 \(\mathrm{F}\) 到 C，也不能直接从 F 到 D)。玩家轮流行动，第一个无法再移动的玩家输。
\end{tcolorbox}

\begin{center}
\includegraphics[max width=0.5\textwidth]{images/bo_d2e5cmbef24c73b03tn0_46_450_1208_742_541_0.jpg}
\end{center}
\hspace*{3em} 

(a) Explain why this game fulfills the ending condition.

\begin{tcolorbox}
(a) 解释为什么这个游戏满足终止条件。
\end{tcolorbox}

(b) Who is winning in the above position? If it is the first player to move, describe all possible winning moves. Justify your answer.

\begin{tcolorbox}
(b) 在上述局面中谁会获胜？如果是先手玩家行动，描述所有可能的必胜步。说明你的答案的合理性。
\end{tcolorbox}

(c) How does the answer to (b) change when the arrow from \(\mathrm{J}\) to \(\mathrm{K}\) is reversed so that it points from \(\mathrm{K}\) to \(\mathrm{J}\) instead?

\begin{tcolorbox}
(c) 当从 \(\mathrm{J}\) 到 \(\mathrm{K}\) 的箭头反转，使其从 \(\mathrm{K}\) 指向 \(\mathrm{J}\) 时，(b) 的答案会如何变化？
\end{tcolorbox}

Exercise 1.10. Consider the game Chomp from Exercise 1.5 of size \(2 \times  4\) , in a game sum with a Nim heap of size 4.

\begin{tcolorbox}
练习 1.10。考虑练习 1.5 中大小为 \(2 \times  4\) 的巧取游戏，与大小为 4 的尼姆堆进行博弈之和。
\end{tcolorbox}

\begin{center}
\includegraphics[max width=0.5\textwidth]{images/bo_d2e5cmbef24c73b03tn0_47_564_191_757_147_0.jpg}
\end{center}
\hspace*{3em} 

What are the winning moves of the starting player I, if any? [Hint: Represent Chomp as a game in the normal play convention (see Exercise 1.5(b), by changing the dot pattern), so that the losing player is not the player who takes the "poisoned cookie", but the player who can no longer move. This will simplify finding the various Nim values.]

\begin{tcolorbox}
先手玩家 I 是否有必胜步？如果有，是哪些？[提示:将巧取游戏(Chomp)表示为正常规则下的博弈(见练习 1.5(b)，通过改变点模式)，这样输的玩家不是拿到“有毒曲奇”的玩家，而是无法继续行动的玩家。这将简化计算各种尼姆值(Nim value)的过程。]
\end{tcolorbox}

Exercise 1.11. In \(m \times  n\) Cram (see Exercise 1.4), a rectangular board of \(m \times  n\) squares is given. The two players alternately place a domino either horizontally or vertically on two unoccupied adjacent squares, which then become occupied. The last player to be able to place a domino wins.

\begin{tcolorbox}
练习 1.11。在 \(m \times  n\) 塞砖棋(Cram)游戏中(见练习 1.4)，给定一个由 \(m \times  n\) 个方格组成的矩形棋盘。两名玩家交替地将一块多米诺骨牌(domino)水平或垂直放置在两个未被占据的相邻方格上，放置后这两个方格就被占据了。最后能够放置多米诺骨牌的玩家获胜。
\end{tcolorbox}

(a) Find the Nim value (size of the equivalent Nim heap) of \(2 \times  3\) Cram.

\begin{tcolorbox}
(a) 求出 \(2 \times  3\) 塞砖棋的尼姆值(等效尼姆堆(Nim heap)的大小)。
\end{tcolorbox}

(b) Find all winning moves, if any, for the game sum of a \(2 \times  3\) Cram game and a \(1 \times  4\) Cram game.

\begin{tcolorbox}
(b) 对于一个 \(2 \times  3\) 塞砖棋游戏和一个 \(1 \times  4\) 塞砖棋游戏的博弈之和(game sum)，如果存在必胜步，找出所有的必胜步。
\end{tcolorbox}

Exercise 1.12. Consider the following variant of Nim called Split-Nim, which is played with heaps of tokens as in Nim. Like in Nim, a player can remove some tokens from one of the heaps, or else split a heap into two (not necessarily equal) new heaps. For example, a heap of size 4 can be reduced to size 3, 2, 1, or 0 as in ordinary Nim, or be split into two heaps of sizes 1 and 3, respectively, or into two heaps of sizes 2 and 2. As usual, the last player to be able to move wins.

\begin{tcolorbox}
练习 1.12。考虑尼姆游戏(Nim)的一个变体，称为分裂尼姆游戏(Split - Nim)，它和尼姆游戏一样使用代币堆进行游戏。和尼姆游戏一样，玩家可以从其中一个代币堆中移除一些代币，或者将一个代币堆分裂成两个(不一定相等)新的代币堆。例如，一个大小为 4 的代币堆可以像普通尼姆游戏那样减少到大小为 3、2、1 或 0，或者分裂成大小分别为 1 和 3 的两个代币堆，或者分裂成大小为 2 和 2 的两个代币堆。和往常一样，最后能够行动的玩家获胜。
\end{tcolorbox}

(a) Find the Nim values (size of equivalent Nim heaps) of the single Split-Nim heaps of size1,2,3, and 4, respectively.

\begin{tcolorbox}
(a) 分别求出大小为 1、2、3 和 4 的单个分裂尼姆代币堆的尼姆值(等效尼姆堆的大小)。
\end{tcolorbox}

(b) Find all winning moves, if any, when Split-Nim is played with three heaps of sizes1,2,3.

\begin{tcolorbox}
(b) 当分裂尼姆游戏使用大小为 1、2、3 的三个代币堆进行时，如果存在必胜步，找出所有的必胜步。
\end{tcolorbox}

(c) Find all winning moves, if any, when Split-Nim is played with three heaps of sizes 1, 2, 4 .

\begin{tcolorbox}
(c) 当分裂尼姆游戏使用大小为 1、2、4 的三个代币堆进行时，如果存在必胜步，找出所有的必胜步。
\end{tcolorbox}

(d) Determine if the following statement is true or false: "Two heaps in Split-Nim are equivalent if and only if they have equal size." Justify your answer.

\begin{tcolorbox}
(d) 判断以下陈述的真假:“在分裂尼姆游戏中，两个代币堆等效当且仅当它们的大小相等。”证明你的答案。
\end{tcolorbox}

Exercise 1.13. The impartial game Gray Hackenbush is played on a figure consisting of nodes and edges that are connected to these nodes or to the ground (the ground is the dashed line in the pictures below). A move is to remove an edge, and with it all the edges that are then no longer connected to the ground. For example, in the leftmost figure below, one can remove any edge in one of the three stalks. Removing the second edge from the stalk that consists of three edges takes the topmost edge with it, leaving only a single edge. All edges are colored gray, which means they can be removed by either player, so this is an impartial game. As usual, players alternate and the last player able to move wins.

\begin{tcolorbox}
练习 1.13。无偏博弈灰色哈肯布什游戏(Gray Hackenbush)在一个由节点(node)和连接到这些节点或地面的边(edge)组成的图形上进行(地面在下面的图片中用虚线表示)。一次行动是移除一条边，同时移除所有不再与地面相连的边。例如，在下面最左边的图形中，玩家可以移除三根茎中的任何一条边。移除由三条边组成的茎中的第二条边时，会同时移除最上面的边，只留下一条边。所有的边都被涂成灰色，这意味着任何玩家都可以移除它们，所以这是一个无偏博弈。和往常一样，玩家交替行动，最后能够行动的玩家获胜。
\end{tcolorbox}

(a) Compute the Nim values for the following three Gray Hackenbush figures (using what you know about Nim, and the mex rule):

\begin{tcolorbox}
(a) 计算以下三个灰色哈肯布什图形的尼姆值(使用你所知道的尼姆游戏和 mex 规则):
\end{tcolorbox}

\begin{center}
\includegraphics[max width=0.8\textwidth]{images/bo_d2e5cmbef24c73b03tn0_48_272_202_1098_205_0.jpg}
\end{center}
\hspace*{3em} 

(b) Compute the Nim values for the following four Gray Hackenbush figures:

\begin{tcolorbox}
(b) 计算以下四个灰色哈肯布什图形的尼姆值:
\end{tcolorbox}

\begin{center}
\includegraphics[max width=0.8\textwidth]{images/bo_d2e5cmbef24c73b03tn0_48_263_474_1112_229_0.jpg}
\end{center}
\hspace*{3em} 

(c) Based on the answer in (b), give a rule to compute the Nim value of a "tree" that is constructed by putting several "branches" with known Nim values on top of a "stem" of size \(n\) , for example \(n = 2\) in the picture below, where three branches of heights3,2, and 4 are put on top of the "stem". Prove the correctness of that rule (you may find it advantageous to glue the branches together at the bottom, as in (b)). Use this rule to find the Nim value of the rightmost picture.

\begin{tcolorbox}
(c) 根据 (b) 的答案，给出一个规则来计算一个“树”的尼姆值，这个“树”是通过将几个已知尼姆值的“分支”放在一个大小为 \(n\) 的“茎”上构建的，例如下面图片中的 \(n = 2\) ，其中高度为 3、2 和 4 的三个分支被放在“茎”的顶部。证明该规则的正确性(你可能会发现像 (b) 中那样将分支在底部粘在一起会更有利)。使用这个规则来求出最右边图片的尼姆值。
\end{tcolorbox}

\begin{center}
\includegraphics[max width=0.9\textwidth]{images/bo_d2e5cmbef24c73b03tn0_48_276_1093_1162_313_0.jpg}
\end{center}
\hspace*{3em} 

Exercise 1.14. This exercise is about partizan games as treated in Section 1.8.

\begin{tcolorbox}
练习 1.14。本练习是关于第 1.8 节中讨论的有偏博弈(partizan game)。
\end{tcolorbox}

(a) Complete the proof of (1.32) by showing that Right as starting player loses.

\begin{tcolorbox}
(a) 通过证明右玩家(Right player)为先手时会输来完成 (1.32) 的证明。
\end{tcolorbox}

The partizan game (Black-White) Hackenbush is played on a figure consisting of nodes and black or white edges that are connected to these nodes or to the ground, as in the pictures below. In a move the Left player removes a black edge and the Right player a white edge; if no such edge exists then the player loses. After such a move, all nodes and edges no longer connected to the ground disappear as well. For example, in a Hackenbush figure with a white edge on top of a black edge, Left can remove the bottom (black) edge and creates the zero game, whereas Right can remove the top (white) edge and leaves a single black edge.

\begin{tcolorbox}
有偏博弈(黑白)哈肯布什游戏((Black - White) Hackenbush)在一个由节点和连接到这些节点或地面的黑色或白色边组成的图形上进行，如下方图片所示。在一次行动中，左玩家(Left player)移除一条黑色边，右玩家移除一条白色边；如果不存在这样的边，那么该玩家输。在这样的行动之后，所有不再与地面相连的节点和边也会消失。例如，在一个黑色边上面有一条白色边的哈肯布什图形中，左玩家可以移除底部(黑色)边并形成零博弈，而右玩家可以移除顶部(白色)边并留下一条黑色边。
\end{tcolorbox}

\begin{center}
\includegraphics[max width=0.8\textwidth]{images/bo_d2e5cmbef24c73b03tn0_48_253_2008_1136_181_0.jpg}
\end{center}
\hspace*{3em} 

(b) Recall that for a combinatorial game \(G\) its negative is a game \(- G\) so that \(G + \left( {-G}\right)  \equiv  0\) . Describe the negative of a Hackenbush figure, and give an example. [Hint: Compare (1.24).]

\begin{tcolorbox}
(b) 回顾一下，对于一个组合博弈 \(G\) ，它的负博弈是一个博弈 \(- G\) ，使得 \(G + \left( {-G}\right)  \equiv  0\) 。描述哈肯布什图形的负博弈，并给出一个例子。[提示:参考 (1.24)。]
\end{tcolorbox}

We consider Hackenbush stalks, which are just paths of edges of either color. It can be shown that any such stalk is a number, with some examples considered here.

\begin{tcolorbox}
我们考虑哈肯布什杆，它们只是由任意一种颜色的边组成的路径。可以证明，任何这样的杆都代表一个数字，这里会考虑一些例子。
\end{tcolorbox}

(c) Draw the game that represents the number 3 (where Left is 3 moves ahead).

\begin{tcolorbox}
(c) 画出表示数字 3 的博弈图(其中先手玩家已经进行了 3 步)。
\end{tcolorbox}

(d) Find the numbers represented by the following Hackenbush stalks. Justify your answers. It helps to study and understand (1.36), and to consider only undominated moves of the players; also, compare the strength of the Left (black) player for these games to get an idea about the answer:

\begin{tcolorbox}
(d) 找出以下哈肯布什杆所代表的数字。证明你的答案。研究并理解 (1.36) 会有帮助，并且只考虑玩家的非劣势移动；此外，比较这些博弈中先手(黑色)玩家的实力，以了解答案:
\end{tcolorbox}

\begin{center}
\includegraphics[max width=0.5\textwidth]{images/bo_d2e5cmbef24c73b03tn0_49_614_742_759_238_0.jpg}
\end{center}
\hspace*{3em} 

where you can use as shorthand the following notation for these games:

\begin{tcolorbox}
对于这些博弈，你可以使用以下简写符号:
\end{tcolorbox}

\begin{center}
\includegraphics[max width=0.5\textwidth]{images/bo_d2e5cmbef24c73b03tn0_49_653_1101_677_122_0.jpg}
\end{center}
\hspace*{3em} 

(e) Construct a Hackenbush game (preferably a single stalk) that represents the number \(\frac{7}{4}\) .

\begin{tcolorbox}
(e) 构造一个哈肯布什博弈(最好是单个杆)，它代表数字 \(\frac{7}{4}\) 。
\end{tcolorbox}

Exercise 1.15. (Albert, Nowakowski, and Wolfe, 2007, Prep Problems 4.1 and 4.2.)

\begin{tcolorbox}
练习 1.15。(Albert、Nowakowski 和 Wolfe，2007 年，预备问题 4.1 和 4.2。)
\end{tcolorbox}

(a) Show that if \(G\) and \(H\) are partizan games and Left wins when moving second in \(G\) and in \(H\) , then Left wins when moving second in \(G + H\) . [Hint: Express what it means for Left to win when moving second in terms of outcome classes, as in (1.31).]

\begin{tcolorbox}
(a) 证明如果 \(G\) 和 \(H\) 是有偏博弈，并且先手玩家在 \(G\) 和 \(H\) 中后手移动时获胜，那么先手玩家在 \(G + H\) 中后手移动时也获胜。[提示:用结果类来表示先手玩家后手移动获胜的含义，如 (1.31) 所示。]
\end{tcolorbox}

(b) Give an example of games \(G\) and \(H\) so that Left wins when moving first in \(G\) and in \(H\) , but cannot win in \(G + H\) , even if Left can choose to move first or second in the game sum \(G + H\) . [Hint: Can both games be numbers? Or impartial games?] 2

\begin{tcolorbox}
(b) 给出博弈 \(G\) 和 \(H\) 的一个例子，使得先手玩家在 \(G\) 和 \(H\) 中先手移动时获胜，但在先手玩家可以选择在博弈之和 \(G + H\) 中先手或后手移动的情况下，也无法在 \(G + H\) 中获胜。[提示:这两个博弈都能是数字博弈吗？或者是无偏博弈吗？] 2
\end{tcolorbox}

\section*{Congestion Games}

\begin{tcolorbox}
\section*{拥塞博弈}
\end{tcolorbox}

Congestion means that a shared resource, such as a road, becomes more costly when more people use it. In a congestion game, multiple players decide on which resource to use, with the aim to minimize their cost. This interaction defines a game because the cost depends on what the other players do.

\begin{tcolorbox}
拥塞是指当更多人使用共享资源(如道路)时，该资源的使用成本会增加。在拥塞博弈中，多个玩家决定使用哪种资源，目的是最小化他们的成本。这种相互作用构成了一个博弈，因为成本取决于其他玩家的选择。
\end{tcolorbox}

We present congestion networks as a model of traffic where individual users choose routes from an origin to a destination. Each edge of the network is congested by creating a cost to each user that weakly increases with the number of users of that edge. The central result (Theorem 2.2) states that every congestion game has an equilibrium where no user can individually reduce her cost by changing her chosen route.

\begin{tcolorbox}
我们将拥塞网络作为交通模型，其中单个用户从起点到终点选择路线。网络的每条边都会产生拥塞，使得每个用户的成本随着使用该边的用户数量的增加而弱增加。核心结果(定理 2.2)表明，每个拥塞博弈都存在一个均衡，在该均衡下，没有用户可以通过改变其选择的路线来单独降低其成本。
\end{tcolorbox}

Like Chapter 1, this is another introductory chapter to an active and diverse area of game theory, where we can quickly show some important and nontrivial results. It also gives an introduction to the game-theoretic concepts of strategies and equilibrium before we develop the general theory.

\begin{tcolorbox}
与第 1 章一样，这是博弈论中一个活跃且多样化领域的另一章导论，我们可以快速展示一些重要且非平凡的结果。在我们发展一般理论之前，它还介绍了博弈论中的策略和均衡概念。
\end{tcolorbox}

Equilibrium is the result of selfish routing by the users who individually minimize their costs. This is typically not the socially optimal outcome, which could have a lower average cost. A simple example is the "Pigou network" (see Section 2.2), named after the English economist Arthur Pigou, who introduced the concept of an externality (such as congestion) to economics. Even more surprisingly, the famous Braess paradox, threated in Section 2.3, shows that adding capacity to a road network can worsen equilibrium congestion.

\begin{tcolorbox}
均衡是用户为了最小化自身成本而进行自私路由的结果。这通常不是社会最优结果，社会最优结果可能具有更低的平均成本。一个简单的例子是“庇古网络”(见第 2.2 节)，它以英国经济学家亚瑟·庇古的名字命名，他将外部性(如拥塞)的概念引入了经济学。更令人惊讶的是，著名的布雷斯悖论(在第 2.3 节讨论)表明，增加道路网络的容量可能会加剧均衡拥塞。
\end{tcolorbox}

In Section 2.4 we give the general definition of congestion networks. In Section 2.5 we prove that every congestion game has an equilibrium. This is proved with a cleverly chosen potential function, which is a single function that simultaneously reflects the change in cost of every user when she changes her strategy. Its minimum over all strategy choices therefore defines an equilibrium.

\begin{tcolorbox}
在第 2.4 节中，我们给出拥塞网络的一般定义。在第 2.5 节中，我们证明每个拥塞博弈都存在一个均衡。这是通过巧妙选择的势函数来证明的，势函数是一个单一的函数，它同时反映了每个用户改变策略时其成本的变化。因此，它在所有策略选择上的最小值定义了一个均衡。
\end{tcolorbox}

In Section 2.6, we explain the wider context of the presented model. It is the discrete model of atomic (non-splittable) flow with finitely many individual users who choose their traffic routes. As illustrated by the considered examples, the resulting equilibria are often not unique when the "last" user can optimally choose between more than one edge. The limit of many users is the continuous model of splittable flow where users can be infinitesmally small fractions of a "mass" of users who want to travel from an origin to a destination. In this model the equilibrium is often unique, but its existence proof is more technical and not given here.

\begin{tcolorbox}
在第 2.6 节中，我们解释所提出模型的更广泛背景。它是原子流(不可拆分)的离散模型，有有限数量的单个用户选择他们的交通路线。正如所考虑的例子所示，当“最后一个”用户可以在多条边上进行最优选择时，得到的均衡通常不是唯一的。大量用户的极限情况是可拆分流的连续模型，其中用户可以是想要从起点到终点旅行的“用户群体”的无穷小部分。在这个模型中，均衡通常是唯一的，但它的存在性证明更具技术性，这里不给出。
\end{tcolorbox}

We also mention the Price of Anarchy that compares the worst equilibrium cost with the socially optimal cost. For further details we give references in Section 2.7.

\begin{tcolorbox}
我们还提到了无政府代价(Price of Anarchy)，它用于比较最差均衡成本与社会最优成本。更多详细信息请参考2.7节中的参考文献。
\end{tcolorbox}

\section*{2.1 Prerequisites and Learning Outcomes}

\begin{tcolorbox}
\section*{2.1 先决条件与学习成果}
\end{tcolorbox}

There are no particular prequisites for this chapter. We use the terminology of directed graphs with nodes and their connecting edges (with possibly several "parallel" edges that connect the same two nodes); they are formally defined in Section 2.4. Standard calculus is used to find the minimum of a function of one variable. In Section 2.5 the flow through a network, given by the collection of flows \({f}_{e}\) through each edge \(e\) , is denoted by the letter \(f\) , which is here a variable.

\begin{tcolorbox}
本章没有特定的先决条件。我们使用有向图的术语，包括节点及其连接的边(可能存在连接相同两个节点的多条“平行”边)；它们在2.4节中有正式定义。使用标准微积分来求单变量函数的最小值。在2.5节中，网络中的流由每条边 \(e\) 上的流 \({f}_{e}\) 集合给出，用字母 \(f\) 表示，这里它是一个变量。
\end{tcolorbox}

After studying this chapter, you should be able to:

\begin{tcolorbox}
学习本章后，你应该能够:
\end{tcolorbox}

\begin{itemize}
\item understand the concept of congestion networks with their cost functions, and how edges in the network have a flow that depends on users' chosen paths;
\end{itemize}

\begin{tcolorbox}
\begin{itemize}
\item 理解拥塞网络(congestion network)及其成本函数的概念，以及网络中的边如何具有取决于用户所选路径的流；
\end{itemize}
\end{tcolorbox}

\begin{itemize}
\item find an equilibrium flow in a network, and the socially optimal flow, and compare them;
\end{itemize}

\begin{tcolorbox}
\begin{itemize}
\item 找出网络中的均衡流(equilibrium flow)和社会最优流，并对它们进行比较；
\end{itemize}
\end{tcolorbox}

\begin{itemize}
\item understand how a potential function helps finding an equilibrium.
\end{itemize}

\begin{tcolorbox}
\begin{itemize}
\item 理解势函数(potential function)如何帮助找到均衡。
\end{itemize}
\end{tcolorbox}

\section*{2.2 Introduction: The Pigou Network}

\begin{tcolorbox}
\section*{2.2 引言:庇古网络(Pigou network)}
\end{tcolorbox}

Getting around at rush hour takes much longer than at less busy times. Commuters take buses where people have to queue to board. There are more buses that fill the bus lanes, and cars are in the way and clog the streets because only a limited number can pass each green light. Commuters can take the bus, walk, or ride a bike, and cars can choose different routes. Congestion, which slows everything down, depends on how many people use certain forms of transport, and which route they choose. We assume that people try to choose an optimal route, but what is optimal depends on what others do. This is an interactive situation, which we call a game. We describe a mathematical model that makes the rules of this game precise.

\begin{tcolorbox}
高峰时段出行比非繁忙时段要花费长得多的时间。通勤者乘坐公交车时需要排队上车。更多的公交车挤满了公交专用道，而汽车也会阻碍交通，使街道堵塞，因为每个绿灯期间只有有限数量的车辆能够通过。通勤者可以选择乘坐公交车、步行或骑自行车，汽车也可以选择不同的路线。拥堵会使一切变慢，它取决于有多少人使用特定的交通方式以及他们选择的路线。我们假设人们会尝试选择最优路线，但什么是最优路线取决于其他人的选择。这是一种交互式的情况，我们称之为博弈。我们将描述一个数学模型，精确地定义这个博弈的规则。
\end{tcolorbox}

Figure 2.1 shows an example of a congestion network. This particular network has two nodes, an origin \(o\) and a destination \(d\) , and two edges that both connect \(o\) to \(d\) . Suppose there are several users of the network who all want to travel from \(o\) to \(d\) and can choose either edge. Each edge has a \(\operatorname{cost}c\left( x\right)\) associated with it, which is a function that describes how costly it is for each user to use that edge when it has a flow or load of \(x\) users. The top edge has constant cost \(c\left( x\right)  = 2\) , and the bottom edge has \(\operatorname{cost}c\left( x\right)  = x\) . The cost could for example be travel time, or incorporate additional monetary costs. The cost is the same for all users of the edge. The cost for each user of the top edge, no matter how many people use it, is always 2, whereas the bottom edge has cost 1 for one user, 2 for two users, and so on (and zero cost for no users, but then there is no one to profit from taking that zero-cost route).

\begin{tcolorbox}
图2.1展示了一个拥塞网络的示例。这个特定的网络有两个节点，一个起点 \(o\) 和一个终点 \(d\) ，以及两条都连接 \(o\) 到 \(d\) 的边。假设网络中有多个用户，他们都想从 \(o\) 前往 \(d\) ，并且可以选择任意一条边。每条边都有一个与之关联的 \(\operatorname{cost}c\left( x\right)\) ，它是一个函数，描述了当该边有 \(x\) 个用户的流或负载时，每个用户使用该边的成本。上边的边具有恒定成本 \(c\left( x\right)  = 2\) ，下边的边具有 \(\operatorname{cost}c\left( x\right)  = x\) 。成本例如可以是旅行时间，或者包含额外的货币成本。对于该边的所有用户来说，成本是相同的。无论有多少人使用上边的边，每个用户的成本始终是2，而下边的边对于一个用户的成本是1，对于两个用户的成本是2，依此类推(对于没有用户的情况，成本为零，但此时没有人能从选择这条零成本路线中受益)。
\end{tcolorbox}

\begin{center}
\includegraphics[max width=0.3\textwidth]{images/bo_d2e5cmbef24c73b03tn0_52_198_201_455_316_0.jpg}
\end{center}
\hspace*{3em} 

Figure 2.1 A congestion network with two nodes \(o\) and \(d\) that define the origin and destination of all users, with two parallel edges between them which have different costs \(c\left( x\right)\) that depend on the usage \(x\) of the respective edge.

\begin{tcolorbox}
图2.1 一个拥塞网络，有两个节点 \(o\) 和 \(d\) ，它们定义了所有用户的起点和终点，节点之间有两条平行边，其成本 \(c\left( x\right)\) 取决于相应边的使用情况 \(x\) 。
\end{tcolorbox}

Suppose the network in Figure 2.1 is used by two users. They can either both use the top edge, or choose a different edge each, or both use the bottom edge. If both use the top edge, both pay cost 2, but if one of them switches to the bottom edge, that user will travel more cheaply and only pay cost 1 on that edge. We say that this situation is not in equilibrium because at least one user can improve her cost by changing her action. Note that we only consider the unilateral deviation of a single user in this scenario. If both would simultaneously switch to the bottom edge, then this edge would be congested with \(x = 2\) and again incur cost 2, which is no improvement to the situation that both use the top edge.

\begin{tcolorbox}
假设图2.1中的网络有两个用户使用。他们可以都使用上边的边，或者各自选择不同的边，或者都使用下边的边。如果他们都使用上边的边，两人都需支付成本2，但如果其中一人切换到下边的边，该用户的出行成本会更低，只需支付成本1。我们说这种情况不是均衡状态，因为至少有一个用户可以通过改变自己的行为来降低成本。注意，在这个场景中我们只考虑单个用户的单方面偏离(deviation)。如果两人同时切换到下边的边，那么这条边会因 \(x = 2\) 而变得拥堵，成本又会变为2，这与两人都使用上边的边的情况相比并没有改善。
\end{tcolorbox}

If the two users use different edges, this is an equilibrium. Namely, the user of the bottom edge currently has cost 1 but would pay cost 2 by switching to the top edge, and clearly she has no incentive to do so. And the user of the top edge, when switching from the top edge, would increase the load on the bottom edge from 1 to 2 with the resulting new cost 2 which is also no improvement for her. This proves the equilibrium property.

\begin{tcolorbox}
如果两个用户使用不同的边，这就是一种均衡状态。具体来说，使用下边边的用户当前成本为1，但如果切换到上边的边，成本将变为2，显然她没有动机这么做。而使用上边边的用户，如果从上边的边切换，会使下边边的负载从1增加到2，新的成本也为2，这对她来说也没有改善。这就证明了均衡性质。
\end{tcolorbox}

Finally, the situation where both users use the bottom edge is also an equilibrium. Both have cost 2, but that is not worse than the cost on the top edge. So this network, with two users, has two equilibria.

\begin{tcolorbox}
最后，两个用户都使用底边的情况也是一种均衡。两者的成本均为2，但这并不比顶边的成本更差。因此，这个有两个用户的网络有两种均衡状态。
\end{tcolorbox}

The underlying behavior of the users in this game is called selfish routing. That is, every user tries to find the route that is least costly for her. Equilibrium is obtained when every user is in a situation where she cannot find a better route on her own.

\begin{tcolorbox}
这个博弈中用户的潜在行为被称为自私路由。也就是说，每个用户都试图找到对自己成本最低的路线。当每个用户都处于无法自行找到更好路线的情况时，就达到了均衡。
\end{tcolorbox}

Is selfish routing always socially optimal, as measured by the average cost per user? In this example, it is better if one user chooses the top edge and the other the bottom edge, with costs 2 and 1, respectively (or 1.5 on average), which is clearly socially optimal. The other equilibrium where both use the bottom edge has the higher cost 2 (paid by both users).

\begin{tcolorbox}
以每个用户的平均成本来衡量，自私路由是否总是符合社会最优呢？在这个例子中，如果一个用户选择顶边，另一个用户选择底边，成本分别为2和1(平均为1.5)，这显然是社会最优的。另一种两个用户都使用底边的均衡状态下，成本更高，为2(两个用户都支付这个成本)。
\end{tcolorbox}

\begin{center}
\includegraphics[max width=0.3\textwidth]{images/bo_d2e5cmbef24c73b03tn0_53_337_332_454_314_0.jpg}
\end{center}
\hspace*{3em} 

Figure 2.2 The same Pigou network as in Figure 2.1, but with a constant cost \(c\left( x\right)  =\) 100 on the top edge. We assume 100 users want to travel from \(o\) to \(d\) .

\begin{tcolorbox}
图2.2 与图2.1相同的庇古网络(Pigou network)，但顶边的固定成本 \(c\left( x\right)  =\) 为100。我们假设100个用户想从 \(o\) 前往 \(d\) 。
\end{tcolorbox}

However, selfish routing may lead to a situation where the average user is worse off than in the socially optimal flow through the network. Figure 2.2 is identical to Figure 2.1 except that the top edge has a constant cost of 100 . It is called the Pigou network. Assume that there are 100 users who want to travel from \(o\) to \(d\) . It is easy to see that there are only two equilibria in this congestion game, namely one where 99 users use the bottom edge and only 1 the top edge, or another where all 100 users use the bottom edge. In all other cases, at most 98 users use the bottom edge, so that any one of the remaining users on the top edge would pay at most 99 by switching to the bottom edge, which improves her cost. In other words, the bottom edge will be "saturated" with users until its usage costs reaches or nearly reaches the cost of the top edge.

\begin{tcolorbox}
然而，自私路由可能导致一种情况，即与通过网络的社会最优流相比，普通用户的情况更糟。图2.2与图2.1相同，只是顶边的固定成本为100。它被称为庇古网络(Pigou network)。假设100个用户想从 \(o\) 前往 \(d\) 。很容易看出，在这个拥塞博弈(congestion game)中只有两种均衡状态，即一种是99个用户使用底边，只有1个用户使用顶边，另一种是100个用户都使用底边。在所有其他情况下，最多98个用户使用底边，因此顶边上剩余的任何一个用户通过切换到底边最多只需支付99，这降低了她的成本。换句话说，底边将被用户“饱和”，直到其使用成本达到或接近顶边的成本。
\end{tcolorbox}

The average cost for each user in these equilibria is found as follows. In the equilibrium where all users use the bottom edge the average cost is 100 . In the equilibrium where one user uses the top edge and 99 the bottom edge the average cost is99.01, computed as \(\frac{99}{100}{99} + \frac{1}{100}{100} = \frac{99}{100}{99} + \frac{1}{100}\left( {{99} + 1}\right)  = \frac{{99} + 1}{100}{99} + \frac{1}{100}\) .

\begin{tcolorbox}
这些均衡状态下每个用户的平均成本计算如下。在所有用户都使用底边的均衡状态下，平均成本为100。在一个用户使用顶边、99个用户使用底边的均衡状态下，平均成本为99.01，计算方式为 \(\frac{99}{100}{99} + \frac{1}{100}{100} = \frac{99}{100}{99} + \frac{1}{100}\left( {{99} + 1}\right)  = \frac{{99} + 1}{100}{99} + \frac{1}{100}\) 。
\end{tcolorbox}

In contrast, if the flow from \(o\) to \(d\) is split equally between the two edges, then 50 users will use the top edge and each pay cost 100, and 50 users will use the bottom edge and each pay cost 50 . So the average cost is \(\frac{50}{100}{100} + \frac{50}{100}{50} = \frac{1}{2}\left( {{100} + {50}}\right)  = {75}\) . We show that this 50-50 split has the optimal average cost. Namely, suppose that \(y\) users use the bottom edge and correspondingly \({100} - y\) use the top edge. Then their average cost is \(\frac{{100} - y}{100}{100} + \frac{y}{100}y\) , which equals \({100} - y + \frac{1}{100}{y}^{2}\) . Assume first that \(y\) can take any real value. Then the derivative of this function of \(y\) is \(- 1 + \frac{2}{100}y\) , which is zero if and only if \(y = {50}\) , which is an integer and gives the unique minimum of the function. Equivalently, \({100} - y + \frac{1}{100}{y}^{2} = {100} + \frac{1}{100}{\left( y - {50}\right) }^{2} - \frac{2500}{100}\) , which has its smallest value 75 for \(y = {50}\) .

\begin{tcolorbox}
相比之下，如果从 \(o\) 到 \(d\) 的流在两条边上平均分配，那么50个用户将使用顶边，每个用户支付成本100，50个用户将使用底边，每个用户支付成本50。因此平均成本为 \(\frac{50}{100}{100} + \frac{50}{100}{50} = \frac{1}{2}\left( {{100} + {50}}\right)  = {75}\) 。我们证明这种50 - 50的分配方式具有最优平均成本。即，假设 \(y\) 个用户使用底边，相应地 \({100} - y\) 个用户使用顶边。那么他们的平均成本为 \(\frac{{100} - y}{100}{100} + \frac{y}{100}y\) ，这等于 \({100} - y + \frac{1}{100}{y}^{2}\) 。首先假设 \(y\) 可以取任何实数值。那么这个关于 \(y\) 的函数的导数为 \(- 1 + \frac{2}{100}y\) ，当且仅当 \(y = {50}\) 时该导数为零， \(y = {50}\) 是一个整数，并且使该函数取得唯一最小值。等价地， \({100} - y + \frac{1}{100}{y}^{2} = {100} + \frac{1}{100}{\left( y - {50}\right) }^{2} - \frac{2500}{100}\) ，当 \(y = {50}\) 时其最小值为75。
\end{tcolorbox}

However, note that the 50-50 split in this example is definitely not an equilibrium, because any user of the top edge would switch to the bottom edge. This would continue until that edge no longer has a lower cost. The socially optimal 50-50 split would have to be enforced in some way (for example by limiting the usage of the bottom edge), or would require new incentives that "change the game", for example a congestion charge.

\begin{tcolorbox}
然而，请注意，在这个例子中 50 - 50 的分流绝对不是一种均衡状态，因为顶边(top edge)的任何用户都会切换到底边(bottom edge)。这种情况会一直持续，直到该边的成本不再更低。社会最优的 50 - 50 分流必须以某种方式来实施(例如，限制底边的使用)，或者需要有新的激励措施来“改变博弈规则”，例如征收拥堵费。
\end{tcolorbox}

\section*{2.3 The Braess Paradox}

\begin{tcolorbox}
\section*{2.3 布雷斯悖论(Braess paradox)}
\end{tcolorbox}

\begin{center}
\includegraphics[max width=0.9\textwidth]{images/bo_d2e5cmbef24c73b03tn0_54_219_362_1213_399_0.jpg}
\end{center}
\hspace*{3em} 

Figure 2.3 The Braess paradox. Each edge has a cost function \(c\left( x\right)\) that depends on the flow \(x\) on that edge (the flow may be different for different edges). The right congestion network has one extra edge \({uv}\) with zero cost compared to the left congestion network. For 100 users, the equilibrium flow in the left network is optimal, whereas in the right network it is worse, due to the extra capacity.

\begin{tcolorbox}
图 2.3 布雷斯悖论(Braess paradox)。每条边(edge)都有一个成本函数 \(c\left( x\right)\) ，该函数取决于该边上的流(flow) \(x\) (不同边上的流可能不同)。与左侧的拥塞网络(congestion network)相比，右侧的拥塞网络多了一条成本为零的边 \({uv}\) 。对于 100 个用户来说，左侧网络中的均衡流(equilibrium flow)是最优的，而在右侧网络中，由于额外的容量，均衡流反而更差。
\end{tcolorbox}

We describe a famous congestion game that shows an unexpected effect of selfish routing. Figure 2.3 shows two congestion networks. In both cases, we assume that there are 100 users who all want to travel from the origin \(o\) to the destination \(d\) . In the left network in Figure 2.3, there are only two routes, across the top and across the bottom. Suppose that \(y\) users take the bottom route via the intermediate node \(u\) , and \({100} - y\) users take the top route via node \(v\) . Both routes use two edges, whose costs are added. With this usage distribution on the edges, the bottom route incurs cost \(y + {100}\) and the top route incurs cost \({100} + \left( {{100} - y}\right)\) . In equilibrium, both costs should be equal or nearly equal in order to not incentivize any user to switch routes. That is, if \(y = {50}\) , then both routes cost 150, and any switch would only lead to a higher cost on the other route. If, for example, \(y = {51}\) , then the bottom route would cost 151 and the top route would cost 149 . Then a user on the bottom route could switch to the top route which has a new, lower cost of 150 . So the only equilibrium is an equal split between the two routes. It is easy to see that this is also the optimal flow through the network.

\begin{tcolorbox}
我们描述一个著名的拥塞博弈(congestion game)，它展示了自私路由的意外效果。图 2.3 展示了两个拥塞网络(congestion network)。在这两种情况下，我们假设共有 100 个用户，他们都想从起点 \(o\) 前往终点 \(d\) 。在图 2.3 的左侧网络中，只有两条路线，即顶部路线和底部路线。假设 \(y\) 个用户选择经过中间节点 \(u\) 的底部路线， \({100} - y\) 个用户选择经过节点 \(v\) 的顶部路线。两条路线都使用两条边，其成本相加。根据边上的这种使用分布，底部路线的成本为 \(y + {100}\) ，顶部路线的成本为 \({100} + \left( {{100} - y}\right)\) 。在均衡状态下，为了不激励任何用户切换路线，两条路线的成本应该相等或近乎相等。也就是说，如果 \(y = {50}\) ，那么两条路线的成本均为 150，并且任何切换只会导致另一条路线的成本更高。例如，如果 \(y = {51}\) ，那么底部路线的成本为 151，顶部路线的成本为 149。那么，底部路线上的一个用户可以切换到顶部路线，此时顶部路线的新成本更低，为 150。因此，唯一的均衡状态是两条路线的用户数量平均分配。很容易看出，这也是通过该网络的最优流(flow)。
\end{tcolorbox}

The right picture in Figure 2.3 has an extra edge \({uv}\) (from \(u\) to \(v\) ). This edge provides a "shortcut" with a new route from \(o\) via \(u\) and \(v\) to \(d\) along the edges \({ou}\) , \({uv},{vd}\) . Suppose that \(z\) users take this new route, that \(y - z\) users for \(y \geq  z\) take the bottom route from \(o\) via \(u\) to \(d\) (so that \(y\) users use the edge \({ou}\) , and \(y - z\) users use the edge \({ud}\) ), and that \({100} - y\) users take the top route from \(o\) to \(v\) to \(d\) . The resulting usage of the edge \({ov}\) is then \({100} - y\) and of the edge \({vd}\) is \({100} - y + z\) . The left diagram in Figure 2.3 corresponds to the case where the edge \({uv}\) is blocked and therefore \(z = 0\) . Now consider the previous equilibrium flow where \(y = {50}\) and \(z = 0\) and suppose that the shortcut \({uv}\) is opened. In this situation, the cost of getting from \(o\) to \(v\) via the shortcut is currently 50, much lower than 100 along the edge \({ov}\) , and similarly getting from \(u\) to \(d\) via \(v\) is also cheaper than via the edge \({ud}\) with cost 100 . This is not in equilibrium. In fact, the only equilibrium flow is given with at most one user on edge \({ov}\) and at most one user on the edge \({ud}\) , and everybody else using the shortcut route, in the worst case with \(y = z = {100}\) where everyone takes the shortcut and nobody uses the edges ov and \({ud}\) with constant costs 100. Then the overall cost of traveling from \(o\) to \(d\) is now \({100} + 0 + {100} = {200}\) , whereas it was 150 beforehand.

\begin{tcolorbox}
图2.3中的右图有一条额外的边 \({uv}\) (从 \(u\) 到 \(v\) )。这条边提供了一条“捷径”，即从 \(o\) 经 \(u\) 和 \(v\) 沿着边 \({ou}\) 、 \({uv},{vd}\) 到 \(d\) 的新路线。假设 \(z\) 个用户选择这条新路线， \(y \geq  z\) 对应的 \(y - z\) 个用户选择从 \(o\) 经 \(u\) 到 \(d\) 的底部路线(这样就有 \(y\) 个用户使用边 \({ou}\) ， \(y - z\) 个用户使用边 \({ud}\) )，并且 \({100} - y\) 个用户选择从 \(o\) 经 \(v\) 到 \(d\) 的顶部路线。那么，边 \({ov}\) 的使用量为 \({100} - y\) ，边 \({vd}\) 的使用量为 \({100} - y + z\) 。图2.3中的左图对应边 \({uv}\) 被阻塞的情况，因此 \(z = 0\) 。现在考虑之前的均衡流，其中 \(y = {50}\) 且 \(z = 0\) ，并假设捷径 \({uv}\) 开通了。在这种情况下，经捷径从 \(o\) 到 \(v\) 的成本目前为50，远低于沿边 \({ov}\) 的100，同样，经 \(v\) 从 \(u\) 到 \(d\) 也比经成本为100的边 \({ud}\) 便宜。这不是均衡状态。实际上，唯一的均衡流是边 \({ov}\) 上最多有一个用户，边 \({ud}\) 上最多有一个用户，其他所有人都使用捷径路线，在最坏的情况下， \(y = z = {100}\) 即每个人都走捷径，没有人使用成本恒定为100的边ov和 \({ud}\) 。那么，现在从 \(o\) 到 \(d\) 的总出行成本为 \({100} + 0 + {100} = {200}\) ，而之前是150。
\end{tcolorbox}

On closer inspection, the new edge \({uv}\) with zero cost creates two Pigou networks as shown on the right, if one identifies the two nodes \(u\) and \(v\) . The first Pigou network goes from \(o\) to destination \(v\) via two routes, namely along edge ov with constant cost 100, and along the edges ou and \({uv}\) with \(\operatorname{cost}c\left( x\right)  = x\) , all of which is incurred on the first edge \({ou}\) . The second Pigou network starts at node \(u\) with destination \(d\) , also via two routes, one along the edge \({ud}\) with cost 100, and the other along edges \({uv}\) and \({vd}\) , of which only the second edge is costly. Like in the Pigou network, neither of the edges ov and \({ud}\) with constant cost is used in selfish routing (except possibly by one user).

\begin{tcolorbox}
仔细观察，如果将两个节点 \(u\) 和 \(v\) 视为同一节点，那么成本为零的新边 \({uv}\) 会形成两个庇古网络(Pigou network)，如右图所示。第一个庇古网络从 \(o\) 经两条路线到达目的地 \(v\) ，即沿成本恒定为100的边ov，以及沿边ou和 \({uv}\) 且 \(\operatorname{cost}c\left( x\right)  = x\) ，所有这些成本都产生在第一条边 \({ou}\) 上。第二个庇古网络从节点 \(u\) 出发，目的地为 \(d\) ，同样有两条路线，一条沿成本为100的边 \({ud}\) ，另一条沿边 \({uv}\) 和 \({vd}\) ，其中只有第二条边有成本。与庇古网络一样，在自私路由中，成本恒定的边ov和 \({ud}\) 都不会被使用(最多可能有一个用户使用)。
\end{tcolorbox}

\begin{center}
\includegraphics[max width=0.5\textwidth]{images/bo_d2e5cmbef24c73b03tn0_55_952_605_639_403_0.jpg}
\end{center}
\hspace*{3em} 

The Pigou network shows that a selfish equilibrium flow may be more costly than a centrally planned socially optimal flow. The added zero-cost edge \({uv}\) in Figure 2.3 goes beyond this observation, because it shows that adding extra capacity to the network can make congestion worse in equilibrium because the new shortcut destroys the balance between the previous two routes. This is called Braess's paradox, after Braess (1968).

\begin{tcolorbox}
庇古网络表明，自私的均衡流可能比中央计划的社会最优流成本更高。图2.3中新增的零成本边 \({uv}\) 进一步说明了这一点，因为它表明在网络中增加额外容量可能会使均衡状态下的拥塞情况恶化，因为新的捷径破坏了之前两条路线之间的平衡。这被称为布雷斯悖论(Braess's paradox)，以布雷斯(Braess，1968)的名字命名。
\end{tcolorbox}

Similarly, closing a road may reduce traffic congestion. This has also been observed in practice. Most evidence on this is anecdotal, but at least some urban planners seem to be aware of this possible effect. In Seoul (Korea), a large expressway over the central Cheonggyecheon riverbed was dismantled in 2002 and replaced by a long recreational park (Hong, Song, and Wu, 2007, p. 235). Traffic improved, although certainly due to many factors in this complex project.

\begin{tcolorbox}
同样，关闭一条道路可能会减少交通拥堵。这在实践中也得到了证实。关于这一点的大多数证据都是轶事性的，但至少一些城市规划者似乎意识到了这种可能的影响。在首尔(韩国)，2002年拆除了清溪川中央河床上的一条大型高速公路，并将其改造成了一个长长的休闲公园(Hong、Song和Wu，2007，第235页)。尽管这个复杂项目的交通改善肯定是由多种因素造成的，但交通状况确实得到了改善。
\end{tcolorbox}

\section*{2.4 Definition of Congestion Games}

\begin{tcolorbox}
\section*{2.4 拥塞博弈(congestion game)的定义}
\end{tcolorbox}

In this section, we give a general definition of congestion games and of the concept of an equilibrium. A congestion network has the following components:

\begin{tcolorbox}
在本节中，我们给出拥塞博弈和均衡概念的一般定义。拥塞网络(congestion network)具有以下组成部分:
\end{tcolorbox}

\begin{itemize}
\item A finite set of nodes.
\end{itemize}

\begin{tcolorbox}
\begin{itemize}
\item 一个有限的节点(node)集合。
\end{itemize}
\end{tcolorbox}

\begin{itemize}
\item A finite collection \(E\) of edges. Each edge \(e\) is an ordered pair, written as \({uv}\) , from some node \(u\) to some node \(v\) , which is graphically drawn as an arrow from \(u\) to \(v\) . Parallel edges (that is, with the same pair \({uv}\) ) are allowed (hence the edges form a "collection" \(E\) rather than a set, which would not allow for such repetitions), as in Figure 2.1.
\end{itemize}

\begin{tcolorbox}
\begin{itemize}
\item 边(edge)的有限集合 \(E\) 。每条边 \(e\) 是一个有序对，记为 \({uv}\) ，从某个节点 \(u\) 到某个节点 \(v\) ，在图形上表示为从 \(u\) 到 \(v\) 的箭头。允许存在平行边(parallel edges)(即具有相同的有序对 \({uv}\) )(因此边构成一个“集合” \(E\) 而非集合，集合不允许有这样的重复)，如图 2.1 所示。
\end{itemize}
\end{tcolorbox}

\begin{itemize}
\item Each edge \(e\) in \(E\) has a cost function \({c}_{e}\) that gives a value \({c}_{e}\left( x\right)\) when there are \(x\) users on edge \(e\) , which describes the same cost to each user for using \(e\) . Each cost function is weakly increasing, that is, \(x \leq  y\) implies \({c}_{e}\left( x\right)  \leq  {c}_{e}\left( y\right)\) .
\end{itemize}

\begin{tcolorbox}
\begin{itemize}
\item 图 \(E\) 中的每条边 \(e\) 都有一个成本函数 \({c}_{e}\) ，当边 \(e\) 上有 \(x\) 个用户时，该函数给出一个值 \({c}_{e}\left( x\right)\) ，它描述了每个用户使用边 \(e\) 的相同成本。每个成本函数都是弱递增的，即 \(x \leq  y\) 意味着 \({c}_{e}\left( x\right)  \leq  {c}_{e}\left( y\right)\) 。
\end{itemize}
\end{tcolorbox}

\begin{itemize}
\item A number \(N\) of users of the network. Each user \(i = 1,\ldots ,N\) has an origin \({o}_{i}\) and destination \({d}_{i}\) , which are two nodes in the network, which may or may not be the same for all users (if they are the same, they are usually called \(o\) and \(d\) as in the above examples).
\end{itemize}

\begin{tcolorbox}
\begin{itemize}
\item 网络的用户数量为 \(N\) 。每个用户 \(i = 1,\ldots ,N\) 都有一个起点 \({o}_{i}\) 和一个终点 \({d}_{i}\) ，它们是网络中的两个节点，对于所有用户来说，这两个节点可能相同也可能不同(如果相同，通常像上述示例中那样称为 \(o\) 和 \(d\) )。
\end{itemize}
\end{tcolorbox}

The underlying structure of nodes and edges is called a directed graph or digraph (where edges are sometimes called "arcs"). In such a digraph, a path \(P\) from \(u\) to \(v\) is a sequence of distinct nodes \({u}_{0},{u}_{1},\ldots ,{u}_{m}\) for \(m \geq  0\) where \({u}_{k}{u}_{k + 1}\) is an edge for \(0 \leq  k < m\) , and \(u = {u}_{0}\) and \(v = {u}_{m}\) . For any such edge \(e = {u}_{k}{u}_{k + 1}\) for \(0 \leq  k < m\) we write \(e \in  P\) . Note that a node may appear at most once in a path. Every user \(i\) chooses a path (which we have earlier also called a "route") from her origin \({o}_{i}\) to her destination \({d}_{i}\) .

\begin{tcolorbox}
节点(node)和边(edge)的底层结构称为有向图(digraph)(边有时也称为“弧”)。在这样的有向图中，从 \(u\) 到 \(v\) 的路径(path) \(P\) 是一个不同节点的序列 \({u}_{0},{u}_{1},\ldots ,{u}_{m}\) ，其中 \(m \geq  0\) ， \({u}_{k}{u}_{k + 1}\) 是一条边， \(0 \leq  k < m\) ，且 \(u = {u}_{0}\) 和 \(v = {u}_{m}\) 。对于任何这样的边 \(e = {u}_{k}{u}_{k + 1}\) ，其中 \(0 \leq  k < m\) ，我们记为 \(e \in  P\) 。注意，一个节点在一条路径中最多出现一次。每个用户 \(i\) 都从她的起点 \({o}_{i}\) 选择一条路径(我们之前也称之为“路线”)到她的终点 \({d}_{i}\) 。
\end{tcolorbox}

\begin{itemize}
\item A strategy of user \(i\) is a path \({P}_{i}\) from \({o}_{i}\) to \({d}_{i}\) .
\end{itemize}

\begin{tcolorbox}
\begin{itemize}
\item 用户 \(i\) 的一个策略是从 \({o}_{i}\) 到 \({d}_{i}\) 的一条路径 \({P}_{i}\) 。
\end{itemize}
\end{tcolorbox}

\begin{itemize}
\item Given a strategy \({P}_{i}\) for each user \(i\) , the load on or flow through an edge \(e\) is defined as \({f}_{e} = \left| \left\{  {i \mid  e \in  {P}_{i}}\right\}  \right|\) , which is the number of chosen paths that contain \(e\) , that is, the number of users on \(e\) . The cost to user \(i\) for her strategy \({P}_{i}\) , given that the other users have chosen their strategies, is then
\end{itemize}

\begin{tcolorbox}
\begin{itemize}
\item 给定每个用户 \(i\) 的一个策略 \({P}_{i}\) ，边 \(e\) 上的负载或流量(flow)定义为 \({f}_{e} = \left| \left\{  {i \mid  e \in  {P}_{i}}\right\}  \right|\) ，即包含边 \(e\) 的所选路径的数量，也就是边 \(e\) 上的用户数量。在其他用户已经选择了他们的策略的情况下，用户 \(i\) 采用她的策略 \({P}_{i}\) 的成本为
\end{itemize}
\end{tcolorbox}

\[
\mathop{\sum }\limits_{{e \in  {P}_{i}}}{c}_{e}\left( {f}_{e}\right) . \tag{2.1}
\]

The congestion game obtained from this network is defined by the strategies of the users (or "players"), and by the resulting cost to each user when all strategies are known. Each user tries, on her own, to minimize her cost. Every edge \(e\) has a certain flow \({f}_{e}\) , which is number of users on \(e\) , and creates the cost \({c}_{e}\left( {f}_{e}\right)\) which is the same for every user of that edge. Users may have different total costs when they use different strategies (that is, different paths; these paths necessarily differ if they do not have the same origin and destination).

\begin{tcolorbox}
从这个网络得到的拥塞博弈(congestion game)由用户(或“玩家”)的策略以及当所有策略都已知时每个用户产生的成本来定义。每个用户都独自尝试最小化她的成本。每条边 \(e\) 都有一定的流量 \({f}_{e}\) ，即边 \(e\) 上的用户数量，并产生成本 \({c}_{e}\left( {f}_{e}\right)\) ，该成本对于使用该边的每个用户都是相同的。当用户采用不同的策略(即不同的路径；如果他们的起点和终点不同，这些路径必然不同)时，他们可能有不同的总成本。
\end{tcolorbox}

Now consider what happens when user \(i\) changes her strategy \({P}_{i}\) to a different strategy \({Q}_{i}\) . Some edges of \({Q}_{i}\) are the same as before, which define the set \({Q}_{i} \cap  {P}_{i}\) . Other edges in \({Q}_{i}\) , those not in \({P}_{i}\) , are new, which define the set \({Q}_{i} \smallsetminus  {P}_{i}\) . Clearly, every edge in \({Q}_{i}\) belongs to exactly one of these two sets. On any edge in \({Q}_{i} \cap  {P}_{i}\) , the flow and hence the cost remains the same as before, but every edge \(e\) in \({Q}_{i} \smallsetminus  {P}_{i}\) has user \(i\) as one additional user and therefore an increased flow of \({f}_{e} + 1\) . We say \({P}_{i}\) is a best response of user \(i\) to the remaining strategies \({P}_{j}\) for \(j \neq  i\) if for every alternative strategy \({Q}_{i}\) of user \(i\)

\begin{tcolorbox}
现在考虑当用户 \(i\) 将其策略 \({P}_{i}\) 改变为不同的策略 \({Q}_{i}\) 时会发生什么。 \({Q}_{i}\) 中的一些边与之前相同，这些边定义了集合 \({Q}_{i} \cap  {P}_{i}\) 。 \({Q}_{i}\) 中其他不在 \({P}_{i}\) 里的边是新的，这些边定义了集合 \({Q}_{i} \smallsetminus  {P}_{i}\) 。显然， \({Q}_{i}\) 中的每条边恰好属于这两个集合之一。在 \({Q}_{i} \cap  {P}_{i}\) 中的任何一条边上，流(flow)以及因此产生的成本(cost)都与之前相同，但 \({Q}_{i} \smallsetminus  {P}_{i}\) 中的每条边 \(e\) 都有用户 \(i\) 作为一个额外的用户，因此流增加了 \({f}_{e} + 1\) 。我们说，如果对于用户 \(i\) 的每个替代策略 \({Q}_{i}\) ， \({P}_{i}\) 是用户 \(i\) 对其余策略 \({P}_{j}\) 的最优反应(best response)
\end{tcolorbox}

\[
\mathop{\sum }\limits_{{e \in  {P}_{i}}}{c}_{e}\left( {f}_{e}\right)  \leq  \mathop{\sum }\limits_{{e \in  {Q}_{i} \cap  {P}_{i}}}{c}_{e}\left( {f}_{e}\right)  + \mathop{\sum }\limits_{{e \in  {Q}_{i} \smallsetminus  {P}_{i}}}{c}_{e}\left( {{f}_{e} + 1}\right) . \tag{2.2}
\]

The left-hand side of this inequality is the cost to user \(i\) when she uses strategy \({P}_{i}\) , and the right-hand side is the possibly changed cost when she deviates to strategy \({Q}_{i}\) . Then \({P}_{i}\) is a best response to the other strategies if no deviation gives rise to a lower cost. (When user \(i\) changes from \({P}_{i}\) to \({Q}_{i}\) , some edges of \({P}_{i}\) may no longer be used in \({Q}_{i}\) and therefore have a reduced flow; this may help other users of those edges, but is irrelevant for user \(i\) .) Note that the concept of best response concerns only deviations of a single user.

\begin{tcolorbox}
这个不等式的左边是用户 \(i\) 使用策略 \({P}_{i}\) 时的成本，右边是她偏离到策略 \({Q}_{i}\) 时可能改变的成本。那么，如果没有偏离能产生更低的成本， \({P}_{i}\) 就是对其他策略的最优反应。(当用户 \(i\) 从 \({P}_{i}\) 改变到 \({Q}_{i}\) 时， \({P}_{i}\) 中的一些边可能在 \({Q}_{i}\) 中不再使用，因此流会减少；这可能对那些边的其他用户有帮助，但与用户 \(i\) 无关。)请注意，最优反应的概念仅涉及单个用户的偏离(deviation)。
\end{tcolorbox}

Definition 2.1. The strategies \({P}_{1},\ldots ,{P}_{N}\) of all \(N\) users define an equilibrium if each strategy is a best response to the other strategies, that is, if (2.2) holds for all \(i = 1,\ldots ,N\) .

\begin{tcolorbox}
定义 2.1。如果每个策略都是对其他策略的最优反应，即如果 (2.2) 对所有 \(i = 1,\ldots ,N\) 都成立，那么所有 \(N\) 个用户的策略 \({P}_{1},\ldots ,{P}_{N}\) 定义了一个均衡(equilibrium)。
\end{tcolorbox}

We have seen examples of equilibria in congestion games in the previous sections. Next, we prove that such an equilibrium always exists.

\begin{tcolorbox}
我们在前面的章节中已经看到了拥塞博弈(congestion game)中均衡的例子。接下来，我们证明这样的均衡总是存在的。
\end{tcolorbox}

\section*{2.5 Existence of Equilibrium in a Congestion Game}

\begin{tcolorbox}
\section*{2.5 拥塞博弈中均衡的存在性}
\end{tcolorbox}

The following is the central theorem of this chapter. It is proved with the help of a potential function \(\Phi\) . The potential function is constructed in such a way that it defines for each edge the increase in cost created by each additional user on the edge, as explained further after the proof.

\begin{tcolorbox}
以下是本章的核心定理。它是借助势函数(potential function) \(\Phi\) 来证明的。势函数的构造方式是，它为每条边定义了该边上每个额外用户所产生的成本增加，证明后会进一步解释。
\end{tcolorbox}

Theorem 2.2. Every congestion game (as obtained from a congestion network) has at least one equilibrium.

\begin{tcolorbox}
定理 2.2。每个拥塞博弈(从拥塞网络(congestion network)得到)至少有一个均衡。
\end{tcolorbox}

Proof. Suppose the \(N\) strategies of the users are \({P}_{1},\ldots ,{P}_{N}\) , which defines a flow \({f}_{e}\) on each edge \(e \in  E\) , namely the number of users \(i\) with \(e \in  {P}_{i}\) . We call this the flow \(f\) induced by these strategies. We now define the following function \(\Phi \left( f\right)\) of this flow by

\begin{tcolorbox}
证明。假设用户的 \(N\) 个策略是 \({P}_{1},\ldots ,{P}_{N}\) ，这在每条边 \(e \in  E\) 上定义了一个流 \({f}_{e}\) ，即具有 \(e \in  {P}_{i}\) 的用户数量 \(i\) 。我们称这为由这些策略诱导的流 \(f\) 。我们现在通过以下方式定义这个流的函数 \(\Phi \left( f\right)\)
\end{tcolorbox}

\[
\Phi \left( f\right)  = \mathop{\sum }\limits_{{e \in  E}}\left( {{c}_{e}\left( 1\right)  + {c}_{e}\left( 2\right)  + \cdots  + {c}_{e}\left( {f}_{e}\right) }\right) . \tag{2.3}
\]

Suppose that user \(i\) changes her path \({P}_{i}\) to \({Q}_{i}\) . We call the resulting new flow \({f}^{{Q}_{i}}\) . We will prove that

\begin{tcolorbox}
假设用户 \(i\) 将其路径 \({P}_{i}\) 改变为 \({Q}_{i}\) 。我们称由此产生的新流为 \({f}^{{Q}_{i}}\) 。我们将证明
\end{tcolorbox}

\[
\Phi \left( {f}^{{Q}_{i}}\right)  - \Phi \left( f\right)  = \mathop{\sum }\limits_{{e \in  {Q}_{i}}}{c}_{e}\left( {f}_{e}^{{Q}_{i}}\right)  - \mathop{\sum }\limits_{{e \in  {P}_{i}}}{c}_{e}\left( {f}_{e}\right) . \tag{2.4}
\]

The right-hand side is the difference in costs to user \(i\) between her strategy \({Q}_{i}\) and her strategy \({P}_{i}\) , according to (2.1). Her cost for the flow \({f}^{{Q}_{i}}\) has been written on the right-hand side of (2.2) in terms of the original flow \(f\) (when she uses \({P}_{i}\) ) as

\begin{tcolorbox}
根据 (2.1)，右边是用户 \(i\) 在其策略 \({Q}_{i}\) 和策略 \({P}_{i}\) 之间的成本差异。她在流 \({f}^{{Q}_{i}}\) 下的成本在 (2.2) 的右边已经根据原始流 \(f\) (当她使用 \({P}_{i}\) 时)写成了
\end{tcolorbox}

\[
\mathop{\sum }\limits_{{e \in  {Q}_{i}}}{c}_{e}\left( {f}_{e}^{{Q}_{i}}\right)  = \mathop{\sum }\limits_{{e \in  {Q}_{i} \cap  {P}_{i}}}{c}_{e}\left( {f}_{e}\right)  + \mathop{\sum }\limits_{{e \in  {Q}_{i} \smallsetminus  {P}_{i}}}{c}_{e}\left( {{f}_{e} + 1}\right) , \tag{2.5}
\]

and her cost for the original flow \(f\) can be expressed similarly, namely as

\begin{tcolorbox}
并且她在原始流 \(f\) 下的成本可以类似地表示，即表示为
\end{tcolorbox}

\[
\mathop{\sum }\limits_{{e \in  {P}_{i}}}{c}_{e}\left( {f}_{e}\right)  = \mathop{\sum }\limits_{{e \in  {P}_{i} \cap  {Q}_{i}}}{c}_{e}\left( {f}_{e}\right)  + \mathop{\sum }\limits_{{e \in  {P}_{i} \smallsetminus  {Q}_{i}}}{c}_{e}\left( {f}_{e}\right) . \tag{2.6}
\]

Hence, the right-hand side of (2.4) is

\begin{tcolorbox}
因此，(2.4) 的右边是
\end{tcolorbox}

\[
\mathop{\sum }\limits_{{e \in  {Q}_{i}}}{c}_{e}\left( {f}_{e}^{{Q}_{i}}\right)  - \mathop{\sum }\limits_{{e \in  {P}_{i}}}{c}_{e}\left( {f}_{e}\right)  = \mathop{\sum }\limits_{{e \in  {Q}_{i} \smallsetminus  {P}_{i}}}{c}_{e}\left( {{f}_{e} + 1}\right)  - \mathop{\sum }\limits_{{e \in  {P}_{i} \smallsetminus  {Q}_{i}}}{c}_{e}\left( {f}_{e}\right) . \tag{2.7}
\]

We claim that because of (2.3) (which is why \(\Phi\) is defined that way), the right-hand side of (2.7) is equal to the left-hand side \(\Phi \left( {f}^{{Q}_{i}}\right)  - \Phi \left( f\right)\) of (2.4). Namely, by changing her path from \({P}_{i}\) to \({Q}_{i}\) , user \(i\) increases the flow on any new edge \(e\) in \({Q}_{i} \smallsetminus  {P}_{i}\) from \({f}_{e}\) to \({f}_{e} + 1\) , and thus adds the term \({c}_{e}\left( {{f}_{e} + 1}\right)\) to the sum in (2.3). Similarly, for any edge \(e\) in \({P}_{i} \smallsetminus  {Q}_{i}\) which is in \({P}_{i}\) but no longer in \({Q}_{i}\) , the flow \({f}^{{Q}_{i}}\left( e\right)\) is reduced from \({f}_{e}\) to \({f}_{e} - 1\) , so that the term \({c}_{e}\left( {f}_{e}\right)\) has to be subtracted from the sum in (2.3). This shows (2.4).

\begin{tcolorbox}
我们断言，由于(2.3)(这就是 \(\Phi\) 如此定义的原因)，(2.7)的右侧等于(2.4)的左侧 \(\Phi \left( {f}^{{Q}_{i}}\right)  - \Phi \left( f\right)\) 。即，用户 \(i\) 将其路径从 \({P}_{i}\) 改变为 \({Q}_{i}\) 时，会使 \({Q}_{i} \smallsetminus  {P}_{i}\) 中任意新边 \(e\) 上的流从 \({f}_{e}\) 增加到 \({f}_{e} + 1\) ，从而在(2.3)的求和项中加上项 \({c}_{e}\left( {{f}_{e} + 1}\right)\) 。类似地，对于 \({P}_{i} \smallsetminus  {Q}_{i}\) 中存在于 \({P}_{i}\) 但不再存在于 \({Q}_{i}\) 中的任意边 \(e\) ，流 \({f}^{{Q}_{i}}\left( e\right)\) 从 \({f}_{e}\) 减少到 \({f}_{e} - 1\) ，因此必须从(2.3)的求和项中减去项 \({c}_{e}\left( {f}_{e}\right)\) 。这就证明了(2.4)。
\end{tcolorbox}

Equation (2.4) states that any change in cost to any user \(i\) when she changes her strategy, and thus the flow, is the same as the change in \(\Phi\) . Consider now a flow that achieves the global minimum of \(\Phi \left( f\right)\) for all possible flows \(f\) that result from users’ strategies \({P}_{1},\ldots ,{P}_{N}\) . There are only finitely many such strategy combinations, so this minimum exists. For such a minimum, no change in flow will reduce \(\Phi \left( f\right)\) any further, and hence no individual player can reduce her cost. Therefore, the players’ strategies for such a minimum of \(\Phi\) define an equilibrium.

\begin{tcolorbox}
方程(2.4)表明，任意用户 \(i\) 改变其策略(从而改变流)时，其成本的任何变化都与 \(\Phi\) 的变化相同。现在考虑一个流，它在所有由用户策略 \({P}_{1},\ldots ,{P}_{N}\) 产生的可能流 \(f\) 中实现了 \(\Phi \left( f\right)\) 的全局最小值。这样的策略组合只有有限多个，所以这个最小值存在。对于这样的最小值，流的任何变化都不会使 \(\Phi \left( f\right)\) 进一步减小，因此没有单个玩家可以降低其成本。因此，玩家对于 \(\Phi\) 这样一个最小值的策略定义了一个均衡。
\end{tcolorbox}

The function \(\Phi\) introduced in the proof of Theorem 2.2 is of special interest. Its definition (2.3) can be thought of being obtained in the following way. Consider an edge \(e\) , and assume that there are no users of the network yet. Then \(c\left( 1\right)\) is the cost when one user travels on edge \(e\) . Next, \(c\left( 2\right)\) is the cost when a second user travels on edge \(e\) , and so on until \(c\left( {f}_{e}\right)\) for the full flow \({f}_{e}\) on edge \(e\) when all users are present (recall that \({f}_{e}\) is the number of users of the edge). However, the sum \({c}_{e}\left( 1\right)  + {c}_{e}\left( 2\right)  + \cdots  + {c}_{e}\left( {f}_{e}\right)\) in (2.3) is not the combined cost to all users of edge \(e\) , because the load \({f}_{e}\) creates the cost \(c\left( {f}_{e}\right)\) for every user of the edge, so their combined cost is \({f}_{e} \cdot  c\left( {f}_{e}\right)\) , which is larger than the sum \({c}_{e}\left( 1\right)  + {c}_{e}\left( 2\right)  + \cdots  + {c}_{e}\left( {f}_{e}\right)\) . Rather, this sum accumulates the additional cost created by each new user of the edge.

\begin{tcolorbox}
定理2.2证明中引入的函数 \(\Phi\) 特别值得关注。它的定义(2.3)可以通过以下方式得到。考虑一条边 \(e\) ，并假设网络中还没有用户。那么 \(c\left( 1\right)\) 是一个用户在边 \(e\) 上通行时的成本。接着， \(c\left( 2\right)\) 是第二个用户在边 \(e\) 上通行时的成本，依此类推，直到所有用户都在时边 \(e\) 上的全流 \({f}_{e}\) 对应的 \(c\left( {f}_{e}\right)\) (回想一下， \({f}_{e}\) 是该边的用户数量)。然而，(2.3)中的和 \({c}_{e}\left( 1\right)  + {c}_{e}\left( 2\right)  + \cdots  + {c}_{e}\left( {f}_{e}\right)\) 并不是边 \(e\) 上所有用户的总费用，因为负载 \({f}_{e}\) 为该边的每个用户产生成本 \(c\left( {f}_{e}\right)\) ，所以他们的总费用是 \({f}_{e} \cdot  c\left( {f}_{e}\right)\) ，这比和 \({c}_{e}\left( 1\right)  + {c}_{e}\left( 2\right)  + \cdots  + {c}_{e}\left( {f}_{e}\right)\) 要大。相反，这个和累积了该边每个新用户所产生的额外成本。
\end{tcolorbox}

Because of this definition, the function \(\Phi\) has the crucial property (2.4) that any change of strategy by a single user incurs a cost change for that user that equals the change of \(\Phi\) . A single function \(\Phi\) achieves this for all users simultaneously. A function \(\Phi\) with this property (2.4) is called a potential function, and a game that has such a potential function is called a potential game. By the preceding proof, any potential game has an equilibrium, which can be found via the minimum of the potential function.

\begin{tcolorbox}
由于这个定义，函数 \(\Phi\) 具有关键性质(2.4)，即单个用户的任何策略改变所导致的该用户的成本变化等于 \(\Phi\) 的变化。单个函数 \(\Phi\) 能同时为所有用户实现这一点。具有这种性质(2.4)的函数 \(\Phi\) 被称为势函数，具有这种势函数的博弈被称为势博弈。根据前面的证明，任何势博弈都存在一个均衡，这个均衡可以通过势函数的最小值找到。
\end{tcolorbox}

The property (2.4) can also be used in the reverse direction to find an equilibrium. Namely, a smaller value of the potential function \(\Phi\) can be found by identifying one user who can reduce her cost by changing her strategy. So we change her strategy accordingly and reduce the value of the potential function. This can be seen as an iterative process where every user (one at a time) adapts her strategy until no user can find further improvements, which then defines an equilibrium. This clearly occurs for any equilibrium. Such an equilibrium may only represent a "local" minimum of the potential function in the sense that it does not allow further reduction by changing a single user's strategy at a time. In the example in Figure 2.1, there are two equilibria that have different values of the potential function.

\begin{tcolorbox}
性质 (2.4) 也可以反向使用来寻找均衡。即，可以通过识别出一个能够通过改变策略降低自身成本的用户，来找到势函数 \(\Phi\) 的更小值。因此，我们相应地改变该用户的策略，从而降低势函数的值。这可以看作一个迭代过程，每个用户(一次一个)调整自己的策略，直到没有用户能找到进一步的改进方法，此时就确定了一个均衡。显然，任何均衡都会出现这种情况。从一次改变单个用户的策略无法进一步降低势函数值的意义上来说，这样的均衡可能仅仅代表势函数的一个“局部”最小值。在图 2.1 的例子中，有两个势函数值不同的均衡。
\end{tcolorbox}

\section*{2.6 Atomic and Splittable Flow, Price of Anarchy}

\begin{tcolorbox}
\section*{2.6 原子流与可拆分流、无政府代价}
\end{tcolorbox}

The presented model of congestion games with finitely many users on a network is only one of many ways of modeling congestion by selfish routing.

\begin{tcolorbox}
所提出的网络上有限多个用户的拥塞博弈模型，只是通过自私路由对拥塞进行建模的众多方法之一。
\end{tcolorbox}

First, the representation of costs in (2.1) is a simplification in the sense that it assumes that user \(i\) creates a "flow" of 1 along its chosen path, and thus contributes to the congestion on every edge on that path. This model does not take any temporal aspects into account, like the fact that the user can be only on one edge at a time but not everywhere on the path at once. Nevertheless, this is not completely unrealistic because everywhere along her chosen route, the user will create congestion at some point in time together with the other users. The flow model is very appropriate for internet routing, for example, with a continuous flow of data packets that consume bandwidth along the taken route.

\begin{tcolorbox}
首先，(2.1) 中成本的表示是一种简化，因为它假设用户 \(i\) 沿着其选择的路径创建了一个大小为 1 的“流”，从而对该路径上的每条边的拥塞都有贡献。这个模型没有考虑任何时间因素，例如用户一次只能在一条边上，而不能同时在路径上的所有位置。然而，这并非完全不现实，因为在用户选择的路线上的任何位置，用户都会在某个时刻与其他用户一起造成拥塞。例如，对于互联网路由，数据分组沿着所选路线持续流动并消耗带宽，这种流模型就非常合适。
\end{tcolorbox}

The specific model that we consider here is called atomic (or non-splittable) flow where single users decide on their paths from their origin to their destination.

\begin{tcolorbox}
我们这里考虑的特定模型称为原子(或不可拆分)流，其中单个用户决定他们从起点到终点的路径。
\end{tcolorbox}

In contrast, splittable flow means that a flow unit is just some "mass" of users of which any parts can be routed along different paths. For example, in Figure 2.1 the flow of 2 from \(o\) to \(d\) could be split into a flow of 0.01 on the top edge and a flow of 1.99 on the bottom edge. This would still not be an equilibrium because any one user from the 0.01 fraction of users on the top edge could improve her cost by moving to the bottom edge (each user is negligibly small). Splittable flows are in some sense simpler because there will be no minor variation of equilibria, such as the distinct equilibria with a flow of either \(y = {100}\) or \(y = {99}\) on the bottom edge in Figure 2.2; the only equilibrium flow would be \(y = {100}\) . The equilibrium condition is then simply that there is no alternative path with currently smaller cost, without taking into account that there is an increased cost by switching to another path, as in the term \({c}_{e}\left( {{f}_{e} + 1}\right)\) in (2.2) where the increase of \({f}_{e}\) by 1 results from that switch. In the splittable flow model, users have negligible size, so there is no such increase. It can be shown that the cost of an equilibrium flow in a splittable flow network is unique (however, there may be several equilibria with that cost, for example with an undetermined flow across two parallel edges with the same constant cost). For splittable flow, the proof of Theorem 2.2 can be amended with a potential function \(\Phi\) defined with an integral \({\int }_{0}^{{f}_{e}}{c}_{e}\left( x\right) {dx}\) that replaces the sum \({c}_{e}\left( 1\right)  + {c}_{e}\left( 2\right)  + \cdots  + {c}_{e}\left( {f}_{e}\right)\) in (2.3); see Roughgarden (2016, section 13.2.2).

\begin{tcolorbox}
相比之下，可拆分流意味着一个流单元只是一些“群体”用户，其任何部分都可以沿着不同的路径进行路由。例如，在图 2.1 中，从 \(o\) 到 \(d\) 的流量为 2 的流可以拆分为顶部边上流量为 0.01 和底部边上流量为 1.99。这仍然不是一个均衡，因为顶部边上 0.01 部分的任何一个用户都可以通过移动到底部边来降低自己的成本(每个用户的影响可以忽略不计)。从某种意义上说，可拆分流更简单，因为不会有均衡的微小变化，例如图 2.2 中底部边上流量为 \(y = {100}\) 或 \(y = {99}\) 的不同均衡；唯一的均衡流量将是 \(y = {100}\) 。那么均衡条件就简单地是，不存在当前成本更低的替代路径，而无需考虑切换到另一条路径会增加成本，就像 (2.2) 中的项 \({c}_{e}\left( {{f}_{e} + 1}\right)\) 那样，其中 \({f}_{e}\) 增加 1 就是由这种切换导致的。在可拆分流模型中，用户的规模可以忽略不计，因此不会有这样的增加。可以证明，可拆分流网络中均衡流的成本是唯一的(然而，可能存在多个具有该成本的均衡，例如在两条具有相同恒定成本的平行边上存在未确定的流量)。对于可拆分流，定理 2.2 的证明可以用一个由积分 \({\int }_{0}^{{f}_{e}}{c}_{e}\left( x\right) {dx}\) 定义的势函数 \(\Phi\) 进行修正，该积分替代了 (2.3) 中的求和 \({c}_{e}\left( 1\right)  + {c}_{e}\left( 2\right)  + \cdots  + {c}_{e}\left( {f}_{e}\right)\) ；参见拉夫加登(2016 年，第 13.2.2 节)。
\end{tcolorbox}

We also mention the Price of Anarchy in this context. This is defined, for a specific congestion network, as the ratio between the worst average cost to a user in an equilibrium and the optimal average cost in the social optimum. The "anarchy" refers to the selfishness in the routing game where each user chooses her optimal route on her own. In the Pigou network, the Price of Anarchy is \(\frac{100}{75} = \frac{4}{3}\) . For splittable flow, it can be shown that this is the worst possible Price of Anarchy when the cost function on each edge is affine, that is, of the form \(c\left( x\right)  = {ax} + b\) for suitable nonnegative real numbers \(a\) and \(b\) (which depend on the edge). (For a proof of this result see Roughgarden and Tardos, 2002, theorem 4.5, or Roughgarden, 2016, chapter 11.) In other words, for affine cost functions, selfish routing is never more than \(1/3\) more costly than what could be achieved with central planning. For non-splittable flow (which has been the topic of this chapter), the Price of Anarchy for affine cost functions can be shown to be at most \(\frac{5}{2}\) (Roughgarden,2016, theorem 12.3). Exercise 2.3 describes a network where this bad ratio is actually achieved, so \(\frac{5}{2}\) as the Price of Anarchy for this model is tight.

\begin{tcolorbox}
我们还会在这个背景下提及无政府代价(Price of Anarchy)。对于特定的拥塞网络(congestion network)，它被定义为均衡(equilibrium)状态下用户的最差平均成本(cost)与社会最优(social optimum)状态下的最优平均成本之比。“无政府”指的是路由博弈中的自私行为，即每个用户自行选择其最优路线。在庇古网络(Pigou network)中，无政府代价为 \(\frac{100}{75} = \frac{4}{3}\) 。对于可拆分流(splittable flow)，可以证明，当每条边(edge)的成本函数为仿射函数，即对于合适的非负实数 \(a\) 和 \(b\) (取决于具体的边)，具有 \(c\left( x\right)  = {ax} + b\) 的形式时，这是可能出现的最差无政府代价。(关于此结果的证明，请参阅Roughgarden和Tardos(2002年)的定理4.5，或Roughgarden(2016年)的第11章。)换句话说，对于仿射成本函数，自私路由的成本永远不会比中央规划所能达到的成本高出 \(1/3\) 以上。对于不可拆分流(本章的主题)，可以证明仿射成本函数的无政府代价至多为 \(\frac{5}{2}\) (Roughgarden，2016年，定理12.3)。练习2.3描述了一个实际上能达到这个糟糕比率的网络，因此在这个模型中 \(\frac{5}{2}\) 作为无政府代价是紧的。
\end{tcolorbox}

\section*{2.7 Further Reading}

\begin{tcolorbox}
\section*{2.7 延伸阅读}
\end{tcolorbox}

We have used material from chapter 13 of the book by Roughgarden (2016), in particular theorem 13.6 and its notation for our proof of Theorem 2.2. Exercise 2.3 is figure 12.3 of that book, rotated so as to display its symmetry. The Pigou network has been described qualitatively by Pigou (1920, p. 194).

\begin{tcolorbox}
我们参考了Roughgarden(2016年)所著书籍第13章的内容，特别是定理13.6及其符号，用于我们对定理2.2的证明。练习2.3是那本书中图12.3旋转后的图形，以展示其对称性。庇古网络已由庇古(1920年，第194页)进行了定性描述。
\end{tcolorbox}

The Braess paradox is due to Braess (1968). Exercise 2.2 discusses the actual network used by Braess. Hong, Song, and Wu (2007) do not mention the Braess paradox in the Cheonggyecheon urban regeneration project, but show some nice "before" and "after" pictures on page 235, and add some Fengshui to this book.

\begin{tcolorbox}
布雷斯悖论(Braess paradox)由布雷斯(1968年)提出。练习2.2讨论了布雷斯实际使用的网络。Hong、Song和Wu(2007年)在清溪川城市再生项目中未提及布雷斯悖论，但在第235页展示了一些不错的“前后”对比图片，并为这本书增添了一些风水元素。
\end{tcolorbox}

The proof of Theorem 2.2 for atomic selfish routing with the help of a potential function \(\Phi\) was originally described by Rosenthal (1973). More general games where an equilibrium is found with the help of a potential function are the potential games studied by Monderer and Shapley (1996).

\begin{tcolorbox}
借助势函数 \(\Phi\) 对原子自私路由的定理2.2进行的证明最初由罗森塔尔(Rosenthal，1973年)提出。在更一般的博弈中，借助势函数找到均衡的博弈是由蒙德勒(Monderer)和夏普利(Shapley，1996年)研究的势博弈。
\end{tcolorbox}

The model of splittable flow is due to Wardrop (1952), and its equilibrium is often called a Wardrop equilibrium.

\begin{tcolorbox}
可拆分流模型由沃德罗普(Wardrop，1952年)提出，其均衡通常被称为用户均衡。
\end{tcolorbox}

\section*{2.8 Exercises for Chapter 2}

\begin{tcolorbox}
\section*{2.8 第2章练习}
\end{tcolorbox}

Exercise 2.1. Consider the following congestion network with origin and destination nodes \(o\) and \(d\) , and an intermediate node \(v\) . There are ten individual users who want to travel from \(o\) to \(d\) .

\begin{tcolorbox}
练习2.1。考虑以下拥塞网络，其起点和终点节点分别为 \(o\) 和 \(d\) ，中间节点为 \(v\) 。有十个个体用户想从 \(o\) 前往 \(d\) 。
\end{tcolorbox}

\begin{center}
\includegraphics[max width=0.4\textwidth]{images/bo_d2e5cmbef24c73b03tn0_61_674_728_579_290_0.jpg}
\end{center}
\hspace*{3em} 

As shown, the cost on the bottom edge is \(c\left( x\right)  = x + 2\) when it has \(x\) users, where \(x + y = {10}\) . The edge from \(o\) to \(v\) has cost function \(c\left( y\right)  = y\) . The top edge from \(v\) to \(d\) has constant cost 3, and the lower edge from \(v\) to \(d\) has cost function \(c\left( z\right)  = z\) if \(z\) of the \(y\) users use that edge, where \(0 \leq  z \leq  y\) .

\begin{tcolorbox}
如图所示，当底部的边有 \(x\) 个用户时，其成本为 \(c\left( x\right)  = x + 2\) ，其中 \(x + y = {10}\) 。从 \(o\) 到 \(v\) 的边的成本函数为 \(c\left( y\right)  = y\) 。从 \(v\) 到 \(d\) 的顶部边的成本恒定为3，从 \(v\) 到 \(d\) 的底部边，如果 \(y\) 个用户中有 \(z\) 个使用该边，其成本函数为 \(c\left( z\right)  = z\) ，其中 \(0 \leq  z \leq  y\) 。
\end{tcolorbox}

(a) Find all equilibria of this congestion game (note that \(x,y\) and \(z\) are integers), and their resulting average cost per user. [Hint: Consider first what \(z\) can be in an equilibrium.]

\begin{tcolorbox}
(a) 找出这个拥塞博弈的所有均衡(注意 \(x,y\) 和 \(z\) 是整数)，以及每个用户的最终平均成本。[提示:首先考虑在均衡状态下 \(z\) 可能的取值。]
\end{tcolorbox}

(b) For the same network, find the equilibrium when \(x,y,z\) can be fractional for a total flow of 10 from \(o\) to \(d\) (so this flow is now splittable). Compare its average cost with the costs of the equilibria in (a).

\begin{tcolorbox}
(b) 对于同一个网络，当从 \(o\) 到 \(d\) 的总流量为10且 \(x,y,z\) 可以取分数值时(即现在的流量是可拆分的)，找出均衡。将其平均成本与(a)中均衡的成本进行比较。
\end{tcolorbox}

Exercise 2.2. Consider the following two congestion networks with the indicated cost functions for a flow \(x\) on the respective edge; for example, \({50} + x\) is short for \(c\left( x\right)  = {50} + x\) . Note that different edges may have different flows \(x\) that depend on the chosen routes by the users.

\begin{tcolorbox}
练习2.2。考虑以下两个拥塞网络，各边的流量 \(x\) 对应的成本函数已标明；例如， \({50} + x\) 是 \(c\left( x\right)  = {50} + x\) 的简写。注意，不同的边可能有不同的流量 \(x\) ，这取决于用户选择的路线。
\end{tcolorbox}

\begin{center}
\includegraphics[max width=0.9\textwidth]{images/bo_d2e5cmbef24c73b03tn0_61_359_1773_1214_393_0.jpg}
\end{center}
\hspace*{3em} 

Suppose there are six users who want to travel from \(o\) to \(d\) .

\begin{tcolorbox}
假设有六个用户想从 \(o\) 前往 \(d\) 。
\end{tcolorbox}

(a) Find the equilibrium flow in the network on the left.

\begin{tcolorbox}
(a) 找出左边网络中的均衡流量。
\end{tcolorbox}

(b) In the network on the right, start from the equilibrium flow in (a) and improve the route of one user at a time until you find an equilibrium flow. Compare the equilibrium costs in both networks.

\begin{tcolorbox}
(b) 在右边的网络中，从(a)中的均衡流量开始，每次改进一个用户的路线，直到找到一个均衡流量。比较两个网络中的均衡成本。
\end{tcolorbox}

Exercise 2.3. Consider the following congestion network with three nodes. As shown, an edge has the cost function \(c\left( x\right)  = x\) if the flow (number of users) on that edge is \(x\) , except for the edges \({vu}\) and \({wu}\) where \(c\left( x\right)  = 0\) . The table on the right shows four users \(i = 1,2,3,4\) with different origins \({o}_{i}\) and destinations \({d}_{i}\) . In this network, each user has two possible routes from her origin to her destination.

\begin{tcolorbox}
练习2.3。考虑以下具有三个节点的拥塞网络。如图所示，如果某条边的流量(用户数量)为 \(x\) ，则该边的成本函数为 \(c\left( x\right)  = x\) ，但边 \({vu}\) 和 \({wu}\) 除外，其中 \(c\left( x\right)  = 0\) 。右边的表格显示了四个用户 \(i = 1,2,3,4\) ，他们有不同的起点 \({o}_{i}\) 和终点 \({d}_{i}\) 。在这个网络中，每个用户从其起点到终点有两条可能的路线。
\end{tcolorbox}

\begin{center}
\includegraphics[max width=0.8\textwidth]{images/bo_d2e5cmbef24c73b03tn0_62_298_740_1051_421_0.jpg}
\end{center}
\hspace*{3em} 

Find two equilibria of this congestion game, one with low cost and one with high cost. Explain why the equilibrium property holds in each case. Compare the costs of the two equilibria. What is the socially optimal (possibly non-equilibrium) flow?

\begin{tcolorbox}
找出这个拥塞博弈的两个均衡，一个是低成本均衡，一个是高成本均衡。解释在每种情况下均衡性质为何成立。比较这两个均衡的成本。社会最优(可能是非均衡)流量是多少？
\end{tcolorbox}

\section*{3}

\section*{Games in Strategic Form}

\begin{tcolorbox}
\section*{策略形式的博弈}
\end{tcolorbox}

In this chapter we start with the systematic development of non-cooperative game theory. Its most basic model is the game in strategic form, the topic of this chapter. The available actions of each player, called strategies, are assumed as given. The players choose their strategies simultaneously and independently, and receive individual payoffs that represent their preferences for strategy profiles (combinations of strategies).

\begin{tcolorbox}
在本章中，我们开始系统地发展非合作博弈论。其最基本的模型是策略形式的博弈，这也是本章的主题。假设每个玩家可用的行动(称为策略)是已知的。玩家们同时且独立地选择他们的策略，并获得代表他们对策略组合(策略的组合)偏好的个体收益。
\end{tcolorbox}

For two players, a game in strategic form is a table where one player chooses a row and the other player a column, with two payoffs in each cell of the table. We present a number of standard games such as the Prisoner's Dilemma, the Quality game, Chicken, the Battle of the Sexes, and Matching Pennies.

\begin{tcolorbox}
对于两个玩家的情况，策略形式的博弈是一个表格，其中一个玩家选择行，另一个玩家选择列，表格的每个单元格中有两个收益。我们介绍一些标准博弈，如囚徒困境、产品质量博弈、胆小鬼博弈、性别战和硬币匹配 (博弈)。
\end{tcolorbox}

The central concept of equilibrium is a profile of strategies that are mutual best responses. A game may have one, several, or no equilibria. As shown in Chapter 6, allowing for mixed (randomized) strategies will ensure that every finite game has an equilibrium, as shown by Nash (1951). An equilibrium without randomization as considered in the present chapter is also known as a pure Nash equilibrium.

\begin{tcolorbox}
均衡的核心概念是一组相互为最优反应(best response)的策略组合。一个博弈可能有一个、多个或没有均衡。正如第6章所示，引入混合(随机化)策略将确保每个有限博弈都有一个均衡，这是由纳什(1951)证明的。本章所考虑的无随机化的均衡也被称为纯纳什均衡(Nash equilibrium)。
\end{tcolorbox}

A strategy dominates another strategy if the player strictly prefers it for any fixed strategies of the other players. Dominated strategies are never played in equilibrium and can therefore be eliminated from the game. If iterated elimination of dominated strategies results in a unique strategy profile, the game is called dominance solvable. We illustrate this with the "Cournot duopoly" of quantity competition. (In Section 4.7 in Chapter 4, we will change this game to a commitment game, known as Stackelberg leadership.)

\begin{tcolorbox}
如果对于其他参与者的任何固定策略，某一参与者严格偏好某一策略而非另一策略，那么该策略就占优于另一策略。被占优策略(dominated strategy)在均衡中永远不会被采用，因此可以从博弈中剔除。如果重复剔除劣势策略(iterated elimination of dominated)得到唯一的策略组合，那么该博弈被称为占优可解(dominance solvable)博弈。我们用产量竞争的“古诺双寡头博弈(Cournot duopoly)”来说明这一点。(在第4章的4.7节中，我们将把这个博弈转变为一个承诺博弈(commitment game)，即斯塔克尔伯格领导博弈。)
\end{tcolorbox}

A strategy weakly dominates another strategy if the player weakly prefers it for any fixed strategies of the other players, and in at least one case strictly prefers it. Eliminating a weakly dominated strategy does not introduce new equilibria, but may lose equilibria, which reduces the understanding of the game. Unless one is interested in finding just one equilibrium of the game, one should therefore not eliminate weakly dominated strategies, nor iterate that process.

\begin{tcolorbox}
如果对于其他参与者的任何固定策略，某一参与者弱偏好某一策略而非另一策略，并且至少在一种情况下严格偏好该策略，那么该策略就弱占优于另一策略。剔除一个弱被占优策略(weakly dominated strategy)不会引入新的均衡，但可能会丢失一些均衡，这会降低对博弈的理解。因此，除非只对找到博弈的一个均衡感兴趣，否则不应该剔除弱被占优策略，也不应该重复这个过程。
\end{tcolorbox}

The final Section 3.8 shows that symmetric \(N\) -player games with two strategies per player always have an equilibrium.

\begin{tcolorbox}
最后的3.8节表明，每个参与者有两种策略的对称(symmetric) \(N\) 人博弈总是存在一个均衡。
\end{tcolorbox}

\section*{3.1 Prerequisites and Learning Outcomes}

\begin{tcolorbox}
\section*{3.1 先决条件和学习成果}
\end{tcolorbox}

This is a first core chapter of the book. The previous Chapter 2 has introduced the concepts of strategies and equilibrium for the special congestion games and is therefore useful but not a prerequisite. We deal mostly with finite sets. You should be thoroughly familiar with Cartesian products \({S}_{1} \times  \cdots  \times  {S}_{N}\) of sets \({S}_{1},\ldots ,{S}_{N}\) . The Cournot game on intervals is analyzed with basic calculus.

\begin{tcolorbox}
这是本书的第一个核心章节。前一章(第2章)介绍了特殊拥塞博弈(congestion game)的策略和均衡概念，因此有帮助但不是先决条件。我们主要处理有限集。你应该熟悉集合 \({S}_{1},\ldots ,{S}_{N}\) 的笛卡尔积 \({S}_{1} \times  \cdots  \times  {S}_{N}\) 。区间上的古诺博弈用基础微积分进行分析。
\end{tcolorbox}

After studying this chapter, you should

\begin{tcolorbox}
学习本章后，你应该
\end{tcolorbox}

\begin{itemize}
\item know the components of a game in strategic form: strategies, strategy profiles, and payoffs;
\end{itemize}

\begin{tcolorbox}
\begin{itemize}
\item 了解策略形式(strategic form)博弈的组成部分:策略、策略组合和收益；
\end{itemize}
\end{tcolorbox}

\begin{itemize}
\item be able to write down and correctly interpret tables for two-player games;
\end{itemize}

\begin{tcolorbox}
\begin{itemize}
\item 能够写出并正确解释两人博弈的表格；
\end{itemize}
\end{tcolorbox}

\begin{itemize}
\item understand the concept of an equilibrium, use this term correctly (and its plural "equilibria"), and know how it relates to strategies, partial strategy profiles, and best responses;
\end{itemize}

\begin{tcolorbox}
\begin{itemize}
\item 理解均衡的概念，正确使用该术语(及其复数形式“equilibria”)，并知道它与策略、部分策略组合(partial strategy profile)和最优反应的关系；
\end{itemize}
\end{tcolorbox}

\begin{itemize}
\item be familiar with common \(2 \times  2\) games such as the Prisoner’s Dilemma or the Stag Hunt, and understand how they differ;
\end{itemize}

\begin{tcolorbox}
\begin{itemize}
\item 熟悉常见的 \(2 \times  2\) 博弈，如囚徒困境(Prisoner’s Dilemma)或猎鹿博弈(Stag Hunt)，并理解它们的区别；
\end{itemize}
\end{tcolorbox}

\begin{itemize}
\item know the difference between dominance and weak dominance and why dominated strategies can be eliminated in a complete equilibrium analysis but weakly dominated strategies cannot;
\end{itemize}

\begin{tcolorbox}
\begin{itemize}
\item 了解占优和弱占优的区别，以及为什么在完整的均衡分析中可以剔除被占优策略，但不能剔除弱被占优策略；
\end{itemize}
\end{tcolorbox}

\begin{itemize}
\item know the Cournot quantity game, also when strategy sets are real intervals;
\end{itemize}

\begin{tcolorbox}
\begin{itemize}
\item 了解古诺产量博弈，包括策略集为实数区间的情况；
\end{itemize}
\end{tcolorbox}

\begin{itemize}
\item understand symmetry in games with two players and more than two players.
\end{itemize}

\begin{tcolorbox}
\begin{itemize}
\item 理解两人及多人博弈中的对称性。
\end{itemize}
\end{tcolorbox}

\section*{3.2 Games in Strategic Form}

\begin{tcolorbox}
\section*{3.2 策略形式的博弈}
\end{tcolorbox}

A game in strategic form is the fundamental model of non-cooperative game theory. The game has \(N\) players, \(N \geq  1\) , and each player \(i = 1,\ldots ,N\) has a nonempty set \({S}_{i}\) of strategies. If each player \(i\) chooses a strategy \({s}_{i}\) from \({S}_{i}\) , the resulting \(N\) -tuple \(s = \left( {{s}_{1},\ldots ,{s}_{n}}\right)\) is called a strategy profile. The game is specified by assigning to each strategy profile \(s\) a real-valued payoff \({u}_{i}\left( s\right)\) to each player \(i\) .

\begin{tcolorbox}
策略形式的博弈是非合作博弈论的基本模型。该博弈有 \(N\) 个参与者， \(N \geq  1\) ，并且每个参与者 \(i = 1,\ldots ,N\) 都有一个非空的策略集 \({S}_{i}\) 。如果每个参与者 \(i\) 从 \({S}_{i}\) 中选择一个策略 \({s}_{i}\) ，那么得到的 \(N\) 元组 \(s = \left( {{s}_{1},\ldots ,{s}_{n}}\right)\) 被称为一个策略组合。该博弈通过为每个策略组合 \(s\) 给每个参与者 \(i\) 分配一个实值收益 \({u}_{i}\left( s\right)\) 来确定。
\end{tcolorbox}

The payoffs represent each player’s preference. For two strategy profiles \(s\) and \(\widehat{s}\) , player \(i\) strictly prefers \(s\) to \(\widehat{s}\) if \({u}_{i}\left( s\right)  > {u}_{i}\left( \widehat{s}\right)\) , and is indifferent between \(s\) and \(\widehat{s}\) if \({u}_{i}\left( s\right)  = {u}_{i}\left( \widehat{s}\right)\) . If \({u}_{i}\left( s\right)  \geq  {u}_{i}\left( \widehat{s}\right)\) then player \(i\) weakly prefers \(s\) to \(\widehat{s}\) . Player \(i\) is only interested in maximizing his own payoff, and not interested in the payoffs to other players (other than in trying to anticipate their actions); any "social" concern a player has about a particular outcome has to be (and could be) built into his own payoff. All players know the available strategies and payoffs of the other players, and know that they know them, etc. (that is, the game is "common knowledge").

\begin{tcolorbox}
收益代表了每个玩家的偏好。对于两个策略组合 \(s\) 和 \(\widehat{s}\) ，如果 \({u}_{i}\left( s\right)  > {u}_{i}\left( \widehat{s}\right)\) ，玩家 \(i\) 严格偏好 \(s\) 而非 \(\widehat{s}\) ；如果 \({u}_{i}\left( s\right)  = {u}_{i}\left( \widehat{s}\right)\) ，玩家 \(i\) 对 \(s\) 和 \(\widehat{s}\) 无差异。如果 \({u}_{i}\left( s\right)  \geq  {u}_{i}\left( \widehat{s}\right)\) ，则玩家 \(i\) 弱偏好 \(s\) 而非 \(\widehat{s}\) 。玩家 \(i\) 只关心最大化自己的收益，而不关心其他玩家的收益(除了试图预测他们的行动)；玩家对特定结果的任何“社会”考量都必须(也可以)纳入他自己的收益中。所有玩家都知道其他玩家可用的策略和收益，并且知道他们都知道这些，以此类推(也就是说，该博弈是“共同知识”)。
\end{tcolorbox}

The game is played as follows: The players choose their strategies simultaneously (without knowing what the other players choose), and receive their respective payoffs for the resulting strategy profile. They cannot enter into any binding agreements about what they should play (which is why the game is called "noncooperative"). Furthermore, the game is assumed to be played only once, and therefore also called a one-shot game. Playing the same game (or a varying game) many times leads to the much more advanced theory of repeated games.

\begin{tcolorbox}
该博弈的进行方式如下:玩家同时选择他们的策略(不知道其他玩家的选择)，并根据所产生的策略组合获得各自的收益。他们不能就应该采取的行动达成任何有约束力的协议(这就是为什么该博弈被称为“非合作博弈”)。此外，假设该博弈只进行一次，因此也被称为一次性博弈。多次进行相同的博弈(或不同的博弈)会引出更为高深的重复博弈理论。
\end{tcolorbox}

Much of this book is concerned with two-player games. The players are typically named 1 and 2 or I and II. Then the strategic form is conveniently represented by a table. The rows of the table represent the strategies of player I (also called the row player), and the columns represent the strategies of player II (the column player). A strategy profile is a strategy pair, that is, a row and a column, with a corresponding cell of the table that contains two payoffs, one for player I and the other for player II.

\begin{tcolorbox}
本书大部分内容涉及两人博弈。玩家通常被命名为 1 和 2 或 I 和 II。然后，策略形式可以方便地用一个表格来表示。表格的行代表玩家 I(也称为行玩家)的策略，列代表玩家 II(列玩家)的策略。一个策略组合是一个策略对，即一行和一列，表格中对应的单元格包含两个收益，一个是玩家 I 的，另一个是玩家 II 的。
\end{tcolorbox}

If \(m\) and \(n\) are positive integers, then an \(m \times  n\) game is a two-player game in strategic form with \(m\) strategies for player I (the rows of the table) and \(n\) strategies for player II (the columns of the table).

\begin{tcolorbox}
如果 \(m\) 和 \(n\) 是正整数，那么一个 \(m \times  n\) 博弈是一个策略形式的两人博弈，玩家 I 有 \(m\) 个策略(表格的行)，玩家 II 有 \(n\) 个策略(表格的列)。
\end{tcolorbox}

Many interesting game-theoretic observations apply already to \(2 \times  2\) games. Figure 3.1(a) shows the famous Prisoner's Dilemma game. Each player has two strategies, called \(C\) and \(D\) , which stand for "cooperate" and "defect". The payoffs are as follows: If both players choose \(C\) , then both get payoff 2. If both players choose \(D\) , then both get payoff 1 . If the two players choose different strategies, then the player who chooses \(C\) gets payoff 0 and the player who chooses \(D\) gets the highest possible payoff 3.

\begin{tcolorbox}
许多有趣的博弈论观察结果已经适用于 \(2 \times  2\) 博弈。图 3.1(a) 展示了著名的囚徒困境博弈。每个玩家有两个策略，分别称为 \(C\) 和 \(D\) ，分别代表“合作”和“背叛”。收益情况如下:如果两个玩家都选择 \(C\) ，则两人都获得收益 2。如果两个玩家都选择 \(D\) ，则两人都获得收益 1。如果两个玩家选择不同的策略，那么选择 \(C\) 的玩家获得收益 0，选择 \(D\) 的玩家获得最高可能的收益 3。
\end{tcolorbox}

(a)

\begin{center}
\includegraphics[max width=0.8\textwidth]{images/bo_d2e5cmbef24c73b03tn0_66_289_1504_1062_362_0.jpg}
\end{center}
\hspace*{3em} 

Figure 3.1 Prisoner's Dilemma game (a), its symmetry shown by reflection along the dotted diagonal line in (b).

\begin{tcolorbox}
图 3.1 囚徒困境博弈 (a)，其对称性在 (b) 中通过沿虚线对角线反射显示。
\end{tcolorbox}

In each cell of the payoff table, the payoff to the row player is shown in the lower left corner and the payoff to the column player in the upper right corner. There are several very good reasons for this display:

\begin{tcolorbox}
在收益表的每个单元格中，行玩家的收益显示在左下角，列玩家的收益显示在右上角。这种显示方式有几个很好的理由:
\end{tcolorbox}

\begin{itemize}
\item There is no ambiguity to whom the payoff belongs. The payoff close to the left is near the strategy labels for the rows, the payoff close to the top is near the labels for the columns; this subdivision is also shown at the top left of the table that shows the two player names, here I and II.
\end{itemize}

\begin{tcolorbox}
\begin{itemize}
\item 收益归属明确。靠近左边的收益靠近行的策略标签，靠近顶部的收益靠近列的标签；这种划分也在表格的左上角显示，那里显示了两个玩家的名称，这里是 I 和 II。
\end{itemize}
\end{tcolorbox}

\begin{itemize}
\item It is immediately apparent if the game is symmetric (that is, it stays the same game when exchanging the players), by reflecting the table along the diagonal line which exchanges the players, their strategies (rows and columns), and the payoffs in each cell, as shown in Figure 3.1(b). This reflection works because the rows are listed from top to bottom and the columns from left to right.
\end{itemize}

\begin{tcolorbox}
\begin{itemize}
\item 通过沿对角线反射表格(交换玩家、他们的策略(行和列)以及每个单元格中的收益)，可以立即看出该博弈是否是对称的(即，交换玩家后博弈保持不变)，如图 3.1(b) 所示。这种反射可行是因为行是从上到下列出的，列是从左到右列出的。
\end{itemize}
\end{tcolorbox}

\begin{itemize}
\item An \(m \times  n\) game is specified by two \(m \times  n\) matrices \(A\) and \(B\) which contain the payoffs to players I and II. Because the payoffs are staggered in each cell, these two matrices are easy to identify separately so that one can focus on the payoffs of a single player.
\end{itemize}

\begin{tcolorbox}
\begin{itemize}
\item 一个 \(m \times  n\) 博弈由两个 \(m \times  n\) 矩阵 \(A\) 和 \(B\) 确定，它们分别包含玩家 I 和玩家 II 的收益。由于每个单元格中的收益是错开的，这两个矩阵很容易分别识别，这样就可以专注于单个玩家的收益。
\end{itemize}
\end{tcolorbox}

The only unfamiliar aspect of putting the row player's payoff at the bottom left of a cell is that it is below the column player's payoff, where the usual convention is to consider rows first and columns second. However, the above reasons do not depend on this convention and treat rows and columns symmetrically. Another common display (usually for typesetting reasons) that uses the convention "rows first, columns second" is

\begin{tcolorbox}
将行参与者的收益置于单元格左下角唯一不熟悉的方面在于，它位于列参与者的收益下方，而通常的惯例是先考虑行，再考虑列。然而，上述原因并不依赖于这一惯例，而是对称地(symmetric)对待行和列。另一种采用“先考虑行，再考虑列”这一惯例的常见显示方式(通常出于排版原因)是
\end{tcolorbox}

(3.1)where the payoffs are not staggered and therefore harder to separate, and where the symmetry of a game is less apparent. We will always use the staggered payoff display as in Figure 3.1.

\begin{tcolorbox}
在这种显示方式中，收益没有错开，因此更难区分，而且博弈的对称性也不太明显。我们将始终使用如图3.1所示的错开收益显示方式。
\end{tcolorbox}

\begin{center}
\adjustbox{max width=\textwidth}{
\begin{tabular}{|c|c|c|}
\hline
\multirow{3}{*}{AI \(C\)  \(D\)} & \(C\) & \(D\) \\
\cline{2-3}
 & 2, 2 & 0,3 \\
\cline{2-3}
 & 3,0 & 1,1 \\
\cline{1-3}
\hline
\end{tabular}
}
\end{center}

\begin{tcolorbox}

\begin{center}
\adjustbox{max width=\textwidth}{
\begin{tabular}{|c|c|c|}
\hline
\multirow{3}{*}{人工智能(AI) \(C\)  \(D\)} & \(C\) & \(D\) \\
\cline{2-3}
 & 2, 2 & 0,3 \\
\cline{2-3}
 & 3,0 & 1,1 \\
\cline{1-3}
\hline
\end{tabular}
}
\end{center}

\end{tcolorbox}

The story behind the name "Prisoner's Dilemma" is that of two prisoners held suspect of a serious crime. There is no judicial evidence for this crime except if one of the prisoners testifies against the other. If one of them testifies (choosing " \(D\) "), he will be rewarded with immunity from prosecution (payoff 3), whereas the other will serve a long prison sentence (payoff 0 ). If both testify, their punishment will be less severe (payoff 1 for each). If they both "cooperate" with each other by not testifying at all (choosing " \(C\) "), they will only be imprisoned briefly for some minor charge that can be held against them (payoff 2 for each). The "defection" from that mutually beneficial outcome is to testify, which gives a higher payoff no matter what the other prisoner does. However, the resulting payoff is lower to both. This constitutes their "dilemma".

\begin{tcolorbox}
“囚徒困境”这一名称背后的故事涉及两名因涉嫌严重犯罪而被关押的囚犯。除了其中一名囚犯指证另一名囚犯之外，没有其他司法证据证明这一罪行。如果其中一人指证(选择“ \(D\) ”)，他将获得免于起诉的奖励(收益为3)，而另一人将面临长期监禁(收益为0)。如果两人都指证，他们的刑罚会较轻(每人收益为1)。如果他们通过完全不指证而相互“合作”(选择“ \(C\) ”)，他们只会因一些可指控的轻微罪行而被短期监禁(每人收益为2)。从这种互利结果的“背叛”行为就是指证，无论另一名囚犯怎么做，指证都会带来更高的收益。然而，最终两人的收益都会降低。这就构成了他们的“困境”。
\end{tcolorbox}

In the Prisoner's Dilemma, both players always have an incentive to defect, so that the strategy pair(D, D)will be the predicted outcome of playing the game, even though both players are worse off than if they played(C, C). The reason is that the non-cooperative model assumes that the players cannot enforce an agreement to play(C, C)(if they could, for example by signing a contract of some sort, this would have to be a separate strategy in the game). The Prisoner's Dilemma shows that individual incentives may jeopardize a common good. For example, in an arms race between two nations, \(C\) and \(D\) may represent the choices between restricted and unlimited arms build-up. The Prisoner's Dilemma can also be extended to many players. For example, it may be in every fisher's interest not to overfish so that fish stocks do not decline, by introducing a fishing quota that limits each fisher's catch. If adopted collectively, this would benefit all. However, without an enforcement of fishing quotas, each single fisher who exceeds his quota will barely contribute to overfishing but strongly benefit for himself. Collectively, they destroy their fish stock this way. This is also called the "tragedy of the commons". It applies similarly to many types of actions that collectively harm the environment (and each participant) but are individually beneficial.

\begin{tcolorbox}
在囚徒困境中，两名玩家总是有背叛的动机，因此策略组合(D，D)将是这场博弈的预期结果，尽管如果他们选择(C，C)，两名玩家的情况都会更好。原因在于非合作模型假设玩家无法执行达成的选择(C，C)的协议(例如，如果他们可以通过签署某种合同来达成协议，那么这在博弈中必须是一种单独的策略)。囚徒困境表明个人动机可能会损害公共利益。例如，在两个国家的军备竞赛中， \(C\) 和 \(D\) 可能代表限制军备建设和无限制军备建设之间的选择。囚徒困境也可以扩展到多个玩家。例如，通过引入限制每个渔民捕捞量的捕捞配额，不进行过度捕捞以避免鱼类资源减少可能符合每个渔民的利益。如果集体采纳这一措施，将对所有人都有益。然而，如果不执行捕捞配额，每个超过配额的渔民对过度捕捞的影响微乎其微，但自己却能获得巨大利益。从总体上看，他们这样做会破坏鱼类资源。这也被称为“公地悲剧”。这同样适用于许多对环境(以及每个参与者)造成集体损害但对个人有益的行为。
\end{tcolorbox}

\section*{3.3 Best Responses and Equilibrium}

\begin{tcolorbox}
\section*{3.3 最优反应与均衡}
\end{tcolorbox}

Consider the game in Figure 3.2(a). This is the Prisoner's Dilemma game, with further graphical information added to help the analysis of the game. Namely, certain payoffs are surrounded by boxes that indicate that they are best-response payoffs.

\begin{tcolorbox}
考虑图3.2(a)中的博弈。这是囚徒困境博弈，并添加了更多图形信息以帮助分析该博弈。具体来说，某些收益被方框包围，表明它们是最优反应收益。
\end{tcolorbox}

A best response is a strategy of a player that is optimal for that player assuming it is known what all other players do. In a two-player game, there is only one other player. Hence, for the row player, a best response is found for each column. In the Prisoner’s Dilemma game, the best response of the row player against column \(C\) is row \(D\) because his payoff when playing \(D\) is 3 rather than 2 when playing \(C\) (so in the cell for row \(D\) and column \(C\) , the payoff 3 to player \(\mathrm{I}\) is put in a box), and the best response against column \(D\) is also row \(D\) (with payoff 1, in a box, which is larger than payoff 0 for row \(C\) ). Similarly, the best response of the column player against row \(C\) is column \(D\) , and against row \(D\) is column \(D\) , as shown by the boxes around the payoffs for the column player. Because the game is symmetric, the best-response payoffs are also symmetric.

\begin{tcolorbox}
最优反应是指在已知其他所有玩家的行动的情况下，对某个玩家而言最优的策略。在两人博弈中，只有另一名玩家。因此，对于行玩家来说，要为每一列找到最优反应。在囚徒困境博弈中，行玩家针对列 \(C\) 的最优反应是行 \(D\) ，因为他选择 \(D\) 时的收益是3，而选择 \(C\) 时的收益是2(因此在行 \(D\) 和列 \(C\) 的单元格中，玩家 \(\mathrm{I}\) 的收益3被放在方框中)，并且针对列 \(D\) 的最优反应也是行 \(D\) (收益为1，放在方框中，大于行 \(C\) 的收益0)。同样，列玩家针对行 \(C\) 的最优反应是列 \(D\) ，针对行 \(D\) 的最优反应也是列 \(D\) ，列玩家的收益周围的方框显示了这一点。由于该博弈是对称的，最优反应收益也是对称的。
\end{tcolorbox}

The game in Figure 3.2(b) has payoffs that are nearly identical to those of the Prisoner's Dilemma game, except for the payoff to player II for the top right cell, which is changed from 3 to 1 . Because the game is not symmetric, we name the strategies of the players differently, here \(T\) and \(B\) for the row player (for "top" and "bottom") and \(l\) and \(r\) for the column player ("left" and "right"). To distinguish them more easily, we often write the strategies of player I in upper case and those of player II in lower case.

\begin{tcolorbox}
图3.2(b)中的博弈收益与囚徒困境博弈的收益几乎相同，除了右上角单元格中玩家II的收益从3变为1。由于该博弈不是对称的，我们对玩家的策略进行了不同的命名，这里行玩家的策略是 \(T\) 和 \(B\) (分别表示“上”和“下”)，列玩家的策略是 \(l\) 和 \(r\) (分别表示“左”和“右”)。为了更方便地区分它们，我们通常将玩家I的策略用大写字母表示，将玩家II的策略用小写字母表示。
\end{tcolorbox}

\begin{center}
\includegraphics[max width=0.8\textwidth]{images/bo_d2e5cmbef24c73b03tn0_69_383_194_1101_363_0.jpg}
\end{center}
\hspace*{3em} 

Figure 3.2 Prisoner's Dilemma game (a) with best-response payoffs shown by boxes, and the same done for the Quality game (b).

\begin{tcolorbox}
图3.2 囚徒困境博弈(a)，其中最优反应收益用方框表示，产品质量博弈(b)也是如此。
\end{tcolorbox}

Figure 3.2(b) is called the Quality game. It may describe the situation of player I as a restaurant owner, who can provide food of good quality (strategy \(T\) ) or bad quality(B), and a potential customer, player II, who may decide to eat there(l)or not(r). The customer prefers \(l\) to \(r\) only if the quality is good. However, whatever player II does, player I is better off by choosing \(B\) . These preferences are shown by the best-response boxes. The best response of player II to row \(T\) is different from her best response to row \(B\) . The best response of player I to either column is always row \(B\) .

\begin{tcolorbox}
图3.2(b)被称为产品质量博弈。它可以描述参与者I作为餐厅老板的情况，老板可以提供高质量的食物(策略 \(T\) )或低质量的食物(B)，而潜在客户，即参与者II，可以决定在那里就餐(l)或不在那里就餐(r)。只有当食物质量好时，客户才会偏好 \(l\) 而非 \(r\) 。然而，无论参与者II怎么做，参与者I选择 \(B\) 都会更好。这些偏好通过最优反应框显示出来。参与者II对行 \(T\) 的最优反应与她对行 \(B\) 的最优反应不同。参与者I对任意列的最优反应始终是行 \(B\) 。
\end{tcolorbox}

In both games in Figure 3.2, consider the bottom right cell of the table, which corresponds to the strategy pair(D, D)in the Prisoner’s Dilemma game and to (B, r)in the Quality game. In each game, this is the only cell that has both payoffs surrounded by a box. This means that the two strategies that define that cell are mutual best responses. Such a strategy pair is called an equilibrium, which is the central solution concept for non-cooperative games.

\begin{tcolorbox}
在图3.2的两个博弈中，考虑表格的右下角单元格，在囚徒困境博弈中对应策略组合(D，D)，在产品质量博弈中对应(B，r)。在每个博弈中，这个单元格是唯一两个收益都被框起来的单元格。这意味着定义该单元格的两个策略是相互的最优反应。这样的策略组合被称为均衡，它是非合作博弈的核心解概念。
\end{tcolorbox}

We define the concepts of best response and equilibrium for a general \(N\) -player game.

\begin{tcolorbox}
我们为一般的 \(N\) 人博弈定义最优反应和均衡的概念。
\end{tcolorbox}

Definition 3.1. Consider a game in strategic form with \(N\) players where player \(i = 1,\ldots ,N\) has the strategy set \({S}_{i}\) and receives payoff \({u}_{i}\left( s\right)\) for each strategy profile \(s = \left( {{s}_{1},\ldots ,{s}_{N}}\right)\) in \({S}_{1} \times  \cdots  \times  {S}_{N}\) . For player \(i\) , consider a partial profile

\begin{tcolorbox}
定义3.1。考虑一个具有 \(N\) 个参与者的策略形式博弈，其中参与者 \(i = 1,\ldots ,N\) 的策略集为 \({S}_{i}\) ，并且对于 \({S}_{1} \times  \cdots  \times  {S}_{N}\) 中的每个策略组合 \(s = \left( {{s}_{1},\ldots ,{s}_{N}}\right)\) ，参与者 \(i = 1,\ldots ,N\) 获得收益 \({u}_{i}\left( s\right)\) 。对于参与者 \(i\) ，考虑其他参与者的部分策略组合
\end{tcolorbox}

\[
{s}_{-i} = \left( {{s}_{1},\ldots ,{s}_{i - 1},{s}_{i + 1},\ldots ,{s}_{N}}\right)  \tag{3.2}
\]

of strategies of the other players, given by their strategies \({s}_{j}\) for \(j \neq  i\) . For any strategy \({s}_{i}\) in \({S}_{i}\) , this partial profile is extended to a full strategy profile denoted by \(s = \left( {{s}_{i},{s}_{-i}}\right)\) . Then \({s}_{i}\) is a best response to \({s}_{-i}\) if

\begin{tcolorbox}
由其他参与者在 \(j \neq  i\) 时的策略 \({s}_{j}\) 给出。对于 \({S}_{i}\) 中的任何策略 \({s}_{i}\) ，这个部分策略组合可以扩展为一个完整的策略组合，记为 \(s = \left( {{s}_{i},{s}_{-i}}\right)\) 。那么，如果满足以下条件，则 \({s}_{i}\) 是对 \({s}_{-i}\) 的最优反应
\end{tcolorbox}

\[
{u}_{i}\left( s\right)  = {u}_{i}\left( {{s}_{i},{s}_{-i}}\right)  \geq  {u}_{i}\left( {{\bar{s}}_{i},{s}_{-i}}\right) \;\text{ for all }{\bar{s}}_{i} \in  {S}_{i}. \tag{3.3}
\]

An equilibrium of the game is a strategy profile \(s\) where each strategy is a best response to the other strategies in that profile, that is, (3.3) holds for all players \(i = 1,\ldots ,N\) .

\begin{tcolorbox}
博弈的一个均衡是一个策略组合 \(s\) ，其中每个策略都是对该组合中其他策略的最优反应，即，对于所有参与者 \(i = 1,\ldots ,N\) ，式(3.3)都成立。
\end{tcolorbox}

The notation (3.2) for a partial profile \({s}_{-i}\) of strategies of the players other than player \(i\) is very common. For a two-player game \(\left( {N = 2}\right)\) , it is just a strategy of the other player. The notation \(\left( {{s}_{i},{s}_{-i}}\right)\) creates a profile \(s\) for all players, where \({s}_{i}\) is assumed to inserted at the \(i\) th position (as the strategy of player \(i\) ) to define the strategy profile \(s\) . The condition (3.3) states that \({s}_{i}\) is a best response to \({s}_{-i}\) because player \(i\) cannot get a better payoff by changing to a different strategy \({\bar{s}}_{i}\) . That payoff may stay the same (in which case \({\bar{s}}_{i}\) is also a best response to \({s}_{-i}\) ). Moreover, only unilateral changes by a single player are considered, not joint changes by several players, as for example, a change from(D, D)to(C, C)in the Prisoner’s Dilemma. An equilibrium, as a strategy profile of mutual best responses, is therefore a profile where no player can profitably gain by unilaterally changing his strategy, that is, with the remaining strategies kept fixed.

\begin{tcolorbox}
除参与者 \(i\) 之外的其他参与者的部分策略组合 \({s}_{-i}\) 的表示法(3.2)非常常见。对于两人博弈 \(\left( {N = 2}\right)\) ，它只是另一个参与者的一个策略。表示法 \(\left( {{s}_{i},{s}_{-i}}\right)\) 为所有参与者创建了一个策略组合 \(s\) ，其中假设 \({s}_{i}\) 被插入到第 \(i\) 个位置(作为参与者 \(i\) 的策略)来定义策略组合 \(s\) 。条件(3.3)表明 \({s}_{i}\) 是对 \({s}_{-i}\) 的最优反应，因为参与者 \(i\) 通过改变为不同的策略 \({\bar{s}}_{i}\) 无法获得更好的收益。收益可能保持不变(在这种情况下， \({\bar{s}}_{i}\) 也是对 \({s}_{-i}\) 的最优反应)。此外，只考虑单个参与者的单方面改变，而不考虑多个参与者的联合改变，例如，在囚徒困境中从(D，D)变为(C，C)。因此，作为相互最优反应的策略组合，均衡是一个没有参与者可以通过单方面改变其策略而获利的组合，即，其余策略保持不变。
\end{tcolorbox}

Some notes on terminology: An equilibrium is often called Nash equilibrium, named after John Nash (1951) who proved that an "equilibrium point" always exists for finite strategy sets that are extended to include randomized or "mixed" strategies (see Chapter 6). In distinction to such a "mixed equilibrium", an equilibrium according to Definition 3.1 is often called a pure equilibrium where players choose their strategies deterministically. The plural of "equilibrium" is equilibria.

\begin{tcolorbox}
关于术语的一些说明:均衡通常被称为纳什均衡(Nash equilibrium)，以约翰·纳什(John Nash，1951)的名字命名，他证明了对于扩展到包括随机或“混合”策略的有限策略集(见第6章)，“均衡点”总是存在的。与这种“混合均衡(mixed equilibrium)”不同，根据定义3.1的均衡通常被称为纯策略均衡(pure equilibrium)，其中参与者确定性地选择他们的策略。“equilibrium”的复数形式是“equilibria”。
\end{tcolorbox}

The congestion games studied in Chapter 2 are \(N\) -player games, where the strategy set of each player \(i\) is the set of paths from her origin \({o}_{i}\) to her destination \({d}_{i}\) , and a player's payoff is the negative of her cost as it results from the cost of using the edges on such a path, depending on their usage by all players. Exercise 3.2 considers a three-player game where the strategic form is given directly by a suitable "three-dimensional" table.

\begin{tcolorbox}
第2章中研究的拥塞博弈(congestion game)是 \(N\) 人博弈，其中每个参与者 \(i\) 的策略集是从其起点 \({o}_{i}\) 到其终点 \({d}_{i}\) 的路径集合，而参与者的收益是其成本的负值，该成本源于使用这条路径上的边的成本，这取决于所有参与者对这些边的使用情况。练习3.2考虑了一个三人博弈，其策略形式(strategic form)由一个合适的“三维”表格直接给出。
\end{tcolorbox}

Why is equilibrium such an important concept? The analysis of a game should describe what players will (or should) do when following their preferences, taking into account that the other players do the same. A recommendation of a strategy to each player defines a strategy profile. This profile should be an equilibrium because otherwise this recommendation would be self-defeating, that is, at least one player would have an incentive to deviate from her recommended strategy to some other strategy. In that sense, equilibrium is a necessary condition of "stability" for a game-theoretic "solution". This does not mean that an equilibrium will always exist, or that the players will always reach an equilibrium when they adapt their strategies to be best responses. In practice, people may not play an equilibrium because the game is only a simplified model of the interactive situation, and even in a game with clear rules (such as Chess) it may be too difficult to play perfectly.

\begin{tcolorbox}
为什么均衡(equilibrium)是如此重要的概念呢？对博弈的分析应该描述参与者在遵循自身偏好的情况下会(或应该)怎么做，同时要考虑到其他参与者也会这么做。为每个参与者推荐一种策略就定义了一个策略组合。这个组合应该是一个均衡，否则这个推荐就会自相矛盾，也就是说，至少有一个参与者会有动机从推荐给她的策略转向其他策略。从这个意义上说，均衡是博弈论“解”的“稳定性”的必要条件。这并不意味着均衡总是存在，也不意味着参与者在调整策略以达到最优反应(best response)时总能达到均衡。在实践中，人们可能不会达到均衡，因为博弈只是对互动情况的简化模型，即使在规则明确的博弈(如国际象棋)中，要完美地进行博弈也可能太难了。
\end{tcolorbox}

Much of this book is about determining the equilibria of a game, as the first step in a game-theoretic analysis. A separate argument is then which of these equilibria, if any, make sense as a practical recommendation of how to play. Another insight may be that an equilibrium has undesirable properties, like in the Prisoner's Dilemma or in the Braess paradox in Figure 2.3.

\begin{tcolorbox}
本书的大部分内容是关于确定博弈的均衡，这是博弈论分析的第一步。然后另一个问题是，如果存在多个均衡，那么其中哪些均衡可以作为如何进行博弈的实际建议。另一个见解可能是，一个均衡可能有不理想的性质，比如在囚徒困境或图2.3中的布雷斯悖论(Braess paradox)中。
\end{tcolorbox}

\section*{3.4 Games with Multiple Equilibria}

\begin{tcolorbox}
\section*{3.4 具有多个均衡的博弈}
\end{tcolorbox}

We consider several well-known \(2 \times  2\) games that have more than one equilibrium. The game of Chicken is shown in Figure 3.3. It is symmetric, as shown by the dotted line, and we have directly marked the best-response payoffs with boxes; recall that neither markup is part of the game description but already part of its analysis. The two strategies \(A\) and \(C\) stand for "aggressive" and "cautious" behavior, for example of two car drivers that drive towards each other on a narrow road. The aggressive strategy is only advantageous (with payoff 2) if the other player is cautious but leads to a crash (payoff 0 ) if the other player is aggressive, whereas a cautious strategy always gives payoff 1 to the player who uses it. This game has two equilibria,(C, A)and(A, C), because the best response to aggressive behavior is to be cautious and vice versa.

\begin{tcolorbox}
我们考虑几个著名的具有多个均衡的对称博弈(symmetric game)。胆小鬼博弈(Chicken)如图3.3所示。它是对称的，如虚线所示，我们已经直接用方框标记出了最优反应的收益；请记住，这两种标记都不是博弈描述的一部分，而是其分析的一部分。两种策略 \(A\) 和 \(C\) 分别代表“激进”和“谨慎”的行为，例如两个司机在狭窄道路上相向行驶的情况。如果另一个参与者采取谨慎策略，激进策略才是有利的(收益为2)，但如果另一个参与者也采取激进策略，就会导致碰撞(收益为0)，而谨慎策略总是给采取该策略的参与者带来收益1。这个博弈有两个均衡，(C, A)和(A, C)，因为对激进行为的最优反应是谨慎，反之亦然。
\end{tcolorbox}

\begin{center}
\includegraphics[max width=0.2\textwidth]{images/bo_d2e5cmbef24c73b03tn0_71_397_997_353_357_0.jpg}
\end{center}
\hspace*{3em} 

Figure 3.3 The symmetric game of Chicken.

\begin{tcolorbox}
图3.3 对称的胆小鬼博弈。
\end{tcolorbox}

The game known as the Battle of the Sexes is shown in Figure 3.4(a); its rather antiquated gender stereotypes should not be taken too seriously. In this scenario, player I and player II are a couple who each decide (simultaneously and independently, which is rather unrealistic) whether to go to a concert(C)or to a sports event(S). The players have different payoffs arising from which event they go to, but that payoff is zero if they have to attend the event alone. The game has two equilibria:(C, C)where they both go to the concert, or(S, S)where they both go to the sports event. Their preferences between these events differ, however.

\begin{tcolorbox}
被称为性别战(Battle of the Sexes)的博弈如图3.4(a)所示；其相当过时的性别刻板印象不应被太当真。在这个场景中，参与者I和参与者II是一对情侣，他们各自决定(同时且独立地，这相当不现实)是去听音乐会(C)还是去看体育赛事(S)。参与者从他们去的活动中获得不同的收益，但如果他们必须独自参加活动，收益就是零。这个博弈有两个均衡:(C, C)，即他们都去听音乐会；或者(S, S)，即他们都去看体育赛事。然而，他们对这些活动的偏好不同。
\end{tcolorbox}

The Battle of the Sexes game is not symmetric when written as in Figure 3.4(a), because the payoffs for the strategy pairs(C, C)and(S, S)on the diagonal are not the same for both players, which is clearly necessary for symmetry. However, changing the order of the strategies of one player, for example of player I as shown in Figure 3.4(b), makes this a game with a symmetric payoff structure; for a true symmetric game, one would also have to exchange the strategy names \(C\) and \(S\) of the strategies of player I, which would not represent the actions that player I takes. The game in Figure 3.4(b) is very similar to Chicken. However, because of the meaning of the strategies, the Battle of the Sexes is a coordination game (both players benefit from choosing the same action), whereas Chicken is an "anti-coordination" game.

\begin{tcolorbox}
如图3.4(a)所示写出的性别战博弈不是对称的，因为对角线上的策略对(C, C)和(S, S)的收益对两个参与者来说并不相同，而这显然是对称的必要条件。然而，改变一个参与者的策略顺序，例如如图3.4(b)所示改变参与者I的策略顺序，会使这个博弈具有对称的收益结构；对于真正的对称博弈，还必须交换参与者I的策略名称 \(C\) 和 \(S\) ，但这就不能代表参与者I实际采取的行动了。图3.4(b)中的博弈与胆小鬼博弈非常相似。然而，由于策略的含义，性别战是一个协调博弈(coordination game)(两个参与者都从选择相同的行动中受益)，而胆小鬼博弈是一个“反协调”博弈。
\end{tcolorbox}

\begin{center}
\includegraphics[max width=0.8\textwidth]{images/bo_d2e5cmbef24c73b03tn0_72_244_195_1105_362_0.jpg}
\end{center}
\hspace*{3em} 

Figure 3.4 The Battle of the Sexes game (a), whose payoff structure is symmetric if the strategies of one player are exchanged, as in (b).

\begin{tcolorbox}
图3.4 性别战博弈(a)，如果交换一个参与者的策略，其收益结构就会对称，如图(b)所示。
\end{tcolorbox}

The Stag Hunt game in Figure 3.5 models a conflict between cooperation and safety. Two hunters individually decide between \(S\) , to hunt a stag (a male deer), or \(H\) , to hunt a hare. Each can catch the hare, without the help of the other player, and then receives payoff 1 . In order to catch the stag they have to cooperate - if the other hunter goes for the hare then the stag hunter's payoff is zero. If they both hunt the stag they succeed and each get payoff 3 .

\begin{tcolorbox}
图3.5中的猎鹿博弈(Stag Hunt)模型展示了合作与安全之间的冲突。两名猎人各自要在 \(S\) (猎鹿，雄鹿)和 \(H\) (猎兔)之间做出选择。每个猎人无需他人帮助就能捕获野兔，从而获得收益1。要捕获鹿，他们必须合作 — 如果另一个猎人去猎兔，那么猎鹿者的收益为零。如果两人都猎鹿，他们就能成功，且各自获得收益3。
\end{tcolorbox}

\begin{center}
\includegraphics[max width=0.2\textwidth]{images/bo_d2e5cmbef24c73b03tn0_72_258_1187_352_357_0.jpg}
\end{center}
\hspace*{3em} 

Figure 3.5 The Stag Hunt game.

\begin{tcolorbox}
图3.5 猎鹿博弈(Stag Hunt)。
\end{tcolorbox}

This symmetric game has two symmetric equilibria,(S, S)and(H, H). The first equilibrium has a much higher payoff to both, but it requires trusting that the other player will also play \(S\) , because playing \(H\) is completely safe no matter what the other player does. The Stag Hunt game is sometimes called the "Trust Dilemma" for that reason.

\begin{tcolorbox}
这个对称博弈(symmetric game)有两个对称均衡(symmetric equilibrium)，即(S, S)和(H, H)。第一个均衡为双方带来的收益要高得多，但这需要信任对方也会选择 \(S\) ，因为无论对方怎么做，选择 \(H\) 都是完全安全的。出于这个原因，猎鹿博弈有时也被称为“信任困境(Trust Dilemma)”。
\end{tcolorbox}

It would be useful if game theory gave guidance of how to select between multiple equilibria. In the Stag Hunt game, it seems very reasonable to recommend (S, S), because both players get more than when playing(H, H), and it is strictly better to play \(S\) if the other player also plays \(S\) . However, that reasoning alone is tenuous, because it also applies in the following variation of the game. Suppose each payoff 1 in Figure 3.5 is replaced by 2.9 and these payoffs are in millions of dollars, so a player may get 3 million dollars by playing \(S\) (and risk getting nothing if the other player plays \(H\) ) or safely get 2.9 million dollars by playing \(H\) . Here the risk-free choice seems clearly better.

\begin{tcolorbox}
如果博弈论能为如何在多个均衡中进行选择提供指导，那将非常有用。在猎鹿博弈中，推荐选择(S, S)似乎非常合理，因为与选择(H, H)相比，双方的收益都更高，而且如果对方也选择 \(S\) ，那么选择 \(S\) 会更好。然而，仅靠这种推理是站不住脚的，因为它同样适用于该博弈的以下变体。假设图3.5中每个收益1都被替换为290万，且这些收益以百万美元为单位，那么玩家选择 \(S\) 可能获得300万美元(但如果对方选择 \(H\) ，则有一无所获的风险)，或者选择 \(H\) 安全地获得290万美元。在这里，无风险的选择显然更好。
\end{tcolorbox}

There are a number of different approaches to the problem of equilibrium selection. Typically, an equilibrium is meant to fulfill additional refinement criteria which are meant to be generally applicable (for example that there should not be another equilibrium with higher payoff to both players, although this may conflict with the inherent "riskiness" of an equilibrium, as the Stag Hunt game demonstrates). This generality is often debated, because it may require that the players use the same theoretical thought process to arrive at the supposedly "right" equilibrium. We leave this problem open, and are content with just identifying when a game has multiple equilibria.

\begin{tcolorbox}
对于均衡选择(equilibrium selection)问题，有多种不同的解决方法。通常，一个均衡需要满足额外的精炼标准，这些标准应具有普遍适用性(例如，不应存在另一个能让双方玩家获得更高收益的均衡，不过，正如猎鹿博弈所示，这可能与均衡固有的“风险性”相冲突)。这种普遍性常常引发争议，因为这可能要求玩家使用相同的理论思维过程来达成所谓“正确”的均衡。我们暂不讨论这个问题，只满足于识别一个博弈何时存在多个均衡。
\end{tcolorbox}

\section*{3.5 Dominated Strategies}

\begin{tcolorbox}
\section*{3.5 被占优策略(dominated strategy)}
\end{tcolorbox}

In the Quality game in Figure 3.2(b), strategy \(B\) gives a better payoff to the row player than strategy \(T\) , no matter what the other player does. We say \(B\) dominates \(T\) , by which we always mean \(B\) strictly dominates \(T\) . The following definition states the concepts of strict dominance, weak dominance, and payoff equivalence for a general \(N\) -player game.

\begin{tcolorbox}
在图3.2(b)的产品质量博弈(Quality game)中，无论对方玩家选择什么策略，行玩家选择策略 \(B\) 的收益都比选择策略 \(T\) 更高。我们说 \(B\) 占优 \(T\) ，这里我们始终指的是 \(B\) 严格占优 \(T\) 。以下定义阐述了一般 \(N\) 人博弈中严格占优、弱占优和收益等价(payoff equivalent)的概念。
\end{tcolorbox}

Definition 3.2. Consider an \(N\) -player game with strategy sets \({S}_{i}\) and payoff functions \({u}_{i}\) for each player \(i = 1,\ldots ,N\) , let \({s}_{i}\) and \({t}_{i}\) be two strategies of some player \(i\) , and let \({S}_{-i} = \mathop{\sum }\limits_{{j = 1,j \neq  i}}^{N}{S}_{j}\) be the set of partial profiles of the other players. Then

\begin{tcolorbox}
定义3.2。考虑一个 \(N\) 人博弈，每个玩家 \(i = 1,\ldots ,N\) 的策略集为 \({S}_{i}\) ，收益函数为 \({u}_{i}\) ，设 \({s}_{i}\) 和 \({t}_{i}\) 是某个玩家 \(i\) 的两个策略， \({S}_{-i} = \mathop{\sum }\limits_{{j = 1,j \neq  i}}^{N}{S}_{j}\) 是其他玩家的部分策略组合集合。那么
\end{tcolorbox}

\begin{itemize}
\item \({t}_{i}\) dominates (or strictly dominates) \({s}_{i}\) if
\end{itemize}

\begin{tcolorbox}
\begin{itemize}
\item 如果满足以下条件，则 \({t}_{i}\) 占优(或严格占优) \({s}_{i}\)
\end{itemize}
\end{tcolorbox}

\[
{u}_{i}\left( {{t}_{i},{s}_{-i}}\right)  > {u}_{i}\left( {{s}_{i},{s}_{-i}}\right) \;\text{ for all }{s}_{-i} \in  {S}_{-i}, \tag{3.4}
\]

\begin{itemize}
\item \({t}_{i}\) is payoff equivalent to \({s}_{i}\) if
\end{itemize}

\begin{tcolorbox}
\begin{itemize}
\item 如果满足以下条件，则 \({t}_{i}\) 与 \({s}_{i}\) 收益等价
\end{itemize}
\end{tcolorbox}

\[
{u}_{i}\left( {{t}_{i},{s}_{-i}}\right)  = {u}_{i}\left( {{s}_{i},{s}_{-i}}\right) \;\text{ for all }{s}_{-i} \in  {S}_{-i},
\]

\begin{itemize}
\item \({t}_{i}\) weakly dominates \({s}_{i}\) if \({t}_{i}\) and \({s}_{i}\) are not payoff equivalent and
\end{itemize}

\begin{tcolorbox}
\begin{itemize}
\item 如果 \({t}_{i}\) 和 \({s}_{i}\) 收益不等价，且满足以下条件，则 \({t}_{i}\) 弱占优 \({s}_{i}\)
\end{itemize}
\end{tcolorbox}

\[
{u}_{i}\left( {{t}_{i},{s}_{-i}}\right)  \geq  {u}_{i}\left( {{s}_{i},{s}_{-i}}\right) \;\text{ for all }{s}_{-i} \in  {S}_{-i}. \tag{3.5}
\]

The condition (3.4) of strict dominance is particularly easy to check in a two-player game. For example, if \(i\) is the column player, then \({s}_{i}\) and \({t}_{i}\) are two columns, and \({s}_{-i}\) is an arbitrary row, and we look at the payoffs of the column player. Then (3.4) states that, row by row, each entry in column \({t}_{i}\) is greater than the respective entry in column \({s}_{i}\) . For the columns \(D\) and \(C\) in the Prisoner’s Dilemma game in Figure 3.2(a), this can be written as \(\left( \begin{array}{l} 3 \\  1 \end{array}\right)  > \left( \begin{array}{l} 2 \\  0 \end{array}\right)\) , which is true. However, for the Quality game in Figure 3.2(b), the two payoff columns for \(r\) and \(l\) are \(\left( \begin{array}{l} 1 \\  1 \end{array}\right)\) and \(\left( \begin{array}{l} 2 \\  0 \end{array}\right)\) where this does not hold, and neither \(r\) dominates \(l\) nor \(l\) dominates \(r\) . However, for both games (which have the same payoffs to the row player) the bottom row dominates the top row because \(\left( \begin{array}{ll} 3 & 1 \end{array}\right)  > \left( \begin{array}{ll} 2 & 0 \end{array}\right)\) (we compare two vectors component by component).

\begin{tcolorbox}
在两人博弈中，严格占优的条件(3.4)特别容易验证。例如，如果 \(i\) 是列参与者，那么 \({s}_{i}\) 和 \({t}_{i}\) 是两列， \({s}_{-i}\) 是任意一行，我们观察列参与者的收益。那么(3.4)表明，逐行来看，列 \({t}_{i}\) 中的每个元素都大于列 \({s}_{i}\) 中的相应元素。对于图3.2(a)中囚徒困境博弈里的列 \(D\) 和 \(C\) ，这可以写成 \(\left( \begin{array}{l} 3 \\  1 \end{array}\right)  > \left( \begin{array}{l} 2 \\  0 \end{array}\right)\) ，这是成立的。然而，对于图3.2(b)中的产品质量博弈， \(r\) 和 \(l\) 的两个收益列分别是 \(\left( \begin{array}{l} 1 \\  1 \end{array}\right)\) 和 \(\left( \begin{array}{l} 2 \\  0 \end{array}\right)\) ，这种情况不成立， \(r\) 既不占优 \(l\) ， \(l\) 也不占优 \(r\) 。不过，对于这两个博弈(行参与者的收益相同)，底行占优顶行，因为 \(\left( \begin{array}{ll} 3 & 1 \end{array}\right)  > \left( \begin{array}{ll} 2 & 0 \end{array}\right)\) (我们逐分量比较两个向量)。
\end{tcolorbox}

Note that " \({t}_{i}\) dominates \({s}_{i}\) " does not mean that \({t}_{i}\) is "always better" than \({s}_{i}\) . In the Prisoner’s Dilemma, \(D\) dominates \(C\) , but playing \(D\) can result in payoff 1 and playing \(C\) in payoff 2, just not for the same strategy of the other player.

\begin{tcolorbox}
请注意，“ \({t}_{i}\) 占优 \({s}_{i}\) ”并不意味着 \({t}_{i}\) 总是比 \({s}_{i}\) “更好”。在囚徒困境中， \(D\) 占优 \(C\) ，但选择 \(D\) 可能得到收益1，选择 \(C\) 可能得到收益2，只是这不是在另一个参与者采用相同策略的情况下。
\end{tcolorbox}

Weak dominance is essentially the same as strict dominance where the strict inequality in (3.4) is replaced by a weak inequality in (3.5). However, it should not be equality throughout (which would mean payoff equivalence). That is, \({t}_{i}\) weakly dominates \({s}_{i}\) if (3.5) holds and \({u}_{i}\left( {{t}_{i},{s}_{-i}}\right)  > {u}_{i}\left( {{s}_{i},{s}_{-i}}\right)\) for at least one partial profile \({s}_{-i} \in  {S}_{-i}\) .

\begin{tcolorbox}
弱占优本质上与严格占优相同，只是(3.4)中的严格不等式被(3.5)中的弱不等式所取代。然而，不能始终取等号(这意味着收益等价)。也就是说，如果(3.5)成立，并且对于至少一个部分策略组合 \({s}_{-i} \in  {S}_{-i}\) 有 \({u}_{i}\left( {{t}_{i},{s}_{-i}}\right)  > {u}_{i}\left( {{s}_{i},{s}_{-i}}\right)\) ，那么 \({t}_{i}\) 弱占优 \({s}_{i}\) 。
\end{tcolorbox}

\begin{center}
\includegraphics[max width=0.2\textwidth]{images/bo_d2e5cmbef24c73b03tn0_74_258_1052_354_358_0.jpg}
\end{center}
\hspace*{3em} 

Figure 3.6 A game where column \(r\) weakly dominates column \(l\) . The weakly dominated strategy \(l\) is part of the equilibrium(T, l).

\begin{tcolorbox}
图3.6 一个列 \(r\) 弱占优列 \(l\) 的博弈。弱被占优策略 \(l\) 是均衡(T，l)的一部分。
\end{tcolorbox}

Figure 3.6 shows a game where the strategy \(T\) of player I has two best responses, \(l\) and \(r\) , because they result in the same payoff for player II (in fact for player I as well, but player I does not choose between columns). Against row \(B\) , player II’s payoff for \(l\) is worse than for \(r\) , so \(l\) is weakly dominated by \(r\) . Nevertheless, \(l\) is part of an equilibrium, namely(T, l), in fact with a higher payoff to player II than the equilibrium(B, r), so player II arguably has an advantage when "threatening" to play \(l\) . Therefore, one should not disregard weakly dominated strategies when analyzing a game, where one is typically interested in finding all its equilibria.

\begin{tcolorbox}
图3.6展示了一个博弈，其中参与者I的策略 \(T\) 有两个最优反应， \(l\) 和 \(r\) ，因为它们给参与者II带来相同的收益(实际上给参与者I也带来相同收益，但参与者I不选择列)。对于行 \(B\) ，参与者II选择 \(l\) 的收益比选择 \(r\) 的收益差，所以 \(l\) 被 \(r\) 弱占优。然而， \(l\) 是一个均衡，即(T，l)的一部分，实际上给参与者II带来的收益比均衡(B，r)更高，所以可以说参与者II在“威胁”选择 \(l\) 时具有优势。因此，在分析博弈时，不应该忽视弱被占优策略，因为通常人们感兴趣的是找出博弈的所有均衡。
\end{tcolorbox}

In contrast, in an equilibrium a player never plays a strictly dominated strategy, because he can play instead the strategy that dominates it. The following proposition considers what happens when a weakly or strictly dominated strategy is removed from the game.

\begin{tcolorbox}
相比之下，在均衡中，参与者永远不会选择严格被占优策略，因为他可以选择占优该策略的策略。下面的命题考虑了从博弈中移除弱被占优或严格被占优策略时会发生什么。
\end{tcolorbox}

Proposition 3.3. Consider an \(N\) -player game \(G\) in strategic form with strategy sets \({S}_{1},\ldots ,{S}_{N}\) . Let \({s}_{i}\) and \({t}_{i}\) be strategies of player \(i\) . Suppose \({t}_{i}\) weakly dominates or is payoff equivalent to \({s}_{i}\) . Consider the game \({G}^{\prime }\) with the same payoff functions as \(G\) but where \({S}_{i}\) is replaced by \({S}_{i} - \left\{  {s}_{i}\right\}\) . Then:

\begin{tcolorbox}
命题3.3。考虑一个具有策略集 \({S}_{1},\ldots ,{S}_{N}\) 的 \(N\) 人策略形式博弈 \(G\) 。设 \({s}_{i}\) 和 \({t}_{i}\) 是参与者 \(i\) 的策略。假设 \({t}_{i}\) 弱占优于 \({s}_{i}\) 或与 \({s}_{i}\) 收益等价。考虑与 \(G\) 具有相同收益函数，但 \({S}_{i}\) 被 \({S}_{i} - \left\{  {s}_{i}\right\}\) 替代的博弈 \({G}^{\prime }\) 。那么:
\end{tcolorbox}

(a) Any equilibrium of \({G}^{\prime }\) is an equilibrium of \(G\) .

\begin{tcolorbox}
(a) \({G}^{\prime }\) 的任何均衡都是 \(G\) 的均衡。
\end{tcolorbox}

(b) If \({t}_{i}\) dominates \({s}_{i}\) , then \(G\) and \({G}^{\prime }\) have the same equilibria.

\begin{tcolorbox}
(b) 如果 \({t}_{i}\) 占优于 \({s}_{i}\) ，那么 \(G\) 和 \({G}^{\prime }\) 具有相同的均衡。
\end{tcolorbox}

Proof. To show (a), consider an equilibrium of \({G}^{\prime }\) , which is a valid strategy profile in \(G\) . If player \(i\) could profitably deviate to strategy \({s}_{i}\) in \(G\) , then, clearly, player \(i\) could also profitably deviate to \({t}_{i}\) in \({G}^{\prime }\) , which violates the equilibrium property. The equilibrium property holds for the players other than player \(i\) because they have the same strategies and payoffs in \(G\) and \({G}^{\prime }\) .

\begin{tcolorbox}
证明。为了证明 (a)，考虑 \({G}^{\prime }\) 的一个均衡，它是 \(G\) 中的一个有效策略组合。如果参与者 \(i\) 在 \(G\) 中可以通过偏离到策略 \({s}_{i}\) 获得收益，那么显然，参与者 \(i\) 在 \({G}^{\prime }\) 中也可以通过偏离到 \({t}_{i}\) 获得收益，这违反了均衡性质。除参与者 \(i\) 之外的其他参与者的均衡性质成立，因为他们在 \(G\) 和 \({G}^{\prime }\) 中有相同的策略和收益。
\end{tcolorbox}

To show (b), suppose \({t}_{i}\) dominates \({s}_{i}\) . Then the dominated strategy \({s}_{i}\) is never part of an equilibrium \(s = \left( {{s}_{i},{s}_{-i}}\right)\) because of player \(i\) ’s profitable deviation to \({t}_{i}\) according to (3.4). Any strategy profile in \(G\) that does not contain \({s}_{i}\) is also a strategy profile of \({G}^{\prime }\) . If that profile is an equilibrium of \(G\) , then it is also an equilibrium of \({G}^{\prime }\) , which has fewer deviation possibilities. Any equilibrium in \({G}^{\prime }\) is an equilibrium in \(G\) by (a). Hence, \(G\) and \({G}^{\prime }\) have the same equilibria, as claimed.

\begin{tcolorbox}
为了证明 (b)，假设 \({t}_{i}\) 占优于 \({s}_{i}\) 。根据 (3.4)，由于参与者 \(i\) 可以通过偏离到 \({t}_{i}\) 获得收益，被占优策略 \({s}_{i}\) 永远不会是均衡 \(s = \left( {{s}_{i},{s}_{-i}}\right)\) 的一部分。 \(G\) 中任何不包含 \({s}_{i}\) 的策略组合也是 \({G}^{\prime }\) 的一个策略组合。如果该组合是 \(G\) 的一个均衡，那么它也是 \({G}^{\prime }\) 的一个均衡，因为 \({G}^{\prime }\) 的偏离可能性更少。根据 (a)， \({G}^{\prime }\) 中的任何均衡也是 \(G\) 中的均衡。因此，如所声称的， \(G\) 和 \({G}^{\prime }\) 具有相同的均衡。
\end{tcolorbox}

As explained before with Figure 3.6, eliminating a weakly dominated strategy may result in losing equilibria of a game. Proposition 3.3(a) shows that this does not introduce new equilibria. Hence, deleting a weakly dominated strategy may be considered acceptable when trying to find only one equilibrium of a game. However, this normally distorts the analysis of the game.

\begin{tcolorbox}
如之前结合图3.6所解释的，剔除一个弱被占优策略可能会导致博弈的均衡丢失。命题3.3(a)表明这不会引入新的均衡。因此，当试图仅找到博弈的一个均衡时，剔除一个弱被占优策略可以被认为是可接受的。然而，这通常会扭曲对博弈的分析。
\end{tcolorbox}

Proposition 3.3(b) shows that removing a (strictly) dominated strategy from a game does not change its equilibria. Because a payoff-maximizing player should not play a dominated strategy, it does make sense to eliminate a dominated strategy from the game by removing it from the player's strategy set. This changes the game, and the smaller game may have new dominated strategies that were not dominated before. Removing such a dominated strategy again, and continuing this process until all remaining strategies are undominated, is called iterated elimination of dominated strategies.

\begin{tcolorbox}
命题3.3(b)表明，从博弈中剔除一个(严格)被占优策略不会改变其均衡。因为一个追求收益最大化的参与者不应该采用被占优策略，所以从参与者的策略集中剔除被占优策略从而将其从博弈中消除是有意义的。这会改变博弈，并且更小的博弈可能会有之前未被占优的新的被占优策略。再次剔除这样的被占优策略，并持续这个过程直到所有剩余策略都未被占优，这被称为重复剔除劣势策略。
\end{tcolorbox}

In contrast, eliminating weakly dominated strategies is problematic because it may lose equilibria, and introduces further issues when iterated. As shown in Exercise 3.1, the order in which weakly dominated strategies are eliminated generally matters. The process may get stuck because weak domination turns into payoff equivalence, unless one also removes a payoff equivalent strategy, with a rather arbitrary decision as to which strategy should be eliminated. In addition, even if the elimination of weakly dominated strategies results in a single strategy profile, that profile may depend on the elimination order.

\begin{tcolorbox}
相比之下，剔除弱被占优策略是有问题的，因为它可能会丢失均衡，并且在重复操作时会引入更多问题。如练习3.1所示，剔除弱被占优策略的顺序通常是重要的。这个过程可能会陷入停滞，因为弱占优会变成收益等价，除非同时剔除一个收益等价的策略，而关于应该剔除哪个策略的决策相当随意。此外，即使剔除弱被占优策略得到了一个单一的策略组合，该组合也可能依赖于剔除顺序。
\end{tcolorbox}

Figure 3.7 demonstrates iterated elimination of dominated strategies for the Quality game. In the first step, the top row \(T\) is dominated by \(B\) and removed to create a \(1 \times  2\) game. In that game, column \(r\) dominates \(l\) (which is not true in the original \(2 \times  2\) game). After eliminating \(l\) , a single cell remains that corresponds to the strategy profile(B, r)which is also the unique equilibrium of the original game. The following proposition shows that this holds generally.

\begin{tcolorbox}
图3.7展示了产品质量博弈的重复剔除劣势策略过程。第一步，顶行 \(T\) 被 \(B\) 占优并被剔除，形成一个 \(1 \times  2\) 博弈。在那个博弈中，列 \(r\) 占优于 \(l\) (在原始的 \(2 \times  2\) 博弈中并非如此)。剔除 \(l\) 后，剩下一个单元格，对应策略组合 (B, r)，它也是原始博弈的唯一均衡。以下命题表明这通常是成立的。
\end{tcolorbox}

\begin{center}
\includegraphics[max width=0.9\textwidth]{images/bo_d2e5cmbef24c73b03tn0_76_214_233_1224_365_0.jpg}
\end{center}
\hspace*{3em} 

Figure 3.7 Iterated elimination of dominated strategies in the Quality game.

\begin{tcolorbox}
图3.7 产品质量博弈中的重复剔除劣势策略。
\end{tcolorbox}

Proposition 3.4. Consider a game G that by a finite iterated elimination of dominated strategies is reduced to a single strategy profile \(s\) (that is, every player has only one strategy). Then \(s\) is the unique equilibrium of \(G\) .

\begin{tcolorbox}
命题3.4。考虑一个博弈G，通过有限次重复剔除劣势策略将其简化为单一的策略组合 \(s\) (即每个参与者只有一种策略)。那么 \(s\) 是 \(G\) 的唯一均衡。
\end{tcolorbox}

Proof. As assumed, eliminate a dominated strategy from \(G\) , and continue this process in finitely many steps to arrive at a single strategy for each player. The final game only has one strategy profile \(s\) , which is trivially an equilibrium because no player can deviate. By Proposition 3.3, the previous game has the same unique equilibrium \(s\) , and by going backwards in the produced sequence of games, this eventually holds for the original game \(G\) .

\begin{tcolorbox}
证明。根据假设，从 \(G\) 中剔除一个(被占优)策略，并在有限步骤内继续此过程，为每个参与者得到单一策略。最终的博弈只有一个策略组合 \(s\) ，这显然是一个均衡，因为没有参与者可以偏离。根据命题3.3，前一个博弈有相同的唯一均衡 \(s\) ，通过在生成的博弈序列中逆向推导，这最终对原始博弈 \(G\) 也成立。
\end{tcolorbox}

If finite iterated elimination of dominated strategies ends in a unique strategy profile, the game is said to be dominance solvable, with the final strategy profile as its solution, which is also the unique equilibrium of the game. The following lemma states the reassuring fact that the order in which dominated strategies are eliminated does not matter, and that several dominated strategies can be removed simultaneously.

\begin{tcolorbox}
如果有限次重复剔除劣势策略以唯一的策略组合结束，则称该博弈是占优可解的，最终的策略组合即为其解，这也是该博弈的唯一均衡。以下引理表明了一个令人放心的事实，即剔除(被占优)策略的顺序无关紧要，并且可以同时剔除多个(被占优)策略。
\end{tcolorbox}

Lemma 3.5. Consider an \(N\) -player game \(G\) in strategic form with strategy sets \({S}_{1},\ldots\) , \({S}_{N}\) , let \({s}_{i}\) and \({\widehat{s}}_{j}\) be two dominated strategies of player \(i\) and player \(j\) , respectively, possibly the same player \(\left( {i = j}\right)\) . Let \({G}^{\prime }\) be the (otherwise unchanged) game obtained from \(G\) by removing \({s}_{i}\) from \({S}_{i}\) , and let \(\widehat{G}\) be the game obtained from \(G\) by removing \({\widehat{s}}_{j}\) from \({S}_{j}\) . Then \({\widehat{s}}_{j}\) is dominated in \({G}^{\prime }\) , and \({s}_{i}\) is dominated in \(\widehat{G}\) , and removing \({\widehat{s}}_{j}\) from \({G}^{\prime }\) and removing \({s}_{i}\) from \(\widehat{G}\) results in the same game obtained by removing directly \({s}_{i}\) and \({\widehat{s}}_{j}\) from \(G\) .

\begin{tcolorbox}
引理3.5。考虑一个具有策略集 \({S}_{1},\ldots\) 、 \({S}_{N}\) 的 \(N\) 个参与者的策略形式博弈 \(G\) ，设 \({s}_{i}\) 和 \({\widehat{s}}_{j}\) 分别是参与者 \(i\) 和参与者 \(j\) 的两个(被占优)策略，可能是同一个参与者 \(\left( {i = j}\right)\) 。设 \({G}^{\prime }\) 是从 \(G\) 中通过从 \({S}_{i}\) 移除 \({s}_{i}\) 得到的(其他方面不变)博弈，设 \(\widehat{G}\) 是从 \(G\) 中通过从 \({S}_{j}\) 移除 \({\widehat{s}}_{j}\) 得到的博弈。那么 \({\widehat{s}}_{j}\) 在 \({G}^{\prime }\) 中是(被占优)策略， \({s}_{i}\) 在 \(\widehat{G}\) 中是(被占优)策略，并且从 \({G}^{\prime }\) 移除 \({\widehat{s}}_{j}\) 和从 \(\widehat{G}\) 移除 \({s}_{i}\) 得到的博弈与直接从 \(G\) 移除 \({s}_{i}\) 和 \({\widehat{s}}_{j}\) 得到的博弈相同。
\end{tcolorbox}

Proof. Suppose \({s}_{i}\) is dominated by \({t}_{i}\) and \({\widehat{s}}_{j}\) is dominated by \({\widehat{t}}_{j}\) . If \({\widehat{t}}_{j} \neq  {s}_{i}\) , then \({\widehat{s}}_{j}\) is also dominated in \({G}^{\prime }\) by \({\widehat{t}}_{j}\) because any partial profile \({s}_{-j}\) in \({G}^{\prime }\) is also a partial profile in \(G\) where the dominating inequality holds. If \({\widehat{t}}_{j} = {s}_{i}\) (which requires \(i = j\) ), then \({\widehat{s}}_{j}\) is also dominated in \({G}^{\prime }\) (by \({t}_{i}\) , because domination is clearly transitive), as claimed. Similarly, if \({t}_{i} \neq  {\widehat{s}}_{j}\) , then \({s}_{i}\) is dominated in \(\widehat{G}\) by \({t}_{i}\) , and if \({t}_{i} = {\widehat{s}}_{j}\) , then \({s}_{i}\) is dominated in \(\widehat{G}\) by \({\widehat{t}}_{j}\) . Removing the dominated strategies clearly results in the same game with both \({s}_{i}\) and \({\widehat{s}}_{j}\) removed directly from \(G\) .

\begin{tcolorbox}
证明。假设 \({s}_{i}\) 被 \({t}_{i}\) 占优， \({\widehat{s}}_{j}\) 被 \({\widehat{t}}_{j}\) 占优。如果 \({\widehat{t}}_{j} \neq  {s}_{i}\) ，那么 \({\widehat{s}}_{j}\) 在 \({G}^{\prime }\) 中也被 \({\widehat{t}}_{j}\) 占优，因为 \({G}^{\prime }\) 中的任何部分策略组合 \({s}_{-j}\) 也是 \(G\) 中的部分策略组合，在其中占优不等式成立。如果 \({\widehat{t}}_{j} = {s}_{i}\) (这需要 \(i = j\) )，那么 \({\widehat{s}}_{j}\) 在 \({G}^{\prime }\) 中也被 \({t}_{i}\) 占优(因为占优显然是传递的)，如所声称的那样。类似地，如果 \({t}_{i} \neq  {\widehat{s}}_{j}\) ，那么 \({s}_{i}\) 在 \(\widehat{G}\) 中被 \({t}_{i}\) 占优，如果 \({t}_{i} = {\widehat{s}}_{j}\) ，那么 \({s}_{i}\) 在 \(\widehat{G}\) 中被 \({\widehat{t}}_{j}\) 占优。显然，剔除(被占优)策略会得到与直接从 \(G\) 剔除 \({s}_{i}\) 和 \({\widehat{s}}_{j}\) 相同的博弈。
\end{tcolorbox}

Iterated elimination of strictly dominated strategies has two uses. First, it may simplify the game, ideally until only one strategy profile remains, which is the unique equilibrium of the game (which of course only works if the game has only one equilibrium). Second, playing an equilibrium strategy assumes that the other players also play their equilibrium strategies. However, a player will always prefer a dominating strategy over the dominated strategy irrespective of the strategies of the other players. In that sense, choosing a dominating strategy is more "robust" in terms of assumptions about what the other players do. This applies certainly to a strategy that is dominating in the original game; iterated elimination of dominated strategies does require the "rationality" assumption that the other players maximize their payoff (in Figure 3.7, player II chooses \(r\) under the assumption that player I will not play \(T\) ).

\begin{tcolorbox}
重复剔除严格被占优策略有两个用途。首先，它可以简化博弈，理想情况下，直到只剩下一个策略组合，这就是该博弈的唯一均衡(当然，这只有在博弈只有一个均衡时才有效)。其次，选择均衡策略假设其他参与者也会选择他们的均衡策略。然而，无论其他参与者选择什么策略，参与者总是会更喜欢占优策略而不是(被占优)策略。从这个意义上说，选择占优策略在对其他参与者行为的假设方面更“稳健”。这当然适用于在原博弈中占优的策略；重复剔除劣势策略确实需要“理性”假设，即其他参与者会最大化他们的收益(在图3.7中，参与者II假设参与者I不会选择 \(r\) 而选择 \(T\) )。
\end{tcolorbox}

\section*{3.6 The Cournot Duopoly of Quantity Competition}

\begin{tcolorbox}
\section*{3.6 产量竞争的古诺双寡头博弈}
\end{tcolorbox}

In this section, we discuss an economic model of competition between two competing firms who choose the quantity of their product, where the price decreases with the total quantity on the market. It is also called a Cournot duopoly named after Augustin Cournot (1838) who proposed this model, nowadays described as a game with an equilibrium. In this game, the strategy sets of the players are infinite. Finite versions of the game are dominance solvable, and in principle also the infinite version, as we describe at the end of this section.

\begin{tcolorbox}
在本节中，我们讨论一个经济模型，该模型涉及两家相互竞争的企业之间的竞争，它们各自选择产品的产量，而市场价格会随着市场上的总产量的增加而下降。这个模型也被称为古诺双寡头博弈(Cournot duopoly)，以奥古斯丁·古诺(Augustin Cournot，1838 年)的名字命名，他提出了这个模型，如今该模型被描述为具有均衡的博弈。在这个博弈中，参与者的策略集是无限的。该博弈的有限版本是占优可解的，原则上无限版本也是如此，我们将在本节末尾进行描述。
\end{tcolorbox}

The players I and II are two firms who choose a nonnegative quantity of producing some good up to some upper bound \(M\) , say, so \({S}_{1} = {S}_{2} = \left\lbrack  {0,M}\right\rbrack\) . Let \(x\) and \(y\) be the strategies chosen by player I and II, respectively. For simplicity there are no costs of production, and the total quantity \(x + y\) is sold at a price \({12} - \left( {x + y}\right)\) per unit, which is also the firm’s profit per unit, so the payoffs \(a\left( {x,y}\right)\) and \(b\left( {x,y}\right)\) to players I and II are given by

\begin{tcolorbox}
参与者 I 和 II 是两家企业，它们选择生产某种商品的非负产量，最高可达某个上限 \(M\) ，比如说，因此 \({S}_{1} = {S}_{2} = \left\lbrack  {0,M}\right\rbrack\) 。设 \(x\) 和 \(y\) 分别是参与者 I 和 II 选择的策略。为简单起见，假设没有生产成本，总产量 \(x + y\) 以每单位 \({12} - \left( {x + y}\right)\) 的价格出售，这也是企业每单位的利润，因此参与者 I 和 II 的收益 \(a\left( {x,y}\right)\) 和 \(b\left( {x,y}\right)\) 由下式给出
\end{tcolorbox}

\[
\text{ payoff to }\mathrm{I} : \;a\left( {x,y}\right)  = x \cdot  \left( {{12} - y - x}\right) \text{ , } \tag{3.6}
\]

\[
\text{ payoff to II : }\;b\left( {x,y}\right)  = y \cdot  \left( {{12} - x - y}\right) \text{ . }
\]

\(\Rightarrow\) Exercise 3.3 describes an extension that incorporates production costs.

\begin{tcolorbox}
\(\Rightarrow\) 练习3.3描述了一个包含生产成本的扩展模型。
\end{tcolorbox}

The game (3.6) is clearly symmetric in \(x\) and \(y\) . Figure 3.8 shows a finite version of this game where the players' strategies are restricted to the four quantities 0 ,

\begin{tcolorbox}
博弈(3.6)在 \(x\) 和 \(y\) 上显然是对称的。图3.8展示了这个博弈的一个有限版本，其中参与者的策略被限制在四个产量水平上:0 ，
\end{tcolorbox}

\begin{center}
\includegraphics[max width=0.4\textwidth]{images/bo_d2e5cmbef24c73b03tn0_78_196_198_617_621_0.jpg}
\end{center}
\hspace*{3em} 

Figure 3.8 Symmetric duopoly game between two firms who both have four strategies 0,3 , 4, 6 with payoffs (3.6) if player I chooses \(x\) and player II chooses \(y\) .

\begin{tcolorbox}
图3.8 两家企业之间的对称双寡头博弈，两家企业都有四个策略:0、3、4、6，收益由(3.6)式确定。如果参与者 I 选择 \(x\) ，参与者 II 选择 \(y\) 。
\end{tcolorbox}

3,4, or 6 . The payoffs are determined by (3.6). For example, for \(x = 3\) and \(y = 4\) player I gets payoff 15 and player II gets 20. Best-response payoffs, as determined by this payoff table, are marked by boxes. The game has the unique equilibrium (4,4)where each player gets payoff 16 .

\begin{tcolorbox}
3、4 或 6。收益由(3.6)式确定。例如，当 \(x = 3\) 和 \(y = 4\) 时，参与者 I 的收益为15，参与者 II 的收益为20。根据这个收益表确定的最优反应收益用方框标记。该博弈有唯一的均衡(4,4)，此时每个参与者的收益为16。
\end{tcolorbox}

\begin{center}
\includegraphics[max width=0.9\textwidth]{images/bo_d2e5cmbef24c73b03tn0_78_207_1111_1213_491_0.jpg}
\end{center}
\hspace*{3em} 

Figure 3.9 The game in Figure 3.8 after eliminating for both players the dominated strategy 0, then strategy 6 (dominated by 4), then strategy 3 (dominated by 4), which shows that the game is dominance solvable.

\begin{tcolorbox}
图3.9 图3.8中的博弈，在为两个参与者剔除(被占优)策略 0，然后剔除策略 6(被 4 占优)，接着剔除策略 3(被 4 占优)之后，这表明该博弈是占优可解的。
\end{tcolorbox}

Figure 3.9 shows that the game can be solved by iterated elimination of dominated strategies. In the \(4 \times  4\) game, strategy 0 is dominated by 3 or 4, and after eliminating it directly for both players one obtains the \(3 \times  3\) game shown on the left in Figure 3.9. In that game, strategy 6 is dominated by 4, and after eliminating this strategy for both players the middle \(2 \times  2\) game is reached. This game has the structure of a Prisoner's Dilemma game because 4 dominates 3 (so that 3 will be eliminated), but the final strategy pair of undominated strategies (4,4) gives both players a lower payoff than the strategy pair(3,3). Here the Prisoner’s Dilemma arises in an economic context: The two firms could cooperate by equally splitting the optimal "monopoly" quantity 6, but in response the other player would "defect" and choose 4 rather than 3 . If both players do this, the total quantity 8 reduces the price to 4 and both players have an overall lower payoff.

\begin{tcolorbox}
图3.9表明，该博弈可以通过重复剔除劣势策略来求解。在 \(4 \times  4\) 博弈中，策略 0 被 3 或 4 占优，在为两个参与者直接剔除该策略后，得到图3.9左侧所示的 \(3 \times  3\) 博弈。在那个博弈中，策略 6 被 4 占优，在为两个参与者剔除该策略后，得到中间的 \(2 \times  2\) 博弈。这个博弈具有囚徒困境博弈的结构，因为 4 占优 3(因此 3 将被剔除)，但最终未被占优的策略对(4,4)给两个参与者带来的收益低于策略对(3,3)。这里的囚徒困境出现在经济背景中:两家企业可以通过平均分配最优的“垄断”产量 6 来进行合作，但作为回应，另一个参与者会“背叛”并选择 4 而不是 3。如果两个参与者都这样做，总产量 8 会使价格降至 4，两个参与者的总体收益都会降低。
\end{tcolorbox}

The game can also be analyzed when each player has the real interval \(\left\lbrack  {0,M}\right\rbrack\) as a strategy set. We can assume that \(M = {12}\) because if a player chooses a higher quantity than 12 then the price \({12} - x - y\) would be negative and the player would be better off by producing nothing. Player I’s best response to \(y\) is the quantity \(x\) that maximizes \(a\left( {x,y}\right)\) , where

\begin{tcolorbox}
当每个玩家将实数区间 \(\left\lbrack  {0,M}\right\rbrack\) 作为策略集时，也可以对该博弈进行分析。我们可以假设 \(M = {12}\) ，因为如果一个玩家选择的产量高于 12，那么价格 \({12} - x - y\) 将为负，此时该玩家不生产会更好。玩家 I 对 \(y\) 的最优反应是使 \(a\left( {x,y}\right)\) 最大化的产量 \(x\) ，其中
\end{tcolorbox}

\[
a\left( {x,y}\right)  = x \cdot  \left( {{12} - y - x}\right)  =  - {\left( x - \frac{{12} - y}{2}\right) }^{2} + {\left( \frac{{12} - y}{2}\right) }^{2} \tag{3.7}
\]

which clearly has its maximum at \(x = \frac{{12} - y}{2} = 6 - y/2\) , where \(x \geq  0\) because \(y \leq  {12}\) . (When just looking at \(x \cdot  \left( {{12} - y - x}\right)\) , this is a parabola with its apex at the midpoint between its zeros at \(x = 0\) and \(x = {12} - y\) . The maximum is also found by solving \(\left. {\frac{d}{dx}a\left( {x,y}\right)  = 0\text{ . }}\right)\) Note that the best response to \(y\) is unique. Because the game is symmetric, player II’s best response \(y\) to \(x\) is given by \(y = 6 - x/2\) .

\begin{tcolorbox}
显然，它在 \(x = \frac{{12} - y}{2} = 6 - y/2\) 处取得最大值，其中 \(x \geq  0\) ，因为 \(y \leq  {12}\) 。(仅考虑 \(x \cdot  \left( {{12} - y - x}\right)\) 时，这是一条抛物线，其顶点位于零点 \(x = 0\) 和 \(x = {12} - y\) 的中点。也可以通过求解 \(\left. {\frac{d}{dx}a\left( {x,y}\right)  = 0\text{ . }}\right)\) 来找到最大值。注意，对 \(y\) 的最优反应是唯一的。由于该博弈是对称的，玩家 II 对 \(x\) 的最优反应 \(y\) 由 \(y = 6 - x/2\) 给出。
\end{tcolorbox}

\begin{center}
\includegraphics[max width=0.5\textwidth]{images/bo_d2e5cmbef24c73b03tn0_79_344_1070_662_668_0.jpg}
\end{center}
\hspace*{3em} 

Figure 3.10 Best-response functions \(x = 6 - y/2\) and \(y = 6 - x/2\) of players \(\mathrm{I}\) and II for \(x,y \in  \left\lbrack  {0,{12}}\right\rbrack\) in the duopoly game (3.6).

\begin{tcolorbox}
图 3.10 在双头垄断博弈 (3.6) 中，玩家 \(\mathrm{I}\) 和玩家 II 对于 \(x,y \in  \left\lbrack  {0,{12}}\right\rbrack\) 的最优反应函数 \(x = 6 - y/2\) 和 \(y = 6 - x/2\) 。
\end{tcolorbox}

In order to obtain an equilibrium(x, y), the strategies \(x\) and \(y\) have to be mutual best responses, that is, \(x = 6 - y/2\) and \(y = 6 - x/2\) . This pair of linear equations has the unique solution \(x = y = 4\) , which is the unique equilibrium of the game, shown as the intersection point(4,4)of the best-response lines in Figure 3.10. It coincides with the equilibrium of the \(4 \times  4\) game in Figure 3.8 where 4 was one of the considered four strategies.

\begin{tcolorbox}
为了得到一个均衡 (x, y)，策略 \(x\) 和 \(y\) 必须是相互的最优反应，即 \(x = 6 - y/2\) 且 \(y = 6 - x/2\) 。这对线性方程有唯一解 \(x = y = 4\) ，它是该博弈的唯一均衡，如图 3.10 中最优反应线的交点 (4,4) 所示。它与图 3.8 中 \(4 \times  4\) 博弈的均衡一致，在图 3.8 中，4 是所考虑的四个策略之一。
\end{tcolorbox}

In the remainder of this section, we explain how even the game with infinite strategy sets \(\left\lbrack  {0,M}\right\rbrack\) is (to an arbitarily close approximation) dominance solvable, by iteratively eliminating whole intervals of dominated strategies from the strategy sets. The first step already happened when we chose \(M = {12}\) : In principle, player \(\mathrm{I}\) could choose \({x}^{\prime } > {12}\) , but then \(a\left( {{x}^{\prime },y}\right)  = {x}^{\prime } \cdot  \left( {{12} - y - {x}^{\prime }}\right)  < 0\) because \(y \geq  0\) , which is worse than \(a\left( {0,y}\right)\) for any \(y\) . That is, if \({x}^{\prime } > {12}\) then \({x}^{\prime }\) is strictly dominated by \(x = 0\) . The same applies to the other player, so we can assume that both strategy sets are equal to the interval \(\left\lbrack  {0,{12}}\right\rbrack\) .

\begin{tcolorbox}
在本节的剩余部分，我们将解释如何通过从策略集中迭代消除整个区间的(被占优)策略，使得即使是具有无限策略集 \(\left\lbrack  {0,M}\right\rbrack\) 的博弈也是(任意接近地近似)占优可解的。当我们选择 \(M = {12}\) 时，第一步就已经完成了:原则上，玩家 \(\mathrm{I}\) 可以选择 \({x}^{\prime } > {12}\) ，但此时 \(a\left( {{x}^{\prime },y}\right)  = {x}^{\prime } \cdot  \left( {{12} - y - {x}^{\prime }}\right)  < 0\) ，因为 \(y \geq  0\) ，对于任何 \(y\) 来说，这都比 \(a\left( {0,y}\right)\) 更差。也就是说，如果 \({x}^{\prime } > {12}\) ，那么 \({x}^{\prime }\) 被 \(x = 0\) 严格占优。同样的情况也适用于另一个玩家，所以我们可以假设两个策略集都等于区间 \(\left\lbrack  {0,{12}}\right\rbrack\) 。
\end{tcolorbox}

The following proposition states that when player II is assumed to play \(y\) from an interval \(\left\lbrack  {A,B}\right\rbrack\) , then there is another interval \(\left\lbrack  {{A}^{\prime },{B}^{\prime }}\right\rbrack   = \left\lbrack  {6 - \frac{B}{2},6 - \frac{A}{2}}\right\rbrack\) that contains all possible best responses of player I (condition (a)), and so that every strategy \({x}^{\prime }\) of player I outside the interval \(\left\lbrack  {{A}^{\prime },{B}^{\prime }}\right\rbrack\) is dominated by either \({B}^{\prime }\) or \({A}^{\prime }\) (conditions (b) and (c)). After eliminating these dominated strategies, we can assume that player I only plays strategies from \(\left\lbrack  {{A}^{\prime },{B}^{\prime }}\right\rbrack\) . If both players are initially considering strategies from the same interval \(\left\lbrack  {A,B}\right\rbrack\) , then condition (d) shows when this leads to a sequence of successively narrower intervals, which we will discuss afterwards.

\begin{tcolorbox}
以下命题表明，当假设玩家 II 从区间 \(y\) 中选择策略 \(\left\lbrack  {A,B}\right\rbrack\) 时，存在另一个区间 \(\left\lbrack  {{A}^{\prime },{B}^{\prime }}\right\rbrack   = \left\lbrack  {6 - \frac{B}{2},6 - \frac{A}{2}}\right\rbrack\) 包含玩家 I 的所有可能的最优反应(条件 (a))，并且区间 \(\left\lbrack  {{A}^{\prime },{B}^{\prime }}\right\rbrack\) 之外玩家 I 的每个策略 \({x}^{\prime }\) 都被 \({B}^{\prime }\) 或 \({A}^{\prime }\) 占优(条件 (b) 和 (c))。在消除这些(被占优)策略后，我们可以假设玩家 I 仅从 \(\left\lbrack  {{A}^{\prime },{B}^{\prime }}\right\rbrack\) 中选择策略。如果两个玩家最初都在考虑来自同一区间 \(\left\lbrack  {A,B}\right\rbrack\) 的策略，那么条件 (d) 表明何时会导致一系列逐渐变窄的区间，我们将在后面讨论。
\end{tcolorbox}

Proposition 3.6. Consider the game with payoffs (3.6) where \(y \in  \left\lbrack  {A,B}\right\rbrack\) with \(0 \leq  A \leq\)  \(B \leq  {12}\) . Let \({A}^{\prime } = 6 - \frac{B}{2}\) and \({B}^{\prime } = 6 - \frac{A}{2}\) . Then

\begin{tcolorbox}
命题 3.6. 考虑具有收益 (3.6) 的博弈，其中 \(y \in  \left\lbrack  {A,B}\right\rbrack\) 满足 \(0 \leq  A \leq\)  \(B \leq  {12}\) 。设 \({A}^{\prime } = 6 - \frac{B}{2}\) 且 \({B}^{\prime } = 6 - \frac{A}{2}\) 。那么
\end{tcolorbox}

(a) Each \(x\) in \(\left\lbrack  {{A}^{\prime },{B}^{\prime }}\right\rbrack\) is player \({I}^{\prime }\) s unique best response to some \(y\) in \(\left\lbrack  {A,B}\right\rbrack\) .

\begin{tcolorbox}
(a) 区间 \(\left\lbrack  {{A}^{\prime },{B}^{\prime }}\right\rbrack\) 中的每个 \(x\) 都是玩家 \({I}^{\prime }\) 对区间 \(\left\lbrack  {A,B}\right\rbrack\) 中某个 \(y\) 的唯一最优反应。
\end{tcolorbox}

(b) If \({x}^{\prime } > {B}^{\prime }\) then strategy \({x}^{\prime }\) of player \(I\) is strictly dominated by \(x = {B}^{\prime }\) .

\begin{tcolorbox}
(b) 如果 \({x}^{\prime } > {B}^{\prime }\) ，那么玩家 \(I\) 的策略 \({x}^{\prime }\) 被 \(x = {B}^{\prime }\) 严格占优。
\end{tcolorbox}

(c) If \({x}^{\prime } < {A}^{\prime }\) then strategy \({x}^{\prime }\) of player \(I\) is strictly dominated by \(x = {A}^{\prime }\) .

\begin{tcolorbox}
(c) 如果 \({x}^{\prime } < {A}^{\prime }\) ，那么玩家 \(I\) 的策略 \({x}^{\prime }\) 被 \(x = {A}^{\prime }\) 严格占优。
\end{tcolorbox}

(d) The interval \(\left\lbrack  {{A}^{\prime },{B}^{\prime }}\right\rbrack\) has half the length of the interval \(\left\lbrack  {A,B}\right\rbrack\) , and \(\left\lbrack  {{A}^{\prime },{B}^{\prime }}\right\rbrack   \subseteq  \left\lbrack  {A,B}\right\rbrack\) if and only if

\begin{tcolorbox}
(d) 区间 \(\left\lbrack  {{A}^{\prime },{B}^{\prime }}\right\rbrack\) 的长度是区间 \(\left\lbrack  {A,B}\right\rbrack\) 长度的一半，并且当且仅当 \(\left\lbrack  {{A}^{\prime },{B}^{\prime }}\right\rbrack   \subseteq  \left\lbrack  {A,B}\right\rbrack\)
\end{tcolorbox}

\[
4 \in  \left\lbrack  {A + \frac{1}{3}\left( {B - A}\right) ,A + \frac{2}{3}\left( {B - A}\right) }\right\rbrack  , \tag{3.8}
\]

and if (3.8) holds then it also holds with \({A}^{\prime },{B}^{\prime }\) instead of \(A,B\) . \(y = A\) and \(y = B\) (solid curves) and an intermediate case (dashed curve). Player I’s best response to \(y\) is \(x = \frac{{12} - y}{2}\) and therefore between \(\frac{{12} - B}{2}\) and \(\frac{{12} - A}{2}\) , that is, in the interval \(\left\lbrack  {{A}^{\prime },{B}^{\prime }}\right\rbrack\) , shown by the long solid vertical lines in Figure 3.11. The mapping \(\left\lbrack  {A,B}\right\rbrack   \rightarrow  \left\lbrack  {6 - \frac{B}{2},6 - \frac{A}{2}}\right\rbrack  ,y \mapsto  6 - \frac{y}{2}\) is a bijection, which shows (a).

\begin{tcolorbox}
并且如果 (3.8) 成立，那么用 \({A}^{\prime },{B}^{\prime }\) 代替 \(A,B\) 时它也成立。 \(y = A\) 和 \(y = B\) (实线曲线)以及一个中间情况(虚线曲线)。玩家 I 对 \(y\) 的最优反应是 \(x = \frac{{12} - y}{2}\) ，因此在 \(\frac{{12} - B}{2}\) 和 \(\frac{{12} - A}{2}\) 之间，即在区间 \(\left\lbrack  {{A}^{\prime },{B}^{\prime }}\right\rbrack\) 内，如图 3.11 中的长实线垂直线所示。映射 \(\left\lbrack  {A,B}\right\rbrack   \rightarrow  \left\lbrack  {6 - \frac{B}{2},6 - \frac{A}{2}}\right\rbrack  ,y \mapsto  6 - \frac{y}{2}\) 是一个双射，这表明了 (a)。
\end{tcolorbox}

\begin{center}
\includegraphics[max width=0.6\textwidth]{images/bo_d2e5cmbef24c73b03tn0_80_203_1479_791_555_0.jpg}
\end{center}
\hspace*{3em} 

Figure 3.11 The dashed curve shows the payoff \(x \cdot  \left( {{12} - y - x}\right)\) to player I in Proposition 3.6 (negative payoffs are not shown). Its maximum is at \(\frac{{12} - y}{2}\) . Here \({A}^{\prime } = \frac{{12} - B}{2}\) and \({B}^{\prime } = \frac{{12} - A}{2}\)

\begin{tcolorbox}
图3.11 虚线曲线展示了命题3.6中参与者I的收益 \(x \cdot  \left( {{12} - y - x}\right)\) (负收益未显示)。其最大值位于 \(\frac{{12} - y}{2}\) 处。这里 \({A}^{\prime } = \frac{{12} - B}{2}\) 和 \({B}^{\prime } = \frac{{12} - A}{2}\)
\end{tcolorbox}

Proof. Assume that player II plays some \(y\) in \(\left\lbrack  {A,B}\right\rbrack\) . According to (3.7), the payoff to player I as a function of \(x\) is a parabola. Figure 3.11 shows the two extreme cases

\begin{tcolorbox}
证明。假设参与者II在 \(\left\lbrack  {A,B}\right\rbrack\) 中选择某个 \(y\) 。根据式(3.7)，参与者I的收益作为 \(x\) 的函数是一条抛物线。图3.11展示了两种极端情况
\end{tcolorbox}

The case \(y = A\) is best possible for player I, which is the largest parabola in Figure 3.11. The figure shows that if \({x}^{\prime } > {B}^{\prime } = \frac{{12} - A}{2}\) then player \(\mathrm{I}\) is certainly better off playing \(x = {B}^{\prime }\) because for each \(y \geq  A\) the dashed parabola slopes downwards (possibly even with negative values) from \({B}^{\prime }\) to \({x}^{\prime }\) . This intuition for (b) is formally shown as follows. Let \({x}^{\prime } > {B}^{\prime }\) . We want to show that \(a\left( {{B}^{\prime },y}\right)  > a\left( {{x}^{\prime },y}\right)\) , which by (3.7) is equivalent to

\begin{tcolorbox}
情况 \(y = A\) 对参与者I来说是最有利的，这对应图3.11中最大的抛物线。该图表明，如果 \({x}^{\prime } > {B}^{\prime } = \frac{{12} - A}{2}\) ，那么参与者 \(\mathrm{I}\) 选择 \(x = {B}^{\prime }\) 肯定更有利，因为对于每个 \(y \geq  A\) ，虚线抛物线从 \({B}^{\prime }\) 到 \({x}^{\prime }\) 是向下倾斜的(甚至可能取负值)。对于(b)的这种直观理解正式证明如下。设 \({x}^{\prime } > {B}^{\prime }\) 。我们要证明 \(a\left( {{B}^{\prime },y}\right)  > a\left( {{x}^{\prime },y}\right)\) ，根据式(3.7)，这等价于
\end{tcolorbox}

\[
- {\left( {B}^{\prime } - \frac{{12} - y}{2}\right) }^{2} >  - {\left( {x}^{\prime } - \frac{{12} - y}{2}\right) }^{2}\text{ . } \tag{3.9}
\]

With

\begin{tcolorbox}
有了
\end{tcolorbox}

\[
{x}^{\prime } = {B}^{\prime } + c\;\left( {c > 0}\right) ,\;y = A + d\;\left( {d \geq  0}\right)
\]

and therefore \(\frac{{12} - y}{2} = 6 - \frac{A + d}{2} = {B}^{\prime } - \frac{d}{2}\) the inequality (3.9) is equivalent to

\begin{tcolorbox}
因此 \(\frac{{12} - y}{2} = 6 - \frac{A + d}{2} = {B}^{\prime } - \frac{d}{2}\) ，不等式(3.9)等价于
\end{tcolorbox}

\[
{\left( {B}^{\prime } - \left( {B}^{\prime } - \frac{d}{2}\right) \right) }^{2} < {\left( {B}^{\prime } + c - \left( {B}^{\prime } - \frac{d}{2}\right) \right) }^{2},
\]

that is, to \({\left( d/2\right) }^{2} < {\left( c + d/2\right) }^{2}\) which is true. This shows (b).

\begin{tcolorbox}
即 \({\left( d/2\right) }^{2} < {\left( c + d/2\right) }^{2}\) ，这是成立的。这证明了(b)。
\end{tcolorbox}

Similarly, the case \(y = B\) is worst possible for player \(\mathrm{I}\) , which is the smallest parabola in Figure 3.11. Then if \({x}^{\prime } < {A}^{\prime } = \frac{{12} - B}{2}\) then player I is better off playing \(x = {A}^{\prime }\) because for each \(y \leq  B\) the parabola slopes upwards from \({x}^{\prime }\) to \({A}^{\prime }\) . Formally, let \({x}^{\prime } < {A}^{\prime }\) . We want to show that \(a\left( {{A}^{\prime },y}\right)  > a\left( {{x}^{\prime },y}\right)\) , that is,

\begin{tcolorbox}
类似地，情况 \(y = B\) 对参与者 \(\mathrm{I}\) 来说是最不利的，这对应图3.11中最小的抛物线。那么如果 \({x}^{\prime } < {A}^{\prime } = \frac{{12} - B}{2}\) ，参与者I选择 \(x = {A}^{\prime }\) 更有利，因为对于每个 \(y \leq  B\) ，抛物线从 \({x}^{\prime }\) 到 \({A}^{\prime }\) 是向上倾斜的。正式地，设 \({x}^{\prime } < {A}^{\prime }\) 。我们要证明 \(a\left( {{A}^{\prime },y}\right)  > a\left( {{x}^{\prime },y}\right)\) ，即
\end{tcolorbox}

\[
- {\left( {A}^{\prime } - \frac{{12} - y}{2}\right) }^{2} >  - {\left( {x}^{\prime } - \frac{{12} - y}{2}\right) }^{2}\text{ . } \tag{3.10}
\]

Now we set

\begin{tcolorbox}
现在我们设
\end{tcolorbox}

\[
{x}^{\prime } = {A}^{\prime } - c\;\left( {c > 0}\right) ,\;y = B - d\;\left( {d \geq  0}\right)
\]

and therefore \(\frac{{12} - y}{2} = 6 - \frac{B - d}{2} = {A}^{\prime } + \frac{d}{2}\) . Then (3.10) is equivalent to

\begin{tcolorbox}
因此 \(\frac{{12} - y}{2} = 6 - \frac{B - d}{2} = {A}^{\prime } + \frac{d}{2}\) 。那么式(3.10)等价于
\end{tcolorbox}

\[
{\left( {A}^{\prime } - \left( {A}^{\prime } + \frac{d}{2}\right) \right) }^{2} < {\left( {A}^{\prime } - c - \left( {A}^{\prime } + \frac{d}{2}\right) \right) }^{2},
\]

that is, to \({\left( -d/2\right) }^{2} < {\left( -c - d/2\right) }^{2}\) or again \({\left( d/2\right) }^{2} < {\left( c + d/2\right) }^{2}\) which is true. This shows (c).

\begin{tcolorbox}
即 \({\left( -d/2\right) }^{2} < {\left( -c - d/2\right) }^{2}\) 或者再次得到 \({\left( d/2\right) }^{2} < {\left( c + d/2\right) }^{2}\) ，这是成立的。这证明了(c)。
\end{tcolorbox}

To show (d), clearly \({B}^{\prime } - {A}^{\prime } = \left( {B - A}\right) /2\) . The condition \(\left\lbrack  {{A}^{\prime },{B}^{\prime }}\right\rbrack   \subseteq  \left\lbrack  {A,B}\right\rbrack\) is equivalent to \(A \leq  {A}^{\prime } = 6 - \frac{B}{2}\) and \({B}^{\prime } = 6 - \frac{A}{2} \leq  B\) , which has the following equivalent statements:

\begin{tcolorbox}
为了证明(d)，显然 \({B}^{\prime } - {A}^{\prime } = \left( {B - A}\right) /2\) 。条件 \(\left\lbrack  {{A}^{\prime },{B}^{\prime }}\right\rbrack   \subseteq  \left\lbrack  {A,B}\right\rbrack\) 等价于 \(A \leq  {A}^{\prime } = 6 - \frac{B}{2}\) 和 \({B}^{\prime } = 6 - \frac{A}{2} \leq  B\) ，它有以下等价表述:
\end{tcolorbox}

\[
A + \frac{B}{2} \leq  6 \leq  B + \frac{A}{2},
\]

\[
A + \frac{A + \left( {B - A}\right) }{2} \leq  6 \leq  A + \left( {B - A}\right)  + \frac{A}{2},
\]

\[
\frac{3}{2}A + \frac{\left( B - A\right) }{2} \leq  6 \leq  \frac{3}{2}A + \left( {B - A}\right) ,
\]

that is,

\begin{tcolorbox}
即
\end{tcolorbox}

\[
A + \frac{1}{3}\left( {B - A}\right)  \leq  4 \leq  A + \frac{2}{3}\left( {B - A}\right)  \tag{3.11}
\]

which is stated in (3.8). Assuming that (3.11) holds, we want to show the same inequalities with \({A}^{\prime },{B}^{\prime }\) instead of \(A,B\) , that is,

\begin{tcolorbox}
这在(3.8)中有所陈述。假设(3.11)成立，我们想证明用 \({A}^{\prime },{B}^{\prime }\) 代替 \(A,B\) 时的相同不等式，即
\end{tcolorbox}

\[
{A}^{\prime } + \frac{1}{3}\left( {{B}^{\prime } - {A}^{\prime }}\right)  \leq  4 \leq  {A}^{\prime } + \frac{2}{3}\left( {{B}^{\prime } - {A}^{\prime }}\right) ,
\]

\[
{B}^{\prime } - \frac{2}{3}\left( {{B}^{\prime } - {A}^{\prime }}\right)  \leq  4 \leq  {B}^{\prime } - \frac{1}{3}\left( {{B}^{\prime } - {A}^{\prime }}\right) ,
\]

\[
6 - \frac{A}{2} - \frac{2}{3}\left( {{B}^{\prime } - {A}^{\prime }}\right)  \leq  4 \leq  6 - \frac{A}{2} - \frac{1}{3}\left( {{B}^{\prime } - {A}^{\prime }}\right)
\]

or

\begin{tcolorbox}
或者
\end{tcolorbox}

\[
2 \leq  \frac{A}{2} + \frac{2}{3}\left( {{B}^{\prime } - {A}^{\prime }}\right) ,\;\frac{A}{2} + \frac{1}{3}\left( {{B}^{\prime } - {A}^{\prime }}\right)  \leq  2
\]

which when multiplied by 2 and using \(2\left( {{B}^{\prime } - {A}^{\prime }}\right)  = B - A\) is equivalent to (3.11) which we assumed to be true. This completes the proof of (d).

\begin{tcolorbox}
将其乘以 2 并使用 \(2\left( {{B}^{\prime } - {A}^{\prime }}\right)  = B - A\) 后，等价于我们假设为真的(3.11)。这就完成了(d)的证明。
\end{tcolorbox}

By Proposition 3.6, we can apply iterated elimination of dominated strategies to the game (3.6) as follows. Initially, assume that both players consider any strategy in the interval \(\left\lbrack  {A,B}\right\rbrack   = \left\lbrack  {0,{12}}\right\rbrack\) . Then \(\left\lbrack  {{A}^{\prime },{B}^{\prime }}\right\rbrack   = \left\lbrack  {0,6}\right\rbrack\) . Any strategy \({x}^{\prime }\) of player I outside \(\left\lbrack  {{A}^{\prime },{B}^{\prime }}\right\rbrack\) (which here means only \({x}^{\prime } > 6\) ) is dominated by \(x = {B}^{\prime }\) . (However, as stated in (a), any strategy in \(\left\lbrack  {{A}^{\prime },{B}^{\prime }}\right\rbrack\) is not dominated because it is a best response to some \(y\) in \(\left\lbrack  {A,B}\right\rbrack\) , which shows that the elimination at this stage is best possible.) Similarly, because the game is symmetric, any strategy \({y}^{\prime }\) of player II outside \(\left\lbrack  {{A}^{\prime },{B}^{\prime }}\right\rbrack\) is dominated by \(y = {B}^{\prime }\) . We eliminate these dominated strategies, so that the new strategy sets will be \(\left\lbrack  {{A}^{\prime },{B}^{\prime }}\right\rbrack\) . Condition (3.8) states that 4 is in the middle third of the interval \(\left\lbrack  {A,B}\right\rbrack\) , which is (just about) true if \(\left\lbrack  {A,B}\right\rbrack   = \left\lbrack  {0,{12}}\right\rbrack\) where that middle third is \(\left\lbrack  {4,8}\right\rbrack\) . Furthermore,(d) states that then the same holds for \(\left\lbrack  {{A}^{\prime },{B}^{\prime }}\right\rbrack\) (indeed,4 is in the middle third \(\left\lbrack  {2,4}\right\rbrack\) of the interval \(\left\lbrack  {0,6}\right\rbrack\) ). We can now repeat the elimination of dominated strategies with \({A}^{\prime },{B}^{\prime }\) instead of \(A,B\) , and each time the length \(B - A\) of the considered interval is reduced by half. The strategy \(x = 4\) and \(y = 4\) belongs to all these intervals. In fact, any other strategy will sooner or later be outside a sufficiently small interval and will therefore be eliminated, although there is no bound on the number of elimination rounds for strategies that are very close to 4 . In the end, only the strategy pair(4,4)survives, which is the unique equilibrium of the game.

\begin{tcolorbox}
根据命题 3.6，我们可以对博弈(3.6)应用重复剔除劣势策略(iterated elimination of dominated strategies)，具体如下。最初，假设两个参与者都考虑区间 \(\left\lbrack  {A,B}\right\rbrack   = \left\lbrack  {0,{12}}\right\rbrack\) 内的任何策略。然后 \(\left\lbrack  {{A}^{\prime },{B}^{\prime }}\right\rbrack   = \left\lbrack  {0,6}\right\rbrack\) 。参与者 I 在 \(\left\lbrack  {{A}^{\prime },{B}^{\prime }}\right\rbrack\) 之外的任何策略 \({x}^{\prime }\) (这里仅指 \({x}^{\prime } > 6\) )都被 \(x = {B}^{\prime }\) 占优。(然而，如(a)中所述， \(\left\lbrack  {{A}^{\prime },{B}^{\prime }}\right\rbrack\) 内的任何策略都不被占优，因为它是对 \(\left\lbrack  {A,B}\right\rbrack\) 中某个 \(y\) 的最优反应(best response)，这表明此阶段的剔除是最优的。)类似地，由于该博弈是对称的(symmetric)，参与者 II 在 \(\left\lbrack  {{A}^{\prime },{B}^{\prime }}\right\rbrack\) 之外的任何策略 \({y}^{\prime }\) 都被 \(y = {B}^{\prime }\) 占优。我们剔除这些被占优策略(dominated strategy)，使得新的策略集将是 \(\left\lbrack  {{A}^{\prime },{B}^{\prime }}\right\rbrack\) 。条件(3.8)表明 4 在区间 \(\left\lbrack  {A,B}\right\rbrack\) 的中间三分之一处，如果 \(\left\lbrack  {A,B}\right\rbrack   = \left\lbrack  {0,{12}}\right\rbrack\) 成立，这(几乎)是正确的，其中中间三分之一是 \(\left\lbrack  {4,8}\right\rbrack\) 。此外，(d)表明对于 \(\left\lbrack  {{A}^{\prime },{B}^{\prime }}\right\rbrack\) 同样成立(实际上，4 在区间 \(\left\lbrack  {0,6}\right\rbrack\) 的中间三分之一 \(\left\lbrack  {2,4}\right\rbrack\) 处)。我们现在可以用 \({A}^{\prime },{B}^{\prime }\) 代替 \(A,B\) 重复剔除劣势策略，并且每次所考虑区间的长度 \(B - A\) 都会减半。策略 \(x = 4\) 和 \(y = 4\) 属于所有这些区间。实际上，任何其他策略迟早都会处于足够小的区间之外，因此会被剔除，尽管对于非常接近 4 的策略，剔除轮数没有上限。最终，只有策略对(4,4)留存下来，这是该博弈的唯一均衡(equilibrium)。
\end{tcolorbox}

This shows that the quantity duopoly game (3.6) is dominance solvable (if one allows for an unlimited number of elimination rounds, otherwise up to any desired degree of accuracy). Dominance solvability is considered a stronger solution concept than equilibrium. The equilibrium itself, of course, can be computed directly as the solution(x, y)to the two equations \(x = 6 - y/2\) and \(y = 6 - x/2\) .

\begin{tcolorbox}
这表明产量双头垄断博弈(3.6)是占优可解的(dominance solvable)(如果允许进行无限轮次的剔除，否则可达到任何期望的精度)。占优可解性被认为是比均衡更强的解概念。当然，均衡本身可以直接作为两个方程 \(x = 6 - y/2\) 和 \(y = 6 - x/2\) 的解(x, y)来计算。
\end{tcolorbox}

\section*{3.7 Games without a Pure-Strategy Equilibrium}

\begin{tcolorbox}
\section*{3.7 没有纯策略均衡的博弈}
\end{tcolorbox}

Not every game has an equilibrium in pure strategies. Figure 3.12 shows two well-known examples. In Matching Pennies, the two players reveal a penny which can show Heads(H)or Tails(T). If the pennies match, then player I wins the other player's penny, otherwise player II. No strategy pair can be stable because the losing player would always deviate.

\begin{tcolorbox}
并非每个博弈都存在纯策略均衡。图3.12展示了两个著名的例子。在硬币匹配(博弈)中，两名玩家各出示一枚硬币，硬币可能正面(H)朝上或反面(T)朝上。如果两枚硬币的正反面相同，则玩家I赢得对方的硬币；否则，玩家II赢得硬币。不存在稳定的策略组合，因为输的玩家总会改变策略。
\end{tcolorbox}

\begin{center}
\includegraphics[max width=0.8\textwidth]{images/bo_d2e5cmbef24c73b03tn0_83_416_550_1084_490_0.jpg}
\end{center}
\hspace*{3em} 

Figure 3.12 Matching Pennies (left) and Rock-Paper-Scissors (right).

\begin{tcolorbox}
图3.12 硬币匹配(左)和石头剪刀布(右)。
\end{tcolorbox}

Rock-Paper-Scissors is a \(3 \times  3\) game where both players choose simultaneously one of their three strategies Rock(R), Paper(P), or Scissors(S). Rock loses to Paper, Paper loses to Scissors, and Scissors lose to Rock, and it is a draw otherwise. No two strategies are best responses to each other. Like Matching Pennies, this is a zero-sum game because the payoffs in any cell of the table sum to zero. Unlike Matching Pennies, Rock-Paper-Scissors is symmetric. Hence, when both players play the same strategy (the cells on the diagonal), they get the same payoff, which is zero because the game is zero-sum.

\begin{tcolorbox}
石头剪刀布是一种零和博弈，两名玩家同时从石头(R)、布(P)、剪刀(S)这三种策略中选择其一。石头输给布，布输给剪刀，剪刀输给石头；若双方选择相同，则为平局。没有两种策略能互为最优反应。和硬币匹配(博弈)一样，这是一个零和博弈，因为表格中任意单元格的收益之和为零。与硬币匹配(博弈)不同的是，石头剪刀布是对称博弈。因此，当两名玩家采用相同的策略(对角线上的单元格)时，他们的收益相同，由于这是零和博弈，所以收益为零。
\end{tcolorbox}

The game-theoretic recommendation is to play randomly in games like Matching Pennies or Rock-Paper-Scissors that have no equilibrium, according to certain probabilities that depend on the payoffs. As we will see in Chapter 6, any finite game has an equilibrium when players are allowed to use randomized strategies.

\begin{tcolorbox}
博弈论的建议是，在像硬币匹配(博弈)或石头剪刀布这类不存在均衡的博弈中，玩家应根据取决于收益的特定概率随机选择策略。正如我们将在第6章看到的，当允许玩家使用随机策略时，任何有限博弈都存在均衡。
\end{tcolorbox}

\section*{3.8 Symmetric Games with Two Strategies per Player}

\begin{tcolorbox}
\section*{3.8 每个玩家有两种策略的对称博弈}
\end{tcolorbox}

In this section, we consider \(N\) -player games that, somewhat surprisingly, always have a pure equilibrium, namely symmetric games where each player has only two strategies.

\begin{tcolorbox}
在本节中，我们考虑n人博弈，有些令人惊讶的是，这类博弈总是存在纯策略均衡，即每个玩家只有两种策略的对称博弈。
\end{tcolorbox}

An \(N\) -player game is symmetric if each player has the same set of strategies, and if the game stays the same after any permutation (shuffling) of the players and, correspondingly, their payoffs. For two players, this means that the game stays the same when exchanging the players and their payoffs, visualized by reflection along the diagonal.

\begin{tcolorbox}
如果每个玩家都有相同的策略集，并且在对玩家及其相应的收益进行任何排列(洗牌)后，博弈保持不变，那么该n人博弈就是对称的。对于两名玩家来说，这意味着当交换玩家及其收益时，博弈保持不变，这可以通过沿对角线反射来直观表示。
\end{tcolorbox}

We now consider symmetric \(N\) -player games where each player has two strategies, which we call 0 and 1 . Normally, any combination of these strategies defines a separate strategy profile, so there are \({2}^{N}\) profiles, and each of them specifies a payoff to each player, so the game is defined by \(N \cdot  {2}^{N}\) payoffs. If the game is symmetric, vastly fewer payoffs are needed. Then a strategy profile is determined by how many players choose 1, say \(k\) players (where \(0 \leq  k \leq  N\) ), and then the remaining \(N - k\) players choose 0, so the profile can be written as

\begin{tcolorbox}
现在我们考虑对称的n人博弈，其中每个玩家有两种策略，我们称之为0和1。通常，这些策略的任何组合都定义了一个单独的策略组合，因此有n + 1个组合，并且每个组合都为每个玩家指定了一个收益，所以该博弈由2(n + 1)个收益来定义。如果该博弈是对称的，则需要的收益数量会大大减少。那么，一个策略组合由选择1的玩家数量决定，假设为k名玩家(其中0 ≤ k ≤ n)，然后剩下的n - k名玩家选择0，所以该组合可以写成
\end{tcolorbox}

\[
\left( {\underset{k}{\underbrace{1,\ldots ,1}},\underset{N - k}{\underbrace{0,\ldots ,0}}}\right) . \tag{3.12}
\]

Because the game is symmetric, any profile where \(k\) players choose 1 has to give the same payoff as (3.12) to any player who chooses 1, and a second payoff to any player who chooses 0 . Hence, we need only two payoffs for these profiles (3.12) when \(1 \leq  k \leq  N - 1\) . When \(k = 0\) then the profile is \(\left( {0,\ldots ,0}\right)\) and all players play the same strategy 0 and only one payoff is needed, and similarly when \(k = N\) where all players play 1. Therefore, a symmetric \(N\) -player game with two strategies per player is specified by only \({2N}\) payoffs.

\begin{tcolorbox}
因为该博弈是对称的，任何有k名玩家选择1的组合，对于选择1的任何玩家来说，其收益必须与(3.12)中的收益相同，对于选择0的任何玩家来说，有另一个收益。因此，当0 < k < n时，对于这些组合(3.12)，我们只需要两个收益。当k = 0时，组合为(0, …, 0)，所有玩家都采用相同的策略0，只需要一个收益；同样，当k = n时，所有玩家都采用策略1。因此，每个玩家有两种策略的对称n人博弈仅由n + 2个收益来定义。
\end{tcolorbox}

Proposition 3.7. Consider a symmetric \(N\) -player game where each player has two strategies, 0 and 1 . Then this game has a pure equilibrium. The strategy profile \(\left( {1,1,\ldots ,1}\right)\) is the unique equilibrium of the game if and only if 1 dominates 0 .

\begin{tcolorbox}
命题3.7。考虑一个对称的n人博弈，其中每个玩家有两种策略，0和1。那么这个博弈存在纯策略均衡。当且仅当策略1占优于策略0时，策略组合(1, …, 1)是该博弈的唯一均衡。
\end{tcolorbox}

Proof. If 1 dominates 0, then \(\left( {1,1,\ldots ,1}\right)\) is clearly the unique equilibrium of the game, so suppose this not the case. Because 0 is not dominated and there are only two strategies, for some profile the payoff when playing 0 is greater than or equal to the payoff when playing 1, that is, 0 is a best response to the remaining strategies. Consider such a profile (3.12) where 0 is a best response, with the smallest number \(k\) of players who play 1, where \(0 \leq  k < N\) . Then this profile is an equilibrium: By assumption and symmetry, 0 is a best response to the strategies of the other players for every player who plays 0 . If 1 was not a best response for every player who plays 1, then such a player would obtain a higher payoff when changing his strategy to 0 . This change would result in a profile where \(k - 1\) players play 1, the remaining players play 0, and 0 is a best response, which contradicts the smallest of choice of \(k\) . This shows that the game has an equilibrium.

\begin{tcolorbox}
证明。如果策略1占优于策略0，那么 \(\left( {1,1,\ldots ,1}\right)\) 显然是该博弈的唯一均衡，因此假设并非如此。由于策略0未被占优，且只有两种策略，对于某些策略组合而言，选择策略0的收益大于或等于选择策略1的收益，即策略0是其余策略的最优反应。考虑这样一个策略组合(3.12)，其中策略0是最优反应，且选择策略1的玩家数量 \(k\) 最少，其中 \(0 \leq  k < N\) 。那么这个策略组合就是一个均衡:根据假设和对称性，对于每个选择策略0的玩家来说，策略0是对其他玩家策略的最优反应。如果对于每个选择策略1的玩家来说，策略1不是最优反应，那么这样的玩家将策略改为0时会获得更高的收益。这种改变会导致一个新的策略组合，其中有 \(k - 1\) 个玩家选择策略1，其余玩家选择策略0，并且策略0是最优反应，这与 \(k\) 的最小选择相矛盾。这表明该博弈存在一个均衡。
\end{tcolorbox}

We consider this proposition for symmetric games when \(N = 2\) . The symmetric Prisoner's Dilemma (Figure 3.1) has a single equilibrium, and it indeed uses a dominating strategy. The equilibrium is symmetric, that is, both players play the same strategy. The symmetric Stag Hunt game (Figure 3.5) has two symmetric equilibria. The symmetric Chicken game (Figure 3.3) has two equilibria which are not symmetric. Clearly, in any symmetric two-player game, non-symmetric equilibria come in pairs, here(A, C)and(C, A), obtained by exchanging the players and their strategies. Chicken does not have a symmetric pure equilibrium. The \(3 \times  3\) game of Rock-Paper-Scissors (Figure 3.12) is symmetric but does not have a pure equilibrium. This shows that Proposition 3.7 cannot be generalized to symmetric games with more than two strategies per player.

\begin{tcolorbox}
我们考虑当 \(N = 2\) 时对称博弈的这一命题。对称的囚徒困境(图3.1)有一个单一均衡，并且确实使用了占优策略。该均衡是对称的，即两个玩家选择相同的策略。对称的猎鹿博弈(图3.5)有两个对称均衡。对称的胆小鬼博弈(图3.3)有两个非对称的均衡。显然，在任何对称的两人博弈中，非对称均衡是成对出现的，这里是(A，C)和(C，A)，通过交换玩家及其策略得到。胆小鬼博弈没有对称的纯策略均衡。 \(3 \times  3\) 的石头剪刀布博弈(图3.12)是对称的，但没有纯策略均衡。这表明命题3.7不能推广到每个玩家有两种以上策略的对称博弈。
\end{tcolorbox}

\(\Rightarrow\) Exercise 3.4 describes a game related to Proposition 3.7.

\begin{tcolorbox}
\(\Rightarrow\) 练习3.4描述了一个与命题3.7相关的博弈。
\end{tcolorbox}

\section*{3.9 Further Reading}

\begin{tcolorbox}
\section*{3.9 延伸阅读}
\end{tcolorbox}

The display of staggered payoffs in the lower left and upper right of each cell in the payoff table is due to Thomas Schelling. In 2005, he received, together with Robert Aumann, the Nobel memorial prize in Economic Sciences (officially: The Sveriges Riksbank Prize in Economic Sciences in Memory of Alfred Nobel) "for having enhanced our understanding of conflict and cooperation through game-theory analysis." According to Dixit and Nalebuff (1991, p. 90), he said with excessive modesty: "If I am ever asked whether I ever made a contribution to game theory, I shall answer yes. Asked what it was, I shall say the invention of staggered payoffs in a matrix."

\begin{tcolorbox}
收益表中每个单元格左下角和右上角交错显示收益的方式归功于托马斯·谢林(Thomas Schelling)。2005年，他与罗伯特·奥曼(Robert Aumann)共同获得了诺贝尔经济学奖(正式名称:瑞典中央银行纪念阿尔弗雷德·诺贝尔经济学奖)，“因为他们通过博弈论分析增进了我们对冲突与合作的理解”。根据迪克西特(Dixit)和奈尔伯夫(Nalebuff)(1991年，第90页)的说法，他过分谦虚地说:“如果有人问我是否对博弈论做出过贡献，我会回答是的。如果问我贡献是什么，我会说是发明了矩阵中的交错收益。”
\end{tcolorbox}

The strategic form is taught in every course on non-cooperative game theory (it is sometimes called the "normal form", now less used in game theory because "normal" is an overworked term in mathematics). Osborne (2004) gives careful historical explanations of the games considered in this chapter, including the original duopoly model of Cournot (1838), and many others. Our Exercise 3.4 is taken from that book. Gibbons (1992) shows that the Cournot game is dominance solvable, with less detail than our proof of Proposition 3.6. Both Osborne and Gibbons disregard Schelling and use comma-separated payoffs as in (3.1).

\begin{tcolorbox}
策略形式在每一门非合作博弈论课程中都会讲授(它有时被称为“标准型”，现在在博弈论中较少使用，因为“标准”在数学中是一个被过度使用的术语)。奥斯本(Osborne，2004年)对本章所考虑的博弈进行了细致的历史解释，包括古诺(Cournot，1838年)最初的双头垄断模型以及许多其他模型。我们的练习3.4取自那本书。吉本斯(Gibbons，1992年)表明古诺博弈是占优可解的，但细节不如我们对命题3.6的证明。奥斯本和吉本斯都忽略了谢林，使用如(3.1)中用逗号分隔的收益。
\end{tcolorbox}

The Cournot game in Section 3.6 is also a potential game with a strictly concave potential function, which has a unique maximum and therefore a unique equilibrium. Neyman (1997) showed that it is also a unique correlated equilibrium (see Chapter 12). Potential games (Monderer and Shapley, 1996) generalize games with a potential function such as the congestion games considered in Section 2.5.

\begin{tcolorbox}
3.6节中的古诺博弈也是一个具有严格凹势函数的势博弈，该势函数有唯一的最大值，因此有唯一的均衡。内曼(Neyman，1997年)表明它也是唯一的相关均衡(见第12章)。势博弈(蒙德勒(Monderer)和沙普利(Shapley)，1996年)推广了具有势函数的博弈，如2.5节中考虑的拥塞博弈。
\end{tcolorbox}

A classic survey of equilibrium refinements is van Damme (1987). Proposition 3.7 seems to have been shown first by Cheng, Reeves, Vorobeychik, and Wellman (2004, theorem 1).

\begin{tcolorbox}
关于均衡精炼的经典综述是范·达姆(van Damme，1987年)的著作。命题3.7似乎最早由程(Cheng)、里夫斯(Reeves)、沃罗贝奇克(Vorobeychik)和韦尔曼(Wellman)(2004年，定理1)证明。
\end{tcolorbox}

\section*{3.10 Exercises for Chapter 3}

\begin{tcolorbox}
\section*{3.10 第3章练习}
\end{tcolorbox}

Exercise 3.1. Consider the following \(3 \times  3\) game.

\begin{tcolorbox}
练习3.1。考虑以下 \(3 \times  3\) 博弈。
\end{tcolorbox}

\begin{center}
\includegraphics[max width=0.3\textwidth]{images/bo_d2e5cmbef24c73b03tn0_86_575_351_488_483_0.jpg}
\end{center}
\hspace*{3em} 

(a) Identify all pairs of strategies where one strategy weakly dominates the other.

\begin{tcolorbox}
(a)找出所有一个策略弱占优于另一个策略的策略对。
\end{tcolorbox}

(b) Assume you are allowed to remove a weakly dominated strategy of some player. Do so, and repeat this process (of iterated elimination of weakly dominated strategies) until you find a single strategy pair of the original game. (Remember that two payoff equivalent strategies do not weakly dominate each other!)

\begin{tcolorbox}
(b) 假设允许你剔除某个参与者的一个弱被占优策略。进行此操作，并重复这个过程(即迭代剔除弱被占优策略)，直到你找到原博弈的一个单一策略组合。(记住，两个收益等价的策略不会相互弱占优！)
\end{tcolorbox}

(c) Find such an iterated elimination of weakly dominated strategies that results in a strategy pair other than the one found in (b), where both strategies, and the payoffs to the players, are different.

\begin{tcolorbox}
(c) 找出一种迭代剔除弱被占优策略的方法，使其得到一个与 (b) 中不同的策略组合，其中两个策略以及参与者的收益都不同。
\end{tcolorbox}

(d) What are the equilibria (in pure strategies) of the game?

\begin{tcolorbox}
(d) 该博弈的(纯策略)均衡是什么？
\end{tcolorbox}

Exercise 3.2. Consider the following three-player game in strategic form.

\begin{tcolorbox}
练习 3.2。考虑以下策略形式的三人博弈。
\end{tcolorbox}

\begin{center}
\includegraphics[max width=0.2\textwidth]{images/bo_d2e5cmbef24c73b03tn0_86_435_1415_356_351_0.jpg}
\end{center}
\hspace*{3em} 

\begin{center}
\includegraphics[max width=0.2\textwidth]{images/bo_d2e5cmbef24c73b03tn0_86_862_1414_353_351_0.jpg}
\end{center}
\hspace*{3em} 

III: \(L\) III: \(R\)

\begin{tcolorbox}
参与者 III: \(L\) 参与者 III: \(R\)
\end{tcolorbox}

Each player has two strategies: Player I has the strategies \(T\) and \(B\) (the top and bottom row), player II has the strategies \(l\) and \(r\) (the left and right column in each \(2 \times  2\) panel), and player III has the strategies \(L\) and \(R\) (the right or left panel). The payoffs to the players in each cell are given as triples of numbers to players I, II, III.

\begin{tcolorbox}
每个参与者有两个策略:参与者 I 有策略 \(T\) 和 \(B\) (顶行和底行)，参与者 II 有策略 \(l\) 和 \(r\) (每个 \(2 \times  2\) 面板中的左列和右列)，参与者 III 有策略 \(L\) 和 \(R\) (右面板或左面板)。每个单元格中参与者的收益以参与者 I、II、III 的数字三元组形式给出。
\end{tcolorbox}

(a) Identify all pairs of strategies where one strategy strictly, or weakly, dominates the other. [Hint: Make sure you understand what dominance means for more than two players. Be careful to consider the correct payoffs.]

\begin{tcolorbox}
(a) 找出所有一个策略严格或弱占优另一个策略的策略对。[提示:确保你理解多于两个参与者时占优的含义。注意考虑正确的收益。]
\end{tcolorbox}

(b) Apply iterated elimination of strictly dominated strategies to this game. What are the equilibria of the game?

\begin{tcolorbox}
(b) 对这个博弈应用迭代剔除严格被占优策略的方法。该博弈的均衡是什么？
\end{tcolorbox}

Exercise 3.3. Consider the duopoly game as studied in Section 3.6 where player I and II produce nonnegative quantities \(x\) and \(y\) and the unit price on the market is \({12} - x - y\) . In addition, the players now have a cost of production, which is a cost of 1 per unit for player I and a cost of 2 per unit for player II. Write down the payoff to the players as a modification of (3.6), find their best response functions, and determine the equilibrium of this game. What are the players' profits in equilibrium?

\begin{tcolorbox}
练习 3.3。考虑第 3.6 节中研究的双头垄断博弈，其中参与者 I 和 II 生产非负产量 \(x\) 和 \(y\) ，市场上的单位价格是 \({12} - x - y\) 。此外，现在参与者有生产成本，参与者 I 的单位成本为 1，参与者 II 的单位成本为 2。将参与者的收益写成 (3.6) 的修正形式，找出他们的最优反应函数，并确定这个博弈的均衡。均衡时参与者的利润是多少？
\end{tcolorbox}

Exercise 3.4. (Osborne, 2004, Exercise 33.1, Contributing to a public good.)

\begin{tcolorbox}
练习 3.4。(奥斯本，2004 年，练习 33.1，为公共品做贡献。)
\end{tcolorbox}

Each of \(N\) people chooses whether or not to contribute a fixed amount towards the provision of a public good. The good is provided if and only if at least \(K\) people contribute, where \(2 \leq  K \leq  N\) . If the good is not provided, contributions are not refunded. Each person strictly ranks outcomes from best to worst as follows:

\begin{tcolorbox}
\(N\) 个人中的每个人都选择是否为提供一种公共品贡献一笔固定的金额。当且仅当至少 \(K\) 个人做出贡献时，该公共品才会被提供，其中 \(2 \leq  K \leq  N\) 。如果公共品未被提供，贡献不会退还。每个人对结果的严格排序从最好到最差如下:
\end{tcolorbox}

(i) any outcome in which the good is provided and she does not contribute,

\begin{tcolorbox}
(i) 公共品被提供且她没有做出贡献的任何结果，
\end{tcolorbox}

(ii) any outcome in which the good is provided and she contributes,

\begin{tcolorbox}
(ii) 公共品被提供且她做出贡献的任何结果，
\end{tcolorbox}

(iii) any outcome in which the good is not provided and she does not contribute,

\begin{tcolorbox}
(iii) 公共品未被提供且她没有做出贡献的任何结果，
\end{tcolorbox}

(iv) any outcome in which the good is not provided and she contributes.

\begin{tcolorbox}
(iv) 公共品未被提供且她做出贡献的任何结果。
\end{tcolorbox}

Formulate this situation as a strategic-form game and find its equilibria. You need to know the number of players, the strategies that a player has available, and the players' preferences for strategy profiles. Is there a symmetry here? Is there an equilibrium in which more than \(K\) people contribute? One in which \(K\) people contribute? One in which fewer than \(K\) people contribute? 4

\begin{tcolorbox}
将这种情况表述为一个策略形式的博弈并找出其均衡。你需要知道参与者的数量、参与者可用的策略以及参与者对策略组合的偏好。这里有对称性吗？是否存在一个均衡，其中超过 \(K\) 个人做出贡献？一个 \(K\) 个人做出贡献的均衡？一个少于 \(K\) 个人做出贡献的均衡？4
\end{tcolorbox}

\section*{Game Trees with Perfect Information}

\begin{tcolorbox}
\section*{完美信息博弈树}
\end{tcolorbox}

This chapter considers game trees, the second main way for defining a noncooperative game in addition to the strategic form. In a game tree, players move sequentially and (in the case of perfect information studied in this chapter) are aware of the previous moves of the other players. In contrast, in a strategic-form game players move simultaneously. In this "dynamic" setting, a play means a specific run of the game given by a sequence of actions of the players.

\begin{tcolorbox}
本章将探讨博弈树，它是除策略形式之外定义非合作博弈的第二种主要方式。在博弈树中，玩家按顺序行动，并且(在本章研究的完美信息情况下)知晓其他玩家之前的行动。相比之下，在策略形式的博弈中，玩家同时行动。在这种“动态”设定下，一次博弈指的是由玩家的一系列行动所构成的特定博弈过程。
\end{tcolorbox}

A game tree is described by a tree whose nodes represent possible states of play. Some nodes are chance nodes where the next node is determined randomly according to a known probability distribution. A decision node belongs to a player who makes a move to determine the next node. Play ends at a terminal node or leaf of the game tree where each player receives a payoff.

\begin{tcolorbox}
博弈树由一棵树来描述，其节点代表可能的博弈状态。一些节点是机会节点，在这些节点上，下一个节点根据已知的概率分布随机确定。决策节点属于某个玩家，该玩家通过行动来确定下一个节点。博弈在博弈树的终端节点或叶节点结束，此时每个玩家会获得一个收益。
\end{tcolorbox}

Game trees can be solved by backward induction, where one finds optimal moves for each player, given that all future moves have already been determined. Backward induction starts with the decision nodes closest to the leaves, and assumes that a player always makes a payoff-maximizing move. Such a move is in general not unique, so that the players' moves determined by backward induction may be interdependent (see Figure 4.4).

\begin{tcolorbox}
博弈树可以通过逆向归纳法求解，即在已知所有未来行动都已确定的情况下，为每个玩家找到最优行动。逆向归纳法从最接近叶节点的决策节点开始，并假设玩家总是做出使收益最大化的行动。一般来说，这样的行动并非唯一，因此通过逆向归纳法确定的玩家行动可能相互依赖(见图4.4)。
\end{tcolorbox}

A strategy in a game tree is a derived concept. A strategy defines a move at every decision node of the player and therefore a plan of action for every possible state of play. Any strategy profile therefore defines a payoff to each player (which may be an expected payoff if there are chance moves), which defines the strategic form of the game.

\begin{tcolorbox}
博弈树中的策略是一个派生概念。策略定义了玩家在每个决策节点的行动，因此是针对每种可能博弈状态的行动计划。因此，任何策略组合都为每个玩家定义了一个收益(如果存在机会行动，可能是期望效用)，这就定义了博弈的策略形式。
\end{tcolorbox}

Because every strategy is a combination of moves, the number of strategies in a game tree is very large. Reduced strategies reduce this number to some extent. In a reduced strategy, moves at decision nodes that cannot be reached due to an earlier own of the player are left unspecified.

\begin{tcolorbox}
由于每个策略都是行动的组合，博弈树中的策略数量非常多。精简策略在一定程度上减少了这个数量。在精简策略中，由于玩家先前的行动而无法到达的决策节点上的行动不做具体规定。
\end{tcolorbox}

Backward induction defines a strategy profile which is an equilibrium of the game (Theorem 4.6). This is called a subgame-perfect equilibrium (SPE) because it defines an equilibrium in every subgame. (In a game tree with perfect information, a subgame is just a subtree of the game. In games with imperfect information, treated in Chapter 10, an SPE can in general not be found by backward induction.)

\begin{tcolorbox}
逆向归纳法定义了一个策略组合，它是博弈的一个均衡(定理4.6)。这被称为子博弈完美均衡(SPE)，因为它在每个子博弈中都定义了一个均衡。(在具有完美信息的博弈树中，子博弈只是博弈的一个子树。在第10章讨论的不完美信息博弈中，一般不能通过逆向归纳法找到子博弈完美均衡。)
\end{tcolorbox}

In Section 4.7, we consider commitment games, which are two-player games where a given strategic-form game is changed to a game tree where one player moves first and the other player is informed about the first player's move and can react to that move. It is assumed that the second player always chooses a best response, as in an SPE. This changes the game, typically to the advantage of the first mover. We study the commitment game for the Quality game and for the Cournot duopoly game of quantity competition from Section 3.6, where the commitment game is known as a Stackelberg leadership game.

\begin{tcolorbox}
在4.7节中，我们将考虑承诺博弈，这是一种两人博弈，其中给定的策略形式博弈被转变为一个博弈树，其中一个玩家先行动，另一个玩家知晓第一个玩家的行动并可以对此做出反应。假设第二个玩家总是选择最优反应，就像在子博弈完美均衡中一样。这改变了博弈，通常对先行动者有利。我们研究产品质量博弈和第3.6节中的古诺双寡头产量竞争博弈的承诺博弈，其中承诺博弈被称为斯塔克尔伯格领导博弈。
\end{tcolorbox}

Every game tree can be converted to strategic form. A general game in strategic form can only be represented by a tree with extra components that model imperfect information. Game trees with imperfect information will be treated in Chapter 10 .

\begin{tcolorbox}
每个博弈树都可以转换为策略形式。一般的策略形式博弈只能由带有模拟不完美信息的额外组件的树来表示。不完美信息的博弈树将在第10章讨论。
\end{tcolorbox}

\section*{4.1 Prerequisites and Learning Outcomes}

\begin{tcolorbox}
\section*{4.1 先决条件和学习成果}
\end{tcolorbox}

Chapter 3 is a prerequisite for this chapter. Chapter 1 is not, although combinatorial games are sequential games with perfect information as considered here, and therefore illustrative examples. We will recall the mathematical concepts of directed graphs and trees, and of the expected value of a random variable, as background material.

\begin{tcolorbox}
第3章是本章的先决条件。虽然组合博弈也是这里所考虑的具有完美信息的序贯博弈，是具有说明性的例子，但第1章不是先决条件。我们将回顾有向图与树图以及随机变量的期望值等数学概念，作为背景材料。
\end{tcolorbox}

After studying this chapter, you should be able to

\begin{tcolorbox}
学习本章后，你应该能够
\end{tcolorbox}

\begin{itemize}
\item interpret game trees with their moves and payoffs;
\end{itemize}

\begin{tcolorbox}
\begin{itemize}
\item 解读博弈树及其行动和收益；
\end{itemize}
\end{tcolorbox}

\begin{itemize}
\item create the strategies of a player as combinations of moves, and know how to count them;
\end{itemize}

\begin{tcolorbox}
\begin{itemize}
\item 将玩家的行动组合成策略，并知道如何计算策略的数量；
\end{itemize}
\end{tcolorbox}

\begin{itemize}
\item combine strategies into reduced strategies, which identify moves at decision nodes that are unreachable due to an earlier own move of the player, and understand where the corresponding "wildcard" symbol "*" is placed in the move list that represents the strategy;
\end{itemize}

\begin{tcolorbox}
\begin{itemize}
\item 将策略组合成精简策略，识别由于玩家先前的行动而无法到达的决策节点上的行动，并理解在表示策略的行动列表中，相应的“通配符”符号“*”应放在何处；
\end{itemize}
\end{tcolorbox}

\begin{itemize}
\item write down the strategic form of a game tree, with players' payoffs in the correct place in each cell, and with computed expected payoffs if needed;
\end{itemize}

\begin{tcolorbox}
\begin{itemize}
\item 写出博弈树的策略形式，将玩家的收益正确地放置在每个单元格中，并在需要时计算期望效用；
\end{itemize}
\end{tcolorbox}

\begin{itemize}
\item explain backward induction and why it requires full (unreduced) strategies, which may not be unique;
\end{itemize}

\begin{tcolorbox}
\begin{itemize}
\item 解释逆向归纳法(backward induction)以及为什么它需要完整(未简化)的策略，而这些策略可能并非唯一；
\end{itemize}
\end{tcolorbox}

\begin{itemize}
\item construct any equilibrium, even if not subgame-perfect, directly from the game tree, also for more than two players, as in Exercise 4.5;
\end{itemize}

\begin{tcolorbox}
\begin{itemize}
\item 直接从博弈树(game tree)构建任何均衡(equilibrium)，即使不是子博弈完美均衡，也适用于两个以上的参与者，如练习4.5所示；
\end{itemize}
\end{tcolorbox}

\begin{itemize}
\item create the commitment game from a game in strategic form, and compare its subgame-perfect equilibria with the equilibria in the strategic form.
\end{itemize}

\begin{tcolorbox}
\begin{itemize}
\item 从策略形式(strategic form)的博弈中构建承诺博弈(commitment game)，并将其子博弈完美均衡与策略形式中的均衡进行比较。
\end{itemize}
\end{tcolorbox}

\section*{4.2 Definition of Game Trees}

\begin{tcolorbox}
\section*{4.2 博弈树的定义}
\end{tcolorbox}

Figure 4.1 shows an example of a game tree. We always draw trees downwards, with the starting node, called the root, at the top. (Conventions on drawing game trees vary. Sometimes trees are drawn from the bottom upwards, sometimes from left to right, and sometimes with the root at the center with edges in any direction.)

\begin{tcolorbox}
图4.1展示了一个博弈树的示例。我们通常将树向下绘制，起始节点(称为根节点)位于顶部。(绘制博弈树的约定各不相同。有时树是从下向上绘制，有时是从左向右绘制，有时根节点位于中心，边可以指向任何方向。)
\end{tcolorbox}

\begin{center}
\includegraphics[max width=0.6\textwidth]{images/bo_d2e5cmbef24c73b03tn0_90_204_533_775_575_0.jpg}
\end{center}
\hspace*{3em} 

Figure 4.1 Example of a game tree. The square node indicates a chance move. At a leaf of the tree, the top payoff is to player I, the bottom payoff to player II.

\begin{tcolorbox}
图4.1 博弈树示例。方形节点表示随机行动。在树的叶节点(leaf)处，顶部的收益属于参与者I，底部的收益属于参与者II。
\end{tcolorbox}

The nodes of the tree denote states of play (which have been called "positions" in the combinatorial games considered in Chapter 1). Nodes are connected by lines, called edges. An edge from a node \(u\) to a child node \(v\) (where \(v\) is drawn below \(u\) ) indicates a possible move in the game. This may be a move of a "personal" player, for example move \(X\) of player I in Figure 4.1. Then \(u\) is also called a decision node. Alternatively, \(u\) is a chance node, like the node \(u\) that follows move \(b\) of player II in Figure 4.1. We draw decision nodes as small filled circles and chance nodes as squares. After a chance node \(u\) , the next node \(v\) is determined by a random choice made with the probability associated with the edge that leads from \(u\) to \(v\) . In Figure 4.1, these probabilities are \(\frac{1}{3}\) for the left move and \(\frac{2}{3}\) for the right move.

\begin{tcolorbox}
树的节点(node)表示博弈状态(在第1章所考虑的组合博弈中，这些状态被称为“位置”)。节点通过称为边(edge)的线连接。从节点 \(u\) 到子节点 \(v\) (其中 \(v\) 绘制在 \(u\) 下方)的边表示博弈中的一个可能行动。这可能是“个人”参与者的行动，例如图4.1中参与者I的行动 \(X\) 。此时 \(u\) 也称为决策节点。或者， \(u\) 是一个随机节点，如图4.1中参与者II行动 \(b\) 之后的节点 \(u\) 。我们将决策节点绘制为小实心圆，将随机节点绘制为方形。在随机节点 \(u\) 之后，下一个节点 \(v\) 由与从 \(u\) 到 \(v\) 的边相关联的概率随机确定。在图4.1中，向左行动的概率为 \(\frac{1}{3}\) ，向右行动的概率为 \(\frac{2}{3}\) 。
\end{tcolorbox}

At a terminal node or leaf of the game tree, every player gets a payoff. In Figure 4.1, leaves are not explicitly drawn, but the payoffs given instead, with the top payoff to player I and the bottom payoff to player II.

\begin{tcolorbox}
在博弈树的终端节点(terminal node)或叶节点处，每个参与者都会获得一个收益。在图4.1中，叶节点未明确绘制，而是直接给出了收益，顶部的收益属于参与者I，底部的收益属于参与者II。
\end{tcolorbox}

It does not matter how the tree is drawn, only how the nodes are connected by edges, as summarized in the background material on directed graphs and trees. The following is the formal definition of a game tree.

\begin{tcolorbox}
树的绘制方式并不重要，重要的是节点如何通过边连接，这在有向图与树图(directed graphs and trees)的背景材料(background material)中已有总结。以下是博弈树的正式定义。
\end{tcolorbox}

Definition 4.2. A game tree is a tree with the following additional information:

\begin{tcolorbox}
定义4.2. 博弈树是具有以下额外信息的树:
\end{tcolorbox}

\begin{itemize}
\item A finite set of \(N\) players.
\end{itemize}

\begin{tcolorbox}
\begin{itemize}
\item 有限个 \(N\) 参与者。
\end{itemize}
\end{tcolorbox}

\begin{itemize}
\item Each non-terminal node of the tree is either a decision node and belongs to one of the \(N\) players, or is a chance node. Background material: Directed graphs and trees
\end{itemize}

\begin{tcolorbox}
\begin{itemize}
\item 树的每个非终端节点要么是决策节点，属于 \(N\) 个参与者之一，要么是随机节点。背景材料:有向图与树图
\end{itemize}
\end{tcolorbox}

A directed graph (or digraph) is given by a finite set of nodes and a finite set of edges. Every edge is an ordered pair of nodes \(u\) and \(v\) , written as \({uv}\) . A walk from \(u\) to \(v\) is a sequence of nodes \({u}_{0},{u}_{1},\ldots ,{u}_{m}\) for \(m \geq  0\) where \({u}_{k}{u}_{k + 1}\) is an edge for \(0 \leq  k < m\) , and \(u = {u}_{0}\) and \(v = {u}_{m}\) . If these nodes are all distinct, the walk is called a path.

\begin{tcolorbox}
有向图(digraph)由有限个节点和有限条边组成。每条边是节点 \(u\) 和 \(v\) 的有序对，记为 \({uv}\) 。从 \(u\) 到 \(v\) 的行走(walk)是节点序列 \({u}_{0},{u}_{1},\ldots ,{u}_{m}\) ，其中 \(m \geq  0\) 且 \({u}_{k}{u}_{k + 1}\) 是边， \(0 \leq  k < m\) ， \(u = {u}_{0}\) 且 \(v = {u}_{m}\) 。如果这些节点都不同，则该行走称为路径(path)。
\end{tcolorbox}

A (directed) tree is a digraph with a distinguished node, the root, so that for every node \(u\) there is a unique walk from the root to \(u\) . If \({uv}\) is an edge of the tree, then \(v\) is called a child of \(u\) and \(u\) is called the parent of \(v\) . In a tree, a node \(v\) is a descendant of a node \(u\) if there is a nonempty walk from \(u\) to \(v\) . A node without a child is called a leaf or terminal node of the tree.

\begin{tcolorbox}
(有向)树是具有一个特殊节点(根节点)的有向图，使得对于每个节点 \(u\) ，从根节点到 \(u\) 都有唯一的行走。如果 \({uv}\) 是树的一条边，则 \(v\) 称为 \(u\) 的子节点， \(u\) 称为 \(v\) 的父节点。在树中，如果从节点 \(u\) 到节点 \(v\) 存在非空行走，则 \(v\) 是 \(u\) 的后代节点(descendant)。没有子节点的节点称为树的叶节点或终端节点。
\end{tcolorbox}

We represent trees graphically with the root at the top and edges drawn downwards from a parent to a child.

\begin{tcolorbox}
我们以图形方式表示树，根节点位于顶部，边从父节点向下绘制到子节点。
\end{tcolorbox}

Lemma 4.1. In a tree, the unique walk from the root to a node is a path. Every node has a unique parent, except for the root which has no parent.

\begin{tcolorbox}
引理4.1. 在树中，从根节点到节点的唯一行走是一条路径。每个节点都有唯一的父节点，除了根节点没有父节点。
\end{tcolorbox}

Proof. If the walk \({u}_{0},{u}_{1},\ldots ,{u}_{m}\) from the root \({u}_{0}\) to the node \({u}_{m}\) had a repeated node \({u}_{i} = {u}_{j}\) for some \(i < j\) , we could remove the nodes \({u}_{i + 1},\ldots ,{u}_{j}\) from the walk and obtain a second walk from \({u}_{0}\) to \({u}_{m}\) , which is not allowed. Hence, the walk has no repeated nodes and is therefore a path.

\begin{tcolorbox}
证明。如果从根节点 \({u}_{0}\) 到节点 \({u}_{m}\) 的行走 \({u}_{0},{u}_{1},\ldots ,{u}_{m}\) 存在某个 \(i < j\) 使得有重复节点 \({u}_{i} = {u}_{j}\) ，我们可以从该行走中移除节点 \({u}_{i + 1},\ldots ,{u}_{j}\) ，并得到从 \({u}_{0}\) 到 \({u}_{m}\) 的第二条行走路径，这是不允许的。因此，该行走没有重复节点，所以它是一条路径。
\end{tcolorbox}

Consider a tree node \(u\) which is not the root, and let \({vu}\) be the last edge on the unique path from the root to \(u\) . If there was a second edge \({wu}\) with \(w \neq  v\) , then the path from the root to \(w\) extended by \(u\) would give a second path from the root to \(u\) , which is not allowed. Hence, \(v\) is the unique parent of \(u\) . If the root had a parent \(v\) , the path from the root to \(v\) could be continued to a walk from the root to itself, which would be a second walk in addition to the empty path that consists of a single node. Again, this is not allowed.

\begin{tcolorbox}
考虑一个非根节点的树节点 \(u\) ，设 \({vu}\) 是从根节点到 \(u\) 的唯一路径上的最后一条边。如果存在第二条边 \({wu}\) 满足 \(w \neq  v\) ，那么从根节点到 \(w\) 再经过 \(u\) 的路径将构成从根节点到 \(u\) 的第二条路径，这是不允许的。因此， \(v\) 是 \(u\) 的唯一父节点。如果根节点有一个父节点 \(v\) ，从根节点到 \(v\) 的路径可以继续延伸为从根节点到其自身的一条行走路径，这将是除了由单个节点组成的空路径之外的第二条行走路径。同样，这也是不允许的。
\end{tcolorbox}

\begin{itemize}
\item For a decision node \(u\) , every edge \({uv}\) to a child \(v\) of \(u\) is assigned a different move. These moves are called the moves at \(u\) .
\end{itemize}

\begin{tcolorbox}
\begin{itemize}
\item 对于一个决策节点 \(u\) ，从该节点到其子节点 \(v\) 的每条边 \({uv}\) 都被赋予一个不同的行动。这些行动被称为在节点 \(u\) 处的行动。
\end{itemize}
\end{tcolorbox}

\begin{itemize}
\item For a chance node \(u\) , the edges \({uv}\) to the children \(v\) of \(u\) are assigned probabilities (that is, nonnegative reals that sum to 1).
\end{itemize}

\begin{tcolorbox}
\begin{itemize}
\item 对于一个机会节点 \(u\) ，从该节点到其子节点 \(v\) 的边 \({uv}\) 被赋予概率(即非负实数且其和为 1)。
\end{itemize}
\end{tcolorbox}

\begin{itemize}
\item Every leaf of the game tree has an \(N\) -tuple of real-valued payoffs, one for each player.
\end{itemize}

\begin{tcolorbox}
\begin{itemize}
\item 博弈树的每个叶节点(leaf)都有一个 \(N\) 元实值收益向量，每个玩家对应一个收益值。
\end{itemize}
\end{tcolorbox}

The game tree, with its decision nodes, moves, chance probabilities and payoffs, is known to the players, and defines the game completely. The game is played by starting at the root. At a decision node, the respective player chooses a move, which determines the next node. At a chance node, the move is made randomly according to the given probabilities. Play ends when a leaf is reached, where all players receive their payoffs.

\begin{tcolorbox}
博弈树及其决策节点、行动、机会概率和收益情况对玩家是已知的，并且完全定义了该博弈。博弈从根节点开始进行。在决策节点处，相应的玩家选择一个行动，该行动决定了下一个节点。在机会节点处，根据给定的概率随机进行行动。当到达叶节点时，博弈结束，此时所有玩家获得他们的收益。
\end{tcolorbox}

Players are interested in maximizing their own payoffs. If the outcome of playing the game is random, then the players are assumed to be interested in maximizing their expected payoffs.

\begin{tcolorbox}
玩家们关心的是最大化他们自己的收益。如果博弈的结果是随机的，那么假设玩家们关心的是最大化他们的期望效用(expected payoff)。
\end{tcolorbox}

\section*{Background material: Expected value of a finite random variable}

\begin{tcolorbox}
\section*{背景材料(background material):有限随机变量的期望值(expected value)}
\end{tcolorbox}

A random variable \(U\) takes real values that depend on a random event. Let \(U\) take finitely many values \({u}_{1},\ldots ,{u}_{k}\) that occur with probabilities \({p}_{1},\ldots ,{p}_{k}\) , respectively. Then the expected value or expectation of \(U\) is defined as

\begin{tcolorbox}
随机变量 \(U\) 取依赖于随机事件的实数值。设 \(U\) 取有限个值 \({u}_{1},\ldots ,{u}_{k}\) ，这些值分别以概率 \({p}_{1},\ldots ,{p}_{k}\) 出现。那么 \(U\) 的期望值或期望定义为
\end{tcolorbox}

\[
\mathbb{E}\left\lbrack  U\right\rbrack   = {p}_{1}{u}_{1} + \cdots  + {p}_{k}{u}_{k}. \tag{4.1}
\]

Because the probabilities \({p}_{1},\ldots ,{p}_{k}\) sum to 1, the expected value is the weighted average of the values \({u}_{i}\) with their probabilities \({p}_{i}\) as weights.

\begin{tcolorbox}
因为概率 \({p}_{1},\ldots ,{p}_{k}\) 之和为 1，所以期望值是以概率 \({p}_{i}\) 为权重的 \({u}_{i}\) 值的加权平均值。
\end{tcolorbox}

As a special case, if the random variable is deterministic and takes only a single value \({u}_{1}\) (with probability \({p}_{1} = 1\) ), then that value is the expected value, \(\mathbb{E}\left\lbrack  U\right\rbrack   = {u}_{1}.\)

\begin{tcolorbox}
作为一种特殊情况，如果随机变量是确定性的，并且只取单个值 \({u}_{1}\) (概率为 \({p}_{1} = 1\) )，那么该值就是期望值 \(\mathbb{E}\left\lbrack  U\right\rbrack   = {u}_{1}.\)
\end{tcolorbox}

In Figure 4.1, the expected payoff to player I after the chance move is \(\frac{1}{3}0 + \frac{2}{3}3 = 2\) , and to player II it is \(\frac{1}{3}{15} + \frac{2}{3}0 = 5\) . In this game, the chance node could therefore be replaced by a leaf with payoff 2 for player I and payoff 5 for player II. These expected payoffs reflect the players' preferences for random outcomes; Chapter 5 discusses in detail the assumptions that underly the use of such payoffs.

\begin{tcolorbox}
在图 4.1 中，机会行动之后玩家 I 的期望效用(expected payoff)是 \(\frac{1}{3}0 + \frac{2}{3}3 = 2\) ，玩家 II 的期望效用是 \(\frac{1}{3}{15} + \frac{2}{3}0 = 5\) 。在这个博弈中，机会节点因此可以用一个叶节点(leaf)来代替，该叶节点给玩家 I 的收益为 2，给玩家 II 的收益为 5。这些期望效用反映了玩家对随机结果的偏好(preference)；第 5 章详细讨论了使用此类收益的假设条件。
\end{tcolorbox}

In a two-player game, we normally name the players I and II and refer to player I as "he" and to player II as "she". Moves of player I are typically in upper case and of player II in lower case.

\begin{tcolorbox}
在两人博弈中，我们通常将玩家命名为玩家 I 和玩家 II，并将玩家 I 称为“他”，将玩家 II 称为“她”。玩家 I 的行动通常用大写字母表示，玩家 II 的行动通常用小写字母表示。
\end{tcolorbox}

Figure 4.2 demonstrates the tree property that every node is reached by a unique path from the root. This fails to hold in the left picture (a) for the middle leaf. Such a structure may make sense in a game. For example, the figure could represent two people in a pub where first player I chooses to pay for the drinks (move \(P\) ) or to accept (move \(A\) ) that the first round is paid by player II. In the second round, player II may then decide to either accept not to pay (move \(a\) ) or to pay(p). Then the players may only care how often, but not when, they have paid a round, with the middle payoff pair(1,1)for two possible "histories" of moves. The tree property prohibits states of play with more than one history, because the history is represented by the unique path from the root. Even if certain differences in the game history do not matter, these states of play are distinguished, mostly as a matter of mathematical convenience. The correct representation of the above game is shown in the right picture, Figure 4.2(b).

\begin{tcolorbox}
图4.2展示了博弈树(game tree)的一个特性，即从根节点出发，每个节点都有唯一的路径可达。在左图(a)中，中间的叶节点(leaf)不满足这一特性。这样的结构在博弈中可能是合理的。例如，该图可以表示酒吧里的两个人，首先玩家I可以选择为酒水买单(行动 \(P\) )，或者接受(行动 \(A\) )由玩家II支付第一轮的费用。在第二轮，玩家II可以决定要么接受不买单(行动 \(a\) )，要么买单(p)。然后玩家可能只关心他们买单的次数，而不关心买单的时间，中间的收益对(1,1)对应着两种可能的行动“历史”。博弈树的特性禁止出现具有多个历史的游戏状态，因为历史是由从根节点出发的唯一路径表示的。即使博弈历史中的某些差异无关紧要，这些游戏状态也会被区分开来，这主要是为了数学处理的方便。上述博弈的正确表示如图4.2(b)右图所示。
\end{tcolorbox}

\begin{center}
\includegraphics[max width=0.8\textwidth]{images/bo_d2e5cmbef24c73b03tn0_93_391_218_1138_465_0.jpg}
\end{center}
\hspace*{3em} 

Figure 4.2 The left picture (a) is not a tree because the leaf where both players receive payoff 1 is reachable by two different paths. The correct game tree representation is shown on the right in (b).

\begin{tcolorbox}
图4.2 左图(a)不是一棵博弈树，因为两个玩家收益都为1的叶节点可以通过两条不同的路径到达。正确的博弈树表示如右图(b)所示。
\end{tcolorbox}

Game trees are also called extensive games with perfect information. Perfect information means that a player always knows the state of play and therefore the complete history of play up to then. Game trees can be enhanced with an additional structure that represents "imperfect information", which is the topic of Chapter 10. Imperfect information is required to represent simultaneous moves. In the game trees considered in the present chapter, players move sequentially and observe each other's moves.

\begin{tcolorbox}
博弈树也被称为具有完美信息(perfect information)的扩展式博弈。完美信息意味着玩家总是知道游戏的状态，因此也知道到那时为止的完整游戏历史。博弈树可以通过额外的结构来表示“不完美信息”，这是第10章的主题。不完美信息用于表示同时行动。在本章所考虑的博弈树中，玩家按顺序行动并观察彼此的行动。
\end{tcolorbox}

Note how game trees differ from the combinatorial games in Chapter 1:

\begin{tcolorbox}
注意博弈树与第1章中的组合博弈有何不同:
\end{tcolorbox}

\begin{itemize}
\item A combinatorial game can be described very compactly, in particular when it is given as a sum of games. For general game trees, such game sums are usually not considered.
\end{itemize}

\begin{tcolorbox}
\begin{itemize}
\item 组合博弈可以非常简洁地描述，特别是当它表示为博弈之和(game sum)时。对于一般的博弈树，通常不考虑这种博弈之和。
\end{itemize}
\end{tcolorbox}

\begin{itemize}
\item The "rules" in a game tree are much more flexible: more than two players and chance moves are allowed, players do not have to alternate, and payoffs do not just stand for "win" or "lose".
\end{itemize}

\begin{tcolorbox}
\begin{itemize}
\item 博弈树中的“规则”更加灵活:允许有两个以上的玩家和随机行动，玩家不必轮流行动，收益也不只是代表“赢”或“输”。
\end{itemize}
\end{tcolorbox}

\begin{itemize}
\item The flexibility of game trees comes at a price, though: The game description is much longer than a combinatorial game. For example, a simple instance of Nim may require a huge game tree. Regularities like the mex rule do not apply to general game trees.
\end{itemize}

\begin{tcolorbox}
\begin{itemize}
\item 然而，博弈树的灵活性是有代价的:博弈的描述比组合博弈长得多。例如，一个简单的尼姆游戏(Nim)可能需要一个巨大的博弈树。像mex规则(mex rule)这样的规律并不适用于一般的博弈树。
\end{itemize}
\end{tcolorbox}

\section*{4.3 Backward Induction}

\begin{tcolorbox}
\section*{4.3 逆向归纳法(backward induction)}
\end{tcolorbox}

Which moves should the players choose in a game tree? "Optimal" play should maximize a player's payoff. This can be decided irrespective of other players' actions when the player is the last player to move. In the game in Figure 4.3(a), player II maximizes her payoff by move \(r\) . Going backward in time, player I makes his move \(T\) or \(B\) at the root of the game tree, where he will receive either 1 or 2, assuming the described future behavior of player II. Consequently, he will choose \(B\) , and play ends with payoff 2 to each player. These chosen moves are shown in Figure 4.3(b) with arrows on the edges; similar to the boxes that we put around best-response payoffs in a strategic-form game, this is additional information to help the analysis of the game and not part of the game specification.

\begin{tcolorbox}
在博弈树中，玩家应该选择哪些行动呢？“最优”行动应该使玩家的收益最大化。当玩家是最后行动的玩家时，这可以不考虑其他玩家的行动来决定。在图4.3(a)的博弈中，玩家II通过行动 \(r\) 使自己的收益最大化。逆向推导，玩家I在博弈树的根节点选择行动 \(T\) 或 \(B\) ，假设玩家II按上述方式行动，他将获得1或2的收益。因此，他会选择 \(B\) ，游戏结束时每个玩家的收益为2。这些选择的行动在图4.3(b)中用边上的箭头表示；类似于我们在战略式博弈中给最优反应收益加上的方框，这是帮助分析博弈的额外信息，而不是博弈规范的一部分。
\end{tcolorbox}

\begin{center}
\includegraphics[max width=0.8\textwidth]{images/bo_d2e5cmbef24c73b03tn0_94_285_533_1056_463_0.jpg}
\end{center}
\hspace*{3em} 

Figure 4.3 A two-player game tree (a) with backward induction moves (b).

\begin{tcolorbox}
图4.3 一个两人博弈树(a)及逆向归纳法行动(b)。
\end{tcolorbox}

This process is called backward induction: Starting with the decision nodes closest to the leaves, a player's move is chosen that maximizes that player's payoff at the node. In general, a move is chosen in this way for each decision node provided all subsequent moves have already been decided. Eventually, this will determine a move for every decision node, and hence for the entire game.

\begin{tcolorbox}
这个过程称为逆向归纳法:从最接近叶节点的决策节点开始，为每个决策节点选择一个能使该节点上玩家收益最大化的行动。一般来说，只要所有后续行动都已经确定，就以这种方式为每个决策节点选择一个行动。最终，这将为每个决策节点确定一个行动，从而为整个博弈确定行动。
\end{tcolorbox}

\begin{center}
\includegraphics[max width=0.5\textwidth]{images/bo_d2e5cmbef24c73b03tn0_94_203_1399_672_624_0.jpg}
\end{center}
\hspace*{3em} 

Figure 4.4 The game tree in Figure 4.1 with the chance move replaced by its expected payoffs, and backward induction moves that depend on player II's choice between \(b\) (white arrow, then two possible choices \(X\) or \(Y\) by player I) and \(c\) (black arrow, unique choice \(X\) by player \(\mathrm{I}\) ).

\begin{tcolorbox}
图4.4 图4.1中的博弈树，随机行动被其期望效用(expected payoff)所取代，逆向归纳法行动取决于玩家II在 \(b\) (白色箭头，然后玩家I有两种可能的选择 \(X\) 或 \(Y\) )和 \(c\) (黑色箭头，玩家 \(\mathrm{I}\) 唯一的选择 \(X\) )之间的选择。
\end{tcolorbox}

The move selected by backward induction is not necessarily unique, if there is more than one move that gives maximal payoff to the player. This may influence the choice of moves earlier in the tree. In the game in Figure 4.4, backward induction chooses either move \(b\) or move \(c\) for player II, both of which give her payoff 5 (which is an expected payoff for move \(b\) in the original game in Figure 4.1) that exceeds her payoff 4 for move \(a\) . At the rightmost node, player I chooses \(Q\) . This determines the preceding move \(d\) by player II which gives her the higher payoff 3 as opposed to 2 (via move \(Q\) ). In turn, this means that player I, when choosing between \(X,Y\) , or \(Z\) at the root of the game tree, will get payoff 2 for \(Y\) and payoff 1 for \(Z\) . The payoff when he chooses \(X\) depends on the choice of player II: if that is \(b\) , then player I gets 2, and he can choose either \(X\) or \(Y\) , both of which give him maximal payoff 2. These choices are shown with white arrows in Figure 4.4. If player II chooses \(c\) , however, then the payoff to player I is 4 when choosing \(X\) , so this is the unique optimal move.

\begin{tcolorbox}
如果有不止一种行动能为玩家带来最大效用，那么逆向归纳法(backward induction)所选择的行动不一定是唯一的。这可能会影响博弈树(game tree)中早期的行动选择。在图4.4的博弈中，逆向归纳法为玩家II选择行动 \(b\) 或行动 \(c\) ，这两种行动都能为她带来效用5(这是图4.1原始博弈中行动 \(b\) 的期望效用(expected payoff))，超过了她选择行动 \(a\) 时的效用4。在最右侧的节点(node)，玩家I选择 \(Q\) 。这决定了玩家II的前一个行动 \(d\) ，该行动为她带来了更高的效用3，而不是2(通过行动 \(Q\) )。反过来，这意味着玩家I在博弈树的根节点选择 \(X,Y\) 或 \(Z\) 时，选择 \(Y\) 能得到效用2，选择 \(Z\) 能得到效用1。他选择 \(X\) 时的效用取决于玩家II的选择:如果玩家II选择 \(b\) ，那么玩家I得到2，他可以选择 \(X\) 或 \(Y\) ，这两种选择都能为他带来最大效用2。这些选择在图4.4中用白色箭头表示。然而，如果玩家II选择 \(c\) ，那么玩家I选择 \(X\) 时的效用是4，所以这是唯一的最优行动。
\end{tcolorbox}

To summarize, the possible combinations of moves that can arise in Figure 4.1 by backward induction are, by listing the moves for each player: \(\left( {{XQ},{bd}}\right) ,\left( {{XQ},{cd}}\right)\) , and(YQ, bd). Note that \(Q\) and \(d\) are always chosen, but that \(Y\) can only be chosen in combination with the move \(b\) by player II.

\begin{tcolorbox}
综上所述，通过逆向归纳法(backward induction)在图4.1中可能出现的行动组合，按每个玩家的行动列出为: \(\left( {{XQ},{bd}}\right) ,\left( {{XQ},{cd}}\right)\) 和(YQ, bd)。请注意， \(Q\) 和 \(d\) 总是会被选择，但 \(Y\) 只能与玩家II的行动 \(b\) 组合选择。
\end{tcolorbox}

We do not try to argue that player II "should not" choose \(b\) because of the "risk" that player I might choose \(Y\) . This would be a "refinement" of backward induction that we do not pursue.

\begin{tcolorbox}
我们并不试图争辩说，由于参与者I可能选择 \(b\) 的“风险”，参与者II“不应该”选择 \(Y\) 。这将是我们不探讨的逆向归纳法(backward induction)的一种“精炼”。
\end{tcolorbox}

The moves determined by backward induction are therefore, in general, not unique, and possibly interdependent.

\begin{tcolorbox}
因此，一般来说，逆向归纳法(backward induction)所确定的行动并非唯一，并且可能相互依赖。
\end{tcolorbox}

Backward induction gives a unique recommendation to each player if there is always only one move that gives maximal payoff. This applies to generic games. A generic game is a game where the payoff parameters are real numbers with no special dependency between them (like two payoffs being equal). In particular, it should be allowed to replace them with values nearby. For example, the payoffs may be given in some practical scenario which has some "noise" that effects the precise value of each payoff. Then two such payoffs are equal with probability zero, and so this case can be disregarded. In generic games, the optimal move is always unique, so that backward induction gives a unique result.

\begin{tcolorbox}
如果总是只有一种行动能带来最大效用，那么逆向归纳法(backward induction)会为每个参与者提供唯一的建议。这适用于通有博弈(generic game)。通有博弈是指效用参数为实数且它们之间没有特殊依赖关系(如两个效用相等)的博弈。特别是，应该允许用附近的值替换它们。例如，效用可能是在某个存在影响每个效用精确值的“噪声”的实际场景中给出的。那么，两个这样的效用相等的概率为零，因此这种情况可以忽略不计。在通有博弈中，最优行动总是唯一的，因此逆向归纳法(backward induction)会给出唯一的结果。
\end{tcolorbox}

\section*{4.4 Strategies in Game Trees}

\begin{tcolorbox}
\section*{4.4 博弈树(game tree)中的策略}
\end{tcolorbox}

In a game in strategic form, each player's strategies are assumed as given. For a game tree, a strategy is a derived concept. The best way to think of a strategy in a game tree is as a plan of action for every decision node.

\begin{tcolorbox}
在策略形式(strategic form)的博弈中，每个参与者的策略被视为给定的。对于博弈树(game tree)，策略是一个派生概念。将博弈树(game tree)中的策略视为每个决策节点(node)的行动计划是最好的理解方式。
\end{tcolorbox}

Definition 4.3. In a game tree, a strategy of a player specifies a move for every decision node of that player.

\begin{tcolorbox}
定义4.3. 在博弈树(game tree)中，参与者的策略为该参与者的每个决策节点(node)指定一个行动。
\end{tcolorbox}

Hence, if a player has \(k\) decision nodes, a strategy of that player is a \(k\) -tuple of moves. For simplicity of notation, we often write down a strategy as a list of moves, one for each decision node of the player (and not as a comma-separated tuple surrounded by parentheses). In the game tree in Figure 4.1, the strategies of player I are \({XP},{XQ},{YP},{YQ},{ZP},{ZQ}\) . The strategies of player II are \({ad},{ae},{bd}\) , be, \({cd},{ce}\) . When specifying a strategy in that way, this must be done with respect to a fixed order of the decision nodes in the tree, in order to identify each move uniquely. This matters when a move name appears more than once, for example in Figure 4.2(b). In that tree, the strategies of player II are \({aa},{ap},{pa},{pp}\) , with the understanding that the first move refers to the left decision node and the second move to the right decision node of player II.

\begin{tcolorbox}
因此，如果一个参与者有 \(k\) 个决策节点(node)，那么该参与者的策略就是一个由 \(k\) 个行动组成的元组。为了符号表示的简便，我们通常将策略写成一个行动列表，每个行动对应参与者的一个决策节点(node)(而不是用括号括起来、用逗号分隔的元组)。在图4.1的博弈树(game tree)中，参与者I的策略是 \({XP},{XQ},{YP},{YQ},{ZP},{ZQ}\) 。参与者II的策略是 \({ad},{ae},{bd}\) 、be、 \({cd},{ce}\) 。以这种方式指定策略时，必须按照树中决策节点(node)的固定顺序进行，以便唯一地识别每个行动。当一个行动名称出现多次时，这一点很重要，例如在图4.2(b)中。在该树中，参与者II的策略是 \({aa},{ap},{pa},{pp}\) ，要理解为第一个行动对应参与者II的左决策节点(node)，第二个行动对应右决策节点(node)。
\end{tcolorbox}

Backward induction defines a move for every decision node of the game tree, and therefore for every decision node of each player, which in turn gives a strategy for each player. The result of backward induction is therefore a strategy profile.

\begin{tcolorbox}
逆向归纳法(backward induction)为博弈树(game tree)的每个决策节点(node)定义一个行动，因此也为每个参与者的每个决策节点(node)定义一个行动，进而为每个参与者给出一个策略。因此，逆向归纳法(backward induction)的结果是一个策略组合。
\end{tcolorbox}

With a strategy for each player, we can define the strategic form of the game tree, which is determined as follows. Consider a strategy profile, and assume that players move according to their strategies. If there are no chance moves, their play leads to a unique leaf. If there are chance moves, a strategy profile may lead to a probability distribution on the leaves of the game tree, with resulting expected payoffs. For chance moves that occur earlier in the tree, the computation of expected payoffs is more complicated than in Figure 4.1; see Figure 4.8 for an example. In general, any strategy profile defines an expected payoff to each player; this includes the deterministic case where payoffs and expected payoffs conincide, as mentioned after (4.1).

\begin{tcolorbox}
有了每个参与者的策略，我们就可以定义博弈树(game tree)的策略形式(strategic form)，其确定方式如下。考虑一个策略组合，并假设参与者按照他们的策略行动。如果没有随机行动，他们的行动会导致一个唯一的叶节点(leaf)。如果有随机行动，一个策略组合可能会导致博弈树(game tree)的叶节点(leaf)上出现一个概率分布，从而产生期望效用(expected payoff)。对于在树中较早出现的随机行动，期望效用(expected payoff)的计算比图4.1中的情况更复杂；见图4.8的示例。一般来说，任何策略组合都会为每个参与者定义一个期望效用(expected payoff)；这包括如(4.1)之后所提到的效用和期望效用(expected payoff)一致的确定性情况。
\end{tcolorbox}

Definition 4.4. The strategic form of a game tree is defined by the set of strategies for each player according to Definition 4.3, and the expected payoff to each player resulting from each strategy profile.

\begin{tcolorbox}
定义4.4. 博弈树(game tree)的策略形式(strategic form)由根据定义4.3为每个参与者设定的策略集以及每个策略组合为每个参与者带来的期望效用(expected payoff)来定义。
\end{tcolorbox}

We first consider the simple game tree in Figure 4.3, which is shown in Figure 4.5 together with its strategic form (already encountered in Figure 3.6), which is a \(2 \times  2\) game. In this game, each player has only one decision node, so a strategy is the same as a move. (In general, a strategy specifies several moves, so do not use the terms "strategy" and "move" interchangeably!) Note also that we write the payoff to player I at the top at each leaf of the game tree, but at the bottom left of a cell in the payoff table.

\begin{tcolorbox}
我们首先考虑图4.3中的简单博弈树(game tree)，它与它的策略形式(已经在图3.6中见过)一起显示在图4.5中，这是一个 \(2 \times  2\) 博弈。在这个博弈中，每个参与者只有一个决策节点(node)，因此策略与行动相同。(一般来说，策略指定多个行动，因此不要将“策略”和“行动”这两个术语互换使用！)还要注意，我们在博弈树(game tree)的每个叶节点(leaf)的顶部写出参与者I的效用，但在效用表中一个单元格的左下角写出。
\end{tcolorbox}

The game tree in Figure 4.3 has only three leaves (and no chance moves) but the strategic form has four cells. Therefore, the payoffs for one of the leaves appear twice, namely for the strategy pairs(T, l)and(T, r). The reason is that the leftmost leaf is reached immediately after move \(T\) of player \(\mathrm{I}\) , and then the move \(l\) or \(r\) of player II does not matter.

\begin{tcolorbox}
图4.3中的博弈树(game tree)只有三个叶节点(leaf)(且没有随机行动)，但策略形式有四个单元格。因此，其中一个叶节点(leaf)的效用出现了两次，即对于策略对(T, l)和(T, r)。原因是，在参与者 \(\mathrm{I}\) 采取行动 \(T\) 后立即到达最左边的叶节点(leaf)，然后参与者II的行动 \(l\) 或 \(r\) 就无关紧要了。
\end{tcolorbox}

\begin{center}
\includegraphics[max width=0.8\textwidth]{images/bo_d2e5cmbef24c73b03tn0_97_422_221_1083_459_0.jpg}
\end{center}
\hspace*{3em} 

Figure 4.5 Game tree (a) of the Threat game with backward induction moves, and its strategic form (b) with best responses. Note: In the game tree, the payoffs to player I are shown at the top, but in the strategic form they appear at the bottom left of each cell.

\begin{tcolorbox}
图4.5 带有逆向归纳法步骤的威胁博弈的博弈树(a)及其带有最优反应的策略形式(b)。注意:在博弈树中，玩家I的收益显示在顶部，但在策略形式中，它们出现在每个单元格的左下角。
\end{tcolorbox}

The game tree and the strategic form specify the same game, but are played differently. In the game tree, play proceeds sequentially. In the strategic form, the players choose their strategies simultaneously, and then they play according to these "plans of action".

\begin{tcolorbox}
博弈树和策略形式描述的是同一个博弈，但玩法不同。在博弈树中，博弈按顺序进行。在策略形式中，玩家同时选择他们的策略，然后根据这些“行动计划”进行博弈。
\end{tcolorbox}

In the game in Figure 4.5, backward induction specifies the unique strategy profile(B, r), which is also an equilibrium (which holds generally, see Theorem 4.6 below). The game has a second equilibrium(T, l), which can also be seen from the game tree: If player II chooses \(l\) , then player I’s best response is \(T\) . Player II can threaten this suboptimal move because she never has to make it when player I chooses \(T\) , as he is incentivized to do. The game is also known as the Threat game for that reason. Player II gets a higher payoff in this equilibrium(T, r). We will distinguish equilibria that are obtained by backward induction and those that are not, as discussed in Section 4.6.

\begin{tcolorbox}
在图4.5的博弈中，逆向归纳法确定了唯一的策略组合(B，r)，这也是一个均衡(一般情况下都成立，见下面的定理4.6)。该博弈还有第二个均衡(T，l)，这也可以从博弈树中看出:如果玩家II选择 \(l\) ，那么玩家I的最优反应是 \(T\) 。玩家II可以威胁采取这种次优行动，因为当玩家I选择 \(T\) 时(他有动机这么做)，她永远不必采取这一行动。出于这个原因，该博弈也被称为威胁博弈。在这个均衡(T，r)中，玩家II获得了更高的收益。我们将区分通过逆向归纳法得到的均衡和未通过逆向归纳法得到的均衡，如4.6节所述。
\end{tcolorbox}

Figure 4.6 shows the strategic form of the game tree in Figure 4.1. In that game, players I and II both have two decision nodes with three and two moves, respectively, and therefore \(3 \times  2\) strategies each. The strategic form has 36 cells, which have many repetitions because the game tree has only eight (or, after contracting the two chance outcomes as in Figure 4.4, seven) leaves.

\begin{tcolorbox}
图4.6显示了图4.1中博弈树的策略形式。在该博弈中，玩家I和玩家II分别有两个决策节点，分别有三种和两种行动，因此每个玩家有 \(3 \times  2\) 种策略。策略形式有36个单元格，其中有很多重复，因为博弈树只有八个叶节点(或者，如图4.4那样合并两个随机结果后，有七个叶节点)。
\end{tcolorbox}

\section*{4.5 Reduced Strategies}

\begin{tcolorbox}
\section*{4.5 精简策略}
\end{tcolorbox}

In Figure 4.6, the two rows \({XP}\) and \({XQ}\) have identical payoffs for both players, and similarly the two rows \({YP}\) and \({YQ}\) . The reason is that the strategies \({XP}\) and \({XQ}\) specify that player I chooses \(X\) at his first decision node and either \(P\) or \(Q\) at his second decision node. However, after choosing \(X\) , the second decision node cannot be reached. The second decision node can only be reached if player I chooses \(Z\) as his first move, and then the two strategies \({ZP}\) and \({ZQ}\) may lead to different leaves (if player II chooses \(e\) , as in her three strategies \({ae},{be}\) , or \({ce}\) ). Because player I himself chooses between \(X,Y\) , or \(Z\) , it makes sense to replace both strategies \({XP}\) and \({XQ}\) by a less specific "plan of action" \(X *\) that precribes only move \(X\) . Similarly, the two strategies \({YP}\) and \({YQ}\) are collapsed into \(Y\) *. We always use the star " \(*\) " as a placeholder. It stands for any unspecified move at the respective unreachable decision node, to identify the node where this move is left unspecified.

\begin{tcolorbox}
在图4.6中，两行 \({XP}\) 和 \({XQ}\) 对两个玩家来说收益相同，同样，两行 \({YP}\) 和 \({YQ}\) 也是如此。原因是策略 \({XP}\) 和 \({XQ}\) 规定玩家I在他的第一个决策节点选择 \(X\) ，在他的第二个决策节点选择 \(P\) 或 \(Q\) 。然而，在选择 \(X\) 之后，无法到达第二个决策节点。只有当玩家I将 \(Z\) 作为他的第一步行动时，才能到达第二个决策节点，然后两种策略 \({ZP}\) 和 \({ZQ}\) 可能会导致不同的叶节点(如果玩家II选择 \(e\) ，就像她的三种策略 \({ae},{be}\) 或 \({ce}\) 那样)。因为玩家I自己在 \(X,Y\) 或 \(Z\) 之间进行选择，所以用一个不太具体的“行动计划” \(X *\) 来取代策略 \({XP}\) 和 \({XQ}\) 是有意义的，该“行动计划”只规定行动 \(X\) 。同样，两种策略 \({YP}\) 和 \({YQ}\) 合并为 \(Y\) *。我们总是使用星号“ \(*\) ”作为占位符。它代表在相应无法到达的决策节点上的任何未指定的行动，以确定该行动未被指定的节点。
\end{tcolorbox}

\begin{center}
\includegraphics[max width=0.7\textwidth]{images/bo_d2e5cmbef24c73b03tn0_98_355_195_909_889_0.jpg}
\end{center}
\hspace*{3em} 

Figure 4.6 Strategic form of the extensive game in Figure 4.1. The less redundant reduced strategic form is shown in Figure 4.7 below.

\begin{tcolorbox}
图4.6 图4.1中扩展式博弈的策略形式。冗余较少的精简策略型如图4.7所示。
\end{tcolorbox}

Leaving moves at unreachable nodes unspecified in this manner defines a reduced strategy according to the following definition. Because the resulting expected payoffs remain uniquely defined, tabulating these reduced strategies and the payoff for the resulting reduced strategy profiles gives the reduced strategic form of the game. The reduced strategic form of the game tree of Figure 4.1 is shown in Figure 4.7, with indicated best-response payoffs.

\begin{tcolorbox}
以这种方式对无法到达的节点上的行动不做具体规定，根据以下定义确定了一个精简策略。由于由此产生的期望效用仍然是唯一确定的，将这些精简策略和由此产生的精简策略组合的收益制成表格，就得到了博弈的精简策略型。图4.1中博弈树的精简策略型如图4.7所示，并标明了最优反应收益。
\end{tcolorbox}

Definition 4.5. In a game tree, a reduced strategy of a player specifies a move for every decision node of that player, except for those decision nodes that are unreachable due to an earlier own move, where the move is replaced by "*". A reduced strategy profile is a tuple of reduced strategies, one for each player of the game. The reduced strategic form of a game tree lists all reduced strategies for each player, and tabulates the expected payoff to each player for each reduced strategy profile.

\begin{tcolorbox}
定义4.5。在博弈树中，玩家的精简策略为该玩家的每个决策节点指定一个行动，但由于玩家自己先前的行动而无法到达的决策节点除外，在这些节点上，行动用“*”代替。精简策略组合是一个精简策略的元组，每个玩家对应一个。博弈树的精简策略型列出了每个玩家的所有精简策略，并将每个精简策略组合下每个玩家的期望效用制成表格。
\end{tcolorbox}

\begin{center}
\includegraphics[max width=0.7\textwidth]{images/bo_d2e5cmbef24c73b03tn0_99_503_197_898_622_0.jpg}
\end{center}
\hspace*{3em} 

Figure 4.7 Reduced strategic form of the extensive game in Figure 4.1. The star * stands for an arbitrary move at the second decision node of player I, which is not reachable after move \(X\) or \(Y\) .

\begin{tcolorbox}
图4.7 图4.1中扩展式博弈的精简策略型。星号*代表玩家I在第二个决策节点上的任意行动，在行动 \(X\) 或 \(Y\) 之后无法到达该节点。
\end{tcolorbox}

The preceding definition generalizes Definitions 4.3 and 4.4. It is important that the only moves that are left unspecified are at decision nodes which are unreachable due to an earlier own move of the player. A reduced strategy must not disregard a move because another player may not move there, because that possibility cannot be excluded by looking only at the player's own moves. In the game tree in Figure 4.1, for example, no reduction is possible for the strategies of player II, because neither of her moves at one decision node precludes a move at her other decision node. Therefore, the reduced strategies of player II in that game are the same as her (unreduced, original) strategies.

\begin{tcolorbox}
上述定义对定义 4.3 和 4.4 进行了推广。重要的是，唯一未明确指定的行动是在由于玩家自身先前的行动而无法到达的决策节点处。一个精简策略不能因为另一个玩家可能不会在某个节点行动就忽略该行动，因为仅通过查看玩家自身的行动无法排除这种可能性。例如，在图 4.1 的博弈树中，玩家 II 的策略无法进行精简，因为她在一个决策节点的任何行动都不会排除她在另一个决策节点的行动。因此，在该博弈中玩家 II 的精简策略与她的(未精简的、原始的)策略相同。
\end{tcolorbox}

In Definition 4.5, reduced strategies are defined without any reference to payoffs, only by the structure of the game tree. One could derive the reduced strategic form in Figure 4.7 from the strategic form in Figure 4.6 by identifying the strategies that have identical payoffs for both players and replacing them by a single reduced strategy. Some authors define the reduced strategic form in terms of this elimination of duplicate strategies in the strategic form. (Sometimes, dominated strategies are also eliminated when defining the reduced strategic form.) We define the reduced strategic form without any reference to payoffs because strategies may have identical payoffs by accident, due to special payoffs at the leaves of the game (this does not occur in generic games). Moreover, we prefer that reduced strategies refer to the structure of the game tree, not to some relationship of the payoffs, which is a different aspect of the game.

\begin{tcolorbox}
在定义 4.5 中，精简策略的定义与收益无关，仅取决于博弈树的结构。可以通过识别对两个玩家都有相同收益的策略，并用一个单一的精简策略替换它们，从图 4.6 的策略形式推导出图 4.7 的精简策略型。一些作者根据策略形式中重复策略的消除来定义精简策略型。(有时，在定义精简策略型时也会消除(被占优) 策略。)我们在定义精简策略型时不考虑收益，因为由于博弈叶节点的特殊收益，策略可能偶然具有相同的收益(这在通有博弈中不会发生)。此外，我们更倾向于让精简策略参考博弈树的结构，而不是收益的某种关系，这是博弈的不同方面。
\end{tcolorbox}

The following is a recipe for filling, reasonably quickly, the payoff table for the reduced strategic form (or, if one wishes, the full strategic form) of a game tree like Figure 4.1:

\begin{tcolorbox}
以下是一个相对快速地填写类似图 4.1 这样的博弈树的精简策略型(或者，如果愿意，完整策略型)收益表的方法:
\end{tcolorbox}

\begin{itemize}
\item For each player, choose a fixed order of the decision nodes, where nodes that come earlier in the tree should come first. If the player has \(k\) decision nodes, a (reduced) strategy has \(k\) "slots" to define a move, one for each decision node.
\end{itemize}

\begin{tcolorbox}
\begin{itemize}
\item 为每个玩家选择一个固定的决策节点顺序，树中较早出现的节点应排在前面。如果玩家有 \(k\) 个决策节点，一个(精简)策略有 \(k\) 个“空位”来定义一个行动，每个决策节点对应一个空位。
\end{itemize}
\end{tcolorbox}

\begin{itemize}
\item Generate all (reduced) strategies. A reduced strategy defines for each slot a move, or an unspecified move " \(*\) " if the respective decision node is unreachable due to an earlier move of the same player. For two players, list the (reduced) strategies of player I as rows and of player II as columns.
\end{itemize}

\begin{tcolorbox}
\begin{itemize}
\item 生成所有(精简)策略。一个精简策略为每个空位定义一个行动，如果相应的决策节点由于同一玩家先前的行动而无法到达，则定义为未指定的行动“ \(*\) ”。对于两个玩家，将玩家 I 的(精简)策略列为行，玩家 II 的(精简)策略列为列。
\end{itemize}
\end{tcolorbox}

\begin{itemize}
\item For each leaf of the game tree, consider the path from the root to the leaf. Assume first that there are no chance moves on that path. For each player, consider the moves of that player on that path. For example, the leftmost leaf in Figure 4.1 with payoffs(3,4)has move \(X\) for player I and move \(a\) for player II.
\end{itemize}

\begin{tcolorbox}
\begin{itemize}
\item 对于博弈树的每个叶节点，考虑从根节点到该叶节点的路径。首先假设该路径上没有随机行动。对于每个玩家，考虑该玩家在该路径上的行动。例如，图 4.1 中最左边的叶节点，收益为(3,4)，玩家 I 的行动为 \(X\) ，玩家 II 的行动为 \(a\) 。
\end{itemize}
\end{tcolorbox}

\begin{itemize}
\item For each player, identify all strategies that agree with the moves on the path to the leaf. For the leftmost leaf in Figure 4.1 and the reduced strategic form in Figure 4.7, these strategies are given by the reduced strategy \(X *\) for player \(\mathrm{I}\) (or the unreduced strategies \({XP}\) and \({XQ}\) in Figure 4.6), and the strategies \({ad}\) and ae for player II. For all these strategy pairs, enter the payoff pair, here (3,4), in the respective cells. Repeat this for each leaf.
\end{itemize}

\begin{tcolorbox}
\begin{itemize}
\item 对于每个玩家，识别所有与到达叶节点路径上的行动一致的策略。对于图 4.1 中最左边的叶节点和图 4.7 中的精简策略型，这些策略对于玩家 \(\mathrm{I}\) 由精简策略 \(X *\) 给出(或者在图 4.6 中是未精简策略 \({XP}\) 和 \({XQ}\) )，对于玩家 II 是策略 \({ad}\) 和 ae 。对于所有这些策略对，将收益对(这里是(3,4))填入相应的单元格。对每个叶节点重复此操作。
\end{itemize}
\end{tcolorbox}

As another example, the leaf with payoffs(2,3)that directly follows move \(Y\) is reached by move \(Y\) of player I (which defines row \(Y *\) in the reduced strategic form in Figure 4.7), and no moves of player II, so all strategies of player II agree with these moves; that is, the entire row \(Y *\) is filled with the payoff pair(2,3).

\begin{tcolorbox}
再举一个例子，直接跟随行动 \(Y\) 且收益为(2,3)的叶节点，是由玩家 I 的行动 \(Y\) 到达的(该行动在图 4.7 的精简策略型中定义了行 \(Y *\) )，玩家 II 没有行动，所以玩家 II 的所有策略都与这些行动一致；也就是说，整行 \(Y *\) 都填入收益对(2,3)。
\end{tcolorbox}

\begin{itemize}
\item If there are chance moves on the path from the root to the leaf, let \(p\) be the product of chance probabilities on that path. Multiply each payoff at the leaf with \(p\) , and add the result to the current payoffs in the respective cells that are determined as before (assuming the payoffs in each cell are initially zero). For an example see Figure 4.8.
\end{itemize}

\begin{tcolorbox}
\begin{itemize}
\item 如果从根节点到叶节点的路径上有随机行动，令 \(p\) 为该路径上随机概率的乘积。将叶节点的每个收益乘以 \(p\) ，并将结果加到如前所述确定的相应单元格的当前收益中(假设每个单元格的初始收益为零)。示例见图 4.8。
\end{itemize}
\end{tcolorbox}

This recipe for creating the strategic form of a game tree is much less laborious than considering each strategy pair and identifying to which leaf it leads (or leaves, in case there are chance moves).

\begin{tcolorbox}
这种创建博弈树策略型的方法比考虑每对策略并确定它通向哪个叶节点(或者，如果有随机行动，则是哪些叶节点)要省力得多。
\end{tcolorbox}

However, it should be clear that the strategic form is typically substantially larger than the game tree. It is therefore often easier to directly analyze the game tree rather than the strategic form.

\begin{tcolorbox}
然而，应该清楚的是，策略型通常比博弈树大得多。因此，直接分析博弈树通常比分析策略型更容易。
\end{tcolorbox}

Strategies are combinations of moves, so for every additional decision node of a player, each move at that node can be combined with move combinations already considered. The number of move combinations therefore grows exponentially with

\begin{tcolorbox}
策略是行动的组合，因此对于玩家的每个额外决策节点，该节点的每个行动都可以与已经考虑过的行动组合相结合。因此，行动组合的数量会随着
\end{tcolorbox}

\begin{center}
\includegraphics[max width=0.8\textwidth]{images/bo_d2e5cmbef24c73b03tn0_101_390_219_1139_459_0.jpg}
\end{center}
\hspace*{3em} 

Figure 4.8 Creating the strategic form by iterating through tree leaves when there are chance moves, which lead to expected payoffs.

\begin{tcolorbox}
图4.8 当存在导致期望效用的随机行动时，通过遍历博弈树的叶节点来创建策略形式。
\end{tcolorbox}

\begin{center}
\includegraphics[max width=0.5\textwidth]{images/bo_d2e5cmbef24c73b03tn0_101_369_844_644_379_0.jpg}
\end{center}
\hspace*{3em} 

Figure 4.9 Game tree (without payoffs) with \(m\) moves \({C}_{1},\ldots ,{C}_{m}\) of player I at the root of the game tree. After each move \({C}_{i}\) , player II has two moves \({l}_{i}\) and \({r}_{i}\) . Player II has \({2}^{m}\) strategies.

\begin{tcolorbox}
图4.9 博弈树(无效用值)，玩家I在博弈树的根节点有 \(m\) 个行动 \({C}_{1},\ldots ,{C}_{m}\) 。在玩家I的每次行动 \({C}_{i}\) 之后，玩家II有两个行动 \({l}_{i}\) 和 \({r}_{i}\) 。玩家II有 \({2}^{m}\) 个策略。
\end{tcolorbox}

the number of decision nodes of the player, because it is the product of the numbers of moves at each node. In the game in Figure 4.9, for example, where \(m\) possible initial moves of player I are followed each time by two possible moves of player II, the number of strategies of player II is \({2}^{m}\) . Moreover, this is also the number of reduced strategies of player II because no reduction is possible. This shows that even the reduced strategic form can be exponentially large in the size of the game tree. If a player's move is preceded by an own earlier move, the reduced strategic form is smaller because then that move can be left unspecified in any reduced strategy that does not choose the preceding own move.

\begin{tcolorbox}
玩家的决策节点数量，因为它是每个节点的行动数量的乘积。例如，在图4.9的博弈中，玩家I有 \(m\) 种可能的初始行动，每次之后玩家II都有两种可能的行动，玩家II的策略数量是 \({2}^{m}\) 。此外，这也是玩家II的精简策略的数量，因为无法进行精简。这表明，即使是精简策略型在博弈树的规模上也可能呈指数级增长。如果一个玩家的行动之前有自己的早期行动，那么精简策略型会更小，因为在任何不选择前面自己行动的精简策略中，该行动可以不被确定。
\end{tcolorbox}

\(\Rightarrow\) Exercise 4.4 is about reduced strategies.

\begin{tcolorbox}
\(\Rightarrow\) 练习4.4是关于精简策略的。
\end{tcolorbox}

In the reduced strategic form, an equilibrium is defined, as in Definition 3.1, as a profile of reduced strategies, each of which is a best response to the others. In this context, we will for brevity sometimes refer to reduced strategies simply as "strategies". This is justified because, when looking at the reduced strategic form, the concepts of dominance and equilibrium can be applied directly to the reduced strategic form, for example to the table that defines a two-player game. Then "strategy" means simply a row or column of that table. The term "reduced strategy" is only relevant when referring to the game tree.

\begin{tcolorbox}
在精简策略型中，如同定义3.1一样，均衡被定义为一个精简策略的组合，其中每个策略都是对其他策略的最优反应。在这种情况下，为了简洁起见，我们有时会简单地将精简策略称为“策略”。这是合理的，因为在考虑精简策略型时，占优和均衡的概念可以直接应用于精简策略型，例如应用于定义两人博弈的表格。那么“策略”就简单地意味着该表格中的一行或一列。“精简策略”这个术语仅在提及博弈树时才相关。
\end{tcolorbox}

The equilibria in Figure 4.7 are identified as those pairs of (reduced) strategies that are best responses to each other, with both payoffs surrounded by a box in the respective cell of the table. These equilibria in reduced strategies are \(\left( {X*,{bd}}\right)\) , \(\left( {X*,{cd}}\right) ,\left( {X*,{ce}}\right)\) , and \(\left( {Y*,{bd}}\right)\) .

\begin{tcolorbox}
图4.7中的均衡被确定为那些(精简)策略对，它们彼此是最优反应，并且在表格的相应单元格中，两个效用值都被框起来。这些精简策略的均衡是 \(\left( {X*,{bd}}\right)\) 、 \(\left( {X*,{cd}}\right) ,\left( {X*,{ce}}\right)\) 和 \(\left( {Y*,{bd}}\right)\) 。
\end{tcolorbox}

\section*{4.6 Subgame-Perfect Equilibrium (SPE)}

\begin{tcolorbox}
\section*{4.6 子博弈完美均衡(SPE)}
\end{tcolorbox}

In this section, we consider the relationship between backward induction and equilibrium. The concept of equilibrium is based on the strategic form because it applies to a strategy profile. In a game tree, the strategies in that profile are assumed to be chosen by the players before play starts, and the concept of best response assumes that the other players play the moves in the strategy profile.

\begin{tcolorbox}
在本节中，我们考虑逆向归纳法和均衡之间的关系。均衡的概念基于策略形式，因为它适用于一个策略组合。在博弈树中，假设该组合中的策略是玩家在游戏开始前选择的，并且最优反应的概念假设其他玩家按照策略组合中的行动进行游戏。
\end{tcolorbox}

With the game tree as the given specification of the game, it is often desirable to keep its sequential interpretation. The strategies chosen by the players should therefore also express some "sequential rationality" as expressed by backward induction. That is, the moves in a strategy profile should be optimal for any part of the game, including subtrees that cannot be reached due to earlier moves, possibly by other players, like the little tree that starts with the decision node of player II in the Threat game in Figure 4.3(a).

\begin{tcolorbox}
以博弈树作为博弈的给定规范，通常希望保留其顺序解释。因此，玩家选择的策略也应该表达出逆向归纳法所表达的某种“顺序理性”。也就是说，策略组合中的行动对于博弈的任何部分都应该是最优的，包括由于早期行动(可能是其他玩家的行动)而无法到达的子树，就像图4.3(a)中威胁博弈中以玩家II的决策节点开始的小树一样。
\end{tcolorbox}

In a game tree, a subgame is any subtree of the game, given by a node of the tree as the root of the subtree and all its descendants. (Note: This definition of a subgame applies only to games with perfect information, which are the game trees considered so far. In extensive games with imperfect information, which we consider in Chapter 10, a subtree is a subgame only if all players know that they are in that subtree; see Section 10.10.) A strategy profile that defines an equilibrium for every subgame is called a subgame-perfect equilibrium or SPE.

\begin{tcolorbox}
在博弈树中，子博弈是博弈的任何子树，由树的一个节点作为子树的根节点及其所有后代节点给出。(注意:这个子博弈的定义仅适用于具有完美信息的博弈，这是我们目前考虑的博弈树。在第10章中我们将考虑的不完美信息扩展式博弈中，一个子树只有在所有玩家都知道他们处于该子树时才是一个子博弈；见第10.10节。)一个为每个子博弈定义均衡的策略组合被称为子博弈完美均衡或SPE。
\end{tcolorbox}

\section*{Theorem 4.6. Backward induction defines an SPE.}

\begin{tcolorbox}
\section*{定理4.6. 逆向归纳法定义了一个子博弈完美均衡。}
\end{tcolorbox}

Proof. Recall the process of backward induction: Starting with the nodes closest to the leaves, consider a decision node \(u\) , say, with the assumption that all moves at descendants of \(u\) have already been selected. Among the moves at \(u\) , select a move that maximizes the expected payoff to the player who moves at \(u\) . (Expected payoffs must be regarded if there are chance moves after \(u\) .) In that manner, a move is selected for every decision node, which determines an entire strategy profile.

\begin{tcolorbox}
证明。回顾逆向归纳法的过程:从最接近叶节点的节点开始，考虑一个决策节点 \(u\) ，假设 \(u\) 的所有后代节点的行动都已经被选择。在 \(u\) 的行动中，选择一个能使在 \(u\) 行动的玩家的期望效用最大化的行动。(如果在 \(u\) 之后有随机行动，则必须考虑期望效用。)以这种方式，为每个决策节点选择一个行动，从而确定一个完整的策略组合。
\end{tcolorbox}

We prove the theorem inductively (see Figure 4.10): Consider a non-terminal node \(u\) of the game tree, which may be a decision node (as in the backward induction process), or a chance node. Suppose that the moves at \(u\) are \({c}_{1},{c}_{2},\ldots ,{c}_{k}\) , which lead to subtrees \({T}_{1},{T}_{2},\ldots ,{T}_{k}\) of the game tree. If \(u\) is a chance node, then instead of the moves we have probabilities \({p}_{1},{p}_{2},\ldots ,{p}_{k}\) . Assume, as inductive hypothesis, that the moves selected so far define an SPE in each of these trees. (As the "base case" of the induction, this is certainly true for each subtree \({T}_{1},{T}_{2},\ldots ,{T}_{k}\) that is just a leaf of the tree; if all subtrees are leaves, then \(u\) is a "last" decision node considered first in the backward induction process.) The induction step is completed if one shows that, by selecting the move at \(u\) , one obtains an SPE for the subgame with root \(u\) .

\begin{tcolorbox}
我们用归纳法证明该定理(见图4.10):考虑博弈树(game tree)的一个非终端节点(non-terminal node) \(u\) ，它可能是一个决策节点(如在逆向归纳法(backward induction)过程中)，也可能是一个机会节点。假设在节点 \(u\) 的行动为 \({c}_{1},{c}_{2},\ldots ,{c}_{k}\) ，这些行动会导向博弈树的子树(subtree) \({T}_{1},{T}_{2},\ldots ,{T}_{k}\) 。如果 \(u\) 是一个机会节点，那么我们有的不是行动，而是概率 \({p}_{1},{p}_{2},\ldots ,{p}_{k}\) 。作为归纳假设，假定到目前为止所选择的行动在每一棵子树中都定义了一个子博弈完美均衡(SPE)。(作为归纳的“基础情形”，对于每一棵只是树的叶节点(leaf)的子树 \({T}_{1},{T}_{2},\ldots ,{T}_{k}\) ，这肯定是成立的；如果所有子树都是叶节点，那么 \(u\) 就是逆向归纳法过程中首先考虑的“最后”决策节点。)如果能证明通过在节点 \(u\) 选择行动，可以得到以 \(u\) 为根节点的子博弈(subgame)的一个子博弈完美均衡，那么归纳步骤就完成了。
\end{tcolorbox}

\begin{center}
\includegraphics[max width=0.8\textwidth]{images/bo_d2e5cmbef24c73b03tn0_103_399_219_1136_459_0.jpg}
\end{center}
\hspace*{3em} 

Figure 4.10 Induction step for proving Theorem 4.6 when the next node \(u\) belongs to chance (left) or to player \(i\) (right).

\begin{tcolorbox}
图4.10 当下一个节点 \(u\) 属于机会(左图)或属于参与者 \(i\) (右图)时，证明定理4.6的归纳步骤。
\end{tcolorbox}

First, suppose that \(u\) is a chance node, so that the next node is chosen according to the fixed probabilities specified in the game tree. Then backward induction does not select a move for \(u\) , and the inductive step holds trivially: For every player, the payoff in the subgame that starts at \(u\) is the expectation of the payoffs for the subgames \({T}_{l}\) weighted with the probabilities \({p}_{l}\) for \(1 \leq  l \leq  k\) , and if a player could improve on that payoff, she would have to do so by changing her moves within at least one subtree \({T}_{l}\) , which, by inductive hypothesis, she cannot.

\begin{tcolorbox}
首先，假设 \(u\) 是一个机会节点，这样下一个节点是根据博弈树中指定的固定概率来选择的。那么逆向归纳法不会为 \(u\) 选择行动，并且归纳步骤显然成立:对于每个参与者来说，从 \(u\) 开始的子博弈中的收益是子博弈 \({T}_{l}\) 的收益以概率 \({p}_{l}\) 对 \(1 \leq  l \leq  k\) 进行加权后的期望值，而且如果一个参与者能够提高该收益，她就必须通过改变至少一棵子树 \({T}_{l}\) 内的行动来实现，根据归纳假设，她无法做到这一点。
\end{tcolorbox}

Second, suppose that \(u\) is a decision node, let \(i\) be the player to move at \(u\) , and consider a player \(j\) other than player \(i\) . Again, for player \(j\) , the moves in the subgame that starts at \(u\) are completely specified. Irrespective of what move player \(i\) selects at \(u\) , player \(j\) cannot improve her payoff because that would mean she could improve her payoff already in some subtree \({T}_{l}\) .

\begin{tcolorbox}
其次，假设 \(u\) 是一个决策节点，设 \(i\) 是在节点 \(u\) 行动的参与者，并考虑除参与者 \(i\) 之外的另一个参与者 \(j\) 。同样，对于参与者 \(j\) 来说，从 \(u\) 开始的子博弈中的行动是完全确定的。无论参与者 \(i\) 在节点 \(u\) 选择什么行动，参与者 \(j\) 都无法提高她的收益，因为这意味着她已经可以在某棵子树 \({T}_{l}\) 中提高她的收益了。
\end{tcolorbox}

Third, consider player \(i\) who moves at \(u\) . Backward induction selects a move for \(u\) that is best for player \(i\) , given the remaining moves. Suppose player \(i\) could improve his payoff by choosing some move \({c}_{l}\) and additionally change his moves in the subtree \({T}_{l}\) . But the resulting improved payoff would only be the improved payoff in \({T}_{l}\) , that is, the player could already get a better payoff in \({T}_{l}\) itself. This contradicts the inductive assumption that the moves selected so far define an SPE for \({T}_{l}\) . This completes the induction.

\begin{tcolorbox}
第三，考虑在节点 \(u\) 行动的参与者 \(i\) 。逆向归纳法会为节点 \(u\) 选择一个在其余行动给定的情况下对参与者 \(i\) 来说最优的行动。假设参与者 \(i\) 可以通过选择某个行动 \({c}_{l}\) 并额外改变他在子树 \({T}_{l}\) 中的行动来提高他的收益。但由此得到的提高后的收益仅仅是子树 \({T}_{l}\) 中提高后的收益，也就是说，该参与者本来就可以在子树 \({T}_{l}\) 中获得更好的收益。这与到目前为止所选择的行动为子树 \({T}_{l}\) 定义了一个子博弈完美均衡的归纳假设相矛盾。这样就完成了归纳。
\end{tcolorbox}

Theorem 4.6 has two important consequences. In backward induction, each move can be chosen deterministically, so that backward induction determines a profile of pure strategies. Theorem 4.6 therefore implies that game trees have equilibria, and it is not necessary to consider randomized strategies. Secondly, subgame-perfect equilibria exist. For game trees, we can use "SPE" synonymously with "strategy profile obtained by backward induction".

\begin{tcolorbox}
定理4.6有两个重要的推论。在逆向归纳法中，每个行动都可以确定性地选择，因此逆向归纳法确定了一个纯策略组合。因此，定理4.6意味着博弈树存在均衡(equilibrium)，而且没有必要考虑随机策略。其次，子博弈完美均衡是存在的。对于博弈树，我们可以将“子博弈完美均衡(SPE)”与“通过逆向归纳法得到的策略组合”同义使用。
\end{tcolorbox}

Corollary 4.7. Every game tree with perfect information has a pure-strategy equilibrium.

\begin{tcolorbox}
推论4.7:每个具有完美信息(perfect information)的博弈树都有一个纯策略均衡。
\end{tcolorbox}

Corollary 4.8. Every game tree with perfect information has an SPE.

\begin{tcolorbox}
推论4.8:每个具有完美信息的博弈树都有一个子博弈完美均衡。
\end{tcolorbox}

An SPE can only be described as a profile of unreduced (that is, fully specified) strategies. Recall that in the game in Figure 4.1, the SPE are \(\left( {{XQ},{bd}}\right) ,\left( {{XQ},{cd}}\right)\) , and(YQ, bd). The reduced strategic form of the game in Figure 4.7 shows the equilibria \(\left( {X*,{bd}}\right) ,\left( {X*,{cd}}\right)\) , and \(\left( {Y*,{bd}}\right)\) . However, we cannot call any of these an SPE because they leave the second move of player I unspecified, as indicated by the \(*\) symbol. In this case, replacing \(*\) by \(Q\) results in all cases in an SPE. The full strategies are necessary to determine if they define an equilibrium in every subgame. As seen in Figure 4.7, the game has, in addition, the equilibrium \(\left( {X*,{ce}}\right)\) . This is not subgame-perfect because it prescribes the move \(e\) , which is not part of an equilibrium in the subgame that starts with the decision node where player II chooses between \(d\) and \(e\) : Replacing \(*\) by \(P\) would not be the best response by player I to \(e\) in that subgame, and when replacing \(*\) by \(Q\) , player II’s best response would not be \(e\) . This property can only be seen in the game tree, not even in the unreduced strategic form in Figure 4.6, because it concerns an unreached part of the game tree due to the earlier move \(X\) of player \(\mathrm{I}\) .

\begin{tcolorbox}
子博弈完美均衡(SPE)只能被描述为未简化(即完全明确)策略的组合。回顾图4.1中的博弈，子博弈完美均衡是 \(\left( {{XQ},{bd}}\right) ,\left( {{XQ},{cd}}\right)\) 和(YQ, bd)。图4.7中博弈的精简策略型展示了均衡 \(\left( {X*,{bd}}\right) ,\left( {X*,{cd}}\right)\) 和 \(\left( {Y*,{bd}}\right)\) 。然而，我们不能将这些中的任何一个称为子博弈完美均衡，因为它们没有明确玩家I的第二步行动，如 \(*\) 符号所示。在这种情况下，在所有情况下将 \(*\) 替换为 \(Q\) 会得到一个子博弈完美均衡。要确定它们是否在每个子博弈中都定义了一个均衡，完整的策略是必要的。如图4.7所示，该博弈还有均衡 \(\left( {X*,{ce}}\right)\) 。这不是子博弈完美的，因为它规定了行动 \(e\) ，而这不是以玩家II在 \(d\) 和 \(e\) 之间进行选择的决策节点开始的子博弈中的均衡的一部分:在该子博弈中，将 \(*\) 替换为 \(P\) 不是玩家I对 \(e\) 的最优反应，并且当将 \(*\) 替换为 \(Q\) 时，玩家II的最优反应不会是 \(e\) 。这个性质只能在博弈树中看到，即使在图4.6的未简化策略形式中也看不到，因为它涉及到由于玩家 \(\mathrm{I}\) 更早的行动 \(X\) 而未到达的博弈树部分。
\end{tcolorbox}

\(\Rightarrow\) Exercise 4.3 is very instructive for understanding subgame-perfect equilibria. Exercise 4.5 studies a three-player game.

\begin{tcolorbox}
\(\Rightarrow\) 练习4.3对于理解子博弈完美均衡非常有启发性。练习4.5研究了一个三人博弈。
\end{tcolorbox}

\section*{4.7 Commitment Games}

\begin{tcolorbox}
\section*{4.7 承诺博弈}
\end{tcolorbox}

The main difference between game trees and games in strategic form is that in a game tree the players move sequentially and are aware of any previous moves of the other players. In contrast, in a strategic-form game the players move simultaneously. The difference between these two descriptions becomes striking when changing a two-player game in strategic form to a commitment game.

\begin{tcolorbox}
博弈树和策略形式的博弈之间的主要区别在于，在博弈树中，玩家按顺序行动，并且知道其他玩家之前的任何行动。相比之下，在策略形式的博弈中，玩家同时行动。当将一个两人的策略形式博弈转变为承诺博弈时，这两种描述之间的差异变得非常明显。
\end{tcolorbox}

The commitment game is a new game obtained from a two-player game in strategic form. One of the players (we always choose player I) moves first by choosing one of his strategies. Then, unlike in the strategic form, player II is informed about the move of player I and moves second. Player II chooses one of her original strategies, which can depend on the move of player I. The resulting payoffs to both players are as specified in the original game.

\begin{tcolorbox}
承诺博弈是从一个两人的策略形式博弈中得到的新博弈。其中一个玩家(我们总是选择玩家I)先行动，选择他的一个策略。然后，与策略形式不同的是，玩家II会得知玩家I的行动，然后再行动。玩家II选择她原来的一个策略，这个策略可以取决于玩家I的行动。两个玩家最终获得的收益与原博弈中规定的相同。
\end{tcolorbox}

An example is shown in Figure 4.11. The game in strategic form in (a) is the Quality game in Figure 3.2(b). The commitment game, shown in (b) and (c), is a game tree with perfect information, where player I moves first and player II moves second. The original strategies \(T\) and \(B\) of player \(\mathrm{I}\) become his moves \(T\) and \(B\) . In Figure 4.11(b), we have shown the subsequent choices of player II as two \(1 \times  2\) games which are the two rows of the strategic form in (b), which depend on the move \(T\) or \(B\) of player \(\mathrm{I}\) , with the same payoffs as in (a). Player \(\mathrm{{II}}\) can then separately choose between \(l\) and \(r\) . We assume that player II chooses her best move, shown by the best-response payoffs, which is \(l\) in response to \(T\) and \(r\) in response to \(B\) . We assume that player I anticipates that move of player II and therefore chooses \(T\) , as shown by the backward induction arrow.

\begin{tcolorbox}
图4.11展示了一个例子。(a)中的策略形式的博弈是图3.2(b)中的产品质量博弈。(b)和(c)中展示的承诺博弈是一个具有完美信息的博弈树，其中玩家I先行动，玩家II后行动。玩家 \(\mathrm{I}\) 原来的策略 \(T\) 和 \(B\) 变成了他的行动 \(T\) 和 \(B\) 。在图4.11(b)中，我们将玩家II随后的选择展示为两个 \(1 \times  2\) 博弈，它们是(b)中策略形式的两行，这取决于玩家 \(\mathrm{I}\) 的行动 \(T\) 或 \(B\) ，收益与(a)中相同。玩家 \(\mathrm{{II}}\) 然后可以分别在 \(l\) 和 \(r\) 之间进行选择。我们假设玩家II选择她的最优行动，由最优反应收益所示，即对 \(T\) 的反应是 \(l\) ，对 \(B\) 的反应是 \(r\) 。我们假设玩家I预期到了玩家II的这一行动，因此选择 \(T\) ，如逆向归纳箭头所示。
\end{tcolorbox}

\begin{center}
\includegraphics[max width=0.9\textwidth]{images/bo_d2e5cmbef24c73b03tn0_105_320_195_1260_926_0.jpg}
\end{center}
\hspace*{3em} 

Figure 4.11 The simultaneous Quality game (a) from Figure 3.2(b), converted to its sequential commitment game (b) where player I moves first and player II can react to player I's move. The conventional representation of (b) is the extensive game (c). The arrows in (b) and (c) show the backward induction outcome. The strategic form of (c) with best-response payoffs is shown in (d).

\begin{tcolorbox}
图4.11 图3.2(b)中的同时进行的产品质量博弈(a)，转换为顺序的承诺博弈(b)，其中玩家I先行动，玩家II可以对玩家I的行动做出反应。(b)的传统表示形式是扩展式博弈(c)。(b)和(c)中的箭头显示了逆向归纳的结果。(c)带有最优反应收益的策略形式在(d)中展示。
\end{tcolorbox}

The graphical description in Figure 4.11(b) just illustrates how the commitment game is obtained from the original game in strategic form. It represents a game tree, which is shown in (c), where each move of player I leads to a separate decision node for player II, whose moves are the original strategies of player II. In order to emphasize that these moves \(l\) and \(r\) are separate and depend on the decision node, we have in (c) given them subscripts, and call them \({l}_{T}\) and \({r}_{T}\) after move \(T\) , and \({l}_{B}\) and \({r}_{B}\) after move \(B\) . Figure 4.11(d) shows the resulting strategic form of the commitment game. It shows that the original strategic-form game and the commitment game are very different.

\begin{tcolorbox}
图4.11(b)中的图形描述仅说明了如何从原始的策略形式博弈中得到承诺博弈。它表示一个博弈树，如图(c)所示，其中参与者I的每一步行动都会为参与者II产生一个单独的决策节点，参与者II的行动就是其原始策略。为了强调这些行动 \(l\) 和 \(r\) 是相互独立的，并且依赖于决策节点，我们在图(c)中给它们加上了下标，并将在行动 \(T\) 之后的行动称为 \({l}_{T}\) 和 \({r}_{T}\) ，在行动 \(B\) 之后的行动称为 \({l}_{B}\) 和 \({r}_{B}\) 。图4.11(d)展示了承诺博弈的策略形式。它表明原始的策略形式博弈和承诺博弈有很大的不同。
\end{tcolorbox}

In the commitment game, we then look for subgame-perfect equilibria, that is, apply backward induction. Figure 4.11(b) and (c) show that the resulting unique SPE is \(\left( {T,{l}_{T}{r}_{B}}\right)\) , which means that (going backwards) player II when offered high quality \(T\) chooses \({l}_{T}\) and when offered low quality \(B\) chooses \({r}_{B}\) . In consequence, player I will then choose \(T\) because that gives him the higher payoff 2 rather than just 1 with \(B\) .

\begin{tcolorbox}
在承诺博弈中，我们寻找子博弈完美均衡，即应用逆向归纳法。图4.11(b)和(c)表明，得到的唯一子博弈完美均衡是 \(\left( {T,{l}_{T}{r}_{B}}\right)\) ，这意味着(逆向推导)当参与者II面对高质量 \(T\) 时选择 \({l}_{T}\) ，面对低质量 \(B\) 时选择 \({r}_{B}\) 。因此，参与者I会选择 \(T\) ，因为这会给他带来更高的收益2，而选择 \(B\) 只能得到1。
\end{tcolorbox}

Compared to the equilibrium in the original Quality game, the commitment game has a higher payoff (in the SPE) to both players. This holds because the game has changed. It is important that player I has the power to commit to \(T\) and cannot change later back to \(B\) after seeing that player II chose \(l\) . Without that commitment power of player I, the players would not act as described and the game would not be accurately reflected by the game tree.

\begin{tcolorbox}
与原始产品质量博弈中的均衡相比，承诺博弈(在子博弈完美均衡中)给两个参与者都带来了更高的收益。这是因为博弈发生了变化。重要的是，参与者I有能力承诺选择 \(T\) ，并且在看到参与者II选择 \(l\) 后不能再改回 \(B\) 。如果参与者I没有这种承诺能力，参与者的行为就不会如上述描述，博弈树也就不能准确反映该博弈。
\end{tcolorbox}

The commitment game in Figure 4.11 also demonstrates the difference between an equilibrium, which is a strategy profile, and the equilibrium path, which is described by the moves that are actually played in the game when players use their equilibrium strategies. (If the game has no chance moves, the equilibrium path is just a path in the game tree from the root to a leaf; if the game has chance moves, the path may branch.) In contrast, a strategy profile defines a move in every part of the game tree. Here, the SPE \(\left( {T,{l}_{T}{r}_{B}}\right)\) specifies moves \({l}_{T}\) and \({r}_{B}\) for the two decision nodes of player II. On the equilibrium path, the moves are \(T\) of player I followed by move \({l}_{T}\) of player II. However, it is not sufficient to simply call this equilibrium \(\left( {T,{l}_{T}}\right)\) or(T, l), because player II’s move \({r}_{B}\) must be specified to know that \(T\) is player I’s best response.

\begin{tcolorbox}
图4.11中的承诺博弈还展示了均衡(一种策略组合)和均衡路径(当参与者使用其均衡策略时，博弈中实际采取的行动所描述的路径)之间的区别。(如果博弈没有随机行动，均衡路径就是博弈树中从根节点到叶节点的一条路径；如果博弈有随机行动，路径可能会分叉。)相比之下，一个策略组合定义了博弈树中每个部分的行动。这里，子博弈完美均衡 \(\left( {T,{l}_{T}{r}_{B}}\right)\) 为参与者II的两个决策节点指定了行动 \({l}_{T}\) 和 \({r}_{B}\) 。在均衡路径上，参与者I的行动是 \(T\) ，随后是参与者II的行动 \({l}_{T}\) 。然而，仅仅称这个均衡为 \(\left( {T,{l}_{T}}\right)\) 或(T, l)是不够的，因为必须指定参与者II的行动 \({r}_{B}\) ，才能知道 \(T\) 是参与者I的最优反应。
\end{tcolorbox}

We only consider SPE when analysing commitment games. For example, the equilibrium(B, r)of the original Quality game in Figure 4.11(a) can also be found in (d) where player II always chooses the move \(r\) that corresponds to her equilibrium strategy in (a), which defines the strategy \({r}_{T}{r}_{B}\) in the commitment game. This determines the equilibrium \(\left( {B,{r}_{T}{r}_{B}}\right)\) in (d), because \(B\) is a best response to \({r}_{T}{r}_{B}\) , and \({r}_{T}{r}_{B}\) is a best response to \(B\) . However, this equilibrium of the commitment game is not subgame-perfect, because it prescribes the suboptimal move \({r}_{T}\) off the equilibrium path; this does not affect the equilibrium property because \(T\) is not chosen by player I. (This holds as a general argument: consider any equilibrium (x, y)of the original game, and in the commitment game let player I commit to \(x\) and let player II always respond with \(y\) ; then this is an equilibrium of the commitment game, which, however, is in general not subgame-perfect.)

\begin{tcolorbox}
我们在分析承诺博弈时只考虑子博弈完美均衡。例如，图4.11(a)中原始产品质量博弈的均衡(B, r)也可以在图(d)中找到，其中参与者II总是选择与她在图(a)中的均衡策略相对应的行动 \(r\) ，这在承诺博弈中定义了策略 \({r}_{T}{r}_{B}\) 。这就确定了图(d)中的均衡 \(\left( {B,{r}_{T}{r}_{B}}\right)\) ，因为 \(B\) 是对 \({r}_{T}{r}_{B}\) 的最优反应，而 \({r}_{T}{r}_{B}\) 是对 \(B\) 的最优反应。然而，这个承诺博弈的均衡不是子博弈完美的，因为它规定了在均衡路径之外采取次优行动 \({r}_{T}\) ；这并不影响均衡性质，因为参与者I不会选择 \(T\) 。(这是一个通用的论点:考虑原始博弈的任何均衡(x, y)，在承诺博弈中让参与者I承诺选择 \(x\) ，让参与者II总是以 \(y\) 回应；那么这就是承诺博弈的一个均衡，但一般来说它不是子博弈完美的。)
\end{tcolorbox}

In order to compare a strategic-form game with its commitment game in an interesting way, we consider only subgame-perfect equilibria of the commitment game.

\begin{tcolorbox}
为了以一种有趣的方式比较策略形式博弈和其承诺博弈，我们只考虑承诺博弈的子博弈完美均衡。
\end{tcolorbox}

In practical applications, it is often of interest to "change the rules of the game" and see what happens (with a proper game-theoretic analysis). Changing the Quality game to its commitment game shows how both players benefit.

\begin{tcolorbox}
在实际应用中，人们通常会对“改变博弈规则”并观察结果(通过适当的博弈论分析)感兴趣。将产品质量博弈转变为承诺博弈，能展示出双方参与者如何从中受益。
\end{tcolorbox}

Commitment games are also called leadership games where the first player to move is called the leader and the second player the follower. The leader's ability to commit irrevocably to his strategy is important, and is also called commitment power. The ability to commit typically confers a first-mover advantage, which states that the leader is not worse off than in the original game where the players act simultaneously (as shown in Exercise 4.7, this holds definitely when best responses are always unique).

\begin{tcolorbox}
承诺博弈也被称为领导博弈，在这种博弈中，先行动的参与者被称为领导，后行动的参与者被称为跟随者。领导不可撤销地坚持其策略的能力非常重要，这也被称为承诺能力。承诺能力通常会带来先行优势，即领导的处境不会比参与者同时行动的原始博弈更差(如练习4.7所示，当最优反应总是唯一时，这一点肯定成立)。
\end{tcolorbox}

This statement must be interpreted carefully. For example, if either player has the power to commit, then it is not necessarily best to go first. For example, consider changing the game in Figure 4.11(a) so that player II can commit to her strategy, and player I moves second. Then player I will always respond by choosing \(B\) because this is his dominant strategy in (a). Backward induction would then amount to player II choosing \(r\) , and player I choosing \({B}_{l}{B}_{r}\) , with the low payoff 1 to both. Then player II is not worse off than in the simultaneous-choice game, as asserted by the first-mover advantage, but does not gain anything either. In contrast, making player I the first mover as in Figure 4.11(b) is beneficial to both.

\begin{tcolorbox}
必须谨慎解读这一说法。例如，如果任何一方参与者都有承诺的能力，那么先行行动并不一定是最优选择。例如，考虑改变图4.11(a)中的博弈，让参与者II能够承诺其策略，而参与者I后行动。那么参与者I总是会选择 \(B\) 作为回应，因为这是他在(a)中的占优策略。然后，通过逆向归纳法，参与者II会选择 \(r\) ，参与者I会选择 \({B}_{l}{B}_{r}\) ，双方的收益都很低，仅为1。此时，正如先行优势所断言的，参与者II的处境并不比同时选择的博弈更差，但也没有获得任何好处。相比之下，如图4.11(b)所示，让参与者I先行动对双方都有利。
\end{tcolorbox}

Commitment games are also known as Stackelberg games, after the economist Heinrich von Stackelberg who applied them to the Cournot duopoly of quantity competition such as the game that we considered in Section 3.6. We study this commitment game first for its finite version in Figure 3.8 and then its infinite version (3.6) for arbitrary nonnegative \(x\) and \(y\) . We will see that in this game, the leader is better off than in the simultaneous game, but the follower is worse off.

\begin{tcolorbox}
承诺博弈也被称为斯塔克尔伯格博弈，这是以经济学家海因里希·冯·斯塔克尔伯格(Heinrich von Stackelberg)的名字命名的，他将这种博弈应用于古诺双寡头产量竞争，比如我们在第3.6节中考虑的博弈。我们首先研究图3.8中该承诺博弈的有限版本，然后研究任意非负 \(x\) 和 \(y\) 情况下的无限版本(3.6)。我们会发现，在这个博弈中，领导的处境比同时行动的博弈更好，但跟随者的处境更差。
\end{tcolorbox}

The commitment game for the game in Figure 3.8 is shown in Figure 4.12, just as Figure 4.11(b) is obtained from Figure 4.11(a). The best responses of player II are given by her (original) best responses to the first move0,3,4,6of player \(\mathrm{I}\) , which are6,4,4,3, respectively, with payoff pairs \(\left( {0,{36}}\right) ,\left( {{15},{20}}\right) ,\left( {{16},{16}}\right)\) , and(18,9)to the two players. Among these, 18 is the best for player I. Hence, by backward induction, player I will commit to his rightmost move 6 . The corresponding payoffs are 18 for player I and 9 for player II, compared to 16 and 16 in the simultaneous game in Figure 3.8. In this symmetric game, the ability of one player to credibly commit himself gives a clear first-mover advantage to the leader which, unlike in the Quality game, is here at the expense of the follower.

\begin{tcolorbox}
图3.8中博弈的承诺博弈如图4.12所示，就像图4.11(b)是从图4.11(a)推导出来的一样。参与者II的最优反应由她对参与者 \(\mathrm{I}\) 的首次行动0、3、4、6的(原始)最优反应给出，分别为6、4、4、3，收益对为 \(\left( {0,{36}}\right) ,\left( {{15},{20}}\right) ,\left( {{16},{16}}\right)\) ，两个参与者的收益为(18, 9)。其中，18对参与者I来说是最好的。因此，通过逆向归纳法，参与者I会承诺选择最右边的行动6。相应的收益是参与者I为18，参与者II为9，而在图3.8的同时行动博弈中，双方的收益均为16。在这个对称博弈中，一方参与者可信地承诺自己策略的能力为领导带来了明显的先行优势，与产品质量博弈不同的是，这里是以牺牲跟随者的利益为代价的。
\end{tcolorbox}

We now consider the commitment game for the game with payoffs \(a\left( {x,y}\right)\) and \(b\left( {x,y}\right)\) in (3.6) to players I and II when their strategies \(x\) and \(y\) are arbitrary reals in \(\left\lbrack  {0,{12}}\right\rbrack\) (they will not choose \(x > {12}\) or \(y > {12}\) which would give them negative payoffs). We have seen that the best response of player II to \(x\) is \(y\left( x\right)  = 6 - x/2\) . In the commitment game, it is indeed useful to see player II's choice as such a function \(y\left( x\right)\) of the commitment \(x\) of player \(\mathrm{I}\) .

\begin{tcolorbox}
现在我们考虑在(3.6)中，当参与者I和II的策略 \(x\) 和 \(y\) 是 \(\left\lbrack  {0,{12}}\right\rbrack\) 中的任意实数时(他们不会选择 \(x > {12}\) 或 \(y > {12}\) ，因为这会给他们带来负收益)，该博弈的承诺博弈。我们已经知道，参与者II对 \(x\) 的最优反应是 \(y\left( x\right)  = 6 - x/2\) 。在承诺博弈中，将参与者II的选择看作是参与者 \(\mathrm{I}\) 的承诺 \(x\) 的函数 \(y\left( x\right)\) 确实很有用。
\end{tcolorbox}

\begin{center}
\includegraphics[max width=0.8\textwidth]{images/bo_d2e5cmbef24c73b03tn0_108_254_256_1021_960_0.jpg}
\end{center}
\hspace*{3em} 

Figure 4.12 Commitment game for the duopoly game in Figure 3.8.

\begin{tcolorbox}
图4.12 图3.8中双寡头博弈的承诺博弈。
\end{tcolorbox}

As always in a commitment game, we look for an SPE and assume that the follower chooses her best response, even off the equilibrium path. The backward induction outcome is therefore the solution \(x\) to the problem of maximizing \(a\left( {x,y\left( x\right) }\right)\) , that is, of maximizing \(x \cdot  \left( {{12} - y\left( x\right)  - x}\right)  = x \cdot  \left( {6 - x/2}\right)  = {6x} - {x}^{2}/2\) . The derivative of this function with respect to \(x\) is \(6 - x\) , which is zero for \(x = 6\) . That is, player I commits to 6 and player II responds with \(y\left( x\right)  = 6 - x/2\) for any commitment \(x\) (which defines the full strategy of player II in this SPE), in particular with \(y\left( 6\right)  = 3\) on the equilibrium path. The resulting payoffs are 18 for player I and 9 for player II. The strategies 6 and 3 are part of the \(4 \times  4\) game in Figure 3.8, along with the equilibrium strategy 4 of both players in that simultaneous game.

\begin{tcolorbox}
与承诺博弈中一贯的做法一样，我们寻找子博弈完美均衡(SPE)，并假设跟随者即使在非均衡路径上也会选择其最优反应。因此，逆向归纳法的结果是问题最大化 \(x\) 的解，即最大化 \(a\left( {x,y\left( x\right) }\right)\) 的解。该函数关于 \(x\) 的导数为 \(6 - x\) ，当 \(x = 6\) 时其值为零。也就是说，参与者 I 承诺选择 6，而参与者 II 对于任何承诺 \(x\) 都会以 \(y\left( x\right)  = 6 - x/2\) 作为回应(这定义了参与者 II 在这个子博弈完美均衡中的完整策略)，特别是在均衡路径上以 \(y\left( 6\right)  = 3\) 回应。最终结果是参与者 I 的收益为 18，参与者 II 的收益为 9。策略 6 和 3 是图 3.8 中 \(4 \times  4\) 博弈的一部分，同时还有该同时博弈中两个参与者的均衡策略 4。
\end{tcolorbox}

\(\Rightarrow\) Exercise 4.7 studies a game with infinite strategy sets like (3.6), and makes interesting observations, in particular in (d), on the payoffs to the two players in the commitment game compared to the original simultaneous game.

\begin{tcolorbox}
\(\Rightarrow\) 练习 4.7 研究了一个像(3.6)那样具有无限策略集的博弈，并做出了有趣的观察，特别是在(d)中，比较了承诺博弈与原始同时博弈中两个参与者的收益。
\end{tcolorbox}

\section*{4.8 Further Reading}

\begin{tcolorbox}
\section*{4.8 进一步阅读}
\end{tcolorbox}

Backward induction is also known as Zermelo's algorithm, attributed to an article of Zermelo (1913) on Chess. However, Zermelo was not concerned with backward induction but instead proved that if the white player can force a win, then he can do so in a bounded number of moves (Schwalbe and Walker, 2001). For that reason, "backward induction" is the better term.

\begin{tcolorbox}
逆向归纳法也被称为策梅洛算法，这一名称源于策梅洛(1913 年)关于国际象棋的一篇文章。然而，策梅洛关注的并非逆向归纳法，而是证明了如果白方能够迫使获胜，那么他可以在有限步数内做到这一点(施瓦尔贝和沃克，2001 年)。因此，“逆向归纳法”是更合适的术语。
\end{tcolorbox}

Reduced strategies are occasionally defined in terms of the strategic form by identifying strategies that have identical payoffs for all players. We define this reduction independently of payoffs only via the structure of the game tree, which is clearer. We also give the * symbol (which is sometimes omitted) for each unspecified move and each decision node (there may be several such decision nodes) in order to ensure that even a reduced strategy has \(k\) "slots" of moves when a player has \(k\) decision nodes.

\begin{tcolorbox}
精简策略有时是根据策略形式来定义的，即通过识别对所有参与者具有相同收益的策略。我们仅通过博弈树的结构独立于收益来定义这种精简，这样更清晰。我们还为每个未指定的行动和每个决策节点(可能有多个这样的决策节点)给出 * 符号(有时会省略)，以确保即使是精简策略，当参与者有 \(k\) 个决策节点时也有 \(k\) 个“行动槽”。
\end{tcolorbox}

The economist von Stackelberg (1934) analyzed various ways how firms can act in a duopoly situation, one of which is the leadership game where one firm commits to a certain output and the other firm reacts with a best response. We have described it with concepts of modern game theory (see also Gibbons, 1992), which were not around in 1934.

\begin{tcolorbox}
经济学家冯·斯塔克尔伯格(1934 年)分析了企业在双头垄断情况下的各种行为方式，其中之一就是领导博弈，即一家企业承诺一定的产量，而另一家企业以最优反应做出回应。我们用现代博弈论的概念描述了这一情况(另见吉本斯，1992 年)，这些概念在 1934 年还不存在。
\end{tcolorbox}

The commitment game of the game (3.6) treated at the end of Section 4.7 is a symmetric game where players choose their strategies from an interval, which has a unique symmetric equilibrium, and where the players' best response functions are unique and monotonic in the other player's choice. Because the game is symmetric, it suffices to analyze the game when player \(\mathrm{I}\) is the leader in the corresponding commmitment game, which has an SPE payoff to the leader and to the follower. The leader's payoff is at least as high as in the simultaneous game (see Exercise 4.6). Curiously, in such games the payoff to the follower is either worse than in the simultanenous game, or even higher than the payoff to the leader (as in Exercise 4.7). That is, the seemingly natural case that both players profit from sequential play compared to simultaneous play, but the leader more so than the follower, does not occur. The proof is short and uses only some monotonicity conditions (von Stengel, 2010, theorem 1).

\begin{tcolorbox}
在 4.7 节末尾讨论的博弈(3.6)的承诺博弈是一个对称博弈，参与者从一个区间中选择他们的策略，该博弈有唯一的对称均衡，并且参与者的最优反应函数是唯一的，并且随着另一个参与者的选择单调变化。由于该博弈是对称的，在相应的承诺博弈中，当参与者 \(\mathrm{I}\) 是领导时，分析该博弈就足够了，该博弈为领导和跟随者都带来了子博弈完美均衡收益。领导的收益至少与同时博弈中的收益一样高(见练习 4.6)。奇怪的是，在这类博弈中，跟随者的收益要么比同时博弈中的收益差，要么甚至比领导的收益还高(如练习 4.7)。也就是说，与同时博弈相比，两个参与者都从序贯博弈中获利，但领导比跟随者获利更多这种看似自然的情况并不会出现。证明过程很简短，只使用了一些单调性条件(冯·施滕格尔，2010 年，定理 1)。
\end{tcolorbox}

\section*{4.9 Exercises for Chapter 4}

\begin{tcolorbox}
\section*{4.9 第 4 章练习}
\end{tcolorbox}

Exercise 4.1. Consider the following game tree. As always, the top payoffs at a leaf are for player I and bottom payoffs for player II.

\begin{tcolorbox}
练习 4.1。考虑以下博弈树。和往常一样，叶节点处的上方收益是参与者 I 的，下方收益是参与者 II 的。
\end{tcolorbox}

(a) What is the number of strategies of player I and of player II?

\begin{tcolorbox}
(a) 参与者 I 和参与者 II 的策略数量分别是多少？
\end{tcolorbox}

\begin{center}
\includegraphics[max width=0.5\textwidth]{images/bo_d2e5cmbef24c73b03tn0_110_784_447_665_629_0.jpg}
\end{center}
\hspace*{3em} 

(b) How many reduced strategies do they have?

\begin{tcolorbox}
(b) 他们各有多少个精简策略？
\end{tcolorbox}

(c) Give the reduced strategic form of the game.

\begin{tcolorbox}
(c) 给出该博弈的精简策略型。
\end{tcolorbox}

(d) What are the equilibria of the game in reduced strategies?

\begin{tcolorbox}
(d) 该博弈在精简策略下的均衡是什么？
\end{tcolorbox}

(e) What are the subgame-perfect equilibria of the game?

\begin{tcolorbox}
(e) 该博弈的子博弈完美均衡是什么？
\end{tcolorbox}

Exercise 4.2. Consider the following game tree.

\begin{tcolorbox}
练习 4.2。考虑以下博弈树。
\end{tcolorbox}

(a) What is the number of strategies of player I and of player II? How many reduced strategies does each of the players have?

\begin{tcolorbox}
(a) 参与者 I 和参与者 II 的策略数量分别是多少？每个参与者各有多少个精简策略？
\end{tcolorbox}

\begin{center}
\includegraphics[max width=0.4\textwidth]{images/bo_d2e5cmbef24c73b03tn0_110_795_1319_616_622_0.jpg}
\end{center}
\hspace*{3em} 

(b) Give the reduced strategic form of the game.

\begin{tcolorbox}
(b) 给出该博弈的精简策略型。
\end{tcolorbox}

(c) What are the equilibria of the game in reduced strategies? What are the subgame-perfect equilibria of the game?

\begin{tcolorbox}
(c) 该博弈在精简策略下的均衡是什么？该博弈的子博弈完美均衡是什么？
\end{tcolorbox}

(d) Identify every pair of reduced strategies where one strategy weakly or strictly dominates the other, and indicate if the dominance is weak or strict.

\begin{tcolorbox}
(d) 找出每一对精简策略，其中一个策略弱占优或严格占优于另一个策略，并指出这种占优是弱占优还是严格占优。
\end{tcolorbox}

Exercise 4.3. Consider the following game trees.

\begin{tcolorbox}
练习4.3。考虑以下博弈树。
\end{tcolorbox}

\begin{center}
\includegraphics[max width=0.7\textwidth]{images/bo_d2e5cmbef24c73b03tn0_111_482_281_978_459_0.jpg}
\end{center}
\hspace*{3em} 

(a) Find all equilibria for the game tree on the left. Which of these are subgame-perfect?

\begin{tcolorbox}
(a) 找出左边博弈树的所有均衡。其中哪些是子博弈完美均衡？
\end{tcolorbox}

(b) In the game tree on the right, the payoffs \(a,b,c,d\) are positive real numbers. For each of the following statements (i), (ii), (iii), decide if it is true or false, justifying your answer with an argument or counterexample; you may refer to any standard results. For any \(a,b,c,d > 0\) ,

\begin{tcolorbox}
(b) 在右边的博弈树中，收益 \(a,b,c,d\) 是正实数。对于以下陈述 (i)、(ii)、(iii) 中的每一个，判断其是真还是假，并用论据或反例证明你的答案；你可以引用任何标准结果。对于任何 \(a,b,c,d > 0\) ，
\end{tcolorbox}

(i) the game always has a subgame-perfect equilibrium (SPE);

\begin{tcolorbox}
(i) 该博弈始终存在子博弈完美均衡(SPE)；
\end{tcolorbox}

(ii) the payoff to player II in any SPE is always at least as high as her payoff in any equilibrium;

\begin{tcolorbox}
(ii) 在任何子博弈完美均衡中，玩家 II 的收益总是至少不低于她在任何均衡中的收益；
\end{tcolorbox}

(iii) the payoff to player I in any SPE is always at least as high as his payoff in any equilibrium.

\begin{tcolorbox}
(iii) 在任何子博弈完美均衡中，玩家 I 的收益总是至少不低于他在任何均衡中的收益。
\end{tcolorbox}

Exercise 4.4. Consider the following two game trees (a) and (b). Payoffs have been omitted because they are not relevant for the question. In each case, how many strategies does each player have? How many reduced strategies?

\begin{tcolorbox}
练习 4.4。考虑以下两个博弈树 (a) 和 (b)。由于收益与问题无关，因此省略。在每种情况下，每个玩家有多少种策略？有多少种精简策略？
\end{tcolorbox}

(a)

\begin{center}
\includegraphics[max width=0.3\textwidth]{images/bo_d2e5cmbef24c73b03tn0_111_1069_1562_505_495_0.jpg}
\end{center}
\hspace*{3em} 

\begin{center}
\includegraphics[max width=0.5\textwidth]{images/bo_d2e5cmbef24c73b03tn0_111_342_1630_643_342_0.jpg}
\end{center}
\hspace*{3em} 

Exercise 4.5. Consider the following three-player game tree. At a leaf, the topmost payoff is to player I, the middle payoff is to player II, and the bottom payoff is to player III.

\begin{tcolorbox}
练习 4.5。考虑以下三人博弈树。在叶节点处，最上面的收益属于玩家 I，中间的收益属于玩家 II，最下面的收益属于玩家 III。
\end{tcolorbox}

(a) How many strategy profiles does this game have?

\begin{tcolorbox}
(a) 这个博弈有多少个策略组合？
\end{tcolorbox}

\begin{center}
\includegraphics[max width=0.4\textwidth]{images/bo_d2e5cmbef24c73b03tn0_112_906_385_539_499_0.jpg}
\end{center}
\hspace*{3em} 

(b) Identify all pairs of strategies where one strategy strictly, or weakly, dominates the other.

\begin{tcolorbox}
(b) 找出所有一个策略严格或弱占优另一个策略的策略对。
\end{tcolorbox}

(c) Find all equilibria (in pure strategies). Which of these are subgame-perfect?

\begin{tcolorbox}
(c) 找出所有(纯策略)均衡。其中哪些是子博弈完美均衡？
\end{tcolorbox}

Exercise 4.6. Consider a game \(G\) in strategic form. Recall that the commitment game derived from \(G\) is defined by letting player I choose one of his strategies \(x\) , which is then announced to player II, who can then in each case choose one of her strategies in \(G\) as a response to \(x\) . The resulting payoffs are as in the original game \(G\) .

\begin{tcolorbox}
练习 4.6。考虑一个策略形式的博弈 \(G\) 。回顾一下，从 \(G\) 导出的承诺博弈的定义是，让玩家 I 选择他的一个策略 \(x\) ，然后将其告知玩家 II，玩家 II 随后可以在每种情况下选择 \(G\) 中的一个策略作为对 \(x\) 的回应。由此产生的收益与原始博弈 \(G\) 中的收益相同。
\end{tcolorbox}

(a) If \(G\) is an \(m \times  n\) game, how many strategies do player \(\mathrm{I}\) and player \(\mathrm{{II}}\) have, respectively, in the commitment game?

\begin{tcolorbox}
(a) 如果 \(G\) 是一个 \(m \times  n\) 博弈，在承诺博弈中，玩家 \(\mathrm{I}\) 和玩家 \(\mathrm{{II}}\) 分别有多少种策略？
\end{tcolorbox}

For each of the following statements, determine whether they are true or false; justify your answer by a short argument or counterexample.

\begin{tcolorbox}
对于以下每个陈述，判断它们是真还是假；通过简短的论证或反例来证明你的答案。
\end{tcolorbox}

(b) In an SPE of the commitment game, player I never commits to a strategy that is strictly dominated in \(G\) .

\begin{tcolorbox}
(b) 在承诺博弈的子博弈完美均衡中，玩家 I 永远不会选择一个在 \(G\) 中严格被占优的策略。
\end{tcolorbox}

(c) In an SPE of the commitment game, player II never chooses a move that is a strictly dominated strategy in \(G\) .

\begin{tcolorbox}
(c) 在承诺博弈的子博弈完美均衡中，玩家 II 永远不会选择一个在 \(G\) 中严格被占优的策略。
\end{tcolorbox}

Exercise 4.7. Let \(G\) be the following game: Player I chooses a non-negative real number \(x\) , and simultaneously player II chooses a non-negative real number \(y\) . The resulting (symmetric) payoffs are

\begin{tcolorbox}
练习 4.7。设 \(G\) 为以下博弈:玩家 I 选择一个非负实数 \(x\) ，同时玩家 II 选择一个非负实数 \(y\) 。由此产生的(对称)收益为
\end{tcolorbox}

\[
x \cdot  \left( {4 + y - x}\right) \;\text{ for player }\mathrm{I}\text{ , }
\]

\[
y \cdot  \left( {4 + x - y}\right) \;\text{ for player II. }
\]

(a) Given \(x\) , determine player II’s best response \(y\left( x\right)\) (which is a function of \(x\) ), and player I’s best response \(x\left( y\right)\) to \(y\) . Find an equilibrium, and give the payoffs to the two players.

\begin{tcolorbox}
(a) 给定 \(x\) ，确定玩家 II 的最优反应 \(y\left( x\right)\) (它是 \(x\) 的函数)，以及玩家 I 对 \(y\) 的最优反应 \(x\left( y\right)\) 。找出一个均衡，并给出两个玩家的收益。
\end{tcolorbox}

(b) Find an SPE of the commitment game (where player I moves first), and give the payoffs to the two players.

\begin{tcolorbox}
(b) 找出承诺博弈(玩家 I 先行动)的一个子博弈完美均衡，并给出两个玩家的收益。
\end{tcolorbox}

(c) Are the equilibria in (a) and (b) unique?

\begin{tcolorbox}
(c) (a)和(b)中的均衡是否唯一？
\end{tcolorbox}

(d) Let \(G\) be a game where the best response \(y\left( x\right)\) of player II to any strategy \(x\) of player I is always unique. Show that in any SPE of the commitment game (where player I moves first), the payoff to player I is at least as large as his payoff in any equilibrium of the original game \(G\) .

\begin{tcolorbox}
(d) 设 \(G\) 为一个博弈，其中玩家 II 对玩家 I 的任何策略 \(x\) 的最优反应 \(y\left( x\right)\) 总是唯一的。证明在承诺博弈(玩家 I 先行动)的任何子博弈完美均衡(SPE)中，玩家 I 的收益至少不低于原博弈 \(G\) 中任何均衡下他的收益。
\end{tcolorbox}

\section*{Expected Utility}

\begin{tcolorbox}
\section*{期望效用}
\end{tcolorbox}

In a non-cooperative game, every outcome is associated with a payoff to each player that the player wants to maximize. The payoff represents the player's preference for the outcome. The games considered so far often have a pure-strategy equilibrium, where the preference applies to deterministic outcomes and is usually straightforward. From Chapter 6 onwards we will consider mixed strategies, which are actively randomized actions. The game is then analyzed under the assumption that players faced with random outcomes maximize their expected payoffs. The present chapter is about justifications for this assumption.

\begin{tcolorbox}
在非合作博弈中，每个结果都对应着每个玩家希望最大化的收益。收益代表了玩家对结果的偏好。到目前为止所考虑的博弈通常有纯策略均衡，此时偏好适用于确定性结果，通常较为直接。从第 6 章开始，我们将考虑混合策略，即主动随机化的行动。然后在玩家面对随机结果时最大化其期望效用的假设下分析博弈。本章正是关于这一假设的合理性证明。
\end{tcolorbox}

A player’s payoff is often called utility, with a utility function \(u : X \rightarrow  \mathbb{R}\) where \(X\) is the set of outcomes of a game. Suppose the outcomes \({x}_{1},\ldots ,{x}_{n}\) have probabilities \({p}_{1},\ldots ,{p}_{n}\) , respectively, which is called a lottery over the set of outcomes. The expected utility for this lottery \(P\) is \(\mathbb{E}\left( {u,P}\right)  = {p}_{1}u\left( {x}_{1}\right)  + \cdots  + {p}_{n}u\left( {x}_{n}\right)\) . An expected-utility function represents a preference of the player for lotteries in the sense that the player strictly prefers one lottery over another if it has a higher expected utility, and is indifferent between two lotteries if their expected utility is the same.

\begin{tcolorbox}
玩家的收益通常称为效用，其效用函数为 \(u : X \rightarrow  \mathbb{R}\) ，其中 \(X\) 是博弈结果的集合。假设结果 \({x}_{1},\ldots ,{x}_{n}\) 的概率分别为 \({p}_{1},\ldots ,{p}_{n}\) ，这被称为结果集合上的一个彩票。该彩票 \(P\) 的期望效用为 \(\mathbb{E}\left( {u,P}\right)  = {p}_{1}u\left( {x}_{1}\right)  + \cdots  + {p}_{n}u\left( {x}_{n}\right)\) 。期望效用函数代表了玩家对彩票的偏好，即如果一个彩票的期望效用更高，玩家会严格偏好该彩票；如果两个彩票的期望效用相同，玩家对它们无差异。
\end{tcolorbox}

One of the "rationality" assumptions of game theory is that players are expected-utility maximizers. This seems like a rather strong assumption in the sense that the player is "risk neutral" with respect to the received utility or payoff. However, as we explain in this chapter, this is a misconception. The payoff does not mean money, but is an abstract and much more general representation of a player's preference for lotteries, as long as this preference fulfills certain consistency conditions.

\begin{tcolorbox}
博弈论的“理性”假设之一是玩家是期望效用最大化者。从玩家在收到的效用或收益方面是“风险中性”的意义上来说，这似乎是一个相当强的假设。然而，正如我们在本章中所解释的，这是一种误解。收益并不意味着金钱，而是玩家对彩票偏好的一种抽象且更一般的表示，只要这种偏好满足某些一致性条件。
\end{tcolorbox}

The consistency conditions about the player's preference between lotteries are known as the axioms (basic assumptions) of von Neumann and Morgenstern (1947). The two main axioms are

\begin{tcolorbox}
关于玩家对彩票偏好的一致性条件被称为冯·诺伊曼和摩根斯坦(1947)公理(基本假设)。两个主要公理是
\end{tcolorbox}

\begin{itemize}
\item the preference is total (any two lotteries are comparable) and transitive;
\end{itemize}

\begin{tcolorbox}
\begin{itemize}
\item 偏好是完全的(任何两个彩票都是可比的)且具有传递性；
\end{itemize}
\end{tcolorbox}

\begin{itemize}
\item any preference or indifference between two lotteries is preserved when one of the lotteries is substituted for the other in an otherwise identical larger lottery (independence axiom).
\end{itemize}

\begin{tcolorbox}
\begin{itemize}
\item 当在一个其他方面相同的更大彩票中，用一个彩票替换另一个彩票时，两个彩票之间的任何偏好或无差异关系都保持不变(独立性公理)。
\end{itemize}
\end{tcolorbox}

The theorem of von Neumann and Morgenstern (1947) states that these axioms (plus a more technical axiom) allow the construction of an expected-utility function.

\begin{tcolorbox}
冯·诺伊曼和摩根斯坦(1947)定理指出，这些公理(加上一个更具技术性的公理)允许构建一个期望效用函数。
\end{tcolorbox}

We derive this theorem in this chapter. The main idea is relatively straightforward and explained in the first summary Section 5.2, under the assumption (which we relax later) that there is a least preferred outcome \({x}_{0}\) and a most preferred outcome \({x}_{1}\) . With these "reference outcomes", every other outcome \(z\) is replaced by an equally preferred lottery between \({x}_{0}\) and \({x}_{1}\) , where \({x}_{1}\) is chosen with probability \(\alpha\) . The utility for \(z\) can then be defined as \(u\left( z\right)  = \alpha\) . With this representation "utility = probability", the expected-utility theorem holds because an expected probability is again a probability (see Figure 5.1).

\begin{tcolorbox}
我们在本章中推导这个定理。主要思想相对直接，在第一个总结部分 5.2 中进行了解释，假设(我们稍后会放宽这个假设)存在一个最不偏好的结果 \({x}_{0}\) 和一个最偏好的结果 \({x}_{1}\) 。有了这些“参考结果”，每个其他结果 \(z\) 都被替换为 \({x}_{0}\) 和 \({x}_{1}\) 之间一个同样受偏好的彩票，其中 \({x}_{1}\) 以概率 \(\alpha\) 被选中。然后 \(z\) 的效用可以定义为 \(u\left( z\right)  = \alpha\) 。通过这种“效用 = 概率”的表示，期望效用定理成立，因为期望概率仍然是一个概率(见图 5.1)。
\end{tcolorbox}

The remaining sections of this "foundational" chapter give in essence a brief introduction to decision theory, which studies how an individual or a group can make a decision between several given alternatives (see Section 5.3). We focus on a single player who faces decisions under risk, that is, the alternatives are lotteries with known probabilities for outcomes (Section 5.4). As part of a decision analysis, a utility function is often used as a tool for simplifying complex decisions, by eliciting the decision maker's preferences from artificial simpler scenarios.

\begin{tcolorbox}
这一“基础”章节的其余部分本质上简要介绍了决策理论，该理论研究个人或群体如何在几个给定的替代方案之间做出决策(见第 5.3 节)。我们关注面临风险决策的单个玩家，即替代方案是已知结果概率的彩票(第 5.4 节)。作为决策分析的一部分，效用函数通常被用作简化复杂决策的工具，通过从人为的简单场景中引出决策者的偏好。
\end{tcolorbox}

A lottery with only one possible outcome means to get that outcome with certainty. Any expected-utility function \(u\) therefore also represents preferences under certainty, via the values \(u\left( x\right)\) for outcomes \(x\) . If the decision is only between alternatives that have certain outcomes, the utility function is called ordinal and can be changed without changing the preference by applying any strictly increasing transformation (Section 5.5).

\begin{tcolorbox}
只有一个可能结果的彩票意味着肯定会得到该结果。因此，任何期望效用函数 \(u\) 也通过结果 \(x\) 的值 \(u\left( x\right)\) 来表示确定性条件下的偏好。如果决策仅在具有确定结果的选项之间进行，那么该效用函数被称为序数效用函数，并且可以通过应用任何严格递增变换来改变它，而不改变偏好(第5.5节)。
\end{tcolorbox}

For decisions between risky alternatives, the expected-utility function is cardinal, which means that the preference is only preserved when applying positive-affine transformations like for a temperature scale (Section 5.6). Then the utility function \(u\) is uniquely determined by fixing the "utility scale" by choosing an origin \({x}_{0}\) so that \(u\left( {x}_{0}\right)  = 0\) , and a unit, that is, a more preferred outcome \({x}_{1}\) so that \(u\left( {x}_{1}\right)  = 1\) . All other values of \(u\) are then already determined by the player’s preference for simple lotteries, which have only two outcomes \(x\) and \(y\) , say, where \(x\) occurs with probability \(1 - \alpha\) and \(y\) with probability \(\alpha\) . We use the notation \(x{ \land  }_{\alpha }y\) (see (5.16), unique to this book) for this simple lottery where \({ \land  }_{\alpha }\) is meant to represent a little "probability tree" where the right branch has probability \(\alpha\) .

\begin{tcolorbox}
对于有风险的选项之间的决策，期望效用函数是基数效用函数，这意味着只有在应用类似于温度刻度的正仿射变换时，偏好才会保持不变(第5.6节)。然后，通过选择一个原点 \({x}_{0}\) 来固定“效用尺度”，使得 \(u\left( {x}_{0}\right)  = 0\) ，并选择一个单位，即一个更受偏好的结果 \({x}_{1}\) ，使得 \(u\left( {x}_{1}\right)  = 1\) ，效用函数 \(u\) 就被唯一确定了。 \(u\) 的所有其他值则已经由玩家对简单彩票的偏好所确定，这些简单彩票只有两个结果 \(x\) 和 \(y\) ，比如说，其中 \(x\) 以概率 \(1 - \alpha\) 发生， \(y\) 以概率 \(\alpha\) 发生。我们使用符号 \(x{ \land  }_{\alpha }y\) (见(5.16)，本书特有的表示法)来表示这个简单彩票，其中 \({ \land  }_{\alpha }\) 表示一个小的“概率树”，其右分支的概率为 \(\alpha\) 。
\end{tcolorbox}

Simple lotteries specify fully how the values \(u\left( z\right)\) of the utility function result from probabilities, namely as \(u\left( z\right)  = \alpha\) if the player is indifferent between \(z\) and the simple lottery \(x{ \land  }_{\alpha }y\) , where \(u\left( x\right)  = 0\) and \(u\left( y\right)  = 1\) . This also means that the player's preference for lotteries can be constructed from these simple lotteries as soon as one accepts the consistency axioms. In this way, constructing a utility function is a useful decision aid, because lotteries with many outcomes and associated probabilities are rather complex prospects. In practice, finding a probability \(\alpha\) so that the player is indifferent between \(z\) and \(x{ \land  }_{\alpha }y\) is conceptually easier if \(x\) and \(y\) are not extreme like the worst and best possible outcomes that can occur. For that reason, we do not even assume (as in Section 5.2) that such worst and best possible outcomes exist for our main Theorem 5.4. The proof of the theorem then requires adapting the origin and unit of the utility function to the lottery under consideration, which leads to somewhat lengthy but straightforward manipulations with positive-affine transformations.

\begin{tcolorbox}
简单彩票完全规定了效用函数的值 \(u\left( z\right)\) 是如何由概率产生的，即如果玩家对 \(z\) 和简单彩票 \(x{ \land  }_{\alpha }y\) 无差异，那么 \(u\left( z\right)  = \alpha\) ，其中 \(u\left( x\right)  = 0\) 且 \(u\left( y\right)  = 1\) 。这也意味着，一旦接受了一致性公理，玩家对彩票的偏好就可以从这些简单彩票中构建出来。通过这种方式，构建效用函数是一种有用的决策辅助工具，因为具有多个结果和相关概率的彩票是相当复杂的前景。在实践中，如果 \(x\) 和 \(y\) 不是像可能出现的最差和最好结果那样极端，那么从概念上讲，找到一个概率 \(\alpha\) ，使得玩家对 \(z\) 和 \(x{ \land  }_{\alpha }y\) 无差异会更容易。出于这个原因，在我们的主要定理5.4中，我们甚至不假设(如第5.2节那样)存在这样的最差和最好可能结果。那么，该定理的证明就需要将效用函数的原点和单位调整为所考虑的彩票，这会导致使用正仿射变换进行一些冗长但直接的操作。
\end{tcolorbox}

In Section 5.7 we explain the consistency axioms of independence and continuity that imply the existence of an expected-utility function, which is proved in full in Section 5.8. In Section 5.9, outcomes are themselves real numbers (such as amounts of money) where a lottery has an expected value. Risk aversion means that getting this expected value for certain is always preferred to the (risky) lottery, which is equivalent to a concave utility function. In the final Section 5.10 we discuss the "rationality" of the utility axioms (as a short "philosophical" discussion that we generally avoid in this book), and give further references.

\begin{tcolorbox}
在第5.7节中，我们解释了独立性和连续性的一致性公理，这些公理意味着存在一个期望效用函数，这将在第5.8节中进行完整证明。在第5.9节中，结果本身是实数(如金额)，其中彩票有一个期望值。风险厌恶意味着确定地获得这个期望值总是比(有风险的)彩票更受偏好，这等价于一个凹效用函数。在最后一节5.10中，我们讨论了效用公理的“理性”(作为一个简短的“哲学”讨论，我们在本书中通常避免此类讨论)，并给出了进一步的参考文献。
\end{tcolorbox}

\section*{5.1 Prerequisites and Learning Outcomes}

\begin{tcolorbox}
\section*{5.1 先决条件和学习成果}
\end{tcolorbox}

We need the concept of a binary relation, like that of an order (see Definition 1.2). Expected values have been defined in (4.1).

\begin{tcolorbox}
我们需要二元关系的概念，比如序关系的概念(见定义1.2)。期望值已在(4.1)中定义。
\end{tcolorbox}

After studying this chapter, you should be able to

\begin{tcolorbox}
学习本章后，你应该能够
\end{tcolorbox}

\begin{itemize}
\item understand a player's preference as a binary relation on the set of lotteries over outcomes, and its expected-utility representation;
\end{itemize}

\begin{tcolorbox}
\begin{itemize}
\item 将玩家的偏好理解为结果的彩票集合上的二元关系，以及其期望效用表示；
\end{itemize}
\end{tcolorbox}

\begin{itemize}
\item distinguish ordinal functions for decisions under certainty and cardinal utility functions for decisions under risk;
\end{itemize}

\begin{tcolorbox}
\begin{itemize}
\item 区分确定性决策下的序数效用函数和风险决策下的基数效用函数；
\end{itemize}
\end{tcolorbox}

\begin{itemize}
\item understand the role of a utility function for creating a total and transitive preference, for example an additive function for outcomes with multiple attributes (Exercise 5.1);
\end{itemize}

\begin{tcolorbox}
\begin{itemize}
\item 理解效用函数在创建完全且传递的偏好中的作用，例如对于具有多个属性的结果的加性函数(练习5.1)；
\end{itemize}
\end{tcolorbox}

\begin{itemize}
\item explain the lexicographic order and why it is problematic both in practice, see the example (5.10), and theory, see Figure 5.3;
\end{itemize}

\begin{tcolorbox}
\begin{itemize}
\item 解释字典序，以及为什么它在实践(见示例(5.10))和理论(见图5.3)中都存在问题；
\end{itemize}
\end{tcolorbox}

\begin{itemize}
\item apply positive-affine transformations to payoffs and know why they represent the same preference for lotteries;
\end{itemize}

\begin{tcolorbox}
\begin{itemize}
\item 对收益应用正仿射变换，并知道为什么它们表示对彩票的相同偏好；
\end{itemize}
\end{tcolorbox}

\begin{itemize}
\item use simple lotteries and their notation \(x{ \land  }_{\alpha }y\) ;
\end{itemize}

\begin{tcolorbox}
\begin{itemize}
\item 使用简单彩票及其符号 \(x{ \land  }_{\alpha }y\) ；
\end{itemize}
\end{tcolorbox}

\begin{itemize}
\item state the von Neumann-Morgenstern consistency axioms and their use;
\end{itemize}

\begin{tcolorbox}
\begin{itemize}
\item 阐述冯·诺伊曼 - 摩根斯坦一致性公理及其用途；
\end{itemize}
\end{tcolorbox}

\begin{itemize}
\item construct an expected-utility function via probabilities;
\end{itemize}

\begin{tcolorbox}
\begin{itemize}
\item 通过概率构建期望效用函数；
\end{itemize}
\end{tcolorbox}

\begin{itemize}
\item apply concepts of risk aversion and risk neutrality to outcomes that are themselves numbers.
\end{itemize}

\begin{tcolorbox}
\begin{itemize}
\item 将风险厌恶和风险中性的概念应用于本身为数值的结果。
\end{itemize}
\end{tcolorbox}

\section*{5.2 Summary}

\begin{tcolorbox}
\section*{5.2 总结}
\end{tcolorbox}

This section is a kind of "executive summary" (with more details later) of the main theorem of this chapter: If the player's preference for lotteries fulfills certain consistency axioms, then this preference can be represented by an expected-utility function.

\begin{tcolorbox}
本节是本章主要定理的一种“内容提要”(后续会有更多细节):如果玩家对彩票的偏好满足某些一致性公理，那么这种偏好可以用一个期望效用函数来表示。
\end{tcolorbox}

We study decisions under risk of a single decision maker (here still called player for brevity). Risk means the player faces lotteries with outcomes \({x}_{1},\ldots ,{x}_{n}\) from a set \(X\) , which depend on the lottery and occur with known probabilities \({p}_{1},\ldots ,{p}_{n}\) . For simplicity, we consider only lotteries with finitely many outcomes, but the set \(X\) of possible outcomes may be infinite, for example a real interval.

\begin{tcolorbox}
我们研究单个决策者(此处为简便仍称为玩家)在风险下的决策。风险意味着玩家面临结果 \({x}_{1},\ldots ,{x}_{n}\) 来自集合 \(X\) 的彩票，这些结果取决于彩票且以已知概率 \({p}_{1},\ldots ,{p}_{n}\) 出现。为简单起见，我们仅考虑具有有限个结果的彩票，但可能结果的集合 \(X\) 可能是无限的，例如一个实数区间。
\end{tcolorbox}

The set of lotteries over \(X\) is denoted by \(\Delta \left( X\right)\) . We start with a binary preference relation \(\precsim\) on \(\Delta \left( X\right)\) , where \(P \precsim  Q\) means that the player weakly prefers lottery \(Q\) to lottery \(P\) . If \(P \precsim  Q\) and \(Q \precsim  P\) hold then this is written as \(P \sim  Q\) and means that the player is indifferent between \(P\) and \(Q\) . If \(P \precsim  Q\) but not \(Q \precsim  P\) then this defines the strict preference \(P \prec  Q\) .

\begin{tcolorbox}
关于 \(X\) 的彩票集合记为 \(\Delta \left( X\right)\) 。我们从 \(\Delta \left( X\right)\) 上的二元偏好关系 \(\precsim\) 开始，其中 \(P \precsim  Q\) 表示玩家弱偏好彩票 \(Q\) 甚于彩票 \(P\) 。如果 \(P \precsim  Q\) 和 \(Q \precsim  P\) 都成立，那么记为 \(P \sim  Q\) ，表示玩家对 \(P\) 和 \(Q\) 无差异。如果 \(P \precsim  Q\) 成立但 \(Q \precsim  P\) 不成立，那么这就定义了严格偏好 \(P \prec  Q\) 。
\end{tcolorbox}

The goal is to study conditions on the player’s preference relation \(\precsim\) on \(\Delta \left( X\right)\) so that it can be represented by an expected-utility function. This is a function \(u : X \rightarrow  \mathbb{R}\) so that for \(P,Q \in  \Delta \left( X\right)\) ,

\begin{tcolorbox}
目标是研究玩家在 \(\Delta \left( X\right)\) 上的偏好关系 \(\precsim\) 应满足的条件，以便能用一个期望效用函数来表示它。这是一个函数 \(u : X \rightarrow  \mathbb{R}\) ，使得对于 \(P,Q \in  \Delta \left( X\right)\) ，
\end{tcolorbox}

\[
P \precsim  Q\; \Leftrightarrow  \;\mathbb{E}\left( {u,P}\right)  \leq  \mathbb{E}\left( {u,Q}\right)  \tag{5.1}
\]

where \(\mathbb{E}\left( {u,P}\right)\) is the expected value of \(u\) under \(P\) . In a decision situation (and later in a game) the player therefore wants to maximize his expected utility.

\begin{tcolorbox}
其中 \(\mathbb{E}\left( {u,P}\right)\) 是 \(u\) 在 \(P\) 下的期望值。因此，在决策情形(以及之后的博弈中)玩家希望最大化其期望效用。
\end{tcolorbox}

If a utility function exists, then it imposes a number of conditions on the preference \(\precsim\) of the player. Suitable conditions called consistency axioms will be sufficient for the existence of an expected-utility function.

\begin{tcolorbox}
如果存在一个效用函数，那么它会对玩家的偏好 \(\precsim\) 施加一些条件。合适的条件(称为一致性公理)将足以保证期望效用函数的存在性。
\end{tcolorbox}

The first axiom states that \(\precsim\) is total and transitive. The only difference to a total order (see Definition 1.2) is that \(P \sim  Q\) is allowed to hold if \(P \neq  Q\) .

\begin{tcolorbox}
第一个公理表明 \(\precsim\) 是全序且传递的。与全序(见定义 1.2)的唯一区别在于，当 \(P \neq  Q\) 时允许 \(P \sim  Q\) 成立。
\end{tcolorbox}

The lotteries in \(\Delta \left( X\right)\) include deterministic lotteries that choose a single outcome \(x\) with probability 1, which are denoted by \({\mathbf{1}}_{x}\) . The preference relation \(\precsim\) therefore applies to deterministic outcomes \(x\) and \(y\) where \(x \precsim  y\) stands for \({\mathbf{1}}_{x} \precsim  {\mathbf{1}}_{y}\) , which holds if and only if \(u\left( x\right)  \leq  u\left( y\right)\) for an expected-utility function \(u\) . This preference for deterministic outcomes is preserved if \(u\left( x\right)\) is replaced by \(\phi \left( {u\left( x\right) }\right)\) for a strictly increasing function \(\phi\) , for example \(\phi \left( u\right)  = {e}^{u}\) . That is, the utility function for deterministic outcomes is ordinal in the sense that more preferred outcomes have higher values of the utility function, and this is all that matters.

\begin{tcolorbox}
\(\Delta \left( X\right)\) 中的彩票包括确定性彩票，即它们以概率 1 选择单个结果 \(x\) ，记为 \({\mathbf{1}}_{x}\) 。因此，偏好关系 \(\precsim\) 适用于确定性结果 \(x\) 和 \(y\) ，其中 \(x \precsim  y\) 代表 \({\mathbf{1}}_{x} \precsim  {\mathbf{1}}_{y}\) ，当且仅当对于期望效用函数 \(u\) 有 \(u\left( x\right)  \leq  u\left( y\right)\) 时成立。如果对于严格递增函数 \(\phi\) ，将 \(u\left( x\right)\) 替换为 \(\phi \left( {u\left( x\right) }\right)\) ，例如 \(\phi \left( u\right)  = {e}^{u}\) ，那么对确定性结果的这种偏好得以保留。也就是说，确定性结果的效用函数是序数意义上的，即更受偏好的结果具有更高的效用函数值，而这就是全部要点。
\end{tcolorbox}

Whereas a utility function for decisions under certainty is ordinal, an expected-utility function \(u\) for decisions under risk is cardinal. This means that the only admitted transformations of \(u\) are positive-affine, given by \({au} + b\) for constants \(a\) and \(b\) with \(a > 0\) , for example as between the temperature scales Celsius and Fahrenheit. This uniqueness holds because an expected-utility function evaluates lotteries via their expected value. It already applies when considering only simple lotteries that have only two outcomes (Proposition 5.2). The "cardinal" property of an expected-utility function refers to its uniqueness up to positive-affine transformations, not to an intrinsic "strength" of preferences for outcomes.

\begin{tcolorbox}
确定性条件下决策的效用函数是序数的，而风险条件下决策的期望效用函数 \(u\) 是基数的。这意味着 \(u\) 唯一允许的变换是正仿射变换，由 \({au} + b\) 给出，其中 \(a\) 和 \(b\) 为常数且 \(a > 0\) ，例如温度刻度中的摄氏度和华氏度之间的变换。这种唯一性成立是因为期望效用函数通过彩票的期望值来评估彩票。当仅考虑只有两种结果的简单彩票时，这一点就已经适用(命题5.2)。期望效用函数的“基数”性质是指它在正仿射变换下的唯一性，而不是指对结果偏好的内在“强度”。
\end{tcolorbox}

By applying a positive-affine transformation, for two arbitrary outcomes \({x}_{0}\) and \({x}_{1}\) so that \({x}_{0} \prec  {x}_{1}\) , one can choose the expected-utility function \(u\) so that \(u\left( {x}_{0}\right)  = 0\) and \(u\left( {x}_{1}\right)  = 1\) , which fixes the origin and unit of the utility scale. Consider now any outcome \(z\) so that \({x}_{0} \precsim  z\) and \(z \precsim  {x}_{1}\) (which we abbreviate as \({x}_{0} \precsim  z \precsim  {x}_{1}\) using that \(\precsim\) is transitive). Then \(0 \leq  u\left( z\right)  \leq  1\) , and the precise value of \(u\left( z\right)\) is given by the probability \(\alpha\) so that the player is indifferent between \(z\) and the simple lottery that chooses \({x}_{0}\) with probability \(1 - \alpha\) and \({x}_{1}\) with probability \(\alpha\) (which we write as \({x}_{0}{ \land  }_{\alpha }{x}_{1}\) , see (5.16), where \({ \land  }_{\alpha }\) is meant to suggest a small probability tree). This holds because the expected utility \(\left( {1 - \alpha }\right) u\left( {x}_{0}\right)  + {\alpha u}\left( {x}_{1}\right)\) for this simple lottery is \(\alpha\) .

\begin{tcolorbox}
通过应用正仿射变换，对于任意两个结果 \({x}_{0}\) 和 \({x}_{1}\) ，使得 \({x}_{0} \prec  {x}_{1}\) ，可以选择期望效用函数 \(u\) ，使得 \(u\left( {x}_{0}\right)  = 0\) 且 \(u\left( {x}_{1}\right)  = 1\) ，这就确定了效用尺度的原点和单位。现在考虑任意结果 \(z\) ，使得 \({x}_{0} \precsim  z\) 且 \(z \precsim  {x}_{1}\) (由于 \(\precsim\) 是传递的，我们将其缩写为 \({x}_{0} \precsim  z \precsim  {x}_{1}\) )。那么 \(0 \leq  u\left( z\right)  \leq  1\) ，并且 \(u\left( z\right)\) 的精确值由概率 \(\alpha\) 给出，使得玩家在 \(z\) 和以概率 \(1 - \alpha\) 选择 \({x}_{0}\) 、以概率 \(\alpha\) 选择 \({x}_{1}\) 的简单彩票之间无差异(我们将其写为 \({x}_{0}{ \land  }_{\alpha }{x}_{1}\) ，见(5.16)，其中 \({ \land  }_{\alpha }\) 表示一个小概率树)。这是因为这个简单彩票的期望效用 \(\left( {1 - \alpha }\right) u\left( {x}_{0}\right)  + {\alpha u}\left( {x}_{1}\right)\) 为 \(\alpha\) 。
\end{tcolorbox}

With this scaling \(u\left( {x}_{0}\right)  = 0\) and \(u\left( {x}_{1}\right)  = 1\) , each value \(u\left( z\right)\) of the expected-utility function is therefore just the probability \(\alpha\) of choosing \({x}_{1}\) (and \({x}_{0}\) otherwise) that makes the player indifferent between this lottery \({x}_{0}{ \land  }_{\alpha }{x}_{1}\) and getting \(z\) with certainty. This insight has profound consequences. Assume for simplicity that \({x}_{0}\) is least preferred and \({x}_{1}\) is most preferred among all outcomes, so that \({x}_{0} \precsim  z \precsim  {x}_{1}\) for all outcomes \(z\) . Define \(u\left( z\right)\) in \(\left\lbrack  {0,1}\right\rbrack\) as the probability \(\alpha\) so that \(z \sim  {x}_{0}{ \land  }_{\alpha }{x}_{1}\) . Then \(\alpha\) is unique, and \(u\) is therefore unique up to positive-affine transformations.

\begin{tcolorbox}
通过这种缩放 \(u\left( {x}_{0}\right)  = 0\) 和 \(u\left( {x}_{1}\right)  = 1\) ，期望效用函数的每个值 \(u\left( z\right)\) 因此仅仅是选择 \({x}_{1}\) 的概率 \(\alpha\) (否则选择 \({x}_{0}\) )，该概率使得玩家在这个彩票 \({x}_{0}{ \land  }_{\alpha }{x}_{1}\) 和确定获得 \(z\) 之间无差异。这一见解具有深远的影响。为简单起见，假设在所有结果中 \({x}_{0}\) 是最不被偏好的， \({x}_{1}\) 是最被偏好的，使得对于所有结果 \(z\) 都有 \({x}_{0} \precsim  z \precsim  {x}_{1}\) 。在 \(\left\lbrack  {0,1}\right\rbrack\) 中，将 \(u\left( z\right)\) 定义为概率 \(\alpha\) ，使得 \(z \sim  {x}_{0}{ \land  }_{\alpha }{x}_{1}\) 。那么 \(\alpha\) 是唯一的，因此 \(u\) 在正仿射变换下是唯一的。
\end{tcolorbox}

\begin{center}
\includegraphics[max width=0.9\textwidth]{images/bo_d2e5cmbef24c73b03tn0_118_202_1497_1246_379_0.jpg}
\end{center}
\hspace*{3em} 

Figure 5.1 Equivalent probability trees for the lottery \(P\) with probabilities \({p}_{1},\ldots ,{p}_{n}\) for outcomes \({z}_{1},\ldots ,{z}_{n}\) and probability \(p = \mathop{\sum }\limits_{{i = 1}}^{n}{p}_{i}{\alpha }_{i}\) for outcome \({x}_{1}\) .

\begin{tcolorbox}
图5.1 彩票 \(P\) 的等价概率树，结果 \({z}_{1},\ldots ,{z}_{n}\) 的概率为 \({p}_{1},\ldots ,{p}_{n}\) ，结果 \({x}_{1}\) 的概率为 \(p = \mathop{\sum }\limits_{{i = 1}}^{n}{p}_{i}{\alpha }_{i}\) 。
\end{tcolorbox}

Most importantly, the preference for an arbitrary lottery \(P\) is represented by its expectation \(\mathbb{E}\left( {u,P}\right)\) . Namely, suppose that \(P\) chooses the outcomes \({z}_{1},\ldots ,{z}_{n}\) with probabilities \({p}_{1},\ldots ,{p}_{n}\) , respectively. Let \({z}_{i} \sim  {x}_{0}{ \land  }_{{\alpha }_{i}}{x}_{1}\) where \({\alpha }_{i} = u\left( {z}_{i}\right)\) for \(1 \leq  i \leq  n\) . Figure 5.1 shows on the left the probability tree for \(P\) . In the middle tree, each outcome \({z}_{i}\) is replaced by its equally preferred simple lottery \({x}_{0}{ \land  }_{{\alpha }_{i}}{x}_{1}\) between \({x}_{0}\) and \({x}_{1}\) . In that tree, the only outcomes are \({x}_{0}\) and \({x}_{1}\) . Along the \(i\) th edge from the root, outcome \({x}_{1}\) is reached with probability \({p}_{i}{\alpha }_{i}\) because independent probabilities are multiplied. Summed over \(i = 1,\ldots ,n\) , this gives the overall probability

\begin{tcolorbox}
最重要的是，对任意彩票 \(P\) 的偏好由其期望 \(\mathbb{E}\left( {u,P}\right)\) 表示。即，假设 \(P\) 分别以概率 \({p}_{1},\ldots ,{p}_{n}\) 选择结果 \({z}_{1},\ldots ,{z}_{n}\) 。设 \({z}_{i} \sim  {x}_{0}{ \land  }_{{\alpha }_{i}}{x}_{1}\) ，其中对于 \(1 \leq  i \leq  n\) 有 \({\alpha }_{i} = u\left( {z}_{i}\right)\) 。图 5.1 左侧展示了 \(P\) 的概率树。在中间的树中，每个结果 \({z}_{i}\) 都被其在 \({x}_{0}\) 和 \({x}_{1}\) 之间等偏好的简单彩票 \({x}_{0}{ \land  }_{{\alpha }_{i}}{x}_{1}\) 所替代。在该树中，唯一的结果是 \({x}_{0}\) 和 \({x}_{1}\) 。沿着从根节点出发的第 \(i\) 条边，以概率 \({p}_{i}{\alpha }_{i}\) 到达结果 \({x}_{1}\) ，因为独立概率是相乘的。对 \(i = 1,\ldots ,n\) 求和，这给出了总体概率
\end{tcolorbox}

\[
p = \mathop{\sum }\limits_{{i = 1}}^{n}{p}_{i}{\alpha }_{i} \tag{5.2}
\]

for \({x}_{1}\) and \(1 - p\) for \({x}_{0}\) . This is shown in the right tree in Figure 5.1 that aggregates the probabilities for \({x}_{0}\) and \({x}_{1}\) . Hence, the player is indifferent between \(P\) and \({x}_{0}{ \land  }_{p}{x}_{1}\) where \(p = \mathbb{E}\left( {u,P}\right)\) by (5.2) because \({\alpha }_{i} = u\left( {z}_{i}\right)\) for \(1 \leq  i \leq  n\) . The comparison between different lotteries \(P\) and \(Q\) therefore becomes a comparison between their corresponding probabilities for \({x}_{1}\) : If \(p = \mathbb{E}\left( {u,P}\right)\) as described and in a similar way \(q = \mathbb{E}\left( {u,Q}\right)\) , then \(P \precsim  Q\) if and only if \(p \leq  q\) , which is exactly the property (5.1) of an expected-utility function.

\begin{tcolorbox}
对于 \({x}_{1}\) 和 \(1 - p\) 对于 \({x}_{0}\) 。这在图 5.1 右侧的树中展示，该树汇总了 \({x}_{0}\) 和 \({x}_{1}\) 的概率。因此，玩家在 \(P\) 和 \({x}_{0}{ \land  }_{p}{x}_{1}\) 之间无差异，其中根据 (5.2) 有 \(p = \mathbb{E}\left( {u,P}\right)\) ，因为对于 \(1 \leq  i \leq  n\) 有 \({\alpha }_{i} = u\left( {z}_{i}\right)\) 。因此，不同彩票 \(P\) 和 \(Q\) 之间的比较就变成了它们对应 \({x}_{1}\) 的概率之间的比较:如果如所述有 \(p = \mathbb{E}\left( {u,P}\right)\) ，并且以类似的方式有 \(q = \mathbb{E}\left( {u,Q}\right)\) ，那么当且仅当 \(p \leq  q\) 时 \(P \precsim  Q\) ，这正是期望效用函数(expected-utility function)的性质 (5.1)。
\end{tcolorbox}

In summary, assume "reference outcomes" \({x}_{0}\) and \({x}_{1}\) with \(u\left( {x}_{0}\right)  = 0\) and \(u\left( {x}_{1}\right)  = 1\) so that \({x}_{0} \precsim  {z}_{i} \precsim  {x}_{1}\) for the outcomes \({z}_{1},\ldots ,{z}_{n}\) of the considered lottery \(P\) . Because the values \(u\left( {z}_{i}\right)\) of the utility function are probabilities \({\alpha }_{i}\) for simple lotteries \({x}_{0}{ \land  }_{{\alpha }_{i}}{x}_{1}\) between these reference outcomes, the expected utility of \(P\) as in Figure 5.1 is just the aggregated probability \(p\) for \({x}_{1}\) as in (5.2), because probabilities can be multiplied.

\begin{tcolorbox}
总之，假设“参考结果” \({x}_{0}\) 和 \({x}_{1}\) ，且 \(u\left( {x}_{0}\right)  = 0\) 和 \(u\left( {x}_{1}\right)  = 1\) ，使得对于所考虑的彩票 \(P\) 的结果 \({z}_{1},\ldots ,{z}_{n}\) 有 \({x}_{0} \precsim  {z}_{i} \precsim  {x}_{1}\) 。因为效用函数的值 \(u\left( {z}_{i}\right)\) 是这些参考结果之间简单彩票 \({x}_{0}{ \land  }_{{\alpha }_{i}}{x}_{1}\) 的概率 \({\alpha }_{i}\) ，如图 5.1 所示， \(P\) 的期望效用就是如 (5.2) 中所示 \({x}_{1}\) 的汇总概率 \(p\) ，因为概率可以相乘。
\end{tcolorbox}

Figure 5.1 is therefore a derivation of an expected-utility function. Crucially, this argument can be made in terms of the preference relation \(\precsim\) alone, provided the following hold:

\begin{tcolorbox}
因此，图 5.1 是期望效用函数(expected-utility function)的一种推导。关键在于，只要满足以下条件，这个论证就可以仅用偏好关系(preference relation) \(\precsim\) 来进行:
\end{tcolorbox}

\begin{itemize}
\item If \({x}_{0} \prec  z \prec  {x}_{1}\) for outcomes \({x}_{0},{x}_{1},z\) , then \(z \sim  {x}_{0}{ \land  }_{\alpha }{x}_{1}\) for some \(\alpha  \in  \left( {0,1}\right)\) .
\end{itemize}

\begin{tcolorbox}
\begin{itemize}
\item 如果对于结果 \({x}_{0},{x}_{1},z\) 有 \({x}_{0} \prec  z \prec  {x}_{1}\) ，那么对于某个 \(\alpha  \in  \left( {0,1}\right)\) 有 \(z \sim  {x}_{0}{ \land  }_{\alpha }{x}_{1}\) 。
\end{itemize}
\end{tcolorbox}

\begin{itemize}
\item If \(z \sim  {x}_{0}{ \land  }_{\alpha }{x}_{1}\) , then substituting in any lottery the outcome \(z\) with \({x}_{0}{ \land  }_{\alpha }{x}_{1}\) creates an equally preferred lottery.
\end{itemize}

\begin{tcolorbox}
\begin{itemize}
\item 如果 \(z \sim  {x}_{0}{ \land  }_{\alpha }{x}_{1}\) ，那么在任何彩票中用 \({x}_{0}{ \land  }_{\alpha }{x}_{1}\) 替换结果 \(z\) 会得到一个同样受偏好的彩票。
\end{itemize}
\end{tcolorbox}

\begin{itemize}
\item In combined lotteries, only the probability for an outcome matters, not the structure of the lottery.
\end{itemize}

\begin{tcolorbox}
\begin{itemize}
\item 在组合彩票中，只有某个结果的概率重要，而不是彩票的结构。
\end{itemize}
\end{tcolorbox}

The first of these conditions is called the continuity axiom. The second condition, together with the corresponding condition with \(\sim\) replaced by \(<\) , is called the independence axiom. These axioms are assumed to hold for general lotteries and not just outcomes. The third condition is not considered a constraint on the preference relation because decisions under risk deal with "objective" known probabilities.

\begin{tcolorbox}
这些条件中的第一个被称为连续性公理(continuity axiom)。第二个条件，连同将 \(\sim\) 替换为 \(<\) 后的相应条件，被称为独立性公理(independence axiom)。假定这些公理适用于一般的彩票，而不仅仅是结果。第三个条件不被视为对偏好关系的约束，因为风险决策处理的是“客观”已知的概率。
\end{tcolorbox}

Hence, any preference relation \(\precsim\) on lotteries that is total, transitive, and fulfills the continuity and independence axioms can be represented by an expected-utility function. This is the main Theorem 5.4. Its full proof does not assume that "extreme" outcomes \({x}_{0}\) and \({x}_{1}\) exist so that \({x}_{0} \precsim  z \precsim  {x}_{1}\) for all outcomes \(z\) to define \(u\left( z\right)  = \alpha\) via \(z \sim  {x}_{0}{ \land  }_{\alpha }{x}_{1}\) . Instead, \({x}_{0}\) and \({x}_{1}\) can be chosen depending on the considered lottery to define the utility function \(u\left( z\right)\) "locally", which is then extended to all outcomes \(z\) using related lotteries and positive-affine transformations.

\begin{tcolorbox}
因此，彩票上任何完全、传递且满足连续性和独立性公理的偏好关系 \(\precsim\) 都可以用一个期望效用函数(expected-utility function)来表示。这就是主要的定理 5.4。其完整证明并不假定存在“极端”结果 \({x}_{0}\) 和 \({x}_{1}\) ，使得对于所有结果 \(z\) 有 \({x}_{0} \precsim  z \precsim  {x}_{1}\) ，从而通过 \(z \sim  {x}_{0}{ \land  }_{\alpha }{x}_{1}\) 来定义 \(u\left( z\right)  = \alpha\) 。相反，可以根据所考虑的彩票选择 \({x}_{0}\) 和 \({x}_{1}\) 来“局部地”定义效用函数 \(u\left( z\right)\) ，然后使用相关彩票和正仿射变换(positive-affine transformation)将其扩展到所有结果 \(z\) 。
\end{tcolorbox}

The remainder of this chapter explains in more detail what decision theory is about. For the development of game theory it can be skipped, except for the background material on line segments in Figure 5.4, which is needed for the geometric concept of convexity in the following chapters. Section 5.9 defines concave utility functions, which are used in Chapter 11 on bargaining.

\begin{tcolorbox}
本章的其余部分更详细地解释了决策理论(decision theory)是关于什么的。对于博弈论的发展来说，除了图 5.4 中关于线段(line segments)的背景材料(background material)(这在后续章节的凸性几何概念中是需要的)之外，这部分可以跳过。第 5.9 节定义了凹效用函数(concave utility function)，这在第 11 章的讨价还价中会用到。
\end{tcolorbox}

\section*{5.3 Decisions Under Risk}

\begin{tcolorbox}
\section*{5.3 风险决策}
\end{tcolorbox}

The following situation illustrates a decision problem under risk. Suppose you have booked a rental car for a week. Your liability for damage to the car is 1,000 Euros. At the rental desk you are offered additional insurance at the cost of 20 Euros per day that would reduce this excess from 1,000 Euros to zero, which would cost you an extra 140 Euros. Should you buy this insurance or not? The choice is between being 140 Euros out of pocket for sure, or risking with some probability \(p\) that you pay up to 1,000 Euros in case you damage the car.

\begin{tcolorbox}
以下情况说明了一个风险决策问题。假设你已经预订了一辆租车，租期为一周。你对汽车损坏的赔偿责任是 1000 欧元。在租车柜台，你可以以每天 20 欧元的成本购买额外保险，这将把你的自付额从 1000 欧元降至零，这意味着你要多花 140 欧元。你应该购买这份保险吗？选择是要么确定地多花 140 欧元，要么以一定概率 \(p\) 承担风险，即如果你损坏了汽车，你最多要支付 1000 欧元。
\end{tcolorbox}

There are a number of further complications to this decision:

\begin{tcolorbox}
这个决策还有一些其他的复杂情况:
\end{tcolorbox}

\begin{itemize}
\item This is a one-off decision - you are not renting cars all the time. If you did, you could think for yourself whether during 50 days of car rental you are likely to create a damage of 1,000 Euros, because that is what you would pay for sure if you took out this insurance all the time.
\end{itemize}

\begin{tcolorbox}
\begin{itemize}
\item 这是一次性决策——你不是一直租车。如果你是，你可以自己想想，在 50 天的租车期间，你是否可能造成 1000 欧元的损坏，因为如果你一直购买这份保险，这就是你肯定要支付的金额。
\end{itemize}
\end{tcolorbox}

\begin{itemize}
\item If there is damage, maybe it is just a scratch that costs less than 1,000 Euros. This means that the risk is defined by multiple outcomes with different probabilities.
\end{itemize}

\begin{tcolorbox}
\begin{itemize}
\item 如果有损坏，也许只是一个划痕，费用不到 1000 欧元。这意味着风险是由具有不同概率的多个结果定义的。
\end{itemize}
\end{tcolorbox}

\begin{itemize}
\item Buying the insurance will give you peace of mind so you do not have to worry about damaging the car, which is an extra benefit. (On the other hand, would you drive less safely?)
\end{itemize}

\begin{tcolorbox}
\begin{itemize}
\item 购买保险会让你安心，这样你就不必担心损坏汽车，这是额外的好处。(另一方面，你开车会更不安全吗？)
\end{itemize}
\end{tcolorbox}

\begin{itemize}
\item Maybe the car rental itself costs only 200 Euros, which is is raised significantly by an extra 140 Euros. (The "independence axiom" says that this should not matter.)
\end{itemize}

\begin{tcolorbox}
\begin{itemize}
\item 也许租车本身只需要 200 欧元，而额外的 140 欧元会使费用大幅增加。(“独立性公理”说这应该无关紧要。)
\end{itemize}
\end{tcolorbox}

\begin{itemize}
\item If you decline the insurance and then have to pay for damage, you will regret your decision. Here the decision theorist will tell you: The risky choice means by definition that you do not know the future. If you took the risk, accept your decision that you did so.
\end{itemize}

\begin{tcolorbox}
\begin{itemize}
\item 如果你拒绝购买保险，然后不得不为损坏买单，你会后悔你的决定。在这里，决策理论家会告诉你:风险选择从定义上来说意味着你不知道未来。如果你冒了这个险，就接受你这样做的决定。
\end{itemize}
\end{tcolorbox}

\begin{itemize}
\item Paying 1,000 Euros might put you in severe financial difficulties. Hopefully, it doesn't. However, this consideration may be a good guide in other situations, for example whether you should pay another \({80}\mathrm{{GBP}}\) for insuring your new 400-GBP refrigerator against "accidental breakdown" for three years after purchase. (Such insurances are popular in the UK but tend to be over-priced.)
\end{itemize}

\begin{tcolorbox}
\begin{itemize}
\item 支付 1000 欧元可能会让你陷入严重的财务困境。希望不会这样。然而，这种考虑在其他情况下可能是一个很好的指导，例如你是否应该再花 \({80}\mathrm{{GBP}}\) 为你新买的价值 400 英镑的冰箱购买三年的“意外故障”保险。(这种保险在英国很受欢迎，但往往定价过高。)
\end{itemize}
\end{tcolorbox}

\begin{itemize}
\item You do not know the probability \(p\) that your car gets a 1,000-Euro damage. Thinking of whether you will create such a damage during 50 days of driving a rental car may help estimating that probability. This also depends on how safely you can drive in the place that you are visiting.
\end{itemize}

\begin{tcolorbox}
\begin{itemize}
\item 你不知道你的汽车遭受 1000 欧元损坏的概率 \(p\) 。想想在 50 天的租车驾驶中你是否会造成这样的损坏可能有助于估计这个概率。这也取决于你在访问地的驾驶安全程度。
\end{itemize}
\end{tcolorbox}

Concerning the probability \(p\) , the expected cost if you do not buy insurance is \(1,{000} \cdot  p\) Euros, which exceeds 140 Euros if \(p > {0.14}\) . However, there is no rule that says "buy the insurance if \(p > {0.14}^{\prime \prime }\) because you may well prefer to buy the insurance if \(p = {0.1}\) , for example. (The actual \(p\) is probably much lower than suggested by the prices of the rental car companies - separate "car rental excess" insurers cover the excess for 60 Euros per year for multiple rentals, so (a) this is an indication that the risk is low, (b) this is the kind of insurance you might prefer to buy in advance.)

\begin{tcolorbox}
关于概率 \(p\) ，如果您不购买保险，预期成本为 \(1,{000} \cdot  p\) 欧元，若 \(p > {0.14}\) ，该成本将超过 140 欧元。然而，并没有“若 \(p > {0.14}^{\prime \prime }\) 就购买保险”的规则，例如，若 \(p = {0.1}\) ，您很可能更愿意购买保险。(实际的 \(p\) 可能远低于租车公司价格所暗示的水平——单独的“租车超额”保险公司每年收取 60 欧元，为多次租车提供超额保障，所以 (a) 这表明风险较低，(b) 这是一种您可能会提前购买的保险。)
\end{tcolorbox}

Assume that in this car rental problem, you have an estimate of the probability \(p\) that you will damage the car (perhaps with a rather high \(p\) to be on the safe side). Then the problem whether you should buy the extra insurance for 140 Euros is called a decision problem under risk, which is one of the problems of decision theory and which is the one we study.

\begin{tcolorbox}
假设在这个租车问题中，您对自己损坏汽车的概率 \(p\) 有一个估计(为了安全起见， \(p\) 可能估计得相当高)。那么，您是否应该花 140 欧元购买额外保险的问题就被称为风险决策问题，这是决策理论中的问题之一，也是我们要研究的问题。
\end{tcolorbox}

Decision theory deals with situations that are distinguished in the following respects:

\begin{tcolorbox}
决策理论处理的情况在以下方面有所不同:
\end{tcolorbox}

\begin{itemize}
\item decisions of an individual or of a group;
\end{itemize}

\begin{tcolorbox}
\begin{itemize}
\item 个人决策或群体决策；
\end{itemize}
\end{tcolorbox}

\begin{itemize}
\item decisions under certainty, under risk, or under uncertainty;
\end{itemize}

\begin{tcolorbox}
\begin{itemize}
\item 确定性决策、风险决策或不确定性决策；
\end{itemize}
\end{tcolorbox}

\begin{itemize}
\item outcomes with single or multiple attributes.
\end{itemize}

\begin{tcolorbox}
\begin{itemize}
\item 单属性结果或多属性结果。
\end{itemize}
\end{tcolorbox}

The car rental problem has only a single decision maker. The problem of group decisions is to aggregate the preferences of multiple individuals. For example, there may be three candidates \(x,y,z\) for a job that is decided by a committee of three people \(A,B,C\) who have the following preferences:

\begin{tcolorbox}
租车问题只有一个决策者。群体决策问题是要汇总多个个体的偏好。例如，有三个候选人 \(x,y,z\) 竞争一份工作，由一个三人委员会 \(A,B,C\) 来决定，他们的偏好如下:
\end{tcolorbox}

\[
A : x \prec  y \prec  z
\]

\[
B : \;y \prec  z \prec  x \tag{5.3}
\]

\[
C : z \prec  x \prec  y
\]

which means that \(A\) prefers \(y\) over \(x\) and \(z\) over both \(x\) and \(y\) , and correspondingly for \(B\) and \(C\) . One possible decision rule could be a pairwise comparison by majority vote. In this case, \(A\) and \(B\) (the majority) prefer \(z\) over \(y\) , and \(A\) and \(C\) prefer \(y\) over \(x\) , and \(B\) and \(C\) prefer \(x\) over \(z\) . Hence, even though each committee member has a strict order of the candidates, this pairwise comparison yields a cycle and therefore no clear winner. This is known as the Condorcet Paradox and one of the problems of group decisions, as part of the area of Social Choice. In this chapter, we only consider an individual decision maker and his preferences, whom we still call player for brevity.

\begin{tcolorbox}
这意味着 \(A\) 对 \(y\) 的偏好高于 \(x\) ，对 \(z\) 的偏好高于 \(x\) 和 \(y\) ， \(B\) 和 \(C\) 也有相应的偏好。一种可能的决策规则是通过多数人规则进行两两比较。在这种情况下， \(A\) 和 \(B\) (多数人)对 \(z\) 的偏好高于 \(y\) ， \(A\) 和 \(C\) 对 \(y\) 的偏好高于 \(x\) ， \(B\) 和 \(C\) 对 \(x\) 的偏好高于 \(z\) 。因此，即使每个委员会成员对候选人都有严格序，但这种两两比较会产生循环，因此没有明确的获胜者。这就是孔多塞悖论(Condorcet Paradox)，是群体决策问题之一，属于社会选择(Social Choice)领域。在本章中，我们只考虑单个决策者及其偏好，为简洁起见，我们仍称其为参与者。
\end{tcolorbox}

In decisions under certainty, the outcome of the decision is known for sure. Decisions under risk deal with situations where the outcomes may be random but have probabilities that (as accurately as necessary) are known to the player. Decisions under uncertainty concern situations where the actual probability of an outcome is unknown. One approach of decision theory (which we will not cover) is to find conditions on preferences that amount to a player's (possibly implicitly defined) subjective probabilities for such outcomes that can be used like ordinary probabilities. We consider the situation of risk where probabilities are known.

\begin{tcolorbox}
在确定性决策中，决策的结果是确定已知的。风险决策处理的情况是，结果可能是随机的，但参与者知道(必要时尽可能准确地)这些结果的概率。不确定性决策涉及的情况是，结果的实际概率未知。决策理论的一种方法(我们不讨论)是找到偏好的条件，这些条件相当于参与者对这些结果的(可能是隐含定义的)主观概率，这些主观概率可以像普通概率一样使用。我们考虑已知概率的风险情况。
\end{tcolorbox}

Some decision problems involve only a single attribute like money, for which the preference is typically easy to decide when the outcome is certain (the preference is to pay less rather than more). Many decisions involve multiple attributes where the decision is no longer easy, such as whether going to a good but expensive restaurant or a less good but cheaper restaurant, with the two attributes of quality and price. The attribute structure of outcomes will not be relevant for us, but is useful for demonstrating complications that can arise with seemingly natural decision rules, as in the example (5.10) below.

\begin{tcolorbox}
一些决策问题只涉及单一属性，如金钱，当结果确定时，通常很容易确定偏好(偏好是支付更少而不是更多)。许多决策涉及多个属性，此时决策就不再容易了，例如是去一家高档但昂贵的餐厅，还是去一家稍差但便宜的餐厅，这里涉及质量和价格两个属性。结果的属性结构对我们来说并不重要，但它有助于说明看似自然的决策规则可能出现的复杂情况，如下例 (5.10) 所示。
\end{tcolorbox}

\section*{5.4 Preferences for Lotteries}

\begin{tcolorbox}
\section*{5.4 对彩票的偏好}
\end{tcolorbox}

We consider the following setup of a single player in a decision situation. A set \(X\) contains the possible outcomes, and the set \(\Delta \left( X\right)\) of (finite) lotteries over \(X\) is formally defined as follows.

\begin{tcolorbox}
我们考虑单个参与者在决策情况下的以下设定。集合 \(X\) 包含可能的结果， \(X\) 上的(有限)彩票集合 \(\Delta \left( X\right)\) 正式定义如下。
\end{tcolorbox}

Definition 5.1. Let \(X\) be a nonempty set of outcomes. A lottery over \(X\) is given by finitely many elements \({x}_{1},\ldots ,{x}_{n}\) of \(X\) (also called the outcomes of \(P\) ) with corresponding probabilities \({p}_{1},\ldots ,{p}_{n}\) , which are nonnegative numbers that sum to 1 . The set of lotteries over \(X\) is denoted by \(\Delta \left( X\right)\) .

\begin{tcolorbox}
定义 5.1。设 \(X\) 是一个非空的结果集合。 \(X\) 上的一个彩票由 \(X\) 中的有限个元素 \({x}_{1},\ldots ,{x}_{n}\) (也称为 \(P\) 的结果)以及相应的概率 \({p}_{1},\ldots ,{p}_{n}\) 给出，这些概率是非负实数，且它们的和为 1。 \(X\) 上的彩票集合记为 \(\Delta \left( X\right)\) 。
\end{tcolorbox}

If we compare two lotteries \(P\) and \(Q\) we can assume that they assign probabilities \({p}_{1},\ldots ,{p}_{n}\) and \({q}_{1},\ldots ,{q}_{n}\) to the same outcomes \({x}_{1},\ldots ,{x}_{n}\) , if necessary by letting some probabilities be zero if an outcome has positive probability for only one of the lotteries. This is simplest if \(X\) is itself finite, in which case we can take \(n = \left| X\right|\) , but we may also consider cases where \(X\) is infinite.

\begin{tcolorbox}
如果我们比较两个彩票 \(P\) 和 \(Q\) ，我们可以假设它们对相同的结果 \({x}_{1},\ldots ,{x}_{n}\) 分配概率 \({p}_{1},\ldots ,{p}_{n}\) 和 \({q}_{1},\ldots ,{q}_{n}\) ，如有必要，若某个结果仅在其中一个彩票中有正概率，则可令另一些概率为零。如果 \(X\) 本身是有限的，这是最简单的情况，此时我们可以取 \(n = \left| X\right|\) ，但我们也可以考虑 \(X\) 为无限的情况。
\end{tcolorbox}

We assume that the player has a preference for any two lotteries \(P\) and \(Q\) . If, when faced with the choice between \(P\) and \(Q\) , the player is willing to accept \(Q\) rather than \(P\) , we write this as

\begin{tcolorbox}
我们假设玩家对任意两个彩票 \(P\) 和 \(Q\) 有偏好。如果当面对 \(P\) 和 \(Q\) 之间的选择时，玩家愿意接受 \(Q\) 而非 \(P\) ，我们将其记为
\end{tcolorbox}

\[
P \lesssim  Q
\]

and also say the player (weakly) prefers \(Q\) to \(P\) . We then define the indifference relation \(\sim\) by

\begin{tcolorbox}
并且也说玩家(弱)偏好 \(Q\) 甚于 \(P\) 。然后我们通过以下方式定义无差异关系 \(\sim\)
\end{tcolorbox}

\[
P \sim  Q\; \Leftrightarrow  \;P \precsim  Q\text{ and }Q \precsim  P \tag{5.4}
\]

where \(P \sim  Q\) means the player is indifferent between \(P\) and \(Q\) . The strict preference relation \(<\) is defined by

\begin{tcolorbox}
其中 \(P \sim  Q\) 表示玩家对 \(P\) 和 \(Q\) 无差异。严格偏好关系 \(<\) 定义为
\end{tcolorbox}

\[
P \prec  Q\; \Leftrightarrow  \;P \precsim  Q\text{ and }\operatorname{not}Q \precsim  P \tag{5.5}
\]

where \(P \prec  Q\) means the player strictly prefers \(Q\) to \(P\) ; it is clearly equivalent to

\begin{tcolorbox}
其中 \(P \prec  Q\) 表示玩家严格偏好 \(Q\) 甚于 \(P\) ；显然它等价于
\end{tcolorbox}

\[
P \prec  Q\; \Leftrightarrow  \;P \precsim  Q\text{ and }\operatorname{not}P \sim  Q.
\]

We will derive a number of conditions on the weak preference relation \(\precsim\) on \(\Delta \left( X\right)\) that makes this relation mathematically tractable for its use in decision theory and game theory. These will be restrictions on \(\precsim\) that may not be descriptive of how people actually behave in certain decision situations. Instead, these restrictions are often called normative, that is, they state how people should behave, for example in order to come to a conclusive decision. In that way, the theory may be used as a tool in a decision analysis.

\begin{tcolorbox}
我们将推导出 \(\Delta \left( X\right)\) 上的弱偏好关系 \(\precsim\) 的一些条件，这些条件使该关系在决策理论和博弈论中的应用在数学上易于处理。这些将是对 \(\precsim\) 的限制，它们可能无法描述人们在某些决策情况下的实际行为。相反，这些限制通常被称为规范的，也就是说，它们说明了人们应该如何行为，例如为了得出一个决定性的决策。这样，该理论可以用作决策分析的工具。
\end{tcolorbox}

The goal will be to find conditions for the existence of a function \(u : X \rightarrow  \mathbb{R}\) , which is called a utility function for the player, so that the player strictly prefers a lottery \(Q\) over lottery \(P\) if \(Q\) has higher expected utility under \(Q\) than under \(P\) , and is indifferent between two lotteries if they have the same expected utility.

\begin{tcolorbox}
目标是找到一个函数 \(u : X \rightarrow  \mathbb{R}\) 存在的条件，该函数被称为玩家的效用函数，使得如果彩票 \(Q\) 在 \(Q\) 下的期望效用高于在 \(P\) 下的期望效用，玩家会严格偏好彩票 \(Q\) 而非彩票 \(P\) ；如果两张彩票的期望效用相同，玩家对它们无差异。
\end{tcolorbox}

Expected utility is the expected value of the utility function for the given lottery. Suppose the lottery \(P\) has probabilities \({p}_{1},\ldots ,{p}_{n}\) for the outcomes \({x}_{1},\ldots ,{x}_{n}\) , respectively. Then the utility function \(u\) applied to these outcomes defines a random variable. In extension of (4.1), we denote its expectation by \(\mathbb{E}\left( {u,P}\right)\) , given by

\begin{tcolorbox}
期望效用是给定彩票的效用函数的期望值。假设彩票 \(P\) 对于结果 \({x}_{1},\ldots ,{x}_{n}\) 分别有概率 \({p}_{1},\ldots ,{p}_{n}\) 。那么应用于这些结果的效用函数 \(u\) 定义了一个随机变量。作为 (4.1) 的扩展，我们用 \(\mathbb{E}\left( {u,P}\right)\) 表示其期望值，表达式为
\end{tcolorbox}

\[
\mathbb{E}\left( {u,P}\right)  = {p}_{1}u\left( {x}_{1}\right)  + \cdots  + {p}_{n}u\left( {x}_{n}\right) . \tag{5.6}
\]

Then the utility function \(u\) represents the player’s preference \(\precsim\) on \(\Delta \left( X\right)\) if, as stated in (5.1), for any two lotteries \(P\) and \(Q\) in \(\Delta \left( X\right)\) we have \(P \precsim  Q\) if and only if \(\mathbb{E}\left( {u,P}\right)  \leq  \mathbb{E}\left( {u,Q}\right)\) .

\begin{tcolorbox}
那么，如果如 (5.1) 所述，对于 \(\Delta \left( X\right)\) 中的任意两张彩票 \(P\) 和 \(Q\) ，当且仅当 \(\mathbb{E}\left( {u,P}\right)  \leq  \mathbb{E}\left( {u,Q}\right)\) 时 \(P \precsim  Q\) ，则效用函数 \(u\) 表示玩家在 \(\Delta \left( X\right)\) 上的偏好 \(\precsim\) 。
\end{tcolorbox}

If a utility function exists, then it imposes a number of conditions on the preference \(\precsim\) of the player, which are necessary conditions for a utility function. A selected set of these conditions is also sufficient for the existence of a utility function. These conditions are considered by many as sufficiently weak and plausible to justify them as a normative assumptions about the player's preference.

\begin{tcolorbox}
如果存在一个效用函数，那么它会对玩家的偏好 \(\precsim\) 施加一些条件，这些条件是效用函数存在的必要条件。一组经过挑选的这些条件对于效用函数的存在也是充分的。许多人认为这些条件足够弱且合理，可以将它们作为关于玩家偏好的规范假设。
\end{tcolorbox}

A first necessary condition that follows from (5.1) is that \(\precsim\) is a total relation, because the order \(\leq\) on \(\mathbb{R}\) is total: For any two lotteries \(P\) and \(Q\) ,

\begin{tcolorbox}
从 (5.1) 得出的第一个必要条件是 \(\precsim\) 是一个全关系，因为 \(\mathbb{R}\) 上的序 \(\leq\) 是全序:对于任意两张彩票 \(P\) 和 \(Q\) ，
\end{tcolorbox}

\[
P \precsim  Q\text{ or }Q \precsim  P. \tag{5.7}
\]

In a sense, this seems already problematic because lotteries represent complex scenarios, with multiple outcomes and their probabilities, so a person may already face difficulties in deciding between one or the other. In fact, the representation (5.1) with a utility \(u\) function may help in such a decision, once \(u\) has been determined by lotteries that are not very complex. We will see that lotteries that have only two outcomes will suffice for that purpose.

\begin{tcolorbox}
从某种意义上说，这似乎已经存在问题，因为彩票代表着复杂的场景，有多个结果及其概率，所以一个人可能已经在决定选择这张或那张彩票时面临困难。事实上，一旦通过不太复杂的彩票确定了效用 \(u\) 函数，(5.1) 中的表示法可能有助于做出这样的决策。我们将看到，只有两个结果的彩票就足以达到这个目的。
\end{tcolorbox}

\section*{5.5 Ordinal Preferences for Decisions Under Certainty}

\begin{tcolorbox}
\section*{5.5 确定性条件下决策的序数偏好}
\end{tcolorbox}

As a first step, we ignore uncertainties and look only at lotteries that have only a single outcome \(x\) , which has necessarily probability 1 . Let \({\mathbf{1}}_{x}\) be the lottery in \(\Delta \left( X\right)\) that assigns probability 1 to \(x\) . Then clearly \(\mathbb{E}\left( {u,{\mathbf{1}}_{x}}\right)  = u\left( x\right)\) . If \(x\) and \(y\) are two outcomes, we write \(x \precsim  y\) if \({\mathbf{1}}_{x} \precsim  {\mathbf{1}}_{y}\) which is simply the player’s weak preference of \(y\) over \(x\) as a deterministic outcome (technically, the preference \(\lesssim\) is a relation on the set \(\Delta \left( X\right)\) of lotteries). Therefore, if \(u\) is an expected-utility function according to (5.1), then for any \(x,y \in  X\) ,

\begin{tcolorbox}
作为第一步，我们忽略不确定性，只考虑只有单一结果 \(x\) 的彩票，该结果的概率必然为 1。设 \({\mathbf{1}}_{x}\) 是 \(\Delta \left( X\right)\) 中赋予 \(x\) 概率为 1 的彩票。那么显然 \(\mathbb{E}\left( {u,{\mathbf{1}}_{x}}\right)  = u\left( x\right)\) 。如果 \(x\) 和 \(y\) 是两个结果，当 \({\mathbf{1}}_{x} \precsim  {\mathbf{1}}_{y}\) 时我们记为 \(x \precsim  y\) ，这仅仅是玩家对 \(y\) 相对于 \(x\) 作为确定性结果的弱偏好(从技术上讲，偏好 \(\lesssim\) 是彩票集合 \(\Delta \left( X\right)\) 上的一个关系)。因此，如果 \(u\) 是根据 (5.1) 定义的期望效用函数，那么对于任意 \(x,y \in  X\) ，
\end{tcolorbox}

\[
x \precsim  y\; \Leftrightarrow  \;u\left( x\right)  \leq  u\left( y\right) , \tag{5.8}
\]

that is, the utility function represents the player's preference over outcomes that occur with certainty.

\begin{tcolorbox}
即，效用函数表示玩家对确定发生的结果的偏好。
\end{tcolorbox}

In a number of decision situations, in particular when the outcome has only a single attribute, the preference for receiving the outcome for certain is clear. For example, if \(x\) is the amount of money saved, then \(u\left( x\right)\) is strictly increasing in \(x\) , or if \(x\) is money to be paid, then \(u\left( x\right)\) is strictly decreasing in \(x\) .

\begin{tcolorbox}
在许多决策场景中，特别是当结果只有单一属性时，对确定获得该结果的偏好是明确的。例如，如果 \(x\) 是节省的金额，那么 \(u\left( x\right)\) 关于 \(x\) 严格递增；或者如果 \(x\) 是要支付的金额，那么 \(u\left( x\right)\) 关于 \(x\) 严格递减。
\end{tcolorbox}

An immediate consequence of the representation (5.1) is that the preference relation \(\precsim\) is transitive,

\begin{tcolorbox}
表示法 (5.1) 的一个直接结果是偏好关系 \(\precsim\) 是传递的，
\end{tcolorbox}

\[
x \precsim  y\text{ and }y \precsim  z \Rightarrow  x \precsim  z. \tag{5.9}
\]

This is typically not problematic for single-attribute decisions, but already for two attributes some decision rules can lead to a non-transitive preference relation, as the following example demonstrates. Suppose there are three candidates \(x,y,z\) for a job. A candidate \(c\) is described by a pair of attributes \(\left( {{c}_{1},{c}_{2}}\right)\) , where the attribute "intelligence" is represented by an IQ score \({c}_{1}\) from a cognitive test, and the attribute "experience" by \({c}_{2}\) years of previous work in the profession. Their attributes are as follows:

\begin{tcolorbox}
对于单属性决策而言，这通常不会有问题，但对于双属性决策，一些决策规则就可能导致非传递的偏好关系，正如下面的例子所示。假设一份工作有三名候选人 \(x,y,z\) 。候选人 \(c\) 由一对属性 \(\left( {{c}_{1},{c}_{2}}\right)\) 描述，其中“智力”属性由认知测试的智商得分 \({c}_{1}\) 表示，“经验”属性由该职业的 \({c}_{2}\) 年工作经验表示。他们的属性如下:
\end{tcolorbox}

(5.10)

\begin{center}
\adjustbox{max width=\textwidth}{
\begin{tabular}{|c|c|c|}
\hline
candidate \(\left( {{c}_{1},{c}_{2}}\right)\) & IQ score \({c}_{1}\) & years experience \({c}_{2}\) \\
\cline{1-3}
\(x = \left( {{x}_{1},{x}_{2}}\right)\) & 121 & 5 \\
\cline{1-3}
\(y = \left( {{y}_{1},{y}_{2}}\right)\) & 124 & 4 \\
\cline{1-3}
\(z = \left( {{z}_{1},{z}_{2}}\right)\) & 128 & 2 \\
\cline{1-3}
\hline
\end{tabular}
}
\end{center}

\begin{tcolorbox}

\begin{center}
\adjustbox{max width=\textwidth}{
\begin{tabular}{|c|c|c|}
\hline
候选人 \(\left( {{c}_{1},{c}_{2}}\right)\) & 智商得分 \({c}_{1}\) & 年工作经验 \({c}_{2}\) \\
\cline{1-3}
\(x = \left( {{x}_{1},{x}_{2}}\right)\) & 121 & 5 \\
\cline{1-3}
\(y = \left( {{y}_{1},{y}_{2}}\right)\) & 124 & 4 \\
\cline{1-3}
\(z = \left( {{z}_{1},{z}_{2}}\right)\) & 128 & 2 \\
\cline{1-3}
\hline
\end{tabular}
}
\end{center}

\end{tcolorbox}

The hiring manager considers intelligence as more important, and considers experience as relevant only for candidates with equal intelligence. This amounts to the lexicographic order \({ \leq  }_{\text{ lex }}\) where two attribute pairs \(\left( {{c}_{1},{c}_{2}}\right)\) and \(\left( {{d}_{1},{d}_{2}}\right)\) are compared according to

\begin{tcolorbox}
招聘经理认为智力更为重要，并且仅对智力相当的候选人考虑经验因素。这相当于字典序 \({ \leq  }_{\text{ lex }}\) ，其中两个属性对 \(\left( {{c}_{1},{c}_{2}}\right)\) 和 \(\left( {{d}_{1},{d}_{2}}\right)\) 按照
\end{tcolorbox}

\[
\left( {{c}_{1},{c}_{2}}\right) { \leq  }_{\text{ lex }}\left( {{d}_{1},{d}_{2}}\right) \; \Leftrightarrow  \;{c}_{1} \leq  {d}_{1}\text{ or }\left( {{c}_{1} = {d}_{1}\text{ and }{c}_{2} \leq  {d}_{2}}\right) . \tag{5.11}
\]

Both attributes are real numbers. Because \(\leq\) is a total order (see Definition 1.2) on \(\mathbb{R}\) , it is easy to see that \({ \leq  }_{\text{ lex }}\) is also a total order on \({\mathbb{R}}^{2}\) . If \({ \leq  }_{\text{ lex }}\) defines the preference relation, then \(z\) is the most preferred candidate in (5.10), because \({x}_{1} < {z}_{1}\) and \({y}_{1} < {z}_{1}\) .

\begin{tcolorbox}
两个属性均为实数。由于 \(\leq\) 是 \(\mathbb{R}\) 上的全序(参见定义 1.2)，很容易看出 \({ \leq  }_{\text{ lex }}\) 也是 \({\mathbb{R}}^{2}\) 上的全序。如果 \({ \leq  }_{\text{ lex }}\) 定义了偏好关系，那么 \(z\) 是(5.10)中最受偏好的候选人，因为 \({x}_{1} < {z}_{1}\) 且 \({y}_{1} < {z}_{1}\) 。
\end{tcolorbox}

However, the lexicographic order \({ \leq  }_{\text{ lex }}\) in (5.11) means that whenever \({c}_{1} < {d}_{1}\) then \(\left( {{c}_{1},{c}_{2}}\right)\) is always strictly preferred to \(\left( {{d}_{1},{d}_{2}}\right)\) no matter how close \({c}_{1}\) and \({d}_{1}\) are, even if \({d}_{2}\) is vastly better than \({c}_{2}\) ; the order between \({c}_{2}\) and \({d}_{2}\) matters only when \({c}_{1} = {d}_{1}\) .

\begin{tcolorbox}
然而，(5.11)中的字典序 \({ \leq  }_{\text{ lex }}\) 意味着，只要 \({c}_{1} < {d}_{1}\) ，那么无论 \({c}_{1}\) 和 \({d}_{1}\) 多么接近， \(\left( {{c}_{1},{c}_{2}}\right)\) 总是严格优于 \(\left( {{d}_{1},{d}_{2}}\right)\) ，即使 \({d}_{2}\) 比 \({c}_{2}\) 好得多；仅当 \({c}_{1} = {d}_{1}\) 时， \({c}_{2}\) 和 \({d}_{2}\) 之间的顺序才起作用。
\end{tcolorbox}

Suppose that, for that reason, the manager adopts the following rule: If the IQ score between two candidates differs by less than 5, then the candidate with higher experience is preferred. Then in (5.10) \(x\) is preferred to \(y\) because their IQ scores differ by only 3 and \(x\) has more experience; similarly, \(y\) is preferred to \(z\) . However, \(z\) is preferred to \(x\) because their IQ difference is high enough to count. The rule therefore leads to an intransitivity of the strict preference.

\begin{tcolorbox}
假设出于这个原因，经理采用以下规则:如果两名候选人的智商得分相差小于 5，那么经验更丰富的候选人更受偏好。那么在(5.10)中， \(x\) 比 \(y\) 更受偏好，因为他们的智商得分仅相差 3，且 \(x\) 经验更丰富；类似地， \(y\) 比 \(z\) 更受偏好。然而， \(z\) 比 \(x\) 更受偏好，因为他们的智商差异足够大，可以产生影响。因此，该规则导致了严格偏好的非传递性。
\end{tcolorbox}

While the rule itself is plausible (for example because the IQ result is considered "noisy"), the intransitivity creates a cycle. A similar cycle resulted from the pairwise comparison rule in the Condorcet paradox (5.3) (where the three rankings by the committee members \(A,B,C\) could also be considered as three "attributes" of the candidates \(x,y,z\) ). However, a cyclic preference, for whatever reason it occurs, cannot result in a satisfactory choice of one of the outcomes in the cycle because there is always a more preferred outcome. If the decision is then made in some arbitrary way, the cycle represents in effect an indifference among the outcomes in the cycle. It is therefore useful to impose the transitivity rule (5.9) as one of the normative conditions on the preference relation \(\lesssim\) .

\begin{tcolorbox}
虽然该规则本身是合理的(例如，因为智商结果被认为“有误差”)，但非传递性会产生一个循环。类似的循环也出现在孔多塞悖论(5.3)的两两比较规则中(其中委员会成员的三个排名 \(A,B,C\) 也可以被视为候选人 \(x,y,z\) 的三个“属性”)。然而，无论出于何种原因出现的循环偏好，都无法在循环中的结果中做出令人满意的选择，因为总会有更受偏好的结果。如果随后以某种任意方式做出决策，那么这个循环实际上代表了循环中各结果之间的无差异。因此，将传递性规则(5.9)作为偏好关系 \(\lesssim\) 的规范条件之一是有用的。
\end{tcolorbox}

The rule "if the IQ scores differ by less than 5, go by experience" is intransitive because "differing by less than 5", which replaces the condition \({c}_{1} = {d}_{1}\) in (5.11), is not a transitive relation. One way out of this would be to round the IQ score to the nearest multiple of 5 and then apply the lexicographic rule (5.11). In (5.10), the IQ scores of \(x,y,z\) would then be rounded to120,125,130and create the (transitive) ranking \(x \prec  y \prec  z\) . However, this rounding process is sensitive to small changes: if \(z\) had an IQ score of 127, this would be rounded to 125 and considered equal to the rounded IQ score of \(y\) , with the preference switching to \(x \prec  z \prec  y\) . Rather than using the lexicographic rule, it is probably better to realize that there may be a tradeoff between the two attributes. That tradeoff that is better represented by weights in an additive representation, such as " 2 IQ points have the same weight as 1 year job experience", which is represented by a utility function such as \(u\left( {{c}_{1},{c}_{2}}\right)  = {c}_{1} + 2{c}_{2}\) . Such an additive utility function may be in practice a good approximation of a player's preference, and choosing its weights may help the player to clarify the necessary tradeoffs. (Exercise 5.1 relates to such an "additive" utility function.)

\begin{tcolorbox}
规则“如果智商分数相差小于5分，则依据经验来判断”是非传递的，因为取代了(5.11)中条件 \({c}_{1} = {d}_{1}\) 的“相差小于5分”并非传递关系。解决这个问题的一种方法是将智商分数四舍五入到最接近的5的倍数，然后应用字典序规则(5.11)。在(5.10)中， \(x,y,z\) 的智商分数将被四舍五入为120、125、130，并形成(传递的)排序 \(x \prec  y \prec  z\) 。然而，这种四舍五入过程对微小变化很敏感:如果 \(z\) 的智商分数为127，它将被四舍五入为125，并被视为与 \(y\) 四舍五入后的智商分数相等，偏好就会转向 \(x \prec  z \prec  y\) 。与其使用字典序规则，不如认识到这两个属性之间可能存在权衡。这种权衡可以用加法表示中的权重来更好地体现，例如“2分智商相当于1年工作经验”，这可以用像 \(u\left( {{c}_{1},{c}_{2}}\right)  = {c}_{1} + 2{c}_{2}\) 这样的效用函数来表示。这样的加法效用函数在实践中可能是对玩家偏好的一个很好的近似，选择其权重可能有助于玩家明确必要的权衡。(练习5.1与这种“加法”效用函数相关。)
\end{tcolorbox}

A utility function \(u\) that represents a preference for deterministic outcomes as in (5.8) has a lot of freedom for its shape. Namely, for the comparison of outcomes that occur with certainty, only the ordinal comparison between values of \(u\) matters, that is, whether \(u\left( x\right)\) is less than, greater than, or equal to \(u\left( y\right)\) . This implies that \(u\left( x\right)\) can also be replaced by a \(\phi \left( {u\left( x\right) }\right)\) for a strictly increasing function \(\phi\) (also called a transformation of the utility function). For example, if \(x\) is "personal wealth" and always \(x > 0\) , then the preference \(x \precsim  y\) can either be represented by the wealth itself, that is, if and only if \(x \leq  y\) , or, for example, by the logarithm of that wealth, namely \(\log x \leq  \log y\) , where the strictly increasing transformation is given by \(\phi  = \log\) . The represented preference (a higher personal wealth is strictly preferred) is exactly the same. A utility function for decisions under certainty is therefore also called an ordinal utility function.

\begin{tcolorbox}
如(5.8)中那样表示确定性结果偏好的效用函数 \(u\) 在形状上有很大的自由度。也就是说，对于确定会发生的结果的比较，只有 \(u\) 值之间的序数比较才重要，即 \(u\left( x\right)\) 是小于、大于还是等于 \(u\left( y\right)\) 。这意味着 \(u\left( x\right)\) 也可以被一个严格递增函数 \(\phi\) 的 \(\phi \left( {u\left( x\right) }\right)\) 所取代(也称为效用函数的变换)。例如，如果 \(x\) 是“个人财富”且总是 \(x > 0\) ，那么偏好 \(x \precsim  y\) 既可以用财富本身来表示，即当且仅当 \(x \leq  y\) 时，也可以用该财富的对数来表示，即 \(\log x \leq  \log y\) ，其中严格递增变换由 \(\phi  = \log\) 给出。所表示的偏好(严格偏好更高的个人财富)是完全相同的。因此，用于确定性决策的效用函数也被称为序数效用函数。
\end{tcolorbox}

\section*{5.6 Cardinal Utility Functions and Simple Lotteries}

\begin{tcolorbox}
\section*{5.6 基数效用函数与简单彩票}
\end{tcolorbox}

As soon as a utility function represents a preference for risky choices, its values cannot be transformed by an arbitrary strictly increasing transformation. Suppose a player can choose between the following two lotteries \(P\) and \(Q\) . Lottery \(P\) gives with probability \(\frac{1}{2}\) either 100 dollars or 400 dollars, whereas lottery \(Q\) gives 240 dollars for certain. If the utility for receiving \(x\) dollars is just \(x\) , then the expected utility for \(P\) is \(\frac{1}{2}{100} + \frac{1}{2}{400} = {250}\) whereas for \(Q\) it is 240, so the player strictly prefers \(P\) to \(Q\) . However, if the utility for receiving \(x\) dollars is \(\sqrt{x}\) , then the expected utility for \(P\) is \(\frac{1}{2}\sqrt{100} + \frac{1}{2}\sqrt{400} = \frac{1}{2}{10} + \frac{1}{2}{20} = {15}\) whereas for \(Q\) it is \(\sqrt{240} \approx  {15.5}\) , so the player strictly prefers \(Q\) to \(P\) . There is no particular reason to assume a special utility function like \(u\left( x\right)  = \sqrt{x}\) , but the player should be able to express a preference between the lotteries \(P\) and \(Q\) , which should be either \(Q \prec  P\) , or \(P \prec  Q\) , or the indifference \(P \sim  Q\) . That is, an expected-utility function that represents a preference for lotteries is no longer just ordinal but cardinal.

\begin{tcolorbox}
一旦效用函数表示对风险选择的偏好，其值就不能通过任意的严格递增变换来转换。假设一个玩家可以在以下两种彩票 \(P\) 和 \(Q\) 之间进行选择。彩票 \(P\) 以概率 \(\frac{1}{2}\) 给出 100 美元或 400 美元，而彩票 \(Q\) 确定给出 240 美元。如果获得 \(x\) 美元的效用恰好是 \(x\) ，那么彩票 \(P\) 的期望效用是 \(\frac{1}{2}{100} + \frac{1}{2}{400} = {250}\) ，而彩票 \(Q\) 的期望效用是 240，所以玩家严格偏好 \(P\) 而非 \(Q\) 。然而，如果获得 \(x\) 美元的效用是 \(\sqrt{x}\) ，那么彩票 \(P\) 的期望效用是 \(\frac{1}{2}\sqrt{100} + \frac{1}{2}\sqrt{400} = \frac{1}{2}{10} + \frac{1}{2}{20} = {15}\) ，而彩票 \(Q\) 的期望效用是 \(\sqrt{240} \approx  {15.5}\) ，所以玩家严格偏好 \(Q\) 而非 \(P\) 。没有特别的理由去假设像 \(u\left( x\right)  = \sqrt{x}\) 这样的特殊效用函数，但玩家应该能够表达对彩票 \(P\) 和 \(Q\) 的偏好，这种偏好应该是 \(Q \prec  P\) ，或者 \(P \prec  Q\) ，或者是无差异 \(P \sim  Q\) 。也就是说，代表对彩票偏好的期望效用函数不再仅仅是序数的，而是基数的。
\end{tcolorbox}

By definition, a cardinal utility function is unique up to positive affine transformations. That is, if \(u : X \rightarrow  \mathbb{R}\) is a utility function, then the function \(v : X \rightarrow  \mathbb{R}\) defined by

\begin{tcolorbox}
根据定义，基数效用函数在正仿射变换下是唯一的。也就是说，如果 \(u : X \rightarrow  \mathbb{R}\) 是一个效用函数，那么由下式定义的函数 \(v : X \rightarrow  \mathbb{R}\)
\end{tcolorbox}

\[
v\left( x\right)  = a \cdot  u\left( x\right)  + b \tag{5.12}
\]

for real constants \(a\) and \(b\) where \(a > 0\) is also a utility function, and these are the only allowed transformations. More formally, call two functions \(u,v : X \rightarrow  \mathbb{R}\) scale-equivalent if (5.12) holds for some reals \(a\) and \(b\) with \(a > 0\) , for all \(x\) in \(X\) . This is clearly an equivalence relation because it is reflexive (take \(a = 1,b = 0\) ), symmetric because (5.12) is equivalent to \(u\left( x\right)  = \frac{1}{a}v\left( x\right)  - \frac{b}{a}\) , and transitive because if \(w\left( x\right)  = {cv}\left( x\right)  + d\) with \(c > 0\) then \(w\left( x\right)  = \left( {ca}\right) u\left( x\right)  + \left( {{cb} + d}\right)\) . We will show that any two expected-utility functions are scale-equivalent.

\begin{tcolorbox}
对于实常数 \(a\) 和 \(b\) ，其中 \(a > 0\) 也是一个效用函数，并且这些是唯一允许的变换。更正式地说，如果对于某些实数 \(a\) 和 \(b\) ，且 \(a > 0\) ，对于 \(X\) 中的所有 \(x\) ，(5.12) 式成立，那么称两个函数 \(u,v : X \rightarrow  \mathbb{R}\) 是尺度等价的。这显然是一种关系的等价性，因为它是自反的(取 \(a = 1,b = 0\) )，对称的，因为 (5.12) 式等价于 \(u\left( x\right)  = \frac{1}{a}v\left( x\right)  - \frac{b}{a}\) ，并且是传递的，因为如果 \(w\left( x\right)  = {cv}\left( x\right)  + d\) 且 \(c > 0\) ，那么 \(w\left( x\right)  = \left( {ca}\right) u\left( x\right)  + \left( {{cb} + d}\right)\) 。我们将证明任意两个期望效用函数是尺度等价的。
\end{tcolorbox}

When replacing \(u\left( x\right)\) with \(v\left( x\right)  = {au}\left( x\right)  + b\) , a different factor \(a\) changes the unit of the utility function (as if changing a monetary payoff from dollars to cents), and the additive constant \(b\) changes its origin (where the utility function takes value zero). This change of origin in a positive-affine transformation is familiar from the temperature scales Celcius and Fahrenheit, which are related by the positive-affine transformation \(F = {1.8C} + {32}\) to convert a temperature of \(C\) degrees Celsius to \(F\) degrees Fahrenheit; for example, a temperature of 20 degrees Celsius equals 68 degrees Fahrenheit. The temperature conversion from Celsius to Kelvin keeps the unit and only changes the origin, where degrees in Kelvin are degrees in Celsius plus 273.15 (so that the physically lowest possible temperature of "absolute zero" is represented by zero degrees Kelvin).

\begin{tcolorbox}
当用 \(v\left( x\right)  = {au}\left( x\right)  + b\) 替换 \(u\left( x\right)\) 时，不同的因子 \(a\) 会改变效用函数的单位(就像将货币收益从美元转换为美分一样)，而加法常数 \(b\) 会改变其原点(即效用函数取值为零的点)。正仿射变换中这种原点的改变在摄氏度和华氏度这两种温度刻度中很常见，它们通过正仿射变换 \(F = {1.8C} + {32}\) 相关联，用于将 \(C\) 摄氏度的温度转换为 \(F\) 华氏度；例如，20 摄氏度等于 68 华氏度。从摄氏度到开尔文温度的转换保持单位不变，只改变原点，开尔文温度等于摄氏温度加上 273.15(这样，物理上可能的最低温度“绝对零度”就用 0 开尔文表示)。
\end{tcolorbox}

We first show that any expected-utility function \(u\) that fulfills (5.1) can be replaced by the scale-equivalent function \(v\) in (5.12), which we also write as \({au} + b\) , for any fixed \(a,b \in  \mathbb{R}\) with \(a > 0\) . Namely, by (5.6),

\begin{tcolorbox}
我们首先证明，任何满足(5.1)的期望效用函数 \(u\) 都可以在(5.12)中被尺度等价的函数 \(v\) 所替代，我们也将其写为 \({au} + b\) ，其中 \(a,b \in  \mathbb{R}\) 为任意固定值且 \(a > 0\) 。即，根据(5.6)，
\end{tcolorbox}

\[
\mathbb{E}\left( {v,P}\right)  = \mathbb{E}\left( {{au} + b,P}\right)
\]

\[
= {p}_{1}\left( {{au}\left( {x}_{1}\right)  + b}\right)  + \cdots  + {p}_{n}\left( {{au}\left( {x}_{n}\right)  + b}\right)
\]

\[
= a\left( {{p}_{1}u\left( {x}_{1}\right)  + \cdots  + {p}_{n}u\left( {x}_{n}\right) }\right)  + \left( {{p}_{1} + \cdots  + {p}_{n}}\right) b
\]

\[
= a\mathbb{E}\left( {u,P}\right)  + b \tag{5.13}
\]

so that for any lotteries \(P\) and \(Q\) we have the equivalent statements

\begin{tcolorbox}
因此，对于任意彩票 \(P\) 和 \(Q\) ，我们有等价的陈述
\end{tcolorbox}

\[
P \precsim  Q \Leftrightarrow  \mathbb{E}\left( {u,P}\right) \; \leq  \;\mathbb{E}\left( {u,Q}\right)
\]

\[
\Leftrightarrow  a\mathbb{E}\left( {u,P}\right)  + b \leq  a\mathbb{E}\left( {u,Q}\right)  + b
\]

\[
\Leftrightarrow  \mathbb{E}\left( {v,P}\right)  \leq  \mathbb{E}\left( {v,Q}\right)  \tag{5.14}
\]

which proves the claim. While (5.14) is immediate from the simple observation (5.13), it emphasizes that "utility" is an abstract scale to represent a player's preference with an origin that can be arbitrarily shifted (by adding \(b\) ) and whose units can be scaled by multiplication by some positive factor \(a\) . Despite its name, "utility" is not an intrinsic "usefulness" of an outcome, but instead represents this preference (which is why we use the more neutral term "payoff" for the games in the other chapters of this book).

\begin{tcolorbox}
这就证明了该断言。虽然(5.14)可直接从简单观察(5.13)得出，但它强调了“效用”是一种抽象尺度，用于表示玩家的偏好，其原点可以任意移动(通过加上 \(b\) )，并且其单位可以通过乘以某个正因子 \(a\) 来缩放。尽管名为“效用”，但它并非结果的内在“有用性”，而是代表这种偏好(这就是为什么我们在本书其他章节的博弈中使用更中性的术语“收益”)。
\end{tcolorbox}

Our second claim is that an expected-utility function only admits the positive-affine transformations in (5.12) in order to represent a given preference for lotteries. For that purpose, it will suffice to look at lotteries that have only one or two outcomes.

\begin{tcolorbox}
我们的第二个断言是，期望效用函数只有通过(5.12)中的正仿射变换才能表示对彩票的给定偏好。为此，只需考虑只有一个或两个结果的彩票即可。
\end{tcolorbox}

The following operation combines two existing lotteries. Suppose \(\alpha  \in  \left\lbrack  {0,1}\right\rbrack\) and \(P\) and \(Q\) are two lotteries over \(X\) . Then the combined lottery \(\left( {1 - \alpha }\right) P + {\alpha Q}\) is a new lottery that can be thought of a two-stage process as in the following picture, which we call a probability tree (a game tree that has only chance moves),

\begin{tcolorbox}
以下操作将两个现有的彩票组合起来。假设 \(\alpha  \in  \left\lbrack  {0,1}\right\rbrack\) 和 \(P\) 以及 \(Q\) 是关于 \(X\) 的两个彩票。那么组合彩票 \(\left( {1 - \alpha }\right) P + {\alpha Q}\) 是一个新的彩票，可以看作是一个两阶段过程，如下图所示，我们称之为概率树(一种只有机会行动的博弈树)，
\end{tcolorbox}

(5.15)

\begin{center}
\includegraphics[max width=0.2\textwidth]{images/bo_d2e5cmbef24c73b03tn0_128_701_351_236_244_0.jpg}
\end{center}
\hspace*{3em} 

where at the first stage \(P\) is chosen with probability \(1 - \alpha\) and \(Q\) with probability \(\alpha\) , and then the chosen lottery is played with its respective probabilities for outcomes (as always, different random chance moves are independent). If these outcomes are \({x}_{i}\) with probabilities \({p}_{i}\) under \(P\) and \({q}_{i}\) under \(Q\) for \(1 \leq  i \leq  n\) , then the combined lottery \(\left( {1 - \alpha }\right) P + {\alpha Q}\) can also be considered as a single-stage process that chooses \({x}_{i}\) with probability \(\left( {1 - \alpha }\right) {p}_{i} + \alpha {q}_{i}\) for \(1 \leq  i \leq  n\) , and is therefore a lottery in the usual sense.

\begin{tcolorbox}
在第一阶段，以概率 \(1 - \alpha\) 选择 \(P\) ，以概率 \(\alpha\) 选择 \(Q\) ，然后以各自的结果概率来进行所选的彩票(和往常一样，不同的随机机会行动是相互独立的)。如果这些结果是 \({x}_{i}\) ，在 \(P\) 下的概率为 \({p}_{i}\) ，在 \(Q\) 下的概率为 \({q}_{i}\) ，对于 \(1 \leq  i \leq  n\) ，那么组合彩票 \(\left( {1 - \alpha }\right) P + {\alpha Q}\) 也可以被视为一个单阶段过程，以概率 \(\left( {1 - \alpha }\right) {p}_{i} + \alpha {q}_{i}\) 选择 \({x}_{i}\) ，对于 \(1 \leq  i \leq  n\) ，因此它是通常意义上的彩票。
\end{tcolorbox}

We always write a combined lottery as \(\left( {1 - \alpha }\right) P + {\alpha Q}\) so that the probability \(\alpha\) chooses the second lottery \(Q\) , because for \(\alpha  \in  \left\lbrack  {0,1}\right\rbrack\) the case \(\alpha  = 0\) as the left endpoint of the interval \(\left\lbrack  {0,1}\right\rbrack\) naturally corresponds to the left branch in (5.15) which chooses \(P\) , and the case \(\alpha  = 1\) as the right endpoint of the interval \(\left\lbrack  {0,1}\right\rbrack\) corresponds to the right branch which chooses \(Q\) . This correspondence would be lost if (as it is often done) we wrote such a combination as \({\beta P} + \left( {1 - \beta }\right) Q\) for \(\beta  \in  \left\lbrack  {0,1}\right\rbrack\) (which with \(\beta  = 1 - \alpha\) is of course mathematically equivalent).

\begin{tcolorbox}
我们总是将组合彩票写为 \(\left( {1 - \alpha }\right) P + {\alpha Q}\) ，这样概率 \(\alpha\) 选择第二个彩票 \(Q\) ，因为对于 \(\alpha  \in  \left\lbrack  {0,1}\right\rbrack\) ，情况 \(\alpha  = 0\) 作为区间 \(\left\lbrack  {0,1}\right\rbrack\) 的左端点自然对应于(5.15)中选择 \(P\) 的左分支，而情况 \(\alpha  = 1\) 作为区间 \(\left\lbrack  {0,1}\right\rbrack\) 的右端点对应于选择 \(Q\) 的右分支。如果(像通常那样)我们将这样的组合写为 \({\beta P} + \left( {1 - \beta }\right) Q\) 对于 \(\beta  \in  \left\lbrack  {0,1}\right\rbrack\) (当然，当 \(\beta  = 1 - \alpha\) 时在数学上是等价的)，这种对应关系就会丢失。
\end{tcolorbox}

A simple lottery is a lottery that chooses outcome \(x\) with probability \(1 - \alpha\) and outcome \(y\) with probability \(\alpha\) , which we denote by \(x{ \land  }_{\alpha }y\) (as in a little probability tree). Recall that the lottery that chooses the certain outcome \(x\) is written as \({\mathbf{1}}_{x}\) , so that by definition

\begin{tcolorbox}
简单彩票是一种以概率 \(1 - \alpha\) 选择结果 \(x\) ，以概率 \(\alpha\) 选择结果 \(y\) 的彩票，我们用 \(x{ \land  }_{\alpha }y\) 表示(就像一个小概率树)。回想一下，选择确定结果 \(x\) 的彩票写为 \({\mathbf{1}}_{x}\) ，因此根据定义
\end{tcolorbox}

\[
x{ \land  }_{\alpha }y = \left( {1 - \alpha }\right) {\mathbf{1}}_{x} + \alpha {\mathbf{1}}_{y}. \tag{5.16}
\]

Note that the simple lottery \(x{ \land  }_{\alpha }y\) is a lottery with one or two outcomes (it has only one outcome if \(x = y\) or if \(\alpha  = 0\) or \(\alpha  = 1\) ). Clearly,

\begin{tcolorbox}
注意，简单彩票 \(x{ \land  }_{\alpha }y\) 是一种有一个或两个结果的彩票(如果 \(x = y\) 或者 \(\alpha  = 0\) 或者 \(\alpha  = 1\) ，则它只有一个结果)。显然，
\end{tcolorbox}

\[
\mathbb{E}\left( {u,x{ \land  }_{\alpha }y}\right)  = \left( {1 - \alpha }\right) u\left( x\right)  + {\alpha u}\left( y\right) . \tag{5.17}
\]

In the following proposition, we assume that the player's preference for simple lotteries is represented by their expected value of a function \(u : X \rightarrow  \mathbb{R}\) as in (5.1) and (5.17), which we call an "expected-utility function for simple lotteries". The proposition states that \(u\) is then unique up to scale-equivalence. Once the player’s preference \(\precsim\) for simple lotteries is determined in this way, his preference for general lotteries is automatically determined (provided it fulfills certain assumptions which we will consider later).

\begin{tcolorbox}
在下面的命题中，我们假设玩家对简单彩票的偏好由他们对函数 \(u : X \rightarrow  \mathbb{R}\) 的期望值来表示，如(5.1)和(5.17)所示，我们将其称为“简单彩票的期望效用函数”。该命题表明， \(u\) 在尺度等价意义下是唯一的。一旦玩家对简单彩票的偏好 \(\precsim\) 以这种方式确定，他对一般彩票的偏好就会自动确定(前提是它满足我们稍后会考虑的某些假设)。
\end{tcolorbox}

Proposition 5.2. Let \(S\) be a subset of \(\Delta \left( X\right)\) that includes all simple lotteries. Let \(u : X \rightarrow  \mathbb{R}\) be a function that fulfills (5.1) for all \(P,Q \in  S\) . Then if \(v\) is another such function, then \(u\) and \(v\) are scale-equivalent, that is,(5.12) holds for some \(a,b\) with \(a > 0\) for all \(x \in  X\) .

\begin{tcolorbox}
命题 5.2。设 \(S\) 是 \(\Delta \left( X\right)\) 的一个包含所有简单彩票的子集。设 \(u : X \rightarrow  \mathbb{R}\) 是一个对所有 \(P,Q \in  S\) 都满足(5.1)的函数。那么，如果 \(v\) 是另一个这样的函数，则 \(u\) 和 \(v\) 是尺度等价的，即，对于某些 \(a,b\) 且对于所有 \(x \in  X\) 有 \(a > 0\) ，(5.12)成立。
\end{tcolorbox}

Proof. In this proof, let "utility function" be a function \(u\) that fulfills (5.1) for all simple lotteries \(P,Q \in  S\) . Consider a given utility function \(\widehat{u} : X \rightarrow  \mathbb{R}\) . If the decision maker is indifferent between all outcomes, and hence between all lotteries, then any utility function is constant and all utility functions are scale-equivalent (by adding suitable constants). Hence, we can assume there are two outcomes \({x}_{0}\) and \({x}_{1}\) with \({x}_{0} \prec  {x}_{1}\) so that \(\widehat{u}\left( {x}_{0}\right)  < \widehat{u}\left( {x}_{1}\right)\) . Define the function \(u : X \rightarrow  \mathbb{R}\) by

\begin{tcolorbox}
证明。在这个证明中，设“效用函数”是一个对所有简单彩票 \(P,Q \in  S\) 都满足(5.1)的函数 \(u\) 。考虑一个给定的效用函数 \(\widehat{u} : X \rightarrow  \mathbb{R}\) 。如果决策者对所有结果无差异，因此对所有彩票也无差异，那么任何效用函数都是常数，并且所有效用函数都是尺度等价的(通过添加合适的常数)。因此，我们可以假设存在两个结果 \({x}_{0}\) 和 \({x}_{1}\) ，使得 \({x}_{0} \prec  {x}_{1}\) ，从而 \(\widehat{u}\left( {x}_{0}\right)  < \widehat{u}\left( {x}_{1}\right)\) 。定义函数 \(u : X \rightarrow  \mathbb{R}\) 为
\end{tcolorbox}

\[
u\left( x\right)  = \frac{1}{\widehat{u}\left( {x}_{1}\right)  - \widehat{u}\left( {x}_{0}\right) }\widehat{u}\left( x\right)  - \frac{\widehat{u}\left( {x}_{0}\right) }{\widehat{u}\left( {x}_{1}\right)  - \widehat{u}\left( {x}_{0}\right) }
\]

which is a utility function that is scale-equivalent to \(\widehat{u}\) and which fulfills

\begin{tcolorbox}
这是一个与 \(\widehat{u}\) 尺度等价且满足
\end{tcolorbox}

\[
u\left( {x}_{0}\right)  = 0,\;u\left( {x}_{1}\right)  = 1. \tag{5.18}
\]

We show that any other utility function \(v : X \rightarrow  \mathbb{R}\) is scale-equivalent to \(u\) , via

\begin{tcolorbox}
我们证明任何其他效用函数 \(v : X \rightarrow  \mathbb{R}\) 都与 \(u\) 尺度等价，通过
\end{tcolorbox}

\[
v\left( x\right)  = {au}\left( x\right)  + b\;\text{ with }\;a = v\left( {x}_{1}\right)  - v\left( {x}_{0}\right) ,\;b = v\left( {x}_{0}\right) . \tag{5.19}
\]

Here, \(a\) and \(b\) are uniquely determined by (5.18) and

\begin{tcolorbox}
这里， \(a\) 和 \(b\) 由(5.18)唯一确定，并且
\end{tcolorbox}

\[
v\left( {x}_{0}\right)  = {au}\left( {x}_{0}\right)  + b,\;v\left( {x}_{1}\right)  = {au}\left( {x}_{1}\right)  + b. \tag{5.20}
\]

Consider any outcome \(z\) . Then one of the following three cases applies:

\begin{tcolorbox}
考虑任何结果 \(z\) 。那么以下三种情况之一适用:
\end{tcolorbox}

\[
z \prec  {x}_{0},\;{x}_{0} \precsim  z \precsim  {x}_{1},\;\text{ or }\;{x}_{1} \prec  z. \tag{5.21}
\]

We first consider the middle case \({x}_{0} \lesssim  z \lesssim  {x}_{1}\) in (5.21). Then \(0 = u\left( {x}_{0}\right)  \leq  u\left( z\right)  \leq\)  \(u\left( {x}_{1}\right)  = 1\) . Let \(\alpha  = u\left( z\right)\) , and consider the simple lottery \(P = {x}_{0}{ \land  }_{\alpha }{x}_{1}\) , so that by (5.17)

\begin{tcolorbox}
我们首先考虑(5.21)中的中间情况 \({x}_{0} \lesssim  z \lesssim  {x}_{1}\) 。那么 \(0 = u\left( {x}_{0}\right)  \leq  u\left( z\right)  \leq\)  \(u\left( {x}_{1}\right)  = 1\) 。设 \(\alpha  = u\left( z\right)\) ，并考虑简单彩票 \(P = {x}_{0}{ \land  }_{\alpha }{x}_{1}\) ，使得根据(5.17)
\end{tcolorbox}

\[
\mathbb{E}\left( {u,P}\right)  = \mathbb{E}\left( {u,{x}_{0}{ \land  }_{\alpha }{x}_{1}}\right)  = \left( {1 - \alpha }\right) u\left( {x}_{0}\right)  + {\alpha u}\left( {x}_{1}\right)  = \alpha  = u\left( z\right)
\]

and therefore \(z \sim  P\) . The utility function \(v\) represents the same indifference, so that by (5.20) and (5.18)

\begin{tcolorbox}
因此 \(z \sim  P\) 。效用函数 \(v\) 表示相同的无差异，因此根据(5.20)和(5.18)
\end{tcolorbox}

\[
v\left( z\right)  = \mathbb{E}\left( {v,P}\right)  = \mathbb{E}\left( {v,{x}_{0}{ \land  }_{\alpha }{x}_{1}}\right)
\]

\[
= \left( {1 - \alpha }\right) v\left( {x}_{0}\right)  + {\alpha v}\left( {x}_{1}\right)
\]

\[
= \left( {1 - \alpha }\right) \left( {{au}\left( {x}_{0}\right)  + b}\right)  + \alpha \left( {{au}\left( {x}_{1}\right)  + b}\right)
\]

\[
= \alpha  \cdot  a + b = {au}\left( z\right)  + b
\]

which shows (5.19) for \(x = z\) .

\begin{tcolorbox}
这表明对于 \(x = z\) (5.19)成立。
\end{tcolorbox}

Next, consider the case \(z \prec  {x}_{0}\) , which is similar but slightly more involved. Because \(z \prec  {x}_{0} \prec  {x}_{1}\) , we now choose a simple lottery \(Q = z{ \land  }_{\beta }{x}_{1}\) so that \(Q \sim  {x}_{0}\) , that is,

\begin{tcolorbox}
接下来，考虑情况 \(z \prec  {x}_{0}\) ，这种情况类似但稍微复杂一些。因为 \(z \prec  {x}_{0} \prec  {x}_{1}\) ，我们现在选择一个简单彩票 \(Q = z{ \land  }_{\beta }{x}_{1}\) ，使得 \(Q \sim  {x}_{0}\) ，即，
\end{tcolorbox}

\[
\left( {1 - \beta }\right) u\left( z\right)  + {\beta u}\left( {x}_{1}\right)  = u\left( {x}_{0}\right)
\]

which holds if \(\beta \left( {u\left( {x}_{1}\right)  - u\left( z\right) }\right)  = u\left( {x}_{0}\right)  - u\left( z\right)\) or \(\beta  = \frac{-u\left( z\right) }{1 - u\left( z\right) } = 1 - \frac{1}{1 - u\left( z\right) }\) where \(0 < \beta  < 1\) because \(u\left( z\right)  < u\left( {x}_{0}\right)  = 0\) . Furthermore,

\begin{tcolorbox}
如果 \(\beta \left( {u\left( {x}_{1}\right)  - u\left( z\right) }\right)  = u\left( {x}_{0}\right)  - u\left( z\right)\) 或 \(\beta  = \frac{-u\left( z\right) }{1 - u\left( z\right) } = 1 - \frac{1}{1 - u\left( z\right) }\) 成立，则该式成立，其中 \(0 < \beta  < 1\) ，因为 \(u\left( z\right)  < u\left( {x}_{0}\right)  = 0\) 。此外，
\end{tcolorbox}

\[
u\left( z\right)  = \frac{-\beta }{1 - \beta }. \tag{5.22}
\]

Because \(Q \sim  {x}_{0}\) , we have

\begin{tcolorbox}
因为 \(Q \sim  {x}_{0}\) ，我们有
\end{tcolorbox}

\[
\mathbb{E}\left( {v,Q}\right)  = v\left( {x}_{0}\right)
\]

\[
\left( {1 - \beta }\right) v\left( z\right)  + {\beta v}\left( {x}_{1}\right)  = v\left( {x}_{0}\right)
\]

\[
\left( {1 - \beta }\right) v\left( z\right)  + \beta \left( {{au}\left( {x}_{1}\right)  + b}\right)  = {au}\left( {x}_{0}\right)  + b
\]

\[
\left( {1 - \beta }\right) v\left( z\right)  + \beta \left( {a + b}\right)  = b
\]

\[
\left( {1 - \beta }\right) v\left( z\right)  =  - {\beta a} + \left( {1 - \beta }\right) b
\]

\[
v\left( z\right)  = \frac{-\beta }{1 - \beta }a + b = {au}\left( z\right)  + b \tag{5.23}
\]

which shows (5.19) for \(x = z\) if \(z \prec  {x}_{0}\) . The third case \({x}_{1} \prec  z\) in (5.21) is similar, by choosing a simple lottery \({x}_{0}{ \land  }_{\gamma }z\) so that the player is indifferent between that lottery and \({x}_{1}\) , and again using (5.20) (see Exercise 5.2). This proves the claim.

\begin{tcolorbox}
这表明，如果 \(z \prec  {x}_{0}\) ，则对于 \(x = z\) 有 (5.19) 式成立。(5.21) 中的第三种情况 \({x}_{1} \prec  z\) 是类似的，通过选择一个简单彩票 \({x}_{0}{ \land  }_{\gamma }z\) ，使得玩家对该彩票和 \({x}_{1}\) 无差异，然后再次使用 (5.20) 式(见练习 5.2)。这就证明了该命题。
\end{tcolorbox}

The proof of Proposition 5.2 highlights the meaning of the values of the utility function: Consider two outcomes \({x}_{0}\) and \({x}_{1}\) so that \({x}_{0} \prec  {x}_{1}\) . First of all, the utility function \(u\) can always be scaled so that, as in (5.18), the origin of the utility scale is \(u\left( {x}_{0}\right)  = 0\) , with \(u\left( {x}_{1}\right)  = u\left( {x}_{1}\right)  - u\left( {x}_{0}\right)  = 1\) as its unit. Second, consider some outcome \(z\) so that \({x}_{0} \prec  z \prec  {x}_{1}\) . This condition indicates that there is some probability \(\alpha\) so that the player is indifferent between \(z\) and the simple lottery \({x}_{0}{ \land  }_{\alpha }{x}_{1}\) . That is, we have the equivalent statements

\begin{tcolorbox}
命题 5.2 的证明凸显了效用函数(utility function)值的含义:考虑两个结果 \({x}_{0}\) 和 \({x}_{1}\) ，使得 \({x}_{0} \prec  {x}_{1}\) 。首先，效用函数 \(u\) 总是可以进行缩放，使得如 (5.18) 式所示，效用尺度(utility scale)的原点为 \(u\left( {x}_{0}\right)  = 0\) ，以 \(u\left( {x}_{1}\right)  = u\left( {x}_{1}\right)  - u\left( {x}_{0}\right)  = 1\) 为单位。其次，考虑某个结果 \(z\) ，使得 \({x}_{0} \prec  z \prec  {x}_{1}\) 。这个条件表明存在某个概率 \(\alpha\) ，使得玩家对 \(z\) 和简单彩票 \({x}_{0}{ \land  }_{\alpha }{x}_{1}\) 无差异。也就是说，我们有等价的表述
\end{tcolorbox}

\[
z \sim  {x}_{0}{ \land  }_{\alpha }{x}_{1},\;u\left( z\right)  = \left( {1 - \alpha }\right) u\left( {x}_{0}\right)  + {\alpha u}\left( {x}_{1}\right) ,
\]

and

\begin{tcolorbox}
和
\end{tcolorbox}

\[
\alpha  = \frac{u\left( z\right)  - u\left( {x}_{0}\right) }{u\left( {x}_{1}\right)  - u\left( {x}_{0}\right) } \tag{5.24}
\]

where by (5.18) \(\alpha  = u\left( z\right)\) . In other words, the utility \(u\left( z\right)\) for the outcome \(z\) that is between the less and more preferred outcomes \({x}_{0}\) and \({x}_{1}\) is just the probability for getting the more preferred outcome \({x}_{1}\) in the simple lottery \({x}_{0}{ \land  }_{\alpha }{x}_{1}\) , assuming the scaling (5.18). Without that scaling, \(\alpha\) is the more general quotient (5.24), which follows from the observation that for any real numbers \(A,B,C\) with \(A < C\) ,

\begin{tcolorbox}
其中根据 (5.18) 式有 \(\alpha  = u\left( z\right)\) 。换句话说，对于介于较不偏好和较偏好结果 \({x}_{0}\) 和 \({x}_{1}\) 之间的结果 \(z\) ，其效用 \(u\left( z\right)\) 就是在简单彩票 \({x}_{0}{ \land  }_{\alpha }{x}_{1}\) 中获得较偏好结果 \({x}_{1}\) 的概率，假设采用了 (5.18) 式的缩放方式。如果不进行这种缩放， \(\alpha\) 就是更一般的商 (5.24) ，这是基于以下观察得到的:对于任意实数 \(A,B,C\) ，且 \(A < C\) ，
\end{tcolorbox}

\[
A < B < C,\;B = \left( {1 - \alpha }\right) A + {\alpha C}\; \Leftrightarrow  \;\alpha  = \frac{B - A}{C - A},\;0 < \alpha  < 1. \tag{5.25}
\]

Taking two outcomes \({x}_{0}\) and \({x}_{1}\) and identifying the probability \(\alpha\) so that \(z \sim  {x}_{0}{ \land  }_{\alpha }{x}_{1}\) is therefore a possible way to construct the values \(u\left( z\right)\) of the utility function \(u\) for outcomes \(z\) that fulfill \({x}_{0} \prec  z \prec  {x}_{1}\) . If \({x}_{0}\) is least preferred and \({x}_{1}\) most preferred among all outcomes, then this determines the utility function uniquely (where if \(z \sim  {x}_{0}\) then \(u\left( z\right)  = u\left( {x}_{0}\right)\) and if \(z \sim  {x}_{1}\) then \(u\left( z\right)  = u\left( {x}_{1}\right)\) ). However, even if such extreme outcomes exist (for example if \(X\) is finite), in practice they might be not as useful for constructing the utility function because the player may be uncomfortable to consider simple lotteries between the worst and best possible outcomes. More realistic outcomes \({x}_{0}\) and \({x}_{1}\) are typically more useful instead. Then this leads to the other cases \(z \prec  {x}_{0}\) and \({x}_{1} \prec  z\) in (5.21), where one can construct other simple lotteries as in the proof of Proposition 5.2. The probability in that lottery is then used to find the unknown utility value \(u\left( z\right)\) using the general expression (5.24), but with permuted roles of the outcomes \(z,{x}_{0},{x}_{1}\) according to the player’s preference for \(z\) relative to \({x}_{0}\) and \({x}_{1}\) .

\begin{tcolorbox}
选取两个结果 \({x}_{0}\) 和 \({x}_{1}\) ，并确定概率 \(\alpha\) ，使得 \(z \sim  {x}_{0}{ \land  }_{\alpha }{x}_{1}\) 成为一种为满足 \({x}_{0} \prec  z \prec  {x}_{1}\) 的结果 \(z\) 构建效用函数 \(u\) 的值 \(u\left( z\right)\) 的可能方法。如果 \({x}_{0}\) 是所有结果中最不被偏好的，而 \({x}_{1}\) 是最被偏好的，那么这将唯一地确定效用函数(即如果 \(z \sim  {x}_{0}\) 则 \(u\left( z\right)  = u\left( {x}_{0}\right)\) ，如果 \(z \sim  {x}_{1}\) 则 \(u\left( z\right)  = u\left( {x}_{1}\right)\) )。然而，即使存在这样的极端结果(例如，如果 \(X\) 是有限的)，在实践中它们对于构建效用函数可能并不那么有用，因为参与者可能不愿意考虑在最差和最好的可能结果之间进行简单彩票选择。更现实的结果 \({x}_{0}\) 和 \({x}_{1}\) 通常更有用。这就引出了 (5.21) 中的其他情况 \(z \prec  {x}_{0}\) 和 \({x}_{1} \prec  z\) ，在这种情况下，可以像命题 5.2 的证明中那样构建其他简单彩票。然后使用该彩票中的概率，通过一般表达式 (5.24) 来找到未知的效用值 \(u\left( z\right)\) ，但要根据参与者对 \(z\) 相对于 \({x}_{0}\) 和 \({x}_{1}\) 的偏好来交换结果 \(z,{x}_{0},{x}_{1}\) 的角色。
\end{tcolorbox}

In summary, an expected-utility function is unique up to positive-affine transformations already if it represents the player's preference for simple lotteries, that is, lotteries that have two outcomes, with an arbitrary probability \(\alpha\) for the second outcome (and probability \(1 - \alpha\) for the first outcome). Finding a suitable \(\alpha\) so that the player is indifferent between this lottery and another outcome can be used to construct the utility function. The utility function is then useful for deciding the player's preference for more complicated lotteries.

\begin{tcolorbox}
总之，如果一个期望效用函数(expected-utility function)能表示参与者对简单彩票(即有两个结果，第二个结果的概率为任意 \(\alpha\) ，第一个结果的概率为 \(1 - \alpha\) 的彩票)的偏好，那么它在正仿射变换(positive-affine transformation)下是唯一的。找到一个合适的 \(\alpha\) ，使得参与者对该彩票和另一个结果无差异(indifference)，就可以用来构建效用函数(utility function)。然后，该效用函数可用于确定参与者对更复杂彩票的偏好(preference)。
\end{tcolorbox}

Next, we will consider only the preference between lotteries and certain consistency axioms for this preference. We then show that a preference relation that fulfills these axioms can be represented by an expected-utility function.

\begin{tcolorbox}
接下来，我们将只考虑彩票之间的偏好以及这种偏好的某些一致性公理。然后我们将证明，满足这些公理的偏好关系(preference relation)可以用一个期望效用函数来表示。
\end{tcolorbox}

\section*{5.7 Consistency Axioms}

\begin{tcolorbox}
\section*{5.7 一致性公理}
\end{tcolorbox}

Consider a preference relation \(\precsim\) on the set \(\Delta \left( X\right)\) of lotteries over the outcome set \(X\) , and its derived relations \(\sim\) and \(\prec\) in (5.4) and (5.5). In order to prove that \(\precsim\) can be represented by an expected-utility function according to (5.1), we consider a number of sufficient assumptions, called axioms. These are reasonable, normative conditions about how a player's preference should look like.

\begin{tcolorbox}
考虑结果集 \(X\) 上的彩票集 \(\Delta \left( X\right)\) 上的偏好关系 \(\precsim\) ，以及它在 (5.4) 和 (5.5) 中导出的关系 \(\sim\) 和 \(\prec\) 。为了证明 \(\precsim\) 可以根据 (5.1) 用一个期望效用函数来表示，我们考虑一些充分的假设，称为公理。这些是关于参与者的偏好应该是什么样的合理、规范(normative)条件。
\end{tcolorbox}

The first condition, already stated in (5.7) and (5.9), is that \(\precsim\) is total and transitive. As argued above, this is a useful normative assumption because otherwise the player would sometimes not be able to make up his mind; this may result in an arbitrary choice and factual indifference, which is typically not what is intended.

\begin{tcolorbox}
第一个条件，已在 (5.7) 和 (5.9) 中陈述，即 \(\precsim\) 是完全的且具有传递性(transitive)。如前所述，这是一个有用的规范假设，因为否则参与者有时将无法下定决心；这可能导致任意选择和事实上的无差异，而这通常不是预期的情况。
\end{tcolorbox}

The following conditions on \(\precsim\) affect the combination of lotteries as in (5.15). First of all, it is assumed that only the eventual probabilities for outcomes matter and not whether these probabilities result from a multi-stage process.

\begin{tcolorbox}
以下关于 \(\precsim\) 的条件会影响如 (5.15) 中彩票的组合。首先，假设只有结果的最终概率是重要的，而这些概率是否来自多阶段过程并不重要。
\end{tcolorbox}

Let \(P,Q,R\) be lotteries and \(\alpha\) be a probability. The following condition states that indifference between lotteries is preserved when they are substituted into other lotteries:

\begin{tcolorbox}
设 \(P,Q,R\) 为彩票， \(\alpha\) 为概率。以下条件表明，当彩票被代入其他彩票时，彩票之间的无差异关系得以保留:
\end{tcolorbox}

\[
P \sim  Q\; \Rightarrow  \;\left( {1 - \alpha }\right) R + {\alpha P} \sim  \left( {1 - \alpha }\right) R + {\alpha Q}. \tag{5.26}
\]

Here the player is indifferent between \(P\) and \(Q\) , and stays indifferent if these lotteries occur with some probability \(\alpha\) in a combined lottery that otherwise chooses some other lottery \(R\) . A similar condition states the same for strict preference, where \(\alpha  > 0\) is required:

\begin{tcolorbox}
这里，参与者对 \(P\) 和 \(Q\) 无差异，并且如果这些彩票以某个概率 \(\alpha\) 出现在一个组合彩票中，而该组合彩票在其他情况下选择另一个彩票 \(R\) ，参与者仍然保持无差异。类似的条件对于严格偏好(strict preference)也成立，其中需要 \(\alpha  > 0\) :
\end{tcolorbox}

\[
P \prec  Q,\;0 < \alpha  \leq  1\; \Rightarrow  \;\left( {1 - \alpha }\right) R + {\alpha P} \prec  \left( {1 - \alpha }\right) R + {\alpha Q}. \tag{5.27}
\]

The two conditions (5.26) and (5.27) are also known as independence of irrelevant alternatives or just independence axioms. The "irrelevant alternative" is the lottery \(R\) which should not affect the indifference \(P \sim  Q\) or the strict preference \(P \prec  Q\) .

\begin{tcolorbox}
条件(5.26)和(5.27)也被称为无关选项独立性(independence of irrelevant alternatives)，或简称为独立性公理(independence axioms)。“无关选项”是彩票 \(R\) ，它不应影响无差异(indifference)关系 \(P \sim  Q\) 或严格偏好(strict preference)关系 \(P \prec  Q\) 。
\end{tcolorbox}

\begin{center}
\includegraphics[max width=0.9\textwidth]{images/bo_d2e5cmbef24c73b03tn0_132_203_805_1233_377_0.jpg}
\end{center}
\hspace*{3em} 

Figure 5.2 The Allais paradox, which is an empirical violation of (5.27) for three hypothetical outcomes \(x,y,z\) of receiving 1,5, or 0 million dollars, \(P = y{ \land  }_{0.1}z\) and \(Q = {\mathbf{1}}_{x}\) , and where \(R\) is either \({\mathbf{1}}_{x}\) or \({\mathbf{1}}_{z}\) .

\begin{tcolorbox}
图5.2展示了阿莱悖论(Allais paradox)，这是对条件(5.27)的一个实证反例，涉及三种假设结果 \(x,y,z\) ，即获得100万美元、500万美元或0美元， \(P = y{ \land  }_{0.1}z\) 和 \(Q = {\mathbf{1}}_{x}\) ，其中 \(R\) 可以是 \({\mathbf{1}}_{x}\) 或 \({\mathbf{1}}_{z}\) 。
\end{tcolorbox}

Figure 5.2 shows an example, known as the Allais paradox, where in certain hypothetical situations people may violate condition (5.27). Consider three outcomes \(x,y,z\) where \(x\) means "receive one million dollars", \(y\) means "receive five million dollars", and \(z\) means "receive zero dollars", let \(P\) be the simple lottery \(y{ \land  }_{0.1}z\) and let \(Q = {\mathbf{1}}_{x}\) . The preference \(P \prec  Q\) means that the player would rather receive one million dollars for sure than five million dollars with 90 percent probability and otherwise nothing, which may be the player's choice. Consider now the lottery \(R = {\mathbf{1}}_{z}\) and let \(\alpha  = {0.1}\) . Then the combined lottery \(\left( {1 - \alpha }\right) R + {\alpha P}\) , shown on the right in Figure 5.2, means that the player receives nothing with 91 percent probability and five million dollars with 9 percent probability, and \(\left( {1 - \alpha }\right) R + {\alpha Q}\) that he receives nothing with 90 percent probability and one million dollars with 10 percent probability. Given that the player very likely receives nothing anyhow, his preference may now be reversed, that is, \(\left( {1 - \alpha }\right) R + {\alpha Q} \prec  \left( {1 - \alpha }\right) R + {\alpha P}\) , which contradicts condition (5.27). On the other hand, if \(R = {\mathbf{1}}_{x}\) , shown on the left in Figure 5.2, then \(\left( {1 - \alpha }\right) R + {\alpha P}\) means that the player receives one million dollars with 90 percent probability, nothing with 1 percent probability, and five million dollars with 9 percent probability, whereas \(\left( {1 - \alpha }\right) R + {\alpha Q}\) means that he receives one million with certainty, which may well preserve his original preference \(P \prec  Q\) also in combination with \(R\) , in agreement with (5.27).

\begin{tcolorbox}
图5.2给出了一个被称为阿莱悖论(Allais paradox)的例子，在某些假设情形下，人们可能会违背条件(5.27)。考虑三种结果 \(x,y,z\) ，其中 \(x\) 表示“获得100万美元”， \(y\) 表示“获得500万美元”， \(z\) 表示“获得0美元”。设 \(P\) 为简单彩票(simple lottery) \(y{ \land  }_{0.1}z\) ，并设 \(Q = {\mathbf{1}}_{x}\) 。偏好关系 \(P \prec  Q\) 意味着玩家宁愿确定地获得100万美元，也不愿以90\%的概率获得500万美元，否则一无所获，这可能是玩家的选择。现在考虑彩票 \(R = {\mathbf{1}}_{z}\) ，并设 \(\alpha  = {0.1}\) 。那么，图5.2右侧所示的组合彩票 \(\left( {1 - \alpha }\right) R + {\alpha P}\) 意味着玩家有91\%的概率一无所获，有9\%的概率获得500万美元；而 \(\left( {1 - \alpha }\right) R + {\alpha Q}\) 意味着玩家有90\%的概率一无所获，有10\%的概率获得100万美元。鉴于玩家很可能无论如何都一无所获，他的偏好现在可能会逆转，即 \(\left( {1 - \alpha }\right) R + {\alpha Q} \prec  \left( {1 - \alpha }\right) R + {\alpha P}\) ，这与条件(5.27)相矛盾。另一方面，如果是图5.2左侧所示的 \(R = {\mathbf{1}}_{x}\) ，那么 \(\left( {1 - \alpha }\right) R + {\alpha P}\) 意味着玩家有90\%的概率获得100万美元，有1\%的概率一无所获，有9\%的概率获得500万美元；而 \(\left( {1 - \alpha }\right) R + {\alpha Q}\) 意味着玩家确定地获得100万美元，这很可能使他在与 \(R\) 组合时也保持原来的偏好 \(P \prec  Q\) ，符合条件(5.27)。
\end{tcolorbox}

The suggested remedy for resolving the inconsistent preference in Figure 5.2, which violates (5.27), is that the player should (a) decide whether his preference is \(P \prec  Q\) or \(Q \prec  P\) (or \(P \sim  Q\) ), which ignores the irrelevant alternative \({\mathbf{1}}_{z}\) or \({\mathbf{1}}_{x}\) , and (b) consider how much weight he should give his regret when the underdesirable outcome \({\mathbf{1}}_{z}\) occurs in the lottery \(P = {\mathbf{1}}_{z}{ \land  }_{0.9}{\mathbf{1}}_{y}\) and whether it matters that this regret is "invisible" in the combined lottery \(\left( {1 - \alpha }\right) {\mathbf{1}}_{z} + {\alpha P}\) that he prefers in Figure 5.2. If the regret matters, it should be made part of the outcome \(z\) which then takes two forms "receiving zero dollars when I could have gotten one million", and "receiving zero dollars but that was the likely outcome anyhow". In summary, the inconsistency in Figure 5.2 may well be observed in practice, but the independence axiom (5.27) may help resolve it by highlighting that the basic decision should be between \(P\) and \(Q\) .

\begin{tcolorbox}
为解决图5.2中违背条件(5.27)的不一致偏好问题，建议的解决办法是，玩家应该:(a)确定自己的偏好是 \(P \prec  Q\) 还是 \(Q \prec  P\) (或者 \(P \sim  Q\) )，这需要忽略无关选项 \({\mathbf{1}}_{z}\) 或 \({\mathbf{1}}_{x}\) ；(b)考虑当彩票 \(P = {\mathbf{1}}_{z}{ \land  }_{0.9}{\mathbf{1}}_{y}\) 中出现不理想的结果 \({\mathbf{1}}_{z}\) 时，自己应给予后悔(regret)情绪多大的权重，以及在图5.2中他偏好的组合彩票 \(\left( {1 - \alpha }\right) {\mathbf{1}}_{z} + {\alpha P}\) 中这种后悔情绪“不可见”是否重要。如果后悔情绪很重要，那么它应该成为结果 \(z\) 的一部分，此时结果 \(z\) 有两种形式:“本可以获得100万美元却只得到0美元”和“得到0美元，但这本来就是很可能出现的结果”。总之，图5.2中的不一致情况在实践中很可能会出现，但独立性公理(5.27)可以通过强调基本决策应在 \(P\) 和 \(Q\) 之间进行，来帮助解决这一问题。
\end{tcolorbox}

One consequence of (5.27) is the following, rather natural, monotonicity property, which states that increasing the probability for a more preferred part \(Q\) of a combined lottery makes the combined lottery more preferable.

\begin{tcolorbox}
(5.27)的一个推论是以下相当自然的单调性性质，该性质表明，增加组合彩票中更受偏好部分 \(Q\) 的概率会使该组合彩票更受偏好。
\end{tcolorbox}

Lemma 5.3. Consider a preference relation \(\precsim\) on \(\Delta \left( X\right)\) that fulfills (5.27). Then for \(P,Q\) in \(\Delta \left( X\right)\) ,

\begin{tcolorbox}
引理5.3。考虑 \(\Delta \left( X\right)\) 上满足(5.27)的偏好关系 \(\precsim\) 。那么对于 \(\Delta \left( X\right)\) 中的 \(P,Q\) ，
\end{tcolorbox}

\[
P \prec  Q,\;0 \leq  \alpha  < \beta  \leq  1\; \Rightarrow  \;\left( {1 - \alpha }\right) P + {\alpha Q} \prec  \left( {1 - \beta }\right) P + {\beta Q}. \tag{5.28}
\]

Proof. Let \(P \prec  Q\) and \(0 \leq  \alpha  < \beta  \leq  1\) , which implies that \(0 \leq  \frac{\alpha }{\beta } < 1\) and therefore \(0 < 1 - \frac{\alpha }{\beta } \leq  1\) . The following probability trees convey the proof:

\begin{tcolorbox}
证明。设 \(P \prec  Q\) 和 \(0 \leq  \alpha  < \beta  \leq  1\) ，这意味着 \(0 \leq  \frac{\alpha }{\beta } < 1\) ，因此 \(0 < 1 - \frac{\alpha }{\beta } \leq  1\) 。以下概率树说明了该证明过程:
\end{tcolorbox}

\begin{center}
\includegraphics[max width=0.9\textwidth]{images/bo_d2e5cmbef24c73b03tn0_133_342_1365_1243_382_0.jpg}
\end{center}
\hspace*{3em} 

We use (5.27) twice, first to show that \(\frac{\alpha }{\beta }Q + \left( {1 - \frac{\alpha }{\beta }}\right) P \prec  \frac{\alpha }{\beta }Q + \left( {1 - \frac{\alpha }{\beta }}\right) Q\) , and then by substituting either combined lottery for \({Q}^{\prime }\) in \(\left( {1 - \beta }\right) P + \beta {Q}^{\prime }\) , which shows (5.28).

\begin{tcolorbox}
我们两次使用(5.27)，首先证明 \(\frac{\alpha }{\beta }Q + \left( {1 - \frac{\alpha }{\beta }}\right) P \prec  \frac{\alpha }{\beta }Q + \left( {1 - \frac{\alpha }{\beta }}\right) Q\) ，然后在 \(\left( {1 - \beta }\right) P + \beta {Q}^{\prime }\) 中用任一组合彩票替换 \({Q}^{\prime }\) ，这就证明了(5.28)。
\end{tcolorbox}

The next condition is called the continuity axiom, where \(P,Q,R\) are lotteries and \(\alpha\) is a probability:

\begin{tcolorbox}
下一个条件称为连续性公理(continuity axiom)，其中 \(P,Q,R\) 是彩票， \(\alpha\) 是一个概率:
\end{tcolorbox}

\[
P \prec  Q \prec  R\; \Rightarrow  \;\exists \alpha  \in  \left( {0,1}\right)  : Q \sim  \left( {1 - \alpha }\right) P + {\alpha R}. \tag{5.29}
\]

To see why condition (5.29) may be considered problematic, let \(P\) be "die in the next minute", \(Q\) be "receive nothing", and \(R\) "receive one dollar". Then \(Q \sim  \left( {1 - \alpha }\right) P + {\alpha R}\) means that the player would risk his life with positive probability \(1 - \alpha\) just to receive one dollar with probability \(\alpha\) , which seems unreasonable. The counterargument to this is that \(1 - \alpha\) can be arbitrarily small as long as it is positive, for example \({2}^{-{500}}\) as when 500 flips of a fair coin come out "heads" 500 times in a row. This is a minimal risk, and people do risk their lives with small probabilities for small gains, for example when cycling without a helmet. Condition (5.29) is a technical, not a practical restriction.

\begin{tcolorbox}
为了理解为什么条件(5.29)可能被认为有问题，设 \(P\) 为“在接下来的一分钟内死亡”， \(Q\) 为“一无所获”， \(R\) 为“获得一美元”。那么 \(Q \sim  \left( {1 - \alpha }\right) P + {\alpha R}\) 意味着玩家会以正概率 \(1 - \alpha\) 冒生命危险，只是为了以概率 \(\alpha\) 获得一美元，这似乎不合理。对此的反驳观点是，只要 \(1 - \alpha\) 为正，它可以任意小，例如 \({2}^{-{500}}\) ，就像一枚公平的硬币连续500次都掷出“正面”一样。这是极小的风险，而且人们确实会为了小收益而以小概率冒生命危险，例如骑自行车不戴头盔。条件(5.29)是一个技术性限制，而非实际限制。
\end{tcolorbox}

The continuity axiom (5.29) holds necessarily when the preference \(<\) is represented by an expected-utility function \(u\) , because \(P \prec  Q \prec  R\) implies \(\mathbb{E}\left( {u,P}\right)  < \mathbb{E}\left( {u,Q}\right)  < \mathbb{E}\left( {u,R}\right)\) , so that by (5.25) with \(\alpha  = \frac{\mathbb{E}\left( {u,Q}\right)  - \mathbb{E}\left( {u,P}\right) }{\mathbb{E}\left( {u,R}\right)  - \mathbb{E}\left( {u,P}\right) }\) we have \(0 < \alpha  < 1\) and

\begin{tcolorbox}
当偏好 \(<\) 由期望效用函数(expected-utility function) \(u\) 表示时，连续性公理(5.29)必然成立，因为 \(P \prec  Q \prec  R\) 意味着 \(\mathbb{E}\left( {u,P}\right)  < \mathbb{E}\left( {u,Q}\right)  < \mathbb{E}\left( {u,R}\right)\) ，所以根据(5.25)，当 \(\alpha  = \frac{\mathbb{E}\left( {u,Q}\right)  - \mathbb{E}\left( {u,P}\right) }{\mathbb{E}\left( {u,R}\right)  - \mathbb{E}\left( {u,P}\right) }\) 时，我们有 \(0 < \alpha  < 1\) 且
\end{tcolorbox}

\[
\mathbb{E}\left( {u,Q}\right)  = \left( {1 - \alpha }\right) \mathbb{E}\left( {u,P}\right)  + \alpha \mathbb{E}\left( {u,R}\right)  = \mathbb{E}\left( {u,\left( {1 - \alpha }\right) P + {\alpha R}}\right)  \tag{5.30}
\]

and therefore \(Q \sim  \left( {1 - \alpha }\right) P + {\alpha R}\) .

\begin{tcolorbox}
因此 \(Q \sim  \left( {1 - \alpha }\right) P + {\alpha R}\) 。
\end{tcolorbox}

Because a utility function \(u : X \rightarrow  \mathbb{R}\) takes real values, there are sets \(X\) with a preference relation \(\lesssim\) , already restricted to deterministic outcomes, that cannot be represented with a utility function according to (5.8). The classic example for this is \(X = {\left\lbrack  0,1\right\rbrack  }^{2}\) with \(\precsim\) as the lexicographic order \({ \leq  }_{\text{ lex }}\) defined in (5.11), illustrated in Figure 5.3.

\begin{tcolorbox}
因为效用函数(utility function) \(u : X \rightarrow  \mathbb{R}\) 取实数值，所以存在具有偏好关系 \(\lesssim\) 的集合 \(X\) (已限制为确定性结果)，根据(5.8)不能用效用函数来表示。经典的例子是 \(X = {\left\lbrack  0,1\right\rbrack  }^{2}\) 与 \(\precsim\) 构成的字典序(lexicographic order) \({ \leq  }_{\text{ lex }}\) ，如(5.11)所定义，如图5.3所示。
\end{tcolorbox}

\begin{center}
\includegraphics[max width=0.4\textwidth]{images/bo_d2e5cmbef24c73b03tn0_134_251_1168_599_470_0.jpg}
\end{center}
\hspace*{3em} 

Figure 5.3 The set \(X = {\left\lbrack  0,1\right\rbrack  }^{2}\) with \(\precsim\) as the lexicographic order \({ \leq  }_{\text{ lex }}\) . A utility function \(u\left( {{x}_{1},{x}_{2}}\right)\) would have to map the uncountably many vertical strips of \(X\) to nonempty disjoint intervals in \(\mathbb{R}\) , which is not possible.

\begin{tcolorbox}
图5.3 集合 \(X = {\left\lbrack  0,1\right\rbrack  }^{2}\) ，其字典序(lexicographic order)为 \({ \leq  }_{\text{ lex }}\) ，以 \(\precsim\) 表示。一个效用函数(utility function) \(u\left( {{x}_{1},{x}_{2}}\right)\) 必须将 \(X\) 中不可数的垂直条带映射到 \(\mathbb{R}\) 中的非空不相交区间，这是不可能的。
\end{tcolorbox}

Suppose there was a utility function \(u : {\left\lbrack  0,1\right\rbrack  }^{2} \rightarrow  \mathbb{R}\) so that, according to (5.11),

\begin{tcolorbox}
假设存在一个效用函数(utility function) \(u : {\left\lbrack  0,1\right\rbrack  }^{2} \rightarrow  \mathbb{R}\) ，使得根据式(5.11)，
\end{tcolorbox}

\[
\left( {{x}_{1},{x}_{2}}\right)  \prec  \left( {{y}_{1},{y}_{2}}\right)  \Leftrightarrow  \left( {{x}_{1},{x}_{2}}\right) { \leq  }_{\text{ lex }}\left( {{y}_{1},{y}_{2}}\right) \;\text{ and }\;\left( {{x}_{1},{x}_{2}}\right)  \neq  \left( {{y}_{1},{y}_{2}}\right)
\]

\[
\Leftrightarrow  {x}_{1} < {y}_{1}\;\text{ or }\;\left( {{x}_{1} = {y}_{1}\text{ and }{x}_{2} < {y}_{2}}\right)
\]

\[
\Leftrightarrow  u\left( {{x}_{1},{x}_{2}}\right)  < u\left( {{y}_{1},{y}_{2}}\right) \text{ . } \tag{5.31}
\]

Consider now a "vertical strip" of the form \(\left\{  {x}_{1}\right\}   \times  \left\lbrack  {0,1}\right\rbrack\) , a subset of \(X\) , of which three are shown as bold lines in Figure 5.3. For any two points \(\left( {{x}_{1},{x}_{2}}\right)\) and \(\left( {{x}_{1},{y}_{2}}\right)\) on such a strip so that \({x}_{2} < {y}_{2}\) (for example, the endpoints \(\left( {{x}_{1},0}\right)\) and \(\left( {{x}_{1},1}\right)\) ) we have \(u\left( {{x}_{1},{x}_{2}}\right)  < u\left( {{x}_{1},{y}_{2}}\right)\) . For two such strips \(\left\{  {x}_{1}\right\}   \times  \left\lbrack  {0,1}\right\rbrack\) and \(\left\{  {y}_{1}\right\}   \times  \left\lbrack  {0,1}\right\rbrack\) with \({x}_{1} < {y}_{1}\) the lexicographic order dictates \(u\left( {{x}_{1},1}\right)  < u\left( {{y}_{1},0}\right)\) . Therefore, the images under \(u\) of these strips must not overlap, and they are subsets of pairwise disjoint intervals in the range \(\mathbb{R}\) of \(u\) . This would work fine if there were only finitely many such strips (as when they have a certain width as drawn in the diagram), but mathematically they have no width, and therefore the images of these strips "cannot fit" disjointly into \(\mathbb{R}\) . Technically, for any \({x}_{1} \in  \left\lbrack  {0,1}\right\rbrack\) we can choose a rational number \(r\left( {x}_{1}\right)\) so that \(u\left( {{x}_{1},0}\right)  < r\left( {x}_{1}\right)  < u\left( {{x}_{1},1}\right)\) . If \({x}_{1} < {y}_{1}\) then \(u\left( {{x}_{1},1}\right)  < u\left( {{y}_{1},0}\right)\) and therefore \(r\left( {x}_{1}\right)  < r\left( {y}_{1}\right)\) . Hence, the mapping \({x}_{1} \mapsto  r\left( {x}_{1}\right)\) defines an injection from the set \(\left\lbrack  {0,1}\right\rbrack\) to the set \(\mathbb{Q}\) of rationals, which implies that \(\left\lbrack  {0,1}\right\rbrack\) has at most the set-theoretic cardinality of \(\mathbb{Q}\) . This is a contradiction because \(\mathbb{Q}\) is a countable set but \(\left\lbrack  {0,1}\right\rbrack\) is not. Hence, a real-valued utility function \(u : {\left\lbrack  0,1\right\rbrack  }^{2} \rightarrow  \mathbb{R}\) cannot represent the lexicographic order \({ \leq  }_{\text{ lex }}\) as in (5.31).

\begin{tcolorbox}
现在考虑形式为 \(\left\{  {x}_{1}\right\}   \times  \left\lbrack  {0,1}\right\rbrack\) 的“垂直条带”，它是 \(X\) 的一个子集，图5.3中用粗线展示了其中三条。对于这样一个条带上的任意两点 \(\left( {{x}_{1},{x}_{2}}\right)\) 和 \(\left( {{x}_{1},{y}_{2}}\right)\) ，使得 \({x}_{2} < {y}_{2}\) (例如，端点 \(\left( {{x}_{1},0}\right)\) 和 \(\left( {{x}_{1},1}\right)\) )，我们有 \(u\left( {{x}_{1},{x}_{2}}\right)  < u\left( {{x}_{1},{y}_{2}}\right)\) 。对于两个这样的条带 \(\left\{  {x}_{1}\right\}   \times  \left\lbrack  {0,1}\right\rbrack\) 和 \(\left\{  {y}_{1}\right\}   \times  \left\lbrack  {0,1}\right\rbrack\) ，且 \({x}_{1} < {y}_{1}\) ，字典序(lexicographic order)要求 \(u\left( {{x}_{1},1}\right)  < u\left( {{y}_{1},0}\right)\) 。因此，这些条带在 \(u\) 下的像不能重叠，它们是 \(u\) 的值域 \(\mathbb{R}\) 中两两不相交区间的子集。如果只有有限个这样的条带(就像图中画出的具有一定宽度的条带那样)，这是可行的，但从数学角度来看，它们没有宽度，因此这些条带的像无法不相交地“放入” \(\mathbb{R}\) 中。从技术上讲，对于任意 \({x}_{1} \in  \left\lbrack  {0,1}\right\rbrack\) ，我们可以选择一个有理数 \(r\left( {x}_{1}\right)\) ，使得 \(u\left( {{x}_{1},0}\right)  < r\left( {x}_{1}\right)  < u\left( {{x}_{1},1}\right)\) 。如果 \({x}_{1} < {y}_{1}\) ，那么 \(u\left( {{x}_{1},1}\right)  < u\left( {{y}_{1},0}\right)\) ，因此 \(r\left( {x}_{1}\right)  < r\left( {y}_{1}\right)\) 。因此，映射 \({x}_{1} \mapsto  r\left( {x}_{1}\right)\) 定义了一个从集合 \(\left\lbrack  {0,1}\right\rbrack\) 到有理数集 \(\mathbb{Q}\) 的单射，这意味着 \(\left\lbrack  {0,1}\right\rbrack\) 的集合论基数至多与 \(\mathbb{Q}\) 相同。这是一个矛盾，因为 \(\mathbb{Q}\) 是可数集，而 \(\left\lbrack  {0,1}\right\rbrack\) 不是。因此，一个实值效用函数(utility function) \(u : {\left\lbrack  0,1\right\rbrack  }^{2} \rightarrow  \mathbb{R}\) 不能像式(5.31)那样表示字典序(lexicographic order) \({ \leq  }_{\text{ lex }}\) 。
\end{tcolorbox}

The extreme "sensitivity" to changes in the first component is a mathematical reason not to use a lexicographic rule as a preference relation. For the same reason, the rule is also behaviorally questionable, as we discussed earlier after (5.11) and after (5.29) for the lexicographic preference of "safety to your life" over "monetary gain".

\begin{tcolorbox}
对第一个分量变化的极端“敏感性”是不将字典序规则用作偏好关系的一个数学原因。出于同样的原因，该规则在行为上也存在疑问，正如我们之前在(5.11)和(5.29)之后针对“生命安全”相对于“金钱收益”的字典序偏好所讨论的那样。
\end{tcolorbox}

As shown in (5.30), the existence of an expected-utility function implies the continuity axiom (5.29). If the preference \(\lesssim\) restricted to deterministic outcomes is the lexicographic order \({ \leq  }_{\text{ lex }}\) on \(X = {\left\lbrack  0,1\right\rbrack  }^{2}\) , then it violates the continuity axiom, and an expected-utility function does not exist. In the next section, we show how the continuity axiom (5.29) is used in the construction of an expected-utility function.

\begin{tcolorbox}
如(5.30)所示，期望效用函数的存在意味着连续性公理(5.29)。如果限制在确定性结果上的偏好 \(\lesssim\) 是 \(X = {\left\lbrack  0,1\right\rbrack  }^{2}\) 上的字典序 \({ \leq  }_{\text{ lex }}\) ，那么它违反了连续性公理，并且期望效用函数不存在。在下一节中，我们将展示连续性公理(5.29)如何用于构建期望效用函数。
\end{tcolorbox}

\section*{5.8 Existence of an Expected-Utility Function}

\begin{tcolorbox}
\section*{5.8 期望效用函数的存在性}
\end{tcolorbox}

The following is the central theorem of this chapter. The main idea of the proof has been given in Section 5.2. The proof given here is longer because it does not assume that there exists a least preferred outcome \({x}_{0}\) and most preferred outcome \({x}_{1}\) . Instead, any two outcomes \({x}_{0}\) and \({x}_{1}\) can be chosen to define the origin and unit of the utility function. Then the utility function can be defined for further outcomes \(z\) (for example, with \(z \prec  {x}_{0}\) ) as in Proposition 5.2. The expected-utility property is shown to hold for arbitrary lotteries \(P\) for a function \(v\) that is scale-equivalent to \(u\) and scaled with the least and most preferred outcomes of \(P\) (where we use that \(P\) has only finitely many outcomes).

\begin{tcolorbox}
以下是本章的核心定理。证明的主要思路已在5.2节给出。这里给出的证明更长，因为它不假设存在最不偏好的结果 \({x}_{0}\) 和最偏好的结果 \({x}_{1}\) 。相反，可以选择任意两个结果 \({x}_{0}\) 和 \({x}_{1}\) 来定义效用函数的原点和单位。然后，可以像命题5.2那样为更多的结果 \(z\) (例如，对于 \(z \prec  {x}_{0}\) )定义效用函数。对于与 \(u\) 尺度等价且以 \(P\) 的最不偏好和最偏好结果进行缩放的函数 \(v\) (其中我们利用 \(P\) 只有有限多个结果这一事实)，期望效用性质被证明对于任意彩票 \(P\) 都成立。
\end{tcolorbox}

Theorem 5.4. Let \(\Delta \left( X\right)\) be the set of lotteries on the set of outcomes \(X\) and let \(\precsim\) be a preference relation on \(\Delta \left( X\right)\) . Assume that \(\lesssim\) is total and transitive, and fulfills the independence axioms (5.26) and (5.27) and the continuity axiom (5.29). Then there exists an expected-utility function \(u : X \rightarrow  \mathbb{R}\) that represents \(\lesssim\) according to (5.1). Proof. The fact that \(\precsim\) is total and transitive implies that \(\sim\) and \(\prec\) as defined in (5.4) and (5.5) are also total and transitive, which we use all the time, for example when writing \({x}_{0} \precsim  z \precsim  {x}_{1}\) as short for \({x}_{0} \precsim  z\) and \(z \precsim  {x}_{1}\) . We will make repeated use of the continuity axiom. Note that if \(P \prec  Q \prec  R\) and \(Q \sim  \left( {1 - \alpha }\right) P + {\alpha R}\) as in (5.29), then \(\alpha\) is unique, because if \(\alpha  < \beta\) then \(\left( {1 - \alpha }\right) P + {\alpha R} \prec  \left( {1 - \beta }\right) P + {\beta R}\) by Lemma 5.3, so we cannot also have \(Q \sim  \left( {1 - \beta }\right) P + {\beta R}\) .

\begin{tcolorbox}
定理5.4。设 \(\Delta \left( X\right)\) 是结果集 \(X\) 上的彩票集合，设 \(\precsim\) 是 \(\Delta \left( X\right)\) 上的偏好关系。假设 \(\lesssim\) 是完全的且传递的，并且满足独立性公理(5.26)和(5.27)以及连续性公理(5.29)。那么存在一个期望效用函数 \(u : X \rightarrow  \mathbb{R}\) ，它根据(5.1)表示 \(\lesssim\) 。证明。 \(\precsim\) 是完全的且传递的这一事实意味着(5.4)和(5.5)中定义的 \(\sim\) 和 \(\prec\) 也是完全的且传递的，我们一直会用到这一点，例如当将 \({x}_{0} \precsim  z \precsim  {x}_{1}\) 简写为 \({x}_{0} \precsim  z\) 和 \(z \precsim  {x}_{1}\) 时。我们将反复使用连续性公理。注意，如果如(5.29)中的 \(P \prec  Q \prec  R\) 和 \(Q \sim  \left( {1 - \alpha }\right) P + {\alpha R}\) ，那么 \(\alpha\) 是唯一的，因为如果 \(\alpha  < \beta\) ，那么根据引理5.3有 \(\left( {1 - \alpha }\right) P + {\alpha R} \prec  \left( {1 - \beta }\right) P + {\beta R}\) ，所以我们不能同时有 \(Q \sim  \left( {1 - \beta }\right) P + {\beta R}\) 。
\end{tcolorbox}

If the player is indifferent between all outcomes, then he is also indifferent between all lotteries over outcomes, and a constant utility function suffices. Hence, assume there are two outcomes \({x}_{0}\) and \({x}_{1}\) so that \({x}_{0} \prec  {x}_{1}\) . As in the proof of Proposition 5.2, we use \({x}_{0}\) and \({x}_{1}\) to fix the origin and unit of the utility function \(u\) as in (5.18), that is, \(u\left( {x}_{0}\right)  = 0\) and \(u\left( {x}_{1}\right)  = 1\) . For any outcome \(z\) we then choose the value \(u\left( z\right)\) of the utility function with the help of a suitable simple lottery:

\begin{tcolorbox}
如果玩家对所有结果都无差异，那么他对所有基于结果的彩票也无差异，此时一个常数效用函数就足够了。因此，假设存在两个结果 \({x}_{0}\) 和 \({x}_{1}\) ，使得 \({x}_{0} \prec  {x}_{1}\) 。正如在命题 5.2 的证明中那样，我们使用 \({x}_{0}\) 和 \({x}_{1}\) 来确定效用函数 \(u\) 的原点和单位，如式 (5.18) 所示，即 \(u\left( {x}_{0}\right)  = 0\) 且 \(u\left( {x}_{1}\right)  = 1\) 。对于任何结果 \(z\) ，我们随后借助一个合适的简单彩票来选择效用函数的值 \(u\left( z\right)\) :
\end{tcolorbox}

\begin{itemize}
\item If \(z \sim  {x}_{0}\) then let \(u\left( z\right)  = 0\) , and if \(z \sim  {x}_{1}\) then let \(u\left( z\right)  = 1\) .
\end{itemize}

\begin{tcolorbox}
\begin{itemize}
\item 如果 \(z \sim  {x}_{0}\) ，则令 \(u\left( z\right)  = 0\) ；如果 \(z \sim  {x}_{1}\) ，则令 \(u\left( z\right)  = 1\) 。
\end{itemize}
\end{tcolorbox}

\begin{itemize}
\item If \({x}_{0} < z < {x}_{1}\) then find \(\alpha\) so that \(z \sim  \left( {1 - \alpha }\right) {\mathbf{1}}_{{x}_{0}} + \alpha {\mathbf{1}}_{{x}_{1}} = {x}_{0}{ \land  }_{\alpha }{x}_{1}\) according to (5.29) and let \(u\left( z\right)  = \alpha\) .
\end{itemize}

\begin{tcolorbox}
\begin{itemize}
\item 如果 \({x}_{0} < z < {x}_{1}\) ，则根据式 (5.29) 找到 \(\alpha\) 使得 \(z \sim  \left( {1 - \alpha }\right) {\mathbf{1}}_{{x}_{0}} + \alpha {\mathbf{1}}_{{x}_{1}} = {x}_{0}{ \land  }_{\alpha }{x}_{1}\) ，并令 \(u\left( z\right)  = \alpha\) 。
\end{itemize}
\end{tcolorbox}

\begin{itemize}
\item If \(z \prec  {x}_{0} \prec  {x}_{1}\) then find \(\beta\) so that \({x}_{0} \sim  z{ \land  }_{\beta }{x}_{1}\) , again by (5.29), and let \(u\left( z\right)  = \frac{-\beta }{1 - \beta }\) as in (5.22).
\end{itemize}

\begin{tcolorbox}
\begin{itemize}
\item 如果 \(z \prec  {x}_{0} \prec  {x}_{1}\) ，同样根据式 (5.29) 找到 \(\beta\) 使得 \({x}_{0} \sim  z{ \land  }_{\beta }{x}_{1}\) ，并如式 (5.22) 那样令 \(u\left( z\right)  = \frac{-\beta }{1 - \beta }\) 。
\end{itemize}
\end{tcolorbox}

\begin{itemize}
\item If \({x}_{0} \prec  {x}_{1} \prec  z\) then find \(\gamma\) so that \({x}_{1} \sim  {x}_{0}{ \land  }_{\gamma }z\) by (5.29) and let \(u\left( z\right)  = \frac{1}{\gamma }\) .
\end{itemize}

\begin{tcolorbox}
\begin{itemize}
\item 如果 \({x}_{0} \prec  {x}_{1} \prec  z\) ，根据式 (5.29) 找到 \(\gamma\) 使得 \({x}_{1} \sim  {x}_{0}{ \land  }_{\gamma }z\) ，并令 \(u\left( z\right)  = \frac{1}{\gamma }\) 。
\end{itemize}
\end{tcolorbox}

In all cases, \(u\left( z\right)\) is determined by the player’s indifference between the certain outcome and the simple lottery. For example, in the last case, \(u\left( {x}_{1}\right)  = 1 =\)  \(\mathbb{E}\left( {u,{x}_{0}{ \land  }_{\gamma }z}\right)  = \left( {1 - \gamma }\right) u\left( {x}_{0}\right)  + {\gamma u}\left( z\right)  = {\gamma u}\left( z\right)\) , so that \(u\left( z\right)  = \frac{1}{\gamma }\) .

\begin{tcolorbox}
在所有情况下， \(u\left( z\right)\) 由玩家对确定结果和简单彩票的无差异来确定。例如，在最后一种情况下， \(u\left( {x}_{1}\right)  = 1 =\)  \(\mathbb{E}\left( {u,{x}_{0}{ \land  }_{\gamma }z}\right)  = \left( {1 - \gamma }\right) u\left( {x}_{0}\right)  + {\gamma u}\left( z\right)  = {\gamma u}\left( z\right)\) ，从而 \(u\left( z\right)  = \frac{1}{\gamma }\) 。
\end{tcolorbox}

With \(u : X \rightarrow  \mathbb{R}\) determined in this way, we now show that \(u\) fulfills (5.1), that is, \(u\) represents the player’s preference between any two lotteries \(P\) and \(Q\) . Let \({z}_{i}\) for \(1 \leq  i \leq  n\) be the outcomes of \(P\) or \(Q\) where \({z}_{i}\) has probability \({p}_{i}\) under \(P\) and \({q}_{i}\) under \(Q\) . We first assume that all these outcomes fulfill \({x}_{0} \precsim  {z}_{i} \precsim  {x}_{1}\) . Let \({\alpha }_{i} = u\left( {z}_{i}\right)\) , where by the construction of \(u\) we have \({z}_{i} \sim  {x}_{0}{ \land  }_{{\alpha }_{i}}{x}_{1}\) . We now replace, by the successive substitution of equivalent lotteries according to (5.26), each outcome \({z}_{i}\) by its equivalent lottery \({x}_{0}{ \land  }_{{\alpha }_{i}}{x}_{1}\) . That is (see Figure 5.1),

\begin{tcolorbox}
以这种方式确定 \(u : X \rightarrow  \mathbb{R}\) 后，我们现在证明 \(u\) 满足 (5.1)，即 \(u\) 表示玩家在任意两个彩票 \(P\) 和 \(Q\) 之间的偏好。设 \({z}_{i}\) ( \(1 \leq  i \leq  n\) )为 \(P\) 或 \(Q\) 的结果，其中 \({z}_{i}\) 在 \(P\) 下的概率为 \({p}_{i}\) ，在 \(Q\) 下的概率为 \({q}_{i}\) 。我们首先假设所有这些结果都满足 \({x}_{0} \precsim  {z}_{i} \precsim  {x}_{1}\) 。设 \({\alpha }_{i} = u\left( {z}_{i}\right)\) ，根据 \(u\) 的构造，我们有 \({z}_{i} \sim  {x}_{0}{ \land  }_{{\alpha }_{i}}{x}_{1}\) 。现在，根据 (5.26) 依次用等价彩票替换，我们将每个结果 \({z}_{i}\) 替换为其等价彩票 \({x}_{0}{ \land  }_{{\alpha }_{i}}{x}_{1}\) 。即(见图 5.1)
\end{tcolorbox}

\[
P = \mathop{\sum }\limits_{{i = 1}}^{n}{p}_{i}{\mathbf{1}}_{{z}_{i}} \sim  \mathop{\sum }\limits_{{i = 1}}^{n}{p}_{i}\left( {{x}_{0}{ \land  }_{{\alpha }_{i}}{x}_{1}}\right)
\]

\[
= \mathop{\sum }\limits_{{i = 1}}^{n}{p}_{i}\left( {\left( {1 - {\alpha }_{i}}\right) {\mathbf{1}}_{{x}_{0}} + {\alpha }_{i}{\mathbf{1}}_{{x}_{1}}}\right)
\]

\[
= \left( {1 - \mathop{\sum }\limits_{{i = 1}}^{n}{p}_{i}{\alpha }_{i}}\right) {\mathbf{1}}_{{x}_{0}} + \left( {\mathop{\sum }\limits_{{i = 1}}^{n}{p}_{i}{\alpha }_{i}}\right) {\mathbf{1}}_{{x}_{1}}
\]

\[
= \left( {1 - p}\right) {\mathbf{1}}_{{x}_{0}} + p{\mathbf{1}}_{{x}_{1}} = {x}_{0}{ \land  }_{p}{x}_{1} \tag{5.32}
\]

with

\begin{tcolorbox}
与
\end{tcolorbox}

\[
p = \mathop{\sum }\limits_{{i = 1}}^{n}{p}_{i}{\alpha }_{i} = \mathop{\sum }\limits_{{i = 1}}^{n}{p}_{i}u\left( {z}_{i}\right)  = \mathbb{E}\left( {u,P}\right) . \tag{5.33}
\]

Similarly,

\begin{tcolorbox}
类似地
\end{tcolorbox}

\[
Q \sim  {x}_{0}{ \land  }_{q}{x}_{1}\;\text{ with }\;q = \mathop{\sum }\limits_{{i = 1}}^{n}{q}_{i}{\alpha }_{i} = \mathop{\sum }\limits_{{i = 1}}^{n}{q}_{i}u\left( {z}_{i}\right)  = \mathbb{E}\left( {u,Q}\right) . \tag{5.34}
\]

The derivation (5.32) is the core of the argument:

\begin{tcolorbox}
推导式 (5.32) 是该论证的核心:
\end{tcolorbox}

\begin{itemize}
\item By construction, the utility \(u\left( {z}_{i}\right)\) of the outcome \({z}_{i}\) is the probability \({\alpha }_{i}\) so that the player is indifferent between \({z}_{i}\) and the simple lottery \({x}_{0}{ \land  }_{{\alpha }_{i}}{x}_{1}\) that chooses the better outcome \({x}_{1}\) with probability \({\alpha }_{i}\) .
\end{itemize}

\begin{tcolorbox}
\begin{itemize}
\item 根据构造，结果 \({z}_{i}\) 的效用 \(u\left( {z}_{i}\right)\) 是概率 \({\alpha }_{i}\) ，使得玩家对 \({z}_{i}\) 和以概率 \({\alpha }_{i}\) 选择更好结果 \({x}_{1}\) 的简单彩票 \({x}_{0}{ \land  }_{{\alpha }_{i}}{x}_{1}\) 无差异。
\end{itemize}
\end{tcolorbox}

\begin{itemize}
\item Each outcome \({z}_{i}\) is replaced by its simple lottery \({x}_{0}{ \land  }_{{\alpha }_{i}}{x}_{1}\) in \(P\) , which by repeated application of (5.26) gives an equally preferred lottery.
\end{itemize}

\begin{tcolorbox}
\begin{itemize}
\item 在 \(P\) 中，每个结果 \({z}_{i}\) 都被其简单彩票 \({x}_{0}{ \land  }_{{\alpha }_{i}}{x}_{1}\) 替换，通过重复应用 (5.26) 得到一个同等偏好的彩票。
\end{itemize}
\end{tcolorbox}

\begin{itemize}
\item The resulting overall lottery has only the two outcomes \({x}_{0}\) and \({x}_{1}\) , with probability \(p = \mathop{\sum }\limits_{{i = 1}}^{n}{p}_{i}{\alpha }_{i}\) for \({x}_{1}\) because independent probabilities can be multiplied.
\end{itemize}

\begin{tcolorbox}
\begin{itemize}
\item 最终得到的总体彩票只有两个结果 \({x}_{0}\) 和 \({x}_{1}\) ， \({x}_{1}\) 的概率为 \(p = \mathop{\sum }\limits_{{i = 1}}^{n}{p}_{i}{\alpha }_{i}\) ，因为独立概率可以相乘。
\end{itemize}
\end{tcolorbox}

\begin{itemize}
\item The sum \(\mathop{\sum }\limits_{{i = 1}}^{n}{p}_{i}{\alpha }_{i}\) is equal to the expectation \(\mathbb{E}\left( {u,P}\right)\) by (5.33). The corresponding statement for \(Q\) is (5.34).
\end{itemize}

\begin{tcolorbox}
\begin{itemize}
\item 根据 (5.33)，和 \(\mathop{\sum }\limits_{{i = 1}}^{n}{p}_{i}{\alpha }_{i}\) 等于期望 \(\mathbb{E}\left( {u,P}\right)\) 。对于 \(Q\) 的相应表述是 (5.34)。
\end{itemize}
\end{tcolorbox}

With the help of the monotonicity Lemma 5.3, we show that \(P \precsim  Q\) holds if and only if \(\mathbb{E}\left( {u,P}\right)  \leq  \mathbb{E}\left( {u,Q}\right)\) , that is, if and only if \(p \leq  q\) . Suppose \(P \precsim  Q\) . If \(q < p\) then, because \({x}_{0} < {x}_{1},\left( {5.28}\right)\) (with \({x}_{0},{x}_{1},q,p\) for \(P,Q,\alpha ,\beta\) ) implies \(Q \sim  {x}_{0}{ \land  }_{q}{x}_{1} \prec  {x}_{0}{ \land  }_{p}{x}_{1} \sim  P\) , that is, \(Q \prec  P\) , a contradiction. This shows that \(P \precsim  Q\) implies \(p \leq  q\) . To show the converse, let \(p \leq  q\) . If \(p = q\) then clearly \(P \sim  Q\) , and if \(p < q\) then \(P \sim  {x}_{0}{ \land  }_{p}{x}_{1} \prec  {x}_{0}{ \land  }_{q}{x}_{1} \sim  Q\) , again by (5.28), that is, \(P \prec  Q\) . This shows that \(p \leq  q\) implies \(P \precsim  Q\) , as required.

\begin{tcolorbox}
借助单调性引理5.3，我们证明 \(P \precsim  Q\) 成立当且仅当 \(\mathbb{E}\left( {u,P}\right)  \leq  \mathbb{E}\left( {u,Q}\right)\) 成立，即当且仅当 \(p \leq  q\) 成立。假设 \(P \precsim  Q\) 成立。如果 \(q < p\) 成立，那么，因为 \({x}_{0} < {x}_{1},\left( {5.28}\right)\) (其中 \({x}_{0},{x}_{1},q,p\) 对应 \(P,Q,\alpha ,\beta\) )意味着 \(Q \sim  {x}_{0}{ \land  }_{q}{x}_{1} \prec  {x}_{0}{ \land  }_{p}{x}_{1} \sim  P\) 成立，也就是 \(Q \prec  P\) 成立，这产生了矛盾。这表明 \(P \precsim  Q\) 蕴含 \(p \leq  q\) 。为了证明其逆命题，设 \(p \leq  q\) 成立。如果 \(p = q\) 成立，那么显然 \(P \sim  Q\) 成立；如果 \(p < q\) 成立，那么根据(5.28)， \(P \sim  {x}_{0}{ \land  }_{p}{x}_{1} \prec  {x}_{0}{ \land  }_{q}{x}_{1} \sim  Q\) 成立，即 \(P \prec  Q\) 成立。这表明 \(p \leq  q\) 蕴含 \(P \precsim  Q\) ，符合要求。
\end{tcolorbox}

It remains to show (5.1) for lotteries \(P\) and \(Q\) that have outcomes \({z}_{1},\ldots ,{z}_{n}\) with \({z}_{i} \prec  {x}_{0}\) or \({x}_{1} \prec  {z}_{i}\) for some \(i\) . Let \({z}_{1}\) be the least preferred and \({z}_{n}\) be the most preferred of all outcomes \({z}_{1},\ldots ,{z}_{n}\) . We adopt the same approach as before, but use two new "reference outcomes" \({y}_{0}\) and \({y}_{1}\) instead of \({x}_{0}\) and \({x}_{1}\) :

\begin{tcolorbox}
接下来需要证明对于彩票 \(P\) 和 \(Q\) ，(5.1) 式成立，这些彩票的结果 \({z}_{1},\ldots ,{z}_{n}\) 满足对于某些 \(i\) 有 \({z}_{i} \prec  {x}_{0}\) 或者 \({x}_{1} \prec  {z}_{i}\) 。设 \({z}_{1}\) 是所有结果 \({z}_{1},\ldots ,{z}_{n}\) 中最不被偏好的， \({z}_{n}\) 是最被偏好的。我们采用与之前相同的方法，但使用两个新的“参考结果” \({y}_{0}\) 和 \({y}_{1}\) 来代替 \({x}_{0}\) 和 \({x}_{1}\) :
\end{tcolorbox}

\begin{itemize}
\item If \({z}_{1} \prec  {x}_{0}\) then let \({y}_{0} = {z}_{1}\) , otherwise let \({y}_{0} = {x}_{0}\) . If \({x}_{1} \prec  {z}_{n}\) then let \({y}_{1} = {z}_{n}\) , otherwise let \({y}_{1} = {x}_{1}\) . Then \({y}_{0} \precsim  {z}_{i} \precsim  {y}_{1}\) for all outcomes \({z}_{i}\) of the lotteries \(P\) and \(Q\) .
\end{itemize}

\begin{tcolorbox}
\begin{itemize}
\item 如果 \({z}_{1} \prec  {x}_{0}\) 成立，那么令 \({y}_{0} = {z}_{1}\) ；否则，令 \({y}_{0} = {x}_{0}\) 。如果 \({x}_{1} \prec  {z}_{n}\) 成立，那么令 \({y}_{1} = {z}_{n}\) ；否则，令 \({y}_{1} = {x}_{1}\) 。那么对于彩票 \(P\) 和 \(Q\) 的所有结果 \({z}_{i}\) ，有 \({y}_{0} \precsim  {z}_{i} \precsim  {y}_{1}\) 成立。
\end{itemize}
\end{tcolorbox}

\begin{itemize}
\item Define the function \(v : X \rightarrow  \mathbb{R}\) by
\end{itemize}

\begin{tcolorbox}
\begin{itemize}
\item 定义函数 \(v : X \rightarrow  \mathbb{R}\) 如下
\end{itemize}
\end{tcolorbox}

\[
v\left( x\right)  = \frac{u\left( x\right)  - u\left( {y}_{0}\right) }{u\left( {y}_{1}\right)  - u\left( {y}_{0}\right) },
\]

which is scale-equivalent to \(u\) and fulfills

\begin{tcolorbox}
该函数与 \(u\) 是尺度等价的(scale - equivalent)，并且满足
\end{tcolorbox}

\[
v\left( {y}_{0}\right)  = 0,\;v\left( {y}_{1}\right)  = 1.
\]

\begin{itemize}
\item Replace every outcome \({z}_{i}\) with an equally preferred simple lottery between \({y}_{0}\) and \({y}_{1}\) for some probability \({\beta }_{i}\) (which is possible because \({y}_{0} \precsim  {z}_{i} \precsim  {y}_{1}\) ),
\end{itemize}

\begin{tcolorbox}
\begin{itemize}
\item 用 \({y}_{0}\) 和 \({y}_{1}\) 之间的一个等偏好简单彩票(simple lottery)来替换每个结果 \({z}_{i}\) ，对于某个概率 \({\beta }_{i}\) (这是可行的，因为 \({y}_{0} \precsim  {z}_{i} \precsim  {y}_{1}\) )
\end{itemize}
\end{tcolorbox}

\[
{z}_{i} \sim  {y}_{0}{ \land  }_{{\beta }_{i}}{y}_{1}\;\left( {1 \leq  i \leq  n}\right) . \tag{5.35}
\]

\begin{itemize}
\item Substitute in \(P\) and \(Q\) the simple lotteries in (5.35) for the outcomes \({z}_{i}\) , using repeatedly (5.26). The player is then indifferent between \(P\) and a suitable simple lottery with the two outcomes \({y}_{0}\) and \({y}_{1}\) , and similarly for \(Q\) ,
\end{itemize}

\begin{tcolorbox}
\begin{itemize}
\item 在 \(P\) 和 \(Q\) 中，使用(5.26)式反复替换结果 \({z}_{i}\) 为(5.35)式中的简单彩票。然后，玩家对于 \(P\) 和一个具有两个结果 \({y}_{0}\) 和 \({y}_{1}\) 的合适简单彩票是无差异的，对于 \(Q\) 也是类似的。
\end{itemize}
\end{tcolorbox}

\[
P \sim  {y}_{0}{ \land  }_{\widehat{p}}{y}_{1},\;\widehat{p} = \mathop{\sum }\limits_{{i = 1}}^{n}{p}_{i}{\beta }_{i},\;Q \sim  {y}_{0}{ \land  }_{\widehat{q}}{y}_{1},\;\widehat{q} = \mathop{\sum }\limits_{{i = 1}}^{n}{q}_{i}{\beta }_{i}.
\]

Then we are done if we can show that

\begin{tcolorbox}
那么，如果我们能证明
\end{tcolorbox}

\[
{\beta }_{i} = v\left( {z}_{i}\right) \;\left( {0 \leq  i \leq  n}\right)  \tag{5.36}
\]

which implies that \(\widehat{p} = \mathbb{E}\left( {v,P}\right)\) and \(\widehat{q} = \mathbb{E}\left( {v,Q}\right)\) as in (5.33) and (5.34), where we argue as before to prove that \(P \precsim  Q\) if and only if \(\widehat{p} \leq  \widehat{q}\) . This shows that \(v\) is an expected-utility function, and so is the scale-equivalent function \(u\) .

\begin{tcolorbox}
这意味着如 (5.33) 和 (5.34) 中的 \(\widehat{p} = \mathbb{E}\left( {v,P}\right)\) 和 \(\widehat{q} = \mathbb{E}\left( {v,Q}\right)\) ，我们如前所述证明当且仅当 \(\widehat{p} \leq  \widehat{q}\) 时 \(P \precsim  Q\) 。这表明 \(v\) 是一个期望效用函数(expected-utility function)，尺度等价的函数 \(u\) 也是。
\end{tcolorbox}

The probabilities \({\beta }_{i}\) are uniquely determined by the indifference conditions (5.35), but we cannot assume that (5.36) holds because \(v\) , although scale-equivalent to \(u\) , is not yet proved to be an expected-utility function because we have not yet shown this for \(u\) (for the general case that \({z}_{i} \prec  {x}_{0}\) or \({x}_{1} \prec  {z}_{i}\) for some \(i\) ).

\begin{tcolorbox}
概率 \({\beta }_{i}\) 由无差异条件(indifference conditions)(5.35) 唯一确定，但我们不能假设 (5.36) 成立，因为 \(v\) 尽管与 \(u\) 尺度等价，但尚未被证明是一个期望效用函数，因为我们还没有针对 \(u\) 证明这一点(对于某些 \(i\) 有 \({z}_{i} \prec  {x}_{0}\) 或 \({x}_{1} \prec  {z}_{i}\) 的一般情况)。
\end{tcolorbox}

Hence, it remains to prove (5.36). We do this, using (5.26), with the help of suitable indifferences between simple lotteries that involve \({y}_{0},{x}_{0},{x}_{1},{y}_{1}\) and \({z}_{i}\) for \(0 \leq  i \leq  n\) , using the definition of \(u\left( {z}_{i}\right)\) . The computations are straightforward applications of affine transformations but somewhat lengthy.

\begin{tcolorbox}
因此，仍需证明 (5.36)。我们利用 (5.26)，借助涉及 \({y}_{0},{x}_{0},{x}_{1},{y}_{1}\) 和 \({z}_{i}\) (对于 \(0 \leq  i \leq  n\) )的简单彩票(simple lottery)之间的合适无差异关系来完成这一点，使用 \(u\left( {z}_{i}\right)\) 的定义。计算过程是仿射变换(affine transformation)的直接应用，但有些冗长。
\end{tcolorbox}

The goal is to represent all of \({z}_{1},\ldots ,{z}_{n},{x}_{0},{x}_{1}\) as simple lotteries between \({y}_{0}\) and \({y}_{1}\) . We first do this for \({x}_{0}\) and \({x}_{1}\) , with

\begin{tcolorbox}
目标是将所有的 \({z}_{1},\ldots ,{z}_{n},{x}_{0},{x}_{1}\) 表示为 \({y}_{0}\) 和 \({y}_{1}\) 之间的简单彩票(simple lottery)。我们首先对 \({x}_{0}\) 和 \({x}_{1}\) 进行此操作，其中
\end{tcolorbox}

\[
{x}_{0} \sim  {y}_{0}{ \land  }_{{\eta }_{0}}{y}_{1},\;{x}_{1} \sim  {y}_{0}{ \land  }_{{\eta }_{1}}{y}_{1}, \tag{5.37}
\]

where we claim that

\begin{tcolorbox}
我们声称
\end{tcolorbox}

\[
{\eta }_{0} = v\left( {x}_{0}\right)  = \frac{-u\left( {y}_{0}\right) }{u\left( {y}_{1}\right)  - u\left( {y}_{0}\right) },\;{\eta }_{1} = v\left( {x}_{1}\right)  = \frac{1 - u\left( {y}_{0}\right) }{u\left( {y}_{1}\right)  - u\left( {y}_{0}\right) }. \tag{5.38}
\]

To show (5.38), we have \({y}_{0} \precsim  {x}_{0} \prec  {x}_{1}\) and therefore

\begin{tcolorbox}
为了证明 (5.38)，我们有 \({y}_{0} \precsim  {x}_{0} \prec  {x}_{1}\) ，因此
\end{tcolorbox}

\[
{x}_{0} \sim  {y}_{0}{ \land  }_{\beta }{x}_{1},\;u\left( {y}_{0}\right)  = \frac{-\beta }{1 - \beta } \tag{5.39}
\]

according to the construction of \(u\left( {y}_{0}\right)\) (which applies also if \({y}_{0} = {x}_{0}\) where \(\beta  = 0\) ). With an obvious extension of the notation (5.16) to combinations of lotteries, we therefore have

\begin{tcolorbox}
根据 \(u\left( {y}_{0}\right)\) 的构造(如果 \({y}_{0} = {x}_{0}\) 也适用，其中 \(\beta  = 0\) )。因此，通过将符号表示(5.16)明显扩展到彩票(lottery)组合，我们有
\end{tcolorbox}

\[
{x}_{0} \sim  {y}_{0}{ \land  }_{\beta }{x}_{1} \sim  {y}_{0}{ \land  }_{\beta }\left( {{y}_{0}{ \land  }_{{\eta }_{1}}{y}_{1}}\right)  = {y}_{0}{ \land  }_{{\eta }_{0}}{y}_{1}
\]

which shows that

\begin{tcolorbox}
这表明
\end{tcolorbox}

\[
{\eta }_{0} = \beta {\eta }_{1}. \tag{5.40}
\]

Similarly, using that \({x}_{0} \prec  {x}_{1} \precsim  {y}_{1}\) and the construction of \(u\left( {y}_{1}\right)\) ,

\begin{tcolorbox}
同样，利用 \({x}_{0} \prec  {x}_{1} \precsim  {y}_{1}\) 和 \(u\left( {y}_{1}\right)\) 的构造
\end{tcolorbox}

\[
{x}_{1} \sim  {x}_{0}{ \land  }_{\gamma }{y}_{1},\;u\left( {y}_{1}\right)  = \frac{1}{\gamma } \tag{5.41}
\]

(also if \({y}_{1} = {x}_{1}\) where \(\gamma  = 1\) ). Therefore,

\begin{tcolorbox}
(如果 \({y}_{1} = {x}_{1}\) 也适用，其中 \(\gamma  = 1\) )。因此，
\end{tcolorbox}

\[
{x}_{1} \sim  {x}_{0}{ \land  }_{\gamma }{y}_{1} \sim  \left( {{y}_{0}{ \land  }_{{\eta }_{0}}{y}_{1}}\right) { \land  }_{\gamma }{y}_{1} = {y}_{0}{ \land  }_{{\eta }_{1}}{y}_{1}
\]

which shows that

\begin{tcolorbox}
这表明
\end{tcolorbox}

\[
{\eta }_{1} = \left( {1 - \gamma }\right) {\eta }_{0} + \gamma  \tag{5.42}
\]

We solve (5.40) and (5.42) for \({\eta }_{0}\) and \({\eta }_{1}\) in terms of \(\beta\) and \(\gamma\) , namely

\begin{tcolorbox}
我们用 \(\beta\) 和 \(\gamma\) 来求解(5.40)和(5.42)中的 \({\eta }_{0}\) 和 \({\eta }_{1}\) ，即
\end{tcolorbox}

\[
{\eta }_{1} = \left( {1 - \gamma }\right) \beta {\eta }_{1} + \gamma ,\;{\eta }_{1} = \frac{\gamma }{1 - \left( {1 - \gamma }\right) \beta },
\]

which agrees with the expression for \({\eta }_{1}\) in (5.38) because by (5.39) and (5.41)

\begin{tcolorbox}
这与(5.38)中 \({\eta }_{1}\) 的表达式一致，因为根据(5.39)和(5.41)
\end{tcolorbox}

\[
\frac{1 - u\left( {y}_{0}\right) }{u\left( {y}_{1}\right)  - u\left( {y}_{0}\right) } = \frac{1 + \frac{\beta }{1 - \beta }}{\frac{1}{\gamma } + \frac{\beta }{1 - \beta }} = \frac{\gamma \left( {1 - \beta }\right)  + {\gamma \beta }}{1 - \beta  + {\gamma \beta }} = \frac{\gamma }{1 - \left( {1 - \gamma }\right) \beta }.
\]

Furthermore, by (5.40)

\begin{tcolorbox}
此外，根据(5.40)
\end{tcolorbox}

\[
\frac{-u\left( {y}_{0}\right) }{u\left( {y}_{1}\right)  - u\left( {y}_{0}\right) } = \frac{\frac{\beta }{1 - \beta }}{\frac{1}{\gamma } + \frac{\beta }{1 - \beta }} = \frac{\gamma \beta }{1 - \beta  + {\gamma \beta }} = \beta {\eta }_{1} = {\eta }_{0}
\]

which shows (5.38).

\begin{tcolorbox}
这表明了(5.38)。
\end{tcolorbox}

To show (5.36), we make the usual case distinctions for \({z}_{i}\) . First, if \({x}_{0} \precsim  {z}_{i} \precsim  {x}_{1}\) then \({z}_{i} \sim  {x}_{0}{ \land  }_{{\alpha }_{i}}{x}_{1}\) with \(u\left( {z}_{i}\right)  = {\alpha }_{i}\) . We substitute (5.37) and obtain

\begin{tcolorbox}
为了证明(5.36)，我们对 \({z}_{i}\) 进行通常的情况区分。首先，如果 \({x}_{0} \precsim  {z}_{i} \precsim  {x}_{1}\) ，那么 \({z}_{i} \sim  {x}_{0}{ \land  }_{{\alpha }_{i}}{x}_{1}\) 且 \(u\left( {z}_{i}\right)  = {\alpha }_{i}\) 。我们代入(5.37)并得到
\end{tcolorbox}

\[
{z}_{i} \sim  {x}_{0}{ \land  }_{{\alpha }_{i}}{x}_{1} \sim  \left( {{y}_{0}{ \land  }_{{\eta }_{0}}{y}_{1}}\right) { \land  }_{{\alpha }_{i}}\left( {{y}_{0}{ \land  }_{{\eta }_{1}}{y}_{1}}\right)  = {y}_{0}{ \land  }_{{\beta }_{i}}{y}_{1}
\]

which shows

\begin{tcolorbox}
这表明
\end{tcolorbox}

\[
{\beta }_{i} = \left( {1 - {\alpha }_{i}}\right) {\eta }_{0} + {\alpha }_{i}{\eta }_{1} = {\eta }_{0} + {\alpha }_{i}\left( {{\eta }_{1} - {\eta }_{0}}\right)  = {\eta }_{0} + u\left( {z}_{i}\right) \left( {{\eta }_{1} - {\eta }_{0}}\right)
\]

\[
= \frac{-u\left( {y}_{0}\right) }{u\left( {y}_{1}\right)  - u\left( {y}_{0}\right) } + u\left( {z}_{i}\right) \frac{1}{u\left( {y}_{1}\right)  - u\left( {y}_{0}\right) } = v\left( {z}_{i}\right) , \tag{5.43}
\]

which again shows (5.36).

\begin{tcolorbox}
这再次表明了(5.36)。
\end{tcolorbox}

If \({y}_{0} \precsim  {z}_{i} \prec  {x}_{0}\) , then by construction of \(u\left( {z}_{i}\right)\) we have \({x}_{0} \sim  {z}_{i}{ \land  }_{{\delta }_{i}}{x}_{1}\) with

\begin{tcolorbox}
如果 \({y}_{0} \precsim  {z}_{i} \prec  {x}_{0}\) ，那么根据 \(u\left( {z}_{i}\right)\) 的构造，我们有 \({x}_{0} \sim  {z}_{i}{ \land  }_{{\delta }_{i}}{x}_{1}\) 且
\end{tcolorbox}

\[
u\left( {z}_{i}\right)  = \frac{-{\delta }_{i}}{1 - {\delta }_{i}} = 1 - \frac{1}{1 - {\delta }_{i}},\;\frac{1}{1 - {\delta }_{i}} = 1 - u\left( {z}_{i}\right) . \tag{5.44}
\]

Then \({z}_{i} \sim  {y}_{0}{ \land  }_{{\beta }_{i}}{y}_{1}\) and thus

\begin{tcolorbox}
那么 \({z}_{i} \sim  {y}_{0}{ \land  }_{{\beta }_{i}}{y}_{1}\) ，因此
\end{tcolorbox}

\[
{x}_{0} \sim  {z}_{i}{ \land  }_{{\delta }_{i}}{x}_{1} \sim  \left( {{y}_{0}{ \land  }_{{\beta }_{i}}{y}_{1}}\right) { \land  }_{{\delta }_{i}}\left( {{y}_{0}{ \land  }_{{\eta }_{1}}{y}_{1}}\right)  = {y}_{0}{ \land  }_{{\eta }_{0}}{y}_{1},
\]

that is, \({\eta }_{0} = \left( {1 - {\delta }_{i}}\right) {\beta }_{i} + {\delta }_{i}{\eta }_{1}\) and therefore by (5.44)

\begin{tcolorbox}
即， \({\eta }_{0} = \left( {1 - {\delta }_{i}}\right) {\beta }_{i} + {\delta }_{i}{\eta }_{1}\) ，因此根据(5.44)
\end{tcolorbox}

\[
{\beta }_{i} = \frac{{\eta }_{0} - {\delta }_{i}{\eta }_{1}}{1 - {\delta }_{i}} = {\eta }_{0}\left( {1 - u\left( {z}_{i}\right) }\right)  + u\left( {z}_{i}\right) {\eta }_{1} = {\eta }_{0} + u\left( {z}_{i}\right) \left( {{\eta }_{1} - {\eta }_{0}}\right)  = v\left( {z}_{i}\right)
\]

as in (5.43).

\begin{tcolorbox}
如(5.43)所示。
\end{tcolorbox}

Finally, if \({x}_{1} \prec  {z}_{i} \precsim  {y}_{1}\) , then by construction of \(u\left( {z}_{i}\right)\) we have \({x}_{1} \sim  {x}_{0}{ \land  }_{{\gamma }_{i}}{z}_{i}\) with \(u\left( {z}_{i}\right)  = 1/{\gamma }_{i}\) . Then \({z}_{i} \sim  {y}_{0}{ \land  }_{{\beta }_{i}}{y}_{1}\) and thus

\begin{tcolorbox}
最后，如果 \({x}_{1} \prec  {z}_{i} \precsim  {y}_{1}\) ，那么根据 \(u\left( {z}_{i}\right)\) 的构造，我们有 \({x}_{1} \sim  {x}_{0}{ \land  }_{{\gamma }_{i}}{z}_{i}\) 且 \(u\left( {z}_{i}\right)  = 1/{\gamma }_{i}\) 。那么 \({z}_{i} \sim  {y}_{0}{ \land  }_{{\beta }_{i}}{y}_{1}\) ，因此
\end{tcolorbox}

\[
{x}_{1} \sim  {x}_{0}{ \land  }_{{\gamma }_{i}}{z}_{i} \sim  \left( {{y}_{0}{ \land  }_{{\eta }_{0}}{y}_{1}}\right) { \land  }_{{\gamma }_{i}}\left( {{y}_{0}{ \land  }_{{\beta }_{i}}{y}_{1}}\right)  = {y}_{0}{ \land  }_{{\eta }_{1}}{y}_{1},
\]

that is, \({\eta }_{1} = \left( {1 - {\gamma }_{i}}\right) {\eta }_{0} + {\gamma }_{i}{\beta }_{i}\) and therefore

\begin{tcolorbox}
也就是说， \({\eta }_{1} = \left( {1 - {\gamma }_{i}}\right) {\eta }_{0} + {\gamma }_{i}{\beta }_{i}\) ，所以
\end{tcolorbox}

\[
{\beta }_{i} = \frac{{\eta }_{1} - \left( {1 - {\gamma }_{i}}\right) {\eta }_{0}}{{\gamma }_{i}} = \frac{{\eta }_{1}}{{\gamma }_{i}} - \frac{{\eta }_{0}}{{\gamma }_{i}} + {\eta }_{0} = u\left( {z}_{i}\right) \left( {{\eta }_{1} - {\eta }_{0}}\right)  + {\eta }_{0} = v\left( {z}_{i}\right)
\]

again as in (5.43). This shows (5.36) for all \({z}_{i}\) as claimed, and therefore the theorem.

\begin{tcolorbox}
再次如 (5.43) 式。这表明对于所有的 \({z}_{i}\) ，(5.36) 式成立，从而证明了该定理。
\end{tcolorbox}

\section*{5.9 Risk Aversion}

\begin{tcolorbox}
\section*{5.9 风险厌恶}
\end{tcolorbox}

So far, we have considered an abstract set \(X\) of outcomes and lotteries for these outcomes, and have shown that a utility function \(u : X \rightarrow  \mathbb{R}\) exists that represents the player's preference for these lotteries if this preference satisfies certain consistency axioms.

\begin{tcolorbox}
到目前为止，我们考虑了一个由结果和这些结果对应的彩票组成的抽象集合 \(X\) ，并且证明了如果玩家对这些彩票的偏好满足某些一致性公理，那么存在一个效用函数 \(u : X \rightarrow  \mathbb{R}\) 来表示玩家的偏好。
\end{tcolorbox}

In this section, we assume that \(X\) is a real interval. Each outcome is therefore a real number. This is typically a payment in money which the player wants to maximize. Is there any relationship between \(x\) and \(u\left( x\right)\) , apart from \(u\) being monotonically increasing? The concept of risk aversion means that if \(P\) is a lottery over \(X\) with an expected monetary amount \(\widehat{x}\) , then the player prefers to get the certain amount \(\widehat{x}\) to the lottery \(P\) . As an example, suppose your current wealth is \(\widehat{x}\) dollars, and you are offered a lottery where you win or lose with equal probability 100 dollars. This lottery has outcomes \(\widehat{x} - {100}\) and \(\widehat{x} + {100}\) with probability \(\frac{1}{2}\) each, and its expected value is clearly \(\widehat{x}\) dollars. If you prefer this safe amount (i.e., not play the lottery), then you are risk averse (for this specific lottery). A risk neutral player is always indifferent between \(\widehat{x}\) and \(P\) , and a risk seeking player would prefer the lottery. None of these concepts necessarily has to apply in general, because the preference may depend on \(P\) . However, in practice a player is typically risk averse. We will see that this is equivalent to the utility function being concave. In order to define concavity, consider the background material and Figure 5.4 on line segments, which provide a useful geometric view of probabilities.

\begin{tcolorbox}
在本节中，我们假设 \(X\) 是一个实数区间。因此，每个结果都是一个实数。这通常是玩家想要最大化的货币支付。除了 \(u\) 是单调递增的之外， \(x\) 和 \(u\left( x\right)\) 之间还有其他关系吗？风险厌恶的概念是指，如果 \(P\) 是一个关于 \(X\) 的彩票，其期望货币金额为 \(\widehat{x}\) ，那么玩家更倾向于获得确定的金额 \(\widehat{x}\) ，而不是参与彩票 \(P\) 。例如，假设你当前的财富是 \(\widehat{x}\) 美元，有人提供给你一个彩票，你有相等的概率赢得或输掉 100 美元。这个彩票的结果是 \(\widehat{x} - {100}\) 和 \(\widehat{x} + {100}\) ，每个结果的概率都是 \(\frac{1}{2}\) ，其期望值显然是 \(\widehat{x}\) 美元。如果你更倾向于这个安全的金额(即不参与彩票)，那么你就是风险厌恶的(针对这个特定的彩票)。一个风险中性的玩家总是对 \(\widehat{x}\) 和 \(P\) 无差异，而一个风险偏好的玩家会更倾向于彩票。这些概念不一定普遍适用，因为偏好可能取决于 \(P\) 。然而，在实践中，玩家通常是风险厌恶的。我们将看到，这等价于效用函数是凹的。为了定义凹性，请参考关于线段的背景材料和图 5.4，它们为概率提供了一个有用的几何视角。
\end{tcolorbox}

A function \(u\) is called concave if any two points on the graph of the function are connected by a line segment that is nowhere above the graph of the function.

\begin{tcolorbox}
如果函数 \(u\) 图像上的任意两点都由一条线段连接，且该线段处处不高于函数图像，则称该函数为凹函数。
\end{tcolorbox}

Background material: Line segments

\begin{tcolorbox}
背景材料:线段
\end{tcolorbox}

Figure 5.4 shows two points \(x\) and \(y\) , here in the plane, but the picture may also be regarded as a suitable view of the situation in a higher-dimensional space. The (infinite) line that goes through the points \(x\) and \(y\) is obtained by adding to the point \(x\) , regarded as a vector, any multiple of the difference vector \(y - x\) . The resulting vector \(x + \lambda  \cdot  \left( {y - x}\right)\) , for \(\lambda  \in  \mathbb{R}\) , is \(x\) when \(\lambda  = 0\) and \(y\) when \(\lambda  = 1\) . Figure 5.4 shows some examples \(a,b,c\) of other points on that line. When \(0 \leq  \lambda  \leq  1\) , like for point \(a\) , the resulting points give the line segment that joins \(x\) and \(y\) . If \(\lambda  > 1\) , then one obtains points on the line through \(x\) and \(y\) on the other side of \(y\) relative to \(x\) , like the point \(b\) in Figure 5.4. For \(\lambda  < 0\) , the corresponding point, like \(c\) in Figure 5.4, is on the other side of \(x\) relative to \(y\) .

\begin{tcolorbox}
图5.4展示了两个点 \(x\) 和 \(y\) ，这里它们位于平面上，但该图也可被视为高维度(dimension)空间中此情形的合适视图。经过点 \(x\) 和 \(y\) 的(无限)直线是通过将点 \(x\) (视为一个向量)加上差向量 \(y - x\) 的任意倍数得到的。对于 \(\lambda  \in  \mathbb{R}\) ，所得向量 \(x + \lambda  \cdot  \left( {y - x}\right)\) 在 \(\lambda  = 0\) 时为 \(x\) ，在 \(\lambda  = 1\) 时为 \(y\) 。图5.4展示了该直线上其他点的一些示例 \(a,b,c\) 。当 \(0 \leq  \lambda  \leq  1\) 时，如对于点 \(a\) ，所得点构成连接 \(x\) 和 \(y\) 的线段(line segment)。如果 \(\lambda  > 1\) ，则会得到经过 \(x\) 和 \(y\) 的直线上相对于 \(x\) 在 \(y\) 另一侧的点，如图5.4中的点 \(b\) 。对于 \(\lambda  < 0\) ，相应的点，如图5.4中的 \(c\) ，相对于 \(y\) 在 \(x\) 的另一侧。
\end{tcolorbox}

The expression \(x + \lambda \left( {y - x}\right)\) can be re-written as

\begin{tcolorbox}
表达式 \(x + \lambda \left( {y - x}\right)\) 可以重写为
\end{tcolorbox}

\[
\left( {1 - \lambda }\right) x + {\lambda y} \tag{5.45}
\]

where the given points \(x\) and \(y\) appear only once. This expression (5.45) (with \(1 - \lambda\) as the coefficient of the first point \(x\) and \(\lambda\) as the coefficient of the second point \(y\) ) shows how the real interval \(\left\lbrack  {0,1}\right\rbrack\) naturally maps to the line segment that joins \(x\) and \(y\) via the possible values of \(\lambda  \in  \left\lbrack  {0,1}\right\rbrack\) , with \(0 \mapsto  x\) and \(1 \mapsto  y\) for the endpoints. See also the right diagram in Figure 5.4. That is, for all \(x\) and \({x}^{\prime }\) in the domain of \(u\) ,

\begin{tcolorbox}
其中给定的点 \(x\) 和 \(y\) 仅出现一次。这个表达式(5.45)(以 \(1 - \lambda\) 作为第一个点 \(x\) 的系数，以 \(\lambda\) 作为第二个点 \(y\) 的系数)展示了实数区间 \(\left\lbrack  {0,1}\right\rbrack\) 如何通过 \(\lambda  \in  \left\lbrack  {0,1}\right\rbrack\) 的可能取值自然地映射到连接 \(x\) 和 \(y\) 的线段(line segment)上，端点分别对应 \(0 \mapsto  x\) 和 \(1 \mapsto  y\) 。另见图5.4中的右图。也就是说，对于 \(u\) 定义域内的所有 \(x\) 和 \({x}^{\prime }\) ，
\end{tcolorbox}

\[
\left( {1 - \lambda }\right) u\left( x\right)  + {\lambda u}\left( {x}^{\prime }\right)  \leq  u\left( {\left( {1 - \lambda }\right) x + \lambda {x}^{\prime }}\right) \;\text{ for all }\lambda  \in  \left\lbrack  {0,1}\right\rbrack  . \tag{5.46}
\]

\begin{center}
\includegraphics[max width=0.9\textwidth]{images/bo_d2e5cmbef24c73b03tn0_141_339_1184_1241_442_0.jpg}
\end{center}
\hspace*{3em} 

Figure 5.4 Left diagram: The line through the points \(x\) and \(y\) is given by the points \(x + \lambda \left( {y - x}\right)\) where \(\lambda  \in  \mathbb{R}\) . Examples are point \(a\) for \(\lambda  = {0.6}\) , point \(b\) for \(\lambda  = {1.5}\) , and point \(c\) for \(\lambda  =  - {0.4}\) . The line segment that connects \(x\) and \(y\) results when \(\lambda\) is restricted to \(0 \leq  \lambda  \leq  1\) . The right diagram shows how the real interval \(\left\lbrack  {0,1}\right\rbrack\) is naturally mapped to that line segment via the map \(\lambda  \mapsto  \left( {1 - \lambda }\right) x + {\lambda y}\) .

\begin{tcolorbox}
图5.4 左图:经过点 \(x\) 和 \(y\) 的直线由点 \(x + \lambda \left( {y - x}\right)\) 给出，其中 \(\lambda  \in  \mathbb{R}\) 。示例包括当 \(\lambda  = {0.6}\) 时的点 \(a\) 、当 \(\lambda  = {1.5}\) 时的点 \(b\) 以及当 \(\lambda  =  - {0.4}\) 时的点 \(c\) 。当 \(\lambda\) 限制在 \(0 \leq  \lambda  \leq  1\) 时得到连接 \(x\) 和 \(y\) 的线段(line segment)。右图展示了实数区间 \(\left\lbrack  {0,1}\right\rbrack\) 如何通过映射 \(\lambda  \mapsto  \left( {1 - \lambda }\right) x + {\lambda y}\) 自然地映射到该线段上。
\end{tcolorbox}

As an example, Figure 5.5 shows the function \(u\left( x\right)  = \sqrt{x}\) defined on \(\left\lbrack  {0,1}\right\rbrack\) . On the graph of the function (as a subset of \({\mathbb{R}}^{2}\) ) we show three white dots that are the points \(\left( {x,u\left( x\right) }\right) ,\left( {\left( {1 - \lambda }\right) x + \lambda {x}^{\prime },u\left( {\left( {1 - \lambda }\right) x + \lambda {x}^{\prime }}\right) }\right)\) , and \(\left( {{x}^{\prime },u\left( {x}^{\prime }\right) }\right)\) . The black dot is the convex combination \(\left( {1 - \lambda }\right)  \cdot  \left( {x,u\left( x\right) }\right)  + \lambda  \cdot  \left( {{x}^{\prime },u\left( {x}^{\prime }\right) }\right)\) of the outer white dots, which is below the middle white dot, as asserted in (5.46). This holds generally and shows that \(u\) is concave.

\begin{tcolorbox}
例如，图5.5展示了定义在 \(u\left( x\right)  = \sqrt{x}\) 上的函数 \(\left\lbrack  {0,1}\right\rbrack\) 。在该函数的图(作为 \({\mathbb{R}}^{2}\) 的子集)上，我们展示了三个白色点，分别是点 \(\left( {x,u\left( x\right) }\right) ,\left( {\left( {1 - \lambda }\right) x + \lambda {x}^{\prime },u\left( {\left( {1 - \lambda }\right) x + \lambda {x}^{\prime }}\right) }\right)\) 和 \(\left( {{x}^{\prime },u\left( {x}^{\prime }\right) }\right)\) 。黑色点是外部白色点的凸组合 \(\left( {1 - \lambda }\right)  \cdot  \left( {x,u\left( x\right) }\right)  + \lambda  \cdot  \left( {{x}^{\prime },u\left( {x}^{\prime }\right) }\right)\) ，正如(5.46)中所断言的，它位于中间白色点的下方。这一情况普遍成立，表明 \(u\) 是凹函数。
\end{tcolorbox}

If we have equality in (5.46), then that part of the graph of the function is itself a line segment. If we never have equality, that is, (5.46) holds as a strict inequality whenever \(x \neq  {x}^{\prime }\) and \(0 < \lambda  < 1\) , then the function \(u\) is said to be strictly concave.

\begin{tcolorbox}
如果(5.46)式取等号，那么函数图的那部分本身就是一条线段。如果永远不取等号，即只要 \(x \neq  {x}^{\prime }\) 和 \(0 < \lambda  < 1\) 时，(5.46)式以严格不等式成立，那么函数 \(u\) 被称为严格凹函数。
\end{tcolorbox}

\begin{center}
\includegraphics[max width=0.6\textwidth]{images/bo_d2e5cmbef24c73b03tn0_142_275_562_859_643_0.jpg}
\end{center}
\hspace*{3em} 

Figure 5.5 Example of a concave function \(u\left( x\right)  = \sqrt{x}\) on the interval \(\left\lbrack  {0,1}\right\rbrack\) .

\begin{tcolorbox}
图5.5 区间 \(\left\lbrack  {0,1}\right\rbrack\) 上凹函数 \(u\left( x\right)  = \sqrt{x}\) 的示例。
\end{tcolorbox}

It is easy to see that a function that is twice differentiable is concave if and only if its second derivative is always nonnegative, and strictly concave if its second derivative is always positive (as in the case of the square root function \(\sqrt{x}\) for \(x > 0\) ). Furthermore, it can be shown that, except at the endpoints of the interval that is the domain of the function, concavity implies continuity.

\begin{tcolorbox}
很容易看出，一个二次可微的函数是凹函数当且仅当它的二阶导数始终是非负的；如果它的二阶导数始终为正(例如对于 \(x > 0\) 时的平方根函数 \(\sqrt{x}\) )，则该函数是严格凹函数。此外，可以证明，除了函数定义域区间的端点外，凹性意味着连续性。
\end{tcolorbox}

The following observation states that risk aversion is equivalent to the concavity of the player's utility function. For simple lotteries its proof is almost immediate.

\begin{tcolorbox}
以下观察表明，风险厌恶等价于玩家的效用函数是凹函数。对于简单彩票，其证明几乎是直接的。
\end{tcolorbox}

Proposition 5.5. Let \(X\) be a real interval and let \(u : X \rightarrow  \mathbb{R}\) be an expected-utility function. Consider any lottery \(P\) over \(X\) with expected value \(\widehat{x} = \mathbb{E}\left( P\right)\) . Then the player is risk averse, that is, always \(P \precsim  \widehat{x}\) , if and only if \(u\) is concave.

\begin{tcolorbox}
命题5.5。设 \(X\) 是一个实数区间， \(u : X \rightarrow  \mathbb{R}\) 是一个期望效用函数。考虑 \(X\) 上的任意彩票 \(P\) ，其期望值为 \(\widehat{x} = \mathbb{E}\left( P\right)\) 。那么玩家是风险厌恶的，即总是 \(P \precsim  \widehat{x}\) ，当且仅当 \(u\) 是凹函数。
\end{tcolorbox}

Proof. Because \(X\) is a real interval, any lottery \(P\) over \(X\) is a random variable that has an expectation \(\widehat{x} = \mathbb{E}\left( P\right)\) , where \(\widehat{x} \in  X\) . Consider first (see (5.16)) a simple lottery \(P = x{ \land  }_{\lambda }{x}^{\prime }\) for \(\lambda  \in  \left\lbrack  {0,1}\right\rbrack\) and \(x,{x}^{\prime } \in  X\) , where \(\widehat{x} = \mathbb{E}\left( P\right)  = \left( {1 - \lambda }\right) x + \lambda {x}^{\prime }\) . Then \(P \precsim  \widehat{x}\) is equivalent to \(\mathbb{E}\left( {u,P}\right)  = \left( {1 - \lambda }\right) u\left( x\right)  + {\lambda u}\left( {x}^{\prime }\right)  \leq  u\left( \widehat{x}\right)\) , which is exactly the concavity condition (5.46).

\begin{tcolorbox}
证明。因为 \(X\) 是一个实数区间， \(X\) 上的任何彩票 \(P\) 都是一个具有期望值 \(\widehat{x} = \mathbb{E}\left( P\right)\) 的随机变量，其中 \(\widehat{x} \in  X\) 。首先考虑(见(5.16))一个简单彩票 \(P = x{ \land  }_{\lambda }{x}^{\prime }\) ，其中 \(\lambda  \in  \left\lbrack  {0,1}\right\rbrack\) 且 \(x,{x}^{\prime } \in  X\) ，这里 \(\widehat{x} = \mathbb{E}\left( P\right)  = \left( {1 - \lambda }\right) x + \lambda {x}^{\prime }\) 。那么 \(P \precsim  \widehat{x}\) 等价于 \(\mathbb{E}\left( {u,P}\right)  = \left( {1 - \lambda }\right) u\left( x\right)  + {\lambda u}\left( {x}^{\prime }\right)  \leq  u\left( \widehat{x}\right)\) ，这正是凹性条件(5.46)。
\end{tcolorbox}

It remains to show that if \(u\) is concave then \(P \precsim  \mathbb{E}\left( P\right)\) for any lottery \(P\) that assigns positive probabilities \({p}_{1},\ldots ,{p}_{n}\) to outcomes \({x}_{1},\ldots ,{x}_{n}\) for \(n \geq  3\) . As inductive hypothesis, assume this is true for \(n - 1\) instead of \(n\) (we have just shown it for \(n = 2\) ). Write \(P\) as \(Q{ \land  }_{{p}_{n}}{x}_{n}\) where \(Q\) assigns probability \(\frac{{p}_{i}}{1 - {p}_{n}}\) to \({x}_{i}\) for \(1 \leq  i \leq  n - 1\) . By inductive hypothesis, \(\mathbb{E}\left( {Q,u}\right)  \leq  u\left( {\mathbb{E}\left( Q\right) }\right)\) . Therefore (where the second inequality holds because \(u\) is concave),

\begin{tcolorbox}
接下来需要证明，如果 \(u\) 是凹的，那么对于任何将正概率 \({p}_{1},\ldots ,{p}_{n}\) 分配给结果 \({x}_{1},\ldots ,{x}_{n}\) (其中 \(n \geq  3\) )的彩票 \(P\) ，有 \(P \precsim  \mathbb{E}\left( P\right)\) 。作为归纳假设，假设对于 \(n - 1\) 而非 \(n\) 这是成立的(我们刚刚已经证明了对于 \(n = 2\) 是成立的)。将 \(P\) 写为 \(Q{ \land  }_{{p}_{n}}{x}_{n}\) ，其中 \(Q\) 将概率 \(\frac{{p}_{i}}{1 - {p}_{n}}\) 分配给 \({x}_{i}\) (对于 \(1 \leq  i \leq  n - 1\) )。根据归纳假设， \(\mathbb{E}\left( {Q,u}\right)  \leq  u\left( {\mathbb{E}\left( Q\right) }\right)\) 。因此(其中第二个不等式成立是因为 \(u\) 是凹的)
\end{tcolorbox}

\[
\mathbb{E}\left( {P,u}\right)  = \mathop{\sum }\limits_{{i = 1}}^{{n - 1}}{p}_{i}u\left( {x}_{i}\right)  + {p}_{n}{x}_{n}
\]

\[
= \left( {1 - {p}_{n}}\right) \mathbb{E}\left( {Q,u}\right)  + {p}_{n}{x}_{n}
\]

\[
\leq  \left( {1 - {p}_{n}}\right) u\left( {\mathbb{E}\left( Q\right) }\right)  + {p}_{n}{x}_{n}
\]

\[
\leq  u\left( {\left( {1 - {p}_{n}}\right) \mathbb{E}\left( Q\right)  + {p}_{n}{x}_{n}}\right)  = u\left( {\mathbb{E}\left( P\right) }\right) ,
\]

which completes the proof.

\begin{tcolorbox}
这就完成了证明。
\end{tcolorbox}

We stress again that risk aversion requires real-valued outcomes, such as money, where the expectation of a lottery is itself an outcome which can be compared with the lottery.

\begin{tcolorbox}
我们再次强调，风险厌恶要求结果是实值的，例如金钱，在这种情况下，彩票的期望本身就是一个可以与彩票进行比较的结果。
\end{tcolorbox}

Beyond concavity, risk aversion is not easy to quantify. If the player has a wealth of \(\widehat{x}\) dollars and is risk averse, he will prefer \(\widehat{x}\) to an equal gamble between \(\widehat{x} - {100}\) and \(\widehat{x} + {100}\) dollars. The player should therefore be willing to pay a "risk premium" \(\pi\) so that he is indifferent between safely receiving \(\widehat{x} - \pi\) and the gamble. For the utility function \(u\left( x\right)  =  - {e}^{-{cx}}\) for some \(c > 0\) , it is easy to see that \(\pi\) does not depend on \(\widehat{x}\) , and that this is up to positive-affine transformations the only concave utility function with this property. If the values of \(\widehat{x}\) are realistic, and the lottery itself (like winning or losing 100 dollars) is comparable to the risks under consideration, then estimating the risk premium \(\pi\) may be a useful approach to constructing the utility function.

\begin{tcolorbox}
除了凹性之外，风险厌恶并不容易量化。如果玩家拥有 \(\widehat{x}\) 美元的财富并且是风险厌恶的，他会更喜欢确定的 \(\widehat{x}\) 美元，而不是在 \(\widehat{x} - {100}\) 美元和 \(\widehat{x} + {100}\) 美元之间进行等概率的赌博。因此，玩家应该愿意支付一个“风险溢价” \(\pi\) ，使得他对安全地获得 \(\widehat{x} - \pi\) 美元和进行赌博无差异。对于效用函数 \(u\left( x\right)  =  - {e}^{-{cx}}\) (对于某个 \(c > 0\) )，很容易看出 \(\pi\) 不依赖于 \(\widehat{x}\) ，并且在正仿射变换下，这是唯一具有此性质的凹效用函数。如果 \(\widehat{x}\) 的值是现实的，并且彩票本身(例如赢或输 100 美元)与所考虑的风险具有可比性，那么估计风险溢价 \(\pi\) 可能是构建效用函数的一种有用方法。
\end{tcolorbox}

\section*{5.10 Discussion and Further Reading}

\begin{tcolorbox}
\section*{5.10 讨论与进一步阅读}
\end{tcolorbox}

The consistency axioms that imply the existence of an expected-utility function are sometimes called "rationality axioms". A player is called "rational" when he acts in accordance with his own consistent preferences. In game theory, further assumptions are that all players know the entire game, including the preferences of all players, that everybody is rational, and often that this is common knowledge (everybody knows that everybody knows, and so on). For more detailed discussions see Aumann and Brandenburger (1995), Aumann and Dreze (2008), and the book by Perea (2012).

\begin{tcolorbox}
暗示期望效用函数存在的一致性公理有时被称为“理性公理”。当一个玩家按照自己一致的偏好行动时，他被称为“理性的”。在博弈论中，进一步的假设是所有玩家都了解整个博弈，包括所有玩家的偏好，每个人都是理性的，并且通常这是共同知识(每个人都知道每个人都知道，依此类推)。更详细的讨论请参见奥曼和布兰登伯格(1995 年)、奥曼和德雷兹(2008 年)以及佩雷亚(2012 年)的书。
\end{tcolorbox}

Game theory, and economic theory in general, is often criticized for these idealized assumptions. It should be clear that the assumptions are part of the mathematical model, which may not be particularly realistic, and need to be questioned when considering the conclusions of a game-theoretic analysis. These conclusions already pose difficulties when a solution (such as equilibrium) is not unique. But even if unique, the conclusions are normative (stating what players should do) rather than descriptive (stating what players actually do).

\begin{tcolorbox}
博弈论以及一般的经济理论常常因这些理想化的假设而受到批评。应该明确的是，这些假设是数学模型的一部分，可能并不特别现实，在考虑博弈论分析的结论时需要对其进行质疑。当一个解(如均衡)不唯一时，这些结论就已经存在困难。但即使解是唯一的，这些结论也是规范的(陈述玩家应该做什么)而非描述性的(陈述玩家实际做了什么)。
\end{tcolorbox}

A further source of misunderstanding game theory is the use of words for mathematical concepts that have strong connotations elsewhere. It is clear that a game tree is not an actual tree but describes a sequence of decision "branches", so "tree" can be used as a neutral word. However, the word "utility" suggests an intrinsic strength of preference associated with the "usefulness" of an outcome. As argued in this chapter, it should be considered as a mathematical tool for representing a preference for risky outcomes.

\begin{tcolorbox}
对博弈论产生误解的另一个原因是，用于数学概念的词汇在其他地方有强烈的内涵。显然，博弈树并不是真正的树，而是描述了一系列决策“分支”，因此“树”可以作为一个中性的词汇。然而，“效用”这个词暗示了与结果的“有用性”相关的内在偏好强度。正如本章所论述的，它应该被视为一种用于表示对风险结果偏好的数学工具。
\end{tcolorbox}

Because of the connotations of the word "utility", we instead use the more neutral and active term "payoff" in the context of games in this book. However, it is useful to remember that payoffs represent a preference and can therefore be scaled with positive-affine transformations.

\begin{tcolorbox}
由于“效用”这个词的内涵，在本书的博弈背景下，我们改用更中性和积极的术语“收益”。然而，记住收益代表一种偏好，因此可以通过正仿射变换进行缩放是很有用的。
\end{tcolorbox}

A number of classic books on game theory, such as Luce and Raiffa (1957) and Myerson (1991), start with the rationality assumptions that justify the use of expected payoffs. Unusually, Osborne and Rubinstein (1994) even define games in terms of preference relations rather than payoffs.

\begin{tcolorbox}
许多关于博弈论的经典著作，如卢斯和雷法(1957 年)以及迈尔森(1991 年)的著作，都始于证明使用期望效用合理性的理性假设。不同寻常的是，奥斯本和鲁宾斯坦(1994 年)甚至用偏好关系而非效用来定义博弈。
\end{tcolorbox}

The concept of expected utility is generally attributed to Daniel Bernoulli (1738), who proposed the logarithm as a utility function for lotteries that involve money. For an informative and entertaining historical account see Szpiro (2020), and also Szpiro (2010) on approaches to group decision making, the Condorcet Paradox, and social choice in general. The consistency axioms for an expected-utility function are due to, and often named after, von Neumann and Morgenstern (1947). The paradox in Figure 5.2 has been described by Allais (1953). For a short history of the discussion around these axioms, and further references, see Moscati (2016).

\begin{tcolorbox}
期望效用的概念通常归功于丹尼尔·伯努利(1738 年)，他提出用对数作为涉及货币的彩票的效用函数。关于这一概念的信息丰富且有趣的历史记载，可参见斯皮罗(2020 年)的著作，关于群体决策方法、孔多塞悖论和一般社会选择的内容也可参见斯皮罗(2010 年)的著作。期望效用函数的一致性公理由冯·诺依曼和摩根斯坦(1947 年)提出，并且通常以他们的名字命名。图 5.2 中的悖论由阿莱(1953 年)描述。关于围绕这些公理的讨论的简短历史及更多参考文献，可参见莫斯卡蒂(2016 年)的著作。
\end{tcolorbox}

A classic book on utility theory, which connects axioms about preferences and their representation with a utility function, is Fishburn (1970). An abstract approach to preferences, due to Herstein and Milnor (1953), is called a "mixture set" where two lotteries (including safe outcomes) \(P\) and \(Q\) and a probability \(\alpha\) are combined into a new lottery \(\left( {1 - \alpha }\right) P + {\alpha Q}\) as in (5.15), with suitable axioms about these mixtures of \(P,Q,\alpha\) . Rather than considering such abstract "mixtures", we have kept their meaning as lotteries, so that the algebraic manipulations, for example as used in the proof of Theorem 5.4, stay concrete.

\begin{tcolorbox}
一本关于效用理论的经典著作是菲什伯恩(1970 年)的作品，它将关于偏好的公理及其用效用函数表示的内容联系起来。赫斯坦和米尔诺(1953 年)提出的一种关于偏好的抽象方法被称为“混合集”，其中两个彩票(包括确定结果) \(P\) 和 \(Q\) 以及一个概率 \(\alpha\) 按照(5.15)式组合成一个新的彩票 \(\left( {1 - \alpha }\right) P + {\alpha Q}\) ，并且有关于这些 \(P,Q,\alpha\) 混合的合适公理。我们没有考虑这种抽象的“混合”，而是保留了它们作为彩票的含义，这样代数运算，例如在定理 5.4 的证明中使用的运算，就保持具体。
\end{tcolorbox}

Utility theory is not a good descriptive theory of human decision making, as documented and extensively studied by the psychologists Daniel Kahneman and Amos Tversky. Their collaboration and research is vividly described in the popular-science book by Lewis (2016). The example (5.10) of intransitive preferences is due to Tversky (1969).

\begin{tcolorbox}
效用理论并不是一个很好的描述人类决策的理论，心理学家丹尼尔·卡尼曼和阿莫斯·特沃斯基对此进行了记录和广泛研究。刘易斯(2016 年)的科普书籍生动地描述了他们的合作与研究。非传递偏好的例子(5.10)由特沃斯基(1969 年)提出。
\end{tcolorbox}

An important paper on risk aversion is Pratt (1964).

\begin{tcolorbox}
一篇关于风险厌恶的重要论文是普拉特(1964 年)的作品。
\end{tcolorbox}

\section*{5.11 Exercises for Chapter 5}

\begin{tcolorbox}
\section*{5.11 第五章练习}
\end{tcolorbox}

Exercise 5.1. Find positive weights \(w\) in a utility function of the form \(u\left( {{c}_{1},{c}_{2}}\right)  =\)  \({c}_{1} + w \cdot  {c}_{2}\) to create at least four different strict rankings among the candidates \(x,y,z\) in (5.10). [Hint: Experiment with some values of \(w\) and vary them. What are the values of \(u\left( {{c}_{1},{c}_{2}}\right)\) for the three candidates as a function of \(w\) ?]

\begin{tcolorbox}
练习 5.1。在形式为 \(u\left( {{c}_{1},{c}_{2}}\right)  =\)  \({c}_{1} + w \cdot  {c}_{2}\) 的效用函数中找到正权重 \(w\) ，以使(5.10)中的候选人 \(x,y,z\) 之间至少有四种不同的严格排序。[提示:尝试一些 \(w\) 的值并改变它们。作为 \(w\) 的函数，三个候选人的 \(u\left( {{c}_{1},{c}_{2}}\right)\) 值是多少？]
\end{tcolorbox}

Exercise 5.2. Complete the proof of Proposition 5.2 for the third case \({x}_{1} \prec  z\) in (5.21), to show that \(v\left( z\right)  = {au}\left( z\right)  + b\) . As stated in the proof, choose a simple lottery \(Q = {x}_{0}{ \land  }_{\gamma }z\) so that the player is indifferent between that lottery and \({x}_{1}\) , and proceed similarly to (5.23) using (5.18) and (5.20).

\begin{tcolorbox}
练习 5.2。完成命题 5.2 中(5.21)式第三种情况 \({x}_{1} \prec  z\) 的证明，以表明 \(v\left( z\right)  = {au}\left( z\right)  + b\) 。如证明中所述，选择一个简单彩票 \(Q = {x}_{0}{ \land  }_{\gamma }z\) ，使得玩家对该彩票和 \({x}_{1}\) 无差异，然后使用(5.18)式和(5.20)式，按照(5.23)式的方式进行推导。
\end{tcolorbox}

Exercise 5.3.

\begin{tcolorbox}
练习 5.3。
\end{tcolorbox}

(a) Prove the equivalence (5.25), including the inequalities, for real numbers \(A,B,C\) that fulfill \(A < C\) .

\begin{tcolorbox}
(a) 对于满足 \(A < C\) 的实数 \(A,B,C\) ，证明等价关系(5.25)，包括其中的不等式。
\end{tcolorbox}

(b) Prove \(\left( {1 - \alpha }\right) \mathbb{E}\left( {u,P}\right)  + \alpha \mathbb{E}\left( {u,R}\right)  = \mathbb{E}\left( {u,\left( {1 - \alpha }\right) P + {\alpha R}}\right)\) for two lotteries \(P\) and \(R\) , which is the second equation in (5.30).

\begin{tcolorbox}
(b) 证明两个彩票 \(P\) 和 \(R\) 的 \(\left( {1 - \alpha }\right) \mathbb{E}\left( {u,P}\right)  + \alpha \mathbb{E}\left( {u,R}\right)  = \mathbb{E}\left( {u,\left( {1 - \alpha }\right) P + {\alpha R}}\right)\) 关系，这是(5.30)式中的第二个等式。
\end{tcolorbox}

\section*{Mixed Equilibrium}

\begin{tcolorbox}
\section*{混合均衡}
\end{tcolorbox}

A game in strategic form does not always have an equilibrium in which each player chooses her strategy deterministically. As we describe in this chapter, Nash (1951) showed that any finite strategic-form game has an equilibrium if players are allowed to use mixed strategies. A mixed strategy is an actively randomized choice of the player's given "pure" strategies according to certain probabilities. A profile of mixed strategies is called a mixed equilibrium if no player can improve her expected payoff by unilaterally changing her strategy.

\begin{tcolorbox}
策略形式的博弈并不总是存在每个玩家都确定性地选择其策略的均衡。正如我们在本章中所描述的，纳什(1951 年)证明了，如果允许玩家使用混合策略，任何有限策略形式的博弈都存在均衡。混合策略是玩家根据一定概率对给定的“纯”策略进行主动随机选择。如果没有玩家可以通过单方面改变其策略来提高其期望效用，那么一个混合策略组合就被称为混合均衡。
\end{tcolorbox}

In this chapter, we briefly recall from Chapter 5 how a player's payoffs represent her preference for random outcomes via an expected-utility function. This is illustrated with a single-player decision problem whether to comply or not with a legal requirement, which is then turned into a game. The game is known as an Inspection game between an inspector as player I and an inspectee as player II. This game has no equilibrium in pure strategies. However, active randomization comes up naturally in this game, because the inspector can and will choose to inspect only some of the time. The mixed equilibrium in this game demonstrates many aspects of general mixed equilibria, in particular that a player's randomization probabilities depend on the payoffs of the other player.

\begin{tcolorbox}
在本章中，我们简要回顾第5章的内容，即玩家的收益如何通过期望效用函数来表示她对随机结果的偏好。这通过一个单人决策问题来说明，即是否遵守一项法律要求，然后将其转化为一个博弈。这个博弈被称为监管博弈，玩家I是监管者，玩家II是被监管者。这个博弈在纯策略中不存在均衡。然而，在这个博弈中自然会出现主动随机化，因为监管者可以而且只会选择在某些时候进行检查。这个博弈中的混合均衡展示了一般混合均衡的许多方面，特别是一个玩家的随机化概率取决于另一个玩家的收益。
\end{tcolorbox}

We then turn to general two-player games in strategic form, already studied in Chapter 3, here called bimatrix games. The payoff matrices for the two players are convenient for representing expected payoffs via multiplication with vectors of probabilities that are the players' mixed strategies. These mixed strategies, as vectors of probabilities, have a geometric interpretation, which we study in detail (and even more so in Chapter 9). Mixing with probabilities is geometrically equivalent to taking convex combinations.

\begin{tcolorbox}
然后我们转向策略形式的一般两人博弈，在第3章中已经研究过，这里称为双矩阵博弈。两个玩家的收益矩阵便于通过与玩家混合策略的概率向量相乘来表示期望收益。这些作为概率向量的混合策略具有几何解释，我们将详细研究(在第9章会研究得更多)。从几何角度看，用概率进行混合等同于取凸组合。
\end{tcolorbox}

In a two-player game we usually name the players I and II, and denote their pure strategies as rows \(i = 1,\ldots ,m\) and columns \(j = 1,\ldots ,n\) , respectively. In \(N\) -player games, which we only consider for general definitions, the players are numbered as players \(i = 1,\ldots ,N\) , which is the most commonly used notation.

\begin{tcolorbox}
在两人博弈中，我们通常将玩家命名为I和II，并分别将他们的纯策略表示为行 \(i = 1,\ldots ,m\) 和列 \(j = 1,\ldots ,n\) 。在 \(N\) 人博弈中，我们仅考虑一般定义，玩家编号为 \(i = 1,\ldots ,N\) ，这是最常用的表示法。
\end{tcolorbox}

The best-response condition gives a convenient, finite condition for when a mixed strategy is a best response to the mixed strategy of the other player (or against the mixed strategies of all other players in a game with more than two players). Namely, only the pure strategies that are best responses are allowed to be played with positive probability by a best-response mixed strategy. These pure best responses must have maximal, and hence equal, expected payoff. This is stated as a "trivial condition" by Nash (1951, p. 287), and it is indeed not hard to prove (see Proposition 6.1). However, the best-response condition is central for understanding and computing mixed-strategy equilibria.

\begin{tcolorbox}
最优反应条件为混合策略何时是对另一个玩家的混合策略(或在多人博弈中对所有其他玩家的混合策略)的最优反应提供了一个方便的有限条件。即，最优反应混合策略只允许以正概率采用那些作为最优反应的纯策略。这些纯最优反应必须具有最大的，因此是相等的期望收益。纳什(1951，第287页)将其表述为一个“平凡条件”，而且确实不难证明(见命题6.1)。然而，最优反应条件对于理解和计算混合策略均衡至关重要。
\end{tcolorbox}

In Section 6.5, we state and prove Nash's theorem that every game has an equilibrium. We follow Nash's original proof (for two players, to simplify notation), and state and briefly motivate Brouwer's fixed-point theorem used in that proof. In Chapter 7 we prove the fixed-point theorem itself.

\begin{tcolorbox}
在6.5节中，我们陈述并证明纳什定理，即每个博弈都存在一个均衡。我们遵循纳什的原始证明(为简化表示，针对两人博弈)，并陈述并简要说明该证明中使用的布劳威尔不动点定理。在第7章中，我们将证明不动点定理本身。
\end{tcolorbox}

The important final Sections 6.6-6.8 show how to find equilibria in bimatrix games, in particular when one player has only two strategies.

\begin{tcolorbox}
重要的最后6.6 - 6.8节展示了如何在双矩阵博弈中找到均衡，特别是当一个玩家只有两种策略时。
\end{tcolorbox}

The special and interesting case of zero-sum games is treated in Chapter 8 .

\begin{tcolorbox}
零和博弈这一特殊且有趣的情况将在第8章中讨论。
\end{tcolorbox}

\section*{6.1 Prerequisites and Learning Objectives}

\begin{tcolorbox}
\section*{6.1 先决条件和学习目标}
\end{tcolorbox}

Prerequisites are Chapter 3 and knowing that payoffs are only unique up to positive-affine transformations, as discussed in Chapter 5. Fluency with matrix multiplication is useful, which we will recall as background material. We also need the algebraic representation of line segments explained in Figure 5.4.

\begin{tcolorbox}
先决条件是熟悉第3章的内容，并知道收益仅在正仿射变换下是唯一的，如第5章所讨论的。熟练掌握矩阵乘法会很有用，我们将把它作为背景材料进行回顾。我们还需要图5.4中解释的线段的代数表示。
\end{tcolorbox}

After studying this chapter, you should be able to:

\begin{tcolorbox}
学习本章后，你应该能够:
\end{tcolorbox}

\begin{itemize}
\item explain what mixed strategies and mixed equilibria are;
\end{itemize}

\begin{tcolorbox}
\begin{itemize}
\item 解释什么是混合策略和混合均衡；
\end{itemize}
\end{tcolorbox}

\begin{itemize}
\item state the main theorems about mixed strategies, namely the best-response condition, and Nash's theorem on the existence of mixed equilibria;
\end{itemize}

\begin{tcolorbox}
\begin{itemize}
\item 陈述关于混合策略的主要定理，即最优反应条件和纳什关于混合均衡存在性的定理；
\end{itemize}
\end{tcolorbox}

\begin{itemize}
\item (as an important skill) find all mixed equilibria of small bimatrix games, using the "difference trick" and, where necessary, the upper-envelope method;
\end{itemize}

\begin{tcolorbox}
\begin{itemize}
\item (作为一项重要技能)使用“差分技巧”，必要时使用上包络法，找出小型双矩阵博弈的所有混合均衡；
\end{itemize}
\end{tcolorbox}

\begin{itemize}
\item give the definition of a degenerate game, and find all equilibria of small degenerate games.
\end{itemize}

\begin{tcolorbox}
\begin{itemize}
\item 给出退化博弈的定义，并找出小型退化博弈的所有均衡。
\end{itemize}
\end{tcolorbox}

\section*{6.2 Compliance Inspections}

\begin{tcolorbox}
\section*{6.2 合规检查}
\end{tcolorbox}

In this section, we consider a two-player Inspection game where randomized or mixed strategies occur naturally. Against a mixed strategy, the other player faces a lottery, which is evaluated according to its expected payoffs, as studied in detail in the preceding Chapter 5.

\begin{tcolorbox}
在本节中，我们考虑一个两人监管博弈，其中随机或混合策略自然出现。面对一个混合策略，另一个玩家面临一场彩票，根据其期望收益对其进行评估，如前一章第5章详细研究的那样。
\end{tcolorbox}

To recap this briefly, Figure 6.1 shows a game with a single player who can decide to comply with a legal requirement, like buying a train ticket, or to cheat otherwise. The payoff when she chooses to comply is 0 . Cheating involves a 10 percent chance of getting caught and having to pay a penalty, stated as the negative payoff -90, and otherwise a 90 percent chance of gaining a payoff of 10 . With these numbers, cheating leads to a random outcome with an expected payoff of \({0.9} \cdot  {10} + {0.1} \cdot  \left( {-{90}}\right)\) , which is zero, so that the player is indifferent between her two moves. In fact, assuming that the two outcomes with payoffs 0 and 10 have been set as origin and unit of the utility function, the utility of -90 for "getting caught" is determined by this indifference. If the player is indifferent when the probability of getting caught is only four percent, then -90 has to be replaced by \(- u\) so that

\begin{tcolorbox}
简要回顾一下，图6.1展示了一个单人博弈，该玩家可以选择遵守法律要求(如购买火车票)，或者选择违规。她选择遵守时的收益为0。违规有10\%的概率被抓住并需支付罚款，此时收益记为 -90，另外有90\%的概率获得10的收益。根据这些数值，违规会导致一个随机结果，期望效用为 \({0.9} \cdot  {10} + {0.1} \cdot  \left( {-{90}}\right)\) ，即0，因此玩家对这两种选择无差异。实际上，假设收益为0和10的两种结果已被设定为效用函数的原点和单位，那么“被抓住”的效用 -90 就由这种无差异决定。如果玩家在被抓住的概率仅为4\%时无差异，那么 -90 必须替换为 \(- u\) ，使得
\end{tcolorbox}

\[
0 = {0.96} \cdot  {10} + {0.04} \cdot  \left( {-u}\right) ,
\]

\begin{center}
\includegraphics[max width=0.2\textwidth]{images/bo_d2e5cmbef24c73b03tn0_148_338_201_352_373_0.jpg}
\end{center}
\hspace*{3em} 

Figure 6.1 One-player decision problem between comply and cheat. With the stated payoffs, the player is indifferent.

\begin{tcolorbox}
图6.1 遵守与违规之间的单人决策问题。在给定的收益下，玩家无差异。
\end{tcolorbox}

which is equivalent to \(u = {9.6}/{0.04} = {240}\) . Also, recall that the origin and unit of the utility function are arbitrary and can be changed by a positive-affine transformation. For example, \(0,{10}, - {90}\) could be replaced by \(0,1, - 9\) or \(- 1,0, - {10}\) .

\begin{tcolorbox}
这等同于 \(u = {9.6}/{0.04} = {240}\) 。此外，请记住，效用函数的原点和单位是任意的，可以通过正仿射变换进行改变。例如， \(0,{10}, - {90}\) 可以替换为 \(0,1, - 9\) 或 \(- 1,0, - {10}\) 。
\end{tcolorbox}

In practice, the preference for risk of a player may not be known. A game-theoretic analysis should be carried out for different choices of the payoff parameters in order to test how much they influence the results.

\begin{tcolorbox}
在实践中，玩家的风险偏好可能是未知的。为了测试收益参数对结果的影响程度，应该针对不同的收益参数选择进行博弈论分析。
\end{tcolorbox}

We now consider the Inspection game. For player II, called inspectee, the strategy comply means again to fulfill some legal obligation (such as buying a transport ticket, or paying tax). The inspectee has an incentive to violate this obligation by choosing to cheat. Player I, the inspector, should verify that the inspectee complies, but doing so requires inspections which are costly. If the inspector inspects and catches the inspectee cheating, then the inspectee has to pay a penalty.

\begin{tcolorbox}
现在我们考虑监管博弈。对于玩家II，即被监管者，“遵守”策略再次意味着履行某些法律义务(如购买交通票或纳税)。被监管者有动机通过选择“违规”来违反这项义务。玩家I，即监管者，应该核实被监管者是否遵守规定，但这样做需要进行检查，会产生成本。如果监管者进行检查并抓住被监管者违规，那么被监管者必须支付罚款。
\end{tcolorbox}

Figure 6.2 shows possible payoffs for the Inspection game. The outcome that defines the reference payoff zero to both players is that the inspector chooses Don't Inspect and the inspectee chooses to comply. Without inspection, the inspectee prefers to cheat, represented by her payoff 10, with negative payoff -10 to the inspector. The inspector may also decide to Inspect. If the inspectee complies, inspection leaves her payoff 0 unchanged, while the inspector incurs a cost with a negative payoff of -1 . If the inspectee cheats, however, inspection will result in a penalty (payoff -90 for player II) and still create a certain amount of hassle for player I (payoff -6).

\begin{tcolorbox}
图6.2展示了监管博弈的可能收益情况。为两个玩家定义参考收益为0的结果是，监管者选择“不检查”，被监管者选择“遵守”。如果不进行检查，被监管者更倾向于违规，其收益为10，而监管者的收益为 -10。监管者也可以决定“检查”。如果被监管者遵守规定，检查不会改变她的收益0，而监管者会产生成本，收益为 -1。然而，如果被监管者违规，检查将导致被监管者支付罚款(玩家II的收益为 -90)，并且仍然会给玩家I带来一定的麻烦(收益为 -6)。
\end{tcolorbox}

\begin{center}
\includegraphics[max width=0.4\textwidth]{images/bo_d2e5cmbef24c73b03tn0_149_335_196_542_401_0.jpg}
\end{center}
\hspace*{3em} 

Figure 6.2 Inspection game between the inspector (player I) and inspectee (player II).

\begin{tcolorbox}
图6.2 监管者(玩家I)与被监管者(玩家II)之间的监管博弈。
\end{tcolorbox}

In all cases, player I would strongly prefer if player II complied, but this is outside of player I's control. However, the inspector prefers to inspect if the inspectee cheats (because -6 is better than -10 ), as shown by the best-response box for payoff -6 in Figure 6.2. If the inspector always preferred Don't Inspect, then this would be a dominating strategy and be part of a (unique) equilibrium where the inspectee cheats.

\begin{tcolorbox}
在所有情况下，玩家I非常希望玩家II遵守规定，但这不在玩家I的控制范围内。然而，如图6.2中收益 -6 的最优反应框所示，如果被监管者违规，监管者更倾向于检查(因为 -6 比 -10 好)。如果监管者总是倾向于“不检查”，那么这将是一个占优策略，并且会成为被监管者违规的(唯一)均衡的一部分。
\end{tcolorbox}

The best-response payoffs in Figure 6.2 show that this game has no equilibrium in pure strategies. If any player settles on a deterministic choice (like Don't Inspect by player I), then the best reponse of the other player is unique (here to cheat by player II), to which the original choice is not a best reponse (player I prefers to inspect when the other player cheats, against which player II in turn prefers to comply). The strategies in an equilibrium must be best responses to each other, so in this game this fails for any pair of pure strategies.

\begin{tcolorbox}
图6.2中的最优反应收益表明，这个博弈在纯策略中没有均衡。如果任何一个玩家确定了一个确定性的选择(如玩家I选择“不检查”)，那么另一个玩家的最优反应是唯一的(这里是玩家II选择违规)，而最初的选择并不是对这个最优反应的最优反应(当另一个玩家违规时，玩家I更倾向于检查，而玩家II则更倾向于遵守)。均衡中的策略必须是相互的最优反应，因此在这个博弈中，任何一对纯策略都无法满足这一条件。
\end{tcolorbox}

What should the players do in the game of Figure 6.2? One possibility is that they prepare for the worst, that is, choose a so-called max-min strategy. A max-min strategy maximizes the player's worst payoff against all possible choices of the opponent. The max-min strategy (as a pure strategy) for player \(\mathrm{I}\) is to inspect (where the inspector guarantees himself payoff -6), and for player II it is to comply (which guarantees her a payoff of 0 ). However, this is not an equilibrium and hence not a stable recommendation to the two players, because player I could switch his strategy to Don't Inspect and improve his payoff.

\begin{tcolorbox}
在图6.2的博弈中，玩家应该怎么做呢？一种可能性是他们为最坏的情况做准备，即选择所谓的极大极小策略。极大极小策略是在对手所有可能的选择下，最大化玩家的最差收益。玩家 \(\mathrm{I}\) 的极大极小策略(作为纯策略)是检查(监管者可以保证自己的收益为 -6)，而玩家II的极大极小策略是遵守(这可以保证她的收益为0)。然而，这不是一个均衡，因此也不是给两个玩家的稳定建议，因为玩家I可以将他的策略改为“不检查”来提高他的收益。
\end{tcolorbox}

A mixed strategy of player I in this game is to inspect only with a certain probability. In the context of inspections, randomizing is also a practical approach that reduces costs. Even if an inspection is not certain, a sufficiently high chance of being caught should deter the inspectee from cheating.

\begin{tcolorbox}
在这个博弈中，参与者I的混合策略(mixed strategy)是以一定概率进行检查。在检查的情境下，随机化也是一种降低成本(cost)的实用方法。即使检查不是必然的，被抓住的可能性足够高也应该能阻止被检查者作弊。
\end{tcolorbox}

A mixed strategy of player I is given by a probability \(p\) for Inspect, and thus \(1 - p\) for Don’t Inspect. Instead of his two pure strategies, player I can now choose a probability \(p \in  \left\lbrack  {0,1}\right\rbrack\) . This includes his pure strategies Don’t Inspect when \(p = 0\) and Inspect when \(p = 1\) . A mixed strategy of player II is given by a probability \(q\) for cheat and \(1 - q\) for comply. Her choice between two pure strategies is thereby extended to any choice of \(q \in  \left\lbrack  {0,1}\right\rbrack\) . In effect, this means the players now each have an infinite strategy set \(\left\lbrack  {0,1}\right\rbrack\) with payoffs computed as expected payoffs. This new game is called the mixed extension of the original \(2 \times  2\) game.

\begin{tcolorbox}
参与者I的混合策略(mixed strategy)是为“检查”设定一个概率 \(p\) ，因此“不检查”的概率为 \(1 - p\) 。现在，参与者I可以选择一个概率 \(p \in  \left\lbrack  {0,1}\right\rbrack\) ，而不是他的两种纯策略。这包括当 \(p = 0\) 时的纯策略“不检查”和当 \(p = 1\) 时的纯策略“检查”。参与者II的混合策略(mixed strategy)是为“作弊”设定一个概率 \(q\) ，为“遵守”设定一个概率 \(1 - q\) 。这样，她在两种纯策略之间的选择就扩展到了对 \(q \in  \left\lbrack  {0,1}\right\rbrack\) 的任何选择。实际上，这意味着参与者现在各自有一个无限的策略集 \(\left\lbrack  {0,1}\right\rbrack\) ，其收益按期望效用(expected payoff)计算。这个新的博弈被称为原 \(2 \times  2\) 博弈的混合扩展(mixed extension)。
\end{tcolorbox}

When does the strategy pair \(\left( {p,q}\right)  \in  \left\lbrack  {0,1}\right\rbrack   \times  \left\lbrack  {0,1}\right\rbrack\) define an equilibrium (which will be called a mixed equilibrium)? These strategies have to be best responses to each other. Suppose player I chooses \(p\) and player II has to determine her best response. We first consider her possible pure responses. Her expected payoff for comply is \(\left( {1 - p}\right)  \cdot  0 + p \cdot  0 = 0\) , which is independent of \(p\) . Her expected payoff for cheat is \(\left( {1 - p}\right)  \cdot  {10} + p \cdot  \left( {-{90}}\right)  = {10} - {100p}\) . In effect, player II faces the decision situation in Figure 6.1 with 0.9 and 0.1 replaced by \(1 - p\) and \(p\) . Clearly, she prefers comply when \(0 > {10} - {100p}\) , that is, \(p > {0.1}\) , and cheat when \(0 < {10} - {100p}\) , that is, \(p < {0.1}\) , and she is indifferent if and only if \(0 = {10} - {100p}\) , that is, \(p = {0.1}\) .

\begin{tcolorbox}
策略对 \(\left( {p,q}\right)  \in  \left\lbrack  {0,1}\right\rbrack   \times  \left\lbrack  {0,1}\right\rbrack\) 何时构成一个均衡(equilibrium)(将其称为混合均衡(mixed equilibrium))？这些策略必须是彼此的最优反应(best response)。假设参与者I选择 \(p\) ，参与者II必须确定她的最优反应。我们首先考虑她可能的纯反应。她选择“遵守”的期望效用(expected payoff)是 \(\left( {1 - p}\right)  \cdot  0 + p \cdot  0 = 0\) ，这与 \(p\) 无关。她选择“作弊”的期望效用(expected payoff)是 \(\left( {1 - p}\right)  \cdot  {10} + p \cdot  \left( {-{90}}\right)  = {10} - {100p}\) 。实际上，参与者II面临着图6.1中的决策情形，其中0.9和0.1分别被 \(1 - p\) 和 \(p\) 取代。显然，当 \(0 > {10} - {100p}\) 时，即 \(p > {0.1}\) 时，她更倾向于“遵守”；当 \(0 < {10} - {100p}\) 时，即 \(p < {0.1}\) 时，她更倾向于“作弊”；当且仅当 \(0 = {10} - {100p}\) 时，即 \(p = {0.1}\) 时，她无差异。
\end{tcolorbox}

In the mixed extension, player II can also choose to mix by choosing \(q\) so that \(0 < q < 1\) . We claim that this only happens when she is indifferent. Namely, suppose that \(p\) is very low, for example \(p = {0.01}\) , so that cheat gives an expected payoff of \({0.99} \cdot  {10} + {0.01} \cdot  \left( {-{90}}\right)  = 9\) . With a mixed strategy \(q\) , player II’s expected payoff is \(0 \cdot  \left( {1 - q}\right)  + {9q} = {9q}\) , which is clearly maximized for \(q = 1\) , which is her pure best response cheat. Similarly, if \(p\) is high, for example \(p = {0.2}\) , then cheat gives an expected payoff of \({0.8} \cdot  {10} + {0.2} \cdot  \left( {-{90}}\right)  =  - {10}\) . With a mixed strategy \(q\) , player II’s expected payoff is \(0 \cdot  \left( {1 - q}\right)  - {10q} =  - {10q}\) , which is maximized for \(q = 0\) , which is her pure best response comply. Only when comply and cheat give the same expected payoffs, player II’s expected payoff does not depend on \(q\) and she can in fact choose any \(q \in  \left\lbrack  {0,1}\right\rbrack\) as a best response. Figure 6.3 shows a drawing of the best-response probability \(q\) against \(p\) .

\begin{tcolorbox}
在混合扩展中，玩家 II 也可以通过选择 \(q\) 进行混合，使得 \(0 < q < 1\) 。我们断言，只有当她无差异时才会出现这种情况。也就是说，假设 \(p\) 非常低，例如 \(p = {0.01}\) ，那么作弊的期望效用为 \({0.99} \cdot  {10} + {0.01} \cdot  \left( {-{90}}\right)  = 9\) 。采用混合策略 \(q\) 时，玩家 II 的期望效用为 \(0 \cdot  \left( {1 - q}\right)  + {9q} = {9q}\) ，显然当 \(q = 1\) 时该期望效用达到最大化，这是她的纯最优反应——作弊。类似地，如果 \(p\) 很高，例如 \(p = {0.2}\) ，那么作弊的期望效用为 \({0.8} \cdot  {10} + {0.2} \cdot  \left( {-{90}}\right)  =  - {10}\) 。采用混合策略 \(q\) 时，玩家 II 的期望效用为 \(0 \cdot  \left( {1 - q}\right)  - {10q} =  - {10q}\) ，当 \(q = 0\) 时期望效用达到最大化，这是她的纯最优反应——遵守规则。只有当遵守规则和作弊带来相同的期望效用时，玩家 II 的期望效用才不依赖于 \(q\) ，实际上她可以选择任何 \(q \in  \left\lbrack  {0,1}\right\rbrack\) 作为最优反应。图 6.3 展示了最优反应概率 \(q\) 相对于 \(p\) 的关系图。
\end{tcolorbox}

\begin{center}
\includegraphics[max width=0.4\textwidth]{images/bo_d2e5cmbef24c73b03tn0_150_248_1432_598_587_0.jpg}
\end{center}
\hspace*{3em} 

Figure 6.3 Best-response mixed-strategy probability \(q\) for cheat (black line) of player II in the Inspection game, against the mixed-strategy probability \(p\) for Inspect of player I. The mixed-strategy square is drawn so that its corners match the four cells in Figure 6.2.

\begin{tcolorbox}
图 6.3 监管博弈中玩家 II 作弊的最优反应混合策略概率 \(q\) (黑线)与玩家 I 检查的混合策略概率 \(p\) 的关系图。混合策略方块的绘制使得其角点与图 6.2 中的四个单元格相对应。
\end{tcolorbox}

In summary, the only case where player II can possibly randomize between her pure strategies is if both strategies give her the same expected payoff. As stated and proved formally in Proposition 6.1 below, it is never optimal for a player to assign a positive probability to a pure strategy that is inferior to other pure strategies, for given mixed strategies of the other players.

\begin{tcolorbox}
综上所述，玩家 II 可能在其纯策略之间进行随机选择的唯一情况是，两种策略给她带来相同的期望效用。正如下面命题 6.1 所陈述和正式证明的那样，对于给定的其他玩家的混合策略，玩家对劣于其他纯策略的纯策略赋予正概率永远不是最优的。
\end{tcolorbox}

\begin{center}
\includegraphics[max width=0.9\textwidth]{images/bo_d2e5cmbef24c73b03tn0_151_339_379_1246_589_0.jpg}
\end{center}
\hspace*{3em} 

Figure 6.4 Left: Best-response mixed-strategy probability \(p\) (gray line) of player I against \(q\) . Right: The intersection of the two best-response lines determines the unique mixed equilibrium where \(p = {0.1}\) and \(q = {0.2}\) .

\begin{tcolorbox}
图 6.4 左图:玩家 I 相对于 \(q\) 的最优反应混合策略概率 \(p\) (灰线)。右图:两条最优反应线的交点确定了唯一的混合均衡，其中 \(p = {0.1}\) 且 \(q = {0.2}\) 。
\end{tcolorbox}

Similarly, consider a mixed strategy of player II where she cheats with probability \(q\) . The expected payoffs for player I’s pure strategies are \(- {10q}\) for Don’t Inspect and \(\left( {-1}\right) \left( {1 - q}\right)  - {6q} =  - 1 - {5q}\) for Inspect. If player I plays Inspect with probability \(p\) , then his expected payoff is therefore \(\left( {-{10q}}\right) \left( {1 - p}\right)  + \left( {-1 - {5q}}\right) p =  - {10q} + p\left( {{5q} - 1}\right)\) . This is a linear function of \(p\) , and has its maximum at \(p = 0\) if \({5q} - 1 < 0\) , that is, \(q < {0.2}\) , and at \(p = 1\) if \({5q} - 1 > 0\) , that is, \(q > {0.2}\) , and only if \(q = {0.2}\) does it not depend on \(p\) so that \(p \in  \left\lbrack  {0,1}\right\rbrack\) can then be chosen arbitrarily. This is shown in the left diagram in Figure 6.4.

\begin{tcolorbox}
类似地，考虑玩家 II 的一种混合策略，她以概率 \(q\) 作弊。玩家 I 的纯策略的期望效用分别为:不检查时为 \(- {10q}\) ，检查时为 \(\left( {-1}\right) \left( {1 - q}\right)  - {6q} =  - 1 - {5q}\) 。如果玩家 I 以概率 \(p\) 选择检查，那么他的期望效用为 \(\left( {-{10q}}\right) \left( {1 - p}\right)  + \left( {-1 - {5q}}\right) p =  - {10q} + p\left( {{5q} - 1}\right)\) 。这是关于 \(p\) 的线性函数，若 \({5q} - 1 < 0\) ，即 \(q < {0.2}\) ，则在 \(p = 0\) 处达到最大值；若 \({5q} - 1 > 0\) ，即 \(q > {0.2}\) ，则在 \(p = 1\) 处达到最大值；只有当 \(q = {0.2}\) 时，它才不依赖于 \(p\) ，此时 \(p \in  \left\lbrack  {0,1}\right\rbrack\) 可以任意选择。这在图 6.4 的左图中有所展示。
\end{tcolorbox}

The right diagram in Figure 6.4 shows the intersection of the two best-response lines with the unique mixed equilibrium of the game at \(p = {0.1}\) and \(q = {0.2}\) when the two mixed strategies are mutual best responses. It is similar to the best-response diagram in Figure 3.10 for the Cournot duopoly, except for a quarter clock-wise turn so that the vertical and horizontal parameters \(p\) and \(q\) match the rows and columns in the \(2 \times  2\) game in Figure 6.2.

\begin{tcolorbox}
图6.4中的右图展示了当两个混合策略互为最优反应时，两条最优反应线的交点，该交点对应着博弈在 \(p = {0.1}\) 和 \(q = {0.2}\) 处的唯一混合均衡(mixed equilibrium)。它与图3.10中用于古诺双寡头博弈(Cournot duopoly)的最优反应示意图(best-response diagram)相似，只是顺时针旋转了四分之一圈，以便垂直和水平参数 \(p\) 和 \(q\) 与图6.2中 \(2 \times  2\) 博弈的行和列相匹配。
\end{tcolorbox}

In the unique mixed equilibrium with \(p = {0.1}\) and \(q = {0.2}\) of the Inspection game, the expected payoffs are

\begin{tcolorbox}
在监管博弈(Inspection game)的 \(p = {0.1}\) 和 \(q = {0.2}\) 唯一混合均衡(mixed equilibrium)中，期望效用(expected payoff)为
\end{tcolorbox}

\[
- {10q} =  - 1\left( {1 - q}\right)  - {6q} =  - 2\;\text{ for player }\mathrm{I},
\]

\[
0 = {10}\left( {1 - p}\right)  - {90p} = 0\;\text{ for player II. }
\]

The equilibrium probabilities \(q\) and \(p\) are (independently) determined by these equations. It is noteworthy that they depend on the opponent's payoffs and not on the player's own payoffs (which can therefore be changed as long as the circular preference structure in Figure 6.2 remains intact). For example, one would expect that a more severe penalty than -90 for being caught lowers the probability of cheating in equilibrium. In fact, it does not. What does change is the probability \(p\) of inspection, which is reduced until the inspectee is indifferent.

\begin{tcolorbox}
均衡概率 \(q\) 和 \(p\) 由这些方程(独立地)确定。值得注意的是，它们取决于对手的收益，而非玩家自身的收益(因此，只要图6.2中的循环偏好(preference)结构保持不变，玩家自身的收益就可以改变)。例如，人们可能认为，被抓后处以比 -90 更严厉的惩罚会降低均衡状态下作弊的概率。实际上并非如此。真正改变的是检查的概率 \(p\) ，该概率会降低，直到被检查者感到无差异为止。
\end{tcolorbox}

This dependence on opponent payoffs is the most counter-intuitive property of mixed equilibrium. Moreover, mixing itself seems paradoxical when the player is indifferent in equilibrium. If player II, for example, can equally comply or cheat, why should she randomize? In particular, she could comply and get payoff zero for certain, which is simpler and safer. The answer is that precisely because there is no incentive to choose one strategy over the other, a player can mix, and only in that case can there be an equilibrium. If player II would comply for certain, then the only optimal choice of player \(\mathrm{I}\) is not to inspect, but then complying is not optimal, so this is not an equilibrium.

\begin{tcolorbox}
这种对对手收益的依赖是混合均衡(mixed equilibrium)最违反直觉的特性。此外，当玩家在均衡状态下感到无差异时，混合策略(mixed strategy)本身似乎也自相矛盾。例如，如果玩家 II 可以同样地选择遵守规则或作弊，那她为什么要随机选择呢？特别是，她可以选择遵守规则并确定获得零收益，这样既简单又安全。答案是，正是因为没有动机去选择一种策略而非另一种策略，玩家才可以采用混合策略，而且只有在这种情况下才可能存在均衡。如果玩家 II 确定会遵守规则，那么玩家 \(\mathrm{I}\) 的唯一最优选择就是不进行检查，但这样一来，遵守规则就不是最优选择了，所以这不是一个均衡。
\end{tcolorbox}

\section*{6.3 Bimatrix Games}

\begin{tcolorbox}
\section*{6.3 双矩阵博弈(bimatrix game)}
\end{tcolorbox}

In the following, we discuss mixed equilibria for general games in strategic form. Most ideas can be explained, with much simpler notation, for two-player games. Then the strategies of player I and player II are the rows and columns of a table with two payoffs in each cell, so that the game is defined by two payoff matrices. Some basic facts about matrix manipulation are in the background material.

\begin{tcolorbox}
接下来，我们将讨论策略形式下一般博弈的混合均衡。对于两人博弈，大多数概念可以用更简单的符号来解释。此时，玩家I和玩家II的策略分别对应一个表格的行和列，每个单元格中有两个收益值，因此该博弈由两个收益矩阵定义。有关矩阵运算的一些基本事实可在背景材料中找到。
\end{tcolorbox}

Background material: Matrix multiplication

\begin{tcolorbox}
背景材料:矩阵乘法
\end{tcolorbox}

If \(A\) is an \(m \times  k\) matrix and \(B\) is an \(k \times  n\) matrix, then their matrix product \(A \cdot  B\) (or \({AB}\) for short) is the \(m \times  n\) matrix \(C\) with entries

\begin{tcolorbox}
如果 \(A\) 是一个 \(m \times  k\) 矩阵， \(B\) 是一个 \(k \times  n\) 矩阵，那么它们的矩阵乘积 \(A \cdot  B\) (或简记为 \({AB}\) )是一个 \(m \times  n\) 矩阵 \(C\) ，其元素为
\end{tcolorbox}

\[
{c}_{ij} = \mathop{\sum }\limits_{{s = 1}}^{k}{a}_{is}{b}_{sj}\;\left( {1 \leq  i \leq  m,\;1 \leq  j \leq  n}\right) . \tag{6.1}
\]

When writing down a matrix product \({AB}\) it is assumed to be defined, that is, \(A\) has as many columns as \(B\) has rows. Matrix multiplication is associative, that is, \(\left( {AB}\right) C = A\left( {BC}\right)\) , written as \({ABC}\) , for matrices \(A,B,C\) .

\begin{tcolorbox}
在写下矩阵乘积 \({AB}\) 时，默认该乘积是有定义的，即 \(A\) 的列数等于 \(B\) 的行数。矩阵乘法满足结合律，即对于矩阵 \(A,B,C\) ，有 \(\left( {AB}\right) C = A\left( {BC}\right)\) ，可写成 \({ABC}\) 。
\end{tcolorbox}

When dealing with matrices, we use the following notation throughout. All matrices considered here are real matrices (with entries in \(\mathbb{R}\) ). The set of \(m \times  n\) matrices is \({\mathbb{R}}^{m \times  n}\) . The transpose of an \(m \times  n\) matrix \(A\) is written \({A}^{\top }\) , which is the \(n \times  m\) matrix with entries \({a}_{ji}\) if \({a}_{ij}\) are the entries of \(A\) , for \(1 \leq  i \leq  m,1 \leq  j \leq  n\) .

\begin{tcolorbox}
在处理矩阵时，我们始终使用以下符号。这里考虑的所有矩阵都是实矩阵(元素属于 \(\mathbb{R}\) )。 \(m \times  n\) 矩阵的集合记为 \({\mathbb{R}}^{m \times  n}\) 。 \(m \times  n\) 矩阵 \(A\) 的转置记为 \({A}^{\top }\) ，它是一个 \(n \times  m\) 矩阵，若 \(A\) 的元素为 \({a}_{ij}\) ，则 \({A}^{\top }\) 的元素为 \({a}_{ji}\) ，其中 \(1 \leq  i \leq  m,1 \leq  j \leq  n\) 。
\end{tcolorbox}

All vectors are column vectors, so if \(x \in  {\mathbb{R}}^{n}\) then \(x\) is an \(n \times  1\) matrix and \({x}^{\top }\) is the corresponding row vector in \({\mathbb{R}}^{1 \times  n}\) , and row vectors are easily identified by this transposition sign. The components of a vector \(x \in  {\mathbb{R}}^{n}\) are \({x}_{1},\ldots ,{x}_{n}\) . We write the scalar product of two vectors \(x,y\) in \({\mathbb{R}}^{n}\) as the matrix product \({x}^{\top }y\) . In pictures (with a square for each matrix entry, here for \(n = 4\) ):

\begin{tcolorbox}
所有向量均为列向量，因此如果 \(x \in  {\mathbb{R}}^{n}\) ，那么 \(x\) 是一个 \(n \times  1\) 矩阵， \({x}^{\top }\) 是 \({\mathbb{R}}^{1 \times  n}\) 中对应的行向量，通过这个转置符号可以很容易识别出行向量。向量 \(x \in  {\mathbb{R}}^{n}\) 的分量为 \({x}_{1},\ldots ,{x}_{n}\) 。我们将 \({\mathbb{R}}^{n}\) 中两个向量 \(x,y\) 的标量积写成矩阵乘积 \({x}^{\top }y\) 。用图形表示(每个矩阵元素用一个方块表示，这里是 \(n = 4\) ):
\end{tcolorbox}

(6.2)

\begin{center}
\includegraphics[max width=0.3\textwidth]{images/bo_d2e5cmbef24c73b03tn0_153_753_327_424_155_0.jpg}
\end{center}
\hspace*{3em} 

We use matrix multiplication where possible. In particular, scalars are treated like \(1 \times  1\) matrices. A column vector \(x\) is multiplied with a scalar \(\lambda\) from the right, and a row vector \({x}^{\top }\) is multiplied with a scalar \(\lambda\) from the left,

\begin{tcolorbox}
我们尽可能使用矩阵乘法。特别地，标量被视为 \(1 \times  1\) 矩阵。列向量 \(x\) 从右侧与标量 \(\lambda\) 相乘，行向量 \({x}^{\top }\) 从左侧与标量 \(\lambda\) 相乘。
\end{tcolorbox}

(6.3)

\begin{center}
\includegraphics[max width=0.7\textwidth]{images/bo_d2e5cmbef24c73b03tn0_153_513_657_901_158_0.jpg}
\end{center}
\hspace*{3em} 

For more details see the background material on visualizing matrix multiplication.

\begin{tcolorbox}
更多详细信息请参阅矩阵乘法可视化的背景材料。
\end{tcolorbox}

\(\Rightarrow\) Exercise 6.1 compares \({x}^{\top }y\) and \(x{y}^{\top }\) .

\begin{tcolorbox}
\(\Rightarrow\) 练习6.1比较了 \({x}^{\top }y\) 和 \(x{y}^{\top }\) 。
\end{tcolorbox}

Furthermore, we use the special vectors0and1with all components equal to 0 and 1, respectively, where their dimension depends on the context. Inequalities like \(x \geq  \mathbf{0}\) hold for all components; note that the resulting order \(\leq\) between vectors is only a partial order.

\begin{tcolorbox}
此外，我们使用特殊向量0和1，它们的所有分量分别等于0和1，其维度取决于具体情境。像 \(x \geq  \mathbf{0}\) 这样的不等式对所有分量都成立；注意，向量之间由此产生的序 \(\leq\) 只是一个偏序。
\end{tcolorbox}

Recall that a game in strategic form is specified by a finite set of "pure" strategies for each player, and a payoff for each player for each strategy profile, which is a tuple of strategies, one for each player. The game is played by each player independently and simultaneously choosing one strategy, whereupon the players receive their respective payoffs.

\begin{tcolorbox}
回顾一下，策略形式的博弈由每个玩家的一组有限“纯”策略以及每个策略组合(即每个玩家各选一个策略所构成的元组)下每个玩家的收益来确定。该博弈通过每个玩家独立且同时选择一个策略来进行，然后玩家获得各自的收益。
\end{tcolorbox}

In this chapter we focus on two-player games. Unless specified otherwise, we assume that player I has \(m\) strategies and player II has \(n\) strategies. The \(m\) strategies of the row player I are denoted by \(i = 1,\ldots ,m\) , and the \(n\) strategies of the column player II are denoted by \(j = 1,\ldots ,n\) .

\begin{tcolorbox}
在本章中，我们主要关注两人博弈。除非另有说明，我们假设玩家I有 \(m\) 个策略，玩家II有 \(n\) 个策略。行玩家I的 \(m\) 个策略用 \(i = 1,\ldots ,m\) 表示，列玩家II的 \(n\) 个策略用 \(j = 1,\ldots ,n\) 表示。
\end{tcolorbox}

An \(m \times  n\) two-player game in strategic form is also called a bimatrix game (A, B). Here, \(A\) and \(B\) are the two \(m \times  n\) payoff matrices. The \(m\) rows are the pure strategies \(i\) of player I and the \(n\) columns are the pure strategies \(j\) of player II. For a row \(i\) and column \(j\) , the matrix entry \({a}_{ij}\) of \(A\) is the payoff to player \(\mathrm{I}\) , and the matrix entry \({b}_{ij}\) of \(B\) is the payoff to player II. The staggered payoffs in a table such as Figure 3.1 for the Prisoner’s Dilemma show \(A\) as the matrix of lower-left payoffs in each cell, and \(B\) as the matrix of upper-right payoffs in each cell, here

\begin{tcolorbox}
一个处于策略形式(strategic form)的 \(m \times  n\) 两人博弈也被称为双矩阵博弈(bimatrix game)(A, B)。这里， \(A\) 和 \(B\) 是两个 \(m \times  n\) 收益矩阵。 \(m\) 行是玩家 I 的纯策略 \(i\) ， \(n\) 列是玩家 II 的纯策略 \(j\) 。对于第 \(i\) 行和第 \(j\) 列， \(A\) 的矩阵元素 \({a}_{ij}\) 是玩家 \(\mathrm{I}\) 的收益， \(B\) 的矩阵元素 \({b}_{ij}\) 是玩家 II 的收益。像图 3.1 中囚徒困境(Prisoner’s Dilemma)表格里的交错收益(staggered payoffs)显示， \(A\) 是每个单元格左下角收益的矩阵， \(B\) 是每个单元格右上角收益的矩阵，这里
\end{tcolorbox}

\[
A = \left( \begin{array}{ll} 2 & 0 \\  3 & 1 \end{array}\right) ,\;B = \left( \begin{array}{ll} 2 & 3 \\  0 & 1 \end{array}\right) .
\]

This bimatrix game is symmetric because \(B = {A}^{\top }\) .

\begin{tcolorbox}
这个双矩阵博弈是对称的(symmetric)，因为 \(B = {A}^{\top }\) 。
\end{tcolorbox}

Background material: Visualizing matrix multiplication

\begin{tcolorbox}
背景材料(background material):矩阵乘法的可视化(visualizing matrix multiplication)
\end{tcolorbox}

In a matrix product \(C = {AB}\) , the entry \({c}_{ij}\) of \(C\) in row \(i\) and column \(j\) is the scalar product of the \(i\) th row of \(A\) with the \(j\) th column of \(B\) , according to (6.2) and (6.1). There are two useful ways of visualizing the product \({AB}\) . Let \(A \in  {\mathbb{R}}^{m \times  k}\) and \(x \in  {\mathbb{R}}^{k}\) . Then \({Ax} \in  {\mathbb{R}}^{m}\) . Let \(A = \left\lbrack  {{A}_{1}\cdots {A}_{k}}\right\rbrack\) , that is, \({A}_{1},\ldots ,{A}_{k}\) are the columns of \(A\) . Then \({Ax} = {A}_{1}{x}_{1} + \cdots  + {A}_{k}{x}_{k}\) , that is, \({Ax}\) is the linear combination of the columns of \(A\) with the components \({x}_{1},\ldots ,{x}_{k}\) of \(x\) as coefficients. In pictures for \(m = 3,k = 4\) :

\begin{tcolorbox}
在矩阵乘积 \(C = {AB}\) 中，根据(6.2)和(6.1)， \(C\) 中位于第 \(i\) 行和第 \(j\) 列的元素 \({c}_{ij}\) 是 \(A\) 的第 \(i\) 行与 \(B\) 的第 \(j\) 列的标量积(scalar product)。有两种直观理解乘积 \({AB}\) 的有用方法。设 \(A \in  {\mathbb{R}}^{m \times  k}\) 和 \(x \in  {\mathbb{R}}^{k}\) 。那么 \({Ax} \in  {\mathbb{R}}^{m}\) 。设 \(A = \left\lbrack  {{A}_{1}\cdots {A}_{k}}\right\rbrack\) ，即 \({A}_{1},\ldots ,{A}_{k}\) 是 \(A\) 的列。那么 \({Ax} = {A}_{1}{x}_{1} + \cdots  + {A}_{k}{x}_{k}\) ，即 \({Ax}\) 是以 \(x\) 的分量 \({x}_{1},\ldots ,{x}_{k}\) 为系数的 \(A\) 的列的线性组合。对于 \(m = 3,k = 4\) 的图示如下:
\end{tcolorbox}

(6.4)

\begin{center}
\includegraphics[max width=0.4\textwidth]{images/bo_d2e5cmbef24c73b03tn0_154_517_623_616_165_0.jpg}
\end{center}
\hspace*{3em} 

Let \(B \in  {\mathbb{R}}^{k \times  n}\) and \(B = \left\lbrack  {{B}_{1}\cdots {B}_{n}}\right\rbrack\) . Then the \(j\) th column of \({AB}\) is the linear combination \(A{B}_{j}\) of the columns of \(A\) with the components of \({B}_{j}\) as coefficients. We can visualize the columns of \({AB}\) as follows (for \(n = 2\) ):

\begin{tcolorbox}
设 \(B \in  {\mathbb{R}}^{k \times  n}\) 和 \(B = \left\lbrack  {{B}_{1}\cdots {B}_{n}}\right\rbrack\) 。那么 \({AB}\) 的第 \(j\) 列是以 \({B}_{j}\) 的分量为系数的 \(A\) 的列的线性组合 \(A{B}_{j}\) 。我们可以如下直观理解 \({AB}\) 的列(对于 \(n = 2\) ):
\end{tcolorbox}

\begin{center}
\includegraphics[max width=0.6\textwidth]{images/bo_d2e5cmbef24c73b03tn0_154_395_970_859_155_0.jpg}
\end{center}
\hspace*{3em} 

The second view of \({AB}\) is the same but using rows. Let \(y \in  {\mathbb{R}}^{k}\) , so that \({y}^{\top }B\) is a row vector in \({\mathbb{R}}^{1 \times  n}\) , given as the linear combination \({y}_{1}{b}_{1}^{\top } + \cdots  + {y}_{k}{b}_{k}^{\top }\) of the rows \({b}_{1}^{\top },\ldots ,{b}_{k}^{\top }\) of \(B\) (which we can write as \({B}^{\top } = \left\lbrack  {{b}_{1}\cdots {b}_{k}}\right\rbrack\) ),

\begin{tcolorbox}
对 \({AB}\) 的第二种理解方式相同，但使用行。设 \(y \in  {\mathbb{R}}^{k}\) ，使得 \({y}^{\top }B\) 是 \({\mathbb{R}}^{1 \times  n}\) 中的一个行向量，它表示为 \(B\) 的行 \({b}_{1}^{\top },\ldots ,{b}_{k}^{\top }\) 的线性组合 \({y}_{1}{b}_{1}^{\top } + \cdots  + {y}_{k}{b}_{k}^{\top }\) (我们可以写成 \({B}^{\top } = \left\lbrack  {{b}_{1}\cdots {b}_{k}}\right\rbrack\) )，
\end{tcolorbox}

(6.5)

\begin{center}
\includegraphics[max width=0.5\textwidth]{images/bo_d2e5cmbef24c73b03tn0_154_456_1305_740_171_0.jpg}
\end{center}
\hspace*{3em} 

Let \({A}^{\top } = \left\lbrack  {{a}_{1}\cdots {a}_{m}}\right\rbrack\) . Then the \(i\) th row of \({AB}\) is the linear combination \({a}_{i}^{\top }B\) of the rows of \(B\) with the components of the \(i\) th row \({a}_{i}^{\top }\) of \(A\) as coefficients:

\begin{tcolorbox}
设 \({A}^{\top } = \left\lbrack  {{a}_{1}\cdots {a}_{m}}\right\rbrack\) 。那么 \({AB}\) 的第 \(i\) 行是以 \(A\) 的第 \(i\) 行 \({a}_{i}^{\top }\) 的分量为系数的 \(B\) 的行的线性组合 \({a}_{i}^{\top }B\) :
\end{tcolorbox}

\begin{center}
\includegraphics[max width=0.6\textwidth]{images/bo_d2e5cmbef24c73b03tn0_154_398_1610_856_155_0.jpg}
\end{center}
\hspace*{3em} 

It is useful to acquire some routine in these manipulations. Most common is multiplication of a matrix with a column vector from the right as in (6.4) or with a row vector from the left as in (6.5).

\begin{tcolorbox}
熟悉这些操作是很有用的。最常见的是矩阵与列向量右乘，如式(6.4)所示；或者矩阵与行向量左乘，如式(6.5)所示。
\end{tcolorbox}

A mixed strategy is a randomized strategy of a player. It is defined as a lottery (probability distribution) on the set of pure strategies of that player. This is played as an "active randomization", where the player selects each pure strategy according to its probability. When the other player considers a best response to such a mixed strategy, she is assumed to know the probabilities but not the outcome of the random choice, and bases her decision on the resulting expected payoffs.

\begin{tcolorbox}
混合策略(mixed strategy)是玩家的一种随机策略。它被定义为该玩家纯策略集上的一种彩票(概率分布)。这是一种“主动随机化”玩法，玩家根据概率选择每种纯策略。当另一位玩家考虑对这种混合策略的最优反应(best response)时，假定她知道概率但不知道随机选择的结果，并根据由此产生的期望效用(expected payoff)来做出决策。
\end{tcolorbox}

A pure strategy is considered as a special mixed strategy, which chooses the pure strategy with probability one (in the same way as a deterministic outcome is considered a "random" outcome that occurs with probability one).

\begin{tcolorbox}
纯策略被视为一种特殊的混合策略，它以概率1选择该纯策略(就像确定性结果被视为以概率1发生的“随机”结果一样)。
\end{tcolorbox}

A mixed strategy is determined by the probabilities that it assigns to the player’s pure strategies. For player I, a mixed strategy \(x\) is therefore identified with the \(m\) -tuple of probabilities \({x}_{1},{x}_{2},\ldots ,{x}_{m}\) that it assigns to the pure strategies \(1,2,\ldots ,m\) of player I, which we write as the (column) vector \(x \in  {\mathbb{R}}^{m}\) . Similarly, a mixed strategy \(y\) of player II is an \(n\) -tuple of probabilities \({y}_{1},\ldots ,{y}_{n}\) for playing the pure strategies \(j = 1,\ldots ,n\) , written as \(y \in  {\mathbb{R}}^{n}\) .

\begin{tcolorbox}
混合策略由它分配给玩家纯策略的概率决定。对于玩家I，混合策略 \(x\) 因此由它分配给玩家I的纯策略 \(1,2,\ldots ,m\) 的 \(m\) 维(dimension)概率元组 \({x}_{1},{x}_{2},\ldots ,{x}_{m}\) 确定，我们将其写为(列)向量 \(x \in  {\mathbb{R}}^{m}\) 。类似地，玩家II的混合策略 \(y\) 是玩纯策略 \(j = 1,\ldots ,n\) 的 \(n\) 维概率元组 \({y}_{1},\ldots ,{y}_{n}\) ，写为 \(y \in  {\mathbb{R}}^{n}\) 。
\end{tcolorbox}

The sets of mixed strategies of player I and player II are denoted by \(X\) and \(Y\) . The probabilities \({x}_{i}\) in a mixed strategy \(x\) are nonnegative and sum to 1, where we can write the conditions \({x}_{i} \geq  0\) for \(1 \leq  i \leq  m\) and \({x}_{1} + \cdots  + {x}_{m} = 1\) as \(x \geq  \mathbf{0}\) and \({\mathbf{1}}^{\top }x = 1\) , where in this context \(\mathbf{1} \in  {\mathbb{R}}^{m}\) to match the dimension of \(x\) ; when we similarly write \({y}_{1} + \cdots  + {y}_{n} = 1\) as \({\mathbf{1}}^{\top }y = 1\) then \(\mathbf{1} \in  {\mathbb{R}}^{n}\) . That is,

\begin{tcolorbox}
玩家I和玩家II的混合策略集分别用 \(X\) 和 \(Y\) 表示。混合策略 \(x\) 中的概率 \({x}_{i}\) 是非负的，且总和为1，我们可以将 \(1 \leq  i \leq  m\) 和 \({x}_{1} + \cdots  + {x}_{m} = 1\) 的条件 \({x}_{i} \geq  0\) 写为 \(x \geq  \mathbf{0}\) 和 \({\mathbf{1}}^{\top }x = 1\) ，在这种情况下 \(\mathbf{1} \in  {\mathbb{R}}^{m}\) 以匹配 \(x\) 的维度；当我们类似地将 \({y}_{1} + \cdots  + {y}_{n} = 1\) 写为 \({\mathbf{1}}^{\top }y = 1\) 时，则 \(\mathbf{1} \in  {\mathbb{R}}^{n}\) 。即，
\end{tcolorbox}

\[
X = \left\{  {x \in  {\mathbb{R}}^{m} \mid  x \geq  \mathbf{0},{\mathbf{1}}^{\top }x = 1}\right\}  ,\;Y = \left\{  {y \in  {\mathbb{R}}^{n} \mid  y \geq  \mathbf{0},{\mathbf{1}}^{\top }y = 1}\right\}  . \tag{6.6}
\]

For \(x \in  X\) and \(y \in  Y\) , the expected payoff to player \(\mathrm{I}\) is \({x}^{\top }{Ay}\) , and the expected payoff to player II is \({x}^{\top }{By}\) . It is best to think of \({x}^{\top }{Ay}\) as \({x}^{\top }\left( {Ay}\right)\) , that is, as the scalar product of \(x\) with \({Ay}\) . For each row \(i\) , the entry of \({Ay}\) in row \(i\) is \({\left( Ay\right) }_{i}\) , given by

\begin{tcolorbox}
对于 \(x \in  X\) 和 \(y \in  Y\) ，玩家 \(\mathrm{I}\) 的期望效用是 \({x}^{\top }{Ay}\) ，玩家II的期望效用是 \({x}^{\top }{By}\) 。最好将 \({x}^{\top }{Ay}\) 看作 \({x}^{\top }\left( {Ay}\right)\) ，即 \(x\) 与 \({Ay}\) 的标量积(scalar product)。对于每一行 \(i\) ， \({Ay}\) 中第 \(i\) 行的元素是 \({\left( Ay\right) }_{i}\) ，由下式给出
\end{tcolorbox}

\[
{\left( Ay\right) }_{i} = \mathop{\sum }\limits_{{j = 1}}^{n}{a}_{ij}{y}_{j}\;\text{ for }1 \leq  i \leq  m. \tag{6.7}
\]

That is, \({\left( Ay\right) }_{i}\) is the expected payoff to player \(I\) when he plays row \(i\) . Therefore, \({x}^{\top }{Ay}\) is the expected payoff to player I when the players use \(x\) and \(y\) , because

\begin{tcolorbox}
即， \({\left( Ay\right) }_{i}\) 是玩家 \(I\) 选择第 \(i\) 行时的期望效用。因此， \({x}^{\top }{Ay}\) 是当玩家使用 \(x\) 和 \(y\) 时玩家I的期望效用，因为
\end{tcolorbox}

\[
{x}^{\top }\left( {Ay}\right)  = \mathop{\sum }\limits_{{i = 1}}^{m}{x}_{i}{\left( Ay\right) }_{i} = \mathop{\sum }\limits_{{i = 1}}^{m}{x}_{i}\mathop{\sum }\limits_{{j = 1}}^{n}{a}_{ij}{y}_{j} = \mathop{\sum }\limits_{{i = 1}}^{m}\mathop{\sum }\limits_{{j = 1}}^{n}{x}_{i}{a}_{ij}{y}_{j}. \tag{6.8}
\]

Because the players choose their pure strategies \(i\) and \(j\) independently, the probability that they choose the pure strategy pair(i, j)is the product \({x}_{i}{y}_{j}\) of these probabilities, which is the coefficient of the payoff \({a}_{ij}\) in (6.8).

\begin{tcolorbox}
由于玩家独立地选择他们的纯策略 \(i\) 和 \(j\) ，他们选择纯策略对(i, j)的概率是这些概率的乘积 \({x}_{i}{y}_{j}\) ，这就是式(6.8)中收益 \({a}_{ij}\) 的系数。
\end{tcolorbox}

Analogously, \({x}^{\top }{By}\) is the expected payoff to player II when the players use the mixed strategies \(x\) and \(y\) . Here, it is best to read this as \(\left( {{x}^{\top }B}\right) y\) . The row vector \({x}^{\top }B\) is the vector of expected payoffs \({\left( {x}^{\top }B\right) }_{j}\) to player II for her columns \(j\) . Multiplied with the column probabilities \({y}_{j}\) , the sum over all columns \(j\) is the expected payoff to player II, which in analogy to (6.8) is given by

\begin{tcolorbox}
类似地， \({x}^{\top }{By}\) 是当玩家使用混合策略 \(x\) 和 \(y\) 时玩家 II 的期望效用。在这里，最好将其理解为 \(\left( {{x}^{\top }B}\right) y\) 。行向量 \({x}^{\top }B\) 是玩家 II 对于她的列 \(j\) 的期望效用 \({\left( {x}^{\top }B\right) }_{j}\) 的向量。与列概率 \({y}_{j}\) 相乘后，所有列 \(j\) 的总和就是玩家 II 的期望效用，类似于 (6.8)，其由下式给出
\end{tcolorbox}

\[
\left( {{x}^{\top }B}\right) y = \mathop{\sum }\limits_{{j = 1}}^{n}{\left( {x}^{\top }B\right) }_{j}{y}_{j} = \mathop{\sum }\limits_{{j = 1}}^{n}\left( {\mathop{\sum }\limits_{{i = 1}}^{m}{x}_{i}{b}_{ij}}\right) {y}_{j} = \mathop{\sum }\limits_{{j = 1}}^{n}\mathop{\sum }\limits_{{i = 1}}^{m}{x}_{i}{b}_{ij}{y}_{j}.
\]

\section*{6.4 The Best-Response Condition}

\begin{tcolorbox}
\section*{6.4 最优反应条件}
\end{tcolorbox}

A mixed-strategy equilibrium is a profile of mixed strategies such that no player can improve his expected payoff by unilaterally changing his own strategy. In a two-player game, an equilibrium is a pair(x, y)of mixed strategies such that \(x\) is a best response to \(y\) and vice versa. It does not seem easy to decide if \(x\) is a best response to \(y\) among all possible mixed strategies, that is, if \(x\) maximizes \({x}^{\top }{Ay}\) for all \(x\) in \(X\) , because \(X\) is an infinite set. However, the following proposition, known as the best-response condition, shows how to recognise this. This proposition is not difficult but important. We discuss it afterwards.

\begin{tcolorbox}
混合策略均衡是一种混合策略组合，使得没有玩家可以通过单方面改变自己的策略来提高其期望效用。在两人博弈中，均衡是一对混合策略 (x, y)，使得 \(x\) 是对 \(y\) 的最优反应，反之亦然。要判断 \(x\) 是否是所有可能的混合策略中对 \(y\) 的最优反应，即 \(x\) 是否能使 \({x}^{\top }{Ay}\) 对于 \(X\) 中的所有 \(x\) 都达到最大，似乎并不容易，因为 \(X\) 是一个无限集。然而，下面这个被称为最优反应条件的命题展示了如何识别这一点。这个命题并不难但很重要。我们随后会讨论它。
\end{tcolorbox}

Proposition 6.1 (Best-response condition, two players). Consider an \(m \times  n\) bimatrix game(A, B)and mixed strategies \(x \in  X\) and \(y \in  Y\) . Let

\begin{tcolorbox}
命题 6.1(两人博弈的最优反应条件)。考虑一个 \(m \times  n\) 双矩阵博弈 (A, B) 以及混合策略 \(x \in  X\) 和 \(y \in  Y\) 。设
\end{tcolorbox}

\[
u = \max \left\{  {{\left( Ay\right) }_{i} \mid  1 \leq  i \leq  m}\right\}  ,\;v = \max \left\{  {{\left( {x}^{\top }B\right) }_{j} \mid  1 \leq  j \leq  n}\right\}  . \tag{6.9}
\]

Then

\begin{tcolorbox}
那么
\end{tcolorbox}

\[
\left( {\forall \bar{x} \in  X : {x}^{\top }{Ay} \geq  {\bar{x}}^{\top }{Ay}}\right) \; \Leftrightarrow  \;{x}^{\top }{Ay} = u, \tag{6.10}
\]

that is, \(x\) is a best response to \(y\) if and only if \({x}^{\top }{Ay} = u\) . This is equivalent to

\begin{tcolorbox}
也就是说， \(x\) 是对 \(y\) 的最优反应，当且仅当 \({x}^{\top }{Ay} = u\) 。这等价于
\end{tcolorbox}

\[
{x}_{i} > 0 \Rightarrow  {\left( Ay\right) }_{i} = u\;\left( {1 \leq  i \leq  m}\right) . \tag{6.11}
\]

Similarly, \(y\) is a best response to \(x\) if and only if \({x}^{\top }{By} = v\) , which is equivalent to

\begin{tcolorbox}
类似地， \(y\) 是对 \(x\) 的最优反应，当且仅当 \({x}^{\top }{By} = v\) ，这等价于
\end{tcolorbox}

\[
{y}_{j} > 0 \Rightarrow  {\left( {x}^{\top }B\right) }_{j} = v\;\left( {1 \leq  j \leq  n}\right) . \tag{6.12}
\]

Proof. Recall that \({\left( Ay\right) }_{i}\) is the \(i\) th component of \({Ay}\) , which is the expected payoff to player I when he plays row \(i\) , as in (6.7). The maximum of these expected payoffs for the pure strategies of player \(\mathrm{I}\) is \(u\) . Then

\begin{tcolorbox}
证明。回顾 \({\left( Ay\right) }_{i}\) 是 \({Ay}\) 的第 \(i\) 个分量，如 (6.7) 所示，它是玩家 I 选择第 \(i\) 行时的期望效用。玩家 \(\mathrm{I}\) 的纯策略的这些期望效用的最大值是 \(u\) 。那么
\end{tcolorbox}

\[
{x}^{\top }{Ay} = \mathop{\sum }\limits_{{i = 1}}^{m}{x}_{i}{\left( Ay\right) }_{i} = \mathop{\sum }\limits_{{i = 1}}^{m}{x}_{i}\left( {u - \left( {u - {\left( Ay\right) }_{i}}\right) }\right.
\]

\[
= \mathop{\sum }\limits_{{i = 1}}^{m}{x}_{i}u - \mathop{\sum }\limits_{{i = 1}}^{m}{x}_{i}\left( {u - {\left( Ay\right) }_{i}}\right)
\]

\[
= u - \mathop{\sum }\limits_{{i = 1}}^{m}{x}_{i}\left( {u - {\left( Ay\right) }_{i}}\right) . \tag{6.13}
\]

Because \({x}_{i} \geq  0\) and \(u - {\left( Ay\right) }_{i} \geq  0\) for any row \(i\) , we have \(\mathop{\sum }\limits_{{i = 1}}^{m}{x}_{i}\left( {u - {\left( Ay\right) }_{i}}\right)  \geq  0\) and therefore \({x}^{\top }{Ay} \leq  u\) . The expected payoff \({x}^{\top }{Ay}\) achieves the maximum \(u\) if and only if \(\mathop{\sum }\limits_{{i = 1}}^{m}{x}_{i}\left( {u - {\left( Ay\right) }_{i}}\right)  = 0\) , that is, if \({x}_{i} > 0\) implies \({\left( Ay\right) }_{i} = u\) , which shows (6.11). Condition (6.12) is shown analogously.

\begin{tcolorbox}
因为对于任何行 \(i\) 都有 \({x}_{i} \geq  0\) 和 \(u - {\left( Ay\right) }_{i} \geq  0\) ，我们有 \(\mathop{\sum }\limits_{{i = 1}}^{m}{x}_{i}\left( {u - {\left( Ay\right) }_{i}}\right)  \geq  0\) ，因此 \({x}^{\top }{Ay} \leq  u\) 。期望效用 \({x}^{\top }{Ay}\) 达到最大值 \(u\) 当且仅当 \(\mathop{\sum }\limits_{{i = 1}}^{m}{x}_{i}\left( {u - {\left( Ay\right) }_{i}}\right)  = 0\) ，即如果 \({x}_{i} > 0\) 意味着 \({\left( Ay\right) }_{i} = u\) ，这就证明了 (6.11)。条件 (6.12) 的证明类似。
\end{tcolorbox}

We repeat this argument using vector notation. First, the scalar product of two nonnegative vectors is nonnegative: for \(x,w \in  {\mathbb{R}}^{m}\) ,

\begin{tcolorbox}
我们用向量符号重复这个论证。首先，两个非负向量的标量积是非负的:对于 \(x,w \in  {\mathbb{R}}^{m}\) ，
\end{tcolorbox}

\[
x \geq  \mathbf{0},\;w \geq  \mathbf{0}\; \Rightarrow  \;{x}^{\top }w \geq  0. \tag{6.14}
\]

Furthermore,(6.9) implies \(\mathbf{1}u \geq  {Ay}\) , that is, \(w = \mathbf{1}u - {Ay} \geq  \mathbf{0}\) . Hence,

\begin{tcolorbox}
此外，(6.9) 意味着 \(\mathbf{1}u \geq  {Ay}\) ，即 \(w = \mathbf{1}u - {Ay} \geq  \mathbf{0}\) 。因此，
\end{tcolorbox}

\[
{x}^{\top }{Ay} = {x}^{\top }\left( {\mathbf{1}u - \left( {\mathbf{1}u - {Ay}}\right) }\right)  = {x}^{\top }\mathbf{1}u - {x}^{\top }\left( {\mathbf{1}u - {Ay}}\right)
\]

\[
= u - {x}^{\top }\left( {\mathbf{1}u - {Ay}}\right) \; \leq  u \tag{6.15}
\]

(note \({x}^{\top }\mathbf{1}u = \left( {{x}^{\top }\mathbf{1}}\right) u = 1 \cdot  u = u\) , which shows the advantage of using matrix products consistently). Second, (6.14) can be extended to

\begin{tcolorbox}
(注意 \({x}^{\top }\mathbf{1}u = \left( {{x}^{\top }\mathbf{1}}\right) u = 1 \cdot  u = u\) ，这显示了始终使用矩阵乘积的优势)。其次，(6.14) 可以扩展到
\end{tcolorbox}

\[
x \geq  \mathbf{0},\;w \geq  \mathbf{0}\; \Rightarrow  \;\left( {{x}^{\top }w = 0\; \Leftrightarrow  \;{x}_{i}{w}_{i} = 0\;\text{ for }1 \leq  i \leq  m}\right)
\]

so that \({x}^{\top }{Ay} = u\) in (6.15) if and only if \({x}_{i} = 0\) or \(u = {\left( Ay\right) }_{i}\) for \(1 \leq  i \leq  m\) , which says the same as (6.11).

\begin{tcolorbox}
因此，当且仅当对于 \(1 \leq  i \leq  m\) 有 \({x}_{i} = 0\) 或 \(u = {\left( Ay\right) }_{i}\) 时，(6.15) 式中的 \({x}^{\top }{Ay} = u\) 成立，这与 (6.11) 式表述相同。
\end{tcolorbox}

The mixed strategy \(x\) is a best response to \(y\) if \(x\) gives maximum expected payoff \({x}^{\top }{Ay}\) to player I among all mixed strategies in \(X\) , as in the left statement in (6.10). In contrast, the payoff \(u\) in (6.9) is the maximum expected payoff \({\left( Ay\right) }_{i}\) among all pure strategies \(i\) of player I. According to (6.10), these payoffs are the same, so a player cannot improve his payoff by mixing. This is intuitive because mixing creates a weighted average \({x}^{\top }\left( {Ay}\right)\) of the expected payoffs \({\left( Ay\right) }_{i}\) with the probabilities \({x}_{i}\) as weights as in (6.8).

\begin{tcolorbox}
如果混合策略 \(x\) 在 \(X\) 中的所有混合策略里，能为参与者 I 带来最大期望效用 \({x}^{\top }{Ay}\) ，如 (6.10) 式左侧表述的那样，那么 \(x\) 就是对 \(y\) 的最优反应。相比之下，(6.9) 式中的收益 \(u\) 是参与者 I 在所有纯策略 \(i\) 中的最大期望效用 \({\left( Ay\right) }_{i}\) 。根据 (6.10) 式，这些收益是相同的，所以参与者无法通过混合策略来提高其收益。这是符合直觉的，因为混合策略会如 (6.8) 式那样，以概率 \({x}_{i}\) 为权重，对期望效用 \({\left( Ay\right) }_{i}\) 进行加权平均 \({x}^{\top }\left( {Ay}\right)\) 。
\end{tcolorbox}

Furthermore,(6.11) states that \(x\) is a best response to \(y\) if and only if \(x\) only chooses pure best responses with positive probability \({x}_{i}\) . The condition whether a pure strategy is a best response is easy to check, because one only has to compute the \(m\) expected payoffs \({\left( Ay\right) }_{i}\) for \(i = 1,\ldots ,m\) . For example, if player I has three pure strategies \(\left( {m = 3}\right)\) , and the expected payoffs in (6.7) are \({\left( Ay\right) }_{1} = 4,{\left( Ay\right) }_{2} = 4\) , and \({\left( Ay\right) }_{3} = 3\) , then only the first two strategies are pure best responses. If these expected payoffs are 3, 5, and 3, then only the second strategy is a best response. Clearly, at least one pure best response exists, because the numbers \({\left( Ay\right) }_{i}\) in (6.11) take their maximum \(u\) for at least one \(i\) .

\begin{tcolorbox}
此外，(6.11) 式表明，当且仅当 \(x\) 仅以正概率 \({x}_{i}\) 选择纯最优反应时， \(x\) 才是对 \(y\) 的最优反应。判断一个纯策略是否为最优反应的条件很容易验证，因为只需要计算 \(i = 1,\ldots ,m\) 的 \(m\) 个期望效用 \({\left( Ay\right) }_{i}\) 即可。例如，如果参与者 I 有三个纯策略 \(\left( {m = 3}\right)\) ，且 (6.7) 式中的期望效用分别为 \({\left( Ay\right) }_{1} = 4,{\left( Ay\right) }_{2} = 4\) 和 \({\left( Ay\right) }_{3} = 3\) ，那么只有前两个策略是纯最优反应。如果这些期望效用分别为 3、5 和 3，那么只有第二个策略是最优反应。显然，至少存在一个纯最优反应，因为 (6.11) 式中的数字 \({\left( Ay\right) }_{i}\) 至少会在一个 \(i\) 处取得最大值 \(u\) 。
\end{tcolorbox}

The best-response condition is the central tool for finding mixed-strategy equilibria, which we describe here for an example; a related approach is described in Section 6.6. Consider the following \(2 \times  3\) game(A, B), where we already marked the best-response payoffs and give the two payoff matrices (as determined by the table):

\begin{tcolorbox}
最优反应条件是寻找混合策略均衡的核心工具，我们在此通过一个例子来进行说明；相关方法将在 6.6 节介绍。考虑以下 \(2 \times  3\) 博弈 (A, B)，我们已经标记出了最优反应收益，并给出了两个收益矩阵(由表格确定):
\end{tcolorbox}

\[
A = \left( \begin{array}{lll} 0 & 4 & 0 \\  2 & 2 & 1 \end{array}\right) ,\;B = \left( \begin{array}{lll} 3 & 0 & 2 \\  1 & 6 & 4 \end{array}\right) . \tag{6.16}
\]

\begin{center}
\includegraphics[max width=0.3\textwidth]{images/bo_d2e5cmbef24c73b03tn0_157_359_1713_500_353_0.jpg}
\end{center}
\hspace*{3em} 

This game has no equilibrium in pure strategies. For any pure strategy, its unique best response is a pure strategy of the other player, which does not result in an equilibrium, so that both players have to play mixed strategies that are not pure. We try to find a mixed equilibrium(x, y), were \(x = {\left( {x}_{1},{x}_{2}\right) }^{\top }\) and \(y = {\left( {y}_{1},{y}_{2},{y}_{3}\right) }^{\top }\) . (When writing down specific vectors, we will from now on omit the transposition sign for brevity and just write \(x = \left( {{x}_{1},{x}_{2}}\right)\) , for example.)

\begin{tcolorbox}
这个博弈不存在纯策略均衡。对于任何纯策略，其唯一的最优反应是另一个参与者的纯策略，这不会导致均衡，因此两个参与者都必须采用非纯的混合策略。我们尝试寻找一个混合均衡 (x, y)，其中 \(x = {\left( {x}_{1},{x}_{2}\right) }^{\top }\) 且 \(y = {\left( {y}_{1},{y}_{2},{y}_{3}\right) }^{\top }\) 。(为了简洁起见，从现在起，在写下具体向量时，我们将省略转置符号，例如只写 \(x = \left( {{x}_{1},{x}_{2}}\right)\) 。)
\end{tcolorbox}

Because player I has only two pure strategies, both have to be played with positive probability, that is, \({x}_{1} > 0\) and \({x}_{2} > 0\) . According to the best-response condition (6.11), player I has to be indifferent between his two rows, that is,

\begin{tcolorbox}
由于参与者 I 只有两个纯策略，这两个策略都必须以正概率被采用，即 \({x}_{1} > 0\) 且 \({x}_{2} > 0\) 。根据最优反应条件 (6.11)，参与者 I 必须对他的两行策略无差异，即
\end{tcolorbox}

\[
{\left( Ay\right) }_{1} = {\left( Ay\right) }_{2} = u,\;4{y}_{2} = 2{y}_{1} + 2{y}_{2} + {y}_{3} = u,\;2{y}_{2} = 2{y}_{1} + {y}_{3}. \tag{6.17}
\]

We have already seen that this cannot be achieved when \(y\) is a pure strategy, so we now consider any pair of strategies that are played with positive probability under \(y\) . Suppose first that these are the first two columns, that is, \({y}_{1} > 0\) and \({y}_{2} > 0\) and \({y}_{3} = 0\) . Then the unique solution to (6.17) (and \({y}_{1} + {y}_{2} = 1\) ) is \({y}_{1} = {y}_{2} = \frac{1}{2}\) with \(u = 2\) . In turn, by (6.12) player II has to be indifferent between her first two columns, that is,

\begin{tcolorbox}
我们已经看到，当 \(y\) 是纯策略时，这无法实现，因此我们现在考虑在 \(y\) 下以正概率采用的任意一对策略。首先假设这些策略是前两列，即 \({y}_{1} > 0\) 、 \({y}_{2} > 0\) 和 \({y}_{3} = 0\) 。那么，(6.17)(以及 \({y}_{1} + {y}_{2} = 1\) )的唯一解是 \({y}_{1} = {y}_{2} = \frac{1}{2}\) ，其中 \(u = 2\) 。继而，根据(6.12)，玩家 II 必须对她的前两列无差异，即
\end{tcolorbox}

\[
{\left( {x}^{\top }B\right) }_{1} = {\left( {x}^{\top }B\right) }_{2} = v,\;3{x}_{1} + {x}_{2} = 6{x}_{2},\;3{x}_{1} = 5{x}_{2} \tag{6.18}
\]

which has the unique solution \({x}_{1} = \frac{5}{8},{x}_{2} = \frac{3}{8}\) and \({\left( {x}^{\top }B\right) }_{1} = {\left( {x}^{\top }B\right) }_{2} = \frac{18}{8} = \frac{9}{4}\) . However, we also have to make sure that \(\frac{9}{4}\) is the best-response payoff to player II, a consideration that is not necessary in \(2 \times  2\) games where it suffices that both players are indifferent between their two strategies. Here we have \({\left( {x}^{\top }B\right) }_{3} = \frac{22}{8} = \frac{11}{4} > \frac{9}{4}\) . That is, even though player II is indifferent between her first two columns, which is necessary because \({y}_{1} > 0\) and \({y}_{2} > 0\) , these are not pure best responses against \(x = \left( {\frac{5}{8},\frac{3}{8}}\right)\) as required by (6.12). So there is no mixed equilibrium where player \(\mathrm{L}\) plays her first two columns.

\begin{tcolorbox}
其唯一解为 \({x}_{1} = \frac{5}{8},{x}_{2} = \frac{3}{8}\) 和 \({\left( {x}^{\top }B\right) }_{1} = {\left( {x}^{\top }B\right) }_{2} = \frac{18}{8} = \frac{9}{4}\) 。然而，我们还必须确保 \(\frac{9}{4}\) 是玩家 II 的最优反应收益，在 \(2 \times  2\) 博弈中这一考虑并非必要，在这类博弈中，只要两个玩家对他们的两种策略无差异即可。这里我们有 \({\left( {x}^{\top }B\right) }_{3} = \frac{22}{8} = \frac{11}{4} > \frac{9}{4}\) 。也就是说，尽管玩家 II 对她的前两列无差异，这是因为 \({y}_{1} > 0\) 和 \({y}_{2} > 0\) 所必需的，但根据(6.12)的要求，这些并非针对 \(x = \left( {\frac{5}{8},\frac{3}{8}}\right)\) 的纯最优反应。因此，不存在玩家 \(\mathrm{L}\) 采用她的前两列的混合均衡。
\end{tcolorbox}

Next, suppose player II chooses her first and third column, with \({y}_{1} > 0,{y}_{2} = 0\) and \({y}_{3} > 0\) . However, then the right equation in (6.17) states \(0 = 2{y}_{1} + {y}_{3}\) , which is not possible. Intuitively, when player II plays only the first and third column in (6.16), then the bottom row \(D\) is the unique best response and player I cannot be made indifferent.

\begin{tcolorbox}
接下来，假设玩家 II 选择她的第一列和第三列，其中 \({y}_{1} > 0,{y}_{2} = 0\) 和 \({y}_{3} > 0\) 。然而，此时(6.17)中的右侧方程表明 \(0 = 2{y}_{1} + {y}_{3}\) ，这是不可能的。直观地说，当玩家 II 在(6.16)中仅采用第一列和第三列时，那么底行 \(D\) 是唯一的最优反应，并且无法使玩家 I 无差异。
\end{tcolorbox}

This leaves the final possibility that \({y}_{1} = 0\) and \({y}_{2} > 0\) and \({y}_{3} > 0\) where (6.17) becomes \(2{y}_{2} = {y}_{3}\) and therefore \({y}_{2} = \frac{1}{3},{y}_{3} = \frac{2}{3}\) . Similar to (6.18), we now need

\begin{tcolorbox}
这就只剩下最后一种可能性，即 \({y}_{1} = 0\) 、 \({y}_{2} > 0\) 和 \({y}_{3} > 0\) ，此时(6.17)变为 \(2{y}_{2} = {y}_{3}\) ，因此 \({y}_{2} = \frac{1}{3},{y}_{3} = \frac{2}{3}\) 。与(6.18)类似，我们现在需要
\end{tcolorbox}

\[
{\left( {x}^{\top }B\right) }_{2} = {\left( {x}^{\top }B\right) }_{3} = v,\;6{x}_{2} = 2{x}_{1} + 4{x}_{2},\;2{x}_{1} = 2{x}_{2} \tag{6.19}
\]

which has the unique solution \({x}_{1} = {x}_{2} = \frac{1}{2}\) with \({\left( {x}^{\top }B\right) }_{2} = {\left( {x}^{\top }B\right) }_{3} = v = 3\) where \(v\) is indeed the maximum of expected column payoffs because \({\left( {x}^{\top }B\right) }_{1} = 2 < 3\) . We have found the mixed equilibrium \(x = \left( {\frac{1}{2},\frac{1}{2}}\right) ,y = \left( {0,\frac{1}{3},\frac{2}{3}}\right)\) . There is no equilibrium where player II chooses all three columns with positive probability because this implies the column indifferences (6.18) and (6.19), which have contradictory solutions \(x\) . The mixed equilibrium is therefore unique.

\begin{tcolorbox}
其唯一解为 \({x}_{1} = {x}_{2} = \frac{1}{2}\) ，其中 \({\left( {x}^{\top }B\right) }_{2} = {\left( {x}^{\top }B\right) }_{3} = v = 3\) ，由于 \({\left( {x}^{\top }B\right) }_{1} = 2 < 3\) ， \(v\) 确实是预期列收益的最大值。我们已经找到了混合均衡 \(x = \left( {\frac{1}{2},\frac{1}{2}}\right) ,y = \left( {0,\frac{1}{3},\frac{2}{3}}\right)\) 。不存在玩家 II 以正概率选择所有三列的均衡，因为这意味着列无差异条件(6.18)和(6.19)，而它们的解 \(x\) 相互矛盾。因此，混合均衡是唯一的。
\end{tcolorbox}

\(\Rightarrow\) Exercise 6.2 asks you to find all mixed equilibria of a \(3 \times  2\) game, which can be done similarly and illustrates the best-response condition. It also motivates methods for finding mixed equilibria that we develop in Section 6.6.

\begin{tcolorbox}
\(\Rightarrow\) 练习 6.2 要求你找出一个 \(3 \times  2\) 博弈的所有混合均衡，这可以用类似的方法完成，并能说明最优反应条件。它还为我们在第 6.6 节中开发的寻找混合均衡的方法提供了动机。
\end{tcolorbox}

The best-response condition applies also to a game with more than two players, which we explain in the remainder of this section. In Definition 3.1 we defined a commonly used notation and the concepts of best response and equilibrium for an \(N\) -player game. The following definition defines its mixed extension. For the finite strategy set \({S}_{i}\) of player \(i\) , we denote the set of mixed strategies over \({S}_{i}\) by \(\Delta \left( {S}_{i}\right)\) , in agreement with Definition 5.1.

\begin{tcolorbox}
最优反应条件同样适用于有两个以上参与者的博弈，我们将在本节的剩余部分对此进行解释。在定义3.1中，我们为一个 \(N\) 参与者博弈定义了常用符号以及最优反应和均衡的概念。以下定义给出了它的混合扩展。对于参与者 \(i\) 的有限策略集 \({S}_{i}\) ，根据定义5.1，我们用 \(\Delta \left( {S}_{i}\right)\) 表示 \({S}_{i}\) 上的混合策略集。
\end{tcolorbox}

Definition 6.2 (Mixed extension). Consider an \(N\) -player game in strategic form with finite strategy set \({S}_{i}\) for player \(i = 1,\ldots ,N\) and payoff \({u}_{i}\left( s\right)\) for each strategy profile \(s = \left( {{s}_{1},\ldots ,{s}_{N}}\right)\) in \(S = {S}_{1} \times  \cdots  \times  {S}_{N}\) . Let \({X}_{i} = \Delta \left( {S}_{i}\right)\) be the set of mixed strategies \({\sigma }_{i}\) of player \(i\) , where \({\sigma }_{i}\left( {s}_{i}\right)\) is the probability that \({\sigma }_{i}\) selects the pure strategy \({s}_{i} \in  {S}_{i}\) . (If \({\sigma }_{i}\left( {s}_{i}\right)  = 1\) for some \({s}_{i}\) then \({\sigma }_{i}\) is considered the same as the pure strategy \({s}_{i}\) .) Then the mixed extension of this game is the \(N\) -player game in strategic form with the (infinite) strategy sets \({X}_{1},\ldots ,{X}_{N}\) for the \(N\) players, where for a mixed-strategy profile \(\sigma  = \left( {{\sigma }_{1},\ldots ,{\sigma }_{N}}\right)\) player \(i\) receives the expected payoff

\begin{tcolorbox}
定义6.2(混合扩展)。考虑一个策略形式的 \(N\) 参与者博弈，其中参与者 \(i = 1,\ldots ,N\) 的有限策略集为 \({S}_{i}\) ，对于 \(S = {S}_{1} \times  \cdots  \times  {S}_{N}\) 中的每个策略组合 \(s = \left( {{s}_{1},\ldots ,{s}_{N}}\right)\) ，其支付为 \({u}_{i}\left( s\right)\) 。设 \({X}_{i} = \Delta \left( {S}_{i}\right)\) 为参与者 \(i\) 的混合策略 \({\sigma }_{i}\) 的集合，其中 \({\sigma }_{i}\left( {s}_{i}\right)\) 是 \({\sigma }_{i}\) 选择纯策略 \({s}_{i} \in  {S}_{i}\) 的概率。(如果对于某个 \({s}_{i}\) 有 \({\sigma }_{i}\left( {s}_{i}\right)  = 1\) ，则 \({\sigma }_{i}\) 被视为与纯策略 \({s}_{i}\) 相同。)那么，该博弈的混合扩展是一个策略形式的 \(N\) 参与者博弈，其(无限)策略集为 \({X}_{1},\ldots ,{X}_{N}\) ，对于所有 \(N\) 个参与者，在混合策略组合 \(\sigma  = \left( {{\sigma }_{1},\ldots ,{\sigma }_{N}}\right)\) 下，参与者 \(i\) 获得期望效用
\end{tcolorbox}

\[
{U}_{i}\left( \sigma \right)  = \mathop{\sum }\limits_{{s \in  S}}{u}_{i}\left( s\right) \mathop{\prod }\limits_{{k = 1}}^{N}{\sigma }_{k}\left( {s}_{k}\right) .
\]

An equilibrium (or Nash equilibrium) \(\sigma\) of the mixed extension is called a mixed equilibrium of the original game.

\begin{tcolorbox}
混合扩展的一个均衡(或纳什均衡) \(\sigma\) 被称为原博弈的混合均衡。
\end{tcolorbox}

The following proposition extends the best-response condition to an \(N\) -player game. As in the case of two players, it allows us to verify that a mixed-strategy profile is an equilibrium by the condition that a player's mixed strategy only selects pure best responses with positive probability.

\begin{tcolorbox}
以下命题将最优反应条件扩展到了一个 \(N\) 参与者博弈中。与两个参与者的情况一样，它使我们能够通过“参与者的混合策略仅以正概率选择纯最优反应”这一条件来验证一个混合策略组合是否为均衡。
\end{tcolorbox}

Proposition 6.3 (Best-response condition, \(N\) players). Consider the mixed extension of an \(N\) -player game with the notation in Definition 6.2. Let \({\sigma }_{i} \in  {X}_{i}\) be a mixed strategy of player \(i\) , let \({\sigma }_{-i}\) be a partial profile of mixed strategies of the other \(N - 1\) players, and let

\begin{tcolorbox}
命题6.3(最优反应条件， \(N\) 个参与者)。考虑一个 \(N\) 参与者博弈的混合扩展，采用定义6.2中的符号。设 \({\sigma }_{i} \in  {X}_{i}\) 为参与者 \(i\) 的混合策略，设 \({\sigma }_{-i}\) 为其他 \(N - 1\) 个参与者的混合策略的部分组合，并且设
\end{tcolorbox}

\[
{v}_{i} = \max \left\{  {{U}_{i}\left( {{s}_{i},{\sigma }_{-i}}\right)  \mid  {s}_{i} \in  {S}_{i}}\right\}
\]

be the best possible expected payoff against \({\sigma }_{-i}\) for a pure strategy \({s}_{i}\) of player \(i\) . Then \({\sigma }_{i}\) is a best response to \({\sigma }_{-i}\) if and only if

\begin{tcolorbox}
为参与者 \(i\) 的纯策略 \({s}_{i}\) 针对 \({\sigma }_{-i}\) 所能获得的最大期望效用。那么，当且仅当
\end{tcolorbox}

\[
{\sigma }_{i}\left( {s}_{i}\right)  > 0 \Rightarrow  {U}_{i}\left( {{s}_{i},{\sigma }_{-i}}\right)  = {v}_{i}\;\left( {{s}_{i} \in  {S}_{i}}\right) ,
\]

that is, if \({\sigma }_{i}\) only selects pure best responses \({s}_{i}\) against \({\sigma }_{i}\) with positive probability. Proof. The proof is the same as for Proposition 6.1 with the pure strategy \({s}_{i}\) (where \(i\) denotes a player) instead of row \(i\) (where \(i\) denotes a pure strategy) and probability \({\sigma }_{i}\left( {s}_{i}\right)\) instead of \({x}_{i}\) in (6.11). The crucial condition is that the expected payoff to player \(i\) is linear in his mixed-strategy probabilities \({\sigma }_{i}\left( {s}_{i}\right)\) , because for a mixed-strategy profile \(\sigma\) ,

\begin{tcolorbox}
即，如果 \({\sigma }_{i}\) 仅以正概率选择针对 \({\sigma }_{i}\) 的纯最优反应 \({s}_{i}\) 。证明:该证明与命题6.1的证明相同，只是将(6.11)中的行 \(i\) (其中 \(i\) 表示一个纯策略)替换为纯策略 \({s}_{i}\) (其中 \(i\) 表示一个参与者)，并将 \({x}_{i}\) 替换为概率 \({\sigma }_{i}\left( {s}_{i}\right)\) 。关键条件是，参与者 \(i\) 的期望效用关于其混合策略概率 \({\sigma }_{i}\left( {s}_{i}\right)\) 是线性的，因为对于一个混合策略组合 \(\sigma\) ，
\end{tcolorbox}

\[
{U}_{i}\left( \sigma \right)  = {U}_{i}\left( {{\sigma }_{i},{\sigma }_{-i}}\right)  = \mathop{\sum }\limits_{{s \in  S}}{u}_{i}\left( s\right) \mathop{\prod }\limits_{{k = 1}}^{N}{\sigma }_{k}\left( {s}_{k}\right)
\]

\[
= \mathop{\sum }\limits_{{{s}_{i} \in  {S}_{i}}}\mathop{\sum }\limits_{{{s}_{-i} \in  {S}_{-i}}}{u}_{i}\left( {{s}_{i},{s}_{-i}}\right) \mathop{\prod }\limits_{{k = 1}}^{N}{\sigma }_{k}\left( {s}_{k}\right)
\]

\[
= \mathop{\sum }\limits_{{{s}_{i} \in  {S}_{i}}}\left( {\mathop{\sum }\limits_{{{s}_{-i} \in  {S}_{-i}}}{u}_{i}\left( {{s}_{i},{s}_{-i}}\right) \mathop{\prod }\limits_{{k = 1,k \neq  i}}^{N}{\sigma }_{k}\left( {s}_{k}\right) }\right) {\sigma }_{i}\left( {s}_{i}\right)
\]

\[
= \mathop{\sum }\limits_{{{s}_{i} \in  {S}_{i}}}{U}_{i}\left( {{s}_{i},{\sigma }_{-i}}\right) {\sigma }_{i}\left( {s}_{i}\right) . \tag{6.20}
\]

That is, the expected payoff \({U}_{i}\left( \sigma \right)  = {U}_{i}\left( {{\sigma }_{i},{\sigma }_{-i}}\right)\) is exactly the expectation, with probabilities \({\sigma }_{i}\left( {s}_{i}\right)\) , of the expected payoffs \({U}_{i}\left( {{s}_{i},{\sigma }_{-i}}\right)\) for the "rows" \({s}_{i} \in  {S}_{i}\) , by considering player \(i\) as the row player. In particular, if \(N = 2\) and \(i\) is player 1, then \({\sigma }_{-i}\) is the mixed strategy \(y\) of player 2, and \({U}_{i}\left( {{s}_{i},{\sigma }_{-i}}\right)\) is the expected payoff \({\left( Ay\right) }_{{s}_{i}}\) for row \({s}_{i}\) . The proof in (6.13) with \({\sigma }_{i}\) instead of \(x\) and \({v}_{i}\) instead of \(u\) applies as before.

\begin{tcolorbox}
也就是说，期望效用 \({U}_{i}\left( \sigma \right)  = {U}_{i}\left( {{\sigma }_{i},{\sigma }_{-i}}\right)\) 恰好是以概率 \({\sigma }_{i}\left( {s}_{i}\right)\) 对“行” \({s}_{i} \in  {S}_{i}\) 的期望效用 \({U}_{i}\left( {{s}_{i},{\sigma }_{-i}}\right)\) 取期望，这里将玩家 \(i\) 视为行玩家。特别地，如果 \(N = 2\) 且 \(i\) 是玩家 1，那么 \({\sigma }_{-i}\) 是玩家 2 的混合策略 \(y\) ，而 \({U}_{i}\left( {{s}_{i},{\sigma }_{-i}}\right)\) 是行 \({s}_{i}\) 的期望效用 \({\left( Ay\right) }_{{s}_{i}}\) 。用 \({\sigma }_{i}\) 代替 \(x\) ，用 \({v}_{i}\) 代替 \(u\) 后，(6.13) 中的证明与之前相同。
\end{tcolorbox}

\section*{\(\Rightarrow\) Exercise 6.3 is about a game with three players.}

\begin{tcolorbox}
\section*{\(\Rightarrow\) 练习 6.3 是关于一个有三个玩家的博弈。}
\end{tcolorbox}

We continue to focus on two-player games, partly because of the simpler notation. More importantly, for two players the best-response condition compares expected payoffs (such as \({\left( Ay\right) }_{i}\) for the rows \(i\) of the row player, say) that depend linearly on the mixed-strategy probabilities \({y}_{j}\) of a single other player. In contrast, the corresponding expression \({U}_{i}\left( {{s}_{i},{\sigma }_{-i}}\right)\) in (6.20) for \(N\) players with \(N > 2\) is nonlinear in the other players' mixed-strategy probabilities, which makes it much harder to find a mixed equilibrium in a general \(N\) -player game.

\begin{tcolorbox}
我们继续关注两人博弈，部分原因是其符号表示更简单。更重要的是，对于两个玩家，最优反应条件比较的是期望效用(例如，对于行玩家的行 \(i\) 的 \({\left( Ay\right) }_{i}\) )，这些期望效用线性依赖于另一个玩家的混合策略概率 \({y}_{j}\) 。相比之下，(6.20) 中对于 \(N\) 个玩家且 \(N > 2\) 的相应表达式 \({U}_{i}\left( {{s}_{i},{\sigma }_{-i}}\right)\) 关于其他玩家的混合策略概率是非线性的，这使得在一般的 \(N\) 人博弈中找到混合均衡变得困难得多。
\end{tcolorbox}

\section*{6.5 Existence of Mixed Equilibria}

\begin{tcolorbox}
\section*{6.5 混合均衡的存在性}
\end{tcolorbox}

In this section, we give the original proof of John Nash (1951) that shows that any finite game has a mixed equilibrium. This proof uses Brouwer's fixed-point theorem (Theorem 6.4 below) which is about continuous functions on compact convex sets. A subset of \({\mathbb{R}}^{d}\) is compact if it is closed and bounded.

\begin{tcolorbox}
在本节中，我们给出约翰·纳什(1951 年)的原始证明，该证明表明任何有限博弈都存在混合均衡。这个证明使用了布劳威尔不动点定理(见下面的定理 6.4)，该定理是关于紧致凸集上的连续函数的。 \({\mathbb{R}}^{d}\) 的一个子集如果是闭集且有界，则它是紧致的。
\end{tcolorbox}

A subset \(C\) of \({\mathbb{R}}^{d}\) is convex if it contains with any two points \(x\) and \(y\) the line segment that connects \(x\) and \(y\) . As described in Figure 5.4 and (5.45), this is the case if and only for all \(x,y \in  C\) and \(\lambda  \in  \left\lbrack  {0,1}\right\rbrack\) also \(x\left( {1 - \lambda }\right)  + {y\lambda } \in  C\) , where we now multiply the column vectors \(x\) and \(y\) from the right with the scalars \(\left( {1 - \lambda }\right)\) and \(\lambda\) so that they are matrix products as in (6.3). For \(\lambda  \in  \left\lbrack  {0,1}\right\rbrack\) , the term \(x\left( {1 - \lambda }\right)  + {y\lambda }\) is called the convex combination of \(x\) and \(y\) (which is any point on the line segment that connects \(x\) and \(y\) ). In general, a convex combination of points \({z}_{1},{z}_{2},\ldots ,{z}_{n}\) in \({\mathbb{R}}^{d}\) is a linear combination \({z}_{1}{\lambda }_{1} + {z}_{2}{\lambda }_{2} + \cdots  + {z}_{n}{\lambda }_{n}\) where the coefficients \({\lambda }_{1},\ldots ,{\lambda }_{n}\) are nonnegative and sum to one.

\begin{tcolorbox}
\({\mathbb{R}}^{d}\) 的一个子集 \(C\) 是凸集，如果它包含连接任意两点 \(x\) 和 \(y\) 的线段。如图 5.4 和 (5.45) 所述，当且仅当对于所有 \(x,y \in  C\) 和 \(\lambda  \in  \left\lbrack  {0,1}\right\rbrack\) 也有 \(x\left( {1 - \lambda }\right)  + {y\lambda } \in  C\) 时才是这种情况，这里我们现在将列向量 \(x\) 和 \(y\) 从右边与标量 \(\left( {1 - \lambda }\right)\) 和 \(\lambda\) 相乘，使得它们像 (6.3) 中那样是矩阵乘积。对于 \(\lambda  \in  \left\lbrack  {0,1}\right\rbrack\) ，项 \(x\left( {1 - \lambda }\right)  + {y\lambda }\) 称为 \(x\) 和 \(y\) 的凸组合(它是连接 \(x\) 和 \(y\) 的线段上的任意一点)。一般来说， \({\mathbb{R}}^{d}\) 中各点 \({z}_{1},{z}_{2},\ldots ,{z}_{n}\) 的凸组合是一个线性组合 \({z}_{1}{\lambda }_{1} + {z}_{2}{\lambda }_{2} + \cdots  + {z}_{n}{\lambda }_{n}\) ，其中系数 \({\lambda }_{1},\ldots ,{\lambda }_{n}\) 是非负的且和为 1。
\end{tcolorbox}

\begin{center}
\includegraphics[max width=1.0\textwidth]{images/bo_d2e5cmbef24c73b03tn0_161_329_538_1270_1091_0.jpg}
\end{center}
\hspace*{3em} 

The players’ mixed-strategy sets are special convex sets. For an \(m \times  n\) game, the mixed-strategy sets of players I and II are \(X\) and \(Y\) as in (6.6). The \(i\) th unit vector \({e}_{i}\) in \({\mathbb{R}}^{m}\) , for \(1 \leq  i \leq  m\) , is the vector \({\left( 0,\ldots ,0,1,0,\ldots ,0\right) }^{\top }\) where the \(i\) th component is one and every other component is zero. This vector belongs to \(X\) and represents the pure strategy \(i\) of player I. For any mixed strategy \(x = {\left( {x}_{1},\ldots ,{x}_{m}\right) }^{\top }\) in \(X\) , clearly \(x = {e}_{1}{x}_{1} + {e}_{2}{x}_{2} + \cdots  + {e}_{n}{x}_{m}\) , that is, \(x\) is the convex combination of the unit vectors \({e}_{1},\ldots ,{e}_{m}\) with the mixed-strategy probabilities \({x}_{1},\ldots ,{x}_{m}\) as coefficients (because they are nonnegative and sum to one). Similarly, any mixed strategy \(y \in  Y\) is the convex combination of the unit vectors in \({\mathbb{R}}^{n}\) , which represent the pure strategies of player II.

\begin{tcolorbox}
玩家的混合策略集是特殊的凸集。对于一个 \(m \times  n\) 博弈，玩家 I 和玩家 II 的混合策略集分别为 (6.6) 式中的 \(X\) 和 \(Y\) 。 \({\mathbb{R}}^{m}\) 中的第 \(i\) 个单位向量 \({e}_{i}\) (其中 \(1 \leq  i \leq  m\) )是向量 \({\left( 0,\ldots ,0,1,0,\ldots ,0\right) }^{\top }\) ，其第 \(i\) 个分量为 1，其他分量均为 0。该向量属于 \(X\) ，表示玩家 I 的纯策略 \(i\) 。对于 \(X\) 中的任何混合策略 \(x = {\left( {x}_{1},\ldots ,{x}_{m}\right) }^{\top }\) ，显然有 \(x = {e}_{1}{x}_{1} + {e}_{2}{x}_{2} + \cdots  + {e}_{n}{x}_{m}\) ，即 \(x\) 是以混合策略概率 \({x}_{1},\ldots ,{x}_{m}\) 为系数的单位向量 \({e}_{1},\ldots ,{e}_{m}\) 的凸组合(因为这些概率是非负的且总和为 1)。类似地，任何混合策略 \(y \in  Y\) 是 \({\mathbb{R}}^{n}\) 中单位向量的凸组合，这些单位向量表示玩家 II 的纯策略。
\end{tcolorbox}

Examples of \(X\) are shown in Figure 6.5 for \(m = 2,3,4\) where \(X\) is a line segment, triangle, and tetrahedron, respectively. In general, the convex hull of the \(m\) unit vectors in \({\mathbb{R}}^{m}\) is called the unit simplex. It is easy to see that \(X\) and \(Y\) , and their Cartesian product \(X \times  Y\) , are convex sets (see (6.23) below). In order to prove the existence of a mixed equilibrium, we will use the following theorem for a suitable continuous function \(f : C \rightarrow  C\) on the set \(C = X \times  Y\) of mixed strategy profiles.

\begin{tcolorbox}
图 6.5 展示了 \(X\) 在 \(m = 2,3,4\) 情况下的示例，其中 \(X\) 分别为线段、三角形和四面体。一般来说， \({\mathbb{R}}^{m}\) 中 \(m\) 个单位向量的凸包称为单位单纯形。容易看出， \(X\) 和 \(Y\) 以及它们的笛卡尔积 \(X \times  Y\) 都是凸集(见下面的 (6.23) 式)。为了证明混合均衡的存在性，我们将对混合策略组合集 \(C = X \times  Y\) 上的合适连续函数 \(f : C \rightarrow  C\) 使用以下定理。
\end{tcolorbox}

Theorem 6.4 (Brouwer’s fixed-point theorem). Let \(C\) be a nonempty subset of \({\mathbb{R}}^{d}\) that is convex and compact, and let \(f\) be a continuous function from \(C\) to \(C\) . Then \(f\) has at least one fixed point, that is, a point \(z\) in \(C\) so that \(f\left( z\right)  = z\) .

\begin{tcolorbox}
定理 6.4(布劳威尔不动点定理)。设 \(C\) 是 \({\mathbb{R}}^{d}\) 的一个非空子集，它是凸集且是紧集，设 \(f\) 是从 \(C\) 到 \(C\) 的连续函数。那么 \(f\) 至少有一个不动点，即 \(C\) 中存在一点 \(z\) ，使得 \(f\left( z\right)  = z\) 。
\end{tcolorbox}

\begin{center}
\includegraphics[max width=0.3\textwidth]{images/bo_d2e5cmbef24c73b03tn0_162_277_700_461_444_0.jpg}
\end{center}
\hspace*{3em} 

Figure 6.6 Proof of Brouwer's fixed-point theorem for \(C = \left\lbrack  {0,1}\right\rbrack\) .

\begin{tcolorbox}
图 6.6 是 \(C = \left\lbrack  {0,1}\right\rbrack\) 情况下布劳威尔不动点定理的证明。
\end{tcolorbox}

Figure 6.6 illustrates Theorem 6.4 for the interval \(C = \left\lbrack  {0,1}\right\rbrack\) . If the continuous function \(f : C \rightarrow  C\) fulfills \(f\left( 0\right)  = 0\) , then 0 is a fixed point of \(f\) , and similarly if \(f\left( 1\right)  = 1\) then 1 is a fixed point of \(f\) . Hence, we can assume that \(f\left( 0\right)  > 0\) and \(f\left( 1\right)  < 1\) . Then the graph of the function \(f\) crosses the line of the identity function \(x \mapsto  x\) at least once for some point \(z\) where \(f\left( z\right)  = z\) , which is the desired fixed point. More formally, the continuous function \(x \mapsto  x - f\left( x\right)\) is negative for \(x = 0\) and positive for \(x = 1\) , and therefore zero for an intermediate value \(z\) with \(0 < z < 1\) by the intermediate value theorem, that is, \(z - f\left( z\right)  = 0\) or \(f\left( z\right)  = z\) .

\begin{tcolorbox}
图 6.6 说明了区间 \(C = \left\lbrack  {0,1}\right\rbrack\) 上的定理 6.4。如果连续函数 \(f : C \rightarrow  C\) 满足 \(f\left( 0\right)  = 0\) ，那么 0 是 \(f\) 的一个不动点；类似地，如果 \(f\left( 1\right)  = 1\) ，那么 1 是 \(f\) 的一个不动点。因此，我们可以假设 \(f\left( 0\right)  > 0\) 且 \(f\left( 1\right)  < 1\) 。那么函数 \(f\) 的图至少与恒等函数 \(x \mapsto  x\) 的线相交一次，交点为 \(z\) ，此时 \(f\left( z\right)  = z\) ，这就是我们所需要的不动点。更正式地说，连续函数 \(x \mapsto  x - f\left( x\right)\) 在 \(x = 0\) 处为负，在 \(x = 1\) 处为正，根据介值定理，对于介于两者之间的值 \(z\) (其中 \(0 < z < 1\) )，该函数值为 0，即 \(z - f\left( z\right)  = 0\) 或 \(f\left( z\right)  = z\) 。
\end{tcolorbox}

The following simple examples show that the assumptions in Theorem 6.4 are necessary. Consider the function \(f : \mathbb{R} \rightarrow  \mathbb{R}\) defined by \(f\left( x\right)  = x + 1\) . This is clearly a continuous function, and the set \(\mathbb{R}\) of all real numbers is closed and convex, but \(f\) has no fixed point. Some assumption of the fixed-point theorem must be violated, and in this case it is compactness, because the set \(\mathbb{R}\) is not bounded. Consider another function \(f : \left\lbrack  {0,1}\right\rbrack   \rightarrow  \left\lbrack  {0,1}\right\rbrack\) , given by \(f\left( x\right)  = {x}^{2}\) . The fixed points of this function are 0 and 1 . If we consider the function \(f\left( x\right)  = {x}^{2}\) as a function on the open interval(0,1), then this function has no longer any fixed points. In this case, the missing condition is that the function is not defined on a closed set, which is therefore not compact. Another function is \(f\left( x\right)  = 1 - x\) for \(x \in  \{ 0,1\}\) , where the domain of this function has just two elements, so this is a compact set. This function has no fixed points, which is possible because its domain is not convex. Finally, the function on \(\left\lbrack  {0,1}\right\rbrack\) defined by \(f\left( x\right)  = 1\) for \(0 \leq  x \leq  \frac{1}{2}\) and \(f\left( x\right)  = 0\) for \(\frac{1}{2} < x \leq  1\) has no fixed point, which is possible in this case because the function is not continuous. These examples demonstrate why the assumptions of Theorem 6.4 are necessary.

\begin{tcolorbox}
以下简单示例表明，定理6.4中的假设是必要的。考虑由 \(f\left( x\right)  = x + 1\) 定义的函数 \(f : \mathbb{R} \rightarrow  \mathbb{R}\) 。这显然是一个连续函数，并且所有实数的集合 \(\mathbb{R}\) 是闭集且凸集，但 \(f\) 没有不动点。不动点定理的某些假设必定不成立，在这种情况下，不满足的是紧性，因为集合 \(\mathbb{R}\) 是无界的。再考虑另一个函数 \(f : \left\lbrack  {0,1}\right\rbrack   \rightarrow  \left\lbrack  {0,1}\right\rbrack\) ，由 \(f\left( x\right)  = {x}^{2}\) 给出。这个函数的不动点是0和1。如果我们将函数 \(f\left( x\right)  = {x}^{2}\) 看作是定义在开区间(0,1)上的函数，那么这个函数就不再有任何不动点。在这种情况下，缺失的条件是函数不是定义在闭集上，因此该集合不是紧集。另一个函数是当 \(x \in  \{ 0,1\}\) 时的 \(f\left( x\right)  = 1 - x\) ，该函数的定义域只有两个元素，所以这是一个紧集。这个函数没有不动点，这是可能的，因为它的定义域不是凸集。最后，定义在 \(\left\lbrack  {0,1}\right\rbrack\) 上的函数，当 \(0 \leq  x \leq  \frac{1}{2}\) 时为 \(f\left( x\right)  = 1\) ，当 \(\frac{1}{2} < x \leq  1\) 时为 \(f\left( x\right)  = 0\) ，它没有不动点，在这种情况下这是可能的，因为该函数不连续。这些示例说明了为什么定理6.4的假设是必要的。
\end{tcolorbox}

Theorem 6.4 is a major theorem in topology and much harder to prove for general dimension \(d\) . We will give a full proof in Chapter 7.

\begin{tcolorbox}
定理6.4是拓扑学中的一个重要定理，对于一般维度 \(d\) 而言，其证明要困难得多。我们将在第7章给出完整的证明。
\end{tcolorbox}

The following theorem states the existence of a mixed equilibrium. We prove it using Brouwer's fixed-point theorem along the lines of the original proof by Nash (1951), for notational simplicity for the case of two players.

\begin{tcolorbox}
以下定理阐述了混合均衡的存在性。我们使用布劳威尔不动点定理，按照纳什(1951)原始证明的思路来证明它，为简化符号，这里考虑两个参与者的情况。
\end{tcolorbox}

Theorem 6.5 (Nash, 1951). Every finite game has at least one mixed equilibrium.

\begin{tcolorbox}
定理6.5(纳什，1951)。每个有限博弈至少有一个混合均衡。
\end{tcolorbox}

Proof. We will give the proof for an \(m \times  n\) two-player game(A, B), to simplify notation. It extends in the same manner to any finite number of players. The set \(C\) that is used in the present context is the product of the sets of mixed strategies of the players. Let \(X\) and \(Y\) be the sets of mixed strategies of player I and player II as in (6.6), and let \(C = X \times  Y\) , which is clearly convex and compact.

\begin{tcolorbox}
证明。为简化符号，我们将给出 \(m \times  n\) 两人博弈(A, B)的证明。它可以以同样的方式推广到任意有限数量的参与者。在当前上下文中使用的集合 \(C\) 是参与者混合策略集合的乘积。设 \(X\) 和 \(Y\) 分别是参与者I和参与者II的混合策略集合，如(6.6)所示，并设 \(C = X \times  Y\) ，显然它是凸集且是紧集。
\end{tcolorbox}

Then the function \(f : C \rightarrow  C\) that we are going to construct maps a pair of mixed strategies(x, y)to another pair \(f\left( {x,y}\right)  = \left( {\bar{x},\bar{y}}\right)\) . Intuitively, a mixed-strategy probability \({x}_{i}\) of player I (and similarly \({y}_{j}\) of player II) is changed to \({\bar{x}}_{i}\) such that it will increase if the pure strategy \(i\) does better than the current expected payoff \({x}^{\top }{Ay}\) . If the pure strategy \(i\) does worse, then \({x}_{i}\) will decrease unless \({x}_{i}\) is already zero. In equilibrium, \(x\) is a best response to \(y\) and \({x}^{\top }{Ay} = u\) for the pure best-response payoff \(u\) according to (6.10), and all sub-optimal pure strategies have probability zero. This means that the mixed strategy \(x\) does not change, \(\bar{x} = x\) . Similarly, \(\bar{y} = y\) if \(y\) is a best response to \(x\) . Therefore, the equilibrium property is equivalent to the fixed-point property \(\left( {x,y}\right)  = \left( {\bar{x},\bar{y}}\right)  = f\left( {x,y}\right)\) .

\begin{tcolorbox}
然后，我们要构建的函数 \(f : C \rightarrow  C\) 将一对混合策略 (x, y) 映射到另一对 \(f\left( {x,y}\right)  = \left( {\bar{x},\bar{y}}\right)\) 。直观地说，玩家 I 的混合策略概率 \({x}_{i}\) (玩家 II 的 \({y}_{j}\) 同理)会变为 \({\bar{x}}_{i}\) ，如果纯策略 \(i\) 的表现优于当前期望效用 \({x}^{\top }{Ay}\) ，则 \({x}_{i}\) 会增加。如果纯策略 \(i\) 的表现较差，那么除非 \({x}_{i}\) 已经为零，否则 \({x}_{i}\) 会减小。在均衡状态下，根据 (6.10)， \(x\) 是对 \(y\) 和 \({x}^{\top }{Ay} = u\) 的最优反应，对应纯最优反应效用 \(u\) ，并且所有次优纯策略的概率为零。这意味着混合策略 \(x\) 不会改变，即 \(\bar{x} = x\) 。类似地，如果 \(y\) 是对 \(x\) 的最优反应，那么 \(\bar{y} = y\) 。因此，均衡性质等价于不动点性质 \(\left( {x,y}\right)  = \left( {\bar{x},\bar{y}}\right)  = f\left( {x,y}\right)\) 。
\end{tcolorbox}

In order to define \(f\) as described, consider the following functions \(\chi  : X \times  Y \rightarrow\)  \({\mathbb{R}}^{m}\) and \(\psi  : X \times  Y \rightarrow  {\mathbb{R}}^{n}\) . For each pure strategy \(i\) of player I, let \({\chi }_{i}\left( {x,y}\right)\) be the \(i\) th component of \(\chi \left( {x,y}\right)\) , and for each pure strategy \(j\) of player II, let \({\psi }_{j}\left( {x,y}\right)\) be the \(j\) th component of \(\psi \left( {x,y}\right)\) . The functions \(\chi\) and \(\psi\) are defined by

\begin{tcolorbox}
为了如前文所述地定义 \(f\) ，考虑以下函数 \(\chi  : X \times  Y \rightarrow\) 、 \({\mathbb{R}}^{m}\) 和 \(\psi  : X \times  Y \rightarrow  {\mathbb{R}}^{n}\) 。对于玩家 I 的每个纯策略 \(i\) ，令 \({\chi }_{i}\left( {x,y}\right)\) 为 \(\chi \left( {x,y}\right)\) 的第 \(i\) 个分量；对于玩家 II 的每个纯策略 \(j\) ，令 \({\psi }_{j}\left( {x,y}\right)\) 为 \(\psi \left( {x,y}\right)\) 的第 \(j\) 个分量。函数 \(\chi\) 和 \(\psi\) 定义如下
\end{tcolorbox}

\[
{\chi }_{i}\left( {x,y}\right)  = \max \left\{  {0,{\left( Ay\right) }_{i} - {x}^{\top }{Ay}}\right\}  \;\left( {1 \leq  i \leq  m}\right) , \tag{6.21}
\]

\[
{\psi }_{j}\left( {x,y}\right)  = \max \left\{  {0,{\left( {x}^{\top }B\right) }_{j} - {x}^{\top }{By}}\right\}  \;\left( {1 \leq  j \leq  n}\right) .
\]

Recall that \({\left( Ay\right) }_{i}\) is the expected payoff to player I against \(y\) when he uses the pure strategy \(i\) , and that \({\left( {x}^{\top }B\right) }_{j}\) is the expected payoff to player II against \(x\) when she uses the pure strategy \(j\) , and that \({x}^{\top }{Ay}\) and \({x}^{\top }{By}\) are the expected payoffs to player \(\mathrm{I}\) and player II. Hence, the difference \({\left( Ay\right) }_{i} - {x}^{\top }{Ay}\) is positive if the pure strategy \(i\) gives more than the average \({x}^{\top }{Ay}\) against \(y\) , zero if it gives the same payoff, and negative if it gives less. The term \({\chi }_{i}\left( {x,y}\right)\) is this difference, except that it is replaced by zero if the difference is negative. The term \({\psi }_{j}\left( {x,y}\right)\) is defined analogously. Thus, \(\chi \left( {x,y}\right)\) is a nonnegative vector in \({\mathbb{R}}^{m}\) , and \(\psi \left( {x,y}\right)\) is a nonnegative vector in \({\mathbb{R}}^{n}\) . The functions \(\chi\) and \(\psi\) are continuous.

\begin{tcolorbox}
回顾一下， \({\left( Ay\right) }_{i}\) 是玩家 I 使用纯策略 \(i\) 对抗 \(y\) 时的期望效用， \({\left( {x}^{\top }B\right) }_{j}\) 是玩家 II 使用纯策略 \(j\) 对抗 \(x\) 时的期望效用， \({x}^{\top }{Ay}\) 和 \({x}^{\top }{By}\) 分别是玩家 \(\mathrm{I}\) 和玩家 II 的期望效用。因此，如果纯策略 \(i\) 对抗 \(y\) 时的收益高于平均收益 \({x}^{\top }{Ay}\) ，则差值 \({\left( Ay\right) }_{i} - {x}^{\top }{Ay}\) 为正；如果收益相同，则差值为零；如果收益较低，则差值为负。项 \({\chi }_{i}\left( {x,y}\right)\) 就是这个差值，但如果差值为负，则该项被替换为零。项 \({\psi }_{j}\left( {x,y}\right)\) 的定义类似。因此， \(\chi \left( {x,y}\right)\) 是 \({\mathbb{R}}^{m}\) 中的非负向量， \(\psi \left( {x,y}\right)\) 是 \({\mathbb{R}}^{n}\) 中的非负向量。函数 \(\chi\) 和 \(\psi\) 是连续的。
\end{tcolorbox}

The pair of vectors(x, y)is now changed by replacing \(x\) by \(x + \chi \left( {x,y}\right)\) in order to get \(\bar{x}\) , and \(y\) by \(y + \psi \left( {x,y}\right)\) to get \(\bar{y}\) . Both sums are nonnegative. However, in general these new vectors are no longer mixed strategies because their components do not sum to one. For that purpose, they are "re-normalized" by the following continuous functions \(r\) and \(s\) :

\begin{tcolorbox}
现在，通过将 \(x\) 替换为 \(x + \chi \left( {x,y}\right)\) 来改变向量对 (x, y) 以得到 \(\bar{x}\) ，将 \(y\) 替换为 \(y + \psi \left( {x,y}\right)\) 以得到 \(\bar{y}\) 。两个和都是非负的。然而，一般来说，这些新向量不再是混合策略，因为它们的分量之和不为 1。为此，通过以下连续函数 \(r\) 和 \(s\) 对它们进行“重新归一化”:
\end{tcolorbox}

\[
r\left( x\right)  = x\frac{1}{{\mathbf{1}}^{\top }x},\;s\left( y\right)  = y\frac{1}{{\mathbf{1}}^{\top }y}
\]

for nonzero nonnegative vectors \(x \in  {\mathbb{R}}^{m}\) and \(y \in  {\mathbb{R}}^{n}\) . Then \(r\left( x\right)  \in  X\) because \({\mathbf{1}}^{\top }r\left( x\right)  = {\mathbf{1}}^{\top }x\frac{1}{{\mathbf{1}}^{\top }x} = 1\) . Analogously, \(s\left( y\right)  \in  Y\) . The function \(f : C \rightarrow  C\) is now defined by

\begin{tcolorbox}
对于非零非负向量 \(x \in  {\mathbb{R}}^{m}\) 和 \(y \in  {\mathbb{R}}^{n}\) 。那么 \(r\left( x\right)  \in  X\) ，因为 \({\mathbf{1}}^{\top }r\left( x\right)  = {\mathbf{1}}^{\top }x\frac{1}{{\mathbf{1}}^{\top }x} = 1\) 。类似地， \(s\left( y\right)  \in  Y\) 。函数 \(f : C \rightarrow  C\) 现在定义为
\end{tcolorbox}

\[
f\left( {x,y}\right)  = \left( {r\left( {x + \chi \left( {x,y}\right) }\right) ,s\left( {y + \psi \left( {x,y}\right) }\right) }\right) .
\]

We claim that

\begin{tcolorbox}
我们断言
\end{tcolorbox}

\[
x = r\left( {x + \chi \left( {x,y}\right) }\right) \; \Leftrightarrow  \;x\text{ is a best response to }y, \tag{6.22}
\]

\[
y = s\left( {y + \psi \left( {x,y}\right) }\right) \; \Leftrightarrow  \;y\text{ is a best response to }x,
\]

which proves that(x, y)is a fixed point of \(f\) if and only if(x, y)is an equilibrium of the game(A, B). By symmetry, it suffices to prove the first equivalence in (6.22), where " \(\Leftarrow\) " is easy: If \(x\) is a best response to \(y\) , then the best-response condition (6.10) implies \(\chi \left( {x,y}\right)  = \mathbf{0} \in  {\mathbb{R}}^{m}\) in (6.21) and therefore \(x = r\left( {x + \chi \left( {x,y}\right) }\right)\) .

\begin{tcolorbox}
这证明了 (x, y) 是 \(f\) 的一个不动点当且仅当 (x, y) 是博弈 (A, B) 的一个均衡。由对称性，只需证明 (6.22) 中的第一个等价关系，其中“ \(\Leftarrow\) ”很容易:如果 \(x\) 是对 \(y\) 的最优反应，那么最优反应条件 (6.10) 意味着 (6.21) 中的 \(\chi \left( {x,y}\right)  = \mathbf{0} \in  {\mathbb{R}}^{m}\) ，因此 \(x = r\left( {x + \chi \left( {x,y}\right) }\right)\) 。
\end{tcolorbox}

Conversely, let \(x = r\left( {x + \chi \left( {x,y}\right) }\right)\) . In theory, this could hold if \(\chi \left( {x,y}\right)  \neq  \mathbf{0}\) (so that \(x\) is not a best response to \(y\) ) due to the re-normalization with \(r\) , for example if \(x = \left( {\frac{1}{3},\frac{1}{3},\frac{1}{3}}\right)\) and \(\chi \left( {x,y}\right)  = \left( {\frac{1}{5},\frac{1}{5},\frac{1}{5}}\right)\) . However, this is not possible. Let

\begin{tcolorbox}
反之，设 \(x = r\left( {x + \chi \left( {x,y}\right) }\right)\) 。理论上，如果 \(\chi \left( {x,y}\right)  \neq  \mathbf{0}\) (使得 \(x\) 不是对 \(y\) 的最优反应)，由于用 \(r\) 进行重新归一化，这可能成立，例如如果 \(x = \left( {\frac{1}{3},\frac{1}{3},\frac{1}{3}}\right)\) 且 \(\chi \left( {x,y}\right)  = \left( {\frac{1}{5},\frac{1}{5},\frac{1}{5}}\right)\) 。然而，这是不可能的。设
\end{tcolorbox}

\[
u = \max \left\{  {{\left( Ay\right) }_{i} \mid  1 \leq  i \leq  m}\right\}   = {\left( Ay\right) }_{\ell },
\]

\[
t = \min \left\{  {{\left( Ay\right) }_{i} \mid  1 \leq  i \leq  m,{x}_{i} > 0}\right\}   = {\left( Ay\right) }_{k}\;\left( {\text{ with }{x}_{k} > 0}\right) .
\]

Shown analogously to (6.13), \(t \leq  {x}^{\top }{Ay}\) , and therefore \({\chi }_{k}\left( {x,y}\right)  = 0\) . If \(x\) is not a best response to \(y\) , then \({x}^{\top }{Ay} < u\) , thus \({\chi }_{\ell }\left( {x,y}\right)  > 0\) and \({\mathbf{1}}^{\top }\left( {x + \chi \left( {x,y}\right) }\right)  > 1\) , that is, \(\frac{1}{{\mathbf{1}}^{\top }\left( {x + \chi \left( {x,y}\right) }\right) } < 1\) . Because \({\chi }_{k}\left( {x,y}\right)  = 0\) , this means that the re-normalized vector \(\bar{x} = r\left( {x + \chi \left( {x,y}\right) }\right)\) shrinks in its positive \(k\) th component, \({\bar{x}}_{k} < {x}_{k}\) , which contradicts the assumption \(\bar{x} = x\) . This proves " \(\Rightarrow\) " in (6.22).

\begin{tcolorbox}
与 (6.13) 类似地证明， \(t \leq  {x}^{\top }{Ay}\) ，因此 \({\chi }_{k}\left( {x,y}\right)  = 0\) 。如果 \(x\) 不是对 \(y\) 的最优反应，那么 \({x}^{\top }{Ay} < u\) ，因此 \({\chi }_{\ell }\left( {x,y}\right)  > 0\) 且 \({\mathbf{1}}^{\top }\left( {x + \chi \left( {x,y}\right) }\right)  > 1\) ，即 \(\frac{1}{{\mathbf{1}}^{\top }\left( {x + \chi \left( {x,y}\right) }\right) } < 1\) 。因为 \({\chi }_{k}\left( {x,y}\right)  = 0\) ，这意味着重新归一化后的向量 \(\bar{x} = r\left( {x + \chi \left( {x,y}\right) }\right)\) 在其正的第 \(k\) 个分量 \({\bar{x}}_{k} < {x}_{k}\) 上缩小，这与假设 \(\bar{x} = x\) 矛盾。这证明了 (6.22) 中的“ \(\Rightarrow\) ”。
\end{tcolorbox}

As claimed, the fixed points(x, y)are therefore exactly the equilibria of the game, and an equilibrium exists by Brouwer's fixed-point Theorem 6.4.

\begin{tcolorbox}
如所断言，因此不动点 (x, y) 恰好是博弈的均衡，并且根据布劳威尔不动点定理 6.4 存在一个均衡。
\end{tcolorbox}

We give the proof of Brouwer's fixed-point theorem in Chapter 7.

\begin{tcolorbox}
我们将在第 7 章给出布劳威尔不动点定理的证明。
\end{tcolorbox}

\section*{6.6 Finding Mixed Equilibria in Small Games}

\begin{tcolorbox}
\section*{6.6 在小型博弈中寻找混合均衡}
\end{tcolorbox}

How can we find all equilibria of a two-player game in strategic form? It is easy to determine the pure-strategy equilibria, because these are the cells in the payoff table where both payoffs are marked with a box as best-response payoffs. We now describe how to find mixed-strategy equilibria, using the best-response condition Proposition 6.1.

\begin{tcolorbox}
我们如何找到策略形式的两人博弈的所有均衡呢？确定纯策略均衡很容易，因为这些是收益表中两个收益都被框起来作为最优反应收益的单元格。我们现在描述如何使用最优反应条件命题 6.1 来找到混合策略均衡。
\end{tcolorbox}

We first consider \(2 \times  2\) games. Our first example above, the Inspection game in Figure 6.2, has no pure-strategy equilibrium. There we have determined the mixed-strategy probabilities of a player so as to make the other player indifferent between her pure strategies, because only then will she mix between these strategies. This is a consequence of Proposition 6.1: Only pure strategies that get maximum, and hence equal, expected payoff can be played with positive probability in equilibrium.

\begin{tcolorbox}
我们首先考虑 \(2 \times  2\) 博弈。我们上面的第一个例子，即图6.2中的监管博弈，没有纯策略均衡。在那里，我们确定了一个玩家的混合策略概率，以便让另一个玩家在她的纯策略之间无差异，因为只有这样她才会在这些策略之间进行混合。这是命题6.1的一个推论:在均衡状态下，只有那些获得最大(因此相等)期望效用的纯策略才有可能以正概率被采用。
\end{tcolorbox}

\begin{center}
\includegraphics[max width=0.7\textwidth]{images/bo_d2e5cmbef24c73b03tn0_165_456_667_979_493_0.jpg}
\end{center}
\hspace*{3em} 

Figure 6.7 Battle of the Sexes game (left) and a \(3 \times  2\) game (right), which is the same game with an extra strategy \(B\) for player \(\mathrm{I}\) .

\begin{tcolorbox}
图6.7 性别战博弈(左)和一个 \(3 \times  2\) 博弈(右)，后者是玩家 \(\mathrm{I}\) 增加了一个额外策略 \(B\) 的相同博弈。
\end{tcolorbox}

\begin{center}
\includegraphics[max width=0.4\textwidth]{images/bo_d2e5cmbef24c73b03tn0_165_403_1330_566_595_0.jpg}
\end{center}
\hspace*{3em} 

Figure 6.8 Best-response mixed strategies(1 - p, p) (gray) and(1 - q, q)(black) for the \(2 \times  2\) Battle of the Sexes game on the left in Figure 6.7. The game has the two pure equilibria \(\;\left( {\left( {1,0}\right) ,\left( {1,0}\right) }\right) \;\) and \(\left( {\left( {0,1}\right) ,\left( {0,1}\right) }\right)\) and the mixed equilibrium \(\left( {\left( {\frac{1}{3},\frac{2}{3}}\right) ,\left( {\frac{2}{3},\frac{1}{3}}\right) }\right)\) .

\begin{tcolorbox}
图6.8 图6.7左侧性别战 \(2 \times  2\) 博弈的最优反应混合策略 (1 - p, p)(灰色)和 (1 - q, q)(黑色)。该博弈有两个纯策略均衡 \(\;\left( {\left( {1,0}\right) ,\left( {1,0}\right) }\right) \;\) 和 \(\left( {\left( {0,1}\right) ,\left( {0,1}\right) }\right)\) 以及一个混合均衡 \(\left( {\left( {\frac{1}{3},\frac{2}{3}}\right) ,\left( {\frac{2}{3},\frac{1}{3}}\right) }\right)\) 。
\end{tcolorbox}

A mixed-strategy equilibrium can also exist in \(2 \times  2\) games that have pure-strategy equilibria. As an example, consider the Battle of the Sexes game on the left in Figure 6.7. The best-response payoffs show that it has the pure-strategy equilibria(C, c)and(S, s). Suppose now that player I plays the mixed strategy \(x = \left( {1 - p,p}\right)\) , and player II plays the mixed strategy \(y = \left( {1 - q,q}\right)\) . Similar to Figure 6.4 for the Inspection game, Figure 6.8 shows the best responses of the player against the mixed strategy of the other player. For example, if \(p = \frac{1}{2}\) then player II’s expected payoffs are 1 for column \(c\) and \(\frac{1}{2}\) for column \(s\) and her best response is \(c\) , which is the mixed strategy(1,0)where \(q = 0\) . In general, the payoffs for \(c\) and \(s\) are \(2\left( {1 - p}\right)\) and \(p\) , so player II weakly prefers \(c\) to \(s\) if \(2\left( {1 - p}\right)  \geq  p\) or equivalently \(\frac{2}{3} \geq  p\) , that is, \(p \in  \left\lbrack  {0,\frac{2}{3}}\right\rbrack\) . For \(p = \frac{2}{3}\) player II is indifferent and any \(q \in  \left\lbrack  {0,1}\right\rbrack\) defines a best response(1 - q, q). For \(p > \frac{2}{3}\) player II strictly prefers \(s\) to \(c\) which defines the unique best response(0,1). Similarly, the best-response mixed strategy(1 - p, p)of player \(\mathrm{I}\) against the mixed strategy(1 - q, q)is \(p = 0\) for \(q \in  \left\lbrack  {0,\frac{1}{3}}\right)\) , any \(p \in  \left\lbrack  {0,1}\right\rbrack\) for \(q = \frac{1}{3}\) , and \(p = 1\) for \(q \in  \left( {\frac{1}{3},1}\right\rbrack\) (an example to check this quickly is \(q = \frac{1}{2}\) , where strategy \(S\) gives player \(\mathrm{I}\) a higher expected payoff). Apart from the pure-strategy equilibria, this gives the mixed equilibrium \(\left( {\left( {\frac{1}{3},\frac{2}{3}}\right) ,\left( {\frac{2}{3},\frac{1}{3}}\right) }\right)\) .

\begin{tcolorbox}
在具有纯策略均衡的 \(2 \times  2\) 博弈中，也可能存在混合策略均衡。例如，考虑图6.7左侧的性别战(Battle of the Sexes)博弈。最优反应效用表明，该博弈存在纯策略均衡 (C, c) 和 (S, s)。现在假设参与者I采用混合策略 \(x = \left( {1 - p,p}\right)\) ，参与者II采用混合策略 \(y = \left( {1 - q,q}\right)\) 。与图6.4中的监管博弈(Inspection game)类似，图6.8展示了参与者针对对方混合策略的最优反应。例如，如果 \(p = \frac{1}{2}\) ，那么参与者II在列 \(c\) 的期望效用(expected payoff)为1，在列 \(s\) 的期望效用为 \(\frac{1}{2}\) ，她的最优反应是 \(c\) ，即混合策略 (1,0)，其中 \(q = 0\) 。一般来说， \(c\) 和 \(s\) 的效用分别为 \(2\left( {1 - p}\right)\) 和 \(p\) ，因此如果 \(2\left( {1 - p}\right)  \geq  p\) 或者等价地 \(\frac{2}{3} \geq  p\) ，即 \(p \in  \left\lbrack  {0,\frac{2}{3}}\right\rbrack\) ，参与者II弱偏好 \(c\) 而非 \(s\) 。对于 \(p = \frac{2}{3}\) ，参与者II无差异，任何 \(q \in  \left\lbrack  {0,1}\right\rbrack\) 都定义了一个最优反应 (1 - q, q)。对于 \(p > \frac{2}{3}\) ，参与者II严格偏好 \(s\) 而非 \(c\) ，这定义了唯一的最优反应 (0,1)。类似地，参与者 \(\mathrm{I}\) 针对混合策略 (1 - q, q) 的最优反应混合策略 (1 - p, p) 为:当 \(q \in  \left\lbrack  {0,\frac{1}{3}}\right)\) 时为 \(p = 0\) ；当 \(q = \frac{1}{3}\) 时，任何 \(p \in  \left\lbrack  {0,1}\right\rbrack\) 均可；当 \(q \in  \left( {\frac{1}{3},1}\right\rbrack\) 时为 \(p = 1\) (快速验证的一个例子是 \(q = \frac{1}{2}\) ，此时策略 \(S\) 给参与者 \(\mathrm{I}\) 带来更高的期望效用)。除了纯策略均衡之外，这给出了混合均衡(mixed equilibrium) \(\left( {\left( {\frac{1}{3},\frac{2}{3}}\right) ,\left( {\frac{2}{3},\frac{1}{3}}\right) }\right)\) 。
\end{tcolorbox}

In summary, the rule to find a mixed-strategy equilibrium in a \(2 \times  2\) game is to make the other player indifferent, because only in that case the other player can mix. Then the expected payoffs for the two opponent strategies have to be equal, and the resulting equation determines the player's own probability.

\begin{tcolorbox}
总之，在 \(2 \times  2\) 博弈中寻找混合策略均衡的规则是让对方参与者无差异，因为只有在这种情况下，对方参与者才会采用混合策略。然后，两个对手策略的期望效用必须相等，由此得到的方程决定了参与者自身的概率。
\end{tcolorbox}

The mixed-strategy probabilities in the Battle of the Sexes game can be seen relatively quickly by looking at the payoff matrices: For example, player II must give twice as much weight to the left column \(c\) as to the right column \(s\) , because player I’s top row \(C\) only gets payoff 1 when player II plays \(c\) , whereas his bottom row \(S\) gets payoff 2 when she plays \(s\) , and the payoffs for the strategy pairs(C, s) and(S, c)are zero. Hence, the two rows \(C\) and \(S\) have the same expected payoff only when player II mixes \(c\) and \(s\) with probabilities \(\frac{2}{3}\) and \(\frac{1}{3}\) , respectively.

\begin{tcolorbox}
通过观察效用矩阵，可以相对快速地看出性别战(Battle of the Sexes)博弈中的混合策略概率:例如，参与者II对左列 \(c\) 的权重必须是右列 \(s\) 的两倍，因为当参与者II选择 \(c\) 时，参与者I的顶行 \(C\) 只能获得效用1，而当她选择 \(s\) 时，他的底行 \(S\) 能获得效用2，并且策略组合 (C, s) 和 (S, c) 的效用为零。因此，只有当参与者II分别以概率 \(\frac{2}{3}\) 和 \(\frac{1}{3}\) 混合 \(c\) 和 \(s\) 时，两行 \(C\) 和 \(S\) 才具有相同的期望效用。
\end{tcolorbox}

\begin{center}
\includegraphics[max width=0.8\textwidth]{images/bo_d2e5cmbef24c73b03tn0_166_287_1467_1069_416_0.jpg}
\end{center}
\hspace*{3em} 

Figure 6.9 The "difference trick" for finding mixed equilibrium probabilities in a \(2 \times  2\) game. The left figure shows the game, and the difference in payoffs to the other player for each pure strategy. As shown on the right, these differences are assigned to the respective other own strategy and are re-normalized to become probabilities. The fractions \(\frac{11}{5}\) and \(\frac{38}{10}\) are the resulting equal expected payoffs to player II and player I, respectively.

\begin{tcolorbox}
图6.9 用于在 \(2 \times  2\) 博弈中寻找混合均衡(mixed equilibrium)概率的“差分技巧”(difference trick)。左图展示了该博弈，以及每个纯策略下另一个参与者的收益差异。如右图所示，这些差异被分配给各自的另一个自身策略，并重新归一化为概率。分数 \(\frac{11}{5}\) 和 \(\frac{38}{10}\) 分别是参与者II和参与者I的期望效用(expected payoff)相等的结果。
\end{tcolorbox}

We now explain a quick method, the "difference trick", for finding the probabilities in a mixed equilibrium of a general \(2 \times  2\) game. Consider the game on the left in Figure 6.9. When player I plays \(T\) , then \(l\) is a best response, and when player I plays \(B\) , then \(r\) is a best response. Consequently, there must be a way to mix \(T\) and \(B\) so that player II is indifferent between \(l\) and \(r\) . We consider now the difference \(\Delta\) in payoffs to the other player for both rows: When player \(\mathrm{I}\) plays \(T\) , then the difference between the two payoffs to player II is \(\Delta  = 2 - 1 = 1\) , and when player I plays \(B\) , then that difference, in absolute value, is \(\Delta  = 7 - 3 = 4\) , as shown on the side of the game. (Note that these differences, when considered as differences between the payoffs for \(l\) versus \(r\) , have opposite sign because \(l\) is preferred to \(r\) against \(T\) , and the other way around against \(B\) ; otherwise, player II would always prefer the same strategy and could not be made indifferent.) Now the payoff difference \(\Delta\) is assigned as a probability weight to the respective other strategy of the row player, meaning \(T\) is given weight 4 (which is the \(\Delta\) computed for \(B\) ), and \(B\) is given weight 1 (which is the \(\Delta\) for \(T\) ). The probabilities for \(T\) and \(B\) are then chosen proportional to these weights, so one has to divide each weight by 5 (which is the sum of the weights) in order to obtain probabilities. This is shown on the right in Figure 6.9. The expected payoffs to player II are also shown there, at the bottom, and are for \(l\) given by \(\left( {4 \cdot  2 + 1 \cdot  3}\right) /5 = \frac{11}{5}\) , and for \(r\) by \(\left( {4 \cdot  1 + 1 \cdot  7}\right) /5 = \frac{11}{5}\) , so they are indeed equal as claimed. Similarly, the two payoff differences for player I in columns \(l\) and \(r\) are 3 and 7, respectively, so \(l\) and \(r\) should be played with probabilities that are proportional to 7 and 3, respectively. With the resulting probabilities \(\frac{7}{10}\) and \(\frac{3}{10}\) , the two rows \(T\) and \(B\) get the same expected payoff \(\frac{38}{10}\) .

\begin{tcolorbox}
我们现在解释一种快速方法，即“差分技巧”(difference trick)，用于在一般的 \(2 \times  2\) 博弈中寻找混合均衡(mixed equilibrium)的概率。考虑图6.9左侧的博弈。当参与者I选择 \(T\) 时， \(l\) 是最优反应(best response)；当参与者I选择 \(B\) 时， \(r\) 是最优反应。因此，必然存在一种方法来混合 \(T\) 和 \(B\) ，使得参与者II在 \(l\) 和 \(r\) 之间无差异。我们现在考虑两行中另一个参与者的收益差异 \(\Delta\) :当参与者 \(\mathrm{I}\) 选择 \(T\) 时，参与者II的两个收益之间的差异是 \(\Delta  = 2 - 1 = 1\) ；当参与者I选择 \(B\) 时，该差异的绝对值是 \(\Delta  = 7 - 3 = 4\) ，如博弈图旁边所示。(注意，当将这些差异视为 \(l\) 与 \(r\) 的收益之间的差异时，它们的符号相反，因为在面对 \(T\) 时， \(l\) 优于 \(r\) ，而在面对 \(B\) 时则相反；否则，参与者II将总是偏好相同的策略，并且无法使其无差异。)现在，收益差异 \(\Delta\) 被作为概率权重分配给行参与者的各自另一个策略，即 \(T\) 的权重为4(这是为 \(B\) 计算的 \(\Delta\) )， \(B\) 的权重为1(这是为 \(T\) 计算的 \(\Delta\) )。然后， \(T\) 和 \(B\) 的概率选择与这些权重成比例，因此必须将每个权重除以5(这是权重之和)以获得概率。这在图6.9的右侧展示。参与者II的期望效用也显示在那里，在底部， \(l\) 的期望效用由 \(\left( {4 \cdot  2 + 1 \cdot  3}\right) /5 = \frac{11}{5}\) 给出， \(r\) 的期望效用由 \(\left( {4 \cdot  1 + 1 \cdot  7}\right) /5 = \frac{11}{5}\) 给出，因此它们确实如所声称的那样相等。类似地，参与者I在列 \(l\) 和 \(r\) 中的两个收益差异分别为3和7，因此 \(l\) 和 \(r\) 的选择概率应分别与7和3成比例。使用得到的概率 \(\frac{7}{10}\) 和 \(\frac{3}{10}\) ，两行 \(T\) 和 \(B\) 获得相同的期望效用 \(\frac{38}{10}\) 。
\end{tcolorbox}

\(\Rightarrow\) Exercise 6.4 asks you to prove that this "difference trick" always works.

\begin{tcolorbox}
\(\Rightarrow\) 练习6.4要求你证明这种“差分技巧”(difference trick)总是有效的。
\end{tcolorbox}

Next, we consider games where one player has two and the other more than two strategies. The \(3 \times  2\) game on the right in Figure 6.7 is like the Battle of the Sexes game, except that player I has an additional strategy B. Essentially, such a game can be analyzed by considering the \(2 \times  2\) games obtained by restricting both players to two strategies only, and checking if the resulting equilibria carry over to the whole game. First, the pure-strategy equilibrium(C, c)of the original Battle of the Sexes game is also an equilibrium of the larger game, but(S, s)is not because \(S\) is no longer a best response to \(s\) because \(B\) gives a larger payoff to player I. Second, consider the mixed-strategy equilibrium of the smaller game where player I chooses \(C\) and \(S\) with probabilities \(\frac{1}{3}\) and \(\frac{2}{3}\) (and \(B\) with probability zero), so that player II is indifferent between \(c\) and \(s\) . Hence, the mixed strategy \(\left( {\frac{2}{3},\frac{1}{3}}\right)\) of player II is still a best response, against which \(C\) and \(S\) both give the same expected payoff \(\frac{2}{3}\) . However, this is not enough to guarantee an equilibrium, because \(C\) and \(S\) have to be best responses, that is, their payoff must be at least as large as that for the additional strategy \(B\) . That payoff is \(\left( {-1}\right)  \cdot  \frac{2}{3} + 3 \cdot  \frac{1}{3} = \frac{1}{3}\) , so it is indeed not larger than the payoff for \(C\) and \(S\) . In other words, the mixed equilibrium of the smaller game is also a mixed equilibrium of the larger game, given by the pair of mixed strategies \(\left( {\left( {\frac{1}{3},\frac{2}{3},0}\right) ,\left( {\frac{2}{3},\frac{1}{3}}\right) }\right)\) .

\begin{tcolorbox}
接下来，我们考虑一方玩家有两种策略而另一方玩家有超过两种策略的博弈。图6.7右侧的 \(3 \times  2\) 博弈与性别战(Battle of the Sexes)博弈类似，不同之处在于玩家I有一个额外的策略B。本质上，此类博弈可以通过考虑仅限制双方玩家使用两种策略而得到的 \(2 \times  2\) 博弈来进行分析，并检查由此得到的均衡是否适用于整个博弈。首先，原始性别战博弈的纯策略均衡(C, c)也是这个更大博弈的一个均衡，但(S, s)不是，因为 \(S\) 不再是对 \(s\) 的最优反应(best response)，因为 \(B\) 给玩家I带来了更高的收益。其次，考虑较小博弈的混合策略均衡(mixed equilibrium)，其中玩家I以概率 \(\frac{1}{3}\) 和 \(\frac{2}{3}\) 选择 \(C\) 和 \(S\) (以概率零选择 \(B\) )，从而使玩家II在 \(c\) 和 \(s\) 之间无差异。因此，玩家II的混合策略 \(\left( {\frac{2}{3},\frac{1}{3}}\right)\) 仍然是最优反应，在这种情况下， \(C\) 和 \(S\) 都给出相同的期望效用(expected payoff) \(\frac{2}{3}\) 。然而，这并不足以保证达到均衡，因为 \(C\) 和 \(S\) 必须是最优反应，即它们的收益必须至少与额外策略 \(B\) 的收益一样大。该收益为 \(\left( {-1}\right)  \cdot  \frac{2}{3} + 3 \cdot  \frac{1}{3} = \frac{1}{3}\) ，所以它确实不大于 \(C\) 和 \(S\) 的收益。换句话说，较小博弈的混合均衡也是较大博弈的混合均衡，由混合策略对 \(\left( {\left( {\frac{1}{3},\frac{2}{3},0}\right) ,\left( {\frac{2}{3},\frac{1}{3}}\right) }\right)\) 给出。
\end{tcolorbox}

Other "restricted" \(2 \times  2\) games are obtained by letting player I play only \(C\) and \(B\) , or only \(S\) and \(B\) . In the first case, the difference trick gives the mixed strategy \(\left( {\frac{3}{5},\frac{2}{5}}\right)\) of player II (for playing \(c,s\) ) so that player I is indifferent between \(C\) and \(B\) , where both strategies receive expected payoff \(\frac{3}{5}\) . However, the expected payoff to the third strategy \(S\) is then \(\frac{4}{5}\) , which is higher, so that \(C\) and \(B\) are not best responses, which means we cannot have an equilibrium where only \(C\) and \(B\) are played with positive probability by player I. Another reason why no equilibrium strategy of player II mixes \(C\) and \(B\) is that against both \(C\) and \(B\) , player II’s best response is always \(c\) , so that player I could not make player II indifferent by playing only \(C\) and \(B\) .

\begin{tcolorbox}
其他“受限”的 \(2 \times  2\) 博弈可通过让玩家I仅采用 \(C\) 和 \(B\) 策略，或者仅采用 \(S\) 和 \(B\) 策略来获得。在第一种情况下，差分技巧(difference trick)给出了玩家II的混合策略 \(\left( {\frac{3}{5},\frac{2}{5}}\right)\) (用于采用 \(c,s\) )，使得玩家I在 \(C\) 和 \(B\) 之间无差异，此时这两种策略的期望效用均为 \(\frac{3}{5}\) 。然而，第三种策略 \(S\) 的期望效用为 \(\frac{4}{5}\) ，这更高，因此 \(C\) 和 \(B\) 不是最优反应，这意味着我们无法得到一个玩家I仅以正概率采用 \(C\) 和 \(B\) 的均衡。玩家II没有混合采用 \(C\) 和 \(B\) 的均衡策略的另一个原因是，对于 \(C\) 和 \(B\) ，玩家II的最优反应始终是 \(c\) ，因此玩家I无法通过仅采用 \(C\) 和 \(B\) 使玩家II无差异。
\end{tcolorbox}

Finally, consider \(S\) and \(B\) as pure strategies that player I mixes in a possible equilibrium. This requires, via the difference trick, the mixed strategy \(\left( {0,\frac{3}{4},\frac{1}{4}}\right)\) of player I (as a probability vector for playing his three pure strategies \(C,S,B\) ) so that player II is indifferent between \(c\) and \(s\) with payoff 1 for both strategies. Then if player II uses the mixed strategy \(\left( {\frac{1}{2},\frac{1}{2}}\right)\) (in order to make player I indifferent between \(S\) and \(B\) ), the expected payoffs for the three rows \(C,S,B\) are \(\frac{1}{2},1\) , and 1, respectively, so that indeed player I only uses best responses with positive probability, which gives a third equilibrium of the game.

\begin{tcolorbox}
最后，考虑将 \(S\) 和 \(B\) 作为玩家I在可能的均衡中混合采用的纯策略。通过差分技巧，这需要玩家I采用混合策略 \(\left( {0,\frac{3}{4},\frac{1}{4}}\right)\) (作为采用他的三个纯策略 \(C,S,B\) 的概率向量)，从而使玩家II在 \(c\) 和 \(s\) 之间无差异，且两种策略的收益均为1。然后，如果玩家II采用混合策略 \(\left( {\frac{1}{2},\frac{1}{2}}\right)\) (以使玩家I在 \(S\) 和 \(B\) 之间无差异)，则三行策略 \(C,S,B\) 的期望效用分别为 \(\frac{1}{2},1\) 和1，因此玩家I确实仅以正概率采用最优反应，这给出了该博弈的第三个均衡。
\end{tcolorbox}

Why is there no mixed equilibrium of the \(3 \times  2\) game in Figure 6.7 where player I mixes between his three pure strategies? The reason is that player II has only a single probability at her disposal (which determines the complementary probability for the other pure strategy), which does not give enough freedom to satisfy two equations, namely indifference between \(C\) and \(S\) as well as between \(S\) and \(B\) (the third indifference between \(C\) and \(B\) would then hold automatically). We have already computed the probabilities \(\left( {\frac{2}{3},\frac{1}{3}}\right)\) for \(c,s\) that are needed to give the same expected payoff for \(C\) and \(S\) , against which the payoff for row \(B\) is different, namely \(\frac{1}{3}\) . This alone suffices to show that it is not possible to make player I indifferent between all three pure strategies, so there is no equilibrium where player I mixes between all of them. In certain games, which are called "degenerate" and which are treated in Section 6.8, it is indeed the case that player II, say, mixes only between two pure strategies, but where "by accident" three strategies of player I have the same optimal payoff. Equilibria of degenerate games are more complicated to determine, as discussed in Section 6.8.

\begin{tcolorbox}
为什么图6.7中的 \(3 \times  2\) 博弈不存在参与者I在其三种纯策略之间进行混合的混合均衡(mixed equilibrium)呢？原因在于参与者II只能使用单一概率(该概率决定了另一种纯策略的互补性(complementary)概率)，这不足以满足两个方程，即 \(C\) 和 \(S\) 之间的无差异(indifference)以及 \(S\) 和 \(B\) 之间的无差异(第三个无差异，即 \(C\) 和 \(B\) 之间的无差异会自动成立)。我们已经计算出了使得 \(C\) 和 \(S\) 具有相同期望效用(expected payoff)所需的 \(c,s\) 的概率 \(\left( {\frac{2}{3},\frac{1}{3}}\right)\) ，而在此情况下，行 \(B\) 的效用不同，即为 \(\frac{1}{3}\) 。仅此一点就足以表明，不可能使参与者I在所有三种纯策略之间保持无差异，因此不存在参与者I在所有这些策略之间进行混合的均衡。在某些被称为“退化博弈(degenerate game)”的博弈中(将在6.8节中讨论)，确实会出现这样的情况，例如参与者II仅在两种纯策略之间进行混合，但“偶然地”参与者I的三种策略具有相同的最优效用。如6.8节所述，确定退化博弈的均衡更为复杂。
\end{tcolorbox}

\(\Rightarrow\) Attempt Exercise 6.2 if you have not done this earlier. Compare your solution with the approach for the \(3 \times  2\) game in Figure 6.7.

\begin{tcolorbox}
\(\Rightarrow\) 如果你之前没有做过，请尝试练习6.2。将你的解决方案与图6.7中 \(3 \times  2\) 博弈的方法进行比较。
\end{tcolorbox}

\section*{6.7 The Upper-Envelope Method}

\begin{tcolorbox}
\section*{6.7 上包络(upper envelope)法}
\end{tcolorbox}

The topic of this section is the upper-envelope diagram for finding mixed equilibria, in particular in games where one player has only two strategies.

\begin{tcolorbox}
本节的主题是用于寻找混合均衡(mixed equilibrium)的上包络图，特别是在其中一个参与者只有两种策略的博弈中。
\end{tcolorbox}

\begin{center}
\includegraphics[max width=0.9\textwidth]{images/bo_d2e5cmbef24c73b03tn0_169_338_452_1244_595_0.jpg}
\end{center}
\hspace*{3em} 

Figure 6.10 Upper envelope of expected payoffs to player I against the mixed strategy of player II for the two games in Figure 6.7. The white circles indicate mixed strategies of player II that are part of an equilibrium.

\begin{tcolorbox}
图6.10 针对图6.7中两个博弈，参与者II的混合策略下参与者I的期望效用的上包络。白色圆圈表示作为均衡一部分的参与者II的混合策略。
\end{tcolorbox}

The graphs in Figure 6.10 show the upper-envelope diagrams for the two games in Figure 6.7. This diagram is a plot of the expected payoff of one player against the mixed strategy of the other player. We use this typically when that mixed strategy is determined by a single probability, which applies when the player has only two pure strategies. In this example, player II has in both games only the two strategies \(c\) and \(s\) , and her mixed strategy is(1 - q, q). The horizontal axis represents the probability \(q = \operatorname{prob}\left( s\right)\) for playing the second pure strategy \(s\) of player II. In the vertical direction, we plot the resulting expected payoffs to the other player, here player I, for each of his pure strategies.

\begin{tcolorbox}
图6.10中的图展示了图6.7中两个博弈的上包络图。该图绘制了一个参与者的期望效用与另一个参与者的混合策略之间的关系。当混合策略由单一概率决定时，我们通常会使用这种方法，这适用于参与者只有两种纯策略的情况。在这个例子中，在两个博弈中参与者II都只有两种策略 \(c\) 和 \(s\) ，她的混合策略是(1 - q, q)。横轴表示参与者II采用第二种纯策略 \(s\) 的概率 \(q = \operatorname{prob}\left( s\right)\) 。在纵轴方向上，我们绘制了另一个参与者(这里是参与者I)对于他的每种纯策略所得到的期望效用。
\end{tcolorbox}

Consider first the left game, where player I has only the two strategies \(C\) and \(S\) . For \(C\) , his expected payoff against(1 - q, q)is \(1 - q\) , and for \(S\) it is \({2q}\) . The upper envelope is the maximum of the resulting two lines, shown in bold.

\begin{tcolorbox}
首先考虑左边的博弈，其中参与者I只有两种策略 \(C\) 和 \(S\) 。对于 \(C\) ，他针对(1 - q, q)的期望效用是 \(1 - q\) ，对于 \(S\) ，期望效用是 \({2q}\) 。上包络是得到的两条线中的最大值，用粗线表示。
\end{tcolorbox}

The right game has an additional strategy \(B\) for player \(\mathrm{I}\) , with payoff -1 against the left column \(c\) where \(q = 0\) , and payoff 3 against the right column \(s\) where \(q = 1\) . For \(q \in  \left\lbrack  {0,1}\right\rbrack\) , the resulting expected payoff is \(- \left( {1 - q}\right)  + {3q}\) , that is, \({4q} - 1\) . This expected payoff becomes a new part of the upper envelope for \(q \geq  \frac{1}{2}\) , because then the expected payoffs for rows \(S\) and \(B\) are equal to 1 when \(q = \frac{1}{2}\) , where the two lines cross, and for \(q > \frac{1}{2}\) the expected payoff for row \(B\) is strictly better than for the other rows.

\begin{tcolorbox}
右边的博弈中参与者 \(\mathrm{I}\) 有一个额外的策略 \(B\) ，在左列 \(c\) (其中 \(q = 0\) )下的效用为 -1，在右列 \(s\) (其中 \(q = 1\) )下的效用为3。对于 \(q \in  \left\lbrack  {0,1}\right\rbrack\) ，得到的期望效用是 \(- \left( {1 - q}\right)  + {3q}\) ，即 \({4q} - 1\) 。当 \(q \geq  \frac{1}{2}\) 时，这个期望效用成为上包络的新部分，因为当 \(q = \frac{1}{2}\) 时，行 \(S\) 和 \(B\) 的期望效用等于1，此时两条线相交，并且当 \(q > \frac{1}{2}\) 时，行 \(B\) 的期望效用严格优于其他行。
\end{tcolorbox}

The lines that describe the expected payoffs are particularly easy to obtain graphically, according to the following goalpost method for constructing the upper envelope:

\begin{tcolorbox}
根据以下用于构建上包络的球门柱法(goalpost method)，描述期望效用的线特别容易通过图形方式获得:
\end{tcolorbox}

\begin{itemize}
\item Identify a player who has only two pure strategies (here player II).
\end{itemize}

\begin{tcolorbox}
\begin{itemize}
\item 确定一个只有两种纯策略的参与者(这里是参与者II)。
\end{itemize}
\end{tcolorbox}

\begin{itemize}
\item Plot the probability of the second strategy of the player along a horizontal axis from 0 to 1, here \(q = \operatorname{prob}\left( s\right)\) . This defines the mixed strategy(1 - q, q)of player II. In that way, the left end 0 of the interval \(\left\lbrack  {0,1}\right\rbrack\) for \(q\) corresponds to the left strategy of player II, and the right end 1 to her right strategy in the payoff table.
\end{itemize}

\begin{tcolorbox}
\begin{itemize}
\item 沿水平轴从0到1绘制玩家第二个策略的概率，这里 \(q = \operatorname{prob}\left( s\right)\) 。这就定义了玩家II的混合策略(1 - q, q)。这样，区间 \(\left\lbrack  {0,1}\right\rbrack\) 中对应 \(q\) 的左端点0对应玩家II在收益表中的左策略，右端点1对应其右策略。
\end{itemize}
\end{tcolorbox}

\begin{itemize}
\item Erect a vertical line (a "goalpost") at each endpoint of the interval \(\left\lbrack  {0,1}\right\rbrack\) .
\end{itemize}

\begin{tcolorbox}
\begin{itemize}
\item 在区间 \(\left\lbrack  {0,1}\right\rbrack\) 的每个端点处竖立一条垂直线(“球门柱”)。
\end{itemize}
\end{tcolorbox}

\begin{itemize}
\item For each pure strategy of the other player (here player I), do the following: At the left goalpost, mark the payoff to player I against the left strategy of player II as a height on the goalpost, and similarly at the right goalpost for the right strategy of player II; then connect these two heights by a straight line. For row \(C\) , the left and right heights are 1 and 0 ; for row \(S\) , they are 0 and 2 ; and for row \(B\) (right game only), they are -1 and 3 . These are exactly the entries in the rows of the payoff matrix to player I.
\end{itemize}

\begin{tcolorbox}
\begin{itemize}
\item 对于另一个玩家(这里是玩家I)的每个纯策略，执行以下操作:在左球门柱上，将玩家I对抗玩家II左策略的收益标记为球门柱上的一个高度，同样在右球门柱上标记玩家I对抗玩家II右策略的收益；然后用一条直线连接这两个高度。对于行 \(C\) ，左右高度分别为1和0；对于行 \(S\) ，分别为0和2；对于行 \(B\) (仅适用于右侧博弈)，分别为 -1和3。这些恰好是收益矩阵中玩家I所在行的元素。
\end{itemize}
\end{tcolorbox}

\section*{The maximum of the plotted lines is the upper envelope.}

\begin{tcolorbox}
\section*{所绘制直线的最大值就是上包络。}
\end{tcolorbox}

The fact that this construction works, in particular that one obtains straight lines for the expected payoffs, is due to the algebraic description of line segments explained in Figure 5.4 above. Namely, consider some row with payoffs \(a\) and \(b\) to player \(\mathrm{I}\) , like \(a =  - 1\) and \(b = 3\) for row \(B\) . Then his expected payoff against the mixed strategy(1 - q, q)is \(\left( {1 - q}\right) a + {qb}\) . This expected payoff is the vertical coordinate ("height") in the goalpost diagram, whereas the horizontal coordinate is just \(q\) . The two endpoints on the goalposts have coordinates(0, a)and(1, b), and the line segment that is drawn to connect them consists of the points \(\left( {1 - q}\right) \left( {0,a}\right)  + q\left( {1,b}\right)\) for \(q \in  \left\lbrack  {0,1}\right\rbrack\) , which are the points \(\left( {q,\left( {1 - q}\right) a + {qb}}\right)\) , as claimed.

\begin{tcolorbox}
这种构造方法可行，特别是能得到期望效用的直线，这是由于上述图5.4中对线段的代数描述。即，考虑某一行，玩家 \(\mathrm{I}\) 的收益为 \(a\) 和 \(b\) ，例如行 \(B\) 的 \(a =  - 1\) 和 \(b = 3\) 。那么他对抗混合策略(1 - q, q)的期望效用为 \(\left( {1 - q}\right) a + {qb}\) 。这个期望效用是球门柱图中的垂直坐标(“高度”)，而水平坐标就是 \(q\) 。球门柱上的两个端点坐标为(0, a)和(1, b)，连接它们所绘制的线段由 \(\left( {1 - q}\right) \left( {0,a}\right)  + q\left( {1,b}\right)\) (其中 \(q \in  \left\lbrack  {0,1}\right\rbrack\) )组成，也就是所声称的 \(\left( {q,\left( {1 - q}\right) a + {qb}}\right)\) 这些点。
\end{tcolorbox}

In the left diagram in Figure 6.10, the two lines of expected payoffs for rows \(C\) and \(S\) obviously intersect exactly where player \(\mathrm{I}\) is indifferent between these two rows, which happens when \(q = \frac{1}{3}\) . However, the diagram not only gives information about this indifference, but about player I's preference in general: Any pure strategy that is a best response of player I is obviously the topmost line segment of all the lines describing the expected payoffs. Here, row \(C\) is a best response whenever \(0 \leq  q \leq  \frac{1}{3}\) , and row \(S\) is a best response whenever \(\frac{1}{3} \leq  q \leq  1\) , with indifference between \(C\) and \(S\) exactly when \(q = \frac{1}{3}\) . The upper envelope of expected payoffs is the maximum of the lines that describe the expected payoffs for the different pure strategies, and it is indicated by the bold line segments in the diagram. Marking this upper envelope in bold is therefore the final step of the goalpost method.

\begin{tcolorbox}
在图6.10的左图中，行 \(C\) 和 \(S\) 的两条期望效用线显然恰好相交于玩家 \(\mathrm{I}\) 对这两行无差异的地方，即 \(q = \frac{1}{3}\) 时。然而，该图不仅给出了关于这种无差异的信息，还给出了玩家I的一般偏好信息:玩家I的任何最优反应纯策略显然是描述期望效用的所有直线中最上方的线段。这里，当 \(0 \leq  q \leq  \frac{1}{3}\) 时，行 \(C\) 是最优反应；当 \(\frac{1}{3} \leq  q \leq  1\) 时，行 \(S\) 是最优反应，而当 \(q = \frac{1}{3}\) 时，玩家I对 \(C\) 和 \(S\) 无差异。期望效用的上包络是描述不同纯策略期望效用的直线的最大值，在图中用粗线段表示。因此，用粗线标记这个上包络是球门柱法的最后一步。
\end{tcolorbox}

The upper envelope is of particular use when a player has more than two pure strategies, because mixed equilibria of a \(2 \times  2\) are anyhow quickly determined with the "difference trick". Consider the upper-envelope diagram on the right of Figure 6.10, for the \(3 \times  2\) game in Figure 6.7. The extra strategy \(B\) of player \(\mathrm{I}\) gives the line connecting height -1 on the left goalpost to height 3 on the right goalpost, which shows that \(B\) is a best response for all \(q\) so that \(\frac{1}{2} \leq  q \leq  1\) , with indifference between \(S\) and \(B\) when \(q = \frac{1}{2}\) . The upper envelope consists of three line segments for the three pure strategies \(C,S\) , and \(B\) of player I. Moreover, there are only two points where player \(\mathrm{I}\) is indifferent between any two such strategies, namely between \(C\) and \(S\) when \(q = \frac{1}{3}\) , and between \(S\) and \(B\) when \(q = \frac{1}{2}\) . These two points, indicated by small circles on the upper envelope, give the only two mixed-strategy probabilities of player II where player I can mix between two pure strategies (the small circle for \(q = 0\) corresponds to the pure equilibrium strategy \(c\) ). There is a third indifference between \(C\) and \(B\) , for \(q = \frac{2}{5}\) , where the two lines for \(C\) and \(B\) intersect, but this intersection point is below the line for \(S\) and hence not on the upper envelope. Consequently, this intersection point does not have to be investigated as a possible mixed-strategy equilibrium strategy where player I mixes between the two rows \(C\) and \(B\) because they are not best responses.

\begin{tcolorbox}
当一个玩家有两个以上的纯策略时，上包络特别有用，因为退化博弈(degenerate game)的混合均衡无论如何都可以通过“差分技巧(difference trick)”快速确定。考虑图6.10右侧的上包络图，它对应图6.7中的退化博弈。玩家 \(\mathrm{I}\) 的额外策略 \(B\) 给出了一条从左端点高度 -1 连接到右端点高度 3 的线段，这表明对于所有的 \(q\) ， \(B\) 都是最优反应(best response)，使得 \(\frac{1}{2} \leq  q \leq  1\) ，并且当 \(q = \frac{1}{2}\) 时，玩家在 \(S\) 和 \(B\) 之间无差异(indifference)。上包络由玩家 I 的三个纯策略 \(C,S\) 和 \(B\) 对应的三条线段(line segments)组成。此外，玩家 \(\mathrm{I}\) 在任意两个这样的策略之间无差异的点只有两个，即当 \(q = \frac{1}{3}\) 时，在 \(C\) 和 \(S\) 之间；当 \(q = \frac{1}{2}\) 时，在 \(S\) 和 \(B\) 之间。这两个点在上包络上用小圆圈表示，它们给出了玩家 II 的仅有的两个混合策略概率，在这些概率下玩家 I 可以在两个纯策略之间进行混合( \(q = 0\) 对应的小圆圈对应纯均衡(pure equilibrium)策略 \(c\) )。在 \(C\) 和 \(B\) 之间还有第三种无差异情况，即当 \(q = \frac{2}{5}\) 时， \(C\) 和 \(B\) 对应的两条线相交，但这个交点在 \(S\) 对应的线下方，因此不在上包络上。因此，这个交点不需要作为玩家 I 在两行 \(C\) 和 \(B\) 之间进行混合的可能的混合策略均衡(equilibrium)策略来研究，因为它们不是最优反应。
\end{tcolorbox}

Given these candidates of mixtures for player I, we can quickly find if and how the two respective pure strategies can be mixed in order to make player II indifferent using the difference trick, as described in the previous section.

\begin{tcolorbox}
给定玩家 I 的这些混合策略候选，我们可以如前一节所述，使用差分技巧快速找出这两个相应的纯策略是否可以以及如何混合，以使玩家 II 无差异。
\end{tcolorbox}

Hence, the upper envelope restricts the pairs of strategies that need to be tested as possible strategies that are mixed in equilibrium. Exercise 6.5, for example, describes a \(2 \times  5\) game. Here it is player I who has two strategies, and the upper envelope describes the expected payoffs to player II against player I's mixed strategy (1 - p, p). Because player II has five pure strategies, the upper envelope may consist of up to five line segments; there may be fewer line segments if some pure strategy of player II is never a best response. Consequently, there are up to four intersection points where these line segments join up on the upper envelope, where player II is indifferent between her corresponding two strategies and can mix between them. Without the upper envelope, there would not be four but ten (which are the different ways to pick two out of five) pure strategy pairs to be checked as possible pure strategies that player II mixes in equilibrium, which would be much more laborious to investigate.

\begin{tcolorbox}
因此，上包络限制了需要作为可能的均衡混合策略进行测试的策略对。例如，练习 6.5 描述了一个退化博弈。在这里，玩家 I 有两个策略，上包络描述了玩家 II 针对玩家 I 的混合策略 (1 - p, p) 的期望效用(expected payoff)。因为玩家 II 有五个纯策略，上包络最多可能由五条线段组成；如果玩家 II 的某些纯策略永远不是最优反应，线段可能会更少。因此，在上包络上线段连接的交点最多有四个，在这些点上玩家 II 在她对应的两个策略之间无差异，并且可以在它们之间进行混合。如果没有上包络，需要检查的可能是玩家 II 在均衡中混合的纯策略对不是四个，而是十个(即从五个中选两个的不同方式)，这将需要进行更繁琐的研究。
\end{tcolorbox}

\(\Rightarrow\) Exercise 6.2 can also be solved with the upper-envelope method. This method is of particular help for larger games as in Exercise 6.5.

\begin{tcolorbox}
\(\Rightarrow\) 练习 6.2 也可以用上包络方法解决。对于像练习 6.5 中的较大博弈，这种方法特别有帮助。
\end{tcolorbox}

In principle, the upper-envelope diagram is also defined for the expected payoffs against a mixed strategy that mixes between more than two pure strategies. This diagram is harder to draw and visualize in this case because it requires higher dimensions. An example is shown in Figure 9.2 in Chapter 9, where we develop the upper-envelope method for finding equilibria further.

\begin{tcolorbox}
原则上，上包络图也定义了针对在两个以上纯策略之间进行混合的混合策略的期望效用。在这种情况下，这个图更难绘制和可视化，因为它需要更高的维度(dimension)。第 9 章的图 9.2 展示了一个例子，我们将进一步开发上包络方法来寻找均衡。
\end{tcolorbox}

\(\Rightarrow\) The simple Exercise 6.6 helps understanding dominated strategies with the upper-envelope diagram.

\begin{tcolorbox}
\(\Rightarrow\) 简单的练习 6.6 有助于通过上包络图理解(被占优)策略(dominated strategy)。
\end{tcolorbox}

\section*{6.8 Degenerate Games}

\begin{tcolorbox}
\section*{6.8 退化博弈}
\end{tcolorbox}

In this section, we treat games that have to be analyzed more carefully in order to find all their equilibria. We completely analyze an example, and then give the definition of a degenerate game that applies to \(2 \times  2\) games like in this example. We then discuss larger games, and give the general Definition 6.7 for when an \(m \times  n\) game is degenerate.

\begin{tcolorbox}
在本节中，我们将处理那些需要更仔细分析才能找到其所有均衡的博弈。我们将完全分析一个例子，然后给出适用于像这个例子中的退化博弈的定义。然后我们讨论更大的博弈，并给出当一个退化博弈何时适用的一般定义 6.7。
\end{tcolorbox}

Figure 6.11 is the strategic form, with best response payoffs, of the Threat game in Figure 4.3. It is similar to the Battle of the Sexes game in that it has two pure equilibria, here(T, l)and(B, r). As in the Battle of the Sexes game, we should expect an additional mixed-strategy equilibrium, which indeed exists. Assume this equilibrium is given by the mixed strategy(1 - p, p)for player \(\mathrm{I}\) and(1 - q, q) for player II. The difference trick gives \(q = \frac{1}{2}\) so that player \(\mathrm{I}\) is indifferent between \(T\) and \(B\) and can mix, and for player I the difference trick gives \(p = 0\) so that player II is indifferent between \(l\) and \(r\) .

\begin{tcolorbox}
图6.11是图4.3中威胁博弈的策略形式，其中包含了最优反应的效用。它与性别战博弈类似，都有两个纯策略均衡，这里是(T, l)和(B, r)。和性别战博弈一样，我们预计会存在一个额外的混合策略均衡，实际上也确实存在。假设这个均衡由参与者 \(\mathrm{I}\) 的混合策略(1 - p, p)和参与者II的混合策略(1 - q, q)给出。通过差分技巧得到 \(q = \frac{1}{2}\) ，使得参与者 \(\mathrm{I}\) 在 \(T\) 和 \(B\) 之间无差异，从而可以进行混合；对于参与者I，差分技巧得到 \(p = 0\) ，使得参与者II在 \(l\) 和 \(r\) 之间无差异。
\end{tcolorbox}

In more detail, if player \(I\) plays \(T\) and \(B\) with probabilities \(1 - p\) and \(p\) , then the expected payoff to player II is \(3\left( {1 - p}\right)\) for her strategy \(l\) , and \(3\left( {1 - p}\right)  + {2p}\) for \(r\) , so that she is indifferent between \(l\) and \(r\) only when \(p = 0\) . Indeed, because \(r\) weakly dominates \(l\) , and player II’s payoff for \(r\) is larger when player I plays \(B\) , any positive probability for \(B\) would make \(r\) the unique best response of player II (against which \(B\) is the unique best response of player \(\mathrm{I}\) ). Hence, the only equilibrium where \(B\) is played with positive probability is the pure-strategy equilibrium(B, r).

\begin{tcolorbox}
更详细地说，如果参与者 \(I\) 分别以概率 \(1 - p\) 和 \(p\) 选择 \(T\) 和 \(B\) ，那么参与者II选择策略 \(l\) 的期望效用是 \(3\left( {1 - p}\right)\) ，选择 \(r\) 的期望效用是 \(3\left( {1 - p}\right)  + {2p}\) ，所以只有当 \(p = 0\) 时，她在 \(l\) 和 \(r\) 之间才无差异。实际上，由于 \(r\) 弱占优 \(l\) ，并且当参与者I选择 \(B\) 时，参与者II选择 \(r\) 的效用更大，所以 \(B\) 以任何正概率出现都会使 \(r\) 成为参与者II的唯一最优反应(而对于 \(r\) ， \(B\) 是参与者 \(\mathrm{I}\) 的唯一最优反应)。因此， \(B\) 以正概率出现的唯一均衡是纯策略均衡(B, r)。
\end{tcolorbox}

\begin{center}
\includegraphics[max width=0.2\textwidth]{images/bo_d2e5cmbef24c73b03tn0_172_259_1580_354_359_0.jpg}
\end{center}
\hspace*{3em} 

Figure 6.11 Strategic form of the Threat game, which is a degenerate game.

\begin{tcolorbox}
图6.11 威胁博弈的策略形式，这是一个退化博弈。
\end{tcolorbox}

Consequently, a mixed-strategy equilibrium of the game is \(\left( {\left( {1,0}\right) ,\left( {\frac{1}{2},\frac{1}{2}}\right) }\right)\) . However, in this equilibrium only player II uses a properly mixed strategy and player I plays the pure strategy \(T\) . The best-response condition requires that player II is indifferent between her pure strategies \(l\) and \(r\) because both are played with positive probability, and this implies that \(T\) is played with probability one, as just described. For player I, however, the best-response condition requires that \(T\) has maximal expected payoff, which does not have to be the same as the expected payoff for \(B\) because \(B\) does not have positive probability! That is to say, player \({\mathrm{I}}^{\prime }\mathrm{s}\) indifference between \(T\) and \(B\) , which we have used to determine player II’s mixed strategy(1 - q, q)as \(q = \frac{1}{2}\) , is too strong a requirement. All we need is that \(T\) is a best response to this mixed strategy, so the expected payoff 1 to player I when he plays \(T\) has to be at least as large as his expected payoff \({2q}\) when he plays \(B\) , that is, \(1 \geq  {2q}\) . In other words, the equation \(1 = {2q}\) has to be replaced by an inequality.

\begin{tcolorbox}
因此，该博弈的一个混合策略均衡是 \(\left( {\left( {1,0}\right) ,\left( {\frac{1}{2},\frac{1}{2}}\right) }\right)\) 。然而，在这个均衡中，只有参与者II使用了恰当的混合策略，而参与者I选择了纯策略 \(T\) 。最优反应条件要求参与者II在她的纯策略 \(l\) 和 \(r\) 之间无差异，因为这两个策略都以正概率出现，这意味着正如前面所述， \(T\) 以概率1出现。然而，对于参与者I来说，最优反应条件要求 \(T\) 的期望效用最大，这不一定与 \(B\) 的期望效用相同，因为 \(B\) 出现的概率为零！也就是说，我们用来确定参与者II的混合策略(1 - q, q)为 \(q = \frac{1}{2}\) 的参与者 \({\mathrm{I}}^{\prime }\mathrm{s}\) 在 \(T\) 和 \(B\) 之间无差异这一条件太强了。我们只需要 \(T\) 是对这个混合策略的最优反应，所以参与者I选择 \(T\) 时期望效用为1，至少要和他选择 \(B\) 时期望效用 \({2q}\) 一样大，即 \(1 \geq  {2q}\) 。换句话说，等式 \(1 = {2q}\) 必须用一个不等式来代替。
\end{tcolorbox}

The inequality \(1 \geq  {2q}\) is equivalent to \(0 \leq  q \leq  \frac{1}{2}\) . With these possible values for \(q\) , we obtain an infinite set of equilibria \(\left( {\left( {1,0}\right) ,\left( {1 - q,q}\right) }\right)\) . The two extreme cases for \(q\) , namely \(q = 0\) and \(q = \frac{1}{2}\) , give the pure-strategy equilibrium(T, l)and the mixed equilibrium \(\left( {\left( {1,0}\right) ,\left( {\frac{1}{2},\frac{1}{2}}\right) }\right)\) that we found earlier.

\begin{tcolorbox}
不等式 \(1 \geq  {2q}\) 等价于 \(0 \leq  q \leq  \frac{1}{2}\) 。对于 \(q\) 的这些可能取值，我们得到了一个无限的均衡集合 \(\left( {\left( {1,0}\right) ,\left( {1 - q,q}\right) }\right)\) 。 \(q\) 的两个极端情况，即 \(q = 0\) 和 \(q = \frac{1}{2}\) ，分别给出了纯策略均衡(T, l)和我们之前找到的混合均衡 \(\left( {\left( {1,0}\right) ,\left( {\frac{1}{2},\frac{1}{2}}\right) }\right)\) 。
\end{tcolorbox}

Note: An equilibrium where at least one player uses a mixed strategy that is not a pure strategy is always called a mixed equilibrium. Also, when we want to find "all mixed equilibria" of a game, we usually include the pure-strategy equilibria among them (which are at any rate easy to find).

\begin{tcolorbox}
注意:至少有一个参与者使用非纯策略的混合策略的均衡总是被称为混合均衡。此外，当我们想要找到一个博弈的“所有混合均衡”时，我们通常会把纯策略均衡也包括在内(无论如何，纯策略均衡很容易找到)。
\end{tcolorbox}

The described set of equilibria (where player II uses a mixed-strategy probability \(q\) from some interval) has an intuition in terms of the possible "threat" in this game that player II plays \(l\) (with an undesirable outcome after player I has chosen \(B\) , in the game tree in Figure 4.3(a)). First, \(r\) weakly dominates \(l\) , so if there is any positive probability that player I ignores the threat and plays \(B\) , then the threat of playing \(l\) at all is not sustainable as a best response, so any equilibrium where player II does play \(l\) requires that player I’s probability \(p\) for playing \(B\) is zero. Second, if player I plays \(T\) , then player II is indifferent between \(l\) and \(r\) and she can in principle play anything, but in order to maintain the threat, \(T\) must be a best response, which requires that the probability \(1 - q\) for \(l\) is sufficiently high. Player I is indifferent between \(T\) and \(B\) (and therefore \(T\) still optimal) when \(1 - q = \frac{1}{2}\) . The threat works whenever \(1 - q \geq  \frac{1}{2}\) . In other words,"threats" can work even when the threatened action with the undesirable outcome is uncertain, as long as its probability is sufficiently high.

\begin{tcolorbox}
所描述的均衡集(其中玩家 II 从某个区间中选择混合策略概率 \(q\) )可以从这个博弈中玩家 II 选择 \(l\) 所构成的可能“威胁”来直观理解(在图 4.3(a) 的博弈树中，玩家 I 选择 \(B\) 后会出现不理想的结果)。首先， \(r\) 弱占优 \(l\) ，因此如果玩家 I 有任何正概率忽视该威胁并选择 \(B\) ，那么玩家 II 选择 \(l\) 的威胁就无法作为最优反应持续下去，所以任何玩家 II 选择 \(l\) 的均衡都要求玩家 I 选择 \(B\) 的概率 \(p\) 为零。其次，如果玩家 I 选择 \(T\) ，那么玩家 II 在 \(l\) 和 \(r\) 之间无差异，原则上她可以选择任何策略，但为了维持威胁， \(T\) 必须是最优反应，这要求选择 \(l\) 的概率 \(1 - q\) 足够高。当 \(1 - q = \frac{1}{2}\) 时，玩家 I 在 \(T\) 和 \(B\) 之间无差异(因此 \(T\) 仍然是最优的)。只要 \(1 - q \geq  \frac{1}{2}\) ，威胁就会起作用。换句话说，只要产生不理想结果的受威胁行动的概率足够高，即使该行动不确定，“威胁”也能奏效。
\end{tcolorbox}

The complications in this game arise because the game in Figure 6.11 is "degenerate" according to the following definition.

\begin{tcolorbox}
根据以下定义，图 6.11 中的博弈是“退化的”，这就导致了该博弈的复杂性。
\end{tcolorbox}

Definition 6.6. A \(2 \times  2\) game is called degenerate if some player has a pure strategy with two pure best responses of the other player.

\begin{tcolorbox}
定义 6.6:如果某个玩家的一个纯策略对应另一个玩家的两个纯最优反应，则称该 \(2 \times  2\) 博弈为退化博弈。
\end{tcolorbox}

A degenerate game is complicated to analyze because one player can use a pure strategy (like \(T\) in the threat game) so that the other player can mix between her two best responses, but the mixed-strategy probabilities do not have to fulfill the equation that the other player is indifferent between his strategies. Instead of that equation, it suffices to fulfill the inequality that the first player's pure strategy is a best response. Note that the mentioned equation, as well as the inequality, is a consequence of the best-response condition, Proposition 6.1.

\begin{tcolorbox}
退化博弈分析起来很复杂，因为一个玩家可以使用纯策略(如威胁博弈中的 \(T\) )，使得另一个玩家可以在她的两个最优反应之间进行混合，但混合策略的概率不必满足另一个玩家在其策略之间无差异的等式。取而代之的是，只需满足第一个玩家的纯策略是最优反应的不等式即可。请注意，上述等式以及不等式都是最优反应条件(命题 6.1)的结果。
\end{tcolorbox}

\(\Rightarrow\) Exercise 6.7 is about a degenerate game.

\begin{tcolorbox}
\(\Rightarrow\) 练习 6.7 是关于一个退化博弈的。
\end{tcolorbox}

\begin{center}
\includegraphics[max width=0.9\textwidth]{images/bo_d2e5cmbef24c73b03tn0_174_191_493_1251_432_0.jpg}
\end{center}
\hspace*{3em} 

Figure 6.12 Degenerate \(2 \times  3\) game and upper envelope of expected payoffs to player II.

\begin{tcolorbox}
图 6.12 退化 \(2 \times  3\) 博弈以及玩家 II 的期望效用上包络。
\end{tcolorbox}

Degeneracy occurs for larger games as well, which requires a more general definition. Consider the \(2 \times  3\) game in Figure 6.12. Each pure strategy has a unique best response, so the condition in Definition 6.6 does not hold. This game does not have a pure-strategy equilibrium, so both players have to use mixed strategies. We apply the upper-envelope method, plotting the payoff to player II as a function of the probability that player I plays \(B\) , shown on the right in Figure 6.12. Here, the lines of expected payoffs cross in a single point, and all three pure strategies of player II are best responses when player I uses the mixed strategy \(\left( {\frac{1}{2},\frac{1}{2}}\right)\) .

\begin{tcolorbox}
退化现象在更大规模的博弈中也会出现，这需要更一般的定义。考虑图 6.12 中的 \(2 \times  3\) 博弈。每个纯策略都有唯一的最优反应，因此定义 6.6 中的条件不成立。这个博弈没有纯策略均衡，所以两个玩家都必须使用混合策略。我们应用上包络法，将玩家 II 的收益绘制为玩家 I 选择 \(B\) 的概率的函数，如图 6.12 右侧所示。在这里，期望效用线在一个点相交，当玩家 I 使用混合策略 \(\left( {\frac{1}{2},\frac{1}{2}}\right)\) 时，玩家 II 的所有三个纯策略都是最优反应。
\end{tcolorbox}

How should we analyze this game? Clearly, we only have an equilibrium if both players mix, which requires the mixed strategy \(\left( {\frac{1}{2},\frac{1}{2}}\right)\) for player \(\mathrm{I}\) . Then player I has to be made indifferent, and in order to achieve this, player II can mix between any of her three best responses \(l,c\) , and \(r\) . We now have the same problem as earlier in the Threat game, namely too much freedom: Three probabilities for player II, call them \({y}_{l},{y}_{c}\) , and \({y}_{r}\) , have to be chosen so that \({y}_{l} + {y}_{c} + {y}_{r} = 1\) , and so that player \(\mathrm{I}\) is indifferent between his pure strategies \(T\) and \(B\) , which gives the equation \(3{y}_{c} + 2{y}_{r} = {y}_{l} + 2{y}_{c}\) . These are three unknowns subject to two equations, with an underdetermined solution.

\begin{tcolorbox}
我们应该如何分析这个博弈呢？显然，只有当两个玩家都采用混合策略时才会达到均衡，这要求玩家 \(\mathrm{I}\) 使用混合策略 \(\left( {\frac{1}{2},\frac{1}{2}}\right)\) 。然后必须使玩家 I 无差异，为了实现这一点，玩家 II 可以在她的三个最优反应 \(l,c\) 和 \(r\) 之间进行混合。现在我们遇到了与之前威胁博弈中相同的问题，即自由度太大:玩家 II 的三个概率，设为 \({y}_{l},{y}_{c}\) 和 \({y}_{r}\) ，必须选择使得 \({y}_{l} + {y}_{c} + {y}_{r} = 1\) ，并且使得玩家 \(\mathrm{I}\) 在他的纯策略 \(T\) 和 \(B\) 之间无差异，这给出了等式 \(3{y}_{c} + 2{y}_{r} = {y}_{l} + 2{y}_{c}\) 。这是三个未知数受两个方程约束，解是不确定的。
\end{tcolorbox}

We give a method for finding out all possible probabilities, which can be generalized to the situation that three probabilities are subject to only two linear equations, like here. First, sort the coefficients of the probabilities in the equation \(3{y}_{c} + 2{y}_{r} = {y}_{l} + 2{y}_{c}\) , which describes the indifference of the expected payoffs to the other player, so that each probability appears only once and so that its coefficient is positive, which here gives the equation \({y}_{c} + 2{y}_{r} = {y}_{l}\) . Of the three probabilities, normally two appear on one side of the equation and one on the other side. The "extreme solutions" of this equation are obtained by setting either of the probabilities on the side of the equation with two probabilities to zero, here either \({y}_{c} = 0\) or \({y}_{r} = 0\) . In the former case, this gives the solution \(\left( {\frac{2}{3},0,\frac{1}{3}}\right)\) for \(\left( {{y}_{l},{y}_{c},{y}_{r}}\right)\) , in the latter case the solution \(\left( {\frac{1}{2},\frac{1}{2},0}\right)\) . This implies that in any equilibrium \(0 \leq  {y}_{c} \leq  \frac{1}{2}\) (and similarly \(\left. {0 \leq  {y}_{r} \leq  \frac{1}{3}}\right)\) . By chosing \({y}_{c} \in  \left\lbrack  {0,\frac{1}{2}}\right\rbrack\) as the free parameter, the remaining probabilities are given by \({y}_{l} = {y}_{c} + 2{y}_{r} = {y}_{c} + 2\left( {1 - {y}_{l} - {y}_{c}}\right)  = 2 - 2{y}_{l} - {y}_{c}\) or \({y}_{l} = \frac{2}{3} - \frac{1}{3}{y}_{c}\) , and \({y}_{r} = 1 - {y}_{l} - {y}_{c} = 1 - \left( {\frac{2}{3} - \frac{1}{3}{y}_{c}}\right)  - {y}_{c} = \frac{1}{3} - \frac{2}{3}{y}_{c}\) .

\begin{tcolorbox}
我们给出一种找出所有可能概率的方法，该方法可以推广到像这里这样三个概率仅受两个线性方程约束的情形。首先，对描述另一个参与者期望效用无差异的方程 \(3{y}_{c} + 2{y}_{r} = {y}_{l} + 2{y}_{c}\) 中概率的系数进行整理，使得每个概率仅出现一次且其系数为正，在此得到方程 \({y}_{c} + 2{y}_{r} = {y}_{l}\) 。在这三个概率中，通常有两个出现在方程的一侧，另一个出现在另一侧。通过将方程中包含两个概率那一侧的任意一个概率设为零，即可得到该方程的“极端解”，在此即 \({y}_{c} = 0\) 或 \({y}_{r} = 0\) 。在前一种情况下，对于 \(\left( {{y}_{l},{y}_{c},{y}_{r}}\right)\) 可得到解 \(\left( {\frac{2}{3},0,\frac{1}{3}}\right)\) ；在后一种情况下，可得到解 \(\left( {\frac{1}{2},\frac{1}{2},0}\right)\) 。这意味着在任何均衡中 \(0 \leq  {y}_{c} \leq  \frac{1}{2}\) (同理 \(\left. {0 \leq  {y}_{r} \leq  \frac{1}{3}}\right)\) )。通过选择 \({y}_{c} \in  \left\lbrack  {0,\frac{1}{2}}\right\rbrack\) 作为自由参数，其余概率由 \({y}_{l} = {y}_{c} + 2{y}_{r} = {y}_{c} + 2\left( {1 - {y}_{l} - {y}_{c}}\right)  = 2 - 2{y}_{l} - {y}_{c}\) 或 \({y}_{l} = \frac{2}{3} - \frac{1}{3}{y}_{c}\) 以及 \({y}_{r} = 1 - {y}_{l} - {y}_{c} = 1 - \left( {\frac{2}{3} - \frac{1}{3}{y}_{c}}\right)  - {y}_{c} = \frac{1}{3} - \frac{2}{3}{y}_{c}\) 给出。
\end{tcolorbox}

Another way to find the extreme solutions to the underdetermined mixed-strategy probabilities of player II is to simply ignore one of the three best responses and apply the difference trick: Assume that player II uses only his best responses \(l\) and \(c\) . Then the difference trick, to make player I indifferent, gives probabilities \(\left( {\frac{1}{2},\frac{1}{2},0}\right)\) for \(l,c,r\) . If player II uses only his best responses \(l\) and \(r\) , the difference trick gives the mixed strategy \(\left( {\frac{2}{3},0,\frac{1}{3}}\right)\) . These are exactly the two mixed strategies just described. If player II uses only his best responses \(c\) and \(r\) , then the difference trick does not work because player I’s best response to both \(c\) and \(r\) is always \(T\) .

\begin{tcolorbox}
找出参与者 II 未完全确定的混合策略概率的极端解的另一种方法是，简单地忽略三个最优反应中的一个，然后应用差分技巧:假设参与者 II 仅使用其最优反应 \(l\) 和 \(c\) 。那么，为了使参与者 I 达到无差异，差分技巧给出 \(l,c,r\) 的概率为 \(\left( {\frac{1}{2},\frac{1}{2},0}\right)\) 。如果参与者 II 仅使用其最优反应 \(l\) 和 \(r\) ，差分技巧给出混合策略 \(\left( {\frac{2}{3},0,\frac{1}{3}}\right)\) 。这正是刚刚描述的两种混合策略。如果参与者 II 仅使用其最优反应 \(c\) 和 \(r\) ，那么差分技巧就不起作用了，因为参与者 I 对 \(c\) 和 \(r\) 的最优反应始终是 \(T\) 。
\end{tcolorbox}

The following is a definition of degeneracy for general \(m \times  n\) games.

\begin{tcolorbox}
以下是一般 \(m \times  n\) 博弈的退化定义。
\end{tcolorbox}

Definition 6.7 (Degenerate game). A two-player game is called degenerate if some player has a mixed strategy that assigns positive probability to exactly \(k\) pure strategies so that the other player has more than \(k\) pure best responses to that mixed strategy.

\begin{tcolorbox}
定义 6.7(退化博弈)。如果某个参与者有一个混合策略，该策略恰好对 \(k\) 个纯策略赋予正概率，使得另一个参与者对该混合策略有超过 \(k\) 个纯最优反应，那么这个两人博弈就被称为退化博弈。
\end{tcolorbox}

In a degenerate game, a mixed strategy with "too many" best responses creates the difficulty that we have described with the above examples. Consider, in contrast, the "normal" situation that the game is nondegenerate.

\begin{tcolorbox}
在退化博弈中，具有“过多”最优反应的混合策略会产生我们在上述例子中所描述的困难。相比之下，考虑博弈是非退化的“正常”情形。
\end{tcolorbox}

Proposition 6.8. In any equilibrium of a two-player game that is not degenerate, both players use mixed strategies that mix the same number of pure strategies.

\begin{tcolorbox}
命题 6.8。在任何非退化的两人博弈的均衡中，两个参与者使用的混合策略所混合的纯策略数量相同。
\end{tcolorbox}

\(\Rightarrow\) You are asked to give the easy and instructive proof of Proposition 6.8 in Exercise 6.8.

\begin{tcolorbox}
\(\Rightarrow\) 你需要在练习 6.8 中给出命题 6.8 的简单且有启发性的证明。
\end{tcolorbox}

In other words, in a nondegenerate game we have pure-strategy equilibria (both players use exactly one pure strategy), or mixed-strategy equilibria where both players mix between exactly two strategies, or equilibria where both player mix between exactly three strategies, and so on. In general, in a nondegenerate game both players mix between exactly \(k\) pure strategies. Each of these strategies has a probability, and these probabilities are subject to \(k\) equations in order to make the other player indifferent between his or her pure strategies: One of these equations states that the probabilities sum to one, and the other \(k - 1\) equations state the equality between the first and second, second and third, etc., up to the (k - 1)st and \(k\) th strategy of the other player. These equations have unique solutions. (It can be shown that non-unique solutions can only occur in a degenerate game.) These solutions have to be checked for the equilibrium property: The resulting probabilities have to be nonnegative, and the unused pure strategies of the other player must not have a higher payoff. Finding mixed-strategy probabilities by equating the expected payoffs to the other player no longer gives unique solutions in degenerate games, as we have demonstrated.

\begin{tcolorbox}
换句话说，在非退化博弈(nondegenerate game)中，我们有纯策略均衡(两个参与者都恰好使用一种纯策略)，或者混合策略均衡(两个参与者都恰好混合使用两种策略)，或者均衡中两个参与者都恰好混合使用三种策略，依此类推。一般来说，在非退化博弈中，两个参与者都恰好混合使用 \(k\) 种纯策略。这些策略中的每一种都有一个概率，并且这些概率要满足 \(k\) 个方程，以使另一个参与者在其纯策略之间无差异:其中一个方程表示概率之和为 1，另外 \(k - 1\) 个方程表示另一个参与者的第一种和第二种、第二种和第三种，依此类推，直到第 (k - 1) 种和第 \(k\) 种策略之间的期望效用相等。这些方程有唯一解。(可以证明，非唯一解只可能出现在退化博弈(degenerate game)中。)需要检查这些解是否具有均衡性质:得到的概率必须是非负的，并且另一个参与者未使用的纯策略的期望效用不能更高。正如我们所证明的，在退化博弈中，通过使另一个参与者的期望效用相等来找到混合策略的概率不再能得到唯一解。
\end{tcolorbox}

Should we or should we not care about degenerate games? In a certain sense, degenerate games can be ignored when considering games that are given in strategic form, and where each cell of the payoff table arises from a different circumstance of the interactive situation that is modeled by the game. In that case, it is "unlikely" that two payoffs are identical, which is necessary to get two pure best responses to a pure strategy, which is the only case where a \(2 \times  2\) game can be degenerate. A small change of the payoff will result in a unique preference of the other player. Similarly, it is unlikely that more than two lines that define the upper envelope of expected payoffs cross in one point, or that more than three planes of expected payoffs against a mixed strategy that mixes three pure strategies (like in the two planes in Figure 9.2) cross in one point, and so on. In other words, "generic" (or "almost all") games in strategic form are not degenerate. On the other hand, degeneracy is very likely when looking at a game tree, because a payoff of the game tree will occur repeatedly in the strategic form, as demonstrated by the Threat game.

\begin{tcolorbox}
我们是否应该关注退化博弈(degenerate game)呢？从某种意义上说，当考虑以策略形式(strategic form)给出的博弈，并且收益表的每个单元格都源于该博弈所建模的交互情境的不同情况时，可以忽略退化博弈。在这种情况下，两个收益相同的可能性“很小”，而这是对一种纯策略有两种纯最优反应(best response)的必要条件，而这也是退化博弈可能出现的唯一情况。收益的微小变化会导致另一个参与者有唯一的偏好。同样，定义期望效用上包络(upper envelope)的两条以上的线在一点相交的可能性很小，或者针对混合三种纯策略的混合策略(如在图 9.2 中的两个平面)的三个以上的期望效用平面在一点相交的可能性也很小，依此类推。换句话说，策略形式下的“一般”(或“几乎所有”)博弈是非退化的。另一方面，当观察博弈树(game tree)时，退化的可能性非常大，因为博弈树的一个收益会在策略形式中反复出现，如威胁博弈所示。
\end{tcolorbox}

We conclude with a useful proposition that can be used when looking for all equilibria of a degenerate game. In the Threat game in Figure 6.11, there are two equilibria \(\left( {\left( {1,0}\right) ,\left( {1,0}\right) }\right)\) and \(\left( {\left( {1,0}\right) ,\left( {\frac{1}{2},\frac{1}{2}}\right) }\right)\) that are obtained, essentially, by "ignoring" the degeneracy. These two equilibria have the same strategy(1,0)of player I. Similarly, the game in Figure 6.12 has two mixed equilibria \(\left( {\left( {\frac{1}{2},\frac{1}{2}}\right) ,\left( {\frac{1}{2},\frac{1}{2},0}\right) }\right)\) and \(\left( {\left( {\frac{1}{2},\frac{1}{2}}\right) ,\left( {\frac{2}{3},0,\frac{1}{3}}\right) }\right)\) , which again share the same strategy of player \(\mathrm{I}\) .

\begin{tcolorbox}
我们以一个有用的命题作为结论，该命题可用于寻找退化博弈(degenerate game)的所有均衡。在图 6.11 的威胁博弈中，有两个均衡 \(\left( {\left( {1,0}\right) ,\left( {1,0}\right) }\right)\) 和 \(\left( {\left( {1,0}\right) ,\left( {\frac{1}{2},\frac{1}{2}}\right) }\right)\) ，本质上是通过“忽略”退化得到的。这两个均衡中参与者 I 的策略相同，均为 (1,0)。类似地，图 6.12 中的博弈有两个混合均衡 \(\left( {\left( {\frac{1}{2},\frac{1}{2}}\right) ,\left( {\frac{1}{2},\frac{1}{2},0}\right) }\right)\) 和 \(\left( {\left( {\frac{1}{2},\frac{1}{2}}\right) ,\left( {\frac{2}{3},0,\frac{1}{3}}\right) }\right)\) ，它们同样具有参与者 \(\mathrm{I}\) 的相同策略。
\end{tcolorbox}

The following proposition states that in a two-player game, two mixed equilibria that have the same strategy of one player can be combined with convex combinations to obtain further equilibria, as it is the case in the above examples. In the proposition player II’s strategy \(y\) is fixed, but it holds similarly when player I’s strategy is fixed. Fixing strategies is important here, because the set of equilibria is in general not convex.

\begin{tcolorbox}
以下命题指出，在两人博弈中，具有一个参与者相同策略的两个混合均衡可以通过凸组合(convex combination)相结合，以得到更多的均衡，就像上述例子中的情况一样。在该命题中，参与者 II 的策略 \(y\) 是固定的，但当参与者 I 的策略固定时，情况类似。在这里固定策略很重要，因为均衡集一般不是凸集。
\end{tcolorbox}

Proposition 6.9. Consider a bimatrix game(A, B)with mixed-strategy sets \(X\) and \(Y\) as in (6.6). Let(x, y)and \(\left( {\widehat{x},y}\right)\) be equilibria of the game, where \(x,\widehat{x} \in  X\) and \(y \in  Y\) . Then \(\left( {x\left( {1 - p}\right)  + \widehat{x}p}\right) ,y)\) is also an equilibrium of the game, for any \(p \in  \left\lbrack  {0,1}\right\rbrack\) .

\begin{tcolorbox}
命题6.9。考虑一个双矩阵博弈(A, B)，其混合策略集为 \(X\) 和 \(Y\) ，如(6.6)所示。设(x, y)和 \(\left( {\widehat{x},y}\right)\) 是该博弈的均衡，其中 \(x,\widehat{x} \in  X\) 且 \(y \in  Y\) 。那么对于任意 \(p \in  \left\lbrack  {0,1}\right\rbrack\) ， \(\left( {x\left( {1 - p}\right)  + \widehat{x}p}\right) ,y)\) 也是该博弈的一个均衡。
\end{tcolorbox}

Proof. Let \(\bar{x} = x\left( {1 - p}\right)  + \widehat{x}p\) . Clearly, \(\bar{x} \in  X\) because \({\bar{x}}_{i} \geq  0\) for all pure strategies \(i\) of player I, and

\begin{tcolorbox}
证明。设 \(\bar{x} = x\left( {1 - p}\right)  + \widehat{x}p\) 。显然， \(\bar{x} \in  X\) ，因为对于参与者I的所有纯策略 \(i\) ，都有 \({\bar{x}}_{i} \geq  0\) ，并且
\end{tcolorbox}

\[
{\mathbf{1}}^{\top }\bar{x} = {\mathbf{1}}^{\top }\left( {x\left( {1 - p}\right)  + \widehat{x}p}\right)  = {\mathbf{1}}^{\top }x\left( {1 - p}\right)  + {\mathbf{1}}^{\top }\widehat{x}p = \left( {1 - p}\right)  + p = 1. \tag{6.23}
\]

The pair \(\left( {\bar{x},y}\right)\) is an equilibrium if \(\bar{x}\) is a best response to \(y\) and vice versa. For any mixed strategy \(\widetilde{x}\) of player \(\mathrm{I}\) , we have \({x}^{\top }{Ay} \geq  {\widetilde{x}}^{\top }{Ay}\) and \({\widehat{x}}^{\top }{Ay} \geq  {\widetilde{x}}^{\top }{Ay}\) , and consequently

\begin{tcolorbox}
如果 \(\bar{x}\) 是对 \(y\) 的最优反应，反之亦然，那么 \(\left( {\bar{x},y}\right)\) 这一对就是一个均衡。对于参与者 \(\mathrm{I}\) 的任何混合策略 \(\widetilde{x}\) ，我们有 \({x}^{\top }{Ay} \geq  {\widetilde{x}}^{\top }{Ay}\) 和 \({\widehat{x}}^{\top }{Ay} \geq  {\widetilde{x}}^{\top }{Ay}\) ，因此
\end{tcolorbox}

\[
{\bar{x}}^{\top }{Ay} = \left( {1 - p}\right) {x}^{\top }{Ay} + p{\widehat{x}}^{\top }{Ay} \geq  \left( {1 - p}\right) {\widetilde{x}}^{\top }{Ay} + p{\widetilde{x}}^{\top }{Ay} = {\widetilde{x}}^{\top }{Ay},
\]

which shows that \(\bar{x}\) is a best response to \(y\) . In other words, these inequalities hold because they are preserved under convex combinations. Similarly, \(y\) is a best response to \(x\) and \(\widehat{x}\) , that is, for any \(\widetilde{y} \in  Y\) we have \({x}^{\top }{By} \geq  {x}^{\top }B\widetilde{y}\) and \({\widehat{x}}^{\top }{By} \geq  {\widehat{x}}^{\top }B\widetilde{y}\) . Again, taking convex combinations shows \({\bar{x}}^{\top }{By} \geq  {\bar{x}}^{\top }B\widetilde{y}\) . Hence, \(\left( {\bar{x},y}\right)\) is an equilibrium, as claimed.

\begin{tcolorbox}
这表明 \(\bar{x}\) 是对 \(y\) 的最优反应。换句话说，这些不等式成立是因为它们在凸组合下保持不变。类似地， \(y\) 是对 \(x\) 和 \(\widehat{x}\) 的最优反应，即对于任何 \(\widetilde{y} \in  Y\) ，我们有 \({x}^{\top }{By} \geq  {x}^{\top }B\widetilde{y}\) 和 \({\widehat{x}}^{\top }{By} \geq  {\widehat{x}}^{\top }B\widetilde{y}\) 。再次，取凸组合表明 \({\bar{x}}^{\top }{By} \geq  {\bar{x}}^{\top }B\widetilde{y}\) 。因此，如所声称的， \(\left( {\bar{x},y}\right)\) 是一个均衡。
\end{tcolorbox}

We come back to the example of Figure 6.12 to illustrate Proposition 6.9: This game has two equilibria(x, y)and \(\left( {x,\widehat{y}}\right)\) where \(x = \left( {\frac{1}{2},\frac{1}{2}}\right) ,y = \left( {\frac{1}{2},\frac{1}{2},0}\right)\) , and \(\widehat{y} = \left( {\frac{2}{3},0,\frac{1}{3}}\right)\) , so this is the situation of Proposition 6.9, except that the two equilibria have the same mixed strategy of player I rather than of player II. Consequently, \(\left( {x,\bar{y}}\right)\) is an equilibrium for any convex combination \(\bar{y}\) of \(y\) and \(\widehat{y}\) . The equilibria (x, y)and \(\left( {x,\widehat{y}}\right)\) are "extreme" in the sense that \(y\) and \(\widehat{y}\) have as many zero probabilities as possible, which means that \(y\) and \(\widehat{y}\) are the endpoints of a line segment of equilibrium strategies \(\bar{y} = y\left( {1 - p}\right)  + \widehat{y}p\) (that is, \(\left( {x,\bar{y}}\right)\) is an equilibrium).

\begin{tcolorbox}
我们回到图6.12的例子来说明命题6.9:这个博弈有两个均衡(x, y)和 \(\left( {x,\widehat{y}}\right)\) ，其中 \(x = \left( {\frac{1}{2},\frac{1}{2}}\right) ,y = \left( {\frac{1}{2},\frac{1}{2},0}\right)\) 且 \(\widehat{y} = \left( {\frac{2}{3},0,\frac{1}{3}}\right)\) ，所以这就是命题6.9的情况，只是这两个均衡中参与者I的混合策略相同，而不是参与者II的。因此，对于 \(y\) 和 \(\widehat{y}\) 的任何凸组合 \(\bar{y}\) ， \(\left( {x,\bar{y}}\right)\) 都是一个均衡。均衡(x, y)和 \(\left( {x,\widehat{y}}\right)\) 是“极端的”，因为 \(y\) 和 \(\widehat{y}\) 的概率尽可能多地为零，这意味着 \(y\) 和 \(\widehat{y}\) 是均衡策略线段 \(\bar{y} = y\left( {1 - p}\right)  + \widehat{y}p\) 的端点(即 \(\left( {x,\bar{y}}\right)\) 是一个均衡)。
\end{tcolorbox}

\(\Rightarrow\) Express the set of equilibria for the Threat game in Figure 6.11 with the help of Proposition 6.9.

\begin{tcolorbox}
\(\Rightarrow\) 借助命题6.9，求出图6.11中威胁博弈的均衡集。
\end{tcolorbox}

\section*{6.9 Further Reading}

\begin{tcolorbox}
\section*{6.9 进一步阅读}
\end{tcolorbox}

In Section 6.2, the Inspection game is used to motivate mixed strategies. Inspection games are an applied area of game theory; see, for example, Avenhaus and Canty (1996).

\begin{tcolorbox}
在6.2节中，监管博弈被用于引入混合策略。监管博弈是博弈论的一个应用领域；例如，可参见Avenhaus和Canty(1996)。
\end{tcolorbox}

The seminal paper by Nash (1951) defined "equilibrium points" in finite games and proved their existence. As mentioned, the best-response condition is stated as "trivial" by Nash (1951, p. 287), but it is central for finding mixed equilibria. Nash's "equilibrium point" is now universally called "Nash equilibrium", which is the equilibrium concept considered throughout this book, except for the generalization of correlated equilibrium in Chapter 12. For textbook treatments of mixed-strategy equilibria see Osborne (2004, Chapter 4) or Maschler, Solan, and Zamir (2013, Chapter 5).

\begin{tcolorbox}
纳什(1951)的开创性论文定义了有限博弈中的“均衡点”并证明了它们的存在性。如前所述，纳什(1951，第287页)称最优反应条件为“显而易见的”，但它对于寻找混合均衡至关重要。纳什的“均衡点”现在被普遍称为“纳什均衡”，这是本书通篇所考虑的均衡概念，第12章中相关均衡的推广情况除外。关于混合策略均衡的教科书式处理，可参阅奥斯本(2004，第4章)或马斯勒、索兰和扎米尔(2013，第5章)。
\end{tcolorbox}

The use of the upper envelope seems to be folklore, and is at any rate straightforward. For Definition 6.7 see van Damme (1987, p. 52), and for a detailed comparison of conditions of degeneracy for two-player games (which are all equivalent to this definition) see von Stengel (2021).

\begin{tcolorbox}
上包络的使用似乎是业内常识，而且无论如何都很直接。关于定义6.7，可参阅范·达默(1987，第52页)；关于两人博弈退化条件的详细比较(所有这些条件都与此定义等价)，可参阅冯·施滕格尔(2021)。
\end{tcolorbox}

\section*{6.10 Exercises for Chapter 6}

\begin{tcolorbox}
\section*{6.10 第6章练习}
\end{tcolorbox}

Exercise 6.1. Let \(x \in  {\mathbb{R}}^{m}\) and \(y \in  {\mathbb{R}}^{n}\) and remember that all vectors are column vectors. If \(m = n\) , what is the difference between \({x}^{\top }y\) and \(x{y}^{\top }\) ? What if \(m \neq  n\) ?

\begin{tcolorbox}
练习6.1。设 \(x \in  {\mathbb{R}}^{m}\) 和 \(y \in  {\mathbb{R}}^{n}\) ，并记住所有向量都是列向量。如果 \(m = n\) ， \({x}^{\top }y\) 和 \(x{y}^{\top }\) 之间的差异是什么？如果 \(m \neq  n\) 呢？
\end{tcolorbox}

Exercise 6.2. Find all mixed equilibria (which always includes any pure equilibria) of this \(3 \times  2\) game.

\begin{tcolorbox}
练习6.2。找出这个 \(3 \times  2\) 博弈的所有混合均衡(这总是包括任何纯策略均衡)。
\end{tcolorbox}

\begin{center}
\includegraphics[max width=0.2\textwidth]{images/bo_d2e5cmbef24c73b03tn0_178_1100_442_353_488_0.jpg}
\end{center}
\hspace*{3em} 

Exercise 6.3. Consider the three-player game tree in Exercise 4.5. Recall that some or all of the strategies in a mixed equilibrium may be pure strategies.

\begin{tcolorbox}
练习6.3。考虑练习4.5中的三人博弈树。回想一下，混合均衡中的部分或所有策略可能是纯策略。
\end{tcolorbox}

(a) Is the following statement true or false? It is easiest to argue directly with the game tree.

\begin{tcolorbox}
(a) 以下陈述是对还是错？直接用博弈树来论证是最容易的。
\end{tcolorbox}

For each of the players I, II, or III, the game has a mixed equilibrium in which that player plays a mixed strategy that is not a pure strategy.

\begin{tcolorbox}
对于玩家I、II或III中的每一个，该博弈都有一个混合均衡，其中该玩家采用的混合策略不是纯策略。
\end{tcolorbox}

(b) Is the following statement true or false? Explain.

\begin{tcolorbox}
(b) 以下陈述是对还是错？请解释。
\end{tcolorbox}

In every mixed equilibrium of the game, at least one player I, II, or III plays a pure strategy.

\begin{tcolorbox}
在该博弈的每个混合均衡中，玩家I、II或III中至少有一个采用纯策略。
\end{tcolorbox}

Exercise 6.4. In this \(2 \times  2\) game, \(A,B,C,D\) are the payoffs to player I, which are real numbers, no two of which are equal. Similarly, \(a,b,c,d\) are the payoffs to player II, which are real numbers, also no two of which are equal.

\begin{tcolorbox}
练习6.4。在这个 \(2 \times  2\) 博弈中， \(A,B,C,D\) 是玩家I的收益，它们是实数，且任意两个都不相等。类似地， \(a,b,c,d\) 是玩家II的收益，它们也是实数，且任意两个也不相等。
\end{tcolorbox}

\begin{center}
\includegraphics[max width=0.3\textwidth]{images/bo_d2e5cmbef24c73b03tn0_178_1038_1481_415_356_0.jpg}
\end{center}
\hspace*{3em} 

(a) Under which conditions does this game have a mixed equilibrium which is not a pure-strategy equilibrium?

\begin{tcolorbox}
(a) 在什么条件下，这个博弈有一个不是纯策略均衡的混合均衡？
\end{tcolorbox}

[Hint: Consider the possible patterns of best responses and resulting possible dominance relations, and express this by comparing the payoffs, as, for example, \(A > C\) .]

\begin{tcolorbox}
[提示:考虑最优反应的可能模式以及由此产生的可能的占优关系，并通过比较收益来表达，例如 \(A > C\) 。]
\end{tcolorbox}

(b) Under which conditions in (a) is this the only equilibrium of the game?

\begin{tcolorbox}
(b) 在(a)中的什么条件下，这是该博弈的唯一均衡？
\end{tcolorbox}

(c) Consider one of the situations in (b) and compute the probabilities \(1 - p\) and \(p\) for playing Top and Bottom, respectively, and \(1 - q\) and \(q\) for playing left and right, respectively, that hold in equilibrium.

\begin{tcolorbox}
(c) 考虑(b)中的一种情况，分别计算在均衡状态下选择“上”和“下”的概率 \(1 - p\) 和 \(p\) ，以及选择“左”和“右”的概率 \(1 - q\) 和 \(q\) 。
\end{tcolorbox}

(d) For the solution in (c), give a simple formula for the quotients \(\left( {1 - p}\right) /p\) and \(\left( {1 - q}\right) /q\) in terms of the payoff parameters. Try to write this formula such that the denominator and numerator are both positive.

\begin{tcolorbox}
(d) 对于(c)中的解，用收益参数给出商 \(\left( {1 - p}\right) /p\) 和 \(\left( {1 - q}\right) /q\) 的简单公式。尽量使该公式的分母和分子均为正数。
\end{tcolorbox}

(e) Show that the mixed equilibrium strategies do not change if player I's payoffs \(A\) and \(C\) are replaced by \(A + E\) and \(C + E\) , respectively, for some constant \(E\) , and similarly if player II’s payoffs \(a\) and \(b\) are replaced by \(a + e\) and \(b + e\) , respectively, for some constant \(e\) .

\begin{tcolorbox}
(e) 证明:若玩家I的收益 \(A\) 和 \(C\) 分别被替换为 \(A + E\) 和 \(C + E\) (其中 \(E\) 为某个常数)，且若玩家II的收益 \(a\) 和 \(b\) 分别被替换为 \(a + e\) 和 \(b + e\) (其中 \(e\) 为某个常数)，则混合均衡(mixed equilibrium)策略不会改变。
\end{tcolorbox}

Exercise 6.5. Consider the following \(2 \times  5\) game:

\begin{tcolorbox}
练习6.5. 考虑以下 \(2 \times  5\) 博弈:
\end{tcolorbox}

\begin{center}
\adjustbox{max width=\textwidth}{
\begin{tabular}{|c|c|c|c|c|c|}
\hline
II I & \(a\) & \(b\) & \(C\) & \(d\) & \(e\) \\
\cline{1-6}
\(T\) & 2 0 & 4 2 & 3 1 & 5 0 & 0 3 \\
\cline{1-6}
\(B\) & 7 1 & 4 0 & 5 4 & 0 1 & 8 0 \\
\cline{1-6}
\hline
\end{tabular}
}
\end{center}

\begin{tcolorbox}

\begin{center}
\adjustbox{max width=\textwidth}{
\begin{tabular}{|c|c|c|c|c|c|}
\hline
二 一 & \(a\) & \(b\) & \(C\) & \(d\) & \(e\) \\
\cline{1-6}
\(T\) & 2 0 & 4 2 & 3 1 & 5 0 & 0 3 \\
\cline{1-6}
\(B\) & 7 1 & 4 0 & 5 4 & 0 1 & 8 0 \\
\cline{1-6}
\hline
\end{tabular}
}
\end{center}

\end{tcolorbox}

(a) Draw the expected payoffs to player II for all her strategies \(a,b,c,d,e\) , in terms of the probability \(p\) , say, that player I plays strategy \(B\) . Indicate the best responses of player II, depending on that probability \(p\) .

\begin{tcolorbox}
(a) 绘制玩家 II 所有策略 \(a,b,c,d,e\) 的期望效用(期望效用, expected payoff)，以玩家 I 选择策略 \(B\) 的概率 \(p\) 来表示。根据该概率 \(p\) 指出玩家 II 的最优反应(最优反应, best response)。
\end{tcolorbox}

(b) Using the diagram in (a), find all mixed (including pure) equilibria of the game.

\begin{tcolorbox}
(b) 利用 (a) 中的图表，找出该博弈的所有混合(包括纯策略)均衡。
\end{tcolorbox}

Exercise 6.6. Consider the Quality game in Figure 3.2(b).

\begin{tcolorbox}
练习 6.6。考虑图 3.2(b) 中的产品质量博弈(产品质量博弈, Quality game)。
\end{tcolorbox}

(a) Use the "goalpost" method twice, to draw the upper envelope of best-response payoffs to player I against the mixed strategy of player II, and vice versa.

\begin{tcolorbox}
(a) 两次使用“门柱”方法，绘制玩家 I 针对玩家 II 混合策略的最优反应收益的上包络(上包络, upper envelope)，反之亦然。
\end{tcolorbox}

(b) Explain which of these diagrams shows that a pure strategy dominates another pure strategy.

\begin{tcolorbox}
(b) 解释这些图表中哪一个表明一个纯策略占优于另一个纯策略。
\end{tcolorbox}

(c) Explain how that diagram shows that if a pure strategy \(s\) dominates a strategy \(t\) , then \(t\) is never a best response, even against a mixed strategy of the other player.

\begin{tcolorbox}
(c) 解释该图表如何表明，如果一个纯策略 \(s\) 占优于策略 \(t\) ，那么 \(t\) 永远不是最优反应(最优反应, best response)，即使是针对另一个玩家的混合策略。
\end{tcolorbox}

Exercise 6.7. Find all equilibria of this degenerate game, which is a variant of the Inspection game in Figure 6.2. The difference is that player II, the inspectee, derives no gain from acting illegally (playing \(r\) ) even if she is not inspected (player \(\mathrm{I}\) choosing \(T\) ).

\begin{tcolorbox}
练习 6.7。找出这个退化博弈(退化博弈, degenerate game)的所有均衡，它是图 6.2 中监管博弈(监管博弈, Inspection game)的一个变体。不同之处在于，玩家 II，即被监管者，即使未被监管(玩家 \(\mathrm{I}\) 选择 \(T\) )，采取非法行为(选择 \(r\) )也不会获得任何收益。
\end{tcolorbox}

\begin{center}
\adjustbox{max width=\textwidth}{
\begin{tabular}{|c|c|c|}
\hline
II I & \(l\) & \(r\) \\
\cline{1-3}
\(T\) & 0 0 & 0 -10 \\
\cline{1-3}
\(B\) & 0 -1 & -90 -6 \\
\cline{1-3}
\hline
\end{tabular}
}
\end{center}

\begin{tcolorbox}

\begin{center}
\adjustbox{max width=\textwidth}{
\begin{tabular}{|c|c|c|}
\hline
二 一 & \(l\) & \(r\) \\
\cline{1-3}
\(T\) & 0 0 & 0 -10 \\
\cline{1-3}
\(B\) & 0 -1 & -90 -6 \\
\cline{1-3}
\hline
\end{tabular}
}
\end{center}

\end{tcolorbox}

\section*{Brouwer's Fixed-Point Theorem}

\begin{tcolorbox}
\section*{布劳威尔不动点定理}
\end{tcolorbox}

As we showed in Section 6.5, Nash (1951) proved the existence of a mixed equilibrium in a finite game with the help of Brouwer's fixed-point Theorem 6.4. Fixed-point theorems are powerful tools for proving the existence of many equilibrium concepts in economics. Brouwer's theorem is the first and most important of these.

\begin{tcolorbox}
正如我们在6.5节中所展示的，纳什(1951年)借助布劳威尔不动点定理6.4证明了有限博弈中混合均衡的存在性。不动点定理是证明经济学中许多均衡概念存在性的有力工具。布劳威尔定理是其中第一个也是最重要的一个。
\end{tcolorbox}

This chapter is about proving Brouwer's fixed-point theorem. We present a standard elementary proof (the most complex proof in this book), without the machinery of algebraic topology, along the following route:

\begin{tcolorbox}
本章将证明布劳威尔不动点定理。我们将沿着以下思路给出一个标准的初等证明(本书中最复杂的证明)，不使用代数拓扑的工具:
\end{tcolorbox}

\begin{itemize}
\item The compact convex domain of the considered continuous function is first assumed to be the unit simplex \(\Delta\) . In a simplex, every point is the unique convex combination of the vertices of the simplex, which for \(\Delta\) are the unit vectors.
\end{itemize}

\begin{tcolorbox}
\begin{itemize}
\item 首先假设所考虑的连续函数的紧致凸域是单位单纯形 \(\Delta\) 。在单纯形中，每个点都是该单纯形顶点的唯一凸组合，对于 \(\Delta\) 来说，其顶点就是单位向量。
\end{itemize}
\end{tcolorbox}

\begin{itemize}
\item A finite number of "grid points" of \(\Delta\) is considered, which are the vertices in a so-called simplicial subdivision or triangulation of \(\Delta\) . Although such triangulations have a clear intuition in low dimensions, we carefully prove all required properties about them.
\end{itemize}

\begin{tcolorbox}
\begin{itemize}
\item 考虑 \(\Delta\) 的有限个“网格点”，它们是 \(\Delta\) 的所谓单纯剖分或三角剖分中的顶点。尽管这种三角剖分在低维度下有清晰的直观理解，但我们会仔细证明关于它们的所有必要性质。
\end{itemize}
\end{tcolorbox}

\begin{itemize}
\item Suitable labels identify "completely labeled sub-simplices" of the triangulation as approximate fixed points. Sperner's lemma (Theorem 7.10) states that a completely labeled sub-simplex exists.
\end{itemize}

\begin{tcolorbox}
\begin{itemize}
\item 合适的标签将三角剖分中的“完全标记子单纯形”识别为近似不动点。斯佩纳引理(定理7.10)表明存在完全标记子单纯形。
\end{itemize}
\end{tcolorbox}

\begin{itemize}
\item The KKM lemma (Theorem 7.12) of Knaster, Kuratowski, and Mazurkiewicz (1929) is proved via a limit process from Sperner's lemma and implies the existence of a fixed point of a continuous function on \(\Delta\) .
\end{itemize}

\begin{tcolorbox}
\begin{itemize}
\item 克纳斯特、库拉托夫斯基和马祖尔克维奇(1929年)的KKM引理(定理7.12)通过对斯佩纳引理进行极限过程得到证明，并由此推出 \(\Delta\) 上连续函数不动点的存在性。
\end{itemize}
\end{tcolorbox}

\begin{itemize}
\item The nontrivial part of this limit process is to find finer and finer triangulations, for which we present (at the end of the chapter) the "Freudenthal" triangulation.
\end{itemize}

\begin{tcolorbox}
\begin{itemize}
\item 这个极限过程的非平凡部分是找到越来越精细的三角剖分，为此我们在本章末尾介绍“弗罗伊登塔尔”三角剖分。
\end{itemize}
\end{tcolorbox}

\begin{itemize}
\item A general compact convex set \(C\) is "embedded" into the unit simplex in order to prove the fixed-point theorem for a continuous function on \(C\) .
\end{itemize}

\begin{tcolorbox}
\begin{itemize}
\item 为了证明 \(C\) 上连续函数的不动点定理，将一般的紧致凸集 \(C\) “嵌入”到单位单纯形中。
\end{itemize}
\end{tcolorbox}

The main step is Sperner's lemma. In an ingenious induction proof, Sperner (1928) counted the completely labeled sub-simplices in order to show that their number is odd, and therefore not zero. At its core, Sperner's argument is due to a Background material: Undirected graphs

\begin{tcolorbox}
主要步骤是斯佩纳引理。斯佩纳(1928年)在一个巧妙的归纳证明中，对完全标记子单纯形进行计数，以表明它们的数量是奇数，因此不为零。斯佩纳论证的核心基于一个背景材料:无向图
\end{tcolorbox}

A graph \(G\) is given by a finite set of nodes and a set of edges. Every edge is an unordered pair \(\{ u,v\}\) of nodes \(u\) and \(v\) , which are called the endpoints of the edge, where \(u \neq  v\) . The degree of a node \(u\) is the number of its neighbors in \(G\) , that is, the number of nodes \(v\) so that \(\{ u,v\}\) is an edge. A path in \(G\) of length \(m \geq  0\) is a sequence of nodes \({u}_{0},{u}_{1},\ldots ,{u}_{m}\) so that \(\left\{  {{u}_{k},{u}_{k + 1}}\right\}\) is an edge for \(0 \leq  k < m\) . The end nodes of the path are \({u}_{0}\) and \({u}_{m}\) . If \({u}_{0} = {u}_{m}\) then the path is called a cycle. simple observation from graph theory (see the background material on graphs), namely that the number of odd-degree nodes in a graph is even. This parity argument holds because the sum of the degrees of all nodes is twice the number of edges (and hence an even number), because each edge contributes to the degree of both its endpoints.

\begin{tcolorbox}
图 \(G\) 由有限个节点集和边集给出。每条边是节点 \(u\) 和 \(v\) 的无序对 \(\{ u,v\}\) ， \(u\) 和 \(v\) 称为该边的端点，其中 \(u \neq  v\) 。节点 \(u\) 的度是它在 \(G\) 中的邻居数量，即满足 \(\{ u,v\}\) 是边的节点 \(v\) 的数量。 \(G\) 中长度为 \(m \geq  0\) 的路径是节点序列 \({u}_{0},{u}_{1},\ldots ,{u}_{m}\) ，使得对于 \(0 \leq  k < m\) ， \(\left\{  {{u}_{k},{u}_{k + 1}}\right\}\) 是边。路径的端点是 \({u}_{0}\) 和 \({u}_{m}\) 。如果 \({u}_{0} = {u}_{m}\) ，则该路径称为一个环。这是图论中的一个简单观察(见关于图的背景材料)，即图中奇数度节点的数量是偶数。这个奇偶性论证成立是因为所有节点度的总和是边数量的两倍(因此是偶数)，因为每条边对其两个端点的度都有贡献。
\end{tcolorbox}

In addition to the parity-counting by Sperner (1928), we give a clean constructive proof of Sperner's lemma. This alternative step in the proof combines the graphs across different dimensions in such a way that there is a single path from a vertex of the simplex to a completely labeled sub-simplex (see Figure 7.10). The dimension of the sub-simplices that define the path may go up and down multiple times along the path.

\begin{tcolorbox}
除了斯佩纳(1928年)的奇偶性计数方法外，我们还给出了斯佩纳引理的一个简洁构造性证明。证明中的这个替代步骤以一种特殊方式将不同维度的图组合起来，使得从单纯形的一个顶点到一个完全标记子单纯形存在一条单一路径(见图7.10)。定义该路径的子单纯形的维度可能会沿着路径多次上升和下降。
\end{tcolorbox}

Section 7.2 introduces the concept of a label of a point in the unit simplex, and illustrates its use for the fixed-point theorem when the simplex is an interval. Simplices and triangulations are introduced in Section 7.3. Section 7.4 proves Sperner's lemma. The KKM lemma is presented in Section 7.5.

\begin{tcolorbox}
7.2节介绍单位单纯形中一个点的标签概念，并说明当单纯形是一个区间时它在不动点定理中的应用。7.3节介绍单纯形和三角剖分。7.4节证明斯佩纳引理。7.5节介绍KKM引理。
\end{tcolorbox}

In Section 7.6, we prove Brouwer's fixed-point theorem on a general compact convex set \(C\) by first "embedding" \(C\) with an injection into the unit simplex \(\Delta\) so that we can replace \(C\) by its image \(D\) , which is now a subset of \(\Delta\) . Secondly, we extend the given continuous function \(f\) on \(D\) to a function on \(\Delta\) whose fixed points are exactly the fixed points of \(f\) . For points \(x\) in \(\Delta\) but not in \(D\) , the new function simply chooses the point \(p\left( x\right)\) in \(D\) that is closest to \(x\) . We carefully prove that this "projection" function \(p\) is continuous, in fact "contractive" (see Theorem 7.18), and emphasize the underlying geometric intuition. Even though these considerations are simple, they produce, in passing, important results from convex analysis such as the theorem of the "separating hyperplane" (Corollary 7.16 and Figure 7.13).

\begin{tcolorbox}
在7.6节中，我们通过首先用一个单射将一般紧凸集 \(C\) “嵌入”到单位单纯形 \(\Delta\) 中，来证明布劳威尔不动点定理，这样我们就可以用 \(C\) 的像 \(D\) 来代替 \(C\) ，而 \(D\) 现在是 \(\Delta\) 的一个子集。其次，我们将给定的在 \(D\) 上的连续函数 \(f\) 扩展为在 \(\Delta\) 上的一个函数，其不动点恰好就是 \(f\) 的不动点。对于 \(\Delta\) 中但不在 \(D\) 中的点 \(x\) ，新函数简单地选择 \(D\) 中离 \(x\) 最近的点 \(p\left( x\right)\) 。我们仔细证明了这个“投影”函数 \(p\) 是连续的，实际上是“压缩的”(见定理7.18)，并强调了其背后的几何直觉。尽管这些考虑很简单，但它们顺带产生了凸分析中的重要结果，比如“分离超平面”定理(推论7.16和图7.13)。
\end{tcolorbox}

Section 7.7 treats in detail the nontrivial problem of finding finer and finer triangulations of a simplex in arbitrary dimension, which has been postponed in the proof. We present the elegant triangulation due to Freudenthal (1942), which has two components: the subdivision of the unit cube into small "cubelets", which works well in any dimension \(d\) , and the subdivision of the cube and each cubelet into \(d\) ! simplices obtained from the possible "monotone paths" on the cube.

\begin{tcolorbox}
7.7节详细处理了在任意维度下寻找单纯形越来越精细的三角剖分这一非平凡问题，该问题在证明中被推迟了。我们介绍了弗赖登塔尔(1942年)提出的优雅的三角剖分方法，它有两个部分:将单位立方体细分为小的“子立方体”，这在任何维度 \(d\) 下都适用；以及将立方体和每个子立方体细分为 \(d\) ! 个单纯形，这些单纯形是从立方体上可能的“单调路径”得到的。
\end{tcolorbox}

\section*{7.1 Prerequisites and Learning Outcomes}

\begin{tcolorbox}
\section*{7.1 先决条件和学习成果}
\end{tcolorbox}

This chapter continues from Section 6.5, where we demonstrated Brouwer's fixed-point theorem on an interval and the necessity of its assumptions. The main geometric tool is convexity, with line segments defined in Figure 5.4 and convex combinations defined in Section 6.5. We assume basics of linear algebra when explaining affine combinations, for example in Lemma 7.2.

\begin{tcolorbox}
本章从6.5节继续展开，在6.5节中我们证明了布劳威尔不动点定理在区间上的情况以及其假设的必要性。主要的几何工具是凸性，其中线段在图5.4中定义，凸组合在6.5节中定义。在解释仿射组合时，例如在引理7.2中，我们会假设一些线性代数的基础知识。
\end{tcolorbox}

We also assume familiarity with continuity of functions in \({\mathbb{R}}^{d}\) , with the notions of closed and compact sets, and that a continuous real-valued function on a compact domain assumes its minimum and maximum. We will recall as background material the Euclidean norm and distance, and the Bolzano-Weierstrass theorem on the existence of converging subsequences.

\begin{tcolorbox}
我们还假设熟悉 \({\mathbb{R}}^{d}\) 中函数的连续性、闭集和紧集的概念，以及在紧域上的连续实值函数能取到其最小值和最大值。我们将回顾欧几里得范数与距离以及关于收敛子列存在性的波尔查诺 - 魏尔施特拉斯定理作为背景材料。
\end{tcolorbox}

After studying this chapter, you should be able to:

\begin{tcolorbox}
学习完本章后，你应该能够:
\end{tcolorbox}

\begin{itemize}
\item understand affine and convex combinations, simplices, and triangulations;
\end{itemize}

\begin{tcolorbox}
\begin{itemize}
\item 理解仿射组合和凸组合、单纯形和三角剖分；
\end{itemize}
\end{tcolorbox}

\begin{itemize}
\item define the dimension of a convex set, and the faces of a simplex;
\end{itemize}

\begin{tcolorbox}
\begin{itemize}
\item 定义凸集的维度和单纯形的面；
\end{itemize}
\end{tcolorbox}

\begin{itemize}
\item state the main theorems of this chapter: Sperner's lemma, the KKM lemma, Brouwer's fixed-point theorem both on the unit simplex and on a general compact convex set, and explain the difference between these domains;
\end{itemize}

\begin{tcolorbox}
\begin{itemize}
\item 陈述本章的主要定理:斯佩纳引理、KKM引理、单位单纯形和一般紧凸集上的布劳威尔不动点定理，并解释这些定义域之间的区别；
\end{itemize}
\end{tcolorbox}

\begin{itemize}
\item understand the significance of graphs with maximum node degree two in the proof of Sperner's lemma, and their connection to Lemma 7.6;
\end{itemize}

\begin{tcolorbox}
\begin{itemize}
\item 理解在斯佩纳引理的证明中最大节点度为2的图的重要性，以及它们与引理7.6的联系；
\end{itemize}
\end{tcolorbox}

\begin{itemize}
\item explain the role of the projection onto a closed convex set and why it is contractive and therefore continuous (Theorem 7.18);
\end{itemize}

\begin{tcolorbox}
\begin{itemize}
\item 解释投影到闭凸集上的作用，以及为什么它是压缩的，因此是连续的(定理7.18)；
\end{itemize}
\end{tcolorbox}

\begin{itemize}
\item explain the Freudenthal triangulation.
\end{itemize}

\begin{tcolorbox}
\begin{itemize}
\item 解释弗赖登塔尔三角剖分。
\end{itemize}
\end{tcolorbox}

\section*{7.2 Labels}

\begin{tcolorbox}
\section*{7.2 标签}
\end{tcolorbox}

In this chapter we give a proof of Brouwer's fixed-point Theorem 6.4 for a continuous function \(f\) on a compact and convex set \(C\) . The main effort will be to prove the theorem when \(C\) is the unit simplex \(\Delta\) in \({\mathbb{R}}^{d}\) defined by

\begin{tcolorbox}
在本章中，我们将证明对于紧凸集 \(C\) 上的连续函数 \(f\) 的布劳威尔不动点定理(定理6.4)。主要的工作将是证明当 \(C\) 是由下式定义的 \({\mathbb{R}}^{d}\) 中的单位单纯形 \(\Delta\) 时的定理。
\end{tcolorbox}

\[
\Delta  = \left\{  {x \in  {\mathbb{R}}^{d} \mid  x \geq  \mathbf{0},{\mathbf{1}}^{\top }x = 1}\right\}  . \tag{7.1}
\]

The mixed strategy sets \(X\) and \(Y\) in (6.6) are unit simplices, as shown in Figure 6.5. (However, \(X \times  Y\) is not a unit simplex.)

\begin{tcolorbox}
(6.6)中的混合策略集 \(X\) 和 \(Y\) 是单位单纯形，如图6.5所示。(然而， \(X \times  Y\) 不是单位单纯形。)
\end{tcolorbox}

Recall that the \(i\) th unit vector \({e}_{i}\) in \({\mathbb{R}}^{d}\) has its \(i\) th component equal to one and all other components equal to zero. The unit vectors are the vertices (extreme points or "corners") of the unit simplex \(\Delta\) .

\begin{tcolorbox}
回想一下， \({\mathbb{R}}^{d}\) 中的第 \(i\) 个单位向量 \({e}_{i}\) 的第 \(i\) 个分量等于1，其他所有分量都等于0。单位向量是单位单纯形 \(\Delta\) 的顶点(极点或“角点”)。
\end{tcolorbox}

Let \(f : \Delta  \rightarrow  \Delta\) and \(x \in  \Delta\) , and write \(f\left( x\right)  = {\left( {f}_{1}\left( x\right) ,\ldots ,{f}_{d}\left( x\right) \right) }^{\top }\) . For a point in \({\mathbb{R}}^{d}\) , a label is an element of \(\{ 1,\ldots ,d\}\) , and depending on \(f\) we assign a set of labels

\begin{tcolorbox}
设 \(f : \Delta  \rightarrow  \Delta\) 和 \(x \in  \Delta\) ，并记 \(f\left( x\right)  = {\left( {f}_{1}\left( x\right) ,\ldots ,{f}_{d}\left( x\right) \right) }^{\top }\) 。对于 \({\mathbb{R}}^{d}\) 中的一个点，一个标签是 \(\{ 1,\ldots ,d\}\) 中的一个元素，并且根据 \(f\) 我们为每个点分配一组标签
\end{tcolorbox}

\begin{center}
\includegraphics[max width=0.6\textwidth]{images/bo_d2e5cmbef24c73b03tn0_183_331_209_851_619_0.jpg}
\end{center}
\hspace*{3em} 

Figure 7.1 Illustration of (7.2). The dashed lines show points \(\;\left( {{z}_{1},{z}_{2},{z}_{3}}\right)\) where \({z}_{i} = {x}_{i}\) . Here \(x\) has labels 1 and 2 because \({x}_{1} \geq  {f}_{1}\left( x\right)\) and \({x}_{2} \geq  {f}_{2}\left( x\right)\) .

\begin{tcolorbox}
图 7.1 对 (7.2) 的说明。虚线表示 \(\;\left( {{z}_{1},{z}_{2},{z}_{3}}\right)\) 点，其中 \({z}_{i} = {x}_{i}\) 。这里 \(x\) 有标签 1 和 2，因为 \({x}_{1} \geq  {f}_{1}\left( x\right)\) 且 \({x}_{2} \geq  {f}_{2}\left( x\right)\) 。
\end{tcolorbox}

to each \(x \in  \Delta\) , according to

\begin{tcolorbox}
根据以下规则为每个 \(x \in  \Delta\) 分配
\end{tcolorbox}

\[
x\text{ has label }i \Leftrightarrow  {x}_{i} \geq  {f}_{i}\left( x\right) \;\left( {1 \leq  i \leq  d}\right) \text{ . } \tag{7.2}
\]

Any label \(i\) of \(x = \left( {{x}_{1},\ldots ,{x}_{d}}\right)\) identifies a component \({x}_{i}\) of \(x\) that in \(f\left( x\right)\) is not increased (increased means being moved closer to the vertex \({e}_{i}\) of \(\Delta\) ); see Figure 7.1.

\begin{tcolorbox}
\(x = \left( {{x}_{1},\ldots ,{x}_{d}}\right)\) 的任何标签 \(i\) 都标识了 \(x\) 的一个分量 \({x}_{i}\) ，在 \(f\left( x\right)\) 中该分量不会增加(增加意味着向 \(\Delta\) 的顶点 \({e}_{i}\) 移动)；见图 7.1。
\end{tcolorbox}

A point \(x \in  \Delta\) is completely labeled if \(x\) has every label \(i \in  \{ 1,\ldots ,d\}\) . This condition characterizes the fixed points of \(f\) .

\begin{tcolorbox}
如果 \(x\) 拥有每个标签 \(i \in  \{ 1,\ldots ,d\}\) ，则点 \(x \in  \Delta\) 是完全标记的。这个条件刻画了 \(f\) 的不动点。
\end{tcolorbox}

Lemma 7.1. Let \(f : \Delta  \rightarrow  \Delta\) and \(x \in  \Delta\) . Then \(f\left( x\right)  = x\) if and only if \(x\) is completely labeled.

\begin{tcolorbox}
引理 7.1。设 \(f : \Delta  \rightarrow  \Delta\) 和 \(x \in  \Delta\) 。那么 \(f\left( x\right)  = x\) 当且仅当 \(x\) 是完全标记的。
\end{tcolorbox}

Proof. If \(f\left( x\right)  = x\) then \({f}_{i}\left( x\right)  = {x}_{i}\) for all \(i = 1,\ldots ,d\) and \(x\) is therefore completely labeled. Conversely, suppose \(x\) is completely labeled, \({x}_{i} \geq  {f}_{i}\left( x\right)\) for all \(i\) . If \({x}_{i} > {f}_{i}\left( x\right)\) for some \(i\) then \(\mathop{\sum }\limits_{{k = 1}}^{d}{x}_{k} > \mathop{\sum }\limits_{{k = 1}}^{d}{f}_{k}\left( x\right)  = 1\) , which contradicts \(x \in  \Delta\) . Hence, \({x}_{i} = {f}_{i}\left( x\right)\) for all \(i\) , that is, \(f\left( x\right)  = x\) , as claimed.

\begin{tcolorbox}
证明。如果 \(f\left( x\right)  = x\) ，那么对于所有 \(i = 1,\ldots ,d\) 有 \({f}_{i}\left( x\right)  = {x}_{i}\) ，因此 \(x\) 是完全标记的。反之，假设 \(x\) 是完全标记的，对于所有 \(i\) 有 \({x}_{i} \geq  {f}_{i}\left( x\right)\) 。如果对于某个 \(i\) 有 \({x}_{i} > {f}_{i}\left( x\right)\) ，那么 \(\mathop{\sum }\limits_{{k = 1}}^{d}{x}_{k} > \mathop{\sum }\limits_{{k = 1}}^{d}{f}_{k}\left( x\right)  = 1\) ，这与 \(x \in  \Delta\) 矛盾。因此，对于所有 \(i\) 有 \({x}_{i} = {f}_{i}\left( x\right)\) ，即 \(f\left( x\right)  = x\) ，得证。
\end{tcolorbox}

Our strategy for proving Brouwer's fixed-point theorem is as follows:

\begin{tcolorbox}
我们证明布劳威尔不动点定理的策略如下:
\end{tcolorbox}

\begin{itemize}
\item Discretize \(\Delta\) with finitely many points on some kind of "grid".
\end{itemize}

\begin{tcolorbox}
\begin{itemize}
\item 在某种“网格”上用有限个点对 \(\Delta\) 进行离散化。
\end{itemize}
\end{tcolorbox}

\begin{itemize}
\item For simplicity, assign to each gridpoint a single label rather than a set of labels as in (7.2), so that in (7.2)," \(\Leftrightarrow\) " is replaced by " \(\Rightarrow\) ".
\end{itemize}

\begin{tcolorbox}
\begin{itemize}
\item 为简单起见，为每个网格点分配一个单一的标签，而不是像 (7.2) 中那样分配一组标签，这样在 (7.2) 中，“ \(\Leftrightarrow\) ”被替换为“ \(\Rightarrow\) ”。
\end{itemize}
\end{tcolorbox}

\begin{itemize}
\item Show the existence of a set of "neighboring gridpoints" that together have all labels, which will locate an "approximate fixed point".
\end{itemize}

\begin{tcolorbox}
\begin{itemize}
\item 证明存在一组“相邻网格点”，它们共同拥有所有标签，这将定位一个“近似不动点”。
\end{itemize}
\end{tcolorbox}

\begin{itemize}
\item By choosing a finer and finer grid, a suitable limit of such approximate fixed points defines a fixed point of the considered function.
\end{itemize}

\begin{tcolorbox}
\begin{itemize}
\item 通过选择越来越精细的网格，这种近似不动点的合适极限定义了所考虑函数的一个不动点。
\end{itemize}
\end{tcolorbox}

We explain this strategy for the example in Figure 6.6 where \(f\) is defined on \(\left\lbrack  {0,1}\right\rbrack\) . The interval \(\left\lbrack  {0,1}\right\rbrack\) corresponds to \(\Delta\) in (7.1) for \(d = 2\) , which is the line segment that connects the two points(1,0)and(0,1)in \({\mathbb{R}}^{2}\) , because the elements \(\left( {{x}_{1},{x}_{2}}\right)\) of \(\Delta\) are uniquely given by \(\left( {1 - {x}_{2},{x}_{2}}\right)\) for \({x}_{2} \in  \left\lbrack  {0,1}\right\rbrack\) . Similarly, we can identify a function \(f : \Delta  \rightarrow  \Delta\) with its second component \({f}_{2}\left( {{x}_{1},{x}_{2}}\right)\) because \({f}_{1}\left( {{x}_{1},{x}_{2}}\right)  = 1 - {f}_{2}\left( {{x}_{1},{x}_{2}}\right)\) . Then \(\left( {{x}_{1},{x}_{2}}\right)\) has label 1 if \({x}_{1} \geq  {f}_{1}\left( {{x}_{1},{x}_{2}}\right)\) or equivalently \({x}_{2} \leq  {f}_{2}\left( {1 - {x}_{2},{x}_{2}}\right)\) , and label 2 if \({x}_{2} \geq  {f}_{2}\left( {1 - {x}_{2},{x}_{2}}\right)\) . Furthermore, we can assume

\begin{tcolorbox}
我们以图6.6中的例子来说明这一策略，其中 \(f\) 定义在 \(\left\lbrack  {0,1}\right\rbrack\) 上。区间 \(\left\lbrack  {0,1}\right\rbrack\) 对应于(7.1)中 \(d = 2\) 时的 \(\Delta\) ，它是 \({\mathbb{R}}^{2}\) 中连接两点 (1,0) 和 (0,1) 的线段，因为 \(\Delta\) 中的元素 \(\left( {{x}_{1},{x}_{2}}\right)\) 由 \(\left( {1 - {x}_{2},{x}_{2}}\right)\) 唯一给出，其中 \({x}_{2} \in  \left\lbrack  {0,1}\right\rbrack\) 。类似地，我们可以将函数 \(f : \Delta  \rightarrow  \Delta\) 与其第二个分量 \({f}_{2}\left( {{x}_{1},{x}_{2}}\right)\) 等同起来，因为 \({f}_{1}\left( {{x}_{1},{x}_{2}}\right)  = 1 - {f}_{2}\left( {{x}_{1},{x}_{2}}\right)\) 。那么，如果 \({x}_{1} \geq  {f}_{1}\left( {{x}_{1},{x}_{2}}\right)\) 或者等价地 \({x}_{2} \leq  {f}_{2}\left( {1 - {x}_{2},{x}_{2}}\right)\) ，则 \(\left( {{x}_{1},{x}_{2}}\right)\) 的标签为1；如果 \({x}_{2} \geq  {f}_{2}\left( {1 - {x}_{2},{x}_{2}}\right)\) ，则标签为2。此外，我们可以假设
\end{tcolorbox}

\[
\left( {1,0}\right) \text{ has label }1,\;\left( {0,1}\right) \text{ has label 2 . } \tag{7.3}
\]

\begin{center}
\includegraphics[max width=0.9\textwidth]{images/bo_d2e5cmbef24c73b03tn0_184_241_202_1190_447_0.jpg}
\end{center}
\hspace*{3em} 

Figure 7.2 Two examples of a labeled subdivision of \(\Delta\) for \(d = 2\) . A completely labeled sub-segment contains a fixed point of \(f\) .

\begin{tcolorbox}
图7.2 \(d = 2\) 时 \(\Delta\) 的标记剖分的两个例子。一个完全标记的子线段包含 \(f\) 的一个不动点。
\end{tcolorbox}

The left diagram in Figure 7.2 shows the example from Figure 6.6 with the line segment \(\Delta\) shown as the interval \(\left\lbrack  {0,1}\right\rbrack\) that contains \({x}_{2}\) . A number of "gridpoints" in \(\Delta\) are shown as white dots. Each is given a label according to the sign of \({x}_{2} - {f}_{2}\left( {1 - {x}_{2},{x}_{2}}\right)\) . These points subdivide \(\Delta\) into sub-segments, each of which has a labeled endpoint. A sub-segment that has two endpoints with both labels 1 and 2 is called \({completelylabeled}\) . Such a sub-segment contains a fixed point \(z\) of \(f\) because the sign of \({x}_{2} - {f}_{2}\left( {1 - {x}_{2},{x}_{2}}\right)\) changes when \({x}_{2}\) is in that sub-interval. In Figure 7.2, condition (7.3) holds automatically because(1,0)has only label 1 and (0,1)has only label 2. In general,(7.3) implies that there is at least one completely labeled sub-segment because when changing \({x}_{2}\) from 0 to 1, the label of a gridpoint \(\left( {1 - {x}_{2},{x}_{2}}\right)\) changes at some point from 1 to 2 .

\begin{tcolorbox}
图7.2的左图展示了图6.6中的例子，线段 \(\Delta\) 表示为包含 \({x}_{2}\) 的区间 \(\left\lbrack  {0,1}\right\rbrack\) 。 \(\Delta\) 中的一些“格点”用白点表示。每个格点根据 \({x}_{2} - {f}_{2}\left( {1 - {x}_{2},{x}_{2}}\right)\) 的符号被赋予一个标签。这些点将 \(\Delta\) 细分为子线段，每个子线段都有一个标记的端点。一个子线段的两个端点分别带有标签1和2，则称其为 \({completelylabeled}\) 。这样的子线段包含 \(f\) 的一个不动点 \(z\) ，因为当 \({x}_{2}\) 处于该子区间时， \({x}_{2} - {f}_{2}\left( {1 - {x}_{2},{x}_{2}}\right)\) 的符号会发生变化。在图7.2中，条件(7.3)自动成立，因为 (1,0) 只有标签1，而 (0,1) 只有标签2。一般来说，(7.3)意味着至少存在一个完全标记的子线段，因为当 \({x}_{2}\) 从0变化到1时，某个格点 \(\left( {1 - {x}_{2},{x}_{2}}\right)\) 的标签会在某一点从1变为2。
\end{tcolorbox}

In the right diagram in Figure 7.2, the function \(f\) has three fixed points, which are again identified by the completely labeled sub-segments. In both diagrams, the fixed points would be located more accurately with a "finer" grid.

\begin{tcolorbox}
在图7.2的右图中，函数 \(f\) 有三个不动点，同样由完全标记的子线段确定。在两个图中，使用“更精细”的网格可以更精确地定位不动点。
\end{tcolorbox}

\section*{7.3 Simplices and Triangulations}

\begin{tcolorbox}
\section*{7.3 单纯形与三角剖分}
\end{tcolorbox}

In Figure 7.2, the unit simplex \(\Delta\) in \({\mathbb{R}}^{d}\) is a line segment which is subdivided into smaller sub-segments. The generalization of this, for general \(d\) , is a subdivision of \(\Delta\) into smaller sub-simplices, called a simplicial subdivision or triangulation, which will take the role of the indicated "grid".

\begin{tcolorbox}
在图7.2中， \({\mathbb{R}}^{d}\) 中的单位单纯形 \(\Delta\) 是一个线段，它被细分为更小的子线段。对于一般的 \(d\) ，这种情况的推广是将 \(\Delta\) 细分为更小的子单纯形，称为单纯剖分或三角剖分，它将起到所指示的“网格”的作用。
\end{tcolorbox}

The basic operation from geometry that we need is drawing a line through two points \(x\) and \(y\) in \({\mathbb{R}}^{d}\) , as shown in Figure 5.4, and iterating that process (where a third point not on that line together with any point on the line creates a plane, and so on). The points on the line through \(x\) and \(y\) are the special linear combinations \(x\left( {1 - \lambda }\right)  + {y\lambda }\) for \(\lambda  \in  \mathbb{R}\) called affine combinations. In general, an affine combination of points \({v}^{1},\ldots ,{v}^{n}\) in \({\mathbb{R}}^{d}\) is of the form \(\mathop{\sum }\limits_{{i = 1}}^{n}{v}^{i}{\lambda }_{i}\) where \({\lambda }_{1},\ldots ,{\lambda }_{n}\) are reals with \(\mathop{\sum }\limits_{{i = 1}}^{n}{\lambda }_{i} = 1\) . It is called a convex combination if \({\lambda }_{i} \geq  0\) for \(1 \leq  i \leq  n\) . A set of points is convex if it is closed under forming convex combinations. The convex hull conv \(V\) of a set of points \(V\) is the (with respect to set inclusion) smallest convex set \(C\) so that \(V \subseteq  C\) .

\begin{tcolorbox}
我们需要的几何学基本操作是，在 \({\mathbb{R}}^{d}\) 中过两点 \(x\) 和 \(y\) 画一条直线，如图5.4所示，并重复这个过程(不在该直线上的第三点与直线上的任意一点可确定一个平面，依此类推)。过 \(x\) 和 \(y\) 的直线上的点是形如 \(x\left( {1 - \lambda }\right)  + {y\lambda }\) (其中 \(\lambda  \in  \mathbb{R}\) )的特殊线性组合，称为仿射组合(affine combination)。一般地， \({\mathbb{R}}^{d}\) 中各点 \({v}^{1},\ldots ,{v}^{n}\) 的仿射组合形式为 \(\mathop{\sum }\limits_{{i = 1}}^{n}{v}^{i}{\lambda }_{i}\) ，其中 \({\lambda }_{1},\ldots ,{\lambda }_{n}\) 为实数且满足 \(\mathop{\sum }\limits_{{i = 1}}^{n}{\lambda }_{i} = 1\) 。若对于 \(1 \leq  i \leq  n\) 有 \({\lambda }_{i} \geq  0\) ，则称其为凸组合(convex combination)。如果一个点集在形成凸组合时是封闭的，那么该点集就是凸集(convex set)。点集 \(V\) 的凸包(convex hull)conv \(V\) 是(关于集合包含关系)包含 \(V\) 的最小凸集 \(C\) ，即 \(V \subseteq  C\) 。
\end{tcolorbox}

Given points are affinely independent if none of these points is an affine combination of the others. This has an easy equivalent characterization in terms of linear independence in one higher dimension.

\begin{tcolorbox}
若给定的点集中没有一个点是其他点的仿射组合，则这些点是仿射无关的。这可以通过高一维的线性无关性给出一个简单的等价刻画。
\end{tcolorbox}

Lemma 7.2. The vectors \({v}^{1},\ldots ,{v}^{n}\) in \({\mathbb{R}}^{d}\) are affinely independent if and only if the

\begin{tcolorbox}
引理7.2. \({\mathbb{R}}^{d}\) 中的向量 \({v}^{1},\ldots ,{v}^{n}\) 仿射无关当且仅当
\end{tcolorbox}

vectors

\begin{tcolorbox}
向量
\end{tcolorbox}

\[
\left( \begin{array}{l} 1 \\  {v}^{1} \end{array}\right) ,\ldots ,\left( \begin{matrix} 1 \\  {v}^{n} \end{matrix}\right)  \tag{7.4}
\]

in \({\mathbb{R}}^{1 + d}\) are linearly independent.

\begin{tcolorbox}
在 \({\mathbb{R}}^{1 + d}\) 中线性无关。
\end{tcolorbox}

Proof. Suppose \({v}^{1},\ldots ,{v}^{n}\) in \({\mathbb{R}}^{d}\) are not affinely independent. That is, some vector, say \({v}^{1}\) , is an affine combination of the others, with reals \({\lambda }_{2},\ldots ,{\lambda }_{n}\) so that

\begin{tcolorbox}
证明. 假设 \({\mathbb{R}}^{d}\) 中的 \({v}^{1},\ldots ,{v}^{n}\) 不是仿射无关的。也就是说，某个向量，比如 \({v}^{1}\) ，是其他向量的仿射组合，存在实数 \({\lambda }_{2},\ldots ,{\lambda }_{n}\) 使得
\end{tcolorbox}

\[
1 = {\lambda }_{2} + \cdots  + {\lambda }_{n}
\]

\[
{v}^{1} = {v}^{2}{\lambda }_{2} + \cdots  + {v}^{n}{\lambda }_{n}.
\]

This is equivalent to

\begin{tcolorbox}
这等价于
\end{tcolorbox}

\[
\left( \begin{matrix} 1 \\  {v}^{1} \end{matrix}\right) \left( {-1}\right)  + \left( \begin{matrix} 1 \\  {v}^{2} \end{matrix}\right) {\lambda }_{2} + \cdots  + \left( \begin{matrix} 1 \\  {v}^{n} \end{matrix}\right) {\lambda }_{n} = \left( \begin{array}{l} 0 \\  \mathbf{0} \end{array}\right)  \tag{7.5}
\]

which shows that the vectors in (7.4) are not linearly independent. Conversely, if they are not linearly independent, with

\begin{tcolorbox}
这表明(7.4)中的向量不是线性无关的。反之，如果它们不是线性无关的，即
\end{tcolorbox}

\[
\left( \begin{matrix} 1 \\  {v}^{1} \end{matrix}\right) {\lambda }_{1} + \left( \begin{matrix} 1 \\  {v}^{2} \end{matrix}\right) {\lambda }_{2} + \cdots  + \left( \begin{matrix} 1 \\  {v}^{n} \end{matrix}\right) {\lambda }_{n} = \left( \begin{array}{l} 0 \\  \mathbf{0} \end{array}\right)  \tag{7.6}
\]

where \({\lambda }_{1},\ldots ,{\lambda }_{n}\) are not all zero, say \({\lambda }_{1} \neq  0\) , then multiplication of (7.6) with \(- 1/{\lambda }_{1}\) gives (7.5) and therefore the affine dependence of \({v}^{1},\ldots ,{v}^{n}\) . (The additional first row with entries 1 in (7.4) is needed for linear independence, for example if some \({v}^{i}\) is the zero vector, which has no special role in affine combinations.)

\begin{tcolorbox}
其中 \({\lambda }_{1},\ldots ,{\lambda }_{n}\) 不全为零，不妨设 \({\lambda }_{1} \neq  0\) ，那么用 \(- 1/{\lambda }_{1}\) 乘以(7.6)式就得到(7.5)式，从而说明 \({v}^{1},\ldots ,{v}^{n}\) 是仿射相关的。((7.4)中额外的第一行元素都为1是为了保证线性无关性，例如当某个 \({v}^{i}\) 是零向量时，零向量在仿射组合中没有特殊作用。)
\end{tcolorbox}

A convex set has dimension \(n\) if and only if it has \(n + 1\) , but no more, affinely independent points. A polytope is the convex hull of a finite set of points in some Euclidean space \({\mathbb{R}}^{d}\) . If these points are affinely independent, then the polytope is called a simplex.

\begin{tcolorbox}
一个凸集的维度(dimension)为 \(n\) 当且仅当它有 \(n + 1\) 个但不会更多的仿射无关点。有界多面体(polytope)是某个欧几里得空间 \({\mathbb{R}}^{d}\) 中有限点集的凸包。如果这些点是仿射无关的，那么这个有界多面体就称为单纯形(simplex)。
\end{tcolorbox}

The unit simplex \(\Delta\) in (7.1) is the convex hull of the unit vectors \({e}_{1},\ldots ,{e}_{d}\) in \({\mathbb{R}}^{d}\) . They are affinely independent, so the unit simplex has dimension \(d - 1\) , and for that reason is sometimes denoted by \({\Delta }_{d - 1}\) . For example (see Figure 7.1), \({\Delta }_{2}\) is the convex hull of the three unit vectors in \({\mathbb{R}}^{3}\) and a triangle, which has dimension two (even though it is a subset of \({\mathbb{R}}^{3}\) , because all \({\left( {x}_{1},{x}_{2},{x}_{3}\right) }^{\top }\) in \({\Delta }_{2}\) fulfill the equation \(\left. {{x}_{1} + {x}_{2} + {x}_{3} = 1}\right)\) .

\begin{tcolorbox}
式(7.1)中的单位单纯形 \(\Delta\) 是 \({\mathbb{R}}^{d}\) 中单位向量 \({e}_{1},\ldots ,{e}_{d}\) 的凸包。它们仿射无关，因此单位单纯形的维度为 \(d - 1\) ，出于这个原因，它有时用 \({\Delta }_{d - 1}\) 表示。例如(见图7.1)， \({\Delta }_{2}\) 是 \({\mathbb{R}}^{3}\) 中三个单位向量的凸包，是一个三角形，其维度为2(尽管它是 \({\mathbb{R}}^{3}\) 的子集，因为 \({\Delta }_{2}\) 中的所有 \({\left( {x}_{1},{x}_{2},{x}_{3}\right) }^{\top }\) 都满足方程 \(\left. {{x}_{1} + {x}_{2} + {x}_{3} = 1}\right)\) )。
\end{tcolorbox}

In agreement with the definition of dimension, the empty set has dimension -1. A set of dimension zero is any singleton. A simplex of dimension one is the convex hull of two distinct points and therefore a line segment. A simplex of dimension two is a triangle. A simplex of dimension three is a tetrahedron. Every simplex of a given dimension (in any Euclidean space of possibly much larger dimension) is in an affine (and hence continuous) bijection with the unit simplex of the same dimension, as follows.

\begin{tcolorbox}
根据维度的定义，空集的维度为 -1。维度为零的集合是任何单元素集。一维单纯形是两个不同点的凸包，因此是一条线段。二维单纯形是一个三角形。三维单纯形是一个四面体。给定维度的每个单纯形(在任何可能维度大得多的欧几里得空间中)都与相同维度的单位单纯形存在仿射(因此是连续的)双射，如下所示。
\end{tcolorbox}

Lemma 7.3. Let \(S\) be a simplex of dimension \(d - 1\) given as the convex hull of the affinely independent points \({v}^{1},\ldots ,{v}^{d}\) in some \({\mathbb{R}}^{k}\) . Then there is a linear map \({\mathbb{R}}^{d} \rightarrow  {\mathbb{R}}^{k}\) that restricted to \(\Delta\) in (7.1) is a bijection \(\Delta  \rightarrow  S\) . Its inverse \(J : S \rightarrow  \Delta\) is the restriction of an affine map \({\mathbb{R}}^{k} \rightarrow  {\mathbb{R}}^{d}\) .

\begin{tcolorbox}
引理7.3。设 \(S\) 是一个维度为 \(d - 1\) 的单纯形，它是某个 \({\mathbb{R}}^{k}\) 中仿射无关点 \({v}^{1},\ldots ,{v}^{d}\) 的凸包。那么存在一个线性映射 \({\mathbb{R}}^{d} \rightarrow  {\mathbb{R}}^{k}\) ，其在式(7.1)中的 \(\Delta\) 上的限制是一个双射 \(\Delta  \rightarrow  S\) 。其逆 \(J : S \rightarrow  \Delta\) 是一个仿射映射 \({\mathbb{R}}^{k} \rightarrow  {\mathbb{R}}^{d}\) 的限制。
\end{tcolorbox}

Proof. If \(x \in  S\) , then there are nonnegative reals \({\lambda }_{1},\ldots ,{\lambda }_{d}\) such that

\begin{tcolorbox}
证明。如果 \(x \in  S\) ，那么存在非负实数 \({\lambda }_{1},\ldots ,{\lambda }_{d}\) 使得
\end{tcolorbox}

\[
\left( \begin{array}{l} 1 \\  x \end{array}\right)  = \left( \begin{matrix} 1 \\  {v}^{1} \end{matrix}\right) {\lambda }_{1} + \cdots  + \left( \begin{matrix} 1 \\  {v}^{d} \end{matrix}\right) {\lambda }_{d} \tag{7.7}
\]

By affine independence and Lemma 7.2, the vector \(\lambda  = {\left( {\lambda }_{1},\ldots ,{\lambda }_{d}\right) }^{\top }\) of coefficients in (7.7) is unique, and different coefficient vectors define different points \(x \in  S\) . These coefficient vectors \(\lambda\) belong to the unit simplex \(\Delta\) , which is therefore in a bijection \(\Delta  \rightarrow  S\) via the linear map \(\lambda  \mapsto  x\) in (7.7).

\begin{tcolorbox}
由仿射无关性和引理7.2可知，式(7.7)中的系数向量 \(\lambda  = {\left( {\lambda }_{1},\ldots ,{\lambda }_{d}\right) }^{\top }\) 是唯一的，不同的系数向量定义不同的点 \(x \in  S\) 。这些系数向量 \(\lambda\) 属于单位单纯形 \(\Delta\) ，因此通过式(7.7)中的线性映射 \(\lambda  \mapsto  x\) 与 \(\Delta  \rightarrow  S\) 存在双射。
\end{tcolorbox}

The inverse of this bijection is seen to be the restriction of an affine map \({\mathbb{R}}^{k} \rightarrow  {\mathbb{R}}^{d}\) as follows. We first construct a linear map \({\mathbb{R}}^{1 + k} \rightarrow  {\mathbb{R}}^{d}\) , which is uniquely defined by its (arbitrary) images on a basis of \({\mathbb{R}}^{1 + k}\) . For this purpose, extend \(\left( \begin{matrix} 1 \\  {v}^{1} \end{matrix}\right) ,\ldots ,\left( \begin{matrix} 1 \\  {v}^{d} \end{matrix}\right)\) to such a basis, map \(\left( \begin{matrix} 1 \\  {v}^{i} \end{matrix}\right)\) to the unit vector \({e}_{i}\) in \({\mathbb{R}}^{d}\) for \(1 \leq  i \leq  d\) , and each remaining basis vector to0. Because this is a linear map, it can be written as \(\left( \begin{array}{l} \alpha \\  x \end{array}\right)  \mapsto  \lambda  = A\left( \begin{array}{l} \alpha \\  x \end{array}\right)\) for \(\alpha  \in  \mathbb{R}\) and \(x \in  {\mathbb{R}}^{k}\) for some \(d \times  \left( {1 + k}\right)\) matrix \(A\) . By (7.7), \(\lambda  \in  \Delta\) if \(\alpha  = 1\) and \(x \in  S\) . For \(\alpha  = 1\) we obtain an affine map \({\mathbb{R}}^{k} \rightarrow  {\mathbb{R}}^{d}\) , \(x \mapsto  \lambda  = A\left( \begin{array}{l} 1 \\  x \end{array}\right)\) . The restriction of this map to \(S\) is the desired affine bijection \(J : S \rightarrow  \Delta\) . (If \(x \notin  S\) then \(A\left( \begin{array}{l} 1 \\  x \end{array}\right)\) depends on the choice of the additional basis vectors, so \(A\) itself is not unique, but this is irrelevant here; the restriction \(J : S \rightarrow  \Delta\) is unique.)

\begin{tcolorbox}
可以看出，这个双射的逆是如下仿射映射 \({\mathbb{R}}^{k} \rightarrow  {\mathbb{R}}^{d}\) 的限制。我们首先构造一个线性映射 \({\mathbb{R}}^{1 + k} \rightarrow  {\mathbb{R}}^{d}\) ，它由其在 \({\mathbb{R}}^{1 + k}\) 的一组基上的(任意)像唯一确定。为此，将 \(\left( \begin{matrix} 1 \\  {v}^{1} \end{matrix}\right) ,\ldots ,\left( \begin{matrix} 1 \\  {v}^{d} \end{matrix}\right)\) 扩充为这样一组基，对于 \(1 \leq  i \leq  d\) ，将 \(\left( \begin{matrix} 1 \\  {v}^{i} \end{matrix}\right)\) 映射到 \({\mathbb{R}}^{d}\) 中的单位向量 \({e}_{i}\) ，并将其余每个基向量映射到 0。因为这是一个线性映射，所以对于某个 \(d \times  \left( {1 + k}\right)\) 矩阵 \(A\) ，它可以写成 \(\left( \begin{array}{l} \alpha \\  x \end{array}\right)  \mapsto  \lambda  = A\left( \begin{array}{l} \alpha \\  x \end{array}\right)\) (其中 \(\alpha  \in  \mathbb{R}\) 且 \(x \in  {\mathbb{R}}^{k}\) )。根据 (7.7)，若 \(\alpha  = 1\) 且 \(x \in  S\) ，则 \(\lambda  \in  \Delta\) 。对于 \(\alpha  = 1\) ，我们得到一个仿射映射 \({\mathbb{R}}^{k} \rightarrow  {\mathbb{R}}^{d}\) ，即 \(x \mapsto  \lambda  = A\left( \begin{array}{l} 1 \\  x \end{array}\right)\) 。这个映射在 \(S\) 上的限制就是所需的仿射双射 \(J : S \rightarrow  \Delta\) 。(若 \(x \notin  S\) ，则 \(A\left( \begin{array}{l} 1 \\  x \end{array}\right)\) 依赖于额外基向量的选取，所以 \(A\) 本身不是唯一的，但这在这里无关紧要；限制映射 \(J : S \rightarrow  \Delta\) 是唯一的。)
\end{tcolorbox}

In (7.7), the unique coefficients \({\lambda }_{1},\ldots ,{\lambda }_{d}\) are called the barycentric coordinates of \(x\) in \(S\) (which depend on the chosen order of the points \({v}^{1},\ldots ,{v}^{d}\) that define the simplex \(S\) ). The reason for the name "barycentric" is that \(x\) is the center of weight (barycenter) of the weights \({\lambda }_{i}\) placed at the points \({v}^{i}\) for \(1 \leq  i \leq  d\) .

\begin{tcolorbox}
在 (7.7) 中，唯一的系数 \({\lambda }_{1},\ldots ,{\lambda }_{d}\) 被称为 \(x\) 在 \(S\) 中的重心坐标(它依赖于定义单纯形 \(S\) 的点 \({v}^{1},\ldots ,{v}^{d}\) 的选取顺序)。“重心”这个名称的由来是，对于 \(1 \leq  i \leq  d\) ， \(x\) 是放置在点 \({v}^{i}\) 处权重为 \({\lambda }_{i}\) 的加权重心(重心)。
\end{tcolorbox}

Consider a simplex \(S\) of dimension \(d - 1\) given as the convex hull of the \(d\) affinely independent points \({v}^{1},\ldots ,{v}^{d}\) . Then any subset of these points (including the empty set) is also affinely independent, and its convex hull is also a simplex. Each such simplex is a called a face of \(S\) . Any singleton face \(\left\{  {v}^{i}\right\}\) , or \({v}^{i}\) itself, is called a vertex of \(S\) . A face of dimension 1 is called an edge. A face of dimension \(d - 2\) of \(S\) is called a facet of \(S\) . A vertex \(v\) of a general polytope \(P\) is an extreme point of \(P\) , that is, \(v \in  P\) and \(v\) is not a convex combination of other points of \(P\) . Exercise 7.2 shows that each vertex of a simplex is indeed such an extreme point.

\begin{tcolorbox}
考虑一个维度为 \(d - 1\) 的单纯形 \(S\) ，它是 \(d\) 个仿射无关点 \({v}^{1},\ldots ,{v}^{d}\) 的凸包。那么这些点的任何子集(包括空集)也是仿射无关的，并且其凸包也是一个单纯形。每个这样的单纯形都称为 \(S\) 的一个面。任何单元素面 \(\left\{  {v}^{i}\right\}\) 或 \({v}^{i}\) 本身都称为 \(S\) 的一个顶点。维度为 1 的面称为一条边。 \(S\) 中维度为 \(d - 2\) 的面称为 \(S\) 的一个极面。一般有界多面体 \(P\) 的一个顶点 \(v\) 是 \(P\) 的一个极点，即 \(v \in  P\) 且 \(v\) 不是 \(P\) 中其他点的凸组合。练习 7.2 表明，单纯形的每个顶点确实是这样一个极点。
\end{tcolorbox}

Definition 7.4. Let \(P\) be a polytope. A triangulation or simplicial subdivision of \(P\) is a finite set \(\sum  = \left\{  {{S}_{1},\ldots ,{S}_{N}}\right\}\) of simplices (called sub-simplices) \({S}_{i}\) for \(i = 1,\ldots ,N\) so that (a) \(P = \mathop{\bigcup }\limits_{{i = 1}}^{N}{S}_{i}\) ,(b) no sub-simplex is a subset of another, and (c) the intersection of any two sub-simplices is a face of both. A vertex in \(\sum\) is any vertex of any sub-simplex in \(\sum\) . A sub-facet in \(\sum\) is any facet of any sub-simplex in \(\sum\) .

\begin{tcolorbox}
定义7.4。设 \(P\) 为一个有界多面体(polytope)。 \(P\) 的一个三角剖分(triangulation)或单纯剖分(simplicial subdivision)是一个由单纯形(simplex，称为子单纯形，sub - simplex) \({S}_{i}\) 组成的有限集合 \(\sum  = \left\{  {{S}_{1},\ldots ,{S}_{N}}\right\}\) ，其中 \(i = 1,\ldots ,N\) 满足:(a) \(P = \mathop{\bigcup }\limits_{{i = 1}}^{N}{S}_{i}\) ；(b) 没有一个子单纯形是另一个子单纯形的子集；(c) 任意两个子单纯形的交集是这两个子单纯形的一个面(face)。 \(\sum\) 中的一个顶点(vertex)是 \(\sum\) 中任意一个子单纯形的任意一个顶点。 \(\sum\) 中的一个子极面(sub - facet)是 \(\sum\) 中任意一个子单纯形的任意一个极面(facet)。
\end{tcolorbox}

\begin{center}
\includegraphics[max width=0.9\textwidth]{images/bo_d2e5cmbef24c73b03tn0_187_340_1381_1248_325_0.jpg}
\end{center}
\hspace*{3em} 

Figure 7.3 Left: \(\left\{  {{S}_{1},{S}_{2},{S}_{3}}\right\}\) is not a triangulation of the triangle because \({S}_{1} \cap  {S}_{3}\) is not a face of \({S}_{3}\) . Center and right: triangulations of the triangle.

\begin{tcolorbox}
图7.3 左图: \(\left\{  {{S}_{1},{S}_{2},{S}_{3}}\right\}\) 不是该三角形的一个三角剖分，因为 \({S}_{1} \cap  {S}_{3}\) 不是 \({S}_{3}\) 的一个面。中图和右图:该三角形的三角剖分。
\end{tcolorbox}

In Definition 7.4(c), the intersection of two sub-simplices may be the empty set, which is a face of any polytope. Figure 7.3 illustrates the concept of simplicial subdivisions of a triangle. The left picture fulfills conditions (a) and (b) of Definition 7.4 but not (c), which does hold in the middle picture. The right picture shows that the sub-simplices in a triangulation may vary in shape and size.

\begin{tcolorbox}
在定义7.4(c)中，两个子单纯形的交集可能是空集，空集是任意有界多面体的一个面。图7.3展示了三角形的单纯剖分的概念。左图满足定义7.4的条件(a)和(b)，但不满足条件(c)，而中图满足条件(c)。右图表明，三角剖分中的子单纯形的形状和大小可以不同。
\end{tcolorbox}

The following lemmas state some, rather intuitive, properties of triangulations.

\begin{tcolorbox}
以下引理陈述了三角剖分的一些相当直观的性质。
\end{tcolorbox}

Lemma 7.5. Let \(S\) be a simplex with vertices \({v}^{1},\ldots ,{v}^{d}\) and let \(\sum\) be a triangulation of \(S\) . Then all sub-simplices \({S}_{i}\) in \(\sum\) have the same dimension \(d - 1\) as \(S\) .

\begin{tcolorbox}
引理7.5。设 \(S\) 是一个顶点为 \({v}^{1},\ldots ,{v}^{d}\) 的单纯形， \(\sum\) 是 \(S\) 的一个三角剖分。那么 \(\sum\) 中的所有子单纯形 \({S}_{i}\) 与 \(S\) 具有相同的维度(dimension) \(d - 1\) 。
\end{tcolorbox}

Proof. The vertices of any sub-simplex belong to \(S\) , which has dimension \(d - 1\) , so no sub-simplex can have higher dimension. Suppose one sub-simplex, call it \({S}_{1}\) , has lower dimension, with vertices \({w}^{1},\ldots ,{w}^{k}\) and \(k < d\) . Let \(x = \mathop{\sum }\limits_{{i = 1}}^{k}{w}^{i}\frac{1}{k}\) be the barycenter of \({S}_{1}\) . At least one vertex \({v}^{i}\) of \(S\) is not an affine combination of \({w}^{1},\ldots ,{w}^{k}\) . For \(\varepsilon  \in  \mathbb{R}\) , let

\begin{tcolorbox}
证明。任意子单纯形的顶点都属于 \(S\) ， \(S\) 的维度为 \(d - 1\) ，所以没有子单纯形的维度可以更高。假设一个子单纯形，称它为 \({S}_{1}\) ，维度更低，其顶点为 \({w}^{1},\ldots ,{w}^{k}\) 和 \(k < d\) 。设 \(x = \mathop{\sum }\limits_{{i = 1}}^{k}{w}^{i}\frac{1}{k}\) 是 \({S}_{1}\) 的重心(barycenter)。 \(S\) 中至少有一个顶点 \({v}^{i}\) 不是 \({w}^{1},\ldots ,{w}^{k}\) 的仿射组合(affine combination)。对于 \(\varepsilon  \in  \mathbb{R}\) ，设
\end{tcolorbox}

\[
z\left( \varepsilon \right)  = x\left( {1 - \varepsilon }\right)  + {v}^{i}\varepsilon . \tag{7.8}
\]

Geometrically, the points \(z\left( \varepsilon \right)\) lie on the line through \(x\) and \({v}^{i}\) , see Figure 7.4.

\begin{tcolorbox}
从几何角度看，点 \(z\left( \varepsilon \right)\) 位于过 \(x\) 和 \({v}^{i}\) 的直线上，见图7.4。
\end{tcolorbox}

\begin{center}
\includegraphics[max width=0.4\textwidth]{images/bo_d2e5cmbef24c73b03tn0_188_204_724_609_595_0.jpg}
\end{center}
\hspace*{3em} 

Figure 7.4 Illustration of (7.8) for \(d = 4,k = 3\) , and \(i = 2\) in (7.9).

\begin{tcolorbox}
图7.4 展示了(7.8)式对于 \(d = 4,k = 3\) 以及(7.9)式中 \(i = 2\) 的情况。
\end{tcolorbox}

Then for \(\varepsilon  \neq  0\) we have \({v}^{i} = z\left( \varepsilon \right) \frac{1}{\varepsilon } + x\left( {1 - \frac{1}{\varepsilon }}\right)\) , which shows that \(z\left( \varepsilon \right)\) is not an affine combination of \({w}^{1},\ldots ,{w}^{k}\) because otherwise \({v}^{i}\) would be. That is,

\begin{tcolorbox}
那么对于 \(\varepsilon  \neq  0\) ，我们有 \({v}^{i} = z\left( \varepsilon \right) \frac{1}{\varepsilon } + x\left( {1 - \frac{1}{\varepsilon }}\right)\) ，这表明 \(z\left( \varepsilon \right)\) 不是 \({w}^{1},\ldots ,{w}^{k}\) 的仿射组合，因为否则 \({v}^{i}\) 也会是。也就是说，
\end{tcolorbox}

\[
z\left( \varepsilon \right)  \notin  {S}_{1}\;\left( {\varepsilon  \neq  0}\right) . \tag{7.9}
\]

Consider a decreasing sequence \({\varepsilon }_{1},{\varepsilon }_{2},\ldots\) in \((0,1\rbrack\) that tends to 0, for example \({\varepsilon }_{n} = \frac{1}{n}\) , where \(z\left( {\varepsilon }_{n}\right)  \in  S\) by (7.8) because \(S\) is convex. Because \(S\) is the union of the sub-simplices in \(\sum\) , there is for each \(n\) some sub-simplex that contains \(z\left( {\varepsilon }_{n}\right)\) . Because \(\sum\) is finite, one sub-simplex, call it \({S}_{2}\) , contains infinitely many of the points \(z\left( {\varepsilon }_{n}\right)\) (and in fact all points in between in that sequence because they are on a line segment and \({S}_{2}\) is convex). Because \({\varepsilon }_{n} \rightarrow  0\) as \(n \rightarrow  \infty\) , the points \(z\left( {\varepsilon }_{n}\right)\) converge to \(x\) , which implies \(x \in  {S}_{2}\) because \({S}_{2}\) is closed. Hence, \(x \in  {S}_{1} \cap  {S}_{2}\) . By condition (c) in Definition 7.4, the intersection \({S}_{1} \cap  {S}_{2}\) is a face of both \({S}_{1}\) and \({S}_{2}\) , which includes all vertices \({w}^{1},\ldots ,{w}^{k}\) of \({S}_{1}\) because they are used with positive weight \(\frac{1}{k}\) in the representation \(x = \mathop{\sum }\limits_{{i = 1}}^{k}{w}^{i}\frac{1}{k}\) . Hence, \({S}_{1} = {S}_{1} \cap  {S}_{2}\) , that is, \({S}_{1} \subseteq  {S}_{2}\) . Moreover, \(z\left( \varepsilon \right)  \in  {S}_{2} - {S}_{1}\) for some \(\varepsilon  > 0\) , which implies \({S}_{1} \neq  {S}_{2}\) . This contradicts Definition 7.4(b). Hence, all sub-simplices have the same dimension \(d - 1\) as \(S\) , as claimed.

\begin{tcolorbox}
考虑在 \((0,1\rbrack\) 中的一个递减序列 \({\varepsilon }_{1},{\varepsilon }_{2},\ldots\) ，它趋于 0，例如 \({\varepsilon }_{n} = \frac{1}{n}\) ，根据 (7.8) 式，因为 \(S\) 是凸集，所以 \(z\left( {\varepsilon }_{n}\right)  \in  S\) 。由于 \(S\) 是 \(\sum\) 中子单纯形(sub - simplex)的并集，对于每个 \(n\) ，存在某个子单纯形包含 \(z\left( {\varepsilon }_{n}\right)\) 。因为 \(\sum\) 是有限的，所以有一个子单纯形，记为 \({S}_{2}\) ，包含序列中无限多个点 \(z\left( {\varepsilon }_{n}\right)\) (实际上，该序列中介于这些点之间的所有点也在其中，因为它们在一条线段(line segment)上，且 \({S}_{2}\) 是凸集)。由于当 \(n \rightarrow  \infty\) 时 \({\varepsilon }_{n} \rightarrow  0\) ，点 \(z\left( {\varepsilon }_{n}\right)\) 收敛于 \(x\) ，这意味着 \(x \in  {S}_{2}\) ，因为 \({S}_{2}\) 是闭集。因此， \(x \in  {S}_{1} \cap  {S}_{2}\) 。根据定义 7.4 中的条件 (c)，交集 \({S}_{1} \cap  {S}_{2}\) 是 \({S}_{1}\) 和 \({S}_{2}\) 的一个面(face)，它包含 \({S}_{1}\) 的所有顶点(vertex) \({w}^{1},\ldots ,{w}^{k}\) ，因为在表示式 \(x = \mathop{\sum }\limits_{{i = 1}}^{k}{w}^{i}\frac{1}{k}\) 中它们的权重 \(\frac{1}{k}\) 为正。因此， \({S}_{1} = {S}_{1} \cap  {S}_{2}\) ，即 \({S}_{1} \subseteq  {S}_{2}\) 。此外，对于某个 \(\varepsilon  > 0\) 有 \(z\left( \varepsilon \right)  \in  {S}_{2} - {S}_{1}\) ，这意味着 \({S}_{1} \neq  {S}_{2}\) 。这与定义 7.4(b) 矛盾。因此，如所声称的，所有子单纯形的维度(dimension) \(d - 1\) 都与 \(S\) 相同。
\end{tcolorbox}

Lemma 7.6. Let \(S\) be a simplex, let \(\sum\) be a triangulation of \(S\) , and let \(F\) be a sub-facet in \(\sum\) . Then exactly one of the following two cases occurs:

\begin{tcolorbox}
引理 7.6。设 \(S\) 为一个单纯形(simplex)， \(\sum\) 为 \(S\) 的一个三角剖分(triangulation)， \(F\) 为 \(\sum\) 中的一个子极面(sub - facet)。则以下两种情况恰好发生其一:
\end{tcolorbox}

(a) \(F\) is a subset of a facet of \(S\) , and \(F\) is the facet of exactly one sub-simplex in \(\sum\) .

\begin{tcolorbox}
(a) \(F\) 是 \(S\) 的一个极面(facet)的子集，且 \(F\) 是 \(\sum\) 中恰好一个子单纯形的极面。
\end{tcolorbox}

(b) \(F\) is not a subset of a facet of \(S\) , and \(F\) is the facet of exactly two sub-simplices in \(\sum\) .

\begin{tcolorbox}
(b) \(F\) 不是 \(S\) 的一个极面的子集，且 \(F\) 是 \(\sum\) 中恰好两个子单纯形的极面。
\end{tcolorbox}

Proof. Let \(S\) have dimension \(d - 1\) and vertices \({v}^{1},\ldots ,{v}^{d}\) . By Lemma 7.5, every sub-simplex also has dimension \(d - 1\) . Let \(F\) be a facet of some sub-simplex \(T\) in \(\sum\) and let \({w}^{1},\ldots ,{w}^{d - 1}\) be the vertices of \(F\) . Let \(x = \mathop{\sum }\limits_{{i = 1}}^{{d - 1}}{w}^{i}\frac{1}{d - 1}\) be the barycenter of \(F\) , and let \({w}^{d}\) be the remaining vertex of \(T\) . Consider the line through \(x\) and \({w}^{d}\) given by the points

\begin{tcolorbox}
证明。设 \(S\) 的维度为 \(d - 1\) ，顶点为 \({v}^{1},\ldots ,{v}^{d}\) 。根据引理 7.5，每个子单纯形的维度也为 \(d - 1\) 。设 \(F\) 是 \(\sum\) 中某个子单纯形 \(T\) 的一个极面， \({w}^{1},\ldots ,{w}^{d - 1}\) 是 \(F\) 的顶点。设 \(x = \mathop{\sum }\limits_{{i = 1}}^{{d - 1}}{w}^{i}\frac{1}{d - 1}\) 是 \(F\) 的重心(barycenter)， \({w}^{d}\) 是 \(T\) 剩下的那个顶点。考虑过 \(x\) 和 \({w}^{d}\) 的直线，该直线上的点为
\end{tcolorbox}

\[
y\left( \varepsilon \right)  = x\left( {1 - \varepsilon }\right)  + {w}^{d}\varepsilon
\]

for \(\varepsilon  \in  \mathbb{R}\) , which for \(\varepsilon  \in  \left\lbrack  {0,1}\right\rbrack\) belong to \(T\) . Suppose \(F\) is the facet of another sub-simplex \({T}^{\prime }\) , which has another \(d\) th vertex \({w}_{0}\) , different from \({w}_{d}\) . The \(d\) vertices \({w}^{0},{w}^{1},\ldots ,{w}^{d - 1}\) of \({T}^{\prime }\) are affinely independent. Hence, there are unique \({\lambda }_{0},{\lambda }_{1},\ldots ,{\lambda }_{d - 1}\) with \({w}^{d} = \mathop{\sum }\limits_{{i = 0}}^{{d - 1}}{w}^{i}{\lambda }_{i}\) , where \({\lambda }_{0} \neq  0\) because \({w}^{d} \notin  F\) . We claim that \({\lambda }_{0} < 0\) , that is, \({w}_{0}\) and \({w}_{d}\) are on opposite sides of the facet \(F\) . Namely,

\begin{tcolorbox}
对于 \(\varepsilon  \in  \mathbb{R}\) ，其中对于 \(\varepsilon  \in  \left\lbrack  {0,1}\right\rbrack\) 属于 \(T\) 。假设 \(F\) 是另一个子单纯形 \({T}^{\prime }\) 的极面，该子单纯形有另一个第 \(d\) 个顶点 \({w}_{0}\) ，不同于 \({w}_{d}\) 。 \({T}^{\prime }\) 的 \(d\) 个顶点 \({w}^{0},{w}^{1},\ldots ,{w}^{d - 1}\) 仿射无关。因此，存在唯一的 \({\lambda }_{0},{\lambda }_{1},\ldots ,{\lambda }_{d - 1}\) 满足 \({w}^{d} = \mathop{\sum }\limits_{{i = 0}}^{{d - 1}}{w}^{i}{\lambda }_{i}\) ，其中 \({\lambda }_{0} \neq  0\) 是因为 \({w}^{d} \notin  F\) 。我们断言 \({\lambda }_{0} < 0\) ，即 \({w}_{0}\) 和 \({w}_{d}\) 在极面 \(F\) 的两侧。即
\end{tcolorbox}

\[
y\left( \varepsilon \right)  = x\left( {1 - \varepsilon }\right)  + \mathop{\sum }\limits_{{i = 0}}^{{d - 1}}{w}^{i}{\lambda }_{i}\varepsilon
\]

\[
= {w}^{0}{\lambda }_{0}\varepsilon  + \mathop{\sum }\limits_{{i = 1}}^{{d - 1}}{w}^{i}\left( {\frac{1}{d - 1}\left( {1 - \varepsilon }\right)  + {\lambda }_{i}\varepsilon }\right)
\]

\[
= {w}^{0}{\lambda }_{0}\varepsilon  + \mathop{\sum }\limits_{{i = 1}}^{{d - 1}}{w}^{i}\left( {\frac{1}{d - 1} + \left( {{\lambda }_{i} - \frac{1}{d - 1}}\right) \varepsilon }\right) . \tag{7.10}
\]

If \({\lambda }_{0} > 0\) , this would imply \(y\left( \varepsilon \right)  \in  \left( {T \cap  {T}^{\prime }}\right)  - F\) for sufficiently small positive \(\varepsilon\) , which is not possible because \(T\) and \({T}^{\prime }\) intersect in a common facet which can only be \(F\) . Hence, \({\lambda }_{0} < 0\) as claimed.

\begin{tcolorbox}
如果 \({\lambda }_{0} > 0\) ，对于足够小的正 \(\varepsilon\) ，这将意味着 \(y\left( \varepsilon \right)  \in  \left( {T \cap  {T}^{\prime }}\right)  - F\) ，但这是不可能的，因为 \(T\) 和 \({T}^{\prime }\) 在一个公共极面相交，而这个公共极面只能是 \(F\) 。因此，如所断言的 \({\lambda }_{0} < 0\) 。
\end{tcolorbox}

The two cases (a) and (b) are shown as follows. If \(F\) is a subset of a facet of \(S\) , then \(y\left( \varepsilon \right)  \notin  S\) for all \(\varepsilon  < 0\) : Otherwise, for \(\varepsilon  < 0\) so that \(y\left( \varepsilon \right)\) and \(y\left( {-\varepsilon }\right)\) are in \(S\) , we have \(x = y\left( \varepsilon \right) \frac{1}{2} + y\left( {-\varepsilon }\right) \frac{1}{2} = \mathop{\sum }\limits_{{i = 1}}{v}^{i}{\mu }_{i}\) where \({\mu }_{i} > 0\) for all \(i = 1,\ldots ,d\) , that is, \(x\) is in the interior of \(S\) instead of on a facet of \(S\) . This implies that there is no second sub-simplex \({T}^{\prime }\) with facet \(F\) as considered above, because \({T}^{\prime }\) (and hence \(S\) ) would contain such \(y\left( \varepsilon \right)\) with \(\varepsilon  < 0\) and \({\lambda }_{0} < 0\) by (7.10). So the sub-simplex \(T\) is unique as stated in (a).

\begin{tcolorbox}
以下展示了情况 (a) 和 (b)。如果 \(F\) 是 \(S\) 的一个极面的子集，那么对于所有 \(\varepsilon  < 0\) 有 \(y\left( \varepsilon \right)  \notin  S\) :否则，对于 \(\varepsilon  < 0\) 使得 \(y\left( \varepsilon \right)\) 和 \(y\left( {-\varepsilon }\right)\) 在 \(S\) 中，我们有 \(x = y\left( \varepsilon \right) \frac{1}{2} + y\left( {-\varepsilon }\right) \frac{1}{2} = \mathop{\sum }\limits_{{i = 1}}{v}^{i}{\mu }_{i}\) ，其中对于所有 \(i = 1,\ldots ,d\) 有 \({\mu }_{i} > 0\) ，即 \(x\) 在 \(S\) 的内部而不是在 \(S\) 的一个极面上。这意味着不存在如上述所考虑的以 \(F\) 为极面的第二个子单纯形 \({T}^{\prime }\) ，因为 \({T}^{\prime }\) (进而 \(S\) )会包含这样的 \(y\left( \varepsilon \right)\) ，其中有 \(\varepsilon  < 0\) 和 \({\lambda }_{0} < 0\) (根据 (7.10))。所以如 (a) 中所述，子单纯形 \(T\) 是唯一的。
\end{tcolorbox}

On the other hand, if \(F\) is not a subset of a facet of \(S\) , then \(x\) is in the interior of \(S\) , and there are \(y\left( \varepsilon \right)\) in \(S\) for \(\varepsilon  < 0\) and therefore there is a second sub-simplex \({T}^{\prime }\) with facet \(F\) . By the same reasoning as applied to \(T\) , there is no other sub-simplex with a vertex on the same side of \(F\) as the vertex \({w}^{0}\) of \({T}^{\prime }\) . This shows case (b), and obviously exactly one of the two cases (a) or (b) applies.

\begin{tcolorbox}
另一方面，如果 \(F\) 不是 \(S\) 的一个极面的子集，那么 \(x\) 在 \(S\) 的内部，并且对于 \(\varepsilon  < 0\) 存在 \(y\left( \varepsilon \right)\) 在 \(S\) 中，因此存在以 \(F\) 为极面的第二个子单纯形 \({T}^{\prime }\) 。通过对 \(T\) 应用的相同推理，不存在其他子单纯形，其顶点与 \({T}^{\prime }\) 的顶点 \({w}^{0}\) 在 \(F\) 的同一侧。这展示了情况 (b)，并且显然情况 (a) 或 (b) 中恰好有一种情况成立。
\end{tcolorbox}

The following lemma will be used in an inductive proof in the next section.

\begin{tcolorbox}
以下引理将用于下一节的归纳证明。
\end{tcolorbox}

Lemma 7.7. Let \(n \geq  2\) , let \(S\) be a simplex with vertices \({v}^{1},\ldots ,{v}^{n}\) , let \(F\) be the facet of \(S\) with vertices \({v}^{1},\ldots ,{v}^{n - 1}\) , let \(\sum\) be a triangulation of \(S\) , and let \({\sum }^{\prime }\) be the set of sub-facets in \(\sum\) that are subsets of \(F\) . Then \({\sum }^{\prime }\) is a triangulation of \(F\) .

\begin{tcolorbox}
引理7.7。设 \(n \geq  2\) ，设 \(S\) 为一个顶点为 \({v}^{1},\ldots ,{v}^{n}\) 的单纯形，设 \(F\) 为 \(S\) 中顶点为 \({v}^{1},\ldots ,{v}^{n - 1}\) 的极面，设 \(\sum\) 为 \(S\) 的一个三角剖分，设 \({\sum }^{\prime }\) 为 \(\sum\) 中作为 \(F\) 子集的子极面的集合。那么 \({\sum }^{\prime }\) 是 \(F\) 的一个三角剖分。
\end{tcolorbox}

Proof. Let \({\sum }_{1} = \{ T \cap  F \mid  T \in  \sum \}\) . Then \(F = \mathop{\bigcup }\limits_{{{T}^{\prime } \in  {\sum }_{1}}}{T}^{\prime }\) because every \(x\) in \(F\) (as an element of \(S\) ) belongs to some sub-simplex \(T\) in \(\sum\) . For every \(T \in  \sum\) , we claim that \(T \cap  F\) is a face of \(T\) . Namely, every vertex \(v\) of \(T\) is a unique convex combination \(v = \mathop{\sum }\limits_{{i = 1}}^{n}{v}^{i}{\lambda }_{i}\) of the vertices \({v}^{i}\) of \(S\) , where \(v \in  F\) if and only if \({\lambda }_{n} = 0\) . Then \(T \cap  F\) is the convex hull of these vertices \(v\) that belong to \(F\) . Therefore, \(T \cap  F\) is a face of \(T\) as claimed, and is a simplex of dimension at most \(n - 2\) .

\begin{tcolorbox}
证明。设 \({\sum }_{1} = \{ T \cap  F \mid  T \in  \sum \}\) 。那么 \(F = \mathop{\bigcup }\limits_{{{T}^{\prime } \in  {\sum }_{1}}}{T}^{\prime }\) ，因为 \(F\) 中的每个 \(x\) (作为 \(S\) 的一个元素)都属于 \(\sum\) 中的某个子单纯形 \(T\) 。对于每个 \(T \in  \sum\) ，我们断言 \(T \cap  F\) 是 \(T\) 的一个面。即， \(T\) 的每个顶点 \(v\) 都是 \(S\) 的顶点 \({v}^{i}\) 的唯一凸组合 \(v = \mathop{\sum }\limits_{{i = 1}}^{n}{v}^{i}{\lambda }_{i}\) ，其中当且仅当 \({\lambda }_{n} = 0\) 时 \(v \in  F\) 。那么 \(T \cap  F\) 是这些属于 \(F\) 的顶点 \(v\) 的凸包。因此，如所断言， \(T \cap  F\) 是 \(T\) 的一个面，并且是一个维度至多为 \(n - 2\) 的单纯形。
\end{tcolorbox}

Let \({\sum }_{2}\) be the set of those simplices in \({\sum }_{1}\) that have dimension \(n - 2\) . Because all simplices in \(\sum\) have dimension \(n - 1\) by Lemma 7.5, \({\sum }_{2} = {\sum }^{\prime }\) .

\begin{tcolorbox}
设 \({\sum }_{2}\) 为 \({\sum }_{1}\) 中维度为 \(n - 2\) 的那些单纯形的集合。根据引理7.5，由于 \(\sum\) 中的所有单纯形的维度都是 \(n - 1\) ，所以 \({\sum }_{2} = {\sum }^{\prime }\) 。
\end{tcolorbox}

We show that every \(x \in  F\) belongs to some simplex \({T}^{\prime }\) in \({\sum }^{\prime }\) , that is, in \({\sum }_{2}\) . Suppose there is some \(x \in  F\) where that is not the case. Then \(x \in  {S}_{1}\) for some simplex \({S}_{1} \in  {\sum }_{1}\) where \({S}_{1}\) has dimension less than \(n - 2\) . Let \({S}_{1}\) be the simplex in \({\sum }_{1}\) of largest dimension (by assumption at most \(n - 3\) ) that contains \(x\) . By the same argument as used in the proof of Lemma 7.5 with \(d = n - 1\) , there is some \({S}_{2} \in  {\sum }_{1}\) which is a superset of \({S}_{1}\) and has at least one further vertex and higher dimension than \({S}_{1}\) , a contradiction. This shows that \(F = \mathop{\bigcup }\limits_{{{T}^{\prime } \in  {\sum }^{\prime }}}{T}^{\prime }\) as claimed.

\begin{tcolorbox}
我们证明每个 \(x \in  F\) 都属于 \({\sum }^{\prime }\) 中的某个单纯形 \({T}^{\prime }\) ，即属于 \({\sum }_{2}\) 。假设存在某个 \(x \in  F\) 不满足该情况。那么对于某个单纯形 \({S}_{1} \in  {\sum }_{1}\) 有 \(x \in  {S}_{1}\) ，其中 \({S}_{1}\) 的维度小于 \(n - 2\) 。设 \({S}_{1}\) 为 \({\sum }_{1}\) 中包含 \(x\) 的最大维度(根据假设至多为 \(n - 3\) )的单纯形。通过与引理7.5证明中使用 \(d = n - 1\) 时相同的论证，存在某个 \({S}_{2} \in  {\sum }_{1}\) 是 \({S}_{1}\) 的超集，且比 \({S}_{1}\) 至少多一个顶点且维度更高，这产生矛盾。这表明如所断言 \(F = \mathop{\bigcup }\limits_{{{T}^{\prime } \in  {\sum }^{\prime }}}{T}^{\prime }\) 。
\end{tcolorbox}

For no two sub-facets in \({\sum }^{\prime }\) is one a subset of the other because they have the same dimension but different sets of vertices. Moreover, the intersection of two such sub-facets is determined by the intersection of their vertices, and hence a face of both these sub-facets. This shows that \({\sum }^{\prime }\) is a triangulation of \(F\) .

\begin{tcolorbox}
在 \({\sum }^{\prime }\) 中没有两个子极面使得一个是另一个的子集，因为它们具有相同的维度但顶点集不同。此外，两个这样的子极面的交集由它们顶点的交集确定，因此是这两个子极面的一个面。这表明 \({\sum }^{\prime }\) 是 \(F\) 的一个三角剖分。
\end{tcolorbox}

\begin{center}
\includegraphics[max width=0.4\textwidth]{images/bo_d2e5cmbef24c73b03tn0_190_197_1414_614_542_0.jpg}
\end{center}
\hspace*{3em} 

Figure 7.5 Triangulation of a regular triangle into small similar triangles.

\begin{tcolorbox}
图7.5 正三角形划分为小的相似三角形的三角剖分。
\end{tcolorbox}

The triangulation of a simplex \(S\) takes the role of the "grid" mentioned earlier, with its vertices as gridpoints. The goal is to be able to create for any \(\varepsilon  > 0\) an " \(\varepsilon\) -fine grid" where the vertices in every sub-simplex are no more than \(\varepsilon\) apart. Figure 7.5 shows a good way to do this when \(S\) is a triangle. However, triangulating a higher-dimensional simplex to achieve this is not as easy. For example, consider a regular tetrahedron where each of its facets, which are triangles, is subdivided into four regular triangles obtained by joining the midpoints of each edge of the tetrahedron. Then "chopping off" the small tetrahedron at each corner leaves a central polytope that is not a tetrahedron but an octahedron. This octahedron has to be subdivided further into tetrahedra in some new way. In addition, it is not clear how such a construction would extend to higher dimensions. For the moment, we assume that in any dimension we can find a sufficiently "fine" triangulation of a simplex. A specific construction that achieves this is described in Section 7.7.

\begin{tcolorbox}
单纯形 \(S\) 的三角剖分起到了前面提到的“网格”的作用，其顶点即为网格点。目标是能够为任意 \(\varepsilon  > 0\) 创建一个“ \(\varepsilon\) -精细网格”，使得每个子单纯形中的顶点之间的距离不超过 \(\varepsilon\) 。图7.5展示了当 \(S\) 为三角形时实现这一目标的一种好方法。然而，对高维单纯形进行三角剖分以达到这一要求并非易事。例如，考虑一个正四面体，它的每个极面(即三角形)被细分为四个正三角形，这些正三角形是通过连接四面体每条边的中点得到的。然后，“切掉”每个角上的小四面体后，会留下一个中心有界多面体，它不是四面体，而是八面体。这个八面体必须以某种新的方式进一步细分为四面体。此外，目前尚不清楚这种构造如何扩展到更高维度。目前，我们假设在任何维度下，我们都能找到单纯形的足够“精细”的三角剖分。实现这一目标的具体构造方法将在7.7节中介绍。
\end{tcolorbox}

\section*{7.4 Sperner's Lemma}

\begin{tcolorbox}
\section*{7.4 斯佩纳引理}
\end{tcolorbox}

In this section, we prove Sperner's lemma, which is often called a combinatorial version of Brouwer's fixed-point theorem, and which is a crucial step in its proof. Sperner’s lemma applies to a simplex \(S\) with \(d\) vertices and triangulation \(\sum\) of \(S\) , where each vertex in \(\sum\) is given a label in \(\{ 1,\ldots ,d\}\) , which fulfills the following condition.

\begin{tcolorbox}
在本节中，我们将证明斯佩纳引理，它常被称为布劳威尔不动点定理的组合版本，并且是证明该定理的关键一步。斯佩纳引理适用于具有 \(d\) 个顶点的单纯形 \(S\) 以及 \(S\) 的三角剖分 \(\sum\) ，其中 \(\sum\) 中的每个顶点都被赋予了 \(\{ 1,\ldots ,d\}\) 中的一个标签，且满足以下条件。
\end{tcolorbox}

Definition 7.8 (Sperner labeling). Let \(S\) be a simplex with vertices \({v}^{1},\ldots ,{v}^{d}\) , let \(\sum\) be a triangulation of \(S\) , and let \(V\) be the set of vertices in \(\sum\) . Then a function \(V \rightarrow  \{ 1,\ldots ,d\}\) that assigns to each \(x\) in \(V\) a label is said to fulfill the Sperner condition, and is called a Sperner labeling, if whenever \(x = \mathop{\sum }\limits_{{k = 1}}^{d}{v}^{k}{\lambda }_{k}\) and \(x\) has label \(i\) , then \({\lambda }_{i} > 0\) .

\begin{tcolorbox}
定义7.8(斯佩纳标号)。设 \(S\) 是一个顶点为 \({v}^{1},\ldots ,{v}^{d}\) 的单纯形， \(\sum\) 是 \(S\) 的三角剖分， \(V\) 是 \(\sum\) 中的顶点集。那么，一个为 \(V\) 中的每个 \(x\) 分配标签的函数 \(V \rightarrow  \{ 1,\ldots ,d\}\) ，若满足每当 \(x = \mathop{\sum }\limits_{{k = 1}}^{d}{v}^{k}{\lambda }_{k}\) 且 \(x\) 的标签为 \(i\) 时，有 \({\lambda }_{i} > 0\) ，则称该函数满足斯佩纳条件，并称其为斯佩纳标号。
\end{tcolorbox}

\begin{center}
\includegraphics[max width=0.5\textwidth]{images/bo_d2e5cmbef24c73b03tn0_191_334_1446_746_724_0.jpg}
\end{center}
\hspace*{3em} 

Figure 7.6 A triangulation of the triangle \(S\) with a Sperner labeling. The vertices \({v}^{1},{v}^{2},{v}^{3}\) of \(S\) are uniquely identified by their respective labels 1, 2, 3.

\begin{tcolorbox}
图7.6 带有斯佩纳标号的三角形 \(S\) 的三角剖分。 \(S\) 的顶点 \({v}^{1},{v}^{2},{v}^{3}\) 通过它们各自的标签1、2、3唯一标识。
\end{tcolorbox}

Note that in Definition 7.8, every point \(x\) in the simplex \(S\) has the unique barycentric coordinates \({\lambda }_{1},\ldots ,{\lambda }_{d}\) as in (7.7). The Sperner condition implies that each vertex \({v}^{i}\) of \(S\) has label \(i\) (and therefore every vertex has a different label). In general, the Sperner condition states that if \(x\) is on a face \(F\) of \(S\) (and no smaller face), then \(x\) is only allowed to have one of the labels of the vertices of \(F\) . For example, if \(x\) is on the edge that connects \({v}^{1}\) and \({v}^{2}\) , then \(x\) can only have labels 1 or 2 . For \(d = 2\) the Sperner condition is stated in (7.3). For \(d = 3\) the simplex \(S\) is a triangle, with an example of a triangulation with a Sperner labeling shown in Figure 7.6.

\begin{tcolorbox}
请注意，在定义7.8中，单纯形 \(S\) 中的每个点 \(x\) 都具有如(7.7)所示的唯一重心坐标 \({\lambda }_{1},\ldots ,{\lambda }_{d}\) 。斯佩纳条件意味着 \(S\) 的每个顶点 \({v}^{i}\) 都有标签 \(i\) (因此每个顶点都有不同的标签)。一般来说，斯佩纳条件规定，如果 \(x\) 位于 \(S\) 的一个面 \(F\) 上(而不是更小的面上)，那么 \(x\) 只能具有 \(F\) 的顶点的标签之一。例如，如果 \(x\) 位于连接 \({v}^{1}\) 和 \({v}^{2}\) 的边上，那么 \(x\) 只能有标签1或2。对于 \(d = 2\) ，斯佩纳条件在(7.3)中给出。对于 \(d = 3\) ，单纯形 \(S\) 是一个三角形，图7.6展示了一个带有斯佩纳标号的三角剖分示例。
\end{tcolorbox}

For a function \(f : \Delta  \rightarrow  \Delta\) on the unit simplex \(\Delta\) in (7.1) and \(x \in  \Delta\) , one of the labels of \(x\) defined in (7.2) can always be chosen to fulfill the Sperner condition.

\begin{tcolorbox}
对于 (7.1) 中单位单纯形 \(f : \Delta  \rightarrow  \Delta\) 上的函数 \(\Delta\) 以及 \(x \in  \Delta\) ，(7.2) 中定义的 \(x\) 的标签之一总能被选来满足斯佩纳条件。
\end{tcolorbox}

Lemma 7.9. Consider the(d - 1)-dimensional unit simplex \(\Delta\) in (7.1) and a function \(f : \Delta  \rightarrow  \Delta\) . Then every point \(x \in  \Delta\) can be given a label \(i\) so that \({x}_{i} \geq  {f}_{i}\left( x\right)\) and the Sperner condition holds.

\begin{tcolorbox}
引理 7.9。考虑 (7.1) 中 (d - 1) 维单位单纯形 \(\Delta\) 以及函数 \(f : \Delta  \rightarrow  \Delta\) 。那么每个点 \(x \in  \Delta\) 都可以被赋予一个标签 \(i\) ，使得 \({x}_{i} \geq  {f}_{i}\left( x\right)\) 且斯佩纳条件成立。
\end{tcolorbox}

Proof. Because \(\Delta\) is the convex hull of the unit vectors \({e}_{1},\ldots ,{e}_{d}\) , the barycentric coordinates of \(x \in  \Delta\) are just \({x}_{1},\ldots ,{x}_{d}\) . Let \(I = \left\{  {i \mid  {x}_{i} > 0}\right\}\) . If \({f}_{i}\left( x\right)  > {x}_{i}\) for all \(i \in  I\) , then

\begin{tcolorbox}
证明。因为 \(\Delta\) 是单位向量 \({e}_{1},\ldots ,{e}_{d}\) 的凸包， \(x \in  \Delta\) 的重心坐标就是 \({x}_{1},\ldots ,{x}_{d}\) 。设 \(I = \left\{  {i \mid  {x}_{i} > 0}\right\}\) 。如果对于所有 \(i \in  I\) 都有 \({f}_{i}\left( x\right)  > {x}_{i}\) ，那么
\end{tcolorbox}

\[
\mathop{\sum }\limits_{{i = 1}}^{d}{f}_{i}\left( x\right)  \geq  \mathop{\sum }\limits_{{i \in  I}}{f}_{i}\left( x\right)  > \mathop{\sum }\limits_{{i \in  I}}{x}_{i} = \mathop{\sum }\limits_{{i = 1}}^{d}{x}_{i} = 1,
\]

which contradicts \(f\left( x\right)  \in  \Delta\) . So there is at least one \(i \in  I\) with \({x}_{i} \geq  {f}_{i}\left( x\right)\) , that is, \({x}_{i} > 0\) and \(x\) has label \(i\) , as required.

\begin{tcolorbox}
这与 \(f\left( x\right)  \in  \Delta\) 矛盾。所以至少存在一个 \(i \in  I\) 使得 \({x}_{i} \geq  {f}_{i}\left( x\right)\) ，即 \({x}_{i} > 0\) 且 \(x\) 具有标签 \(i\) ，符合要求。
\end{tcolorbox}

Sperner's lemma asserts that a triangulation of a simplex with a Sperner labeling has at least one completely labeled sub-simplex. A suitable induction proof shows the stronger statement that the number of completely labeled sub-simplices is odd. Figure 7.7 shows this for the subdivided triangle considered earlier.

\begin{tcolorbox}
斯佩纳引理断言，带有斯佩纳标号的单纯形的三角剖分至少有一个完全标记的子单纯形。一个合适的归纳证明可以得出更强的结论，即完全标记的子单纯形的数量是奇数。图 7.7 展示了前面所考虑的细分三角形的这种情况。
\end{tcolorbox}

Theorem 7.10 (Sperner’s lemma). Let \(S\) be a simplex with vertices \({v}^{1},\ldots ,{v}^{d}\) , let \(\sum\) be a triangulation of \(S\) , and consider a Sperner labeling of the vertices in \(V\) . Then \(\sum\) has an odd number of sub-simplices \(T\) that are completely labeled (that is, every label in \(\{ 1,\ldots ,d\}\) is the label of some vertex of \(T\) ).

\begin{tcolorbox}
定理 7.10(斯佩纳引理)。设 \(S\) 是一个顶点为 \({v}^{1},\ldots ,{v}^{d}\) 的单纯形，设 \(\sum\) 是 \(S\) 的一个三角剖分，并考虑 \(V\) 中顶点的斯佩纳标号。那么 \(\sum\) 有奇数个完全标记的子单纯形 \(T\) (即， \(\{ 1,\ldots ,d\}\) 中的每个标签都是 \(T\) 的某个顶点的标签)。
\end{tcolorbox}

Proof. In the statement of the theorem, the vertices \({v}^{1},\ldots ,{v}^{d}\) of \(S\) can be elements of some space \({\mathbb{R}}^{k}\) for any \(k \geq  d - 1\) . Consider the following statement:

\begin{tcolorbox}
证明。在定理的陈述中， \(S\) 的顶点 \({v}^{1},\ldots ,{v}^{d}\) 可以是任意 \(k \geq  d - 1\) 时某个空间 \({\mathbb{R}}^{k}\) 的元素。考虑以下陈述:
\end{tcolorbox}

(i) Let \(1 \leq  n \leq  d\) , let \({F}_{n}\) be the simplex with vertices \({v}^{1},\ldots ,{v}^{n}\) , and let \({\sum }_{n} = \{ F \mid  F\) is a face of dimension \(n - 1\) of some \(\left. {T \in  \sum \text{ and }F \subseteq  {F}_{n}}\right\}\) . Then \({\sum }_{n}\) is a triangulation of \({F}_{n}\) .

\begin{tcolorbox}
(i) 设 \(1 \leq  n \leq  d\) ，设 \({F}_{n}\) 是顶点为 \({v}^{1},\ldots ,{v}^{n}\) 的单纯形，且设 \({\sum }_{n} = \{ F \mid  F\) 是某个 \(\left. {T \in  \sum \text{ and }F \subseteq  {F}_{n}}\right\}\) 的 \(n - 1\) 维面。那么 \({\sum }_{n}\) 是 \({F}_{n}\) 的一个三角剖分。
\end{tcolorbox}

By assumption,(i) is true for \(n = d\) . If (i) holds for some \(n > 1\) , then Lemma 7.7 (with \(d = n\) ) implies that \({\sum }_{n - 1}\) is a triangulation of \({F}_{n - 1}\) . Hence, by induction,(i) is true for \(n = d,d - 1,\ldots ,1\) , as claimed. We now prove by induction:

\begin{tcolorbox}
根据假设，(i) 对于 \(n = d\) 成立。如果 (i) 对于某个 \(n > 1\) 成立，那么引理 7.7(对于 \(d = n\) )意味着 \({\sum }_{n - 1}\) 是 \({F}_{n - 1}\) 的一个三角剖分。因此，通过归纳法，(i) 对于 \(n = d,d - 1,\ldots ,1\) 成立，如所声称的。我们现在通过归纳法证明:
\end{tcolorbox}

\begin{center}
\includegraphics[max width=0.5\textwidth]{images/bo_d2e5cmbef24c73b03tn0_193_335_202_744_721_0.jpg}
\end{center}
\hspace*{3em} 

Figure 7.7 Illustration of Sperner's lemma (Theorem 7.10). The black (full) dots mark completely labeled sub-simplices, and the white (hollow) dots mark completely labeled sub-facets, as used in the induction proof of Theorem 7.10. Both white and black dots are odd in number.

\begin{tcolorbox}
图 7.7 斯佩纳引理(定理 7.10)的图示。黑色(实心)点标记完全标记的子单纯形，白色(空心)点标记完全标记的子极面，如定理 7.10 的归纳证明中所使用的。白色和黑色点的数量都是奇数。
\end{tcolorbox}

(ii) For \(n = 1,2,\ldots ,d\) , let \({\sum }_{n}\) be the triangulation of \({F}_{n}\) as defined in (i). Call \(T \in  {\sum }_{n}\) completely labeled if the vertices of \(T\) have all labels \(1,\ldots ,n\) . Then \({\sum }_{n}\) has an odd number of completely labeled simplices \(T\) .

\begin{tcolorbox}
(ii) 对于 \(n = 1,2,\ldots ,d\) ，设 \({\sum }_{n}\) 为 (i) 中所定义的 \({F}_{n}\) 的三角剖分。若 \(T\) 的顶点具有所有标签 \(1,\ldots ,n\) ，则称 \(T \in  {\sum }_{n}\) 为完全标记的。那么 \({\sum }_{n}\) 有奇数个完全标记的单纯形 \(T\) 。
\end{tcolorbox}

Statement (ii) is trivial for \(n = 1\) where \({F}_{1} = \left\{  {v}_{1}\right\}\) . To prove the induction step, assume (ii) holds for some \(n \geq  1\) where \(n < d\) , and we want to prove (ii) for \(n + 1\) . For any face \(F\) of a simplex \(T\) in \(\sum\) , let

\begin{tcolorbox}
对于 \(n = 1\) (其中 \({F}_{1} = \left\{  {v}_{1}\right\}\) )，命题 (ii) 是显然的。为了证明归纳步骤，假设 (ii) 对于某个 \(n \geq  1\) (其中 \(n < d\) )成立，我们要证明 (ii) 对于 \(n + 1\) 也成立。对于 \(\sum\) 中单纯形 \(T\) 的任意面 \(F\) ，设
\end{tcolorbox}

\[
L\left( F\right)  = \{ i \mid  i\text{ is a label of some vertex of }F\} .
\]

Consider the following sets:

\begin{tcolorbox}
考虑以下集合:
\end{tcolorbox}

\[
C = \left\{  {F \in  {\sum }_{n} \mid  L\left( F\right)  = \{ 1,\ldots ,n\} }\right\}  ,
\]

\[
{C}^{\prime } = \left\{  {F \mid  F\text{ is a facet of some }T \in  {\sum }_{n + 1},L\left( F\right)  = \{ 1,\ldots ,n\} ,F \notin  {\sum }_{n}}\right\}  , \tag{7.11}
\]

\[
D = \left\{  {T \in  {\sum }_{n + 1} \mid  L\left( T\right)  = \{ 1,\ldots ,n + 1\} }\right\}  ,
\]

\[
A = \left\{  {T \in  {\sum }_{n + 1} \mid  L\left( T\right)  = \{ 1,\ldots ,n\} }\right\}  .
\]

The case \(n = 2\) is shown in Figures 7.7 and 7.8. The set \({F}_{3} = {F}_{n + 1}\) is the triangle considered earlier, with a triangulation \({\sum }_{3}\) of \({F}_{3}\) into triangles; the vertices in \({\sum }_{3}\) may have labels 1,2, and 3 . The completely labeled elements of \({\sum }_{3}\) define the set \(D\) and are marked with black dots. The set \({F}_{2} = {F}_{n}\) is the bottom facet of the triangle \({F}_{3}\) , with a triangulation \({\sum }_{2}\) of \({F}_{2}\) into line segments; the vertices in \({\sum }_{2}\) may have labels 1 and 2 . The completely labeled elements of \({\sum }_{2}\) define the set \(C\) and are marked with white dots. By inductive assumption, \(\left| C\right|\) is odd.

\begin{tcolorbox}
情况 \(n = 2\) 如图 7.7 和 7.8 所示。集合 \({F}_{3} = {F}_{n + 1}\) 是前面考虑的三角形，它有一个将 \({F}_{3}\) 三角剖分成三角形的三角剖分 \({\sum }_{3}\) ； \({\sum }_{3}\) 中的顶点可能具有标签 1、2 和 3。 \({\sum }_{3}\) 中完全标记的元素定义了集合 \(D\) ，并用黑点标记。集合 \({F}_{2} = {F}_{n}\) 是三角形 \({F}_{3}\) 的底部极面，它有一个将 \({F}_{2}\) 三角剖分成线段的三角剖分 \({\sum }_{2}\) ； \({\sum }_{2}\) 中的顶点可能具有标签 1 和 2。 \({\sum }_{2}\) 中完全标记的元素定义了集合 \(C\) ，并用白点标记。根据归纳假设， \(\left| C\right|\) 是奇数。
\end{tcolorbox}

The set \({C}^{\prime }\) is the set of sub-facets \(F\) whose \(n\) vertices have all labels \(1,\ldots ,n\) and which are not in \({\sum }_{n}\) ; in Figure 7.8 they are also marked with white dots. (We use the terminology when \(n + 1 = d\) and call the elements of \({\sum }_{n + 1}\) "sub-simplices" and any facet of a sub-simplex a "sub-facet", which includes the elements of \({\sum }_{n}\) .) The set \(A\) is the set of almost completely labeled sub-simplices where for any \(T \in  A\) the \(n + 1\) vertices of \(T\) have all labels \(1,\ldots ,n\) but not label \(n + 1\) .

\begin{tcolorbox}
集合 \({C}^{\prime }\) 是子极面 \(F\) 的集合，其 \(n\) 个顶点具有所有标签 \(1,\ldots ,n\) 且不在 \({\sum }_{n}\) 中；在图 7.8 中，它们也用白点标记。(当 \(n + 1 = d\) 时，我们使用这样的术语，称 \({\sum }_{n + 1}\) 的元素为“子单纯形”，子单纯形的任何极面为“子极面”，这包括 \({\sum }_{n}\) 的元素。)集合 \(A\) 是几乎完全标记的子单纯形的集合，对于任意 \(T \in  A\) ， \(T\) 的 \(n + 1\) 个顶点具有所有标签 \(1,\ldots ,n\) 但不具有标签 \(n + 1\) 。
\end{tcolorbox}

\begin{center}
\includegraphics[max width=0.5\textwidth]{images/bo_d2e5cmbef24c73b03tn0_194_199_200_741_726_0.jpg}
\end{center}
\hspace*{3em} 

Figure 7.8 The sets defined in (7.11) for \(n = 2\) . The elements of \(C\) (white dots) are in \({\sum }_{2}\) with labels 1,2, and the elements of \(D\) (black dots) are sub-simplices in \({\sum }_{3}\) with labels1,2,3. They are connected by paths with further nodes in \({C}^{\prime }\) (white dots, sub-facets inside the triangle with labels 1, 2), which define the graph \({G}_{3}\) . Hence, \(\left| {C \cup  D}\right|\) is even, and thus \(\left| D\right|\) is odd because \(\left| C\right|\) is odd.

\begin{tcolorbox}
图 7.8 为 \(n = 2\) 定义的集合 (7.11)。 \(C\) 的元素(白点)在 \({\sum }_{2}\) 中，标签为 1、2， \(D\) 的元素(黑点)是 \({\sum }_{3}\) 中的子单纯形，标签为 1、2、3。它们通过路径相连，路径上还有 \({C}^{\prime }\) 中的其他节点(白点，三角形内标签为 1、2 的子极面)，这些定义了图 \({G}_{3}\) 。因此， \(\left| {C \cup  D}\right|\) 是偶数，由于 \(\left| C\right|\) 是奇数，所以 \(\left| D\right|\) 是奇数。
\end{tcolorbox}

We now define a new graph \({G}_{n + 1}\) with node set \(C \cup  {C}^{\prime } \cup  D\) , where any two nodes are joined by an edge if they are subsets of the same sub-simplex, which belongs to \(A\) or \(D\) . Figure 7.8 shows the nodes of the graph \({G}_{3}\) as white or black dots, and its edges as thick gray lines. The following cases can occur for such an edge of \({G}_{n + 1}\) :

\begin{tcolorbox}
我们现在定义一个新的图 \({G}_{n + 1}\) ，其节点集为 \(C \cup  {C}^{\prime } \cup  D\) ，如果任意两个节点是同一个子单纯形的子集，且该子单纯形属于 \(A\) 或 \(D\) ，则这两个节点由一条边相连。图 7.8 中，图 \({G}_{3}\) 的节点用白点或黑点表示，其边用粗灰线表示。对于 \({G}_{n + 1}\) 的这样一条边，可能出现以下情况:
\end{tcolorbox}

Case A. One endpoint of the edge is a sub-facet \(F \in  C\) , so that \(F\) is a subset of the facet \({F}_{n}\) of \({F}_{n + 1}\) . Hence, by Lemma 7.6(a), \(F\) is the face of exactly one sub-simplex \(T\) in \(D \cup  A\) . Then:

\begin{tcolorbox}
情况A。该边的一个端点是子极面 \(F \in  C\) ，使得 \(F\) 是 \({F}_{n + 1}\) 的极面 \({F}_{n}\) 的子集。因此，根据引理7.6(a)， \(F\) 是 \(D \cup  A\) 中恰好一个子单纯形 \(T\) 的面。那么:
\end{tcolorbox}

(a) If \(T \in  D\) , then \(T\) is completely labeled, and has exactly one facet with labels \(1,\ldots ,n\) , which is \(F\) . The edge of the graph is \(\{ F,T\}\) , and both \(F\) and \(T\) have degree one. The edge itself defines a path of length one in the graph. An example of such a single-edge path is shown in Figure 7.8 at the bottom in the fourth triangle from the left.

\begin{tcolorbox}
(a) 如果 \(T \in  D\) ，则 \(T\) 是完全标记的，并且恰好有一个带有标签 \(1,\ldots ,n\) 的极面，即 \(F\) 。该图的边是 \(\{ F,T\}\) ，并且 \(F\) 和 \(T\) 的度均为1。该边本身在图中定义了一条长度为1的路径。图7.8中左侧起第四个三角形底部展示了这样一条单边路径的示例。
\end{tcolorbox}

(b) If \(T \in  A\) , then \(T\) has \(n + 1\) vertices that together have labels \(1,\ldots ,n\) , so one of these labels occurs twice. Hence, \(T\) has exactly two sub-facets that have labels \(1,\ldots ,n\) , one of which is \(F\) and the second, say \({F}^{\prime }\) , belongs to \({C}^{\prime }\) . The edge of the graph is \(\left\{  {F,{F}^{\prime }}\right\}\) , with \({F}^{\prime }\) as the only neighbor of \(F\) , so \(F\) again has degree one. In Figure 7.8, these are all the other cases where \(F \in  {\sum }_{2}\) (that is, \(F\) is a white dot at the bottom of the triangle), where \(T\) (which represents in effect the edge of the graph) has either labels1,2,1or1,2,2.

\begin{tcolorbox}
(b) 如果 \(T \in  A\) ，则 \(T\) 有 \(n + 1\) 个顶点，这些顶点一起具有标签 \(1,\ldots ,n\) ，因此这些标签中有一个出现了两次。因此， \(T\) 恰好有两个具有标签 \(1,\ldots ,n\) 的子极面，其中一个是 \(F\) ，另一个，设为 \({F}^{\prime }\) ，属于 \({C}^{\prime }\) 。该图的边是 \(\left\{  {F,{F}^{\prime }}\right\}\) ， \({F}^{\prime }\) 是 \(F\) 的唯一邻居，因此 \(F\) 的度再次为1。在图7.8中，这些都是 \(F \in  {\sum }_{2}\) 的所有其他情况(即 \(F\) 是三角形底部的一个白点)，其中 \(T\) (实际上代表图的边)的标签为1,2,1或1,2,2。
\end{tcolorbox}

Case B. One endpoint of the edge of the graph is a sub-facet \(F \in  {C}^{\prime }\) . Then \(F \notin  {\sum }_{n}\) by definition of \({C}^{\prime }\) , and \(F\) is therefore not a subset of any facet of \({F}_{n + 1}\) because of the Sperner condition: the vertices of \(F\) have the labels \(1,\ldots ,n\) , and every facet of \({F}_{n + 1}\) other than \({F}_{n}\) is missing one of these labels. Hence, by Lemma 7.6(b), such a sub-facet \(F\) in \({C}^{\prime }\) is the facet of exactly two sub-simplices in \(D \cup  A\) . In the graph, \(F\) has degree two. For each of the two sub-simplices \(T \in  D \cup  A\) that have \(F\) as a facet, either (a) or (b) applies as before: (a) if \(T \in  D\) , then \(F\) is the unique neighbor of \(T\) , the edge of the graph is \(\{ F,T\}\) , and \(T\) has degree one in the graph; (b) if \(T \in  A\) , then \(A\) has exactly one other facet \({F}^{\prime }\) with labels \(1,\ldots ,n\) , where \({F}^{\prime } \in  C \cup  {C}^{\prime }\) , and the edge of the graph is \(\left\{  {F,{F}^{\prime }}\right\}\) .

\begin{tcolorbox}
情况B。该图的边的一个端点是子极面 \(F \in  {C}^{\prime }\) 。那么根据 \({C}^{\prime }\) 的定义有 \(F \notin  {\sum }_{n}\) ，并且由于斯佩纳条件， \(F\) 因此不是 \({F}_{n + 1}\) 的任何极面的子集: \(F\) 的顶点具有标签 \(1,\ldots ,n\) ，并且 \({F}_{n + 1}\) 除 \({F}_{n}\) 之外的每个极面都缺少这些标签中的一个。因此，根据引理7.6(b)， \({C}^{\prime }\) 中的这样一个子极面 \(F\) 是 \(D \cup  A\) 中恰好两个子单纯形的极面。在图中， \(F\) 的度为2。对于以 \(F\) 为极面的两个子单纯形 \(T \in  D \cup  A\) 中的每一个，如前所述，要么适用情况(a)，要么适用情况(b):(a) 如果 \(T \in  D\) ，则 \(F\) 是 \(T\) 的唯一邻居，该图的边是 \(\{ F,T\}\) ，并且 \(T\) 在图中的度为1；(b) 如果 \(T \in  A\) ，则 \(A\) 恰好有另一个带有标签 \(1,\ldots ,n\) 的极面 \({F}^{\prime }\) ，其中 \({F}^{\prime } \in  C \cup  {C}^{\prime }\) ，并且该图的边是 \(\left\{  {F,{F}^{\prime }}\right\}\) 。
\end{tcolorbox}

Case C. One endpoint of the edge of the graph is a completely labeled sub-simplex \(T \in  D\) . This case has been covered before: \(T\) has a unique facet \(F\) with labels \(1,\ldots ,n\) , which is the unique neighbor of \(T\) , so that \(T\) has degree one in the graph, and \(F \in  C \cup  {C}^{\prime }\) ; the edge of the graph is \(\{ T,F\}\) .

\begin{tcolorbox}
情况C。该图的边的一个端点是完全标记的子单纯形 \(T \in  D\) 。这种情况之前已经讨论过: \(T\) 有一个唯一的带有标签 \(1,\ldots ,n\) 的极面 \(F\) ，它是 \(T\) 的唯一邻居，因此 \(T\) 在图中的度为1，并且 \(F \in  C \cup  {C}^{\prime }\) ；该图的边是 \(\{ T,F\}\) 。
\end{tcolorbox}

In summary, the graph \({G}_{n + 1}\) has nodes in \(C\) and \(D\) which have degree one, and nodes in \({C}^{\prime }\) that have degree two, and no other nodes. Clearly, such a graph is a set of paths and cycles (as can be seen in Figure 7.8; note the cycle of length four on the left). The end nodes of the paths are exactly the degree-one nodes, which are the elements of \(C \cup  D\) (the cycles have no end nodes). Because every path has two end nodes, the number \(\left| C\right|  + \left| D\right|\) of these end nodes is even. (Alternatively, as stated at the beginning of this chapter, every graph has an even number of odd-degree nodes, which are here the nodes of degree one.) Because \(\left| C\right|\) is odd, this implies that \(\left| D\right|\) is also odd, which completes the induction proof for (ii). For \(n = d\) ,(ii) states the theorem.

\begin{tcolorbox}
总之，图 \({G}_{n + 1}\) 在 \(C\) 和 \(D\) 中有度数为一的节点，在 \({C}^{\prime }\) 中有度数为二的节点，没有其他节点。显然，这样的图是一组路径和环(如图 7.8 所示；注意左边长度为四的环)。路径的端点恰好是度数为一的节点，即 \(C \cup  D\) 的元素(环没有端点)。因为每条路径有两个端点，所以这些端点的数量 \(\left| C\right|  + \left| D\right|\) 是偶数。(或者，正如本章开头所述，每个图的奇数度节点数量为偶数，这里度数为一的节点就是奇数度节点。)因为 \(\left| C\right|\) 是奇数，这意味着 \(\left| D\right|\) 也是奇数，这就完成了对 (ii) 的归纳证明。对于 \(n = d\) ，(ii) 表述了该定理。
\end{tcolorbox}

The fact that the number of end nodes of paths is even is called a parity argument. The original proof in Sperner (1928), and very similarly in Knaster, Kuratowski, and Mazurkiewicz (1929), uses also a parity argument without the graphs \({G}_{n}\) , namely the following equation:

\begin{tcolorbox}
路径端点数量为偶数这一事实被称为奇偶性论证。施佩纳(Sperner，1928 年)的原始证明，以及非常类似地，克纳斯特(Knaster)、库拉托夫斯基(Kuratowski)和马祖尔克维奇(Mazurkiewicz，1929 年)的证明，也使用了奇偶性论证，但没有用到图 \({G}_{n}\) ，即以下等式:
\end{tcolorbox}

\[
\left| D\right|  + 2\left| A\right|  = \left| C\right|  + 2\left| {C}^{\prime }\right| . \tag{7.12}
\]

On the left-hand side in (7.12), each completely labeled sub-facet in \(C \cup  {C}^{\prime }\) is counted as either the unique facet of a sub-simplex \(T\) in \(D\) , or as one of the two facets of a sub-simplex \(T\) in \(A\) ; that is, we look at all the sub-simplices in \(D \cup  A\) and count for each one its (one or two) sub-facets that have the labels \(1,\ldots ,n\) . On the right-hand side of (7.12), the sub-facet is either an element of \(C\) and then the facet of exactly one sub-simplex in \(D \cup  A\) , or an element of \({C}^{\prime }\) , and then it is counted twice because it belongs to two sub-simplices in \(D \cup  A\) ; here we have used Lemma 7.6. Clearly,(7.12) implies \(\left| D\right|  = \left| C\right|  + 2\left| {C}^{\prime }\right|  - 2\left| A\right|\) , and therefore that \(\left| D\right|\) is odd because \(\left| C\right|\) is odd, which again completes the induction step for (ii).

\begin{tcolorbox}
在 (7.12) 的左边， \(C \cup  {C}^{\prime }\) 中的每个完全标记的子极面要么被计为 \(D\) 中子单纯形 \(T\) 的唯一极面，要么被计为 \(A\) 中子单纯形 \(T\) 的两个极面之一；也就是说，我们查看 \(D \cup  A\) 中的所有子单纯形，并为每个子单纯形计算其具有标签 \(1,\ldots ,n\) 的(一个或两个)子极面。在 (7.12) 的右边，子极面要么是 \(C\) 的元素，此时它是 \(D \cup  A\) 中恰好一个子单纯形的极面，要么是 \({C}^{\prime }\) 的元素，此时它被计为两次，因为它属于 \(D \cup  A\) 中的两个子单纯形；这里我们使用了引理 7.6。显然，(7.12) 意味着 \(\left| D\right|  = \left| C\right|  + 2\left| {C}^{\prime }\right|  - 2\left| A\right|\) ，因此 \(\left| D\right|\) 是奇数，因为 \(\left| C\right|\) 是奇数，这再次完成了对 (ii) 的归纳步骤。
\end{tcolorbox}

We use the graphs \({G}_{n}\) for \(n = 2,\ldots ,d\) because the picture in Figure 7.8 shows the parity argument immediately via the end nodes of paths, compared to the more abstract equation (7.12) (although the reason is the same). Moreover, the graphs lead to a constructive proof of Theorem 7.10, as we show next.

\begin{tcolorbox}
我们为 \(n = 2,\ldots ,d\) 使用图 \({G}_{n}\) ，是因为图 7.8 中的图示通过路径的端点能立即展示奇偶性论证，与更抽象的等式 (7.12) 相比(尽管原理相同)。此外，正如我们接下来要展示的，这些图能为定理 7.10 提供一个构造性证明。
\end{tcolorbox}

We have completed the proof of Sperner's lemma in two ways, where the second way based on (7.12) seems much shorter. In the remainder of this section on Sperner’s lemma, we combine the graphs \({G}_{2},\ldots ,{G}_{d}\) to show that they can actually be used to find a completely labeled sub-simplex in the given triangulation.

\begin{tcolorbox}
我们用两种方法完成了施佩纳引理的证明，其中基于 (7.12) 的第二种方法似乎要短得多。在关于施佩纳引理这一节的剩余部分，我们将结合图 \({G}_{2},\ldots ,{G}_{d}\) 来展示它们实际上可用于在给定的三角剖分中找到一个完全标记的子单纯形。
\end{tcolorbox}

At first sight, the use of these graphs does not seem constructive, and only proves the induction step of claim (ii) in the proof. For example, consider the triangulation \({\sum }_{2}\) of the bottom edge of the triangle in Figure 7.8. The first completely labeled line segment on the left starts a path that ends at another such line segment, not at a completely labeled triangle. It is easy to continue along the edge (which is just a sequence of line segments) until finding a path that ends at a black dot inside the triangle (there are three such paths in the figure). However, the triangle may just be the next face \({F}_{3}\) in a larger simplex, for example a tetrahedron, and some other black dots (for example near the top right) may be needed to find a completely labeled sub-simplex in the next triangulation \({\sum }_{4}\) .

\begin{tcolorbox}
乍一看，使用这些图似乎不是构造性的，只是证明了证明中命题 (ii) 的归纳步骤。例如，考虑图 7.8 中三角形底边的三角剖分 \({\sum }_{2}\) 。左边第一个完全标记的线段开启了一条路径，该路径终止于另一个这样的线段，而不是一个完全标记的三角形。很容易沿着边(它只是一系列线段)继续，直到找到一条终止于三角形内部黑点的路径(图中有三条这样的路径)。然而，这个三角形可能只是一个更大单纯形(例如四面体)中的下一个面 \({F}_{3}\) ，可能需要其他一些黑点(例如右上角附近的)来在下一步三角剖分 \({\sum }_{4}\) 中找到一个完全标记的子单纯形。
\end{tcolorbox}

\begin{center}
\includegraphics[max width=0.6\textwidth]{images/bo_d2e5cmbef24c73b03tn0_196_384_1023_870_409_0.jpg}
\end{center}
\hspace*{3em} 

Figure 7.9 Triangulation of \({F}_{2}\) , which is the bottom edge of the triangle in Figures 7.6-7.8, with line segments as sub-simplices, and the vertices in \({\sum }_{2}\) , marked by small dots, as sub-facets. The Sperner labeling is (a), analogous to Figure 7.6. The completely labeled sub-simplices of \({\sum }_{2}\) are marked with white dots in (b), analogous to the black dots in Figure 7.7; the completely labeled sub-facet is the leftmost vertex marked in gray, analogous to any white dot in Figure 7.7. The graph \({G}_{2}\) is shown in (c), analogous to Figure 7.8, and the graph alone in (d), which in this dimension is only a set of paths.

\begin{tcolorbox}
图7.9 \({F}_{2}\) 的三角剖分， \({F}_{2}\) 是图7.6 - 7.8中三角形的底边，以线段作为子单纯形， \({\sum }_{2}\) 中的顶点(用小点标记)作为子极面。斯佩纳标号如(a)所示，与图7.6类似。 \({\sum }_{2}\) 的完全标记子单纯形在(b)中用白点标记，类似于图7.7中的黑点；完全标记的子极面是最左边用灰色标记的顶点，类似于图7.7中的任何一个白点。图 \({G}_{2}\) 如(c)所示，类似于图7.8，(d)中单独展示了该图，在这个维度中它只是一组路径。
\end{tcolorbox}

A more systematic approach is needed, for which we consider the first induction step from \(n = 1\) to \(n = 2\) of claim (ii) in the proof, shown in Figure 7.9. The facet \({F}_{1}\) consists only of a single vertex \({v}^{1}\) , marked as a gray node on the left. This defines the graph \({G}_{1}\) , which has no edges, but one node (and hence an odd number of nodes) given by \({v}^{1}\) , which is also the single sub-simplex of \({\sum }_{1}\) and completely labeled; the graph \({G}_{1}\) is the only graph which does not consist of paths and cycles. In \({\sum }_{2}\) , the sub-simplices are line segments, and their sub-facets are single vertices. The set \({C}^{\prime }\) in (7.12) consists of those vertices that have label 1, marked in gray in Figure 7.9, in addition to the leftmost gray node as the unique element of \(C\) . This is one of the end nodes of the paths in \({G}_{2}\) . All other end nodes are the completely labeled sub-simplices in \({\sum }_{2}\) , marked as before with white dots. (There are no cycles because \({F}_{2}\) is just a line segment.)

\begin{tcolorbox}
需要一种更系统的方法，为此我们考虑证明中命题(ii)从 \(n = 1\) 到 \(n = 2\) 的第一步归纳，如图7.9所示。极面 \({F}_{1}\) 仅由单个顶点 \({v}^{1}\) 组成，在左边标记为灰色节点。这定义了图 \({G}_{1}\) ，它没有边，但有一个由 \({v}^{1}\) 给出的节点(因此节点数为奇数)， \({v}^{1}\) 也是 \({\sum }_{1}\) 的唯一子单纯形且是完全标记的；图 \({G}_{1}\) 是唯一不由路径和环组成的图。在 \({\sum }_{2}\) 中，子单纯形是线段，它们的子极面是单个顶点。(7.12)中的集合 \({C}^{\prime }\) 由那些标签为1的顶点组成，在图7.9中用灰色标记，此外最左边的灰色节点是 \(C\) 的唯一元素。这是 \({G}_{2}\) 中路径的一个端点。所有其他端点是 \({\sum }_{2}\) 中完全标记的子单纯形，如之前一样用白点标记。(因为 \({F}_{2}\) 只是一条线段，所以没有环。)
\end{tcolorbox}

We now state our constructive version of Sperner's lemma, obtained by just taking the union of the considered graphs, which concatenates their paths.

\begin{tcolorbox}
现在我们陈述斯佩纳引理的构造性版本，它是通过取所考虑图的并集得到的，该并集将它们的路径连接起来。
\end{tcolorbox}

Proposition 7.11. Consider a triangulation \(\sum\) of a simplex \(S\) with vertices \({v}^{1},\ldots ,{v}^{d}\) and a Sperner labeling as in Theorem 7.10, \(d \geq  2\) . Let \({G}_{1},\ldots ,{G}_{d}\) be the graphs defined in the proof of the theorem, and define the graph \(G\) by

\begin{tcolorbox}
命题7.11。考虑单纯形 \(S\) 的三角剖分 \(\sum\) ，其顶点为 \({v}^{1},\ldots ,{v}^{d}\) ，并采用如定理7.10中的斯佩纳标号 \(d \geq  2\) 。设 \({G}_{1},\ldots ,{G}_{d}\) 是定理证明中定义的图，并按如下方式定义图 \(G\)
\end{tcolorbox}

\[
G = {G}_{1} \cup  \cdots  \cup  {G}_{d} \tag{7.13}
\]

where the union is taken separately for the sets of nodes and edges, respectively. Then every node in \(G\) has degree one or two, so that \(G\) consists of paths and cycles. The end nodes of the paths are given by \({v}^{1}\) and the completely labeled sub-simplices of \(\sum\) .

\begin{tcolorbox}
其中并集分别对节点集和边集取。那么 \(G\) 中的每个节点的度数为1或2，因此 \(G\) 由路径和环组成。路径的端点由 \({v}^{1}\) 和 \(\sum\) 的完全标记子单纯形给出。
\end{tcolorbox}

Proof. Let \(1 \leq  n < d\) . According to the construction in (7.12), the common nodes of \({G}_{n}\) and \({G}_{n + 1}\) are the sub-simplices in \({\sum }_{n}\) with \(n\) vertices that have all labels \(1,\ldots ,n\) . For \(n > 1\) , these nodes have degree one in both \({G}_{n}\) and \({G}_{n + 1}\) , and therefore degree two in \({G}_{n} \cup  {G}_{n + 1}\) and in \(G\) . Hence, the degree-one nodes in \(G\) are given by \({v}^{1}\) (which is the only node in \({G}_{1}\) , where it has degree zero, and which has degree one in \({G}_{2}\) ), and the completely labeled sub-simplices in \({G}_{d}\) , as claimed.

\begin{tcolorbox}
证明。设 \(1 \leq  n < d\) 。根据(7.12)中的构造， \({G}_{n}\) 和 \({G}_{n + 1}\) 的公共节点是 \({\sum }_{n}\) 中具有 \(n\) 个顶点且所有顶点标签为 \(1,\ldots ,n\) 的子单纯形。对于 \(n > 1\) ，这些节点在 \({G}_{n}\) 和 \({G}_{n + 1}\) 中的度数均为1，因此在 \({G}_{n} \cup  {G}_{n + 1}\) 和 \(G\) 中的度数为2。因此，如所声称的， \(G\) 中度数为1的节点由 \({v}^{1}\) (它是 \({G}_{1}\) 中的唯一节点，在其中度数为0，在 \({G}_{2}\) 中度数为1)和 \({G}_{d}\) 中完全标记的子单纯形给出。
\end{tcolorbox}

\begin{center}
\includegraphics[max width=0.6\textwidth]{images/bo_d2e5cmbef24c73b03tn0_197_524_1329_865_382_0.jpg}
\end{center}
\hspace*{3em} 

Figure 7.10 Example of \(G\) in (7.13) for \(d = 4\) . The graph \({G}_{2}\) is from Figure 7.9. The graph \({G}_{3}\) from Figure 7.8 is drawn as a "layer" between the completely labeled elements of \({\sum }_{2}\) and \({\sum }_{3}\) , and the graph \({G}_{4}\) as another such layer. The triangles indicate sub-facets in \({\sum }_{4}\) with labels1,2,3(apart from the black dots, which also belong to \({G}_{3}\) ), and the squares are sub-simplices in \({\sum }_{4}\) with all labels1,2,3,4.

\begin{tcolorbox}
图7.10 (7.13)中 \(G\) 在 \(d = 4\) 情况下的示例。图 \({G}_{2}\) 取自图7.9。图7.8中的图 \({G}_{3}\) 被绘制为 \({\sum }_{2}\) 和 \({\sum }_{3}\) 的完全标记元素之间的一个“层”，图 \({G}_{4}\) 则作为另一个这样的层。三角形表示 \({\sum }_{4}\) 中标签为1、2、3的子极面(除黑点外，黑点也属于 \({G}_{3}\) )，正方形表示 \({\sum }_{4}\) 中标签为1、2、3、4的子单纯形。
\end{tcolorbox}

An example of \(G\) is shown in Figure 7.10, for a triangulated tetrahedron which has \({F}_{3}\) in Figure 7.8 as its bottom facet. The concatenated paths from \({G}_{2},\ldots ,{G}_{d}\) may result in new cycles, as on the bottom right. The path that starts at \({v}^{1}\) ends at a completely labeled sub-simplex in \(\sum\) , which is therefore found by following the path. The example shows that the path may move up and down through the different dimensions, such as visiting a triangle in \({G}_{4}\) and coming back to a single vertex in \({G}_{2}\) . The construction of \(G\) in (7.13) is very clean because it does not need any further addition to the triangulation \(\sum\) and its labeled vertices.

\begin{tcolorbox}
图7.10展示了 \(G\) 的一个示例，针对一个三角剖分的四面体，其底面极面是图7.8中的 \({F}_{3}\) 。来自 \({G}_{2},\ldots ,{G}_{d}\) 的拼接路径可能会形成新的循环，如右下角所示。从 \({v}^{1}\) 开始的路径终止于 \(\sum\) 中的一个完全标记子单纯形，因此通过沿着该路径可以找到它。该示例表明，路径可能会在不同维度间上下移动，例如访问 \({G}_{4}\) 中的一个三角形，然后回到 \({G}_{2}\) 中的一个顶点。(7.13)中 \(G\) 的构造非常简洁，因为它不需要对三角剖分 \(\sum\) 及其标记顶点进行任何进一步的添加。
\end{tcolorbox}

\section*{7.5 The Knaster-Kuratowski-Mazurkiewicz Lemma}

\begin{tcolorbox}
\section*{7.5 克纳斯特 - 库拉托夫斯基 - 马祖尔克维奇引理}
\end{tcolorbox}

In this section, we use Sperner's lemma to prove a theorem due to Knaster, Kuratowski, and Mazurkiewicz (1929) known as the "KKM lemma" after the initials of the authors. The KKM lemma considers a simplex \(S\) with vertices \({v}^{1},\ldots ,{v}^{d}\) and closed sets \({A}_{1},\ldots ,{A}_{d}\) . The assumption, which is closely related to the Sperner condition, is that for any set of vertices, the face of \(S\) with these vertices \({v}^{i}\) is covered by the corresponding closed sets \({A}_{i}\) , as illustrated in Figure 7.11. The conclusion is that then these sets have a common point. In some contexts, such closed sets can be defined more naturally than a function for which one tries to find a fixed point, so that the KKM lemma is a useful tool in its own right. At any rate, it leads almost immediately to a proof of Brouwer's fixed-point theorem on the unit simplex (Theorem 7.13). As background material, we recall the concept of Euclidean norm and distance, and the classic theorem of Bolzano-Weierstrass on the existence of converging subsequences, which is used in the proof of the KKM lemma.

\begin{tcolorbox}
在本节中，我们使用斯佩纳引理来证明一个由克纳斯特(Knaster)、库拉托夫斯基(Kuratowski)和马祖尔克维奇(Mazurkiewicz)(1929年)提出的定理，该定理以作者姓氏首字母命名为“KKM引理”。KKM引理考虑一个单纯形 \(S\) ，其顶点为 \({v}^{1},\ldots ,{v}^{d}\) ，以及闭集 \({A}_{1},\ldots ,{A}_{d}\) 。与斯佩纳条件密切相关的假设是，对于任意顶点集， \(S\) 中以这些顶点 \({v}^{i}\) 构成的面被相应的闭集 \({A}_{i}\) 覆盖，如图7.11所示。结论是这些集合有一个公共点。在某些情况下，这样的闭集比试图寻找不动点的函数更自然地被定义，因此KKM引理本身就是一个有用的工具。无论如何，它几乎可以立即导出单位单纯形上布劳威尔不动点定理的证明(定理7.13)。作为背景材料，我们回顾欧几里得范数与距离的概念，以及关于收敛子列存在性的经典波尔查诺 - 魏尔斯特拉斯定理，该定理用于KKM引理的证明。
\end{tcolorbox}

Background material: Euclidean norm and distance

\begin{tcolorbox}
背景材料:欧几里得范数与距离
\end{tcolorbox}

For \(a = {\left( {a}_{1},\ldots ,{a}_{d}\right) }^{\top } \in  {\mathbb{R}}^{d}\) , the Euclidean norm \(\parallel a\parallel\) of \(a\) is defined as

\begin{tcolorbox}
对于 \(a = {\left( {a}_{1},\ldots ,{a}_{d}\right) }^{\top } \in  {\mathbb{R}}^{d}\) ， \(a\) 的欧几里得范数 \(\parallel a\parallel\) 定义为
\end{tcolorbox}

\[
\parallel a\parallel  = \sqrt{{a}^{\top }a} = \sqrt{{a}_{1}^{2} + \cdots  + {a}_{d}^{2}}.
\]

The Euclidean distance between two points \(a,b \in  {\mathbb{R}}^{d}\) is defined as \(\parallel a - b\parallel\) .

\begin{tcolorbox}
两点 \(a,b \in  {\mathbb{R}}^{d}\) 之间的欧几里得距离定义为 \(\parallel a - b\parallel\) 。
\end{tcolorbox}

Theorem 7.12 (The KKM lemma). Let \(S\) be a simplex with vertices \({v}^{1},\ldots ,{v}^{d}\) , and let \({A}_{1},\ldots ,{A}_{d}\) be closed sets such that for all \(I \subseteq  \{ 1,\ldots ,d\}\) the face \({F}_{I} = \operatorname{conv}\left\{  {{v}^{i} \mid  i \in  I}\right\}\) of \(S\) is a subset of \(\mathop{\bigcup }\limits_{{i \in  I}}{A}_{i}\) . Then \(\mathop{\bigcap }\limits_{{i = 1}}^{d}{A}_{i} \neq  \varnothing\) .

\begin{tcolorbox}
定理7.12(KKM引理)。设 \(S\) 是一个顶点为 \({v}^{1},\ldots ,{v}^{d}\) 的单纯形，设 \({A}_{1},\ldots ,{A}_{d}\) 是闭集，使得对于所有 \(I \subseteq  \{ 1,\ldots ,d\}\) ， \(S\) 中以这些顶点构成的面 \({F}_{I} = \operatorname{conv}\left\{  {{v}^{i} \mid  i \in  I}\right\}\) 是相应闭集 \(\mathop{\bigcup }\limits_{{i \in  I}}{A}_{i}\) 的子集。那么 \(\mathop{\bigcap }\limits_{{i = 1}}^{d}{A}_{i} \neq  \varnothing\) 。
\end{tcolorbox}

Proof. Let \(\varepsilon  > 0\) , and consider a triangulation \({\sum }_{\varepsilon }\) of \(S\) so that for any sub-simplex \(T\) in \({\sum }_{\varepsilon }\) and any \(x,y \in  T\) we have \(\parallel x - y\parallel  \leq  \varepsilon\) ; we will show in Section 7.7 that such an \(\varepsilon\) -fine triangulation exists for any \(\varepsilon  > 0\) . For any vertex \(x\) in \({\sum }_{\varepsilon }\) , let \(I \subseteq  \{ 1,\ldots ,d\}\) be inclusion-minimal so that \(x \in  {F}_{I}\) . By assumption, \({F}_{I} \subseteq  \mathop{\bigcup }\limits_{{i \in  I}}{A}_{i}\) and therefore \(x \in  {A}_{i}\) for some \(i \in  I\) , where we choose \(i\) as a label of \(x\) . By construction, this is a Sperner labeling. According to Sperner's lemma (Theorem 7.10), there is a completely labeled sub-simplex in \({\sum }_{\varepsilon }\) with vertices \({w}_{\varepsilon }^{1},\ldots ,{w}_{\varepsilon }^{d}\) so that \({w}_{\varepsilon }^{i}\) has label \(i\) and therefore belongs to \({A}_{i}\) for \(1 \leq  i \leq  d\) .

\begin{tcolorbox}
证明。设 \(\varepsilon  > 0\) ，并考虑 \(S\) 的一个三角剖分 \({\sum }_{\varepsilon }\) ，使得对于 \({\sum }_{\varepsilon }\) 中的任何子单纯形 \(T\) 和任何 \(x,y \in  T\) ，我们有 \(\parallel x - y\parallel  \leq  \varepsilon\) ；我们将在 7.7 节中证明，对于任何 \(\varepsilon  > 0\) ，这样的 \(\varepsilon\) - 精细三角剖分是存在的。对于 \({\sum }_{\varepsilon }\) 中的任何顶点 \(x\) ，设 \(I \subseteq  \{ 1,\ldots ,d\}\) 是包含关系下的极小元，使得 \(x \in  {F}_{I}\) 。根据假设， \({F}_{I} \subseteq  \mathop{\bigcup }\limits_{{i \in  I}}{A}_{i}\) ，因此对于某个 \(i \in  I\) ，有 \(x \in  {A}_{i}\) ，我们选择 \(i\) 作为 \(x\) 的标签。根据构造，这是一个斯佩纳标号。根据斯佩纳引理(定理 7.10)， \({\sum }_{\varepsilon }\) 中存在一个完全标记的子单纯形，其顶点为 \({w}_{\varepsilon }^{1},\ldots ,{w}_{\varepsilon }^{d}\) ，使得 \({w}_{\varepsilon }^{i}\) 具有标签 \(i\) ，因此对于 \(1 \leq  i \leq  d\) ，它属于 \({A}_{i}\) 。
\end{tcolorbox}

Background material: Every bounded sequence has a converging subsequence

\begin{tcolorbox}
背景材料:每个有界序列都有一个收敛子序列
\end{tcolorbox}

The theorem of Bolzano-Weierstrass states: Let \({x}^{\left( 1\right) },{x}^{\left( 2\right) },\ldots\) be a bounded sequence of points in \({\mathbb{R}}^{d}\) , that is, there is some \(M\) so that \(\begin{Vmatrix}{x}^{\left( n\right) }\end{Vmatrix} \leq  M\) for all \(n\) . Then this sequence has a converging subsequence.

\begin{tcolorbox}
波尔查诺 - 魏尔斯特拉斯定理指出:设 \({x}^{\left( 1\right) },{x}^{\left( 2\right) },\ldots\) 是 \({\mathbb{R}}^{d}\) 中的一个有界点列，即存在某个 \(M\) ，使得对于所有的 \(n\) ，都有 \(\begin{Vmatrix}{x}^{\left( n\right) }\end{Vmatrix} \leq  M\) 。那么这个序列有一个收敛子序列。
\end{tcolorbox}

This important theorem has a nice proof that we recall here. We first consider a bounded sequence of real numbers \({y}^{\left( 1\right) },{y}^{\left( 2\right) },\ldots\) where \(\left| {y}^{\left( n\right) }\right|  \leq  M\) for all \(n\) . We show that this sequence has a monotonic subsequence \({y}^{\left( {k}_{1}\right) },{y}^{\left( {k}_{2}\right) },\ldots\) , that is, either \({y}^{\left( {k}_{1}\right) } \geq  {y}^{\left( {k}_{2}\right) } \geq  \cdots\) or \({y}^{\left( {k}_{1}\right) } \leq  {y}^{\left( {k}_{2}\right) } \leq  \cdots\) holds. Imagine the real numbers \({y}^{\left( 1\right) },{y}^{\left( 2\right) },\ldots\) as the heights of "hotels at the seaside". From the top of each hotel with height \({y}^{\left( k\right) }\) you can look beyond the subsequent hotels with heights \({y}^{\left( k + 1\right) },{y}^{\left( k + 2\right) },\ldots\) as long as they have lower height, and see the sea at infinity if these are all lower. In other words, a hotel has "seaview" if it belongs to the set \(K = \left\{  {k \mid  {y}^{\left( k\right) } > {y}^{\left( j\right) }\text{ for all }j > k}\right\}\) (presumably, these are very expensive hotels). If \(K\) is infinite, then we take the elements of \(K\) in ascending order as the superscripts \({k}_{1},{k}_{2},{k}_{3},\ldots\) that give our subsequence \({y}^{\left( {k}_{1}\right) },{y}^{\left( {k}_{2}\right) },{y}^{\left( {k}_{3}\right) },\ldots\) , which is clearly decreasing. If, however, \(K\) is finite with maximal element \(m\) (take \(m = 0\) if \(K\) is empty), then for each \(k > m\) we have \(k \notin  K\) and hence there exists some \(j > k\) with \({y}^{\left( k\right) } \leq  {y}^{\left( j\right) }\) . Starting with \({y}^{\left( {k}_{1}\right) }\) for \({k}_{1} = m + 1\) we let \({k}_{2} > {k}_{1}\) with \({y}^{\left( {k}_{1}\right) } \leq  {y}^{\left( {k}_{2}\right) }\) . Then find another \({k}_{3} > {k}_{2}\) so that \({y}^{\left( {k}_{2}\right) } \leq  {y}^{\left( {k}_{3}\right) }\) , and so on, which gives a weakly increasing subsequence \({y}^{\left( {k}_{1}\right) } \leq  {y}^{\left( {k}_{2}\right) } \leq  {y}^{\left( {k}_{3}\right) } \leq  \cdots\) . In either case, the original sequence has a monotonic subsequence.

\begin{tcolorbox}
这个重要的定理有一个很好的证明，我们在此回顾一下。我们首先考虑一个有界的实数序列 \({y}^{\left( 1\right) },{y}^{\left( 2\right) },\ldots\) ，其中对于所有的 \(n\) 都有 \(\left| {y}^{\left( n\right) }\right|  \leq  M\) 。我们证明这个序列有一个单调子序列 \({y}^{\left( {k}_{1}\right) },{y}^{\left( {k}_{2}\right) },\ldots\) ，即要么 \({y}^{\left( {k}_{1}\right) } \geq  {y}^{\left( {k}_{2}\right) } \geq  \cdots\) 成立，要么 \({y}^{\left( {k}_{1}\right) } \leq  {y}^{\left( {k}_{2}\right) } \leq  \cdots\) 成立。把实数 \({y}^{\left( 1\right) },{y}^{\left( 2\right) },\ldots\) 想象成“海边酒店”的高度。从每个高度为 \({y}^{\left( k\right) }\) 的酒店顶部，只要后续高度为 \({y}^{\left( k + 1\right) },{y}^{\left( k + 2\right) },\ldots\) 的酒店高度更低，你就可以看到它们之外的景色，如果这些酒店的高度都更低，你还能看到无穷远处的大海。换句话说，如果一家酒店属于集合 \(K = \left\{  {k \mid  {y}^{\left( k\right) } > {y}^{\left( j\right) }\text{ for all }j > k}\right\}\) ，那么它就有“海景”(想必，这些都是非常昂贵的酒店)。如果 \(K\) 是无限集，那么我们按升序取 \(K\) 中的元素作为下标 \({k}_{1},{k}_{2},{k}_{3},\ldots\) ，从而得到我们的子序列 \({y}^{\left( {k}_{1}\right) },{y}^{\left( {k}_{2}\right) },{y}^{\left( {k}_{3}\right) },\ldots\) ，它显然是递减的。然而，如果 \(K\) 是有限集，其最大元素为 \(m\) (如果 \(K\) 为空集，则取 \(m = 0\) )，那么对于每个 \(k > m\) ，我们有 \(k \notin  K\) ，因此存在某个 \(j > k\) 使得 \({y}^{\left( k\right) } \leq  {y}^{\left( j\right) }\) 。从 \({k}_{1} = m + 1\) 对应的 \({y}^{\left( {k}_{1}\right) }\) 开始，我们令 \({k}_{2} > {k}_{1}\) 且 \({y}^{\left( {k}_{1}\right) } \leq  {y}^{\left( {k}_{2}\right) }\) 。然后找到另一个 \({k}_{3} > {k}_{2}\) 使得 \({y}^{\left( {k}_{2}\right) } \leq  {y}^{\left( {k}_{3}\right) }\) ，依此类推，这样就得到一个弱递增子序列 \({y}^{\left( {k}_{1}\right) } \leq  {y}^{\left( {k}_{2}\right) } \leq  {y}^{\left( {k}_{3}\right) } \leq  \cdots\) 。无论哪种情况，原序列都有一个单调子序列。
\end{tcolorbox}

This monotonic subsequence \({y}^{\left( {k}_{1}\right) },{y}^{\left( {k}_{2}\right) },\ldots\) is bounded because the original sequence is bounded. It is easy to show that the subsequence converges to the infimum of its members if it is decreasing, or to the supremum if it is weakly increasing. This shows the Bolzano-Weierstrass theorem for a sequence of real numbers in \({\mathbb{R}}^{1}\) .

\begin{tcolorbox}
这个单调子序列 \({y}^{\left( {k}_{1}\right) },{y}^{\left( {k}_{2}\right) },\ldots\) 是有界的，因为原序列是有界的。很容易证明，如果该子序列是递减的，它收敛于其元素的下确界；如果它是弱递增的，则收敛于其上确界。这就证明了实数序列在 \({\mathbb{R}}^{1}\) 中的波尔查诺 - 魏尔施特拉斯定理(Bolzano - Weierstrass theorem)。
\end{tcolorbox}

To show the theorem for a sequence of points \({x}^{\left( n\right) } = \left( {{x}_{1}^{\left( n\right) },\ldots ,{x}_{d}^{\left( n\right) }}\right)\) in \({\mathbb{R}}^{d}\) for \(n = 1,2,3,\ldots\) , note that the components \({x}_{i}^{\left( n\right) }\) for \(i = 1,\ldots ,d\) of these points are bounded because \(\left| {x}_{i}^{\left( n\right) }\right|  \leq  \begin{Vmatrix}{x}^{\left( n\right) }\end{Vmatrix} \leq  M\) . First, find a converging subsequence of the first components \({x}_{1}^{\left( 1\right) },{x}_{1}^{\left( 2\right) },{x}_{1}^{\left( 3\right) },\ldots\) with limit \({z}_{1}\) , say. For ease of notation, replace the original sequence of points with this subsequence, which now converges in its first component. Second, find a subsequence so that the second components \({x}_{2}^{\left( 1\right) },{x}_{2}^{\left( 2\right) },{x}_{2}^{\left( 3\right) },\ldots\) have limit \({z}_{2}\) , so that now both the first and second components converge. By continuing in this manner, we obtain a subsequence with limit \(\left( {{z}_{1},{z}_{2},\ldots ,{z}_{d}}\right)\) , which completes the proof.

\begin{tcolorbox}
为了证明该定理对于 \(n = 1,2,3,\ldots\) 时 \({\mathbb{R}}^{d}\) 中的点列 \({x}^{\left( n\right) } = \left( {{x}_{1}^{\left( n\right) },\ldots ,{x}_{d}^{\left( n\right) }}\right)\) 成立，注意到这些点的 \(i = 1,\ldots ,d\) 对应的分量 \({x}_{i}^{\left( n\right) }\) 是有界的，因为 \(\left| {x}_{i}^{\left( n\right) }\right|  \leq  \begin{Vmatrix}{x}^{\left( n\right) }\end{Vmatrix} \leq  M\) 。首先，找出第一个分量 \({x}_{1}^{\left( 1\right) },{x}_{1}^{\left( 2\right) },{x}_{1}^{\left( 3\right) },\ldots\) 的一个收敛子列，设其极限为 \({z}_{1}\) 。为了方便表示，用这个子列替换原来的点列，此时该子列的第一个分量收敛。其次，找出一个子列，使得第二个分量 \({x}_{2}^{\left( 1\right) },{x}_{2}^{\left( 2\right) },{x}_{2}^{\left( 3\right) },\ldots\) 有极限 \({z}_{2}\) ，这样一来，第一个和第二个分量都收敛。按照这种方式继续下去，我们得到一个极限为 \(\left( {{z}_{1},{z}_{2},\ldots ,{z}_{d}}\right)\) 的子列，从而完成了证明。
\end{tcolorbox}

Consider a sequence of positive reals \(\varepsilon \left( 1\right) ,\varepsilon \left( 2\right) ,\ldots\) that tends to 0 (for example, \(\varepsilon \left( n\right)  = \frac{1}{n}\) ), an \(\varepsilon \left( n\right)\) -fine triangulation \({\sum }_{\varepsilon \left( n\right) }\) of \(S\) , and the corresponding sequence of points \({w}_{\varepsilon \left( 1\right) }^{1},{w}_{\varepsilon \left( 2\right) }^{1},\ldots\) in the simplex \(S\) that have label 1 considered earlier. Because

\begin{tcolorbox}
考虑一个趋于 0 的正实数序列 \(\varepsilon \left( 1\right) ,\varepsilon \left( 2\right) ,\ldots\) (例如， \(\varepsilon \left( n\right)  = \frac{1}{n}\) )， \(S\) 的一个 \(\varepsilon \left( n\right)\) - 精细三角剖分 \({\sum }_{\varepsilon \left( n\right) }\) ，以及前面所考虑的单纯形 \(S\) 中标签为 1 的对应点列 \({w}_{\varepsilon \left( 1\right) }^{1},{w}_{\varepsilon \left( 2\right) }^{1},\ldots\) 。因为
\end{tcolorbox}

\begin{center}
\includegraphics[max width=0.9\textwidth]{images/bo_d2e5cmbef24c73b03tn0_200_198_208_1251_980_0.jpg}
\end{center}
\hspace*{3em} 

Figure 7.11 Illustration of the assumptions of the KKM lemma for \(d = 3\) , for the sets \(I = \{ 1,2\} ,\{ 1,3\} ,\{ 2,3\}\) , and \(\{ 1,2,3\}\) .

\begin{tcolorbox}
图 7.11 展示了对于 \(d = 3\) 、集合 \(I = \{ 1,2\} ,\{ 1,3\} ,\{ 2,3\}\) 和 \(\{ 1,2,3\}\) 的 KKM 引理的假设。
\end{tcolorbox}

\(S\) is bounded and closed, this sequence of points has a converging subsequence with limit \(z\) in \(S\) . For all \(i = 1,\ldots ,d\) , we have \(\begin{Vmatrix}{{w}_{\varepsilon \left( n\right) }^{1} - {w}_{\varepsilon \left( n\right) }^{i}}\end{Vmatrix} \leq  \varepsilon \left( n\right)\) , so that the points \({w}_{\varepsilon \left( n\right) }^{i}\) in that subsequence with label \(i\) also converge to \(z\) as \(n \rightarrow  \infty\) , where \({w}_{\varepsilon \left( n\right) }^{i} \in  {A}_{i}\) . Because \({A}_{i}\) is closed, \(z \in  {A}_{i}\) . Therefore, the intersection \(\mathop{\bigcap }\limits_{{i = 1}}^{d}{A}_{i}\) contains \(z\) and is not empty, as claimed.

\begin{tcolorbox}
由于 \(S\) 是有界闭集，这个点列有一个收敛子列，其极限 \(z\) 属于 \(S\) 。对于所有的 \(i = 1,\ldots ,d\) ，我们有 \(\begin{Vmatrix}{{w}_{\varepsilon \left( n\right) }^{1} - {w}_{\varepsilon \left( n\right) }^{i}}\end{Vmatrix} \leq  \varepsilon \left( n\right)\) ，使得该子列中标签为 \(i\) 的点 \({w}_{\varepsilon \left( n\right) }^{i}\) 当 \(n \rightarrow  \infty\) 时也收敛到 \(z\) ，其中 \({w}_{\varepsilon \left( n\right) }^{i} \in  {A}_{i}\) 。因为 \({A}_{i}\) 是闭集，所以 \(z \in  {A}_{i}\) 。因此，交集 \(\mathop{\bigcap }\limits_{{i = 1}}^{d}{A}_{i}\) 包含 \(z\) 且非空，正如所声称的那样。
\end{tcolorbox}

The following theorem is an immediate consequence of the KKM lemma, applied to the unit simplex \(\Delta\) . The closed sets \({A}_{i}\) are those points \(x\) in \(\Delta\) that have label \(i\) according to our original definition (7.2).

\begin{tcolorbox}
以下定理是将 KKM 引理应用于单位单纯形 \(\Delta\) 的直接结果。闭集 \({A}_{i}\) 是根据我们最初的定义(7.2)在 \(\Delta\) 中标签为 \(i\) 的那些点 \(x\) 。
\end{tcolorbox}

Theorem 7.13 (Brouwer's fixed-theorem on the unit simplex). A continuous function \(f : \Delta  \rightarrow  \Delta\) on the(d - 1)-dimensional unit simplex \(\Delta\) in (7.1) has a fixed point.

\begin{tcolorbox}
定理7.13(单位单纯形上的布劳威尔不动点定理)。式(7.1)中 (d - 1) 维度的单位单纯形 \(\Delta\) 上的连续函数 \(f : \Delta  \rightarrow  \Delta\) 有一个不动点。
\end{tcolorbox}

Proof. For \(i = 1,\ldots ,d\) , let \({A}_{i} = \left\{  {x \in  \Delta  \mid  {x}_{i} \geq  {f}_{i}\left( x\right) }\right\}\) . Then \({A}_{i}\) is a closed set because \(f\) and hence \({f}_{i}\) is continuous. Let \(x \in  \Delta\) and \(I = \left\{  {i \mid  {x}_{i} > 0}\right\}\) . Because the vertices of \(\Delta\) are the unit vectors, the face \({F}_{I}\) in Theorem 7.12 is the smallest facet of \(\Delta\) that contains \(x\) . As shown in Lemma 7.9, there is some \(i\) so that \({x}_{i} > 0\) and \(x \in  {A}_{i}\) . Because \(x\) is arbitrary, this shows \({F}_{I} \subseteq  \mathop{\bigcup }\limits_{{i \in  I}}{A}_{i}\) for any \(I \subseteq  \{ 1,\ldots ,d\}\) . By Theorem 7.12, there is some \(z \in  \mathop{\bigcap }\limits_{{i = 1}}^{d}{A}_{i}\) , where \(f\left( z\right)  = z\) by Lemma 7.1.

\begin{tcolorbox}
证明。对于 \(i = 1,\ldots ,d\) ，设 \({A}_{i} = \left\{  {x \in  \Delta  \mid  {x}_{i} \geq  {f}_{i}\left( x\right) }\right\}\) 。那么 \({A}_{i}\) 是一个闭集，因为 \(f\) ，因此 \({f}_{i}\) 是连续的。设 \(x \in  \Delta\) 且 \(I = \left\{  {i \mid  {x}_{i} > 0}\right\}\) 。因为 \(\Delta\) 的顶点是单位向量，定理7.12中的面 \({F}_{I}\) 是 \(\Delta\) 中包含 \(x\) 的最小极面。如引理7.9所示，存在某个 \(i\) 使得 \({x}_{i} > 0\) 且 \(x \in  {A}_{i}\) 。由于 \(x\) 是任意的，这表明对于任意 \(I \subseteq  \{ 1,\ldots ,d\}\) 有 \({F}_{I} \subseteq  \mathop{\bigcup }\limits_{{i \in  I}}{A}_{i}\) 。根据定理7.12，存在某个 \(z \in  \mathop{\bigcap }\limits_{{i = 1}}^{d}{A}_{i}\) ，根据引理7.1有 \(f\left( z\right)  = z\) 。
\end{tcolorbox}

\section*{7.6 Brouwer's Fixed-Point Theorem on a General Compact Convex Set}

\begin{tcolorbox}
\section*{7.6 一般紧致凸集上的布劳威尔不动点定理}
\end{tcolorbox}

Theorem 7.13 states that a continuous function on the unit simplex \(\Delta\) has a fixed point. In this section, we prove this for a continuous function \(f\) on a compact convex set \(C\) that is a subset of some \({\mathbb{R}}^{k}\) , which is the general form of Brouwer’s fixed-point theorem. The first step is to embed \(C\) into \(\Delta\) with a bijective linear map. Assuming that \(C\) is a subset of \(\Delta\) , the second step is to construct a continuous function \(p : \Delta  \rightarrow  C\) so that \(p\left( x\right)  = x\) for \(x \in  C\) . Then the continuous function \(h : \Delta  \rightarrow  \Delta\) defined by \(h\left( x\right)  = f\left( {p\left( x\right) }\right)\) for \(x \in  \Delta\) has a fixed point, which is easily seen to be a fixed point of the given function \(f : C \rightarrow  C\) .

\begin{tcolorbox}
定理7.13表明单位单纯形 \(\Delta\) 上的连续函数有一个不动点。在本节中，我们证明对于某个 \({\mathbb{R}}^{k}\) 的子集——紧致凸集 \(C\) 上的连续函数 \(f\) 也成立，这是布劳威尔不动点定理的一般形式。第一步是通过一个双射线性映射将 \(C\) 嵌入到 \(\Delta\) 中。假设 \(C\) 是 \(\Delta\) 的一个子集，第二步是构造一个连续函数 \(p : \Delta  \rightarrow  C\) 使得对于 \(x \in  C\) 有 \(p\left( x\right)  = x\) 。那么由对于 \(x \in  \Delta\) 定义为 \(h\left( x\right)  = f\left( {p\left( x\right) }\right)\) 的连续函数 \(h : \Delta  \rightarrow  \Delta\) 有一个不动点，容易看出这个不动点也是给定函数 \(f : C \rightarrow  C\) 的不动点。
\end{tcolorbox}

The following lemma is about the embedding of \(C\) into \(\Delta\) .

\begin{tcolorbox}
以下引理是关于将 \(C\) 嵌入到 \(\Delta\) 中的。
\end{tcolorbox}

Lemma 7.14. Let \(C \subset  {\mathbb{R}}^{k}\) and let \(C\) be a compact and convex set of dimension \(d - 1\) . Then there is an injective affine map \(J : C \rightarrow  \Delta\) from \(C\) to the(d - 1)-dimensional unit simplex \(\Delta\) in (7.1).

\begin{tcolorbox}
引理7.14。设 \(C \subset  {\mathbb{R}}^{k}\) ，且设 \(C\) 是一个维度为 \(d - 1\) 的紧致凸集。那么存在一个从 \(C\) 到式(7.1)中 (d - 1) 维度单位单纯形 \(\Delta\) 的单射仿射映射 \(J : C \rightarrow  \Delta\) 。
\end{tcolorbox}

Proof. Because \(C\) has dimension \(d - 1\) , it has \(d\) , but no more, affinely independent points \({w}^{1},\ldots ,{w}^{d}\) . If \(d = 1\) then \(C\) is just a point, so we can assume \(d > 1\) . By Lemma 7.2, each point \(x \in  C\) is the affine combination \(x = \mathop{\sum }\limits_{{i = 1}}^{d}{w}^{i}{\lambda }_{i}\) where this equation and the equation \(\mathop{\sum }\limits_{{i = 1}}^{d}{\lambda }_{i} = 1\) uniquely determine the coefficients \({\lambda }_{1},\ldots ,{\lambda }_{d}\) . Hence, these coefficients are defined by a linear map from \({\mathbb{R}}^{k}\) (which contains \(C\) ) to \({\mathbb{R}}^{d}\) and therefore depend continuously on \(x\) . Their minimum \(\mathop{\min }\limits_{{1 \leq  i \leq  d}}{\lambda }_{i}\) is a continuous function of \(x\) and therefore assumes its minimum \(M\) on the compact set \(C\) , that is,

\begin{tcolorbox}
证明。因为 \(C\) 的维度为 \(d - 1\) ，所以它有 \(d\) 个(但不会更多)仿射无关的点 \({w}^{1},\ldots ,{w}^{d}\) 。如果 \(d = 1\) ，那么 \(C\) 只是一个点，因此我们可以假设 \(d > 1\) 。根据引理7.2，每个点 \(x \in  C\) 都是仿射组合 \(x = \mathop{\sum }\limits_{{i = 1}}^{d}{w}^{i}{\lambda }_{i}\) ，其中这个方程以及方程 \(\mathop{\sum }\limits_{{i = 1}}^{d}{\lambda }_{i} = 1\) 唯一地确定了系数 \({\lambda }_{1},\ldots ,{\lambda }_{d}\) 。因此，这些系数由一个从 \({\mathbb{R}}^{k}\) (包含 \(C\) )到 \({\mathbb{R}}^{d}\) 的线性映射定义，所以它们连续依赖于 \(x\) 。它们的最小值 \(\mathop{\min }\limits_{{1 \leq  i \leq  d}}{\lambda }_{i}\) 是 \(x\) 的连续函数，因此在紧集 \(C\) 上取得其最小值 \(M\) ，即
\end{tcolorbox}

\[
x = \mathop{\sum }\limits_{{i = 1}}^{d}{w}^{i}{\lambda }_{i},\;\mathop{\sum }\limits_{{i = 1}}^{d}{\lambda }_{i} = 1\; \Rightarrow  \;{\lambda }_{i} \geq  M\;\left( {x \in  C,1 \leq  i \leq  d}\right) . \tag{7.14}
\]

For \(x = {w}^{2}\) we have \({\lambda }_{1} = 0\) , so \(M \leq  0\) . Let \(w = \mathop{\sum }\limits_{{i = 1}}^{d}{w}^{i}\frac{1}{d}\) be the barycenter of the points \({w}^{1},\ldots ,{w}^{d}\) . Our aim is to "stretch" the points \({w}^{i}\) to new points \({v}^{i}\) (typically outside of \(C\) ) "away from \({w}^{\prime \prime }\) so that every \(x \in  C\) is a convex and not just affine combination of \({v}^{1},\ldots ,{v}^{d}\) . For some yet to be determined \(\alpha  > 0\) , let

\begin{tcolorbox}
对于 \(x = {w}^{2}\) ，我们有 \({\lambda }_{1} = 0\) ，所以 \(M \leq  0\) 。设 \(w = \mathop{\sum }\limits_{{i = 1}}^{d}{w}^{i}\frac{1}{d}\) 为点 \({w}^{1},\ldots ,{w}^{d}\) 的重心。我们的目标是将点 \({w}^{i}\) “拉伸” 为新的点 \({v}^{i}\) (通常在 \(C\) 之外)，“远离 \({w}^{\prime \prime }\) ”，使得每个 \(x \in  C\) 是 \({v}^{1},\ldots ,{v}^{d}\) 的凸组合，而不仅仅是仿射组合。对于某个有待确定的 \(\alpha  > 0\) ，设
\end{tcolorbox}

\[
{v}^{i} = w + \left( {{w}^{i} - w}\right) \alpha \;\left( {1 \leq  i \leq  d}\right)
\]

and therefore

\begin{tcolorbox}
因此
\end{tcolorbox}

\[
{w}^{i} = w + \left( {{v}^{i} - w}\right) \frac{1}{\alpha } = w\left( {1 - \frac{1}{\alpha }}\right)  + {v}^{i}\frac{1}{\alpha }.
\]

Then \(w\) is also the barycenter of the points \({v}^{1},\ldots ,{v}^{d}\) because

\begin{tcolorbox}
那么 \(w\) 也是点 \({v}^{1},\ldots ,{v}^{d}\) 的重心，因为
\end{tcolorbox}

\[
\mathop{\sum }\limits_{{i = 1}}^{d}{v}^{i}\frac{1}{d} = \mathop{\sum }\limits_{{i = 1}}^{d}\left( {w + \left( {{w}^{i} - w}\right) \alpha }\right) \frac{1}{d} = w.
\]

Hence, for \(x \in  C\) we have

\begin{tcolorbox}
因此，对于 \(x \in  C\) ，我们有
\end{tcolorbox}

\[
x = \mathop{\sum }\limits_{{i = 1}}^{d}{w}^{i}{\lambda }_{i} = \mathop{\sum }\limits_{{i = 1}}^{d}\left( {w\left( {1 - \frac{1}{\alpha }}\right)  + {v}^{i}\frac{1}{\alpha }}\right) {\lambda }_{i} = \mathop{\sum }\limits_{{i = 1}}^{d}{v}_{i}\left( {\frac{1}{d}\left( {1 - \frac{1}{\alpha }}\right)  + {\lambda }_{i}\frac{1}{\alpha }}\right)  \tag{7.15}
\]

where we want to achieve that for \(1 \leq  i \leq  d\) ,

\begin{tcolorbox}
其中我们希望对于 \(1 \leq  i \leq  d\) 能满足
\end{tcolorbox}

\[
\frac{1}{d}\left( {1 - \frac{1}{\alpha }}\right)  + {\lambda }_{i}\frac{1}{\alpha } \geq  0 \tag{7.16}
\]

which holds if \({\lambda }_{i} \geq  \frac{1 - \alpha }{d}\) and therefore if \(M \geq  \frac{1 - \alpha }{d}\) because of (7.14). For this purpose, it suffices to choose \(\alpha  = 1 - {Md}\) (note that \(M \leq  0\) ). In addition, \(\mathop{\sum }\limits_{{i = 1}}^{d}\left( {\frac{1}{d}\left( {1 - \frac{1}{\alpha }}\right)  + {\lambda }_{i}\frac{1}{\alpha }}\right)  = 1\) , so (7.16) and (7.15) imply that \(C\) is a subset of the simplex \(T\) with vertices \({v}^{1},\ldots ,{v}^{d}\) . By Lemma 7.3, there is an affine bijection \(T \rightarrow  \Delta\) , which restricted to \(C\) gives the desired affine injection \(J : C \rightarrow  \Delta\) .

\begin{tcolorbox}
如果 \({\lambda }_{i} \geq  \frac{1 - \alpha }{d}\) 成立，那么该式成立，并且由于(7.14)，如果 \(M \geq  \frac{1 - \alpha }{d}\) 成立则该式也成立。为此，选择 \(\alpha  = 1 - {Md}\) 就足够了(注意 \(M \leq  0\) )。此外， \(\mathop{\sum }\limits_{{i = 1}}^{d}\left( {\frac{1}{d}\left( {1 - \frac{1}{\alpha }}\right)  + {\lambda }_{i}\frac{1}{\alpha }}\right)  = 1\) ，所以(7.16)和(7.15)意味着 \(C\) 是顶点为 \({v}^{1},\ldots ,{v}^{d}\) 的单纯形 \(T\) 的子集。根据引理7.3，存在一个仿射双射 \(T \rightarrow  \Delta\) ，它在 \(C\) 上的限制给出了所需的仿射单射 \(J : C \rightarrow  \Delta\) 。
\end{tcolorbox}

Let \(C\) be a compact convex set and let \(D = J\left( C\right)\) be the image of \(C\) under the affine injection \(J\) in Lemma 7.14. The set \(D\) is a closed convex subset of \(\Delta\) . Next, we construct a continuous function \(p : \Delta  \rightarrow  D\) that restricted to \(D\) is the identity, that is, \(p\left( x\right)  = x\) for \(x \in  D\) . The following function \(p\) is easiest to use.

\begin{tcolorbox}
设 \(C\) 是一个紧凸集，设 \(D = J\left( C\right)\) 是 \(C\) 在引理7.14中的仿射单射 \(J\) 下的像。集合 \(D\) 是 \(\Delta\) 的一个闭凸子集。接下来，我们构造一个连续函数 \(p : \Delta  \rightarrow  D\) ，它在 \(D\) 上的限制是恒等映射，即对于 \(x \in  D\) 有 \(p\left( x\right)  = x\) 。下面的函数 \(p\) 是最容易使用的。
\end{tcolorbox}

Let \(D\) be any nonempty closed convex subset of \({\mathbb{R}}^{d}\) . The projection onto \(D\) is the function \(p : {\mathbb{R}}^{d} \rightarrow  D\) defined for \(x \in  {\mathbb{R}}^{d}\) by

\begin{tcolorbox}
设 \(D\) 为 \({\mathbb{R}}^{d}\) 的任意非空闭凸子集。到 \(D\) 上的投影是为 \(x \in  {\mathbb{R}}^{d}\) 定义的函数 \(p : {\mathbb{R}}^{d} \rightarrow  D\) ，其定义如下
\end{tcolorbox}

\[
p\left( x\right)  = \arg \mathop{\min }\limits_{{z \in  D}}\parallel x - z\parallel . \tag{7.17}
\]

Condition (7.17) means that \(p\left( x\right)\) is the point \(z\) in \(D\) that is closest in Euclidean distance to \(x\) , where clearly \(p\left( x\right)  = x\) if \(x \in  D\) . Because \(D\) is nonempty and closed and the distance is bounded from below by zero, such a closest point in \(D\) exists. We will shortly see that this point is also unique, so that \(p\left( x\right)\) is well defined. We will then show that \(p\) is a continuous function. We first need some properties of the Euclidean norm.

\begin{tcolorbox}
条件 (7.17) 意味着 \(p\left( x\right)\) 是 \(D\) 中与 \(x\) 的欧几里得距离 (Euclidean distance) 最近的点 \(z\) ，显然，如果 \(x \in  D\) ，则 \(p\left( x\right)  = x\) 。因为 \(D\) 是非空且闭的，并且距离有下界零，所以 \(D\) 中存在这样一个最近的点。我们很快会看到这个点也是唯一的，因此 \(p\left( x\right)\) 是良定义的。然后我们将证明 \(p\) 是一个连续函数。我们首先需要一些欧几里得范数 (Euclidean norm) 的性质。
\end{tcolorbox}

General vectors \(a,b \in  {\mathbb{R}}^{d}\) are called orthogonal if \({a}^{\top }b = 0\) . (This includes the case that \(a\) or \(b\) is the zero vector, which is always covered in the following observations, typically in some trivial way.) Then the triangle with corners \(a,\mathbf{0},b\) has a right angle at0, and the theorem of Pythagoras applies,

\begin{tcolorbox}
如果 \({a}^{\top }b = 0\) ，则称一般向量 \(a,b \in  {\mathbb{R}}^{d}\) 是正交的 (orthogonal)。(这包括 \(a\) 或 \(b\) 是零向量的情况，在以下的讨论中通常会以某种平凡的方式涵盖这种情况。)那么，以 \(a,\mathbf{0},b\) 为顶点的三角形在 0 处有一个直角，勾股定理适用
\end{tcolorbox}

\[
{a}^{\top }b = 0\; \Leftrightarrow  \;\parallel a - b{\parallel }^{2} = \parallel a{\parallel }^{2} + \parallel b{\parallel }^{2},
\]

\begin{center}
\includegraphics[max width=0.2\textwidth]{images/bo_d2e5cmbef24c73b03tn0_202_1161_1754_270_162_0.jpg}
\end{center}
\hspace*{3em} 

which holds because for any \(a,b\) we have \({b}^{\top }a = {a}^{\top }b \in  \mathbb{R}\) and

\begin{tcolorbox}
这是因为对于任意 \(a,b\) ，我们有 \({b}^{\top }a = {a}^{\top }b \in  \mathbb{R}\) 且
\end{tcolorbox}

\[
\parallel a - b{\parallel }^{2} = {\left( a - b\right) }^{\top }\left( {a - b}\right)  = {a}^{\top }a - {a}^{\top }b - {b}^{\top }a + {b}^{\top }b = \parallel a{\parallel }^{2} - 2{a}^{\top }b + \parallel b{\parallel }^{2}. \tag{7.18}
\]

Hence,

\begin{tcolorbox}
因此，
\end{tcolorbox}

\[
{a}^{\top }b = 0\; \Rightarrow  \;\parallel a - b\parallel  \geq  \parallel a\parallel ,\;\parallel a - b\parallel  \geq  \parallel b\parallel . \tag{7.19}
\]

The following lemma is concerned with the distance of a point \(x\) to the points on a ray (or halfline) with endpoint \(c\) through some other point \(b\) , given by the set \(\{ c + \left( {b - c}\right) \lambda  \mid  \lambda  \geq  0\}\) . As shown in Figure 7.12, this depends on the sign of \({\left( x - c\right) }^{\top }\left( {b - c}\right)\) , which indicates the angle between the vectors \(x - c\) and \(b - c\) ; the angle is acute if \({\left( x - c\right) }^{\top }\left( {b - c}\right)  > 0\) , obtuse if \({\left( x - c\right) }^{\top }\left( {b - c}\right)  < 0\) , and a right angle if \({\left( x - c\right) }^{\top }\left( {b - c}\right)  = 0\) . Exactly in the latter two cases the closest point to \(x\) on the ray is its endpoint \(c\) .

\begin{tcolorbox}
以下引理涉及点 \(x\) 到由端点 \(c\) 经过另一点 \(b\) 所确定的射线 (ray 或 halfline) 上的点的距离，该射线由集合 \(\{ c + \left( {b - c}\right) \lambda  \mid  \lambda  \geq  0\}\) 给出。如图 7.12 所示，这取决于 \({\left( x - c\right) }^{\top }\left( {b - c}\right)\) 的符号，它表示向量 \(x - c\) 和 \(b - c\) 之间的角 (angle)；如果 \({\left( x - c\right) }^{\top }\left( {b - c}\right)  > 0\) ，则该角为锐角；如果 \({\left( x - c\right) }^{\top }\left( {b - c}\right)  < 0\) ，则为钝角；如果 \({\left( x - c\right) }^{\top }\left( {b - c}\right)  = 0\) ，则为直角。恰好在后两种情况下，射线上距离 \(x\) 最近的点是其端点 \(c\) 。
\end{tcolorbox}

\begin{center}
\includegraphics[max width=0.9\textwidth]{images/bo_d2e5cmbef24c73b03tn0_203_342_703_1252_270_0.jpg}
\end{center}
\hspace*{3em} 

Figure 7.12 Illustration of Lemma 7.15 (a) (left) and (b) (center and right).

\begin{tcolorbox}
图 7.12 引理 7.15 (a)(左)和 (b)(中、右)的图示。
\end{tcolorbox}

Lemma 7.15. Let \(x,b,c \in  {\mathbb{R}}^{d}\) with \(b \neq  c\) , and \({z}_{\lambda } = c + \left( {b - c}\right) \lambda\) for \(\lambda  \in  \mathbb{R}\) . Then

\begin{tcolorbox}
引理 7.15。设 \(x,b,c \in  {\mathbb{R}}^{d}\) 满足 \(b \neq  c\) ，且对于 \(\lambda  \in  \mathbb{R}\) 有 \({z}_{\lambda } = c + \left( {b - c}\right) \lambda\) 。则
\end{tcolorbox}

(a) \({\left( x - c\right) }^{\top }\left( {b - c}\right)  > 0 \Leftrightarrow  \begin{Vmatrix}{x - {z}_{\lambda }}\end{Vmatrix} < \parallel x - c\parallel\) for all sufficiently small

\begin{tcolorbox}
(a) 对于所有足够小的
\end{tcolorbox}

positive \(\lambda\) ,

\begin{tcolorbox}
正的 \(\lambda\) ，
\end{tcolorbox}

(b) \({\left( x - c\right) }^{\top }\left( {b - c}\right)  \leq  0 \Leftrightarrow  \begin{Vmatrix}{x - {z}_{\lambda }}\end{Vmatrix} > \parallel x - c\parallel\) for all \(\lambda  > 0\) .

\begin{tcolorbox}
(b) 对于所有 \(\lambda  > 0\) ，有 \({\left( x - c\right) }^{\top }\left( {b - c}\right)  \leq  0 \Leftrightarrow  \begin{Vmatrix}{x - {z}_{\lambda }}\end{Vmatrix} > \parallel x - c\parallel\) 。
\end{tcolorbox}

Proof. By (7.18),

\begin{tcolorbox}
证明。由 (7.18)，
\end{tcolorbox}

\[
{\begin{Vmatrix}x - {z}_{\lambda }\end{Vmatrix}}^{2} = {\left( x - {z}_{\lambda }\right) }^{\top }\left( {x - {z}_{\lambda }}\right)
\]

\[
= {\left( x - c - \left( b - c\right) \lambda \right) }^{\top }\left( {x - c - \left( {b - c}\right) \lambda }\right)
\]

\[
= \parallel x - c{\parallel }^{2} - 2{\left( x - c\right) }^{\top }\left( {b - c}\right) \lambda  + \parallel b - c{\parallel }^{2}{\lambda }^{2}
\]

\[
= \parallel x - c{\parallel }^{2} - \left( {2{\left( x - c\right) }^{\top }\left( {b - c}\right)  - \parallel b - c{\parallel }^{2}\lambda }\right) \lambda .
\]

Hence, if \({\left( x - c\right) }^{\top }\left( {b - c}\right)  > 0\) , then \(\begin{Vmatrix}{x - {z}_{\lambda }}\end{Vmatrix} < \parallel x - c\parallel\) for sufficiently small positive \(\lambda\) , and if \({\left( x - c\right) }^{\top }\left( {b - c}\right)  \leq  0\) , then \(\begin{Vmatrix}{x - {z}_{\lambda }}\end{Vmatrix} > \parallel x - c\parallel\) for all \(\lambda  > 0\) . This shows " \(\Rightarrow\) " in (a) and (b). To show the converse implication " \(\Leftarrow\) " in (a), \(\begin{Vmatrix}{x - {z}_{\lambda }}\end{Vmatrix} < \parallel x - c\parallel\) for some \(\lambda  > 0\) and \({\left( x - c\right) }^{\top }\left( {b - c}\right)  \leq  0\) would contradict " \(\Rightarrow\) " in (b). Similarly, " \(\Leftarrow\) " in (b) holds because \(\begin{Vmatrix}{x - {z}_{\lambda }}\end{Vmatrix} > \parallel x - c\parallel\) for some sufficiently small \(\lambda  > 0\) and \({\left( x - c\right) }^{\top }\left( {b - c}\right)  > 0\) contradicts " \(\Rightarrow\) " in (a).

\begin{tcolorbox}
因此，如果 \({\left( x - c\right) }^{\top }\left( {b - c}\right)  > 0\) ，那么对于足够小的正 \(\lambda\) 有 \(\begin{Vmatrix}{x - {z}_{\lambda }}\end{Vmatrix} < \parallel x - c\parallel\) ；如果 \({\left( x - c\right) }^{\top }\left( {b - c}\right)  \leq  0\) ，那么对于所有 \(\lambda  > 0\) 有 \(\begin{Vmatrix}{x - {z}_{\lambda }}\end{Vmatrix} > \parallel x - c\parallel\) 。这表明了 (a) 和 (b) 中的 “ \(\Rightarrow\) ”。为了证明 (a) 中的反向蕴含关系 “ \(\Leftarrow\) ”，对于某个 \(\lambda  > 0\) 有 \(\begin{Vmatrix}{x - {z}_{\lambda }}\end{Vmatrix} < \parallel x - c\parallel\) 且 \({\left( x - c\right) }^{\top }\left( {b - c}\right)  \leq  0\) 会与 (b) 中的 “ \(\Rightarrow\) ” 矛盾。类似地，(b) 中的 “ \(\Leftarrow\) ” 成立，因为对于某个足够小的 \(\lambda  > 0\) 有 \(\begin{Vmatrix}{x - {z}_{\lambda }}\end{Vmatrix} > \parallel x - c\parallel\) 且 \({\left( x - c\right) }^{\top }\left( {b - c}\right)  > 0\) 与 (a) 中的 “ \(\Rightarrow\) ” 矛盾。
\end{tcolorbox}

Lemma 7.15 implies that for a nonempty closed convex set \(D\) , the point \(c\) in \(D\) that is closest to \(x\) is unique, so that \(p\left( x\right)  = c\) in (7.17) is well defined: Namely, if \(b\) is any other point in \(D\) and \({\left( x - c\right) }^{\top }\left( {b - c}\right)  > 0\) , then by Lemma 7.15(a) there is a point \({z}_{\lambda }\) that is closer to \(x\) than \(c\) , where for \(0 < \lambda  \leq  1\) this point belongs to \(D\) because \(D\) is convex, so \(c\) is not the closest point to \(x\) in \(D\) , a contradiction; and if \({\left( x - c\right) }^{\top }\left( {b - c}\right)  \leq  0\) , then Lemma 7.15(b) with \(\lambda  = 1\) (and thus \({z}_{\lambda } = b\) ) shows that \(b\) has greater distance \(\parallel x - b\parallel\) from \(x\) than \(c\) . Moreover, we have shown the following.

\begin{tcolorbox}
引理 7.15 表明，对于非空闭凸集 \(D\) ， \(D\) 中距离 \(x\) 最近的点 \(c\) 是唯一的，因此 (7.17) 中的 \(p\left( x\right)  = c\) 是良定义的:即，如果 \(b\) 是 \(D\) 中的任意其他点且 \({\left( x - c\right) }^{\top }\left( {b - c}\right)  > 0\) ，那么根据引理 7.15(a) ，存在一个点 \({z}_{\lambda }\) ，它比 \(c\) 更接近 \(x\) ，其中对于 \(0 < \lambda  \leq  1\) ，这个点属于 \(D\) ，因为 \(D\) 是凸集，所以 \(c\) 不是 \(D\) 中距离 \(x\) 最近的点，这是一个矛盾；如果 \({\left( x - c\right) }^{\top }\left( {b - c}\right)  \leq  0\) ，那么结合 \(\lambda  = 1\) (从而 \({z}_{\lambda } = b\) )应用引理 7.15(b) 表明， \(b\) 到 \(x\) 的距离 \(\parallel x - b\parallel\) 比 \(c\) 到 \(x\) 的距离更大。此外，我们已经证明了以下内容。
\end{tcolorbox}

Corollary 7.16. Let \(D\) be a nonempty closed convex subset of \({\mathbb{R}}^{d}\) and \(x \in  {\mathbb{R}}^{d}\) and let \(c = p\left( x\right)\) be the projection of \(x\) onto \(D\) . Then \({\left( x - c\right) }^{\top }\left( {b - c}\right)  \leq  0\) for all \(b \in  D\) .

\begin{tcolorbox}
推论 7.16。设 \(D\) 是 \({\mathbb{R}}^{d}\) 的非空闭凸子集， \(x \in  {\mathbb{R}}^{d}\) ，并且设 \(c = p\left( x\right)\) 是 \(x\) 在 \(D\) 上的投影。那么对于所有 \(b \in  D\) 有 \({\left( x - c\right) }^{\top }\left( {b - c}\right)  \leq  0\) 。
\end{tcolorbox}

\begin{center}
\includegraphics[max width=0.5\textwidth]{images/bo_d2e5cmbef24c73b03tn0_204_204_604_741_382_0.jpg}
\end{center}
\hspace*{3em} 

Figure 7.13 For \(x \notin  D\) in Corollary 7.16, the hyperplane \(H\) with normal vector \(v = x - c\) and \(\alpha  = {\left( x - c\right) }^{\top }c\) separates \(x\) from \(D\) . (The set \(D\) can be unbounded.)

\begin{tcolorbox}
图 7.13 对于推论 7.16 中的 \(x \notin  D\) ，具有法向量 \(v = x - c\) 且 \(\alpha  = {\left( x - c\right) }^{\top }c\) 的超平面 \(H\) 将 \(x\) 与 \(D\) 分隔开。(集合 \(D\) 可以是无界的。)
\end{tcolorbox}

Corollary 7.16 is also known as the theorem of the separating hyperplane. A hyperplane \(H\) in \({\mathbb{R}}^{d}\) is defined as \(H = \left\{  {z \in  {\mathbb{R}}^{d} \mid  {v}^{\top }z = \alpha }\right\}\) for a nonzero normal vector \(v\) in \({\mathbb{R}}^{d}\) and a real number \(\alpha\) . The corresponding halfspace is defined as \(\left\{  {z \in  {\mathbb{R}}^{d} \mid  {v}^{\top }z \leq  \alpha }\right\}\) . Suppose that in Corollary 7.16 we have \(x \notin  D\) and therefore \(x \neq  c\) and \({\left( x - c\right) }^{\top }\left( {x - c}\right)  = \parallel x - c{\parallel }^{2} > 0\) . With \(v = x - c\) and \(\alpha  = {\left( x - c\right) }^{\top }c\) the corollary states that \(D\) is a subset of the halfspace defined by \(H\) , whereas \(x\) is clearly not in the halfspace and therefore "separated" from \(D\) by \(H\) . See also Figure 7.13. The separating hyperplane theorem is a very useful tool, but one has to ensure that its assumptions hold, in particular that the set \(D\) is closed.

\begin{tcolorbox}
推论7.16也被称为分离超平面定理。在 \({\mathbb{R}}^{d}\) 中的一个超平面 \(H\) 定义为 \(H = \left\{  {z \in  {\mathbb{R}}^{d} \mid  {v}^{\top }z = \alpha }\right\}\) ，其中 \(v\) 是 \({\mathbb{R}}^{d}\) 中的非零法向量， \(\alpha\) 是实数。相应的半空间定义为 \(\left\{  {z \in  {\mathbb{R}}^{d} \mid  {v}^{\top }z \leq  \alpha }\right\}\) 。假设在推论7.16中我们有 \(x \notin  D\) ，因此 \(x \neq  c\) 和 \({\left( x - c\right) }^{\top }\left( {x - c}\right)  = \parallel x - c{\parallel }^{2} > 0\) 。对于 \(v = x - c\) 和 \(\alpha  = {\left( x - c\right) }^{\top }c\) ，该推论表明 \(D\) 是由 \(H\) 定义的半空间的一个子集，而 \(x\) 显然不在这个半空间中，因此被 \(H\) 从 \(D\) 中“分离”出来。另见 图7.13。分离超平面定理是一个非常有用的工具，但必须确保其假设成立，特别是集合 \(D\) 是闭集。
\end{tcolorbox}

A second easy consequence of Lemma 7.15 is that the shortest connecting line segment of a point \(x\) to a line \(L\) is orthogonal to that line.

\begin{tcolorbox}
引理7.15的第二个简单推论是，点 \(x\) 到直线 \(L\) 的最短连接线段与该直线正交。
\end{tcolorbox}

\begin{center}
\includegraphics[max width=0.3\textwidth]{images/bo_d2e5cmbef24c73b03tn0_204_1067_1486_385_202_0.jpg}
\end{center}
\hspace*{3em} 

Lemma 7.17. Let \(L\) be a line in \({\mathbb{R}}^{d}\) , that is, the set of affine combinations of two points in \({\mathbb{R}}^{d}\) , let \(x \in  {\mathbb{R}}^{d}\) , and let \(c = p\left( x\right)\) be the projection of \(x\) onto \(L\) . Then \({\left( x - c\right) }^{\top }\left( {b - c}\right)  = 0\) for all \(b \in  L\) .

\begin{tcolorbox}
引理7.17。设 \(L\) 是 \({\mathbb{R}}^{d}\) 中的一条直线，即 \({\mathbb{R}}^{d}\) 中两个点的仿射组合的集合，设 \(x \in  {\mathbb{R}}^{d}\) ，并设 \(c = p\left( x\right)\) 是 \(x\) 在 \(L\) 上的投影。那么对于所有 \(b \in  L\) 有 \({\left( x - c\right) }^{\top }\left( {b - c}\right)  = 0\) 。
\end{tcolorbox}

Proof. Let \(b\) be any point on \(L\) other than \(c\) . Then \(L = \{ c + \left( {b - c}\right) \lambda  \mid  \lambda  \in  \mathbb{R}\}\) . If \({\left( x - c\right) }^{\top }\left( {b - c}\right)  > 0\) , then \(c\) is not the closest point to \(x\) by Lemma 7.15(a). If \({\left( x - c\right) }^{\top }\left( {b - c}\right)  < 0\) , let \({b}^{\prime } = c + \left( {b - c}\right) \left( {-1}\right)\) so that \({b}^{\prime } - c =  - \left( {b - c}\right)\) and replace \(b\) by \({b}^{\prime }\) to observe the same, again by Lemma 7.15(a). Hence, \({\left( x - c\right) }^{\top }\left( {b - c}\right)  = 0\) as claimed.

\begin{tcolorbox}
证明。设 \(b\) 是 \(L\) 上除 \(c\) 之外的任意一点。那么 \(L = \{ c + \left( {b - c}\right) \lambda  \mid  \lambda  \in  \mathbb{R}\}\) 。如果 \({\left( x - c\right) }^{\top }\left( {b - c}\right)  > 0\) ，那么根据引理7.15(a)， \(c\) 不是到 \(x\) 的最近点。如果 \({\left( x - c\right) }^{\top }\left( {b - c}\right)  < 0\) ，设 \({b}^{\prime } = c + \left( {b - c}\right) \left( {-1}\right)\) 使得 \({b}^{\prime } - c =  - \left( {b - c}\right)\) ，并将 \(b\) 替换为 \({b}^{\prime }\) 来观察同样的情况，同样根据引理7.15(a)。因此，如所声称的那样有 \({\left( x - c\right) }^{\top }\left( {b - c}\right)  = 0\) 。
\end{tcolorbox}

Our goal is to show that the projection map \(p\) in (7.17) is continuous. This is a consequence of the following theorem.

\begin{tcolorbox}
我们的目标是证明(7.17)中的投影映射 \(p\) 是连续的。这是以下定理的一个推论。
\end{tcolorbox}

Theorem 7.18. Let \(D\) be a nonempty closed convex subset of \({\mathbb{R}}^{d}\) , and let \(p : {\mathbb{R}}^{d} \rightarrow  D\) be the projection onto \(D\) defined in (7.17). Then for all \(x,y \in  {\mathbb{R}}^{d}\)

\begin{tcolorbox}
定理7.18。设 \(D\) 是 \({\mathbb{R}}^{d}\) 的一个非空闭凸子集，并设 \(p : {\mathbb{R}}^{d} \rightarrow  D\) 是(7.17)中定义的到 \(D\) 上的投影。那么对于所有 \(x,y \in  {\mathbb{R}}^{d}\)
\end{tcolorbox}

\[
\parallel p\left( x\right)  - p\left( y\right) \parallel  \leq  \parallel x - y\parallel . \tag{7.20}
\]

\begin{center}
\includegraphics[max width=0.5\textwidth]{images/bo_d2e5cmbef24c73b03tn0_205_346_450_714_380_0.jpg}
\end{center}
\hspace*{3em} 

Figure 7.14 Projections \(r\left( x\right)\) and \(r\left( y\right)\) onto a line \(L\) .

\begin{tcolorbox}
图7.14 到直线 \(L\) 上的投影 \(r\left( x\right)\) 和 \(r\left( y\right)\) 。
\end{tcolorbox}

Proof. Let \(x,y \in  {\mathbb{R}}^{d}\) . We first show the claim when \(D\) is a line \(L\) , that is,

\begin{tcolorbox}
证明。设 \(x,y \in  {\mathbb{R}}^{d}\) 。我们首先证明当 \(D\) 是一条直线 \(L\) 时的结论，即
\end{tcolorbox}

\[
\parallel r\left( x\right)  - r\left( y\right) \parallel  \leq  \parallel x - y\parallel  \tag{7.21}
\]

where \(r\left( x\right)\) and \(r\left( y\right)\) denote the projections of \(x\) and \(y\) onto \(L\) . The claim is apparent from Figure 7.14, where we consider the point \(z = r\left( x\right)  + y - r\left( y\right)\) , so that the points \(z,y,r\left( y\right) ,r\left( x\right)\) form a rectangle where \(\parallel z - y\parallel  = \parallel r\left( x\right)  - r\left( y\right) \parallel\) , and the points \(x,z,y\) are three corners of a rectangle whose side length \(\parallel z - y\parallel\) is not longer than the diagonal length \(\parallel x - y\parallel\) . Formally,

\begin{tcolorbox}
其中 \(r\left( x\right)\) 和 \(r\left( y\right)\) 分别表示 \(x\) 和 \(y\) 在 \(L\) 上的投影。这一结论从图7.14中显而易见，在图中我们考虑点 \(z = r\left( x\right)  + y - r\left( y\right)\) ，使得点 \(z,y,r\left( y\right) ,r\left( x\right)\) 构成一个矩形，其中 \(\parallel z - y\parallel  = \parallel r\left( x\right)  - r\left( y\right) \parallel\) ，并且点 \(x,z,y\) 是一个矩形的三个顶点，其边长 \(\parallel z - y\parallel\) 不大于对角线长度 \(\parallel x - y\parallel\) 。形式上，
\end{tcolorbox}

\[
{\left( x - z\right) }^{\top }\left( {y - z}\right)
\]

\[
= {\left( x - r\left( x\right)  - y + r\left( y\right) \right) }^{\top }\left( {r\left( y\right)  - r\left( x\right) }\right)
\]

\[
= {\left( x - r\left( x\right) \right) }^{\top }\left( {r\left( y\right)  - r\left( x\right) }\right)  - {\left( y - r\left( y\right) \right) }^{\top }\left( {r\left( x\right)  - r\left( y\right) }\right)
\]

\[
= 0 - 0
\]

where we use Lemma 7.17 twice, with \(r\left( x\right)\) as the closest point to \(x\) on \(L\) and \(r\left( y\right)\) as the closest point to \(y\) on \(L\) . Then (7.19) implies \(\parallel x - y\parallel  = \parallel x - z - \left( {y - z}\right) \parallel  \geq\)  \(\parallel y - z\parallel  = \parallel y - \left( {r\left( x\right)  + y - r\left( y\right) }\right) \parallel  = \parallel r\left( y\right)  - r\left( x\right) \parallel  = \parallel r\left( x\right)  - r\left( y\right) \parallel\) , which shows (7.21).

\begin{tcolorbox}
这里我们两次使用引理7.17，其中 \(r\left( x\right)\) 是 \(L\) 上距离 \(x\) 最近的点， \(r\left( y\right)\) 是 \(L\) 上距离 \(y\) 最近的点。那么(7.19)式意味着 \(\parallel x - y\parallel  = \parallel x - z - \left( {y - z}\right) \parallel  \geq\)  \(\parallel y - z\parallel  = \parallel y - \left( {r\left( x\right)  + y - r\left( y\right) }\right) \parallel  = \parallel r\left( y\right)  - r\left( x\right) \parallel  = \parallel r\left( x\right)  - r\left( y\right) \parallel\) ，这就证明了(7.21)式。
\end{tcolorbox}

\begin{center}
\includegraphics[max width=0.5\textwidth]{images/bo_d2e5cmbef24c73b03tn0_205_343_1754_685_359_0.jpg}
\end{center}
\hspace*{3em} 

Figure 7.15 Projections \(p\left( x\right)\) and \(p\left( y\right)\) of \(x\) and \(y\) onto a convex set \(D\) , with the line \(L\) through \(p\left( x\right)\) and \(p\left( y\right)\) . The angle between \(x,p\left( x\right)\) , and \(p\left( y\right)\) cannot be acute, so \(r\left( x\right)\) is to the left of (or equal to) \(p\left( x\right)\) . Similarly, \(r\left( y\right)\) is to the right of (or equal to) \(p\left( y\right)\) .

\begin{tcolorbox}
图7.15 \(x\) 和 \(y\) 在凸集 \(D\) 上的投影 \(p\left( x\right)\) 和 \(p\left( y\right)\) ，以及过 \(p\left( x\right)\) 和 \(p\left( y\right)\) 的直线 \(L\) 。 \(x,p\left( x\right)\) 与 \(p\left( y\right)\) 之间的角不可能是锐角，所以 \(r\left( x\right)\) 在 \(p\left( x\right)\) 的左侧(或与之重合)。类似地， \(r\left( y\right)\) 在 \(p\left( y\right)\) 的右侧(或与之重合)。
\end{tcolorbox}

Secondly, consider now a general nonempty closed convex set \(D\) , shown in Figure 7.15 for an example in dimension two. Let \(p\left( x\right)\) and \(p\left( y\right)\) be the projections of \(x\) and \(y\) onto \(D\) , where we can assume \(p\left( x\right)  \neq  p\left( y\right)\) because otherwise (7.20) holds trivially. Let \(L\) be the line through \(p\left( x\right)\) and \(p\left( y\right)\) , and let \(r\left( x\right)\) and \(r\left( y\right)\) be the projections of \(x\) and \(y\) onto \(L\) .

\begin{tcolorbox}
其次，现在考虑一个一般的非空闭凸集 \(D\) ，图7.15给出了二维情形下的一个示例。设 \(p\left( x\right)\) 和 \(p\left( y\right)\) 分别是 \(x\) 和 \(y\) 在 \(D\) 上的投影，这里我们可以假设 \(p\left( x\right)  \neq  p\left( y\right)\) ，因为否则(7.20)式显然成立。设 \(L\) 是过 \(p\left( x\right)\) 和 \(p\left( y\right)\) 的直线，设 \(r\left( x\right)\) 和 \(r\left( y\right)\) 分别是 \(x\) 和 \(y\) 在 \(L\) 上的投影。
\end{tcolorbox}

The line segment that connects \(p\left( x\right)\) and \(p\left( y\right)\) is a subset of \(D\) because \(D\) is convex. We show that it is a subset of the line segment that connects \(r\left( x\right)\) and \(r\left( y\right)\) , and therefore at most as long. Let \(c = p\left( x\right)\) and \(b = p\left( y\right)\) . By Corollary 7.16, \({\left( x - c\right) }^{\top }\left( {b - c}\right)  \leq  0\) , that is, the angle at \(c\) shown in Figure 7.15 is a right or obtuse angle. Let \(r\left( x\right)  = p\left( x\right)  + \left( {p\left( y\right)  - p\left( x\right) }\right) \lambda  = {z}_{\lambda }\) . Then \(\lambda  \leq  0\) , because if \(\lambda  > 0\) then Lemma 7.15(b) would imply \(\parallel x - r\left( x\right) \parallel  > \parallel x - p\left( x\right) \parallel\) , which is a contradiction because \(p\left( x\right)  \in  L\) and \(r\left( x\right)\) is the closest point to \(x\) on \(L\) . Similarly, \(r\left( y\right)  = p\left( y\right)  + \left( {p\left( x\right)  - p\left( y\right) }\right) \mu\) where \(\mu  \leq  0\) . Then

\begin{tcolorbox}
连接 \(p\left( x\right)\) 和 \(p\left( y\right)\) 的线段(line segment)是 \(D\) 的子集，因为 \(D\) 是凸集。我们证明它是连接 \(r\left( x\right)\) 和 \(r\left( y\right)\) 的线段的子集，因此长度至多相同。设 \(c = p\left( x\right)\) 且 \(b = p\left( y\right)\) 。根据推论7.16， \({\left( x - c\right) }^{\top }\left( {b - c}\right)  \leq  0\) ，即图7.15中 \(c\) 处的角(angle)是直角或钝角。设 \(r\left( x\right)  = p\left( x\right)  + \left( {p\left( y\right)  - p\left( x\right) }\right) \lambda  = {z}_{\lambda }\) 。那么 \(\lambda  \leq  0\) ，因为如果 \(\lambda  > 0\) ，则引理7.15(b)将意味着 \(\parallel x - r\left( x\right) \parallel  > \parallel x - p\left( x\right) \parallel\) ，这是矛盾的，因为 \(p\left( x\right)  \in  L\) 且 \(r\left( x\right)\) 是 \(L\) 上距离 \(x\) 最近的点。类似地， \(r\left( y\right)  = p\left( y\right)  + \left( {p\left( x\right)  - p\left( y\right) }\right) \mu\) ，其中 \(\mu  \leq  0\) 。那么
\end{tcolorbox}

\[
\parallel r\left( x\right)  - r\left( y\right) \parallel  = \parallel p\left( x\right)  + \left( {p\left( y\right)  - p\left( x\right) }\right) \lambda  - p\left( y\right)  - \left( {p\left( x\right)  - p\left( y\right) }\right) \mu \parallel
\]

\[
= \parallel \left( {p\left( x\right)  - p\left( y\right) }\right) \left( {1 - \lambda  - \mu }\right) \parallel
\]

\[
= \begin{Vmatrix}\left( {p\left( x\right)  - p\left( y\right) }\right) \end{Vmatrix}\left( {1 - \lambda  - \mu }\right)
\]

\[
\geq  \parallel \left( {p\left( x\right)  - p\left( y\right) }\right) \parallel
\]

because \(\lambda  \leq  0\) and \(\mu  \leq  0\) , which with (7.21) implies (7.20).

\begin{tcolorbox}
因为 \(\lambda  \leq  0\) 和 \(\mu  \leq  0\) ，结合(7.21)可推出(7.20)。
\end{tcolorbox}

Theorem 7.18 states that the projection \(p\) onto \(D\) is "contractive", and therefore indeed continuous:

\begin{tcolorbox}
定理7.18表明，到 \(p\) 上的投影 \(D\) 是“压缩的”，因此确实是连续的:
\end{tcolorbox}

Corollary 7.19. Let \(D\) be a nonempty closed convex subset of \({\mathbb{R}}^{d}\) . Then the projection map \(p\) of \({\mathbb{R}}^{d}\) onto \(D\) in (7.17) is continuous.

\begin{tcolorbox}
推论7.19。设 \(D\) 是 \({\mathbb{R}}^{d}\) 的一个非空闭凸子集。那么，(7.17) 中从 \({\mathbb{R}}^{d}\) 到 \(D\) 的投影映射 \(p\) 是连续的。
\end{tcolorbox}

Proof. Continuity of \(p : {\mathbb{R}}^{d} \rightarrow  D\) means that for all \(x,y \in  {\mathbb{R}}^{d}\)

\begin{tcolorbox}
证明。 \(p : {\mathbb{R}}^{d} \rightarrow  D\) 的连续性意味着对于所有 \(x,y \in  {\mathbb{R}}^{d}\)
\end{tcolorbox}

\[
\forall \varepsilon  > 0\exists \delta  > 0 : \parallel x - y\parallel  < \delta  \Rightarrow  \parallel p\left( x\right)  - p\left( y\right) \parallel  < \varepsilon .
\]

By Theorem 7.18, it suffices to take \(\delta  = \varepsilon\) .

\begin{tcolorbox}
根据定理7.18，取 \(\delta  = \varepsilon\) 就足够了。
\end{tcolorbox}

The convexity of the set \(D\) in Corollary 7.19 is essential. For example, if \(D\) is just a set of two distinct points \(a\) and \(b\) , then the projection \(p\) onto \(D\) is not well defined for points that are equidistant to \(a\) and \(b\) , and at any rate "jumps" from \(a\) to \(b\) near these points. Moreover, the function \(D \rightarrow  D\) that exchanges \(a\) and \(b\) is trivially continuous but has no fixed points, which it would have if \(p\) was continuous, according to the proof of the following theorem. This theorem is Brouwer's fixed-point theorem in full generality.

\begin{tcolorbox}
推论7.19中集合 \(D\) 的凸性是至关重要的。例如，如果 \(D\) 只是由两个不同的点 \(a\) 和 \(b\) 组成的集合，那么对于与 \(a\) 和 \(b\) 等距的点，到 \(D\) 上的投影 \(p\) 是没有明确定义的，并且无论如何，在这些点附近会从 \(a\) “跳跃”到 \(b\) 。此外，交换 \(a\) 和 \(b\) 的函数 \(D \rightarrow  D\) 显然是连续的，但没有不动点。根据以下定理的证明，如果 \(p\) 是连续的，它应该有不动点。这个定理就是一般形式的布劳威尔不动点定理。
\end{tcolorbox}

Theorem 7.20. Let \(C\) be a nonempty compact and convex subset of \({\mathbb{R}}^{k}\) , and let \(f : C \rightarrow  C\) be continuous. Then \(f\) has a fixed point \(z \in  C\) , that is, \(f\left( z\right)  = z\) . Proof. Suppose \(C\) has dimension \(d - 1\) . By Lemma 7.14, there is an affine injective map \(J : C \rightarrow  \Delta\) with image \(D = J\left( C\right)  = \{ J\left( z\right)  \mid  z \in  C\}  \subseteq  \Delta\) , and its inverse function \({J}^{-1} : D \rightarrow  C\) is a linear bijection. The set \(D\) is compact, and convex because \(J\) preserves convex combinations. Define the function \(g : D \rightarrow  D\) by

\begin{tcolorbox}
定理7.20。设 \(C\) 是 \({\mathbb{R}}^{k}\) 的一个非空紧致凸子集，并且 \(f : C \rightarrow  C\) 是连续的。那么 \(f\) 有一个不动点 \(z \in  C\) ，即 \(f\left( z\right)  = z\) 。证明。假设 \(C\) 的维度为 \(d - 1\) 。根据引理7.14，存在一个仿射单射 \(J : C \rightarrow  \Delta\) ，其像为 \(D = J\left( C\right)  = \{ J\left( z\right)  \mid  z \in  C\}  \subseteq  \Delta\) ，并且它的反函数 \({J}^{-1} : D \rightarrow  C\) 是一个线性双射。集合 \(D\) 是紧致的，并且是凸的，因为 \(J\) 保持凸组合。定义函数 \(g : D \rightarrow  D\) 为
\end{tcolorbox}

\[
g\left( {J\left( z\right) }\right)  = J\left( {f\left( z\right) }\right) , \tag{7.22}
\]

which is unambiguous because \(J\) is injective, and which defines \(g\) on all of \(D\) because \(D = J\left( C\right)\) . Then for all \(z \in  C\)

\begin{tcolorbox}
这是明确的，因为 \(J\) 是单射，并且由于 \(D = J\left( C\right)\) ，它在整个 \(D\) 上定义了 \(g\) 。那么对于所有 \(z \in  C\)
\end{tcolorbox}

\[
f\left( z\right)  = z\; \Leftrightarrow  \;g\left( {J\left( z\right) }\right)  = J\left( z\right)  \tag{7.23}
\]

where " \(\Rightarrow\) " is immediate from (7.22), and " \(\Leftarrow\) " holds because \(g\left( {J\left( z\right) }\right)  = J\left( z\right)\) implies

\begin{tcolorbox}
其中“ \(\Rightarrow\) ” 直接由 (7.22) 得出，而 “ \(\Leftarrow\) ” 成立是因为 \(g\left( {J\left( z\right) }\right)  = J\left( z\right)\) 意味着
\end{tcolorbox}

\[
f\left( z\right)  = {J}^{-1}\left( {J\left( {f\left( z\right) }\right) }\right)  = {J}^{-1}\left( {g\left( {J\left( z\right) }\right) }\right)  = {J}^{-1}\left( {J\left( z\right) }\right)  = z.
\]

Condition (7.23) states that the fixed points of \(f\) are in one-to-one correspondence with the fixed points of \(g\) via the bijection \(J\) .

\begin{tcolorbox}
条件 (7.23) 表明， \(f\) 的不动点通过双射 \(J\) 与 \(g\) 的不动点一一对应。
\end{tcolorbox}

Let \(p : \Delta  \rightarrow  D\) be the projection map onto \(D\) as defined in (7.17), which is continuous by Corollary 7.19 and fulfills \(p\left( x\right)  = x\) for all \(x \in  D\) . Define the function

\begin{tcolorbox}
设 \(p : \Delta  \rightarrow  D\) 是如 (7.17) 中所定义的到 \(D\) 上的投影映射，根据推论7.19，它是连续的，并且对于所有 \(x \in  D\) 满足 \(p\left( x\right)  = x\) 。定义函数
\end{tcolorbox}

\[
h : \Delta  \rightarrow  \Delta ,\;h\left( x\right)  = g\left( {p\left( x\right) }\right) \;\left( {x \in  \Delta }\right) ,
\]

which is a continuous function on the unit simplex \(\Delta\) because \(p\) and \(g\) are continuous. By Theorem 7.13, \(h\) has a fixed point \(x \in  \Delta\) , that is, \(h\left( x\right)  = g\left( {p\left( x\right) }\right)  = x\) , where \(h\left( x\right)  \in  D\) because \(g\) takes only values in \(D\) , and therefore \(p\left( x\right)  = x\) , which implies \(g\left( x\right)  = x\) . That is, the fixed points of \(h\) are exactly the fixed points of \(g\) . By (7.23), the desired fixed point of \(f\) is given by \(z = {J}^{-1}\left( x\right)\) .

\begin{tcolorbox}
这是单位单纯形 \(\Delta\) 上的一个连续函数，因为 \(p\) 和 \(g\) 是连续的。根据定理7.13， \(h\) 有一个不动点 \(x \in  \Delta\) ，即 \(h\left( x\right)  = g\left( {p\left( x\right) }\right)  = x\) ，其中 \(h\left( x\right)  \in  D\) 是因为 \(g\) 只在 \(D\) 中取值，因此 \(p\left( x\right)  = x\) ，这意味着 \(g\left( x\right)  = x\) 。也就是说， \(h\) 的不动点恰好是 \(g\) 的不动点。根据 (7.23)， \(f\) 的所需不动点由 \(z = {J}^{-1}\left( x\right)\) 给出。
\end{tcolorbox}

\section*{7.7 The Freudenthal Triangulation}

\begin{tcolorbox}
\section*{7.7 弗赖登塔尔三角剖分(Freudenthal Triangulation)}
\end{tcolorbox}

This section provides the missing piece in the proof of Brouwer's fixed-point theorem, namely an arbitrarily fine triangulation of a simplex, which is used to prove the KKM lemma (Theorem 7.12). As mentioned at the end of Section 7.3, this is only obvious when the simplex has dimension one or two. In the following we describe an elegant triangulation, due to Freudenthal (1942), that works in any dimension. Its construction rests on two main ideas:

\begin{tcolorbox}
本节补充了布劳威尔不动点定理证明中缺失的部分，即单纯形(simplex)的任意精细三角剖分(triangulation)，该三角剖分用于证明KKM引理(定理7.12)。如7.3节末尾所述，仅当单纯形的维度(dimension)为一维或二维时，这一点才显而易见。下面我们将介绍一种由弗赖登塔尔(Freudenthal，1942)提出的精妙三角剖分方法，该方法适用于任意维度。其构造基于两个主要思想:
\end{tcolorbox}

\begin{itemize}
\item The \(d\) -dimensional unit cube
\end{itemize}

\begin{tcolorbox}
\begin{itemize}
\item \(d\) 维单位立方体(unit cube)
\end{itemize}
\end{tcolorbox}

\[
{C}_{d} = {\left\lbrack  0,1\right\rbrack  }^{d} = \left\{  {x \in  {\mathbb{R}}^{d} \mid  0 \leq  {x}_{i} \leq  1,1 \leq  i \leq  d}\right\}
\]

can be easily subdivided into small "cubelets".

\begin{tcolorbox}
可以很容易地细分为小“立方体块”。
\end{tcolorbox}

\begin{itemize}
\item The cube \({C}_{d}\) has a canonical triangulation into \(d\) ! simplices that can in the same way be applied to the cubelets, which gives then a triangulation of one of the simplices in the triangulation of \({C}_{d}\) .
\end{itemize}

\begin{tcolorbox}
\begin{itemize}
\item 立方体 \({C}_{d}\) 有一个规范的三角剖分，可将其划分为 \(d\) ! 个单纯形，同样的方法也可应用于这些立方体块，从而得到 \({C}_{d}\) 三角剖分中某个单纯形的三角剖分。
\end{itemize}
\end{tcolorbox}

We continue to use \(d\) has the dimension of the space \({\mathbb{R}}^{d}\) of the points that we consider. The unit simplex \(\Delta\) in (7.1) is a subset of \({\mathbb{R}}^{d}\) but has dimension \(d - 1\) . For the Freudenthal triangulation, we use a less symmetric simplex \({T}_{d}\) than \(\Delta\) , which does not matter because any two simplices of a given dimension are in an affine bijection via their barycentric coordinates by Lemma 7.3. The simplex \({T}_{d}\) has \(d + 1\) vertices in \({\mathbb{R}}^{d}\) and has therefore dimension \(d\) , which by Lemma 7.2 is maximal. (A polytope with this property is also called full-dimensional.)

\begin{tcolorbox}
我们继续用 \(d\) 表示所考虑点所在空间 \({\mathbb{R}}^{d}\) 的维度。式(7.1)中的单位单纯形(unit simplex) \(\Delta\) 是 \({\mathbb{R}}^{d}\) 的一个子集，但其维度为 \(d - 1\) 。对于弗赖登塔尔三角剖分(Freudenthal triangulation)，我们使用一个比 \(\Delta\) 对称性更低的单纯形 \({T}_{d}\) ，这并不影响结果，因为根据引理7.3，给定维度的任意两个单纯形通过它们的重心坐标(barycentric coordinates)存在仿射双射。单纯形 \({T}_{d}\) 在 \({\mathbb{R}}^{d}\) 中有 \(d + 1\) 个顶点，因此其维度为 \(d\) ，根据引理7.2，这是最大维度。(具有此性质的有界多面体(polytope)也称为满维(full-dimensional)多面体。)
\end{tcolorbox}

The simplex \({T}_{d}\) is the convex hull of certain vertices of the unit cube \({C}_{d}\) . For example, if \(d = 3\) , then \({T}_{d}\) is the convex hull of the vectors

\begin{tcolorbox}
单纯形 \({T}_{d}\) 是单位立方体 \({C}_{d}\) 某些顶点的凸包(convex hull)。例如，如果 \(d = 3\) ，那么 \({T}_{d}\) 是向量的凸包
\end{tcolorbox}

\[
\left( \begin{array}{l} 0 \\  0 \\  0 \end{array}\right) ,\left( \begin{array}{l} 1 \\  0 \\  0 \end{array}\right) ,\left( \begin{array}{l} 1 \\  1 \\  0 \end{array}\right) ,\left( \begin{array}{l} 1 \\  1 \\  1 \end{array}\right)  \tag{7.24}
\]

The first of these vectors is0in \({\mathbb{R}}^{d}\) , and the next vector is always obtained by changing one component from 0 to 1, where the last vector is the all-one vector1 in \({\mathbb{R}}^{d}\) . We then generalize \({T}_{d}\) in two ways:

\begin{tcolorbox}
这些向量中的第一个是 \({\mathbb{R}}^{d}\) 中的0向量，后续向量总是通过将一个分量从0变为1得到，最后一个向量是 \({\mathbb{R}}^{d}\) 中的全1向量1。然后我们从两个方面对 \({T}_{d}\) 进行推广:
\end{tcolorbox}

(a) The different orders in which the components are changed from 0 to 1 generate \(d\) ! different simplices, which together define a triangulation of the unit cube.

\begin{tcolorbox}
(a) 分量从0变为1的不同顺序会生成 \(d\) ! 个不同的单纯形，它们共同构成了单位立方体的一个三角剖分。
\end{tcolorbox}

(b) In a scaled-up version of \({T}_{d}\) which we call \({T}_{d}\left( k\right)\) , each entry 1 in (7.24) is replaced by a positive integer \(k\) . We then consider all points with integer coordinates in \({T}_{d}\left( k\right)\) , and shifted unit cubes with these integer points as vertices. The subdivision of each such cube into simplices as described in (a), as far as they are subsets of \({T}_{d}\left( k\right)\) , then creates the desired triangulation of \({T}_{d}\left( k\right)\) .

\begin{tcolorbox}
(b) 在 \({T}_{d}\) 的放大版本(我们称之为 \({T}_{d}\left( k\right)\) )中，式(7.24)中的每个1都被一个正整数 \(k\) 取代。然后我们考虑 \({T}_{d}\left( k\right)\) 中所有具有整数坐标的点，并以这些整数点为顶点构造平移后的单位立方体。如(a)中所述，将每个这样的立方体细分为单纯形，只要它们是 \({T}_{d}\left( k\right)\) 的子集，就可以得到 \({T}_{d}\left( k\right)\) 所需的三角剖分。
\end{tcolorbox}

The construction in (b) exploits the fact that a cube with edge length \(k\) is easily subdivided into cubes of side length 1, which in turn are then subdivided into sub-simplices as in (a). Relative to the scaled-up simplex \({T}_{d}\left( k\right)\) , the distances of any two points in a sub-simplex are \(\frac{1}{k}\) of the distances in \({T}_{d}\left( k\right)\) , so that by shrinking \({T}_{d}\left( k\right)\) back to \({T}_{d}\) one obtains an arbitrarily fine triangulation of \({T}_{d}\) by choosing \(k\) large enough. (It is easier to deal with integer coordinates than with multiples of \(\left. {\frac{1}{k}\text{ . }}\right)\)

\begin{tcolorbox}
(b)中的构造利用了这样一个事实:边长为 \(k\) 的立方体很容易细分为边长为1的立方体，然后这些小立方体又可以像(a)中那样细分为子单纯形(sub - simplex)。相对于放大后的单纯形 \({T}_{d}\left( k\right)\) ，子单纯形中任意两点之间的距离是 \({T}_{d}\left( k\right)\) 中对应距离的 \(\frac{1}{k}\) ，因此通过将 \({T}_{d}\left( k\right)\) 缩小回 \({T}_{d}\) ，只要选择足够大的 \(k\) ，就可以得到 \({T}_{d}\) 的任意精细三角剖分。(处理整数坐标比处理 \(\left. {\frac{1}{k}\text{ . }}\right)\) 的倍数更容易。
\end{tcolorbox}

Figure 7.16 shows this construction for a triangle \({T}_{d}\left( k\right)\) where \(d = 2\) and \(k = 8\) , which is essentially a "shifted" triangulation of the symmetric triangulation in Figure 7.5.

\begin{tcolorbox}
图7.16展示了三角形 \({T}_{d}\left( k\right)\) 的这种构造，其中 \(d = 2\) 且 \(k = 8\) ，这本质上是图7.5中对称三角剖分的“平移”版本。
\end{tcolorbox}

The general construction of the Freudenthal triangulation is as follows. Fix a dimension \(d\) . Let \(k\) be a positive integer. Consider the \(d + 1\) vectors in \({\mathbb{R}}^{d}\)

\begin{tcolorbox}
弗赖登塔尔三角剖分(Freudenthal triangulation)的一般构造如下。固定一个维度 \(d\) 。设 \(k\) 为正整数。考虑 \({\mathbb{R}}^{d}\) 中的 \(d + 1\) 个向量
\end{tcolorbox}

\[
{v}^{0}\left( k\right)  = \left( \begin{matrix} 0 \\  0 \\  0 \\  \vdots \\  0 \end{matrix}\right) ,{v}^{1}\left( k\right)  = \left( \begin{matrix} k \\  0 \\  0 \\  \vdots \\  0 \end{matrix}\right) ,{v}^{2}\left( k\right)  = \left( \begin{matrix} k \\  k \\  0 \\  \vdots \\  0 \end{matrix}\right) ,\cdots ,{v}^{d}\left( k\right)  = \left( \begin{matrix} k \\  k \\  k \\  \vdots \\  k \end{matrix}\right) . \tag{7.25}
\]

\begin{center}
\includegraphics[max width=0.5\textwidth]{images/bo_d2e5cmbef24c73b03tn0_209_338_199_648_658_0.jpg}
\end{center}
\hspace*{3em} 

Figure 7.16 The Freudenthal triangulation of the two-dimensional simplex \({T}_{2}\left( k\right)\) with vertices \(\left( {0,0}\right) ,\left( {k,0}\right)\) and(k, k) for \(k = 8\) .

\begin{tcolorbox}
图7.16 二维单纯形(simplex) \({T}_{2}\left( k\right)\) 的弗赖登塔尔三角剖分(Freudenthal triangulation)，其顶点为 \(\left( {0,0}\right) ,\left( {k,0}\right)\) 和 (k, k)，其中 \(k = 8\) 。
\end{tcolorbox}

These are affinely independent because the corresponding vectors \({z}^{i} = \left( \begin{matrix} 1 \\  {v}^{i}\left( k\right)  \end{matrix}\right)\) in \({\mathbb{R}}^{1 + d}\) as considered in (7.4) are linearly independent because \(\mathop{\sum }\limits_{{i = 0}}^{d}{z}^{i}{\lambda }_{i} = \mathbf{0}\) implies \({\lambda }_{d} = 0\) , then \({\lambda }_{d - 1} = 0\) , and so on until \({\lambda }_{0} = 0\) . Their convex hull therefore defines a simplex, which we call \({T}_{d}\left( k\right)\) .

\begin{tcolorbox}
这些向量仿射无关，因为在 (7.4) 中考虑的 \({\mathbb{R}}^{1 + d}\) 中对应的向量 \({z}^{i} = \left( \begin{matrix} 1 \\  {v}^{i}\left( k\right)  \end{matrix}\right)\) 线性无关，因为 \(\mathop{\sum }\limits_{{i = 0}}^{d}{z}^{i}{\lambda }_{i} = \mathbf{0}\) 意味着 \({\lambda }_{d} = 0\) ，然后 \({\lambda }_{d - 1} = 0\) ，依此类推，直到 \({\lambda }_{0} = 0\) 。因此，它们的凸包(convex hull)定义了一个单纯形(simplex)，我们称之为 \({T}_{d}\left( k\right)\) 。
\end{tcolorbox}

Let \(k = 1\) . The vertices of \({T}_{d}\left( 1\right)\) are defined by starting from the all-zero vector \({v}^{0}\) and changing one coordinate at a time from 0 to 1 until we arrive at the all-one vector \({v}^{d} = \mathbf{1} \in  {\mathbb{R}}^{d}\) . In order to construct \({T}_{d}\left( 1\right)\) we have chosen the coordinates in their natural order \(1,\ldots ,d\) . As a generalization, consider an arbitrary order \(\pi \left( 1\right) ,\ldots ,\pi \left( d\right)\) of the coordinates, that is, a permutation \(\pi\) of \(1,\ldots ,d\) (a bijection on the set \(\{ 1,\ldots ,d\} )\) . We denote the set of all permutations of \(1,\ldots ,d\) by \(\operatorname{Perm}\left( d\right)\) . For \(\pi  \in  \operatorname{Perm}\left( d\right)\) , let \({v}_{\pi }^{0} = \mathbf{0} \in  {\mathbb{R}}^{d}\) , and for \(i = 1,\ldots ,d\) define

\begin{tcolorbox}
设 \(k = 1\) 。 \({T}_{d}\left( 1\right)\) 的顶点是从全零向量 \({v}^{0}\) 开始定义的，每次将一个坐标从 0 变为 1，直到得到全 1 向量 \({v}^{d} = \mathbf{1} \in  {\mathbb{R}}^{d}\) 。为了构造 \({T}_{d}\left( 1\right)\) ，我们按自然顺序 \(1,\ldots ,d\) 选择坐标。作为推广，考虑坐标的任意顺序 \(\pi \left( 1\right) ,\ldots ,\pi \left( d\right)\) ，即 \(1,\ldots ,d\) 的一个排列(permutation) \(\pi\) (集合 \(\{ 1,\ldots ,d\} )\) 上的双射)。我们用 \(\operatorname{Perm}\left( d\right)\) 表示 \(1,\ldots ,d\) 的所有排列的集合。对于 \(\pi  \in  \operatorname{Perm}\left( d\right)\) ，设 \({v}_{\pi }^{0} = \mathbf{0} \in  {\mathbb{R}}^{d}\) ，对于 \(i = 1,\ldots ,d\) 定义
\end{tcolorbox}

\[
{v}_{\pi }^{i} = {v}_{\pi }^{i - 1} + {e}_{\pi \left( i\right) }. \tag{7.26}
\]

That is, the \(i\) th vector \({v}_{\pi }^{i}\) is obtained from the previous vector by adding the unit vector \({e}_{\pi \left( i\right) }\) , which changes the coordinate at position \(\pi \left( i\right)\) from 0 to 1 . For example, if \(\left( {\pi \left( 1\right) ,\ldots ,\pi \left( d\right) }\right)  = \left( {2,3,1,4}\right)\) , then

\begin{tcolorbox}
也就是说，第 \(i\) 个向量 \({v}_{\pi }^{i}\) 是通过在前一个向量上加上单位向量 \({e}_{\pi \left( i\right) }\) 得到的，这将位置 \(\pi \left( i\right)\) 处的坐标从 0 变为 1 。例如，如果 \(\left( {\pi \left( 1\right) ,\ldots ,\pi \left( d\right) }\right)  = \left( {2,3,1,4}\right)\) ，那么
\end{tcolorbox}

\[
{v}_{\pi }^{0} = \left( \begin{array}{l} 0 \\  0 \\  0 \\  0 \end{array}\right) ,{v}_{\pi }^{1} = \left( \begin{array}{l} 0 \\  1 \\  0 \\  0 \end{array}\right) ,{v}_{\pi }^{2} = \left( \begin{array}{l} 0 \\  1 \\  1 \\  0 \end{array}\right) ,{v}_{\pi }^{3} = \left( \begin{array}{l} 1 \\  1 \\  1 \\  0 \end{array}\right) ,{v}_{\pi }^{4} = \left( \begin{array}{l} 1 \\  1 \\  1 \\  1 \end{array}\right) . \tag{7.27}
\]

This defines a "monotonic path" of vertices of \({C}_{d}\) that starts at the all-zero vector0 and ends at the all-one vector1. (It is easy to see that the vertices of the unit cube \({C}_{d}\) are the \({2}^{d}\) vectors in \({\mathbb{R}}^{d}\) where each component is either 0 or 1 .) The convex hull of these vectors defines the simplex

\begin{tcolorbox}
这定义了 \({C}_{d}\) 的顶点的一条“单调路径(path)”，它从全零向量 0 开始，到全 1 向量 1 结束。(很容易看出，单位立方体(unit cube) \({C}_{d}\) 的顶点是 \({\mathbb{R}}^{d}\) 中的 \({2}^{d}\) 个向量，其中每个分量要么是 0 要么是 1 。)这些向量的凸包(convex hull)定义了单纯形
\end{tcolorbox}

\[
{T}_{d}^{\pi } = \operatorname{conv}\left\{  {{v}_{\pi }^{0},{v}_{\pi }^{1},\ldots ,{v}_{\pi }^{d}}\right\}  .
\]

The following lemma describes this simplex in terms of \(d + 1\) linear inequalities rather than as a convex hull of vertices.

\begin{tcolorbox}
以下引理用 \(d + 1\) 个线性不等式而不是顶点的凸包(convex hull)来描述这个单纯形(simplex)。
\end{tcolorbox}

Lemma 7.21. Let \(x \in  {\mathbb{R}}^{d}\) and \(\pi  \in  \operatorname{Perm}\left( d\right)\) . Then \(x \in  {T}_{d}^{\pi }\) if and only if

\begin{tcolorbox}
引理7.21。设 \(x \in  {\mathbb{R}}^{d}\) 和 \(\pi  \in  \operatorname{Perm}\left( d\right)\) 。那么当且仅当
\end{tcolorbox}

\[
1 \geq  {x}_{\pi \left( 1\right) } \geq  {x}_{\pi \left( 2\right) } \geq  \cdots  \geq  {x}_{\pi \left( d\right) } \geq  0. \tag{7.28}
\]

Proof. Suppose (7.28) holds. Define \({y}_{0},{y}_{1},\ldots ,{y}_{d}\) by

\begin{tcolorbox}
证明。假设 (7.28) 成立。通过以下方式定义 \({y}_{0},{y}_{1},\ldots ,{y}_{d}\)
\end{tcolorbox}

\[
{y}_{0} = 1 - {x}_{\pi \left( 1\right) }
\]

\[
{y}_{1} = {x}_{\pi \left( 1\right) } - {x}_{\pi \left( 2\right) }
\]

\[
\vdots  \tag{7.29}
\]

\[
{y}_{d - 1} = {x}_{\pi \left( {d - 1}\right) } - {x}_{\pi \left( d\right) }
\]

\[
{y}_{d} = {x}_{\pi \left( d\right) }
\]

which implies

\begin{tcolorbox}
这意味着
\end{tcolorbox}

\[
{x}_{\pi \left( j\right) } = \mathop{\sum }\limits_{{i = j}}^{d}{y}_{i}\;\left( {1 \leq  j \leq  d}\right)
\]

and by (7.29) and (7.28)

\begin{tcolorbox}
并且由 (7.29) 和 (7.28)
\end{tcolorbox}

\[
\mathop{\sum }\limits_{{i = 0}}^{d}{y}_{i} = 1,\;{y}_{i} \geq  0\;\left( {1 \leq  i \leq  d}\right) . \tag{7.30}
\]

By (7.26), as the example (7.27) illustrates, the components of the vectors \({v}_{\pi }^{i}\) are given by

\begin{tcolorbox}
由(7.26)式，正如示例(7.27)所示，向量 \({v}_{\pi }^{i}\) 的分量由下式给出
\end{tcolorbox}

\[
{\left( {v}_{\pi }^{i}\right) }_{\pi \left( j\right) } = \left\{  {\begin{array}{ll} 0 & \text{ if }i < j \\  1 & \text{ if }i \geq  j \end{array}\;\left( {0 \leq  i \leq  d,1 \leq  j \leq  d}\right) .}\right.
\]

Hence, by (7.29),

\begin{tcolorbox}
因此，由(7.29)式
\end{tcolorbox}

\[
\mathop{\sum }\limits_{{i = 0}}^{d}{\left( {v}_{\pi }^{i}\right) }_{\pi \left( j\right) }{y}_{i} = \mathop{\sum }\limits_{{i = j}}^{d}{y}_{i} = {x}_{\pi \left( j\right) }\;\left( {1 \leq  j \leq  d}\right)
\]

and therefore

\begin{tcolorbox}
所以
\end{tcolorbox}

\[
x = \mathop{\sum }\limits_{{i = 0}}^{d}{v}_{\pi }^{i}{y}_{i} \tag{7.31}
\]

That is, \(x\) is by (7.30) a convex combination of the vectors \({v}_{\pi }^{0},\ldots ,{v}_{\pi }^{d}\) , which shows that \(x \in  {T}_{d}^{\pi }\) . Moreover, \({y}_{0},\ldots ,{y}_{d}\) in (7.29) are the barycentric coordinates of \(x\) .

\begin{tcolorbox}
也就是说，由(7.30)式可知， \(x\) 是向量 \({v}_{\pi }^{0},\ldots ,{v}_{\pi }^{d}\) 的凸组合，这表明 \(x \in  {T}_{d}^{\pi }\) 。此外，(7.29)式中的 \({y}_{0},\ldots ,{y}_{d}\) 是 \(x\) 的重心坐标。
\end{tcolorbox}

Conversely, any vertex \(x = {v}_{\pi }^{i}\) of \({T}_{d}^{\pi }\) for \(0 \leq  i \leq  d\) fulfills (7.28), where the only inequality that holds strictly (if \(0 < i < d\) ) is when the component of \(x\) changes from 1 to 0 . Therefore, the inequalities (7.28) also hold for any convex combination of these vertices, that is, for any element of \({T}_{d}^{\pi }\) .

\begin{tcolorbox}
反之，对于 \(0 \leq  i \leq  d\) ， \({T}_{d}^{\pi }\) 的任何顶点 \(x = {v}_{\pi }^{i}\) 都满足(7.28)式，其中唯一严格成立的不等式(如果 \(0 < i < d\) )是当 \(x\) 的分量从 1 变为 0 时。因此，不等式(7.28)对于这些顶点的任何凸组合也成立，即对于 \({T}_{d}^{\pi }\) 的任何元素都成立。
\end{tcolorbox}

The simplex \({T}_{d}\left( 1\right)\) is \({T}_{d}^{\pi }\) for the identity permutation \(\pi\) . It is easy to see that \({T}_{d}\left( k\right)  = \left\{  {{xk} \mid  x \in  {T}_{d}\left( 1\right) }\right\}\) and therefore by Lemma 7.21

\begin{tcolorbox}
对于恒等排列 \(\pi\) ，单纯形 \({T}_{d}\left( 1\right)\) 是 \({T}_{d}^{\pi }\) 。容易看出 \({T}_{d}\left( k\right)  = \left\{  {{xk} \mid  x \in  {T}_{d}\left( 1\right) }\right\}\) ，因此根据引理 7.21
\end{tcolorbox}

\[
z \in  {T}_{d}\left( k\right) \; \Leftrightarrow  \;k \geq  {z}_{1} \geq  {z}_{2} \geq  \cdots  \geq  {z}_{d} \geq  0. \tag{7.32}
\]

Lemma 7.22. The set \(\left\{  {{T}_{d}^{\pi } \mid  \pi  \in  \operatorname{Perm}\left( d\right) }\right\}\) is a triangulation of the unit cube \({C}_{d}\) .

\begin{tcolorbox}
引理 7.22. 集合 \(\left\{  {{T}_{d}^{\pi } \mid  \pi  \in  \operatorname{Perm}\left( d\right) }\right\}\) 是单位立方体 \({C}_{d}\) 的一个三角剖分。
\end{tcolorbox}

Proof. Let \(x = {\left( {x}_{1},\ldots ,{x}_{d}\right) }^{\top } \in  {C}_{d}\) , that is, \(1 \geq  {x}_{i} \geq  0\) for \(1 \leq  i \leq  d\) , and let \(\pi  \in  \operatorname{Perm}\left( d\right)\) . Sort the components \({x}_{i}\) of \(x\) so that (7.28) holds (the permutation \(\pi\) is not unique if some components of \(x\) are equal to each other). By Lemma 7.21, \(x \in  {T}_{d}^{\pi }\) . This shows \({C}_{d} = \mathop{\bigcup }\limits_{{\pi  \in  \operatorname{Perm}\left( d\right) }}{T}_{d}^{\pi }\) .

\begin{tcolorbox}
证明. 设 \(x = {\left( {x}_{1},\ldots ,{x}_{d}\right) }^{\top } \in  {C}_{d}\) ，即对于 \(1 \leq  i \leq  d\) 有 \(1 \geq  {x}_{i} \geq  0\) ，并设 \(\pi  \in  \operatorname{Perm}\left( d\right)\) 。对 \(x\) 的分量 \({x}_{i}\) 进行排序，使得(7.28)式成立(如果 \(x\) 的某些分量彼此相等，则排列 \(\pi\) 不唯一)。根据引理 7.21， \(x \in  {T}_{d}^{\pi }\) 。这表明 \({C}_{d} = \mathop{\bigcup }\limits_{{\pi  \in  \operatorname{Perm}\left( d\right) }}{T}_{d}^{\pi }\) 。
\end{tcolorbox}

If \(\pi ,{\pi }^{\prime } \in  \operatorname{Perm}\left( d\right)\) and \(\pi  \neq  {\pi }^{\prime }\) then \({T}_{d}^{\pi } \nsubseteq  {T}_{d}^{{\pi }^{\prime }}\) , by taking some \(x \in  {T}_{d}^{\pi }\) with strict inequalities in (7.28) where clearly \(x \notin  {T}_{d}^{{\pi }^{\prime }}\) . It remains to show that the intersection of any two sub-simplices \({T}_{d}^{\pi }\) and \({T}_{d}^{{\pi }^{\prime }}\) is a face of both. Any \(x\) in \({T}_{d}^{\pi } \cap  {T}_{d}^{{\pi }^{\prime }}\) fulfills (7.28) as well as

\begin{tcolorbox}
如果 \(\pi ,{\pi }^{\prime } \in  \operatorname{Perm}\left( d\right)\) 且 \(\pi  \neq  {\pi }^{\prime }\) ，那么通过选取(7.28)式中某些严格不等式对应的 \(x \in  {T}_{d}^{\pi }\) (显然 \(x \notin  {T}_{d}^{{\pi }^{\prime }}\) )，有 \({T}_{d}^{\pi } \nsubseteq  {T}_{d}^{{\pi }^{\prime }}\) 。还需证明任意两个子单纯形 \({T}_{d}^{\pi }\) 和 \({T}_{d}^{{\pi }^{\prime }}\) 的交集是它们两者的一个面。 \({T}_{d}^{\pi } \cap  {T}_{d}^{{\pi }^{\prime }}\) 中的任何 \(x\) 都满足(7.28)式以及
\end{tcolorbox}

\[
{x}_{{\pi }^{\prime }\left( 1\right) } \geq  {x}_{{\pi }^{\prime }\left( 2\right) } \geq  \cdots  \geq  {x}_{{\pi }^{\prime }\left( d\right) }. \tag{7.33}
\]

Any inequality that holds in reverse order in (7.28) and (7.33) has to hold as an equality. This includes any implied equality; for example, if for \(d = 3\) we have \({x}_{1} \geq  {x}_{2} \geq  {x}_{3}\) and \({x}_{2} \geq  {x}_{3} \geq  {x}_{1}\) , then this implies \({x}_{1} = {x}_{3}\) as well as \({x}_{2} = {x}_{3}\) . The intersection \({T}_{d}^{\pi } \cap  {T}_{d}^{{\pi }^{\prime }}\) is then characterized by (7.28) where some inequalities are replaced by equalities. The barycentric coordinates \({y}_{1},\ldots ,{y}_{d - 1}\) in (7.29) of \(x\) that correspond to these equalities are then set to zero, which means that the corresponding vertex is not used in the representation of \(x\) in (7.31). Hence, \(x\) is only a convex combination of the remaining vertices. These are the vertices of \({T}_{d}^{\pi } \cap  {T}_{d}^{{\pi }^{\prime }}\) , which is therefore a common face of the two sub-simplices \({T}_{d}^{\pi }\) and \({T}_{d}^{{\pi }^{\prime }}\) , as claimed. It may be the case that all inequalities in (7.28) are equalities, so that the common face has just the vertices0and1.

\begin{tcolorbox}
任何在(7.28)和(7.33)中反向成立的不等式都必须以等式形式成立。这包括任何隐含的等式；例如，如果对于 \(d = 3\) 我们有 \({x}_{1} \geq  {x}_{2} \geq  {x}_{3}\) 和 \({x}_{2} \geq  {x}_{3} \geq  {x}_{1}\) ，那么这意味着 \({x}_{1} = {x}_{3}\) 以及 \({x}_{2} = {x}_{3}\) 。交集 \({T}_{d}^{\pi } \cap  {T}_{d}^{{\pi }^{\prime }}\) 随后由(7.28)表征，其中一些不等式被等式取代。在(7.29)中 \(x\) 的与这些等式对应的重心坐标(barycentric coordinates) \({y}_{1},\ldots ,{y}_{d - 1}\) 随后被设为零，这意味着在(7.31)中 \(x\) 的表示不使用相应的顶点。因此， \(x\) 只是剩余顶点的凸组合(convex combination)。这些是 \({T}_{d}^{\pi } \cap  {T}_{d}^{{\pi }^{\prime }}\) 的顶点，因此它是两个子单纯形(sub - simplex) \({T}_{d}^{\pi }\) 和 \({T}_{d}^{{\pi }^{\prime }}\) 的公共面(face)，正如所声称的那样。可能会出现(7.28)中的所有不等式都是等式的情况，使得公共面只有顶点0和1。
\end{tcolorbox}

There are \(d\) ! many permutations \(\pi\) of the numbers \(1,\ldots ,d\) . Any two simplices \({T}_{d}^{\pi }\) and \({T}_{d}^{{\pi }^{\prime }}\) have the same shape and size because they only differ in the order of their coordinates, and they intersect in a set of dimension \(d - 1\) or lower which has volume zero in \({\mathbb{R}}^{d}\) . Because the unit cube \({C}_{d}\) has volume 1, Lemma 7.22 implies that each simplex \({T}_{d}^{\pi }\) has volume \(\frac{1}{d!}\) .

\begin{tcolorbox}
数字 \(1,\ldots ,d\) 有 \(d\) ! 种排列(permutation) \(\pi\) 。任意两个单纯形(simplex) \({T}_{d}^{\pi }\) 和 \({T}_{d}^{{\pi }^{\prime }}\) 具有相同的形状和大小，因为它们仅在坐标顺序上有所不同，并且它们的交集是一个维度(dimension)为 \(d - 1\) 或更低的集合，该集合在 \({\mathbb{R}}^{d}\) 中的体积为零。由于单位立方体(unit cube) \({C}_{d}\) 的体积为1，引理7.22意味着每个单纯形 \({T}_{d}^{\pi }\) 的体积为 \(\frac{1}{d!}\) 。
\end{tcolorbox}

We now use the triangulation of the unit cube \({C}_{d}\) in Lemma 7.22 to construct a triangulation of the simplex \({T}_{d}\left( k\right)\) in (7.25), for some integer \(k \geq  1\) .

\begin{tcolorbox}
现在，我们使用引理7.22中单位立方体 \({C}_{d}\) 的三角剖分(triangulation)来构造(7.25)中单纯形 \({T}_{d}\left( k\right)\) 的三角剖分，其中 \(k \geq  1\) 为某个整数。
\end{tcolorbox}

For a polytope \(P\) in \({\mathbb{R}}^{d}\) such as \({C}_{d}\) or \({T}_{d}^{\pi }\) and an integer "grid" vector \(g =\)  \({\left( {g}_{1},\ldots ,{g}_{d}\right) }^{\top } \in  {\mathbb{Z}}^{d}\) , let \(g + P = \{ g + x \mid  x \in  P\}\) . Let

\begin{tcolorbox}
对于 \({\mathbb{R}}^{d}\) 中的有界多面体(polytope) \(P\) ，例如 \({C}_{d}\) 或 \({T}_{d}^{\pi }\) ，以及一个整数“网格”向量 \(g =\)  \({\left( {g}_{1},\ldots ,{g}_{d}\right) }^{\top } \in  {\mathbb{Z}}^{d}\) ，设 \(g + P = \{ g + x \mid  x \in  P\}\) 。设
\end{tcolorbox}

\[
\sum  = \left\{  {g + {T}_{d}^{\pi } \mid  \pi  \in  \operatorname{Perm}\left( d\right) ,g \in  {\mathbb{Z}}^{d},g + {T}_{d}^{\pi } \subseteq  {T}_{d}\left( k\right) }\right\}  . \tag{7.34}
\]

In other words, we consider any simplex \({T}_{d}^{\pi }\) of the triangulation of \({C}_{d}\) , translate it by adding some integer vector \(g\) , and if the resulting set \(g + {T}_{d}^{\pi }\) is a subset of \({T}_{d}\left( k\right)\) then this is a sub-simplex in \(\sum\) . By (7.32) this can only hold if \(0 \leq  {g}_{i} \leq  k\) for \(1 \leq  i \leq  d\) because otherwise \(g + \mathbf{0} \notin  {T}_{d}\left( k\right)\) (where \(\mathbf{0} \in  {T}_{d}^{\pi }\) ), so \(\sum\) is finite. Figure 7.17 shows \(\sum\) for \(d = 3\) and \(k = 2\) .

\begin{tcolorbox}
换句话说，我们考虑 \({C}_{d}\) 的三角剖分中的任意单纯形 \({T}_{d}^{\pi }\) ，通过加上某个整数向量 \(g\) 对其进行平移，如果得到的集合 \(g + {T}_{d}^{\pi }\) 是 \({T}_{d}\left( k\right)\) 的子集，那么这就是 \(\sum\) 中的一个子单纯形。根据(7.32)，只有当 \(0 \leq  {g}_{i} \leq  k\) 对于 \(1 \leq  i \leq  d\) 成立时才会这样，因为否则 \(g + \mathbf{0} \notin  {T}_{d}\left( k\right)\) (其中 \(\mathbf{0} \in  {T}_{d}^{\pi }\) )，所以 \(\sum\) 是有限的。图7.17展示了 \(\sum\) 在 \(d = 3\) 和 \(k = 2\) 时的情况。
\end{tcolorbox}

\begin{center}
\includegraphics[max width=0.9\textwidth]{images/bo_d2e5cmbef24c73b03tn0_212_193_349_1262_579_0.jpg}
\end{center}
\hspace*{3em} 

Figure 7.17 The Freudenthal triangulation \(\sum\) of the three-dimensional simplex \({T}_{3}\left( k\right)\) with vertices \(\left( {0,0,0}\right) ,\left( {k,0,0}\right) ,\left( {k,k,0}\right)\) and(k, k, k)for \(k = 2\) .

\begin{tcolorbox}
图7.17 三维单纯形 \({T}_{3}\left( k\right)\) 的弗赖登塔尔三角剖分 \(\sum\) ，其顶点为 \(\left( {0,0,0}\right) ,\left( {k,0,0}\right) ,\left( {k,k,0}\right)\) 以及对于 \(k = 2\) 的 (k, k, k)。
\end{tcolorbox}

Theorem 7.23. \(\sum\) as defined in (7.34) is a triangulation of \({T}_{d}\left( k\right)\) .

\begin{tcolorbox}
定理7.23。如 (7.34) 所定义的 \(\sum\) 是 \({T}_{d}\left( k\right)\) 的一个三角剖分。
\end{tcolorbox}

Proof. By (7.34), \(\mathop{\bigcup }\limits_{{S \in  \sum }}S \subseteq  {T}_{d}\left( k\right)\) . To show the converse inclusion, let \(z \in  {T}_{d}\left( k\right)\) . We want to find \(g \in  {\mathbb{Z}}^{d}\) and \(\pi  \in  \operatorname{Perm}\left( d\right)\) so that \(z \in  g + {T}_{d}^{\pi }\) , and crucially, \(g + {T}_{d}^{\pi } \subseteq  {T}_{d}\left( k\right)\) as required in (7.34). If this holds then \(z = g + x\) for some \(x \in  {T}_{d}^{\pi }\) , which implies

\begin{tcolorbox}
证明。由 (7.34) 可知， \(\mathop{\bigcup }\limits_{{S \in  \sum }}S \subseteq  {T}_{d}\left( k\right)\) 。为证明反向包含关系，设 \(z \in  {T}_{d}\left( k\right)\) 。我们要找到 \(g \in  {\mathbb{Z}}^{d}\) 和 \(\pi  \in  \operatorname{Perm}\left( d\right)\) 使得 \(z \in  g + {T}_{d}^{\pi }\) ，并且至关重要的是，满足 (7.34) 中要求的 \(g + {T}_{d}^{\pi } \subseteq  {T}_{d}\left( k\right)\) 。如果这成立，那么对于某个 \(x \in  {T}_{d}^{\pi }\) 有 \(z = g + x\) ，这意味着
\end{tcolorbox}

\[
{z}_{i} = {g}_{i} + {x}_{i}\;\left( {{g}_{i} \in  \mathbb{Z},{x}_{i} \in  \left\lbrack  {0,1}\right\rbrack  ,1 \leq  i \leq  d}\right) . \tag{7.35}
\]

In (7.35), the choice of \({g}_{i}\) and \({x}_{i}\) is not unique if \({z}_{i} \in  \mathbb{Z}\) because then either \({x}_{i} = 0\) or \({x}_{i} = 1\) are possible. Then we normally choose \({x}_{i} = 0\) , except when \({z}_{i} = k\) where we choose \({g}_{i} = k - 1\) and \({x}_{i} = 1\) , to ensure that \(g + {T}_{d}^{\pi } \subseteq  {T}_{d}\left( k\right)\) . With \(\left\lfloor  {z}_{i}\right\rfloor   = \max \left\{  {t \in  \mathbb{Z} \mid  t \leq  {z}_{i}}\right\}\) , this is achieved by

\begin{tcolorbox}
在 (7.35) 中，如果 \({z}_{i} \in  \mathbb{Z}\) ， \({g}_{i}\) 和 \({x}_{i}\) 的选择不是唯一的，因为此时 \({x}_{i} = 0\) 或 \({x}_{i} = 1\) 都是可能的。那么我们通常选择 \({x}_{i} = 0\) ，除非当 \({z}_{i} = k\) 时我们选择 \({g}_{i} = k - 1\) 和 \({x}_{i} = 1\) ，以确保 \(g + {T}_{d}^{\pi } \subseteq  {T}_{d}\left( k\right)\) 。对于 \(\left\lfloor  {z}_{i}\right\rfloor   = \max \left\{  {t \in  \mathbb{Z} \mid  t \leq  {z}_{i}}\right\}\) ，这可通过
\end{tcolorbox}

\[
{g}_{i} = \left\{  {\begin{array}{ll} \left\lfloor  {z}_{i}\right\rfloor  & \text{ if }{z}_{i} < k \\  k - 1 & \text{ if }{z}_{i} = k, \end{array}\;{x}_{i} = {z}_{i} - {g}_{i}\;\left( {1 \leq  i \leq  d}\right) ,}\right.  \tag{7.36}
\]

which clearly implies (7.35).

\begin{tcolorbox}
这显然意味着 (7.35)。
\end{tcolorbox}

Similarly, the integral parts \({g}_{i}\) of \({z}_{i}\) fulfill

\begin{tcolorbox}
类似地， \({z}_{i}\) 的整数部分 \({g}_{i}\) 满足
\end{tcolorbox}

\[
{g}_{i} \geq  {g}_{i + 1}\;\left( {1 \leq  i < d}\right)  \tag{7.37}
\]

because if \({z}_{i} = k\) then \({z}_{i + 1} \leq  k\) and \({g}_{i} = k - 1 \geq  {g}_{i + 1}\) , and if \({z}_{i} < k\) then by (7.36) we have \({x}_{i} < 1\) so that \({g}_{i} < {g}_{i + 1}\) would imply \({z}_{i} = {g}_{i} + {x}_{i} < {g}_{i} + 1 \leq  {g}_{i + 1} \leq  {z}_{i + 1}\) , which contradicts (7.32).

\begin{tcolorbox}
因为如果 \({z}_{i} = k\) ，那么 \({z}_{i + 1} \leq  k\) 且 \({g}_{i} = k - 1 \geq  {g}_{i + 1}\) ，并且如果 \({z}_{i} < k\) ，那么由 (7.36) 我们有 \({x}_{i} < 1\) ，使得 \({g}_{i} < {g}_{i + 1}\) 意味着 \({z}_{i} = {g}_{i} + {x}_{i} < {g}_{i} + 1 \leq  {g}_{i + 1} \leq  {z}_{i + 1}\) ，这与 (7.32) 矛盾。
\end{tcolorbox}

Given (7.35), we want to find \(\pi  \in  \operatorname{Perm}\left( d\right)\) so that \(x \in  {T}_{d}^{\pi }\) and \(g + {T}_{d}^{\pi } \subseteq  {T}_{d}\left( k\right)\) . The second condition is not automatic: For example, if \(d = 2\) and \(k = 1\) then \(\sum  = \left\{  {{T}_{2}\left( 1\right) }\right\}\) and for \(z = x = \left( {\frac{1}{2},\frac{1}{2}}\right)\) both \({T}_{2}\left( 1\right)\) and \({T}_{2}^{\pi }\left( 1\right)\) for \(\left( {\pi \left( 1\right) ,\pi \left( 2\right) }\right)  = \left( {2,1}\right)\) contain \(x\) , but \(\mathbf{0} + {T}_{2}^{\pi }\left( 1\right)  \notin  \sum\) . The following is a safe choice of \(\pi\) . Sort the components \({x}_{i}\) of \(x\) with a permutation \(\pi\) so that (7.28) holds and so that

\begin{tcolorbox}
给定式(7.35)，我们希望找到 \(\pi  \in  \operatorname{Perm}\left( d\right)\) 使得 \(x \in  {T}_{d}^{\pi }\) 且 \(g + {T}_{d}^{\pi } \subseteq  {T}_{d}\left( k\right)\) 。第二个条件并非自动满足:例如，如果 \(d = 2\) 且 \(k = 1\) ，那么 \(\sum  = \left\{  {{T}_{2}\left( 1\right) }\right\}\) ，并且对于 \(z = x = \left( {\frac{1}{2},\frac{1}{2}}\right)\) ， \({T}_{2}\left( 1\right)\) 和 \({T}_{2}^{\pi }\left( 1\right)\) 在 \(\left( {\pi \left( 1\right) ,\pi \left( 2\right) }\right)  = \left( {2,1}\right)\) 时都包含 \(x\) ，但 \(\mathbf{0} + {T}_{2}^{\pi }\left( 1\right)  \notin  \sum\) 。以下是 \(\pi\) 的一个稳妥选择。用一个排列 \(\pi\) 对 \(x\) 的分量 \({x}_{i}\) 进行排序，使得式(7.28)成立，并且使得
\end{tcolorbox}

\[
1 \leq  h < j \leq  d,\;{x}_{\pi \left( h\right) } = {x}_{\pi \left( j\right) }\; \Rightarrow  \;\pi \left( h\right)  < \pi \left( j\right) . \tag{7.38}
\]

This can always be achieved, if necessary by exchanging \(\pi \left( h\right)\) and \(\pi \left( j\right)\) if \({x}_{\pi \left( h\right) } =\)  \({x}_{\pi \left( j\right) }\) , which maintains (7.28). Then \(x \in  {T}_{d}^{\pi }\) by Lemma 7.21. We want to show that \(g + y \in  {T}_{d}\left( k\right)\) for any \(y \in  {T}_{d}^{\pi }\) , that is, according to (7.32),

\begin{tcolorbox}
这总能实现，如有必要，若 \({x}_{\pi \left( h\right) } =\)  \({x}_{\pi \left( j\right) }\) ，可交换 \(\pi \left( h\right)\) 和 \(\pi \left( j\right)\) ，这能保持式(7.28)成立。然后根据引理 7.21 有 \(x \in  {T}_{d}^{\pi }\) 。我们想证明对于任意 \(y \in  {T}_{d}^{\pi }\) 都有 \(g + y \in  {T}_{d}\left( k\right)\) ，即根据式(7.32)
\end{tcolorbox}

\[
{g}_{i} + {y}_{i} \geq  {g}_{i + 1} + {y}_{i + 1}\;\left( {1 \leq  i < d}\right)  \tag{7.39}
\]

(by (7.36), we have \(k \geq  {g}_{1} + {y}_{1}\) and \({g}_{d} + {y}_{d} \geq  0\) ). By (7.37), condition (7.39) holds if \({y}_{i} \geq  {y}_{i + 1}\) . Otherwise, let \(y \in  {T}_{d}^{\pi }\) and suppose that \({y}_{i} < {y}_{i + 1}\) for some \(1 \leq  i < d\) . Let \(i + 1 = \pi \left( h\right)\) and \(i = \pi \left( j\right)\) . Then \({y}_{\pi \left( h\right) } = {y}_{i + 1} > {y}_{i} = {y}_{\pi \left( j\right) }\) , which implies \(h < j\) by Lemma 7.21. Then \({x}_{\pi \left( h\right) } \geq  {x}_{\pi \left( j\right) }\) , where \({x}_{\pi \left( h\right) } = {x}_{\pi \left( j\right) }\) would imply \(i + 1 = \pi \left( h\right)  <\)  \(\pi \left( j\right)  = i\) by (7.38) which is not possible, that is, \({x}_{i + 1} = {x}_{\pi \left( h\right) } > {x}_{\pi \left( j\right) } = {x}_{i}\) . By (7.37), either \({g}_{i} = {g}_{i + 1}\) or \({g}_{i} > {g}_{i + 1}\) . If \({g}_{i} = {g}_{i + 1}\) then \({z}_{i + 1} = {g}_{i + 1} + {x}_{i + 1} > {g}_{i} + {x}_{i} = {z}_{i}\) , which violates (7.32). Hence, \({g}_{i} + {y}_{i} \geq  {g}_{i} \geq  {g}_{i + 1} + 1 \geq  {g}_{i + 1} + {y}_{i + 1}\) . This shows (7.39), and therefore, as claimed, \(g + {T}_{d}^{\pi } \subseteq  {T}_{d}\left( k\right)\) . That is, the sub-simplices in \(\sum\) cover \({T}_{d}\left( k\right)\) .

\begin{tcolorbox}
(由(7.36)，我们有 \(k \geq  {g}_{1} + {y}_{1}\) 和 \({g}_{d} + {y}_{d} \geq  0\) )。由(7.37)，若 \({y}_{i} \geq  {y}_{i + 1}\) ，则条件(7.39)成立。否则，令 \(y \in  {T}_{d}^{\pi }\) ，并假设对于某个 \(1 \leq  i < d\) 有 \({y}_{i} < {y}_{i + 1}\) 。令 \(i + 1 = \pi \left( h\right)\) 且 \(i = \pi \left( j\right)\) 。那么 \({y}_{\pi \left( h\right) } = {y}_{i + 1} > {y}_{i} = {y}_{\pi \left( j\right) }\) ，根据引理7.21这意味着 \(h < j\) 。接着 \({x}_{\pi \left( h\right) } \geq  {x}_{\pi \left( j\right) }\) ，其中 \({x}_{\pi \left( h\right) } = {x}_{\pi \left( j\right) }\) 由(7.38)将意味着 \(i + 1 = \pi \left( h\right)  <\)  \(\pi \left( j\right)  = i\) ，而这是不可能的，即 \({x}_{i + 1} = {x}_{\pi \left( h\right) } > {x}_{\pi \left( j\right) } = {x}_{i}\) 。由(7.37)，要么 \({g}_{i} = {g}_{i + 1}\) ，要么 \({g}_{i} > {g}_{i + 1}\) 。若 \({g}_{i} = {g}_{i + 1}\) ，则 \({z}_{i + 1} = {g}_{i + 1} + {x}_{i + 1} > {g}_{i} + {x}_{i} = {z}_{i}\) ，这与(7.32)矛盾。因此， \({g}_{i} + {y}_{i} \geq  {g}_{i} \geq  {g}_{i + 1} + 1 \geq  {g}_{i + 1} + {y}_{i + 1}\) 。这表明了(7.39)，因此，如所声称的， \(g + {T}_{d}^{\pi } \subseteq  {T}_{d}\left( k\right)\) 。也就是说， \(\sum\) 中的子单形覆盖 \({T}_{d}\left( k\right)\) 。
\end{tcolorbox}

In order to show that \(\sum\) is a triangulation of \({T}_{d}\left( k\right)\) , it remains to show that no sub-simplex is a subset of another, and that the intersection of two sub-simplices is a face of both. Both properties are shown similarly to Lemma 7.22, which implies them for two sub-simplices \(g + {T}_{d}^{\pi }\) and \({g}^{\prime } + {T}_{d}^{{\pi }^{\prime }}\) if \(g = {g}^{\prime }\) . If \(g \neq  {g}^{\prime }\) then the two sub-simplices are subsets of two different cubes \(g + {C}_{d}\) and \({g}^{\prime } + {C}_{d}\) which have disjoint non-empty interiors, so neither sub-simplex is a subset of the other. The intersection of \(g + {T}_{d}^{\pi }\) and \({g}^{\prime } + {T}_{d}^{{\pi }^{\prime }}\) is a subset of \(\left( {g + {C}_{d}}\right)  \cap  \left( {{g}^{\prime } + {C}_{d}}\right)\) , which is non-empty only if \(\left| {{g}_{i} - {g}_{i}^{\prime }}\right|  \leq  1\) for \(1 \leq  i \leq  d\) . By Lemma 7.21, \(z \in  \left( {g + {T}_{d}^{\pi }}\right)  \cap  \left( {{g}^{\prime } + {T}_{d}^{{\pi }^{\prime }}}\right)\) if and only if \(z = g + x = {g}^{\prime } + {x}^{\prime }\) so that (7.28) holds and, analogously,

\begin{tcolorbox}
为了证明 \(\sum\) 是 \({T}_{d}\left( k\right)\) 的一个三角剖分，还需证明没有子单纯形是另一个子单纯形的子集，并且两个子单纯形的交集是两者的一个面。这两个性质的证明与引理7.22类似，若 \(g = {g}^{\prime }\) ，则该引理对两个子单纯形 \(g + {T}_{d}^{\pi }\) 和 \({g}^{\prime } + {T}_{d}^{{\pi }^{\prime }}\) 成立。若 \(g \neq  {g}^{\prime }\) ，则这两个子单纯形是两个不同立方体 \(g + {C}_{d}\) 和 \({g}^{\prime } + {C}_{d}\) 的子集，这两个立方体的非空内部不相交，因此没有一个子单纯形是另一个的子集。 \(g + {T}_{d}^{\pi }\) 和 \({g}^{\prime } + {T}_{d}^{{\pi }^{\prime }}\) 的交集是 \(\left( {g + {C}_{d}}\right)  \cap  \left( {{g}^{\prime } + {C}_{d}}\right)\) 的子集，仅当对于 \(1 \leq  i \leq  d\) 有 \(\left| {{g}_{i} - {g}_{i}^{\prime }}\right|  \leq  1\) 时，该交集才非空。根据引理7.21，当且仅当 \(z = g + x = {g}^{\prime } + {x}^{\prime }\) 时 \(z \in  \left( {g + {T}_{d}^{\pi }}\right)  \cap  \left( {{g}^{\prime } + {T}_{d}^{{\pi }^{\prime }}}\right)\) ，从而 (7.28) 成立，并且类似地
\end{tcolorbox}

\[
1 \geq  {x}_{{\pi }^{\prime }\left( 1\right) }^{\prime } \geq  {x}_{{\pi }^{\prime }\left( 2\right) }^{\prime } \geq  \cdots  \geq  {x}_{{\pi }^{\prime }\left( d\right) }^{\prime } \geq  0. \tag{7.40}
\]

If \({g}_{i} = {g}_{i}^{\prime } + 1\) , then the equation \({z}_{i} = {g}_{i} + {x}_{i} = {g}_{i}^{\prime } + {x}_{i}^{\prime }\) imposes the additional equations \({x}_{i} = 1\) and \({x}_{i}^{\prime } = 0\) (and similarly \({x}_{i} = 0\) and \({x}_{i}^{\prime } = 1\) if \({g}_{i} + 1 = {g}_{i}^{\prime }\) ). For those subscripts \(i\) where \({g}_{i} = {g}_{i}^{\prime }\) we have \({x}_{i} = {x}_{i}^{\prime }\) , which induce further equations in (7.28) and (7.40) if these components of \(x\) and \({x}^{\prime }\) appear in reverse order in these inequalities. As argued in the proof of Lemma 7.22, these equations mean that certain barycentric coordinates \({y}_{0},{y}_{1},\ldots ,{y}_{d}\) in (7.29) are zero (which may now include \({y}_{0}\) if \({x}_{\pi \left( 1\right) } = 1\) and \({y}_{d}\) if \({x}_{\pi \left( d\right) } = 0\) ), which excludes the corresponding vertices of the two sub-simplices from their intersection. The remaining vertices (if any) are then the vertices of the intersection, which is therefore a common face of both sub-simplices. This completes the proof of the triangulation property.

\begin{tcolorbox}
若 \({g}_{i} = {g}_{i}^{\prime } + 1\) ，则方程 \({z}_{i} = {g}_{i} + {x}_{i} = {g}_{i}^{\prime } + {x}_{i}^{\prime }\) 施加了额外的方程 \({x}_{i} = 1\) 和 \({x}_{i}^{\prime } = 0\) (类似地，若 \({g}_{i} + 1 = {g}_{i}^{\prime }\) ，则有 \({x}_{i} = 0\) 和 \({x}_{i}^{\prime } = 1\) )。对于那些满足 \({g}_{i} = {g}_{i}^{\prime }\) 的下标 \(i\) ，我们有 \({x}_{i} = {x}_{i}^{\prime }\) ，如果 \(x\) 和 \({x}^{\prime }\) 的这些分量在这些不等式中以相反顺序出现，那么这会在 (7.28) 和 (7.40) 中引出更多方程。正如在引理7.22的证明中所论述的，这些方程意味着 (7.29) 中某些重心坐标 \({y}_{0},{y}_{1},\ldots ,{y}_{d}\) 为零(如果 \({x}_{\pi \left( 1\right) } = 1\) ，现在可能包括 \({y}_{0}\) ；如果 \({x}_{\pi \left( d\right) } = 0\) ，可能包括 \({y}_{d}\) )，这就将两个子单纯形的相应顶点排除在它们的交集之外。剩下的顶点(如果有的话)就是交集的顶点，因此该交集是两个子单纯形的公共面。这就完成了三角剖分性质的证明。
\end{tcolorbox}

The simplex \({T}_{d}\left( k\right)\) has \(d + 1\) vertices \({v}^{i}\left( k\right)\) in (7.25) that are vertices of a cube in \({\mathbb{R}}^{k}\) with edge length \(k\) . The triangulation \(\sum\) in (7.34) subdivides \({T}_{d}\left( k\right)\) into simplices \({T}_{d}^{\pi }\) (which triangulate the unit cube \({C}_{d}\) ) shifted by integer vectors \(g \in  {\mathbb{Z}}^{d}\) . The goal is to obtain an arbitrarily fine triangulation of the unit simplex \(\Delta\) , as used in the KKM lemma (Theorem 7.12). This achieved as follows: When mapping \({T}_{d}\left( k\right)\) to \(\Delta\) with the affine bijection \(J\) in Lemma 7.3, we obtain a corresponding triangulation \(\{ J\left( T\right)  \mid  T \in  \sum \}\) of \(\Delta\) . For each sub-simplex \(J\left( {g + {T}_{d}^{\pi }}\right)\) of that triangulation, the maximum of the distances between its vertices only depends on \(\pi\) , not on \(g\) . That maximum (over all \(\pi  \in  \operatorname{Perm}\left( d\right)\) and all vertex pairs) is proportional to \(\frac{1}{k}\) because \(J\) is affine and \({T}_{d}\left( k\right)  = \left\{  {{xk} \mid  x \in  {T}_{d}\left( 1\right) }\right\}\) . For any \(\varepsilon  > 0\) , one can therefore choose \(k\) large enough so that no two points in any sub-simplex of the triangulation of \(\Delta\) have distance greater than \(\varepsilon\) , which gives the desired \(\varepsilon\) -fine triangulation.

\begin{tcolorbox}
单纯形 \({T}_{d}\left( k\right)\) 在 (7.25) 中有 \(d + 1\) 个顶点 \({v}^{i}\left( k\right)\) ，这些顶点是 \({\mathbb{R}}^{k}\) 中边长为 \(k\) 的立方体的顶点。(7.34) 中的三角剖分 \(\sum\) 将 \({T}_{d}\left( k\right)\) 细分为单纯形 \({T}_{d}^{\pi }\) (这些单纯形对单位立方体 \({C}_{d}\) 进行三角剖分)，并通过整数向量 \(g \in  {\mathbb{Z}}^{d}\) 进行平移。目标是得到单位单纯形 \(\Delta\) 的任意精细的三角剖分，如 KKM 引理(定理 7.12)中所使用的那样。这可以通过以下方式实现:当使用引理 7.3 中的仿射双射 \(J\) 将 \({T}_{d}\left( k\right)\) 映射到 \(\Delta\) 时，我们得到 \(\Delta\) 的相应三角剖分 \(\{ J\left( T\right)  \mid  T \in  \sum \}\) 。对于该三角剖分的每个子单纯形 \(J\left( {g + {T}_{d}^{\pi }}\right)\) ，其顶点之间距离的最大值仅取决于 \(\pi\) ，而不取决于 \(g\) 。该最大值(在所有 \(\pi  \in  \operatorname{Perm}\left( d\right)\) 和所有顶点对上)与 \(\frac{1}{k}\) 成正比，因为 \(J\) 是仿射的且 \({T}_{d}\left( k\right)  = \left\{  {{xk} \mid  x \in  {T}_{d}\left( 1\right) }\right\}\) 。因此，对于任何 \(\varepsilon  > 0\) ，可以选择足够大的 \(k\) ，使得 \(\Delta\) 的三角剖分的任何子单纯形中的任意两点之间的距离都不大于 \(\varepsilon\) ，这就得到了所需的 \(\varepsilon\) - 精细三角剖分。
\end{tcolorbox}

The proof of Brouwer's fixed-point Theorem 7.20 is thereby complete.

\begin{tcolorbox}
布劳威尔不动点定理 7.20 的证明至此完成。
\end{tcolorbox}

\section*{7.8 Further Reading}

\begin{tcolorbox}
\section*{7.8 进一步阅读}
\end{tcolorbox}

The fixed-point theorem is stated at the end of Brouwer (1911, Satz 4). In this and related papers Brouwer developed and used the concept of "degree" (number of inverse images of a point) of a simplicial approximation of a continuous function. For a detailed discussion see Brouwer (1976). A historical survey of the wide-ranging impact of Brouwer's fixed-point theorem and its subsequent extensions is Park (1999).

\begin{tcolorbox}
不动点定理在布劳威尔(1911 年，命题 4)的论文结尾处被提出。在这篇以及相关论文中，布劳威尔发展并使用了连续函数的单纯逼近的“度”(一个点的原像数量)的概念。详细讨论见布劳威尔(1976 年)。关于布劳威尔不动点定理及其后续扩展的广泛影响的历史综述见帕克(1999 年)。
\end{tcolorbox}

Brouwer's fixed-point theorem is a major theorem in topology. The basic notion of topology is continuity, but more primitively topology can be seen as a qualitative study of "position". The first paper published in that regard is the solution by Euler (1741) to the problem of the "bridges of Königsberg" (see also Biggs, Lloyd, and Wilson, 1976). Euler's paper marks the foundation of both graph theory and topology; it introduces the concept of a graph, and the degree of a node plays a central role. It is fitting that Sperner's lemma relates to the even number of odd-degree nodes in a graph.

\begin{tcolorbox}
布劳威尔不动点定理是拓扑学中的一个重要定理。拓扑学的基本概念是连续性，但更本质地说，拓扑学可以被视为对“位置”的定性研究。在这方面发表的第一篇论文是欧拉(1741 年)对“哥尼斯堡七桥问题”的解答(另见比格斯、劳埃德和威尔逊，1976 年)。欧拉的论文标志着图论和拓扑学的基础；它引入了图的概念，并且节点的度起着核心作用。斯佩纳引理与图中奇数度节点的偶数性相关是恰当的。
\end{tcolorbox}

Theorem 7.10 is due to Sperner (1928), and was used by Knaster, Kuratowski, and Mazurkiewicz (1929) to prove the fixed-point theorem on a simplex with the help of their Theorem 7.12. As we mentioned in the proof of Theorem 7.10, these proofs use the counting modulo two in equation (7.12). The term "parity argument" has been coined by Papadimitriou (1994).

\begin{tcolorbox}
定理 7.10 归功于斯佩纳(1928 年)，并且被克纳斯特、库拉托夫斯基和马祖尔克维奇(1929 年)用于借助他们的定理 7.12 证明单纯形上的不动点定理。正如我们在定理 7.10 的证明中所提到的，这些证明使用了方程 (7.12) 中的模 2 计数。“奇偶性论证”这一术语由帕帕迪米特里欧(1994 年)创造。
\end{tcolorbox}

The path-following argument is due to Cohen (1967). Cohen's proof is not given in enough detail to see that the path may also come back to several dimensions further below, as in the example in Figure 7.10. For a more precise description see Kuhn (1969). We have defined the "internal" nodes of the graph as the sub-facets in \({C}^{\prime }\) rather than as the elements of \(A\) in (7.11). When staying in a fixed dimension, it is more natural to consider only sub-simplices, the elements of \(D\) and \(A\) , as nodes of a graph, which are joined by an edge if they have a common sub-facet in \({C}^{\prime }\) . However, for the construction of \({G}_{2}\) , in particular, the edges are more easily seen as line segments and the nodes as vertices in \({\sum }_{2}\) . A key argument in the construction of \(G\) in (7.13) is that the end nodes of some paths in a graph \({G}_{n}\) have different dimension, and we change to the higher dimension only at the end node.

\begin{tcolorbox}
路径跟踪的论证归功于科恩(Cohen，1967)。科恩的证明没有提供足够详细的内容，让人无法看到路径也可能回到低几个维度的情况，如图7.10中的示例所示。更精确的描述请参阅库恩(Kuhn，1969)。我们将图的“内部”节点定义为 \({C}^{\prime }\) 中的子极面，而不是(7.11)中 \(A\) 的元素。当处于固定维度时，更自然的做法是仅将单纯剖分中的子单纯形，即 \(D\) 和 \(A\) 的元素，视为图的节点，如果它们在 \({C}^{\prime }\) 中有一个公共子极面，则用一条边将它们连接起来。然而，特别是对于 \({G}_{2}\) 的构造，边更容易被看作线段，节点更容易被看作 \({\sum }_{2}\) 中的顶点。(7.13)中 \(G\) 构造的一个关键论证是，图 \({G}_{n}\) 中某些路径的端点具有不同的维度，并且我们仅在端点处转换到更高的维度。
\end{tcolorbox}

The image of "seaview hotels" to prove the Bolzano-Weierstrass theorem is taken from Bryant (1990).

\begin{tcolorbox}
用于证明波尔查诺 - 魏尔施特拉斯定理的“海景酒店”的例子取自布莱恩特(Bryant，1990)。
\end{tcolorbox}

The Freudenthal simplicial subdivision is due to Freudenthal (1942) and has been re-discovered several times, for example by Kuhn (1960).

\begin{tcolorbox}
弗罗伊登塔尔单纯剖分归功于弗罗伊登塔尔(Freudenthal，1942)，并且被多次重新发现，例如库恩(Kuhn，1960)。
\end{tcolorbox}

\section*{7.9 Exercises for Chapter 7}

\begin{tcolorbox}
\section*{7.9 第7章练习}
\end{tcolorbox}

Exercise 7.1. Let \(C = \left\{  {\left( {x,y}\right)  \in  {\mathbb{R}}^{2} \mid  {x}^{2} + {y}^{2} \leq  1}\right\}\) , and consider the following functions \(g,h,f : C \rightarrow  C\) :

\begin{tcolorbox}
练习7.1。设 \(C = \left\{  {\left( {x,y}\right)  \in  {\mathbb{R}}^{2} \mid  {x}^{2} + {y}^{2} \leq  1}\right\}\) ，并考虑以下函数 \(g,h,f : C \rightarrow  C\) :
\end{tcolorbox}

\[
g\left( {x,y}\right)  = \left( {x,\sqrt{1 - {x}^{2}}}\right) ,\;h\left( {x,y}\right)  = \left( {-y,x}\right) ,\;f\left( {x,y}\right)  = h\left( {g\left( {x,y}\right) }\right) .
\]

Let \(G = g\left( C\right)  = \{ g\left( {x,y}\right)  \mid  \left( {x,y}\right)  \in  C\}\) . Describe the geometric shape of \(C\) and \(G\) , and how the functions \(g\) and \(h\) act. Does each of the functions \(g,h,f\) have a fixed point in \(C\) ? Justify why, or why not. Determine all fixed points, if any, of each of \(g,h\) , and \(f\) .

\begin{tcolorbox}
设 \(G = g\left( C\right)  = \{ g\left( {x,y}\right)  \mid  \left( {x,y}\right)  \in  C\}\) 。描述 \(C\) 和 \(G\) 的几何形状，以及函数 \(g\) 和 \(h\) 的作用方式。每个函数 \(g,h,f\) 在 \(C\) 中是否有不动点？说明原因。如果有，确定 \(g,h\) 和 \(f\) 每个函数的所有不动点。
\end{tcolorbox}

Exercise 7.2. For a general polytope \(P\) , a vertex of \(P\) is defined as an extreme point of \(P\) , that is, it is not a convex combination of other points of \(P\) .

\begin{tcolorbox}
练习7.2。对于一般的有界多面体 \(P\) ， \(P\) 的顶点被定义为 \(P\) 的极点，即它不是 \(P\) 中其他点的凸组合。
\end{tcolorbox}

(a) Show that for a simplex \(S\) given as the convex hull of the affinely independent points \({v}^{1},\ldots ,{v}^{d}\) , each point \({v}^{i}\) is indeed a vertex of \(S\) according to this definition.

\begin{tcolorbox}
(a)证明对于由仿射无关点 \({v}^{1},\ldots ,{v}^{d}\) 的凸包给出的单纯形 \(S\) ，根据这个定义，每个点 \({v}^{i}\) 确实是 \(S\) 的一个顶点。
\end{tcolorbox}

(b) Show that the set of vertices of a simplex is unique.

\begin{tcolorbox}
(b)证明单纯形的顶点集是唯一的。
\end{tcolorbox}

Exercise 7.3. Give examples of

\begin{tcolorbox}
练习7.3。给出以下例子
\end{tcolorbox}

(a) two simplices whose intersection is not a simplex;

\begin{tcolorbox}
(a)两个单纯形，它们的交集不是单纯形；
\end{tcolorbox}

(b) a nonempty convex set that has no extreme points;

\begin{tcolorbox}
(b)一个非空凸集，它没有极点；
\end{tcolorbox}

(c) a bounded convex set that has at least one extreme point but that is not equal to the convex hull of its extreme points.

\begin{tcolorbox}
(c)一个有界凸集，它至少有一个极点，但它不等于其极点的凸包。
\end{tcolorbox}

Exercise 7.4. Show that the following generalization of Sperner's lemma (Theorem 7.10) is false: Let \(S\) be a simplex with vertices \({v}^{1},\ldots ,{v}^{d}\) , let \(\sum\) be a triangulation of \(S\) , and consider a labeling of the vertices in \(V\) so that \({v}_{i}\) has label \(i\) for \(1 \leq  i \leq  d\) . Then \(\sum\) has at least one completely labeled sub-simplex.

\begin{tcolorbox}
练习7.4。证明以下斯佩纳引理(定理7.10)的推广是错误的:设 \(S\) 是一个顶点为 \({v}^{1},\ldots ,{v}^{d}\) 的单纯形，设 \(\sum\) 是 \(S\) 的一个三角剖分，并考虑 \(V\) 中顶点的一个标签，使得 \({v}_{i}\) 具有标签 \(i\) ，其中 \(1 \leq  i \leq  d\) 。那么 \(\sum\) 至少有一个完全标记的子单纯形。
\end{tcolorbox}

\section*{8}

\section*{Zero-Sum Games}

\begin{tcolorbox}
\section*{零和博弈}
\end{tcolorbox}

Zero-sum games are games of two players where the interests of the players are directly opposed: One player's loss is the other player's gain. Competitions between two players in sports or in parlor games can be thought of as zero-sum games. The combinatorial games studied in Chapter 1 are also zero-sum games.

\begin{tcolorbox}
零和博弈是两个玩家的博弈，玩家的利益直接对立:一个玩家的损失就是另一个玩家的收益。体育比赛或室内游戏中两个玩家之间的竞争可以被视为零和博弈。第1章中研究的组合博弈也是零和博弈。
\end{tcolorbox}

In a zero-sum game, the sum of payoffs to the players in any outcome of playing the game is zero. An \(m \times  n\) bimatrix game(A, B)is zero-sum if \(B =  - A\) . It is completely specified by the matrix \(A\) of payoffs to the row player, and the game is therefore often called a matrix game. Consequently, we can analyse matrix games just like bimatrix games. However, zero-sum games have additional strong properties that do not hold for general bimatrix games. These properties are, which are the subject of this chapter, are:

\begin{tcolorbox}
在零和博弈中，博弈任何结果下参与者的收益总和为零。若 \(B =  - A\) ，则 \(m \times  n\) 双矩阵博弈 (A, B) 是零和博弈。它完全由行参与者的收益矩阵 \(A\) 确定，因此该博弈通常被称为矩阵博弈。因此，我们可以像分析双矩阵博弈一样分析矩阵博弈。然而，零和博弈具有一般双矩阵博弈所不具备的额外强大性质。本章的主题就是这些性质，它们是:
\end{tcolorbox}

\begin{itemize}
\item In a zero-sum game, an equilibrium strategy is a max-min strategy that maximizes the worst-case expected payoff to the player, and is therefore independent of what the opponent does.
\end{itemize}

\begin{tcolorbox}
 - 在零和博弈中，均衡策略是一种极大极小策略，它能使参与者的最坏情况期望效用最大化，因此与对手的行动无关。
\end{tcolorbox}

\begin{itemize}
\item Every equilibrium has the same unique payoff to the players, which is called the value of the game.
\end{itemize}

\begin{tcolorbox}
 - 每个均衡给参与者带来的收益都是相同且唯一的，这被称为博弈的值。
\end{tcolorbox}

\begin{itemize}
\item If a player has several equilibrium strategies, these are exchangeable and can be used in any equilibrium. Moreover, the set of a player's equilibrium strategies is convex.
\end{itemize}

\begin{tcolorbox}
 - 如果一个参与者有多个均衡策略，这些策略是可交换的，并且可以在任何均衡中使用。此外，参与者的均衡策略集是凸集。
\end{tcolorbox}

Apart from these special properties, we treat zero-sum games in a chapter of its own, because:

\begin{tcolorbox}
除了这些特殊性质外，我们单独用一章来讨论零和博弈，原因如下:
\end{tcolorbox}

\begin{itemize}
\item Historically, the main result about zero-sum games of von Neumann (1928) precedes the equilibrium concept of Nash (1951) for games with general payoffs.
\end{itemize}

\begin{tcolorbox}
 - 从历史上看，冯·诺伊曼(1928 年)关于零和博弈的主要结果早于纳什(1951 年)针对一般收益博弈提出的均衡概念。
\end{tcolorbox}

\begin{itemize}
\item The concept of equilibrium as "no incentive to deviate" has a stronger form in a zero-sum game, namely that a player can make his strategy public without a disadvantage. As the games of Matching Pennies and Rock-Paper-Scissors show (see Figure 3.12), this requires the concept of mixed strategies. (Zero-sum games may not always require mixed strategies, for example if they have perfect information like the games in Chapters 1 and 4. A game where one player has only two strategies is dominance solvable, see Exercise 8.2.)
\end{itemize}

\begin{tcolorbox}
 - “没有偏离动机”这一均衡概念在零和博弈中有更强的形式，即参与者可以公开自己的策略而不会处于劣势。正如硬币匹配 (博弈) 和石头剪刀布博弈所示(见图 3.12)，这需要混合策略的概念。(零和博弈并不总是需要混合策略，例如，如果它们像第 1 章和第 4 章中的博弈那样具有完美信息。如果一个参与者只有两种策略，那么该博弈是占优可解的，见练习 8.2。)
\end{tcolorbox}

\begin{itemize}
\item This chapter is mathematically self-contained, with an independent two-page proof of von Neumann's minimax theorem, which (see Proposition 8.3) is equivalent to the existence of an equilibrium in a zero-sum game. The minimax theorem is less powerful than Brouwer's fixed-point theorem, and slightly less powerful than the duality theorem of linear programming, but will be useful for proving the existence of "correlated equilibria" in Chapter 12.
\end{itemize}

\begin{tcolorbox}
 - 本章在数学上是自足的，独立用两页篇幅证明了冯·诺伊曼的极小极大定理，该定理(见命题 8.3)等价于零和博弈中均衡的存在性。极小极大定理的效力不如布劳威尔不动点定理，也略逊于线性规划的对偶定理，但在第 12 章证明“相关均衡”的存在性时会很有用。
\end{tcolorbox}

In Section 8.2, we analyze the soccer penalty as an introductory example of a zero-sum game, and show how to find a max-min strategy graphically (Figures 8.2 and 8.3).

\begin{tcolorbox}
在 8.2 节中，我们以足球点球为例介绍零和博弈，并展示如何通过图形找到极大极小策略(图 8.2 和 8.3)。
\end{tcolorbox}

Section 8.3 states the important properties of zero-sum games. The central concept is that of a max-min strategy, which is a (typically mixed) strategy that maximizes the minimum expected payoff to the player for all possible responses of the opponent. This (rather paranoid) approach can be applied to any game, but in a zero-sum game means that the other player uses a best response, and is therefore justified.

\begin{tcolorbox}
8.3 节阐述了零和博弈的重要性质。核心概念是极大极小策略，它是一种(通常是混合的)策略，能使参与者在对手所有可能反应下的最小期望效用最大化。这种(相当偏执的)方法可以应用于任何博弈，但在零和博弈中意味着另一个参与者使用最优反应，因此是合理的。
\end{tcolorbox}

The payoffs to the maximizing row player in a matrix game are costs to the column player who wants to minimize them. In terms of these costs (rather than the payoffs to the column player), a max-min strategy is called a min-max strategy of the minimizer. If the max-min payoff equals the min-max cost, then the max-min strategy and min-max strategy define an equilibrium, and an equilibrium is such a pair of strategies (Proposition 8.3). Because finite games have an equilibrium, the maxmin \(=\) minmax equality holds, known as the minimax theorem of von Neumann (1928) (Theorem 8.3). It also defines the unique equilibrium payoff (to the maximizer, and cost to the minimizer) as the value of the game.

\begin{tcolorbox}
矩阵博弈中最大化收益的行参与者的收益，对于想要最小化这些收益的列参与者来说就是成本。就这些成本(而非列参与者的收益)而言，极大极小策略被称为最小值点的极小极大策略。如果极大极小收益等于极小极大成本，那么极大极小策略和极小极大策略就定义了一个均衡，并且均衡就是这样一对策略(命题 8.3)。因为有限博弈存在均衡，所以极大极小 \(=\) 极小极大等式成立，这就是冯·诺伊曼(1928 年)的极小极大定理(定理 8.3)。它还将唯一的均衡收益(对最大值点而言)和成本(对最小值点而言)定义为博弈的值。
\end{tcolorbox}

In Section 8.4 we give a short proof of the minimax theorem that is independent of the existence of an equilibrium. It connects two related optimization problems that are a special case of pair of "dual linear programs". We consider them only for this special case, where we prove directly, by induction, that the two problems have the same optimal value.

\begin{tcolorbox}
在 8.4 节中，我们给出了一个不依赖于均衡存在性的极小极大定理的简短证明。它将两个相关的优化问题联系起来，这两个问题是一对“对偶线性规划”的特殊情况。我们仅针对这个特殊情况考虑它们，在这种情况下，我们通过归纳法直接证明这两个问题具有相同的最优值。
\end{tcolorbox}

In zero-sum game, an equilibrium strategy is the same as a max-min strategy (in terms of the player's own payoff), also called an optimal strategy, which does not depend on what the other player does, a very strong property. Further properties of optimal strategies are discussed in Section 8.5.

\begin{tcolorbox}
在零和博弈中，均衡策略与极大极小策略(就参与者自身的收益而言)相同，也称为最优策略，它不依赖于其他参与者的行动，这是一个非常强的性质。最优策略的其他性质将在 8.5 节中讨论。
\end{tcolorbox}

\section*{8.1 Prerequisites and Learning Outcomes}

\begin{tcolorbox}
\section*{8.1 先决条件和学习成果}
\end{tcolorbox}

As prerequisites we require the definition of bimatrix games and mixed strategies from Chapter 6, and manipulating matrices and linear inequalities between vectors. The proof of the minimax theorem in Section 8.4 uses that a continuous function on a compact domain assumes its minimum and maximum.

\begin{tcolorbox}
作为先决条件，我们需要第6章中双矩阵博弈(bimatrix game)和混合策略(mixed strategy)的定义，以及矩阵运算和向量间的线性不等式知识。8.4节中极小极大定理(minimax theorem)的证明用到了紧域上的连续函数能取到最小值和最大值这一性质。
\end{tcolorbox}

After studying this chapter, you should be able to:

\begin{tcolorbox}
学习本章后，你应该能够:
\end{tcolorbox}

\begin{itemize}
\item Explain the concept of max-min strategies and their importance for zero-sum games, and apply this to given games.
\end{itemize}

\begin{tcolorbox}
\begin{itemize}
\item 解释极大极小策略(max - min strategy)的概念及其对零和博弈(zero - sum game)的重要性，并将其应用于给定的博弈中。
\end{itemize}
\end{tcolorbox}

\begin{itemize}
\item Find max-min and min-max strategies in small matrix games like in Exercise 8.1, and exploit the "maxmin = minmax" theorem to find the (max-min or min-max) strategy that is easier to find via the lower or upper envelope of expected payoffs.
\end{itemize}

\begin{tcolorbox}
\begin{itemize}
\item 像在练习8.1中那样，在小型矩阵博弈(matrix game)中找到极大极小和极小极大策略(min - max strategy)，并利用“极大极小值 = 极小极大值”定理，通过期望效用(expected payoff)的下包络或上包络来找到更容易计算的(极大极小或极小极大)策略。
\end{itemize}
\end{tcolorbox}

\begin{itemize}
\item State and prove the equivalence of existence of an equilibrium and the minimax theorem in a zero-sum game (Proposition 8.3).
\end{itemize}

\begin{tcolorbox}
\begin{itemize}
\item 陈述并证明零和博弈中均衡(equilibrium)的存在性与极小极大定理的等价性(命题8.3)。
\end{itemize}
\end{tcolorbox}

\begin{itemize}
\item State the main properties of zero-sum games: independence of "optimal strategies" from the what the opponent does, and its implications (unique value, exchangeability, and convexity of the set of equilibrium strategies).
\end{itemize}

\begin{tcolorbox}
\begin{itemize}
\item 陈述零和博弈的主要性质:“最优策略”(optimal strategy)与对手的行动无关，以及其影响(唯一值、可交换性和均衡策略集的凸性)。
\end{itemize}
\end{tcolorbox}

\begin{itemize}
\item Explain the main idea behind the induction proof of the minimax theorem in Section 8.3.
\end{itemize}

\begin{tcolorbox}
\begin{itemize}
\item 解释8.3节中极小极大定理归纳证明背后的主要思想。
\end{itemize}
\end{tcolorbox}

\section*{8.2 Example: Soccer Penalty}

\begin{tcolorbox}
\section*{8.2 示例:足球点球}
\end{tcolorbox}

A zero-sum game is a bimatrix game(A, B)with \(B =  - A\) . Because it is completely specified by the matrix \(A\) of payoffs to the row player \(\mathrm{I}\) , it is also called a matrix game (with payoff matrix) \(A\) . When the players choose row \(i\) and column \(j\) , respectively, player II pays the matrix entry \({a}_{ij}\) to player I, which is therefore a cost that player II wants to minimize, also in expectation when player II faces a mixed strategy. Player I is called the maximizer and player II the minimizer.

\begin{tcolorbox}
零和博弈是一个双矩阵博弈(A, B)，其中 \(B =  - A\) 。由于它完全由行玩家 \(\mathrm{I}\) 的收益矩阵 \(A\) 确定，因此它也被称为(具有收益矩阵的)矩阵博弈 \(A\) 。当玩家分别选择行 \(i\) 和列 \(j\) 时，玩家II向玩家I支付矩阵元素 \({a}_{ij}\) ，因此这是玩家II想要最小化的成本(cost)，当玩家II面对混合策略时，在期望意义上也是如此。玩家I被称为最大值点(maximizer)，玩家II被称为最小值点(minimizer)。
\end{tcolorbox}

Figure 8.1 shows the penalty kick in soccer (football in Britain) as a zero-sum game between striker and goalkeeper. For simplicity, assume that the striker's possible strategies are \(L\) and \(R\) , that is, to kick into the left or right corner, and that the goalkeeper’s strategies are \(l,w\) , and \(r\) , where \(l\) and \(r\) mean that he jumps immediately into the left or right corner, respectively, and \(w\) that he waits and sees where the ball goes and jumps afterwards. In Figure 8.2, we give an example of resulting probabilities that the striker scores a goal, depending on the choices of the two players. Clearly, the striker (player I) tries to maximize and the goalkeeper (player II) to minimize that probability, and we only need this payoff to player I to specify the game. We call the maximizer "Max" and the minimizer "min" (where we keep our tradition of using lower case letters for the column player).

\begin{tcolorbox}
图8.1将足球(在英国称为football)中的点球大战展示为前锋和守门员之间的零和博弈。为简单起见，假设前锋的可能策略是 \(L\) 和 \(R\) ，即踢向左侧或右侧角落，而守门员的策略是 \(l,w\) 和 \(r\) ，其中 \(l\) 和 \(r\) 分别表示他立即跳向左侧或右侧角落， \(w\) 表示他等待观察球的去向然后再跳跃。在图8.2中，我们给出了一个根据双方玩家的选择，前锋得分的概率示例。显然，前锋(玩家I)试图最大化该概率，而守门员(玩家II)试图最小化该概率，我们只需要这个给玩家I的收益就能确定这个博弈。我们称最大值点为“Max”，最小值点为“min”(我们沿用用小写字母表示列玩家的传统)。
\end{tcolorbox}

\begin{center}
\includegraphics[max width=0.9\textwidth]{images/bo_d2e5cmbef24c73b03tn0_219_340_202_1242_534_0.jpg}
\end{center}
\hspace*{3em} 

Figure 8.1 Strategies for striker and goalkeeper in the soccer Penalty game. In addition to jumping left or right, the goalkeeper can wait(w). An extra strategy for the striker of shooting down the middle is considered in Exercise 8.1.

\begin{tcolorbox}
图8.1 足球点球大战中前锋和守门员的策略。除了向左或向右跳跃，守门员还可以等待(w)。练习8.1中考虑了前锋向中间射门的额外策略。
\end{tcolorbox}

Because the players care about the probability that a goal is scored, this payoff (to player I, cost to player II) automatically fulfills the condition of an expected-utility function, because an expected probability is just a probability. (This is a nice "practical" example of the construction of the values of an expected-utility function as probabilities, see Figure 5.1.) In a general zero-sum game, payoffs and costs are not necessarily defined via probabilities, but are nevertheless assumed to represent the preference of both players with respect to taking expectations. This is a stronger assumption than in a non-zero-sum game where the payoffs can be adjusted independently to reflect different attitudes to risk.

\begin{tcolorbox}
因为玩家关心进球的概率，所以这个收益(给玩家I，成本给玩家II)自动满足期望效用函数(expected - utility function)的条件，因为期望概率就是一个概率。(这是一个很好的将期望效用函数的值构建为概率的“实际”例子，见图5.1。)在一般的零和博弈中，收益和成本不一定通过概率来定义，但仍然假设它们代表了两个玩家在取期望方面的偏好(preference)。这比非零和博弈中的假设更强，在非零和博弈中，收益可以独立调整以反映不同的风险态度。
\end{tcolorbox}

The top left payoff table in Figure 8.2 shows the game as a bimatrix game with two payoffs per cell, and the pure best-response payoffs marked as usual with boxes. The top right shows the more concise matrix game with only one number per cell. A best-response payoff for the maximizer is shown with a box, a best-response (minimal) cost for the minimizer with a circle. This represents the same information as the two payoffs per cell, but more concisely.

\begin{tcolorbox}
图8.2左上角的收益表将该博弈展示为一个双矩阵博弈，每个单元格有两个收益，纯最优反应(best response)收益像往常一样用方框标记。右上角展示了更简洁的矩阵博弈，每个单元格只有一个数字。最大值点的最优反应收益用方框表示，最小值点的最优反应(最小)成本用圆圈表示。这与每个单元格有两个收益所表示的信息相同，但更简洁。
\end{tcolorbox}

The best responses show that the game has no pure-strategy equilibrium. This is immediate from the sports situation: obviously, the goalkeeper's best response to a kick into the right or left corner is to jump into that corner, whereas the striker would rather kick into the opposite corner to where the goalkeeper jumps. We therefore have to find a mixed equilibrium, and active randomization is obviously advantageous for both players.

\begin{tcolorbox}
最优反应表明该博弈没有纯策略均衡。从体育情境中可以立即看出这一点:显然，守门员对踢向右侧或左侧角落的球的最优反应是跳向那个角落，而前锋更愿意踢向守门员跳跃方向的对角。因此，我们必须找到一个混合均衡(mixed equilibrium)，显然主动随机化对两个玩家都有利。
\end{tcolorbox}

We first find the equilibria of the game as we would do this for a \(2 \times  3\) game, except that we only use the payoffs to player I. It is more than appropriate to use the goalpost method for solving the Penalty game! The bottom picture in Figure 8.2 shows the expected payoffs to player I as a function of his probability, say \(p\) , of choosing \(R\) , for the reponses \(l,w\) , and \(r\) of player II. These payoffs to player I are costs to player II, so the best-response costs to player II are given by the minimum of these lines, which defines the lower envelope shown by bold line segments. There are two intersection points on the lower envelope, of the lines for \(l\) and \(w\) when \(p = \frac{1}{3}\) , and of the lines for \(w\) and \(r\) when \(p = \frac{3}{5}\) . These intersection points are easily found with the difference trick (where it does not matter if one considers costs or payoffs, because the differences are the same in absolute value): Namely, \(l\) and \(w\) have cost difference 0.2 in row \(L\) and 0.4 in row \(R\) , so that \(L\) and \(R\) should be played in the ratio \({0.4} : {0.2}\) , that is, for \(p = \frac{1}{3}\) , to get indifference between \(l\) and \(w\) . The indifference between \(w\) and \(r\) when \(p = \frac{3}{5}\) is found similarly.

\begin{tcolorbox}
我们首先像求解双矩阵博弈(bimatrix game)那样找出该博弈的均衡，只不过我们仅考虑玩家I的效用。使用球门柱法(goalpost method)来求解点球博弈再合适不过了！图8.2的下图展示了玩家I的期望效用，它是玩家I选择 \(R\) 的概率(设为 \(p\) )的函数，对应玩家II的策略 \(l,w\) 和 \(r\) 。玩家I的这些效用对玩家II来说就是成本，因此玩家II的最优反应成本由这些直线的最小值决定，这就定义了由粗线段表示的下包络(lower envelope)。在下包络上有两个交点，当 \(p = \frac{1}{3}\) 时， \(l\) 和 \(w\) 对应的直线相交；当 \(p = \frac{3}{5}\) 时， \(w\) 和 \(r\) 对应的直线相交。使用差分技巧(difference trick)很容易找到这些交点(在这种情况下，考虑成本还是效用并不重要，因为差值的绝对值是相同的):即， \(l\) 和 \(w\) 在第 \(L\) 行的成本差为0.2，在第 \(R\) 行的成本差为0.4，因此应该以 \({0.4} : {0.2}\) 的比例选择 \(L\) 和 \(R\) ，也就是当 \(p = \frac{1}{3}\) 时，使得玩家I在 \(l\) 和 \(w\) 之间无差异(indifference)。同理可找到当 \(p = \frac{3}{5}\) 时， \(w\) 和 \(r\) 之间的无差异情况。
\end{tcolorbox}

\begin{center}
\includegraphics[max width=1.0\textwidth]{images/bo_d2e5cmbef24c73b03tn0_220_188_190_1276_1040_0.jpg}
\end{center}
\hspace*{3em} 

Figure 8.2 Top left: the Penalty game as a zero-sum bimatrix game with scoring probabilities as payoffs to the striker (player I), which are the negative of the payoffs to the goalkeeper (player II). Top right: the same game as a matrix game, which shows only the payoffs to the maximizer, which are costs to the minimizer. A box indicates a best-response payoff to the maximizer, a circle a best-response (minimal) cost to the minimizer. Bottom: lower envelope of best-response costs as a function of the mixed strategy(1 - p, p)of Max. For \(p = \widehat{p} = \frac{1}{3}\) this is his max-min strategy.

\begin{tcolorbox}
图8.2 左上方:点球博弈作为一个零和双矩阵博弈，得分概率作为前锋(玩家I)的效用，而这是守门员(玩家II)效用的负值。右上方:同样的博弈表示为矩阵博弈(matrix game)，只显示了最大值点(maximizer)的效用，而这是最小值点(minimizer)的成本。方框表示最大值点的最优反应效用，圆圈表示最小值点的最优反应(最小)成本。下方:作为最大值点混合策略(1 - p, p)的函数，最优反应成本的下包络。对于 \(p = \widehat{p} = \frac{1}{3}\) ，这是他的极大极小策略(max - min strategy)。
\end{tcolorbox}

Hence, player II’s best response to the mixed strategy(1 - p, p)is \(l\) for \(0 \leq  p \leq  \frac{1}{3}\) , and \(w\) for \(\frac{1}{3} \leq  p \leq  \frac{3}{5}\) , and \(r\) for \(\frac{3}{5} \leq  p \leq  1\) . Only for \(p = \frac{1}{3}\) and for \(p = \frac{3}{5}\) can player II mix, and only in the first case, when \(l\) and \(w\) are best responses, can player I, in turn, be made indifferent. This gives the unique mixed equilibrium of the game, \(\left( {\left( {\frac{2}{3},\frac{1}{3}}\right) ,\left( {\frac{1}{6},\frac{5}{6},0}\right) }\right)\) . The resulting expected costs for the three columns \(l,w,r\) are \(\frac{2}{3},\frac{2}{3},\frac{4}{5}\) , where indeed player II assigns positive probability only to the smallest-cost columns \(l\) and \(w\) . The expected payoffs for both rows are \(\frac{2}{3}\) . (This unique equilibrium payoff of the game is the value of the game defined in Theorem 8.4.)

\begin{tcolorbox}
因此，玩家II对混合策略(1 - p, p)的最优反应是:当 \(0 \leq  p \leq  \frac{1}{3}\) 时为 \(l\) ，当 \(\frac{1}{3} \leq  p \leq  \frac{3}{5}\) 时为 \(w\) ，当 \(\frac{3}{5} \leq  p \leq  1\) 时为 \(r\) 。只有当 \(p = \frac{1}{3}\) 和 \(p = \frac{3}{5}\) 时，玩家II可以采用混合策略，并且只有在第一种情况下，即当 \(l\) 和 \(w\) 是最优反应时，才能让玩家I反过来达到无差异。这就得到了该博弈唯一的混合均衡(mixed equilibrium) \(\left( {\left( {\frac{2}{3},\frac{1}{3}}\right) ,\left( {\frac{1}{6},\frac{5}{6},0}\right) }\right)\) 。三列 \(l,w,r\) 对应的期望成本为 \(\frac{2}{3},\frac{2}{3},\frac{4}{5}\) ，实际上玩家II只对成本最小的列 \(l\) 和 \(w\) 赋予正概率。两行的期望效用均为 \(\frac{2}{3}\) 。(该博弈唯一的均衡效用就是定理8.4中定义的博弈值。)
\end{tcolorbox}

So far, our analysis is not new. We now consider the lower envelope of expected payoffs to player I as a function of his own mixed strategy \(x = \left( {1 - p,p}\right)\) from that player's own perspective: It represents the worst possible expected payoff to player I, given the possible responses of player II. It is easy to see that it suffices to look at pure responses of player II (see also Proposition 8.2 below). That is, against the mixed strategy \(x\) the minimum is found over all pure responses of the minimizer, here columns \(l,w\) , and \(r\) .

\begin{tcolorbox}
到目前为止，我们的分析并无新意。现在我们从玩家I自身的角度，考虑玩家I的期望效用作为其自身混合策略 \(x = \left( {1 - p,p}\right)\) 的函数的下包络:它表示在玩家II可能的反应下，玩家I可能获得的最差期望效用。很容易看出，只需要考虑玩家II的纯策略反应即可(也可参见下面的命题8.2)。也就是说，对于混合策略 \(x\) ，最小值是在最小值点的所有纯策略反应(这里是列 \(l,w\) 和 \(r\) )中找到的。
\end{tcolorbox}

If player I wants to "secure" the maximum of this worst possible payoff, he should maximize over the lower envelope with a suitable choice of \(p\) . This maximum is easily found for \(\widehat{p} = \frac{1}{3}\) , shown by a dashed line in Figure 8.2. The strategy \(\widehat{x} = \left( {1 - \widehat{p},\widehat{p}}\right)\) is called the max-min strategy of the maximizer, which we found here graphically as the highest point on the lower envelope. The "hat" accent \(\widehat{x}\) may be thought of as representing that "peak".

\begin{tcolorbox}
如果玩家I想“确保”这种最差可能收益的最大值，他应该通过合适地选择 \(p\) 来最大化下包络。对于 \(\widehat{p} = \frac{1}{3}\) ，这个最大值很容易找到，如图8.2中的虚线所示。策略 \(\widehat{x} = \left( {1 - \widehat{p},\widehat{p}}\right)\) 被称为最大值点的极大极小策略，我们在这里通过图形将其确定为下包络上的最高点。“帽子”符号 \(\widehat{x}\) 可以被认为代表那个“峰值”。
\end{tcolorbox}

The max-min strategy \(\left( {1 - \widehat{p},\widehat{p}}\right)\) is one of the intersection points on the lower envelope. The intersection point at \(p = \frac{3}{5}\) has a lower payoff to the maximizer because the best-response line \(w\) slopes downwards from \(\frac{1}{3}\) to \(\frac{3}{5}\) . Hence, \(p = \frac{3}{5}\) is not a max-min strategy. As we show in Proposition 8.3 below, in a zero-sum game an equilibrium strategy is the same as a max-min strategy. Therefore, unlike in a bimatrix game, \(p = \frac{3}{5}\) does not need to be checked as a possible equilibrium strategy.

\begin{tcolorbox}
极大极小策略 \(\left( {1 - \widehat{p},\widehat{p}}\right)\) 是下包络上的一个交点。在 \(p = \frac{3}{5}\) 处的交点给最大值点带来的收益较低，因为最优反应线 \(w\) 从 \(\frac{1}{3}\) 到 \(\frac{3}{5}\) 是向下倾斜的。因此， \(p = \frac{3}{5}\) 不是极大极小策略。正如我们在下面的命题8.3中所示，在零和博弈中，均衡策略与极大极小策略是相同的。因此，与双矩阵博弈不同，不需要将 \(p = \frac{3}{5}\) 作为可能的均衡策略进行检验。
\end{tcolorbox}

The max-min payoff is well defined and does not depend on what player II does, because we always take the minimum over player II's possible responses. The maximum (and thus the max-min payoff) is unique, but there may be more than one choice of the max-min strategy \(\widehat{x}\) . An example is shown in Figure 8.3 where both \(L\) and \(R\) have the same payoff 0.6 in column \(w\) (and are therefore both best responses to \(w\) , so that this game is degenerate according to Definition 6.7). Then the lower envelope has its maximum 0.6 anywhere for \(\widehat{p} \in  \left\lbrack  {\frac{1}{5},\frac{2}{3}}\right\rbrack\) , which defines the possible max-min strategies \(\widehat{x} = \left( {1 - \widehat{p},\widehat{p}}\right)\) ; the endpoints of the interval are found with the difference trick for the indifferences between \(l\) and \(w\) , and \(w\) and \(r\) , respectively. On the side of player II, the unique best response against \(\widehat{p} = \frac{1}{2}\) (for example) is the pure strategy \(w\) . Therefore, without any further considerations, \(w\) is the unique equilibrium strategy of player II in any equilibrium; this is a very useful consequence of Proposition 8.3, as we will explain.

\begin{tcolorbox}
极大极小收益是明确定义的，并且不依赖于玩家II的行动，因为我们总是取玩家II可能反应的最小值。最大值(也就是极大极小收益)是唯一的，但极大极小策略 \(\widehat{x}\) 可能有不止一种选择。图8.3展示了一个例子，其中 \(L\) 和 \(R\) 在列 \(w\) 中的收益都是0.6(因此它们都是对 \(w\) 的最优反应，根据定义6.7，这个博弈是退化的)。那么下包络在 \(\widehat{p} \in  \left\lbrack  {\frac{1}{5},\frac{2}{3}}\right\rbrack\) 的任何位置都有最大值0.6，这就确定了可能的极大极小策略 \(\widehat{x} = \left( {1 - \widehat{p},\widehat{p}}\right)\) ；该区间的端点分别通过差分技巧来确定 \(l\) 和 \(w\) 以及 \(w\) 和 \(r\) 之间的无差异来得到。在玩家II这边，针对 \(\widehat{p} = \frac{1}{2}\) (例如)的唯一最优反应是纯策略 \(w\) 。因此，无需进一步考虑， \(w\) 是玩家II在任何均衡中的唯一均衡策略；正如我们将解释的，这是命题8.3的一个非常有用的结果。
\end{tcolorbox}

\begin{center}
\includegraphics[max width=0.9\textwidth]{images/bo_d2e5cmbef24c73b03tn0_222_204_216_1238_574_0.jpg}
\end{center}
\hspace*{3em} 

Figure 8.3 The Penalty game with the same scoring probability 0.6 when the column player chooses \(w\) . Then the game is degenerate, and the max-min strategy \(\left( {1 - \widehat{p},\widehat{p}}\right)\) is not unique, \(\widehat{p} \in  \left\lbrack  {\frac{1}{5},\frac{2}{3}}\right\rbrack\)

\begin{tcolorbox}
图8.3 当列玩家选择 \(w\) 时，罚球博弈的得分概率相同，均为0.6。此时博弈是退化的，极大极小策略 \(\left( {1 - \widehat{p},\widehat{p}}\right)\) 不是唯一的， \(\widehat{p} \in  \left\lbrack  {\frac{1}{5},\frac{2}{3}}\right\rbrack\)
\end{tcolorbox}

\(\Rightarrow\) Exercise 8.1 is a variant of the Penalty game where the striker is more likely to score, and in part (b) has a third strategy \(M\) of shooting down the middle.

\begin{tcolorbox}
\(\Rightarrow\) 练习8.1是罚球博弈的一个变体，其中射手更有可能得分，并且在(b)部分有第三种策略 \(M\) ，即向中间射门。
\end{tcolorbox}

\section*{8.3 Max-Min and Min-Max Strategies}

\begin{tcolorbox}
\section*{8.3 极大极小和极小极大策略}
\end{tcolorbox}

The important new concept for zero-sum games is that of a max-min strategy. This concept applies in principle to any game, where we focus on two-player games. A max-min strategy is a strategy of a player that is defined in terms of the player's own payoffs. It also depends on the strategy set of the player. In the following definition, we assume general (nonempty) strategy sets \(X\) and \(Y\) , which are normally, but don't have to be, the mixed-strategy sets in (6.6). A min-max strategy is a strategy of the other player.

\begin{tcolorbox}
零和博弈的一个重要新概念是极大极小策略。这个概念原则上适用于任何博弈，这里我们关注两人博弈。极大极小策略是根据玩家自身收益来定义的玩家策略。它也依赖于玩家的策略集。在下面的定义中，我们假设一般的(非空)策略集 \(X\) 和 \(Y\) ，它们通常但不必须是(6.6)中的混合策略集。极小极大策略是另一个玩家的策略。
\end{tcolorbox}

Definition 8.1 (Max-min strategy, min-max strategy). Consider a two-player game with strategy set \(X\) for player I and \(Y\) for player II, and payoff \(a\left( {x,y}\right)\) to player I for \(x \in  X\) and \(y \in  Y\) . A max-min strategy \(\widehat{x} \in  X\) of player I fulfills

\begin{tcolorbox}
定义8.1(极大极小策略，极小极大策略)。考虑一个双人博弈，玩家I的策略集为 \(X\) ，玩家II的策略集为 \(Y\) ，当玩家I选择 \(x \in  X\) 且玩家II选择 \(y \in  Y\) 时，玩家I的收益为 \(a\left( {x,y}\right)\) 。玩家I的极大极小策略 \(\widehat{x} \in  X\) 满足
\end{tcolorbox}

\[
\mathop{\min }\limits_{{y \in  Y}}a\left( {\widehat{x},y}\right)  = \mathop{\max }\limits_{{x \in  X}}\mathop{\min }\limits_{{y \in  Y}}a\left( {x,y}\right) , \tag{8.1}
\]

assuming these minima and maxima exist, which defines the max-min payoff to player I. A min-max strategy \(\widehat{y} \in  Y\) of player II fulfills

\begin{tcolorbox}
假设这些最小值和最大值存在，这就定义了玩家I的极大极小收益。玩家II的极小极大策略 \(\widehat{y} \in  Y\) 满足
\end{tcolorbox}

\[
\mathop{\max }\limits_{{x \in  X}}a\left( {x,\widehat{y}}\right)  = \mathop{\min }\limits_{{y \in  Y}}\mathop{\max }\limits_{{x \in  X}}a\left( {x,y}\right) , \tag{8.2}
\]

which defines the min-max payoff to player I.

\begin{tcolorbox}
这就定义了玩家I的极小极大收益。
\end{tcolorbox}

Consider first a finite game, for example the game at the top right in Figure 8.2 with the \(2 \times  3\) payoff matrix \(A\) to the row player \(I\) . With the strategy sets \(X = \{ L,R\}\) and \(Y = \{ l,w,r\}\) in Definition 8.1, the minima and maxima in (8.1) and (8.2) clearly exist. Because player I cannot randomize, he has to consider the minimum payoff against \(L\) (which is 0.5) or \(R\) (which is 0.4, both indicated by the circles), so his max-min strategy is \(L\) with max-min payoff 0.5 . For the min-max strategy, the column player II has to consider the maximum payoff to player I (indicated by the boxes) against \(l,w\) , and \(r\) , which is 1.0,0.7, and 1.0, respectively, so the min-max strategy is \(w\) with min-max payoff 0.7. Here, max-min and min-max payoff are different.

\begin{tcolorbox}
首先考虑一个有限博弈，例如图8.2右上角的博弈，行玩家 \(I\) 的 \(2 \times  3\) 收益矩阵为 \(A\) 。根据定义8.1中的策略集 \(X = \{ L,R\}\) 和 \(Y = \{ l,w,r\}\) ，式(8.1)和(8.2)中的最小值和最大值显然存在。由于玩家I不能随机化策略，他必须考虑针对 \(L\) (收益为0.5)或 \(R\) (收益为0.4，两者都用圆圈表示)的最小收益，因此他的极大极小策略是 \(L\) ，极大极小收益为0.5。对于极小极大策略，列玩家II必须考虑玩家I针对 \(l,w\) 和 \(r\) 的最大收益(用方框表示)，分别为1.0、0.7和1.0，因此极小极大策略是 \(w\) ，极小极大收益为0.7。这里，极大极小收益和极小极大收益不同。
\end{tcolorbox}

We normally consider the mixed extension of a finite game. Then \(X\) and \(Y\) are the sets of mixed strategies, and the payoff to player I in Definition 8.1 is given by \(a\left( {x,y}\right)  = {x}^{\top }{Ay}\) . With the goalpost diagram in Figure 8.2, we have identified the (mixed) max-min strategy of player \(\mathrm{I}\) as \(\widehat{x} = \left( {\frac{2}{3},\frac{1}{3}}\right)\) with max-min payoff \(\frac{2}{3}\) . We will shortly see that the min-max strategy player II is her equilibrium strategy \(\left( {\frac{1}{6},\frac{5}{6},0}\right)\) , with the min-max payoff to player I also given by \(\frac{2}{3}\) .

\begin{tcolorbox}
我们通常考虑有限博弈的混合扩展。此时 \(X\) 和 \(Y\) 是混合策略集，定义8.1中玩家I的收益由 \(a\left( {x,y}\right)  = {x}^{\top }{Ay}\) 给出。根据图8.2中的门柱图，我们已确定玩家 \(\mathrm{I}\) 的(混合)极大极小策略为 \(\widehat{x} = \left( {\frac{2}{3},\frac{1}{3}}\right)\) ，极大极小收益为 \(\frac{2}{3}\) 。我们很快会看到，玩家II的极小极大策略是她的均衡策略 \(\left( {\frac{1}{6},\frac{5}{6},0}\right)\) ，玩家I的极小极大收益同样为 \(\frac{2}{3}\) 。
\end{tcolorbox}

The rationale behind a max-min strategy is that it gives a guaranteed (expected) payoff to the player, irrespective of what the other player does. (A max-min strategy is therefore sometimes called a "security strategy".) If player I plays \(x\) , his worst-case payoff is \(\mathop{\min }\limits_{{y \in  Y}}a\left( {x,y}\right)\) , which is maximized by player I’s max-min strategy \(\widehat{x}\left( {8.1}\right)\) and gives him his max-min payoff. In a game with general payoffs, player II may not necessarily want to play the minimizing strategy \(y\) against \(x\) in \(\mathop{\min }\limits_{{y \in  Y}}a\left( {x,y}\right)\) . However, in a zero-sum game, this strategy \(y\) is player II’s best response to \(x\) , so this is what player I should expect.

\begin{tcolorbox}
极大极小策略背后的基本原理是，无论另一个玩家采取什么行动，它都能为玩家提供有保证的(期望)收益。(因此，极大极小策略有时也被称为“安全策略”。)如果玩家I选择 \(x\) ，他的最坏情况收益是 \(\mathop{\min }\limits_{{y \in  Y}}a\left( {x,y}\right)\) ，玩家I的极大极小策略 \(\widehat{x}\left( {8.1}\right)\) 使该收益最大化，并为他带来极大极小收益。在一般收益的博弈中，玩家II不一定想在 \(\mathop{\min }\limits_{{y \in  Y}}a\left( {x,y}\right)\) 中针对 \(x\) 采用最小化策略 \(y\) 。然而，在零和博弈中，该策略 \(y\) 是玩家II对 \(x\) 的最优反应，因此这是玩家I应该预期到的。
\end{tcolorbox}

In a min-max strategy \(\widehat{y}\) in (8.2), player II chooses \(\widehat{y}\) under the assumption that player I chooses a best response \(x\) to achieve the payoff \(\mathop{\max }\limits_{{x \in  X}}a\left( {x,\widehat{y}}\right)\) . By playing \(\widehat{y}\) , player II minimizes this payoff. In a game with general payoffs, this is a rather malevolent action by player II and not justified by any security considerations in terms of her own payoffs.

\begin{tcolorbox}
在式(8.2)的极小极大策略 \(\widehat{y}\) 中，玩家II在假设玩家I选择最优反应 \(x\) 以实现收益 \(\mathop{\max }\limits_{{x \in  X}}a\left( {x,\widehat{y}}\right)\) 的前提下选择 \(\widehat{y}\) 。通过选择 \(\widehat{y}\) ，玩家II使该收益最小化。在一般收益的博弈中，这是玩家II相当恶意的行为，从她自身收益的安全考虑来看是不合理的。
\end{tcolorbox}

However, in a zero-sum game, a min-max strategy is just a max-min strategy for the minimizer. Namely, suppose the zero-sum game is(A, B)with \(B =  - A\) . Then with \(a\left( {x,y}\right)  = {x}^{\top }{Ay}\) , a min-max strategy \(\widehat{y}\) fulfills (8.2), or equivalently

\begin{tcolorbox}
然而，在零和博弈(zero-sum game)中，极小极大策略(min-max strategy)只是针对最小值点(minimizer)的极大极小策略(max-min strategy)。即，假设零和博弈为(A, B)，其中 \(B =  - A\) 。那么，对于 \(a\left( {x,y}\right)  = {x}^{\top }{Ay}\) ，极小极大策略 \(\widehat{y}\) 满足式(8.2)，或者等价地
\end{tcolorbox}

\[
\mathop{\min }\limits_{{x \in  X}}{x}^{\top }B\widehat{y} = \mathop{\max }\limits_{{y \in  Y}}\mathop{\min }\limits_{{x \in  X}}{x}^{\top }{By}
\]

so \(\widehat{y}\) is a max-min strategy in terms of player II’s payoff matrix \(B\) . In zero-sum games, a min-max strategy is therefore not a different concept, but merely expresses player II’s preferences in terms of her cost matrix \(A\) , which is also the payoff matrix to player I and defines the matrix game.

\begin{tcolorbox}
因此， \(\widehat{y}\) 是根据玩家 II 的收益矩阵 \(B\) 得出的极大极小策略。在零和博弈中，极小极大策略并非不同的概念，而仅仅是用玩家 II 的成本矩阵 \(A\) 来表达玩家 II 的偏好，该矩阵也是玩家 I 的收益矩阵，并定义了矩阵博弈(matrix game)。
\end{tcolorbox}

The minima and maxima in Definition 8.1 are assumed to exist, which holds if \(X\) and \(Y\) are compact and \(a\left( {x,y}\right)\) is continuous in \(x\) and \(y\) . This is certainly the case for our standard assumption that \(X\) and \(Y\) are the mixed-strategy sets for a matrix game \(A\) and \(a\left( {x,y}\right)  = {x}^{\top }{Ay}\) . The following proposition observes that in \(\mathop{\max }\limits_{{x \in  X}}\mathop{\min }\limits_{{y \in  Y}}{x}^{\top }{Ay}\) the minimum is already obtained for pure strategies \(j\) instead

\begin{tcolorbox}
假设定义 8.1 中的最小值和最大值存在，若 \(X\) 和 \(Y\) 是紧集(compact set)，且 \(a\left( {x,y}\right)\) 关于 \(x\) 和 \(y\) 连续，则该假设成立。当我们通常假设 \(X\) 和 \(Y\) 是矩阵博弈 \(A\) 和 \(a\left( {x,y}\right)  = {x}^{\top }{Ay}\) 的混合策略集(mixed-strategy set)时，情况确实如此。下面的命题指出，在 \(\mathop{\max }\limits_{{x \in  X}}\mathop{\min }\limits_{{y \in  Y}}{x}^{\top }{Ay}\) 中，最小值已经可以通过纯策略 \(j\) 而非混合策略 \(X\) 得到(例如，在考虑图 8.2 中的下包络(lower envelope)时就用到了这一点)，类似地，在 \(Y\) 中，最大值已经可以通过纯策略 \(a\left( {x,y}\right)\) 而非混合策略 \(x\) 得到。
\end{tcolorbox}

of mixed strategies \(y \in  Y\) (this has been used in considering the lower envelope in Figure 8.2, for example), and similarly that in \(\mathop{\min }\limits_{{y \in  Y}}\mathop{\max }\limits_{{x \in  X}}{x}^{\top }{Ay}\) the maximum is already obtained for pure strategies \(i\) instead of mixed strategies \(x \in  X\) .

\begin{tcolorbox}
即，在 \(\mathop{\min }\limits_{{y \in  Y}}\mathop{\max }\limits_{{x \in  X}}{x}^{\top }{Ay}\) 中，最大值已经可以通过纯策略 \(i\) 而非混合策略 \(x \in  X\) 得到。
\end{tcolorbox}

Proposition 8.2. For an \(m \times  n\) matrix game \(A\) and \(x \in  X\) ,

\begin{tcolorbox}
命题 8.2. 对于一个 \(m \times  n\) 矩阵博弈 \(A\) 和 \(x \in  X\) ，
\end{tcolorbox}

\[
\mathop{\min }\limits_{{y \in  Y}}{x}^{\top }{Ay} = \min \left\{  {{\left( {x}^{\top }A\right) }_{j} \mid  1 \leq  j \leq  n}\right\}  . \tag{8.3}
\]

For \(y \in  Y\) ,

\begin{tcolorbox}
对于 \(y \in  Y\) ，
\end{tcolorbox}

\[
\mathop{\max }\limits_{{x \in  X}}{x}^{\top }{Ay} = \max \left\{  {{\left( Ay\right) }_{i} \mid  1 \leq  i \leq  m}\right\}  . \tag{8.4}
\]

Moreover, a max-min strategy \(\widehat{x} \in  X\) and a min-max strategy \(\widehat{y} \in  Y\) exist.

\begin{tcolorbox}
此外，存在一个极大极小策略 \(\widehat{x} \in  X\) 和一个极小极大策略 \(\widehat{y} \in  Y\) 。
\end{tcolorbox}

Proof. Let \(y \in  Y\) . The best response condition (6.10) says that a best response to \(y\) can be found among the pure strategies \(i\) of the row player, which implies (8.4). Condition (8.3) holds similarly.

\begin{tcolorbox}
证明。设 \(y \in  Y\) 。最优反应条件(6.10)表明，对 \(y\) 的最优反应(best response)可以在行玩家的纯策略 \(i\) 中找到，这意味着式(8.4)成立。类似地，条件(8.3)也成立。
\end{tcolorbox}

The right-hand side of (8.3) is the minimum of a finite number of terms \({\left( {x}^{\top }A\right) }_{j}\) and continuous in \(x\) . The equal left-hand side \(\mathop{\min }\limits_{{y \in  Y}}{x}^{\top }{Ay}\) is therefore also continuous in \(x\) and assumes its maximum for some \(\widehat{x}\) on the compact set \(X\) . Then \(\widehat{x}\) defines a max-min strategy. Similarly, a min-max strategy \(\widehat{y}\) exists.

\begin{tcolorbox}
式(8.3)的右边是有限项 \({\left( {x}^{\top }A\right) }_{j}\) 的最小值，且关于 \(x\) 连续。因此，相等的左边 \(\mathop{\min }\limits_{{y \in  Y}}{x}^{\top }{Ay}\) 关于 \(x\) 也连续，并且在紧集 \(X\) 上的某个 \(\widehat{x}\) 处取得最大值。那么 \(\widehat{x}\) 定义了一个极大极小策略。类似地，存在一个极小极大策略 \(\widehat{y}\) 。
\end{tcolorbox}

Because max-min payoffs and min-max costs are (in expectation) "guaranteed", the first cannot be higher than the second: With a max-min strategy \(\widehat{x}\) and min-max strategy \(\widehat{y}\) as in Definition 8.1, we have

\begin{tcolorbox}
因为极大极小收益和极小极大成本(在期望意义下)是“有保证的”，所以前者不能高于后者:根据定义 8.1 中的极大极小策略 \(\widehat{x}\) 和极小极大策略 \(\widehat{y}\) ，我们有
\end{tcolorbox}

\[
\mathop{\max }\limits_{{x \in  X}}\mathop{\min }\limits_{{y \in  Y}}a\left( {x,y}\right)  = \mathop{\min }\limits_{{y \in  Y}}a\left( {\widehat{x},y}\right)
\]

\[
\leq  a\left( {\widehat{x},\widehat{y}}\right)
\]

\[
\leq  \mathop{\max }\limits_{{x \in  X}}a\left( {x,\widehat{y}}\right)  = \mathop{\min }\limits_{{y \in  Y}}\mathop{\max }\limits_{{x \in  X}}a\left( {x,y}\right) . \tag{8.5}
\]

This is the near-obvious payoff inequality "maxmin \(\leq\) minmax", which holds for any nonempty sets \(X\) and \(Y\) , assuming that the minima and maxima exist. We have seen that for pure-strategy sets \(X\) and \(Y\) , the inequality may be strict.

\begin{tcolorbox}
这就是几乎显而易见的收益不等式“极大极小 \(\leq\) 极小极大”，对于任何非空集 \(X\) 和 \(Y\) ，假设最小值和最大值存在，则该不等式成立。我们已经看到，对于纯策略集 \(X\) 和 \(Y\) ，该不等式可能是严格不等式。
\end{tcolorbox}

If \(X\) and \(Y\) are the mixed-strategy sets and \(a\left( {x,y}\right)  = {x}^{\top }{Ay}\) , the famous minimax (short for "maxmin = minmax") theorem of von Neumann (1928) states that this inequality holds as an equality. The following proposition states that this theorem is equivalent to \(\left( {\widehat{x},\widehat{y}}\right)\) being an equilibrium of the game.

\begin{tcolorbox}
如果 \(X\) 和 \(Y\) 是混合策略集，且 \(a\left( {x,y}\right)  = {x}^{\top }{Ay}\) ，冯·诺伊曼(1928)著名的极小极大定理(“极大极小 = 极小极大”的简称)指出，该不等式取等号。下面的命题指出，该定理等价于 \(\left( {\widehat{x},\widehat{y}}\right)\) 是该博弈的一个均衡(equilibrium)。
\end{tcolorbox}

Proposition 8.3. Consider an \(m \times  n\) matrix game \(A\) with mixed-strategy sets \(X\) and \(Y\) as in (6.6), and let \(\left( {{x}^{ * },{y}^{ * }}\right)  \in  X \times  Y\) . Then the following are equivalent:

\begin{tcolorbox}
命题8.3. 考虑一个 \(m \times  n\) 矩阵博弈 \(A\) ，其混合策略集为 \(X\) 和 \(Y\) ，如(6.6)所示，并且设 \(\left( {{x}^{ * },{y}^{ * }}\right)  \in  X \times  Y\) 。那么以下表述等价:
\end{tcolorbox}

(a) \(\left( {{x}^{ * },{y}^{ * }}\right)\) is an equilibrium of(A, - A);

\begin{tcolorbox}
(a) \(\left( {{x}^{ * },{y}^{ * }}\right)\) 是(A, - A)的一个均衡；
\end{tcolorbox}

(b) \({x}^{ * }\) is a max-min strategy and \({y}^{ * }\) is a min-max strategy, and

\begin{tcolorbox}
(b) \({x}^{ * }\) 是一个极大极小策略， \({y}^{ * }\) 是一个极小极大策略，并且
\end{tcolorbox}

\[
\mathop{\max }\limits_{{x \in  X}}\mathop{\min }\limits_{{y \in  Y}}{x}^{\top }{Ay} = \mathop{\min }\limits_{{y \in  Y}}\mathop{\max }\limits_{{x \in  X}}{x}^{\top }{Ay}. \tag{8.6}
\]

Proof. Suppose (b) holds. Consider a max-min strategy \(\widehat{x}\) as in (8.1) and a min-max strategy \(\widehat{y}\) as in (8.2). Then (8.6) implies

\begin{tcolorbox}
证明. 假设(b)成立。考虑如(8.1)所示的极大极小策略 \(\widehat{x}\) 和如(8.2)所示的极小极大策略 \(\widehat{y}\) 。那么(8.6)意味着
\end{tcolorbox}

\[
{\widehat{x}}^{\top }A\widehat{y} \geq  \mathop{\min }\limits_{{y \in  Y}}{\widehat{x}}^{\top }{Ay}
\]

\[
= \mathop{\max }\limits_{{x \in  X}}\mathop{\min }\limits_{{y \in  Y}}{x}^{\top }{Ay}
\]

\[
= \mathop{\min }\limits_{{y \in  Y}}\mathop{\max }\limits_{{x \in  X}}{x}^{\top }{Ay}
\]

\[
= \mathop{\max }\limits_{{x \in  X}}{x}^{\top }A\widehat{y} \geq  {\widehat{x}}^{\top }A\widehat{y}
\]

so all these inequalities hold as equalities. Hence,

\begin{tcolorbox}
因此所有这些不等式都取等号。所以，
\end{tcolorbox}

\[
\mathop{\max }\limits_{{x \in  X}}{x}^{\top }A\widehat{y} = {\widehat{x}}^{\top }A\widehat{y} = \mathop{\min }\limits_{{y \in  Y}}{\widehat{x}}^{\top }{Ay}, \tag{8.7}
\]

which says that \(\left( {\widehat{x},\widehat{y}}\right)\) is a equilibrium because (8.7) is equivalent to

\begin{tcolorbox}
这表明 \(\left( {\widehat{x},\widehat{y}}\right)\) 是一个均衡，因为(8.7)等价于
\end{tcolorbox}

\[
{x}^{\top }A\widehat{y} \leq  {\widehat{x}}^{\top }A\widehat{y} \leq  {\widehat{x}}^{\top }{Ay}\;\text{ for all }x \in  X,y \in  Y, \tag{8.8}
\]

where the left inequality in (8.8) states that \(\widehat{x}\) is a best response to \(\widehat{y}\) and the right inequality that \(\widehat{y}\) is a best response to \(\widehat{x}\) . This shows (a).

\begin{tcolorbox}
其中(8.8)中的左不等式表明 \(\widehat{x}\) 是对 \(\widehat{y}\) 的最优反应，右不等式表明 \(\widehat{y}\) 是对 \(\widehat{x}\) 的最优反应。这证明了(a)。
\end{tcolorbox}

Conversely, suppose (a) holds, that is, as in (8.7),

\begin{tcolorbox}
反之，假设(a)成立，即如(8.7)所示，
\end{tcolorbox}

\[
\mathop{\max }\limits_{{x \in  X}}{x}^{\top }A{y}^{ * } = {x}^{*\top }A{y}^{ * } = \mathop{\min }\limits_{{y \in  Y}}{x}^{*\top }{Ay}.
\]

Using (8.5), this implies

\begin{tcolorbox}
利用(8.5)，这意味着
\end{tcolorbox}

\[
{x}^{*\top }A{y}^{ * } = \mathop{\min }\limits_{{y \in  Y}}{x}^{*\top }{Ay}
\]

\[
\leq  \mathop{\max }\limits_{{x \in  X}}\mathop{\min }\limits_{{y \in  Y}}{x}^{\top }{Ay}
\]

\[
\leq  \mathop{\min }\limits_{{y \in  Y}}\mathop{\max }\limits_{{x \in  X}}{x}^{\top }{Ay}
\]

\[
\leq  \mathop{\max }\limits_{{x \in  X}}{x}^{\top }A{y}^{ * } = {x}^{*\top }A{y}^{ * }
\]

so we also have equalities throughout, which state that \({x}^{ * }\) is a max-min strategy and \({y}^{ * }\) is a min-max strategy, and (8.6) holds, which shows (b).

\begin{tcolorbox}
因此我们同样处处取等号，这表明 \({x}^{ * }\) 是一个极大极小策略， \({y}^{ * }\) 是一个极小极大策略，并且(8.6)成立，这证明了(b)。
\end{tcolorbox}

The minimax theorem is therefore a consequence of the Theorem 6.5 of Nash (1951).

\begin{tcolorbox}
因此，极小极大定理是纳什(1951)的定理6.5的一个推论。
\end{tcolorbox}

Theorem 8.4 (The minimax theorem of von Neumann, 1928). In a matrix game with payoff matrix \(A\) to the maximizing row player \(I\) and mixed-strategy sets \(X\) and \(Y\) for players I and II,

\begin{tcolorbox}
定理8.4(冯·诺伊曼的极小极大定理，1928)。在一个矩阵博弈中，对于最大化的行玩家 \(I\) 其收益矩阵为 \(A\) ，玩家I和玩家II的混合策略集分别为 \(X\) 和 \(Y\) ，
\end{tcolorbox}

\[
\mathop{\max }\limits_{{x \in  X}}\mathop{\min }\limits_{{y \in  Y}}{x}^{\top }{Ay} = v = \mathop{\min }\limits_{{y \in  Y}}\mathop{\max }\limits_{{x \in  X}}{x}^{\top }{Ay} \tag{8.9}
\]

where \(v\) is the unique max-min payoff to player I and min-max cost to player II, called the value of the game.

\begin{tcolorbox}
其中 \(v\) 是玩家I的唯一极大极小收益和玩家II的唯一极小极大成本，称为博弈的值。
\end{tcolorbox}

Nash's theorem is proved with the help of Brouwer's fixed-point theorem, which has a complex proof. We give a short self-contained proof of the minimax Theorem 8.4 in the next Section 8.4.

\begin{tcolorbox}
纳什定理是借助布劳威尔不动点定理证明的，该定理的证明很复杂。我们将在接下来的8.4节中给出极小极大定理8.4的一个简短的自包含证明。
\end{tcolorbox}

Equilibrium strategies for a two-player zero-sum game have the following very desirable properties, which usually do not hold in games with general payoffs:

\begin{tcolorbox}
两人零和博弈的均衡策略具有以下非常理想的性质，而这些性质在一般收益的博弈中通常不成立:
\end{tcolorbox}

\begin{itemize}
\item The equilibrium payoff to the players is unique, and called the value of the game (as payoff \(v\) in (8.9) to the maximizer and cost to the minimizer).
\end{itemize}

\begin{tcolorbox}
\begin{itemize}
\item 玩家的均衡收益是唯一的，称为博弈的值(如(8.9)中对最大值点的收益 \(v\) 和对最小值点的成本)。
\end{itemize}
\end{tcolorbox}

\begin{itemize}
\item An equilibrium strategy is the same as a max-min strategy (for the maximizer, and min-max strategy for the minimizer). It is therefore independent of what the other player does. For that reason it is also called an optimal strategy of the player.
\end{itemize}

\begin{tcolorbox}
\begin{itemize}
\item 均衡策略与极大极小策略(对于最大值点而言)以及极小极大策略(对于最小值点而言)相同。因此，它与其他玩家的行为无关。出于这个原因，它也被称为玩家的最优策略。
\end{itemize}
\end{tcolorbox}

"Optimal strategy" assumes that players play as a well as they can. Optimal strategies are often rather "defensive" strategies that give the same expected payoff to many responses of the opponent. For example, a symmetric zero-sum game necessarily has value zero (see Exercise 8.3(c)). Hence, an optimal strategy of the maximizer will guarantee him at least payoff zero for all pure responses of the opponent, even if the opponent does not play optimally. Bad play of the opponent is thereby not exploited, but it can only help the player who plays optimally.

\begin{tcolorbox}
“最优策略”假定玩家会尽可能地发挥出最佳水平。最优策略通常是相当“保守”的策略，对于对手的许多应对方式，这些策略都能带来相同的期望效用。例如，一个对称的零和博弈的价值必然为零(见练习8.3(c))。因此，最大值点的最优策略将保证他在对手采用所有纯策略应对时，至少能获得零收益，即使对手没有采用最优策略。这样一来，对手的糟糕表现并不会被利用，但这只会对采用最优策略的玩家有利。
\end{tcolorbox}

Because an optimal strategy does not depend on the other player, we also have:

\begin{tcolorbox}
因为最优策略不依赖于其他玩家，所以我们还有:
\end{tcolorbox}

\begin{itemize}
\item In a zero-sum game, equilibrium strategies are exchangeable: If(x, y)and \(\left( {\bar{x},\bar{y}}\right)\) are two equilibria of a zero-sum game, then \(\left( {\bar{x},y}\right)\) and \(\left( {x,\bar{y}}\right)\) are also equilibria.
\end{itemize}

\begin{tcolorbox}
\begin{itemize}
\item 在零和博弈中，均衡策略是可交换的:如果(x, y)和 \(\left( {\bar{x},\bar{y}}\right)\) 是零和博弈的两个均衡，那么 \(\left( {\bar{x},y}\right)\) 和 \(\left( {x,\bar{y}}\right)\) 也是均衡。
\end{itemize}
\end{tcolorbox}

\begin{itemize}
\item As stated in Proposition 8.3, in a zero-sum game an equilibrium strategy is the same as a max-min (respectively, min-max) strategy. For example, if the zero-sum game has an equilibrium with a unique pure best response, that pure strategy is also the unique max-min strategy of the maximizer or min-max strategy of the minimizer. The reason is that if the min-max strategy (say) was not unique, it would be a non-unique strategy in any equilibrium. An example is the game in Figure 8.3: The max-min strategies are given by \(\left( {1 - \widehat{p},\widehat{p}}\right)\) for \(p \in  \left\lbrack  {\frac{1}{5},\frac{2}{3}}\right\rbrack\) ; for any value of \(\widehat{p}\) in the interior of that interval (for example \(\widehat{p} = \frac{1}{2}\) ), the unique best response of player II is the pure strategy \(w\) , which is therefore also the unique equilibrium strategy of player II and min-max strategy.
\end{itemize}

\begin{tcolorbox}
\begin{itemize}
\item 如命题8.3所述，在零和博弈中，均衡策略与极大极小(或极小极大)策略相同。例如，如果零和博弈存在一个具有唯一纯最优反应的均衡，那么该纯策略也是最大值点的唯一极大极小策略或最小值点的唯一极小极大策略。原因在于，如果(比如说)极小极大策略不是唯一的，那么它在任何均衡中都将是一个非唯一的策略。图8.3中的博弈就是一个例子:极大极小策略由 \(\left( {1 - \widehat{p},\widehat{p}}\right)\) 给出，其中 \(p \in  \left\lbrack  {\frac{1}{5},\frac{2}{3}}\right\rbrack\) ；对于该区间内部的任何 \(\widehat{p}\) 值(例如 \(\widehat{p} = \frac{1}{2}\) )，玩家II的唯一最优反应是纯策略 \(w\) ，因此这也是玩家II的唯一均衡策略和极小极大策略。
\end{itemize}
\end{tcolorbox}

Further properties are discussed in Section 8.5, namely the convexity of the set of optimal strategies (Proposition 8.6), and that weakly dominated strategies can be removed without affecting the value of the game (Proposition 8.7); however, this removal may affect the set of optimal strategies.

\begin{tcolorbox}
第8.5节将讨论更多性质，即最优策略集的凸性(命题8.6)，以及可以移除弱被占优策略而不影响博弈的价值(命题8.7)；然而，这种移除可能会影响最优策略集。
\end{tcolorbox}

\section*{8.4 A Short Proof of the Minimax Theorem}

\begin{tcolorbox}
\section*{8.4 极小极大定理的简短证明}
\end{tcolorbox}

The minimax Theorem 8.4 of von Neumann (1928) is, via Proposition 8.3, a special case of Nash's Theorem 6.5 on the existence of an equilibrium in a finite game. However, von Neumann's theorem predates Nash's result, and can be proved on its own. The following proof by induction is particularly simple.

\begin{tcolorbox}
冯·诺伊曼(1928)的极小极大定理8.4，通过命题8.3，是纳什关于有限博弈中均衡存在性的定理6.5的一个特例。然而，冯·诺伊曼的定理早于纳什的结果，并且可以独立证明。以下的归纳证明特别简单。
\end{tcolorbox}

Theorem 8.5. For an \(m \times  n\) matrix game \(A\) with mixed-strategy sets \(X\) and \(Y\) in (6.6) there are \(x \in  X,y \in  Y\) and \(v \in  \mathbb{R}\) so that

\begin{tcolorbox}
定理8.5。对于一个 \(m \times  n\) 矩阵博弈 \(A\) ，其混合策略集为(6.6)中的 \(X\) 和 \(Y\) ，存在 \(x \in  X,y \in  Y\) 和 \(v \in  \mathbb{R}\) 使得
\end{tcolorbox}

\[
{Ay} \leq  \mathbf{1}v,\;{x}^{\top }A \geq  v{\mathbf{1}}^{\top }.
\]

Then \(v\) is the unique value of the game and (8.9) holds, and \(x\) and \(y\) are optimal strategies of the players.

\begin{tcolorbox}
那么 \(v\) 是该博弈的唯一价值，并且(8.9)成立， \(x\) 和 \(y\) 是玩家的最优策略。
\end{tcolorbox}

Proof. Consider the two optimization problems (with \(v,u \in  \mathbb{R}\) )

\begin{tcolorbox}
证明。考虑两个优化问题(其中 \(v,u \in  \mathbb{R}\) )
\end{tcolorbox}

\[
\mathop{\operatorname{minimize}}\limits_{{y,v}}v\;\text{ subject to }{Ay} \leq  \mathbf{1}v,\;y \in  Y. \tag{8.10}
\]

and

\begin{tcolorbox}
以及
\end{tcolorbox}

\[
\mathop{\operatorname{maximize}}\limits_{{x,u}}u\text{ subject to }{x}^{\top }A \geq  u{\mathbf{1}}^{\top },\;x \in  X. \tag{8.11}
\]

The inequalities \({Ay} \leq  {1v}\) state that by playing \(y\) player II pays at most \(v\) in every row of the game. Similarly, \({x}^{\top }A \geq  u{\mathbf{1}}^{\top }\) states that by playing \(x\) player I gets at least \(u\) in every column of the game.

\begin{tcolorbox}
不等式 \({Ay} \leq  {1v}\) 表明，玩家II采用 \(y\) 策略时，在博弈的每一行中支付最多为 \(v\) 。类似地， \({x}^{\top }A \geq  u{\mathbf{1}}^{\top }\) 表明，玩家I采用 \(x\) 策略时，在博弈的每一列中至少获得 \(u\) 。
\end{tcolorbox}

Let \(y,v\) be optimal for (8.10). Then at least one of the \(m\) inequalities \({Ay} \leq  \mathbf{1}v\) holds as equality because if all inequalities were strict then \(v\) could be reduced further. Hence, \(v\) is the largest entry \({\left( Ay\right) }_{i}\) of \({Ay}\) (as in any pure best response \(i\) to \(y\) ) which by (8.4) fulfills

\begin{tcolorbox}
设 \(y,v\) 是 (8.10) 的最优解。那么 \(m\) 个不等式 \({Ay} \leq  \mathbf{1}v\) 中至少有一个取等号，因为如果所有不等式都是严格不等式，那么 \(v\) 可以进一步减小。因此， \(v\) 是 \({Ay}\) 的最大元素 \({\left( Ay\right) }_{i}\) (就像对 \(y\) 的任何纯最优反应 \(i\) 一样)，根据 (8.4) 满足
\end{tcolorbox}

\[
v = \mathop{\max }\limits_{{1 \leq  i \leq  m}}{\left( Ay\right) }_{i} = \mathop{\max }\limits_{{y \in  Y}}{x}^{\top }{Ay}
\]

and an optimal choice of \(y\) in (8.10) is a min-max strategy. Similarly, for an optimal pair \(x,u\) in (8.11) we have by (8.3)

\begin{tcolorbox}
并且 (8.10) 中 \(y\) 的最优选择是一个极小极大策略(min-max strategy)。类似地，对于 (8.11) 中的最优对 \(x,u\) ，根据 (8.3) 我们有
\end{tcolorbox}

\[
u = \mathop{\min }\limits_{{1 \leq  j \leq  n}}{\left( {x}^{\top }A\right) }_{j} = \mathop{\min }\limits_{{y \in  Y}}{x}^{\top }{Ay}
\]

where \(x\) is a max-min strategy.

\begin{tcolorbox}
其中 \(x\) 是一个极大极小策略(max-min strategy)。
\end{tcolorbox}

Consider feasible (not necessarily optimal) solutions \(y,v,x,u\) to the problems (8.10) and (8.11), that is,

\begin{tcolorbox}
考虑问题 (8.10) 和 (8.11) 的可行(不一定是最优)解 \(y,v,x,u\) ，即
\end{tcolorbox}

\[
{Ay} \leq  \mathbf{1}v,\;y \in  Y,\;{x}^{\top }A \geq  u{\mathbf{1}}^{\top },\;x \in  X. \tag{8.12}
\]

Then

\begin{tcolorbox}
那么
\end{tcolorbox}

\[
u = u{\mathbf{1}}^{\top }y \leq  {x}^{\top }{Ay} \leq  {x}^{\top }\mathbf{1}v = v. \tag{8.13}
\]

This implies that the objective functions in (8.10) and (8.11) are bounded and (without changing their optimum values) can be restricted to nonempty compact domains, namely to \(\left( {y,v}\right)  \in  Y \times  \left\lbrack  {\bar{u},\bar{v}}\right\rbrack\) and to \(\left( {x,u}\right)  \in  X \times  \left\lbrack  {\bar{u},\bar{v}}\right\rbrack\) for some \(\bar{y},\bar{v},\bar{x},\bar{u}\) that fulfill (8.12). Hence, the optimum values \(v\) and \(u\) in (8.10) and (8.11) exist (and are unique). The claim is that \(u = v\) , which proves "maxmin \(=\) minmax" and (8.9).

\begin{tcolorbox}
这意味着 (8.10) 和 (8.11) 中的目标函数是有界的，并且(在不改变其最优值的情况下)可以限制在非空紧域上，即对于满足 (8.12) 的某个 \(\bar{y},\bar{v},\bar{x},\bar{u}\) ，限制在 \(\left( {y,v}\right)  \in  Y \times  \left\lbrack  {\bar{u},\bar{v}}\right\rbrack\) 和 \(\left( {x,u}\right)  \in  X \times  \left\lbrack  {\bar{u},\bar{v}}\right\rbrack\) 上。因此，(8.10) 和 (8.11) 中的最优值 \(v\) 和 \(u\) 存在(并且是唯一的)。结论是 \(u = v\) ，这证明了“极大极小 ≤ 极小极大”以及 (8.9)。
\end{tcolorbox}

We prove this claim by induction on \(m + n\) . It is trivial for \(m + n = 2\) . Consider a matrix \(A\) of any dimension \(m \times  n\) and let \(y,v\) and \(x,u\) be optimal solutions to (8.10) and (8.11), which are feasible, that is, (8.12) holds. If all inequalities in (8.12) hold as equalities, that is, \({Ay} = \mathbf{1}v\) and \(u{\mathbf{1}}^{\top } = {x}^{\top }A\) , then we also have equality in (8.13) and \(u = v\) as claimed. Hence, we can assume that at least one inequality in \({Ay} \leq  \mathbf{1}v\) is strict (or in \({x}^{\top }A \geq  u{\mathbf{1}}^{\top }\) , where the same reasoning applies). Suppose this is the \(k\) th row, that is, with matrix entries \({a}_{ij}\) of \(A\) ,

\begin{tcolorbox}
我们通过对 \(m + n\) 进行归纳来证明这个结论。对于 \(m + n = 2\) 来说这是显然的。考虑任意维度 \(m \times  n\) 的矩阵 \(A\) ，并设 \(y,v\) 和 \(x,u\) 分别是 (8.10) 和 (8.11) 的最优解，它们是可行的，即 (8.12) 成立。如果 (8.12) 中的所有不等式都取等号，即 \({Ay} = \mathbf{1}v\) 和 \(u{\mathbf{1}}^{\top } = {x}^{\top }A\) ，那么 (8.13) 也取等号，并且如所声称的 \(u = v\) 。因此，我们可以假设 \({Ay} \leq  \mathbf{1}v\) 中至少有一个不等式是严格的(或者在 \({x}^{\top }A \geq  u{\mathbf{1}}^{\top }\) 中，同样的推理适用)。假设这是第 \(k\) 行，即，对于矩阵 \(A\) 的元素 \({a}_{ij}\) ，
\end{tcolorbox}

\[
\mathop{\sum }\limits_{{j = 1}}^{n}{a}_{kj}{y}_{j} < v \tag{8.14}
\]

Let \(\bar{A}\) be the matrix \(A\) with the \(k\) th row deleted (clearly \(m > 1\) ). Let \(\bar{y},\bar{v}\) be the optimal solution to (8.10) with \(\bar{A}\) instead of \(A\) , that is, to

\begin{tcolorbox}
设 \(\bar{A}\) 是去掉 \(k\) 行后的矩阵 \(A\) (显然 \(m > 1\) )。设 \(\bar{y},\bar{v}\) 是用 \(\bar{A}\) 代替 \(A\) 时 (8.10) 的最优解，即
\end{tcolorbox}

\[
\mathop{\operatorname{minimize}}\limits_{{\bar{y},\bar{v}}}\bar{v}\text{ subject to }\bar{A}\bar{y} \leq  \mathbf{1}\bar{v},\;\bar{y} \in  Y. \tag{8.15}
\]

The maximization problem that corresponds to (8.11) with \(\bar{A}\) instead of \(A\) can be written as

\begin{tcolorbox}
用 \(\bar{A}\) 代替 \(A\) 时，与 (8.11) 对应的最大化问题可以写成
\end{tcolorbox}

\[
\mathop{\operatorname{maximize}}\limits_{{\bar{x},\bar{u}}}\bar{u}\text{ subject to }{\bar{x}}^{\top }A \geq  \bar{u}{\mathbf{1}}^{\top },\;\bar{x} \in  X,\;{\bar{x}}_{k} = 0. \tag{8.16}
\]

Let \(\bar{x},\bar{u}\) be an optimal solution to (8.16). Then

\begin{tcolorbox}
设 \(\bar{x},\bar{u}\) 是 (8.16) 的最优解。那么
\end{tcolorbox}

\[
\bar{u} \leq  u,\;\bar{v} \leq  v, \tag{8.17}
\]

because \(\bar{x},\bar{u}\) is feasible for (8.11) (so \(\bar{u} \leq  u\) for the maximum \(u\) in (8.11)) and because \(y,v\) is feasible for (8.15). Intuitively,(8.17) holds because the game \(\bar{A}\) is at least as favorable for player II as game \(A\) . By inductive hypothesis, \(\bar{u} = \bar{v}\) .

\begin{tcolorbox}
因为 \(\bar{x},\bar{u}\) 是 (8.11) 的可行解(所以对于 (8.11) 中的最大值 \(u\) 有 \(\bar{u} \leq  u\) )，并且因为 \(y,v\) 是 (8.15) 的可行解。直观地说，(8.17) 成立是因为博弈 \(\bar{A}\) 对玩家 II 至少和博弈 \(A\) 一样有利。根据归纳假设， \(\bar{u} = \bar{v}\) 。
\end{tcolorbox}

We claim (this is the main point) that \(\bar{v} = v\) . Namely, if \(\bar{v} < v\) , let \(\varepsilon  \in  (0,1\rbrack\) , so \(v\left( {1 - \varepsilon }\right)  + \bar{v}\varepsilon  = v - \varepsilon \left( {v - \bar{v}}\right)  < v\) . Consider the strategy \({y}^{\prime } = y\left( {1 - \varepsilon }\right)  + \bar{y}\varepsilon\) , which belongs to \(Y\) because \(Y\) is convex. Because \(\bar{A}y \leq  \mathbf{1}v\) and by (8.15), \({y}^{\prime }\) fulfills

\begin{tcolorbox}
我们声称(这是关键点) \(\bar{v} = v\) 。即，如果 \(\bar{v} < v\) ，设 \(\varepsilon  \in  (0,1\rbrack\) ，那么 \(v\left( {1 - \varepsilon }\right)  + \bar{v}\varepsilon  = v - \varepsilon \left( {v - \bar{v}}\right)  < v\) 。考虑策略 \({y}^{\prime } = y\left( {1 - \varepsilon }\right)  + \bar{y}\varepsilon\) ，它属于 \(Y\) ，因为 \(Y\) 是凸集。由于 \(\bar{A}y \leq  \mathbf{1}v\) 且根据 (8.15)， \({y}^{\prime }\) 满足
\end{tcolorbox}

\[
\bar{A}{y}^{\prime } = \bar{A}\left( {y\left( {1 - \varepsilon }\right)  + \bar{y}\varepsilon }\right)  \leq  \mathbf{1}\left( {v\left( {1 - \varepsilon }\right)  + \bar{v}\varepsilon }\right)  < \mathbf{1}v.
\]

Similarly,(8.14) implies for sufficiently small \(\varepsilon\)

\begin{tcolorbox}
类似地，(8.14) 意味着对于足够小的 \(\varepsilon\)
\end{tcolorbox}

\[
\mathop{\sum }\limits_{{j = 1}}^{n}{a}_{kj}{y}_{j}^{\prime } = \mathop{\sum }\limits_{{j = 1}}^{n}{a}_{kj}{y}_{j} + \varepsilon \mathop{\sum }\limits_{{j = 1}}^{n}{a}_{kj}\left( {{\bar{y}}_{j} - {y}_{j}}\right)  < v
\]

which implies that there are \({y}^{\prime } \in  Y\) and \({v}^{\prime } < v\) so that \(A{y}^{\prime } \leq  \mathbf{1}{v}^{\prime }\) , in contradiction to the minimality of \(v\) in (8.10). This shows \(v = \bar{v}\) .

\begin{tcolorbox}
这意味着存在 \({y}^{\prime } \in  Y\) 和 \({v}^{\prime } < v\) 使得 \(A{y}^{\prime } \leq  \mathbf{1}{v}^{\prime }\) ，这与 (8.10) 中 \(v\) 的极小性相矛盾。这表明 \(v = \bar{v}\) 。
\end{tcolorbox}

With (8.17) this gives \(\bar{u} \leq  u \leq  v = \bar{v} = \bar{u}\) and therefore \(u = v\) , which was to be shown.

\begin{tcolorbox}
结合 (8.17) 可得 \(\bar{u} \leq  u \leq  v = \bar{v} = \bar{u}\) ，因此 \(u = v\) ，这正是我们要证明的。
\end{tcolorbox}

The idea behind this induction proof can be summarized as follows: If the players have optimal strategies that equalize the payoffs \(v\) and \(u\) for all rows and columns, then \(u = v\) . Otherwise (if needed by exchanging the players), there is at least one row with lower payoff than \(v\) , which will not be chosen by the row player, and by omitting this row from the game the theorem holds by the inductive hypothesis.

\begin{tcolorbox}
这个归纳证明背后的思想可以总结如下:如果玩家拥有能使所有行和列的收益 \(v\) 和 \(u\) 相等的最优策略，那么 \(u = v\) 。否则(如有必要，通过交换玩家)，至少有一行的收益低于 \(v\) ，行玩家不会选择这一行，并且通过从游戏中去掉这一行，根据归纳假设该定理成立。
\end{tcolorbox}

The optimization problem (8.10) is a special case of a linear program which asks to optimize a linear function subject to linear inequalities and equations. Here the variables are \(\left( {y,v}\right)  \in  {\mathbb{R}}^{n} \times  \mathbb{R}\) with the linear objective function \(v\) , to be minimized, and the linear constraints are \({Ay} - \mathbf{1}v \leq  \mathbf{0},y \geq  \mathbf{0},{\mathbf{1}}^{\top }y = 1\) . The maximization problem (8.11) is the so-called dual linear program to (8.10), which uses the same data of the problem but processed "vertically" rather than "horizontally", that is, using the transposed matrix and with the roles of right-hand side and objective function exchanged (we omit the precise definition here; see the references in Section 8.6).

\begin{tcolorbox}
优化问题 (8.10) 是线性规划(linear programming)的一个特殊情况，线性规划要求在满足线性不等式和方程的条件下优化一个线性函数。这里的变量是 \(\left( {y,v}\right)  \in  {\mathbb{R}}^{n} \times  \mathbb{R}\) ，线性目标函数是 \(v\) ，需要求极小值，线性约束是 \({Ay} - \mathbf{1}v \leq  \mathbf{0},y \geq  \mathbf{0},{\mathbf{1}}^{\top }y = 1\) 。最大化问题 (8.11) 是 (8.10) 的所谓对偶线性规划，它使用问题的相同数据，但“垂直”处理而非“水平”处理，即使用转置矩阵(transposed matrix)，并且交换了右侧和目标函数的角色(这里我们省略精确的定义；见 8.6 节的参考文献)。
\end{tcolorbox}

The inequality (8.13) is a special case of the weak duality of linear programming and, like (8.5), nearly immediate. The strong duality theorem of linear programming states that if both the original linear program and its dual are feasible, then their optimal values are not just mutual bounds but in fact equal. This immediately implies the minimax theorem.

\begin{tcolorbox}
不等式 (8.13) 是线性规划的弱对偶性(duality of linear programming)的一个特殊情况，并且和 (8.5) 一样，几乎是显而易见的。线性规划的强对偶定理表明，如果原始线性规划及其对偶规划都是可行的，那么它们的最优值不仅是相互的界，实际上是相等的。这立即推出了极小极大定理(minimax theorem)。
\end{tcolorbox}

Linear programming duality is a beautiful and very useful theory, for which the proof of Theorem 8.5 is a nice starting point. Moreover, algorithms for solving linear programs (which now work for hundreds of thousands, even millions, of variables and constraints) can be used to find optimal strategies of zero-sum games.

\begin{tcolorbox}
线性规划对偶(linear programming duality)是一个优美且非常有用的理论，定理 8.5 的证明是一个很好的起点。此外，求解线性规划的算法(现在可以处理数十万甚至数百万个变量和约束)可以用来寻找零和博弈(zero-sum game)的最优策略(optimal strategy)。
\end{tcolorbox}

\section*{8.5 Further Notes on Zero-Sum Games}

\begin{tcolorbox}
\section*{8.5 关于零和博弈的进一步说明}
\end{tcolorbox}

A max-min strategy, and the max-min payoff, only depends on the player's own payoffs. It can therefore also be defined for a game with general payoffs which are not necessarily zero-sum, and even for an \(N\) -player game. In that case, a max-min strategy represents a rather pessimistic view of the world: Player I chooses a strategy \(\widehat{x}\) that maximizes his payoff against an opponent (or opponents) who only want to harm player I. (In a zero-sum game, this coincides with the opponent's own objectives.) A min-max strategy in (8.2) can also be defined for a general game, with a different interpretation: The opponent chooses a strategy \(\widehat{y}\) with the intent to punish player \(\mathrm{I}\) , who can defend himself against this by subsequently choosing \(x\) . The min-max payoff is the maximum punishment that can be inflicted on player I that way. This also makes sense in a game with more than two players, where the opponents agree on a partial mixed-strategy profile \(\widehat{y}\) that maximally punishes player I. (Such "punishment strategies" are used to enforce certain equilibria in repeated games, which we do not cover in this book.)

\begin{tcolorbox}
极大极小策略(max - min strategy)以及极大极小收益仅取决于玩家自身的收益。因此，它也可以为具有不一定是零和的一般收益的博弈定义，甚至可以为一个 \(N\) 玩家博弈定义。在这种情况下，极大极小策略代表了一种相当悲观的世界观:玩家 I 选择一种策略 \(\widehat{x}\) ，该策略能在面对只想伤害玩家 I 的对手(或多个对手)时最大化其收益。(在零和博弈中，这与对手自身的目标一致。)式(8.2)中的极小极大策略(min - max strategy)也可以为一般博弈定义，但有不同的解释:对手选择一种策略 \(\widehat{y}\) 以惩罚玩家 \(\mathrm{I}\) ，玩家 \(\mathrm{I}\) 随后可以通过选择 \(x\) 来进行自我防御。极小极大收益是通过这种方式可以施加给玩家 I 的最大惩罚。这在有两个以上玩家的博弈中也有意义，在这种博弈中，对手们就一个部分混合策略组合(mixed - strategy profile) \(\widehat{y}\) 达成一致，该组合能最大程度地惩罚玩家 I。(这种“惩罚策略”用于在重复博弈中维持某些均衡，本书不涉及这部分内容。)
\end{tcolorbox}

Here, we consider max-min and min-max strategies only for zero-sum games, which define the same concept and are based on the rational assumption that the opponent plays a best response, against which the player wants to play optimally.

\begin{tcolorbox}
在这里，我们仅考虑零和博弈(zero - sum game)中的极大极小策略和极小极大策略，它们定义了相同的概念，并且基于对手会采取最优反应(best response)这一理性假设，玩家希望针对此做出最优应对。
\end{tcolorbox}

We state some further observations on the max-min strategies of the maximizer and min-max strategies of the minimizer in a zero-sum game.

\begin{tcolorbox}
我们对零和博弈中最大值点(maximizer)的极大极小策略和最小值点(minimizer)的极小极大策略做一些进一步的观察。
\end{tcolorbox}

Proposition 8.6. In a zero-sum game, the set of optimal strategies of each player is convex.

\begin{tcolorbox}
命题8.6. 在零和博弈中，每个参与者的最优策略集是凸集。
\end{tcolorbox}

Proof. This follows from Propositions 8.3 and 6.9.

\begin{tcolorbox}
证明. 这可由命题8.3和命题6.9推出。
\end{tcolorbox}

A weakly dominated strategy in a zero-sum game can be disregarded without affecting the value of the game.

\begin{tcolorbox}
在零和博弈中，可以忽略弱被占优策略而不影响博弈的值。
\end{tcolorbox}

Proposition 8.7. Consider two matrix games \(A\) and \({A}^{\prime }\) , where \({A}^{\prime }\) is obtained from \(A\) by removing a weakly dominated strategy of one of the players. Then any (mixed) equilibrium of \({A}^{\prime }\) is an equilibrium of \(A\) , and \(A\) and \({A}^{\prime }\) have the same game value.

\begin{tcolorbox}
命题8.7. 考虑两个矩阵博弈 \(A\) 和 \({A}^{\prime }\) ，其中 \({A}^{\prime }\) 是通过移除其中一个参与者的弱被占优策略从 \(A\) 得到的。那么 \({A}^{\prime }\) 的任何(混合)均衡都是 \(A\) 的均衡，并且 \(A\) 和 \({A}^{\prime }\) 具有相同的博弈值。
\end{tcolorbox}

Proof. Consider a mixed equilibrium of \({A}^{\prime }\) , which defines a mixed strategy pair in \(A\) . By the best-response condition, it is also an equilibrium in \(A\) if there is no profitable deviation to a pure strategy. The only profitable deviation not available in \({A}^{\prime }\) would be to the weakly dominated strategy, but then deviating to the strategy that dominates it would also be profitable and available in \({A}^{\prime }\) . Hence, the equilibrium of \({A}^{\prime }\) is also an equilibrium of \(A\) (Proposition 3.3(a) is similar but about pure equilibria). By Proposition 8.3, the equilibrium payoff in a zero-sum game is the value of the game, which is therefore the same in both games.

\begin{tcolorbox}
证明. 考虑 \({A}^{\prime }\) 的一个混合均衡，它在 \(A\) 中定义了一个混合策略对。根据最优反应条件，如果不存在向纯策略的有利偏离，那么它也是 \(A\) 中的一个均衡。在 \({A}^{\prime }\) 中不可行的唯一有利偏离是转向弱被占优策略，但此时转向占优该策略的策略也将是有利的，并且在 \({A}^{\prime }\) 中是可行的。因此， \({A}^{\prime }\) 的均衡也是 \(A\) 的均衡(命题3.3(a)类似，但针对的是纯均衡)。根据命题8.3，零和博弈中的均衡收益就是博弈的值，因此这两个博弈中的值是相同的。
\end{tcolorbox}

The following game, with its pure best responses shown on the right, has two pure equilibria(T, l)and(T, r). Hence, any convex combination of \(l\) and \(r\) (which is an arbitrary mixed strategy of the minimizer) is a min-max strategy. The only max-min strategy is \(T\) , as the unique best response to \(l\) . The game has value 1 . This is also the case when omitting the weakly dominated strategies \(B\) and \(r\) from the game. However, \(r\) is thereby lost as a possible min-max strategy. Consequently, if we are only interested in the value of a zero-sum game (which is often the case), then we can remove weakly dominated strategies (and iterate this), but the set of optimal strategies may become smaller that way.

\begin{tcolorbox}
下面这个博弈，其纯最优反应显示在右侧，有两个纯均衡 (T, l) 和 (T, r)。因此， \(l\) 和 \(r\) 的任何凸组合(这是最小值点的任意混合策略)都是极小极大策略。唯一的极大极小策略是 \(T\) ，因为它是对 \(l\) 的唯一最优反应。该博弈的值为1。当从博弈中剔除弱被占优策略 \(B\) 和 \(r\) 时，情况也是如此。然而， \(r\) 因此失去了作为可能的极小极大策略的资格。因此，如果我们只对零和博弈的值感兴趣(通常是这种情况)，那么我们可以移除弱被占优策略(并重复此操作)，但最优策略集可能会因此变小。
\end{tcolorbox}

\(\min {}_{1}\)  \(r\) Max

\begin{tcolorbox}
\(\min {}_{1}\)  \(r\) 最大值
\end{tcolorbox}

T1 1

B0 1

\begin{center}
\includegraphics[max width=0.2\textwidth]{images/bo_d2e5cmbef24c73b03tn0_230_872_1871_379_312_0.jpg}
\end{center}
\hspace*{3em} 

\section*{8.6 Further Reading}

\begin{tcolorbox}
\section*{8.6 进一步阅读}
\end{tcolorbox}

The history of the minimax theorem is discussed in von Neumann and Fréchet (1953). In earlier articles in the same issue of Econometrica, Fréchet explains how the mathematician Emile Borel first suggested the use of active randomizations (that is, mixed strategies) in games, and posed this theorem (for symmetric zero-sum games) as an open problem. The theorem was first proved by von Neumann (1928), without knowledge of Borel's papers and question. Its proof was later substantially simplified.

\begin{tcolorbox}
极小极大定理的历史在冯·诺伊曼和弗雷歇(1953)的文献中有所讨论。在《计量经济学》同一期的早期文章中，弗雷歇解释了数学家埃米尔·博雷尔如何首先提出在博弈中使用主动随机化(即混合策略)，并将该定理(针对对称零和博弈)作为一个开放问题提出。该定理最早由冯·诺伊曼(1928)证明，当时他并不知道博雷尔的论文和问题。后来该定理的证明得到了大幅简化。
\end{tcolorbox}

The short inductive proof of Theorem 8.5 is due to Loomis (1946) (for a more general statement, which we have specialized to zero-sum games). It is very close to the proof by Owen (1967) who removes only one strict inequality to apply the induction hypothesis, not all such inequalities like Loomis (1946).

\begin{tcolorbox}
定理8.5的简短归纳证明归功于卢米斯(1946)(针对一个更一般的陈述，我们将其专门应用于零和博弈)。它与欧文(1967)的证明非常接近，欧文只消除一个严格不等式来应用归纳假设，而不像卢米斯(1946)那样消除所有此类不等式。
\end{tcolorbox}

A geometric proof of the minimax theorem, which is also used in the book of von Neumann and Morgenstern (1947), is based on the theorem of the separating hyperplane (see Corollary 7.16 and Figure 7.13). This theorem applies to the following situation, where it is known as Farkas’s Lemma: Let \(A \in  {\mathbb{R}}^{d \times  n}\) and \(x \in  {\mathbb{R}}^{d}\) and assume that \(x \notin  D = \left\{  {{Ay} \mid  y \in  {\mathbb{R}}^{n},y \geq  \mathbf{0}}\right\}\) . The lemma states that then there is some \(v \in  {\mathbb{R}}^{d}\) so that \({v}^{\top }A \geq  {\mathbf{0}}^{\top }\) and \({v}^{\top }x < 0\) . The set \(D\) is the "convex cone" of nonnegative linear combinations of the columns of \(A\) , and if the vector \(x\) is not in that cone then there is a "separating hyperplane" with normal vector \(v\) so that all columns of \(A\) are on one side of the hyperplane \(\left( {{v}^{\top }A \geq  {\mathbf{0}}^{\top }}\right)\) and \(x\) is strictly on the other side \(\left( {{v}^{\top }x < 0}\right)\) . This is a consequence of Corollary 7.16 if one proves (which is the only missing piece) that \(D\) is closed. Because \(D\) is closed, one can choose \(c\) as the closest point to \(x\) in \(D\) and let \(v = c - x\) (which has opposite sign to \(v\) as described following Corollary 7.16). Then \(D \subseteq  \left\{  {z \in  {\mathbb{R}}^{d} \mid  {v}^{\top }z \geq  \alpha }\right\}\) where \(\alpha  = {v}^{\top }c =  - {\left( x - c\right) }^{\top }c = 0\) because, if not already \(c = \mathbf{0}\) , then similarly to Lemma 7.17 one can see that \(x - c\) is orthogonal to the ray \(\{ {c\beta } \mid  \beta  \geq  0\}\) which is a subset of \(D\) ; this also implies \({v}^{\top }\left( {x - c}\right)  =  - \parallel v{\parallel }^{2} < 0\) and thus \({v}^{\top }x < {v}^{\top }c = 0\) .

\begin{tcolorbox}
极小极大定理(minimax theorem)的一个几何证明(冯·诺伊曼和摩根斯坦在1947年的著作中也使用了该证明)基于分离超平面(separating hyperplane)定理(见推论7.16和图7.13)。该定理适用于以下情况，在此情况下它被称为法卡斯定理(Farkas’s Lemma):设 \(A \in  {\mathbb{R}}^{d \times  n}\) 且 \(x \in  {\mathbb{R}}^{d}\) ，并假设 \(x \notin  D = \left\{  {{Ay} \mid  y \in  {\mathbb{R}}^{n},y \geq  \mathbf{0}}\right\}\) 。该引理表明，存在某个 \(v \in  {\mathbb{R}}^{d}\) 使得 \({v}^{\top }A \geq  {\mathbf{0}}^{\top }\) 且 \({v}^{\top }x < 0\) 。集合 \(D\) 是矩阵 \(A\) 各列的非负线性组合构成的“凸锥”，如果向量 \(x\) 不在该锥内，那么存在一个以 \(v\) 为法向量(normal vector)的“分离超平面”，使得矩阵 \(A\) 的所有列都在超平面 \(\left( {{v}^{\top }A \geq  {\mathbf{0}}^{\top }}\right)\) 的一侧，而向量 \(x\) 严格位于另一侧 \(\left( {{v}^{\top }x < 0}\right)\) 。如果能证明(这是唯一缺失的部分) \(D\) 是闭集，那么这就是推论7.16的一个结果。因为 \(D\) 是闭集，所以可以选择 \(c\) 作为 \(D\) 中距离 \(x\) 最近的点，并令 \(v = c - x\) (根据推论7.16后的描述，它与 \(v\) 符号相反)。那么 \(D \subseteq  \left\{  {z \in  {\mathbb{R}}^{d} \mid  {v}^{\top }z \geq  \alpha }\right\}\) ，其中 \(\alpha  = {v}^{\top }c =  - {\left( x - c\right) }^{\top }c = 0\) ，因为如果还不满足 \(c = \mathbf{0}\) ，那么类似于引理7.17可以看出， \(x - c\) 与射线 \(\{ {c\beta } \mid  \beta  \geq  0\}\) 正交(orthogonal)，而该射线是 \(D\) 的一个子集；这也意味着 \({v}^{\top }\left( {x - c}\right)  =  - \parallel v{\parallel }^{2} < 0\) ，从而 \({v}^{\top }x < {v}^{\top }c = 0\) 。
\end{tcolorbox}

The strong duality theorem of linear programming, and therefore the minimax theorem, can be proved with some further algebraic manipulation from Farkas's Lemma (see Gale, 1960, p.78f). A classic text on linear programming is Dantzig (1963). Accessible introductions are Chvátal (1983) and, in particular for the geometric aspects, Matoušek and Gärtner (2007). A scholarly book with superb historical references is Schrijver (1986). The minimax theorem is easy to obtain from linear programming duality; the converse direction has been studied by Adler (2013).

\begin{tcolorbox}
线性规划(linear programming)的强对偶定理，进而极小极大定理，可以通过对法卡斯定理进行一些进一步的代数运算来证明(见Gale，1960，第78页及以后)。关于线性规划的经典著作是Dantzig(1963)。通俗易懂的介绍有Chvátal(1983)，特别是从几何角度介绍的有Matoušek和Gärtner(2007)。一本具有出色历史参考文献的学术著作是Schrijver(1986)。从线性规划对偶(linear programming duality)很容易得到极小极大定理；相反的方向已由Adler(2013)进行了研究。
\end{tcolorbox}

\section*{8.7 Exercises for Chapter 8}

\begin{tcolorbox}
\section*{8.7 第8章习题}
\end{tcolorbox}

Exercise 8.1. Consider the following Penalty game. It gives the probability of scoring a goal when the row player (the striker) adopts one of the strategies \(L\) (shoot left), \(R\) (shoot right), and the column player (the goalkeeper) uses the strategies \(l\) (jump left), \(w\) (wait then jump), \(r\) (jump right). The row player is interested in maximizing and the column player in minimizing the probability of scoring a goal.

\begin{tcolorbox}
习题8.1. 考虑以下点球博弈。当行玩家(射手)采用策略 \(L\) (向左射门)、 \(R\) (向右射门)之一，而列玩家(守门员)采用策略 \(l\) (向左扑)、 \(w\) (等待后再扑)、 \(r\) (向右扑)时，该博弈给出了进球的概率。行玩家希望最大化进球概率，而列玩家希望最小化进球概率。
\end{tcolorbox}

\begin{center}
\adjustbox{max width=\textwidth}{
\begin{tabular}{|c|c|c|c|}
\hline
\multirow{3}{*}{Max \(L\)  \(R\)} & min \(l\) & \(w\) & r \\
\cline{2-4}
 & 0.6 & 0.7 & 1.0 \\
\cline{2-4}
 & 1.0 & 0.8 & 0.7 \\
\cline{1-4}
\hline
\end{tabular}
}
\end{center}

\begin{tcolorbox}

\begin{center}
\adjustbox{max width=\textwidth}{
\begin{tabular}{|c|c|c|c|}
\hline
\multirow{3}{*}{最大值 \(L\)  \(R\)} & 最小值 \(l\) & \(w\) & r \\
\cline{2-4}
 & 0.6 & 0.7 & 1.0 \\
\cline{2-4}
 & 1.0 & 0.8 & 0.7 \\
\cline{1-4}
\hline
\end{tabular}
}
\end{center}

\end{tcolorbox}

(a) Find all equilibria of this game in mixed (including pure) strategies, and their equilibrium payoffs. Why is the payoff unique?

\begin{tcolorbox}
(a) 找出该博弈在混合策略(包括纯策略)下的所有均衡(equilibrium)及其均衡收益。为什么收益是唯一的？
\end{tcolorbox}

(b) Now suppose that player I has an additional strategy \(M\) (shoot down the middle), so that the payoff matrix is

\begin{tcolorbox}
(b) 现在假设玩家I有一个额外的策略 \(M\) (向中间射门)，使得收益矩阵为
\end{tcolorbox}

\begin{center}
\adjustbox{max width=\textwidth}{
\begin{tabular}{|c|c|c|c|}
\hline
Max & \({\text{ min }}_{1}\) & \(w\) & \(r\) \\
\cline{1-4}
\multirow{3}{*}{\(L\)  \(R\)  \(M\)} & 0.6 & 0.7 & 1.0 \\
\cline{2-4}
 & 1.0 & 0.8 & 0.7 \\
\cline{2-4}
 & 1.0 & 0.0 & 1.0 \\
\cline{1-4}
\hline
\end{tabular}
}
\end{center}

\begin{tcolorbox}

\begin{center}
\adjustbox{max width=\textwidth}{
\begin{tabular}{|c|c|c|c|}
\hline
马克斯 & \({\text{ min }}_{1}\) & \(w\) & \(r\) \\
\cline{1-4}
\multirow{3}{*}{\(L\)  \(R\)  \(M\)} & 0.6 & 0.7 & 1.0 \\
\cline{2-4}
 & 1.0 & 0.8 & 0.7 \\
\cline{2-4}
 & 1.0 & 0.0 & 1.0 \\
\cline{1-4}
\hline
\end{tabular}
}
\end{center}

\end{tcolorbox}

Find an equilibrium of this game, and the equilibrium payoff. [Hint: The result from (a) will be useful.] Exercise 8.2.

\begin{tcolorbox}
找出该博弈的一个均衡以及均衡收益。[提示:(a) 部分的结果会有用。] 练习 8.2。
\end{tcolorbox}

(a) Consider a \(2 \times  n\) zero-sum game where the best response to every pure strategy is unique. Suppose that the top row and leftmost column define a pure-strategy equilibrium of the game. Show that this equilibrium can be found by iterated elimination of strictly dominated strategies. [Hint: Consider first the case \(n = 2\) .]

\begin{tcolorbox}
(a) 考虑一个零和博弈，其中每个纯策略的最优反应是唯一的。假设最上面一行和最左边一列确定了该博弈的一个纯策略均衡。证明这个均衡可以通过迭代消除严格被占优策略来找到。[提示:首先考虑 \(n = 2\) 的情况。]
\end{tcolorbox}

(b) Does the claim in (a) still hold if some pure strategy has more than one best response (and the game is therefore degenerate)?

\begin{tcolorbox}
(b) 如果某些纯策略有不止一个最优反应(因此该博弈是退化的)，(a) 中的论断是否仍然成立？
\end{tcolorbox}

(c) Give an example of a \(3 \times  3\) zero-sum game, where the best response to every pure strategy is unique, which has a pure-strategy equilibrium that cannot be found by iterated elimination of strictly dominated strategies.

\begin{tcolorbox}
(c) 给出一个零和博弈的例子，其中每个纯策略的最优反应是唯一的，但存在一个纯策略均衡无法通过迭代消除严格被占优策略来找到。
\end{tcolorbox}

Exercise 8.3. Consider the following simultaneous game between two players: Each player chooses an integer between 1 and 5 (e.g., some number of fingers shown with your hand), and the higher number wins, except when it is just one larger than the lower number, in which case the lower number wins. So 4 beats 2 and 5 but it loses against 3. Equal numbers are a draw.

\begin{tcolorbox}
练习 8.3。考虑两个玩家之间的以下同时博弈:每个玩家选择一个介于 1 到 5 之间的整数(例如，用手展示的手指数量)，较大的数字获胜，除非它只比较小的数字大 1，在这种情况下较小的数字获胜。所以 4 胜过 2 和 5，但输给 3。数字相等则平局。
\end{tcolorbox}

(a) Write down the matrix of payoffs to the row player of this zero-sum game. A winning player gets payoff 1, a draw gives payoff 0 .

\begin{tcolorbox}
(a) 写出这个零和博弈中行玩家的收益矩阵。获胜的玩家获得收益 1，平局获得收益 0。
\end{tcolorbox}

(b) Show all pairs of strategies of the row player where one strategy weakly or strictly dominates the other, and indicate the type of domination.

\begin{tcolorbox}
(b) 找出行玩家的所有策略对，其中一个策略弱占优或严格占优另一个策略，并指出占优类型。
\end{tcolorbox}

(c) Argue carefully, without calculations, why the value of this game has to be zero.

\begin{tcolorbox}
(c) 不进行计算，仔细论证为什么这个博弈的值必须为零。
\end{tcolorbox}

(d) Find an equilibrium of this game in mixed (including pure) strategies, and explain why it is an equilibrium.

\begin{tcolorbox}
(d) 找出这个博弈在混合(包括纯)策略下的一个均衡，并解释为什么它是一个均衡。
\end{tcolorbox}

Exercise 8.4. Consider the following bimatrix game.

\begin{tcolorbox}
练习 8.4。考虑以下双矩阵博弈。
\end{tcolorbox}

\begin{center}
\includegraphics[max width=0.3\textwidth]{images/bo_d2e5cmbef24c73b03tn0_233_719_1303_486_357_0.jpg}
\end{center}
\hspace*{3em} 

(a) Find an equilibrium of this game in mixed (including pure) strategies.

\begin{tcolorbox}
(a) 找出这个博弈在混合(包括纯)策略下的一个均衡。
\end{tcolorbox}

(b) Find a max-min strategy of player I, which may be a mixed strategy, and the resulting max-min payoff to player I. [Hint: Consider the game as a zero-sum game with payoffs to player I.]

\begin{tcolorbox}
(b) 找出玩家 I 的一个极大极小策略(可能是混合策略)，以及玩家 I 由此获得的极大极小收益。[提示:将该博弈视为一个给玩家 I 带来收益的零和博弈。]
\end{tcolorbox}

(c) Find a max-min strategy of player II, which may be a mixed strategy, and the resulting max-min payoff to player II. [Hint: Consider the game as a zero-sum game with payoffs to player II; note the reversed roles of maximizer and minimizer.]

\begin{tcolorbox}
(c) 找出玩家 II 的一个极大极小策略(可能是混合策略)，以及玩家 II 由此获得的极大极小收益。[提示:将该博弈视为一个给玩家 II 带来收益的零和博弈；注意最大值点和最小值点的角色相反。]
\end{tcolorbox}

(d) Compare the max-min payoffs to player I and II from (b) and (c) with their payoffs in any equilibrium from (a).

\begin{tcolorbox}
(d) 将 (b) 和 (c) 中玩家 I 和玩家 II 的极大极小收益与 (a) 中任何均衡下他们的收益进行比较。
\end{tcolorbox}

\section*{Geometry of Equilibria in Bimatrix Games}

\begin{tcolorbox}
\section*{双矩阵博弈中均衡的几何性质}
\end{tcolorbox}

This chapter is about the geometric structure of equilibria in two-player games in strategic form. It shows how to quickly identify equilibria with qualitative "best-response diagrams". These diagrams (for example, Figure 9.4) are subdivisions of the two mixed-strategy simplices into labeled "best-response regions", and suitable labels of the outsides of each simplex, so that an equilibrium is a "completely labeled" pair(x, y)of points. The power of this method is that it easily allows studying \(3 \times  3\) games, which have a lot more variety than the rather restricted \(2 \times  2\) games and their intersecting best-response lines as in Figure 6.8.

\begin{tcolorbox}
本章讨论策略形式的两人博弈中均衡的几何结构。它展示了如何用定性的“最优反应示意图”快速识别均衡。这些示意图(例如，图 9.4)将两个混合策略单纯形细分为有标签的“最优反应区域”，并对每个单纯形的外部进行合适的标记，这样一个均衡就是一个“完全标记”的点对 (x, y)。这种方法的强大之处在于它可以轻松研究零和博弈，与图 6.8 中相当受限的矩阵博弈及其相交的最优反应线相比，零和博弈有更多的变化。
\end{tcolorbox}

Best-response diagrams have been introduced by Shapley (1974) to illustrate the algorithm of Lemke and Howson (1964) that finds one equilibrium of a bimatrix game. Moreover, the algorithm shows that every nondegenerate bimatrix game has an odd number of equilibria. In this chapter, we explain the best-response regions in Section 9.2 and the principle behind the Lemke-Howson algorithm in Section 9.3. Even a reader who is not particularly interested in the computation of equilibria should look at these two sections, which are one of the things in this book that every game theorist should know.

\begin{tcolorbox}
沙普利(1974 年)引入了最优反应示意图，以阐释莱姆基和豪森(1964 年)提出的算法，该算法可找出双矩阵博弈的一个均衡。此外，该算法表明，每个非退化双矩阵博弈都有奇数个均衡。在本章中，我们将在 9.2 节解释最优反应区域，并在 9.3 节阐述莱姆基 - 豪森算法背后的原理。即使是对均衡计算不太感兴趣的读者，也应该阅读这两节内容，这是本书中每个博弈论学者都应该了解的内容之一。
\end{tcolorbox}

The remainder of this chapter gives a more detailed introduction to the computation of equilibria in two-player games based on this geometric approach. In the introductory \(3 \times  2\) example (9.1), the mixed strategy simplex \(X\) of player \(\mathrm{I}\) is a triangle, and the mixed strategy simplex \(Y\) of player II is a line segment, where one would normally apply the upper envelope method (see Section 6.7). Section 9.4 explains the best-response diagrams for \(3 \times  3\) games, where they are subdivided triangles. Section 9.5 explains strategic equivalence of games, which leaves best responses and hence the diagrams unchanged, which is generally useful.

\begin{tcolorbox}
本章的其余部分将基于这种几何方法，更详细地介绍两人博弈中均衡的计算。在介绍性的 \(3 \times  2\) 示例(9.1)中，玩家 \(\mathrm{I}\) 的混合策略单纯形 \(X\) 是一个三角形，而玩家 II 的混合策略单纯形 \(Y\) 是一条线段，通常会应用上包络法(见 6.7 节)。9.4 节解释了 \(3 \times  3\) 博弈的最优反应示意图，它们是细分的三角形。9.5 节解释了博弈的策略等价性，它不会改变最优反应，因此也不会改变示意图，这通常很有用。
\end{tcolorbox}

While the partition of \(X\) and \(Y\) into best-response regions is convenient to draw in low dimension, it is in general dimension better to keep the payoff to the other player as an additional real variable, as described in Section 9.6. The resulting "best-response" polyhedra represent the upper envelope of expected payoffs in general dimension. The mathematically simpler, equivalent "best-response polytopes" are the basis for an algebraic implementation of the Lemke-Howson algorithm known as "complementary pivoting", described in Section 9.7. In the last Section 9.8, we describe how to extend complementary pivoting to degenerate games. This completes an alternative proof of the existence of an equilibrium in a two-player game that is independent of Brouwer's fixed-point theorem.

\begin{tcolorbox}
虽然将 \(X\) 和 \(Y\) 划分为最优反应区域在低维度下便于绘制，但在一般维度下，如 9.6 节所述，将对另一个玩家的效用作为一个额外的实变量会更好。由此产生的“最优反应”多面体代表了一般维度下期望效用的上包络。数学上更简单、等价的“最优反应有界多面体”是莱姆基 - 豪森算法的代数实现(称为“互补换基法”)的基础，这在 9.7 节中有所描述。在最后一节 9.8 中，我们描述了如何将互补换基法扩展到退化博弈。这完成了一个独立于布劳威尔不动点定理的两人博弈中均衡存在性的替代证明。
\end{tcolorbox}

\section*{9.1 Prerequisites and Learning Outcomes}

\begin{tcolorbox}
\section*{9.1 先决条件和学习成果}
\end{tcolorbox}

This chapter is a continuation of Chapter 6, and assumes the knowledge of the upper-envelope (also known as "goalpost") method and of degenerate games. All geometric definitions will be recalled here. For example, the concept of a simplex has been defined in Section 7.3, but we will only consider the mixed-strategy simplices as in Figure 6.5. Undirected graphs have been defined as background material early in Chapter 7.

\begin{tcolorbox}
本章是第 6 章的延续，假设读者已经了解上包络法(也称为“门柱”法)和退化博弈。这里将回顾所有几何定义。例如，单纯形的概念已在 7.3 节中定义，但我们仅考虑如图 6.5 所示的混合策略单纯形。无向图在第 7 章开头作为背景材料进行了定义。
\end{tcolorbox}

After studying this chapter, you should be able to:

\begin{tcolorbox}
学习本章后，你应该能够:
\end{tcolorbox}

\begin{itemize}
\item draw best-response diagrams for bimatrix games with up to three strategies per player;
\end{itemize}

\begin{tcolorbox}
\begin{itemize}
\item 为每个玩家最多有三种策略的双矩阵博弈绘制最优反应示意图；
\end{itemize}
\end{tcolorbox}

\begin{itemize}
\item identify the equilibria of the game with the help of the labels used in these diagrams, and explain the connection of labels and the best-response condition;
\end{itemize}

\begin{tcolorbox}
\begin{itemize}
\item 借助这些示意图中使用的标签识别博弈的均衡，并解释标签与最优反应条件之间的联系；
\end{itemize}
\end{tcolorbox}

\begin{itemize}
\item show the computation of one equilibrium with the Lemke-Howson algorithm using the diagram;
\end{itemize}

\begin{tcolorbox}
\begin{itemize}
\item 使用示意图展示莱姆基 - 豪森算法计算一个均衡的过程；
\end{itemize}
\end{tcolorbox}

\begin{itemize}
\item explain strategic equivalence, for two- and three-player games;
\end{itemize}

\begin{tcolorbox}
\begin{itemize}
\item 解释两人和三人博弈的策略等价性；
\end{itemize}
\end{tcolorbox}

\begin{itemize}
\item see the connection between best-response polyhedra, the upper envelope of expected payoffs, and the best response-regions;
\end{itemize}

\begin{tcolorbox}
\begin{itemize}
\item 理解最优反应多面体、期望效用的上包络和最优反应区域之间的联系；
\end{itemize}
\end{tcolorbox}

\begin{itemize}
\item be able to perform pivoting on small dictionaries with the complementary pivoting rule.
\end{itemize}

\begin{tcolorbox}
\begin{itemize}
\item 能够根据互补换基规则对小型字典进行换基操作。
\end{itemize}
\end{tcolorbox}

\section*{9.2 Labeled Best-Response Regions}

\begin{tcolorbox}
\section*{9.2 带标签的最优反应区域}
\end{tcolorbox}

Our leading example will be the \(3 \times  2\) bimatrix game(A, B)with

\begin{tcolorbox}
我们的主要示例将是 \(3 \times  2\) 双矩阵博弈 (A, B)，其中
\end{tcolorbox}

\[
A = \left( \begin{array}{ll} 2 & 0 \\  1 & 2 \\  0 & 3 \end{array}\right) ,\;B = \left( \begin{array}{ll} 2 & 0 \\  0 & 1 \\  4 & 1 \end{array}\right) . \tag{9.1}
\]

which is very similar to the right game Figure 6.7 and shown in Figure 9.1.

\begin{tcolorbox}
这与图 6.7 右侧的博弈非常相似，如图 9.1 所示。
\end{tcolorbox}

For an \(m \times  n\) bimatrix game, we number the pure strategies of the two players consecutively,

\begin{tcolorbox}
对于 \(m \times  n\) 双矩阵博弈，我们依次对两个玩家的纯策略进行编号，
\end{tcolorbox}

\begin{center}
\includegraphics[max width=0.9\textwidth]{images/bo_d2e5cmbef24c73b03tn0_236_248_211_1197_677_0.jpg}
\end{center}
\hspace*{3em} 

Figure 9.1 Left: The game in (9.1) with pure strategies numbered as in (9.2). Right: The upper envelope of expected payoffs to player I for a mixed strategy \(y \in  Y\) of player II. The mixed strategy set \(Y\) (bottom right) is subdivided into best-response regions, which are labeled with player I's pure best responses.

\begin{tcolorbox}
图9.1 左:式(9.1)中的博弈，纯策略编号如式(9.2)所示。右:对于玩家II的混合策略 \(y \in  Y\) ，玩家I的期望效用上包络。混合策略集 \(Y\) (右下角)被细分为最优反应区域，这些区域用玩家I的纯最优反应进行标记。
\end{tcolorbox}

pure strategies of player \(\mathrm{I}\) (rows) \(\; : \;1,\ldots ,m\) ,

\begin{tcolorbox}
玩家 \(\mathrm{I}\) 的纯策略(行) \(\; : \;1,\ldots ,m\) ,
\end{tcolorbox}

(9.2)pure strategies of player II (columns) : \(m + 1,\ldots ,m + n\) .

\begin{tcolorbox}
玩家II的纯策略(列): \(m + 1,\ldots ,m + n\) 。
\end{tcolorbox}

In drawings we typically show them as circled numbers, here (1, 2), 3) for player I and ④, ⑤ for player II. In the text, we usually omit the circles.

\begin{tcolorbox}
在绘图中，我们通常将它们表示为带圆圈的数字，这里玩家I是(1, 2, 3)，玩家II是④, ⑤。在文本中，我们通常省略圆圈。
\end{tcolorbox}

The mixed strategy sets \(X\) and \(Y\) of the two players are defined in (6.6). Using (9.2), we subdivide them into best-response regions of the other player according to

\begin{tcolorbox}
两个玩家的混合策略集 \(X\) 和 \(Y\) 在式(6.6)中定义。利用式(9.2)，我们根据以下规则将它们细分为对方玩家的最优反应区域
\end{tcolorbox}

\[
Y\left( i\right)  = \{ y \in  Y \mid  i\text{ is a best response to }y\} \;\left( {1 \leq  i \leq  m}\right) , \tag{9.3}
\]

\[
X\left( j\right)  = \{ x \in  X \mid  j\text{ is a best response to }x\} \;\left( {m + 1 \leq  j \leq  m + n}\right) .
\]

The best-response regions are not disjoint because a mixed strategy may have more than pure best response, but this typically happens only at the "boundary" of a best-response region. Every mixed strategy has at least one best response, so the union of the best-response regions is the respective set of mixed strategies. In Figure 9.1, the subdivision of \(Y\) can be seen from the upper envelope diagram:

\begin{tcolorbox}
最优反应区域并非不相交，因为一个混合策略可能有多个纯最优反应，但这种情况通常只发生在最优反应区域的“边界”上。每个混合策略至少有一个最优反应，因此最优反应区域的并集就是相应的混合策略集。在图9.1中， \(Y\) 的细分可以从期望效用上包络图中看出:
\end{tcolorbox}

\[
Y\left( 1\right)  = \left\{  {\left( {{y}_{1},{y}_{2}}\right)  \in  Y \mid  0 \leq  {y}_{2} \leq  \frac{1}{3}}\right\}  ,
\]

\[
Y\left( 2\right)  = \left\{  {\left( {{y}_{1},{y}_{2}}\right)  \in  Y \mid  \frac{1}{3} \leq  {y}_{2} \leq  \frac{1}{2}}\right\}  ,
\]

\[
Y\left( 3\right)  = \left\{  {\left( {{y}_{1},{y}_{2}}\right)  \in  Y \mid  \frac{1}{2} \leq  {y}_{2} \leq  1}\right\}  \text{ . }
\]

Figure 9.2 demonstrates the subdivision of \(X\) into best-response regions. The strategy simplex \(X\) is a triangle with the unit vectors \({e}_{1},{e}_{2}\) , and \({e}_{3}\) as vertices, which are the pure strategies of player I. A "goalpost" is erected on each vertex, shown in the perspective drawings in Figure 9.2. In the top-left diagram, the heights2,0, and 4 are marked on these posts, which are the payoffs to player II for her left column ④. The resulting expected payoffs define a plane above the triangle through these height markings, which defines the expected payoff to player II for this column as a function of \(x = \left( {{x}_{1},{x}_{2},{x}_{3}}\right)\) , because the expected payoff is just the convex combination of the corresponding "heights", located above the point \(x\) in the triangle \(X\) below. Similarly, the top-right diagram shows for each goalpost the expected payoffs 0, 1, and 1 for the right column (5) of player II, which span the plane for the expected payoff. The bottom-left diagram shows the upper envelope as the maximum of the two expected-payoff planes.

\begin{tcolorbox}
图9.2展示了将 \(X\) 细分为最优反应区域的情况。策略单纯形 \(X\) 是一个以单位向量 \({e}_{1},{e}_{2}\) 和 \({e}_{3}\) 为顶点的三角形，这些顶点是玩家I的纯策略。在每个顶点上竖起一个“球门柱”，如图9.2的透视图所示。在左上图中，在这些柱子上标记了高度2、0和4，这是玩家II选择左列④时的效用。由此产生的期望效用定义了一个位于三角形上方且穿过这些高度标记的平面，该平面将玩家II选择此列的期望效用定义为 \(x = \left( {{x}_{1},{x}_{2},{x}_{3}}\right)\) 的函数，因为期望效用只是相应“高度”的凸组合，这些“高度”位于下方三角形 \(X\) 中 \(x\) 点的上方。类似地，右上图展示了玩家II右列(5)的每个球门柱的期望效用0、1和1，这些期望效用构成了期望效用平面。左下图展示了期望效用上包络，即两个期望效用平面的最大值。
\end{tcolorbox}

\begin{center}
\includegraphics[max width=0.9\textwidth]{images/bo_d2e5cmbef24c73b03tn0_237_333_195_1258_1214_0.jpg}
\end{center}
\hspace*{3em} 

Figure 9.2 Expected payoffs to player II as a function of the mixed strategy of player I for the \(3 \times  2\) game in Figure 9.2. Top left: plane of expected payoffs for the left column ④, with payoffs2,0,4against the pure strategies \({e}_{1},{e}_{2},{e}_{3}\) . Top right: plane of expected payoffs for the right column (5), with payoffs0,1,1. Bottom left: upper envelope (maximum) of expected payoffs. Bottom right: Resulting subdivision of \(X\) into best-response regions \(X\left( 4\right)\) and \(X\left( 5\right)\) .

\begin{tcolorbox}
图9.2 对于图9.2中的 \(3 \times  2\) 博弈，玩家II的期望效用作为玩家I混合策略的函数。左上方:左列④的期望效用平面，针对纯策略 \({e}_{1},{e}_{2},{e}_{3}\) 的效用为2、0、4。右上方:右列(5)的期望效用平面，效用为0、1、1。左下方:期望效用的上包络(最大值)。右下方:将 \(X\) 细分为最优反应区域 \(X\left( 4\right)\) 和 \(X\left( 5\right)\) 的结果。
\end{tcolorbox}

Finally, the bottom-right diagram in Figure 9.2 shows the "projection" of the upper envelope diagram onto \(X\) , by ignoring the vertical coordinate, to represent the subdivision of \(X\) into the best-response regions \(X\left( 4\right)\) and \(X\left( 5\right)\) . Against \(\left( {{x}_{1},{x}_{2},{x}_{3}}\right)\) , the expected payoffs for the left and right column are \(2{x}_{1} + 4{x}_{3}\) and \({x}_{2} + {x}_{3}\) , respectively, so the equation

\begin{tcolorbox}
最后，图9.2的右下图展示了期望效用上包络图在 \(X\) 上的“投影”，即忽略垂直坐标，以表示将 \(X\) 细分为最优反应区域 \(X\left( 4\right)\) 和 \(X\left( 5\right)\) 的情况。针对 \(\left( {{x}_{1},{x}_{2},{x}_{3}}\right)\) ，左列和右列的期望效用分别为 \(2{x}_{1} + 4{x}_{3}\) 和 \({x}_{2} + {x}_{3}\) ，因此方程
\end{tcolorbox}

\[
2{x}_{1} + 4{x}_{3} = {x}_{2} + {x}_{3}\;\text{ or equivalently }\;2{x}_{1} + 3{x}_{3} = {x}_{2} \tag{9.4}
\]

shows where both columns have equal payoff. This defines a line segment in the triangle that belongs to both \(X\left( 4\right)\) and \(X\left( 5\right)\) . The endpoints \(a\) and \(b\) of this line segment, also shown in the diagram, are obtained by setting in equation (9.4) either \({x}_{3} = 0\) , which gives \(a = \left( {\frac{1}{3},\frac{2}{3},0}\right)\) , or \({x}_{1} = 0\) , which gives \(b = \left( {0,\frac{3}{4},\frac{1}{4}}\right)\) . The equation holds also for all convex combinations of \(a\) and \(b\) , which define the line segment itself.

\begin{tcolorbox}
显示了两列具有相等收益的位置。这在三角形中定义了一条线段，该线段同时属于 \(X\left( 4\right)\) 和 \(X\left( 5\right)\) 。此线段的端点 \(a\) 和 \(b\) (图中也有显示)是通过在方程 (9.4) 中令 \({x}_{3} = 0\) (得到 \(a = \left( {\frac{1}{3},\frac{2}{3},0}\right)\) )或 \({x}_{1} = 0\) (得到 \(b = \left( {0,\frac{3}{4},\frac{1}{4}}\right)\) )得到的。该方程对于 \(a\) 和 \(b\) 的所有凸组合也成立，这些凸组合定义了该线段本身。
\end{tcolorbox}

In Figures 9.1 and 9.2, the best-response regions \(Y\left( i\right)\) and \(X\left( j\right)\) in (9.3) are labeled with the respective pure strategies \(i\) of player I and \(j\) of player II. We now introduce a second set of labels for unplayed own strategies via

\begin{tcolorbox}
在图 9.1 和 9.2 中，(9.3) 中的最优反应区域 \(Y\left( i\right)\) 和 \(X\left( j\right)\) 分别用玩家 I 的纯策略 \(i\) 和玩家 II 的纯策略 \(j\) 进行了标注。我们现在通过以下方式为未采用的自身策略引入第二组标签
\end{tcolorbox}

\[
X\left( i\right)  = \left\{  {\left( {{x}_{1},\ldots ,{x}_{m}}\right)  \in  X \mid  \;{x}_{i} = 0}\right\}  \;\left( {1 \leq  i \leq  m}\right) , \tag{9.5}
\]

\[
Y\left( j\right)  = \left\{  {\left( {{y}_{1},\ldots ,{y}_{n}}\right)  \in  Y \mid  {y}_{j - m} = 0}\right\}  \;\left( {m + 1 \leq  j \leq  m + n}\right) .
\]

We refer to \({y}_{j - m}\) because \(j\) denotes a pure strategy in \(\{ m + 1,\ldots ,m + n\}\) of player II, but the components of \(y\) still have subscripts \(1,\ldots ,n\) . Similarly, the pure strategies of player II, which are the vertices of \(Y\) , are the unit vectors \({e}_{1},\ldots ,{e}_{n}\) in \({\mathbb{R}}^{n}\) .

\begin{tcolorbox}
我们提及 \({y}_{j - m}\) 是因为 \(j\) 表示玩家 II 在 \(\{ m + 1,\ldots ,m + n\}\) 中的一个纯策略，但 \(y\) 的分量仍有下标 \(1,\ldots ,n\) 。类似地，玩家 II 的纯策略(即 \(Y\) 的顶点)是 \({\mathbb{R}}^{n}\) 中的单位向量 \({e}_{1},\ldots ,{e}_{n}\) 。
\end{tcolorbox}

Definition 9.1. Consider an \(m \times  n\) bimatrix game(A, B)with mixed strategy sets \(X\) and \(Y\) in (6.6). Let \(k \in  \{ 1,\ldots ,m + n\}\) be any pure strategy of player I or II with \(X\left( k\right)\) and \(Y\left( k\right)\) defined as in (9.3) and (9.5). Then \(x \in  X\) is said to have label \(k\) if \(x \in  X\left( k\right)\) , and \(y \in  Y\) is said to have label \(k\) if \(y \in  Y\left( k\right)\) . If every \(k \in  \{ 1,\ldots ,m + n\}\) is a label of \(x\) or \(y\) (or both), then(x, y)is called \({completelylabeled.\ The\ sets}X\) and \(Y\) annotated with these labels define the best-response diagram for the game.

\begin{tcolorbox}
定义 9.1。考虑一个 \(m \times  n\) 双矩阵博弈 (A, B)，其混合策略集为 (6.6) 中的 \(X\) 和 \(Y\) 。设 \(k \in  \{ 1,\ldots ,m + n\}\) 是玩家 I 或玩家 II 的任意纯策略， \(X\left( k\right)\) 和 \(Y\left( k\right)\) 如 (9.3) 和 (9.5) 所定义。那么，如果 \(x \in  X\left( k\right)\) ，则称 \(x \in  X\) 具有标签 \(k\) ；如果 \(y \in  Y\left( k\right)\) ，则称 \(y \in  Y\) 具有标签 \(k\) 。如果每个 \(k \in  \{ 1,\ldots ,m + n\}\) 都是 \(x\) 或 \(y\) (或两者)的标签，那么 (x, y) 称为 \({completelylabeled.\ The\ sets}X\) ，用这些标签标注的 \(Y\) 定义了该博弈的最优反应示意图。
\end{tcolorbox}

For the game (9.1), the labels for the unplayed own strategies as defined in (9.5) are shown in Figure 9.3. For example, \(X\left( 1\right)\) is the side of the triangle given by the points \(\left( {{x}_{1},{x}_{2},{x}_{3}}\right)\) in \(X\) where \({x}_{1} = 0\) , which are all the convex combinations of the unit vectors \({e}_{2} = \left( {0,1,0}\right)\) and \({e}_{3} = \left( {0,0,1}\right)\) . In general, \(X\left( i\right)\) is the facet of the triangle that is opposite to the vertex given by the unit vector \({e}_{i}\) , for \(1 \leq  i \leq  m\) . The mixed strategy set \(Y\) of player II is a line segment with endpoints \({e}_{1} = \left( {1,0}\right)\) (which is the left column ④ of the game) and \({e}_{2} = \left( {0,1}\right)\) (which is the right column ⑤). Each"side" or facet \(Y\left( 4\right)\) or \(Y\left( 5\right)\) of this low-dimensional simplex \(Y\) is just one of these endpoints, where \({e}_{1}\) has label 5 because player II’s second strategy 5 is played with probability zero (and this is the only element of \(Y\) with that property), and the other endpoint \({e}_{2}\) has label 4 because player II’s first strategy 4 has probability zero.

\begin{tcolorbox}
对于博弈 (9.1)，如 (9.5) 中所定义的未采用的自身策略的标签如图 9.3 所示。例如， \(X\left( 1\right)\) 是由 \(X\) 中 \(\left( {{x}_{1},{x}_{2},{x}_{3}}\right)\) 这些点所构成的三角形的边，其中 \({x}_{1} = 0\) ，这些点都是单位向量 \({e}_{2} = \left( {0,1,0}\right)\) 和 \({e}_{3} = \left( {0,0,1}\right)\) 的所有凸组合。一般来说，对于 \(1 \leq  i \leq  m\) ， \(X\left( i\right)\) 是三角形中与单位向量 \({e}_{i}\) 所对应的顶点相对的极面。玩家 II 的混合策略集 \(Y\) 是一条线段，其端点为 \({e}_{1} = \left( {1,0}\right)\) (即博弈的左列 ④)和 \({e}_{2} = \left( {0,1}\right)\) (即博弈的右列 ⑤)。这个低维单纯形 \(Y\) 的每个“边”或极面 \(Y\left( 4\right)\) 或 \(Y\left( 5\right)\) 就是这些端点之一，其中 \({e}_{1}\) 的标签为 5，因为玩家 II 的第二个策略 5 被采用的概率为零(并且这是 \(Y\) 中具有该性质的唯一元素)，而另一个端点 \({e}_{2}\) 的标签为 4，因为玩家 II 的第一个策略 4 被采用的概率为零。
\end{tcolorbox}

\begin{center}
\includegraphics[max width=0.3\textwidth]{images/bo_d2e5cmbef24c73b03tn0_239_346_382_495_330_0.jpg}
\end{center}
\hspace*{3em} 

\begin{center}
\includegraphics[max width=0.4\textwidth]{images/bo_d2e5cmbef24c73b03tn0_239_1013_441_571_128_0.jpg}
\end{center}
\hspace*{3em} 

Figure 9.3 The mixed strategy simplices \(X\) and \(Y\) for the game in Figure 9.1 with best-response regions labeled with pure strategies of the other player, and sides labeled with unplayed own strategies. This defines the best-response diagram of the game.

\begin{tcolorbox}
图 9.3 展示了图 9.1 中博弈的混合策略单纯形 \(X\) 和 \(Y\) ，其中最优反应区域用另一个玩家的纯策略进行标记，边则用未采用的自身策略进行标记。这就定义了该博弈的最优反应示意图。
\end{tcolorbox}

\begin{center}
\includegraphics[max width=0.9\textwidth]{images/bo_d2e5cmbef24c73b03tn0_239_350_951_1237_306_0.jpg}
\end{center}
\hspace*{3em} 

Figure 9.4 The same picture as Figure 9.3 with all annotations omitted except for the labels, which fully define the best-response diagram. The equilibria as completely labeled vertex pairs are shown with symbols: a white circle for the pure equilibrium \(\left( {{e}_{1},{e}_{1}}\right)\) , a black square for(a, c), and a white square for(b, d).

\begin{tcolorbox}
图 9.4 与图 9.3 相同，只是省略了除标签之外的所有注释，这些标签完全定义了最优反应示意图。作为完全标记的顶点对的均衡用符号表示:白色圆圈表示纯策略均衡 \(\left( {{e}_{1},{e}_{1}}\right)\) ，黑色方块表示 (a, c)，白色方块表示 (b, d)。
\end{tcolorbox}

Any mixed strategy \(x\) or \(y\) may have several labels according to Definition 9.1. The inclusion of the unplayed own strategies among these labels is needed for the following, very useful theorem.

\begin{tcolorbox}
根据定义 9.1，任何混合策略 \(x\) 或 \(y\) 可能有多个标签。在这些标签中包含未采用的自身策略对于以下非常有用的定理是必要的。
\end{tcolorbox}

Theorem 9.2. Under the assumptions of Definition 9.1, a mixed strategy pair(x, y)is an equilibrium of(A, B)if and only if(x, y)is completely labeled.

\begin{tcolorbox}
定理 9.2。在定义 9.1 的假设下，混合策略对 (x, y) 是 (A, B) 的一个均衡，当且仅当 (x, y) 是完全标记的。
\end{tcolorbox}

Proof. The theorem is essentially another way of stating the best-response condition Proposition 6.1. Let \(\left( {x,y}\right)  \in  X \times  Y\) , let \(u\) be the best-response payoff to player I against \(y\) , and let \(v\) be the best-response payoff to player II against \(x\) , as defined in (6.9). Suppose(x, y)is an equilibrium and therefore \(x\) a best response to \(y\) as in (6.11), and \(y\) a best response to \(x\) as in (6.12). These conditions can be rewritten as

\begin{tcolorbox}
证明。该定理本质上是陈述最优反应条件(命题 6.1)的另一种方式。设 \(\left( {x,y}\right)  \in  X \times  Y\) ，设 \(u\) 是玩家 I 针对 \(y\) 的最优反应收益，设 \(v\) 是玩家 II 针对 \(x\) 的最优反应收益，如 (6.9) 中所定义。假设 (x, y) 是一个均衡，因此如 (6.11) 中所示 \(x\) 是对 \(y\) 的最优反应，并且如 (6.12) 中所示 \(y\) 是对 \(x\) 的最优反应。这些条件可以重写为
\end{tcolorbox}

\[
{x}_{i} = 0\;\text{ or }\;{\left( Ay\right) }_{i} = u\;\left( {1 \leq  i \leq  m}\right) , \tag{9.6}
\]

\[
{y}_{j} = 0\;\text{ or }\;{\left( {x}^{\top }B\right) }_{j} = v\;\left( {1 \leq  j \leq  n}\right) ,
\]

where the first condition says that \(x\) has label \(i\) (because \(x \in  X\left( i\right)\) ) or \(y\) has label \(i\) (because \(i\) is a best response to \(y\) and therefore \(y \in  Y\left( i\right)\) ), and the second condition that \(y\) has label \(m + j\) or \(x\) has label \(m + j\) (because \(m + j\) is a best response to \(x\) , with the numbering of pure strategies in (9.2)). In other words, (9.6) states that every label \(i \in  \{ 1,\ldots ,m\}\) or \(m + j \in  \{ m + 1,\ldots ,m + n\}\) appears as a label of \(x\) or \(y\) , and(x, y)is therefore completely labeled.

\begin{tcolorbox}
其中第一个条件表明 \(x\) 有标签 \(i\) (因为 \(x \in  X\left( i\right)\) )或者 \(y\) 有标签 \(i\) (因为 \(i\) 是对 \(y\) 的最优反应，因此 \(y \in  Y\left( i\right)\) )，第二个条件表明 \(y\) 有标签 \(m + j\) 或者 \(x\) 有标签 \(m + j\) (因为根据 (9.2) 中纯策略的编号， \(m + j\) 是对 \(x\) 的最优反应)。换句话说，(9.6) 表明每个标签 \(i \in  \{ 1,\ldots ,m\}\) 或 \(m + j \in  \{ m + 1,\ldots ,m + n\}\) 都作为 \(x\) 或 \(y\) 的标签出现，因此 (x, y) 是完全标记的。
\end{tcolorbox}

Conversely, if(x, y)is completely labeled, then (9.6) holds and \(x\) and \(y\) are mutual best responses, so that(x, y)is an equilibrium, as claimed.

\begin{tcolorbox}
相反，如果 (x, y) 是完全标记的，那么 (9.6) 式成立，且 \(x\) 和 \(y\) 是相互的最优反应，因此 (x, y) 是一个均衡，正如所声称的那样。
\end{tcolorbox}

Another way of quickly seeing Theorem 9.2 is that a missing label from(x, y) would be a pure strategy of a player that is played with positive probability (because it is not an unplayed pure strategy as in (9.5)) and is also not a best response, which is precisely not allowed according to the best-response condition, and violates (9.6). The players' pure strategies are numbered separately as in (9.2) so that they can appear as labels of either \(x\) or \(y\) , and in order to distinguish the two sets of conditions in (9.6).

\begin{tcolorbox}
另一种快速理解定理 9.2 的方法是，(x, y) 中的一个缺失标签意味着某个参与者以正概率采用的一个纯策略(因为它不像 (9.5) 式中那样是一个未采用的纯策略)，并且它也不是一个最优反应，而根据最优反应条件，这恰恰是不允许的，并且违反了 (9.6) 式。参与者的纯策略像 (9.2) 式那样单独编号，这样它们就可以作为 \(x\) 或 \(y\) 的标签出现，并且是为了区分 (9.6) 式中的两组条件。
\end{tcolorbox}

As Figure 9.3 shows, only the points \({e}_{1},{e}_{2},{e}_{3},a,b\) in \(X\) have three labels (which is the maximum number of labels of a point in \(X\) ), and only the points \({e}_{1},{e}_{2},c,d\) in \(Y\) have two labels (which is the maximum number of labels of a point in \(Y\) ). Because an equilibrium(x, y)needs all five labels for the players’ pure strategies, one can quickly identify the equilibria of the game. These are:

\begin{tcolorbox}
如图 9.3 所示，只有 \(X\) 中的点 \({e}_{1},{e}_{2},{e}_{3},a,b\) 有三个标签(这是 \(X\) 中一个点的最大标签数)，只有 \(Y\) 中的点 \({e}_{1},{e}_{2},c,d\) 有两个标签(这是 \(Y\) 中一个点的最大标签数)。因为一个均衡 (x, y) 需要参与者纯策略的所有五个标签，所以可以快速识别出该博弈的均衡。这些均衡是:
\end{tcolorbox}

\begin{itemize}
\item The pure equilibrium \(\left( {{e}_{1},{e}_{2}}\right)\) with labels 2,3,4 in \(X\) (which means 2 and 3 are not played and 4 is a best response) and 5,1 in \(Y\) (which means 5 is not played and 1 is a best response). This equilibrium is marked by white circles in Figure 9.4.
\end{itemize}

\begin{tcolorbox}
\begin{itemize}
\item 纯策略均衡 \(\left( {{e}_{1},{e}_{2}}\right)\) ，在 \(X\) 中有标签 2、3、4(这意味着 2 和 3 未被采用，4 是最优反应)，在 \(Y\) 中有标签 5、1(这意味着 5 未被采用，1 是最优反应)。这个均衡在图 9.4 中用白色圆圈标记。
\end{itemize}
\end{tcolorbox}

\begin{itemize}
\item The mixed equilibrium(a, c)with labels3,4,5for \(a = \left( {\frac{1}{3},\frac{2}{3},0}\right)\) in \(X\) (which means 3 is not played and 4 and 5 are best responses) and 1,2 for \(c = \left( {\frac{2}{3},\frac{1}{3}}\right)\) in \(Y\) (which means both 1 and 2 are best responses). This is marked by black squares in Figure 9.4.
\end{itemize}

\begin{tcolorbox}
\begin{itemize}
\item 混合均衡 (a, c)，对于 \(X\) 中的 \(a = \left( {\frac{1}{3},\frac{2}{3},0}\right)\) 有标签 3、4、5(这意味着 3 未被采用，4 和 5 是最优反应)，对于 \(Y\) 中的 \(c = \left( {\frac{2}{3},\frac{1}{3}}\right)\) 有标签 1、2(这意味着 1 和 2 都是最优反应)。这在图 9.4 中用黑色方块标记。
\end{itemize}
\end{tcolorbox}

\begin{itemize}
\item The mixed equilibrium(b, d)with labels1,4,5for \(b = \left( {0,\frac{3}{4},\frac{1}{4}}\right)\) in \(X\) (which means 1 is not played and 4 and 5 are best responses) and 2,3 for \(c = \left( {\frac{1}{2},\frac{1}{2}}\right)\) in \(Y\) (which means both 2 and 3 are best responses). This is marked by white squares in Figure 9.4.
\end{itemize}

\begin{tcolorbox}
\begin{itemize}
\item 混合均衡 (b, d)，对于 \(X\) 中的 \(b = \left( {0,\frac{3}{4},\frac{1}{4}}\right)\) 有标签 1、4、5(这意味着 1 未被采用，4 和 5 是最优反应)，对于 \(Y\) 中的 \(c = \left( {\frac{1}{2},\frac{1}{2}}\right)\) 有标签 2、3(这意味着 2 和 3 都是最优反应)。这在图 9.4 中用白色方块标记。
\end{itemize}
\end{tcolorbox}

The quickest way to find these equilibria as completely labeled pairs of points is to check the points in \(Y\) that have two labels, where the three other labels required to complete the full set of labels are unique, and may or may not appear as labels of points in \(X\) . The three points \({e}_{1},c,d\) in \(Y\) are part of an equilibrium, but \({e}_{2}\) is not because it has labels 3 and 4 and the required three other labels1,2,5do not appear in \(X\) , for example because the sets \(X\left( 2\right)\) and \(X\left( 5\right)\) have no point in common. This suffices to find all equilibria. A similar, slightly longer method is to check \({e}_{2}\) in \(X\) with labels1,3,5where no point in \(Y\) has the other two labels 2 and 4, and \({e}_{3}\) in \(X\) with labels1,2,4where no point in \(Y\) has the other two labels 3 and 5 .

\begin{tcolorbox}
将这些均衡作为完全标记的点对来找到的最快方法是，检查 \(Y\) 中具有两个标签的点，其中完成完整标签集所需的另外三个标签是唯一的，并且可能会也可能不会作为 \(X\) 中点的标签出现。 \(Y\) 中的三个点 \({e}_{1},c,d\) 是一个均衡的一部分，但 \({e}_{2}\) 不是，因为它有标签 3 和 4，而所需的另外三个标签 1、2、5 没有出现在 \(X\) 中，例如因为集合 \(X\left( 2\right)\) 和 \(X\left( 5\right)\) 没有公共点。这足以找到所有的均衡。一种类似但稍长的方法是检查 \(X\) 中具有标签 1、3、5 的 \({e}_{2}\) ，其中 \(Y\) 中没有点具有另外两个标签 2 和 4，以及检查 \(X\) 中具有标签 1、2、4 的 \({e}_{3}\) ，其中 \(Y\) 中没有点具有另外两个标签 3 和 5。
\end{tcolorbox}

The advantage of the best-response diagram for the game is that it is a qualitative diagram that allows identifying the equilibria, in a much more informative way than the payoffs do. In Figure 9.4 we have only kept the labels as annotations, which nevertheless encode all the necessary information. In particular, \(X\) is the strategy set of player I because its sides are labeled with the pure strategies of player I, and the inside labels are the pure best responses of player II; the reverse holds for \(Y\) . The vertices of \(X\) are unit vectors that are identified by the incident labels for the unplayed pure strategies; for example, \({e}_{1} = \left( {1,0,0}\right)\) is the vertex with labels 2 and 3 . In \(Y\) this holds for \({e}_{1} = \left( {1,0}\right)\) as the first strategy of player II where the unplayed pure strategy is \(5 = m + 2\) as the label of that endpoint of \(Y\) , and for \({e}_{2} = \left( {0,1}\right)\) as the second endpoint of \(Y\) with label \(4 = m + 1\) . The equilibria themselves are identified by their complementary sets of labels.

\begin{tcolorbox}
该博弈的最优反应示意图的优点在于，它是一种定性的示意图，能够以比收益更具信息量的方式识别均衡。在图9.4中，我们仅将标签作为注释保留，不过这些标签编码了所有必要的信息。特别地， \(X\) 是玩家I的策略集，因为其边用玩家I的纯策略标记，而内部标签是玩家II的纯最优反应；对于 \(Y\) ，情况则相反。 \(X\) 的顶点是单位向量，通过未采用的纯策略的关联标签来识别；例如， \({e}_{1} = \left( {1,0,0}\right)\) 是带有标签2和3的顶点。在 \(Y\) 中，对于玩家II的第一个策略 \({e}_{1} = \left( {1,0}\right)\) ，未采用的纯策略是 \(5 = m + 2\) ，它是 \(Y\) 该端点的标签，而对于 \(Y\) 的第二个端点 \({e}_{2} = \left( {0,1}\right)\) ，其标签为 \(4 = m + 1\) 。均衡本身通过其互补的标签集来识别。
\end{tcolorbox}

\section*{9.3 The Lemke-Howson Algorithm}

\begin{tcolorbox}
\section*{9.3 莱姆基 - 豪森算法}
\end{tcolorbox}

Identifying equilibria as completely labeled pairs(x, y)according to Theorem 9.2 does not yet prove that there is such a pair. Figure 9.5 is an extension of Figure 9.3 that illustrates such an existence proof, which is the topic of this section. In general, it provides an algorithm for finding an equilibrium of a nondegenerate \(m \times  n\) bimatrix game, due to Lemke and Howson (1964). In Section 9.7 below, we give a more specific (and more rigorous) description of this algorithm that can also be applied to degenerate games. This provides a new, elementary proof of the existence of an equilibrium in a two-player game, which does not require Brouwer's fixed point theorem. Nevertheless, this proof will have some similarity to the techniques used for proving the fixed point theorem, in particular the "parity argument" used in the proof of Sperner's lemma.

\begin{tcolorbox}
根据定理9.2将均衡识别为完全标记的对(x, y)，尚未证明存在这样的一对。图9.5是图9.3的扩展，它说明了这样一个存在性证明，这是本节的主题。一般来说，它提供了一种用于找到非退化 \(m \times  n\) 双矩阵博弈均衡的算法，该算法由莱姆基和豪森(1964年)提出。在下面的9.7节中，我们将给出该算法更具体(且更严格)的描述，该算法也可应用于退化博弈。这为两人博弈中均衡的存在性提供了一个新的、基本的证明，该证明不需要布劳威尔不动点定理。然而，这个证明将与用于证明不动点定理的技术有一些相似之处，特别是在斯佩纳引理的证明中使用的“奇偶性论证”。
\end{tcolorbox}

The mixed-strategy set \(X\) is a subset of \({\mathbb{R}}^{m}\) , and \(Y\) is a subset of \({\mathbb{R}}^{n}\) . Neither \(X\) nor \(Y\) contains the origin0of \({\mathbb{R}}^{m}\) or \({\mathbb{R}}^{n}\) , respectively. We now add0to \(X\) and connect it with a line segment to each vertex \({e}_{1},\ldots ,{e}_{m}\) of \(X\) , and denote this extension of \(X\) by \(\widetilde{X}\) . We similarly add0to \(Y\) and call the resulting extension \(\widetilde{Y}\) . Figure 9.5 shows \(\widetilde{X}\) and \(\widetilde{Y}\) for our example (9.1).

\begin{tcolorbox}
混合策略集 \(X\) 是 \({\mathbb{R}}^{m}\) 的一个子集，而 \(Y\) 是 \({\mathbb{R}}^{n}\) 的一个子集。 \(X\) 和 \(Y\) 分别不包含 \({\mathbb{R}}^{m}\) 或 \({\mathbb{R}}^{n}\) 的原点0。我们现在将0添加到 \(X\) 中，并通过线段将其连接到 \(X\) 的每个顶点 \({e}_{1},\ldots ,{e}_{m}\) ，并将 \(X\) 的这个扩展记为 \(\widetilde{X}\) 。我们类似地将0添加到 \(Y\) 中，并将得到的扩展称为 \(\widetilde{Y}\) 。图9.5展示了我们的示例(9.1)中的 \(\widetilde{X}\) 和 \(\widetilde{Y}\) 。
\end{tcolorbox}

Then0in \(\widetilde{X}\) has all labels \(1,\ldots ,m\) , and0in \(\widetilde{Y}\) has all labels \(m + 1,\ldots ,m + n\) , according to (9.5). Because the pair(0,0)is therefore completely labeled, we call it the artificial equilibrium of \(\widetilde{X} \times  \widetilde{Y}\) . It does not represent a mixed strategy pair.

\begin{tcolorbox}
然后，根据(9.5)， \(\widetilde{X}\) 中的0具有所有标签 \(1,\ldots ,m\) ， \(\widetilde{Y}\) 中的0具有所有标签 \(m + 1,\ldots ,m + n\) 。因为对(0, 0)因此是完全标记的，我们称其为 \(\widetilde{X} \times  \widetilde{Y}\) 的人工均衡。它不代表一个混合策略对。
\end{tcolorbox}

We now consider \(\widetilde{X}\) and \(\widetilde{Y}\) as graphs, rather like they are drawn in Figure 9.5. The nodes of the graph \(\widetilde{X}\) are the points that have \(m\) labels, and they are connected by an edge if they share \(m - 1\) labels, which are the labels of that edge. For example, the nodes0(with labels1,2,3) and \({e}_{2}\) (with labels1,3,5) are connected by the edge that has labels 1 and 3 . The nodes of the graph \(\widetilde{Y}\) are the points that have \(n\) labels, and are connected by an edge if they have \(n - 1\) labels in common (which in Figure 9.5 is just a single label because \(n - 1 = 1\) ), which are the labels of that edge.

\begin{tcolorbox}
我们现在将 \(\widetilde{X}\) 和 \(\widetilde{Y}\) 视为图，就像图 9.5 中所绘制的那样。图 \(\widetilde{X}\) 的节点是具有 \(m\) 个标签的点，如果它们共享 \(m - 1\) 个标签，则通过一条边相连，这些标签就是该边的标签。例如，节点 0(标签为 1、2、3)和 \({e}_{2}\) (标签为 1、3、5)由具有标签 1 和 3 的边相连。图 \(\widetilde{Y}\) 的节点是具有 \(n\) 个标签的点，如果它们有 \(n - 1\) 个共同标签(在图 9.5 中由于 \(n - 1 = 1\) 所以只有一个标签)，则通过一条边相连，这些标签就是该边的标签。
\end{tcolorbox}

\begin{center}
\includegraphics[max width=0.9\textwidth]{images/bo_d2e5cmbef24c73b03tn0_242_200_209_1214_550_0.jpg}
\end{center}
\hspace*{3em} 

Figure 9.5 The labeled mixed-strategy sets \(X\) and \(Y\) in Figure 9.3, each extended by an extra point0with all components zero, so that(0,0)is a completely labeled "artificial equilibrium". The gray arrows show the two Lemke-Howson paths for the missing label 2. The end nodes of these paths are the artificial equilibrium and the equilibria of the game, so the number of equilibria is odd.

\begin{tcolorbox}
图 9.5 图 9.3 中标记的混合策略集 \(X\) 和 \(Y\) ，每个都通过一个所有分量都为零的额外点 0 进行扩展，使得 (0, 0) 是一个完全标记的“人工均衡”。灰色箭头表示缺失标签 2 的两条莱姆基 - 豪森路径。这些路径的端点是人工均衡和博弈的均衡，因此均衡的数量是奇数。
\end{tcolorbox}

Next, we choose a label in \(\{ 1,\ldots ,m + n\}\) that is allowed to be missing, which is kept fixed throughout. We consider all pairs(x, y)in \(\widetilde{X} \times  \widetilde{Y}\) that have all labels except possibly this missing label, and call such a pair(x, y)almost completely labeled. The almost completely labeled pairs(x, y)belong to two types of sets, which are either of the form \(\{ x\}  \times  E\) where \(x\) is a node of \(\widetilde{X}\) and \(E\) is an edge of \(\widetilde{Y}\) , with \(m\) labels for \(x\) and \(n - 1\) labels for \(y \in  E\) , except for the endpoints of \(E\) which have \(n\) labels; or these sets are of the form \(E \times  \{ y\}\) where \(E\) is an edge of \(\widetilde{X}\) , with \(m - 1\) labels for \(x \in  E\) , except for the endpoints of \(E\) which have \(m\) labels, and a node \(y\) of \(\widetilde{Y}\) . (Technically, these sets are the edges of the product graph of the two graphs \(\widetilde{X}\) and \(\widetilde{Y}\) .)

\begin{tcolorbox}
接下来，我们在 \(\{ 1,\ldots ,m + n\}\) 中选择一个允许缺失的标签，该标签在整个过程中保持固定。我们考虑 \(\widetilde{X} \times  \widetilde{Y}\) 中所有除了可能缺少这个缺失标签之外具有所有标签的对 (x, y)，并将这样的对 (x, y) 称为几乎完全标记的。几乎完全标记的对 (x, y) 属于两种类型的集合，要么是 \(\{ x\}  \times  E\) 的形式，其中 \(x\) 是 \(\widetilde{X}\) 的一个节点， \(E\) 是 \(\widetilde{Y}\) 的一条边， \(x\) 有 \(m\) 个标签， \(y \in  E\) 有 \(n - 1\) 个标签，除了 \(E\) 的端点有 \(n\) 个标签；或者这些集合是 \(E \times  \{ y\}\) 的形式，其中 \(E\) 是 \(\widetilde{X}\) 的一条边， \(x \in  E\) 有 \(m - 1\) 个标签，除了 \(E\) 的端点有 \(m\) 个标签，以及 \(\widetilde{Y}\) 的一个节点 \(y\) 。(从技术上讲，这些集合是两个图 \(\widetilde{X}\) 和 \(\widetilde{Y}\) 的乘积图的边。)
\end{tcolorbox}

In Figure 9.5, the missing label is chosen to be 2. Any completely labeled pair(x, y)is by definition also almost completely labeled. It serves as the unique starting pair of a path of almost completely labeled pairs(x, y). We start with the artificial equilibrium(0,0). With label 2 allowed to be missing, we drop label 2, which means changing \(x\) along the unique edge of \(\widetilde{X}\) that connects 0 to \({e}_{2}\) (shown by an arrow in \(\widetilde{X}\) ), while keeping \(y = \mathbf{0}\) fixed. The endpoint \(x = {e}_{2}\) of that arrow in \(\widetilde{X}\) has a new label 5 which is picked up. Because \(x\) has three labels1,3,5and \(y = \mathbf{0}\) has two labels 4,5 but label 2 is missing, the label 5 that has just been picked up in \(\widetilde{X}\) is now duplicate. Because \(y\) no longer needs to have the duplicate label 5, the next step is to drop label 5 in \(\widetilde{Y}\) , that is, change \(y\) from 0 to \({e}_{2}\) along the edge which has only label 4 . The endpoint \({e}_{2}\) of that edge of \(\widetilde{Y}\) has labels 4 and 3, where label 3 has been picked up. Therefore, the current pair \(\left( {x,y}\right)  = \left( {{e}_{2},{e}_{2}}\right)\) has the duplicate label 3. Correspondingly, we can now drop label 3 in \(\widetilde{X}\) , that is, move \(x\) along the edge with labels 1 and 5 to point \(b\) , where label 4 is picked up. At the current pair \(\left( {x,y}\right)  = \left( {b,{e}_{2}}\right)\) , label 4 is duplicate. Next, we drop label 4 in \(\widetilde{Y}\) by moving \(y\) along the edge with label 3 to reach point \(d\) , where label 2 is picked up. Because 2 is the missing label, the reached pair \(\left( {x,y}\right)  = \left( {b,d}\right)\) is completely labeled. This is the equilibrium that is found at the end of the path. In summary, the path is given by the sequence of points in \(\widetilde{X} \times  \widetilde{Y}\) with corresponding label sets

\begin{tcolorbox}
在图9.5中，缺失的标签被选为2。根据定义，任何完全标记的对 (x, y) 同时也是几乎完全标记的。它作为几乎完全标记的对 (x, y) 路径的唯一起始对。我们从人工均衡 (0,0) 开始。由于允许标签2缺失，我们去掉标签2，这意味着沿着连接0到 \({e}_{2}\) 的 \(\widetilde{X}\) 的唯一边(在 \(\widetilde{X}\) 中用箭头表示)改变 \(x\) ，同时保持 \(y = \mathbf{0}\) 固定。 \(\widetilde{X}\) 中该箭头的端点 \(x = {e}_{2}\) 获得了一个新标签5。因为 \(x\) 有三个标签1、3、5，而 \(y = \mathbf{0}\) 有两个标签4、5，但标签2缺失，所以在 \(\widetilde{X}\) 中刚获得的标签5现在是重复标签。因为 \(y\) 不再需要重复标签5，下一步是在 \(\widetilde{Y}\) 中去掉标签5，即沿着只有标签4的边将 \(y\) 从0改变到 \({e}_{2}\) 。 \(\widetilde{Y}\) 中该边的端点 \({e}_{2}\) 有标签4和3，其中标签3已被获取。因此，当前对 \(\left( {x,y}\right)  = \left( {{e}_{2},{e}_{2}}\right)\) 有重复标签3。相应地，我们现在可以在 \(\widetilde{X}\) 中去掉标签3，即沿着带有标签1和5的边将 \(x\) 移动到点 \(b\) ，在那里获取标签4。在当前对 \(\left( {x,y}\right)  = \left( {b,{e}_{2}}\right)\) 处，标签4是重复标签。接下来，我们通过沿着带有标签3的边移动 \(y\) 在 \(\widetilde{Y}\) 中去掉标签4，以到达点 \(d\) ，在那里获取标签2。因为2是缺失的标签，所以到达的对 \(\left( {x,y}\right)  = \left( {b,d}\right)\) 是完全标记的。这就是在路径末尾找到的均衡。总之，路径由 \(\widetilde{X} \times  \widetilde{Y}\) 中具有相应标签集的点序列给出
\end{tcolorbox}

(9.7)

\begin{center}
\adjustbox{max width=\textwidth}{
\begin{tabular}{|c|c|c|}
\hline
(x, y) & labels & drop label \\
\cline{1-3}
(0,0) & \(\bullet\) 123,45 & 2 in \(\widetilde{X}\) \\
\cline{1-3}
\(\left( {{e}_{2},0}\right)\) & 135,45 & 5 in \(\widetilde{Y}\) \\
\cline{1-3}
\(\left( {{e}_{2},{e}_{2}}\right)\) & 135, 43 & 3 in \(\widetilde{X}\) \\
\cline{1-3}
\(\left( {b,{e}_{2}}\right)\) & 145, 43 & 4 in \(\widetilde{Y}\) \\
\cline{1-3}
(b, d) & \(\square\) 145, 23 & (found 2 in \(\widetilde{Y}\) ). \\
\cline{1-3}
\hline
\end{tabular}
}
\end{center}

\begin{tcolorbox}

\begin{center}
\adjustbox{max width=\textwidth}{
\begin{tabular}{|c|c|c|}
\hline
(x，y) & 标签 & 删除标签 \\
\cline{1-3}
(0,0) & \(\bullet\) 123,45 & 在 \(\widetilde{X}\) 中有2个 \\
\cline{1-3}
\(\left( {{e}_{2},0}\right)\) & 135,45 & 在 \(\widetilde{Y}\) 中有5个 \\
\cline{1-3}
\(\left( {{e}_{2},{e}_{2}}\right)\) & 135, 43 & 在 \(\widetilde{X}\) 中有3个 \\
\cline{1-3}
\(\left( {b,{e}_{2}}\right)\) & 145, 43 & 在 \(\widetilde{Y}\) 中有4个 \\
\cline{1-3}
(b，d) & \(\square\) 145, 23 & (在 \(\widetilde{Y}\) 中找到2个)。 \\
\cline{1-3}
\hline
\end{tabular}
}
\end{center}

\end{tcolorbox}

Hence, by starting from the artificial equilibrium(0,0), the sequence of almost completely labeled node pairs of \(\widetilde{X} \times  \widetilde{Y}\) terminates at an actual equilibrium of the game. In general, this defines an algorithm for finding an equilibrium of a bimatrix game, due to Lemke and Howson (1964). The path in (9.7) is therefore also called the Lemke-Howson path for missing label 2. Note that it alternates between \(\widetilde{X}\) and \(\widetilde{Y}\) by keeping in each step one of the components \(x\) or \(y\) in(x, y)fixed, and changing the other component by traversing an edge in the respective graph \(\widetilde{Y}\) or \(\widetilde{X}\) .

\begin{tcolorbox}
因此，从人工均衡(0,0)开始， \(\widetilde{X} \times  \widetilde{Y}\) 的几乎完全标记的节点对序列在博弈的一个实际均衡处终止。一般来说，这定义了一种由莱姆基和豪森(1964年)提出的用于寻找双矩阵博弈均衡的算法。因此，(9.7)中的路径也被称为缺失标签2的莱姆基 - 豪森路径。请注意，它通过在每一步中固定(x, y)中的一个分量 \(x\) 或 \(y\) ，并通过遍历相应图 \(\widetilde{Y}\) 或 \(\widetilde{X}\) 中的一条边来改变另一个分量，从而在 \(\widetilde{X}\) 和 \(\widetilde{Y}\) 之间交替。
\end{tcolorbox}

A Lemke-Howson path can be started at any completely labeled pair(x, y), not just at the artificial equilibrium. For example, if started from(b, d)in Figure 9.5, then the path (for missing label 2) just traverses the path (9.7) backwards, which will then go back to(0,0). However, the game has further equilibria, for example the pure-strategy equilibrium \(\left( {{e}_{1},{e}_{1}}\right)\) , indicated by the white circles. The Lemke-Howson path for missing label 2 that starts at this equilibrium is also shown in Figure 9.5 and consists of only two steps, which end at the third equilibrium(a, c):

\begin{tcolorbox}
莱姆基 - 豪森路径可以从任何完全标记的对(x, y)开始，而不仅仅是从人工均衡开始。例如，如果从图9.5中的(b, d)开始，那么(缺失标签2的)路径只是反向遍历路径(9.7)，然后会回到(0,0)。然而，该博弈还有其他均衡，例如由白色圆圈表示的纯策略均衡 \(\left( {{e}_{1},{e}_{1}}\right)\) 。从这个均衡开始的缺失标签2的莱姆基 - 豪森路径也在图9.5中显示，它只包含两步，最终到达第三个均衡(a, c):
\end{tcolorbox}

(9.8)

\begin{center}
\adjustbox{max width=\textwidth}{
\begin{tabular}{|c|c|c|c|}
\hline
(x, y) & \phantom{X} & labels & drop label \\
\cline{1-4}
\(\left( {{e}_{1},{e}_{1}}\right)\) & O & 234, 15 & 2 in \(\widetilde{X}\) \\
\cline{1-4}
\(\left( {a,{e}_{1}}\right)\) & \phantom{X} & 345,15 & 5 in \(\widetilde{Y}\) \\
\cline{1-4}
(a, c) & - & 345,12 & (found 2 in \(\widetilde{Y}\) ). \\
\cline{1-4}
\hline
\end{tabular}
}
\end{center}

\begin{tcolorbox}

\begin{center}
\adjustbox{max width=\textwidth}{
\begin{tabular}{|c|c|c|c|}
\hline
(x，y) & \phantom{X} & 标签 & 删除标签 \\
\cline{1-4}
\(\left( {{e}_{1},{e}_{1}}\right)\) & O & 234, 15 & \(\widetilde{X}\) 中的 2 \\
\cline{1-4}
\(\left( {a,{e}_{1}}\right)\) & \phantom{X} & 345,15 & \(\widetilde{Y}\) 中的 5 \\
\cline{1-4}
(a，c) & - & 345,12 & (在 \(\widetilde{Y}\) 中找到 2)。 \\
\cline{1-4}
\hline
\end{tabular}
}
\end{center}

\end{tcolorbox}

As a consequence, for a given missing label, the equilibria of the game plus the artificial equilibrium are the end nodes of the Lemke-Howson paths. Because each path has two end nodes, their total number is even and therefore the number of equilibria of the game is odd. This is a parity argument similar to that used for Sperner's lemma in Section 7.4.

\begin{tcolorbox}
因此，对于给定的缺失标签，博弈的均衡加上人工均衡是莱姆基 - 豪森路径的端点。由于每条路径有两个端点，它们的总数是偶数，因此博弈的均衡数量是奇数。这是一个奇偶性论证，类似于第7.4节中用于斯珀纳引理的论证。
\end{tcolorbox}

In order to make this parity argument for a general bimatrix game, we have to assume that the game is nondegenerate (see Definition 6.7) because a degenerate game does not necessarily have a odd number of equilibria; for example, the degenerate game in Figure 6.12 has an infinite set of equilibria (which is very well understood with the help of the best-response diagram, see Exercise 9.2).

\begin{tcolorbox}
为了对一般的双矩阵博弈进行这种奇偶性论证，我们必须假设该博弈是非退化的(见定义6.7)，因为退化博弈不一定有奇数个均衡；例如，图6.12中的退化博弈有无限个均衡集(借助最优反应示意图可以很好地理解这一点，见练习9.2)。
\end{tcolorbox}

Lemma 9.3. An \(m \times  n\) bimatrix game as in Definition 9.1 is nondegenerate if and only if no point in \(X\) has more than \(m\) labels, and no point in \(Y\) has more than \(n\) labels.

\begin{tcolorbox}
引理9.3。如定义9.1中的 \(m \times  n\) 双矩阵博弈是非退化的，当且仅当 \(X\) 中没有点的标签超过 \(m\) 个，且 \(Y\) 中没有点的标签超过 \(n\) 个。
\end{tcolorbox}

Proof. Let \(x = \left( {{x}_{1},\ldots ,{x}_{m}}\right)  \in  X\) and consider the sets \(I = \left\{  {i \mid  {x}_{i} = 0,1 \leq  i \leq  m}\right\}\) and \(J = \{ j \mid  j\) is a best response to \(x,m + 1 \leq  j \leq  m + n\}\) , so that \(I \cup  J\) is the set of labels of \(x\) . The number of pure strategies played with positive probability by \(x\) is \(m - \left| I\right|\) , and nondegeneracy requires \(\left| J\right|  \leq  m - \left| I\right|\) or equivalently \(\left| I\right|  + \left| J\right|  \leq  m\) , that is, \(x\) has no more than \(m\) labels. Similarly, in a nondegenerate game no \(y \in  Y\) has more than \(n\) labels. In a degenerate game, for some \(x\) or \(y\) this fails to hold.

\begin{tcolorbox}
证明。设 \(x = \left( {{x}_{1},\ldots ,{x}_{m}}\right)  \in  X\) ，考虑集合 \(I = \left\{  {i \mid  {x}_{i} = 0,1 \leq  i \leq  m}\right\}\) 和 \(J = \{ j \mid  j\) 是对 \(x,m + 1 \leq  j \leq  m + n\}\) 的最优反应，因此 \(I \cup  J\) 是 \(x\) 的标签集。 \(x\) 以正概率采用的纯策略数量是 \(m - \left| I\right|\) ，非退化性要求 \(\left| J\right|  \leq  m - \left| I\right|\) ，或者等价地 \(\left| I\right|  + \left| J\right|  \leq  m\) ，即 \(x\) 的标签不超过 \(m\) 个。类似地，在非退化博弈中，任何 \(y \in  Y\) 的标签都不超过 \(n\) 个。在退化博弈中，对于某些 \(x\) 或 \(y\) ，这一条件不成立。
\end{tcolorbox}

For a nondegenerate \(m \times  n\) game, the graphs \(\widetilde{X}\) and \(\widetilde{Y}\) have the following properties (which we state without proof, because we will give an alternative, more rigorous description of the Lemke-Howson method in Section 9.7 below):

\begin{tcolorbox}
对于非退化的 \(m \times  n\) 博弈，图 \(\widetilde{X}\) 和 \(\widetilde{Y}\) 具有以下性质(我们不给出证明，因为我们将在下面的第9.7节中给出莱姆基 - 豪森方法的另一种更严格的描述):
\end{tcolorbox}

\begin{itemize}
\item The nodes of \(\widetilde{X}\) are the finitely many points in \(X\) that have \(m\) labels, plus \(\mathbf{0} \in  {\mathbb{R}}^{m}\) . The nodes of \(\widetilde{Y}\) are the finitely many points in \(Y\) that have \(n\) labels, plus \(\mathbf{0} \in  {\mathbb{R}}^{n}\) .
\end{itemize}

\begin{tcolorbox}
\begin{itemize}
\item \(\widetilde{X}\) 的节点是 \(X\) 中具有 \(m\) 个标签的有限个点，再加上 \(\mathbf{0} \in  {\mathbb{R}}^{m}\) 。 \(\widetilde{Y}\) 的节点是 \(Y\) 中具有 \(n\) 个标签的有限个点，再加上 \(\mathbf{0} \in  {\mathbb{R}}^{n}\) 。
\end{itemize}
\end{tcolorbox}

\begin{itemize}
\item Any two nodes of \(\widetilde{X}\) are joined by an edge if they have \(m - 1\) labels in common, which are the labels of that edge. (As a geometric object, the edge is the set of points that have at least these labels; as an edge of the graph \(\widetilde{X}\) , the edge is just a pair \(\left\{  {x,{x}^{\prime }}\right\}\) of nodes of \(\widetilde{X}\) .) Any two nodes of \(\widetilde{Y}\) are joined by an edge if they have \(n - 1\) labels in common, which are the labels of that edge of \(\widetilde{Y}\) .
\end{itemize}

\begin{tcolorbox}
\begin{itemize}
\item 如果 \(\widetilde{X}\) 的任意两个节点有 \(m - 1\) 个共同的标签，那么它们由一条边连接，这些共同标签就是该边的标签。(作为一个几何对象，这条边是至少具有这些标签的点的集合；作为图 \(\widetilde{X}\) 的一条边，这条边只是 \(\widetilde{X}\) 的一对节点 \(\left\{  {x,{x}^{\prime }}\right\}\) 。)如果 \(\widetilde{Y}\) 的任意两个节点有 \(n - 1\) 个共同的标签，那么它们由一条边连接，这些共同标签就是 \(\widetilde{Y}\) 中该边的标签。
\end{itemize}
\end{tcolorbox}

\begin{itemize}
\item For any node \(x\) in \(\widetilde{X}\) and a label \(k\) of \(x\) , there is a unique edge of \(\widetilde{X}\) with endpoint \(x\) that has all the labels of \(x\) except \(k\) . We say this edge is traversed from \(x\) by dropping label \(k\) (the other endpoint \({x}^{\prime }\) of this edge \(\left\{  {x,{x}^{\prime }}\right\}\) has some other label than \(k\) that is picked up, as described for Figure 9.5 above). The same holds for any node in \(\widetilde{Y}\) .
\end{itemize}

\begin{tcolorbox}
\begin{itemize}
\item 对于 \(\widetilde{X}\) 中的任意节点 \(x\) 以及 \(x\) 的一个标签 \(k\) ， \(\widetilde{X}\) 中存在一条唯一的以 \(x\) 为端点的边，该边包含 \(x\) 除 \(k\) 之外的所有标签。我们称这条边是通过去掉标签 \(k\) 从 \(x\) 遍历得到的(这条边 \(\left\{  {x,{x}^{\prime }}\right\}\) 的另一个端点 \({x}^{\prime }\) 会获取除 \(k\) 之外的其他标签，如上文图 9.5 所述)。同样的情况适用于 \(\widetilde{Y}\) 中的任意节点。
\end{itemize}
\end{tcolorbox}

Let \(\ell  \in  \{ 1,\ldots ,m + n\}\) . Using the above properties, we define the \({LH}\) graph for missing label \(\ell\) as a certain subset of the "product graph" of \(\widetilde{X}\) and \(\widetilde{Y}\) . The nodes of the LH graph are the node pairs(x, y)in \(\widetilde{X} \times  \widetilde{Y}\) that are almost completely labeled, that is, every label in \(\{ 1,\ldots ,m + n\}  - \{ \ell \}\) appears as a label of \(x\) or \(y\) . Any edge of the LH graph is of the form \(\left\{  {\left( {x,y}\right) ,\left( {x,{y}^{\prime }}\right) }\right\}\) where \(\left\{  {y,{y}^{\prime }}\right\}\) is an edge of \(\widetilde{Y}\) , or of the form \(\left\{  {\left( {x,y}\right) ,\left( {{x}^{\prime },y}\right) }\right\}\) where \(\left\{  {x,{x}^{\prime }}\right\}\) is an edge of \(\widetilde{X}\) . If(x, y)is completely labeled, then it has a unique neighbor in the LH graph reached by dropping label \(\ell\) , that is, by following the edge which does not have label \(\ell\) (via the edge in \(\widetilde{X}\) if \(\ell  \in  \{ 1,\ldots ,m\}\) , and the edge in \(\widetilde{Y}\) otherwise). If(x, y)is almost completely labeled and does not have label \(\ell\) , then there is exactly one duplicate label \(k\) that appears as a label of both \(x\) and \(y\) . Such a node in the LH graph has exactly two neighbors, reached by dropping the duplicate label \(k\) in either \(\widetilde{X}\) or \(\widetilde{Y}\) .

\begin{tcolorbox}
设 \(\ell  \in  \{ 1,\ldots ,m + n\}\) 。利用上述性质，我们将缺失标签 \(\ell\) 的 \({LH}\) 图定义为 \(\widetilde{X}\) 和 \(\widetilde{Y}\) 的“乘积图”的某个子集。LH 图的节点是 \(\widetilde{X} \times  \widetilde{Y}\) 中几乎完全标记的节点对 (x, y)，即 \(\{ 1,\ldots ,m + n\}  - \{ \ell \}\) 中的每个标签都作为 \(x\) 或 \(y\) 的标签出现。LH 图的任意边的形式为 \(\left\{  {\left( {x,y}\right) ,\left( {x,{y}^{\prime }}\right) }\right\}\) ，其中 \(\left\{  {y,{y}^{\prime }}\right\}\) 是 \(\widetilde{Y}\) 的一条边，或者形式为 \(\left\{  {\left( {x,y}\right) ,\left( {{x}^{\prime },y}\right) }\right\}\) ，其中 \(\left\{  {x,{x}^{\prime }}\right\}\) 是 \(\widetilde{X}\) 的一条边。如果 (x, y) 是完全标记的，那么它在 LH 图中有一个唯一的邻居，通过去掉标签 \(\ell\) 到达，即沿着不包含标签 \(\ell\) 的边(如果 \(\ell  \in  \{ 1,\ldots ,m\}\) ，则通过 \(\widetilde{X}\) 中的边；否则通过 \(\widetilde{Y}\) 中的边)。如果 (x, y) 几乎完全标记且不包含标签 \(\ell\) ，那么恰好存在一个重复标签 \(k\) ，它同时作为 \(x\) 和 \(y\) 的标签出现。LH 图中的这样一个节点恰好有两个邻居，通过在 \(\widetilde{X}\) 或 \(\widetilde{Y}\) 中去掉重复标签 \(k\) 到达。
\end{tcolorbox}

Because every node(x, y)in the LH graph has degree one or two, the graph is a collection of paths and cycles. The end nodes of paths are the degree-one nodes, whose number is even. One of these end nodes is the artificial equilibrium(0,0), and every other end node is an equilibrium of the game. This completes the parity argument. We summarize it as a theorem (whose complete proof requires proving the above bullet points).

\begin{tcolorbox}
因为 LH 图中的每个节点 (x, y) 的度为 1 或 2，所以该图是由路径和环组成的集合。路径的端点是度为 1 的节点，其数量为偶数。这些端点之一是人工均衡 (0, 0)，其他每个端点都是该博弈的一个均衡。至此奇偶性论证完成。我们将其总结为一个定理(其完整证明需要证明上述要点)。
\end{tcolorbox}

Theorem 9.4. A nondegenerate bimatrix game has an odd number of equilibria.

\begin{tcolorbox}
定理 9.4. 一个非退化的双矩阵博弈有奇数个均衡。
\end{tcolorbox}

\begin{center}
\includegraphics[max width=0.8\textwidth]{images/bo_d2e5cmbef24c73b03tn0_245_395_630_1140_1184_0.jpg}
\end{center}
\hspace*{3em} 

A way to remember the parity argument is the story of a castle whose rooms create a maze (which is not really a maze because one cannot make a mistake when entering it), shown in Figure 9.6. Every room in the castle has only one or two doors in it. Only one door leads outside, every other door leads to another room. Every room that has only one door has a bone in it. A hungry dog arrives at the castle, and easily finds a bone, as follows. The dog enters the castle from the outside door and enters a room that either has a bone in it, or is a room with two doors, in which case the dog goes through the other door of that room, and continues in that manner. In that way the dog passes through a number of rooms in the castle, but cannot enter a room again, because that room would have needed a third door to enter it. The dog can also not get out of the castle because it has only one door to the outside, which is the one the dog entered. So the sequence of rooms (the castle has only a finite number of rooms) must end somewhere, which can only be a room with only one door that contains a bone. In this analogy, a room in a castle is a node of the LH graph, and a door that connects two rooms is an edge of the LH graph. The outside of the castle is the artificial equilibrium. Every room with only one door (and a bone in it) is an equilibrium. Moreover, by Theorem 9.4, the number of bones is odd, because starting from a bone either leads to the outside or to another bone. In addition, there may be cycles of rooms that all have two doors, like at the bottom left in Figure 9.6.

\begin{tcolorbox}
一种记住奇偶性论证的方法是通过一个城堡的故事，城堡里的房间构成了一个迷宫(实际上这并不是真正的迷宫，因为进入时不会走错路)，如图9.6所示。城堡里的每个房间只有一扇或两扇门。只有一扇门通向外面，其他每扇门都通向另一个房间。每个只有一扇门的房间里都有一根骨头。一只饥饿的狗来到城堡，很容易就找到了一根骨头，过程如下。狗从外面的门进入城堡，进入一个要么有骨头，要么有两扇门的房间。如果是有两扇门的房间，狗就穿过那扇门继续走。这样，狗穿过了城堡里的一些房间，但不会再次进入同一个房间，因为要再次进入那个房间就需要第三扇门。狗也无法走出城堡，因为只有一扇门通向外面，而那扇门就是它进来的门。所以房间序列(城堡里的房间数量是有限的)必然会在某个地方结束，而这个地方只能是一个只有一扇门且有骨头的房间。在这个类比中，城堡里的一个房间是LH图的一个节点，连接两个房间的一扇门是LH图的一条边。城堡外面是人工均衡。每个只有一扇门(且里面有骨头)的房间是一个均衡。此外，根据定理9.4，骨头的数量是奇数，因为从一根骨头出发要么通向外面，要么通向另一根骨头。此外，可能存在所有房间都有两扇门的循环，如图9.6左下角所示。
\end{tcolorbox}

\section*{9.4 Using Best-Response Diagrams}

\begin{tcolorbox}
\section*{9.4 使用最优反应示意图}
\end{tcolorbox}

In this section, we describe the use of best-response diagrams for \(3 \times  3\) games. For \(2 \times  2\) games, it is possible to draw all mixed strategy pairs simultaneously in a square that represents \(X \times  Y\) , like in Figures 6.4 and 6.8, so that an equilibrium is identified as a single point in the square. In a \(3 \times  3\) game, the best-response diagram is given by two triangles that represent \(X\) and \(Y\) , and each equilibrium has to be identified by a separate symbol that identifies the equilibrium strategies in \(X\) and \(Y\) , like for the \(3 \times  2\) game (9.1) in Figure 9.4. However, \(3 \times  3\) games have a lot more variety than \(2 \times  2\) games. It is therefore useful that they can be analyzed with the help of two-dimensional diagrams.

\begin{tcolorbox}
在本节中，我们描述最优反应示意图在 \(3 \times  3\) 博弈中的应用。对于 \(2 \times  2\) 博弈，可以在一个表示 \(X \times  Y\) 的正方形中同时绘制所有混合策略对，如图6.4和6.8所示，这样均衡就可以被识别为正方形中的一个点。在一个 \(3 \times  3\) 博弈中，最优反应示意图由两个表示 \(X\) 和 \(Y\) 的三角形给出，每个均衡必须用一个单独的符号来识别，该符号标识了 \(X\) 和 \(Y\) 中的均衡策略，如图9.4中的 \(3 \times  2\) 博弈(9.1)。然而， \(3 \times  3\) 博弈比 \(2 \times  2\) 博弈有更多的变化。因此，借助二维图来分析它们是很有用的。
\end{tcolorbox}

\begin{center}
\includegraphics[max width=0.9\textwidth]{images/bo_d2e5cmbef24c73b03tn0_246_190_1442_1265_379_0.jpg}
\end{center}
\hspace*{3em} 

Figure 9.7 The \(3 \times  3\) coordination game(A, B)and its best-response diagram. Every point in \(X\) with three labels together with the same point in \(Y\) is an equilibrium, so the game has seven equilibria.

\begin{tcolorbox}
图9.7 \(3 \times  3\) 协调博弈(A, B)及其最优反应示意图。 \(X\) 中每个带有三个标签的点与 \(Y\) 中的同一点一起构成一个均衡，因此该博弈有七个均衡。
\end{tcolorbox}

In a \(3 \times  3\) game, the rows (strategies of player I) have labels1,2,3, and the columns (strategies of player II) have labels4,5,6. Each strategy set \(X\) and \(Y\) is a triangle, with (up to) three best-response regions.

\begin{tcolorbox}
在一个 \(3 \times  3\) 博弈中，行(玩家I的策略)的标签为1、2、3，列(玩家II的策略)的标签为4、5、6。每个策略集 \(X\) 和 \(Y\) 是一个三角形，最多有三个最优反应区域。
\end{tcolorbox}

Figure 9.7 shows the game(A, B)where \(A\) and \(B\) are the identity matrix. This coordination game is a symmetric game where both players get payoff one if they choose the same row and column, and zero otherwise. In addition to the pure strategy equilibria \(\left( {{e}_{1},{e}_{1}}\right) ,\left( {{e}_{2},{e}_{2}}\right)\) , and \(\left( {{e}_{3},{e}_{3}}\right)\) , the game has many mixed equilibria, because any subset of pure strategies, chosen equally on both sides, with a uniformly at random chosen strategy from that subset, defines a mixed equilibrium(x, y). Here, these are the mixed equilibrium strategies \(x = y = \left( {\frac{1}{2},\frac{1}{2},0}\right) ,\left( {\frac{1}{2},0,\frac{1}{2}}\right) ,\left( {0,\frac{1}{2},\frac{1}{2}}\right)\) , and the completely mixed equilibrium where \(x = y = \left( {\frac{1}{3},\frac{1}{3},\frac{1}{3}}\right)\) . These are seven equilibria in total, as can be seen from the labels in Figure 9.7. The best-response regions have all the same shape; for example, \(X\left( 4\right)  = \left\{  {\left( {{x}_{1},{x}_{2},{x}_{3}}\right)  \in  X \mid  {x}_{1} \geq  {x}_{2},{x}_{1} \geq  {x}_{3}}\right\}\) .

\begin{tcolorbox}
图9.7展示了博弈(A, B)，其中 \(A\) 和 \(B\) 为单位矩阵(identity matrix)。这个协调博弈(coordination game)是一个对称博弈(symmetric game)，如果两个参与者选择相同的行和列，他们都能获得收益1，否则收益为0。除了纯策略均衡(pure equilibrium) \(\left( {{e}_{1},{e}_{1}}\right) ,\left( {{e}_{2},{e}_{2}}\right)\) 和 \(\left( {{e}_{3},{e}_{3}}\right)\) 之外，该博弈还有许多混合均衡(mixed equilibrium)，因为任何纯策略的子集，只要双方选择的概率相等，并且从该子集中随机均匀选择一个策略，就可以定义一个混合均衡(x, y)。这里，这些是混合均衡策略 \(x = y = \left( {\frac{1}{2},\frac{1}{2},0}\right) ,\left( {\frac{1}{2},0,\frac{1}{2}}\right) ,\left( {0,\frac{1}{2},\frac{1}{2}}\right)\) ，以及完全混合均衡，其中 \(x = y = \left( {\frac{1}{3},\frac{1}{3},\frac{1}{3}}\right)\) 。从图9.7中的标签可以看出，总共有七个均衡。最优反应(best response)区域的形状都相同；例如， \(X\left( 4\right)  = \left\{  {\left( {{x}_{1},{x}_{2},{x}_{3}}\right)  \in  X \mid  {x}_{1} \geq  {x}_{2},{x}_{1} \geq  {x}_{3}}\right\}\) 。
\end{tcolorbox}

\begin{center}
\includegraphics[max width=0.9\textwidth]{images/bo_d2e5cmbef24c73b03tn0_247_336_756_1258_377_0.jpg}
\end{center}
\hspace*{3em} 

Figure 9.8 A game(A, B)with a full cycle(9.9)of pure best responses. Its only equilibrium is completely mixed, marked by a black square.

\begin{tcolorbox}
图9.8 一个博弈(A, B)，具有纯最优反应的完整循环(9.9)。它唯一的均衡是完全混合均衡，用黑色方块标记。
\end{tcolorbox}

Figure 9.8 shows a game(A, B)where \(A\) is the identity matrix and \(B\) is the identity matrix with permuted columns. The best-response regions in \(X\) therefore have the same shape but different labels. As the best-response diagram and the payoff matrices show, there is a full cycle of pure best responses

\begin{tcolorbox}
图9.8展示了一个博弈(A, B)，其中 \(A\) 是单位矩阵， \(B\) 是列经过置换的单位矩阵。因此， \(X\) 中的最优反应区域形状相同但标签不同。正如最优反应示意图(best-response diagram)和收益矩阵所示，在这个博弈中存在一个纯最优反应的完整循环
\end{tcolorbox}

\[
\text{ ① } \rightarrow  \text{ ⑤ } \rightarrow  \text{ ② } \rightarrow  \text{ ⑥ } \rightarrow  \text{ ③ } \rightarrow  \text{ ④ } \rightarrow  \text{ ① } \tag{9.9}
\]

in this game. There is no completely labeled pair of mixed strategies(x, y)except for the completely mixed equilibrium where \(x = y = \left( {\frac{1}{3},\frac{1}{3},\frac{1}{3}}\right)\) . This can be seen from the best-response diagram because there is clearly no pure equilibrium by (9.9), and the point \(x = \left( {\frac{1}{2},\frac{1}{2},0}\right)\) in \(X\) with labels3,5,6requires some \(y\) with labels 1,2,4to be part of an equilibrium, but no point in \(Y\) has labels 1 and 4 ; similarly, labels1,4,6in \(X\) require labels2,3,5in \(Y\) but no point in \(Y\) has labels 2 and 5, and labels2,4,5in \(X\) require labels1,3,6in \(Y\) but no point in \(Y\) has labels 3 and 6 .

\begin{tcolorbox}
在这个博弈中，除了完全混合均衡(其中 \(x = y = \left( {\frac{1}{3},\frac{1}{3},\frac{1}{3}}\right)\) )之外，没有完全标记(completely labeled)的混合策略对(x, y)。从最优反应示意图可以看出这一点，因为根据(9.9)显然不存在纯均衡，并且 \(X\) 中标签为3、5、6的点 \(x = \left( {\frac{1}{2},\frac{1}{2},0}\right)\) 需要 \(y\) 中标签为1、2、4的点成为均衡的一部分，但 \(Y\) 中没有点同时具有标签1和4；类似地， \(X\) 中标签为1、4、6的点需要 \(Y\) 中标签为2、3、5的点，但 \(Y\) 中没有点同时具有标签2和5， \(X\) 中标签为2、4、5的点需要 \(Y\) 中标签为1、3、6的点，但 \(Y\) 中没有点同时具有标签3和6。
\end{tcolorbox}

This raises the general question if the cycle (9.9) of pure best responses (for a nondegenerate game) implies that the game has a unique mixed equilibrium (which is true for \(2 \times  2\) games). That question is easily answered with the help of the best-response diagrams. The pure best responses in the corners of the triangle are determined by (9.9). We start with Figure 9.9, where we first keep \(A\) and therefore the best-response regions \(Y\left( 1\right) ,Y\left( 2\right)\) , and \(Y\left( 3\right)\) fixed. Can \(y = \left( {\frac{1}{2},\frac{1}{2},0}\right)\) in \(Y\) with labels1,2,6be part of an equilibrium(x, y)? This requires \(x\) to have labels3,4,5, which is indeed possible if, along the bottom edge of \(X\) with label 3, label 4 appears between labels 5 and 6 . In the resulting best-response diagram, \(x = \left( {\frac{2}{3},\frac{1}{3},0}\right)\) has these labels3,4,5, and the equilibrium strategies \(x\) and \(y\) are marked with a black square. The best-response diagram shows that this is the only equilibrium. Hence, the game has indeed only one equilibrium, although it is not completely mixed.

\begin{tcolorbox}
这就引出了一个普遍问题:对于一个非退化博弈(nondegenerate game)，纯最优反应(pure best responses)的循环(9.9)是否意味着该博弈有唯一的混合均衡(mixed equilibrium)(这对于 \(2 \times  2\) 博弈是成立的)。借助最优反应示意图(best - response diagram)可以轻松回答这个问题。三角形角上的纯最优反应由(9.9)确定。我们从图9.9开始，首先固定 \(A\) ，因此最优反应区域 \(Y\left( 1\right) ,Y\left( 2\right)\) 和 \(Y\left( 3\right)\) 也固定。 \(Y\) 中带有标签1、2、6的 \(y = \left( {\frac{1}{2},\frac{1}{2},0}\right)\) 能否成为均衡(x，y)的一部分呢？这要求 \(x\) 带有标签3、4、5，如果在带有标签3的 \(X\) 的底边，标签4出现在标签5和6之间，这确实是可能的。在得到的最优反应示意图中， \(x = \left( {\frac{2}{3},\frac{1}{3},0}\right)\) 带有这些标签3、4、5，均衡策略 \(x\) 和 \(y\) 用黑色方块标记。最优反应示意图表明这是唯一的均衡。因此，该博弈确实只有一个均衡，尽管它不是完全混合的。
\end{tcolorbox}

\begin{center}
\includegraphics[max width=0.9\textwidth]{images/bo_d2e5cmbef24c73b03tn0_248_193_564_1263_378_0.jpg}
\end{center}
\hspace*{3em} 

Figure 9.9 Best-response diagram for a game with a full cycle (9.9) of pure best responses and a unique equilibrium(x, y)for \(x = \left( {\frac{2}{3},\frac{1}{3},0}\right)\) and \(y = \left( {\frac{1}{2},\frac{1}{2},0}\right)\) . The payoff matrix \(B\) is constructed in Figure 9.10.

\begin{tcolorbox}
图9.9 一个具有纯最优反应的完整循环(9.9)且对于 \(x = \left( {\frac{2}{3},\frac{1}{3},0}\right)\) 和 \(y = \left( {\frac{1}{2},\frac{1}{2},0}\right)\) 有唯一均衡(x，y)的博弈的最优反应示意图。收益矩阵 \(B\) 在图9.10中构建。
\end{tcolorbox}

We have drawn the best-response diagram first. The construction of a payoff matrix \(B\) to get this subdivision of \(X\) into best-response regions is shown in Figure 9.10.

\begin{tcolorbox}
我们先绘制了最优反应示意图。图9.10展示了构建一个收益矩阵 \(B\) 以使 \(X\) 划分为最优反应区域的过程。
\end{tcolorbox}

\begin{center}
\includegraphics[max width=0.7\textwidth]{images/bo_d2e5cmbef24c73b03tn0_248_192_1287_926_422_0.jpg}
\end{center}
\hspace*{3em} 

\[
B = \left( \begin{array}{lll} 2 & 3 & 0 \\  2 & 0 & 3 \\  2 & 0 & 0 \end{array}\right)
\]

Figure 9.10 Construction of the game matrix \(B\) from the best-response diagram on the left via the upper envelope of suitable payoffs. The first column 4 defines a horizontal plane of constant payoffs 2, and the second and third columns 5 and 6 have higher payoff 3 in the corners \({e}_{1}\) and \({e}_{2}\) , respectively, and 0 otherwise.

\begin{tcolorbox}
图9.10 通过合适收益的上包络(upper envelope)从左侧的最优反应示意图构建博弈矩阵 \(B\) 。第一列4定义了一个常数收益为2的水平面，第二列和第三列5和6分别在角 \({e}_{1}\) 和 \({e}_{2}\) 处有更高的收益3，其他地方为0。
\end{tcolorbox}

Figure 9.11 shows a game where the best-response regions \(Y\left( 1\right) ,Y\left( 2\right) ,Y\left( 3\right)\) , and the payoff matrix \(A\) , are also changed compared to Figure 9.10. The game still has the full cycle (9.9) of pure best responses, but it has three equilibria. Hence, condition (9.9) does not imply a unique equilibrium. The counterexample was easy to construct with the help of the best-response diagrams.

\begin{tcolorbox}
图9.11展示了一个博弈，与图9.10相比，其最优反应区域 \(Y\left( 1\right) ,Y\left( 2\right) ,Y\left( 3\right)\) 和收益矩阵 \(A\) 也发生了变化。该博弈仍然有纯最优反应的完整循环(9.9)，但它有三个均衡。因此，条件(9.9)并不意味着有唯一的均衡。借助最优反应示意图很容易构造出这个反例。
\end{tcolorbox}

\begin{center}
\includegraphics[max width=0.9\textwidth]{images/bo_d2e5cmbef24c73b03tn0_249_332_211_1263_378_0.jpg}
\end{center}
\hspace*{3em} 

Figure 9.11 A game with a full cycle (9.9) of pure best responses and multiple equilibria.

\begin{tcolorbox}
图9.11 一个具有纯最优反应的完整循环(9.9)和多个均衡的博弈。
\end{tcolorbox}

Figure 9.12 shows best-response regions where the subdivision of \(X\) is qualitatively at the transition between Figures 9.8 and 9.9. According to Lemma 9.3, this is a degenerate game because \(x = \left( {\frac{1}{2},\frac{1}{2},0}\right)\) has three pure best responses and four labels3,4,5,6. Consequently, only the two labels 1 and 2 of a point \(y\) are needed for an equilibrium(x, y). These two labels define \(Y\left( 1\right)  \cap  Y\left( 2\right)\) , which is the line segment with endpoints \(\left( {\frac{1}{2},\frac{1}{2},0}\right)\) and \(\left( {\frac{1}{3},\frac{1}{3},\frac{1}{3}}\right)\) in \(Y\) . Any \(y\) in that line segment is part of an equilibrium(x, y), so this game has an infinite set of equilibria.

\begin{tcolorbox}
图9.12展示了最优反应区域，其中 \(X\) 的划分在性质上处于图9.8和图9.9之间的过渡状态。根据引理9.3，这是一个退化博弈(degenerate game)，因为 \(x = \left( {\frac{1}{2},\frac{1}{2},0}\right)\) 有三个纯最优反应和四个标签3、4、5、6。因此，对于均衡(x，y)，一个点 \(y\) 只需要两个标签1和2。这两个标签定义了 \(Y\left( 1\right)  \cap  Y\left( 2\right)\) ，它是 \(Y\) 中端点为 \(\left( {\frac{1}{2},\frac{1}{2},0}\right)\) 和 \(\left( {\frac{1}{3},\frac{1}{3},\frac{1}{3}}\right)\) 的线段(line segment)。该线段中的任何 \(y\) 都是均衡(x，y)的一部分，所以这个博弈有无限个均衡。
\end{tcolorbox}

\begin{center}
\includegraphics[max width=0.9\textwidth]{images/bo_d2e5cmbef24c73b03tn0_249_335_1128_1257_379_0.jpg}
\end{center}
\hspace*{3em} 

Figure 9.12 A degenerate game with a line segment of multiple equilibria.

\begin{tcolorbox}
图9.12 一个具有多个均衡线段的退化博弈。
\end{tcolorbox}

For a nondegenerate game, Proposition 6.8 states that the equilibrium strategies in an equilibrium(x, y)mix the same number of pure strategies. In terms of the best-response diagram, this means that the smallest face of \(X\) that contains \(x\) has the same dimension as the smallest face of \(Y\) that contains \(y\) . This face is either a vertex if(x, y)is a pure-strategy equilibrium, an edge if \(x\) and \(y\) mix only two pure strategies, or a triangle if they mix three strategies (and a face of higher dimension for larger games). This is the case in the diagrams seen so far, but fails for the degenerate game in Figure 9.12 where the equilibrium strategy \(x = \left( {\frac{1}{2},\frac{1}{2},0}\right)\) on the lower edge of \(X\) and any point \(y = \left( {\alpha ,\frac{1}{2} - \frac{\alpha }{2},\frac{1}{2} - \frac{\alpha }{2}}\right)\) for \(0 < \alpha  \leq  \frac{1}{3}\) in the interior of the triangle \(Y\) form an equilibrium(x, y).

\begin{tcolorbox}
对于非退化博弈，命题6.8指出，均衡 \textbackslash ((x, y) \textbackslash )中的均衡策略所混合的纯策略数量相同。就最优反应示意图而言，这意味着包含 \textbackslash (\( X  \textbackslash )\)的最小面与包含 \textbackslash (\( x  \textbackslash )\)的最小面具有相同的维度。如果 \textbackslash ((x, y) \textbackslash )是纯策略均衡，这个面就是一个顶点；如果 \textbackslash (\( x  \textbackslash )\)和 \textbackslash (\( y  \textbackslash )\)仅混合两种纯策略，这个面就是一条边；如果它们混合三种策略，这个面就是一个三角形(对于更大的博弈，这个面的维度更高)。到目前为止所看到的示意图都是这种情况，但图9.12中的退化博弈并非如此，在该图中， \textbackslash (\( X  \textbackslash )\)下边缘的均衡策略 \textbackslash (\( x =  \textbackslash left( \{ \textbackslash frac\{1\}\{2\}, \textbackslash frac\{1\}\{2\},0\} \textbackslash right)  \textbackslash )\)与三角形 \textbackslash (\( Y  \textbackslash )\)内部 \textbackslash (\( 0 <  \textbackslash alpha   \textbackslash leq   \textbackslash frac\{1\}\{3\}  \textbackslash )\)的任意点 \textbackslash (\( y =  \textbackslash left( \{ \textbackslash alpha , \textbackslash frac\{1\}\{2\} -  \textbackslash frac\{ \textbackslash alpha \}\{2\}, \textbackslash frac\{1\}\{2\} -  \textbackslash frac\{ \textbackslash alpha \}\{2\}\} \textbackslash right)  \textbackslash )\)构成一个均衡 \textbackslash ((x, y) \textbackslash )。
\end{tcolorbox}

\section*{9.5 Strategic Equivalence}

\begin{tcolorbox}
\section*{9.5 策略等价}
\end{tcolorbox}

When constructing the best-response diagram of a game, it is often simpler to consider another game that has the same diagram. Two such games are called strategically equivalent, which we define for \(N\) -player games in general.

\begin{tcolorbox}
在构建一个博弈的最优反应示意图时，考虑另一个具有相同示意图的博弈通常会更简单。这样的两个博弈被称为策略等价，我们将对一般的 \textbackslash (\( N  \textbackslash )\)人博弈进行定义。
\end{tcolorbox}

Definition 9.5. Consider an \(N\) -player game \(G\) in strategic form with finite strategy set \({S}_{i}\) for player \(i = 1,\ldots ,N\) , payoff \({u}_{i}\left( s\right)\) for each strategy profile \(s\) , mixed strategies \({\sigma }_{i} \in  {X}_{i} = \Delta \left( {S}_{i}\right)\) , and partial profiles of mixed strategies \({\sigma }_{-i} \in  {X}_{-i} = \mathop{X}\limits_{{k \neq  i}}{X}_{k}\) . Let \(\widehat{G}\) be another game with the same strategy sets, and payoff function \({\widehat{u}}_{i}\) for each player \(i\) . Then \(G\) and \(\widehat{G}\) are called strategically equivalent if for all \(i = 1,\ldots ,N\) and all \({\sigma }_{-i} \in  {X}_{-i}\) , any strategy \({s}_{i} \in  {S}_{i}\) is a best response to \({\sigma }_{-i}\) in \(G\) if and only if \({s}_{i}\) is a best response to \({\sigma }_{-i}\) in \(\widehat{G}\) .

\begin{tcolorbox}
定义9.5。考虑一个处于策略形式的 \textbackslash (\( N  \textbackslash )\)人博弈 \textbackslash (\( G  \textbackslash )\)，其中玩家 \textbackslash (\( i = 1, \textbackslash ldots, N  \textbackslash )\)的有限策略集为 \textbackslash (\( \{S\}\_\{i\}  \textbackslash )\)，每个策略组合 \textbackslash (\( s  \textbackslash )\)的收益为 \textbackslash (\( \{u\}\_\{i\} \textbackslash left( s \textbackslash right)  \textbackslash )\)，混合策略为 \textbackslash (\( \{ \textbackslash sigma \}\_\{i\}  \textbackslash in  \{X\}\_\{i\} =  \textbackslash Delta  \textbackslash left( \{S\}\_\{i\} \textbackslash right)  \textbackslash )\)，混合策略的部分组合为 \textbackslash (\( \{ \textbackslash sigma \}\_\{-i\}  \textbackslash in  \{X\}\_\{-i\} =  \textbackslash mathop\{X\} \textbackslash limits\_\{\{k  \textbackslash neq  i\}\}\{X\}\_\{k\}  \textbackslash )\)。设 \textbackslash (\(  \textbackslash widehat\{G\}  \textbackslash )\)是另一个具有相同策略集的博弈，每个玩家 \textbackslash (\( i  \textbackslash )\)的收益函数为 \textbackslash (\( \{ \textbackslash widehat\{u\}\}\_\{i\}  \textbackslash )\)。那么，如果对于所有的 \textbackslash (\( i = 1, \textbackslash ldots, N  \textbackslash )\)和所有的 \textbackslash (\( \{ \textbackslash sigma \}\_\{-i\}  \textbackslash in  \{X\}\_\{-i\}  \textbackslash )\)，在 \textbackslash (\( G  \textbackslash )\)中，任何策略 \textbackslash (\( \{s\}\_\{i\}  \textbackslash in  \{S\}\_\{i\}  \textbackslash )\)是对 \textbackslash (\( \{ \textbackslash sigma \}\_\{-i\}  \textbackslash )\)的最优反应，当且仅当在 \textbackslash (\(  \textbackslash widehat\{G\}  \textbackslash )\)中， \textbackslash (\( \{s\}\_\{i\}  \textbackslash )\)是对 \textbackslash (\( \{ \textbackslash sigma \}\_\{-i\}  \textbackslash )\)的最优反应，则称 \textbackslash (\( G  \textbackslash )\)和 \textbackslash (\(  \textbackslash widehat\{G\}  \textbackslash )\)是策略等价的。
\end{tcolorbox}

Hence, strategic equivalence of two \(m \times  n\) bimatrix games means that any row \(i\) is a best response to a mixed strategy \(y\) of player II in one game whenever it is a best response to \(y\) in the other game, and similarly any column \(j\) is a best response to a mixed strategy \(x\) of player I in one game whenever it is in the other game. This is equivalent to saying that the two games have the same best-response regions \(Y\left( i\right)\) and \(X\left( j\right)\) in (9.3), that is, the same best-response diagram.

\begin{tcolorbox}
因此，两个 \textbackslash (\( m  \textbackslash times  n  \textbackslash )\)双矩阵博弈的策略等价意味着，在一个博弈中，任何行 \textbackslash (\( i  \textbackslash )\)是对玩家II的混合策略 \textbackslash (\( y  \textbackslash )\)的最优反应，当且仅当在另一个博弈中它也是对 \textbackslash (\( y  \textbackslash )\)的最优反应；类似地，在一个博弈中，任何列 \textbackslash (\( j  \textbackslash )\)是对玩家I的混合策略 \textbackslash (\( x  \textbackslash )\)的最优反应，当且仅当在另一个博弈中它也是对 \textbackslash (\( x  \textbackslash )\)的最优反应。这等价于说，这两个博弈在(9.3)中具有相同的最优反应区域 \textbackslash (\( Y \textbackslash left( i \textbackslash right)  \textbackslash )\)和 \textbackslash (\( X \textbackslash left( j \textbackslash right)  \textbackslash )\)，即具有相同的最优反应示意图。
\end{tcolorbox}

As Figures 6.10 and 9.1 show, the game on the right in Figure 6.7 and in Figure 9.1 are strategically equivalent. The two games have the same payoff matrix \(B\) for player II, and the payoff matrix \(A\) in the second game is obtained from the first by adding the constant 1 to each entry in the first column of \(A\) . This is generally a way to obtain a strategically equivalent game. Consider a bimatrix game(A, B), and add a constant \(\alpha\) to each entry in column \(j\) of \(A\) . Then if player II plays the mixed strategy \(y\) , with probability \({y}_{j}\) for column \(j\) , then the expected payoff in row \(i\) is changed from \({\left( Ay\right) }_{i}\) to \({\left( Ay\right) }_{i} + \alpha {y}_{j}\) . That is, all these expected payoffs change by the same amount \(\alpha {y}_{j}\) . Although this depends on \(y\) , it does not change which rows \(i\) have maximal payoff, and therefore it does not change player I’s best responses to \(y\) . We call such a payoff change of adding a constant to some column of the row player's payoffs a column shift of the game.

\begin{tcolorbox}
如图6.10和图9.1所示，图6.7右侧的博弈与图9.1中的博弈在策略上是等价的。这两个博弈对于玩家II而言具有相同的收益矩阵 \(B\) ，且第二个博弈中的收益矩阵 \(A\) 是通过给 \(A\) 第一列的每个元素加上常数1从第一个博弈的收益矩阵得到的。这通常是一种获得策略等价博弈的方法。考虑一个双矩阵博弈(A，B)，给 \(A\) 的第 \(j\) 列的每个元素加上一个常数 \(\alpha\) 。然后，如果玩家II采用混合策略 \(y\) ，即第 \(j\) 列的概率为 \({y}_{j}\) ，那么第 \(i\) 行的期望效用就从 \({\left( Ay\right) }_{i}\) 变为 \({\left( Ay\right) }_{i} + \alpha {y}_{j}\) 。也就是说，所有这些期望效用都改变了相同的量 \(\alpha {y}_{j}\) 。尽管这依赖于 \(y\) ，但它并不会改变哪些行 \(i\) 具有最大收益，因此也不会改变玩家I对 \(y\) 的最优反应。我们将这种给行玩家收益的某一列加上一个常数的收益变化称为博弈的列移位。
\end{tcolorbox}

We use "column shift" as a general term in an \(N\) -player game, by assuming that the pure strategies of the player under consideration are temporarily arranged as rows, so that any partial profile of strategies of the other \(N - 1\) players is given by a column. For example, the three-player game in Exercise 3.2 has as rows the strategies of player I and as columns the partial pure strategy profiles of players II and III. However, one should remember that therefore in a bimatrix game(A, B), a "column shift" for player II means adding a constant to some row of player II's payoff matrix \(B\) .

\begin{tcolorbox}
在一个 \(N\) 人博弈中，我们将“列移位”作为一个通用术语，假设所考虑玩家的纯策略暂时被排列为行，这样其他 \(N - 1\) 个玩家的任何部分策略组合就由一列表示。例如，练习3.2中的三人博弈将玩家I的策略作为行，将玩家II和玩家III的部分纯策略组合作为列。然而，应该记住，因此在双矩阵博弈(A，B)中，玩家II的“列移位”意味着给玩家II的收益矩阵 \(B\) 的某一行加上一个常数。
\end{tcolorbox}

Definition 9.6. Consider an \(N\) -player game \(G\) as in Definition 9.5. Let \(i = 1,\ldots ,N\) be a player, let \({a}_{i}\left( {s}_{-i}\right)  \in  \mathbb{R}\) be a separate constant for each partial profile \({s}_{-i}\) of pure strategies of the other \(N - 1\) players, and let

\begin{tcolorbox}
定义9.6。考虑如定义9.5中的一个 \(N\) 人博弈 \(G\) 。设 \(i = 1,\ldots ,N\) 为一个玩家，设 \({a}_{i}\left( {s}_{-i}\right)  \in  \mathbb{R}\) 为其他 \(N - 1\) 个玩家的每个纯策略部分组合 \({s}_{-i}\) 的一个独立常数，并且设
\end{tcolorbox}

\[
{\widehat{u}}_{i}\left( {{s}_{i},{s}_{-i}}\right)  = {u}_{i}\left( {{s}_{i},{s}_{-i}}\right)  + {a}_{i}\left( {s}_{-i}\right)
\]

be a new utility function \({\widehat{u}}_{i}\) for player \(i\) . This is called a column shift for player \(i\) .

\begin{tcolorbox}
是玩家 \(i\) 的一个新的效用函数 \({\widehat{u}}_{i}\) 。这被称为玩家 \(i\) 的列移位。
\end{tcolorbox}

Proposition 9.7. Consider an \(N\) -player game \(G\) as in Definition 9.6, with a new utility function \({\widehat{u}}_{i}\) for player \(i\) obtained by a column shift for player \(i\) . Then the game \(\widehat{G}\) with \({\widehat{u}}_{i}\) instead of \({u}_{i}\) is strategically equivalent to \(G\) .

\begin{tcolorbox}
命题9.7。考虑如定义9.6中的一个 \(N\) 人博弈 \(G\) ，玩家 \(i\) 通过列移位获得了一个新的效用函数 \({\widehat{u}}_{i}\) 。那么用 \({\widehat{u}}_{i}\) 替代 \({u}_{i}\) 后的博弈 \(\widehat{G}\) 与 \(G\) 在策略上是等价的。
\end{tcolorbox}

Proof. With the notation in Definition 9.6, let \({s}_{i} \in  {S}_{i}\) be a strategy of player \(i\) and let \({\sigma }_{-i}\) be a partial profile of mixed strategies of the other players. Then the expected payoff to player \(i\) is \({U}_{i}\left( {{s}_{i},{\sigma }_{-i}}\right)\) in \(G\) , and \({U}_{i}\left( {{s}_{i},{\sigma }_{-i}}\right)  + {a}_{i}\left( {s}_{-i}\right) \mathop{\prod }\limits_{{k = 1,k \neq  i}}^{N}{\sigma }_{k}\left( {s}_{k}\right)\) in \(\widehat{G}\) . The added term does not depend on \({s}_{i}\) . Hence, the maximum over \({s}_{i} \in  {S}_{i}\) for determining the best response of player \(i\) to \({\sigma }_{-i}\) is not affected by the column shift. That is, the best responses in \(G\) and \(\widehat{G}\) to \({\sigma }_{-i}\) are the same, as claimed.

\begin{tcolorbox}
证明。采用定义9.6中的符号，设 \({s}_{i} \in  {S}_{i}\) 是玩家 \(i\) 的一个策略，设 \({\sigma }_{-i}\) 是其他玩家的一个混合策略部分组合。那么在 \(G\) 中玩家 \(i\) 的期望效用是 \({U}_{i}\left( {{s}_{i},{\sigma }_{-i}}\right)\) ，在 \(\widehat{G}\) 中是 \({U}_{i}\left( {{s}_{i},{\sigma }_{-i}}\right)  + {a}_{i}\left( {s}_{-i}\right) \mathop{\prod }\limits_{{k = 1,k \neq  i}}^{N}{\sigma }_{k}\left( {s}_{k}\right)\) 。所加的项不依赖于 \({s}_{i}\) 。因此，为确定玩家 \(i\) 对 \({\sigma }_{-i}\) 的最优反应而对 \({s}_{i} \in  {S}_{i}\) 取最大值不受列移位的影响。也就是说，如所声称的， \(G\) 和 \(\widehat{G}\) 中对 \({\sigma }_{-i}\) 的最优反应是相同的。
\end{tcolorbox}

A column shift may mean adding the same constant to all "columns" and hence to all payoffs. This just changes the origin of the utility function of the player and always represents the same preference. In addition, the player's payoffs may be scaled by multiplying them with a positive constant (which cannot be varied with the column), which trivially creates a strategically equivalent game. However, two games may be strategically equivalent even if one game cannot be obtained from the other by scaling and column shifts. Consider the \(3 \times  2\) games(A, B)and \(\left( {\widehat{A},B}\right)\) with

\begin{tcolorbox}
列移位(column shift)可能意味着给所有“列”加上相同的常数，从而给所有收益加上相同的常数。这仅仅改变了参与者效用函数的原点，并且始终代表相同的偏好。此外，参与者的收益可以通过乘以一个正常数(该常数不能随列变化)来进行缩放，这显然会产生一个策略等价(strategic equivalence)的博弈。然而，即使一个博弈不能通过缩放和列移位从另一个博弈得到，两个博弈也可能是策略等价的。考虑 \(3 \times  2\) 博弈 (A, B) 和 \(\left( {\widehat{A},B}\right)\) ，其中
\end{tcolorbox}

\[
A = \left( \begin{array}{ll} 2 & 0 \\  1 & 2 \\  0 & 3 \end{array}\right) ,\;\widehat{A} = \left( \begin{array}{ll} 2 & 0 \\  1 & 2 \\  \frac{1}{2} & \frac{5}{2} \end{array}\right)  \tag{9.10}
\]

and any \(B\) , for example as in (9.1). The upper envelope diagrams for the two games in Figure 9.13 show that the two games are strategically equivalent because they have the same best-response regions \(Y\left( 1\right) ,Y\left( 2\right) ,Y\left( 3\right)\) , in particular the indifference between the two rows 2 and 3, which are both best responses when \(y = \left( {\frac{1}{2},\frac{1}{2}}\right)\) . However, \(\widehat{A}\) is not obtained from \(A\) by scaling and column shifts: The first two entries of the first column of \(A\) and \(\widehat{A}\) are the same, 2 and 1, and fix the "utility scale" for the payoffs in the first column, so that no scaling or additive shift can be applied to that column, but the third entry in that column is different for \(A\) and \(\widehat{A}\) .

\begin{tcolorbox}
以及任何 \(B\) ，例如如 (9.1) 中所示。图 9.13 中两个博弈的上包络(upper envelope)图表明，这两个博弈是策略等价的，因为它们具有相同的最优反应(best response)区域 \(Y\left( 1\right) ,Y\left( 2\right) ,Y\left( 3\right)\) ，特别是第 2 行和第 3 行之间的无差异(indifference)，当 \(y = \left( {\frac{1}{2},\frac{1}{2}}\right)\) 时，这两行都是最优反应。然而， \(\widehat{A}\) 不能通过缩放和列移位从 \(A\) 得到: \(A\) 和 \(\widehat{A}\) 第一列的前两个元素是相同的，即 2 和 1，这确定了第一列收益的“效用尺度(utility scale)”，因此不能对该列应用任何缩放或加法移位，但该列的第三个元素对于 \(A\) 和 \(\widehat{A}\) 是不同的。
\end{tcolorbox}

Note that we have defined strategic equivalence only in terms of best responses, which are relevant for determining an equilibrium. The rows 1 and 3 for \(A\) and \(\widehat{A}\) in (9.10) do not have the same point of indifference, which occurs for \(y = \left( {\frac{3}{5},\frac{2}{5}}\right)\) with \(A\) when both rows have payoff \(\frac{6}{5}\) , and for \(y = \left( {\frac{5}{8},\frac{3}{8}}\right)\) with \(\widehat{A}\) when both rows have payoff \(\frac{5}{4}\) . However, in both cases rows 1 and 3 are not on the upper envelope because row 2 has a higher payoff, as Figure 9.13 shows.

\begin{tcolorbox}
请注意，我们仅根据最优反应来定义策略等价，而最优反应与确定均衡(equilibrium)相关。(9.10) 中 \(A\) 和 \(\widehat{A}\) 的第 1 行和第 3 行没有相同的无差异点，对于 \(A\) ，当 \(y = \left( {\frac{3}{5},\frac{2}{5}}\right)\) 时，这两行的收益为 \(\frac{6}{5}\) ；对于 \(\widehat{A}\) ，当 \(y = \left( {\frac{5}{8},\frac{3}{8}}\right)\) 时，这两行的收益为 \(\frac{5}{4}\) 。然而，如图 9.13 所示，在这两种情况下，第 1 行和第 3 行都不在上包络上，因为第 2 行的收益更高。
\end{tcolorbox}

\begin{center}
\includegraphics[max width=0.9\textwidth]{images/bo_d2e5cmbef24c73b03tn0_252_220_215_1214_528_0.jpg}
\end{center}
\hspace*{3em} 

Figure 9.13 Upper envelopes of expected payoffs to player I against \(\left( {{y}_{1},{y}_{2}}\right)\) for the payoff matrices \(A\) and \(\widehat{A}\) in (9.10), which define strategically equivalent but not strongly strategically equivalent games(A, B)and \(\left( {\widehat{A},B}\right)\) .

\begin{tcolorbox}
图 9.13 针对 (9.10) 中收益矩阵 \(A\) 和 \(\widehat{A}\) ，参与者 I 相对于 \(\left( {{y}_{1},{y}_{2}}\right)\) 的期望效用(expected payoff)的上包络图，这些矩阵定义了策略等价但不是强策略等价(strong strategic equivalence)的博弈 (A, B) 和 \(\left( {\widehat{A},B}\right)\) 。
\end{tcolorbox}

A stronger definition (with the notation in Definition 9.5) is to say that \(G\) and \(\widehat{G}\) are strongly strategically equivalent if for all players \(i = 1,\ldots ,N\) , all \({\sigma }_{-i} \in  {X}_{-i}\) , and any two strategies \({s}_{i}\) and \({s}_{i}^{\prime }\) in \({S}_{i}\) we have

\begin{tcolorbox}
一个更强的定义(采用定义 9.5 中的符号)是，如果对于所有参与者 \(i = 1,\ldots ,N\) 、所有 \({\sigma }_{-i} \in  {X}_{-i}\) 以及 \({S}_{i}\) 中的任意两个策略 \({s}_{i}\) 和 \({s}_{i}^{\prime }\) ，我们有
\end{tcolorbox}

\[
{u}_{i}\left( {{s}_{i},{\sigma }_{-i}}\right)  \geq  {u}_{i}\left( {{s}_{i}^{\prime },{\sigma }_{-i}}\right) \; \Leftrightarrow  \;{\widehat{u}}_{i}\left( {{s}_{i},{\sigma }_{-i}}\right)  \geq  {\widehat{u}}_{i}\left( {{s}_{i}^{\prime },{\sigma }_{-i}}\right) . \tag{9.11}
\]

Clearly, strong strategic equivalence implies strategic equivalence, and is implied by any column shift because the proof of Proposition 9.6 still applies. However, the converse does not hold, as the example (9.10) shows, because (9.11) is violated for \(i\) as the row player, \({\sigma }_{-i} = y = \left( {{0.61},{0.39}}\right)\) , and rows 1 and 3 for \({s}_{i}\) and \({s}_{i}^{\prime }\) .

\begin{tcolorbox}
显然，强策略等价意味着策略等价，并且任何列移位都意味着强策略等价，因为命题 9.6 的证明仍然适用。然而，反之则不成立，如示例 (9.10) 所示，因为当 \(i\) 为行参与者、 \({\sigma }_{-i} = y = \left( {{0.61},{0.39}}\right)\) 以及 \({s}_{i}\) 和 \({s}_{i}^{\prime }\) 的第 1 行和第 3 行时，(9.11) 不成立。
\end{tcolorbox}

As an example of using column shifts to simplify the drawing of best-response diagrams, we consider the \(3 \times  3\) game(A, B)with

\begin{tcolorbox}
作为使用列移位来简化最优反应示意图(best - response diagram)绘制的一个例子，我们考虑 \(3 \times  3\) 博弈 (A, B)，其中
\end{tcolorbox}

\[
A = \left( \begin{array}{rrr} 0 & 3 & 0 \\  1 & 0 & 1 \\   - 3 & 4 & 5 \end{array}\right) ,\;B = \left( \begin{array}{rrr} 0 & 1 &  - 2 \\  2 & 0 & 3 \\  2 & 1 & 0 \end{array}\right) . \tag{9.12}
\]

We add the constant 3 to the first and third column of \(A\) , and the constant 2 to the first row of \(B\) . This produces a new game \(\left( {\widehat{A},\widehat{B}}\right)\) that is strategically equivalent to (A, B)by Proposition 9.7,

\begin{tcolorbox}
我们给 \(A\) 的第一列和第三列加上常数 3，给 \(B\) 的第一行加上常数 2。根据命题 9.7，这产生了一个与 (A, B) 策略等价的新博弈 \(\left( {\widehat{A},\widehat{B}}\right)\)
\end{tcolorbox}

\[
\widehat{A} = \left( \begin{array}{lll} 3 & 3 & 3 \\  4 & 0 & 4 \\  0 & 4 & 8 \end{array}\right) ,\;\widehat{B} = \left( \begin{array}{lll} 2 & 3 & 0 \\  2 & 0 & 3 \\  2 & 1 & 0 \end{array}\right) . \tag{9.13}
\]

The game \(\left( {\widehat{A},\widehat{B}}\right)\) , and hence the game(A, B), has the best-response diagram shown in Figure 9.14. These diagrams can be relatively easily visualized via an upper envelope of three planes of expected payoffs, spanned across three goalposts on the corners of each triangle, because the plane for the first strategy has goalposts at equal heights and is therefore horizontal. This has been shown in Figure 9.10, which produces almost the same subdivision of \(X\) because the matrix \(B\) there is almost the same as \(\widehat{B}\) .

\begin{tcolorbox}
博弈 \(\left( {\widehat{A},\widehat{B}}\right)\) ，进而博弈 (A, B) ，具有如图 9.14 所示的最优反应示意图。这些示意图可以通过期望效用的三个平面的上包络相对容易地可视化，这些平面跨越每个三角形角落的三个基准点，因为第一个策略的平面在相同高度处有基准点，因此是水平的。这已在图 9.10 中展示，该图对 \(X\) 进行的细分几乎相同，因为那里的矩阵 \(B\) 与 \(\widehat{B}\) 几乎相同。
\end{tcolorbox}

\begin{center}
\includegraphics[max width=0.8\textwidth]{images/bo_d2e5cmbef24c73b03tn0_253_418_419_1091_422_0.jpg}
\end{center}
\hspace*{3em} 

Figure 9.14 Best-response diagram for the games in (9.12) and (9.13). The games have the same three equilibria, namely \(\left( {\left( {\frac{2}{3},\frac{1}{3},0}\right) ,\left( {\frac{3}{4},\frac{1}{4},0}\right) }\right)\) (black square), \(\left( {\left( {\frac{1}{2},0,\frac{1}{2}}\right) ,\left( {\frac{1}{4},\frac{3}{4},0}\right) }\right)\) (white square), and \(\left( {\left( {0,\frac{2}{3},\frac{1}{3}}\right) ,\left( {\frac{1}{2},0,\frac{1}{2}}\right) }\right)\) (white circle).

\begin{tcolorbox}
图 9.14 博弈 (9.12) 和 (9.13) 的最优反应示意图。这些博弈有相同的三个均衡，即 \(\left( {\left( {\frac{2}{3},\frac{1}{3},0}\right) ,\left( {\frac{3}{4},\frac{1}{4},0}\right) }\right)\) (黑色方块)、 \(\left( {\left( {\frac{1}{2},0,\frac{1}{2}}\right) ,\left( {\frac{1}{4},\frac{3}{4},0}\right) }\right)\) (白色方块)和 \(\left( {\left( {0,\frac{2}{3},\frac{1}{3}}\right) ,\left( {\frac{1}{2},0,\frac{1}{2}}\right) }\right)\) (白色圆圈)。
\end{tcolorbox}

The subdivision of \(Y\) into its best-response regions \(Y\left( 1\right) ,Y\left( 2\right) ,Y\left( 3\right)\) can be obtained similarly to Figure 9.10 by considering an upper envelope of three planes of expected payoffs. We also describe another method that can be applied directly to the triangle \(Y\) , without considering a third dimension. For the original game (A, B), the best-response regions are defined by comparing the expected payoffs in each row:

\begin{tcolorbox}
将 \(Y\) 细分为其最优反应区域 \(Y\left( 1\right) ,Y\left( 2\right) ,Y\left( 3\right)\) 可以类似于图 9.10 ，通过考虑期望效用的三个平面的上包络来获得。我们还描述了另一种可以直接应用于三角形 \(Y\) 的方法，而无需考虑第三个维度。对于原始博弈 (A, B) ，最优反应区域是通过比较每行的期望效用定义的:
\end{tcolorbox}

\[
Y\left( 1\right)  = \left\{  {\left( {{y}_{1},{y}_{2},{y}_{3}}\right)  \in  Y \mid  3{y}_{2} \geq  {y}_{1} + {y}_{3}}\right\}
\]

\[
\left. {3{y}_{2} \geq   - 3{y}_{1} + 4{y}_{2} + 5{y}_{3}}\right\}
\]

\[
Y\left( 2\right)  = \left\{  {\left( {{y}_{1},{y}_{2},{y}_{3}}\right)  \in  Y \mid  {y}_{1} + {y}_{3} \geq  3{y}_{2}}\right\}   \tag{9.14}
\]

\[
\left. {{y}_{1} + {y}_{3} \geq   - 3{y}_{1} + 4{y}_{2} + 5{y}_{3}}\right\}
\]

\[
Y\left( 3\right)  = \left\{  {\left( {{y}_{1},{y}_{2},{y}_{3}}\right)  \in  Y \mid   - 3{y}_{1} + 4{y}_{2} + 5{y}_{3} \geq  3{y}_{2}}\right\}
\]

\[
\left. {-3{y}_{1} + 4{y}_{2} + 5{y}_{3} \geq  {y}_{1} + {y}_{3}}\right\}  \text{ . }
\]

For the game \(\left( {\widehat{A},\widehat{B}}\right)\) , these conditions are

\begin{tcolorbox}
对于博弈 \(\left( {\widehat{A},\widehat{B}}\right)\) ，这些条件是
\end{tcolorbox}

\[
Y\left( 1\right)  = \left\{  {\left( {{y}_{1},{y}_{2},{y}_{3}}\right)  \in  Y \mid  3{y}_{1} + 3{y}_{2} + 3{y}_{3} \geq  4{y}_{1} + 4{y}_{3},}\right.
\]

\[
\left. {3{y}_{1} + 3{y}_{2} + 3{y}_{3} \geq  4{y}_{2} + 8{y}_{3}}\right\}
\]

\[
Y\left( 2\right)  = \left\{  {\left( {{y}_{1},{y}_{2},{y}_{3}}\right)  \in  Y \mid  4{y}_{1} + 4{y}_{3} \geq  3{y}_{1} + 3{y}_{2} + 3{y}_{3}}\right\}   \tag{9.15}
\]

\[
\left. {4{y}_{1} + 4{y}_{3} \geq  4{y}_{2} + 8{y}_{3}}\right\}
\]

\[
Y\left( 3\right)  = \left\{  {\left( {{y}_{1},{y}_{2},{y}_{3}}\right)  \in  Y \mid  4{y}_{2} + 8{y}_{3} \geq  3{y}_{1} + 3{y}_{2} + 3{y}_{3}}\right\}
\]

\[
\left. {4{y}_{2} + 8{y}_{3} \geq  4{y}_{1} + 4{y}_{3}}\right\}  \text{ . }
\]

In agreement with Proposition 9.7, these two sets of conditions are completely equivalent, because the term \(3{y}_{1} + 3{y}_{3}\) is added on both sides of each inequality; this corresponds to the column shift of adding the constant 3 to the first and third column of \(A\) to obtain \(\widehat{A}\) .

\begin{tcolorbox}
与命题 9.7 一致，这两组条件是完全等价的，因为项 \(3{y}_{1} + 3{y}_{3}\) 被加到每个不等式的两边；这对应于通过给 \(A\) 的第一列和第三列加上常数 3 以得到 \(\widehat{A}\) 的列移位。
\end{tcolorbox}

\begin{center}
\includegraphics[max width=0.3\textwidth]{images/bo_d2e5cmbef24c73b03tn0_254_257_201_420_418_0.jpg}
\end{center}
\hspace*{3em} 

Figure 9.15 Construction of the subdivision of \(Y\) in the plane, without an upper envelope in three dimensions.

\begin{tcolorbox}
图 9.15 在平面中对 \(Y\) 进行细分的构建，无需三维上包络。
\end{tcolorbox}

We now consider the indifferences between each pair of rows, and write them as follows:

\begin{tcolorbox}
我们现在考虑每两行之间的无差异，并将它们写成如下形式:
\end{tcolorbox}

\[
\text{ ① } \sim  \text{ ②: }\;3{y}_{2} = {y}_{1} + {y}_{3}\text{ , }
\]

\[
\text{ ① } \sim  \text{ ③: }\;3{y}_{1} = {y}_{2} + 5{y}_{3}\text{ , } \tag{9.16}
\]

\[
\text{ ② } \sim  \text{ ③ } : 4{y}_{2} + 4{y}_{3} = 4{y}_{1}\text{ . }
\]

The equations in (9.16) are obtained by converting an inequality in (9.14) or (9.15) into an equality and then sorting the coefficients of \({y}_{1},{y}_{2}\) , and \({y}_{3}\) so that each variable appears only on one side of the equation with a positive coefficient (or not at all, if the variable has the same coefficient on both sides, which is not the case here). This has already been done in (9.4) for the game (9.1), where it conveniently identified the common line segment of the two best-response regions \(X\left( 4\right)\) and \(X\left( 5\right)\) .

\begin{tcolorbox}
式 (9.16) 中的方程是通过将式 (9.14) 或 (9.15) 中的不等式转换为等式，然后对 \({y}_{1},{y}_{2}\) 和 \({y}_{3}\) 的系数进行排序得到的，使得每个变量仅以正系数出现在方程的一侧(或者根本不出现，如果该变量在两边的系数相同，但这里不是这种情况)。对于博弈 (9.1) ，在式 (9.4) 中已经进行了这样的处理，在那里它方便地确定了两个最优反应区域 \(X\left( 4\right)\) 和 \(X\left( 5\right)\) 的公共线段。
\end{tcolorbox}

In the present case, the three best-response regions \(Y\left( 1\right) ,Y\left( 2\right)\) , and \(Y\left( 3\right)\) can be found by extending that approach, which is illustrated in Figure 9.15. The first equation in (9.16) shows where rows 1 and 2 have the same expected payoff. All three variables \({y}_{1},{y}_{2},{y}_{3}\) appear in this equation, with \({y}_{1}\) on the left-hand side and \({y}_{2}\) and \({y}_{3}\) on the right-hand side. Setting \({y}_{3} = 0\) then gives the solution \(y = \left( {\frac{3}{4},\frac{1}{4},0}\right)\) , which is shown with a black square in Figure 9.15, and setting \({y}_{1} = 0\) gives another solution \(y = \left( {0,\frac{1}{4},\frac{3}{4}}\right)\) , shown as point \(a\) in the figure. The line segment that connects these two points separates \(Y\left( 1\right)\) and \(Y\left( 2\right)\) . In order to decide on which side of this line segment the labels 1 and 2 should be drawn, it is useful to identify the best responses to the pure strategies of player II, which are row 2 against \({e}_{1}\) and row 3 against both \({e}_{2}\) and \({e}_{3}\) , shown by the labels in the corner of the triangle.

\begin{tcolorbox}
在当前情况下，可以通过扩展该方法找到三个最优反应区域 \(Y\left( 1\right) ,Y\left( 2\right)\) 和 \(Y\left( 3\right)\) ，如图9.15所示。(9.16)中的第一个方程表明了行1和行2具有相同期望效用(expected payoff)的位置。该方程中出现了所有三个变量 \({y}_{1},{y}_{2},{y}_{3}\) ，其中 \({y}_{1}\) 在左侧， \({y}_{2}\) 和 \({y}_{3}\) 在右侧。令 \({y}_{3} = 0\) ，则得到解 \(y = \left( {\frac{3}{4},\frac{1}{4},0}\right)\) ，在图9.15中用黑色方块表示；令 \({y}_{1} = 0\) ，则得到另一个解 \(y = \left( {0,\frac{1}{4},\frac{3}{4}}\right)\) ，在图中表示为点 \(a\) 。连接这两个点的线段(line segment)将 \(Y\left( 1\right)\) 和 \(Y\left( 2\right)\) 分开。为了确定在这条线段的哪一侧应绘制标签1和2，识别对玩家II的纯策略的最优反应是有用的，即针对 \({e}_{1}\) 的最优反应是行2，针对 \({e}_{2}\) 和 \({e}_{3}\) 的最优反应是行3，这由三角形角落的标签表示。
\end{tcolorbox}

Similarly, the "extreme" solutions to the second equation in (9.16) are \(y =\)  \(\left( {\frac{1}{4},\frac{3}{4},0}\right)\) when \({y}_{3} = 0\) , shown as a white square in Figure 9.15, and \(b = \left( {\frac{5}{8},0,\frac{3}{8}}\right)\) when \({y}_{2} = 0\) . For the third equation in (9.16), they are \(c = \left( {\frac{1}{2},\frac{1}{2},0}\right)\) and \(\left( {\frac{1}{2},0,\frac{1}{2}}\right)\) , shown as a white circle (which shows that \({y}_{1} = \frac{1}{2}\) in any solution to this equation). The three best-response regions have a common point in the interior of \(Y\) , which can be computed (although its precise location is not relevant for the qualitative shape of the best-response diagram) as \(\left( {\frac{1}{2},\frac{1}{4},\frac{1}{4}}\right)\) .

\begin{tcolorbox}
类似地，(9.16)中第二个方程的“极端”解是:当 \({y}_{3} = 0\) 时为 \(y =\)  \(\left( {\frac{1}{4},\frac{3}{4},0}\right)\) ，在图9.15中用白色方块表示；当 \({y}_{2} = 0\) 时为 \(b = \left( {\frac{5}{8},0,\frac{3}{8}}\right)\) 。对于(9.16)中的第三个方程，其解为 \(c = \left( {\frac{1}{2},\frac{1}{2},0}\right)\) 和 \(\left( {\frac{1}{2},0,\frac{1}{2}}\right)\) ，用白色圆圈表示(这表明在该方程的任何解中 \({y}_{1} = \frac{1}{2}\) )。三个最优反应区域在 \(Y\) 的内部有一个公共点，该点可以计算得出(尽管其精确位置与最优反应示意图(best-response diagram)的定性形状无关)，计算结果为 \(\left( {\frac{1}{2},\frac{1}{4},\frac{1}{4}}\right)\) 。
\end{tcolorbox}

In order to identify the equilibria of the game (9.12), the best-response diagram in Figure 9.14 need not necessarily be constructed with full precision, only to show where the best-response regions meet in each triangle to identify the points that have three labels. The three equilibria are shown with symbols as before.

\begin{tcolorbox}
为了确定博弈(9.12)的均衡，图9.14中的最优反应示意图不必精确构建，只需显示每个三角形中最优反应区域的交汇位置，以确定具有三个标签的点。三个均衡如前所示用符号表示。
\end{tcolorbox}

\begin{center}
\includegraphics[max width=0.8\textwidth]{images/bo_d2e5cmbef24c73b03tn0_255_418_591_1090_627_0.jpg}
\end{center}
\hspace*{3em} 

Figure 9.16 Lemke-Howson paths for the game (9.12), shown in (9.17) and (9.18), for missing label 1.

\begin{tcolorbox}
图9.16 博弈(9.12)中缺失标签1的莱姆基 - 豪森路径(Lemke-Howson path)，如(9.17)和(9.18)所示。
\end{tcolorbox}

Figure 9.16 shows the Lemke-Howson paths for missing label 1 for the game (9.12) using the best-response diagram in Figure 9.14. As before, the graphs \(\widetilde{X}\) and \(\widetilde{Y}\) are created by adding the point0to \(X\) and \(Y\) , which is connected to the vertices of \(X\) and \(Y\) , respectively, by an edge. The paths are as follows, where we only identify the pairs(x, y)in \(\widetilde{X} \times  \widetilde{Y}\) by their sets of labels, and equilibria by their symbols, so as not to clutter the diagram:

\begin{tcolorbox}
图9.16使用图9.14中的最优反应示意图展示了博弈(9.12)中缺失标签1的莱姆基 - 豪森路径。如前所述，图(graph) \(\widetilde{X}\) 和 \(\widetilde{Y}\) 是通过将点0分别添加到 \(X\) 和 \(Y\) 中创建的，点0分别通过一条边(edge)与 \(X\) 和 \(Y\) 的顶点相连。路径如下，为避免图表杂乱，我们仅通过 \(\widetilde{X} \times  \widetilde{Y}\) 中的点集的标签来识别对(x, y)，并通过符号来识别均衡:
\end{tcolorbox}

(9.17) (9.18)

\begin{center}
\adjustbox{max width=\textwidth}{
\begin{tabular}{|c|c|c|}
\hline
(x, y) & labels & drop label \\
\cline{1-3}
\(\bullet\) & 123,456 & 1 in \(\widetilde{X}\) \\
\cline{1-3}
\phantom{X} & 235,456 & 5 in \(\widetilde{Y}\) \\
\cline{1-3}
\phantom{X} & 235,356 & 3 in \(\widetilde{X}\) \\
\cline{1-3}
\phantom{X} & 245, 356 & 4 in \(\widetilde{Y}\) \\
\cline{1-3}
\(\square\) & 245,136 & (found 1 in \(\widetilde{Y}\) ) \\
\cline{1-3}
\hline
\end{tabular}
}
\end{center}

\begin{tcolorbox}

\begin{center}
\adjustbox{max width=\textwidth}{
\begin{tabular}{|c|c|c|}
\hline
(x，y) & 标签 & 删除标签 \\
\cline{1-3}
\(\bullet\) & 123,456 & 在 \(\widetilde{X}\) 中有 1 个 \\
\cline{1-3}
\phantom{X} & 235,456 & 在 \(\widetilde{Y}\) 中有 5 个 \\
\cline{1-3}
\phantom{X} & 235,356 & 在 \(\widetilde{X}\) 中有 3 个 \\
\cline{1-3}
\phantom{X} & 245, 356 & 在 \(\widetilde{Y}\) 中有 4 个 \\
\cline{1-3}
\(\square\) & 245,136 & (在 \(\widetilde{Y}\) 中找到 1 个) \\
\cline{1-3}
\hline
\end{tabular}
}
\end{center}

\end{tcolorbox}

\begin{center}
\adjustbox{max width=\textwidth}{
\begin{tabular}{|c|c|c|}
\hline
(x, y) & labels & drop label \\
\cline{1-3}
O & 146,235 & 1 in \(\widetilde{X}\) \\
\cline{1-3}
\phantom{X} & 346,235 & 3 in \(\widetilde{Y}\) \\
\cline{1-3}
\phantom{X} & 346, 256 & 6 in \(\widetilde{X}\) \\
\cline{1-3}
\phantom{X} & 345, 256 & 5 in \(\widetilde{Y}\) \\
\cline{1-3}
- & 345, 126 & (found 1 in \(\widetilde{Y}\) ). \\
\cline{1-3}
\hline
\end{tabular}
}
\end{center}

\begin{tcolorbox}

\begin{center}
\adjustbox{max width=\textwidth}{
\begin{tabular}{|c|c|c|}
\hline
(x，y) & 标签 & 删除标签 \\
\cline{1-3}
O & 146,235 & 在 \(\widetilde{X}\) 中有1个 \\
\cline{1-3}
\phantom{X} & 346,235 & 在 \(\widetilde{Y}\) 中有3个 \\
\cline{1-3}
\phantom{X} & 346, 256 & 在 \(\widetilde{X}\) 中有6个 \\
\cline{1-3}
\phantom{X} & 345, 256 & 在 \(\widetilde{Y}\) 中有5个 \\
\cline{1-3}
- & 345, 126 & (在 \(\widetilde{Y}\) 中找到1个)。 \\
\cline{1-3}
\hline
\end{tabular}
}
\end{center}

\end{tcolorbox}

It is a useful exercise to find the Lemke-Howson paths in Figure 9.16 for each missing label. For example, the equilibrium \(\left( {\left( {0,\frac{2}{3},\frac{1}{3}}\right) ,\left( {\frac{1}{2},0,\frac{1}{2}}\right) }\right)\) marked with a white circle is found from the artificial equilibrium(0,0)for missing label 3, so this is different from the equilibrium found for missing label 1 . The LH path that starts at the white circle for missing label 1 again is shown in (9.18) and finds the third equilibrium. By considering all possible missing labels, one can construct a "network" of LH paths that can be traversed back and forth. However, this network may be disconnected, so that not all equilibria are found when starting from the artificial equilibrium; see Exercise 9.1.

\begin{tcolorbox}
针对每个缺失标签，找出图9.16中的莱姆基 - 豪森路径是一项有益的练习。例如，用白色圆圈标记的均衡 \(\left( {\left( {0,\frac{2}{3},\frac{1}{3}}\right) ,\left( {\frac{1}{2},0,\frac{1}{2}}\right) }\right)\) 是从人工均衡(0,0)出发，针对缺失标签3找到的，因此这与针对缺失标签1找到的均衡不同。针对缺失标签1，再次从白色圆圈开始的莱姆基 - 豪森路径如(9.18)所示，并能找到第三个均衡。通过考虑所有可能的缺失标签，可以构建一个能够来回遍历的莱姆基 - 豪森路径“网络”。然而，这个网络可能是不连通的，因此从人工均衡出发时，并非所有的均衡都能被找到；参见练习9.1。
\end{tcolorbox}

\section*{9.6 Best-Response Polyhedra and Polytopes}

\begin{tcolorbox}
\section*{9.6 最优反应多面体和有界多面体}
\end{tcolorbox}

As shown in the previous two sections, best-response diagrams help visualize the structure of equilibria in \(3 \times  3\) games. They are also a good basis for general algorithms to find equilibria in larger bimatrix games. These algorithms, including the Lemke-Howson algorithm, are based on the upper envelope diagram, which has a general definition in terms of suitable "best-response polyhedra", as we explain in this section.

\begin{tcolorbox}
如前两节所示，最优反应示意图有助于直观呈现双矩阵博弈( \(3 \times  3\) )中均衡的结构。它们也是在更大的双矩阵博弈中寻找均衡的通用算法的良好基础。包括莱姆基 - 豪森算法在内的这些算法，都基于上包络图，我们将在本节中说明，该图依据合适的“最优反应多面体”有一个通用定义。
\end{tcolorbox}

As mentioned earlier, a hyperplane \(H\) in \({\mathbb{R}}^{d}\) is a set \(H = \left\{  {z \in  {\mathbb{R}}^{d} \mid  {v}^{\top }z = \alpha }\right\}\) for a nonzero normal vector \(v\) in \({\mathbb{R}}^{d}\) and a real number \(\alpha\) . The corresponding halfspace is defined as \(\left\{  {z \in  {\mathbb{R}}^{d} \mid  {v}^{\top }z \leq  \alpha }\right\}\) . A polyhedron is the intersection of a finite number of halfspaces. If the polyhedron \(S\) is the intersection of \(n\) halfspaces \(\left\{  {z \in  {\mathbb{R}}^{d} \mid  {c}_{i}^{\top }z \leq  {q}_{i}}\right\}\) with \({c}_{i} \in  {\mathbb{R}}^{d}\) and \({q}_{i} \in  \mathbb{R}\) for \(i = 1,\ldots ,n\) , then it can be written as

\begin{tcolorbox}
如前所述， \({\mathbb{R}}^{d}\) 中的超平面 \(H\) 是一个集合 \(H = \left\{  {z \in  {\mathbb{R}}^{d} \mid  {v}^{\top }z = \alpha }\right\}\) ，其中 \({\mathbb{R}}^{d}\) 中的非零法向量为 \(v\) ，实数为 \(\alpha\) 。相应的半空间定义为 \(\left\{  {z \in  {\mathbb{R}}^{d} \mid  {v}^{\top }z \leq  \alpha }\right\}\) 。多面体是有限个半空间的交集。如果多面体 \(S\) 是 \(n\) 个半空间 \(\left\{  {z \in  {\mathbb{R}}^{d} \mid  {c}_{i}^{\top }z \leq  {q}_{i}}\right\}\) 的交集，其中 \({c}_{i} \in  {\mathbb{R}}^{d}\) 且 \({q}_{i} \in  \mathbb{R}\) ，对于 \(i = 1,\ldots ,n\) ，那么它可以写成
\end{tcolorbox}

\[
S = \left\{  {z \in  {\mathbb{R}}^{d} \mid  {Cz} \leq  q}\right\}   \tag{9.19}
\]

for the \(n \times  d\) matrix \(C\) and \(q \in  {\mathbb{R}}^{n}\) where \({C}^{\top } = \left\lbrack  {{c}_{1}\cdots {c}_{n}}\right\rbrack\) , that is, \({c}_{i}^{\top }\) is the \(i\) th row of \(C\) .

\begin{tcolorbox}
对于 \(n \times  d\) 矩阵 \(C\) 和 \(q \in  {\mathbb{R}}^{n}\) ，其中 \({C}^{\top } = \left\lbrack  {{c}_{1}\cdots {c}_{n}}\right\rbrack\) ，即 \({c}_{i}^{\top }\) 是 \(C\) 的第 \(i\) 行。
\end{tcolorbox}

If \(z \in  S\) and \({c}_{i}^{\top }z = {q}_{i}\) , then we say that the \(i\) th inequality \({c}_{i}^{\top }z \leq  {q}_{i}\) in (9.19) is tight for \(z\) (that is, holds as an equality). A face \(F\) of \(S\) is obtained by requiring some of the \(n\) inequalities in (9.19) to be tight (as a result, \(F\) may be the empty set). That is, \(F = \left\{  {z \in  S \mid  {c}_{i}^{\top }z = {q}_{i}}\right.\) for \(\left. {i \in  T}\right\}\) for some set \(T \subseteq  \{ 1,\ldots ,n\}\) . There are \({2}^{n}\) ways of choosing \(T\) (some of which may define the same face \(F\) ), so the number of faces is finite. If the face consists of a single point, this is called a vertex of the polyhedron. A bounded polyhedron is called a polytope. It can be shown that every polytope is the convex hull of its vertices, which agrees with our earlier definition of a polytope as the convex hull of a finite set of points in \({\mathbb{R}}^{d}\) .

\begin{tcolorbox}
如果 \(z \in  S\) 且 \({c}_{i}^{\top }z = {q}_{i}\) ，那么我们称 (9.19) 中的第 \(i\) 个不等式 \({c}_{i}^{\top }z \leq  {q}_{i}\) 对于 \(z\) 是紧的(即，取等号成立)。通过要求 (9.19) 中的某些 \(n\) 不等式为紧不等式(结果是， \(F\) 可能为空集)，可以得到 \(S\) 的一个面 \(F\) 。也就是说，对于某个集合 \(T \subseteq  \{ 1,\ldots ,n\}\) 中的 \(\left. {i \in  T}\right\}\) ，有 \(F = \left\{  {z \in  S \mid  {c}_{i}^{\top }z = {q}_{i}}\right.\) 。选择 \(T\) 有 \({2}^{n}\) 种方式(其中一些可能定义相同的面 \(F\) )，因此面的数量是有限的。如果这个面由单个点组成，那么这个点被称为多面体的顶点。有界多面体称为有界多面体(polytope)。可以证明，每个有界多面体都是其顶点的凸包，这与我们之前将有界多面体定义为 \({\mathbb{R}}^{d}\) 中有限点集的凸包是一致的。
\end{tcolorbox}

The mixed-strategy sets \(X\) and \(Y\left( {6.6}\right)\) are polytopes, because they are clearly bounded, and defined by a number of linear inequalities and one linear equation (any equation \({v}^{\top }z = \alpha\) can be written as two inequalities \({v}^{\top }z \leq  \alpha\) and \(- {v}^{\top }z \leq   - \alpha\) ).

\begin{tcolorbox}
混合策略集 \(X\) 和 \(Y\left( {6.6}\right)\) 是有界多面体，因为它们显然是有界的，并且由多个线性不等式和一个线性方程定义(任何方程 \({v}^{\top }z = \alpha\) 都可以写成两个不等式 \({v}^{\top }z \leq  \alpha\) 和 \(- {v}^{\top }z \leq   - \alpha\) )。
\end{tcolorbox}

As a central concept of this section, we consider the following best-response polyhedra \(\bar{P}\) for player \(\mathrm{I}\) and \(\bar{Q}\) for player \(\mathrm{{II}}\) , for an \(m \times  n\) bimatrix game(A, B):

\begin{tcolorbox}
作为本节的核心概念，对于一个 \(m \times  n\) 双矩阵博弈 (A, B)，我们考虑玩家 \(\mathrm{I}\) 的以下最优反应多面体 \(\bar{P}\) 和玩家 \(\mathrm{{II}}\) 的 \(\bar{Q}\) :
\end{tcolorbox}

\[
\bar{P} = \left\{  {\left( {x,v}\right)  \in  X \times  \mathbb{R} \mid  {B}^{\top }x \leq  \mathbf{1}v}\right\}  , \tag{9.20}
\]

\[
\bar{Q} = \{ \left( {y,u}\right)  \in  Y \times  \mathbb{R} \mid  {Ay} \leq  \mathbf{1}u\} .
\]

For the game (9.1), the inequalities \({B}^{\top }x \leq  \mathbf{1}v\) that (in addition to \(x \in  X\) ) define \(\bar{P}\) are

\begin{tcolorbox}
对于博弈 (9.1)，(除 \(x \in  X\) 之外)定义 \(\bar{P}\) 的不等式 \({B}^{\top }x \leq  \mathbf{1}v\) 为
\end{tcolorbox}

\[
2{x}_{1}\; + 4{x}_{3} \leq  v \tag{9.21}
\]

\[
{x}_{2} + {x}_{3} \leq  v
\]

and the inequalities \({Ay} \leq  {1u}\) that (in addition to \(y \in  Y\) ) define \(\bar{Q}\) are

\begin{tcolorbox}
并且，(除 \(y \in  Y\) 之外)定义 \(\bar{Q}\) 的不等式 \({Ay} \leq  {1u}\) 为
\end{tcolorbox}

\[
2{y}_{1}\; \leq  u,
\]

\[
{y}_{1} + 2{y}_{2} \leq  u \tag{9.22}
\]

\[
3{y}_{2} \leq  u\text{ . }
\]

They are displayed in the upper envelope diagrams in Figures 9.2 and 9.1 where \(v\) and \(u\) are at least the payoff to player II and player I, respectively. (Because \(v\) and \(u\) have no upper bound, \(\bar{P}\) and \(\bar{Q}\) are polyhedra but not polytopes.) As soon as at least one of the inequalities in (9.21) and (9.22) is tight, \(v\) and \(u\) are equal to the best-response payoff because they cannot take a lower value. That is, for given \(x \in  X\) and \(y \in  Y\) ,

\begin{tcolorbox}
它们显示在图 9.2 和图 9.1 的上包络图中，其中 \(v\) 和 \(u\) 分别至少是玩家 II 和玩家 I 的收益。(因为 \(v\) 和 \(u\) 没有上界，所以 \(\bar{P}\) 和 \(\bar{Q}\) 是多面体但不是有界多面体。)一旦 (9.21) 和 (9.22) 中的至少一个不等式是紧的， \(v\) 和 \(u\) 就等于最优反应收益，因为它们不能取更低的值。也就是说，对于给定的 \(x \in  X\) 和 \(y \in  Y\) ，
\end{tcolorbox}

\[
\min \{ v \mid  \left( {x,v}\right)  \in  \bar{P}\}  = \max \left\{  {{\left( {B}^{\top }x\right) }_{j} \mid  1 \leq  j \leq  n}\right\}  ,
\]

\[
\min \{ u \mid  \left( {y,u}\right)  \in  \bar{Q}\}  = \max \left\{  {{\left( Ay\right) }_{i} \mid  1 \leq  i \leq  m}\right\}  ,
\]

which is just the familiar definition of the upper envelope of expected payoffs. Assuming that at least one of the inequalities in \({B}^{\top }x \leq  \mathbf{1}v\) is tight (which defines \(v\) as the best-response payoff to \(x\) ), the tight inequalities in \({B}^{\top }x \leq  \mathbf{1}v\) correspond to the pure best responses to \(x\) . Similarly, the tight inequalities in \({Ay} \leq  \mathbf{1}u\) (if at least one is tight) correspond to the pure best responses to \(y\) .

\begin{tcolorbox}
这正是我们熟悉的期望效用上包络的定义。假设 \({B}^{\top }x \leq  \mathbf{1}v\) 中的至少一个不等式是紧的(这将 \(v\) 定义为对 \(x\) 的最优反应收益)， \({B}^{\top }x \leq  \mathbf{1}v\) 中的紧不等式对应于对 \(x\) 的纯最优反应。类似地， \({Ay} \leq  \mathbf{1}u\) 中的紧不等式(如果至少有一个是紧的)对应于对 \(y\) 的纯最优反应。
\end{tcolorbox}

Hence, we can identify the labels that denote the pure strategies \(1,\ldots ,m + n\) of the players as in (9.2) with the tight inequalities in \(\bar{P}\) and \(\bar{Q}\) . Let \(\left( {x,v}\right)  \in  \bar{P}\) with \(v\) as best-response payoff to \(x\) , let \(\left( {y,u}\right)  \in  \bar{Q}\) with \(u\) as best-response payoff to \(y\) , and \(1 \leq  i \leq  m\) and \(1 \leq  j \leq  n\) . Then, in agreement with Definition 9.1,

\begin{tcolorbox}
因此，我们可以将表示参与者纯策略 \(1,\ldots ,m + n\) 的标签(如 (9.2) 中所示)与 \(\bar{P}\) 和 \(\bar{Q}\) 中的紧不等式相对应。设 \(\left( {x,v}\right)  \in  \bar{P}\) ，其中 \(v\) 是对 \(x\) 的最优响应收益；设 \(\left( {y,u}\right)  \in  \bar{Q}\) ，其中 \(u\) 是对 \(y\) 的最优响应收益，且 \(1 \leq  i \leq  m\) 和 \(1 \leq  j \leq  n\) 。那么，根据定义 9.1，
\end{tcolorbox}

\[
x\text{ has label }i \Leftrightarrow  {x}_{i} = 0\text{ , }
\]

\[
x\text{ has label }m + j\; \Leftrightarrow  \;{\left( {B}^{\top }x\right) }_{j} = v\text{ , } \tag{9.23}
\]

\[
y\text{ has label }i\; \Leftrightarrow  \;{\left( Ay\right) }_{i} = u\text{ , }
\]

\[
y\text{ has label }m + j\; \Leftrightarrow  \;{y}_{j} = 0\text{ . }
\]

We can therefore rephrase Theorem 9.2 as follows.

\begin{tcolorbox}
因此，我们可以将定理 9.2 重新表述如下。
\end{tcolorbox}

Corollary 9.8. Let \(\left( {x,v}\right)  \in  \bar{P}\) and \(\left( {y,u}\right)  \in  \bar{Q}\) . Then(x, y)is an equilibrium of(A, B) with equilibrium payoffs \(u\) and \(v\) to players \(I\) and \({II}\) if and only if(x, y)is completely labeled.

\begin{tcolorbox}
推论 9.8。设 \(\left( {x,v}\right)  \in  \bar{P}\) 和 \(\left( {y,u}\right)  \in  \bar{Q}\) 。那么，(x, y) 是 (A, B) 的一个均衡，参与者 \(I\) 和 \({II}\) 的均衡收益分别为 \(u\) 和 \(v\) ，当且仅当 (x, y) 是完全标记的。
\end{tcolorbox}

In Definition 9.1, labels are defined in terms of the best-response regions \(Y\left( i\right)\) and \(X\left( j\right)\) in (9.3), and the zero faces of \(X\) and \(Y\) in (9.5). The best-response regions are polytopes, for example as in (9.14). While the best-response regions are convenient to draw in low dimension (in particular as subsets of a triangle), the definition as faces of the best-response polyhedra \(\bar{P}\) and \(\bar{Q}\) requires only a single tight inequality per label, by adding one more dimension for the payoff \(v\) or \(u\) to the other player.

\begin{tcolorbox}
在定义 9.1 中，标签是根据 (9.3) 中的最优响应区域 \(Y\left( i\right)\) 和 \(X\left( j\right)\) 以及 (9.5) 中 \(X\) 和 \(Y\) 的零面来定义的。最优响应区域是有界多面体，例如 (9.14) 中的情况。虽然最优响应区域在低维度(特别是作为三角形的子集)中便于绘制，但将其定义为最优响应多面体 \(\bar{P}\) 和 \(\bar{Q}\) 的面时，每个标签只需要一个紧不等式，只需为另一个参与者的收益 \(v\) 或 \(u\) 增加一个维度。
\end{tcolorbox}

We now simplify the best-response polyhedra further by eliminating these payoff variables, and replacing the mixed strategies \(x\) and \(y\) by arbitrary nonnegative vectors. Let

\begin{tcolorbox}
现在，我们通过消除这些收益变量，并将混合策略 \(x\) 和 \(y\) 替换为任意非负向量，进一步简化最优响应多面体。设
\end{tcolorbox}

\[
P = \left\{  {x \in  {\mathbb{R}}^{m} \mid  x \geq  \mathbf{0},{B}^{\top }x \leq  \mathbf{1}}\right\}  , \tag{9.24}
\]

\[
Q = \left\{  {y \in  {\mathbb{R}}^{n} \mid  {Ay} \leq  \mathbf{1},\;y \geq  \mathbf{0}}\right\}  .
\]

For the game (9.1), \(P\) is defined by the inequalities \({x}_{1} \geq  0,{x}_{2} \geq  0,{x}_{3} \geq  0\) , and

\begin{tcolorbox}
对于博弈 (9.1)， \(P\) 由不等式 \({x}_{1} \geq  0,{x}_{2} \geq  0,{x}_{3} \geq  0\) 定义，并且
\end{tcolorbox}

\[
2{x}_{1}\; + 4{x}_{3} \leq  1
\]

\[
{x}_{2} + {x}_{3} \leq  1
\]

so that \(P\) is the intersection of the "wedges" \({x}_{3} \leq  \frac{1}{4} - \frac{1}{2}{x}_{1}\) and \({x}_{3} \leq  1 - {x}_{2}\) shown on the left in Figure 9.17. The constraints for \(Q\) are \({y}_{1} \geq  0,{y}_{2} \geq  0\) , and

\begin{tcolorbox}
因此， \(P\) 是图 9.17 左侧所示的“楔形” \({x}_{3} \leq  \frac{1}{4} - \frac{1}{2}{x}_{1}\) 和 \({x}_{3} \leq  1 - {x}_{2}\) 的交集。 \(Q\) 的约束条件是 \({y}_{1} \geq  0,{y}_{2} \geq  0\) ，并且
\end{tcolorbox}

\[
2{y}_{1}\; \leq  1
\]

\[
{y}_{1} + 2{y}_{2} \leq  1 \tag{9.25}
\]

\[
3{y}_{2} \leq  1
\]

which translate to \({y}_{1} \leq  \frac{1}{2},{y}_{2} \leq  \frac{1}{2} - \frac{1}{2}{y}_{1},{y}_{2} \leq  \frac{1}{3}\) , and define a polygon in \({\mathbb{R}}^{2}\) .

\begin{tcolorbox}
这转化为 \({y}_{1} \leq  \frac{1}{2},{y}_{2} \leq  \frac{1}{2} - \frac{1}{2}{y}_{1},{y}_{2} \leq  \frac{1}{3}\) ，并在 \({\mathbb{R}}^{2}\) 中定义了一个多边形。
\end{tcolorbox}

In general, we want \(P\) and \(Q\) to be polytopes, which is equivalent to \(v > 0\) and \(u > 0\) for any \(\left( {x,v}\right)  \in  \bar{P}\) and \(\left( {y,u}\right)  \in  \bar{Q}\) , according to the following lemma.

\begin{tcolorbox}
一般来说，我们希望 \(P\) 和 \(Q\) 是有界多面体，根据以下引理，这等价于对于任何 \(\left( {x,v}\right)  \in  \bar{P}\) 和 \(\left( {y,u}\right)  \in  \bar{Q}\) 有 \(v > 0\) 和 \(u > 0\) 。
\end{tcolorbox}

Lemma 9.9. Consider an \(m \times  n\) bimatrix game(A, B). Then \(P\) in (9.24) is a polytope if and only if the best-response payoff to any \(x\) in \(X\) is always positive, and \(Q\) in (9.24) is a polytope if and only if the best-response payoff to any \(y\) in \(Y\) is always positive.

\begin{tcolorbox}
引理 9.9。考虑一个 \(m \times  n\) 双矩阵博弈 (A, B)。那么，(9.24) 中的 \(P\) 是有界多面体，当且仅当对 \(X\) 中任何 \(x\) 的最优响应收益总是正的；(9.24) 中的 \(Q\) 是有界多面体，当且仅当对 \(Y\) 中任何 \(y\) 的最优响应收益总是正的。
\end{tcolorbox}

Proof. We prove the statement for \(Q\) ; the proof for \(P\) is analogous. The best-response payoff to any mixed strategy \(y\) is the maximum entry of \({Ay}\) , so this is not always positive if and only if \({Ay} \leq  \mathbf{0}\) for some \(y \in  Y\) . For such a \(y\) we have \(y \geq  0,y \neq  \mathbf{0}\) , and \({y\alpha } \in  Q\) for any \(\alpha  \geq  0\) , which shows that \(Q\) is not bounded. Conversely, suppose the best-response payoff \(u\) to any \(y\) is always positive. Because

\begin{tcolorbox}
证明。我们针对 \(Q\) 来证明该命题；针对 \(P\) 的证明是类似的。对任何混合策略 \(y\) 的最优反应收益是 \({Ay}\) 中的最大元素，因此，当且仅当对于某个 \(y \in  Y\) 有 \({Ay} \leq  \mathbf{0}\) 时，它并非总是正的。对于这样的 \(y\) ，我们有 \(y \geq  0,y \neq  \mathbf{0}\) ，并且对于任何 \(\alpha  \geq  0\) 有 \({y\alpha } \in  Q\) ，这表明 \(Q\) 是无界的。反之，假设对任何 \(y\) 的最优反应收益 \(u\) 总是正的。因为
\end{tcolorbox}

\begin{center}
\includegraphics[max width=0.9\textwidth]{images/bo_d2e5cmbef24c73b03tn0_259_340_209_1254_531_0.jpg}
\end{center}
\hspace*{3em} 

Figure 9.17 The polytopes \(P\) and \(Q\) for the example (9.12).

\begin{tcolorbox}
图 9.17 示例 (9.12) 中的有界多面体(polytope) \(P\) 和 \(Q\) 。
\end{tcolorbox}

\(Y\) is compact and \(\bar{Q}\) is closed, the minimum \({u}^{\prime }\) of \(\{ u \mid  \left( {y,u}\right)  \in  \bar{Q}\}\) exists, \({u}^{\prime } > 0\) , and \(u \geq  {u}^{\prime }\) for all(y, u)in \(\bar{Q}\) . Then the map

\begin{tcolorbox}
由于 \(Y\) 是紧致的且 \(\bar{Q}\) 是闭集， \(\{ u \mid  \left( {y,u}\right)  \in  \bar{Q}\}\) 的最小值 \({u}^{\prime }\) 存在， \({u}^{\prime } > 0\) ，并且对于 \(\bar{Q}\) 中的所有 (y, u) 有 \(u \geq  {u}^{\prime }\) 。那么映射
\end{tcolorbox}

\[
\bar{Q} \rightarrow  Q - \{ \mathbf{0}\} ,\;\left( {y,u}\right)  \mapsto  y\frac{1}{u} \tag{9.26}
\]

is a bijection with inverse \(z \mapsto  \left( {z\frac{1}{{1}^{\top }z},\frac{1}{{1}^{\top }z}}\right)\) for \(z \in  Q - \{ \mathbf{0}\}\) . Here, \(\frac{1}{{1}^{\top }z} \geq  {u}^{\prime }\) and thus \({\mathbf{1}}^{\top }z \leq  \frac{1}{{u}^{\prime }}\) , where \({\mathbf{1}}^{\top }z = \mathop{\sum }\limits_{{j = 1}}^{n}\left| {z}_{j}\right|\) because \(z \geq  \mathbf{0}\) , which proves that \(Q\) is bounded and therefore a polytope.

\begin{tcolorbox}
是一个双射，其逆映射为对于 \(z \in  Q - \{ \mathbf{0}\}\) 的 \(z \mapsto  \left( {z\frac{1}{{1}^{\top }z},\frac{1}{{1}^{\top }z}}\right)\) 。这里， \(\frac{1}{{1}^{\top }z} \geq  {u}^{\prime }\) 且因此 \({\mathbf{1}}^{\top }z \leq  \frac{1}{{u}^{\prime }}\) ，其中 \({\mathbf{1}}^{\top }z = \mathop{\sum }\limits_{{j = 1}}^{n}\left| {z}_{j}\right|\) 是因为 \(z \geq  \mathbf{0}\) ，这证明了 \(Q\) 是有界的，从而是一个有界多面体(polytope)。
\end{tcolorbox}

As a sufficient condition that \(v > 0\) and \(u > 0\) hold for any(x, v)in \(\bar{P}\) and(y, u) in \(\bar{Q}\) , we assume that

\begin{tcolorbox}
作为对于 \(\bar{P}\) 中的任何 (x, v) 和 \(\bar{Q}\) 中的任何 (y, u) 使 \(v > 0\) 和 \(u > 0\) 成立的一个充分条件，我们假设
\end{tcolorbox}

\(A\) and \({B}^{\top }\) are nonnegative and have no zero column

\begin{tcolorbox}
\(A\) 和 \({B}^{\top }\) 是非负的且没有零列
\end{tcolorbox}

(9.27)

(because then \({B}^{\top }x\) and \({Ay}\) are nonnegative and nonzero for any \(x \in  X\) and \(y \in  Y\) ). We could simply assume \(A > 0\) and \(B > 0\) , but it is useful to admit zero entries in a payoff matrix, like in our examples. Note that condition (9.27) is not necessary for positive best-response payoffs (which is still the case, for example, if the lower left entry of \(A\) in (9.1) is negative, as Figure 9.1 shows; see also Exercise 9.5). By adding a suitable positive constant to all payoffs of a player, which preserves the preferences of that player (and where the original payoff is recovered by simply substracting the constant), we can assume (9.27) without loss of generality.

\begin{tcolorbox}
(因为这样对于任何 \(x \in  X\) 和 \(y \in  Y\) ， \({B}^{\top }x\) 和 \({Ay}\) 是非负且非零的)。我们可以简单地假设 \(A > 0\) 和 \(B > 0\) ，但允许收益矩阵中有零元素是有用的，就像我们的示例中那样。注意，条件 (9.27) 对于正的最优反应收益并非必要条件(例如，如图 9.1 所示，如果 (9.1) 中 \(A\) 的左下角元素为负，情况仍然如此；另见练习 9.5)。通过给一个参与者的所有收益加上一个合适的正常数，这会保持该参与者的偏好(并且通过简单地减去这个常数就可以恢复原始收益)，我们可以不失一般性地假设 (9.27)。
\end{tcolorbox}

With positive best-response payoffs, the polytope \(P\) (apart from its vertex0) is obtained from \(\bar{P}\) by dividing each inequality \(\mathop{\sum }\limits_{{i = 1}}^{m}{b}_{ij}{x}_{i} \leq  v\) by \(v\) , which gives \(\mathop{\sum }\limits_{{i = 1}}^{m}{b}_{ij}\left( {{x}_{i}/v}\right)  \leq  1\) , and then treating \({x}_{i}/v\) as a new variable that is again called \({x}_{i}\) in \(P\) . Similarly, \(Q - \{ \mathbf{0}\}\) is obtained from \(\bar{Q}\) by dividing each inequality in \({Ay} \leq  \mathbf{1}u\) by \(u\) . In effect, we have normalized the expected payoffs to be 1, and dropped the conditions \({\mathbf{1}}^{\top }x = 1\) and \({\mathbf{1}}^{\top }y = 1\) . Conversely, nonzero vectors \(x \in  P\) and \(y \in  Q\) are multiplied by \(v = \frac{1}{{1}^{\top }x}\) and \(u = \frac{1}{{1}^{\top }y}\) to turn them into probability vectors. The scaling factors \(v\) and \(u\) are the expected payoffs to the other player.

\begin{tcolorbox}
在最佳响应效用为正的情况下，有界多面体 \(P\) (除了其顶点 0)是通过将 \(\bar{P}\) 中的每个不等式 \(\mathop{\sum }\limits_{{i = 1}}^{m}{b}_{ij}{x}_{i} \leq  v\) 除以 \(v\) 得到的，这会得到 \(\mathop{\sum }\limits_{{i = 1}}^{m}{b}_{ij}\left( {{x}_{i}/v}\right)  \leq  1\) ，然后将 \({x}_{i}/v\) 视为一个新变量，在 \(P\) 中再次称为 \({x}_{i}\) 。类似地， \(Q - \{ \mathbf{0}\}\) 是通过将 \(\bar{Q}\) 中 \({Ay} \leq  \mathbf{1}u\) 的每个不等式除以 \(u\) 得到的。实际上，我们已将期望效用归一化为 1，并去掉了条件 \({\mathbf{1}}^{\top }x = 1\) 和 \({\mathbf{1}}^{\top }y = 1\) 。相反，非零向量 \(x \in  P\) 和 \(y \in  Q\) 分别乘以 \(v = \frac{1}{{1}^{\top }x}\) 和 \(u = \frac{1}{{1}^{\top }y}\) 以将它们转换为概率向量。缩放因子 \(v\) 和 \(u\) 是另一个参与者的期望效用。
\end{tcolorbox}

\begin{center}
\includegraphics[max width=0.9\textwidth]{images/bo_d2e5cmbef24c73b03tn0_260_200_255_1260_713_0.jpg}
\end{center}
\hspace*{3em} 

Figure 9.18 The projective transformation in (9.26), shown on the left for a single component \({y}_{j}\) of \(y\) , on the right for \(\bar{Q}\) as in Figure 9.1 and \(Q\) in (9.25).

\begin{tcolorbox}
图 9.18 展示了 (9.26) 中的射影变换，左侧展示了 \(y\) 的单个分量 \({y}_{j}\) 的情况，右侧展示了如图 9.1 中的 \(\bar{Q}\) 以及 (9.25) 中的 \(Q\) 的情况。
\end{tcolorbox}

Figure 9.18 shows a geometric interpretation of this division by the expected payoff, that is, the bijection (9.26), as a so-called "projective transformation". The left diagram shows the pair \(\left( {{y}_{j},u}\right)\) as part of(y, u)in \(\bar{Q}\) for any component \({y}_{j}\) of \(y\) . The line that connects this pair to(0,0)contains the point \(\left( {{y}_{j}/u,1}\right)\) . Hence, \(\left( {Q-\{ \mathbf{0}\} }\right)  \times  \{ 1\}\) is the intersection of the lines that connect any(y, u)in \(\bar{Q}\) with (0,0)in \({\mathbb{R}}^{n} \times  \mathbb{R}\) with the "horizontal" hyperplane \({\mathbb{R}}^{n} \times  \{ 1\}\) . The vertex0of \(Q\) does not arise as such a projection, but corresponds to \(\bar{Q}\) "at infinity".

\begin{tcolorbox}
图 9.18 展示了这种除以期望效用操作的几何解释，即双射 (9.26)，也就是所谓的“射影变换”。左图展示了对于 \(y\) 的任何分量 \({y}_{j}\) ， \(\left( {{y}_{j},u}\right)\) 作为 \(\bar{Q}\) 中 (y, u) 的一部分。连接该对与 (0, 0) 的直线包含点 \(\left( {{y}_{j}/u,1}\right)\) 。因此， \(\left( {Q-\{ \mathbf{0}\} }\right)  \times  \{ 1\}\) 是连接 \(\bar{Q}\) 中任何 (y, u) 与 \({\mathbb{R}}^{n} \times  \mathbb{R}\) 中 (0, 0) 的直线与“水平”超平面 \({\mathbb{R}}^{n} \times  \{ 1\}\) 的交点。 \(Q\) 的顶点 0 并非这样的投影产生的，而是对应于 \(\bar{Q}\) “在无穷远处”。
\end{tcolorbox}

The maps \(\left( {x,v}\right)  \mapsto  x\frac{1}{v}\) from \(\bar{P}\) to \(P\) and \(\left( {y,u}\right)  \mapsto  y\frac{1}{u}\) from \(\bar{Q}\) to \(Q\) preserve inequalities and when they are tight. Hence, we can label the inequalities in \(P\) and \(Q\) with the players’ pure strategies \(1,\ldots ,m + n\) , which is why we have written them in the order \(x \geq  \mathbf{0},{B}^{\top }x \leq  \mathbf{1}v\) and \({Ay} \leq  \mathbf{1}v,y \geq  \mathbf{0}\) in (9.24). Consequently, we obtain one more time the characterization of equilibria as completely labeled pairs of points, this time in \(P \times  Q - \{ \left( {\mathbf{0},\mathbf{0}}\right) \}\) .

\begin{tcolorbox}
从 \(\bar{P}\) 到 \(P\) 的映射 \(\left( {x,v}\right)  \mapsto  x\frac{1}{v}\) 以及从 \(\bar{Q}\) 到 \(Q\) 的映射 \(\left( {y,u}\right)  \mapsto  y\frac{1}{u}\) 保持不等式及其紧性。因此，我们可以用参与者的纯策略 \(1,\ldots ,m + n\) 为 \(P\) 和 \(Q\) 中的不等式添加标签，这就是为什么我们在 (9.24) 中按照 \(x \geq  \mathbf{0},{B}^{\top }x \leq  \mathbf{1}v\) 和 \({Ay} \leq  \mathbf{1}v,y \geq  \mathbf{0}\) 的顺序书写它们。因此，我们再次得到了将均衡刻画为完全标记的点对的结果，这次是在 \(P \times  Q - \{ \left( {\mathbf{0},\mathbf{0}}\right) \}\) 中。
\end{tcolorbox}

Proposition 9.10. Let(A, B)be an \(m \times  n\) bimatrix game so that (9.27) holds, and consider \(P\) and \(Q\) in (9.24) with their \(m + n\) inequalities labeled \(1,\ldots ,m + n\) , so that any \(x \in  P\) and \(y \in  Q\) have the labels of the tight inequalities. Then \(\left( {\bar{x},\bar{y}}\right)\) is an equilibrium of(A, B) with equilibrium payoffs \(u\) and \(v\) to players \(I\) and \({II}\) if and only if \(\left( {x,y}\right)  \in  P \times  Q - \{ \left( {\mathbf{0},\mathbf{0}}\right) \}\) and(x, y)is completely labeled, \(x = \bar{x}\frac{1}{v},y = \bar{y}\frac{1}{u},u = \frac{1}{{1}^{\top }y}\) , and \(v = \frac{1}{{1}^{\top }x}\) .

\begin{tcolorbox}
命题9.10。设(A, B)为一个双矩阵博弈(bimatrix game)，使得(9.27)式成立，并考虑(9.24)式中的 \(P\) 和 \(Q\) ，它们的紧不等式(tight inequality)标记为 \(1,\ldots ,m + n\) ，从而任意的 \(x \in  P\) 和 \(y \in  Q\) 具有紧不等式的标签(label)。那么 \(\left( {\bar{x},\bar{y}}\right)\) 是(A, B)的一个均衡(equilibrium)，玩家 \(I\) 和 \({II}\) 的均衡收益分别为 \(u\) 和 \(v\) ，当且仅当 \(\left( {x,y}\right)  \in  P \times  Q - \{ \left( {\mathbf{0},\mathbf{0}}\right) \}\) 且(x, y)是完全标记的(completely labeled)， \(x = \bar{x}\frac{1}{v},y = \bar{y}\frac{1}{u},u = \frac{1}{{1}^{\top }y}\) ，以及 \(v = \frac{1}{{1}^{\top }x}\) 。
\end{tcolorbox}

Proof. Condition (9.27) implies \(u > 0\) and \(v > 0\) for all \(\left( {\bar{x},v}\right)  \in  \bar{P}\) and \(\left( {\bar{y},u}\right)  \in  \bar{Q}\) , so that the maps \(\left( {\bar{x},v}\right)  \mapsto  x = \bar{x}\frac{1}{v} \in  P\) and \(\left( {\bar{y},u}\right)  \mapsto  y = \bar{y}\frac{1}{u} \in  Q\) are well defined, and they preserve inequalities and when these inequalities are tight, so that the labels of \(\bar{x}\) are the labels of \(x\) and the labels of \(\bar{y}\) are the labels of \(y\) . Hence, if \(\left( {\bar{x},\bar{y}}\right)\) is an equilibrium, then \(\left( {\bar{x},\bar{y}}\right)\) and therefore(x, y)is completely labeled by Corollary 9.8.

\begin{tcolorbox}
证明。条件(9.27)意味着对于所有的 \(\left( {\bar{x},v}\right)  \in  \bar{P}\) 和 \(\left( {\bar{y},u}\right)  \in  \bar{Q}\) ，有 \(u > 0\) 和 \(v > 0\) ，因此映射 \(\left( {\bar{x},v}\right)  \mapsto  x = \bar{x}\frac{1}{v} \in  P\) 和 \(\left( {\bar{y},u}\right)  \mapsto  y = \bar{y}\frac{1}{u} \in  Q\) 是良定义的，并且它们保持不等式关系，当这些不等式是紧不等式时， \(\bar{x}\) 的标签就是 \(x\) 的标签， \(\bar{y}\) 的标签就是 \(y\) 的标签。因此，如果 \(\left( {\bar{x},\bar{y}}\right)\) 是一个均衡，那么根据推论9.8， \(\left( {\bar{x},\bar{y}}\right)\) 且因此(x, y)是完全标记的。
\end{tcolorbox}

Conversely, suppose \(\left( {x,y}\right)  \in  P \times  Q\) is completely labeled and \(\left( {x,y}\right)  \neq  \left( {\mathbf{0},\mathbf{0}}\right)\) . If \(x = \mathbf{0}\) then \(y \neq  \mathbf{0}\) , so that \({y}_{j} \neq  0\) for some \(j \in  \{ 1,\ldots ,n\}\) and \(y\) does not have label \(m + j\) ; this means that(x, y)is not completely labeled because \(x\) has only the labels \(1,\ldots ,m\) , a contradiction. Hence, \(x \neq  \mathbf{0}\) , and similarly \(y \neq  \mathbf{0}\) . By the same reasoning, \(x\) has some label \(m + j\) , which by (9.23) means that at least one inequality in \({B}^{\top }x \leq  \mathbf{1}\) is tight, so that \(v = \frac{1}{{1}^{\top }x}\) is the best-response payoff to \(\bar{x} = {xv}\) for \(\left( {\bar{x},v}\right)  \in  \bar{P}\) . Similarly, \(y\) in \(P\) has some label \(i \in  \{ 1,\ldots ,m\}\) , which means \({\left( Ay\right) }_{i} = 1\) , so that \(u = \frac{1}{{1}^{\mathrm{\;T}}y}\) is the best-response payoff to \(\bar{y} = {yu}\) for \(\left( {\bar{y},u}\right)  \in  \bar{Q}\) . This proves the claim.

\begin{tcolorbox}
反之，假设 \(\left( {x,y}\right)  \in  P \times  Q\) 是完全标记的且 \(\left( {x,y}\right)  \neq  \left( {\mathbf{0},\mathbf{0}}\right)\) 。如果 \(x = \mathbf{0}\) ，那么 \(y \neq  \mathbf{0}\) ，使得对于某个 \(j \in  \{ 1,\ldots ,n\}\) 和 \(y\) ， \({y}_{j} \neq  0\) 没有标签 \(m + j\) ；这意味着(x, y)不是完全标记的，因为 \(x\) 只有标签 \(1,\ldots ,m\) ，这是一个矛盾。因此， \(x \neq  \mathbf{0}\) ，并且类似地有 \(y \neq  \mathbf{0}\) 。通过同样的推理， \(x\) 有某个标签 \(m + j\) ，根据(9.23)式，这意味着 \({B}^{\top }x \leq  \mathbf{1}\) 中至少有一个不等式是紧不等式，从而 \(v = \frac{1}{{1}^{\top }x}\) 是 \(\left( {\bar{x},v}\right)  \in  \bar{P}\) 对 \(\bar{x} = {xv}\) 的最优反应收益。类似地， \(P\) 中的 \(y\) 有某个标签 \(i \in  \{ 1,\ldots ,m\}\) ，这意味着 \({\left( Ay\right) }_{i} = 1\) ，从而 \(u = \frac{1}{{1}^{\mathrm{\;T}}y}\) 是 \(\left( {\bar{y},u}\right)  \in  \bar{Q}\) 对 \(\bar{y} = {yu}\) 的最优反应收益。这就证明了该命题。
\end{tcolorbox}

The last paragraph in the preceding proof shows that we do not need to worry about the possibilities \({B}^{\top }x < \mathbf{1}\) or \({Ay} < \mathbf{1}\) in \(P\) or \(Q\) (that is, no tight inequalities) when considering completely labeled pairs(x, y)in \(P \times  Q\) . There is one obvious exception to this, namely the pair(0,0)in \(P \times  Q\) , which takes the role of the artificial equilibrium for the Lemke-Howson algorithm in a natural way; it does not have to be added like in \(\widetilde{X} \times  \widetilde{Y}\) . For our example (9.1), Figure 9.19 shows the Lemke-Howson paths on \(P \times  Q\) for missing label 2. They traverse the almost completely labeled vertex pairs of \(P\) and \(Q\) , which we identify by their sets of labels exactly as in (9.7) and (9.8).

\begin{tcolorbox}
上述证明的最后一段表明，当考虑 \(P \times  Q\) 中的完全标记对 (x, y) 时，我们无需担心 \(P\) 或 \(Q\) 中出现 \({B}^{\top }x < \mathbf{1}\) 或 \({Ay} < \mathbf{1}\) 的情况(即不存在紧不等式)。有一个明显的例外，即 \(P \times  Q\) 中的对 (0, 0)，它自然地充当了莱姆基 - 豪森算法的人工均衡点；不像在 \(\widetilde{X} \times  \widetilde{Y}\) 中那样需要额外添加。对于我们的示例 (9.1)，图 9.19 展示了 \(P \times  Q\) 上缺失标签 2 的莱姆基 - 豪森路径。它们遍历了 \(P\) 和 \(Q\) 的几乎完全标记的顶点对，我们通过它们的标签集来精确识别这些顶点对，就像在 (9.7) 和 (9.8) 中那样。
\end{tcolorbox}

\begin{center}
\includegraphics[max width=0.9\textwidth]{images/bo_d2e5cmbef24c73b03tn0_261_338_1222_1253_375_0.jpg}
\end{center}
\hspace*{3em} 

Figure 9.19 The Lemke-Howson paths from Figure 9.5 for missing label 2 on the polytopes \(P\) and \(Q\) .

\begin{tcolorbox}
图 9.19 图 9.5 中在多面体 \(P\) 和 \(Q\) 上缺失标签 2 的莱姆基 - 豪森路径。
\end{tcolorbox}

\section*{9.7 Complementary Pivoting}

\begin{tcolorbox}
\section*{9.7 互补换基法}
\end{tcolorbox}

We now describe how the Lemke-Howson algorithm works algebraically and in any dimension. As always, we consider an \(m \times  n\) bimatrix game(A, B). Let \(C\) be the \(d \times  d\) matrix defined by

\begin{tcolorbox}
现在我们从代数角度描述莱姆基 - 豪森算法在任意维度下的工作原理。和往常一样，我们考虑一个 \(m \times  n\) 双矩阵博弈 (A, B)。设 \(C\) 是由下式定义的 \(d \times  d\) 矩阵
\end{tcolorbox}

\[
d = m + n,\;C = \left( \begin{matrix} 0 & A \\  {B}^{\top } & 0 \end{matrix}\right) . \tag{9.28}
\]

In agreement with (9.27) we assume, without loss of generality, that \(C\) is nonnegative and has no zero column. Then \(S = P \times  Q\) is a polytope, which can be written as

\begin{tcolorbox}
与 (9.27) 一致，不失一般性，我们假设 \(C\) 是非负的且没有零列。那么 \(S = P \times  Q\) 是一个有界多面体，可以写成
\end{tcolorbox}

\[
S = \left\{  {z \in  {\mathbb{R}}^{d} \mid  z \geq  \mathbf{0},{Cz} \leq  \mathbf{1}}\right\}   \tag{9.29}
\]

where \(z \in  S = P \times  Q\) means

\begin{tcolorbox}
其中 \(z \in  S = P \times  Q\) 表示
\end{tcolorbox}

\[
z = \left( {x,y}\right) ,\;{z}_{i} = {x}_{i}\;\left( {1 \leq  i \leq  m}\right) ,\;{z}_{m + j} = {y}_{j}\;\left( {1 \leq  j \leq  n}\right) . \tag{9.30}
\]

We label both the \(d\) inequalities \(z \geq  \mathbf{0}\) and the \(d\) inequalities \({Cz} \leq  \mathbf{1}\) in (9.29) with labels \(1,\ldots ,d\) . As usual, any \(z\) in \(S\) gets the labels of the tight inequalities. The \(d\) labels correspond to the players’ pure strategies as in (9.2) and (9.23). An equilibrium \(\left( {\bar{x},\bar{y}}\right)\) of(A, B)then corresponds to a completely labeled point \(z = \left( {x,y}\right)\) of \(S - \{ \mathbf{0}\}\) (that is, \(z\) has all labels \(1,\ldots ,d\) ), via the normalizations \(\bar{x} = x\frac{1}{{\mathbf{1}}^{\top }x}\) and \(\bar{y} = y\frac{1}{{\mathbf{1}}^{\top }y}\) in Proposition 9.10.

\begin{tcolorbox}
我们用标签 \(1,\ldots ,d\) 标记 (9.29) 中的 \(d\) 个不等式 \(z \geq  \mathbf{0}\) 和 \(d\) 个不等式 \({Cz} \leq  \mathbf{1}\) 。和往常一样， \(S\) 中的任何 \(z\) 都获得紧不等式的标签。 \(d\) 个标签对应于 (9.2) 和 (9.23) 中玩家的纯策略。通过命题 9.10 中的归一化 \(\bar{x} = x\frac{1}{{\mathbf{1}}^{\top }x}\) 和 \(\bar{y} = y\frac{1}{{\mathbf{1}}^{\top }y}\) ，(A, B) 的一个均衡 \(\left( {\bar{x},\bar{y}}\right)\) 对应于 \(S - \{ \mathbf{0}\}\) 的一个完全标记点 \(z = \left( {x,y}\right)\) (即 \(z\) 拥有所有标签 \(1,\ldots ,d\) )。
\end{tcolorbox}

For an algebraic treatment, we write the inequalities \({Cz} \leq  \mathbf{1}\) in (9.29) with nonnegative "slack variables" \(w \in  {\mathbb{R}}^{d}\) as the equality and nonnegativity constraints \({Cz} + w = \mathbf{1},w \geq  \mathbf{0}\) . Then \(z \in  S\) has label \(i \in  \{ 1,\ldots ,d\}\) if and only if \({z}_{i} = 0\) or \({w}_{i} = 0\) , which can be written as \({z}_{i}{w}_{i} = 0\) . If this holds for all \(i = 1,\ldots ,d\) , then the vectors \(z\) and \(w\) are called \(\mathit{c}{omplementary}\) . Then the problem of finding a completely labeled point \(z\) of \(S\) means to find \(z\) and \(w\) in \({\mathbb{R}}^{d}\) so that

\begin{tcolorbox}
为了进行代数处理，我们用非负的“松弛变量” \(w \in  {\mathbb{R}}^{d}\) 将 (9.29) 中的不等式 \({Cz} \leq  \mathbf{1}\) 写成等式和非负性约束 \({Cz} + w = \mathbf{1},w \geq  \mathbf{0}\) 。那么 \(z \in  S\) 具有标签 \(i \in  \{ 1,\ldots ,d\}\) 当且仅当 \({z}_{i} = 0\) 或 \({w}_{i} = 0\) ，这可以写成 \({z}_{i}{w}_{i} = 0\) 。如果对于所有的 \(i = 1,\ldots ,d\) 都成立，那么向量 \(z\) 和 \(w\) 被称为 \(\mathit{c}{omplementary}\) 。那么，在 \(S\) 中找到一个完全标记点 \(z\) 的问题意味着在 \({\mathbb{R}}^{d}\) 中找到 \(z\) 和 \(w\) ，使得
\end{tcolorbox}

\[
{Cz} + w = \mathbf{1},\;z \geq  \mathbf{0},\;w \geq  \mathbf{0},\;{z}_{i}{w}_{i} = 0\;\left( {1 \leq  i \leq  d}\right) . \tag{9.31}
\]

In summary, if \(z \in  S\) then \(w\) is defined by \(w = \mathbf{1} - {Cz}\) , we have \(z \geq  \mathbf{0}\) and \(w \geq  \mathbf{0}\) , and \(z\) is completely labeled if and only if \(z\) and \(w\) are complementary.

\begin{tcolorbox}
总之，如果 \(z \in  S\) ，那么 \(w\) 由 \(w = \mathbf{1} - {Cz}\) 定义，我们有 \(z \geq  \mathbf{0}\) 和 \(w \geq  \mathbf{0}\) ，并且当且仅当 \(z\) 和 \(w\) 互补时， \(z\) 是完全标记的。
\end{tcolorbox}

Recall that \(z\) has label \(i\) if \({z}_{i} = 0\) or \({w}_{i} = 0\) . A missing label therefore means some \(\ell  \in  \{ 1,\ldots ,d\}\) so that \({z}_{\ell } > 0\) and \({w}_{\ell } > 0\) are allowed to hold. If with this exception all other conditions in (9.31) hold, then \(z\) and \(w\) are called almost complementary, which means that \(z\) is an almost completely labeled point of \(S\) , that is, \(z\) has all labels except possibly \(\ell\) .

\begin{tcolorbox}
回顾一下，如果 \({z}_{i} = 0\) 或 \({w}_{i} = 0\) ，则 \(z\) 有标签 \(i\) 。因此，缺失标签意味着存在某个 \(\ell  \in  \{ 1,\ldots ,d\}\) ，使得 \({z}_{\ell } > 0\) 和 \({w}_{\ell } > 0\) 可以成立。如果除了这个例外，(9.31) 中的所有其他条件都成立，那么 \(z\) 和 \(w\) 被称为几乎互补，这意味着 \(z\) 是 \(S\) 的一个几乎完全标记的点，即 \(z\) 拥有除可能的 \(\ell\) 之外的所有标签。
\end{tcolorbox}

The Lemke-Howson algorithm starts with the trivial solution \(z = \mathbf{0},w = \mathbf{1}\) to (9.31), which defines the artificial equilibrium. For our example (9.1), the system \(w = 1 - {Cz}\) states

\begin{tcolorbox}
莱姆基 - 豪森算法从 (9.31) 的平凡解 \(z = \mathbf{0},w = \mathbf{1}\) 开始，该解定义了人工均衡。对于我们的例子 (9.1)，系统 \(w = 1 - {Cz}\) 表述为
\end{tcolorbox}

\[
{w}_{1} = 1\; - 2{z}_{4}
\]

\[
{w}_{2} = 1\; - \;{z}_{4} - 2{z}_{5}
\]

\[
{w}_{3} = 1\; - 3{z}_{5} \tag{9.32}
\]

\[
{w}_{4} = 1 - 2{z}_{1}\; - 4{z}_{3}
\]

\[
{w}_{5} = 1\; - {z}_{2} - {z}_{3}.
\]

This is called a dictionary, which means that certain variables, here \({w}_{1},\ldots ,{w}_{5}\) , are expressed in dependence of the others, here \({z}_{1},\ldots ,{z}_{5}\) . The dependent variables on the left are called basic variables and the independent variables on the right are called nonbasic variables. The corresponding basic solution is obtained by setting all nonbasic variables to zero, so that the basic variables are equal to the constants in the dictionary. If these constants are all nonnegative, then the basic solution is called feasible, which is the case in (9.32) where the basic solution is \(z = \mathbf{0},w = \mathbf{1}\) .

\begin{tcolorbox}
这被称为字典，这意味着某些变量，这里是 \({w}_{1},\ldots ,{w}_{5}\) ，表示为依赖于其他变量，这里是 \({z}_{1},\ldots ,{z}_{5}\) 。左边的依赖变量称为基变量，右边的独立变量称为非基变量。通过将所有非基变量设为零来得到相应的基解，这样基变量就等于字典中的常数。如果这些常数都是非负的，那么该基解被称为可行解，在 (9.32) 中就是这种情况，其中基解是 \(z = \mathbf{0},w = \mathbf{1}\) 。
\end{tcolorbox}

The algorithm proceeds by the operation of pivoting that exchanges a nonbasic variable with a basic variable in such a way that the new dictionary defines the same linear relationship between its variables, and so that the corresponding basic solution remains feasible. In geometric terms, each basic feasible solution defines a vertex of the polytope \(S\) , and pivoting means moving along an edge of \(S\) from one vertex to an adjacent vertex, as in Figure 9.19, as we will show.

\begin{tcolorbox}
该算法通过换基操作进行，该操作以这样一种方式交换一个非基变量和一个基变量:新的字典定义了其变量之间相同的线性关系，并且相应的基解仍然是可行的。从几何角度来看，每个基本可行解定义了有界多面体 \(S\) 的一个顶点，换基操作意味着沿着 \(S\) 的一条边从一个顶点移动到相邻的顶点，如图 9.19 所示，我们将进行说明。
\end{tcolorbox}

A fixed missing label \(\ell  \in  \{ 1,\ldots ,d\}\) is chosen, here \(\ell  = 2\) , which determines the entering variable \({z}_{\ell }\) . In the current basic solution, all nonbasic variables are set to zero, but now the entering variable, which in (9.32) is \({z}_{2}\) , is allowed to assume positive values. The dictionary (9.32) describes how this changes the basic variables on the left, which is here only the variable \({w}_{5}\) because \({z}_{2}\) appears only in the last equation \({w}_{5} = 1 - {z}_{2} - {z}_{3}\) . The constraint \({w}_{5} \geq  0\) limits the increase of \({z}_{2}\) , because \({w}_{5}\) becomes zero when \({z}_{2} = 1\) . This determines \({w}_{5}\) as the leaving variable.

\begin{tcolorbox}
选择一个固定的缺失标签 \(\ell  \in  \{ 1,\ldots ,d\}\) ，这里是 \(\ell  = 2\) ，它确定了入基变量 \({z}_{\ell }\) 。在当前的基解中，所有非基变量都被设为零，但现在入基变量(在 (9.32) 中是 \({z}_{2}\) )被允许取正值。字典 (9.32) 描述了这如何改变左边的基变量，这里只有变量 \({w}_{5}\) ，因为 \({z}_{2}\) 只出现在最后一个方程 \({w}_{5} = 1 - {z}_{2} - {z}_{3}\) 中。约束 \({w}_{5} \geq  0\) 限制了 \({z}_{2}\) 的增加，因为当 \({z}_{2} = 1\) 时 \({w}_{5}\) 变为零。这就确定了 \({w}_{5}\) 为离基变量。
\end{tcolorbox}

The pivoting step is now to let \({z}_{2}\) enter the basis and \({w}_{5}\) leave the basis, which means rewriting the dictionary so that \({z}_{2}\) becomes basic and \({w}_{5}\) nonbasic. Here, this amounts to a simple exchange of \({w}_{5}\) with \({z}_{2}\) , which defines the new dictionary

\begin{tcolorbox}
现在换基步骤是让 \({z}_{2}\) 进入基， \({w}_{5}\) 离开基，这意味着重写字典，使得 \({z}_{2}\) 变为基变量， \({w}_{5}\) 变为非基变量。在这里，这相当于简单地将 \({w}_{5}\) 与 \({z}_{2}\) 交换，从而定义了新的字典
\end{tcolorbox}

\[
{w}_{1} = 1\; - 2{z}_{4}
\]

\[
{w}_{2} = 1\; - \;{z}_{4} - 2{z}_{5}
\]

\[
{w}_{3} = 1\; - 3{z}_{5} \tag{9.33}
\]

\[
{w}_{4} = 1 - 2{z}_{1}\; - 4{z}_{3}
\]

\[
{z}_{2} = 1\; - {w}_{5} - {z}_{3}.
\]

In the basic solution associated with this new dictionary, \({z}_{2} = 1\) and \({w}_{5} = 0\) . In Figure 9.19 this is the endpoint of first arrow in \(P\) at the vertex \(x = \left( {0,1,0}\right)\) , where label 2 has been dropped and label 5 is picked up (because \({w}_{5}\) has become nonbasic). In \(Q\) , we still have \(\left( {{y}_{1},{y}_{2}}\right)  = \left( {{z}_{4},{z}_{5}}\right)  = \left( {0,0}\right)\) .

\begin{tcolorbox}
在与这个新字典相关的基解中， \({z}_{2} = 1\) 和 \({w}_{5} = 0\) 。在图9.19中，这是 \(P\) 中从顶点 \(x = \left( {0,1,0}\right)\) 出发的第一个箭头的终点，其中标签2被去掉，标签5被选取(因为 \({w}_{5}\) 已变为非基变量)。在 \(Q\) 中，我们仍然有 \(\left( {{y}_{1},{y}_{2}}\right)  = \left( {{z}_{4},{z}_{5}}\right)  = \left( {0,0}\right)\) 。
\end{tcolorbox}

The dictionary (9.33) has two nonbasic variables \({z}_{5}\) and \({w}_{5}\) with the same subscript 5, which is the duplicate label. The complementary pivoting rule now states that because \({w}_{5}\) has just left the basis, its complement \({z}_{5}\) is chosen as the next entering variable.

\begin{tcolorbox}
字典(9.33)有两个非基变量 \({z}_{5}\) 和 \({w}_{5}\) ，它们的下标都是5，这就是重复标签。互补换基法规则现在规定，因为 \({w}_{5}\) 刚刚离开基，所以它的互补变量 \({z}_{5}\) 被选为下一个入基变量。
\end{tcolorbox}

That is, starting from the basic solution associated with (9.33), \({z}_{5}\) is allowed to become positive, with all other nonbasic variables staying at zero. Restricted to this increase of \({z}_{5}\) , the affected variables on the left-hand side are \({w}_{2} = 1 - 2{z}_{5}\) and \({w}_{3} = 1 - 3{z}_{5}\) , both of which should remain nonnegative. The stronger of these constraints is \({w}_{3} = 1 - 3{z}_{5} \geq  0\) , where \({w}_{3}\) becomes zero when \({z}_{5} = \frac{1}{3}\) (while \({w}_{2}\) is still positive). Hence, \({w}_{3}\) is the leaving variable. The pivoting step means now to rewrite the third equation (which currently has \({w}_{3}\) on the left) as \({z}_{5} = \frac{1}{3} - \frac{1}{3}{w}_{3}\) and to substitute this expression for \({z}_{5}\) where it appears elsewhere on the right. The only substitution is here \({w}_{2} = 1 - {z}_{4} - 2{z}_{5} = 1 - {z}_{4} - 2\left( {\frac{1}{3} - \frac{1}{3}{w}_{3}}\right)\) , so that the resulting new dictionary is

\begin{tcolorbox}
也就是说，从与(9.33)相关的基解开始，允许 \({z}_{5}\) 变为正数，而所有其他非基变量保持为零。仅考虑 \({z}_{5}\) 的增加，等式左边受影响的变量是 \({w}_{2} = 1 - 2{z}_{5}\) 和 \({w}_{3} = 1 - 3{z}_{5}\) ，它们都应保持非负。这些约束中更强的是 \({w}_{3} = 1 - 3{z}_{5} \geq  0\) ，当 \({z}_{5} = \frac{1}{3}\) 时 \({w}_{3}\) 变为零(而 \({w}_{2}\) 仍然为正)。因此， \({w}_{3}\) 是离基变量。换基操作步骤现在意味着将第三个方程(目前左边是 \({w}_{3}\) )重写为 \({z}_{5} = \frac{1}{3} - \frac{1}{3}{w}_{3}\) ，并将这个表达式代入 \({z}_{5}\) 在右边其他地方出现的位置。唯一的代入是这里的 \({w}_{2} = 1 - {z}_{4} - 2{z}_{5} = 1 - {z}_{4} - 2\left( {\frac{1}{3} - \frac{1}{3}{w}_{3}}\right)\) ，所以得到的新字典是
\end{tcolorbox}

\[
{w}_{1} = 1\; - 2{z}_{4}
\]

\[
{w}_{2} = \frac{1}{3}\; - \;{z}_{4} + \frac{2}{3}{w}_{3}
\]

\[
{z}_{5} = \frac{1}{3}\; - \frac{1}{3}{w}_{3} \tag{9.34}
\]

\[
{w}_{4} = 1 - 2{z}_{1}\; - 4{z}_{3}
\]

\[
{z}_{2} = 1\; - {w}_{5} - {z}_{3}.
\]

In the preceding pivoting step, \({z}_{5}\) entered and \({w}_{3}\) left the basis. In Figure 9.19, label 5 was dropped and label 3 picked up in \(Q\) , with the new basic solution at the vertex \(y = \left( {0,\frac{1}{3}}\right)\) of \(Q\) .

\begin{tcolorbox}
在前面的换基操作步骤中， \({z}_{5}\) 入基， \({w}_{3}\) 离基。在图9.19中，在 \(Q\) 中去掉了标签5，选取了标签3，新的基解位于 \(Q\) 的顶点 \(y = \left( {0,\frac{1}{3}}\right)\) 处。
\end{tcolorbox}

By the complementary pivoting rule, the next entering variable in (9.34) is \({z}_{3}\) as the complement of \({w}_{3}\) . The increase of \({z}_{3}\) creates the two constraints \({w}_{4} = 1 - 4{z}_{3} \geq  0\) and \({z}_{2} = 1 - {z}_{3} \geq  0\) , where the former is stronger, with \({w}_{4}\) becoming zero when \({z}_{4} = \frac{1}{4}\) . Hence, the leaving variable is \({w}_{4}\) , so that we rewrite the fourth equation as \({z}_{3} = \frac{1}{4} - \frac{1}{2}{z}_{1} - \frac{1}{4}{w}_{4}\) and substitute this expression for \({z}_{3}\) :

\begin{tcolorbox}
根据互补换基法规则，(9.34)中的下一个入基变量是 \({z}_{3}\) ，它是 \({w}_{3}\) 的互补变量。 \({z}_{3}\) 的增加产生了两个约束 \({w}_{4} = 1 - 4{z}_{3} \geq  0\) 和 \({z}_{2} = 1 - {z}_{3} \geq  0\) ，其中前者更强，当 \({z}_{4} = \frac{1}{4}\) 时 \({w}_{4}\) 变为零。因此，离基变量是 \({w}_{4}\) ，所以我们将第四个方程重写为 \({z}_{3} = \frac{1}{4} - \frac{1}{2}{z}_{1} - \frac{1}{4}{w}_{4}\) 并代入 \({z}_{3}\) 的表达式:
\end{tcolorbox}

\[
{w}_{1} = 1\; - 2{z}_{4}
\]

\[
{w}_{2} = \frac{1}{3}\; - \;{z}_{4} + \frac{2}{3}{w}_{3}
\]

\[
{z}_{5} = \frac{1}{3}\; - \frac{1}{3}{w}_{3} \tag{9.35}
\]

\[
{z}_{3} = \frac{1}{4} - \frac{1}{2}{z}_{1}\; - \frac{1}{4}{w}_{4}
\]

\[
{z}_{2} = \frac{3}{4} + \frac{1}{2}{z}_{1} - {w}_{5} + \frac{1}{4}{w}_{4}.
\]

In the preceding pivoting step, \({z}_{3}\) entered and \({w}_{4}\) left the basis. This means that in \(P\) , label 3 was dropped and label 4 is picked up, and \(x\) moves to the new vertex \(\left( {0,\frac{3}{4},\frac{1}{4}}\right)\) of \(P\) , which is already the white square in \(P\) in Figure 9.19.

\begin{tcolorbox}
在前面的换基操作步骤中， \({z}_{3}\) 入基， \({w}_{4}\) 离基。这意味着在 \(P\) 中，去掉了标签3，选取了标签4，并且 \(x\) 移动到 \(P\) 的新顶点 \(\left( {0,\frac{3}{4},\frac{1}{4}}\right)\) ，这已经是图9.19中 \(P\) 里的白色方块。
\end{tcolorbox}

As the complement of \({w}_{4}\) that has left the basis, the new entering variable in (9.35) is \({z}_{4}\) . Of the two constraints \({w}_{1} = 1 - 2{z}_{4} \geq  0\) and \({w}_{2} = \frac{1}{3} - {z}_{4} \geq  0\) , the second is stronger, where \({w}_{2}\) becomes zero when \({z}_{4} = \frac{1}{3}\) . We rewrite the second equation as \({z}_{4} = \frac{1}{3} - {w}_{2} + \frac{2}{3}{w}_{3}\) and substitute this expression for \({z}_{4}\) :

\begin{tcolorbox}
作为已离基的 \({w}_{4}\) 的互补变量，(9.35) 中的新入基变量是 \({z}_{4}\) 。在两个约束条件 \({w}_{1} = 1 - 2{z}_{4} \geq  0\) 和 \({w}_{2} = \frac{1}{3} - {z}_{4} \geq  0\) 中，第二个更强，当 \({z}_{4} = \frac{1}{3}\) 时 \({w}_{2}\) 变为零。我们将第二个方程重写为 \({z}_{4} = \frac{1}{3} - {w}_{2} + \frac{2}{3}{w}_{3}\) 并将这个表达式代入 \({z}_{4}\) :
\end{tcolorbox}

\[
{w}_{1} = \frac{1}{3}\; + 2{w}_{2} - \frac{4}{3}{w}_{3}
\]

\[
{z}_{4} = \frac{1}{3}\; - \;{w}_{2} + \frac{2}{3}{w}_{3}
\]

\[
{z}_{5} = \frac{1}{3}\; - \frac{1}{3}{w}_{3} \tag{9.36}
\]

\[
{z}_{3} = \frac{1}{4} - \frac{1}{2}{z}_{1}\; - \frac{1}{4}{w}_{4}
\]

\[
{z}_{2} = \frac{3}{4} + \frac{1}{2}{z}_{1} - {w}_{5} + \frac{1}{4}{w}_{4}.
\]

Because \({z}_{4}\) entered and \({w}_{2}\) left the basis, this means that in \(Q\) label 4 was dropped and label 2 picked up, with the new vertex \(y = \left( {\frac{1}{3},\frac{1}{3}}\right)\) of \(Q\) , the white square in \(Q\) in Figure 9.19.

\begin{tcolorbox}
因为 \({z}_{4}\) 入基且 \({w}_{2}\) 离基，这意味着在 \(Q\) 中标签 4 被去掉，标签 2 被选取， \(Q\) 的新顶点为 \(y = \left( {\frac{1}{3},\frac{1}{3}}\right)\) ，即图 9.19 中 \(Q\) 里的白色方块。
\end{tcolorbox}

The algorithm terminates because the leaving variable has the missing label \(\ell  = 2\) as a subscript, so that there is now only one basic variable with that subscript and not two as in all the intermediate steps. The basic solution associated with the final dictionary (9.36) solves the system (9.31). Here, this solution is \(z = \left( {x,y}\right)\) with \(x = \left( {0,\frac{3}{4},\frac{1}{4}}\right)\) and \(y = \left( {\frac{1}{3},\frac{1}{3}}\right)\) . It defines the equilibrium \(\left( {\bar{x},\bar{y}}\right)\) with \(\bar{x} = x\) and payoff 1 to player II, and \(\bar{y} = \left( {\frac{1}{2},\frac{1}{2}}\right)\) with payoff \(\frac{3}{2}\) to player I.

\begin{tcolorbox}
算法终止，因为离基变量的下标是缺失标签 \(\ell  = 2\) ，所以现在只有一个基变量带有该下标，而不像所有中间步骤那样有两个。与最终字典 (9.36) 相关联的基解解决了方程组 (9.31)。在此，这个解是 \(z = \left( {x,y}\right)\) ，其中 \(x = \left( {0,\frac{3}{4},\frac{1}{4}}\right)\) 且 \(y = \left( {\frac{1}{3},\frac{1}{3}}\right)\) 。它定义了均衡 \(\left( {\bar{x},\bar{y}}\right)\) ，其中 \(\bar{x} = x\) ，玩家 II 的收益为 1，以及 \(\bar{y} = \left( {\frac{1}{2},\frac{1}{2}}\right)\) ，玩家 I 的收益为 \(\frac{3}{2}\) 。
\end{tcolorbox}

In the rest of this section we describe complementary pivoting in general. For the moment, we assume that the polytope \(S\) in (9.29) is nondegenerate, which means that no point in \(S\) has more than \(d\) tight inequalities; this can be seen to be equivalent to the nondegeneracy of the game(A, B).

\begin{tcolorbox}
在本节的其余部分，我们将一般性地描述互补换基法。目前，我们假设 (9.29) 中的有界多面体 \(S\) 是非退化的，这意味着 \(S\) 中没有点的紧不等式超过 \(d\) 个；可以看出这等价于博弈 (A, B) 的非退化性。
\end{tcolorbox}

With the \(d \times  d\) identity matrix \(I\) , we write the equations in (9.31) as \({Cz} + {Iw} = \mathbf{1}\) . This system has as variables the \({2d}\) components of \(z\) and \(w\) , the matrix \(\left\lbrack  {CI}\right\rbrack\) of coefficients, and a right-hand side 1 . Any \(d\) linearly independent columns of the matrix \(\left\lbrack  {CI}\right\rbrack\) define an invertible matrix \(D\) , called the basis matrix. The corresponding basic columns have associated variables \({z}_{i}\) or \({w}_{i}\) called basic variables; the other columns and variables are called nonbasic.

\begin{tcolorbox}
借助 \(d \times  d\) 单位矩阵 \(I\) ，我们将 (9.31) 中的方程写为 \({Cz} + {Iw} = \mathbf{1}\) 。这个方程组的变量是 \(z\) 和 \(w\) 的 \({2d}\) 个分量、系数矩阵 \(\left\lbrack  {CI}\right\rbrack\) 以及右侧常数 1。矩阵 \(\left\lbrack  {CI}\right\rbrack\) 的任意 \(d\) 个线性无关列定义了一个可逆矩阵 \(D\) ，称为基矩阵。对应的基列有关联的变量 \({z}_{i}\) 或 \({w}_{i}\) ，称为基变量；其他列和变量称为非基变量。
\end{tcolorbox}

Multiplying the system \({Cz} + {Iw} = \mathbf{1}\) with the inverse of the basis matrix \(D\) defines the equivalent system

\begin{tcolorbox}
用基矩阵 \(D\) 的逆乘以方程组 \({Cz} + {Iw} = \mathbf{1}\) 定义了等价方程组
\end{tcolorbox}

\[
{D}^{-1}{Cz} + {D}^{-1}w = {D}^{-1}\mathbf{1}\text{ . } \tag{9.37}
\]

This new system is a dictionary as in (9.32)-(9.36) above, except that we have kept both basic and nonbasic variables on the left-hand side of the equation. The reason is that in the system (9.37) the columns of the basic variables define the identity matrix, so that (9.37) shows how the basic variables depend on the nonbasic variables. When the nonbasic variables are set to zero, the basic variables are equal to \({D}^{-1}\mathbf{1}\) , which (together with the zero nonbasic variables) defines the basic solution that corresponds to the system (9.37). This basic solution \(z,w\) is called feasible if it is nonnegative, that is, \({D}^{-1}\mathbf{1} \geq  \mathbf{0}\) . By nondegeneracy, this means \({D}^{-1}\mathbf{1} > \mathbf{0}\) , that is, all basic variables have positive values.

\begin{tcolorbox}
这个新方程组与上述 (9.32) - (9.36) 中的字典形式类似，只是我们将基变量和非基变量都保留在了方程的左侧。原因是在方程组 (9.37) 中，基变量的列定义了单位矩阵，所以 (9.37) 展示了基变量如何依赖于非基变量。当非基变量设为零时，基变量等于 \({D}^{-1}\mathbf{1}\) ，这(连同为零的非基变量)定义了对应于方程组 (9.37) 的基解。如果这个基解 \(z,w\) 是非负的，即 \({D}^{-1}\mathbf{1} \geq  \mathbf{0}\) ，则称其为可行解。根据非退化性，这意味着 \({D}^{-1}\mathbf{1} > \mathbf{0}\) ，即所有基变量都取正值。
\end{tcolorbox}

Pivoting means to choose a nonbasic variable to enter the basis and a basic variable to leave the basis. Suppose the entering variable is \({z}_{j}\) (the process is exactly the same for some \({w}_{j}\) as entering variable). We assume that the current basic solution is feasible, and want to maintain a feasible solution while increasing \({z}_{j}\) from zero. This increase of \({z}_{j}\) changes the basic variables whose current value is equal to the vector \({D}^{-1}\mathbf{1}\) , which we denote by \(q\) . Suppose the vector \(v\) is the nonbasic column in (9.37) that corresponds to \({z}_{j}\) . With the increase of \({z}_{j}\) , the vector of basic variables is given by \(q - v{z}_{j}\) and should stay nonnegative, which means that \({q}_{i} - {v}_{i}{z}_{j} \geq  0\) for all rows \(i = 1,\ldots ,d\) . Because \(q \geq  \mathbf{0}\) holds already, we only need to consider rows \(i\) so that \({v}_{i} > 0\) . If \(v\) had no positive entry, we could increase \({z}_{j}\) indefinitely, which would define an unbounded "ray" in \(S\) , but \(S\) is bounded, so this cannot happen (complementary pivoting is also applied to general matrices \(C\) , and then this "ray termination" of the algorithm can happen).

\begin{tcolorbox}
换基操作(pivoting)是指选择一个非基变量(nonbasic variable)进入基，同时选择一个基变量(basic variable)离开基。假设入基变量(entering variable)是 \({z}_{j}\) (若以某个 \({w}_{j}\) 作为入基变量，过程完全相同)。我们假设当前的基解(basic solution)是可行的，并且希望在将 \({z}_{j}\) 从零开始增大的同时保持解的可行性。 \({z}_{j}\) 的增大改变了当前值等于向量 \({D}^{-1}\mathbf{1}\) 的基变量，我们用 \(q\) 表示这些基变量。假设向量 \(v\) 是 (9.37) 中对应于 \({z}_{j}\) 的非基列。随着 \({z}_{j}\) 的增大，基变量向量由 \(q - v{z}_{j}\) 给出，并且应该保持非负，这意味着对于所有行 \(i = 1,\ldots ,d\) 都有 \({q}_{i} - {v}_{i}{z}_{j} \geq  0\) 。因为 \(q \geq  \mathbf{0}\) 已经成立，所以我们只需要考虑满足 \({v}_{i} > 0\) 的行 \(i\) 。如果 \(v\) 没有正元素，我们可以无限增大 \({z}_{j}\) ，这将在 \(S\) 中定义一条无界的“射线(ray)”，但 \(S\) 是有界的，所以这种情况不会发生(互补换基法(complementary pivoting)也适用于一般矩阵 \(C\) ，此时算法可能会出现“射线终止(ray termination)”的情况)。
\end{tcolorbox}

Hence, the entering column \(v\) has at least one positive component \({v}_{i}\) , and the increase of \({z}_{j}\) is limited by the condition \(q - v{z}_{j} \geq  \mathbf{0}\) , that is, \({q}_{i} \geq  {v}_{i}{z}_{j}\) for all \(i\) such that \({v}_{i} > 0\) , or equivalently

\begin{tcolorbox}
因此，入基列 \(v\) 至少有一个正分量 \({v}_{i}\) ， \({z}_{j}\) 的增大受到条件 \(q - v{z}_{j} \geq  \mathbf{0}\) 的限制，即对于所有满足 \({v}_{i} > 0\) 的 \(i\) 都有 \({q}_{i} \geq  {v}_{i}{z}_{j}\) ，或者等价地
\end{tcolorbox}

\[
{z}_{j} \leq  \min \left\{  {\frac{{q}_{i}}{{v}_{i}} \mid  i = 1,\ldots ,d,{v}_{i} > 0}\right\}  . \tag{9.38}
\]

Exactly when \({z}_{j}\) is equal to the minimum in (9.38), one of the inequalities \({q}_{i} - {v}_{i}{z}_{j} \geq  0\) is tight. Such a row \(i\) determines the leaving variable. Finding the minimum in (9.38) is also called the minimum ratio test.

\begin{tcolorbox}
当 \({z}_{j}\) 恰好等于 (9.38) 中的最小值时，不等式 \({q}_{i} - {v}_{i}{z}_{j} \geq  0\) 中有一个取等号。这样的行 \(i\) 确定了离基变量(leaving variable)。求 (9.38) 中的最小值也称为最小比值检验(minimum ratio test)。
\end{tcolorbox}

Nondegeneracy implies that the minimum in (9.38) and hence the leaving variable is unique. The reason is that if the minimum is taken for two (or more) rows, then only one of the corresponding two basic variables can be chosen to leave the basis, but the other will stay basic and after the pivoting step will have value zero, which means that the new basis is degenerate; we give an example in the next section.

\begin{tcolorbox}
非退化性意味着 (9.38) 中的最小值以及离基变量是唯一的。原因是，如果最小值在两行(或更多行)中取得，那么在对应的两个基变量中只能选择一个离开基，而另一个将保持为基变量，并且在换基步骤之后其值将为零，这意味着新的基是退化的；我们将在下一节给出一个例子。
\end{tcolorbox}

The column of the system (9.37) for the entering variable \({z}_{j}\) , and the row of the leaving variable, determine the pivot element \({v}_{i}\) which by (9.38) is positive. The pivoting operation replaces the leaving basic variable by the entering variable, which means to replace the column of \(D\) for the leaving variable with the entering column. Rather than computing from scratch the inverse \({D}^{-1}\) of the new basic matrix, it suffices to manipulate the system (9.37) with row operations that turn the entering column into the unit vector \({e}_{i}\) . In these row operations, the pivot row is divided by the pivot element \({v}_{i}\) , and suitable multiples of the pivot row are subtracted from the other rows. In our examples (9.32)-(9.35) above, these row operations reflect the algebraic substitutions to get from one dictionary to the next.

\begin{tcolorbox}
系统 (9.37) 中入基变量 \({z}_{j}\) 所在的列以及离基变量所在的行确定了主元 \({v}_{i}\) ，根据 (9.38) 可知主元是正的。换基操作(pivoting)用入基变量替换离基的基变量，这意味着用入基列替换离基变量对应的 \(D\) 列。不必从头计算新基矩阵的逆 \({D}^{-1}\) ，只需对系统 (9.37) 进行行变换，将入基列变为单位向量(unit vector) \({e}_{i}\) 即可。在这些行变换中，主元行除以主元 \({v}_{i}\) ，并从其他行中减去主元行的适当倍数。在上面的例子 (9.32) - (9.35) 中，这些行变换反映了从一个字典(dictionary)到下一个字典的代数替换。
\end{tcolorbox}

Pivoting, as described so far, is a general method of traversing the edges of a polyhedron to get from one vertex to another, where each vertex corresponds to a basic feasible solution. In linear programming, which is the maximization of a linear function over a polyhedron, pivoting is the main operation of the very important simplex algorithm.

\begin{tcolorbox}
到目前为止所描述的换基操作(pivoting)是一种遍历多面体(polyhedron)的边(edge)从一个顶点(vertex)到另一个顶点的通用方法，其中每个顶点对应一个基本可行解。在线性规划(linear programming)中，即在多面体上最大化一个线性函数，换基操作是非常重要的单纯形算法(simplex algorithm)的主要操作。
\end{tcolorbox}

What is special about complementary pivoting is the initialization and the choice of the next entering variable, with the goal to solve the system (9.31), that is, to find a complementary solution \(z,w\) . An initial complementary solution is \(z = \mathbf{0}\) and \(w = \mathbf{1}\) . The freedom to move along a path is provided by relaxing the complementarity condition to allow for a single subscript \(\ell\) (for the missing label) so that both \({z}_{\ell }\) and \({w}_{\ell }\) can take positive values and thus both be basic variables.

\begin{tcolorbox}
互补换基法的特殊之处在于初始化和下一个入基变量的选择，其目标是求解系统 (9.31)，即找到一个互补解 \(z,w\) 。一个初始互补解是 \(z = \mathbf{0}\) 和 \(w = \mathbf{1}\) 。通过放宽互补性条件，允许一个下标 \(\ell\) (对应缺失标签)，使得 \({z}_{\ell }\) 和 \({w}_{\ell }\) 都可以取正值，从而都成为基变量，以此来提供沿着路径移动的自由度。
\end{tcolorbox}

At the initial complementary solution \(z = \mathbf{0}\) and \(w = \mathbf{1}\) , the first entering variable is \({z}_{\ell }\) . Once \({z}_{\ell }\) has entered the basis, a variable \({w}_{i}\) has become nonbasic, which is now nonbasic alongside \({z}_{i}\) . The complementary pivoting rule is to take the complement of the leaving variable as the next entering variable; the complement of \({w}_{i}\) is \({z}_{i}\) and vice versa (the same variable may enter and leave the basis multiple times, so it is also possible that some \({z}_{i}\) leaves the basis and afterwards \({w}_{i}\) enters the basis). The missing label \(\ell\) provides a unique start of the algorithm, and subsequently the unique leaving variable and the complementary pivoting rule imply a unique continuation at each step. Hence, the algorithm progresses in a unique manner, and no almost complementary basis can be revisited. The algorithm terminates when either \({w}_{\ell }\) or \({z}_{\ell }\) leaves the basis, and then reaches a solution to (9.31).

\begin{tcolorbox}
在初始互补解 \(z = \mathbf{0}\) 和 \(w = \mathbf{1}\) 处，第一个入基变量是 \({z}_{\ell }\) 。一旦 \({z}_{\ell }\) 进入基，变量 \({w}_{i}\) 就变成了非基变量，现在它与 \({z}_{i}\) 一起成为非基变量。互补换基规则是将离基变量的互补变量作为下一个入基变量； \({w}_{i}\) 的互补变量是 \({z}_{i}\) ，反之亦然(同一个变量可能多次进入和离开基，因此也有可能某个 \({z}_{i}\) 离开基，之后 \({w}_{i}\) 进入基)。缺失标签 \(\ell\) 为算法提供了唯一的起始点，随后唯一的离基变量和互补换基规则意味着每一步都有唯一的延续。因此，算法以唯一的方式进行，并且不会重新访问几乎互补的基。当 \({w}_{\ell }\) 或 \({z}_{\ell }\) 离开基时，算法终止，然后得到 (9.31) 的解。
\end{tcolorbox}

Geometrically, this describes a path of edges of the polytope \(S\) . The points on that path have all labels in \(\{ 1,\ldots ,d\}  - \{ \ell \}\) , and the end nodes of the path are completely labeled. With \(C\) as in (9.28), this is the LH path for missing label \(\ell\) on \(S = P \times  Q\) , as shown for our example in Figure 9.19.

\begin{tcolorbox}
从几何角度看，这描述了有界多面体 \(S\) 的一条边路径。该路径上的点在 \(\{ 1,\ldots ,d\}  - \{ \ell \}\) 中都有所有标签，并且路径的端点是完全标记的。对于 (9.28) 中的 \(C\) ，这是在 \(S = P \times  Q\) 上缺失标签 \(\ell\) 的莱姆基 - 豪森(LH)路径，如图 9.19 中的示例所示。
\end{tcolorbox}

\section*{9.8 Degeneracy Resolution}

\begin{tcolorbox}
\section*{9.8 退化问题解决}
\end{tcolorbox}

The Lemke-Howson algorithm needs, and has, an extension when the game is degenerate and the leaving variable is not always unique. As an example, consider the modification of our main example (9.1)

\begin{tcolorbox}
当博弈是退化的且离基变量并不总是唯一时，莱姆基 - 豪森算法需要并已经有了一个扩展。例如，考虑对我们的主要示例 (9.1) 进行修改
\end{tcolorbox}

\[
A = \left( \begin{array}{ll} 2 & 0 \\  1 & 2 \\  0 & 3 \end{array}\right) ,\;B = \left( \begin{array}{ll} 2 & 0 \\  1 & 1 \\  4 & 1 \end{array}\right) . \tag{9.39}
\]

where the second row of \(B\) has both columns as best responses, and the game is therefore degenerate. The last two rows of the initial dictionary (9.32) then need to be changed to

\begin{tcolorbox}
其中 \(B\) 的第二行的两列都是最优反应，因此该博弈是退化的。那么初始字典 (9.32) 的最后两行需要改为
\end{tcolorbox}

\[
{w}_{4} = 1 - 2{z}_{1} - {z}_{2} - 4{z}_{3} \tag{9.40}
\]

\[
{w}_{5} = 1\; - {z}_{2} - {z}_{3}.
\]

With the missing label \(\ell  = 2\) , the first entering variable \({z}_{2}\) means both \({w}_{4}\) and \({w}_{5}\) become simultaneously zero when \({z}_{2} = 1\) . Suppose that, as before, we choose \({w}_{5}\) as the leaving variable. Then (9.40) is rewritten as

\begin{tcolorbox}
对于缺失标签 \(\ell  = 2\) ，第一个入基变量 \({z}_{2}\) 意味着当 \({z}_{2} = 1\) 时， \({w}_{4}\) 和 \({w}_{5}\) 同时变为零。假设和之前一样，我们选择 \({w}_{5}\) 作为离基变量。那么 (9.40) 可以重写为
\end{tcolorbox}

\[
{w}_{4} = 0 - 2{z}_{1} + {w}_{5} - 3{z}_{3} \tag{9.41}
\]

\[
{z}_{2} = 1\; - {w}_{5} - {z}_{3},
\]

which replaces the last two rows in (9.33). The next entering variable is \({z}_{5}\) , with \({w}_{3}\) leaving, which produces the dictionary (9.34) except for the replacement of the last two rows by (9.41). The next entering variable is \({z}_{3}\) in (9.41), where the leaving variable is necessarily \({w}_{4}\) because \({z}_{3}\) cannot be increased at all while keeping \({w}_{4}\) nonnegative. The new (partial) dictionary obtained from (9.41) is

\begin{tcolorbox}
这将替换 (9.33) 中的最后两行。下一个入基变量是 \({z}_{5}\) ， \({w}_{3}\) 离基，除了用 (9.41) 替换最后两行外，得到字典 (9.34)。(9.41) 中的下一个入基变量是 \({z}_{3}\) ，其中离基变量必然是 \({w}_{4}\) ，因为在保持 \({w}_{4}\) 非负的情况下， \({z}_{3}\) 根本无法增加。从 (9.41) 得到的新(部分)字典是
\end{tcolorbox}

\[
{z}_{3} = 0 - \frac{2}{3}{z}_{1} + \frac{1}{3}{w}_{5} - \frac{1}{3}{w}_{4} \tag{9.42}
\]

\[
{z}_{2} = 1 + \frac{2}{3}{z}_{1} - \frac{4}{3}{w}_{5} + \frac{1}{3}{w}_{4},
\]

which represents the same basic feasible solution where \({z}_{3} = 0\) and \({w}_{4} = 0\) and hence the same vertex(0,1,0)of \(P\) as the dictionary(9.41), except that the basis has changed. The only difference is that in (9.41) \({w}_{4}\) is basic and \({z}_{3}\) nonbasic, and vice versa in (9.42).

\begin{tcolorbox}
它表示与字典 (9.41) 相同的基本可行解，其中 \({z}_{3} = 0\) 和 \({w}_{4} = 0\) ，因此与 \(P\) 的同一个顶点 (0, 1, 0) 相同，只是基发生了变化。唯一的区别是在 (9.41) 中 \({w}_{4}\) 是基变量， \({z}_{3}\) 是非基变量，而在 (9.42) 中则相反。
\end{tcolorbox}

The equations (9.42) replace the last two rows in (9.35), where the new entering variable is \({z}_{4}\) as the complement of the variable \({w}_{4}\) that has left the basis. This leads to the same final pivoting step as in (9.35) with \({w}_{2}\) leaving, and the final dictionary composed of the first three rows of (9.36) and the two equations in (9.42). The resulting equilibrium \(\left( {\bar{x},\bar{y}}\right)\) of the game (9.39) is \(\bar{x} = \left( {0,1,0}\right)\) and \(\bar{y} = \left( {\frac{1}{2},\frac{1}{2}}\right)\) .

\begin{tcolorbox}
方程(9.42)替换了(9.35)中的最后两行，其中新的入基变量(entering variable)是 \({z}_{4}\) ，它是已离基的变量 \({w}_{4}\) 的补。这导致了与(9.35)中相同的最终换基操作(pivoting)步骤，其中 \({w}_{2}\) 离基(leaving)，最终的字典(dictionary)由(9.36)的前三行和(9.42)中的两个方程组成。博弈(9.39)的所得均衡(equilibrium) \(\left( {\bar{x},\bar{y}}\right)\) 为 \(\bar{x} = \left( {0,1,0}\right)\) 和 \(\bar{y} = \left( {\frac{1}{2},\frac{1}{2}}\right)\) 。
\end{tcolorbox}

Here, the algorithm has still succeeded in finding an equilibrium of the game. However, we cannot guarantee any more that it terminates in general. Because of more than one possible leaving variable, the castle in Figure 9.6 now has rooms with more than two doors, and so it is well possible to reenter such a room through its third door, which could result in an endless loop.

\begin{tcolorbox}
在这里，该算法仍然成功找到了博弈的一个均衡。然而，我们不能再保证它在一般情况下会终止。由于存在不止一个可能的离基变量(leaving variable)，图9.6中的城堡(castle)现在有多于两扇门的房间，因此很有可能通过第三扇门重新进入这样的房间，这可能会导致无限循环。
\end{tcolorbox}

This problem occurs when there are degenerate bases where the corresponding basic feasible solution has basic variables with value zero. The solution to this problem is to make sure, by a perturbation of the right-hand side 1 of the system \({Cz} + {Iw} = 1\) , that basic variables in a basic feasible solution are always positive. The perturbation can be made vanishingly small so that it is in fact fictitious, and only induces a certain deterministic rule for the leaving variable.

\begin{tcolorbox}
当存在退化基(对应基本可行解的基变量(basic variables)值为零)时，就会出现这个问题。解决这个问题的方法是，通过对系统 \({Cz} + {Iw} = 1\) 的右侧1进行微扰，确保基本可行解中的基变量始终为正。微扰可以小到可以忽略不计，因此实际上是虚拟的，它只是为离基变量引入了某种确定性规则。
\end{tcolorbox}

Consider a new parameter \(\varepsilon  > 0\) , which will be taken to be arbitarily small. The modified system is \({Cz} + {Iw} = \mathbf{1} + {\left( \varepsilon ,{\varepsilon }^{2},\ldots ,{\varepsilon }^{d}\right) }^{\top }\) and consequently the system (9.37) for a basis matrix \(D\) is

\begin{tcolorbox}
考虑一个新参数 \(\varepsilon  > 0\) ，它将被设为任意小。修改后的系统为 \({Cz} + {Iw} = \mathbf{1} + {\left( \varepsilon ,{\varepsilon }^{2},\ldots ,{\varepsilon }^{d}\right) }^{\top }\) ，因此对于基矩阵 \(D\) 的系统(9.37)为
\end{tcolorbox}

\[
{D}^{-1}{Cz} + {D}^{-1}w = {D}^{-1}\mathbf{1} + {D}^{-1}{\left( \varepsilon ,{\varepsilon }^{2},\ldots ,{\varepsilon }^{d}\right) }^{\top }. \tag{9.43}
\]

If \(\varepsilon  = 0\) this is just the original system (9.37), where the corresponding basic solution is feasible if \({D}^{-1}\mathbf{1} \geq  0\) . We now want to strengthen this condition so that the basic solution is also positive when \(\varepsilon\) takes sufficiently small positive values. This holds trivially when \({D}^{-1}\mathbf{1} > 0\) , but for a degenerate basis this vector may have zero components. Denote the \(i\) th row of the matrix \(\left\lbrack  {{D}^{-1}\mathbf{1}{D}^{-1}}\right\rbrack\) by \({r}_{i} = \left( {{q}_{i0},{q}_{i1},\ldots ,{q}_{id}}\right)\) . Then the right-hand side in (9.43) is positive for all sufficiently small positive \(\varepsilon\) if and only if the first nonzero entry of \({r}_{i}\) is positive for all \(i = 1,\ldots ,d\) , because the \(i\) th row in (9.43) is \({q}_{i0} + {q}_{i1}\varepsilon  + {q}_{i1}{\varepsilon }^{2} + \cdots  + {q}_{id}{\varepsilon }^{d}\) , and the first nonzero summand, say \({q}_{ij}{\varepsilon }^{j}\) , outweighs for small \(\varepsilon\) all the remaining summands which contain higher powers of \(\varepsilon\) . Note that \({r}_{i}\) cannot be the all-zero row vector because its last \(d\) components are a row of \({D}^{-1}\) . The matrix \(\left\lbrack  \begin{array}{ll} {D}^{-1}\mathbf{1} & {D}^{-1} \end{array}\right\rbrack\) is called lexico-positive if the first nonzero component of each of its rows \({r}_{i}\) is positive. This is clearly the case initially when the basic variables are \(w\) with \(D = I\) , and will be maintained as a condition by the algorithm; it means that the perturbed system (9.43) is never degenerate for sufficiently small positive \(\varepsilon\) .

\begin{tcolorbox}
如果 \(\varepsilon  = 0\) 这就是原始系统 (9.37)，当 \({D}^{-1}\mathbf{1} \geq  0\) 时，相应的基解(basic solution)是可行的。现在我们想加强这个条件，使得当 \(\varepsilon\) 取足够小的正值时，基解也是正的。当 \({D}^{-1}\mathbf{1} > 0\) 时，这显然成立，但对于退化基，这个向量可能有零分量。用 \({r}_{i} = \left( {{q}_{i0},{q}_{i1},\ldots ,{q}_{id}}\right)\) 表示矩阵 \(\left\lbrack  {{D}^{-1}\mathbf{1}{D}^{-1}}\right\rbrack\) 的第 \(i\) 行。那么，当且仅当对于所有的 \(i = 1,\ldots ,d\) ， \({r}_{i}\) 的第一个非零元素为正时，对于所有足够小的正值 \(\varepsilon\) ，(9.43) 式的右侧为正，因为 (9.43) 式的第 \(i\) 行是 \({q}_{i0} + {q}_{i1}\varepsilon  + {q}_{i1}{\varepsilon }^{2} + \cdots  + {q}_{id}{\varepsilon }^{d}\) ，并且对于小的 \(\varepsilon\) ，第一个非零求和项，比如 \({q}_{ij}{\varepsilon }^{j}\) ，会大于所有包含 \(\varepsilon\) 更高次幂的其余求和项。注意， \({r}_{i}\) 不可能是全零行向量，因为它的最后 \(d\) 个分量是 \({D}^{-1}\) 的一行。如果矩阵 \(\left\lbrack  \begin{array}{ll} {D}^{-1}\mathbf{1} & {D}^{-1} \end{array}\right\rbrack\) 的每一行 \({r}_{i}\) 的第一个非零分量为正，则称该矩阵为字典正序(lexico - positive)。显然，当基变量(basic variables)为 \(w\) 且 \(D = I\) 时，初始情况下就是这种情况，并且算法会将其作为一个条件保持；这意味着对于足够小的正值 \(\varepsilon\) ，扰动系统 (9.43) 永远不会退化。
\end{tcolorbox}

Consequently, the leaving variable in any pivoting step is always unique, and determined by the following extension of the minimum ratio test (9.38). Assume the entering variable is \({z}_{j}\) (which could also be, completely equivalently, \({w}_{j}\) ) with entering column \(v\) , which has at least one positive component \({v}_{i}\) . In the perturbed system (9.43), the constraints for the increase of \({z}_{j}\) are now

\begin{tcolorbox}
因此，在任何换基操作(pivoting)步骤中，离基变量(leaving variable)总是唯一的，并且由最小比值检验(minimum ratio test)(9.38) 的以下扩展来确定。假设入基变量(entering variable)是 \({z}_{j}\) (完全等价地，也可以是 \({w}_{j}\) )，其入基列是 \(v\) ，该列至少有一个正分量 \({v}_{i}\) 。在扰动系统 (9.43) 中， \({z}_{j}\) 增加的约束条件现在是
\end{tcolorbox}

\[
{z}_{j} \leq  \frac{{q}_{i0}}{{v}_{i}} + \frac{{q}_{i1}}{{v}_{i}}\varepsilon  + \frac{{q}_{i2}}{{v}_{i}}{\varepsilon }^{2} + \cdots  + \frac{{q}_{id}}{{v}_{i}}{\varepsilon }^{d} \tag{9.44}
\]

for all rows \(i\) where \({v}_{i} > 0\) . The coefficients of \(1,\varepsilon ,{\varepsilon }^{2},\ldots ,{\varepsilon }^{d}\) in (9.44) are given by the row vector \(\frac{1}{{v}_{i}}{r}_{i}\) . We consider the lexicographically smallest of these row vectors, that is, the one with the smallest first component where any of these components differ; for example, among the rows

\begin{tcolorbox}
对于所有满足 \({v}_{i} > 0\) 的行 \(i\) 。(9.44) 式中 \(1,\varepsilon ,{\varepsilon }^{2},\ldots ,{\varepsilon }^{d}\) 的系数由行向量 \(\frac{1}{{v}_{i}}{r}_{i}\) 给出。我们考虑这些行向量按字典序最小的那个，即，在这些分量不同的情况下，第一个分量最小的那个；例如，在这些行中
\end{tcolorbox}

\[
\left( {0,0,2, - 1}\right) ,\;\left( {0,0,1,{1000}}\right) ,\;\left( {0,3, - {100}, - 8}\right) ,
\]

the second row is the lexicographically smallest. The lexicographically smallest row is unique, because if two rows were identical, say \(\frac{1}{{v}_{i}}{r}_{i}\) and \(\frac{1}{{v}_{k}}{r}_{k}\) , this would imply two linearly dependent rows of \({D}^{-1}\) . Choosing the lexico-smallest row \(\frac{1}{{v}_{i}}{r}_{i}\) in (9.44) is called the lexico-minimum ratio test. Because this row is lexico-positive, this means that \({z}_{j}\) , when it enters the basis in the perturbed system, assumes a (possibly very small) positive value in the new basic solution. Furthermore, after the pivoting step, all other basic variables stay positive in the perturbed system, and the new basis is again lexico-positive.

\begin{tcolorbox}
第二行是按字典序最小的。按字典序最小的行是唯一的，因为如果两行相同，比如 \(\frac{1}{{v}_{i}}{r}_{i}\) 和 \(\frac{1}{{v}_{k}}{r}_{k}\) ，这将意味着 \({D}^{-1}\) 有两个线性相关的行。在 (9.44) 式中选择按字典序最小的行 \(\frac{1}{{v}_{i}}{r}_{i}\) 称为字典序最小比率检验(lexico - minimum ratio test)。因为这一行是字典正序的，这意味着 \({z}_{j}\) 在扰动系统中进入基时，在新的基解中取一个(可能非常小的)正值。此外，在换基操作步骤之后，扰动系统中的所有其他基变量保持为正，并且新的基再次是字典正序的。
\end{tcolorbox}

In short, the lexico-minimum ratio test determines the next leaving variable uniquely, and the perturbed system is always nondegenerate. What is nice is that the lexico-minimum ratio test involves only a comparison of \(\frac{1}{{v}_{i}}{r}_{i}\) for the rows \({r}_{i}\) of the matrix \(\left\lbrack  \begin{array}{ll} {D}^{-1}\mathbf{1} & {D}^{-1} \end{array}\right\rbrack\) where \({v}_{i} > 0\) , and there is no actual perturbation at all, which is merely mimicked as if \(\varepsilon\) was positive but vanishingly small. Moreover, all the information about the rows \({r}_{i}\) is already provided by the original dictionary (9.37) where the matrix \({D}^{-1}\) is given by the columns for the variables \(w\) , so no additional information needs to be stored.

\begin{tcolorbox}
简而言之，字典序最小比率检验能唯一确定下一个离基变量，并且摄动系统始终是非退化的。好的方面在于，字典序最小比率检验仅涉及对矩阵 \(\left\lbrack  \begin{array}{ll} {D}^{-1}\mathbf{1} & {D}^{-1} \end{array}\right\rbrack\) 中满足 \({v}_{i} > 0\) 的行 \({r}_{i}\) 的 \(\frac{1}{{v}_{i}}{r}_{i}\) 进行比较，而且根本不需要实际进行摄动，只是模拟 \(\varepsilon\) 为正但无穷小的情况。此外，关于行 \({r}_{i}\) 的所有信息已由原始字典(9.37)提供，其中矩阵 \({D}^{-1}\) 由变量 \(w\) 对应的列给出，因此无需存储额外信息。
\end{tcolorbox}

In summary, complementary pivoting, extended by lexico-minimum ratio test, proves the existence of an equilibrium in a bimatrix game, even if that game is degenerate. In this way, the Lemke-Howson algorithm provides an alternative and constructive equilibrium existence proof, independently of Brouwer's fixed point theorem.

\begin{tcolorbox}
总之，通过字典序最小比率检验扩展后的互补换基法证明了即使双矩阵博弈是退化的，也存在均衡。这样，莱姆基 - 豪森算法独立于布劳威尔不动点定理，提供了一种替代性的、构造性的均衡存在性证明。
\end{tcolorbox}

The Lemke-Howson algorithm finds one equilibrium of a bimatrix game (A, B)but not necessarily all equilibria. The polytopes \(P\) and \(Q\) are also useful for finding all equilibria, with a different algorithm. Basically, such an algorithm computes all nonzero vertices \(x\) of \(P\) and \(y\) of \(Q\) , and outputs those pairs(x, y)that are completely labeled. If the game is degenerate, then some vertices \(x\) may have more than \(m\) labels, or some vertices \(y\) may have more than \(n\) labels. Then fewer labels are required for an equilibrium in the other polytope (as in the example in Figure 9.12), which may define a higher-dimensional face of \(P\) or \(Q\) (or both), all whose elements are equilibrium strategies. That face has a finite description either by the set of its vertices (as already computed in this process), or by its set of tight inequalities. For details of such algorithms see Avis, Rosenberg, Savani, and von Stengel (2010).

\begin{tcolorbox}
莱姆基 - 豪森算法能找到双矩阵博弈 (A, B) 的一个均衡，但不一定能找到所有均衡。有界多面体 \(P\) 和 \(Q\) 也可用于通过不同算法找到所有均衡。基本上，这种算法会计算 \(P\) 的所有非零顶点 \(x\) 和 \(Q\) 的所有非零顶点 \(y\) ，并输出那些完全标记的对 (x, y)。如果博弈是退化的，那么某些顶点 \(x\) 可能有多于 \(m\) 个标签，或者某些顶点 \(y\) 可能有多于 \(n\) 个标签。那么另一个有界多面体中的均衡所需的标签会更少(如图 9.12 中的示例)，这可能会定义 \(P\) 或 \(Q\) (或两者)的一个更高维度的面，该面的所有元素都是均衡策略。这个面可以通过其顶点集(在这个过程中已经计算得出)或其紧不等式集进行有限描述。有关此类算法的详细信息，请参阅 Avis、Rosenberg、Savani 和 von Stengel(2010)。
\end{tcolorbox}

\section*{9.9 Further Reading}

\begin{tcolorbox}
\section*{9.9 进一步阅读}
\end{tcolorbox}

The subdivisions of the players' mixed-strategy simplices into labeled best-response regions, and the best-response diagrams like in Figure 9.16, are due to Shapley (1974). The example in Exercise 9.1 of a game where not all equilibria are found by the Lemke-Howson algorithm when started from the artificial equilibrium is shown in Shapley (1974, Figure 3) and is due to Robert Wilson. A rigorous definition of the LH graph that we described following Lemma 9.3 is given in von Stengel (2002).

\begin{tcolorbox}
将玩家的混合策略单纯形细分为标记的最优反应区域，以及如图 9.16 所示的最优反应示意图，这些都归功于 Shapley(1974)。练习 9.1 中的一个博弈示例，即从人工均衡开始使用莱姆基 - 豪森算法无法找到所有均衡的情况，在 Shapley(1974，图 3)中有所展示，该示例由 Robert Wilson 提出。我们在引理 9.3 之后描述的 LH 图的严格定义在 von Stengel(2002)中给出。
\end{tcolorbox}

The LH paths have an intrinsic direction, defined in terms of the determinants of the payoff matrices for the equilibrium supports, and the end nodes of the paths have an opposite sign called the index of the equilibrium. This was first proved by Shapley (1974); for a streamlined presentation see von Stengel (2021). The direction of the path accounts for the " \(\mathrm{D}\) " in the complexity class "PPAD" defined by Papadimitriou (1994) that contains the computational problem of finding an equilibrium of a bimatrix game. This problem is in fact complete for PPAD (Chen, Deng, and Teng, 2009).

\begin{tcolorbox}
LH 路径具有内在方向，该方向根据均衡支撑的收益矩阵的行列式来定义，并且路径的端点具有相反的符号，称为均衡的指数。这一点最早由 Shapley(1974)证明；关于简化的表述，请参阅 von Stengel(2021)。路径的方向解释了 Papadimitriou(1994)定义的复杂性类 “PPAD” 中的 “ \(\mathrm{D}\) ”，该类包含寻找双矩阵博弈均衡的计算问题。实际上，这个问题对于 PPAD 是完全的(Chen、Deng 和 Teng，2009)。
\end{tcolorbox}

In their original publication, Lemke and Howson (1964) considered arbitrary nonnegative vectors instead of mixed strategies, like we do in (9.24) to define the polytopes \(P\) and \(Q\) . However, the payoff matrices \(A\) and \(B\) are replaced by strictly positive cost matrices \({A}^{\prime } = \mathbf{1}\alpha {\mathbf{1}}^{\top } - A\) and \({B}^{\prime } = \mathbf{1}\alpha {\mathbf{1}}^{\top } - B\) for some sufficiently large \(\alpha  \in  \mathbb{R}\) , to define the polyhedra

\begin{tcolorbox}
在他们的原始论文中，Lemke 和 Howson(1964)考虑的是任意非负向量而非混合策略，就像我们在(9.24)中定义有界多面体 \(P\) 和 \(Q\) 那样。然而，收益矩阵 \(A\) 和 \(B\) 被严格正的成本矩阵 \({A}^{\prime } = \mathbf{1}\alpha {\mathbf{1}}^{\top } - A\) 和 \({B}^{\prime } = \mathbf{1}\alpha {\mathbf{1}}^{\top } - B\) 所取代，其中 \(\alpha  \in  \mathbb{R}\) 足够大，以定义多面体
\end{tcolorbox}

\[
{P}^{\prime } = \left\{  {x \in  {\mathbb{R}}^{m} \mid  \;x \geq  \mathbf{0},{\left( {B}^{\prime }\right) }^{\top }x \geq  \mathbf{1}}\right\}  ,
\]

\[
{Q}^{\prime } = \left\{  {y \in  {\mathbb{R}}^{n} \mid  {A}^{\prime }y \geq  \mathbf{1},\;y \geq  \mathbf{0}}\right\}  .
\]

These are clearly unbounded, but do not have the artificial equilibrium(0,0)as a complementary solution. Instead, the algorithm is started from the "ray" of almost completely labeled pairs \(\left( {x,y}\right)  \in  {P}^{\prime } \times  {Q}^{\prime }\) with missing label \(\ell\) given by \(x = {e}_{\ell }\alpha\) so that \(\alpha  \in  \mathbb{R}\) and \({e}_{\ell }\alpha  \in  {P}^{\prime }\) (if \(1 \leq  \ell  \leq  m\) , otherwise \(y = {e}_{\ell  - m}\alpha\) ) and the pure best response to \(\ell\) as the corresponding vertex \(y\) in \({Q}^{\prime }\) (respectively, \(x\) in \({P}^{\prime }\) ).

\begin{tcolorbox}
这些显然是无界的，但并不将人工均衡(0, 0)作为互补解。相反，算法从几乎完全标记的对 \(\left( {x,y}\right)  \in  {P}^{\prime } \times  {Q}^{\prime }\) 的“射线”开始，其缺失标签为 \(\ell\) ，由 \(x = {e}_{\ell }\alpha\) 给出，使得 \(\alpha  \in  \mathbb{R}\) 和 \({e}_{\ell }\alpha  \in  {P}^{\prime }\) (如果 \(1 \leq  \ell  \leq  m\) ，否则 \(y = {e}_{\ell  - m}\alpha\) )，并且对 \(\ell\) 的纯最优反应作为 \({Q}^{\prime }\) 中的相应顶点 \(y\) (分别地， \({P}^{\prime }\) 中的 \(x\) )。
\end{tcolorbox}

The example (9.10) of strategic equivalence of two payoff matrices without a column shift is similar to Example 2.1 of Liu (1996). Moulin and Vial (1978) showed that strong strategic equivalence, defined in (9.11), only admits scaling and column shifts.

\begin{tcolorbox}
两个没有列移位的收益矩阵的策略等价的示例 (9.10) 与 Liu (1996) 的示例 2.1 类似。Moulin 和 Vial (1978) 表明，(9.11) 中定义的强策略等价仅允许缩放和列移位。
\end{tcolorbox}

The polyhedra \(\bar{P}\) and \(\bar{Q}\) in (9.20) have been first described by Mangasarian (1964). As mentioned at the end of Section 9.8, these polyhedra, or the polytopes \(P\) and \(Q\) , can be used to find all equilibria of a bimatrix game, even for degenerate games; see Avis, Rosenberg, Savani, and von Stengel (2010). The projective transformation in Figure 9.18 has been described in von Stengel (2002), which is a general survey on computing equilibria in two-player games. For an introduction to polyhedra and polytopes see the first chapters of Ziegler (1995).

\begin{tcolorbox}
式(9.20)中的多面体 \(\bar{P}\) 和 \(\bar{Q}\) 最早由曼加萨里安(Mangasarian，1964 年)描述。正如 9.8 节末尾所提到的，这些多面体，或者有界多面体 \(P\) 和 \(Q\) ，可用于找出双矩阵博弈(bimatrix game)的所有均衡(equilibrium)，即使是退化博弈(degenerate game)也是如此；参见阿维斯(Avis)、罗森伯格(Rosenberg)、萨瓦尼(Savani)和冯·施滕格尔(von Stengel，2010 年)的研究。图 9.18 中的射影变换(projective transformation)在冯·施滕格尔(2002 年)的研究中有描述，该研究是关于计算两人博弈均衡的综合综述。关于多面体和有界多面体的介绍，可参阅齐格勒(Ziegler，1995 年)著作的前几章。
\end{tcolorbox}

For a general \(d \times  d\) matrix \(C\) and a vector \(q \in  {\mathbb{R}}^{d}\) instead of1, the conditions (9.31) define a so-called linear complementarity problem. If the right-hand side \(q\) has negative components, then \(w = q\) and \(z = \mathbf{0}\) is not a suitable initial solution. By extending this system with an extra column (usually the vector -1) and extra variable \({z}_{0}\) to

\begin{tcolorbox}
对于一般的 \(d \times  d\) 矩阵 \(C\) 和向量 \(q \in  {\mathbb{R}}^{d}\) (而非 1)，条件(9.31)定义了所谓的线性互补问题(linear complementarity problem)。如果右侧 \(q\) 有负分量，那么 \(w = q\) 和 \(z = \mathbf{0}\) 就不是合适的初始解。通过给这个系统增加一列(通常是向量 -1)和一个额外变量 \({z}_{0}\) 来扩展这个系统
\end{tcolorbox}

\[
{Cz} + w - \mathbf{1}{z}_{0} = q,\;z \geq  \mathbf{0},\;w \geq  \mathbf{0},\;{z}_{i}{w}_{i} = 0\;\left( {1 \leq  i \leq  d}\right) , \tag{9.45}
\]

then this system has a "primary ray" of solutions where \(z = \mathbf{0}\) and \(w = q + \mathbf{1}{z}_{0}\) for sufficiently large \({z}_{0}\) . Minimizing \({z}_{0}\) in (9.45) so that \(z = \mathbf{0}\) and \(w \geq  \mathbf{0}\) then allows starting the complementary pivoting algorithm, invented by Lemke (1965). A large body of work is concerned with whether Lemke's algorithm terminates with a solution to (9.45) where \({z}_{0} = 0\) , depending on properties of the matrix \(C\) (usually denoted by \(- M\) ); see Cottle, Pang, and Stone (1992).

\begin{tcolorbox}
那么这个系统有一个解的“主射线”(primary ray)，其中对于足够大的 \({z}_{0}\) ，有 \(z = \mathbf{0}\) 和 \(w = q + \mathbf{1}{z}_{0}\) 。在(9.45)中最小化 \({z}_{0}\) ，使得 \(z = \mathbf{0}\) 和 \(w \geq  \mathbf{0}\) ，然后就可以启动由莱姆基(Lemke，1965 年)发明的互补换基法(complementary pivoting)算法。大量的研究关注莱姆基算法是否能以(9.45)的一个解终止，其中 \({z}_{0} = 0\) ，这取决于矩阵 \(C\) 的性质(通常用 \(- M\) 表示)；参见科特尔(Cottle)、庞(Pang)和斯通(Stone，1992 年)的研究。
\end{tcolorbox}

Pivoting is a general method to move from one basic feasible solution of a system of linear equations with nonnegative variables to another. It is the central step of the simplex algorithm of Dantzig (1963) for linear programming. Its explanation in terms of dictionaries is due to Chvátal (1983), who also explains (page 34) the lexicographic method for degeneracy resolution. Complementary pivoting is only special in terms of the choice of the entering variable as the complement of the variable that has just left the basis.

\begin{tcolorbox}
换基操作(pivoting)是一种从具有非负变量的线性方程组的一个基本可行解转移到另一个基本可行解的通用方法。它是丹齐格(Dantzig，1963 年)用于线性规划(linear programming)的单纯形算法(simplex algorithm)的核心步骤。用字典(dictionary)来解释它归功于赫瓦塔尔(Chvátal，1983 年)，他还解释了(第 34 页)用于解决退化问题的字典序方法。互补换基法的特殊之处仅在于入基变量(entering variable)的选择是刚刚离开基的变量的互补变量。
\end{tcolorbox}

As shown in this chapter, the equilibrium structure of two-player games is best understood in terms of labeled polyhedra and polytopes. This insight can be used to apply known constructions of polytopes to questions about equilibria of bimatrix games, in particular the "dual cyclic polytopes" which have the largest possible number of vertices for a given dimension and number of inequalities (Ziegler, 1995, p. 10f). In von Stengel (1999), these polytopes are used to construct nondegenerate \(n \times  n\) games with a larger number of mixed equilibria than the coordination game (shown for \(n = 3\) in Figure 9.7). These polytopes are also used in Savani and von Stengel (2006, 2016) to construct games with exponentially long Lemke-Howson paths for any missing label.

\begin{tcolorbox}
正如本章所示，两人博弈的均衡结构用带标签的多面体和有界多面体来理解是最合适的。这种见解可用于将已知的有界多面体构造方法应用于双矩阵博弈均衡的相关问题，特别是“对偶循环有界多面体”，对于给定的维度(dimension)和不等式数量，它具有最大可能的顶点数(齐格勒，1995 年，第 10 页及以后)。在冯·施滕格尔(1999 年)的研究中，这些有界多面体被用于构造非退化(nondegenerate)的 \(n \times  n\) 博弈，其混合均衡的数量比协调博弈(coordination game)更多(如图 9.7 中 \(n = 3\) 所示)。这些有界多面体也被萨瓦尼和冯·施滕格尔(2006 年，2016 年)用于构造对于任何缺失标签(missing label)都有指数级长的莱姆基 - 豪森路径(Lemke - Howson path)的博弈。
\end{tcolorbox}

\section*{9.10 Exercises for Chapter 9}

\begin{tcolorbox}
\section*{9.10 第 9 章练习}
\end{tcolorbox}

Exercise 9.1. Consider the symmetric \(3 \times  3\) game \(\left( {A,{A}^{\top }}\right)\) with \(A\) as in (9.12).

\begin{tcolorbox}
练习 9.1。考虑对称的 \(3 \times  3\) 博弈 \(\left( {A,{A}^{\top }}\right)\) ，其中 \(A\) 如(9.12)所示。
\end{tcolorbox}

(a) Draw the best response diagram for this game. [Hint: You do not have to compute anything, but can use one of the pictures in Figure 9.14 with a suitable change of labels.]

\begin{tcolorbox}
(a)画出这个博弈的最优反应(best response)图。[提示:你不必进行任何计算，但可以使用图 9.14 中的一幅图并进行适当的标签(label)更改。]
\end{tcolorbox}

(b) Identify all equilibria of this game.

\begin{tcolorbox}
(b)找出这个博弈的所有均衡。
\end{tcolorbox}

(c) Write down the Lemke-Howson path that starts at(0,0)for each missing label 1,2,3,4,5,6, as in (9.7), by suitably annotating the nodes in \(\widetilde{X}\) and \(\widetilde{Y}\) . Which equilibrium is found for each missing label? [Hint: You can save half of the work by using the symmetry of the game.]

\begin{tcolorbox}
(c)对于每个缺失标签 1、2、3、4、5、6，按照(9.7)的方式，通过对 \(\widetilde{X}\) 和 \(\widetilde{Y}\) 中的节点(node)进行适当标注，写下从(0,0)开始的莱姆基 - 豪森路径。对于每个缺失标签，能找到哪个均衡？[提示:利用博弈的对称性可以节省一半的工作量。]
\end{tcolorbox}

Exercise 9.2. Consider the degenerate \(2 \times  3\) game in Figure 6.12. Construct the best response diagram (note that \(X\) is only a line segment) and identify the set of all equilibria with the help of the labels.

\begin{tcolorbox}
练习9.2。考虑图6.12中的退化双矩阵博弈(bimatrix game)。构建最优反应图(注意b1只是一条线段)，并借助标签确定所有均衡(equilibrium)的集合。
\end{tcolorbox}

Exercise 9.3. Consider the following \(3 \times  3\) game.

\begin{tcolorbox}
练习9.3。考虑以下双矩阵博弈(bimatrix game)。
\end{tcolorbox}

\begin{center}
\includegraphics[max width=0.3\textwidth]{images/bo_d2e5cmbef24c73b03tn0_272_578_1159_489_491_0.jpg}
\end{center}
\hspace*{3em} 

(a) Find all pure strategy equilibria of this game.

\begin{tcolorbox}
(a) 找出该博弈的所有纯策略均衡(pure equilibrium)。
\end{tcolorbox}

(b) The picture below shows, for this game, the mixed-strategy sets \(X\) and \(Y\) for players I and II, but not which set is \(X\) and which is \(Y\) . It also shows their subdivisions into best-response regions, but without labels, which you are asked to provide. Explain which triangle is \(X\) and which is \(Y\) . [Hint: consider the pure best responses first.] For the three corners marked by \(a,b,e\) and \(p,s,u\) , respectively, indicate which pure strategy \(\left( {1,0,0}\right) ,\left( {0,1,0}\right)\) , or(0,0,1)they represent, and mark the triangles with the correct labels ①, ..., ⑥. Is this game degenerate?

\begin{tcolorbox}
(b) 下图展示了该博弈中玩家I和玩家II的混合策略集(mixed-strategy set)b0和b1，但未标明哪个是b0，哪个是b1。图中还展示了它们细分为最优反应区域的情况，但没有标签，需要你添加。解释哪个三角形是b0，哪个是b1。[提示:先考虑纯最优反应。] 对于分别由b2和b3标记的三个角，指出它们代表的纯策略是b4还是(0, 0, 1)，并为三角形标上正确的标签①, ..., ⑥。该博弈是否退化？
\end{tcolorbox}

\begin{center}
\includegraphics[max width=0.7\textwidth]{images/bo_d2e5cmbef24c73b03tn0_273_529_205_931_450_0.jpg}
\end{center}
\hspace*{3em} 

(c) Based on these labels, identify all equilibria of the game in mixed strategies (using the indicated nodes \(a,b,c,d,e,f,g\) and \(p,q,r,s,t,u,v\) ).

\begin{tcolorbox}
(c) 根据这些标签，确定该博弈在混合策略下的所有均衡(使用标明的节点b0和b1)。
\end{tcolorbox}

(d) Show all Lemke-Howson paths for missing label 1, and how they connect the equilibria of the game.

\begin{tcolorbox}
(d) 展示缺失标签1的所有莱姆基 - 豪森路径(Lemke-Howson path)，以及它们如何连接博弈的均衡。
\end{tcolorbox}

Exercise 9.4. Consider a nondegenerate bimatrix game and an arbitrary pure-strategy equilibrium of the game (assuming one exists). Explain why the Lemke-Howson path started at(0,0)always finds this equilibrium for a suitable missing label. Suppose the game has at least two pure equilibria. Explain how the Lemke-Howson algorithm can be used to find a third (not necessarily pure) equilibrium.

\begin{tcolorbox}
练习9.4。考虑一个非退化(nondegenerate)的双矩阵博弈(bimatrix game)以及该博弈的任意一个纯策略均衡(假设存在)。解释为什么从(0, 0)开始的莱姆基 - 豪森路径(Lemke-Howson path)总能为合适的缺失标签找到这个均衡。假设该博弈至少有两个纯均衡。解释如何使用莱姆基 - 豪森算法(Lemke-Howson algorithm)找到第三个(不一定是纯策略的)均衡。
\end{tcolorbox}

Exercise 9.5. Give an example of a bimatrix game(A, B)where every column of \(A\) has a positive entry but where \(Q\) in (9.24) is not a polytope.

\begin{tcolorbox}
练习9.5。给出一个双矩阵博弈(A, B)的例子，其中b0的每一列都有一个正元素，但(9.24)中的b1不是有界多面体(polytope)。
\end{tcolorbox}

Exercise 9.6. Consider a nondegenerate bimatrix game that has a pure-strategy equilibrium. Show that there is a strategically equivalent game so that in this pure equilibrium, both players get a strictly higher payoff than for any other pair of pure strategies. Explain why the "dilemma" in the Prisoner's Dilemma has "gone away" when considering this equivalent game. Also, give an example that shows that nondegeneracy is needed to get strictly higher payoffs.

\begin{tcolorbox}
练习9.6。考虑一个具有纯策略均衡的非退化(nondegenerate)双矩阵博弈(bimatrix game)。证明存在一个策略上等价的博弈，使得在这个纯均衡中，两个玩家的收益都比任何其他纯策略组合的收益严格更高。解释为什么在考虑这个等价博弈时，“囚徒困境”中的“困境”消失了。此外，给出一个例子说明需要非退化条件才能获得严格更高的收益。
\end{tcolorbox}

\section*{10}

\section*{Game Trees with Imperfect Information}

\begin{tcolorbox}
\section*{具有不完美信息的博弈树}
\end{tcolorbox}

Although it appears late in this book, this is a core chapter about non-cooperative games. It studies extensive games, which are game trees with imperfect information. Typically, players do not always have full access to all the information which is relevant to their choices. Extensive games with imperfect information model exactly what information is available to the players when they make a move. The precise modeling and evaluation of strategic information is one of the strengths of game theory.

\begin{tcolorbox}
尽管这一章在本书中出现较晚，但它是关于非合作博弈的核心章节。它研究扩展式博弈，即具有不完美信息的博弈树(game tree)。通常情况下，玩家并不总是能完全获取与他们的选择相关的所有信息。具有不完美信息的扩展式博弈精确地模拟了玩家在行动时所能获得的信息。对战略信息进行精确建模和评估是博弈论的优势之一。
\end{tcolorbox}

The central concept of information sets is introduced in Section 10.2 by means of a detailed example. The subsequent sections give the general definition of extensive games, and the definition of strategies and reduced strategies, which generalize these concepts as they have been introduced for game trees with perfect information. Essentially, all one has to do is to replace "decision node" by "information set".

\begin{tcolorbox}
第10.2节通过一个详细的例子引入了信息集(information set)这一核心概念。后续章节给出了扩展式博弈的一般定义，以及策略和精简策略(reduced strategy)的定义，这些定义将之前为具有完美信息的博弈树引入的概念进行了推广。本质上，只需将“决策节点”替换为“信息集”即可。
\end{tcolorbox}

The concept of perfect recall, treated in Section 10.6, means that a player always remembers what he knew or did earlier. This is a condition on the structure of the information sets of the player, not about how to interpret the game. In a game with perfect recall, the important Theorem 10.4 of Kuhn (1953) states that a player can replace any mixed strategy by a behavior strategy, which is the topic of Secton 10.8. In a behavior strategy, the player randomizes "locally" between his moves at each information set. In contrast, a mixed strategy first selects "globally" a pure strategy at the beginning of a play of the game, which the player then uses as a plan of action throughout play. Behavior strategies are much less complex than mixed strategies, as we will explain. In Section 10.9, we will illustrate optimal mixed and behavior strategies in a game-theoretic variant of the "Monty Hall" problem (with an opposite conclusion to the standard statistical version of this problem).

\begin{tcolorbox}
第10.6节讨论的完美记忆(perfect recall)概念意味着玩家总是记得他之前知道的信息或采取的行动。这是对玩家信息集结构的一个条件，而不是关于如何解释博弈。在具有完美记忆的博弈中，库恩(1953)的重要定理10.4表明，玩家可以用行为策略(behavior strategy)替代任何混合策略(mixed strategy)，这是第10.8节的主题。在行为策略中，玩家在每个信息集上“局部地”对他的行动进行随机化。相比之下，混合策略在游戏开始时“全局地”选择一个纯策略，然后玩家在整个游戏过程中将其作为行动计划使用。正如我们将解释的，行为策略比混合策略简单得多。在第10.9节，我们将在“蒙提·霍尔”问题的博弈论变体中说明最优混合策略和行为策略(该问题的结论与标准统计版本相反)。
\end{tcolorbox}

The final Section 10.10 treats subgames and subgame-perfect equilibria (SPE).

\begin{tcolorbox}
最后一节10.10讨论子博弈(subgame)和子博弈完美均衡(SPE)。
\end{tcolorbox}

\section*{10.1 Prerequisites and Learning Outcomes}

\begin{tcolorbox}
\section*{10.1 先决条件和学习成果}
\end{tcolorbox}

This chapter relies on Chapter 4 and Chapter 6.

\begin{tcolorbox}
本章依赖于第四章和第六章。
\end{tcolorbox}

After studying this chapter, you should be able to:

\begin{tcolorbox}
学习本章后，你应该能够:
\end{tcolorbox}

\begin{itemize}
\item explain extensive games and the concept of information sets;
\end{itemize}

\begin{tcolorbox}
\begin{itemize}
\item 解释扩展式博弈和信息集的概念；
\end{itemize}
\end{tcolorbox}

\begin{itemize}
\item describe pure and reduced pure strategies for extensive games, and the corresponding strategic form and reduced strategic form of the game;
\end{itemize}

\begin{tcolorbox}
\begin{itemize}
\item 描述扩展式博弈的纯策略和精简纯策略，以及博弈相应的策略形式和精简策略型；
\end{itemize}
\end{tcolorbox}

\begin{itemize}
\item define the concept of perfect recall, and see if this holds for a given game;
\end{itemize}

\begin{tcolorbox}
\begin{itemize}
\item 定义完美记忆的概念，并判断给定博弈是否满足这一条件；
\end{itemize}
\end{tcolorbox}

\begin{itemize}
\item explain the difference between behavior and mixed strategies, and state Kuhn's Theorem 10.4;
\end{itemize}

\begin{tcolorbox}
\begin{itemize}
\item 解释行为策略和混合策略的区别，并阐述库恩定理10.4；
\end{itemize}
\end{tcolorbox}

\begin{itemize}
\item construct extensive games from descriptions of games with imperfect information, like the game in Exercise 10.4;
\end{itemize}

\begin{tcolorbox}
\begin{itemize}
\item 根据不完全信息博弈的描述构建扩展式博弈，如练习10.4中的博弈；
\end{itemize}
\end{tcolorbox}

\begin{itemize}
\item find all equilibria of simple extensive games, and represent the mixed strategies in these equilibria as behavior strategies;
\end{itemize}

\begin{tcolorbox}
\begin{itemize}
\item 找出简单扩展式博弈的所有均衡，并将这些均衡中的混合策略表示为行为策略；
\end{itemize}
\end{tcolorbox}

\begin{itemize}
\item explain the concept of subgames and subgame-perfect equilibria (SPE), and how it generalizes and differs from backward induction. Understand that choosing single moves at a time as in backward induction cannot be applied when the game tree has imperfect information.
\end{itemize}

\begin{tcolorbox}
\begin{itemize}
\item 解释子博弈和子博弈完美均衡(SPE)的概念，以及它如何推广逆向归纳法和与之的区别。理解当博弈树存在不完全信息时，不能像逆向归纳法那样一次选择单个行动。
\end{itemize}
\end{tcolorbox}

\section*{10.2 Information Sets}

\begin{tcolorbox}
\section*{10.2 信息集}
\end{tcolorbox}

Consider the situation faced by a large software firm (player I) after a small startup firm (player II) has announced the deployment of a key new technology. The large firm has a large research and development operation. It is generally known that they have researchers working on a wide variety of innovations. However, only the large firm knows for sure whether or not they have made any progress on a product similar to the startup's new technology. The startup firm believes that there is a 50 percent chance that the large firm has developed the basis for a strong competing product. For brevity, when the large firm has the ability to produce a strong competing product, we will say that the large firm is "strong", and "weak" otherwise.

\begin{tcolorbox}
考虑这样一种情况:一家小型初创公司(参与者II)宣布部署一项关键新技术后，一家大型软件公司(参与者I)面临的局面。大型公司拥有庞大的研发部门。众所周知，他们有研究人员在进行各种各样的创新研究。然而，只有大型公司确切知道他们是否在与初创公司新技术类似的产品上取得了进展。初创公司认为大型公司有50\%的可能性已经开发出强有力的竞争产品的基础。为简洁起见，当大型公司有能力生产强有力的竞争产品时，我们称大型公司“强大”，否则称其“弱小”。
\end{tcolorbox}

The large firm, after the announcement by the small firm, has two choices. It can either announce that it too will release a competing product, or it can cede the market for this product. The large firm will certainly condition its choice upon its private knowledge whether it is strong or weak. If the large firm has announced a product, then the startup firm has two choices: It can either negotiate a buyout to the large firm, or it can remain independent and launch its product. The startup firm does not have access to the large firm's private information on the status of its research. However, it does observe whether or not the large firm announces its own product, and may attempt to infer from that announcement whether the large firm is strong or weak.

\begin{tcolorbox}
大型公司在小型公司宣布后有两个选择。它可以宣布也将推出竞争产品，或者放弃该产品的市场。大型公司肯定会根据其自身是强大还是弱小的私人信息来做出选择。如果大型公司宣布推出产品，那么初创公司有两个选择:它可以与大型公司协商被收购，或者保持独立并推出自己的产品。初创公司无法获取大型公司研发状况的私人信息。然而，它能观察到大型公司是否宣布推出自己的产品，并可能试图从该宣布中推断大型公司是强大还是弱小。
\end{tcolorbox}

\begin{center}
\includegraphics[max width=0.5\textwidth]{images/bo_d2e5cmbef24c73b03tn0_276_215_218_662_732_0.jpg}
\end{center}
\hspace*{3em} 

Figure 10.1 Game with imperfect information for player II, indicated by the information set that contains two decision nodes with the same moves \(a\) and \(b\) . As always, at a leaf of the tree, the top payoff is to player \(I\) , the bottom payoff to player II.

\begin{tcolorbox}
图10.1 参与者II的不完全信息博弈，由包含两个具有相同行动 \(a\) 和 \(b\) 的决策节点的信息集表示。和往常一样，在树的叶节点处，顶部收益属于参与者 \(I\) ，底部收益属于参与者II。
\end{tcolorbox}

When the large firm is weak, the startup prefers to stay in the market over selling out to the large firm. When the large firm is strong, the opposite is true, and the startup is better off by selling out.

\begin{tcolorbox}
当大型公司弱小的时候，初创公司更愿意留在市场而不是被大型公司收购。当大型公司强大的时候，情况则相反，初创公司被收购会更好。
\end{tcolorbox}

Figure 10.1 shows a game tree with imperfect information (also called an "extensive game") that models this situation. From the perspective of the startup firm, it is random whether the large firm is weak or strong. This is modeled by an initial chance move, here assumed to have probability \(\frac{1}{2}\) for both possibilities.

\begin{tcolorbox}
图10.1展示了一个模拟这种情况的不完全信息博弈树(也称为“扩展式博弈”)。从初创公司的角度来看，大型公司是弱小还是强大是随机的。这通过一个初始的机会行动来建模，这里假设两种可能性的概率均为 \(\frac{1}{2}\) 。
\end{tcolorbox}

When the large firm is weak, it can choose to cede the market to the startup, here written as move \(T\) , with immediate payoffs 0 and 16 to players I and II. It can also announce a competing product, modeled by move \(B\) , in the hope that the startup firm, player II, will sell out by choosing move \(a\) , with payoffs 12 and 4 to the two players. However, if player II then decides instead to stay in the market (move \(b\) ), then it will even profit from the increased publicity and get payoff 20, with a negative payoff -4 to player I.

\begin{tcolorbox}
当大型公司弱小的时候，它可以选择将市场让给初创公司，这里记为行动 \(T\) ，参与者I和II的即时收益分别为0和16。它也可以宣布推出竞争产品，用行动 \(B\) 表示，希望初创公司(参与者II)选择行动 \(a\) 被收购，此时两个参与者的收益分别为12和4。然而，如果参与者II随后决定留在市场(行动 \(b\) )，那么它甚至会从增加的知名度中获利，获得收益20，而参与者I的收益为 -4。
\end{tcolorbox}

In contrast, when player I is strong, then it can again cede the market to player II (move \(C\) , with the same payoff pair 0,16 as before), or instead announce its own product (move \(D\) ). It will quickly become apparent that player I should never choose \(C\) , but we keep it as a possible move. After move \(D\) of the "strong" player I, the payoffs to the two players are 12, 4 if the startup sells out (move \(a\) ) and \({20}, - 4\) if the startup stays in (move \(b\) ).

\begin{tcolorbox}
相反，当参与者I实力较强时，它既可以再次将市场让给参与者II(选择行动 \(C\) ，此时的收益组合与之前一样，为0,16)，也可以宣布推出自己的产品(选择行动 \(D\) )。显然，参与者I绝不应该选择 \(C\) ，但我们仍将其视为一种可能的行动。在“实力较强”的参与者I选择行动 \(D\) 之后，如果初创企业出售(选择行动 \(a\) )，那么两个参与者的收益为12, 4；如果初创企业继续留在市场(选择行动 \(b\) )，收益为 \({20}, - 4\) 。
\end{tcolorbox}

As an additional feature compared to a game tree with perfect information, some nodes of the players are enclosed by ovals that define information sets. An information set is a set of decision nodes, subject to certain restrictions that we will shortly describe in generality. The interpretation is that a player cannot distinguish among the nodes in an information set, given his knowledge at the time he makes the move. Because his knowledge at all nodes in an information set is the same, he makes the same choice at each node in that set. Here, the startup firm, player II, must choose between move \(a\) (sell out) and move \(b\) (stay in the market). These are the two choices at player II's information set, which has two nodes according to the different histories of play, which player II cannot distinguish.

\begin{tcolorbox}
与具有完美信息的博弈树相比，还有一个额外的特征:参与者的某些节点被椭圆圈起来，这些椭圆定义了信息集。信息集是一组决策节点，需满足我们稍后将详细描述的特定限制条件。其含义是，在做出行动时，参与者根据自己当时的所知，无法区分信息集中的各个节点。由于信息集中所有节点上参与者的所知相同，因此他在该集合的每个节点上都会做出相同的选择。在这里，初创企业，即参与者II，必须在行动 \(a\) (出售)和行动 \(b\) (留在市场)之间做出选择。这是参与者II在其信息集上的两种选择，根据不同的博弈历史，该信息集有两个节点，而参与者II无法区分这些节点。
\end{tcolorbox}

Because player II is not informed about her position in a play of the game, backward induction can no longer be applied. If player II knew whether player I is weak or strong, then it would be better for her to choose \(b\) when player \(\mathrm{I}\) is weak, and to choose \(a\) when player \(\mathrm{I}\) is strong.

\begin{tcolorbox}
由于参与者II在博弈过程中不清楚自己所处的位置，因此逆向归纳法不再适用。如果参与者II知道参与者I是实力较弱还是较强，那么当参与者 \(\mathrm{I}\) 实力较弱时，她选择 \(b\) 会更好；当参与者 \(\mathrm{I}\) 实力较强时，她选择 \(a\) 会更好。
\end{tcolorbox}

Player I also has information sets, namely one for each of his decision nodes. Each information set contains only a single node, which shows that the player is perfectly informed about the state of play. We do not draw these singleton information sets in the game tree; they are implicit when a decision node is shown without being surrounded by an oval.

\begin{tcolorbox}
参与者I也有信息集，即每个决策节点对应一个信息集。每个信息集仅包含一个节点，这表明参与者对博弈状态有完美的了解。在博弈树中，我们不会画出这些单节点信息集；当一个决策节点没有被椭圆圈起来时，就隐含了这样的单节点信息集。
\end{tcolorbox}

One "optimal move" that can be seen directly from the game tree is move \(D\) when player I is strong. Choosing \(C\) gives player I payoff zero, whereas \(D\) gives him either 12 or 20, depending on whether player II chooses \(a\) or \(b\) , which are both larger than zero, so player I should choose \(D\) at any rate. Then the game reduces to a decision between move \(T\) and \(B\) for player I, and between \(a\) and \(b\) for player II.

\begin{tcolorbox}
从博弈树中可以直接看出的一个“最优行动”是，当参与者I实力较强时选择行动 \(D\) 。选择 \(C\) 会使参与者I的收益为零，而选择 \(D\) ，根据参与者II选择 \(a\) 还是 \(b\) ，他的收益要么是12，要么是20，这两个值都大于零，所以参与者I无论如何都应该选择 \(D\) 。这样，博弈就简化为参与者I在行动 \(T\) 和 \(B\) 之间进行决策，参与者II在 \(a\) 和 \(b\) 之间进行决策。
\end{tcolorbox}

The game tree shows that no pair of these moves for the two players is optimal against the other. For player I, if player II chooses \(b\) , then move \(T\) is better (because 0 is a higher payoff than -4 ). If player I chooses \(T\) , then the left node of the information set of player II is not reached. In that case, when player II has the choice between \(a\) and \(b\) , she knows with certainty that she is at the right node of her information set. Hence, she gets the higher payoff 4 when choosing \(a\) (rather than -4 when choosing \(b\) ). In turn, if player II chooses \(a\) , then even a weak player I will choose \(B\) over \(T\) with its higher payoff 12 compared to 0 . If player I chooses \(B\) , then player II will be with equal probability at either the left or the right node in her information set. Player II then has to use expected payoffs when comparing \(a\) and \(b\) . Clearly, move \(a\) will give her payoff 4, whereas move \(b\) will give her the expected payoff \(\frac{1}{2} \cdot  {20} + \frac{1}{2} \cdot  \left( {-4}\right)  = 8\) . That expected payoff exceeds the payoff for move \(a\) . That is, against \(B\) , player II is better off choosing \(b\) . We have now come full circle: The sequence of best-response moves for the players is \(b \rightarrow  T \rightarrow  a \rightarrow  B \rightarrow  b\) , so there is no deterministic behavior of the players that defines an equilibrium.

\begin{tcolorbox}
博弈树显示，这两个参与者的任何一对行动都不是针对对方的最优行动。对于参与者I来说，如果参与者II选择 \(b\) ，那么行动 \(T\) 更好(因为收益0比 -4 高)。如果参与者I选择 \(T\) ，那么参与者II信息集的左节点就不会被触及。在这种情况下，当参与者II在 \(a\) 和 \(b\) 之间进行选择时，她可以确定自己处于其信息集的右节点。因此，她选择 \(a\) 时能获得较高的收益4(而选择 \(b\) 时收益为 -4 )。反过来，如果参与者II选择 \(a\) ，那么即使是实力较弱的参与者I也会选择 \(B\) 而非 \(T\) ，因为前者的收益12高于后者的收益0。如果参与者I选择 \(B\) ，那么参与者II处于其信息集的左节点或右节点的概率相等。此时，参与者II在比较 \(a\) 和 \(b\) 时必须使用期望效用。显然，行动 \(a\) 会给她带来收益4，而行动 \(b\) 会给她带来期望效用 \(\frac{1}{2} \cdot  {20} + \frac{1}{2} \cdot  \left( {-4}\right)  = 8\) 。该期望效用超过了行动 \(a\) 的收益。也就是说，针对 \(B\) ，参与者II选择 \(b\) 会更好。现在我们又回到了原点:参与者的最优反应行动序列为 \(b \rightarrow  T \rightarrow  a \rightarrow  B \rightarrow  b\) ，因此参与者没有确定的行为来定义一个均衡。
\end{tcolorbox}

\begin{center}
\includegraphics[max width=0.6\textwidth]{images/bo_d2e5cmbef24c73b03tn0_278_189_194_847_629_0.jpg}
\end{center}
\hspace*{3em} 

Figure 10.2 Left: Strategic form of the game in Figure 10.1. Right: The \(2 \times  2\) game that results after eliminating the strictly dominated strategies \({TC}\) and \({BC}\) , with best-response payoffs.

\begin{tcolorbox}
图10.2 左:图10.1中博弈的策略形式。右:消除严格被占优策略 \(2 \times  2\) 和 \({BC}\) 后得到的博弈，带有最优反应效用。
\end{tcolorbox}

The game in Figure 10.1 can be analyzed in terms of the moves alone (we will give such an analysis at the end of this section), because for each player, only two moves matter. However, in general it is useful to consider the strategic form of the game that has already been defined for game trees with perfect information. This involves systematically considering the strategies of the players. The equilibrium concept is also defined in terms of strategies and not in terms of moves.

\begin{tcolorbox}
图10.1中的博弈可以仅从行动的角度进行分析(我们将在本节末尾给出这样的分析)，因为对于每个参与者来说，只有两种行动是重要的。然而，一般来说，考虑已经为具有完美信息的博弈树定义的博弈的策略形式是有用的。这涉及系统地考虑参与者的策略。均衡概念也是根据策略而不是行动来定义的。
\end{tcolorbox}

The left table in Figure 10.2 gives the strategic form of the game in Figure 10.1. As in the case of game trees with perfect information, a strategy of a player defines a move for each of the player's decision nodes. The difference, however, is that such a move is specified only once for all the decision nodes that belong to an information set, because by definition the choice of that move is the same no matter where the player is in the information set. In Figure 10.1, this means that player II has only two moves, which are the two moves \(a\) and \(b\) at her information set. Because player II only has one information set, her strategies are the same as her moves. Player I has two information sets, so a strategy for player I is a pair of moves, one move for the first and one move for the second information set. These pairs of moves are \({TC},{TD},{BC}\) , and \({BD}\) .

\begin{tcolorbox}
图10.2中的左表给出了图10.1中博弈的策略形式。与具有完美信息的博弈树的情况一样，参与者的策略为该参与者的每个决策节点定义了一个行动。然而，不同之处在于，对于属于一个信息集的所有决策节点，这样的行动只指定一次，因为根据定义，无论参与者在信息集中的哪个位置，该行动的选择都是相同的。在图10.1中，这意味着参与者II只有两种行动，即她在信息集处的两种行动 \(a\) 和 \(b\) 。因为参与者II只有一个信息集，所以她的策略与她的行动相同。参与者I有两个信息集，所以参与者I的一个策略是一对行动，一个是第一个信息集的行动，一个是第二个信息集的行动。这些行动对是 \({TC},{TD},{BC}\) 和 \({BD}\) 。
\end{tcolorbox}

The four strategies of player I correspond to the rows on the left in Figure 10.2, whereas the strategies \(a\) and \(b\) of player II correspond to the columns. Each of the eight cells in that table gives the expected payoffs as they are calculated from the payoffs at the leaves and the chance probabilities. For example, the strategy pair (BC, a) gives expected payoff \(\frac{1}{2} \cdot  {12} + \frac{1}{2} \cdot  0 = 6\) to player I and payoff \(\frac{1}{2} \cdot  4 + \frac{1}{2} \cdot  {16} = {10}\) to player II. These expectations come from the possible chance moves, both of which have probability \(\frac{1}{2}\) according to Figure 10.1. When the chance move goes to the left, the strategy pair(BC, a)means the play of the game follows move \(B\) and then move \(a\) and reaches the leaf with payoffs12,4to players I, II. When the chance move goes to the right, the strategy pair(BC, a)means that move \(C\) leads immediately to the leaf with payoffs 0,16 .

\begin{tcolorbox}
参与者I的四种策略对应于图10.2中左侧的行，而参与者II的策略 \(a\) 和 \(b\) 对应于列。该表中的八个单元格中的每一个都给出了根据叶节点的效用和机会概率计算出的期望效用。例如，策略对 (BC, a) 给参与者I带来期望效用 \(\frac{1}{2} \cdot  {12} + \frac{1}{2} \cdot  0 = 6\) ，给参与者II带来效用 \(\frac{1}{2} \cdot  4 + \frac{1}{2} \cdot  {16} = {10}\) 。这些期望来自可能的机会行动，根据图10.1，这两种机会行动的概率都是 \(\frac{1}{2}\) 。当机会行动向左时，策略对 (BC, a) 意味着博弈的进行遵循行动 \(B\) ，然后是行动 \(a\) ，并到达给参与者I、II带来效用12、4的叶节点。当机会行动向右时，策略对 (BC, a) 意味着行动 \(C\) 直接导致给参与者I、II带来效用0、16的叶节点。
\end{tcolorbox}

The \(4 \times  2\) game on the left in Figure 10.2 has two strictly dominated strategies, namely \({TC}\) , which is strictly dominated by \({TD}\) , and \({BC}\) , which is strictly dominated by \({BD}\) . This is in agreement with the earlier observation that for player I, move \(C\) is always worse than move \(D\) , no matter what player II does. However, one has to be careful to conclude that a "worse move" always leads to a strictly dominated strategy (the dominated strategy, for example \({TC}\) instead of \({TD}\) , would be obtained by replacing a given move, like \(D\) as part of the strategy \({TD}\) , by the "worse move", here \(C\) ). Namely, even if the payoffs from the "worse move" (here \(C\) ) are strictly worse than those resulting from the given move(D), one has to make sure that the decision point with the moves in question(CandD)is always reached with positive probability. Only then the expected payoffs from the two strategies will always differ. This is here the case because the decision point of player I with moves \(C\) and \(D\) is reached by a chance move. If, in contrast, the chance move was replaced by a move of a third player, for example, then if that third player does not move to the right, the choice between \(C\) and \(D\) would not affect the payoff to player I. So in this case, we could only conclude that \({TD}\) weakly dominates \({TC}\) . In short, it is safer to look at the strategic form of the game, even though a careful look at the game tree may eliminate some strategies right away.

\begin{tcolorbox}
图10.2左侧的 \(4 \times  2\) 博弈有两个严格被占优策略，即 \({TC}\) ，它被 \({TD}\) 严格占优，以及 \({BC}\) ，它被 \({BD}\) 严格占优。这与我们之前的观察一致，即对于参与者I来说，无论参与者II采取什么行动，行动 \(C\) 总是比行动 \(D\) 差。然而，我们必须谨慎地得出这样的结论:“较差的行动”总是会导致一个严格被占优策略(例如，用“较差的行动”(这里是 \(C\) )替换给定行动(如作为策略 \({TD}\) 一部分的 \(D\) )，就会得到被占优策略(如 \({TC}\) 而非 \({TD}\) )。也就是说，即使“较差的行动”(这里是 \(C\) )带来的收益严格低于给定行动(D)带来的收益，我们也必须确保涉及这些行动的决策点(C和D)总是以正概率被达到。只有这样，两种策略的期望效用才会始终不同。在这种情况下，由于参与者I选择行动 \(C\) 和 \(D\) 的决策点是通过一个随机行动达到的，所以情况就是如此。相反，如果这个随机行动被第三个参与者的行动所取代，例如，如果第三个参与者不向右行动，那么 \(C\) 和 \(D\) 之间的选择就不会影响参与者I的收益。因此，在这种情况下，我们只能得出 \({TD}\) 弱占优 \({TC}\) 的结论。简而言之，查看博弈的策略形式更为稳妥，尽管仔细观察博弈树可能会立即排除一些策略。
\end{tcolorbox}

Elimination of the strictly dominated strategies \({TC}\) and \({BC}\) gives the \(2 \times  2\) game on the right in Figure 10.2. There, each player has two strategies, which correspond to the two moves discussed earlier. The best-response payoffs indicated by boxes show that this game has no pure-strategy equilibrium.

\begin{tcolorbox}
剔除严格被占优策略 \({TC}\) 和 \({BC}\) 后，得到图10.2右侧的 \(2 \times  2\) 博弈。在这个博弈中，每个参与者有两种策略，这对应于我们之前讨论的两种行动。用方框表示的最优反应收益表明，这个博弈没有纯策略均衡。
\end{tcolorbox}

The game has an equilibrium in mixed strategies, which are determined in the familiar way. Player I randomizes with equal probability \(\frac{1}{2}\) between \({TD}\) and \({BD}\) so that the expected payoff to player II is 7 for both \(a\) and \(b\) . Player II chooses \(a\) with probability \(\frac{1}{4}\) and \(b\) with probability \(\frac{3}{4}\) so that the expected payoff to player I is \(\frac{36}{4} = 9\) for both \(T\) and \(B\) .

\begin{tcolorbox}
这个博弈存在混合策略均衡，其确定方式与我们熟悉的方式相同。参与者I以相等的概率 \(\frac{1}{2}\) 在 \({TD}\) 和 \({BD}\) 之间随机选择，使得参与者II选择 \(a\) 和 \(b\) 的期望效用均为7。参与者II以概率 \(\frac{1}{4}\) 选择 \(a\) ，以概率 \(\frac{3}{4}\) 选择 \(b\) ，使得参与者I选择 \(T\) 和 \(B\) 的期望效用均为 \(\frac{36}{4} = 9\) 。
\end{tcolorbox}

The game in Figure 10.1 is a competition between two firms that is more familiar when viewed as a simple case of the card game Poker. We can think of player I being dealt with equal probability a weak or a strong "hand" of cards. Player I knows whether his hand is weak or strong (in a real Poker game, player I knows his hand but not necessarily whether it is better or worse than that of player II). Player I can then either "fold" (moves \(T\) and \(C\) ), or "raise" (moves \(B\) and \(D\) ). In turn, player II can either "pass" (move a) or "meet". The "fold" and "pass" moves entail a fixed outcome where the players receive payoffs that are independent of the hands they hold. Only the combinations of moves "raise" and "meet", corresponding to the strategy pair(BD, b), lead to outcomes where the players' hands matter, as can be seen from Figure 10.1 (in the strategic form, different strategy pairs give different payoffs because of the chance move).

\begin{tcolorbox}
图10.1中的博弈是两家公司之间的竞争，当把它看作纸牌游戏扑克的一个简单例子时，我们会更熟悉。我们可以认为参与者I以相等的概率拿到一手弱牌或强牌。参与者I知道自己的牌是弱还是强(在真正的扑克游戏中，参与者I知道自己的牌，但不一定知道它比参与者II的牌好还是差)。然后，参与者I可以选择“弃牌”(行动 \(T\) 和 \(C\) )，或者“加注”(行动 \(B\) 和 \(D\) )。反过来，参与者II可以选择“放弃”(行动a)或“跟注”。“弃牌”和“放弃”行动会导致一个固定的结果，在这个结果中，参与者获得的收益与他们手中的牌无关。只有“加注”和“跟注”的行动组合，对应于策略组合(BD，b)，才会导致参与者手中的牌起作用的结果，如图10.1所示(在策略形式中，由于随机行动，不同的策略组合会带来不同的收益)。
\end{tcolorbox}

When player I has a weak hand and decides to "raise" (move \(B\) ), this amounts to bluffing, which means pretending that he has a strong hand, because player II does not know this. The purpose of bluffing, however, is not only to make player II "pass". Indeed, player II would not "pass" (move \(a\) ) if player I bluffed all the time, as analyzed before. One main effect of bluffing is to leave player II uncertain as to whether player I has a strong or weak hand. The best possible outcome for player I is if player II meets a strong hand of player I (moves \(D\) and \(b\) ), because this gives player I maximum payoff. If player I raises only if he has a strong hand, then player II would rather pass, which is not in the interest of player I. Above, we have reasoned that no deterministic combination of moves leads to an equilibrium. As described, this is made very vivid in terms of the Poker game interpretation.

\begin{tcolorbox}
当玩家I持有弱牌并决定“加注”(行动 \(B\) )时，这相当于虚张声势(bluff)，即假装他持有强牌，因为玩家II并不知晓这一点。然而，虚张声势的目的不仅是让玩家II“弃牌”。实际上，正如之前分析的，如果玩家I一直虚张声势，玩家II不会“弃牌”(行动 \(a\) )。虚张声势的一个主要效果是让玩家II不确定玩家I持有的是强牌还是弱牌。对玩家I来说，最好的结果是玩家II遇到玩家I的强牌(行动 \(D\) 和 \(b\) )，因为这会给玩家I带来最大收益。如果玩家I仅在持有强牌时加注，那么玩家II宁愿弃牌，这不符合玩家I的利益。综上所述，我们已经推断出，没有任何确定性的行动组合能导致均衡(equilibrium)。如前所述，在扑克(Poker)游戏的解释中，这一点表现得非常明显。
\end{tcolorbox}

To conclude the discussion of this example, we can also deduce in this case directly from the game tree the probabilities for the moves \(T\) and \(B\) , and for \(a\) and \(b\) , that are necessary in order to reach an equilibrium. Even without constructing the strategic form, we have seen that no deterministic choice of these moves gives an equilibrium. Hence, both players will randomize, and therefore must be indifferent between their choices. For player I, this indifference is very easy to determine: If he has to decide between \(T\) and \(B\) , then he gets the certain payoff 0 with \(T\) , so indifference results only when the expected payoff for \(B\) is also 0 . If player II chooses \(a\) and \(b\) with probabilities \(1 - q\) and \(q\) , respectively, that expected payoff to player \(\mathrm{I}\) is \({12}\left( {1 - q}\right)  - {4q}\) , so the equation \(0 = {12}\left( {1 - q}\right)  - {4q} = {12} - {16q}\) gives \(q = \frac{3}{4}\) .

\begin{tcolorbox}
为了结束对这个例子的讨论，我们还可以在这种情况下直接从博弈树(game tree)中推导出为了达到均衡，行动 \(T\) 和 \(B\) 以及 \(a\) 和 \(b\) 所需的概率。即使不构建策略形式(strategic form)，我们也已经看到，对这些行动的任何确定性选择都无法达到均衡。因此，两个玩家都会随机选择行动，从而必须对他们的选择无差异(indifference)。对于玩家I来说，这种无差异很容易确定:如果他必须在 \(T\) 和 \(B\) 之间做出决定，那么选择 \(T\) 他能获得确定的收益0，所以只有当 \(B\) 的期望效用(expected payoff)也为0时，才会出现无差异。如果玩家II分别以概率 \(1 - q\) 和 \(q\) 选择 \(a\) 和 \(b\) ，那么玩家 \(\mathrm{I}\) 的期望效用是 \({12}\left( {1 - q}\right)  - {4q}\) ，因此方程 \(0 = {12}\left( {1 - q}\right)  - {4q} = {12} - {16q}\) 得出 \(q = \frac{3}{4}\) 。
\end{tcolorbox}

In turn, player II has to be indifferent between her moves \(a\) and \(b\) in order to randomize on her part. Here the situation is somewhat more complicated, because the expected payoff to player II depends on the probabilities of being either at the left or at the right node in her information set. Assuming, as we do, that player I chooses \(D\) with certainty, the probability of being at the right node is \(\frac{1}{2}\) , whereas the probability of being at the left node is \(\frac{1}{2} \cdot  p\) , where \(p\) is the probability of player I choosing \(B\) . Because there is a nonzero probability \(1 - p\) for \(T\) , a play of the game ends with probability \(\frac{1}{2} \cdot  \left( {1 - p}\right)\) at the leftmost leaf of the tree, and then the information set of player II will not be reached at all. Technically, what has to be computed is the conditional probability that player II will be at the left node versus being at the right node. Before going into the details, we note that from the description so far, the probability for the left node is \(p\) times the probability for the right node. Only these relative weights matter. Indifference of player II between \(a\) and \(b\) then gives the following equation, with the weighted payoffs for \(a\) and \(b\) on the left and right, respectively:

\begin{tcolorbox}
反过来，玩家II为了随机选择行动，必须对她的行动 \(a\) 和 \(b\) 无差异。这里的情况稍微复杂一些，因为玩家II的期望效用取决于她在信息集(information set)中处于左节点(node)还是右节点的概率。假设，正如我们所做的，玩家I确定地选择 \(D\) ，那么处于右节点的概率是 \(\frac{1}{2}\) ，而处于左节点的概率是 \(\frac{1}{2} \cdot  p\) ，其中 \(p\) 是玩家I选择 \(B\) 的概率。因为 \(T\) 有非零概率 \(1 - p\) ，游戏以概率 \(\frac{1}{2} \cdot  \left( {1 - p}\right)\) 在博弈树的最左叶节点(leaf)结束，那么玩家II的信息集根本不会被触及。从技术上讲，需要计算的是玩家II处于左节点与处于右节点的条件概率(conditional probability)。在详细讨论之前，我们注意到，根据目前的描述，左节点的概率是右节点概率的 \(p\) 倍。只有这些相对权重才重要。玩家II在 \(a\) 和 \(b\) 之间的无差异会得出以下方程，左右两边分别是 \(a\) 和 \(b\) 的加权收益:
\end{tcolorbox}

\[
p \cdot  4 + 1 \cdot  4 = p \cdot  {20} + 1 \cdot  \left( {-4}\right) , \tag{10.1}
\]

which simplifies to \({4p} + 4 = {20p} - 4\) or \(8 = {16p}\) , that is, \(p = \frac{1}{2}\) .

\begin{tcolorbox}
该方程简化为 \({4p} + 4 = {20p} - 4\) 或 \(8 = {16p}\) ，即 \(p = \frac{1}{2}\) 。
\end{tcolorbox}

The conditional probabilities for player II being at the left and right node are computed as follows: The left node has absolute probability \(\frac{1}{2} \cdot  p\) of being reached, the right node absolute probability \(\frac{1}{2}\) . The total probability of being at any of the two nodes (which are disjoint events) is the sum of these two probabilities, namely \(\frac{p + 1}{2}\) . The conditional probability is then the absolute probability divided by the total probability, that is, for the left node, \(\frac{p/2}{\left( {p + 1}\right) /2}\) or \(\frac{p}{p + 1}\) , and for the right node, \(\frac{1/2}{\left( {p + 1}\right) /2}\) or \(\frac{1}{p + 1}\) . Using these conditional probabilities instead of the "weights" in (10.1) means dividing both sides of equation (10.1) by \(p + 1\) . This does not change the solution to this equation, so what we have done is perfectly in order (as long as the total probability of reaching the information set is nonzero).

\begin{tcolorbox}
玩家 II 位于左节点和右节点的条件概率计算如下:左节点被到达的绝对概率为 \(\frac{1}{2} \cdot  p\) ，右节点被到达的绝对概率为 \(\frac{1}{2}\) 。处于这两个节点(互斥事件)中任意一个的总概率是这两个概率之和，即 \(\frac{p + 1}{2}\) 。那么条件概率就是绝对概率除以总概率，也就是说，对于左节点，为 \(\frac{p/2}{\left( {p + 1}\right) /2}\) 或 \(\frac{p}{p + 1}\) ，对于右节点，为 \(\frac{1/2}{\left( {p + 1}\right) /2}\) 或 \(\frac{1}{p + 1}\) 。在 (10.1) 中使用这些条件概率而非“权重”，意味着将方程 (10.1) 的两边同时除以 \(p + 1\) 。这不会改变该方程的解，所以我们所做的完全合理(只要到达该信息集的总概率不为零)。
\end{tcolorbox}

One more time, we interpret the equilibrium probabilities \(p = \frac{1}{2}\) and \(q = \frac{3}{4}\) in terms of the game tree: Player I, by choosing \(B\) with probability \(\frac{1}{2}\) , and \(D\) with certainty, achieves that player II is half as likely to be at the left node of her information set as at her right node. Therefore, the conditional probabilities for left versus right node are \(\frac{1}{3}\) and \(\frac{2}{3}\) , respectively \(\left( \frac{1}{4}\right.\) and \(\frac{1}{2}\) divided by the total probability \(\frac{3}{4}\) of reaching the information set). Exactly in that case, player II is indifferent between \(a\) and \(b\) : The conditional expected payoffs to player II are then (clearly) 4 for move \(a\) , and \(\frac{1}{3} \cdot  {20} + \frac{2}{3} \cdot  \left( {-4}\right)  = \frac{12}{3} = 4\) for move \(b\) . The overall expected payoff to player II has to take the absolute probabilities into account: With probability \(\frac{1}{4}\) (left chance move, player \(\mathrm{I}\) choosing \(T\) ), she gets payoff 16, and with probability \(\frac{3}{4}\) the mentioned conditional payoff 4, overall \(\frac{1}{4} \cdot  {16} + \frac{3}{4} \cdot  4 = 4 + 3 = 7\) . For player I, we have observed the indifference between \(T\) and \(B\) , so when the chance move goes to the left, player I gets payoff zero. When the chance move goes to the right, player I when choosing \(D\) will get \(\left( {1 - q}\right)  \cdot  {12} + q \cdot  {20} = \frac{1}{4} \cdot  {12} + \frac{3}{4} \cdot  {20} = {18}\) . The overall payoff to player \(I\) is \(\frac{1}{2} \cdot  0 + \frac{1}{2} \cdot  {18} = 9\) .

\begin{tcolorbox}
再一次，我们从博弈树的角度来解释均衡概率 \(p = \frac{1}{2}\) 和 \(q = \frac{3}{4}\) :玩家 I 以概率 \(\frac{1}{2}\) 选择 \(B\) ，并确定地选择 \(D\) ，使得玩家 II 位于其信息集左节点的可能性是位于右节点可能性的一半。因此，左节点和右节点的条件概率分别为 \(\frac{1}{3}\) 和 \(\frac{2}{3}\) (即 \(\left( \frac{1}{4}\right.\) 和 \(\frac{1}{2}\) 分别除以到达该信息集的总概率 \(\frac{3}{4}\) )。恰好在这种情况下，玩家 II 在 \(a\) 和 \(b\) 之间无差异:那么玩家 II 选择行动 \(a\) 的条件期望效用显然为 4，选择行动 \(b\) 的条件期望效用为 \(\frac{1}{3} \cdot  {20} + \frac{2}{3} \cdot  \left( {-4}\right)  = \frac{12}{3} = 4\) 。玩家 II 的总体期望效用必须考虑绝对概率:以概率 \(\frac{1}{4}\) (左方随机行动，玩家 \(\mathrm{I}\) 选择 \(T\) )，她获得的效用为 16，以概率 \(\frac{3}{4}\) 获得上述条件效用 4，总体为 \(\frac{1}{4} \cdot  {16} + \frac{3}{4} \cdot  4 = 4 + 3 = 7\) 。对于玩家 I，我们已经观察到其在 \(T\) 和 \(B\) 之间无差异，所以当随机行动指向左方时，玩家 I 获得的效用为零。当随机行动指向右方时，玩家 I 选择 \(D\) 将获得 \(\left( {1 - q}\right)  \cdot  {12} + q \cdot  {20} = \frac{1}{4} \cdot  {12} + \frac{3}{4} \cdot  {20} = {18}\) 。玩家 \(I\) 的总体效用为 \(\frac{1}{2} \cdot  0 + \frac{1}{2} \cdot  {18} = 9\) 。
\end{tcolorbox}

\section*{10.3 Extensive Games}

\begin{tcolorbox}
\section*{10.3 扩展式博弈}
\end{tcolorbox}

In this section, we give the general definition of an extensive game, also known as a game in extensive form, or as a game tree with imperfect information.

\begin{tcolorbox}
在本节中，我们给出扩展式博弈的一般定义，扩展式博弈也被称为扩展形式的博弈，或具有不完美信息的博弈树。
\end{tcolorbox}

Similar to the game trees with perfect information considered in Chapter 4, the basic structure of an extensive game is a tree. As before, it has three types of nodes: Terminal nodes, called leaves, with a payoff for each player; non-terminal nodes which are chance nodes (typically drawn as squares), where the next node is chosen randomly according to a given probability that is specified in the tree; or non-terminal nodes that are decision nodes, which belong to a particular player. The outgoing edges from each decision node are labeled with the names of choices or moves which uniquely identify the respective edges.

\begin{tcolorbox}
与第 4 章中考虑的具有完美信息的博弈树类似，扩展式博弈的基本结构是一棵树。和之前一样，它有三种类型的节点:终端节点，也称为叶节点，为每个玩家提供一个效用；非终端节点中的机会节点(通常画成正方形)，下一个节点根据树中指定的给定概率随机选择；或者非终端节点中的决策节点，这些节点属于特定的玩家。从每个决策节点发出的边都标有选择或行动的名称，这些名称能唯一地标识相应的边。
\end{tcolorbox}

The decision nodes are partitioned into information sets. "Partitioned" means that every decision node belongs to exactly one information set, and that no information set is empty. When drawing the extensive game, information sets are shown as ovals around the nodes that they contain (except when the information set contains only one node); some people draw information sets by joining the nodes in the set with a single dashed line instead.

\begin{tcolorbox}
决策节点被划分为信息集。“划分”意味着每个决策节点恰好属于一个信息集，且没有信息集为空。在绘制扩展式博弈时，信息集用包含其中节点的椭圆表示(除非信息集仅包含一个节点)；有些人会用一条虚线连接信息集中的节点来表示信息集。
\end{tcolorbox}

The interpretation of an information set is that during a play of the game, a player is only told that he is at some node in an information set, but not at which particular node. Consequently, the information sets have to fulfill certain conditions: In an information set, all decision nodes belong to the same player; all decision nodes in the set have the same number of outgoing edges; and the set of moves (given by the labels on the outgoing edges) is the same for each node in the information set. In Figure 10.1, for example, the information set of player II has two moves, \(a\) and \(b\) . When the player has reached an information set, he chooses a move, which is by definition the same move no matter where he is in the information set.

\begin{tcolorbox}
信息集的含义是，在博弈过程中，玩家只被告知自己处于某个信息集中的某个节点，但不知道具体是哪个节点。因此，信息集必须满足某些条件:在一个信息集中，所有决策节点都属于同一个玩家；该集合中的所有决策节点具有相同数量的出边；并且信息集中每个节点的行动集合(由出边上的标签给出)是相同的。例如，在图10.1中，玩家II的信息集有两个行动， \(a\) 和 \(b\) 。当玩家到达一个信息集时，他选择一个行动，根据定义，无论他在信息集中的哪个位置，这个行动都是相同的。
\end{tcolorbox}

These conditions are trivially fulfilled if every information set is a singleton (contains only one node; recall that we usually do not draw these singleton information sets). If all information sets are singletons, then the game has perfect information and the information sets can as well be omitted. This gives a game tree with perfect information as considered earlier. Only for information sets that contain two or more nodes do these conditions matter.

\begin{tcolorbox}
如果每个信息集都是单元素集(仅包含一个节点；请记住，我们通常不绘制这些单元素信息集)，那么这些条件显然是满足的。如果所有信息集都是单元素集，那么该博弈具有完美信息，信息集也可以省略。这就得到了前面所考虑的具有完美信息的博弈树。只有对于包含两个或更多节点的信息集，这些条件才重要。
\end{tcolorbox}

\begin{center}
\includegraphics[max width=0.3\textwidth]{images/bo_d2e5cmbef24c73b03tn0_282_304_1143_487_539_0.jpg}
\end{center}
\hspace*{3em} 

Figure 10.3 Game with an information set where two nodes share a path, which is usually not allowed in an extensive game.

\begin{tcolorbox}
图10.3 一个信息集中两个节点共享一条路径的博弈，这在扩展式博弈中通常是不允许的。
\end{tcolorbox}

In addition to the above conditions, we will impose additional constraints on the structure of the information sets. The first of these is that no two nodes in an information set share a path, as in the example in Figure 10.3. The single-player game in that figure has two decision nodes of player \(\mathrm{I}\) , which belong to a single information set that has the two moves \(L\) and \(R\) . The only other non-terminal node is a chance move. Every condition about extensive games is met. However, the two decision nodes of player I share a path, because the first node is the root, from which there is a path to the second node. Information sets with this feature have a very problematic interpretation, because we require that a move, for example \(L\) , at such an information set is the same no matter where the player is. However, in Figure 10.3, what would it mean for player I to make move \(L\) ? It certainly means going left from the root, but if afterwards the chance move goes to the right so that player I can move again, does he automatically have to make move \(L\) again?

\begin{tcolorbox}
除了上述条件外，我们还将对信息集的结构施加额外的约束。其中第一个约束是，信息集中的任意两个节点不能共享一条路径，如图10.3中的示例所示。该图中的单人博弈有玩家 \(\mathrm{I}\) 的两个决策节点，它们属于一个具有两个行动 \(L\) 和 \(R\) 的信息集。唯一的其他非终端节点是一个机会行动。扩展式博弈的每个条件都得到了满足。然而，玩家I的两个决策节点共享一条路径，因为第一个节点是根节点，从该节点有一条路径通向第二个节点。具有这种特征的信息集的解释存在很大问题，因为我们要求在这样的信息集中，例如行动 \(L\) ，无论玩家处于何处都是相同的。然而，在图10.3中，玩家I采取行动 \(L\) 意味着什么呢？这肯定意味着从根节点向左移动，但如果之后机会行动向右移动，使得玩家I可以再次移动，他是否必须自动再次采取行动 \(L\) 呢？
\end{tcolorbox}

The problem with the information set in Figure 10.3 is that, by the interpretation of information sets, player I must have forgotten if he made a move at all when called to choose between \(L\) and \(R\) , and it is not even clear whether he is allowed to choose again or not. An information set can represent that a player forgets what he knew or did earlier. In that case, the player is said to have imperfect recall. However, players with imperfect recall do not conform to our idea of "rational" players. Moreover, it is not clear how often a player can move if one node in an information set is the descendant of another. We will therefore assume that all players have perfect recall, which is defined, and discussed in detail, in Section 10.6 below. In particular, we will see that perfect recall implies that no two nodes in an information set are connected by a path in the tree (see Proposition 10.2 below), so that a situation like in Figure 10.3 cannot occur.

\begin{tcolorbox}
图10.3中信息集的问题在于，根据信息集的解释，玩家I在被要求在 \(L\) 和 \(R\) 之间进行选择时，必须忘记自己是否采取过行动，甚至不清楚他是否被允许再次选择。信息集可以表示玩家忘记了他之前知道或做过的事情。在这种情况下，称玩家具有不完美记忆。然而，具有不完美记忆的玩家不符合我们对“理性”玩家的定义。此外，如果信息集中的一个节点是另一个节点的后代节点，那么不清楚玩家可以移动多少次。因此，我们将假设所有玩家都具有完美记忆，这将在下面的10.6节中进行定义和详细讨论。特别地，我们将看到完美记忆意味着信息集中的任意两个节点在树中不会通过一条路径相连(见下面的命题10.2)，因此像图10.3这样的情况不会发生。
\end{tcolorbox}

To summarize the definition of game trees with imperfect information: The only new concept is that of an information set, which models a player's lack of information about where he is in the game.

\begin{tcolorbox}
总结一下具有不完美信息的博弈树的定义:唯一的新概念是信息集，它模拟了玩家对自己在博弈中所处位置的信息缺失。
\end{tcolorbox}

\section*{10.4 Strategies for Extensive Games and the Strategic Form}

\begin{tcolorbox}
\section*{10.4 扩展式博弈的策略与策略形式}
\end{tcolorbox}

The definition of a strategy in an extensive game generalizes the definition of a strategy for a game tree with perfect information: A (pure) strategy of a player specifies a move for each information set of the player.

\begin{tcolorbox}
扩展式博弈中策略的定义是对具有完美信息的博弈树中策略定义的推广:玩家的(纯)策略为玩家的每个信息集指定一个行动。
\end{tcolorbox}

A formal definition of strategies requires some notation to refer to the information sets of a player, and to the moves available at each information set. Suppose player \(i\) is a player in the extensive game (like player I or player II in a two-player game). Then \({H}_{i}\) denotes the set of information sets of player \(i\) . A particular information set of player \(i\) is typically called \(h\) . We denote the set of choices or moves at the information set \(h\) by \({C}_{h}\) for \(h \in  {H}_{i}\) . For example, if \(h\) is the left information set of player I in Figure 10.1, then \({C}_{h} = \{ T,B\}\) .

\begin{tcolorbox}
策略的正式定义需要一些符号来表示玩家的信息集(information set)，以及每个信息集下可用的行动。假设玩家 \(i\) 是扩展式博弈中的一个玩家(就像两人博弈中的玩家 I 或玩家 II)。那么 \({H}_{i}\) 表示玩家 \(i\) 的信息集集合。玩家 \(i\) 的一个特定信息集通常称为 \(h\) 。我们用 \({C}_{h}\) 表示信息集 \(h\) 下的选择或行动集合，其中 \(h \in  {H}_{i}\) 。例如，如果 \(h\) 是图 10.1 中玩家 I 的左侧信息集，那么 \({C}_{h} = \{ T,B\}\) 。
\end{tcolorbox}

The set \({\sum }_{i}\) of strategies of player \(i\) is then

\begin{tcolorbox}
玩家 \(i\) 的策略集合 \({\sum }_{i}\) 为
\end{tcolorbox}

\[
{\sum }_{i} = \mathop{\sum }\limits_{{h \in  {H}_{i}}}{C}_{h} \tag{10.2}
\]

that is, the Cartesian product of the choice sets \({C}_{h}\) of player \(i\) . An element of \({\sum }_{i}\) is a tuple (vector) of moves, one move \(c\) from \({C}_{h}\) for each information set \(h\) of player \(i\) .

\begin{tcolorbox}
即玩家 \(i\) 的选择集 \({C}_{h}\) 的笛卡尔积。 \({\sum }_{i}\) 中的一个元素是一个行动元组(向量)，对于玩家 \(i\) 的每个信息集 \(h\) ，从 \({C}_{h}\) 中选取一个行动 \(c\) 。
\end{tcolorbox}

If \(N\) is the set of players, like \(N = \{ \mathrm{I},\mathrm{{II}}\}\) in the notation we use for a two-player game, then the set of pure strategy profiles is \(\mathop{\sum }\limits_{{i \in  N}}{\sum }_{i}\) , so a strategy profile specifies one pure strategy for each player.

\begin{tcolorbox}
如果 \(N\) 是玩家集合，就像我们在两人博弈符号中使用的 \(N = \{ \mathrm{I},\mathrm{{II}}\}\) ，那么纯策略组合的集合是 \(\mathop{\sum }\limits_{{i \in  N}}{\sum }_{i}\) ，因此一个策略组合为每个玩家指定了一个纯策略。
\end{tcolorbox}

\begin{center}
\includegraphics[max width=0.9\textwidth]{images/bo_d2e5cmbef24c73b03tn0_284_195_195_1206_498_0.jpg}
\end{center}
\hspace*{3em} 

Figure 10.4 \(\mathrm{A}3 \times  2\) game and an extensive game that has this game as its strategic form.

\begin{tcolorbox}
图 10.4 \(\mathrm{A}3 \times  2\) 博弈以及一个以该博弈为策略形式(strategic form)的扩展式博弈。
\end{tcolorbox}

The strategic form of the extensive game is given by the set \(N\) of players, the set of strategies \({\sum }_{i}\) for each player \(i\) , and the expected payoff that each player gets for each strategy profile. If the game has no chance moves, then the payoff for each strategy profile is given deterministically by the leaf of the game tree that is reached when the players use their strategies. If the game has chance moves, then more than one leaf may be reached with a given strategy profile, and then expected payoffs have to be calculated, as in the example of Figure 10.2.

\begin{tcolorbox}
扩展式博弈的策略形式由玩家集合 \(N\) 、每个玩家 \(i\) 的策略集合 \({\sum }_{i}\) 以及每个玩家在每个策略组合下获得的期望效用(expected payoff)给出。如果博弈没有随机行动，那么每个策略组合的收益由玩家使用其策略时到达的博弈树(game tree)的叶节点(leaf)确定。如果博弈有随机行动，那么给定一个策略组合可能到达多个叶节点，此时就需要计算期望效用，如图 10.2 的例子所示。
\end{tcolorbox}

With the help of information sets, it is possible to model simultaneous play in an extensive game. This is not possible with a game tree with perfect information where the sequential moves mean that the player who moves second is informed about the move of the first player.

\begin{tcolorbox}
借助信息集，可以在扩展式博弈中对同时行动进行建模。而对于具有完美信息(perfect information)的博弈树来说，这是不可能的，因为顺序行动意味着后行动的玩家知道先行动玩家的行动。
\end{tcolorbox}

In particular, given a game in strategic form, there is an extensive game that has exactly this strategic form. For example, consider an \(m \times  n\) game. It does not matter which player moves first, so we can assume this is player \(\mathrm{I}\) , who has \(m\) moves, say \({r}_{1},{r}_{2},\ldots ,{r}_{m}\) for the \(m\) rows of the game, at the root of the game tree. Each of these moves leads to a decision node of player II. All these nodes belong to a single information set. At this information set, player II has \(n\) moves \({c}_{1},{c}_{2},\ldots ,{c}_{n}\) that correspond to the \(n\) columns of the game. Each move of player II leads to a leaf of the game tree. Because \(m\) possible moves of player I are succeeded by \(n\) moves of player II, the game tree has \({mn}\) leaves. If the given game has payoffs \({a}_{ij}\) and \({b}_{ij}\) to players I and II when they play row \(i\) and column \(j\) , respectively, then these two payoffs are given at the leaf of the game tree that is reached by move \({r}_{i}\) followed by move \({c}_{j}\) . Because each player has only one information set, his strategies coincide with his moves. The strategic form of the constructed extensive game is then the original \(m \times  n\) game. Its rows are marked with the strategies \({r}_{1},\ldots ,{r}_{m}\) and its columns with \({c}_{1},\ldots ,{c}_{n}\) because these are the names of the moves of the two players in the extensive game. Figure 10.4 gives an example.

\begin{tcolorbox}
特别地，给定一个策略形式的博弈，存在一个扩展式博弈恰好具有该策略形式。例如，考虑一个 \(m \times  n\) 博弈。哪个玩家先行动并不重要，所以我们可以假设是玩家 \(\mathrm{I}\) 先行动，他有 \(m\) 个行动，比如对于博弈的 \(m\) 行，在博弈树的根节点处行动为 \({r}_{1},{r}_{2},\ldots ,{r}_{m}\) 。这些行动中的每一个都通向玩家 II 的一个决策节点(node)。所有这些节点都属于同一个信息集。在这个信息集下，玩家 II 有 \(n\) 个行动 \({c}_{1},{c}_{2},\ldots ,{c}_{n}\) ，对应于博弈的 \(n\) 列。玩家 II 的每个行动都通向博弈树的一个叶节点。因为玩家 I 的 \(m\) 个可能行动之后是玩家 II 的 \(n\) 个行动，所以博弈树有 \({mn}\) 个叶节点。如果给定的博弈中，玩家 I 和玩家 II 分别选择第 \(i\) 行和第 \(j\) 列时的收益为 \({a}_{ij}\) 和 \({b}_{ij}\) ，那么这两个收益在博弈树中通过行动 \({r}_{i}\) 之后接着行动 \({c}_{j}\) 到达的叶节点处给出。因为每个玩家只有一个信息集，所以他的策略与他的行动一致。所构建的扩展式博弈的策略形式就是原来的 \(m \times  n\) 博弈。其行用策略 \({r}_{1},\ldots ,{r}_{m}\) 标记，列用 \({c}_{1},\ldots ,{c}_{n}\) 标记，因为这些是扩展式博弈中两个玩家的行动名称。图 10.4 给出了一个例子。
\end{tcolorbox}

\(\Rightarrow\) Exercise 10.1 asks you to repeat this argument but also with player II moving first and player I moving second.

\begin{tcolorbox}
\(\Rightarrow\) 练习 10.1 要求你重复这个论证，但改为玩家 II 先行动，玩家 I 后行动。
\end{tcolorbox}

\section*{10.5 Reduced Strategies}

\begin{tcolorbox}
\section*{10.5 精简策略}
\end{tcolorbox}

Consider the extensive game in Figure 10.5. Player II has three information sets, call them \(h,{h}^{\prime }\) , and \({h}^{\prime \prime }\) , with \({C}_{h} = \{ l,r\} ,{C}_{{h}^{\prime }} = \{ a,b\}\) , and \({C}_{{h}^{\prime \prime }} = \{ c,d\}\) . Consequently, she has \(2 \times  2 \times  2 = 8\) pure strategies. Each pure strategy is a triple of moves, an element of \({\sum }_{\mathrm{{II}}} = {C}_{h} \times  {C}_{{h}^{\prime }} \times  {C}_{{h}^{\prime \prime }}\) , for example lac (we write the triple of moves by just writing the three moves next to each other for brevity). If player I chooses \(T\) , then this strategy lac of player II will lead to the leftmost leaf of the game tree with the payoff pair(1,5), via the moves \(l,T\) , and then \(a\) . Clearly, if player I plays \(T\) and player II plays the strategy lad, then this will lead to the same leaf. In fact, irrespective of the strategy chosen by player \(\mathrm{I}\) , the strategies \({lac}\) and \({lad}\) of player \(\mathrm{{II}}\) always have the same effect. This can be seen directly from the game tree, because when player II makes move \(l\) , she will never reach the information set \({h}^{\prime \prime }\) where she has to decide between \(c\) and \(d\) .

\begin{tcolorbox}
考虑图10.5中的扩展式博弈。玩家II有三个信息集，将它们称为 \(h,{h}^{\prime }\) 、 \({h}^{\prime \prime }\) ，其中 \({C}_{h} = \{ l,r\} ,{C}_{{h}^{\prime }} = \{ a,b\}\) 和 \({C}_{{h}^{\prime \prime }} = \{ c,d\}\) 。因此，她有 \(2 \times  2 \times  2 = 8\) 个纯策略。每个纯策略是一个行动三元组，是 \({\sum }_{\mathrm{{II}}} = {C}_{h} \times  {C}_{{h}^{\prime }} \times  {C}_{{h}^{\prime \prime }}\) 中的一个元素，例如lac(为简洁起见，我们通过将三个行动紧挨着写来表示这个三元组)。如果玩家I选择 \(T\) ，那么玩家II的这个策略lac将通过行动 \(l,T\) ，然后 \(a\) ，导致博弈树最左边的叶节点，收益对为(1,5)。显然，如果玩家I选择 \(T\) ，玩家II选择策略lad，那么这将导致相同的叶节点。事实上，无论玩家 \(\mathrm{I}\) 选择什么策略，玩家 \(\mathrm{{II}}\) 的策略 \({lac}\) 和 \({lad}\) 总是有相同的效果。这可以直接从博弈树中看出，因为当玩家II采取行动 \(l\) 时，她永远不会到达需要在 \(c\) 和 \(d\) 之间做出选择的信息集 \({h}^{\prime \prime }\) 。
\end{tcolorbox}

\begin{center}
\includegraphics[max width=0.6\textwidth]{images/bo_d2e5cmbef24c73b03tn0_285_522_1074_875_576_0.jpg}
\end{center}
\hspace*{3em} 

Figure 10.5 Extensive game with one information set for player I and three information sets \(h,{h}^{\prime },{h}^{\prime \prime }\) for player II. Note that \(h\) is a singleton and not drawn explicitly.

\begin{tcolorbox}
图10.5 玩家I有一个信息集，玩家II有三个信息集 \(h,{h}^{\prime },{h}^{\prime \prime }\) 的扩展式博弈。注意， \(h\) 是一个单元素集，未明确画出。
\end{tcolorbox}

This leads to the concept of a reduced strategy, which we have already met for a game tree with perfect information. A reduced strategy of player \(i\) specifies a move for each information set of player \(i\) , except for those information sets of player \(i\) that are unreachable because of an own earlier move (that is, because of a move of player \(i\) ). When writing down a reduced strategy, we list the moves for each information set (in a fixed order of the information sets), where an unspecified move is denoted by an extra symbol, the star "*", which stands for any move at that unreachable information set.

\begin{tcolorbox}
这引出了精简策略的概念，我们在具有完美信息(perfect information)的博弈树中已经遇到过这个概念。玩家 \(i\) 的精简策略为玩家 \(i\) 的每个信息集指定一个行动，但不包括那些由于玩家 \(i\) 自己先前的行动(即玩家 \(i\) 的一个行动)而无法到达的信息集。在写下精简策略时，我们列出每个信息集的行动(按信息集的固定顺序)，其中未指定的行动用一个额外的符号“*”表示，它代表在那个无法到达的信息集上的任何行动。
\end{tcolorbox}

In Figure 10.5, the information set \({h}^{\prime \prime }\) of player II is unreachable after her own earlier move \(l\) , and her information set \({h}^{\prime }\) is unreachable after her own earlier move \(r\) . Her reduced strategies are therefore \({la} * ,{lb} * ,r * c,r * d\) . Note that the star stands for any unspecified move that would normally be at that place when considering a full (unreduced) strategy. So player II's reduced strategy \({la} *\) can be thought of representing both her pure strategies lac and lad.

\begin{tcolorbox}
在图10.5中，玩家II的信息集 \({h}^{\prime \prime }\) 在她自己先前的行动 \(l\) 之后无法到达，她的信息集 \({h}^{\prime }\) 在她自己先前的行动 \(r\) 之后无法到达。因此，她的精简策略是 \({la} * ,{lb} * ,r * c,r * d\) 。注意，星号代表在考虑完整(未精简)策略时通常会在那个位置的任何未指定行动。所以玩家II的精简策略 \({la} *\) 可以被认为同时代表了她的纯策略lac和lad。
\end{tcolorbox}

The reduced strategic form of an extensive game is the same as its strategic form, except that for each player, only reduced strategies are considered and not pure strategies. By construction, the resulting expected payoffs for each player are well defined, because the unspecified moves never influence which leaves of the game tree are reached during play. Figure 10.6 gives the reduced strategic form of the game in Figure 10.5.

\begin{tcolorbox}
扩展式博弈的精简策略型(reduced strategic form)与它的策略形式(strategic form)相同，只是对于每个玩家，只考虑精简策略而不考虑纯策略。根据构造，每个玩家的期望效用(expected payoff)是明确定义的，因为未指定的行动永远不会影响在博弈过程中到达博弈树的哪些叶节点(leaf)。图10.6给出了图10.5中博弈的精简策略型。
\end{tcolorbox}

\begin{center}
\includegraphics[max width=0.4\textwidth]{images/bo_d2e5cmbef24c73b03tn0_286_197_828_618_361_0.jpg}
\end{center}
\hspace*{3em} 

Figure 10.6 Reduced strategic form of the game in Figure 10.5. In the strategies of player II, a star * indicates an unspecified move at an information set that is unreachable due to an own earlier move. As usual, the numbers in boxes are best-response payoffs.

\begin{tcolorbox}
图10.6 图10.5中博弈的精简策略型。在玩家II的策略中，星号“*”表示在由于自己先前的行动而无法到达的信息集上的未指定行动。像往常一样，框中的数字是最优反应收益。
\end{tcolorbox}

\section*{10.6 Perfect Recall}

\begin{tcolorbox}
\section*{10.6 完美记忆}
\end{tcolorbox}

What do the information sets of player II in Figure 10.5 describe? At the information set \(h\) , player II knows for certain that she is at the root of the game. At the information set \({h}^{\prime }\) , she does not know the move of player I, but she knows that she made the earlier move \(l\) at the root. At the information set \({h}^{\prime \prime }\) , she again does not know the move of player \(\mathrm{I}\) , but she knows that she made the move \(r\) earlier at the root.

\begin{tcolorbox}
图10.5中玩家II的信息集描述了什么？在信息集 \(h\) ，玩家II确定自己处于博弈的根节点。在信息集 \({h}^{\prime }\) ，她不知道玩家I的行动，但她知道自己在根节点做出了先前的行动 \(l\) 。在信息集 \({h}^{\prime \prime }\) ，她再次不知道玩家 \(\mathrm{I}\) 的行动，但她知道自己在根节点早些时候做出了行动 \(r\) 。
\end{tcolorbox}

For comparison, consider the game in Figure 10.7, which has the same tree and the same payoffs as the game in Figure 10.5, except that the information sets \({h}^{\prime }\) and \({h}^{\prime \prime }\) have been united to form a new set \({h}^{\prime \prime \prime }\) which has four nodes. Moreover, in order to conform with the definition of an extensive game, the moves at \({h}^{\prime \prime \prime }\) have to be the same at all nodes; they are here called \(a\) and \(b\) , which replace the earlier moves \(c\) and \(d\) at \({h}^{\prime \prime }\) . When player II is to make a move at \({h}^{\prime \prime \prime }\) , she knows nothing about the state of play except that she has made some move at the beginning, followed by either move \(T\) or \(B\) of player I. In particular, the interpretation of \({h}^{\prime \prime \prime }\) implies that player II forgot her earlier move \(l\) or \(r\) that she made at the root. We say that player II has imperfect recall in the game in Figure 10.7.

\begin{tcolorbox}
为作比较，考虑图10.7中的博弈，它与图10.5中的博弈有相同的博弈树和相同的收益，除了信息集 \({h}^{\prime }\) 和 \({h}^{\prime \prime }\) 合并形成了一个新的包含四个节点的信息集 \({h}^{\prime \prime \prime }\) 。此外，为了符合扩展式博弈的定义，在信息集 \({h}^{\prime \prime \prime }\) 处的行动在所有节点上必须相同；这里将它们称为 \(a\) 和 \(b\) ，它们取代了之前在信息集 \({h}^{\prime \prime }\) 处的行动 \(c\) 和 \(d\) 。当参与者II在信息集 \({h}^{\prime \prime \prime }\) 处进行行动时，她对博弈状态一无所知，只知道她在开始时采取了某个行动，随后是参与者I的行动 \(T\) 或 \(B\) 。特别地，信息集 \({h}^{\prime \prime \prime }\) 的解释意味着参与者II忘记了她在根节点处采取的早期行动 \(l\) 或 \(r\) 。我们说在图10.7的博弈中参与者II具有不完美记忆。
\end{tcolorbox}

Imperfect recall creates difficulties when interpreting the game, and contradicts our idea of "rational" players. All the extensive games that we will analyze have perfect recall, which we now proceed to define formally. For example, Figure 10.5 shows a game with perfect recall. There, player II has two information sets \({h}^{\prime }\) and \({h}^{\prime \prime }\) that have more than one node, so these are information sets where player II has some lack of information. We can test whether this lack of information involves her earlier moves. Namely, each node in an information set, say in \({h}^{\prime }\) , is reached by a unique path from the root of the game tree. Along that path, we look at the moves that the players make, which is move \(l\) of player II, followed by move \(T\) of player I for the left node of \({h}^{\prime }\) , or followed by move \(B\) for the right node of \({h}^{\prime }\) . (In general, this may also involve chance moves.) Among these moves, we ignore all moves made by the other players (which is here only player I) or by chance. Consequently, we see that each node in \({h}^{\prime }\) is preceded by move \(l\) of player II. That is, each node in the information set has the same sequence of own earlier moves. This is the definition of perfect recall.

\begin{tcolorbox}
不完美记忆在解释博弈时会产生困难，并且与我们对“理性”参与者的概念相矛盾。我们将分析的所有扩展式博弈都具有完美记忆，现在我们来正式定义它。例如，图10.5展示了一个具有完美记忆的博弈。在那里，参与者II有两个包含多个节点的信息集 \({h}^{\prime }\) 和 \({h}^{\prime \prime }\) ，所以这些是参与者II存在信息缺失的信息集。我们可以检验这种信息缺失是否涉及她的早期行动。即，信息集中的每个节点，比如信息集 \({h}^{\prime }\) 中的节点，都可以通过博弈树的根节点的唯一路径到达。沿着这条路径，我们观察参与者采取的行动，对于信息集 \({h}^{\prime }\) 的左节点，是参与者II的行动 \(l\) ，随后是参与者I的行动 \(T\) ；对于信息集 \({h}^{\prime }\) 的右节点，是参与者II的行动 \(l\) ，随后是行动 \(B\) 。(一般来说，这可能还涉及随机行动。)在这些行动中，我们忽略其他参与者(这里只有参与者I)或随机采取的所有行动。因此，我们看到信息集 \({h}^{\prime }\) 中的每个节点之前都有参与者II的行动 \(l\) 。也就是说，信息集中的每个节点都有相同的自身早期行动序列。这就是完美记忆的定义。
\end{tcolorbox}

Definition 10.1. Player \(i\) in an extensive game has perfect recall if for every information set \(h\) of player \(i\) , all nodes in \(h\) are preceded by the same sequence of moves of player \(i\) . The extensive game is said to have perfect recall if all its players have perfect recall.

\begin{tcolorbox}
定义10.1。在扩展式博弈中，如果对于参与者 \(i\) 的每个信息集 \(h\) ，信息集 \(h\) 中的所有节点之前都有相同的参与者 \(i\) 的行动序列，那么参与者 \(i\) 具有完美记忆。如果扩展式博弈的所有参与者都具有完美记忆，则称该扩展式博弈具有完美记忆。
\end{tcolorbox}

To repeat, the crucial condition about the nodes in an information set \(h\) of player \(i\) is that they are preceded by the same own earlier moves. We have to disregard the moves of the other players because there must be something that player \(i\) is lacking information about if player \(i\) does not have perfect information.

\begin{tcolorbox}
再次强调，参与者 \(i\) 的信息集 \(h\) 中节点的关键条件是它们之前都有相同的自身早期行动。我们必须忽略其他参与者的行动，因为如果参与者 \(i\) 不具有完美信息，那么参与者 \(i\) 必然在某些方面存在信息缺失。
\end{tcolorbox}

\(\Rightarrow\) Make sure you clearly understand the distinction between perfect information and perfect recall. How are the respective concepts defined?

\begin{tcolorbox}
\(\Rightarrow\) 确保你清楚理解完美信息和完美记忆之间的区别。各自的概念是如何定义的？
\end{tcolorbox}

\begin{center}
\includegraphics[max width=0.6\textwidth]{images/bo_d2e5cmbef24c73b03tn0_287_520_1518_881_578_0.jpg}
\end{center}
\hspace*{3em} 

Figure 10.7 Extensive game where player II does not have perfect recall.

\begin{tcolorbox}
图10.7 玩家II不具备完美记忆的扩展式博弈。
\end{tcolorbox}

We can verify that all players have perfect recall in Figure 10.5. We have already verified the condition of Definition 10.1 for \({h}^{\prime }\) . The other information sets of player II are \(h\) and \({h}^{\prime \prime }\) . For \(h\) , the condition holds trivially because \(h\) has only a single node (so we can make the easy observation that in a game with perfect information, that is, singleton information sets, all players have perfect recall). For \({h}^{\prime \prime }\) , both nodes in \({h}^{\prime \prime }\) are preceded by the sequence of own moves that consists of the single move \(r\) . This shows that player II has perfect recall. Player I has only one information set, and for both nodes in that set the sequence of own earlier moves is empty, because player I has not made a move before. It is useful to consider the empty sequence, often denoted by \(\varnothing\) like the empty set, as a sequence of moves, so that Definition 10.1 can be applied.

\begin{tcolorbox}
我们可以验证，在图10.5中所有玩家都具备完美记忆。我们已经验证了 \({h}^{\prime }\) 满足定义10.1的条件。玩家II的其他信息集是 \(h\) 和 \({h}^{\prime \prime }\) 。对于 \(h\) ，该条件显然成立，因为 \(h\) 仅包含一个节点(因此我们可以轻易观察到，在具有完美信息的博弈中，即信息集为单例集的情况下，所有玩家都具备完美记忆)。对于 \({h}^{\prime \prime }\) ， \({h}^{\prime \prime }\) 中的两个节点之前都有由单一行动 \(r\) 组成的自身行动序列。这表明玩家II具备完美记忆。玩家I只有一个信息集，且该信息集中的两个节点之前的自身早期行动序列均为空，因为玩家I之前没有采取任何行动。将通常用 \(\varnothing\) 表示的空序列(就像空集一样)视为一种行动序列是很有用的，这样就可以应用定义10.1。
\end{tcolorbox}

In Figure 10.7, player II does not have perfect recall because, for example, the leftmost node of her information set \({h}^{\prime \prime \prime }\) is preceded by her own earlier move \(l\) , whereas the rightmost node is preceded by the move \(r\) . So the sequences of own earlier moves are not the same for all nodes in \({h}^{\prime \prime \prime }\) .

\begin{tcolorbox}
在图10.7中，玩家II不具备完美记忆，因为例如，她的信息集 \({h}^{\prime \prime \prime }\) 中最左边的节点之前有她自己的早期行动 \(l\) ，而最右边的节点之前有行动 \(r\) 。因此， \({h}^{\prime \prime \prime }\) 中所有节点的自身早期行动序列并不相同。
\end{tcolorbox}

One way to understand Definition 10.1 is via the interpretation of information sets. They represent a player's lack of knowledge about the state of play, that is, at which of the nodes in the information set the player is at that moment. When the sequence of own earlier moves is the same for all these nodes, then the player can indeed not get additional information about where he is in the information set from knowing what moves he already played (as should normally be the case). If, on the other hand, some nodes in the information set are preceded by different own earlier moves, then the player can only fail to know where he is in the information set when he has forgotten these moves. For that reason, we say he has imperfect recall. So a player with perfect recall cannot have such information sets, because they would not accurately reflect his knowledge about the state of play.

\begin{tcolorbox}
理解定义10.1的一种方法是通过对信息集的解释。信息集代表玩家对博弈状态的认知缺失，即玩家在该时刻处于信息集中的哪个节点。当所有这些节点的自身早期行动序列相同时，玩家确实无法从已知自己已经采取的行动中获得关于自己在信息集中位置的额外信息(通常情况应如此)。另一方面，如果信息集中的某些节点之前有不同的自身早期行动，那么只有当玩家忘记了这些行动时，才会不知道自己在信息集中的位置。出于这个原因，我们说他具有不完美记忆。因此，具有完美记忆的玩家不会有这样的信息集，因为它们不能准确反映他对博弈状态的认知。
\end{tcolorbox}

However, Definition 10.1 is not about how to interpret the game. Perfect recall is a structural property of the information sets. This property is easily verified by considering for each information set the nodes therein and checking that they are reached by the same sequence of earlier own moves, which can be done without interpreting the game at all.

\begin{tcolorbox}
然而，定义10.1并非关于如何解释博弈。完美记忆是信息集的一种结构属性。通过考虑每个信息集中的节点，并检查它们是否通过相同的自身早期行动序列到达，就可以轻松验证这一属性，而完全无需对博弈进行解释。
\end{tcolorbox}

Figure 10.8 shows a game with only one player, who has an information set (with moves \(L\) and \(R\) ) that contains two nodes. Both nodes are preceded by an own earlier move named " \(B\) ", so it seems that player I has perfect recall in this game. However, this would obviously not be the case if the two singleton information sets had moves that are named differently, for example \(T\) and \(B\) at the left information set and \(C\) and \(D\) at the right information set (as in Figure 10.1). Indeed, player I does not have perfect recall in the game in Figure 10.8 because moves at different information sets are always considered distinct. A formal way of saying this is that the choice sets \({C}_{h}\) and \({C}_{{h}^{\prime }}\) are disjoint whenever \(h \neq  {h}^{\prime }\) . This can be assumed without loss of generality, either by re-naming the moves suitably, or by assuming that in addition to the name of the move that is indicated in the game tree, part of the "identity" of a move \(c \in  {C}_{h}\) is the information set \(h\) where that move is made.

\begin{tcolorbox}
图10.8展示了一个只有一个玩家的博弈，该玩家有一个信息集(包含行动 \(L\) 和 \(R\) )，其中包含两个节点。两个节点之前都有一个名为“ \(B\) ”的自身早期行动，因此在这个博弈中玩家I似乎具有完美记忆。然而，如果两个单例信息集的行动名称不同，情况显然就不是这样了，例如，左边信息集的行动为 \(T\) 和 \(B\) ，右边信息集的行动为 \(C\) 和 \(D\) (如图10.1所示)。实际上，在图10.8的博弈中玩家I不具有完美记忆，因为不同信息集的行动总是被视为不同的。用正式的说法是，只要 \(h \neq  {h}^{\prime }\) ，选择集 \({C}_{h}\) 和 \({C}_{{h}^{\prime }}\) 就是不相交的。不失一般性，可以通过适当地重新命名行动，或者假设除了博弈树中显示的行动名称外，行动 \(c \in  {C}_{h}\) 的“身份”部分还包括该行动所在的信息集 \(h\) ，来做出这样的假设。
\end{tcolorbox}

In the game in Figure 10.8, player I has imperfect recall, because he forgot his earlier knowledge, namely the outcome of the chance move. In contrast, player II in Figure 10.7 has not forgotten any earlier knowledge (all she knew earlier was that she makes the first move in the game, which she still knows when she reaches the information set \({h}^{\prime \prime \prime }\) ), but she has forgotten the move that she made. For that reason, it is sometimes said that a player has perfect recall if she never forgets what she "knew or did earlier". Definition 10.1 is a precise way of stating this.

\begin{tcolorbox}
在图10.8的博弈中，参与者I具有不完美记忆，因为他忘记了自己先前掌握的信息，即随机行动的结果。相比之下，图10.7中的参与者II没有忘记任何先前掌握的信息(她先前知道的所有信息就是她在博弈中率先行动，当她到达信息集 \({h}^{\prime \prime \prime }\) 时，她仍然知道这一点)，但她忘记了自己采取的行动。出于这个原因，有时人们会说，如果一个参与者永远不会忘记她“先前知道或做过的事情”，那么她就具有完美记忆。定义10.1精确地阐述了这一点。
\end{tcolorbox}

\(\Rightarrow\) Exercises 10.2 and 10.3 test your understanding of perfect recall.

\begin{tcolorbox}
\(\Rightarrow\) 练习10.2和10.3用于检验你对完美记忆的理解。
\end{tcolorbox}

Perfect recall implies that a situation like in Figure 10.3 cannot occur.

\begin{tcolorbox}
完美记忆意味着像图10.3中的情况不会发生。
\end{tcolorbox}

Proposition 10.2. If a player has perfect recall, then no two nodes in an information set of that player are connected by a path.

\begin{tcolorbox}
命题10.2:如果一个参与者具有完美记忆，那么该参与者的一个信息集中的任意两个节点都不会由一条路径相连。
\end{tcolorbox}

Proof. Suppose otherwise, that is, \(h\) is an information set of a player who has perfect recall, so that there are nodes in \(h\) where one node precedes the other on a path in the game tree. If there are several such nodes, let \(u\) be the earliest such node (that is, there is no other node in \(h\) from which there is a path to \(u\) ), and let \(v\) be a node in \(h\) so that there is a path from \(u\) to \(v\) . (An example is Figure 10.3 with \(u\) as the root of the game tree and \(v\) as the second node in player I’s information set.) On the path from \(u\) to \(v\) , the player makes a move \(c\) at the decision node \(u\) . This move \(c\) precedes \(v\) . However, there is no other node in \(h\) where move \(c\) precedes \(u\) , by the choice of \(u\) , so \(u\) is not preceded by \(c\) . Consequently, the sets of the player’s own moves that precede \(v\) and \(u\) are already different, so the sequences of own earlier moves for the two nodes certainly differ as well. So the player does not have perfect recall, which is a contradiction.

\begin{tcolorbox}
证明:假设情况并非如此，即 \(h\) 是一个具有完美记忆的参与者的信息集，使得在 \(h\) 中有一些节点，在博弈树中一条路径上其中一个节点先于另一个节点。如果有多个这样的节点，令 \(u\) 为最早的这样一个节点(也就是说，在 \(h\) 中没有其他节点存在一条通向 \(u\) 的路径)，并令 \(v\) 为 \(h\) 中的一个节点，使得存在一条从 \(u\) 到 \(v\) 的路径。(例如，在图10.3中， \(u\) 是博弈树的根节点， \(v\) 是参与者I的信息集中的第二个节点。)在从 \(u\) 到 \(v\) 的路径上，参与者在决策节点 \(u\) 采取行动 \(c\) 。这个行动 \(c\) 先于 \(v\) 。然而，根据 \(u\) 的选取，在 \(h\) 中没有其他节点使得行动 \(c\) 先于 \(u\) ，所以 \(u\) 并不先于 \(c\) 。因此，先于 \(v\) 和 \(u\) 的参与者自身行动的集合已经不同，所以这两个节点的先前自身行动序列肯定也不同。所以该参与者不具有完美记忆，这就产生了矛盾。
\end{tcolorbox}

\begin{center}
\includegraphics[max width=0.5\textwidth]{images/bo_d2e5cmbef24c73b03tn0_289_350_1462_651_684_0.jpg}
\end{center}
\hspace*{3em} 

Figure 10.8 Single-player game where player I does not have perfect recall. Moves at different information sets are always considered as distinct, even if they have the same name (for example, here \(B\) ).

\begin{tcolorbox}
图10.8 参与者I不具有完美记忆的单人博弈。不同信息集的行动总是被视为不同的，即使它们具有相同的名称(例如，这里的 \(B\) )。
\end{tcolorbox}

The examples in Figures 10.3 and 10.8 show that a game with a single player probably does not have perfect recall if it has information sets that contain two or more nodes.

\begin{tcolorbox}
图10.3和图10.8中的例子表明，如果单人博弈的信息集包含两个或更多节点，那么该博弈可能不具有完美记忆。
\end{tcolorbox}

\(\Rightarrow\) Construct a one-player game with perfect recall where not all information sets are singletons. Why do you need a chance move?

\begin{tcolorbox}
\(\Rightarrow\) 构建一个具有完美记忆的单人博弈，其中并非所有信息集都是单元素集。为什么你需要一个随机行动？
\end{tcolorbox}

\section*{10.7 Behavior Strategies}

\begin{tcolorbox}
\section*{10.7 行为策略}
\end{tcolorbox}

Our next topic concerns ways of randomizing in an extensive game. Unlike games with perfect information, which always have a equilibrium in pure strategies that can be found by backward induction, games with imperfect information may require randomized strategies, as the example in Figure 10.1 has demonstrated.

\begin{tcolorbox}
我们接下来的话题涉及扩展式博弈中的随机化方法。与具有完美信息的博弈(总是存在一个可以通过逆向归纳法找到的纯策略均衡)不同，具有不完美信息的博弈可能需要随机策略，如图10.1中的例子所示。
\end{tcolorbox}

\begin{center}
\includegraphics[max width=0.5\textwidth]{images/bo_d2e5cmbef24c73b03tn0_290_222_1007_711_598_0.jpg}
\end{center}
\hspace*{3em} 

Figure 10.9 Upper envelope of expected payoffs to player II for the game in Figure 10.6. The white and black dots mark equilibrium strategies of player I.

\begin{tcolorbox}
图10.9 图10.6中博弈的参与者II期望效用的上包络。白色和黑色圆点标记了参与者I的均衡策略。
\end{tcolorbox}

We find all equilibria of the game in Figure 10.5 that has the reduced strategic form shown in Figure 10.6. The only equilibrium in pure strategies is the strategy pair \(\left( {T,{la} * }\right)\) . In order to find further mixed equilibria, if any, we use the upper envelope method. Figure 10.9 shows the expected payoffs to player II as a function of the probability that player I chooses \(B\) . As this figure shows, the only pairs of strategies between which player II can be indifferent are \({la} *\) and \({lb} *\) when \(\operatorname{prob}\left( B\right)  = \frac{1}{3}\) ; or \({lb} *\) and \(r * c\) when \(\operatorname{prob}\left( B\right)  = \frac{1}{2}\) ; or \(r * c\) and \(r * d\) when \(\operatorname{prob}\left( B\right)  = \frac{3}{4}\) . The first of these cannot be played in an equilibrium because against both \({la} *\) and \({lb} *\) , the best response of player \(\mathrm{I}\) is \(T\) , so that player \(\mathrm{I}\) cannot be made indifferent between \(T\) and \(B\) , which is necessary for player I wanting to mix between \(T\) and \(B\) in equilibrium. In the other two cases, however, the pure best responses of player I differ, and it is possible to obtain equilibria. These equilibria are \(\left( {\left( {\frac{1}{2},\frac{1}{2}}\right) ,\left( {0,\frac{2}{5},\frac{3}{5},0}\right) }\right)\) and \(\left( {\left( {\frac{1}{4},\frac{3}{4}}\right) ,\left( {0,0,\frac{1}{3},\frac{2}{3}}\right) }\right)\) .

\begin{tcolorbox}
我们找出图10.5所示博弈的所有均衡，该博弈的精简策略型如图10.6所示。纯策略下唯一的均衡是策略对 \(\left( {T,{la} * }\right)\) 。为了找出更多的混合均衡(如果有的话)，我们使用上包络法。图10.9展示了玩家II的期望效用，它是玩家I选择 \(B\) 的概率的函数。如图所示，玩家II可能无差异的策略对仅为:当 \(\operatorname{prob}\left( B\right)  = \frac{1}{3}\) 时， \({la} *\) 和 \({lb} *\) ；当 \(\operatorname{prob}\left( B\right)  = \frac{1}{2}\) 时， \({lb} *\) 和 \(r * c\) ；当 \(\operatorname{prob}\left( B\right)  = \frac{3}{4}\) 时， \(r * c\) 和 \(r * d\) 。其中第一种情况不可能在均衡中出现，因为针对 \({la} *\) 和 \({lb} *\) ，玩家 \(\mathrm{I}\) 的最优反应都是 \(T\) ，所以玩家 \(\mathrm{I}\) 不可能在 \(T\) 和 \(B\) 之间保持无差异，而这是玩家I在均衡中想要在 \(T\) 和 \(B\) 之间进行混合的必要条件。然而，在另外两种情况下，玩家I的纯最优反应不同，此时有可能得到均衡。这些均衡是 \(\left( {\left( {\frac{1}{2},\frac{1}{2}}\right) ,\left( {0,\frac{2}{5},\frac{3}{5},0}\right) }\right)\) 和 \(\left( {\left( {\frac{1}{4},\frac{3}{4}}\right) ,\left( {0,0,\frac{1}{3},\frac{2}{3}}\right) }\right)\) 。
\end{tcolorbox}

What do these mixed strategies mean in terms of the extensive game? For player \(I\) , the mixed strategy \(\left( {\frac{1}{2},\frac{1}{2}}\right)\) means that he chooses \(T\) and \(B\) with equal probability. For player II, the mixed strategy \(\left( {0,\frac{2}{5},\frac{3}{5},0}\right)\) amounts to choosing \({lb} *\) with probability \(\frac{2}{5}\) , and \(r * c\) with probability \(\frac{3}{5}\) . The interpretation of such a mixed strategy is that each player chooses at the beginning of a play of the game a reduced pure strategy, which is a plan of action of what do at each reachable information set. The player then sticks to this plan of action throughout the play. In particular, if player II has chosen to play \({lb} *\) , then she first makes the move \(l\) with certainty, and then move \(b\) at her second information set \({h}^{\prime }\) . Her reduced pure strategy \(r * c\) is a similar plan of action.

\begin{tcolorbox}
在扩展式博弈中，这些混合策略意味着什么呢？对于玩家 \(I\) ，混合策略 \(\left( {\frac{1}{2},\frac{1}{2}}\right)\) 意味着他以相等的概率选择 \(T\) 和 \(B\) 。对于玩家II，混合策略 \(\left( {0,\frac{2}{5},\frac{3}{5},0}\right)\) 相当于以概率 \(\frac{2}{5}\) 选择 \({lb} *\) ，以概率 \(\frac{3}{5}\) 选择 \(r * c\) 。这种混合策略的解释是，每个玩家在博弈开始时选择一个精简纯策略，这是一个在每个可到达的信息集上的行动计划。然后玩家在整个博弈过程中都遵循这个行动计划。特别地，如果玩家II选择执行 \({lb} *\) ，那么她首先肯定会采取行动 \(l\) ，然后在她的第二个信息集 \({h}^{\prime }\) 采取行动 \(b\) 。她的精简纯策略 \(r * c\) 是一个类似的行动计划。
\end{tcolorbox}

There is a different way of playing such a mixed strategy, which is called a behavior strategy. A behavior strategy is a randomized choice of a move for each information set. That is, rather than using a lottery to determine a (reduced) pure strategy at the beginning of playing the game, the player uses a "local" lottery whenever she reaches an information set. In our example, player II's mixed strategy \(\left( {0,\frac{2}{5},\frac{3}{5},0}\right)\) is then played as follows: At her first information set \(h\) , she plays moves \(l\) and \(r\) with probabilities \(\frac{2}{5}\) and \(\frac{3}{5}\) , respectively. When she reaches her information set \({h}^{\prime }\) (which is only possible when she has made the earlier move \(l\) ), she chooses move \(a\) with probability zero and move \(b\) with probability one; this "behavior" at \({h}^{\prime }\) is deterministic. Similarly, when she reaches her information set \({h}^{\prime \prime }\) (having made move \(r\) earlier), then she chooses move \(c\) with certainty.

\begin{tcolorbox}
有一种不同的方式来执行这种混合策略，称为行为策略。行为策略是针对每个信息集对行动进行随机选择。也就是说，玩家不是在博弈开始时用彩票来确定一个(精简)纯策略，而是每当到达一个信息集时使用一个“局部”彩票。在我们的例子中，玩家II的混合策略 \(\left( {0,\frac{2}{5},\frac{3}{5},0}\right)\) 按如下方式执行:在她的第一个信息集 \(h\) ，她分别以概率 \(\frac{2}{5}\) 和 \(\frac{3}{5}\) 采取行动 \(l\) 和 \(r\) 。当她到达信息集 \({h}^{\prime }\) 时(只有在她之前采取了行动 \(l\) 时才有可能)，她以概率零选择行动 \(a\) ，以概率一选择行动 \(b\) ；在 \({h}^{\prime }\) 处的这种“行为”是确定的。类似地，当她到达信息集 \({h}^{\prime \prime }\) 时(之前采取了行动 \(r\) )，她肯定会选择行动 \(c\) 。
\end{tcolorbox}

Definition 10.3. In an extensive game, a behavior strategy \(\beta\) of player \(i\) is defined by a probability distribution on the set of moves \({C}_{h}\) for each information set \(h\) of player \(i\) . It is given by a probability \(\beta \left( c\right)\) for each \(c \in  {C}_{h}\) , that is, a number \(\beta \left( c\right)  \geq  0\) for each \(c \in  {C}_{h}\) so that \(\mathop{\sum }\limits_{{c \in  {C}_{h}}}\beta \left( c\right)  = 1\) . In a reduced behavior strategy, these numbers are left unspecified if \(h\) is unreachable because all earlier own moves by player \(i\) that allow her to reach \(h\) have probability zero under \(\beta\) .

\begin{tcolorbox}
定义10.3。在扩展式博弈中，参与者 \(i\) 的行为策略 \(\beta\) 由针对参与者 \(i\) 的每个信息集 \(h\) 的行动集合 \({C}_{h}\) 上的概率分布所定义。它由每个 \(c \in  {C}_{h}\) 的概率 \(\beta \left( c\right)\) 给出，即每个 \(c \in  {C}_{h}\) 对应一个数字 \(\beta \left( c\right)  \geq  0\) ，使得 \(\mathop{\sum }\limits_{{c \in  {C}_{h}}}\beta \left( c\right)  = 1\) 成立。在简化行为策略中，如果信息集 \(h\) 无法到达(因为在行为策略 \(\beta\) 下，参与者 \(i\) 所有能使其到达 \(h\) 的前期自身行动的概率均为零)，则这些数字不做具体规定。
\end{tcolorbox}

The concept of a reduced behavior strategy is illustrated by the other mixed equilibrium of the game in Figure 10.5, with the mixed strategy \(\left( {0,0,\frac{1}{3},\frac{2}{3}}\right)\) of player II. This is a randomization between the two pure strategies \(r * c\) and \(r * d\) . It can also be played as a (reduced) behavior strategy. Namely, at her first information set \(h\) , player II chooses \(l\) and \(r\) with probability 0 and 1, respectively. Then the information set \({h}^{\prime }\) is unreachable because it is only reachable when player II makes move \(l\) with positive probability, which is not the case here. Consequently, the behavior probabilities for making moves \(a\) and \(b\) are not specified, as in a reduced behavior strategy. (To obtain an "unreduced" behavior strategy, one would specify some arbitrary probabilities for \(a\) and \(b\) .) At her information set \({h}^{\prime \prime }\) , player II chooses \(c\) and \(d\) with probability \(\frac{1}{3}\) and \(\frac{2}{3}\) , respectively.

\begin{tcolorbox}
简化行为策略的概念可通过图10.5中博弈的另一个混合均衡(即参与者II的混合策略 \(\left( {0,0,\frac{1}{3},\frac{2}{3}}\right)\) )来说明。这是两种纯策略 \(r * c\) 和 \(r * d\) 之间的随机化选择。它也可以作为一种(简化)行为策略来执行。也就是说，在参与者II的第一个信息集 \(h\) 处，她选择 \(l\) 和 \(r\) 的概率分别为0和1。那么信息集 \({h}^{\prime }\) 就无法到达，因为只有当参与者II以正概率选择行动 \(l\) 时该信息集才可到达，而这里并非如此。因此，正如简化行为策略那样，选择行动 \(a\) 和 \(b\) 的行为概率未做具体规定。(若要得到一个“未简化”的行为策略，则需为 \(a\) 和 \(b\) 指定一些任意概率。)在她的信息集 \({h}^{\prime \prime }\) 处，参与者II选择 \(c\) 和 \(d\) 的概率分别为 \(\frac{1}{3}\) 和 \(\frac{2}{3}\) 。
\end{tcolorbox}

A behavior strategy can be played by "delaying" the player's random choice until she has reached an information set, when she makes the next move. However, nothing prevents the player making these random choices in advance for each information set, which will then give her a pure strategy that she plays throughout the game. Then the behavior strategy is used like a mixed strategy. In other words, behavior strategies can be considered as special mixed strategies.

\begin{tcolorbox}
行为策略可以通过“延迟”参与者的随机选择来执行，直到她到达某个信息集并做出下一步行动。然而，没有什么能阻止参与者提前为每个信息集做出这些随机选择，这样她就会得到一个在整个博弈过程中执行的纯策略。此时，行为策略的使用方式就如同混合策略。换句话说，行为策略可以被视为特殊的混合策略。
\end{tcolorbox}

A behavior strategy \(\beta\) of player \(i\) defines a mixed strategy \(\mu\) by defining a suitable probability \(\mu \left( \pi \right)\) for every pure strategy \(\pi\) of player \(i\) , which is found as follows. The pure strategy \(\pi\) defines a move \(\pi \left( h\right)\) for each information set \(h \in  {H}_{i}\) . When player \(i\) uses the behavior strategy \(\beta\) , that move has a certain probability \(\beta \left( {\pi \left( h\right) }\right)\) . Because these random moves at the information sets \(h\) of player \(i\) are made independently, the probability that all the moves specified by \(\pi\) are made is the product of these probabilities. That is, the mixed strategy probability is

\begin{tcolorbox}
参与者 \(i\) 的行为策略 \(\beta\) 通过为参与者 \(i\) 的每个纯策略 \(\pi\) 定义一个合适的概率 \(\mu \left( \pi \right)\) 来定义一个混合策略 \(\mu\) ，具体如下。纯策略 \(\pi\) 为每个信息集 \(h \in  {H}_{i}\) 定义了一个行动 \(\pi \left( h\right)\) 。当参与者 \(i\) 使用行为策略 \(\beta\) 时，该行动有一定的概率 \(\beta \left( {\pi \left( h\right) }\right)\) 。由于参与者 \(i\) 在信息集 \(h\) 处的这些随机行动是独立做出的，所以 \(\pi\) 所指定的所有行动都被执行的概率是这些概率的乘积。也就是说，混合策略的概率为
\end{tcolorbox}

\[
\mu \left( \pi \right)  = \mathop{\prod }\limits_{{h \in  {H}_{i}}}\beta \left( {\pi \left( h\right) }\right) . \tag{10.3}
\]

We have already used equation (10.3) when we interpreted a mixed strategy, for example \(\mu  = \left( {0,0,\frac{1}{3},\frac{2}{3}}\right)\) in Figure 10.5, as a behavior strategy \(\beta\) . We explain how this works even when, like here, the mixed strategy is defined with probabilities for reduced pure strategies. The behavior strategy \(\beta\) was given by \(\beta \left( l\right)  = 0,\beta \left( r\right)  = 1\) , \(\beta \left( c\right)  = \frac{1}{3}\) , and \(\beta \left( d\right)  = \frac{2}{3}\) , so that \(\frac{2}{3} = \mu \left( {r * d}\right)  = \beta \left( r\right)  \cdot  \beta \left( *\right)  \cdot  \beta \left( d\right)\) . This is expressed as a product of three behavior probabilities, one for each information set, because the pure strategies are given by a move for each information set; note that in (10.3), the product is taken over all information sets \(h\) of player \(i\) . Because we consider reduced pure strategies, the term \(\beta \left( *\right)\) represents the behavior probability for making any move at \({h}^{\prime }\) (in this case, as part of the reduced pure strategy \(r * d\) ), which is unspecified. Clearly, that probability is equal to 1 . That is, we take \(\beta \left( *\right)  = 1\) whenever (10.3) is applied to a reduced pure strategy \(\pi\) .

\begin{tcolorbox}
我们在将混合策略(例如图10.5中的 \(\mu  = \left( {0,0,\frac{1}{3},\frac{2}{3}}\right)\) )解释为行为策略 \(\beta\) 时，就已经使用了方程(10.3)。我们将解释即使像这里这样，混合策略是用精简纯策略的概率来定义的，这种解释是如何起作用的。行为策略 \(\beta\) 由 \(\beta \left( l\right)  = 0,\beta \left( r\right)  = 1\) 、 \(\beta \left( c\right)  = \frac{1}{3}\) 和 \(\beta \left( d\right)  = \frac{2}{3}\) 给出，使得 \(\frac{2}{3} = \mu \left( {r * d}\right)  = \beta \left( r\right)  \cdot  \beta \left( *\right)  \cdot  \beta \left( d\right)\) 。这被表示为三个行为概率的乘积，每个信息集(information set)对应一个概率，因为纯策略是由每个信息集的一个行动给出的；注意在(10.3)中，这个乘积是对玩家 \(i\) 的所有信息集 \(h\) 进行的。由于我们考虑的是精简纯策略，项 \(\beta \left( *\right)\) 表示在 \({h}^{\prime }\) 处采取任何行动的行为概率(在这种情况下，作为精简纯策略 \(r * d\) 的一部分)，这个概率是未明确指定的。显然，该概率等于1。也就是说，每当(10.3)应用于精简纯策略 \(\pi\) 时，我们取 \(\beta \left( *\right)  = 1\) 。
\end{tcolorbox}

The convention \(\beta \left( *\right)  = 1\) is easily seen to agree with the use of (10.3) for unreduced pure strategies. For example, consider the unreduced strategic form where \(r * d\) stands for the two pure strategies rad and rbd. Each of these has a certain probability under both \(\mu\) and \(\beta\) (when extended to the unreduced strategic form), where \(\mu\) is the representation of \(\beta\) as a mixed strategy. Whatever these probabilities are (obtained by splitting the probability for \(r * d\) into probabilities for rad and rbd), we certainly have

\begin{tcolorbox}
很容易看出，约定 \(\beta \left( *\right)  = 1\) 与(10.3)对未精简纯策略的使用是一致的。例如，考虑未精简的策略形式(strategic form)，其中 \(r * d\) 代表两个纯策略rad和rbd。在 \(\mu\) 和 \(\beta\) 下(当扩展到未精简的策略形式时)，这两个纯策略都有一定的概率，其中 \(\mu\) 是 \(\beta\) 作为混合策略的表示。无论这些概率是多少(通过将 \(r * d\) 的概率拆分为rad和rbd的概率得到)，我们肯定有
\end{tcolorbox}

\[
\mu \left( {rad}\right)  + \mu \left( {rbd}\right)  = \beta \left( r\right) \beta \left( a\right) \beta \left( d\right)  + \beta \left( r\right) \beta \left( b\right) \beta \left( d\right)  = \beta \left( r\right)  \cdot  \left( {\beta \left( a\right)  + \beta \left( b\right) }\right)  \cdot  \beta \left( d\right) .
\]

In this equation, the left-hand side is the proper interpretation of the mixed strategy probability \(\mu \left( {r * d}\right)\) , and the right-hand side of the probability \(\beta \left( r\right)  \cdot  \beta \left( *\right)  \cdot  \beta \left( d\right)\) , in agreement with \(\beta \left( *\right)  = \beta \left( a\right)  + \beta \left( b\right)  = 1\) .

\begin{tcolorbox}
在这个方程中，左边是混合策略概率 \(\mu \left( {r * d}\right)\) 的正确解释，右边是概率 \(\beta \left( r\right)  \cdot  \beta \left( *\right)  \cdot  \beta \left( d\right)\) 的正确解释，这与 \(\beta \left( *\right)  = \beta \left( a\right)  + \beta \left( b\right)  = 1\) 一致。
\end{tcolorbox}

\begin{center}
\includegraphics[max width=0.3\textwidth]{images/bo_d2e5cmbef24c73b03tn0_293_424_220_492_428_0.jpg}
\end{center}
\hspace*{3em} 

Figure 10.10 Game with four pure strategies of player II. Here, not all mixed strategies of player II are representable as behavior strategies.

\begin{tcolorbox}
图10.10 玩家II有四个纯策略的博弈。在此博弈中，并非玩家II的所有混合策略都可以表示为行为策略。
\end{tcolorbox}

While every behavior strategy can be considered as a mixed strategy, the converse does not hold. Figure 10.10 shows a game where player II has the four pure strategies \({ac},{ad},{bc},{bd}\) , which are also her reduced pure strategies. Consider the mixed strategy \(\mu  = \left( {\frac{1}{2},0,0,\frac{1}{2}}\right)\) of player II, where she chooses the pure strategy \({ac}\) , that is, moves \(a\) and \(c\) , with probability \(\frac{1}{2}\) , otherwise the pure strategy \({bd}\) . Clearly, every move is made with positive probability overall (in fact, probability \(\frac{1}{2}\) for each single move). Then if (10.3) was true, every move combination would have to have positive probability as well! However, this is not the case, because the strategy ad, for example, has probability zero under \(\mu\) . The reason is that when player II uses \(\mu\) , the random choice between the moves \(a\) and \(b\) is correlated with the random choice between \(c\) and \(d\) . That is, knowing that move \(a\) has been made gives some information about whether move \(c\) or \(d\) has been made (here, the correlation is maximal: move \(a\) implies move \(c\) ). If these random choices were made independently, as required in a behavior strategy, then the random choice between \(a\) and \(b\) would not affect the random choice between \(c\) and \(d\) .

\begin{tcolorbox}
虽然每个行为策略都可以被视为一个混合策略，但反之则不成立。图10.10展示了一个博弈，其中玩家II有四种纯策略 \({ac},{ad},{bc},{bd}\) ，这些也是她的简化纯策略。考虑玩家II的混合策略 \(\mu  = \left( {\frac{1}{2},0,0,\frac{1}{2}}\right)\) ，她以概率 \(\frac{1}{2}\) 选择纯策略 \({ac}\) ，即行动 \(a\) 和 \(c\) ，否则选择纯策略 \({bd}\) 。显然，每个行动总体上都有正概率(实际上，每个单独行动的概率为 \(\frac{1}{2}\) )。那么，如果(10.3)成立，每个行动组合也必须有正概率！然而，事实并非如此，因为例如策略ad在 \(\mu\) 下的概率为零。原因是当玩家II使用 \(\mu\) 时，行动 \(a\) 和 \(b\) 之间的随机选择与 \(c\) 和 \(d\) 之间的随机选择是相关的。也就是说，知道已经采取了行动 \(a\) 会提供一些关于是否采取了行动 \(c\) 或 \(d\) 的信息(这里，相关性是最大的:行动 \(a\) 意味着行动 \(c\) )。如果这些随机选择是独立进行的，就像行为策略所要求的那样，那么 \(a\) 和 \(b\) 之间的随机选择就不会影响 \(c\) 和 \(d\) 之间的随机选择。
\end{tcolorbox}

In this example, the "behavior" of player II at the left information set in Figure 10.10 described by \(\mu\) is to choose \(a\) with probability \(\frac{1}{2}\) and \(b\) with probability \(\frac{1}{2}\) , and ditto for moves \(c\) and \(d\) . This defines a behavior strategy \(\beta\) of player II. This behavior strategy defines the mixed strategy \(\left( {\frac{1}{4},\frac{1}{4},\frac{1}{4},\frac{1}{4}}\right)\) because it chooses each pure strategy with equal probability. This mixed strategy is clearly different from \(\mu\) .

\begin{tcolorbox}
在这个例子中，图10.10中左信息集(information set)处玩家II的“行为”由 \(\mu\) 描述，是以概率 \(\frac{1}{2}\) 选择 \(a\) ，以概率 \(\frac{1}{2}\) 选择 \(b\) ，行动 \(c\) 和 \(d\) 也是如此。这定义了玩家II的一个行为策略(behavior strategy) \(\beta\) 。这个行为策略定义了混合策略(mixed strategy) \(\left( {\frac{1}{4},\frac{1}{4},\frac{1}{4},\frac{1}{4}}\right)\) ，因为它以相等的概率选择每个纯策略。这个混合策略显然与 \(\mu\) 不同。
\end{tcolorbox}

A second example for Figure 10.10 is the mixed strategy \({\mu }^{\prime } = \left( {\frac{1}{6},\frac{1}{3},\frac{1}{4},\frac{1}{4}}\right)\) . The corresponding behavior strategy \({\beta }^{\prime }\) has move probabilities \({\beta }^{\prime }\left( a\right)  = {\beta }^{\prime }\left( b\right)  = \frac{1}{2}\) , because \({\beta }^{\prime }\left( a\right)  = {\mu }^{\prime }\left( {ac}\right)  + {\mu }^{\prime }\left( {ad}\right)  = \frac{1}{6} + \frac{1}{3}\) . Similarly, \({\beta }^{\prime }\left( c\right)  = {\mu }^{\prime }\left( {ac}\right)  + {\mu }^{\prime }\left( {bc}\right)  = \frac{1}{6} + \frac{1}{4} = \frac{5}{12}\) and \({\beta }^{\prime }\left( d\right)  = {\mu }^{\prime }\left( {ad}\right)  + {\mu }^{\prime }\left( {bd}\right)  = \frac{1}{3} + \frac{1}{4} = \frac{7}{12}\) . Again, making these moves independently with \({\beta }^{\prime }\) results in different probabilities for the pure strategies, for example in \({\beta }^{\prime }\left( a\right)  \cdot  {\beta }^{\prime }\left( c\right)  = \frac{5}{24}\) for \({ac}\) , which is different from \({\mu }^{\prime }\left( {ac}\right)\) .

\begin{tcolorbox}
图10.10的第二个例子是混合策略 \({\mu }^{\prime } = \left( {\frac{1}{6},\frac{1}{3},\frac{1}{4},\frac{1}{4}}\right)\) 。相应的行为策略 \({\beta }^{\prime }\) 的行动概率为 \({\beta }^{\prime }\left( a\right)  = {\beta }^{\prime }\left( b\right)  = \frac{1}{2}\) ，因为 \({\beta }^{\prime }\left( a\right)  = {\mu }^{\prime }\left( {ac}\right)  + {\mu }^{\prime }\left( {ad}\right)  = \frac{1}{6} + \frac{1}{3}\) 。类似地， \({\beta }^{\prime }\left( c\right)  = {\mu }^{\prime }\left( {ac}\right)  + {\mu }^{\prime }\left( {bc}\right)  = \frac{1}{6} + \frac{1}{4} = \frac{5}{12}\) 和 \({\beta }^{\prime }\left( d\right)  = {\mu }^{\prime }\left( {ad}\right)  + {\mu }^{\prime }\left( {bd}\right)  = \frac{1}{3} + \frac{1}{4} = \frac{7}{12}\) 。同样，用 \({\beta }^{\prime }\) 独立地进行这些行动会导致纯策略的不同概率，例如 \({ac}\) 的概率为 \({\beta }^{\prime }\left( a\right)  \cdot  {\beta }^{\prime }\left( c\right)  = \frac{5}{24}\) ，这与 \({\mu }^{\prime }\left( {ac}\right)\) 不同。
\end{tcolorbox}

What we have done in this example is to take a mixed strategy \(\mu\) of a player and "observe" at each information set of that player the "behavior" resulting from \(\mu\) . This behavior defines a certain behavior strategy \(\beta\) . In a similar way, \({\beta }^{\prime }\) is obtained from \({\mu }^{\prime }\) . The behavior strategy may differ from the original mixed strategy because with the behavior strategy, random choices at different information sets are no longer correlated. However, what does not change are the probabilities of reaching the nodes of game tree, given some (pure or mixed) strategy of the other players. We say that the strategies \(\mu\) and \(\beta\) are realization equivalent; similarly, \({\mu }^{\prime }\) and \({\beta }^{\prime }\) are realization equivalent. This realization equivalence holds in any extensive game with perfect recall, and is the topic of the next section.

\begin{tcolorbox}
在这个例子中，我们所做的是取一个玩家的混合策略 \(\mu\) ，并在该玩家的每个信息集上“观察”由 \(\mu\) 产生的“行为”。这种行为定义了某种行为策略 \(\beta\) 。类似地， \({\beta }^{\prime }\) 是从 \({\mu }^{\prime }\) 得到的。行为策略可能与原始的混合策略不同，因为在行为策略中，不同信息集上的随机选择不再相关。然而，给定其他玩家的某些(纯或混合)策略，到达博弈树节点的概率是不变的。我们说策略 \(\mu\) 和 \(\beta\) 是实现等价的；类似地， \({\mu }^{\prime }\) 和 \({\beta }^{\prime }\) 也是实现等价的。这种实现等价性在任何具有完美记忆的扩展式博弈中都成立，这也是下一节的主题。
\end{tcolorbox}

\section*{10.8 Kuhn's Theorem: Behavior Strategies Suffice}

\begin{tcolorbox}
\section*{10.8 库恩定理:行为策略足矣}
\end{tcolorbox}

Information sets were introduced by Kuhn (1953). His article also defined perfect recall and behavior strategies, and proved an important connection between the two concepts. This connection is the central property of extensive games and is often called "Kuhn's theorem". It says that a player who has perfect recall can always replace a mixed strategy by a realization equivalent behavior strategy.

\begin{tcolorbox}
信息集是由库恩(1953 年)引入的。他的文章还定义了完美记忆和行为策略，并证明了这两个概念之间的重要联系。这种联系是扩展式博弈的核心性质，通常被称为“库恩定理”。该定理指出，具有完美记忆的玩家总是可以用一个实现等价的行为策略来替代混合策略。
\end{tcolorbox}

Why should we even care about whether a player can use a behavior instead of a mixed strategy? The reason is that a behavior strategy is much simpler to describe than a mixed strategy, significantly so for larger games. As an example, consider a simplified Poker game where each player is dealt only a single card, with 13 different possibilities for the rank of that card. A game tree for that game may start with a chance move with 13 possibilities for the card dealt to player \(\mathrm{I}\) , who learns the rank of his card and can then either "fold" or "raise". Irrespective of the rest of the game, this defines already \({2}^{13} = {8192}\) move combinations for player I. Assuming these are all the pure strategies of player I, a mixed strategy for player I would have to specify 8191 probabilities (which determine the remaining probability because these probabilities sum to one). In contrast, a behavior strategy would specify one probability (of folding, say) for each possible card, which are 13 different numbers. So a behavior strategy is much less complex to describe than a mixed strategy.

\begin{tcolorbox}
我们为什么要关心玩家是否可以使用行为策略而不是混合策略呢？原因是行为策略比混合策略更容易描述，对于较大的博弈来说更是如此。例如，考虑一个简化的扑克游戏，每个玩家只发一张牌，这张牌的点数有 13 种不同的可能性。该游戏的博弈树可能从一个机会行动开始，玩家 \(\mathrm{I}\) 拿到的牌有 13 种可能性，玩家 \(\mathrm{I}\) 知道自己牌的点数后，可以选择“弃牌”或“加注”。不管游戏的其余部分如何，这已经为玩家 I 定义了 \({2}^{13} = {8192}\) 种行动组合。假设这些是玩家 I 的所有纯策略，那么玩家 I 的混合策略必须指定 8191 个概率(因为这些概率之和为 1，所以这些概率决定了剩余的概率)。相比之下，行为策略会为每张可能的牌指定一个概率(例如弃牌的概率)，这是 13 个不同的数字。因此，行为策略比混合策略的描述要简单得多。
\end{tcolorbox}

This example also gives an intuition why a behavior strategy suffices. The moves at the different decision nodes of player I can be randomized independently because they concern disjoint parts of the game. There is no need to correlate the moves for different cards because in any one game, only one card is drawn.

\begin{tcolorbox}
这个例子也让我们直观地理解了为什么行为策略就足够了。玩家 I 在不同决策节点上的行动可以独立地随机化，因为它们涉及博弈的不相交部分。没有必要对不同牌的行动进行关联，因为在任何一局游戏中，只抽取一张牌。
\end{tcolorbox}

In general, a player may make moves in sequence and a priori there may be some gain in correlating an earlier move with a subsequent move. This is where the condition of perfect recall comes in. Namely, the knowledge about any earlier move is already captured by the information sets because perfect recall says that the player knows all his earlier moves. So there is no need to "condition" a later move on an earlier move because all earlier moves are known.

\begin{tcolorbox}
一般来说，玩家可能会按顺序进行行动，并且先验地，将早期行动与后续行动关联起来可能会有一些好处。这就是完美记忆条件发挥作用的地方。也就是说，关于任何早期行动的信息已经被信息集所包含，因为完美记忆意味着玩家知道自己所有的早期行动。所以没有必要将后期行动“基于”早期行动，因为所有早期行动都是已知的。
\end{tcolorbox}

For illustration, consider the game in Figure 10.7 which does not have perfect recall. It is easily seen that this game has the same strategic form, shown in Figure 10.6, as the game in Figure 10.5, except that the pure strategies of player II are \({la},{lb},{ra},{rb}\) instead of the reduced strategies \({la} * ,{lb} * ,r * c,r * d\) . Here, the mixed strategy \(\left( {0,\frac{2}{5},\frac{3}{5},0}\right)\) of player II, which selects the pure strategy \({lb}\) with probability \(\frac{2}{5}\) , and \({ra}\) with probability \(\frac{3}{5}\) , cannot be expressed as a behavior strategy. Namely, this behavior strategy would have to choose moves \(l\) and \(r\) with probability \(\frac{2}{5}\) and \(\frac{3}{5}\) , respectively, but the probabilities for moves \(a\) and \(b\) , which would be \(\frac{3}{5}\) and \(\frac{2}{5}\) , are not independent of the earlier move (because otherwise the pure strategy \({la}\) , for example, would not have probability zero). So this mixed strategy, which is part of an equilibrium, cannot be expressed as a behavior strategy. Using a mixed strategy gives player II more possibilities of playing in the game.

\begin{tcolorbox}
为了说明这一点，考虑图 10.7 中不具有完美记忆的博弈。很容易看出，除了玩家 II 的纯策略是 \({la},{lb},{ra},{rb}\) 而不是精简策略 \({la} * ,{lb} * ,r * c,r * d\) 之外，这个博弈与图 10.5 中的博弈具有相同的策略形式，如图 10.6 所示。在这里，玩家 II 的混合策略 \(\left( {0,\frac{2}{5},\frac{3}{5},0}\right)\) ，以概率 \(\frac{2}{5}\) 选择纯策略 \({lb}\) ，以概率 \(\frac{3}{5}\) 选择 \({ra}\) ，不能表示为行为策略。也就是说，这个行为策略必须分别以概率 \(\frac{2}{5}\) 和 \(\frac{3}{5}\) 选择行动 \(l\) 和 \(r\) ，但是行动 \(a\) 和 \(b\) 的概率(分别为 \(\frac{3}{5}\) 和 \(\frac{2}{5}\) )并不独立于早期行动(因为否则，例如纯策略 \({la}\) 的概率就不会为零)。所以这个作为均衡一部分的混合策略不能表示为行为策略。使用混合策略为玩家 II 在游戏中提供了更多的玩法可能性。
\end{tcolorbox}

In contrast, the mixed strategy \(\left( {0,\frac{2}{5},\frac{3}{5},0}\right)\) of player II in the perfect-recall game in Figure 10.5 chooses \({lb} *\) and \(r * c\) with probability \(\frac{2}{5}\) and \(\frac{3}{5}\) , respectively. As shown earlier, this is in effect a behavior strategy, because correlating move \(l\) with move \(b\) , and move \(r\) with move \(c\) , is done via the information sets \({h}^{\prime }\) and \({h}^{\prime \prime }\) , which are separate in Figure 10.5 but merged in Figure 10.7. That is, player II knows that she made move \(l\) when she has to make her move \(b\) at \({h}^{\prime }\) , and she knows that she made move \(r\) when she has to make her move \(c\) at \({h}^{\prime \prime }\) .

\begin{tcolorbox}
相比之下，图10.5所示的具有完美记忆的博弈中，玩家II的混合策略 \(\left( {0,\frac{2}{5},\frac{3}{5},0}\right)\) 分别以概率 \(\frac{2}{5}\) 和 \(\frac{3}{5}\) 选择 \({lb} *\) 和 \(r * c\) 。如前所述，这实际上是一种行为策略，因为将行动 \(l\) 与行动 \(b\) 相关联，以及将行动 \(r\) 与行动 \(c\) 相关联，是通过信息集 \({h}^{\prime }\) 和 \({h}^{\prime \prime }\) 实现的，这些信息集在图10.5中是分开的，但在图10.7中是合并的。也就是说，玩家II知道，当她在 \({h}^{\prime }\) 处必须做出行动 \(b\) 时，她之前做出了行动 \(l\) ；当她在 \({h}^{\prime \prime }\) 处必须做出行动 \(c\) 时，她之前做出了行动 \(r\) 。
\end{tcolorbox}

Before stating Kuhn's theorem, we have to define what it means to say that two randomized strategies \(\mu\) and \(\rho\) of a player are realization equivalent. This is the case if for any fixed strategies of the other players, every node of the game tree is reached ("realized") with the same probability when the player uses \(\mu\) as when he uses \(\rho\) . Two realization equivalent strategies \(\mu\) and \(\rho\) always give every player the same expected payoff because the leaves of the game tree are reached with the same probability no matter whether \(\mu\) or \(\rho\) is used.

\begin{tcolorbox}
在陈述库恩定理之前，我们必须定义一个玩家的两个随机策略 \(\mu\) 和 \(\rho\) 实现等价是什么意思。如果对于其他玩家的任何固定策略，当该玩家使用 \(\mu\) 时和使用 \(\rho\) 时，博弈树的每个节点被到达(“实现”)的概率相同，那么就称这两个策略是实现等价的。两个实现等价的策略 \(\mu\) 和 \(\rho\) 总是给每个玩家相同的期望效用，因为无论使用 \(\mu\) 还是 \(\rho\) ，博弈树的叶节点被到达的概率是相同的。
\end{tcolorbox}

When are two strategies \(\mu\) and \(\rho\) realization equivalent? Suppose these are strategies of player \(i\) . Consider a particular node, say \(u\) , of the game tree. Equivalence of \(\mu\) and \(\rho\) means that, given that the other players play in a certain way, the probability that \(u\) is reached is the same when player \(i\) uses \(\mu\) as when he uses \(\rho\) . In order to reach \(u\) , the play of the game has to make a certain sequence of moves, including chance moves, that lead to \(u\) . This move sequence is unique because we have a tree. In that sequence, some moves are made by chance or by players other than player \(i\) . These moves have certain given probabilities because the chance move probabilities and the strategies adopted by the other players are fixed. We can assume that the overall probability for these moves is not zero, by considering fixed strategies of the other players which allow reaching \(u\) (otherwise \(u\) is trivially reached with the same zero probability with \(\mu\) as with \(\rho\) ). So what matters in order to compare \(\mu\) and \(\rho\) are the moves by player \(i\) on the path to \(u\) . If the probability for playing the sequence of moves of player \(i\) on the path to \(u\) is the same under \(\mu\) and \(\rho\) , and if this holds for every node \(u\) , then \(\mu\) and \(\rho\) are realization equivalent. This is the key ingredient of the proof of Kuhn's theorem.

\begin{tcolorbox}
两个策略 \(\mu\) 和 \(\rho\) 在什么情况下是实现等价的呢？假设这些是玩家 \(i\) 的策略。考虑博弈树的一个特定节点，比如 \(u\) 。 \(\mu\) 和 \(\rho\) 等价意味着，在其他玩家以某种方式行动的情况下，当玩家 \(i\) 使用 \(\mu\) 时和使用 \(\rho\) 时，到达 \(u\) 的概率是相同的。为了到达 \(u\) ，博弈的进行必须做出一系列特定的行动，包括随机行动，这些行动会通向 \(u\) 。由于我们处理的是一棵树，所以这个行动序列是唯一的。在这个序列中，一些行动是由随机因素或玩家 \(i\) 之外的其他玩家做出的。这些行动有一定的给定概率，因为随机行动的概率和其他玩家采用的策略是固定的。通过考虑其他玩家允许到达 \(u\) 的固定策略，我们可以假设这些行动的总概率不为零(否则，无论使用 \(\mu\) 还是 \(\rho\) ， \(u\) 显然都以相同的零概率被到达)。因此，为了比较 \(\mu\) 和 \(\rho\) ，重要的是玩家 \(i\) 在通向 \(u\) 的路径上的行动。如果在 \(\mu\) 和 \(\rho\) 下，玩家 \(i\) 在通向 \(u\) 的路径上做出行动序列的概率相同，并且对于每个节点 \(u\) 都成立，那么 \(\mu\) 和 \(\rho\) 就是实现等价的。这是库恩定理证明的关键要素。
\end{tcolorbox}

Theorem 10.4 (Kuhn, 1953). If player \(i\) in an extensive game has perfect recall, then for any mixed strategy \(\mu\) of player \(i\) there is a realization equivalent behavior strategy \(\beta\) of player \(i\) .

\begin{tcolorbox}
定理10.4(库恩，1953年)。如果扩展式博弈中的玩家 \(i\) 具有完美记忆，那么对于玩家 \(i\) 的任何混合策略 \(\mu\) ，都存在一个实现等价的行为策略 \(\beta\) 。
\end{tcolorbox}

Proof. We illustrate the proof with the extensive game in Figure 10.11, with player II as player \(i\) . The payoffs are omitted; they are not relevant because only the probabilities of reaching the nodes of the game tree matter.

\begin{tcolorbox}
证明。我们用图10.11中的扩展式博弈来说明证明过程，将玩家II作为玩家 \(i\) 。支付被省略了；因为只有到达博弈树节点的概率是重要的，所以支付无关紧要。
\end{tcolorbox}

\begin{center}
\includegraphics[max width=0.4\textwidth]{images/bo_d2e5cmbef24c73b03tn0_296_247_590_587_757_0.jpg}
\end{center}
\hspace*{3em} 

Figure 10.11 Game that illustrates sequences of moves of player II, and the proof of Kuhn's Theorem 10.4. Payoffs are omitted.

\begin{tcolorbox}
图10.11 该博弈说明了玩家II的行动序列以及库恩定理10.4的证明。支付被省略。
\end{tcolorbox}

Given \(\mu\) , we want to construct a realization equivalent behavior strategy \(\beta\) . Essentially, the probability \(\beta \left( c\right)\) for any move \(c\) at an information set \(h\) of player \(i\) will be the "observed" probability that player \(i\) makes move \(c\) when she uses \(\mu\) . In order to obtain this probability, we consider sequences of moves of player \(i\) . Such a sequence is obtained by considering any node \(x\) of the game tree and taking the moves of player \(i\) on the path from the root to \(x\) ; the resulting sequence shall be denoted by \({\sigma }^{i}\left( x\right)\) . For example, if \(x\) is the leaf \(w\) of the game tree in Figure 10.11 and \(i = \mathrm{{II}}\) , then that sequence is \({\sigma }^{\mathrm{{II}}}\left( w\right)  = {lbc}\) , that is, move \(l\) followed by \(b\) followed by \(c\) . If all players are considered, then node \(w\) is reached by the sequence of moves \({lLbc}\) which includes move \(L\) of player \(\mathrm{I}\) , but that move is ignored in the definition of \({\sigma }^{\mathrm{{II}}}\left( w\right)\) ; for player \(\mathrm{I}\) , the sequence of moves defined by \(w\) is just the move \(L\) , that is, \({\sigma }^{\mathrm{I}}\left( w\right)  = L\) .

\begin{tcolorbox}
给定 \(\mu\) ，我们希望构建一个实现等价的行为策略 \(\beta\) 。本质上，对于玩家 \(i\) 在信息集 \(h\) 处的任何行动 \(c\) ，其概率 \(\beta \left( c\right)\) 将是玩家 \(i\) 使用 \(\mu\) 时选择行动 \(c\) 的“观察到的”概率。为了得到这个概率，我们考虑玩家 \(i\) 的行动序列。这样的序列是通过考虑博弈树(game tree)的任何节点 \(x\) ，并取从根节点到 \(x\) 的路径上玩家 \(i\) 的行动得到的；得到的序列用 \({\sigma }^{i}\left( x\right)\) 表示。例如，如果 \(x\) 是图 10.11 中博弈树的叶节点(leaf) \(w\) 且 \(i = \mathrm{{II}}\) ，那么该序列就是 \({\sigma }^{\mathrm{{II}}}\left( w\right)  = {lbc}\) ，即行动 \(l\) 之后接着 \(b\) 再接着 \(c\) 。如果考虑所有玩家，那么节点 \(w\) 是通过行动序列 \({lLbc}\) 到达的，该序列包括玩家 \(\mathrm{I}\) 的行动 \(L\) ，但在 \({\sigma }^{\mathrm{{II}}}\left( w\right)\) 的定义中忽略了该行动；对于玩家 \(\mathrm{I}\) ，由 \(w\) 定义的行动序列就是行动 \(L\) ，即 \({\sigma }^{\mathrm{I}}\left( w\right)  = L\) 。
\end{tcolorbox}

According to Definition 10.1, the condition of perfect recall states that the sequence of own earlier moves is the same for all nodes in an information set. That is,

\begin{tcolorbox}
根据定义 10.1，完美记忆(perfect recall)条件表明，对于信息集中的所有节点，玩家自己先前的行动序列是相同的。也就是说，
\end{tcolorbox}

\[
h \in  {H}_{i} \Rightarrow  \forall u,v \in  h\;{\sigma }^{i}\left( u\right)  = {\sigma }^{i}\left( v\right) .
\]

An example is given in Figure 10.11 by the information set \(h\) where player II has the two moves \(c\) and \(d\) , which has two nodes \(u\) and \(v\) with \({\sigma }^{i}\left( u\right)  = {\sigma }^{i}\left( v\right)  = {lb}\) . Consequently, the sequence of moves leading to \(h\) can be defined uniquely as \({\sigma }_{h} = {\sigma }^{i}\left( u\right)\) for any \(u \in  h\) , where \(i\) is identified as the player who moves at \(h\) , that is, \(h \in  {H}_{i}\) . Furthermore, we can extend the sequence \({\sigma }_{h}\) that leads to \(h\) by any move \(c\) at \(h\) , which we write as the longer sequence \({\sigma }_{h}c\) , to be read as " \({\sigma }_{h}\) followed by \({c}^{\prime \prime }\) . In the example, the sequence that leads to \(h\) is \({\sigma }_{h} = {lb}\) , and it can be extended either by move \(c\) or by move \(d\) , giving the longer sequences \({lbc}\) and \({lbd}\) . Any sequence of moves of player \(i\) is either the empty sequence \(\varnothing\) (for example \({\sigma }^{i}\left( r\right)  = \varnothing\) if \(r\) is the root of the tree), or it has a last move \(c\) made at some information set \(h\) of player \(i\) , and then it can be written as \({\sigma }_{h}c\) . Hence, these sequences \({\sigma }_{h}c\) for \(h \in  {H}_{i}\) , \(c \in  {C}_{h}\) and \(\varnothing\) are all possible sequences of moves of player \(i\) .

\begin{tcolorbox}
图 10.11 给出了一个例子，信息集 \(h\) 中玩家 II 有两个行动 \(c\) 和 \(d\) ，该信息集有两个节点 \(u\) 和 \(v\) 且 \({\sigma }^{i}\left( u\right)  = {\sigma }^{i}\left( v\right)  = {lb}\) 。因此，对于任何 \(u \in  h\) ，导致 \(h\) 的行动序列可以唯一地定义为 \({\sigma }_{h} = {\sigma }^{i}\left( u\right)\) ，其中 \(i\) 被确定为在 \(h\) 处行动的玩家，即 \(h \in  {H}_{i}\) 。此外，我们可以通过 \(h\) 处的任何行动 \(c\) 扩展导致 \(h\) 的序列 \({\sigma }_{h}\) ，我们将其写为更长的序列 \({\sigma }_{h}c\) ，读作“ \({\sigma }_{h}\) 之后接着 \({c}^{\prime \prime }\) ”。在这个例子中，导致 \(h\) 的序列是 \({\sigma }_{h} = {lb}\) ，它可以通过行动 \(c\) 或行动 \(d\) 进行扩展，得到更长的序列 \({lbc}\) 和 \({lbd}\) 。玩家 \(i\) 的任何行动序列要么是空序列 \(\varnothing\) (例如，如果 \(r\) 是树的根节点，则为 \({\sigma }^{i}\left( r\right)  = \varnothing\) )，要么它有在玩家 \(i\) 的某个信息集 \(h\) 处的最后一个行动 \(c\) ，然后它可以写成 \({\sigma }_{h}c\) 。因此，对于 \(h \in  {H}_{i}\) 、 \(c \in  {C}_{h}\) 和 \(\varnothing\) 的这些序列 \({\sigma }_{h}c\) 都是玩家 \(i\) 所有可能的行动序列。
\end{tcolorbox}

We now consider the probability \(\mu \left\lbrack  \sigma \right\rbrack\) that such a sequence \(\sigma\) of player \(i\) is played when player \(i\) uses \(\mu\) . This probability \(\mu \left\lbrack  \sigma \right\rbrack\) is simply the combined probability under \(\mu\) for all the pure strategies that agree with \(\sigma\) in the sense that they prescribe all the moves in \(\sigma\) . For example, if \(\sigma  = {lb}\) in Figure 10.11, then the pure strategies of player \(i\) that agree with \(\sigma\) are the strategies that prescribe the moves \(l\) and \(b\) , which are the pure strategies \({lbc}\) and \({lbd}\) . So

\begin{tcolorbox}
现在我们考虑当玩家 \(i\) 使用 \(\mu\) 时，玩家 \(i\) 的这样一个序列 \(\sigma\) 被执行的概率 \(\mu \left\lbrack  \sigma \right\rbrack\) 。这个概率 \(\mu \left\lbrack  \sigma \right\rbrack\) 仅仅是在 \(\mu\) 下，所有与 \(\sigma\) 一致(即在规定 \(\sigma\) 中的所有行动这一意义上)的纯策略的组合概率。例如，如果图10.11中的 \(\sigma  = {lb}\) ，那么与 \(\sigma\) 一致的玩家 \(i\) 的纯策略就是规定了行动 \(l\) 和 \(b\) 的策略，即纯策略 \({lbc}\) 和 \({lbd}\) 。所以
\end{tcolorbox}

\[
\mu \left\lbrack  {lb}\right\rbrack   = \mu \left( {lbc}\right)  + \mu \left( {lbd}\right) , \tag{10.4}
\]

where \(\mu \left( {lbc}\right)\) and \(\mu \left( {lbd}\right)\) are the mixed strategy probabilities for the pure strategies \({lbc}\) and \({lbd}\) .

\begin{tcolorbox}
其中 \(\mu \left( {lbc}\right)\) 和 \(\mu \left( {lbd}\right)\) 分别是纯策略 \({lbc}\) 和 \({lbd}\) 的混合策略概率。
\end{tcolorbox}

The longer the sequence \(\sigma\) is, the smaller is the set of pure strategies that agree with \(\sigma\) . Conversely, shorter sequences have more strategies that agree with them. The extreme case is the empty sequence \(\varnothing\) , because every pure strategy agrees with that sequence. Consequently, we have

\begin{tcolorbox}
序列 \(\sigma\) 越长，与 \(\sigma\) 一致的纯策略集合就越小。相反，较短的序列有更多与之一致的策略。极端情况是空序列 \(\varnothing\) ，因为每个纯策略都与该序列一致。因此，我们有
\end{tcolorbox}

\[
\mu \left\lbrack  \varnothing \right\rbrack   = 1\text{ . } \tag{10.5}
\]

We now look at a different interpretation of equation (10.4). In this special example, \({lbc}\) and \({lbd}\) are also sequences of moves, namely \({lbc} = {\sigma }_{h}c\) and \({lbd} = {\sigma }_{h}d\) with \({\sigma }_{h} = {lb}\) . So the combined probability \(\mu \left\lbrack  {lbc}\right\rbrack\) of all pure strategies that agree with \({lbc}\) is simply the mixed strategy probability \(\mu \left( {lbc}\right)\) , and similarly \(\mu \left\lbrack  {lbd}\right\rbrack   = \mu \left( {lbd}\right)\) . Clearly, \(\mu \left\lbrack  {lb}\right\rbrack   = \mu \left\lbrack  {lbc}\right\rbrack   + \mu \left\lbrack  {lbd}\right\rbrack\) . In general, considering any information set \(h\) of player \(i\) , we have the following important equation:

\begin{tcolorbox}
现在我们来看方程(10.4)的另一种解释。在这个特殊例子中， \({lbc}\) 和 \({lbd}\) 也是行动序列，即带有 \({\sigma }_{h} = {lb}\) 的 \({lbc} = {\sigma }_{h}c\) 和 \({lbd} = {\sigma }_{h}d\) 。所以所有与 \({lbc}\) 一致的纯策略的组合概率 \(\mu \left\lbrack  {lbc}\right\rbrack\) 仅仅是混合策略概率 \(\mu \left( {lbc}\right)\) ，同理 \(\mu \left\lbrack  {lbd}\right\rbrack   = \mu \left( {lbd}\right)\) 。显然， \(\mu \left\lbrack  {lb}\right\rbrack   = \mu \left\lbrack  {lbc}\right\rbrack   + \mu \left\lbrack  {lbd}\right\rbrack\) 。一般来说，考虑玩家 \(i\) 的任何信息集 \(h\) ，我们有以下重要方程:
\end{tcolorbox}

\[
\mu \left\lbrack  {\sigma }_{h}\right\rbrack   = \mathop{\sum }\limits_{{c \in  {C}_{h}}}\mu \left\lbrack  {{\sigma }_{h}c}\right\rbrack  . \tag{10.6}
\]

This equation holds because because the sequence \({\sigma }_{h}\) that leads to \(h\) is extended by some move \(c\) at \(h\) , giving the longer sequence \({\sigma }_{h}c\) . When considering the pure strategies that agree with these extended sequences, we obtain exactly the pure strategies that agree with \({\sigma }_{h}\) .

\begin{tcolorbox}
这个方程成立是因为通向 \(h\) 的序列 \({\sigma }_{h}\) 在 \(h\) 处通过某个行动 \(c\) 得到了扩展，从而得到了更长的序列 \({\sigma }_{h}c\) 。当考虑与这些扩展序列一致的纯策略时，我们恰好得到了与 \({\sigma }_{h}\) 一致的纯策略。
\end{tcolorbox}

We are now in a position to define the behavior strategy \(\beta\) that is realization equivalent to \(\mu\) . Namely, given an information set \(h\) of player \(i\) so that \(\mu \left\lbrack  {\sigma }_{h}\right\rbrack   > 0\) ,

\begin{tcolorbox}
现在我们可以定义与 \(\mu\) 实现等价的行为策略 \(\beta\) 。即，给定玩家 \(i\) 的一个信息集 \(h\) ，使得 \(\mu \left\lbrack  {\sigma }_{h}\right\rbrack   > 0\) ，
\end{tcolorbox}

let

\begin{tcolorbox}
令
\end{tcolorbox}

\[
\beta \left( c\right)  = \frac{\mu \left\lbrack  {{\sigma }_{h}c}\right\rbrack  }{\mu \left\lbrack  {\sigma }_{h}\right\rbrack  }. \tag{10.7}
\]

Because of (10.6), this defines a probability distribution on the set \({C}_{h}\) of moves at \(h\) , as required by a behavior strategy. In (10.7), the condition of perfect recall is used because otherwise the sequence \({\sigma }_{h}\) that leads to the information set \(h\) would not be uniquely defined.

\begin{tcolorbox}
由于(10.6)，这就在 \(h\) 处的行动集合 \({C}_{h}\) 上定义了一个概率分布，这正是行为策略所要求的。在(10.7)中，使用了完美记忆的条件，因为否则通向信息集 \(h\) 的序列 \({\sigma }_{h}\) 就不会被唯一确定。
\end{tcolorbox}

What about the case \(\mu \left\lbrack  {\sigma }_{h}\right\rbrack   = 0\) ? This means that no pure strategy that agrees with \({\sigma }_{h}\) has positive probability under \(\mu\) . In other words, \(h\) is unreachable when player \(i\) uses \(\mu\) . We will show soon that \(\mu \left\lbrack  \sigma \right\rbrack   = \beta \left\lbrack  \sigma \right\rbrack\) for all sequences \(\sigma\) of player \(i\) . Hence, \(h\) is also unreachable when the player uses \(\beta\) , so that \(\beta \left( c\right)\) for \(c \in  {C}_{h}\) does not need to be defined, that is, \(\beta\) is a reduced behavior strategy according to Definition 10.3. Alternatively, we can define some behavior at \(h\) (although it will never matter because player \(i\) will not move so as to reach \(h\) ) in an arbitrary manner, for example by giving each move \(c\) at \(h\) equal probability, setting \(\beta \left( c\right)  = 1/\left| {C}_{h}\right|\) .

\begin{tcolorbox}
情况 \(\mu \left\lbrack  {\sigma }_{h}\right\rbrack   = 0\) 呢？这意味着在 \(\mu\) 下，与 \({\sigma }_{h}\) 一致的纯策略的概率为零。换句话说，当玩家 \(i\) 使用 \(\mu\) 时， \(h\) 是无法到达的。我们很快会证明，对于玩家 \(i\) 的所有序列 \(\sigma\) ，有 \(\mu \left\lbrack  \sigma \right\rbrack   = \beta \left\lbrack  \sigma \right\rbrack\) 。因此，当玩家使用 \(\beta\) 时， \(h\) 同样无法到达，所以对于 \(c \in  {C}_{h}\) 不需要定义 \(\beta \left( c\right)\) ，也就是说，根据定义 10.3， \(\beta\) 是一个简化的行为策略。或者，我们可以以任意方式定义在 \(h\) 处的某种行为(尽管这无关紧要，因为玩家 \(i\) 不会移动到 \(h\) )，例如给 \(h\) 处的每个行动 \(c\) 赋予相等的概率，即设置 \(\beta \left( c\right)  = 1/\left| {C}_{h}\right|\) 。
\end{tcolorbox}

We claim that \(\beta\) , as defined by (10.7), is realization equivalent to \(\mu\) . We can define \(\beta \left\lbrack  \sigma \right\rbrack\) for a sequence \(\sigma\) of player \(i\) , which is the probability that all the moves in \(\sigma\) are made when the player uses \(\beta\) . Clearly, this is simply the product of the probabilities for the moves in \(\sigma\) . That is, if \(\sigma  = {c}_{1}{c}_{2}\cdots {c}_{n}\) , then \(\beta \left\lbrack  \sigma \right\rbrack   = \beta \left( {c}_{1}\right) \beta \left( {c}_{2}\right) \cdots \beta \left( {c}_{n}\right)\) . When \(n = 0\) , then \(\sigma\) is the empty sequence \(\varnothing\) , in which case we set \(\beta \left\lbrack  \varnothing \right\rbrack   = 1\) . So we get by (10.7)

\begin{tcolorbox}
我们声称，由 (10.7) 定义的 \(\beta\) 与 \(\mu\) 是实现等价的。我们可以为玩家 \(i\) 的序列 \(\sigma\) 定义 \(\beta \left\lbrack  \sigma \right\rbrack\) ，它是玩家使用 \(\beta\) 时 \(\sigma\) 中的所有行动都被采取的概率。显然，这只是 \(\sigma\) 中行动概率的乘积。也就是说，如果 \(\sigma  = {c}_{1}{c}_{2}\cdots {c}_{n}\) ，那么 \(\beta \left\lbrack  \sigma \right\rbrack   = \beta \left( {c}_{1}\right) \beta \left( {c}_{2}\right) \cdots \beta \left( {c}_{n}\right)\) 。当 \(n = 0\) 时， \(\sigma\) 是空序列 \(\varnothing\) ，在这种情况下我们设置 \(\beta \left\lbrack  \varnothing \right\rbrack   = 1\) 。所以根据 (10.7) 我们得到
\end{tcolorbox}

\[
\beta \left\lbrack  {{c}_{1}{c}_{2}\cdots {c}_{n}}\right\rbrack   = \beta \left( {c}_{1}\right) \beta \left( {c}_{2}\right) \cdots \beta \left( {c}_{n}\right)  = \frac{\mu \left\lbrack  {c}_{1}\right\rbrack  }{\mu \left\lbrack  0\right\rbrack  }\frac{\mu \left\lbrack  {{c}_{1}{c}_{2}}\right\rbrack  }{\mu \left\lbrack  {c}_{1}\right\rbrack  }\cdots \frac{\mu \left\lbrack  {{c}_{1}{c}_{2}\cdots {c}_{n - 1}{c}_{n}}\right\rbrack  }{\mu \left\lbrack  {{c}_{1}{c}_{2}\cdots {c}_{n - 1}}\right\rbrack  } = \mu \left\lbrack  {{c}_{1}{c}_{2}\cdots {c}_{n}}\right\rbrack
\]

because all the intermediate terms cancel out, and by (10.5). (This is best understood by considering an example like \(\sigma  = {lbc} = {c}_{1}{c}_{2}{c}_{3}\) in Figure 10.11.) Because \(\beta \left\lbrack  \sigma \right\rbrack   = \mu \left\lbrack  \sigma \right\rbrack\) for all sequences \(\sigma\) of player \(i\) , the strategies \(\beta\) and \(\mu\) are realization equivalent, as explained before Theorem 10.4; the reason is that any node \(u\) of the game tree is reached by a unique sequence of moves of all players and of chance. The moves of the other players have fixed probabilities, so all that matters is the sequence \({\sigma }^{i}\left( u\right)\) of moves of player \(i\) , which is some sequence \(\sigma\) , which has the same probability no matter whether player \(i\) uses \(\beta\) or \(\mu\) .

\begin{tcolorbox}
因为所有中间项都抵消了，并且根据 (10.5)。(通过考虑图 10.11 中像 \(\sigma  = {lbc} = {c}_{1}{c}_{2}{c}_{3}\) 这样的例子可以最好地理解这一点。)因为对于玩家 \(i\) 的所有序列 \(\sigma\) 都有 \(\beta \left\lbrack  \sigma \right\rbrack   = \mu \left\lbrack  \sigma \right\rbrack\) ，根据定理 10.4 之前的解释，策略 \(\beta\) 和 \(\mu\) 是实现等价的；原因是博弈树(game tree)的任何节点(node)都是由所有玩家和机会的唯一行动序列到达的。其他玩家的行动有固定的概率，所以重要的是玩家 \(i\) 的行动序列 \({\sigma }^{i}\left( u\right)\) ，它是某个序列 \(\sigma\) ，无论玩家 \(i\) 使用 \(\beta\) 还是 \(\mu\) ，该序列的概率都是相同的。
\end{tcolorbox}

It is instructive to reproduce the preceding proof for the examples after Definition 10.3 or Figure 10.10 above, which gives also some practice in finding behavior strategies in small games.

\begin{tcolorbox}
对定义 10.3 之后的例子或上面的图 10.10 重现前面的证明是有启发性的，这也能让我们在小型博弈中寻找行为策略(behavior strategy)方面得到一些练习。
\end{tcolorbox}

\section*{10.9 Behavior Strategies in the Monty Hall Problem}

\begin{tcolorbox}
\section*{10.9 蒙提霍尔问题中的行为策略(behavior strategy)}
\end{tcolorbox}

We illustrate the concept of behavior strategies for a game-theoretic, non-standard version of the Monty Hall problem. The original version has been posed as a problem in probability in the context of an early television game show with the showmaster named "Monty Hall". A contestant can win a prize by choosing one of three closed doors. Behind one of these doors is the prize (a car), and behind each of the other two doors is something the contestant does not want (a goat). Once the contestant has selected one of the doors, Monty Hall opens one of the remaining doors, which reveals a goat, and offers the contestant the option to switch to the other closed door. The question is, should she switch doors? It seems that the contestant is now facing equal chances between two options, the originally chosen door and the other unopened door, so that there is no point in switching. However, the correct solution is that the original probability \(\frac{1}{3}\) that the contestant has chosen the correct door has not changed by Monty Hall's action of revealing a goat behind another door, because Monty Hall knows where the prize is and can always choose such a door. That probability stays at \(\frac{1}{3}\) , so that the probability that the other unopened door has the prize is \(\frac{2}{3}\) , and therefore the contestant should indeed switch.

\begin{tcolorbox}
我们通过博弈论中一个非标准版本的蒙提·霍尔问题(Monty Hall problem)来说明行为策略(behavior strategy)的概念。最初的版本是作为一个概率问题提出的，它源于一档早期的电视游戏节目，节目主持人名为“蒙提·霍尔”。参赛者可以通过选择三扇关闭的门之一来赢得奖品。其中一扇门后面是奖品(一辆汽车)，另外两扇门后面是参赛者不想要的东西(一只山羊)。一旦参赛者选择了其中一扇门，蒙提·霍尔会打开剩下两扇门中的一扇，露出一只山羊，然后给参赛者一个机会，让其选择换到另一扇未打开的门。问题是，她应该换门吗？看起来参赛者现在面临着两个选项(最初选择的门和另一扇未打开的门)的同等机会，所以换门似乎没有意义。然而，正确的答案是，参赛者最初选择正确门的概率 \(\frac{1}{3}\) 并不会因为蒙提·霍尔打开另一扇门露出山羊的行为而改变，因为蒙提·霍尔知道奖品在哪里，并且总能选择这样一扇门。这个概率仍然是 \(\frac{1}{3}\) ，所以另一扇未打开的门后面有奖品的概率是 \(\frac{2}{3}\) ，因此参赛者确实应该换门。
\end{tcolorbox}

This analysis runs counter to the spontaneous reaction by many people who rather see the opened door with the goat as confirmation that they had chosen the correct door in the first place (this is incorrect). Moreover, they are suspicious that Monty Hall only tries to lure them away from their correct choice. That suspicion is correct if Monty Hall would only sometimes open another door with a goat behind it and offer the switch to the remaining door. Rather obviously, if Monty Hall only offered the switch when the contestant has chosen the correct door, then switching would mean that she loses the prize, and therefore she should never switch. In short, the original "solution" to the Monty Hall problem is only valid when Monty Hall always offers the switch.

\begin{tcolorbox}
这种分析与许多人的直觉反应相悖，他们更倾向于将打开有山羊的门视为自己最初选择正确门的确认(这是错误的)。此外，他们怀疑蒙提·霍尔只是试图诱使他们放弃正确的选择。如果蒙提·霍尔只是有时打开另一扇有山羊的门并提供换门的机会，那么这种怀疑是正确的。很明显，如果蒙提·霍尔只在参赛者选择了正确的门时才提供换门机会，那么换门就意味着她会失去奖品，因此她永远不应该换门。简而言之，蒙提·霍尔问题的原始“解决方案”仅在蒙提·霍尔总是提供换门机会时才有效。
\end{tcolorbox}

\begin{center}
\includegraphics[max width=0.5\textwidth]{images/bo_d2e5cmbef24c73b03tn0_299_337_1408_658_724_0.jpg}
\end{center}
\hspace*{3em} 

Figure 10.12 The (modified) Monty Hall problem as a game between player I as the contestant who has chosen the correct door with probability \(\frac{1}{3}\) (left chance move), and the television show host Monty Hall as player II who can keep the doors closed (moves \(C\) and \(c\) ), or can open another non-prize door (moves \(O\) and \(o\) ). When a door is opened, player I can keep her original choice(K)or switch to the remaining closed door(S). Payoff 1 means that player I gets the prize, paid by player II, payoff 0 means no prize.

\begin{tcolorbox}
图10.12 将(修改后的)蒙提·霍尔问题视为玩家I(参赛者)和玩家II(电视节目主持人蒙提·霍尔)之间的博弈。参赛者选择正确门的概率为 \(\frac{1}{3}\) (左侧的随机行动)，蒙提·霍尔可以选择保持门关闭(行动 \(C\) 和 \(c\) )，或者打开另一扇没有奖品的门(行动 \(O\) 和 \(o\) )。当一扇门被打开时，玩家I可以保持她最初的选择(K)或换到剩下的未打开的门(S)。收益为1意味着玩家I获得奖品，由玩家II支付；收益为0意味着没有奖品。
\end{tcolorbox}

Figure 10.12 shows a new game between the contestant as player I and Monty Hall as a player II who has choices after the contestant has chosen one of the three doors. We assume that the prize is hidden behind each door with equal probability. Monty Hall knows whether or not the contestant has chosen the correct door, shown by his two decision nodes with the two moves \(C\) and \(O\) at the left node and \(c\) and \(o\) at the right node, which he can choose independently. The moves \(C\) and \(c\) mean that Monty keeps the doors closed, moves \(O\) and \(o\) that he opens another door with a goat behind it. In that case the contestant (who does not know if she has chosen the correct door) can keep to her choice(K)or switch to the other unopened door(S), and then gets the prize or not. We assume that the hosting television station prefers not to the pay the prize, so that this is zero-sum game.

\begin{tcolorbox}
图10.12展示了参赛者作为玩家I和蒙提·霍尔作为玩家II之间的一个新博弈，在参赛者选择了三扇门之一后，蒙提·霍尔有多种选择。我们假设奖品以相等的概率隐藏在每扇门后面。蒙提·霍尔知道参赛者是否选择了正确的门，这由他的两个决策节点表示，左节点有两个行动 \(C\) 和 \(O\) ，右节点有两个行动 \(c\) 和 \(o\) ，他可以独立地做出选择。行动 \(C\) 和 \(c\) 意味着蒙提保持门关闭，行动 \(O\) 和 \(o\) 意味着他打开另一扇后面有山羊的门。在这种情况下，参赛者(不知道自己是否选择了正确的门)可以坚持自己的选择(K)或换到另一扇未打开的门(S)，然后决定是否获得奖品。我们假设主办电视台更不愿意支付奖品，所以这是一个零和博弈(zero-sum game)。
\end{tcolorbox}

This game has the same structure as the introductory game in Figure 10.1 of this chapter, although with reversed roles of the players, and different payoffs and chance probabilities. It has the following strategic form:

\begin{tcolorbox}
这个博弈与本章图10.1中的介绍性博弈具有相同的结构，尽管玩家的角色相反，并且收益和随机概率不同。它具有以下策略形式(strategic form):
\end{tcolorbox}

(10.8)Although the game is zero-sum, we show the expected payoffs to both players because the best responses are easier to see. If the contestant chooses \(K\) , which is the top row in (10.8), then her expected payoff is always \(\frac{1}{3}\) , and all pure strategies of player II are best responses. This is therefore a degenerate game. Payoff \(\frac{1}{3}\) also results when Monty Hall never opens a door (strategy \({Cc}\) ), where player I does not even have the opportunity to switch doors. Hence, \(K\) and \({Cc}\) are mutual best responses with the pure equilibrium(K, Cc), and \(\frac{1}{3}\) as the value of the game.

\begin{tcolorbox}
尽管这个博弈是零和博弈，但我们展示了两个玩家的期望效用(expected payoff)，因为这样更容易看出最优反应(best response)。如果参赛者选择 \(K\) ，即在(10.8)中的顶行，那么她的期望效用总是 \(\frac{1}{3}\) ，玩家II的所有纯策略都是最优反应。因此，这是一个退化博弈(degenerate game)。当蒙提·霍尔从不打开门(策略 \({Cc}\) )时，也会得到收益 \(\frac{1}{3}\) ，此时玩家I甚至没有机会换门。因此， \(K\) 和 \({Cc}\) 是具有纯策略均衡(pure equilibrium)(K，Cc)的相互最优反应，并且 \(\frac{1}{3}\) 是博弈的值。
\end{tcolorbox}

\begin{center}
\includegraphics[max width=0.5\textwidth]{images/bo_d2e5cmbef24c73b03tn0_300_484_893_690_395_0.jpg}
\end{center}
\hspace*{3em} 

We now study the set of all equilibrium strategies of the players. A second pure equilibrium is(K, Oc), where Monty Hall only opens another door if the contestant has chosen the correct door in the first place (what we called "trying to lure her to switch"). In this equilibrium, the unique best response to \({Oc}\) is \(K\) . Hence, player I’s unique max-min strategy (and strategy in every equilibrium) is \(K\) , by Proposition 8.3 about zero-sum games.

\begin{tcolorbox}
我们现在研究参与者的所有均衡策略集合。第二个纯策略均衡是 (K, Oc)，即蒙提·霍尔仅在参赛者首次选择了正确的门时才打开另一扇门(我们称之为“试图诱使她换门”)。在这个均衡中，对 \({Oc}\) 的唯一最优反应是 \(K\) 。因此，根据关于零和博弈(zero - sum game)的命题 8.3，参与者 I 的唯一极大极小策略(max - min strategy)(也是每个均衡中的策略)是 \(K\) 。
\end{tcolorbox}

On the other hand, there are many equilibrium strategies of player II, because all of his pure strategies are best responses to \(K\) , and even the strategies \({Co}\) and \({Oo}\) can potentially have positive probability in a mixed equilibrium. One such mixed strategy is \(y = \left( {0,0,\frac{1}{2},\frac{1}{2}}\right)\) , that is, to play \({Oc}\) and \({Oo}\) with probability \(\frac{1}{2}\) each. Then the top row \(K\) has expected payoff \(- \frac{1}{3}\) to player II (it always has), and the bottom row \(S\) has also expected payoff \(- \frac{1}{3}\) , so \(y\) is also a max-min strategy (in terms of his own payoffs) of player II. A second max-min strategy is \({y}^{\prime } = \left( {0,\frac{1}{3},\frac{2}{3},0}\right)\) where the same applies. With \({y}^{\prime }\) , Monty Hall plays with probability \(\frac{1}{3}\) the strategy \({Co}\) , with the seemingly strange offer to the contestant to switch only if she initially chose the wrong door, so that switching would give her the prize for sure; however, \({y}^{\prime }\) also chooses with probability \(\frac{2}{3}\) the "opposite" strategy \({Oc}\) . The strategy \({Oo}\) , chosen with probability \(\frac{1}{2}\) by \(y\) , is the strategy of always opening another door with a goat behind it, as in the original Monty Hall problem, and payoff \(\frac{2}{3}\) when player \(\mathrm{I}\) chooses \(S\) . Nevertheless, in the game the contestant should never switch, because doing so with any positive probability is not a max-min strategy (with a worse payoff than \(\frac{1}{3}\) against \({Oc}\) ).

\begin{tcolorbox}
另一方面，参与者 II 有许多均衡策略，因为他的所有纯策略都是对 \(K\) 的最优反应，甚至策略 \({Co}\) 和 \({Oo}\) 在混合均衡(mixed equilibrium)中都可能有正概率。一种这样的混合策略是 \(y = \left( {0,0,\frac{1}{2},\frac{1}{2}}\right)\) ，即分别以概率 \(\frac{1}{2}\) 采用 \({Oc}\) 和 \({Oo}\) 。那么，顶行 \(K\) 给参与者 II 的期望效用(expected payoff)是 \(- \frac{1}{3}\) (始终如此)，底行 \(S\) 的期望效用也是 \(- \frac{1}{3}\) ，所以 \(y\) 也是参与者 II 的一个极大极小策略(就他自己的收益而言)。第二个极大极小策略是 \({y}^{\prime } = \left( {0,\frac{1}{3},\frac{2}{3},0}\right)\) ，同样适用上述情况。对于 \({y}^{\prime }\) ，蒙提·霍尔以概率 \(\frac{1}{3}\) 采用策略 \({Co}\) ，即只有当参赛者最初选错门时才向她提出换门的看似奇怪的提议，这样换门肯定能让她获得奖品；然而， \({y}^{\prime }\) 也以概率 \(\frac{2}{3}\) 选择“相反”的策略 \({Oc}\) 。策略 \({Oo}\) ， \(y\) 以概率 \(\frac{1}{2}\) 选择它，是始终打开另一扇后面有山羊的门的策略，就像原始的蒙提·霍尔问题一样，并且当参与者 \(\mathrm{I}\) 选择 \(S\) 时期望效用为 \(\frac{2}{3}\) 。然而，在这个游戏中，参赛者永远不应该换门，因为以任何正概率换门都不是极大极小策略(与 \({Oc}\) 对抗时，收益比 \(\frac{1}{3}\) 更差)。
\end{tcolorbox}

It can be shown that all optimal mixed strategies of player II in this game are convex combinations of the described max-min strategies \({Cc},{Oc},y\) , and \({y}^{\prime }\) . Except for \({y}^{\prime }\) , these are all behavior strategies. As pure strategies, \({Cc}\) and \({Co}\) are trivially behavior strategies. With the strategy \(y\) , the behavior of player \(\Pi\) is to always choose \(O\) at the left decision node and to randomize between \(c\) and \(o\) with probability \(\frac{1}{2}\) at the right decision node. The strategy \({y}^{\prime }\) is not a behavior strategy, as its distribution on the move pairs \({Cc},{Co},{Oc},{Oo}\) in the following table shows:

\begin{tcolorbox}
可以证明，在这个游戏中，参与者 II 的所有最优混合策略(mixed strategy)都是所描述的极大极小策略 \({Cc},{Oc},y\) 和 \({y}^{\prime }\) 的凸组合(convex combination)。除了 \({y}^{\prime }\) ，这些都是行为策略(behavior strategy)。作为纯策略， \({Cc}\) 和 \({Co}\) 显然是行为策略。采用策略 \(y\) 时，参与者 \(\Pi\) 的行为是在左决策节点(node)始终选择 \(O\) ，在右决策节点以概率 \(\frac{1}{2}\) 在 \(c\) 和 \(o\) 之间随机选择。策略 \({y}^{\prime }\) 不是行为策略，如下表中它在行动对 \({Cc},{Co},{Oc},{Oo}\) 上的分布所示:
\end{tcolorbox}

(10.9)

\begin{center}
\adjustbox{max width=\textwidth}{
\begin{tabular}{|c|c|c|c|c|c|}
\hline
\multirow{3}{*}{left node \(C\) Og} & right node \(C\) & 0 & left node & right node \(C\) & 0 \\
\cline{2-6}
 & 0 & \(\frac{1}{3}\) & \(C\) & OIN & ON \\
\cline{2-6}
 & WIN & 0 & Og & one & ON \\
\cline{1-6}
\hline
\end{tabular}
}
\end{center}

\begin{tcolorbox}

\begin{center}
\adjustbox{max width=\textwidth}{
\begin{tabular}{|c|c|c|c|c|c|}
\hline
\multirow{3}{*}{左节点 \(C\) 奥格(Og)} & 右节点 \(C\) & 0 & 左节点 & 右节点 \(C\) & 0 \\
\cline{2-6}
 & 0 & \(\frac{1}{3}\) & \(C\) & OIN & ON \\
\cline{2-6}
 & WIN & 0 & 奥格(Og) & 一 & ON \\
\cline{1-6}
\hline
\end{tabular}
}
\end{center}

\end{tcolorbox}

The right table in (10.9) shows the behavior strategy \({\beta }^{\prime }\) that, by Kuhn’s theorem, is realization equivalent to \({y}^{\prime }\) . It randomizes between \(C\) and \(O\) with probabilities \(\frac{1}{3}\) and \(\frac{2}{3}\) , and between \(c\) and \(o\) with probabilities \(\frac{2}{3}\) and \(\frac{1}{3}\) . Translated back to probabilities for the pure strategies \({Cc},{Co},{Oc},{Oo}\) , this behavior strategy \({\beta }^{\prime }\) is the mixed strategy \(\left( {\frac{2}{9},\frac{1}{9},\frac{4}{9},\frac{2}{9}}\right)\) , which gives expected payoff \(- \frac{1}{3}\) to player II in row \(S\) .

\begin{tcolorbox}
根据库恩定理，(10.9)中的右表展示了与混合策略 \({y}^{\prime }\) 实现等价的行为策略 \({\beta }^{\prime }\) 。该行为策略以概率 \(\frac{1}{3}\) 和 \(\frac{2}{3}\) 在 \(C\) 和 \(O\) 之间进行随机选择，以概率 \(\frac{2}{3}\) 和 \(\frac{1}{3}\) 在 \(c\) 和 \(o\) 之间进行随机选择。将此行为策略 \({\beta }^{\prime }\) 转换回纯策略 \({Cc},{Co},{Oc},{Oo}\) 的概率，它就是混合策略 \(\left( {\frac{2}{9},\frac{1}{9},\frac{4}{9},\frac{2}{9}}\right)\) ，该混合策略为第 \(S\) 行的参与者 II 带来期望效用 \(- \frac{1}{3}\) 。
\end{tcolorbox}

Figure 10.13 gives a geometric view of these max-min strategies of player II. The left picture shows the mixed-strategy simplex, and the right picture the square that represents the pairs of behavior strategy probabilities. If \(p\) is the probability that he plays \(O\) (vertical coordinate of the square, going downwards) and \(q\) is the probability that he plays \(o\) (horizontal coordinate), then his four pure strategies \({Cc},{Co},{Oc},{Oo}\) are played with probabilities \(\left( {1 - p}\right) \left( {1 - q}\right) ,\left( {1 - p}\right) q,p\left( {1 - q}\right)\) , and \({pq}\) , respectively. For \(p,q \in  \left\{  {\frac{1}{4},\frac{1}{2},\frac{3}{4}}\right\}\) , they define a grid of lines which is shown in both pictures. However, as a subset of the mixed-strategy simplex, this grid is "twisted" and not convex because a convex combination of behavior strategies is a general mixed strategy and usually not a behavior strategy. The square is convex, and so is the set of max-min behavior strategies as a subset of that square. The behavior strategy \({\beta }^{\prime }\) belongs to that set, and is a convex combination of \({Cc}\) and \(y\) .

\begin{tcolorbox}
图 10.13 从几何角度展示了参与者 II 的这些极大极小策略。左图展示了混合策略单纯形，右图展示了代表行为策略概率对的正方形。如果 \(p\) 是他选择 \(O\) 的概率(正方形的垂直坐标，向下)， \(q\) 是他选择 \(o\) 的概率(水平坐标)，那么他的四种纯策略 \({Cc},{Co},{Oc},{Oo}\) 分别以概率 \(\left( {1 - p}\right) \left( {1 - q}\right) ,\left( {1 - p}\right) q,p\left( {1 - q}\right)\) 和 \({pq}\) 被选择。对于 \(p,q \in  \left\{  {\frac{1}{4},\frac{1}{2},\frac{3}{4}}\right\}\) ，它们定义了一个线网格，在两幅图中均有展示。然而，作为混合策略单纯形的一个子集，这个网格是“扭曲的”且非凸的，因为行为策略的凸组合是一般的混合策略，通常不是行为策略。正方形是凸的，作为该正方形子集的极大极小行为策略集也是凸的。行为策略 \({\beta }^{\prime }\) 属于该集合，并且是 \({Cc}\) 和 \(y\) 的凸组合。
\end{tcolorbox}

\begin{center}
\includegraphics[max width=0.9\textwidth]{images/bo_d2e5cmbef24c73b03tn0_302_196_202_1261_521_0.jpg}
\end{center}
\hspace*{3em} 

Figure 10.13 Left: Simplex of mixed strategies (always for player II) for the game (10.8), and in gray its subset of max-min strategies with vertices \({Cc},{Oc},y,{y}^{\prime }\) . The "twisted grid" shows the behavior strategies for move probabilities that are multiples of \(\frac{1}{4}\) . Right: The square of behavior strategies for the extensive game in Figure 10.12, with the behavior strategy \({\beta }^{\prime }\) that is realization equivalent to \({y}^{\prime }\) , and the same grid. The shaded triangle is the set of max-min behavior strategies.

\begin{tcolorbox}
图 10.13 左图:博弈 (10.8) 的混合策略单纯形(始终针对参与者 II)，灰色部分是顶点为 \({Cc},{Oc},y,{y}^{\prime }\) 的极大极小策略子集。“扭曲的网格”展示了移动概率为 \(\frac{1}{4}\) 倍数的行为策略。右图:图 10.12 中扩展式博弈的行为策略正方形，其中行为策略 \({\beta }^{\prime }\) 与 \({y}^{\prime }\) 实现等价，并且有相同的网格。阴影三角形是极大极小行为策略集。
\end{tcolorbox}

The pure strategies are the four corners of the square (in the same positions as the cells in (10.9)), and are also the vertices of the mixed-strategy simplex.

\begin{tcolorbox}
纯策略是正方形的四个角(与 (10.9) 中的单元格位置相同)，也是混合策略单纯形的顶点。
\end{tcolorbox}

\section*{10.10 Subgames and Subgame-Perfect Equilibria}

\begin{tcolorbox}
\section*{10.10 子博弈与子博弈完美均衡}
\end{tcolorbox}

In an extensive game with perfect information as studied in Chapter 4, a subgame is any subtree of the game, given by a node of the tree as the root of the subtree and all its descendants. This definition has to be refined when considering games with imperfect information, where every player has to know that they have entered the subtree.

\begin{tcolorbox}
在第 4 章研究的具有完美信息的扩展式博弈中，子博弈是博弈树的任意子树，由树中的一个节点作为子树的根节点及其所有后代节点构成。在考虑具有不完美信息的博弈时，这个定义需要细化，此时每个参与者都必须知道他们已经进入了子树。
\end{tcolorbox}

Definition 10.5. In an extensive game, a subgame is any subtree of the game tree so that every information set is either a subset of or disjoint from the nodes of the subtree.

\begin{tcolorbox}
定义 10.5:在扩展式博弈中，子博弈是博弈树的任意子树，使得每个信息集要么是子树节点的子集，要么与子树节点不相交。
\end{tcolorbox}

In other words, if any information set \(h\) contains some node of the subtree, then \(h\) must be completely contained in the subtree. This means that the player to move at \(h\) knows that he or she is in the subgame. Information sets that have no nodes in the subtree do not matter because they are disjoint from the subtree.

\begin{tcolorbox}
换句话说，如果任意信息集 \(h\) 包含子树的某个节点，那么 \(h\) 必须完全包含在子树中。这意味着在 \(h\) 处行动的参与者知道他或她处于子博弈中。信息集中在子树中没有节点的情况无关紧要，因为它们与子树不相交。
\end{tcolorbox}

In Figure 10.14, the subtree with player I's decision node as root defines a subgame of the game. On the other hand, a node in player II's second information set is not the root of a subgame, as shown in the figure for the node that follows move \(B\) of player \(\mathrm{I}\) , because this information set is neither contained in nor disjoint from the subtree with that root.

\begin{tcolorbox}
在图 10.14 中，以参与者 I 的决策节点为根的子树定义了该博弈的一个子博弈。另一方面，参与者 II 的第二个信息集中的一个节点不是子博弈的根节点，如图中参与者 \(\mathrm{I}\) 选择行动 \(B\) 之后的节点所示，因为这个信息集既不包含在以该节点为根的子树中，也不与该子树不相交。
\end{tcolorbox}

\begin{center}
\includegraphics[max width=1.0\textwidth]{images/bo_d2e5cmbef24c73b03tn0_303_332_216_1273_609_0.jpg}
\end{center}
\hspace*{3em} 

Figure 10.14 Extensive game, with its subgame, and the reduced strategic form of the game. Only one equilibrium of the game is subgame-perfect, see Figure 10.15.

\begin{tcolorbox}
图10.14 扩展式博弈及其子博弈，以及该博弈的精简策略型。该博弈只有一个均衡是子博弈完美均衡，见图10.15。
\end{tcolorbox}

A trivial subgame is either the game itself, or any leaf of the game tree. Some games have no other subgames, for example the games in Figure 10.1 or 10.5.

\begin{tcolorbox}
平凡子博弈要么是博弈本身，要么是博弈树的任意叶节点。有些博弈没有其他子博弈，例如图10.1或图10.5中的博弈。
\end{tcolorbox}

If every player has perfect recall, then any mixed strategy can be replaced by a realization equivalent behavior strategy by Kuhn's Theorem 10.4. The strategies in any equilibrium can therefore be assumed to be behavior strategies. A behavior strategy can easily be restricted to a subgame, by ignoring the behavior probabilities for the information sets that are disjoint from the subgame. This restriction of a strategy would be much more difficult to define for a profile of mixed strategies.

\begin{tcolorbox}
如果每个参与者都有完美记忆，那么根据库恩定理10.4，任何混合策略都可以用一个实现等价的行为策略来替代。因此，任何均衡中的策略都可以假定为行为策略。通过忽略与子博弈不相交的信息集的行为概率，行为策略可以很容易地限制在子博弈中。对于混合策略组合来说，定义策略的这种限制要困难得多。
\end{tcolorbox}

Definition 10.6. In an extensive game with perfect recall, a subgame-perfect equilibrium (SPE) is a profile of behavior strategies that defines an equilibrium for every subgame of the game.

\begin{tcolorbox}
定义10.6。在具有完美记忆的扩展式博弈中，子博弈完美均衡(SPE)是一个行为策略组合，它为博弈的每个子博弈定义了一个均衡。
\end{tcolorbox}

In Figure 10.14, the reduced strategic form of the extensive game is shown on the right. Figure 10.15 shows the non-trivial subgame that starts with the move of player I, and its strategic form. This subgame has a unique mixed equilibrium, where player I chooses \(T,B\) with probabilities \(\frac{2}{3},\frac{1}{3}\) , and player II chooses \(a,b\) with probabilities \(\frac{1}{2},\frac{1}{2}\) . The resulting payoffs to players I and II are 2 and 4, which are substituted for the subgame to obtain a smaller game where player II chooses between \(l\) and \(r\) .

\begin{tcolorbox}
在图10.14中，扩展式博弈的精简策略型显示在右侧。图10.15显示了从参与者I的行动开始的非平凡子博弈及其策略型。这个子博弈有一个唯一的混合均衡，其中参与者I以概率 \(\frac{2}{3},\frac{1}{3}\) 选择 \(T,B\) ，参与者II以概率 \(\frac{1}{2},\frac{1}{2}\) 选择 \(a,b\) 。参与者I和II的最终收益分别为2和4，用这些收益替代子博弈，得到一个较小的博弈，其中参与者II在 \(l\) 和 \(r\) 之间进行选择。
\end{tcolorbox}

This smaller game is shown on the right in Figure 10.15. Because this small game has perfect information, an SPE is found by backward induction, which here means that player II chooses \(l\) with certainty. In the overall resulting behavior strategies, player I mixes \(T,B\) with probabilities \(\frac{2}{3},\frac{1}{3}\) , and player II chooses \(l\) for sure and mixes \(a,b\) with probabilities \(\frac{1}{2},\frac{1}{2}\) . This profile of behavior strategies is an SPE.

\begin{tcolorbox}
这个较小的博弈显示在图10.15的右侧。因为这个小博弈具有完美信息，所以通过逆向归纳法可以找到一个子博弈完美均衡，这里的逆向归纳法意味着参与者II肯定会选择 \(l\) 。在最终的整体行为策略中，参与者I以概率 \(\frac{2}{3},\frac{1}{3}\) 混合 \(T,B\) ，参与者II肯定会选择 \(l\) ，并以概率 \(\frac{1}{2},\frac{1}{2}\) 混合 \(a,b\) 。这个行为策略组合是一个子博弈完美均衡。
\end{tcolorbox}

\begin{center}
\includegraphics[max width=0.9\textwidth]{images/bo_d2e5cmbef24c73b03tn0_304_197_216_1249_465_0.jpg}
\end{center}
\hspace*{3em} 

Figure 10.15 The subgame of the game in Figure 10.14, and its unique mixed equilibrium with payoffs 2 and 4. In the full game, the subgame is replaced by these payoffs, as shown on the right, so that the optimal move of player II at the root is \(l\) . This defines the SPE.

\begin{tcolorbox}
图10.15 图10.14中博弈的子博弈，以及其唯一的混合均衡，收益为2和4。在整个博弈中，如图右侧所示，子博弈被这些收益所替代，因此参与者II在根节点的最优行动是 \(l\) 。这就定义了子博弈完美均衡。
\end{tcolorbox}

The step of computing an equilibrium of a subgame and substituting the resulting equilibrium payoffs for the subgame generalizes the backward induction step in games with perfect information. In games with perfect information, this backward induction step amounts to finding an optimal move of a single player in a very simple subgame where each move leads directly to a payoff. The optimal move defines an equilibrium in that subgame. In a game with imperfect information, the subgame may involve more than one player, and the equilibrium is usually more difficult to find.

\begin{tcolorbox}
计算子博弈的均衡并用得到的均衡收益替代子博弈这一步骤，推广了具有完美信息博弈中的逆向归纳步骤。在具有完美信息的博弈中，这个逆向归纳步骤相当于在一个非常简单的子博弈中找到单个参与者的最优行动，在这个子博弈中每个行动都直接导致一个收益。最优行动定义了该子博弈中的一个均衡。在具有不完美信息的博弈中，子博弈可能涉及多个参与者，并且通常更难找到均衡。
\end{tcolorbox}

The described SPE of the extensive game in Figure 10.10 is the mixed-strategy equilibrium \(\left( {\left( {\frac{2}{3},\frac{1}{3}}\right) ,\left( {\frac{1}{2},\frac{1}{2},0}\right) }\right)\) of the reduced strategic form, which also defines a pair of behavior strategies. Another equilibrium is the pair \(\left( {B,r * }\right)\) of reduced pure strategies. The strategy \(r *\) is only a reduced behavior strategy for player II, because it does not specify the probabilities for the moves \(a\) and \(b\) at her second information set (which is unreachable when she makes move \(r\) ). However, this equilibrium is not subgame-perfect, because we have already determined the SPE, which in this game is unique.

\begin{tcolorbox}
图10.10中扩展式博弈所描述的子博弈完美均衡是精简策略型的混合策略均衡 \(\left( {\left( {\frac{2}{3},\frac{1}{3}}\right) ,\left( {\frac{1}{2},\frac{1}{2},0}\right) }\right)\) ，它也定义了一对行为策略。另一个均衡是精简纯策略对 \(\left( {B,r * }\right)\) 。策略 \(r *\) 只是参与者II的一个精简行为策略，因为它没有指定她在第二个信息集(当她采取行动 \(r\) 时该信息集无法到达)的行动 \(a\) 和 \(b\) 的概率。然而，这个均衡不是子博弈完美均衡，因为我们已经确定了该博弈中唯一的子博弈完美均衡。
\end{tcolorbox}

In general, just as for games with perfect information, an equilibrium of an extensive game can only be called subgame-perfect if each player's behavior (or pure) strategy is not reduced, but is fully specified and defines a probability distribution (or deterministic move) for each information set of the player.

\begin{tcolorbox}
一般来说，就像具有完美信息的博弈一样，只有当每个参与者的行为(或纯)策略不是精简的，而是完全确定的，并且为参与者的每个信息集定义了一个概率分布(或确定性行动)时，扩展式博弈的一个均衡才能被称为子博弈完美均衡。
\end{tcolorbox}

We described for the example in Figure 10.14 how to find an SPE, which generalizes the backward induction process. We describe this generally in the proof of the following theorem.

\begin{tcolorbox}
我们针对图10.14中的例子描述了如何找到子博弈完美均衡，这推广了逆向归纳过程。我们将在下面定理的证明中对其进行一般性描述。
\end{tcolorbox}

Theorem 10.7. Any extensive game with perfect recall has an SPE in behavior strategies.

\begin{tcolorbox}
定理10.7。任何具有完美记忆的扩展式博弈在行为策略中都存在子博弈完美均衡。
\end{tcolorbox}

Proof. If the game tree consists of a single node, which is simultaneously the root of the game tree and its only leaf with a payoff to each player, then the theorem holds trivially because every player has only one strategy (doing nothing) and this defines an equilibrium of the game, which is also trivially subgame-perfect.

\begin{tcolorbox}
证明。如果博弈树仅由一个节点组成，该节点同时是博弈树的根节点和唯一的叶节点，且每个玩家都有一个收益，那么该定理显然成立，因为每个玩家只有一种策略(什么都不做)，这就定义了博弈的一个均衡，并且这个均衡显然也是子博弈完美均衡。
\end{tcolorbox}

Hence, assume that the game tree does not just consist of a single node. A SPE of the game is then found inductively as as follows. First, find a subgame of the game with root \(u\) that is not just a leaf and that has itself no further nontrivial subgames. Such a subgame exists, because if this does not hold for the root of the entire game, then consider a nontrivial subgame with root \(u\) instead, and continue in that manner until a subgame that has no further subgames is found.

\begin{tcolorbox}
因此，假设博弈树不仅仅由一个节点组成。然后按如下方式归纳地找到该博弈的一个子博弈完美均衡(SPE)。首先，找到该博弈的一个以 \(u\) 为根节点的子博弈，该子博弈不仅仅是一个叶节点，并且它自身没有进一步的非平凡子博弈。这样的子博弈是存在的，因为如果整个博弈的根节点不满足这个条件，那么考虑一个以 \(u\) 为根节点的非平凡子博弈，然后继续这样操作，直到找到一个没有进一步子博弈的子博弈。
\end{tcolorbox}

Second, consider an equilibrium of this subgame, which exists by Nash's Theorem 6.5, and which can be represented using behavior strategies by Kuhn's Theorem 10.4. This is an SPE of the subgame because the subgame has itself no subgames. In the entire game, replace the root \(u\) of the subgame with a leaf whose payoffs are the expected payoffs in the equilibrium of the subgame. The resulting game tree then has fewer nodes, so that this process can be repeated, until \(u\) is the root of the original game. The resulting behavior strategies are then combined to define a profile of behavior strategies for the original game that by construction is subgame-perfect.

\begin{tcolorbox}
其次，考虑这个子博弈的一个均衡，根据纳什定理6.5，这个均衡是存在的，并且根据库恩定理10.4，可以用行为策略来表示。这是该子博弈的一个子博弈完美均衡，因为这个子博弈自身没有子博弈。在整个博弈中，用一个叶节点替换该子博弈的根节点 \(u\) ，该叶节点的收益是子博弈均衡中的期望效用。这样得到的博弈树的节点数就会减少，因此这个过程可以重复进行，直到 \(u\) 是原博弈的根节点。然后将得到的行为策略组合起来，为原博弈定义一个行为策略组合，根据构造，这个组合是子博弈完美的。
\end{tcolorbox}

\begin{center}
\includegraphics[max width=0.9\textwidth]{images/bo_d2e5cmbef24c73b03tn0_305_390_1273_1145_551_0.jpg}
\end{center}
\hspace*{3em} 

Figure 10.16 Left: Zero-sum game, with payoffs to player I, and its nontrivial sub-games. Right: Strategic form of the left subgame, and its optimal mixed strategies.

\begin{tcolorbox}
图10.16 左:零和博弈，显示玩家I的收益，以及它的非平凡子博弈。右:左子博弈的策略形式，以及它的最优混合策略。
\end{tcolorbox}

We illustrate the use of subgames for the example in Figure 10.16, which is a zero-sum game; payoffs are to player I. The game has two nontrivial subgames, each of which starts with a move of player II. In the right subgame, in an SPE player II as the minimizer clearly chooses move \(c\) with payoff 0 to player II. The left subgame is a \(2 \times  2\) game with the strategic form shown on the right in Figure 10.16. This game has a unique mixed equilibrium, found with the "difference trick" as \(\left( {\left( {\frac{3}{5},\frac{2}{5}}\right) ,\left( {\frac{1}{5},\frac{4}{5}}\right) }\right)\) and payoff \(\frac{12}{5} = {2.4}\) to player I, and -2.4 to player II. It describes the optimal behavior of the players in the left subgame, and is part of a behavior strategy for the whole game. At the root of the game tree, player I will choose \(L\) , as this gives him the higher payoff 2.4, which is the unique value of the game because the game is zero-sum.

\begin{tcolorbox}
我们以图10.16中的例子来说明子博弈的使用，这是一个零和博弈；收益是针对玩家I的。该博弈有两个非平凡子博弈，每个子博弈都从玩家II的一次行动开始。在右子博弈中，在一个子博弈完美均衡中，作为最小值点的玩家II显然会选择行动 \(c\) ，此时玩家II的收益为0。左子博弈是一个 \(2 \times  2\) 博弈，其策略形式如图10.16右侧所示。这个博弈有一个唯一的混合均衡，使用“差分技巧”得到为 \(\left( {\left( {\frac{3}{5},\frac{2}{5}}\right) ,\left( {\frac{1}{5},\frac{4}{5}}\right) }\right)\) ，玩家I的收益为 \(\frac{12}{5} = {2.4}\) ，玩家II的收益为 -2.4。它描述了玩家在左子博弈中的最优行为，并且是整个博弈行为策略的一部分。在博弈树的根节点，玩家I会选择 \(L\) ，因为这会给他带来更高的收益2.4，由于该博弈是零和博弈，这是该博弈的唯一值。
\end{tcolorbox}

It is also possible to find all equilibria of this game, including those that are not subgame-perfect, by considering the game tree and without creating the strategic form. Because the game is zero-sum, its payoff is the max-min payoff to each player in any equilibrium, and therefore 2.4 for player I and -2.4 for player II. If \(R\) is played with positive probability, then the minimizing player II has to play \(c\) because otherwise her expected payoff would be lower. With player II choosing \(c\) , this would lower player I’s payoff compared to playing \(L\) and then the max-min strategy that mixes \(S\) and \(T\) as described above. Hence, \(L\) is played for certain in any equilibrium, with player I's behavior strategy as described before. However, this means that the right decision node of player II is never reached, and player II can mix between \(c\) and \(d\) as long as the expected payoff to player I does not exceed 2.4, that is, as long as \(\operatorname{prob}\left( d\right)  \leq  {0.6}\) . This behavior at the right node, plus the behavior for mixing \(a,b\) with probabilities \(\frac{1}{5},\frac{4}{5}\) as described above, is the set of behavior strategies of player II in an arbitrary equilibrium of the game.

\begin{tcolorbox}
也可以通过考虑博弈树而不构建策略形式来找到这个博弈的所有均衡，包括那些不是子博弈完美的均衡。因为该博弈是零和博弈，在任何均衡中，它的收益是每个玩家的极大极小收益，因此玩家I的收益是2.4，玩家II的收益是 -2.4。如果以正概率选择 \(R\) ，那么作为最小化者的玩家II必须选择 \(c\) ，否则她的期望效用会更低。如果玩家II选择 \(c\) ，与选择 \(L\) 然后采用上述混合 \(S\) 和 \(T\) 的极大极小策略相比，这会降低玩家I的收益。因此，在任何均衡中一定会选择 \(L\) ，玩家I的行为策略如前所述。然而，这意味着玩家II的右决策节点永远不会被到达，玩家II可以在 \(c\) 和 \(d\) 之间进行混合，只要玩家I的期望效用不超过2.4，即只要 \(\operatorname{prob}\left( d\right)  \leq  {0.6}\) 。在右节点的这种行为，加上上述以概率 \(\frac{1}{5},\frac{4}{5}\) 混合 \(a,b\) 的行为，就是该博弈任意均衡中玩家II的行为策略集。
\end{tcolorbox}

\section*{10.11 Further Reading}

\begin{tcolorbox}
\section*{10.11 进一步阅读}
\end{tcolorbox}

The example of the competition between software firms in Section 10.2 is adapted from Turocy and von Stengel (2002).

\begin{tcolorbox}
10.2节中软件公司之间竞争的例子改编自图罗西和冯·施滕格尔(2002年)的研究。
\end{tcolorbox}

Information sets were introduced by Harold Kuhn (1953). Realization equivalent strategies are just called "equivalent" by Kuhn. His Definition 17 states that player \(i\) has perfect recall if for every pure strategy \(\pi\) of player \(i\) and every information set \(h\) of player \(i\) that is reachable ("relevant") when playing \(\pi\) , all nodes in \(h\) are reachable ("possible") when playing \(\pi\) . This is equivalent to our Definition 10.1 of perfect recall, which we have adapted from Selten (1975, p. 27).

\begin{tcolorbox}
信息集(information set)由哈罗德·库恩(Harold Kuhn，1953 年)引入。库恩将实现等价策略简称为“等价”策略。他的定义 17 指出，如果对于玩家 \(i\) 的每个纯策略 \(\pi\) 以及玩家 \(i\) 的每个在采用 \(\pi\) 策略时可达(“相关”)的信息集 \(h\) ，在采用 \(\pi\) 策略时， \(h\) 中的所有节点都是可达(“可能”)的，那么玩家 \(i\) 具有完美记忆(perfect recall)。这与我们从塞尔滕(Selten，1975 年，第 27 页)改编而来的完美记忆的定义 10.1 等价。
\end{tcolorbox}

The concept of a reduced behavior strategy is not standard, but we have introduced it here because it corresponds to a behavior strategy in the same way as a reduced pure strategy corresponds to a pure strategy.

\begin{tcolorbox}
约简行为策略(reduced behavior strategy)的概念并非标准概念，但我们在此引入它，是因为它与行为策略的对应关系，就如同约简纯策略与纯策略的对应关系一样。
\end{tcolorbox}

Other textbook treatments of extensive games with imperfect information are Osborne and Rubinstein (1994, Chapter 11) and Gibbons (1992, Section 2.4).

\begin{tcolorbox}
其他关于不完美信息扩展式博弈的教科书处理方式可参考奥斯本和鲁宾斯坦(Osborne and Rubinstein，1994 年，第 11 章)以及吉本斯(Gibbons，1992 年，第 2.4 节)的著作。
\end{tcolorbox}

Our proof of Kuhn's Theorem 10.4 is based on sequences of moves. These sequences play a central role in the sequence form of von Stengel (1996). This is a strategic description that instead of pure strategies considers sequences \(\sigma\) , which are played with probabilities \(\mu \left\lbrack  \sigma \right\rbrack\) that fulfill the equations (10.6) for all \(h \in  {H}_{i}\) and (10.5). Whereas the strategic form may be exponentially larger than the extensive game, the sequence form has the same size. A zero-sum game, for example, can then be solved with the help of a linear program of the same size as the game tree.

\begin{tcolorbox}
我们对库恩定理 10.4 的证明基于行动序列。这些序列在冯·施滕格尔(von Stengel，1996 年)的序列形式(sequence form)中起着核心作用。这是一种策略性描述，它不考虑纯策略，而是考虑序列 \(\sigma\) ，这些序列以概率 \(\mu \left\lbrack  \sigma \right\rbrack\) 被采用，且对于所有 \(h \in  {H}_{i}\) 满足方程(10.6)和(10.5)。虽然策略形式可能比扩展式博弈大指数倍，但序列形式的规模与扩展式博弈相同。例如，一个零和博弈(zero - sum game)可以借助与博弈树(game tree)规模相同的线性规划来求解。
\end{tcolorbox}

The Monty Hall problem was posed in Selvin (1975). Its behavioral assumptions are discussed and tested in Friedman (1998). Even though our game-theoretic analysis suggests that "not switching doors" is very justifiable when Monty Hall does not always offer a switch, this is still considered a "choice anomaly".

\begin{tcolorbox}
蒙提·霍尔问题由塞尔文(Selvin，1975 年)提出。弗里德曼(Friedman，1998 年)对其行为假设进行了讨论和检验。尽管我们的博弈论分析表明，当蒙提·霍尔并非总是提供换门选项时，“不换门”是非常合理的选择，但这仍然被视为一种“选择异常”。
\end{tcolorbox}

\section*{10.12 Exercises for Chapter 10}

\begin{tcolorbox}
\section*{10.12 第 10 章习题}
\end{tcolorbox}

Exercise 10.1. Explain the two ways, with either player I or player II moving first, of representing an \(m \times  n\) bimatrix game as an extensive game that has the given game as its strategic form. Giving an example may be helpful. In both cases, how many decision nodes do the players have? What is the number of terminal nodes of the game tree? Why do we need information sets here?

\begin{tcolorbox}
习题 10.1. 解释将一个 \(m \times  n\) 双矩阵博弈(bimatrix game)表示为扩展式博弈的两种方式，即玩家 I 先行动或玩家 II 先行动，使得该扩展式博弈的策略形式就是给定的双矩阵博弈。举例说明可能会有所帮助。在这两种情况下，玩家各有多少个决策节点(decision node)？博弈树的终端节点(terminal node)数量是多少？为什么这里需要信息集？
\end{tcolorbox}

Exercise 10.2. Which of the following extensive games have perfect recall and if not, why not? For each extensive game with perfect recall, find all its equilibria in mixed (including pure) strategies.

\begin{tcolorbox}
习题 10.2. 以下哪些扩展式博弈具有完美记忆？如果不具有，原因是什么？对于每个具有完美记忆的扩展式博弈，找出其所有混合策略(包括纯策略)均衡(equilibrium)。
\end{tcolorbox}

\begin{center}
\includegraphics[max width=0.9\textwidth]{images/bo_d2e5cmbef24c73b03tn0_307_334_1232_1260_726_0.jpg}
\end{center}
\hspace*{3em} 

Exercise 10.3. Which of the following extensive games have perfect recall and if not, why not? For each extensive game with perfect recall, find all its equilibria in mixed (including pure) strategies.

\begin{tcolorbox}
习题 10.3. 以下哪些扩展式博弈具有完美记忆？如果不具有，原因是什么？对于每个具有完美记忆的扩展式博弈，找出其所有混合策略(包括纯策略)均衡。
\end{tcolorbox}

\begin{center}
\includegraphics[max width=0.9\textwidth]{images/bo_d2e5cmbef24c73b03tn0_308_197_216_1246_590_0.jpg}
\end{center}
\hspace*{3em} 

Exercise 10.4. Consider the following zero-sum game, a simplified version of Poker adapted from Kuhn (1950). A deck has three cards (of rank High, Middle, and Low), and each player is dealt a card. All deals are equally likely, and of course the players get different cards. A player does not know the card dealt to the other player. After seeing his hand, player I has the option to Raise(R)or Fold(F). When he folds, he loses one unit to player II. When he raises, player II has the option to meet(m)or pass(p). When player II chooses "pass", she has to pay one unit to player I. When player II chooses "meet", the higher card wins, and the player with the lower card has to pay two units to the winning player.

\begin{tcolorbox}
习题 10.4. 考虑以下零和博弈，这是一个改编自库恩(Kuhn，1950 年)的简化版扑克(Poker)游戏。一副牌有三张牌(高、中、低等级)，每个玩家发一张牌。所有发牌情况的可能性相同，当然玩家拿到的牌不同。玩家不知道对方拿到的牌。玩家 I 看到自己的牌后，有加注(R)或弃牌(F)的选项。当他弃牌时，他输给玩家 II 一个单位。当他加注时，玩家 II 有跟注(m)或弃牌(p)的选项。当玩家 II 选择“弃牌”时，她要付给玩家 I 一个单位。当玩家 II 选择“跟注”时，牌面大的玩家获胜，牌面小的玩家要付给获胜玩家两个单位。
\end{tcolorbox}

(a) Draw a game in extensive form that models this game, with information sets, and payoffs to player I at the leaves.

\begin{tcolorbox}
(a) 绘制一个具有信息集的扩展式博弈来模拟这个游戏，并在叶节点(leaf)标注玩家 I 的收益。
\end{tcolorbox}

(b) Simplify this game by assuming that at an information set where a player's move is always at least as good as the other move, no matter what the other player's move or the chance move were or will be, then the player will choose that move. Draw the simplified extensive form game.

\begin{tcolorbox}
(b) 通过假设在某个信息集上，无论另一个玩家的行动或随机行动是什么，一个玩家的某个行动至少和另一个行动一样好，那么该玩家会选择那个行动，来简化这个游戏。绘制简化后的扩展式博弈。
\end{tcolorbox}

(c) What does the simplification in (b) have to do with weakly dominated strategies? Why is the simplification legitimate here?

\begin{tcolorbox}
(c) (b) 中的简化与弱被占优策略(weakly dominated strategy)有什么关系？为什么这里的简化是合理的？
\end{tcolorbox}

(d) Give the strategic form of the simplified game from (b) and find an equilibrium of that game. What are the behavior probabilities of the players in that equilibrium? What is the unique payoff to the players in equilibrium?

\begin{tcolorbox}
(d) 给出 (b) 中简化后游戏的策略形式，并找出该游戏的一个均衡。在该均衡中，玩家的行为概率是多少？均衡时玩家的唯一收益是多少？
\end{tcolorbox}

Exercise 10.5. For an extensive game with perfect recall, show that the information set that contains the root (starting node) of any subgame is always a singleton (that is, contains only that node). [Hint: Use Proposition 10.2.]

\begin{tcolorbox}
习题 10.5. 对于一个具有完美记忆的扩展式博弈，证明包含任何子博弈(subgame)的根节点(起始节点)的信息集总是单元素集(即只包含该节点)。[提示:使用命题 10.2。]
\end{tcolorbox}

Exercise 10.6. Find all equilibria in (pure or) mixed strategies of this extensive game. Which of them are subgame-perfect?

\begin{tcolorbox}
练习10.6。找出这个扩展式博弈的所有(纯策略或)混合策略均衡。其中哪些是子博弈完美均衡？
\end{tcolorbox}

\begin{center}
\includegraphics[max width=0.5\textwidth]{images/bo_d2e5cmbef24c73b03tn0_309_942_210_652_1243_0.jpg}
\end{center}
\hspace*{3em} 

Exercise 10.7. Consider this three-player game in extensive form. At a terminal node, the top, middle, and bottom payoffs are to player1,2,3, respectively. Find all subgame-perfect equilibria (SPE) in mixed strategies.

\begin{tcolorbox}
练习10.7。考虑这个扩展式三人博弈。在终端节点处，顶部、中部和底部的收益分别属于玩家1、2、3。找出所有混合策略的子博弈完美均衡(SPE)。
\end{tcolorbox}

Exercise 10.8. Consider this three-player game in extensive form. Find all subgame-perfect equilibria (SPE) in mixed strategies.

\begin{tcolorbox}
练习10.8。考虑这个扩展式三人博弈。找出所有混合策略的子博弈完美均衡(SPE)。
\end{tcolorbox}

\begin{center}
\includegraphics[max width=0.5\textwidth]{images/bo_d2e5cmbef24c73b03tn0_309_952_1473_650_617_0.jpg}
\end{center}
\hspace*{3em} 

\section*{11}

\section*{Bargaining}

\begin{tcolorbox}
\section*{讨价还价}
\end{tcolorbox}

This chapter presents models of bargaining. Bargaining occurs in real life between the buyer and seller of a house, for example. This is a game-theoretic problem because it is useful to think about the situation of the other player. Both players have something to gain by reaching an agreement.

\begin{tcolorbox}
本章介绍讨价还价模型。讨价还价在现实生活中很常见，比如房屋买卖双方之间的讨价还价。这是一个博弈论问题，因为考虑对方的处境很有必要。双方都能通过达成协议获得收益。
\end{tcolorbox}

A basic bargaining model is that of "splitting a pie" in a way that is acceptable to both players, but where within that range a player can only gain at the expense of the other. The question is how to split the pie fairly. In his undergraduate thesis, later published in Nash (1950), John Nash proposed a system of "axioms", or natural conditions, for arriving at a fair "bargaining solution". Nash showed that these axioms always entail a unique solution, which is also easy to calculate. Section 11.2 shows how a bargaining set arises from a bimatrix game, here a variant of the Prisoner's Dilemma, where the players can benefit from cooperation. General bargaining sets and corresponding axioms are presented in Sections 11.3 and 11.4, with a geometric characterization of the bargaining solution in Section 11.5.

\begin{tcolorbox}
一个基本的讨价还价模型是“分馅饼”，要以双方都能接受的方式进行分配，但在这个范围内，一方的收益必然意味着另一方的损失。问题是如何公平地分割馅饼。约翰·纳什(John Nash)在他后来发表于《纳什(1950)》的本科论文中，提出了一套“公理”，即自然条件，以达成公平的“讨价还价解”。纳什证明，这些公理总能得出唯一解，而且这个解也很容易计算。第11.2节展示了如何从双矩阵博弈(这里是囚徒困境的一个变体)中产生讨价还价集，在这个博弈中，玩家可以从合作中受益。第11.3节和第11.4节介绍了一般的讨价还价集和相应的公理，第11.5节对讨价还价解进行了几何刻画。
\end{tcolorbox}

Nash's bargaining model is not a non-cooperative game in strategic or extensive form that we have considered so far. Rather, it belongs to cooperative game theory, where players are assumed to be able to agree in some way and then adhere to their agreement. The model describes which agreements are possible and how the players benefit from them, but does not describe individual actions that lead to this agreement, nor an enforcement mechanism. The typical approach of cooperative game theory is to propose solutions to the game (ideally a unique solution, such as a specific agreement) that have to fulfill certain axioms. The Nash bargaining solution is of this kind.

\begin{tcolorbox}
纳什的讨价还价模型不是我们目前所考虑的策略式或扩展式非合作博弈。相反，它属于合作博弈论，在该理论中，假设玩家能够以某种方式达成协议，然后遵守协议。该模型描述了哪些协议是可行的，以及玩家如何从中受益，但没有描述达成该协议的个体行动，也没有描述执行机制。合作博弈论的典型方法是提出满足某些公理的博弈解(理想情况下是唯一解，如特定协议)。纳什讨价还价解就是这种类型。
\end{tcolorbox}

The later sections of this chapter go back to a non-cooperative bargaining model, due to Rubinstein (1982). This model uses a game tree with perfect information where the players make alternating offers that the other player can accept or reject. The analysis of this game uses the study of commitment games from Section 4.7. The most basic form, described in Section 11.7, is the "Ultimatum game" that consists of a single take-it-or-leave-it offer which can only be accepted or rejected.

\begin{tcolorbox}
本章后面的章节回到了鲁宾斯坦(Rubinstein，1982)提出的非合作讨价还价模型。该模型使用具有完美信息的博弈树，玩家轮流提出报价，另一方可以接受或拒绝。对这个博弈的分析运用了第4.7节中对承诺博弈的研究。第11.7节描述的最基本形式是“最后通牒博弈”，它由一个要么接受要么放弃的单一报价组成，对方只能选择接受或拒绝。
\end{tcolorbox}

With alternating offers over several rounds, described in Sections 11.8 and 11.9, there is a chance at each stage that the game does not proceed to the next round so that players wind up with nothing because they could not agree. The effect of this is a "discount factor" which quantifies to what extent an earlier agreement is better that a later one. By thinking ahead, the optimal strategies of the players define sequences of alternating offers, with agreement in the first round. This can be done with a fixed number of rounds, or with "stationary strategies" for an infinite number of bargaining rounds where, in essence, the game repeats every two rounds when it comes back to the same player to make an offer (see Section 11.10). The nice conclusion of this alternating-offers bargaining model, shown in Section 11.11, is that a discount factor that represents very patient players leads to the same outcome as Nash's axiomatic bargaining solution. Nash's "cooperative" solution concept is thereby justified with a model of self-interested behavior.

\begin{tcolorbox}
在第11.8节和第11.9节描述的多轮交替报价中，每一轮都有可能因为双方无法达成协议而导致博弈无法进入下一轮，最终双方一无所获。这产生了一个“贴现因子”，它量化了早期达成协议比后期达成协议好的程度。通过前瞻性思考，玩家的最优策略定义了交替报价的序列，并在第一轮就达成协议。这可以在固定轮数的情况下实现，也可以通过“平稳策略”在无限轮次的讨价还价中实现，本质上，当轮到同一玩家报价时，博弈每两轮重复一次(见第11.10节)。第11.11节展示的这个交替报价讨价还价模型的一个很好的结论是，代表非常有耐心的玩家的贴现因子会导致与纳什公理讨价还价解相同的结果。这样，纳什的“合作”解概念就可以通过自利行为模型得到证明。
\end{tcolorbox}

\section*{11.1 Prerequisites and Learning Outcomes}

\begin{tcolorbox}
\section*{11.1 先决条件和学习成果}
\end{tcolorbox}

Prerequisites for this chapter are: two-player games in strategic form (Chapter 3), the concept that utility functions admit positive-affine transformations (Section 5.6), concave utility functions (Section 5.9), commitment games (Section 4.7), and max-min strategies (Section 8.2).

\begin{tcolorbox}
本章的先决条件包括:策略式两人博弈(第3章)、效用函数允许正仿射变换的概念(第5.6节)、凹效用函数(第5.9节)、承诺博弈(第4.7节)和极大极小策略(第8.2节)。
\end{tcolorbox}

After studying this chapter, you should be able to:

\begin{tcolorbox}
学习本章后，你应该能够:
\end{tcolorbox}

\begin{itemize}
\item explain the concepts of bargaining sets, bargaining axioms, Nash bargaining solution, threat point, and Nash product;
\end{itemize}

\begin{tcolorbox}
\begin{itemize}
\item 解释讨价还价集、讨价还价公理、纳什讨价还价解、威胁点和纳什积的概念；
\end{itemize}
\end{tcolorbox}

\begin{itemize}
\item draw bargaining sets derived from bimatrix games, and for bargaining over a "unit pie" when players have concave utility functions;
\end{itemize}

\begin{tcolorbox}
\begin{itemize}
\item 绘制从双矩阵博弈导出的讨价还价集，以及当玩家具有凹效用函数时对“单位馅饼”进行讨价还价的讨价还价集；
\end{itemize}
\end{tcolorbox}

\begin{itemize}
\item explain, and solve specific examples of, the model of bargaining over multiple rounds of demands and counter-demands, and of stationary SPE for infinitely many rounds.
\end{itemize}

\begin{tcolorbox}
\begin{itemize}
\item 解释并解决多轮需求和反需求的讨价还价模型的具体例子，以及无限轮次的平稳子博弈完美均衡。
\end{itemize}
\end{tcolorbox}

\section*{11.2 Bargaining Sets}

\begin{tcolorbox}
\section*{11.2 讨价还价集}
\end{tcolorbox}

Consider the variant of the Prisoner's Dilemma game on the left in Figure 11.1. Player I’s strategy \(B\) dominates \(T\) , and \(r\) dominates \(l\) , so that the unique equilibrium of the game is(B, r)with payoff 1 to both players. The outcome(T, l)would give payoffs 2 and 3 to players I and II, which is clearly better.

\begin{tcolorbox}
考虑图11.1左侧的囚徒困境博弈的变体。玩家I的策略 \(B\) 优于 \(T\) ，且 \(r\) 优于 \(l\) ，因此该博弈的唯一均衡是(B，r)，两个玩家的收益均为1。结果(T，l)将使玩家I和玩家II的收益分别为2和3，这显然更好。
\end{tcolorbox}

Suppose now that the players can talk to each other and have agreed on a certain course of action that they intend to follow. We do not model how the players can talk, agree, or enforce their agreement, but instead try to give guidelines about which agreement they should reach.

\begin{tcolorbox}
现在假设玩家可以相互交流，并就他们打算遵循的某种行动方案达成了一致。我们不建模玩家如何交流、达成协议或执行他们的协议，而是试图给出他们应该达成何种协议的指导原则。
\end{tcolorbox}

The first step is to model what players can achieve by an agreement, by describing a "bargaining set" of possible payoff pairs that they can bargain over. Suppose that the players can agree to play any of the cells of the payoff table in Figure 11.1. The resulting utility pairs(u, v)are points in a two-dimensional diagram with horizontal coordinate \(u\) (utility to player I) and vertical coordinate \(v\) (utility to player II). These points are shown on the right in Figure 11.1, marked with the respective strategy pairs. For example, \(T,l\) marks the utility pair(2,3).

\begin{tcolorbox}
第一步是通过描述一个玩家可以讨价还价的可能收益对的“讨价还价集”，来建模玩家通过协议可以实现的目标。假设玩家可以就图11.1收益表中的任何单元格的策略达成一致。由此产生的效用对(u，v)是二维图中的点，其中横坐标 \(u\) (玩家I的效用)，纵坐标 \(v\) (玩家II的效用)。这些点显示在图11.1的右侧，用相应的策略对标记。例如， \(T,l\) 标记效用对(2，3)。
\end{tcolorbox}

\begin{center}
\includegraphics[max width=0.9\textwidth]{images/bo_d2e5cmbef24c73b03tn0_312_186_677_1245_606_0.jpg}
\end{center}
\hspace*{3em} 

Figure 11.1 Left: \(2 \times  2\) game similar to the Prisoner’s Dilemma. Right: corresponding bargaining set \(S\) (dark gray), and its bargaining solution(U, V).

\begin{tcolorbox}
图11.1 左侧: \(2 \times  2\) 类似于囚徒困境的博弈。右侧:相应的讨价还价集 \(S\) (深灰色)及其讨价还价解(U，V)。
\end{tcolorbox}

The convex hull (set of convex combinations) of these utility pairs is shown in light and dark gray in Figure 11.1. Any such convex combination is obtained by a joint lottery that the players have decided to play as a possible result of their agreement. For example, the utility pair(1.5,2)would be the expected utility if the players have decided to toss a fair coin, and, depending on the outcome, to either play(T, l)or else(B, r); we assume that this coin-toss and subsequent joint action is part of the players' agreement and they will follow through with it.

\begin{tcolorbox}
这些效用对的凸包(凸组合的集合)在图11.1中以浅灰色和深灰色显示。任何这样的凸组合都是通过玩家作为协议的可能结果而决定进行的联合彩票获得的。例如，如果玩家决定抛一枚公平的硬币，并根据结果选择玩(T，l)或(B，r)，那么效用对(1.5，2)将是期望效用；我们假设这次抛硬币以及随后的联合行动是玩家协议的一部分，并且他们会执行该协议。
\end{tcolorbox}

Making the bargaining set convex with the help of lotteries is not as strange as it sounds. For example, if two people find something that both would like to have but that cannot be divided, tossing a coin to decide who gets it is often accepted as a fair solution.

\begin{tcolorbox}
借助彩票使讨价还价集成为凸集并不像听起来那么奇怪。例如，如果两个人发现了一件他们都想要但无法分割的东西，抛硬币决定谁得到它通常被认为是一种公平的解决方案。
\end{tcolorbox}

Not every point in the convex hull is "reasonable" because we always assume that the players have the option not to agree to it in advance. For example, player I would not agree that both players choose with certainty the strategy pair(T, r), where player I gets payoff zero. The reason is that without any agreement, player I can use his max-min strategy \(B\) that guarantees him a payoff of 1 irrespective of what the other player does. As shown in Chapter 8 on zero-sum games, a max-min strategy may be a mixed strategy; in the Prisoner's Dilemma game in Figure 11.1 it is a pure strategy. We ask that every agreed outcome of the bargaining situation should give each player at least his max-min payoff (which is, as usual, defined in terms of that player's own payoffs).

\begin{tcolorbox}
并非凸包中的每个点都是“合理的”，因为我们总是假设玩家有权事先不同意某个点。例如，玩家I不会同意两个玩家确定地选择策略对(T，r)，在这种情况下玩家I的收益为零。原因是在没有任何协议的情况下，玩家I可以使用他的极大极小策略 \(B\) ，无论另一个玩家做什么，该策略都能保证他获得1的收益。正如第8章关于零和博弈中所示，极大极小策略可能是一种混合策略；在图11.1的囚徒困境博弈中，它是一种纯策略。我们要求讨价还价情况的每个商定结果都应至少给每个玩家提供他的极大极小收益(通常根据该玩家自己的收益来定义)。
\end{tcolorbox}

The bargaining set \(S\) derived from a two-player game in strategic form is thus defined as the convex hull of utility pairs (as specified in the game), with the additional constraint that \(u \geq  {u}_{0}\) and \(v \geq  {v}_{0}\) for all(u, v)in \(S\) , where \({u}_{0}\) is the max-min payoff to player I and \({v}_{0}\) is the max-min payoff to player II. With \({u}_{0} = {v}_{0} = 1\) , the resulting set \(S\) is shown in dark gray in Figure 11.1.

\begin{tcolorbox}
因此，从策略形式的两人博弈中导出的讨价还价集 \(S\) 被定义为效用对的凸包(如博弈中所指定)，并附加约束条件:对于 \(S\) 中的所有(u，v)，有 \(u \geq  {u}_{0}\) 和 \(v \geq  {v}_{0}\) ，其中 \({u}_{0}\) 是玩家I的极大极小收益， \({v}_{0}\) 是玩家II的极大极小收益。通过 \({u}_{0} = {v}_{0} = 1\) ，得到的集合 \(S\) 在图11.1中以深灰色显示。
\end{tcolorbox}

In this example, the "Nash bargaining solution" (which is not to be confused with "Nash equilibrium", a very different concept) will give the utility pair \(\left( {U,V}\right)  = \left( {{2.5},{2.5}}\right)\) , as we will explain (the fact that \(U = V\) is a coincidence). This expected payoff pair is a convex combination of the utility pairs(2,3)and(5,0). It is achieved by a joint lottery that chooses(T, l)with probability \(\frac{5}{6}\) and(B, l)with probability \(\frac{1}{6}\) . In effect, player II always plays \(l\) and player I chooses between \(T\) and \(B\) with probabilities \(\frac{5}{6}\) and \(\frac{1}{6}\) . However, this is not a "mixed strategy" but rather a lottery performed by an independent device, with an outcome that both players have in advance decided to accept.

\begin{tcolorbox}
在这个例子中，“纳什讨价还价解”(不要将其与“纳什均衡”相混淆，这是一个截然不同的概念)将给出效用对 \(\left( {U,V}\right)  = \left( {{2.5},{2.5}}\right)\) ，我们稍后会解释( \(U = V\) 只是巧合)。这个期望效用对是效用对 (2,3) 和 (5,0) 的凸组合。它是通过一个联合彩票实现的，该彩票以概率 \(\frac{5}{6}\) 选择 (T, l) ，以概率 \(\frac{1}{6}\) 选择 (B, l) 。实际上，玩家 II 总是选择 \(l\) ，而玩家 I 以概率 \(\frac{5}{6}\) 和 \(\frac{1}{6}\) 在 \(T\) 和 \(B\) 之间进行选择。然而，这不是“混合策略”，而是由一个独立设备进行的彩票，其结果是两个玩家事先就决定接受的。
\end{tcolorbox}

\section*{11.3 Bargaining Axioms}

\begin{tcolorbox}
\section*{11.3 讨价还价公理}
\end{tcolorbox}

We consider bargaining situations modeled as bargaining sets, and try to obtain a "solution" for any such situation. The conditions about this model are called axioms and concern both the bargaining set and the solution one wants to obtain.

\begin{tcolorbox}
我们将讨价还价情形建模为讨价还价集，并尝试为任何此类情形找到一个“解”。关于这个模型的条件被称为公理，它们既涉及讨价还价集，也涉及我们想要得到的解。
\end{tcolorbox}

Definition 11.1 (Axioms for bargaining set). A bargaining set \(S\) with threat point \(\left( {{u}_{0},{v}_{0}}\right)\) is a nonempty subset of \({\mathbb{R}}^{2}\) so that

\begin{tcolorbox}
定义 11.1(讨价还价集的公理)。具有威胁点 \(\left( {{u}_{0},{v}_{0}}\right)\) 的讨价还价集 \(S\) 是 \({\mathbb{R}}^{2}\) 的一个非空子集，使得
\end{tcolorbox}

(a) \(u \geq  {u}_{0},v \geq  {v}_{0}\) for all \(\left( {u,v}\right)  \in  S\) ;

\begin{tcolorbox}
(a) 对于所有 \(\left( {u,v}\right)  \in  S\) ， \(u \geq  {u}_{0},v \geq  {v}_{0}\) ；
\end{tcolorbox}

(b) \(S\) is compact;

\begin{tcolorbox}
(b) \(S\) 是紧致的；
\end{tcolorbox}

(c) \(S\) is convex.

\begin{tcolorbox}
(c) \(S\) 是凸的。
\end{tcolorbox}

Each pair(u, v)in a bargaining set \(S\) represents a utility \(u\) for player \(\mathrm{I}\) and a utility \(v\) for player II, which the players can achieve by reaching a specific agreement. At least one agreement is possible because \(S\) is nonempty. The threat point \(\left( {{u}_{0},{v}_{0}}\right)\) is the pair of utilities that each player can (in expectation) guarantee for himself. Axiom (b) in Definition 11.1 states that \(S\) is bounded and closed, with the aim of defining a "bargaining solution" by some sort of maximization process. The convexity axiom (c) means that players have, if necessary, access to joint lotteries as part of their agreement.

\begin{tcolorbox}
讨价还价集 \(S\) 中的每一对 (u, v) 代表玩家 \(\mathrm{I}\) 的效用 \(u\) 和玩家 II 的效用 \(v\) ，玩家们可以通过达成特定协议来实现这些效用。因为 \(S\) 是非空的，所以至少有一种协议是可能的。威胁点 \(\left( {{u}_{0},{v}_{0}}\right)\) 是每个玩家(在期望意义上)可以为自己保证的效用对。定义 11.1 中的公理 (b) 表明 \(S\) 是有界且闭的，其目的是通过某种最大化过程来定义一个“讨价还价解”。凸性公理 (c) 意味着，如果有必要，玩家们可以在他们的协议中使用联合彩票。
\end{tcolorbox}

In Section 11.2, we have constructed a bargaining set \(S\) from a bimatrix game (A, B). We show that this set fulfills the axioms in Definition 11.1. This is clear for (b) and (c). By construction,(a) holds, but we have to show that \(S\) is not empty. The threat point \(\left( {{u}_{0},{v}_{0}}\right)\) is obtained with \({u}_{0}\) as the max-min payoff to player I (in terms of his payoffs) and \({v}_{0}\) as the max-min payoff to player II (in terms of her payoffs). Thus, if player I plays his max-min strategy \(\widehat{x}\) and player II plays her max-min strategy \(\widehat{y}\) , then the payoff pair \(\left( {u,v}\right)  = \left( {{\widehat{x}}^{\top }A\widehat{y},{\widehat{x}}^{\top }B\widehat{y}}\right)\) that results in the game is an element of \(S\) because \(u \geq  {u}_{0}\) and \(v \geq  {v}_{0}\) , so \(S\) is not empty.

\begin{tcolorbox}
在第 11.2 节中，我们从一个双矩阵博弈 (A, B) 构建了一个讨价还价集 \(S\) 。我们证明这个集合满足定义 11.1 中的公理。对于 (b) 和 (c) 这是显然的。根据构造，(a) 成立，但我们必须证明 \(S\) 是非空的。威胁点 \(\left( {{u}_{0},{v}_{0}}\right)\) 是这样得到的: \({u}_{0}\) 是玩家 I 的极大极小收益(就他的收益而言)， \({v}_{0}\) 是玩家 II 的极大极小收益(就她的收益而言)。因此，如果玩家 I 采用他的极大极小策略 \(\widehat{x}\) ，玩家 II 采用她的极大极小策略 \(\widehat{y}\) ，那么博弈中得到的收益对 \(\left( {u,v}\right)  = \left( {{\widehat{x}}^{\top }A\widehat{y},{\widehat{x}}^{\top }B\widehat{y}}\right)\) 是 \(S\) 的一个元素，因为 \(u \geq  {u}_{0}\) 和 \(v \geq  {v}_{0}\) ，所以 \(S\) 是非空的。
\end{tcolorbox}

However, in the construction of a bargaining set from a bimatrix game, the threat point \(\left( {{u}_{0},{v}_{0}}\right)\) need not itself be an element of \(S\) , as shown in Exercise 11.2(ii). For this reason, the threat point \(\left( {{u}_{0},{v}_{0}}\right)\) is explicitly stated together with the bargaining set. If the threat point belongs to \(S\) and (a) holds, then there can be no other threat point. The axioms in Definition 11.2 below would be easier to state if the threat point was included in \(S\) , but then we would have to take additional steps to obtain such a bargaining set from a bimatrix game. In our examples and drawings we will typically consider sets \(S\) that include the threat point.

\begin{tcolorbox}
然而，在从双矩阵博弈(bimatrix game)构建讨价还价集时，威胁点 \(\left( {{u}_{0},{v}_{0}}\right)\) 本身不必是 \(S\) 的元素，如练习11.2(ii)所示。出于这个原因，威胁点 \(\left( {{u}_{0},{v}_{0}}\right)\) 会与讨价还价集一起明确给出。如果威胁点属于 \(S\) 且条件(a)成立，那么就不可能有其他威胁点。如果威胁点包含在 \(S\) 中，下面定义11.2中的公理表述起来会更简单，但那样我们就必须采取额外的步骤才能从双矩阵博弈中得到这样一个讨价还价集。在我们的示例和图示中，我们通常会考虑包含威胁点的集合 \(S\) 。
\end{tcolorbox}

From now on, we simply assume that the bargaining set \(S\) fulfills these axioms. We will later consider bargaining sets for other situations, which are specified directly and not from a bimatrix game.

\begin{tcolorbox}
从现在起，我们简单地假设讨价还价集 \(S\) 满足这些公理。我们稍后会考虑其他情况的讨价还价集，这些情况是直接指定的，而非从双矩阵博弈得出。
\end{tcolorbox}

In the following definition, the axioms (d)-(h) continue (a)-(c) from Definition 11.1. We discuss them afterwards.

\begin{tcolorbox}
在下面的定义中，公理(d) - (h)是定义11.1中(a) - (c)的延续。我们随后会对它们进行讨论。
\end{tcolorbox}

Definition 11.2 (Axioms for bargaining solution). For a given bargaining set \(S\) with threat point \(\left( {{u}_{0},{v}_{0}}\right)\) , a bargaining solution \(N\left( S\right)\) is a pair(U, V)so that

\begin{tcolorbox}
定义11.2(讨价还价解(bargaining solution)的公理)。对于给定的具有威胁点 \(\left( {{u}_{0},{v}_{0}}\right)\) 的讨价还价集 \(S\) ，讨价还价解 \(N\left( S\right)\) 是一个二元组(U, V)，使得
\end{tcolorbox}

(d) \(\left( {U,V}\right)  \in  S\) ;

(e) the solution(U, V)is Pareto-optimal, that is, for all \(\left( {u,v}\right)  \in  S\) , if \(u \geq  U\) and \(v \geq  V\) , then \(\left( {u,v}\right)  = \left( {U,V}\right)\) ;

\begin{tcolorbox}
(e) 解(U, V)是帕累托最优(Pareto - optimal)的，即对于所有 \(\left( {u,v}\right)  \in  S\) ，如果 \(u \geq  U\) 且 \(v \geq  V\) ，那么 \(\left( {u,v}\right)  = \left( {U,V}\right)\) ；
\end{tcolorbox}

(f) it is invariant under utility scaling, that is, if \(a,c > 0\) and \(b,d \in  \mathbb{R}\) and \({S}^{\prime }\) is the bargaining set \(\{ \left( {{au} + b,{cv} + d}\right)  \mid  \left( {u,v}\right)  \in  S\}\) with threat point \(\left( {a{u}_{0} + b,c{v}_{0} + d}\right)\) , then \(N\left( {S}^{\prime }\right)  = \left( {{aU} + b,{cV} + d}\right)\) ;

\begin{tcolorbox}
(f) 它具有效用缩放不变性(invariant under utility scaling)，即如果 \(a,c > 0\) 且 \(b,d \in  \mathbb{R}\) ，并且 \({S}^{\prime }\) 是具有威胁点 \(\left( {a{u}_{0} + b,c{v}_{0} + d}\right)\) 的讨价还价集 \(\{ \left( {{au} + b,{cv} + d}\right)  \mid  \left( {u,v}\right)  \in  S\}\) ，那么 \(N\left( {S}^{\prime }\right)  = \left( {{aU} + b,{cV} + d}\right)\) ；
\end{tcolorbox}

(g) it preserves symmetry, that is, if \({u}_{0} = {v}_{0}\) and \(\left( {u,v}\right)  \in  S\) implies \(\left( {v,u}\right)  \in  S\) , then \(U = V\)

\begin{tcolorbox}
(g) 它保持对称性，即如果 \({u}_{0} = {v}_{0}\) 且 \(\left( {u,v}\right)  \in  S\) 意味着 \(\left( {v,u}\right)  \in  S\) ，那么 \(U = V\)
\end{tcolorbox}

(h) it is independent of irrelevant alternatives: If \(S,T\) are bargaining sets with the same threat point and \(S \subset  T\) , then either \(N\left( T\right)  \notin  S\) or \(N\left( T\right)  = N\left( S\right)\) .

\begin{tcolorbox}
(h) 它与无关选项(irrelevant alternatives)无关:如果 \(S,T\) 是具有相同威胁点的讨价还价集，且 \(S \subset  T\) ，那么要么 \(N\left( T\right)  \notin  S\) ，要么 \(N\left( T\right)  = N\left( S\right)\) 。
\end{tcolorbox}

The notation \(N\left( S\right)\) is short for "Nash bargaining solution". The solution is a pair(U, V)of utilities for the two players that belongs to the bargaining set \(S\) by axiom (d).

\begin{tcolorbox}
符号 \(N\left( S\right)\) 是“纳什讨价还价解(Nash bargaining solution)”的简写。根据公理(d)，该解是属于讨价还价集 \(S\) 的两个参与者的效用二元组(U, V)。
\end{tcolorbox}

The Pareto-optimality stated in (e) means that it should not be possible to improve the solution for one player without hurting the other player. Graphically, the solution is therefore on the upper-right (or "north-east") border of the bargaining set, which is also called the Pareto-frontier of the set.

\begin{tcolorbox}
(e)中所述的帕累托最优性意味着，不可能在不损害另一个参与者的情况下改善一个参与者的解。从图上看，解因此位于讨价还价集的右上角(或“东北”)边界上，该边界也称为该集合的帕累托前沿。
\end{tcolorbox}

Invariance under utility scaling as stated in (f) means the following. The utility function of a player can be changed by changing its "scale" and "origin" without changing its meaning, just like a temperature scale (see Chapter 5). Such a change for player I means that instead of \(u\) we consider \({au} + b\) where \(a\) and \(b\) are fixed real numbers and \(a > 0\) . Similarly, for player II we can replace \(v\) by \({cv} + d\) for fixed \(c,d\) and \(c > 0\) . When drawing the bargaining set \(S\) , this means that (by adding \(b\) ) the first coordinate of any point in that set is moved by the amount \(b\) (to the right or left depending on whether \(b\) is positive or negative), and stretched by the factor \(a\) (if \(a > 1\) , and shrunk if \(a < 1\) ). In (f), the set \({S}^{\prime }\) and the new threat point are obtained from \(S\) by applying such a "change of scale" for both players.

\begin{tcolorbox}
(f)中所述的效用缩放不变性意味着以下内容。玩家的效用函数可以通过改变其“刻度”和“原点”来改变，而不改变其含义，就像温度刻度一样(见第5章)。对于玩家I，这样的改变意味着我们考虑的不是 \(u\) ，而是 \({au} + b\) ，其中 \(a\) 和 \(b\) 是固定的实数，且 \(a > 0\) 。类似地，对于玩家II，我们可以用 \({cv} + d\) 替换 \(v\) ，其中 \(c,d\) 和 \(c > 0\) 是固定的。在绘制讨价还价集 \(S\) 时，这意味着(通过加上 \(b\) )该集合中任何点的第一个坐标移动了 \(b\) 的量(根据 \(b\) 的正负向右或向左移动)，并按 \(a\) 的因子拉伸(如果 \(a > 1\) )，或者收缩(如果 \(a < 1\) )。在(f)中，集合 \({S}^{\prime }\) 和新的威胁点是通过对两个玩家应用这样的“刻度变化”从 \(S\) 得到的。
\end{tcolorbox}

Because this scaling does not change the meaning of the utility values (they still represent the same preference), the bargaining solution should not be affected either, as stated in axiom (f). In a sense, this means the bargaining solution should not depend on a direct comparison of the players' utilities. For example, if you want to split a bar of chocolate with a friend you cannot claim more because you like chocolate "ten times more than she does" (which you could represent by changing your utility from \(u\) to \({10u}\) ), nor can she claim that you should get less because you already have 5 bars of chocolate in your drawer (represented by a change from \(u\) to \(u + 5\) ).

\begin{tcolorbox}
因为这种缩放不会改变效用值的含义(它们仍然代表相同的偏好)，所以如公理(f)所述，讨价还价解也不应受到影响。从某种意义上说，这意味着讨价还价解不应依赖于对玩家效用的直接比较。例如，如果你想和一个朋友分一块巧克力，你不能因为你“比她喜欢巧克力多十倍”就要求多分(你可以通过将你的效用从 \(u\) 改变为 \({10u}\) 来表示)，她也不能因为你抽屉里已经有5块巧克力就要求你少分(通过从 \(u\) 改变为 \(u + 5\) 来表示)。
\end{tcolorbox}

Preservation of symmetry as stated in (g) means that a symmetric bargaining set \(S\) , where for all(u, v)we have \(\left( {u,v}\right)  \in  S \Leftrightarrow  \left( {v,u}\right)  \in  S\) , and with a threat point of the form \(\left( {{u}_{0},{u}_{0}}\right)\) , should also have a bargaining solution of the form \(N\left( S\right)  = \left( {U,U}\right)\) .

\begin{tcolorbox}
(g)中所述的对称性保持意味着，一个对称的讨价还价集 \(S\) (对于所有的(u, v)，我们有 \(\left( {u,v}\right)  \in  S \Leftrightarrow  \left( {v,u}\right)  \in  S\) )，且威胁点的形式为 \(\left( {{u}_{0},{u}_{0}}\right)\) ，其讨价还价解的形式也应该为 \(N\left( S\right)  = \left( {U,U}\right)\) 。
\end{tcolorbox}

The independence of irrelevant alternatives stated in (h) is the most complicated axiom. As illustrated in Figure 11.2, it means that if a bargaining set \(S\) is enlarged to a larger set \(T\) that has additional points(u, v), but without changing the threat point, then either \(T\) has a bargaining solution \(N\left( T\right)\) that is one of these new points, or else it is the same solution as the solution \(N\left( S\right)\) . It is not allowed that the solution for \(T\) moves to a different point, like \(P\) in Figure 11.2, that still belongs to \(S\) .

\begin{tcolorbox}
(h)中所述的无关选项独立性是最复杂的公理。如图11.2所示，这意味着如果一个讨价还价集 \(S\) 扩展为一个更大的集合 \(T\) ，该集合有额外的点(u, v)，但不改变威胁点，那么要么 \(T\) 的讨价还价解 \(N\left( T\right)\) 是这些新点之一，要么它与解 \(N\left( S\right)\) 相同。不允许 \(T\) 的解移动到另一个不同的点，如图11.2中的 \(P\) ，该点仍然属于 \(S\) 。
\end{tcolorbox}

\begin{center}
\includegraphics[max width=0.4\textwidth]{images/bo_d2e5cmbef24c73b03tn0_315_355_1577_597_578_0.jpg}
\end{center}
\hspace*{3em} 

Figure 11.2 Bargaining sets \(S\) (dark gray) and \(T\) (dark and light gray). Independence of irrelevant alternatives states that if \(S\) is a subset of \(T\) , then \(N\left( T\right)\) is either not in \(S\) , like here, or equal to \(N\left( S\right)\) , but cannot be another point in \(S\) , like \(P\) .

\begin{tcolorbox}
图11.2 讨价还价集 \(S\) (深灰色)和 \(T\) (深灰色和浅灰色)。无关选项独立性表明，如果 \(S\) 是 \(T\) 的子集，那么 \(N\left( T\right)\) 要么不在 \(S\) 中，就像这里一样，要么等于 \(N\left( S\right)\) ，但不能是 \(S\) 中的另一个点，如 \(P\) 。
\end{tcolorbox}

The stated axioms (a)-(h) for a bargaining set and a possible bargaining solution are reasonable requirements that one can defend as general principles.

\begin{tcolorbox}
所陈述的讨价还价集和可能的讨价还价解的公理(a)-(h)是合理的要求，可以作为一般原则来辩护。
\end{tcolorbox}

\section*{11.4 The Nash Bargaining Solution}

\begin{tcolorbox}
\section*{11.4 纳什讨价还价解}
\end{tcolorbox}

This section shows that the bargaining axioms (a)-(h) in Definitions 11.1 and 11.2 imply that a bargaining solution exists, which is, moreover, unique.

\begin{tcolorbox}
本节表明，定义11.1和11.2中的讨价还价公理(a)-(h)意味着存在一个讨价还价解，而且这个解是唯一的。
\end{tcolorbox}

Theorem 11.3. Under the Nash bargaining axioms, every bargaining set \(S\) that contains a point(u, v)with \(u > {u}_{0}\) and \(v > {v}_{0}\) has a unique solution \(N\left( S\right)  = \left( {U,V}\right)\) , which is obtained as the unique point(u, v)that maximizes the product (also called the "Nash product") \(\left( {u - {u}_{0}}\right) \left( {v - {v}_{0}}\right)\) for \(\left( {u,v}\right)  \in  S\) .

\begin{tcolorbox}
定理11.3. 在纳什讨价还价公理下，每个包含点(u, v)且 \(u > {u}_{0}\) 和 \(v > {v}_{0}\) 的讨价还价集 \(S\) 都有一个唯一的解 \(N\left( S\right)  = \left( {U,V}\right)\) ，该解是使乘积(也称为“纳什乘积”) \(\left( {u - {u}_{0}}\right) \left( {v - {v}_{0}}\right)\) 在 \(\left( {u,v}\right)  \in  S\) 时达到最大值的唯一的点(u, v)。
\end{tcolorbox}

Proof. First, consider the set \({S}^{\prime } = \left\{  {\left( {u - {u}_{0},v - {v}_{0}}\right)  \mid  \left( {u,v}\right)  \in  S}\right\}\) , which is the set \(S\) translated so that the threat point \(\left( {{u}_{0},{v}_{0}}\right)\) is the origin(0,0), as shown in Figure 11.3(a) and (b). Maximizing the Nash product then amounts to maximizing \({uv}\) for \(\left( {u,v}\right)  \in  {S}^{\prime }\) . Let(U, V)be a utility pair(u, v)where this happens, where by assumption \({UV} > 0\) ; Figure 11.3(b) shows the hyperbola \({uv} = c\) for the maximum number \(c\) so that the hyperbola still intersects \({S}^{\prime }\) . Next, we re-scale the utilities so that \(\left( {U,V}\right)  = \left( {1,1}\right)\) , by replacing \(S\) with the set \({S}^{\prime \prime } = \left\{  {\left. \left( {\frac{1}{U}u,\frac{1}{V}v}\right) \right| \;\left( {u,v}\right)  \in  {S}^{\prime }}\right\}\) , shown in Figure 11.3(c). To simplify notation, we rename \({S}^{\prime \prime }\) to \(S\) .

\begin{tcolorbox}
证明。首先，考虑集合 \({S}^{\prime } = \left\{  {\left( {u - {u}_{0},v - {v}_{0}}\right)  \mid  \left( {u,v}\right)  \in  S}\right\}\) ，它是集合 \(S\) 经过平移后得到的，使得威胁点 \(\left( {{u}_{0},{v}_{0}}\right)\) 为原点(0,0)，如图11.3(a)和(b)所示。那么，最大化纳什积就相当于在 \(\left( {u,v}\right)  \in  {S}^{\prime }\) 的条件下最大化 \({uv}\) 。设(U, V)是出现这种情况时的一个效用对(u, v)，根据假设 \({UV} > 0\) ；图11.3(b)展示了对应最大数值 \(c\) 的双曲线 \({uv} = c\) ，使得该双曲线仍与 \({S}^{\prime }\) 相交。接下来，我们对效用进行重新缩放，使得 \(\left( {U,V}\right)  = \left( {1,1}\right)\) ，方法是用集合 \({S}^{\prime \prime } = \left\{  {\left. \left( {\frac{1}{U}u,\frac{1}{V}v}\right) \right| \;\left( {u,v}\right)  \in  {S}^{\prime }}\right\}\) 替换 \(S\) ，如图11.3(c)所示。为简化符号，我们将 \({S}^{\prime \prime }\) 重命名为 \(S\) 。
\end{tcolorbox}

Next, consider the set \(T = \{ \left( {u,v}\right)  \mid  u \geq  0,v \geq  0,u + v \leq  2\}\) , see Figure 11.3(d). For the bargaining set \(T\) , the solution is \(N\left( T\right)  = \left( {1,1}\right)\) , because \(T\) is a symmetric set, and(1,1)is the only symmetric point on the Pareto-frontier of \(T\) . We claim that \(S \subseteq  T\) . Then \(N\left( T\right)  = \left( {1,1}\right)  \in  S\) , so that the independence of irrelevant alternatives implies \(N\left( S\right)  = N\left( T\right)\) .

\begin{tcolorbox}
接下来，考虑集合 \(T = \{ \left( {u,v}\right)  \mid  u \geq  0,v \geq  0,u + v \leq  2\}\) ，见图11.3(d)。对于讨价还价集 \(T\) ，解为 \(N\left( T\right)  = \left( {1,1}\right)\) ，因为 \(T\) 是一个对称的集合，并且(1,1)是 \(T\) 的帕累托前沿上唯一的对称点。我们断言 \(S \subseteq  T\) 。那么 \(N\left( T\right)  = \left( {1,1}\right)  \in  S\) ，因此无关选项的独立性意味着 \(N\left( S\right)  = N\left( T\right)\) 。
\end{tcolorbox}

We prove \(S \subseteq  T\) using the convexity of \(S\) . Namely, suppose there was a point \(\left( {\bar{u},\bar{v}}\right)\) in \(S\) which is not in \(T\) . Because \(\bar{u} \geq  0,\bar{v} \geq  0\) , this means \(\bar{u} + \bar{v} > 2\) . Then, as shown in Figure 11.3(d), even if the Nash product \(\bar{u} \cdot  \bar{v}\) is not larger than 1 (which is the Nash product \({UV}\) ), then the Nash product \({uv}\) of a suitable convex combination \(\left( {u,v}\right)  = \left( {1 - \varepsilon }\right) \left( {U,V}\right)  + \varepsilon \left( {\bar{u},\bar{v}}\right)\) for some small positive \(\varepsilon\) is larger than 1 .

\begin{tcolorbox}
我们利用 \(S\) 的凸性来证明 \(S \subseteq  T\) 。具体来说，假设 \(S\) 中存在一个点 \(\left( {\bar{u},\bar{v}}\right)\) 不在 \(T\) 中。因为 \(\bar{u} \geq  0,\bar{v} \geq  0\) ，这意味着 \(\bar{u} + \bar{v} > 2\) 。然后，如图11.3(d)所示，即使纳什积 \(\bar{u} \cdot  \bar{v}\) 不大于1(这是纳什积 \({UV}\) )，那么对于某个小的正数 \(\varepsilon\) ，一个合适的凸组合 \(\left( {u,v}\right)  = \left( {1 - \varepsilon }\right) \left( {U,V}\right)  + \varepsilon \left( {\bar{u},\bar{v}}\right)\) 的纳什积 \({uv}\) 也会大于1。
\end{tcolorbox}

This is illustrated in Figure 11.3(d): The hyperbola \(\{ \left( {u,v}\right)  \mid  {uv} = 1\}\) touches the set \(T\) at the point(1,1). The tangent to the hyperbola is the line through points (2,0)and(0,2), which is the Pareto-frontier of \(T\) . Because the point \(\left( {\bar{u},\bar{v}}\right)\) is to the right of that tangent, the line segment that connects(1,1)and \(\left( {\bar{u},\bar{v}}\right)\) intersects the hyperbola and therefore contains points(u, v)to the right of the hyperbola, where \({uv} > 1\) . This is a contradiction, because \(\left( {u,v}\right)  \in  S\) since \(S\) is convex, but the maximum Nash product of a point in \(S\) is 1 .

\begin{tcolorbox}
这在图11.3(d)中有所说明:双曲线 \(\{ \left( {u,v}\right)  \mid  {uv} = 1\}\) 在点(1,1)处与集合 \(T\) 相切。双曲线的切线是过点(2,0)和(0,2)的直线，该直线是 \(T\) 的帕累托前沿。因为点 \(\left( {\bar{u},\bar{v}}\right)\) 在那条切线的右侧，连接(1,1)和 \(\left( {\bar{u},\bar{v}}\right)\) 的线段与双曲线相交，因此包含位于双曲线右侧的点(u, v)，其中 \({uv} > 1\) 。这是一个矛盾，因为 \(\left( {u,v}\right)  \in  S\) ，由于 \(S\) 是凸集，但 \(S\) 中一点的最大纳什积为1。
\end{tcolorbox}

\begin{center}
\includegraphics[max width=0.9\textwidth]{images/bo_d2e5cmbef24c73b03tn0_317_325_206_1260_1124_0.jpg}
\end{center}
\hspace*{3em} 

Figure 11.3 Proof steps for Theorem 11.3. The dashed line in (d) shows that the point \(\left( {\bar{u},\bar{v}}\right)\) , which is outside the triangle \(T\) (dark and light gray), cannot belong to \(S\) (dark gray), which shows that \(S \subseteq  T\) .

\begin{tcolorbox}
图11.3 定理11.3的证明步骤。(d)中的虚线表明，位于三角形 \(\left( {\bar{u},\bar{v}}\right)\) (深灰色和浅灰色)外部的点 \(T\) ，不可能属于 \(S\) (深灰色)，这表明 \(S \subseteq  T\) 。
\end{tcolorbox}

To prove this formally, consider the Nash product for the convex combination (u, v)of \(\left( {U,V}\right)  = \left( {1,1}\right)\) and \(\left( {\bar{u},\bar{v}}\right)\) :

\begin{tcolorbox}
为了正式证明这一点，考虑 \(\left( {U,V}\right)  = \left( {1,1}\right)\) 和 \(\left( {\bar{u},\bar{v}}\right)\) 的凸组合(u, v)的纳什积:
\end{tcolorbox}

\[
{uv} = \left\lbrack  {\left( {1 - \varepsilon }\right) U + \varepsilon \bar{u}}\right\rbrack  \left\lbrack  {\left( {1 - \varepsilon }\right) V + \varepsilon \bar{v}}\right\rbrack
\]

\[
= \left\lbrack  {U + \varepsilon \left( {\bar{u} - U}\right) }\right\rbrack  \left\lbrack  {V + \varepsilon \left( {\bar{v} - V}\right) }\right\rbrack
\]

\[
= {UV} + \varepsilon \left\lbrack  {\bar{u} + \bar{v} - \left( {U + V}\right)  + \varepsilon \left( {\bar{u} - U}\right) \left( {\bar{v} - V}\right) }\right\rbrack  . \tag{11.1}
\]

Now, because \(\bar{u} + \bar{v} > 2\) , but \(U + V = 2\) , the term \(\bar{u} + \bar{v} - \left( {U + V}\right)\) is positive, and \(\varepsilon \left( {\bar{u} - U}\right) \left( {\bar{v} - V}\right)\) can be made smaller than that term in absolute value, for suitably small \(\varepsilon  > 0\) , even if \(\left( {\bar{u} - U}\right) \left( {\bar{v} - V}\right)\) is any negative number. Consequently, \({uv} > 1 = {UV}\) for sufficiently small positive \(\varepsilon\) , which shows that \({UV}\) was not the maximum Nash product in \(S\) , which contradicts the construction of(U, V). This proves that, indeed, \(S \subseteq  T\) , and hence \(N\left( S\right)  = N\left( T\right)\) as claimed.

\begin{tcolorbox}
现在，因为 \(\bar{u} + \bar{v} > 2\) ，但 \(U + V = 2\) ，项 \(\bar{u} + \bar{v} - \left( {U + V}\right)\) 为正，并且对于足够小的 \(\varepsilon  > 0\) ，即使 \(\left( {\bar{u} - U}\right) \left( {\bar{v} - V}\right)\) 是任意负数， \(\varepsilon \left( {\bar{u} - U}\right) \left( {\bar{v} - V}\right)\) 的绝对值也可以小于该项。因此，对于足够小的正 \(\varepsilon\) ， \({uv} > 1 = {UV}\) ，这表明 \({UV}\) 不是 \(S\) 中的最大纳什积，这与(U, V)的构造相矛盾。这证明了，实际上， \(S \subseteq  T\) ，因此如所声称的 \(N\left( S\right)  = N\left( T\right)\) 。
\end{tcolorbox}

In order to show that the point \(\left( {U,V}\right)  = \left( {1,1}\right)\) is unique, we can again look at the hyperbola, and argue similarly as before. Namely, suppose that \(\left( {u,\frac{1}{u}}\right)\) is another point in \(S\) on the hyperbola with \(u \neq  1\) . Because \(S\) is convex, the line segment that connects(1,1)and \(\left( {u,\frac{1}{u}}\right)\) is also in \(S\) , but its interior points have a larger Nash product. For example, the mid-point \(\frac{1}{2} \cdot  \left( {1,1}\right)  + \frac{1}{2} \cdot  \left( {u,\frac{1}{u}}\right)\) has the Nash product \(\left( {1 + u + \frac{1}{u} + u\frac{1}{u}}\right) /4\) , which is larger than 1 if and only \(u + \frac{1}{u} > 2\) , which holds because this is equivalent to \({\left( u - 1\right) }^{2} > 0\) , which holds because \(u \neq  1\) . Hence, again because \(S\) is convex, the Nash product has a unique maximum in \(S\) .

\begin{tcolorbox}
为了证明点 \(\left( {U,V}\right)  = \left( {1,1}\right)\) 是唯一的，我们可以再次观察双曲线，并进行与之前类似的论证。即，假设 \(\left( {u,\frac{1}{u}}\right)\) 是双曲线 \(S\) 上的另一点，且 \(u \neq  1\) 。因为 \(S\) 是凸集，连接(1,1)和 \(\left( {u,\frac{1}{u}}\right)\) 的线段也在 \(S\) 中，但其中间点的纳什积更大。例如，中点 \(\frac{1}{2} \cdot  \left( {1,1}\right)  + \frac{1}{2} \cdot  \left( {u,\frac{1}{u}}\right)\) 的纳什积为 \(\left( {1 + u + \frac{1}{u} + u\frac{1}{u}}\right) /4\) ，当且仅当 \(u + \frac{1}{u} > 2\) 时，该值大于1，这成立是因为这等价于 \({\left( u - 1\right) }^{2} > 0\) ，而这成立是因为 \(u \neq  1\) 。因此，再次因为 \(S\) 是凸集， \(S\) 中的纳什积有唯一的最大值。
\end{tcolorbox}

The assumption that \(S\) has a point(u, v)with \(u > {u}_{0}\) and \(v > {v}_{0}\) has the following purpose: Otherwise, for \(\left( {u,v}\right)  \in  S\) , all utilities \(u\) of player I or \(v\) of player II (or both) fulfill \(u = {u}_{0}\) or \(v = {v}_{0}\) , respectively. Then the Nash product \(\left( {u - {u}_{0}}\right) \left( {v - {v}_{0}}\right)\) for \(\left( {u,v}\right)  \in  S\) would always be zero, so it has no unique maximum unless \(S\) is just the singleton \(\left\{  \left( {{u}_{0},{v}_{0}}\right) \right\}\) . Then the bargaining set \(S\) is a singleton or a line segment \(\left\{  {u}_{0}\right\}   \times  \left\lbrack  {{v}_{0},V}\right\rbrack\) or \(\left\lbrack  {{u}_{0},U}\right\rbrack   \times  \left\{  {v}_{0}\right\}\) . In this rather trivial case, the Pareto-frontier consists of a single point(U, V), which defines the unique Nash solution \(N\left( S\right)\) .

\begin{tcolorbox}
假设 \(S\) 有一个点 (u, v)，且满足 \(u > {u}_{0}\) 和 \(v > {v}_{0}\) ，其目的如下:否则，对于 \(\left( {u,v}\right)  \in  S\) ，玩家 I 的所有效用 \(u\) 或玩家 II 的所有效用 \(v\) (或两者)分别满足 \(u = {u}_{0}\) 或 \(v = {v}_{0}\) 。那么对于 \(\left( {u,v}\right)  \in  S\) ，纳什积 \(\left( {u - {u}_{0}}\right) \left( {v - {v}_{0}}\right)\) 总是为零，所以除非 \(S\) 只是单元素集 \(\left\{  \left( {{u}_{0},{v}_{0}}\right) \right\}\) ，否则它没有唯一的最大值。那么议价集 \(S\) 是一个单元素集或一条线段 \(\left\{  {u}_{0}\right\}   \times  \left\lbrack  {{v}_{0},V}\right\rbrack\) 或 \(\left\lbrack  {{u}_{0},U}\right\rbrack   \times  \left\{  {v}_{0}\right\}\) 。在这种相当简单的情况下，帕累托前沿由一个单点 (U, V) 组成，该点定义了唯一的纳什解 \(N\left( S\right)\) 。
\end{tcolorbox}

\(\Rightarrow\) If the bargaining set is a line segment as discussed in the preceding paragraph, why is its Pareto-frontier a single point? Is this always so if the line segment is not parallel to one of the coordinate axes?

\begin{tcolorbox}
\(\Rightarrow\) 如果议价集是上一段中讨论的线段，为什么它的帕累托前沿是一个单点？如果线段不平行于任何一个坐标轴，情况总是如此吗？
\end{tcolorbox}

It can be shown relatively easily that the maximization of the Nash product as described in Theorem 11.3 gives a solution(U, V)that fulfills the axioms in Definitions 11.1 and 11.2.

\begin{tcolorbox}
可以相对容易地证明，定理 11.3 中描述的纳什积最大化给出了一个满足定义 11.1 和 11.2 中公理的解 (U, V)。
\end{tcolorbox}

\(\Rightarrow\) In Exercise 11.1, you are asked to verify this for axiom (f).

\begin{tcolorbox}
\(\Rightarrow\) 在练习 11.1 中，要求你验证公理 (f) 的情况。
\end{tcolorbox}

We describe how to find the bargaining solution(U, V)for the bargaining set \(S\) in Figure 11.1. The threat point \(\left( {{u}_{0},{v}_{0}}\right)\) is(1,1). The Pareto-frontier consists of two line segments: a first line segment that joins(1,3.5)to(2,3), and a second line segment that joins(2,3)to(4,1), where the latter is part of a longer line segment that extends from(2,3)to(5,0)(the first line segment is also part of the longer segment that connects(4,0)to(2,3), but this will not be needed). The Nash product \(\left( {u - {u}_{0}}\right) \left( {v - {v}_{0}}\right)\) for a point(u, v)on the Pareto-frontier is the size of the area of the rectangle between the threat point \(\left( {{u}_{0},{v}_{0}}\right)\) and the point(u, v). In this case, the diagram shows that the second line segment creates the larger rectangle in that way. A point on the second line segment is of the form \(\left( {1 - p}\right) \left( {2,3}\right)  + p\left( {5,0}\right)\) for \(0 \leq  p \leq  1\) , that is, \(\left( {2 + {3p},3 - {3p}}\right)\) . Because \(\left( {{u}_{0},{v}_{0}}\right)  = \left( {1,1}\right)\) , the resulting Nash product is \(\left( {1 + {3p}}\right) \left( {2 - {3p}}\right)\) or \(2 + {3p} - 9{p}^{2}\) . The unique maximum over \(p\) of this expression is obtained for \(p = \frac{1}{6}\) because this is the zero of its derivative \(3 - {18p}\) , with negative second derivative, so this is a maximum. The point(U, V) is therefore given as(2.5,2.5). (In general, such a maximum can result for \(p < 0\) or \(p > 1\) , which implies that the left or right endpoint of the line segment with \(p = 0\) or \(p = 1\) , respectively, gives the largest value of the Nash product. Then one has to check the adjoining line segment for a possibly larger Nash product, which may occur at the corner point.) This value of \(p\) allows representing the Nash bargaining solution as a convex combination of payoff pairs that correspond to cells in the bimatrix game, namely via a lottery that chooses(T, l)with payoff pair(2,3)with probability \(1 - p = \frac{5}{6}\) and(B, r)with payoff pair(5,0)with probability \(p = \frac{1}{6}\) . Both players have agreed in advance to accept the outcome of this lottery (even though with probability \(\frac{1}{6}\) the pair(5,0)gives payoff zero to player II, which is less than at the threat point where she gets \({v}_{0} = 1\) ; however, her expected payoff 2.5 for the lottery is higher than \({v}_{0}\) ).

\begin{tcolorbox}
我们描述如何找到图11.1中讨价还价集 \(S\) 的讨价还价解(U, V)。威胁点 \(\left( {{u}_{0},{v}_{0}}\right)\) 为(1,1)。帕累托前沿由两条线段组成:第一条线段连接(1,3.5)和(2,3)，第二条线段连接(2,3)和(4,1)，其中后者是从(2,3)延伸到(5,0)的更长线段的一部分(第一条线段也是连接(4,0)和(2,3)的更长线段的一部分，但这里不需要用到)。帕累托前沿上一点(u, v)的纳什积 \(\left( {u - {u}_{0}}\right) \left( {v - {v}_{0}}\right)\) 是威胁点 \(\left( {{u}_{0},{v}_{0}}\right)\) 和点(u, v)之间矩形的面积大小。在这种情况下，图中显示第二条线段以这种方式形成的矩形更大。第二条线段上的点形式为 \(\left( {1 - p}\right) \left( {2,3}\right)  + p\left( {5,0}\right)\) ，其中 \(0 \leq  p \leq  1\) ，即 \(\left( {2 + {3p},3 - {3p}}\right)\) 。因为 \(\left( {{u}_{0},{v}_{0}}\right)  = \left( {1,1}\right)\) ，所以得到的纳什积为 \(\left( {1 + {3p}}\right) \left( {2 - {3p}}\right)\) 或 \(2 + {3p} - 9{p}^{2}\) 。该表达式在 \(p\) 上的唯一最大值在 \(p = \frac{1}{6}\) 时取得，因为这是其导数 \(3 - {18p}\) 的零点，且二阶导数为负，所以这是一个最大值。因此，点(U, V)为(2.5,2.5)。(一般来说，这样的最大值可能在 \(p < 0\) 或 \(p > 1\) 时出现，这意味着线段左端点或右端点分别对应 \(p = 0\) 或 \(p = 1\) 时，能得到纳什积的最大值。然后必须检查相邻线段是否可能有更大的纳什积，这可能出现在角点处。) \(p\) 的这个值允许将纳什讨价还价解表示为双矩阵博弈中对应单元格的收益对的凸组合，即通过一个彩票，以概率 \(1 - p = \frac{5}{6}\) 选择收益对为(2,3)的(T, l)，以概率 \(p = \frac{1}{6}\) 选择收益对为(5,0)的(B, r)。双方玩家事先已同意接受这个彩票的结果(尽管以概率 \(\frac{1}{6}\) 出现的收益对(5,0)会使玩家II的收益为零，这比她在威胁点获得的 \({v}_{0} = 1\) 要少；然而，她从彩票中获得的期望效用2.5高于 \({v}_{0}\) )。
\end{tcolorbox}

\section*{11.5 Geometry of the Bargaining Solution}

\begin{tcolorbox}
\section*{11.5 讨价还价解的几何性质}
\end{tcolorbox}

In this section, we give a geometric characterization of the bargaining solution. This characterization helps finding the bargaining solution quickly, and is needed later in Section 11.11. We first state an easy lemma.

\begin{tcolorbox}
在本节中，我们给出讨价还价解的几何特征。这一特征有助于快速找到讨价还价解，并且在11.11节中会用到。我们首先陈述一个简单的引理。
\end{tcolorbox}

Lemma 11.4. Consider a triangle \(T\) with vertices \(\left( {0,0}\right) ,\left( {a,0}\right) ,\left( {0,b}\right)\) with \(a,b > 0\) . Then the rectangle of maximal area contained in \(T\) with vertices \(\left( {0,0}\right) ,\left( {U,0}\right) ,\left( {U,V}\right)\) and(0, V) is given by \(U = a/2\) and \(V = b/2\) . Furthermore, if \(S \subseteq  T\) and \(\left( {U,V}\right)  \in  S\) , then(U, V)is also the point in \(S\) that maximizes the product \({UV}\) .

\begin{tcolorbox}
引理11.4。考虑一个顶点为 \(\left( {0,0}\right) ,\left( {a,0}\right) ,\left( {0,b}\right)\) 且满足 \(a,b > 0\) 的三角形 \(T\) 。那么包含在 \(T\) 中，顶点为 \(\left( {0,0}\right) ,\left( {U,0}\right) ,\left( {U,V}\right)\) 和(0, V)的最大面积矩形由 \(U = a/2\) 和 \(V = b/2\) 确定。此外，如果 \(S \subseteq  T\) 且 \(\left( {U,V}\right)  \in  S\) ，那么(U, V)也是 \(S\) 中使乘积 \({UV}\) 最大的点。
\end{tcolorbox}

\begin{center}
\includegraphics[max width=0.3\textwidth]{images/bo_d2e5cmbef24c73b03tn0_319_397_1276_437_551_0.jpg}
\end{center}
\hspace*{3em} 

Figure 11.4 Illustration of Lemma 11.4. The product \({UV}\) , which is the area of the gray rectangle inscribed in the triangle \(T\) , is maximal for \(U = a/2\) and \(V = b/2\) .

\begin{tcolorbox}
图11.4 引理11.4的图示。乘积 \({UV}\) 是内接于三角形 \(T\) 的灰色矩形的面积，当 \(U = a/2\) 和 \(V = b/2\) 时该面积最大。
\end{tcolorbox}

Proof. The product \({UV}\) is the area of the rectangle inscribed in \(T\) as shown in Figure 11.4. The maximum value of \({uv}\) for \(\left( {u,v}\right)  \in  T\) clearly results when(u, v)is on the line that joins(b,0)and(a,0), shown as a dashed line in Figure 11.4, that is, when \(v = b - \frac{b}{a}u\) . The derivative of the function \(u \mapsto  u \cdot  \left( {b - \frac{b}{a}u}\right)\) is \(b - 2\frac{b}{a}u\) , which is zero when \(u = a/2\) ; this is a maximum of the function because its second derivative is negative. Hence, \({UV}\) is maximal for \(U = a/2\) and \(V = b/2\) as claimed.

\begin{tcolorbox}
证明。乘积 \({UV}\) 是如图11.4所示内接于 \(T\) 的矩形的面积。对于 \(\left( {u,v}\right)  \in  T\) ， \({uv}\) 的最大值显然在(u, v)位于连接(b,0)和(a,0)的直线上时取得，该直线在图11.4中用虚线表示，即当 \(v = b - \frac{b}{a}u\) 时。函数 \(u \mapsto  u \cdot  \left( {b - \frac{b}{a}u}\right)\) 的导数为 \(b - 2\frac{b}{a}u\) ，当 \(u = a/2\) 时其值为零；这是该函数的最大值，因为其二阶导数为负。因此，如所声称的，当 \(U = a/2\) 和 \(V = b/2\) 时 \({UV}\) 最大。
\end{tcolorbox}

If \(S \subseteq  T\) , then the maximum of \({uv}\) for \(\left( {u,v}\right)  \in  S\) is clearly at most the maximum of \({uv}\) for \(\left( {u,v}\right)  \in  T\) , and therefore attained for \(\left( {u,v}\right)  = \left( {U,V}\right)\) if \(\left( {U,V}\right)  \in  S\) .

\begin{tcolorbox}
如果 \(S \subseteq  T\) ，那么对于 \(\left( {u,v}\right)  \in  S\) ， \({uv}\) 的最大值显然至多是对于 \(\left( {u,v}\right)  \in  T\) 时 \({uv}\) 的最大值，因此如果 \(\left( {U,V}\right)  \in  S\) ，则在 \(\left( {u,v}\right)  = \left( {U,V}\right)\) 时取得。
\end{tcolorbox}

The following proposition gives a useful condition for finding the bargaining solution geometrically. The bargaining set is assumed to be convex as per axiom (c), to have threat point(0,0), and to have a Pareto-frontier of the form \(\left( {u,f\left( u\right) }\right)\) for a decreasing function \(f\) where, for definiteness, \(0 \leq  u \leq  1\) ; any other right endpoint of the interval for \(u\) is also possible.

\begin{tcolorbox}
以下命题给出了一个从几何角度寻找讨价还价解(bargaining solution)的有用条件。根据公理(c)，假设讨价还价集是凸集，威胁点(threat point)为(0,0)，并且其帕累托前沿(Pareto - frontier)的形式为 \(\left( {u,f\left( u\right) }\right)\) ，其中 \(f\) 是一个递减函数，为明确起见， \(0 \leq  u \leq  1\) ； \(u\) 区间的任何其他右端点也是可能的。
\end{tcolorbox}

Proposition 11.5. Suppose that the Pareto-frontier of a bargaining set \(S\) with threat point (0,0)is given by \(\{ \left( {u,f\left( u\right) }\right)  \mid  0 \leq  u \leq  1\}\) for a decreasing and continuous function \(f\) with \(f\left( 0\right)  > 0\) and \(f\left( 1\right)  = 0\) . Then the bargaining solution(U, V)is the unique point \(\left( {U,f\left( U\right) }\right)\) where the bargaining set has a tangent line with slope \(- f\left( U\right) /U\) . If \(f\) is differentiable, then this slope is the derivative \({f}^{\prime }\left( U\right)\) of \(f\) at \(U\) , that is, \({f}^{\prime }\left( U\right)  =  - f\left( U\right) /U\) .

\begin{tcolorbox}
命题11.5。假设威胁点为(0,0)的讨价还价集 \(S\) 的帕累托前沿由 \(\{ \left( {u,f\left( u\right) }\right)  \mid  0 \leq  u \leq  1\}\) 给出，其中 \(f\) 是一个递减且连续的函数，满足 \(f\left( 0\right)  > 0\) 和 \(f\left( 1\right)  = 0\) 。那么讨价还价解(bargaining solution)(U, V)是讨价还价集有一条斜率为 \(- f\left( U\right) /U\) 的切线的唯一一点 \(\left( {U,f\left( U\right) }\right)\) 。如果 \(f\) 可微，那么该斜率是 \(f\) 在 \(U\) 处的导数 \({f}^{\prime }\left( U\right)\) ，即 \({f}^{\prime }\left( U\right)  =  - f\left( U\right) /U\) 。
\end{tcolorbox}

\begin{center}
\includegraphics[max width=0.4\textwidth]{images/bo_d2e5cmbef24c73b03tn0_320_194_865_565_741_0.jpg}
\end{center}
\hspace*{3em} 

Figure 11.5 Illustration of Proposition 11.5 and its proof. The dashed line is the tangent to the hyperbola \(\{ \left( {u,v}\right)  \mid  {uv} = c = {UV}\}\) at the point (U, V), which is also a tangent to the bargaining set \(S\) at(U, V). The endpoints(a,0)and(0, b)of this tangent fulfill \(a = {2U}\) and \(b = {2V}\) with angles \(\alpha\) and \(\beta\) as shown.

\begin{tcolorbox}
图11.5 命题11.5及其证明的图示。虚线是双曲线 \(\{ \left( {u,v}\right)  \mid  {uv} = c = {UV}\}\) 在点 (U, V) 处的切线，该切线也是讨价还价集 \(S\) 在点 (U, V) 处的切线。此切线的端点 (a,0) 和 (0, b) 满足 \(a = {2U}\) 和 \(b = {2V}\) ，夹角分别为 \(\alpha\) 和 \(\beta\) ，如图所示。
\end{tcolorbox}

Proof. Let(U, V)be the bargaining solution for the bargaining set \(S\) . It is on the Pareto-frontier, so \(V = f\left( U\right)\) . Let \(c = {UV}\) and consider the hyperbola \(\{ \left( {u,v}\right)  \mid  {uv} = c\}\) , which intersects \(S\) in the unique point(U, V)by Theorem 11.3. Because(U, V)maximizes the Nash product, \(c\) is the maximal value of \(c = {uv}\) for \(\left( {u,v}\right)  \in  S\) . The function \(u \mapsto  c/u\) is differentiable with derivative \(- c/{u}^{2}\) , which for \(u = U\) is equal to \(- V/U\) . Consider the line through(U, V)with slope \(- V/U\) , which is the tangent to the hyperbola at(U, V), shown as a dashed line in Figure 11.5. This line intersects the \(u\) -axis at the point \(\left( {a,0}\right)  = \left( {{2U},0}\right)\) and the \(v\) -axis at \(\left( {0,b}\right)  = \left( {0,{2V}}\right)\) , and defines a triangle \(T\) with vertices \(\left( {0,0}\right) ,\left( {a,0}\right)\) , and (0, b). Figure 11.5 corresponds to Figure 11.3(b), and if we re-scale the axes by replacing \(u\) with \(u/U\) and \(v\) with \(v/V\) then we obtain Figure 11.3(d) where \(T\) is the triangle with vertices \(\left( {0,0}\right) ,\left( {2,0}\right)\) , and(0,2). With the help of(11.1)we showed \(S \subseteq  T\) , where the assumption \(\left( {U,V}\right)  = \left( {1,1}\right)\) was not used, so this applies also to Figure 11.5, without re-scaling the axes. Hence, the dashed line is also a tangent to the set \(S\) , as claimed.

\begin{tcolorbox}
证明。设 (U, V) 是讨价还价集 \(S\) 的讨价还价解。它位于帕累托前沿上，因此 \(V = f\left( U\right)\) 。设 \(c = {UV}\) 并考虑双曲线 \(\{ \left( {u,v}\right)  \mid  {uv} = c\}\) ，根据定理11.3，该双曲线与 \(S\) 仅在点 (U, V) 处相交。因为 (U, V) 使纳什积最大化，所以 \(c\) 是 \(c = {uv}\) 在 \(\left( {u,v}\right)  \in  S\) 条件下的最大值。函数 \(u \mapsto  c/u\) 可微，其导数为 \(- c/{u}^{2}\) ，当 \(u = U\) 时，该导数等于 \(- V/U\) 。考虑过点 (U, V) 且斜率为 \(- V/U\) 的直线，此直线是双曲线在点 (U, V) 处的切线，如图11.5中的虚线所示。该直线与 \(u\) 轴相交于点 \(\left( {a,0}\right)  = \left( {{2U},0}\right)\) ，与 \(v\) 轴相交于点 \(\left( {0,b}\right)  = \left( {0,{2V}}\right)\) ，并定义了一个顶点为 \(\left( {0,0}\right) ,\left( {a,0}\right)\) 和 (0, b) 的三角形 \(T\) 。图11.5对应于图11.3(b)，如果我们通过将 \(u\) 替换为 \(u/U\) 以及将 \(v\) 替换为 \(v/V\) 来重新缩放坐标轴，那么我们将得到图11.3(d)，其中 \(T\) 是顶点为 \(\left( {0,0}\right) ,\left( {2,0}\right)\) 和 (0,2) 的三角形。借助 (11.1) 我们证明了 \(S \subseteq  T\) ，其中未使用假设 \(\left( {U,V}\right)  = \left( {1,1}\right)\) ，因此这也适用于图11.5，无需重新缩放坐标轴。因此，虚线也是集合 \(S\) 的切线，正如所声称的那样。
\end{tcolorbox}

If the function \(f\) is not differentiable, then there may be more than one tangent to the set \(S\) at a point. In that case it remains to prove that the described construction of(U, V)is unique. Suppose that a tangent of \(S\) touches \(S\) at the point(U, V)and has the downward slope \(- V/U\) , and intersects the \(u\) -axis at(a,0)and the \(v\) -axis at (0, b)as in Figure 11.5 (where we now ignore the hyperbola). The slope \(- V/U\) of the tangent implies that \(a = {2U}\) and \(b = {2V}\) . Being a tangent to the set \(S\) means that \(S\) is a subset of the triangle \(T\) with vertices \(\left( {0,0}\right) ,\left( {a,0}\right)\) , and(0, b). By Lemma 11.4, \({UV}\) is the maximum value of \({uv}\) for \(\left( {u,v}\right)  \in  S\) . This shows that(U, V)is unique.

\begin{tcolorbox}
如果函数 \(f\) 不可微，那么在某一点处集合 \(S\) 可能存在多条切线。在这种情况下，仍需证明所描述的 (U, V) 的构造是唯一的。假设 \(S\) 的一条切线在点 (U, V) 处与 \(S\) 相切，且斜率为向下的 \(- V/U\) ，并与 \(u\) 轴相交于 (a,0)，与 \(v\) 轴相交于 (0, b)，如图11.5所示(现在我们忽略双曲线)。切线的斜率 \(- V/U\) 意味着 \(a = {2U}\) 和 \(b = {2V}\) 。作为集合 \(S\) 的切线意味着 \(S\) 是顶点为 \(\left( {0,0}\right) ,\left( {a,0}\right)\) 和 (0, b) 的三角形 \(T\) 的子集。根据引理11.4， \({UV}\) 是 \({uv}\) 在 \(\left( {u,v}\right)  \in  S\) 条件下的最大值。这表明 (U, V) 是唯一的。
\end{tcolorbox}

If the function \(f\) is differentiable, then the tangent to the set \(S\) is unique at every point, and we can argue directly as follows. The bargaining solution maximizes the Nash product \(u \cdot  f\left( u\right)\) on the Pareto-frontier. This requires the derivative with respect to \(u\) to be zero, that is, \(f\left( u\right)  + u \cdot  {f}^{\prime }\left( u\right)  = 0\) or \({f}^{\prime }\left( u\right)  =  - f\left( u\right) /u\) , so \(U\) has to solve this equation, as claimed.

\begin{tcolorbox}
如果函数 \(f\) 可微，那么集合 \(S\) 在每一点处的切线都是唯一的，我们可以直接进行如下论证。讨价还价解(bargaining solution)在帕累托前沿上使纳什积 \(u \cdot  f\left( u\right)\) 最大化。这要求关于 \(u\) 的导数为零，即 \(f\left( u\right)  + u \cdot  {f}^{\prime }\left( u\right)  = 0\) 或 \({f}^{\prime }\left( u\right)  =  - f\left( u\right) /u\) ，因此 \(U\) 必须满足这个方程，正如所声称的那样。
\end{tcolorbox}

In Figure 11.5, the point(U, V)is the midpoint of the long side of the right-angled triangle with vertices \(\left( {0,0}\right) ,\left( {a,0}\right)\) , and(0, b). The line from(0,0)to(U, V) divides that triangle into two isosceles triangles. The first triangle has apex (U, V)and equal (mirrored) angles \(\alpha\) at its vertices(0,0)and(a,0). These angles correspond to the slopes \(V/U\) and \(- V/U\) , equal except for their sign, of the two lines that connect these vertices to(U, V). The second triangle has apex(U, V)and equal (mirrored) angles \(\beta\) at its vertices(0,0)and(0, b).

\begin{tcolorbox}
在图 11.5 中，点 (U, V) 是直角三角形长边的中点，该直角三角形的顶点为 \(\left( {0,0}\right) ,\left( {a,0}\right)\) 和 (0, b)。从 (0, 0) 到 (U, V) 的直线将该三角形分成两个等腰三角形。第一个三角形的顶点为 (U, V)，在其顶点 (0, 0) 和 (a, 0) 处有相等(镜像)的角 \(\alpha\) 。这些角对应于连接这些顶点与 (U, V) 的两条直线的斜率 \(V/U\) 和 \(- V/U\) ，除了符号外，这两个斜率相等。第二个三角形的顶点为 (U, V)，在其顶点 (0, 0) 和 (0, b) 处有相等(镜像)的角 \(\beta\) 。
\end{tcolorbox}

Figure 11.6 shows bargaining sets \(S\) with a Pareto-frontier given by line segments, as it arises when \(S\) is derived from a bimatrix game as in Figure 11.1 (which with the threat point shifted to(0,0)defines the set shown in Figure 11.6(b)). Then the Nash bargaining solution is found easily with the help of Lemma 11.4: Find a tangent to the bargaining set \(S\) with endpoints(a,0)and(0, b)so that the midpoint(U, V)of that tangent belongs to \(S\) (Proposition 11.5 asserts that this tangent exists). Because the boundary of \(S\) consists of line segments, the tangent at an interior point of the line segment has the same slope as that line segment. Only finitely many slopes need to be checked. In Figure 11.6(a), the middle line segment contains the midpoint(U, V)of the corresponding tangent.

\begin{tcolorbox}
图 11.6 展示了讨价还价集 \(S\) ，其帕累托前沿由线段(line segments)给出，当 \(S\) 从如图 11.1 所示的双矩阵博弈(bimatrix game)中推导得出时就会出现这种情况(将威胁点(threat point)移至 (0, 0) 后定义了图 11.6(b) 中所示的集合)。然后，借助引理 11.4 可以轻松找到纳什讨价还价解:找到与端点为 (a, 0) 和 (0, b) 的讨价还价集 \(S\) 相切的切线，使得该切线的中点 (U, V) 属于 \(S\) (命题 11.5 断言这条切线存在)。因为 \(S\) 的边界由线段组成，所以线段内部点处的切线与该线段具有相同的斜率。只需检查有限多个斜率。在图 11.6(a) 中，中间的线段包含相应切线的中点 (U, V)。
\end{tcolorbox}

The case of a Pareto-frontier with only two line segments is particularly easy. In Figure 11.6(b), the right line segment contains the midpoint(U, V)of the corresponding tangent. That tangent and the line segment have the same right endpoint(a,0), which implies \(U = \frac{a}{2}\) . Similarly, in Figure 11.6(c), the top line segment contains the midpoint(U, V)of the corresponding tangent. That tangent and the line segment have the same left endpoint(0, b), which implies \(V = \frac{b}{2}\) . In Figure 11.6(d), neither line segment defines a tangent whose midpoint belongs to \(S\) , because the right line segment is too steep and the top line segment too shallow. Then the vertex(U, V)of the bargaining set \(S\) is the Nash bargaining solution; this can be verified by considering the dashed line through that vertex with slope \(- V/U\) , which is the negative of the slope \(V/U\) of the line from(0,0)to(U, V), with angle \(\alpha\) with the \(u\) -axis. That line should be a tangent to \(S\) , as it is here.

\begin{tcolorbox}
帕累托前沿仅由两条线段组成的情况特别简单。在图 11.6(b) 中，右侧的线段包含相应切线的中点 (U, V)。该切线和线段具有相同的右端点 (a, 0)，这意味着 \(U = \frac{a}{2}\) 。类似地，在图 11.6(c) 中，顶部的线段包含相应切线的中点 (U, V)。该切线和线段具有相同的左端点 (0, b)，这意味着 \(V = \frac{b}{2}\) 。在图 11.6(d) 中，没有一条线段能定义出其中点属于 \(S\) 的切线，因为右侧的线段太陡，而顶部的线段太平缓。那么讨价还价集 \(S\) 的顶点 (U, V) 就是纳什讨价还价解；这可以通过考虑过该顶点且斜率为 \(- V/U\) 的虚线来验证，该斜率是从 (0, 0) 到 (U, V) 的直线斜率 \(V/U\) 的负值，与 \(u\) 轴的夹角为 \(\alpha\) 。这条线应该是 \(S\) 的切线，就像这里的情况一样。
\end{tcolorbox}

\begin{center}
\includegraphics[max width=0.9\textwidth]{images/bo_d2e5cmbef24c73b03tn0_322_196_206_1256_1018_0.jpg}
\end{center}
\hspace*{3em} 

Figure 11.6 Various bargaining sets \(S\) with a piecewiese linear Pareto-frontier. By Proposition 11.5, the Nash bargaining solution(U, V)is the midpoint of the tangent to \(S\) (dashed line) intersected with the positive quadrant, with endpoints(a,0) and(0, b). In (d) the point(U, V)is a vertex of \(S\) , where there are many tangents to \(S\) through that point. One of these tangents has the same angle \(\alpha\) with the \(u\) -axis as the upwards line from(0,0)to(U, V).

\begin{tcolorbox}
图 11.6 展示了具有分段线性帕累托前沿的各种讨价还价集 \(S\) 。根据命题 11.5，纳什讨价还价解 (U, V) 是与 \(S\) 相切的切线(虚线)与正象限相交部分的中点，切线的端点为 (a, 0) 和 (0, b)。在 (d) 中，点 (U, V) 是 \(S\) 的一个顶点，过该点有许多与 \(S\) 相切的切线。其中一条切线与 \(u\) 轴的夹角 \(\alpha\) 与从 (0, 0) 向上到 (U, V) 的直线的夹角相同。
\end{tcolorbox}

\section*{11.6 Splitting a Unit Pie}

\begin{tcolorbox}
\section*{11.6 分割一个单位馅饼}
\end{tcolorbox}

In the remainder of this chapter, we consider a bargaining situation where player I and player II have to agree to split a "pie" into an amount \(x\) for player I and \(y\) for player II. The total amount to be split is normalized to be one unit, so this is called

\begin{tcolorbox}
在本章的其余部分，我们考虑一个讨价还价的情形，其中参与者I和参与者II必须就如何将一个“馅饼”分割达成一致，分给参与者I的数量为 \(x\) ，分给参与者II的数量为 \(y\) 。待分割的总量被规范化为一个单位，因此这被称为
\end{tcolorbox}

"a unit pie". The possible splits(x, y)have to fulfill \(x \geq  0,y \geq  0\) , and \(x + y \leq  1\) . If the two players cannot agree, they both receive zero, which defines the threat point(0,0). Splits(x, y)so that \(x + y < 1\) are not Pareto-optimal, but we admit such splits in order to obtain a convex set.

\begin{tcolorbox}
“单位馅饼”。可能的分割方案 (x, y) 必须满足 \(x \geq  0,y \geq  0\) 以及 \(x + y \leq  1\) 。如果两个参与者无法达成一致，他们都将得到零，这就定义了威胁点 (0,0) 。满足 \(x + y < 1\) 的分割方案 (x, y) 不是帕累托最优的，但为了得到一个凸集，我们允许这样的分割方案存在。
\end{tcolorbox}

\begin{center}
\includegraphics[max width=0.9\textwidth]{images/bo_d2e5cmbef24c73b03tn0_323_343_432_1241_492_0.jpg}
\end{center}
\hspace*{3em} 

Figure 11.7 Bargaining over a unit pie, with bargaining set \(S\) as defined in (11.2). In (a), the utility functions are \(u\left( x\right)  = x\) and \(v\left( y\right)  = y\) , with bargaining solution \(\left( {U,V}\right)  = \left( {\frac{1}{2},\frac{1}{2}}\right)\) . In (b), they are \(u\left( x\right)  = \sqrt{x}\) and \(v\left( y\right)  = y\) , with bargaining solution \(\left( {U,V}\right)  = \left( {1/\sqrt{3},\frac{2}{3}}\right)  \approx  \left( {{0.577},{0.667}}\right)\) .

\begin{tcolorbox}
图11.7 关于单位馅饼的讨价还价，讨价还价集 \(S\) 如 (11.2) 中所定义。在 (a) 中，效用函数为 \(u\left( x\right)  = x\) 和 \(v\left( y\right)  = y\) ，讨价还价解为 \(\left( {U,V}\right)  = \left( {\frac{1}{2},\frac{1}{2}}\right)\) 。在 (b) 中，它们是 \(u\left( x\right)  = \sqrt{x}\) 和 \(v\left( y\right)  = y\) ，讨价还价解为 \(\left( {U,V}\right)  = \left( {1/\sqrt{3},\frac{2}{3}}\right)  \approx  \left( {{0.577},{0.667}}\right)\) 。
\end{tcolorbox}

If player I gets utility \(x\) and player II gets utility \(y\) , then the bargaining set is the triangle of all pairs(x, y)with \(x \geq  0,y \geq  0\) and \(x + y \leq  1\) , shown in Figure 11.7(a). The split(x, y)is Pareto-optimal if \(x + y = 1\) , so the Pareto-frontier is given by the splits(x,1 - x), shown as the bold line in the diagram. Because the threat point is(0,0), the Nash product is \(x\left( {1 - x}\right)\) , which is maximized for \(x = \frac{1}{2}\) . The Nash bargaining solution \(\left( {U,V}\right)  = \left( {\frac{1}{2},\frac{1}{2}}\right)\) results also from the symmetry axiom in Definition 11.2(g).

\begin{tcolorbox}
如果参与者I获得效用 \(x\) ，参与者II获得效用 \(y\) ，那么讨价还价集就是所有满足 \(x \geq  0,y \geq  0\) 和 \(x + y \leq  1\) 的数对 (x, y) 所构成的三角形，如图11.7(a) 所示。如果 \(x + y = 1\) ，则分割方案 (x, y) 是帕累托最优的，因此帕累托前沿由分割方案 (x,1 - x) 给出，在图中用粗线表示。因为威胁点是 (0,0) ，所以纳什积为 \(x\left( {1 - x}\right)\) ，当 \(x = \frac{1}{2}\) 时该纳什积达到最大值。纳什讨价还价解 \(\left( {U,V}\right)  = \left( {\frac{1}{2},\frac{1}{2}}\right)\) 也可由定义11.2(g) 中的对称性公理得出。
\end{tcolorbox}

More generally, we consider the split of a unit pie where the players have utility functions \(u\) and \(v\) that are not the identity function. We assume that the utility functions are \(u : \left\lbrack  {0,1}\right\rbrack   \rightarrow  \left\lbrack  {0,1}\right\rbrack\) for player \(\mathrm{I}\) and \(v : \left\lbrack  {0,1}\right\rbrack   \rightarrow  \left\lbrack  {0,1}\right\rbrack\) for player \(\mathrm{{II}}\) , which are normalized so that \(u\left( 0\right)  = v\left( 0\right)  = 0\) and \(u\left( 1\right)  = v\left( 1\right)  = 1\) . These utility functions are assumed to be increasing, concave (see (5.46) and Figure 5.5), and continuous (it can be shown that concavity implies continuity except at endpoints of the domain of the function). The bargaining set of resulting utility pairs to the two players is given by

\begin{tcolorbox}
更一般地，我们考虑单位馅饼的分割问题，其中参与者的效用函数 \(u\) 和 \(v\) 不是恒等函数。我们假设参与者 \(\mathrm{I}\) 的效用函数为 \(u : \left\lbrack  {0,1}\right\rbrack   \rightarrow  \left\lbrack  {0,1}\right\rbrack\) ，参与者 \(\mathrm{{II}}\) 的效用函数为 \(v : \left\lbrack  {0,1}\right\rbrack   \rightarrow  \left\lbrack  {0,1}\right\rbrack\) ，并且这些效用函数经过了规范化处理，使得 \(u\left( 0\right)  = v\left( 0\right)  = 0\) 和 \(u\left( 1\right)  = v\left( 1\right)  = 1\) 。假设这些效用函数是递增的、凹的(见 (5.46) 和图5.5)且连续的(可以证明，除了函数定义域的端点外，凹性意味着连续性)。两个参与者所得到的效用对的讨价还价集由
\end{tcolorbox}

\[
S = \{ \left( {u\left( x\right) ,v\left( y\right) }\right)  \mid  x \geq  0,y \geq  0,x + y \leq  1\} . \tag{11.2}
\]

The following proposition asserts that \(S\) fulfills the axioms in Definition 11.1.

\begin{tcolorbox}
以下命题断言 \(S\) 满足定义11.1中的公理。
\end{tcolorbox}

Proposition 11.6. Let \(u,v : \left\lbrack  {0,1}\right\rbrack   \rightarrow  \mathbb{R}\) be increasing, concave, and continuous utility functions of player \(I\) and \({II}\) , respectively, with \(u\left( 0\right)  = v\left( 0\right)  = 0\) and \(u\left( 1\right)  = v\left( 1\right)  = 1\) . Then the bargaining set \(S\) defined (11.2) has threat point(0,0)and is compact and convex. Moreover, \(S\) is downward closed, that is,

\begin{tcolorbox}
命题11.6。设 \(u,v : \left\lbrack  {0,1}\right\rbrack   \rightarrow  \mathbb{R}\) 分别为参与者 \(I\) 和 \({II}\) 的递增、凹且连续的效用函数(utility function)，其中 \(u\left( 0\right)  = v\left( 0\right)  = 0\) 且 \(u\left( 1\right)  = v\left( 1\right)  = 1\) 。那么，由(11.2)定义的议价集 \(S\) 的威胁点(threat point)为(0,0)，并且是紧致且凸的。此外， \(S\) 是向下封闭的，即
\end{tcolorbox}

\[
\left( {\bar{u},\bar{v}}\right)  \in  S,\;0 \leq  {u}^{\prime } \leq  \bar{u},\;0 \leq  {v}^{\prime } \leq  \bar{v}\; \Rightarrow  \;\left( {{u}^{\prime },{v}^{\prime }}\right)  \in  S. \tag{11.3}
\]

Proof. We first show (11.3). Let \(0 \leq  {u}^{\prime } \leq  \bar{u}\) and \(0 \leq  {v}^{\prime } \leq  \bar{v}\) with \(\left( {\bar{u},\bar{v}}\right)  \in  S\) , that is, \(\bar{u} = u\left( \bar{x}\right)\) and \(\bar{v} = u\left( \bar{y}\right)\) for some \(\bar{x} \geq  0\) and \(\bar{y} \geq  0\) with \(\bar{x} + \bar{y} \leq  1\) . Because \(u\left( 0\right)  = 0 \leq  {u}^{\prime } \leq  u\left( \bar{x}\right)\) and \(u\) is continuous, we have \(u\left( {x}^{\prime }\right)  = {u}^{\prime }\) for some \({x}^{\prime } \in  \left\lbrack  {0,\bar{x}}\right\rbrack\) by the intermediate value theorem, and similarly \(v\left( {y}^{\prime }\right)  = {y}^{\prime }\) for some \({y}^{\prime } \in  \left\lbrack  {0,\bar{y}}\right\rbrack\) , which implies \({x}^{\prime } + {y}^{\prime } \leq  \bar{x} + \bar{y} \leq  1\) and thus \(\left( {{u}^{\prime },{v}^{\prime }}\right)  \in  S\) , which shows (11.3).

\begin{tcolorbox}
证明。我们首先证明(11.3)。设 \(0 \leq  {u}^{\prime } \leq  \bar{u}\) 且 \(0 \leq  {v}^{\prime } \leq  \bar{v}\) ，其中 \(\left( {\bar{u},\bar{v}}\right)  \in  S\) ，即对于某些满足 \(\bar{x} + \bar{y} \leq  1\) 的 \(\bar{x} \geq  0\) 和 \(\bar{y} \geq  0\) ，有 \(\bar{u} = u\left( \bar{x}\right)\) 且 \(\bar{v} = u\left( \bar{y}\right)\) 。因为 \(u\left( 0\right)  = 0 \leq  {u}^{\prime } \leq  u\left( \bar{x}\right)\) 且 \(u\) 是连续的，根据介值定理，对于某些 \({x}^{\prime } \in  \left\lbrack  {0,\bar{x}}\right\rbrack\) 有 \(u\left( {x}^{\prime }\right)  = {u}^{\prime }\) ，类似地，对于某些 \({y}^{\prime } \in  \left\lbrack  {0,\bar{y}}\right\rbrack\) 有 \(v\left( {y}^{\prime }\right)  = {y}^{\prime }\) ，这意味着 \({x}^{\prime } + {y}^{\prime } \leq  \bar{x} + \bar{y} \leq  1\) ，从而 \(\left( {{u}^{\prime },{v}^{\prime }}\right)  \in  S\) ，这就证明了(11.3)。
\end{tcolorbox}

The set \(S\) is the image of a compact set under the continuous function \(\left( {x,y}\right)  \mapsto\)  \(\left( {u\left( x\right) ,v\left( y\right) }\right)\) and therefore is also compact. To show that \(S\) is convex, consider a convex combination \(\left( {{u}^{\prime },{v}^{\prime }}\right)  = \left( {1 - \lambda }\right) \left( {u\left( x\right) ,v\left( y\right) }\right)  + \lambda \left( {u\left( \widetilde{x}\right) ,v\left( \widetilde{y}\right) }\right)\) of two points in \(S\) , where \(\lambda  \in  \left\lbrack  {0,1}\right\rbrack\) . Then \({u}^{\prime } \leq  u\left( {\left( {1 - \lambda }\right) x + \lambda \widetilde{x}}\right)  = \bar{u}\) and \({v}^{\prime } \leq  v\left( {\left( {1 - \lambda }\right) y + \lambda \widetilde{y}}\right)  = \bar{v}\) because \(u\) and \(v\) are concave functions, where clearly \(\left( {\bar{u},\bar{v}}\right)  \in  S\) . Moreover, \({u}^{\prime } \geq  0\) and \({v}^{\prime } \geq  0\) . Hence,(11.3) implies \(\left( {{u}^{\prime },{v}^{\prime }}\right)  \in  S\) , which shows that \(S\) is convex.

\begin{tcolorbox}
集合 \(S\) 是紧集(compact set)在连续函数 \(\left( {x,y}\right)  \mapsto\)  \(\left( {u\left( x\right) ,v\left( y\right) }\right)\) 下的像，因此它也是紧集。为了证明 \(S\) 是凸集，考虑 \(S\) 中两点的凸组合(convex combination) \(\left( {{u}^{\prime },{v}^{\prime }}\right)  = \left( {1 - \lambda }\right) \left( {u\left( x\right) ,v\left( y\right) }\right)  + \lambda \left( {u\left( \widetilde{x}\right) ,v\left( \widetilde{y}\right) }\right)\) ，其中 \(\lambda  \in  \left\lbrack  {0,1}\right\rbrack\) 。那么 \({u}^{\prime } \leq  u\left( {\left( {1 - \lambda }\right) x + \lambda \widetilde{x}}\right)  = \bar{u}\) 且 \({v}^{\prime } \leq  v\left( {\left( {1 - \lambda }\right) y + \lambda \widetilde{y}}\right)  = \bar{v}\) ，因为 \(u\) 和 \(v\) 是凹函数(concave function)，显然 \(\left( {\bar{u},\bar{v}}\right)  \in  S\) 。此外， \({u}^{\prime } \geq  0\) 且 \({v}^{\prime } \geq  0\) 。因此，(11.3) 式意味着 \(\left( {{u}^{\prime },{v}^{\prime }}\right)  \in  S\) ，这表明 \(S\) 是凸集。
\end{tcolorbox}

For \(u\left( x\right)  = \sqrt{x}\) and \(v\left( y\right)  = y\) , the set \(S\) is shown in Figure 11.7(b). The Pareto-frontier is given by the points \(\left( {u\left( x\right) ,v\left( {1 - x}\right) }\right)\) for \(0 \leq  x \leq  1\) . Because \(u = \sqrt{x}\) and \(v = 1 - x = 1 - {u}^{2}\) , these are pairs \(\left( {u,1 - {u}^{2}}\right)\) , so \(v\) is easy to draw as the function \(1 - {u}^{2}\) of \(u\) . In general, one may think of traversing the points \(\left( {u\left( x\right) ,v\left( {1 - x}\right) }\right)\) on the Pareto-frontier when changing the parameter \(x\) , from the upper endpoint(0,1) of the frontier when \(x = 0\) , to its right endpoint(1,0)when \(x = 1\) .

\begin{tcolorbox}
对于 \(u\left( x\right)  = \sqrt{x}\) 和 \(v\left( y\right)  = y\) ，集合 \(S\) 如图 11.7(b) 所示。帕累托前沿(Pareto-frontier)由满足 \(0 \leq  x \leq  1\) 的点 \(\left( {u\left( x\right) ,v\left( {1 - x}\right) }\right)\) 给出。因为 \(u = \sqrt{x}\) 和 \(v = 1 - x = 1 - {u}^{2}\) ，这些是成对的 \(\left( {u,1 - {u}^{2}}\right)\) ，所以将 \(v\) 画成 \(u\) 的函数 \(1 - {u}^{2}\) 很容易。一般来说，可以想象当改变参数 \(x\) 时，遍历帕累托前沿上的点 \(\left( {u\left( x\right) ,v\left( {1 - x}\right) }\right)\) ，从 \(x = 0\) 时前沿的上端点 (0,1)，到 \(x = 1\) 时前沿的右端点 (1,0)。
\end{tcolorbox}

The Nash bargaining solution(U, V)is found by maximizing the Nash product \(u\left( x\right) v\left( {1 - x}\right)\) on the Pareto-frontier. In Figure 11.7(b), the Nash product is \(\sqrt{x}\left( {1 - x}\right)\) or \({x}^{1/2} - {x}^{3/2}\) with derivative \(\frac{1}{2}{x}^{-1/2} - \frac{3}{2}{x}^{1/2}\) or \(\frac{1}{2}{x}^{-1/2}\left( {1 - {3x}}\right)\) , which is zero when \(x = \frac{1}{3}\) , with a negative second derivative, so this is a maximum. Hence, the pie is split into \(\frac{1}{3}\) for player I and \(\frac{2}{3}\) for player II. The corresponding pair(U, V)of utilities is \(\left( {1/\sqrt{3},\frac{2}{3}}\right)\) or about(0.577,0.667). For a similar calculation with different utility functions see Exercise 11.4.

\begin{tcolorbox}
纳什讨价还价解(Nash bargaining solution)(U, V) 是通过在帕累托前沿上最大化纳什积 \(u\left( x\right) v\left( {1 - x}\right)\) 得到的。在图 11.7(b) 中，纳什积为 \(\sqrt{x}\left( {1 - x}\right)\) 或 \({x}^{1/2} - {x}^{3/2}\) ，其导数为 \(\frac{1}{2}{x}^{-1/2} - \frac{3}{2}{x}^{1/2}\) 或 \(\frac{1}{2}{x}^{-1/2}\left( {1 - {3x}}\right)\) ，当 \(x = \frac{1}{3}\) 时导数为零，且二阶导数为负，所以这是一个最大值。因此，馅饼被分给参与者 I \(\frac{1}{3}\) ，分给参与者 II \(\frac{2}{3}\) 。相应的效用对 (U, V) 为 \(\left( {1/\sqrt{3},\frac{2}{3}}\right)\) ，约为 (0.577,0.667)。关于使用不同效用函数(utility function)的类似计算，见练习 11.4。
\end{tcolorbox}

\section*{11.7 The Ultimatum Game}

\begin{tcolorbox}
\section*{11.7 最后通牒博弈(Ultimatum game)}
\end{tcolorbox}

Assume that two players have to split a unit pie, as described in the previous section. In the following, we will develop the "alternating offers" bargaining game applied to this situation. This is a non-cooperative game that models how players may interact in a bargaining situation. The solution to this non-cooperative game, in its most elaborate version of "stationary strategies" (see Section 11.10), will approximate the Nash bargaining solution. The model justifies the Nash solution via a different approach, which is not axiomatic, but instead uses the concept of subgame-perfect equilibria in game trees.

\begin{tcolorbox}
假设如前一节所述，两名参与者必须分割一个单位的馅饼。接下来，我们将构建适用于这种情况的“交替出价”讨价还价博弈。这是一个非合作博弈，用于模拟参与者在讨价还价情境中可能的互动方式。这个非合作博弈在其最详尽的“平稳策略”版本(见第11.10节)下的解，将逼近纳什讨价还价解。该模型通过一种不同的方法来证明纳什解的合理性，这种方法不是公理化的，而是使用了博弈树中的子博弈完美均衡概念。
\end{tcolorbox}

The model is that the pie is split according to the suggestion of one of the players. Suppose player I suggests splitting the pie into \(x\) for player I and \(1 - x\) for player II, so that the players receive utility \(u\left( x\right)\) and \(v\left( {1 - x}\right)\) , respectively. The amount \(x\) proposed by player I is also called his demand because that is the share of the pie that he receives in the proposed split.

\begin{tcolorbox}
该模型是，馅饼根据其中一名参与者的提议进行分割。假设参与者I提议将馅饼分割为参与者I得到 \(x\) ，参与者II得到 \(1 - x\) ，这样参与者分别获得效用 \(u\left( x\right)\) 和 \(v\left( {1 - x}\right)\) 。参与者I提出的数量 \(x\) 也被称为他的需求，因为这是他在提议的分割方案中所获得的馅饼份额。
\end{tcolorbox}

\begin{center}
\includegraphics[max width=0.6\textwidth]{images/bo_d2e5cmbef24c73b03tn0_325_517_493_887_496_0.jpg}
\end{center}
\hspace*{3em} 

Figure 11.8 Discrete version of the Ultimatum game. The black and white arrows denote the two subgame-perfect equilibria, which are the same except for player II’s response to the full demand \(x = 1\) , and player I’s best response of \(x = {.99}\) or \(x = 1\) .

\begin{tcolorbox}
图11.8 最后通牒博弈的离散版本。黑白箭头表示两个子博弈完美均衡，除了参与者II对全额需求 \(x = 1\) 的回应，以及参与者I对 \(x = {.99}\) 或 \(x = 1\) 的最优反应外，这两个均衡是相同的。
\end{tcolorbox}

In the most basic form, this model is called the Ultimatum game, where player II can either accept the split or reject it. If player II accepts, the payoff pair is \(\left( {u\left( x\right) ,v\left( {1 - x}\right) }\right)\) , and if she rejects, it is(0,0), in which case neither player receives anything. Figure 11.8 shows a discretized version of this game where the possible demand \(x\) is a multiple of 0.01 . Hence, player I has 101 possible actions for \(x \in  \left\{  {\left. \frac{i}{100}\right| \;i = 0,1,\ldots ,{100}}\right\}\) , and for each proposed \(x\) player II can either choose \(A\) (accept) or \(R\) (reject). Player II therefore has \({2}^{101}\) pure strategies, given by the different combinations of \(A\) and \(R\) for each possible \(x\) .

\begin{tcolorbox}
在最基本的形式下，这个模型被称为最后通牒博弈，在这个博弈中，参与者II可以接受或拒绝分割方案。如果参与者II接受，支付对为 \(\left( {u\left( x\right) ,v\left( {1 - x}\right) }\right)\) ；如果她拒绝，支付对为(0,0)，在这种情况下，两名参与者都得不到任何东西。图11.8展示了这个博弈的离散版本，其中可能的需求 \(x\) 是0.01的倍数。因此，参与者I对于 \(x \in  \left\{  {\left. \frac{i}{100}\right| \;i = 0,1,\ldots ,{100}}\right\}\) 有101种可能的行动，并且对于每个提议的 \(x\) ，参与者II可以选择 \(A\) (接受)或 \(R\) (拒绝)。因此，参与者II有 \({2}^{101}\) 种纯策略，这些策略由每个可能的 \(x\) 对应的 \(A\) 和 \(R\) 的不同组合给出。
\end{tcolorbox}

The game is shown in Figure 11.8, and we look for a subgame-perfect equilibrium (SPE) of the game, that is, because the game has perfect information, we apply backward induction. Whenever \(1 - x > 0\) , the utility \(v\left( {1 - x}\right)\) to player II is positive, because otherwise

\begin{tcolorbox}
该博弈如图11.8所示，我们要寻找该博弈的子博弈完美均衡(SPE)，也就是说，因为该博弈具有完美信息，我们应用逆向归纳法。只要 \(1 - x > 0\) ，参与者II的效用 \(v\left( {1 - x}\right)\) 就是正的，否则
\end{tcolorbox}

\[
0 = v\left( {1 - x}\right)  = v\left( {x \cdot  0 + \left( {1 - x}\right)  \cdot  1}\right)  < 1 - x = x \cdot  v\left( 0\right)  + \left( {1 - x}\right)  \cdot  v\left( 1\right) ,
\]

which contradicts concavity of \(v\) (and is also seen graphically). Hence, the SPE condition implies that player II accepts any positive amount \(1 - x\) offered by player I. These choices, determined by backward induction, are shown as arrows in Figure 11.8, and they are unique except when \(x = 1\) where player I demands the whole pie. Then player II is indifferent between accepting and rejecting because she will get nothing in either case. Consequently, in terms of pure strategies, both \({AA}\cdots {AA}\) and \({AA}\cdots {AR}\) are SPE strategies of player II. Here, the first 100 \({A}^{\prime }\) s denote acceptance for each \(x = 0,{0.01},\ldots ,{0.99}\) , and the last \(A\) or \(R\) is the choice when \(x = 1\) , shown by a white or black arrow, respectively, in Figure 11.8. Given this strategy of player II, player I's best response is to maximize his payoff, which he does as shown by the corresponding white or black arrow for player I. If player II accepts the demand \(x = 1\) , then player I demands \(x = 1\) , and the pair \(\left( {1,{AA}\cdots {AA}}\right)\) is an SPE, and if player II rejects the demand \(x = 1\) , then player I demands \(x = {0.99}\) (because for \(x = 1\) he would get nothing), and the resulting SPE is \(\left( {{0.99},{AA}\cdots {AR}}\right)\) . These are two subgame-perfect equilibria in pure strategies. If \(u\left( {0.99}\right)  < 1\) , they are also unique; if \(u\left( {0.99}\right)  = 1\) then the demand \(x = {.99}\) is also a best response by player I to the strategy \({AA}\cdots {AA}\) of player II that accepts every demand.

\begin{tcolorbox}
这与 \(v\) 的凹性相矛盾(从图上也能看出)。因此，子博弈完美均衡(SPE)条件意味着玩家 II 会接受玩家 I 提供的任何正数量 \(1 - x\) 。这些通过逆向归纳法确定的选择在图 11.8 中用箭头表示，并且除了 \(x = 1\) (即玩家 I 要求获得整个“馅饼”)的情况外，这些选择是唯一的。此时，玩家 II 在接受和拒绝之间无差异，因为在这两种情况下她都得不到任何东西。因此，就纯策略而言， \({AA}\cdots {AA}\) 和 \({AA}\cdots {AR}\) 都是玩家 II 的子博弈完美均衡策略。这里，前 100 个 \({A}^{\prime }\) 表示对每个 \(x = 0,{0.01},\ldots ,{0.99}\) 的接受，最后一个 \(A\) 或 \(R\) 是当 \(x = 1\) 时的选择，分别在图 11.8 中用白色或黑色箭头表示。给定玩家 II 的这个策略，玩家 I 的最优反应是最大化他的收益，这通过玩家 I 对应的白色或黑色箭头来表示。如果玩家 II 接受要求 \(x = 1\) ，那么玩家 I 要求 \(x = 1\) ，并且组合 \(\left( {1,{AA}\cdots {AA}}\right)\) 是一个子博弈完美均衡；如果玩家 II 拒绝要求 \(x = 1\) ，那么玩家 I 要求 \(x = {0.99}\) (因为对于 \(x = 1\) 他将得不到任何东西)，结果得到的子博弈完美均衡是 \(\left( {{0.99},{AA}\cdots {AR}}\right)\) 。这些是纯策略中的两个子博弈完美均衡。如果 \(u\left( {0.99}\right)  < 1\) ，它们也是唯一的；如果 \(u\left( {0.99}\right)  = 1\) ，那么要求 \(x = {.99}\) 也是玩家 I 对玩家 II 接受所有要求的策略 \({AA}\cdots {AA}\) 的最优反应。
\end{tcolorbox}

In essence, player II receives virtually nothing in the Ultimatum game because of the backward induction assumption. This does not reflect people's behavior in real bargaining situations of this kind, as confirmed in laboratory experiments where people play the Ultimatum game. Typically, such an experiment is conducted where player I states, unknown to player II, his demand \(x\) , and player II declares simultaneously, unknown to player I, her own "reserve demand" \(z\) ; the two players may even be matched anonymously from a group of many subjects that take the respective roles of players I and II. Whenever \(1 - x \geq  z\) , that is, the offer to player II is at least as large as her reserve demand, the two players get \(x\) and \(1 - x\) , otherwise nothing. (So player II may declare a small reserve demand like \(z = {0.01}\) but still receive 0.4 because player I’s demand is \(x = {0.6}\) ; both can be regarded as reasonable choices, but they are not mutual best responses.) In such experiments, player II often makes positive reserve demands like \(z = {0.4}\) , and player I makes cautious demands like \(x = {0.5}\) , which contradict the SPE assumption. A better description of these experiments may therefore be equilibria that are not subgame-perfect, with "threats" by player II to reject offers \(1 - x\) that are too low. Nevertheless, we assume SPE because they allow definite conclusions, and because the model will be refined, and made more realistic, by allowing rejections to be followed by "counter-demands". (Exercise 11.3 asks you to find all equilibria of the discrete Ultimatum game.)

\begin{tcolorbox}
本质上，由于逆向归纳法假设，在最后通牒博弈中玩家 II 实际上得不到任何东西。这并不反映人们在这类实际讨价还价情况中的行为，正如在人们进行最后通牒博弈的实验室实验中所证实的那样。通常，进行这样的实验:玩家 I 在玩家 II 不知情的情况下声明他的要求 \(x\) ，玩家 II 在玩家 I 不知情的情况下同时声明她自己的“保留要求” \(z\) ；这两个玩家甚至可以从众多分别扮演玩家 I 和玩家 II 角色的受试者群体中匿名配对。只要 \(1 - x \geq  z\) ，即给玩家 II 的提议至少和她的保留要求一样大，两个玩家就分别得到 \(x\) 和 \(1 - x\) ，否则什么都得不到。(所以玩家 II 可能声明一个小的保留要求，比如 \(z = {0.01}\) ，但仍然能得到 0.4，因为玩家 I 的要求是 \(x = {0.6}\) ；这两者都可以被视为合理的选择，但它们不是相互的最优反应。)在这样的实验中，玩家 II 经常提出像 \(z = {0.4}\) 这样的正保留要求，玩家 I 提出像 \(x = {0.5}\) 这样谨慎的要求，这与子博弈完美均衡假设相矛盾。因此，对这些实验更好的描述可能是非子博弈完美的均衡，其中玩家 II“威胁”拒绝过低的提议 \(1 - x\) 。然而，我们假设子博弈完美均衡，因为它们能得出明确的结论，并且因为通过允许拒绝后有“反要求”，模型将得到改进并变得更现实。(练习 11.3 要求你找出离散最后通牒博弈的所有均衡。)
\end{tcolorbox}

The described two SPE of the discrete Ultimatum game are found whenever the unit pie can be split into \(M + 1\) possible splits(x,1 - x)for \(x = \frac{i}{M}\) , for \(i = 0,1,\ldots ,M - 1,M\) , like \(M = {100}\) above. Of her \({2}^{M + 1}\) pure strategies, player II can play only two in an SPE, namely always choosing \(A\) whenever \(x < 1\) , and choosing either \(A\) or \(R\) in response to \(x = 1\) . The respective demand by player \(\mathrm{I}\) is \(x = 1\) or \(x = 1 - \frac{1}{M}\) . With a very fine discretization where \(M\) becomes very large, these two SPE are nearly identical, with player I demanding all (or nearly all) of the pie, and player II receiving almost nothing.

\begin{tcolorbox}
每当单位馅饼可以被分割成 \(M + 1\) 种可能的分配方式 (x, 1 - x)，其中 \(x = \frac{i}{M}\) ，对于 \(i = 0,1,\ldots ,M - 1,M\) ，就像上面的 \(M = {100}\) 一样，就可以找到离散最后通牒博弈所描述的两种子博弈完美均衡(SPE)。在她的 \({2}^{M + 1}\) 种纯策略中，玩家 II 在子博弈完美均衡中只能采用两种策略，即每当 \(x < 1\) 时总是选择 \(A\) ，以及针对 \(x = 1\) 选择 \(A\) 或 \(R\) 。玩家 \(\mathrm{I}\) 相应的要求是 \(x = 1\) 或 \(x = 1 - \frac{1}{M}\) 。当离散化非常精细，即 \(M\) 变得非常大时，这两种子博弈完美均衡几乎相同，玩家 I 要求得到全部(或几乎全部)馅饼，而玩家 II 几乎一无所获。
\end{tcolorbox}

\begin{center}
\includegraphics[max width=0.4\textwidth]{images/bo_d2e5cmbef24c73b03tn0_327_392_210_524_474_0.jpg}
\end{center}
\hspace*{3em} 

Figure 11.9 The Continuous Ultimatum game. Player I makes any demand \(x \in  \left\lbrack  {0,1}\right\rbrack\) , which player II either accepts (A) or rejects(R). Only one decision point of player II is shown.

\begin{tcolorbox}
图 11.9 连续最后通牒博弈。玩家 I 提出任何要求 \(x \in  \left\lbrack  {0,1}\right\rbrack\) ，玩家 II 可以选择接受(A)或拒绝(R)。这里只展示了玩家 II 的一个决策点。
\end{tcolorbox}

In the Continuous Ultimatum game, the possible demand \(x\) by player \(\mathrm{I}\) is any number in the interval \(\left\lbrack  {0,1}\right\rbrack\) , which defines an infinite number of strategies for player I. A strategy of player II is then given by a function \(f : \left\lbrack  {0,1}\right\rbrack   \rightarrow  \{ A,R\}\) , where \(f\left( x\right)  = A\) means that player II accepts the demand \(x\) , and \(f\left( x\right)  = R\) means that player II rejects it. (Equivalently, a strategy of player II is an arbitrary subset of \(\left\lbrack  {0,1}\right\rbrack\) that consists exactly of those demands \(x\) that she accepts.) This infinite game is indicated in Figure 11.9. The infinitely many choices \(x\) are indicated by the triangle, whose baseline represents the interval \(\left\lbrack  {0,1}\right\rbrack\) . Every point of that baseline has a separate decision point for player II where she can decide between \(A\) and \(R\) .

\begin{tcolorbox}
在连续最后通牒博弈中，玩家 \(\mathrm{I}\) 可能提出的要求 \(x\) 是区间 \(\left\lbrack  {0,1}\right\rbrack\) 内的任意数字，这为玩家 I 定义了无限种策略。玩家 II 的一种策略由一个函数 \(f : \left\lbrack  {0,1}\right\rbrack   \rightarrow  \{ A,R\}\) 给出，其中 \(f\left( x\right)  = A\) 表示玩家 II 接受要求 \(x\) ， \(f\left( x\right)  = R\) 表示玩家 II 拒绝该要求。(等价地，玩家 II 的一种策略是 \(\left\lbrack  {0,1}\right\rbrack\) 的任意一个子集，该子集恰好由她接受的那些要求 \(x\) 组成。)这个无限博弈如图 11.9 所示。无限多种选择 \(x\) 由三角形表示，其底边代表区间 \(\left\lbrack  {0,1}\right\rbrack\) 。底边的每个点都对应玩家 II 的一个单独决策点，她可以在 \(A\) 和 \(R\) 之间做出选择。
\end{tcolorbox}

We want to show that the Continuous Ultimatum game has a unique SPE. This requires an additional assumption, namely that the utility function \(u\) of player \(\mathrm{I}\) is strictly increasing. This is the case if \({x}^{\prime } = 1\) in the following lemma.

\begin{tcolorbox}
我们要证明连续最后通牒博弈有唯一的子博弈完美均衡。这需要一个额外的假设，即玩家 \(\mathrm{I}\) 的效用函数 \(u\) 是严格递增的。如果满足下面引理中的 \({x}^{\prime } = 1\) ，就会出现这种情况。
\end{tcolorbox}

Lemma 11.7. Let \(u : \left\lbrack  {0,1}\right\rbrack   \rightarrow  \mathbb{R}\) be an increasing, concave, and continuous function with \(u\left( 0\right)  = 0\) and \(u\left( 1\right)  = 1\) . Then \(u\) is strictly increasing on an interval \(\left\lbrack  {0,{x}^{\prime }}\right\rbrack\) , where \(u\left( x\right)  = 1\) for all \(x \in  \left\lbrack  {{x}^{\prime },1}\right\rbrack\) .

\begin{tcolorbox}
引理 11.7。设 \(u : \left\lbrack  {0,1}\right\rbrack   \rightarrow  \mathbb{R}\) 是一个递增、凹且连续的函数，且满足 \(u\left( 0\right)  = 0\) 和 \(u\left( 1\right)  = 1\) 。那么 \(u\) 在区间 \(\left\lbrack  {0,{x}^{\prime }}\right\rbrack\) 上是严格递增的，其中对于所有的 \(x \in  \left\lbrack  {{x}^{\prime },1}\right\rbrack\) 有 \(u\left( x\right)  = 1\) 。
\end{tcolorbox}

Proof. The set \(I = \{ x \in  \left\lbrack  {0,1}\right\rbrack   \mid  u\left( x\right)  = 1\}\) is not empty (it contains 1) and closed because \(u\) is continuous. Let \({x}^{\prime }\) be the smallest element of \(I\) , with \(u\left( {x}^{\prime }\right)  = u\left( 1\right)  = 1\) . Then \(I = \left\lbrack  {{x}^{\prime },1}\right\rbrack\) because \({x}^{\prime } \leq  x \leq  1\) implies \(1 = u\left( {x}^{\prime }\right)  \leq  u\left( x\right)  \leq  u\left( 1\right)  = 1\) . We have \({x}^{\prime } > 0\) because \(u\left( 0\right)  = 0 < 1\) . To see that \(u\) is strictly increasing on the interval \(\left\lbrack  {0,{x}^{\prime }}\right\rbrack\) , suppose there were \(\bar{x}\) and \(\widehat{x}\) in that interval with \(\bar{x} < \widehat{x}\) but \(u\left( \bar{x}\right)  = u\left( \widehat{x}\right)  < u\left( {x}^{\prime }\right)  = 1\) . Then it is easy to see (a drawing will help) that \(u\left( \widehat{x}\right)\) is below the line that connects \(u\left( \bar{x}\right)\) and \(u\left( {x}^{\prime }\right)\) , which contradicts the concavity of \(u\) .

\begin{tcolorbox}
证明。集合 \(I = \{ x \in  \left\lbrack  {0,1}\right\rbrack   \mid  u\left( x\right)  = 1\}\) 非空(它包含 1)且是闭集，因为 \(u\) 是连续的。设 \({x}^{\prime }\) 为 \(I\) 中的最小元素，且满足 \(u\left( {x}^{\prime }\right)  = u\left( 1\right)  = 1\) 。那么 \(I = \left\lbrack  {{x}^{\prime },1}\right\rbrack\) ，因为 \({x}^{\prime } \leq  x \leq  1\) 意味着 \(1 = u\left( {x}^{\prime }\right)  \leq  u\left( x\right)  \leq  u\left( 1\right)  = 1\) 。我们有 \({x}^{\prime } > 0\) ，因为 \(u\left( 0\right)  = 0 < 1\) 。为了说明 \(u\) 在区间 \(\left\lbrack  {0,{x}^{\prime }}\right\rbrack\) 上严格递增，假设在该区间内存在 \(\bar{x}\) 和 \(\widehat{x}\) ，满足 \(\bar{x} < \widehat{x}\) 但 \(u\left( \bar{x}\right)  = u\left( \widehat{x}\right)  < u\left( {x}^{\prime }\right)  = 1\) 。那么很容易看出(画个图会有帮助) \(u\left( \widehat{x}\right)\) 位于连接 \(u\left( \bar{x}\right)\) 和 \(u\left( {x}^{\prime }\right)\) 的直线下方，这与 \(u\) 的凹性相矛盾。
\end{tcolorbox}

If \({x}^{\prime } < 1\) in Lemma 11.7, then player \(\mathrm{I}\) is perfectly satisifed by receiving only \({x}^{\prime }\) , which is less than a full share of the pie. Then any demand \(x\) with \({x}^{\prime } \leq  x < 1\) is accepted by player II because she receives \(v\left( {1 - x}\right)  > 0\) , and \(x\) is a possible strategy of player I in an SPE. In that case, we do not get a unique SPE. We therefore require that \(u\) is strictly increasing, or equivalently, by Lemma 11.7, that \(u\left( x\right)  = 1\) implies \(x = 1\) . Hence, we assume that the utility function \(u\) of player \(\mathrm{I}\) , and similarly the utility function \(v\) of player II, is strictly increasing.

\begin{tcolorbox}
如果引理 11.7 中的 \({x}^{\prime } < 1\) 成立，那么玩家 \(\mathrm{I}\) 仅获得 \({x}^{\prime }\) 就会完全满意，而 \({x}^{\prime }\) 小于整个“馅饼”的份额。那么任何满足 \({x}^{\prime } \leq  x < 1\) 的需求 \(x\) 都会被玩家 II 接受，因为她能得到 \(v\left( {1 - x}\right)  > 0\) ，并且 \(x\) 是玩家 I 在子博弈完美均衡(SPE)中的一个可能策略。在这种情况下，我们无法得到唯一的子博弈完美均衡。因此，我们要求 \(u\) 严格递增，或者等价地，根据引理 11.7， \(u\left( x\right)  = 1\) 意味着 \(x = 1\) 。因此，我们假设玩家 \(\mathrm{I}\) 的效用函数 \(u\) ，以及类似地玩家 II 的效用函数 \(v\) ，都是严格递增的。
\end{tcolorbox}

The following proposition states that the only SPE of the Continuous Ultimatum game is where player I demands the whole pie, with \(x = 1\) , and where player II accepts any demand, including the demand \(x = 1\) on the equilibrium path where she is indifferent between accepting and rejecting. The reasoning is more subtle than the SPE just being a "limit" of the discretized Ultimatum game.

\begin{tcolorbox}
以下命题表明，连续最后通牒博弈(Continuous Ultimatum game)的唯一子博弈完美均衡是玩家 I 要求整个“馅饼”，即 \(x = 1\) ，并且玩家 II 接受任何需求，包括在均衡路径(equilibrium path)上的需求 \(x = 1\) ，在该均衡路径上她对接受和拒绝无差异。其推理比子博弈完美均衡仅仅是离散化最后通牒博弈(Ultimatum game)的“极限”更为微妙。
\end{tcolorbox}

Proposition 11.8. Let \(u,v : \left\lbrack  {0,1}\right\rbrack   \rightarrow  \mathbb{R}\) be strictly increasing, concave, and continuous utility functions of player \(I\) and \({II}\) , respectively, with \(u\left( 0\right)  = v\left( 0\right)  = 0\) and \(u\left( 1\right)  = v\left( 1\right)  = 1\) . Then the Continuous Ultimatum game has the unique SPE where the demand \(x\) of player I is \(x = 1\) and player II always chooses \(A\) .

\begin{tcolorbox}
命题 11.8。设 \(u,v : \left\lbrack  {0,1}\right\rbrack   \rightarrow  \mathbb{R}\) 分别是玩家 \(I\) 和 \({II}\) 的严格递增、凹且连续的效用函数，且满足 \(u\left( 0\right)  = v\left( 0\right)  = 0\) 和 \(u\left( 1\right)  = v\left( 1\right)  = 1\) 。那么连续最后通牒博弈(Continuous Ultimatum game)有唯一的子博弈完美均衡，其中玩家 I 的需求 \(x\) 为 \(x = 1\) ，并且玩家 II 总是选择 \(A\) 。
\end{tcolorbox}

Proof. The described strategies define an SPE. To show uniqueness, first observe that player II has to accept whenever \(x < 1\) because then she receives a positive payoff \(v\left( {1 - x}\right)\) . Suppose there is an SPE where player I’s demand is \(x\) so that \(x < 1\) . Because \(u\) is strictly increasing, player I would get more by demanding a little more, but still less than 1, for example \(\frac{x + 1}{2}\) , which player II still accepts. Hence, \(x = 1\) is the only strategy by player I in an SPE. On the equilibrium path, player II is indifferent, and can choose either \(A\) or \(R\) in response to \(x = 1\) , and in fact could randomize between \(A\) and \(R\) . Suppose that \(R\) is played with any positive probability \(\varepsilon  > 0\) , giving player I expected payoff \(\left( {1 - \varepsilon }\right) u\left( 1\right)  + {\varepsilon u}\left( 0\right)\) , which is \(1 - \varepsilon\) . But then player I’s best response to that strategy of player II would not be \(x = 1\) . Instead, player I could improve his payoff, for example by demanding \(1 - \frac{\varepsilon }{2}\) because then player II would accept for sure and player I receives \(u\left( {1 - \frac{\varepsilon }{2}}\right)  \geq  1 - \frac{\varepsilon }{2} > 1 - \varepsilon\) . However, no demand less than 1 is possible in an SPE. Hence, if player II does accept the demand \(x = 1\) with certainty, then this reaction of player II does not define an SPE, so the described SPE is the only SPE.

\begin{tcolorbox}
证明。所描述的策略定义了一个子博弈完美均衡(SPE)。为了证明其唯一性，首先观察到，只要 \(x < 1\) ，玩家 II 就必须接受，因为此时她能获得正收益 \(v\left( {1 - x}\right)\) 。假设存在一个子博弈完美均衡，其中玩家 I 的要求是 \(x\) ，使得 \(x < 1\) 。由于 \(u\) 是严格递增的，玩家 I 通过多要求一点，但仍少于 1，例如 \(\frac{x + 1}{2}\) ，能获得更多收益，而玩家 II 仍然会接受。因此， \(x = 1\) 是玩家 I 在子博弈完美均衡中的唯一策略。在均衡路径上，玩家 II 无差异，可以选择以 \(A\) 或 \(R\) 来回应 \(x = 1\) ，实际上也可以在 \(A\) 和 \(R\) 之间随机选择。假设以正概率 \(\varepsilon  > 0\) 选择 \(R\) ，这使得玩家 I 的期望效用为 \(\left( {1 - \varepsilon }\right) u\left( 1\right)  + {\varepsilon u}\left( 0\right)\) ，即 \(1 - \varepsilon\) 。但这样一来，玩家 I 针对玩家 II 该策略的最优反应就不会是 \(x = 1\) 。相反，玩家 I 可以提高自己的收益，例如要求 \(1 - \frac{\varepsilon }{2}\) ，因为此时玩家 II 肯定会接受，玩家 I 能获得 \(u\left( {1 - \frac{\varepsilon }{2}}\right)  \geq  1 - \frac{\varepsilon }{2} > 1 - \varepsilon\) 。然而，在子博弈完美均衡中，要求少于 1 是不可能的。因此，如果玩家 II 确定会接受要求 \(x = 1\) ，那么玩家 II 的这种反应并不能定义一个子博弈完美均衡，所以所描述的子博弈完美均衡是唯一的子博弈完美均衡。
\end{tcolorbox}

In summary, with demands that can be chosen continuously, the demand of a player in an SPE makes the other player indifferent between accepting and rejecting, which the other player nevertheless accepts with certainty. This will also be the guiding principle in SPE of more complicated games, which we study next.

\begin{tcolorbox}
总之，当要求可以连续选择时，在子博弈完美均衡中，一个玩家的要求会使另一个玩家在接受和拒绝之间无差异，但另一个玩家仍然会确定地接受。这也将是我们接下来要研究的更复杂博弈的子博弈完美均衡的指导原则。
\end{tcolorbox}

\section*{11.8 Alternating Offers Over Two Rounds}

\begin{tcolorbox}
\section*{11.8 两轮交替报价}
\end{tcolorbox}

Next, we extend the bargaining model so that once a player has rejected the demand of the other player, she can make a counter-demand that, in turn, the other player can accept or reject. In order to find an SPE of this game by backward induction, these alternating offers have to end after some time. The simplest case is bargaining over two rounds, discussed in this section. In the first round, player I makes a demand \(x\) , which player II can accept or reject. In the latter case, she can make a counter-demand \(y\) that represents the share of the pie she claims for herself, so the pie is split into \(1 - y\) for player I and \(y\) for player II. If player I accepts that demand, he receives utility \(u\left( {1 - y}\right)\) and player II receives utility \(v\left( y\right)\) . If player \(\mathrm{I}\) rejects the final demand, both players receive zero.

\begin{tcolorbox}
接下来，我们扩展讨价还价模型，使得一旦一个玩家拒绝了另一个玩家的要求，她可以提出反要求，而另一个玩家又可以接受或拒绝该反要求。为了通过逆向归纳法找到这个博弈的子博弈完美均衡，这些交替报价必须在一段时间后结束。最简单的情况是两轮讨价还价，本节将对此进行讨论。在第一轮中，玩家 I 提出要求 \(x\) ，玩家 II 可以接受或拒绝。如果拒绝，她可以提出反要求 \(y\) ，这代表她为自己争取的份额，因此馅饼被分割为玩家 I 得到 \(1 - y\) ，玩家 II 得到 \(y\) 。如果玩家 I 接受该要求，他获得效用 \(u\left( {1 - y}\right)\) ，玩家 II 获得效用 \(v\left( y\right)\) 。如果玩家 \(\mathrm{I}\) 拒绝最终要求，两个玩家的收益都为零。
\end{tcolorbox}

\begin{center}
\includegraphics[max width=1.0\textwidth]{images/bo_d2e5cmbef24c73b03tn0_329_325_340_1271_893_0.jpg}
\end{center}
\hspace*{3em} 

Figure 11.10 Two-round bargaining with demand \(x\) by player I in round one and counter-demand \(y\) by player II in round two. Left: A second round takes place with probability \(\delta\) , otherwise termination with payoffs(0,0). Right: Simplified, equivalent game with \(\delta\) as a discount factor applied to the second-round utility values.

\begin{tcolorbox}
图 11.10 两轮讨价还价，第一轮玩家 I 提出要求 \(x\) ，第二轮玩家 II 提出反要求 \(y\) 。左图:第二轮以概率 \(\delta\) 发生，否则以收益 (0, 0) 结束博弈。右图:简化的等价博弈，其中 \(\delta\) 作为贴现因子应用于第二轮的效用值。
\end{tcolorbox}

In this game, player II's demand in the final round is final. Hence, this subgame is the Ultimatum game with the roles of the players exchanged, where by Proposition 11.8 player II can demand the whole pie. Hence, there seems to be no incentive for her to accept anything offered to her in the first round. For that reason, the model has an additional feature shown in Figure 11.10. Namely, when player II rejects the first demand, there is a positive probability \(1 - \delta\) that bargaining terminates with no agreement, where we assume \(0 < \delta  < 1\) . This is the unsuccessful outcome, with payoff zero to both, given as the result of a chance move. The game proceeds to player II’s counter-demand only with probability \(\delta\) .

\begin{tcolorbox}
在这个博弈中，玩家 II 在最后一轮的要求是最终要求。因此，这个子博弈是玩家角色互换的最后通牒博弈，根据命题 11.8，玩家 II 可以要求整个馅饼。因此，她似乎没有动力在第一轮接受提供给她的任何东西。出于这个原因，该模型有一个如图 11.10 所示的额外特征。即，当玩家 II 拒绝第一个要求时，有一个正概率 \(1 - \delta\) 使得讨价还价以未达成协议而结束，我们假设 \(0 < \delta  < 1\) 。这是不成功的结果，双方收益都为零，这是由一个随机事件决定的。博弈只有以概率 \(\delta\) 进入玩家 II 的反要求阶段。
\end{tcolorbox}

By computing expected values, the chance move can be eliminated from the game, which gives the game shown on the right in Figure 11.10. In that game, the pair \(\left( {u\left( {1 - y}\right) ,v\left( y\right) }\right)\) of utilities when agreement is reached in the second round is multiplied by \(\delta\) , which is a reduction of the payoffs because \(\delta  < 1\) . The usual interpretation of \(\delta\) is that of a discount factor, which expresses that later payoffs are worth less than earlier payoffs.

\begin{tcolorbox}
通过计算期望值，可以将博弈中的机会节点消除，从而得到图11.10右侧所示的博弈。在该博弈中，第二轮达成协议时的效用对 \(\left( {u\left( {1 - y}\right) ,v\left( y\right) }\right)\) 乘以 \(\delta\) ，这是对收益的一种折减，因为 \(\delta  < 1\) 。通常将 \(\delta\) 解释为贴现因子，它表示后期收益的价值低于前期收益。
\end{tcolorbox}

A discount factor \(\delta\) applied to future payoffs is a realistic assumption. In monetary terms, it can express lost interest payments (for money today as opposed to money in the future). More importantly, future payoffs are less secure because, for example, the buyer or seller in a deal may change her mind, or because of other unforeseen circumstances. This is modeled with the chance move in the left game in Figure 11.10. In addition, such a chance move justifies why both players should have the same discount factor. Alternatively, one may use different discount factors \({\delta }_{\mathrm{I}}\) and \({\delta }_{\mathrm{{II}}}\) for the two players, which give second-round payoffs \({\delta }_{\mathrm{I}}u\left( {1 - y}\right)\) and \({\delta }_{\mathrm{{II}}}v\left( y\right)\) . If these discount factors are different, a higher discount factor represents a more patient player whose future payoff differs less from the present payoff. However, the resulting game would be more complicated to analyze, and would not lead to a direct connection with the Nash bargaining solution. We therefore use only a single discount factor, as represented by the model where \(\delta\) is the probability that bargaining continues into the next round.

\begin{tcolorbox}
对未来收益应用贴现因子 \(\delta\) 是一个现实的假设。从货币角度来看，它可以表示损失的利息支付(今天的钱与未来的钱相比)。更重要的是，未来收益的安全性较低，例如，交易中的买方或卖方可能会改变主意，或者由于其他不可预见的情况。这在图11.10左侧博弈的机会节点中得到了体现。此外，这样的机会节点解释了为什么两个参与者应该使用相同的贴现因子。或者，也可以为两个参与者使用不同的贴现因子 \({\delta }_{\mathrm{I}}\) 和 \({\delta }_{\mathrm{{II}}}\) ，它们会产生第二轮收益 \({\delta }_{\mathrm{I}}u\left( {1 - y}\right)\) 和 \({\delta }_{\mathrm{{II}}}v\left( y\right)\) 。如果这些贴现因子不同，较高的贴现因子代表更有耐心的参与者，其未来收益与当前收益的差异较小。然而，由此产生的博弈分析起来会更复杂，并且不会与纳什讨价还价解建立直接联系。因此，我们仅使用一个贴现因子，如 \(\delta\) 是讨价还价进入下一轮的概率的模型所示。
\end{tcolorbox}

An SPE of the game is found as follows. We use backward induction and analyze the last round first. When player II makes her demand \(y\) to player \(\mathrm{I}\) , he can only accept or reject. This is the Ultimatum game in Proposition 11.8 (with exchanged players). Hence, player II demands the whole pie with \(y = 1\) , and player I accepts any offer made to him, in particular this demand \(y = 1\) (recall that player I’s strategy is a function of the demand \(y\) that he faces), and player I will get nothing. Here, it does not matter if the last stage represents the subgame that starts with player II’s demand \(y\) in the left or right game in Figure 11.10, because the payoffs in the latter are merely scaled with the positive scalar \(\delta\) . These subgames in the two games in Figure 11.10 are strategically equivalent and have the same subgame-perfect equilibria. With this assumption, if player II rejects player I's demand \(x\) in the first round and chooses \(R\) , then player I gets payoff zero, and player II gets \({\delta v}\left( 1\right)\) , which is \(\delta\) .

\begin{tcolorbox}
该博弈的子博弈完美均衡(SPE)的求解如下。我们使用逆向归纳法，首先分析最后一轮。当参与者II向参与者 \(\mathrm{I}\) 提出需求 \(y\) 时，他只能接受或拒绝。这是命题11.8中的最后通牒博弈(参与者角色互换)。因此，参与者II以 \(y = 1\) 的方式要求得到整个“蛋糕”，参与者I会接受向他提出的任何提议，特别是这个需求 \(y = 1\) (回想一下，参与者I的策略是他所面临的需求 \(y\) 的函数)，并且参与者I将一无所获。在这里，最后阶段代表的是从参与者II在图11.10左侧或右侧博弈中提出需求 \(y\) 开始的子博弈并不重要，因为后者的收益只是用正标量 \(\delta\) 进行了缩放。图11.10中两个博弈的这些子博弈在策略上是等价的，并且具有相同的子博弈完美均衡。在这个假设下，如果参与者II在第一轮拒绝了参与者I的需求 \(x\) 并选择 \(R\) ，那么参与者I的收益为零，参与者II的收益为 \({\delta v}\left( 1\right)\) ，即 \(\delta\) 。
\end{tcolorbox}

By backward induction, what should player I do in the first round? Player II's reaction is unique except when she is indifferent. This means that she will accept the proposed split if she is offered a utility \(v\left( {1 - x}\right)\) (now undiscounted, because this is the first round) that exceeds the amount \(\delta\) that she can expect to get by rejection; she is indifferent if \(v\left( {1 - x}\right)  = \delta\) ; and she will reject the split if \(v\left( {1 - x}\right)  < \delta\) . By the same reasoning as for Proposition 11.8, the demand \(x\) by player I in the SPE is chosen so that \(v\left( {1 - x}\right)  = \delta\) , and player II accepts this on the equilibrium path. (If she did not accept, player I would offer her a little bit more and make her accept. However, the only SPE is where player I makes player II indifferent.)

\begin{tcolorbox}
通过逆向归纳法，参与者I在第一轮应该怎么做呢？除了参与者II无差异的情况外，她的反应是唯一的。这意味着，如果给她提供的效用 \(v\left( {1 - x}\right)\) (现在未贴现，因为这是第一轮)超过了她通过拒绝所能期望得到的数量 \(\delta\) ，她将接受提议的分配方案；如果 \(v\left( {1 - x}\right)  = \delta\) ，她将无差异；如果 \(v\left( {1 - x}\right)  < \delta\) ，她将拒绝该分配方案。根据与命题11.8相同的推理，在子博弈完美均衡中，参与者I的需求 \(x\) 被选择为使得 \(v\left( {1 - x}\right)  = \delta\) ，并且参与者II在均衡路径上接受这一需求。(如果她不接受，参与者I会给她多一点，让她接受。然而，唯一的子博弈完美均衡是参与者I让参与者II处于无差异状态。)
\end{tcolorbox}

The resulting SPE is unique. Player I demands \(x\) so that \(v\left( {1 - x}\right)  = \delta\) , which is accepted immediately by player II, so that the game terminates with player I receiving \(u\left( x\right)\) and player II receiving \(v\left( {1 - x}\right)  = \delta\) . The full equilibrium strategies are (with this \(x\) ): Player II accepts any demand \({x}^{\prime }\) so that \(0 \leq  {x}^{\prime } \leq  x\) , and rejects any \({x}^{\prime } > x\) . Any counter-demand by player II in the second round is \(y = 1\) , and player I accepts any amount in the second round. The strategies have to specify the actions of the players in the second stage in order to perform the backward induction analysis.

\begin{tcolorbox}
由此产生的子博弈完美均衡(SPE)是唯一的。参与者I提出的要求为 \(x\) ，使得 \(v\left( {1 - x}\right)  = \delta\) ，该要求会被参与者II立即接受，因此博弈结束时参与者I获得 \(u\left( x\right)\) ，参与者II获得 \(v\left( {1 - x}\right)  = \delta\) 。完整的均衡策略(针对这个 \(x\) )为:参与者II接受任何满足 \(0 \leq  {x}^{\prime } \leq  x\) 的要求 \({x}^{\prime }\) ，并拒绝任何 \({x}^{\prime } > x\) 。参与者II在第二轮提出的任何反要求为 \(y = 1\) ，并且参与者I在第二轮接受任何数额。为了进行逆向归纳法分析，这些策略必须明确参与者在第二阶段的行动。
\end{tcolorbox}

\begin{center}
\includegraphics[max width=0.3\textwidth]{images/bo_d2e5cmbef24c73b03tn0_331_342_390_477_491_0.jpg}
\end{center}
\hspace*{3em} 

Figure 11.11 Graphical solution of the two-round bargaining game in Figure 11.10, for \(\delta  = \frac{3}{4}\) . The inner curve is the set of discounted payoff pairs for the second round.

\begin{tcolorbox}
图11.11 图11.10中两轮讨价还价博弈在 \(\delta  = \frac{3}{4}\) 情况下的图形解法。内曲线是第二轮贴现收益对的集合。
\end{tcolorbox}

Player I’s demand \(x\) in the first round is easily found graphically as shown in Figure 11.11. This shows, as the outer curve, the Pareto-frontier of the bargaining set, here for the utility functions \(u\left( x\right)  = \sqrt{x}\) and \(v\left( y\right)  = y\) used earlier. Recall that for \(0 \leq  x \leq  1\) , one obtains any point \(\left( {u\left( x\right) ,v\left( {1 - x}\right) }\right)\) on that Pareto-frontier. Multiplying any such pair of real numbers with the factor \(\delta\) gives a "shrunk" version of the Pareto-frontier, shown as the curve with its lower endpoint \(\left( {\delta ,0}\right)\) and upper endpoint \(\left( {0,\delta }\right)\) . This second curve is traversed when looking at the points \(\left( {{\delta u}\left( {1 - y}\right) ,{\delta v}\left( y\right) }\right)\) when \(y\) is changed from 0 to 1, which are the discounted payoffs after agreement in the second round. In the subgame given by the second round, player I can only accept or reject, so this Ultimatum game has the upper endpoint \(\left( {0,\delta }\right)\) of the curve, denoted by \(A\) in Figure 11.11, as its payoff pair. Player I’s consideration in round one is to find a point \(B\) on the curve of first-round payoffs so that player II is indifferent between accepting and rejecting (where she will accept in the SPE). So this point \(B\) has the same vertical coordinate \(\delta\) as \(A\) , that is, \(B\) is of the form \(\left( {u\left( x\right) ,\delta }\right)\) . As shown by the arrow, \(B\) is found by moving horizontally from point \(A\) until hitting the outer curve. In the example, \(x\) is the solution to the equation \(v\left( {1 - x}\right)  = \delta\) , which is simply \(x = 1 - \delta\) because \(v\) is the identity function. The resulting utility \(u\left( x\right)\) to player \(\mathrm{I}\) is \(\sqrt{1 - \delta }\) . Figure 11.11 shows this for \(\delta  = \frac{3}{4}\) , where \(u\left( x\right)  = \frac{1}{2}\) and \(v\left( {1 - x}\right)  = \frac{3}{4}\) , that is, \(B = \left( {\frac{1}{2},\frac{3}{4}}\right)\) .

\begin{tcolorbox}
如图11.11所示，通过图形很容易找到参与者I在第一轮提出的要求 \(x\) 。图中，外曲线表示讨价还价集的帕累托前沿，这里使用的是前面提到的效用函数 \(u\left( x\right)  = \sqrt{x}\) 和 \(v\left( y\right)  = y\) 。回想一下，对于 \(0 \leq  x \leq  1\) ，可以在该帕累托前沿上得到任意点 \(\left( {u\left( x\right) ,v\left( {1 - x}\right) }\right)\) 。将任意这样的实数对乘以因子 \(\delta\) ，就得到了帕累托前沿的“缩小版”，如图中下端点为 \(\left( {\delta ,0}\right)\) 、上端点为 \(\left( {0,\delta }\right)\) 的曲线所示。当 \(y\) 从0变化到1时，观察点 \(\left( {{\delta u}\left( {1 - y}\right) ,{\delta v}\left( y\right) }\right)\) 就会遍历这条第二条曲线，这些点是第二轮达成协议后的贴现收益。在由第二轮构成的子博弈中，参与者I只能选择接受或拒绝，因此这个最后通牒博弈的收益对是曲线的上端点 \(\left( {0,\delta }\right)\) ，在图11.11中用 \(A\) 表示。参与者I在第一轮的考虑是在第一轮收益曲线上找到一个点 \(B\) ，使得参与者II在接受和拒绝之间无差异(在子博弈完美均衡中她会选择接受)。因此，这个点 \(B\) 与 \(A\) 具有相同的纵坐标 \(\delta\) ，即 \(B\) 的形式为 \(\left( {u\left( x\right) ,\delta }\right)\) 。如图中箭头所示，从点 \(A\) 水平移动直到与外曲线相交，就找到了 \(B\) 。在这个例子中， \(x\) 是方程 \(v\left( {1 - x}\right)  = \delta\) 的解，由于 \(v\) 是恒等函数，所以该解就是 \(x = 1 - \delta\) 。参与者 \(\mathrm{I}\) 由此获得的效用 \(u\left( x\right)\) 为 \(\sqrt{1 - \delta }\) 。图11.11展示了在 \(\delta  = \frac{3}{4}\) 情况下的这种情况，其中 \(u\left( x\right)  = \frac{1}{2}\) 且 \(v\left( {1 - x}\right)  = \frac{3}{4}\) ，即 \(B = \left( {\frac{1}{2},\frac{3}{4}}\right)\) 。
\end{tcolorbox}

\section*{11.9 Alternating Offers Over Several Rounds}

\begin{tcolorbox}
\section*{11.9 多轮交替出价}
\end{tcolorbox}

The bargaining game with three rounds of alternating offers is shown in Figure 11.12. Player I makes a demand \(x\) in the first round, which can be accepted or rejected by player II. If she rejects, she can make a counter-demand \(y\) , which in turn player I can accept or reject. If player I accepts, both players' payoffs are discounted, giving the expected payoff pair \(\left( {{\delta u}\left( {1 - y}\right) ,{\delta v}\left( y\right) }\right)\) . If player I rejects, then he can make in the third and final round a counter-counter-demand \(s\) , which if accepted gives the pair of twice-discounted payoffs \(\left( {{\delta }^{2}u\left( s\right) ,{\delta }^{2}v\left( {1 - s}\right) }\right)\) . If player II rejects the final demand \(s\) , both players get zero. These expected payoffs result because any rejection independently incurs a risk that, with probability \(1 - \delta\) , the game terminates with no agreement and payoff pair(0,0), and continues with probability \(\delta\) into the next round. Figure 11.12 shows directly the discounted payoffs and not the chance moves between the bargaining rounds, which would have to be included in a detailed model like the left game in Figure 11.10.

\begin{tcolorbox}
三轮交替报价的讨价还价博弈如图11.12所示。参与者I在第一轮提出一个要求 \(x\) ，参与者II可以接受或拒绝。如果她拒绝，她可以提出一个反要求 \(y\) ，而参与者I又可以接受或拒绝这个反要求。如果参与者I接受，两个参与者的收益都会被折现，得到期望效用对 \(\left( {{\delta u}\left( {1 - y}\right) ,{\delta v}\left( y\right) }\right)\) 。如果参与者I拒绝，那么他可以在第三轮也是最后一轮提出一个再反要求 \(s\) ，如果这个再反要求被接受，就得到经过两次折现的收益对 \(\left( {{\delta }^{2}u\left( s\right) ,{\delta }^{2}v\left( {1 - s}\right) }\right)\) 。如果参与者II拒绝最终要求 \(s\) ，两个参与者的收益都为零。这些期望效用之所以会出现，是因为任何一次拒绝都会独立地带来一种风险，即以概率 \(1 - \delta\) ，博弈会以未达成协议而结束，收益对为(0,0)，并以概率 \(\delta\) 进入下一轮。图11.12直接展示了折现后的收益，而没有展示讨价还价轮次之间的随机事件，在像图11.10左侧博弈那样的详细模型中，这些随机事件是必须包含在内的。
\end{tcolorbox}

\begin{center}
\includegraphics[max width=0.6\textwidth]{images/bo_d2e5cmbef24c73b03tn0_332_192_208_808_1192_0.jpg}
\end{center}
\hspace*{3em} 

Figure 11.12 Bargaining over three rounds.

\begin{tcolorbox}
图11.12 三轮讨价还价。
\end{tcolorbox}

The SPE of the three-round game is found by backward induction. The third and final round is an Ultimatum game where player I demands \(s = 1\) , which is accepted by player II. The set of twice discounted payoffs \(\left( {{\delta }^{2}u\left( s\right) ,{\delta }^{2}v\left( {1 - s}\right) }\right)\) is the inner curve on the left in Figure 11.13. Player I’s demand \(s = 1\) gives the payoff pair \(\left( {{\delta }^{2},0}\right)\) , shown as point \(A\) in the figure. In the previous round, player II maximizes her payoff \({\delta v}\left( y\right)\) by making player I indifferent between accepting and rejecting.

\begin{tcolorbox}
三轮博弈的子博弈完美均衡(SPE)可以通过逆向归纳法找到。第三轮也是最后一轮是一个最后通牒博弈，参与者I提出要求 \(s = 1\) ，参与者II接受这个要求。经过两次折现的收益集 \(\left( {{\delta }^{2}u\left( s\right) ,{\delta }^{2}v\left( {1 - s}\right) }\right)\) 是图11.13左侧的内曲线。参与者I的要求 \(s = 1\) 给出了收益对 \(\left( {{\delta }^{2},0}\right)\) ，在图中表示为点 \(A\) 。在上一轮，参与者II通过让参与者I在接受和拒绝之间无差异来最大化她的收益 \({\delta v}\left( y\right)\) 。
\end{tcolorbox}

\begin{center}
\includegraphics[max width=0.9\textwidth]{images/bo_d2e5cmbef24c73b03tn0_333_338_206_1244_562_0.jpg}
\end{center}
\hspace*{3em} 

Figure 11.13 Graphical solution of the bargaining game over three rounds (left) and four rounds (right).

\begin{tcolorbox}
图11.13 三轮(左)和四轮(右)讨价还价博弈的图形解法。
\end{tcolorbox}

Hence, \(y\) is the solution to the equation \({\delta u}\left( {1 - y}\right)  = {\delta }^{2}\) . This is point \(B\) in the figure, which is the point on the curve of second-round payoff pairs \(\left( {{\delta u}\left( {1 - y}\right) ,{\delta v}\left( y\right) }\right)\) with the same horizontal coordinate as \(A\) . With \(y\) determined like that, player \({\mathrm{I}}^{\prime }\mathrm{s}\) demand \(x\) in the first round is the solution to the equation \(v\left( {1 - x}\right)  = {\delta v}\left( y\right)\) because this makes player II indifferent between accepting and rejecting in the first round. The corresponding point \(C = \left( {u\left( x\right) ,v\left( {1 - x}\right) }\right)\) on the set of first-round payoff pairs has the same vertical coordinate as \(B\) . On the equilibrium path, player I demands \(x\) in the first round, which is immediately accepted by player II. The complete equilibrium strategies are as before, where a player accepts any demand up to and including the demand in the SPE, and otherwise rejects and makes the described counter- or counter-counter-demand in the next round.

\begin{tcolorbox}
因此， \(y\) 是方程 \({\delta u}\left( {1 - y}\right)  = {\delta }^{2}\) 的解。这在图中是点 \(B\) ，它是第二轮收益对曲线 \(\left( {{\delta u}\left( {1 - y}\right) ,{\delta v}\left( y\right) }\right)\) 上与 \(A\) 具有相同横坐标的点。在这样确定了 \(y\) 之后，参与者 \({\mathrm{I}}^{\prime }\mathrm{s}\) 在第一轮的要求 \(x\) 是方程 \(v\left( {1 - x}\right)  = {\delta v}\left( y\right)\) 的解，因为这使得参与者II在第一轮接受和拒绝之间无差异。第一轮收益对集合上对应的点 \(C = \left( {u\left( x\right) ,v\left( {1 - x}\right) }\right)\) 与 \(B\) 具有相同的纵坐标。在均衡路径上，参与者I在第一轮提出要求 \(x\) ，参与者II立即接受这个要求。完整的均衡策略和之前一样，即参与者会接受直至并包括子博弈完美均衡中的要求在内的任何要求，否则就拒绝并在下一轮提出上述的反要求或再反要求。
\end{tcolorbox}

In our example with utility functions \(u\left( x\right)  = \sqrt{x}\) and \(v\left( y\right)  = y\) , the SPE of the game with three rounds is found as follows (according to the left picture in Figure 11.13). The vertical arrow from point \(A\) to point \(B\) gives the equation \({\delta u}\left( {1 - y}\right)  = {\delta }^{2}\) or \(u\left( {1 - y}\right)  = \sqrt{1 - y} = \delta\) , that is, \(y = 1 - {\delta }^{2}\) . The horizontal arrow from \(B\) to \(C\) represents the equation \(v\left( {1 - x}\right)  = {\delta v}\left( y\right)\) , that is, \(1 - x = \delta \left( {1 - {\delta }^{2}}\right)  = \delta  - {\delta }^{3}\) , which has the solution \(x = 1 - \delta  + {\delta }^{3}\) . This is player I’s demand in the first round in the SPE; his payoff is \(u\left( x\right)  = \sqrt{1 - \delta  + {\delta }^{3}}\) . For \(\delta  = \frac{3}{4}\) , as in Figure 11.13, we have \(y = \frac{7}{16}\) and \(x = \frac{43}{64}\) with \(u\left( x\right)  = \sqrt{43}/8\) and \(v\left( {1 - x}\right)  = \frac{21}{64}\) . These are the coordinates of point \(C\) , approximately(0.82,0.33).

\begin{tcolorbox}
在我们具有效用函数 \(u\left( x\right)  = \sqrt{x}\) 和 \(v\left( y\right)  = y\) 的例子中，根据图11.13中的左图，该三轮博弈的子博弈完美均衡(SPE)求解如下。从点 \(A\) 到点 \(B\) 的垂直箭头给出方程 \({\delta u}\left( {1 - y}\right)  = {\delta }^{2}\) 或 \(u\left( {1 - y}\right)  = \sqrt{1 - y} = \delta\) ，即 \(y = 1 - {\delta }^{2}\) 。从 \(B\) 到 \(C\) 的水平箭头表示方程 \(v\left( {1 - x}\right)  = {\delta v}\left( y\right)\) ，即 \(1 - x = \delta \left( {1 - {\delta }^{2}}\right)  = \delta  - {\delta }^{3}\) ，其解为 \(x = 1 - \delta  + {\delta }^{3}\) 。这是子博弈完美均衡中玩家I在第一轮的需求；他的收益是 \(u\left( x\right)  = \sqrt{1 - \delta  + {\delta }^{3}}\) 。对于 \(\delta  = \frac{3}{4}\) ，如图11.13所示，我们有 \(y = \frac{7}{16}\) 和 \(x = \frac{43}{64}\) ，其中 \(u\left( x\right)  = \sqrt{43}/8\) 且 \(v\left( {1 - x}\right)  = \frac{21}{64}\) 。这些是点 \(C\) 的坐标，约为(0.82, 0.33)。
\end{tcolorbox}

The game with four bargaining rounds has the same structure as the game with three rounds, except that if player II rejects player I’s demand \(s\) in round three, she can make a final demand \(t\) of her own in the fourth and final round, which player I can only accept or reject. If player I accepts, the players receive the three-times discounted payoffs \(\left( {{\delta }^{3}u\left( {1 - t}\right) ,{\delta }^{3}v\left( t\right) }\right)\) . The corresponding SPE is solved graphically as shown on the right in Figure 11.13. The last round is an Ultimatum game where player II can demand the whole pie with \(t = 1\) , which determines the payoff pair \(\left( {0,{\delta }^{3}}\right)\) , which is point \(A\) in the figure. In the previous round, the third round of the game, player \(\mathrm{I}\) demands \(s\) so that player \(\mathrm{{II}}\) is indifferent between accepting and rejecting, which gives point \(B\) , which is the pair of discounted payoffs \(\left( {{\delta }^{2}u\left( s\right) ,{\delta }^{2}v\left( {1 - s}\right) }\right)\) so that \({\delta }^{2}v\left( {1 - s}\right)  = {\delta }^{3}\) . This determines \(s\) . In the round before that, player II makes player I indifferent between accepting and rejecting with the demand \(y\) so that \({\delta u}\left( {1 - y}\right)  = {\delta }^{2}u\left( s\right)\) , shown as point \(C\) . This equation determines \(y\) . In the first round, player I demands \(x\) so as to make player II indifferent, according to \(v\left( {1 - x}\right)  = {\delta v}\left( y\right)\) , which defines point \(D = \left( {u\left( x\right) ,v\left( {1 - x}\right) }\right)\) . This is also the equilibrium outcome because that demand \(x\) is accepted by player II.

\begin{tcolorbox}
四轮讨价还价博弈与三轮博弈结构相同，不同之处在于，如果玩家II在第三轮拒绝了玩家I的需求 \(s\) ，她可以在第四轮(即最后一轮)提出自己的最终需求 \(t\) ，玩家I只能接受或拒绝。如果玩家I接受，玩家将获得三倍贴现后的收益 \(\left( {{\delta }^{3}u\left( {1 - t}\right) ,{\delta }^{3}v\left( t\right) }\right)\) 。相应的子博弈完美均衡如图11.13右侧所示通过图形方式求解。最后一轮是一个最后通牒博弈(Ultimatum game)，玩家II可以以 \(t = 1\) 要求得到整个“馅饼”，这决定了收益对 \(\left( {0,{\delta }^{3}}\right)\) ，即图中的点 \(A\) 。在前一轮，即博弈的第三轮，玩家 \(\mathrm{I}\) 要求 \(s\) ，使得玩家 \(\mathrm{{II}}\) 在接受和拒绝之间无差异，这给出点 \(B\) ，即贴现收益对 \(\left( {{\delta }^{2}u\left( s\right) ,{\delta }^{2}v\left( {1 - s}\right) }\right)\) ，使得 \({\delta }^{2}v\left( {1 - s}\right)  = {\delta }^{3}\) 。这决定了 \(s\) 。再前一轮，玩家II以需求 \(y\) 使玩家I在接受和拒绝之间无差异，使得 \({\delta u}\left( {1 - y}\right)  = {\delta }^{2}u\left( s\right)\) ，如图中点 \(C\) 所示。该方程决定了 \(y\) 。在第一轮，玩家I根据 \(v\left( {1 - x}\right)  = {\delta v}\left( y\right)\) 要求 \(x\) 以使玩家II无差异，这定义了点 \(D = \left( {u\left( x\right) ,v\left( {1 - x}\right) }\right)\) 。这也是均衡结果，因为该需求 \(x\) 会被玩家II接受。
\end{tcolorbox}

For \(u\left( x\right)  = \sqrt{x}\) and \(v\left( y\right)  = y\) , the right picture in Figure 11.13 gives the following SPE. The arrow from \(A\) to \(B\) gives the equation \({\delta }^{2}v\left( {1 - s}\right)  = {\delta }^{3}\) or \(v\left( {1 - s}\right)  = \delta\) , which determines player I’s demand \(s = 1 - \delta\) in round three. The arrow from \(B\) to \(C\) gives the equation \({\delta u}\left( {1 - y}\right)  = {\delta }^{2}u\left( s\right)\) or \(u\left( {1 - y}\right)  = {\delta u}\left( s\right)\) , that is, \(\sqrt{1 - y} = \delta \sqrt{1 - \delta }\) or equivalently \(1 - y = {\delta }^{2}\left( {1 - \delta }\right)  = {\delta }^{2} - {\delta }^{3}\) , which has the solution \(y = 1 - {\delta }^{2} + {\delta }^{3}\) that determines player II’s demand \(y\) in the second round. Finally, the arrow from \(C\) to \(D\) gives the equation \(v\left( {1 - x}\right)  = {\delta v}\left( y\right)\) or \(1 - x = \delta \left( {1 - {\delta }^{2} + {\delta }^{3}}\right)\) , which gives player I’s demand \(x\) in the first round as \(x = 1 - \delta  + {\delta }^{3} - {\delta }^{4}\) .

\begin{tcolorbox}
对于 \(u\left( x\right)  = \sqrt{x}\) 和 \(v\left( y\right)  = y\) ，图11.13中的右图给出了以下子博弈完美均衡(SPE)。从 \(A\) 到 \(B\) 的箭头给出方程 \({\delta }^{2}v\left( {1 - s}\right)  = {\delta }^{3}\) 或 \(v\left( {1 - s}\right)  = \delta\) ，该方程确定了玩家I在第三轮的需求 \(s = 1 - \delta\) 。从 \(B\) 到 \(C\) 的箭头给出方程 \({\delta u}\left( {1 - y}\right)  = {\delta }^{2}u\left( s\right)\) 或 \(u\left( {1 - y}\right)  = {\delta u}\left( s\right)\) ，即 \(\sqrt{1 - y} = \delta \sqrt{1 - \delta }\) 或等价地 \(1 - y = {\delta }^{2}\left( {1 - \delta }\right)  = {\delta }^{2} - {\delta }^{3}\) ，其解为 \(y = 1 - {\delta }^{2} + {\delta }^{3}\) ，该解确定了玩家II在第二轮的需求 \(y\) 。最后，从 \(C\) 到 \(D\) 的箭头给出方程 \(v\left( {1 - x}\right)  = {\delta v}\left( y\right)\) 或 \(1 - x = \delta \left( {1 - {\delta }^{2} + {\delta }^{3}}\right)\) ，该方程得出玩家I在第一轮的需求 \(x\) 为 \(x = 1 - \delta  + {\delta }^{3} - {\delta }^{4}\) 。
\end{tcolorbox}

Note that in the game of bargaining over four rounds, the arrow from \(A\) to \(B\) defines the same equation as in the two-round game in Figure 11.11; essentially, the fourth round in the four-round game describes the same situation, and gives the same picture, as the second round in the two-round game, except that all payoffs are shrunk by the factor \({\delta }^{2}\) .

\begin{tcolorbox}
注意，在四轮讨价还价博弈中，从 \(A\) 到 \(B\) 的箭头所定义的方程与图11.11中的两轮博弈相同；本质上，四轮博弈的第四轮所描述的情况与两轮博弈的第二轮相同，且图形也相同，只是所有的收益都缩小了 \({\delta }^{2}\) 倍。
\end{tcolorbox}

The bargaining game with alternating offers can be defined for any number of rounds. In each odd-numbered round (the first, third, fifth round, etc.), player I makes a demand, which can be accepted or rejected by player II. After rejection, the game progresses with probability \(\delta\) into the next round, and otherwise terminates with payoffs(0,0). Player II can make a demand in each even-numbered round. If the total number of rounds is even, then player II can make the last demand, which is the whole pie, which is advantageous for player II. If the total number of rounds is odd, then player I can make the last demand, and in turn demand the whole pie for himself.

\begin{tcolorbox}
交替出价的讨价还价博弈可以定义为任意轮数。在每一个奇数轮(第一轮、第三轮、第五轮等)，玩家I提出一个需求，玩家II可以接受或拒绝。拒绝后，博弈以概率 \(\delta\) 进入下一轮，否则以收益(0,0)结束。玩家II可以在每一个偶数轮提出需求。如果总轮数为偶数，那么玩家II可以提出最后一个需求，即整个“馅饼”，这对玩家II有利。如果总轮数为奇数，那么玩家I可以提出最后一个需求，进而为自己要求整个“馅饼”。
\end{tcolorbox}

The graphical solutions in Figures 11.11 and 11.13 can be extended to any number of rounds, with payoffs in the \(k\) th round given by \(\left( {{\delta }^{k - 1}u\left( z\right) ,{\delta }^{k - 1}v\left( {1 - z}\right) }\right)\) for some \(z \in  \left\lbrack  {0,1}\right\rbrack\) . These payoff pairs define a curve that is the original Pareto-frontier of the bargaining set shrunk by the factor \({\delta }^{k - 1}\) . The examples for a total number of two, three, or four rounds show that the resulting SPE defines a point \(\left( {u\left( x\right) ,v\left( {1 - x}\right) }\right)\) . This is the pair of payoffs to the two players because they agree in round one.

\begin{tcolorbox}
图11.11和图11.13中的图形解法可以扩展到任意轮数，第 \(k\) 轮的收益由 \(\left( {{\delta }^{k - 1}u\left( z\right) ,{\delta }^{k - 1}v\left( {1 - z}\right) }\right)\) 给出，其中 \(z \in  \left\lbrack  {0,1}\right\rbrack\) 为某个值。这些收益对定义了一条曲线，该曲线是讨价还价集的原始帕累托前沿缩小了 \({\delta }^{k - 1}\) 倍后的曲线。两轮、三轮或四轮的例子表明，得到的子博弈完美均衡(SPE)定义了一个点 \(\left( {u\left( x\right) ,v\left( {1 - x}\right) }\right)\) 。这是因为两个玩家在第一轮就达成了一致，所以这是两个玩家的收益对。
\end{tcolorbox}

In general, if the total number of rounds is odd, player I is always better off than if the total number of rounds is even, which favours player II who can make the last demand. However, with an increasing total number of rounds, the resulting points \(\left( {u\left( x\right) ,v\left( {1 - x}\right) }\right)\) move towards each other. We should expect that with a very large number of rounds, it matters less and less who can make the last demand (because the future discounted payoffs from which the backward induction starts are very small), and expect some convergence. However, we will not study the "limit" of these finite games, but instead consider a new concept, a game with an infinite number of rounds.

\begin{tcolorbox}
一般来说，如果总轮数为奇数，玩家I的情况总是比总轮数为偶数时更好，因为总轮数为偶数有利于能提出最后需求的玩家II。然而，随着总轮数的增加，所得的点 \(\left( {u\left( x\right) ,v\left( {1 - x}\right) }\right)\) 会相互靠近。我们应该预期，当轮数非常多时，谁能提出最后需求就变得越来越不重要了(因为逆向归纳法所依据的未来贴现期望效用非常小)，并且会出现某种收敛。不过，我们不会研究这些有限博弈的“极限”，而是考虑一个新概念，即具有无限轮数的博弈。
\end{tcolorbox}

\section*{11.10 Stationary Strategies}

\begin{tcolorbox}
\section*{11.10 平稳策略}
\end{tcolorbox}

We now analyze the bargaining game with alternating offers with an infinite number of rounds, that is, the game goes on forever. Formally, this is a game \(\Gamma\) with perfect information represented by an infinite tree. It can be defined recursively as shown in Figure 11.14. In the first and second round, player I and player II each make a demand \(x\) and \(y\) , respectively, from the interval \(\left\lbrack  {0,1}\right\rbrack\) , which the other player can accept or reject. After rejection, a chance move terminates the game with probability \(1 - \delta\) where both players receive payoff zero. With probability \(\delta\) , the game continues into the next round where the player who has just rejected the offer can make her counter-demand. With the beginning of the third round, this is again player I, so that this defines the same game \(\Gamma\) that we started with in round one. Correspondingly, \(\Gamma\) is simply appended as the beginning of the subgame that starts after the second chance move that follows player I's rejection of player II's counter-demand \(y\) . (Note that infinitely many copies of \(\Gamma\) are substituted at each node that results from any \(x\) followed by any \(y\) .)

\begin{tcolorbox}
我们现在分析具有无限轮数的交替报价讨价还价博弈，也就是说，这个博弈会永远进行下去。形式上，这是一个由无限树表示的具有完美信息的博弈 \(\Gamma\) 。它可以如图11.14所示进行递归定义。在第一轮和第二轮，玩家I和玩家II分别从区间 \(\left\lbrack  {0,1}\right\rbrack\) 中提出需求 \(x\) 和 \(y\) ，另一名玩家可以接受或拒绝。拒绝后，有一个随机事件会以概率 \(1 - \delta\) 终止博弈，此时两名玩家的收益均为零。以概率 \(\delta\) ，博弈会进入下一轮，此时刚刚拒绝报价的玩家可以提出反需求。从第三轮开始，轮到玩家I，这样就定义了与第一轮开始时相同的博弈 \(\Gamma\) 。相应地， \(\Gamma\) 只是作为子博弈的起始部分被添加进来，该子博弈从玩家I拒绝玩家II的反需求 \(y\) 之后的第二次随机事件开始。(注意，在任何 \(x\) 之后接着任何 \(y\) 所产生的每个节点上，都会代入无限多个 \(\Gamma\) 的副本。)
\end{tcolorbox}

\begin{center}
\includegraphics[max width=0.4\textwidth]{images/bo_d2e5cmbef24c73b03tn0_335_330_1127_589_1022_0.jpg}
\end{center}
\hspace*{3em} 

Figure 11.14 The bargaining game \(\Gamma\) with an infinite number of rounds, which repeats itself after any demand \(x\) of player I in round one and any demand \(y\) of player II in round two after each player has rejected the other's demand. The next round is reached with probability \(\delta\) each time, and otherwise terminates with payoffs(0,0).

\begin{tcolorbox}
图11.14 具有无限轮数的讨价还价博弈 \(\Gamma\) ，在第一轮玩家I提出任何需求 \(x\) 以及第二轮玩家II提出任何需求 \(y\) 且双方都拒绝了对方的需求之后，该博弈会自我重复。每次以概率 \(\delta\) 进入下一轮，否则以收益(0,0)终止。
\end{tcolorbox}

An SPE of this infinite game cannot be defined by backward induction because the game goes on forever, so that there is no final move from which one could start the induction. However, the game is defined recursively and is the same game after two rounds. Hence, it is possible to define strategies of the players that repeat themselves every two rounds. Such strategies are called stationary and mean that player I always demands \(x\) whenever he can make a demand (in each odd-numbered round, starting with the first round), and player II always demands \(y\) (in each even-numbered round, starting with the second round). The SPE condition states that these strategies should define an equilibrium in any subgame of the game. Any such subgame starts either with a demand by one of the players, with the other player's response, or with a chance move.

\begin{tcolorbox}
这个无限博弈的子博弈完美均衡(SPE)不能用逆向归纳法来定义，因为博弈会永远进行下去，所以不存在可以作为归纳起点的最终行动。然而，该博弈是递归定义的，并且每两轮之后都是相同的博弈。因此，有可能定义玩家的策略，使其每两轮重复一次。这样的策略被称为平稳策略，意味着玩家I在他能提出需求时(从第一轮开始的每个奇数轮次)总是要求 \(x\) ，而玩家II总是要求 \(y\) (从第二轮开始的每个偶数轮次)。子博弈完美均衡条件表明，这些策略应该在该博弈的任何子博弈中定义一个均衡。任何这样的子博弈要么从一名玩家的需求开始，要么从另一名玩家的回应开始，要么从一个随机事件开始。
\end{tcolorbox}

In addition to the players's demands in each round, the SPE should also specify whether the other player accepts or rejects the offer. After rejection, the resulting expected payoffs automatically shrink by a factor of \(\delta\) . Consequently, a player’s demand should not be so high as to result in rejection, but should instead make it optimal for the other player to accept. The demand itself is maximal subject to this condition, that is, it is, as before, chosen so as to make the other player indifferent between accepting and rejecting, and the other player accepts in the SPE.

\begin{tcolorbox}
除了玩家在每一轮的需求之外，子博弈完美均衡还应该规定另一名玩家是接受还是拒绝报价。拒绝后，由此产生的期望效用会自动缩小 \(\delta\) 倍。因此，一名玩家的需求不应过高以至于导致被拒绝，而应该使另一名玩家接受报价是最优选择。需求本身在满足这个条件的前提下是最大的，也就是说，和之前一样，需求的选择要使另一名玩家在接受和拒绝之间无差异，并且在子博弈完美均衡中另一名玩家会接受。
\end{tcolorbox}

We now show that an SPE exists, and how to find it, using a graph similar to the left picture in Figure 11.13. Like there, we draw the three curves of payoff pairs \(\left( {u\left( x\right) ,v\left( {1 - x}\right) }\right)\) that result from player I’s demand \(x\) in round one if this demand is accepted by player II, the discounted payoffs \(\delta \left( {u\left( {1 - y}\right) ,{\delta v}\left( y\right) }\right)\) if player I accepts player II’s demand \(y\) in round two, and the third-round twice-discounted payoffs \({\delta }^{2}\left( {u\left( s\right) ,{\delta }^{2}v\left( {1 - s}\right) }\right)\) if player II accepts player I’s demand \(s\) in round three. The demands \(x\) and \(y\) are chosen so as to make the other player indifferent between acceptance and rejection. For the demand \(s\) in round three, we want to fulfill the requirement of stationary strategies, that is, \(s = x\) .

\begin{tcolorbox}
现在我们将使用类似于图11.13左图的一个图(图/graph)来证明子博弈完美均衡(SPE)的存在性，并说明如何找到它。与图中一样，我们绘制收益(payoff)对 \(\left( {u\left( x\right) ,v\left( {1 - x}\right) }\right)\) 的三条曲线，它们分别是:如果参与者II接受参与者I在第一轮提出的需求 \(x\) 所产生的收益对；如果参与者I在第二轮接受参与者II的需求 \(y\) 所产生的贴现收益 \(\delta \left( {u\left( {1 - y}\right) ,{\delta v}\left( y\right) }\right)\) ；如果参与者II在第三轮接受参与者I的需求 \(s\) 所产生的经过两次贴现的第三轮收益 \({\delta }^{2}\left( {u\left( s\right) ,{\delta }^{2}v\left( {1 - s}\right) }\right)\) 。需求 \(x\) 和 \(y\) 的选择要使得另一个参与者在接受和拒绝之间无差异。对于第三轮的需求 \(s\) ，我们希望满足平稳策略(stationary strategy)的要求，即 \(s = x\) 。
\end{tcolorbox}

Figure 11.15 illustrates how to fulfill these requirements. Suppose that player I's demand \(s\) in round three defines the point \(A = \left( {{\delta }^{2}u\left( s\right) ,{\delta }^{2}v\left( {1 - s}\right) }\right)\) , relatively high up on the inner curve as shown in Figure 11.15(a). In the previous round two, point \(B = \left( {{\delta u}\left( {1 - y}\right) ,{\delta v}\left( y\right) }\right)\) has the same horizontal coordinate as point \(A\) because player II tries to make player I indifferent with her demand \(y\) , according to the equation

\begin{tcolorbox}
图11.15说明了如何满足这些要求。假设参与者I在第三轮的需求 \(s\) 确定了点 \(A = \left( {{\delta }^{2}u\left( s\right) ,{\delta }^{2}v\left( {1 - s}\right) }\right)\) ，如图11.15(a)所示，该点位于内曲线上相对较高的位置。在前一轮(第二轮)，点 \(B = \left( {{\delta u}\left( {1 - y}\right) ,{\delta v}\left( y\right) }\right)\) 与点 \(A\) 具有相同的横坐标，因为根据方程，参与者II试图通过她的需求 \(y\) 使参与者I无差异。
\end{tcolorbox}

\[
{\delta u}\left( {1 - y}\right)  = {\delta }^{2}u\left( s\right)  \tag{11.4}
\]

\begin{center}
\includegraphics[max width=0.9\textwidth]{images/bo_d2e5cmbef24c73b03tn0_337_334_203_1215_1134_0.jpg}
\end{center}
\hspace*{3em} 

In the first round, point \(C\) is found similarly, as the point \(\left( {u\left( x\right) ,v\left( {1 - x}\right) }\right)\) that has the same vertical coordinate as \(B\) , because player I demands \(x\) so as to make player II indifferent between accepting and rejecting, according to

\begin{tcolorbox}
在第一轮，类似地找到点 \(C\) ，即与 \(B\) 具有相同纵坐标的点 \(\left( {u\left( x\right) ,v\left( {1 - x}\right) }\right)\) ，因为根据方程，参与者I提出需求 \(x\) 是为了使参与者II在接受和拒绝之间无差异。
\end{tcolorbox}

\[
v\left( {1 - x}\right)  = {\delta v}\left( y\right) . \tag{11.5}
\]

The demand \(x\) is the same as demand \(s\) , as required in stationary strategies, if point \(C\) defines the same relative position on the outer curve as \(A\) on the inner curve. In Figure 11.15, we have shown point \(D\) as \({\delta }^{2}C\) , that is, point \(C\) shrunk by a factor of \({\delta }^{2}\) , which is the expected payoff pair two rounds later. In Figure 11.15(a), \(A\) and \(D\) do not coincide, because \(s < x\) , that is, the third-round demand \(s\) was too small. Figure 11.15(b) shows a point \(A\) on the inner curve corresponding to a demand \(s\) that is too favourable for player \(\mathrm{I}\) , where the arrows from \(A\) to \(B\) to \(C\) according to equations (11.4) and (11.5) give a first-round demand \(x\) with \(x < s\) .

\begin{tcolorbox}
如果点 \(C\) 在外部曲线上所定义的相对位置与 \(A\) 在内曲线上所定义的相对位置相同，那么根据平稳策略的要求，需求 \(x\) 与需求 \(s\) 相同。在图11.15中，我们将点 \(D\) 表示为 \({\delta }^{2}C\) ，即点 \(C\) 缩小了 \({\delta }^{2}\) 倍，这是两轮后的期望效用(expected payoff)对。在图11.15(a)中， \(A\) 和 \(D\) 不重合，因为 \(s < x\) ，即第三轮的需求 \(s\) 太小。图11.15(b)显示了内曲线上对应于需求 \(s\) 的点 \(A\) ，该需求对参与者 \(\mathrm{I}\) 过于有利，根据方程(11.4)和(11.5)从 \(A\) 到 \(B\) 再到 \(C\) 的箭头给出了第一轮需求 \(x\) ，且 \(x < s\) 。
\end{tcolorbox}

In general, the pictures demonstrate that when starting from \(s = 0\) , these equations give a first-round demand \(x\) where \(x > s\) , and when starting from \(s = 1\) , a first-round demand \(x\) where \(x < s\) . Moreover, when starting from any \(s\) in \(\left\lbrack  {0,1}\right\rbrack\) , the resulting demand \(x\) is a continuous function of \(s\) . Because the continuous function \(x - s\) is positive for \(s = 0\) and negative for \(s = 1\) , it is zero for some intermediate value \(s\) , where \(s = x\) , which defines a stationary SPE. (This SPE is also unique, see Binmore, 1992, theorem 5.8.4.)

\begin{tcolorbox}
一般来说，这些图表明，从 \(s = 0\) 开始，这些方程给出的第一轮需求 \(x\) 满足 \(x > s\) ；从 \(s = 1\) 开始，给出的第一轮需求 \(x\) 满足 \(x < s\) 。此外，从 \(\left\lbrack  {0,1}\right\rbrack\) 中的任何 \(s\) 开始，得到的需求 \(x\) 是 \(s\) 的连续函数。因为连续函数 \(x - s\) 在 \(s = 0\) 处为正，在 \(s = 1\) 处为负，所以对于某个中间值 \(s\) 它为零，此时 \(s = x\) ，这就定义了一个平稳的子博弈完美均衡(SPE)。(这个子博弈完美均衡也是唯一的，见宾默尔(Binmore)，1992年，定理5.8.4。)
\end{tcolorbox}

In short, the stationary SPE of the infinite bargaining game is a solution to the equations (11.4),(11.5), and \(s = x\) . These are equivalent to the two, nicely symmetric equations

\begin{tcolorbox}
简而言之，无限期讨价还价博弈的平稳子博弈完美均衡是方程(11.4)、(11.5)和 \(s = x\) 的解。这些方程等价于两个形式优美的对称(symmetric)方程。
\end{tcolorbox}

\[
u\left( {1 - y}\right)  = {\delta u}\left( x\right) ,\;v\left( {1 - x}\right)  = {\delta v}\left( y\right) . \tag{11.6}
\]

The first of these equations expresses that player II makes player I indifferent in every even-numbered round by her demand \(y\) . The second equation in (11.6) states that player I makes player II indifferent by his demand \(x\) in every odd-numbered round. In this stationary SPE, player I’s demand \(x\) in round one is immediately accepted by player II. As before, the full strategies specify acceptance of any lower and rejection of any higher demand. The actual demand of the other player is always accepted, so the later rounds are never reached.

\begin{tcolorbox}
这些方程中的第一个表明，玩家II在每一个偶数轮次中，通过她的要求 \(y\) 让玩家I无差异。式(11.6)中的第二个方程表明，玩家I在每一个奇数轮次中，通过他的要求 \(x\) 让玩家II无差异。在这个平稳子博弈完美均衡(SPE)中，玩家I在第一轮的要求 \(x\) 会被玩家II立即接受。和之前一样，完整的策略规定接受任何更低的要求，拒绝任何更高的要求。另一个玩家提出的实际要求总是会被接受，因此永远不会进入后续轮次。
\end{tcolorbox}

In the example with utility functions \(u\left( x\right)  = \sqrt{x}\) and \(v\left( y\right)  = y\) , the stationary strategies in an SPE are found as follows. They are given by the demands \(x\) and \(y\) that solve (11.6), that is,

\begin{tcolorbox}
在具有效用函数 \(u\left( x\right)  = \sqrt{x}\) 和 \(v\left( y\right)  = y\) 的例子中，子博弈完美均衡(SPE)中的平稳策略可按如下方式找到。它们由满足方程(11.6)的需求 \(x\) 和 \(y\) 给出，即
\end{tcolorbox}

\[
\sqrt{1 - y} = \delta \sqrt{x},\;1 - x = {\delta y}.
\]

The first equation is equivalent to \(y = 1 - {\delta }^{2}x\) , which substituted into the second equation gives \(1 - x = \delta  - {\delta }^{3}x\) , or \(1 - \delta  = x\left( {1 - {\delta }^{3}}\right)\) . Because \(1 - {\delta }^{3} = \left( {1 - \delta }\right) \left( {1 + \delta  + {\delta }^{2}}\right)\) , this gives the solutions for \(x\) and \(y\) as

\begin{tcolorbox}
第一个方程等价于 \(y = 1 - {\delta }^{2}x\) ，将其代入第二个方程得到 \(1 - x = \delta  - {\delta }^{3}x\) ，或 \(1 - \delta  = x\left( {1 - {\delta }^{3}}\right)\) 。因为 \(1 - {\delta }^{3} = \left( {1 - \delta }\right) \left( {1 + \delta  + {\delta }^{2}}\right)\) ，这给出了 \(x\) 和 \(y\) 的解为
\end{tcolorbox}

\[
x = \frac{1 - \delta }{1 - {\delta }^{3}} = \frac{1}{1 + \delta  + {\delta }^{2}},\;y = \frac{1 + \delta }{1 + \delta  + {\delta }^{2}}. \tag{11.7}
\]

The pie is split in the first round into \(x = \frac{1}{1 + \delta  + {\delta }^{2}}\) for player I and \(1 - x = \frac{\delta  + {\delta }^{2}}{1 + \delta  + {\delta }^{2}}\) for player II. The share \(1 - x\) of the pie for player II is lower than her demand \(y\) in the second round, as illustrated by Figure 11.15(c), which shows that the relative position of point \(A\) on the outer curve is higher than that of point \(B\) on the middle curve. This holds because the second-round payoffs are discounted.

\begin{tcolorbox}
在第一轮中，“馅饼”被分割为参与者 I 获得 \(x = \frac{1}{1 + \delta  + {\delta }^{2}}\) ，参与者 II 获得 \(1 - x = \frac{\delta  + {\delta }^{2}}{1 + \delta  + {\delta }^{2}}\) 。参与者 II 获得的“馅饼”份额 \(1 - x\) 低于她在第二轮中的需求 \(y\) ，如图 11.15(c) 所示，该图显示外曲线上点 \(A\) 的相对位置高于中间曲线上点 \(B\) 的相对位置。这是因为第二轮的收益是经过贴现的。
\end{tcolorbox}

Does it matter for our analysis that player I makes the first demand? Of course, if player II made the first demand, one could swap the players and analyze the game with the described method. However, the game where player II makes the first demand, and the corresponding stationary solution, is already analyzed as part of the game \(\Gamma\) where player I moves first. The reason is that the game with player II moving first is the subgame of \(\Gamma\) in Figure 11.14 after the first chance move, which repeats itself as in \(\Gamma\) . All we have to do is to ignore player I’s first demand, and the first chance move and the corresponding discount factor \(\delta\) . That is, we draw the same picture of curves with discounted payoff pairs as in Figure 11.15, except that we omit the outer curve and scale the picture by multiplication by \(\frac{1}{\delta }\) .

\begin{tcolorbox}
参与者 I 先提出需求对我们的分析有影响吗？当然，如果参与者 II 先提出需求，我们可以交换参与者的角色，并用上述方法分析该博弈。然而，参与者 II 先提出需求的博弈以及相应的平稳解，已经作为参与者 I 先行动的博弈 \(\Gamma\) 的一部分进行了分析。原因是参与者 II 先行动的博弈是图 11.14 中参与者 I 先行动的博弈 \(\Gamma\) 在第一次机会行动之后的子博弈，它会像 \(\Gamma\) 中那样重复出现。我们所要做的就是忽略参与者 I 的第一次需求、第一次机会行动以及相应的贴现因子 \(\delta\) 。也就是说，我们绘制与图 11.15 中类似的贴现收益对曲线，只是省略外曲线，并将图形乘以 \(\frac{1}{\delta }\) 进行缩放。
\end{tcolorbox}

For the alternating-offers game with player II moving first, the resulting equations are given as before by (11.6). They have the same stationary equilibrium solution. This solution starts with player II making the first demand \(y\) , followed by player I’s counter-demand \(x\) , and so on in subsequent rounds. Because in this SPE the first demand is accepted, player II gets payoff \(v\left( y\right)\) and player I gets payoff \(u\left( {1 - y}\right)\) , which is \({\delta u}\left( x\right)\) by (11.6). Hence, all that changes when player II makes the first demand is that player II's payoff is no longer discounted by multiplication by \(\delta\) , and player I’s payoff is discounted instead.

\begin{tcolorbox}
对于参与者 II 先行动的交替报价博弈，得到的方程如前文所示由(11.6)给出。它们具有相同的平稳均衡解。这个解从参与者 II 提出第一次需求 \(y\) 开始，接着是参与者 I 的反需求 \(x\) ，依此类推后续轮次。因为在这个子博弈完美均衡中第一次需求被接受，参与者 II 获得收益 \(v\left( y\right)\) ，参与者 I 获得收益 \(u\left( {1 - y}\right)\) ，根据(11.6)，这是 \({\delta u}\left( x\right)\) 。因此，当参与者 II 先提出需求时，唯一的变化是参与者 II 的收益不再乘以 \(\delta\) 进行贴现，而是参与者 I 的收益进行贴现。
\end{tcolorbox}

\section*{11.11 The Nash Bargaining Solution Via Alternating Offers}

\begin{tcolorbox}
\section*{11.11 通过交替报价得到的纳什讨价还价解}
\end{tcolorbox}

In this final section, we show that for very patient players, when the discount factor \(\delta\) approaches 1, the stationary SPE converges to the Nash bargaining solution. This is true for the example that we have used throughout, with utility functions \(u\left( x\right)  = \sqrt{x}\) and \(v\left( y\right)  = y\) , where the split in the Nash bargaining solution is \(\left( {x,1 - x}\right)  = \left( {\frac{1}{3},\frac{2}{3}}\right)\) . This is also the limit of the stationary solution in (11.7) when \(\delta  \rightarrow  1\) (note the convenient cancellation by \(1 - \delta\) in the term \(\frac{1 - \delta }{1 - {\delta }^{3}}\) that represents \(x\) ; this term has an undefined limit \(\frac{0}{0}\) when \(\delta  = 1\) , which without the cancellation requires using l'Hôpital's rule).

\begin{tcolorbox}
在这最后一节中，我们表明对于非常有耐心的参与者，当贴现因子 \(\delta\) 趋近于 1 时，平稳子博弈完美均衡收敛于纳什讨价还价解。对于我们一直使用的具有效用函数 \(u\left( x\right)  = \sqrt{x}\) 和 \(v\left( y\right)  = y\) 的例子来说，这是成立的，其中纳什讨价还价解中的分割比例是 \(\left( {x,1 - x}\right)  = \left( {\frac{1}{3},\frac{2}{3}}\right)\) 。这也是(11.7)中平稳解在 \(\delta  \rightarrow  1\) 时的极限(注意表示 \(x\) 的项 \(\frac{1 - \delta }{1 - {\delta }^{3}}\) 中方便地约去了 \(1 - \delta\) ；当 \(\delta  = 1\) 时，该项的极限 \(\frac{0}{0}\) 是未定义的，如果不约去则需要使用洛必达法则)。
\end{tcolorbox}

\begin{center}
\includegraphics[max width=0.9\textwidth]{images/bo_d2e5cmbef24c73b03tn0_339_340_1319_1242_558_0.jpg}
\end{center}
\hspace*{3em} 

Figure 11.16 Why the player's payoffs in the stationary SPE converge to the Nash bargaining solution(U, V)as \(\delta  \rightarrow  0\) . Point \(C = \left( {u\left( x\right) ,v\left( {1 - x}\right) }\right)\) is the split when player I makes the first demand, point \(E = \left( {u\left( {1 - y}\right) ,v\left( y\right) }\right)\) when player II makes it.

\begin{tcolorbox}
图11.16 为何静态子博弈完美均衡(SPE)中参与者的收益会随着 \(\delta  \rightarrow  0\) 收敛于纳什讨价还价解(U，V)。点 \(C = \left( {u\left( x\right) ,v\left( {1 - x}\right) }\right)\) 是参与者I首先提出需求时的分配点，点 \(E = \left( {u\left( {1 - y}\right) ,v\left( y\right) }\right)\) 是参与者II首先提出需求时的分配点。
\end{tcolorbox}

The stationary solution to the infinite bargaining game of alternating offers is shown in Figure 11.15(c), with points \(B = \left( {{\delta u}\left( {1 - y}\right) ,{\delta v}\left( y\right) }\right)\) and \(C = \left( {u\left( x\right) ,v\left( {1 - x}\right) }\right)\) . These points are shown again in Figure 11.16(a), with a "preview" of the Nash bargaining solution(U, V). Let the point \(E\) be given by \(E = \left( {u\left( {1 - y}\right) ,v\left( y\right) }\right)\) , which is the "undiscounted" version of \(B\) , that is, it has the same relative position on the outer curve as \(B\) on the middle curve. As mentioned at the end of the previous section, these are also the players' payoffs in the stationary SPE when player II moves first with her demand \(y\) . As \(\delta  \rightarrow  1\) , the points \(E\) and \(C\) converge to the same limit point. As shown in Figure 11.16(b), the points \(E\) and \(C\) define a line where the line segment between \(E\) and \(C\) is a "cord" of the Pareto-frontier. In the limit, the line becomes a tangent to the Pareto-frontier. If the Pareto-frontier is differentiable, the tangent is unique, otherwise we obtain in the limit some tangent to the frontier.

\begin{tcolorbox}
交替报价的无限讨价还价博弈的静态解如图11.15(c)所示，包含点 \(B = \left( {{\delta u}\left( {1 - y}\right) ,{\delta v}\left( y\right) }\right)\) 和 \(C = \left( {u\left( x\right) ,v\left( {1 - x}\right) }\right)\) 。这些点在图11.16(a)中再次展示，并“预览”了纳什讨价还价解(U，V)。设点 \(E\) 由 \(E = \left( {u\left( {1 - y}\right) ,v\left( y\right) }\right)\) 给出，它是 \(B\) 的“无贴现”版本，即它在外部曲线上的相对位置与 \(B\) 在中间曲线上的相对位置相同。如前一节末尾所述，当参与者II首先以她的需求 \(y\) 行动时，这些也是静态子博弈完美均衡中参与者的收益。随着 \(\delta  \rightarrow  1\) ，点 \(E\) 和 \(C\) 收敛于同一个极限点。如图11.16(b)所示，点 \(E\) 和 \(C\) 确定了一条直线，其中 \(E\) 和 \(C\) 之间的线段是帕累托前沿的一条“弦”。在极限情况下，这条直线成为帕累托前沿的切线。如果帕累托前沿是可微的，切线是唯一的，否则在极限情况下我们会得到前沿的某条切线。
\end{tcolorbox}

We use Figure 11.16(b) to show that as \(\delta  \rightarrow  1\) , the points \(E\) and \(C\) converge to the Nash bargaining solution. (The picture shows that the midpoint between \(E\) and \(C\) is already very close to(U, V)even when \(\delta\) is not close to 1 .) We use the geometric characterization of the bargaining solution in Proposition 11.5. Consider the points \(E\) and \(C\) as the upper-left and lower-right corner of a rectangle. As shown in Figure 11.16(b), the upper-right corner of that rectangle is the point \(F = \left( {u\left( x\right) ,v\left( y\right) }\right)\) , and its lower-left corner is the point \(G = \left( {u\left( {1 - y}\right) ,v\left( {1 - x}\right) }\right)\) . The equations (11.6) state that \(G = {\delta F}\) , that is, the points \(F\) and \(G\) are on a line through the origin(0,0), where the line segment between \(F\) and \(G\) is a diagonal of the rectangle. This line has slope \(v\left( y\right) /u\left( x\right)\) , shown as the angle \(\alpha\) near the origin. The other diagonal of the rectangle through the points \(E\) and \(C\) has the same (negative) slope, shown as the angle \(\alpha\) when that line intersects the \(u\) -axis.

\begin{tcolorbox}
我们使用图11.16(b)来证明，随着 \(\delta  \rightarrow  1\) ，点 \(E\) 和 \(C\) 收敛于纳什讨价还价解。(图中显示，即使当 \(\delta\) 不接近1时， \(E\) 和 \(C\) 的中点已经非常接近(U，V)。)我们使用命题11.5中讨价还价解的几何特征。将点 \(E\) 和 \(C\) 视为一个矩形的左上角和右下角。如图11.16(b)所示，该矩形的右上角是点 \(F = \left( {u\left( x\right) ,v\left( y\right) }\right)\) ，其左下角是点 \(G = \left( {u\left( {1 - y}\right) ,v\left( {1 - x}\right) }\right)\) 。方程(11.6)表明 \(G = {\delta F}\) ，即点 \(F\) 和 \(G\) 在一条过原点(0，0)的直线上，其中 \(F\) 和 \(G\) 之间的线段是矩形的一条对角线。这条直线的斜率为 \(v\left( y\right) /u\left( x\right)\) ，在原点附近表示为角 \(\alpha\) 。通过点 \(E\) 和 \(C\) 的矩形的另一条对角线具有相同的(负)斜率，当该直线与 \(u\) 轴相交时表示为角 \(\alpha\) 。
\end{tcolorbox}

In the limit, when \(\delta  \rightarrow  1\) , we have \(y \rightarrow  1 - x\) , and we obtain the same picture as in Figure 11.5 of a "roof" with two sides that have, in absolute value, the same slope. The left side of the roof is a line through the origin, and the right side of the roof is a tangent (of correct slope) to the Pareto-frontier of the bargaining set. This is exactly the condition of Proposition 11.5 that characterizes the apex of the roof as the Nash bargaining solution(U, V).

\begin{tcolorbox}
在极限情况下，当 \(\delta  \rightarrow  1\) 时，我们有 \(y \rightarrow  1 - x\) ，并且我们得到与图11.5中相同的“屋顶”图，其两侧的斜率绝对值相同。屋顶的左侧是一条过原点的直线，屋顶的右侧是讨价还价集帕累托前沿的切线(斜率正确)。这正是命题11.5中把屋顶顶点刻画为纳什讨价还价解(U，V)的条件。
\end{tcolorbox}

We summarize the main result of this chapter. The model of alternating offers with a discount factor is an interesting non-cooperative game in its own right. When the number of bargaining rounds is infinite, the game has an SPE in stationary strategies. When the discount factor \(\delta\) gets close to one, this stationary equilibrium gives a payoff pair that approaches the Nash bargaining solution.

\begin{tcolorbox}
我们总结本章的主要结果。具有贴现因子的交替报价模型本身就是一个有趣的非合作博弈。当讨价还价轮数为无限时，该博弈在平稳策略下存在子博弈完美均衡(SPE)。当贴现因子 \(\delta\) 接近 1 时，这个平稳均衡给出的收益对趋近于纳什讨价还价解。
\end{tcolorbox}

The Nash bargaining solution is originally derived from reasonable "axioms" for a solution to the bargaining problem. The "implementation" of this solution is not spelled out in this axiomatic approach, where it is merely assumed that the players reach a binding agreement in some way.

\begin{tcolorbox}
纳什讨价还价解最初是从讨价还价问题解的合理“公理”推导出来的。这种解的“实现”方式在这种公理化方法中并未详细说明，该方法只是假设参与者以某种方式达成了有约束力的协议。
\end{tcolorbox}

The alternating-offers game justifies this solution via a completely different route, namely a model of players that interact by making demands and possible rejections and counter-demands. In the resulting SPE with stationary strategies, each player anticipates the future behavior of the other player. The players immediately reach an agreement, which is the Nash bargaining solution.

\begin{tcolorbox}
交替报价博弈通过一条完全不同的途径证明了这个解的合理性，即一个参与者通过提出要求、可能的拒绝和反要求进行互动的模型。在由此产生的平稳策略子博弈完美均衡中，每个参与者都能预测到另一个参与者的未来行为。参与者立即达成协议，这个协议就是纳什讨价还价解。
\end{tcolorbox}

\section*{11.12 Further Reading}

\begin{tcolorbox}
\section*{11.12 进一步阅读}
\end{tcolorbox}

The axiomatic bargaining model is due to John Nash (1950) (and should not be confused with Nash equilibrium). An authoritative textbook on game theory, and cooperative games in particular, is Maschler, Solan, and Zamir (2013), where chapter 15 treats the Nash bargaining solution. It is also described in Osborne and Rubinstein (1994, chapter 15) and Osborne (2004, section 16.3).

\begin{tcolorbox}
公理化讨价还价模型归功于约翰·纳什(John Nash，1950 年)(不要与纳什均衡混淆)。关于博弈论，特别是合作博弈的权威教科书是马斯勒(Maschler)、索兰(Solan)和扎米尔(Zamir)(2013 年)所著，其中第 15 章讨论了纳什讨价还价解。奥斯本(Osborne)和鲁宾斯坦(Rubinstein)(1994 年，第 15 章)以及奥斯本(2004 年，第 16.3 节)也有相关描述。
\end{tcolorbox}

In favor of cooperative game theory, the assumption that people keep agreements may be quite justified in practice, without worrying too much about "enforcement mechanisms". The Dutch cultural historian Johan Huizinga (1949) argues that "play", the willingness to follow fairly arbitrary rules for their own sake, is an essential element of culture, where even just pretending to follow the rules of the game is much more tolerated than being a spoilsport and questioning the game itself.

\begin{tcolorbox}
支持合作博弈论的观点认为，在实践中，人们遵守协议的假设可能是相当合理的，不必过于担心“执行机制”。荷兰文化史学家约翰·赫伊津哈(Johan Huizinga，1949 年)认为，“游戏”，即出于自身原因愿意遵守相当随意的规则，是文化的一个基本要素，在文化中，即使只是假装遵守游戏规则也比破坏游戏兴致并质疑游戏本身更容易被容忍。
\end{tcolorbox}

For further details on the Ultimatum game see Osborne (2004, section 6.1.1). The infinite iterated-offers bargaining model with its unique stationary SPE that approaches the Nash bargaining solution is due to Rubinstein (1982). A finite multistage game was introduced earlier by Ståhl (1972).

\begin{tcolorbox}
关于最后通牒博弈的更多细节，请参阅奥斯本(2004 年，第 6.1.1 节)。具有唯一平稳子博弈完美均衡且趋近于纳什讨价还价解的无限迭代报价讨价还价模型归功于鲁宾斯坦(Rubinstein，1982 年)。有限多阶段博弈则是由斯塔尔(Ståhl，1972 年)更早提出的。
\end{tcolorbox}

The original inspiration for this chapter was Binmore (1992, chapter 5) and Binmore (2007, chapters 16 and 17). There, the players are assumed to have different "bargaining power", reflected in different discount factors for the players. Then the stationary strategies in the infinite game with very patient players approach a "generalized" Nash bargaining solution. We have simplified this to equal bargaining powers and equal discount factors, so that the stationary strategies approach the original Nash bargaining solution. For a similar textbook treatment see Osborne (2004, chapter 16).

\begin{tcolorbox}
本章的最初灵感来自宾默尔(Binmore，1992 年，第 5 章)和宾默尔(2007 年，第 16 和 17 章)。在那里，假设参与者具有不同的“讨价还价能力”，这反映在参与者不同的贴现因子上。然后，在参与者非常有耐心的无限博弈中，平稳策略趋近于一个“广义”的纳什讨价还价解。我们将其简化为具有相等的讨价还价能力和相等的贴现因子，以便平稳策略趋近于原始的纳什讨价还价解。类似的教科书处理方式可参阅奥斯本(2004 年，第 16 章)。
\end{tcolorbox}

\section*{11.13 Exercises for Chapter 11}

\begin{tcolorbox}
\section*{11.13 第 11 章练习}
\end{tcolorbox}

Exercise 11.1. Consider a bargaining set \(S\) with threat point(0,0). Let \(a > 0,b > 0\) , and let \({S}^{\prime } = \{ \left( {{au},{bv}}\right)  \mid  \left( {u,v}\right)  \in  S\}\) . The set \({S}^{\prime }\) is the set of utility pairs in \(S\) re-scaled by \(a\) for player I and by \(b\) for player II. Recall that one axiomatic property of the bargaining solution is independence from these scale factors.

\begin{tcolorbox}
练习 11.1。考虑一个威胁点为 (0,0) 的讨价还价集 \(S\) 。设 \(a > 0,b > 0\) ，且设 \({S}^{\prime } = \{ \left( {{au},{bv}}\right)  \mid  \left( {u,v}\right)  \in  S\}\) 。集合 \({S}^{\prime }\) 是 \(S\) 中效用对的集合，参与者 I 的效用按 \(a\) 重新缩放，参与者 II 的效用按 \(b\) 重新缩放。回想一下，讨价还价解的一个公理化性质是与这些缩放因子无关。
\end{tcolorbox}

Show that the Nash product fulfills this property, that is, the Nash bargaining solution from \(S\) obtained by maximising the Nash product re-scales to become the solution for \({S}^{\prime }\) .

\begin{tcolorbox}
证明纳什积满足这一性质，即通过最大化纳什积从 \(S\) 得到的纳什讨价还价解重新缩放后成为 \({S}^{\prime }\) 的解。
\end{tcolorbox}

Note: This is extremely easy once you state in the right way what it means that something maximizes something. You may find it useful to consider first the simpler case \(b = 1\) . For each bimatrix game, do the following:

\begin{tcolorbox}
注意:一旦你以正确的方式说明某个东西最大化另一个东西意味着什么，这就极其简单了。你可能会发现先考虑更简单的情况 \(b = 1\) 会很有用。对于每个双矩阵博弈，做以下事情:
\end{tcolorbox}

Exercise 11.2. Consider the following two bimatrix games (i) and (ii):

\begin{tcolorbox}
练习 11.2。考虑以下两个双矩阵博弈(i)和(ii):
\end{tcolorbox}

(i) (ii)

\begin{tcolorbox}
(i)(ii)
\end{tcolorbox}

\begin{center}
\includegraphics[max width=0.2\textwidth]{images/bo_d2e5cmbef24c73b03tn0_342_333_442_354_350_0.jpg}
\end{center}
\hspace*{3em} 

\begin{center}
\includegraphics[max width=0.2\textwidth]{images/bo_d2e5cmbef24c73b03tn0_342_951_443_355_348_0.jpg}
\end{center}
\hspace*{3em} 

(a) Find max-min strategies \(\widehat{x}\) and \(\widehat{y}\) (which may be mixed strategies) and corresponding max-min values \({u}_{0}\) and \({v}_{0}\) for player I and II, respectively.

\begin{tcolorbox}
(a)分别找出参与者 I 和 II 的极大极小策略 \(\widehat{x}\) 和 \(\widehat{y}\) (可能是混合策略)以及相应的极大极小值 \({u}_{0}\) 和 \({v}_{0}\) 。
\end{tcolorbox}

(b) Draw the convex hull of the payoff pairs in the \(2 \times  2\) game, and the bargaining set \(S\) that results from the payoff pairs(u, v)with \(u \geq  {u}_{0}\) and \(v \geq  {v}_{0}\) with the threat point \(\left( {{u}_{0},{v}_{0}}\right)\) found in (a). Indicate the Pareto-frontier of \(S\) .

\begin{tcolorbox}
(b) 绘制 \(2 \times  2\) 博弈中收益对的凸包(convex hull)，以及由满足 \(u \geq  {u}_{0}\) 和 \(v \geq  {v}_{0}\) 的收益对 (u, v) 与 (a) 中找到的威胁点(threat point) \(\left( {{u}_{0},{v}_{0}}\right)\) 所构成的讨价还价集(bargaining set) \(S\) 。指出 \(S\) 的帕累托前沿(Pareto-frontier)。
\end{tcolorbox}

(c) Find the Nash bargaining solution(U, V).

\begin{tcolorbox}
(c) 求纳什讨价还价解(Nash bargaining solution)(U, V)。
\end{tcolorbox}

(d) Show how the players can implement the Nash bargaining solution for \(S\) with a joint lottery over strategy pairs in the \(2 \times  2\) game.

\begin{tcolorbox}
(d) 说明玩家如何通过在 \(2 \times  2\) 博弈的策略对之间进行联合彩票(joint lottery)来实现 \(S\) 的纳什讨价还价解。
\end{tcolorbox}

Exercise 11.3. We repeat the definition of the discretized ultimatum game. Let \(M\) be a positive integer. Player I’s possible actions are to ask for a number \(x\) , called his demand, which is one of the (integer) numbers \(0,1,2,\ldots ,M\) . In response to this demand, player II can either accept (A) or reject (R). When player II accepts, then player I will receive a payoff of \(x\) and player II a payoff of \(M - x\) . When player II rejects, both players receive payoff zero.

\begin{tcolorbox}
练习 11.3。我们重复离散化最后通牒博弈(Ultimatum game)的定义。设 \(M\) 为正整数。玩家 I 的可能行动是提出一个数字 \(x\) ，称为他的要求，该数字是 \(0,1,2,\ldots ,M\) 这些(整数)数字之一。针对这个要求，玩家 II 可以接受(A)或拒绝(R)。当玩家 II 接受时，玩家 I 将获得收益 \(x\) ，玩家 II 将获得收益 \(M - x\) 。当玩家 II 拒绝时，两个玩家的收益均为零。
\end{tcolorbox}

(a) Draw the game tree for this game, as an extensive game with perfect information, for \(M = 3\) .

\begin{tcolorbox}
(a) 为 \(M = 3\) 绘制此博弈的博弈树(game tree)，将其作为具有完美信息(perfect information)的扩展式博弈。
\end{tcolorbox}

(b) What is the number of pure strategies in this game, for general \(M\) , for player \(\mathrm{I}\) and for player II? What is the number of reduced pure strategies in this game, for player I and for player II?

\begin{tcolorbox}
(b) 对于一般的 \(M\) ，玩家 \(\mathrm{I}\) 和玩家 II 在这个博弈中的纯策略数量分别是多少？玩家 I 和玩家 II 在这个博弈中的简化纯策略数量分别是多少？
\end{tcolorbox}

(c) Determine all subgame-perfect equilibria of the game in pure strategies.

\begin{tcolorbox}
(c) 确定该博弈在纯策略下的所有子博弈完美均衡(subgame-perfect equilibrium)。
\end{tcolorbox}

(d) Determine all equilibria (not necessarily subgame-perfect) of this game in pure strategies.

\begin{tcolorbox}
(d) 确定该博弈在纯策略下的所有均衡(equilibrium)(不一定是子博弈完美的)。
\end{tcolorbox}

(e) Determine all subgame-perfect equilibria of this game where both players may use behavior strategies, in addition to those in (c).

\begin{tcolorbox}
(e) 除了 (c) 中的那些，确定该博弈中两个玩家都可以使用行为策略(behavior strategy)的所有子博弈完美均衡。
\end{tcolorbox}

Exercise 11.4. Consider the following bargaining problem. In the usual way, a "unit pie" is split into non-negative amounts \(x\) and \(y\) with \(x + y \leq  1\) . The utility function of player \(\mathrm{I}\) is \(u\left( x\right)  = x\) , the utility function of player \(\mathrm{{II}}\) is \(v\left( y\right)  = 1 - {\left( 1 - y\right) }^{2}\) . The threat point is(0,0).

\begin{tcolorbox}
练习 11.4。考虑以下讨价还价问题。通常，一个“单位馅饼”被分割成非负数量的 \(x\) 和 \(y\) ，满足 \(x + y \leq  1\) 。玩家 \(\mathrm{I}\) 的效用函数(utility function)是 \(u\left( x\right)  = x\) ，玩家 \(\mathrm{{II}}\) 的效用函数是 \(v\left( y\right)  = 1 - {\left( 1 - y\right) }^{2}\) 。威胁点是 (0, 0)。
\end{tcolorbox}

(a) Draw the bargaining set (as set of possible utilities(u, v)) and indicate the Pareto-frontier in your drawing.

\begin{tcolorbox}
(a) 绘制讨价还价集(作为可能的效用 (u, v) 的集合)，并在图中指出帕累托前沿。
\end{tcolorbox}

(b) Find the Nash bargaining solution. How will the pie be split, and what are the utilities of the players?

\begin{tcolorbox}
(b) 求纳什讨价还价解。馅饼将如何分割，玩家的效用分别是多少？
\end{tcolorbox}

Exercise 11.5. Consider the bargaining problem of splitting a unit pie with utility functions \(u,v : \left\lbrack  {0,1}\right\rbrack   \rightarrow  \left\lbrack  {0,1}\right\rbrack\) as in Exercise 11.4, that is, \(u\left( x\right)  = x\) and \(v\left( y\right)  =\)  \(1 - {\left( 1 - y\right) }^{2}\) . Assume now that the bargaining behavior is described by the subgame-perfect equilibrium of the standard alternating-offers bargaining game with \(k\) rounds. In round 1 player I makes a certain demand. If player II rejects this, she can make a certain counter-demand in round 2, and so on. A rejection in the last round \(k\) means that both players get nothing. The last player to make a demand is player I if \(k\) is odd, and player II if \(k\) is even. When agreement takes place in round \(i\) for some \(i\) with \(1 \leq  i \leq  k\) , then the bargaining set has shrunk by a factor of \({\delta }^{i - 1}\) . That is, the utility of what the players can get after each round of rejection is reduced by multiplication by the discount factor \(\delta\) .

\begin{tcolorbox}
练习 11.5。考虑如练习 11.4 中具有效用函数 \(u,v : \left\lbrack  {0,1}\right\rbrack   \rightarrow  \left\lbrack  {0,1}\right\rbrack\) 的分割单位馅饼的讨价还价问题，即 \(u\left( x\right)  = x\) 和 \(v\left( y\right)  =\)  \(1 - {\left( 1 - y\right) }^{2}\) 。现在假设讨价还价行为由具有 \(k\) 轮的标准交替报价讨价还价博弈的子博弈完美均衡描述。在第 1 轮，玩家 I 提出一个特定要求。如果玩家 II 拒绝，她可以在第 2 轮提出一个特定的反要求，依此类推。在最后一轮 \(k\) 拒绝意味着两个玩家都得不到任何东西。如果 \(k\) 是奇数，最后提出要求的玩家是玩家 I；如果 \(k\) 是偶数，最后提出要求的玩家是玩家 II。当在第 \(i\) 轮(对于某个满足 \(1 \leq  i \leq  k\) 的 \(i\) )达成协议时，讨价还价集缩小了 \({\delta }^{i - 1}\) 倍。也就是说，每一轮拒绝后玩家所能获得的效用通过乘以贴现因子 \(\delta\) 而减少。
\end{tcolorbox}

(a) Suppose that the number of rounds is \(k = 2\) . Draw the picture of shrinking bargaining sets (for some choice of \(\delta\) ) in these two rounds. Find the subgame perfect equilibrium, in terms of player I’s demand \(x\) , for general \(\delta\) . What is the demand \(x\) when \(\delta  = \frac{3}{4}\) ? What are the strategies of the two players in this subgame perfect equilibrium?

\begin{tcolorbox}
(a) 假设轮数为 \(k = 2\) 。画出这两轮中(对于 \(\delta\) 的某些选择)不断缩小的谈判集的图形。对于一般的 \(\delta\) ，找出子博弈完美均衡，用参与者 I 的需求 \(x\) 表示。当 \(\delta  = \frac{3}{4}\) 时，需求 \(x\) 是多少？在这个子博弈完美均衡中，两个参与者的策略是什么？
\end{tcolorbox}

(b) Suppose that the number of rounds is \(k = 3\) . Draw the picture of shrinking bargaining sets in these three rounds. Find the subgame perfect equilibrium, for general \(\delta\) , and compute the demand of player I in round 1 (you do not need to specify the players’ strategies). What is player I’s demand \(x\) when \(\delta  = \frac{1}{4}\) ? (Then \(x\) happens to be a rational number.)

\begin{tcolorbox}
(b) 假设轮数为 \(k = 3\) 。画出这三轮中不断缩小的谈判集的图形。对于一般的 \(\delta\) ，找出子博弈完美均衡，并计算参与者 I 在第一轮的需求(你不需要指定参与者的策略)。当 \(\delta  = \frac{1}{4}\) 时，参与者 I 的需求 \(x\) 是多少？(此时 \(x\) 恰好是一个有理数。)
\end{tcolorbox}

(c) Consider now an infinite number of rounds, and look for stationary equilibrium strategies, where player I demands the same amount \(x\) each time he makes a demand, and player II demands the same amount \(y\) every time she makes a demand. What is this demand \(x\) when \(\delta  = \frac{3}{5}\) ? What is the demand \(x\) when \(\delta\) tends to 1? Compare this with the result of Exercise 11.4(b).

\begin{tcolorbox}
(c) 现在考虑有无穷多轮，寻找平稳均衡策略，即参与者 I 每次提出需求时都要求相同的数量 \(x\) ，参与者 II 每次提出需求时都要求相同的数量 \(y\) 。当 \(\delta  = \frac{3}{5}\) 时，这个需求 \(x\) 是多少？当 \(\delta\) 趋近于 1 时，需求 \(x\) 是多少？将此结果与练习 11.4(b) 的结果进行比较。
\end{tcolorbox}

\section*{12}

\section*{Correlated Equilibrium}

\begin{tcolorbox}
\section*{相关均衡}
\end{tcolorbox}

In this final chapter, we explain a new equilibrium concept called correlated equilibrium that is more general than Nash equilibrium. It allows for randomized actions of the players that depend on an external signal (like a traffic light) that is observed by the players (typically in different ways) so that their actions can be correlated.

\begin{tcolorbox}
在这最后一章中，我们将介绍一种新的均衡概念，称为相关均衡(correlated equilibrium)，它比纳什均衡(Nash equilibrium)更具一般性。它允许参与者根据一个外部信号(如交通信号灯)采取随机行动，参与者(通常以不同的方式)可以观察到这个信号，从而使他们的行动能够相互关联。
\end{tcolorbox}

The set of Nash equilibria of a game is a set of fixed points and typically disconnected. The set of correlated equilibria is larger, and mathematically simpler because it is a convex polytope. In Section 12.4, we also explain the concept of coarse correlated equilibrium, where the players commit to letting the correlation device choose their actions for them, rather than just giving them recommendations. This restricts the players' possibility for deviating and leads to an even larger set (with possibly more beneficial outcomes, see Exercise 12.3).

\begin{tcolorbox}
一个博弈的纳什均衡集是一组不动点，通常是不连通的。相关均衡集更大，并且在数学上更简单，因为它是一个凸有界多面体(convex polytope)。在第 12.4 节中，我们还将解释粗相关均衡(coarse correlated equilibrium)的概念，在这种情况下，参与者承诺让相关机制(correlation device)为他们选择行动，而不仅仅是给他们提供建议。这限制了参与者偏离的可能性，并导致一个更大的集合(可能会有更有利的结果，见练习 12.3)。
\end{tcolorbox}

Because the set of correlated equilibria is a convex polytope, it would be nice to prove that it is not empty without relying on a fixed-point theorem. Indeed, Hart and Schmeidler (1989) gave such an existence proof based on the minimax Theorem 8.4, which we present in Section 12.5 (and which is not normally found in textbooks at this level). The minimax theorem leads to elegant proofs of the existence of stationary distributions for Markov chains (Theorem 12.9) and of correlated equilibria (Theorem 12.11).

\begin{tcolorbox}
由于相关均衡集是一个凸有界多面体，若能不依赖不动点定理来证明它非空就好了。事实上，哈特和施迈德勒(Hart and Schmeidler，1989)基于极小极大定理(minimax Theorem 8.4)给出了这样一个存在性证明，我们将在第 12.5 节中介绍(这在这个水平的教科书中通常是找不到的)。极小极大定理为马尔可夫链(Markov chain)平稳分布(stationary distribution)的存在性(定理 12.9)和相关均衡的存在性(定理 12.11)提供了优雅的证明。
\end{tcolorbox}

The concept of correlated equilibrium is used much less often than Nash equilibrium in the analysis of games, but is important in models of learning in games, mentioned in the final reference Section 12.6.

\begin{tcolorbox}
在博弈分析中，相关均衡的概念比纳什均衡使用得少得多，但在第 12.6 节最后参考文献中提到的博弈学习(learning in games)模型中很重要。
\end{tcolorbox}

\section*{12.1 Prerequisites and Learning Outcomes}

\begin{tcolorbox}
\section*{12.1 先决条件和学习成果}
\end{tcolorbox}

Apart from the basic concepts of strategic form and equilibrium in Chapter 3, the main prerequisite is Chapter 8 with the minimax Theorem 8.4.

\begin{tcolorbox}
除了第 3 章中策略形式(strategic form)和均衡(equilibrium)的基本概念外，主要的先决条件是第 8 章中的极小极大定理(minimax Theorem 8.4)。
\end{tcolorbox}

After studying this chapter, you should be able to:

\begin{tcolorbox}
学习本章后，你应该能够:
\end{tcolorbox}

\begin{itemize}
\item explain correlated equilibrium, its difference to Nash equilibrium, and that correlated equilibria are more than just randomizations over Nash equilibria;
\end{itemize}

\begin{tcolorbox}
\begin{itemize}
\item 解释相关均衡，它与纳什均衡的区别，以及相关均衡不仅仅是纳什均衡上的随机化；
\end{itemize}
\end{tcolorbox}

\begin{itemize}
\item write down the incentive constraints of a game that define the set of correlated equilibria;
\end{itemize}

\begin{tcolorbox}
\begin{itemize}
\item 写出定义相关均衡集的博弈的激励约束(incentive constraints)；
\end{itemize}
\end{tcolorbox}

\begin{itemize}
\item solve these incentive constraints for small games, like in Exercise 12.1;
\end{itemize}

\begin{tcolorbox}
\begin{itemize}
\item 像练习 12.1 那样，求解小博弈的这些激励约束；
\end{itemize}
\end{tcolorbox}

\begin{itemize}
\item describe the concept of coarse correlated equilibrium and why it requires a commitment by the players;
\end{itemize}

\begin{tcolorbox}
\begin{itemize}
\item 描述粗相关均衡的概念以及为什么它需要参与者做出承诺；
\end{itemize}
\end{tcolorbox}

\begin{itemize}
\item outline the mathematical argument for the existence of a correlated equilibrium with the help of the minimax theorem, and why this route is overall simpler than via the existence of a Nash equilibrium;
\end{itemize}

\begin{tcolorbox}
\begin{itemize}
\item 借助极小极大定理概述相关均衡存在性的数学论证，以及为什么这种方法总体上比通过纳什均衡的存在性更简单；
\end{itemize}
\end{tcolorbox}

\begin{itemize}
\item describe in more detail the idea behind this existence proof, explained at the end of Section 12.5: the auxiliary zero-sum game with its "correlator" COR and "deviator" DEV, and how COR can "neutralize" any "deviation plan" of DEV.
\end{itemize}

\begin{tcolorbox}
\begin{itemize}
\item 更详细地描述第12.5节末尾解释的这个存在性证明背后的思想:带有“相关者”COR和“偏离者”DEV的辅助零和博弈，以及COR如何“中和”DEV的任何“偏离计划”。
\end{itemize}
\end{tcolorbox}

\section*{12.2 Examples of Correlated Equilibria}

\begin{tcolorbox}
\section*{12.2 相关均衡的例子}
\end{tcolorbox}

Towards the end of the last century I visited the Hebrew University of Jerusalem. In the faculty club, I ran into the game theorists Robert Aumann and Michael Maschler. Aumann described how "a conversation with Professor Maschler" led him to the idea of correlated equilibrium. He said they discussed how one would play a mixed strategy in practice: Let the outcome of some random event determine the action that you take. For example, choose A if the sun shines, and choose B if it does not. The probability is roughly known from the weather forecast. But this approach revealed a new possibility: Using the weather for randomization does not produce a mixed strategy as considered by von Neumann or Nash, where players play randomly by flipping a coin on their own. If players let their actions be decided by the weather, then their random choices will not be independent. The concept of correlated equilibrium, due to Aumann (1974), allows for correlated randomized actions among the players, and generalizes the concept of Nash equilibrium. (In this chapter, we always use the term "Nash equilibrium" to distinguish it from other equilibrium notions.) Correlated equilibria are the topic of this chapter.

\begin{tcolorbox}
上世纪末，我访问了耶路撒冷希伯来大学。在教师俱乐部里，我遇到了博弈论学者罗伯特·奥曼(Robert Aumann)和迈克尔·马施勒(Michael Maschler)。奥曼讲述了“与马施勒教授的一次交谈”如何使他产生了相关均衡的想法。他说，他们讨论了在实践中如何采用混合策略:让某个随机事件的结果决定你采取的行动。例如，如果出太阳就选择A，如果不出太阳就选择B。从天气预报中大致可以知道这种概率。但这种方法揭示了一种新的可能性:利用天气进行随机化并不会产生冯·诺伊曼(von Neumann)或纳什(Nash)所考虑的混合策略，在他们的理论中，玩家通过自己抛硬币来随机行动。如果玩家让自己的行动由天气决定，那么他们的随机选择将不是相互独立的。奥曼(1974)提出的相关均衡概念允许玩家之间采取相关的随机行动，并推广了纳什均衡的概念。(在本章中，我们始终使用“纳什均衡”这一术语来将其与其他均衡概念区分开来。)相关均衡是本章的主题。
\end{tcolorbox}

Consider the variant of the game Chicken in Figure 12.1. Each player has two strategies, here named \(T\) and \(B\) for player I and \(L\) and \(R\) for player II. Compared to the original Chicken game in Figure 3.3, a player's payoff is increased by 3 if the other player plays her second strategy (by Proposition 9.7, or seen directly, this does not affect the best response to any mixed strategy). The game has the two non-symmetric pure Nash equilibria(B, L)with payoffs1,5to players \(\mathrm{I}\) and \(\mathrm{{II}}\) , and (T, R)with payoffs5,1. In the unique mixed equilibrium, both players use the mixed strategy \(\left( {\frac{1}{2},\frac{1}{2}}\right)\) , with payoff \(\frac{5}{2}\) to each player.

\begin{tcolorbox}
考虑图12.1中胆小鬼博弈(Chicken)的变体。每个玩家有两种策略，这里玩家I的策略命名为 \(T\) 和 \(B\) ，玩家II的策略命名为 \(L\) 和 \(R\) 。与图3.3中的原始胆小鬼博弈相比，如果另一个玩家采用她的第二种策略，那么该玩家的收益会增加3(根据命题9.7，或者直接观察可知，这并不影响对任何混合策略的最优反应)。该博弈有两个非对称的纯纳什均衡:(B，L)，玩家 \(\mathrm{I}\) 和 \(\mathrm{{II}}\) 的收益分别为1和5；以及(T，R)，收益分别为5和1。在唯一的混合均衡中，两个玩家都采用混合策略 \(\left( {\frac{1}{2},\frac{1}{2}}\right)\) ，每个玩家的收益为 \(\frac{5}{2}\) 。
\end{tcolorbox}

\begin{center}
\includegraphics[max width=0.9\textwidth]{images/bo_d2e5cmbef24c73b03tn0_346_195_197_1233_356_0.jpg}
\end{center}
\hspace*{3em} 

Figure 12.1 A variant of the game Chicken in Figure 3.3. The three panels on the right show the three Nash equilibria of the game as probability distributions on strategy profiles.

\begin{tcolorbox}
图12.1 图3.3中胆小鬼博弈的一个变体。右侧的三个面板展示了该博弈在策略组合上的概率分布形式的三个纳什均衡。
\end{tcolorbox}

The three panels on the right in Figure 12.1 show the three Nash equilibria as probability distributions on the four strategy pairs. A pure Nash equilibrium selects only one strategy pair. The mixed Nash equilibrium selects each strategy pair with equal probability \(\frac{1}{4}\) .

\begin{tcolorbox}
图12.1右侧的三个面板展示了四个策略对的概率分布形式的三个纳什均衡。一个纯纳什均衡只选择一个策略对。混合纳什均衡以相等的概率 \(\frac{1}{4}\) 选择每个策略对。
\end{tcolorbox}

In a correlated equilibrium, this probability distribution is thought of as a correlation device that randomly selects a pure strategy profile, where for each player the strategy in that profile is privately told to the player as a recommendation of what to play. Each player is assumed to know how the device works, and can infer from his own recommendation what the other players have been recommended to play. The equilibrium condition states that it is optimal for each player to choose their recommended action, assuming that the other players follow their recommendations.

\begin{tcolorbox}
在相关均衡中，这种概率分布被视为一种相关机制，它随机选择一个纯策略组合，并将该组合中每个玩家的策略私下告知该玩家，作为对其行动的建议。假设每个玩家都知道该机制的运作方式，并且可以从自己得到的建议推断出其他玩家得到的建议。均衡条件表明，假设其他玩家遵循他们的建议，那么每个玩家选择被建议的行动是最优的。
\end{tcolorbox}

For the first pure Nash equilibrium(B, L)in Figure 12.1, the device selects this strategy pair, and recommends player I to play \(B\) and player II to play \(L\) . Here, each player knows what the other has been told, and his own recommendation is a best response to that strategy, so this is a correlated equilibrium. The same applies to the second pure equilibrium(T, R).

\begin{tcolorbox}
对于图12.1中的第一个纯纳什均衡(B，L)，该机制选择这个策略对，并建议玩家I采用 \(B\) ，玩家II采用 \(L\) 。在这里，每个玩家都知道对方得到的建议，并且自己得到的建议是对对方策略的最优反应，因此这是一个相关均衡。这同样适用于第二个纯均衡(T，R)。
\end{tcolorbox}

In the third, mixed Nash equilibrium in Figure 12.1, the device chooses each strategy pair with equal probability. This includes strategy pairs that are not themselves Nash equilibria, such as(T, L). If the device selects(T, L), then player \(\mathrm{I}\) is recommended to play \(T\) . However, and this is crucial, player \(\mathrm{I}\) is not told that player II has been recommended to play L. Player I can only infer from the fact that he has been recommended to play \(T\) that player II has been recommended to play \(L\) and \(R\) with equal probability (this is the same when player \(\mathrm{I}\) has been recommended to play \(B\) ). Given this randomized choice of player II, player I is indifferent between his two rows, and therefore it is optimal to follow the recommended action. The same holds for any strategy pair that is chosen by the device. Therefore, this is also a correlated equilibrium.

\begin{tcolorbox}
在图12.1中的第三个混合纳什均衡中，该机制以相等的概率选择每个策略对。这包括本身并非纳什均衡的策略对，例如(T，L)。如果该机制选择(T，L)，那么玩家 \(\mathrm{I}\) 被建议采用 \(T\) 。然而，关键的是，玩家 \(\mathrm{I}\) 并不知道玩家II被建议采用L。玩家I只能从自己被建议采用 \(T\) 这一事实推断出玩家II被建议采用 \(L\) 和 \(R\) 的概率相等(当玩家 \(\mathrm{I}\) 被建议采用 \(B\) 时也是如此)。鉴于玩家II的这种随机选择，玩家I对自己的两行策略无差异，因此遵循被建议的行动是最优的。对于该机制选择的任何策略对都是如此。因此，这也是一个相关均衡。
\end{tcolorbox}

\begin{center}
\includegraphics[max width=0.9\textwidth]{images/bo_d2e5cmbef24c73b03tn0_347_330_202_1258_272_0.jpg}
\end{center}
\hspace*{3em} 

Figure 12.2 Examples of correlated equilibria for the game in Figure 12.1, and their expected payoffs. The general form of a correlated distribution is shown on the right.

\begin{tcolorbox}
图12.2展示了图12.1中博弈的相关均衡(correlated equilibrium)示例及其期望效用(expected payoff)。右侧显示了相关分布的一般形式。
\end{tcolorbox}

In a mixed Nash equilibrium(x, y)of an \(m \times  n\) bimatrix game, a pure strategy pair(i, j)is chosen with probability \({x}_{i}{y}_{j}\) because the two players randomize independently. In a correlated equilibrium for an \(m \times  n\) game, each strategy pair (i, j)has some probability \({z}_{ij}\) , called the joint probability for this pair. Figure 12.2 shows more general probability distributions of this kind. In the left panel, the two strategy pairs(B, L)and(T, L)are chosen with equal probability. No matter which pair is chosen, both players know with certainty the recommendation to the other player. The two pure strategies are best responses to each other, so the equilibrium condition holds. Here, the device chooses the other two strategy pairs (T, L)and(B, R)with probability zero. The joint probability distribution is not independent but correlated, that is, it cannot be expressed in the form \({z}_{ij} = {x}_{i}{y}_{j}\) because, for example, \({z}_{11} = 0\) would imply \({x}_{1} = 0\) or \({y}_{1} = 0\) , where \({x}_{1} = 0\) implies \({z}_{12} = 0\) and \({y}_{1} = 0\) implies \({z}_{21} = 0\) , but both \({z}_{12}\) and \({z}_{21}\) are positive.

\begin{tcolorbox}
在双矩阵博弈(bimatrix game)的混合纳什均衡(Nash equilibrium)(x, y)中，由于两个参与者独立地进行随机化选择，纯策略对(i, j)被选中的概率为 \({x}_{i}{y}_{j}\) 。在双矩阵博弈的相关均衡中，每个策略对 (i, j) 都有一定的概率 \({z}_{ij}\) ，称为该策略对的联合概率。图12.2展示了更一般的此类概率分布。在左图中，两个策略对 (B, L) 和 (T, L) 被选中的概率相等。无论选择哪一对，两个参与者都能确定地知道对另一个参与者的建议。这两个纯策略互为最优反应(best response)，因此均衡条件成立。在这里，装置选择另外两个策略对 (T, L) 和 (B, R) 的概率为零。联合概率分布不是独立的，而是相关的，也就是说，它不能表示为 \({z}_{ij} = {x}_{i}{y}_{j}\) 的形式，因为，例如， \({z}_{11} = 0\) 意味着 \({x}_{1} = 0\) 或 \({y}_{1} = 0\) ，其中 \({x}_{1} = 0\) 意味着 \({z}_{12} = 0\) ， \({y}_{1} = 0\) 意味着 \({z}_{21} = 0\) ，但 \({z}_{12}\) 和 \({z}_{21}\) 都是正数。
\end{tcolorbox}

This correlated equilibrium in the left panel in Figure 12.2 gives an expected payoff of 3 to each player, which is not achieved in any Nash equilibrium. As an equal "mixture" between the pure Nash equilibria, it also seems like a fair way to play the game. It could be implemented by a fair coin toss, whose outcome then determines which of the two pure Nash equilibria should be played.

\begin{tcolorbox}
图12.2左图中的这个相关均衡为每个参与者提供了3的期望效用，这在任何纳什均衡中都无法实现。作为纯纳什均衡之间的平等“混合”，它似乎也是一种公平的博弈方式。这可以通过抛一枚公平的硬币来实现，硬币的结果决定了应该采用哪一个纯纳什均衡。
\end{tcolorbox}

However, the players can achieve even higher expected payoffs with the device shown in the second panel in Figure 12.2. There, the strategy pair(B, R)with payoffs 3 for each player has probability \(\frac{1}{3}\) , the same as the other two pairs(B, L) and(T, R). The top-left strategy pair(T, L)with the low payoffs0,0has probability zero. If the device chooses(B, R), then player I (and similarly player II) only knows that the other player has been recommended either pure strategy with probability \(\frac{1}{2}\) , against which strategy 2 is a best response. If the device chooses (B, L), then player I also infers that player II has been told to play either strategy with probability \(\frac{1}{2}\) , whereas player II knows that player I has been told to play strategy 2 for sure; vice versa, the same applies when the device chooses(T, L). Hence, this is also a correlated equilibrium, with expected payoffs \(\left( {\frac{10}{3},\frac{10}{3}}\right)\) . It can be shown that this correlated equilibrium has the largest sum of expected payoffs to the two players.

\begin{tcolorbox}
然而，通过图12.2第二图所示的装置，参与者可以获得更高的期望效用。在那里，每个参与者收益为3的策略对 (B, R) 的概率为 \(\frac{1}{3}\) ，与另外两个策略对 (B, L) 和 (T, R) 相同。左上角收益为0,0的策略对 (T, L) 的概率为零。如果装置选择 (B, R)，那么参与者I(参与者II同理)只知道另一个参与者被建议采用任一纯策略的概率为 \(\frac{1}{2}\) ，对此策略2是最优反应。如果装置选择 (B, L)，那么参与者I也能推断出参与者II被建议以概率 \(\frac{1}{2}\) 采用任一策略，而参与者II知道参与者I肯定被建议采用策略2；反之，当装置选择 (T, L) 时情况相同。因此，这也是一个相关均衡，期望效用为 \(\left( {\frac{10}{3},\frac{10}{3}}\right)\) 。可以证明，这个相关均衡使两个参与者的期望效用之和最大。
\end{tcolorbox}

However, there can also be correlated equilibria that have worse expected payoffs than any Nash equilibrium, as shown in the third panel in Figure 12.2. There, the device chooses the strategy pair(T, L)with its bad payoffs0,0with probability \(\frac{1}{3}\) , and the strategy pair(B, R)with its good payoffs4,4with probability zero. This is a correlated equilibrium for the same reasons as before: Each player infers from his own recommendation either a deterministic or a \(\left( {\frac{1}{2},\frac{1}{2}}\right)\) random choice of his opponent, against which the own recommendation is a best response. The expected payoff, however, is \(\frac{0 + 5 + 1}{3} = 2\) to each player, a sum of payoffs that is worse than in any Nash equilibrium.

\begin{tcolorbox}
然而，也可能存在相关均衡，其期望效用比任何纳什均衡都差，如图12.2第三图所示。在那里，装置以概率 \(\frac{1}{3}\) 选择收益较差的策略对 (T, L)(收益为0,0)，而以概率零选择收益较好的策略对 (B, R)(收益为4,4)。出于与之前相同的原因，这是一个相关均衡:每个参与者从自己得到的建议中推断出对手是确定性选择还是以 \(\left( {\frac{1}{2},\frac{1}{2}}\right)\) 的概率随机选择，而自己得到的建议对此是最优反应。然而，每个参与者的期望效用为 \(\frac{0 + 5 + 1}{3} = 2\) ，这种收益之和比任何纳什均衡都差。
\end{tcolorbox}

\section*{12.3 Incentive Constraints}

\begin{tcolorbox}
\section*{12.3 激励约束(incentive constraints)}
\end{tcolorbox}

In this section, we formalize the concept of correlated equilibrium, first for two players because the notation is easier. Consider an \(m \times  n\) bimatrix game(A, B)with payoffs \({a}_{ij}\) and \({b}_{ij}\) to players I and II if they play row \(i\) and column \(j\) , respectively. A correlation device is a probability distribution \(z\) on pure strategy pairs, with probability \({z}_{ij}\) for the pair(i, j), where these numbers fulfill

\begin{tcolorbox}
在本节中，我们将形式化相关均衡(correlated equilibrium)的概念，首先考虑两个参与者的情况，因为这种情况下的符号表示更简单。考虑一个 \(m\times n\) 双矩阵博弈(bimatrix game) \((A,B)\) ，如果参与者 I 和参与者 II 分别选择第 \(i\) 行和第 \(j\) 列，他们的收益分别为 \(a_{ij}\) 和 \(b_{ij}\) 。相关机制(correlation device)是纯策略对的一个概率分布 \(p\) ，策略对 \((i,j)\) 的概率为 \(p_{ij}\) ，其中这些数值满足
\end{tcolorbox}

\[
{z}_{ij} \geq  0\;\left( {1 \leq  i \leq  m,1 \leq  j \leq  n}\right) ,\;\mathop{\sum }\limits_{{i = 1}}^{m}\mathop{\sum }\limits_{{j = 1}}^{n}{z}_{ij} = 1. \tag{12.1}
\]

These probabilities are common knowledge to the players. The device selects a strategy pair according to these probabilities, and privately tells each player only his pure strategy in that pair. Suppose player \(\mathrm{I}\) is told strategy \(i\) . Then the \(i\) th row of the device (as in the panels in Figure 12.2) has probabilities \({z}_{i1},\ldots ,{z}_{in}\) for the \(n\) columns \(j = 1,\ldots ,n\) of what player II has been told. Given that player I has been told \(i\) , the conditional probability for column \(j\) is therefore

\begin{tcolorbox}
这些概率是参与者的共同知识(common knowledge)。该机制根据这些概率选择一个策略对，并私下只告知每个参与者他在该策略对中的纯策略。假设参与者 I 被告知选择策略 \(i\) 。那么，该机制的第 \(i\) 行(如图 12.2 中的面板所示)给出了参与者 II 被告知选择各列 \(j\) 的概率 \(p_{ij}\) 。已知参与者 I 被告知选择 \(i\) ，那么参与者 II 选择第 \(j\) 列的条件概率(conditional probability)为
\end{tcolorbox}

\[
\operatorname{prob}\left( {j \mid  i}\right)  = \frac{{z}_{ij}}{\mathop{\sum }\limits_{{l = 1}}^{n}{z}_{il}}. \tag{12.2}
\]

This expression is well defined whenever row \(i\) is chosen with positive probability by \(z\) , because then \({z}_{il} > 0\) for at least one column \(l\) , so that the denominator in (12.2) is positive. This conditional probability distribution is what player I infers (when told to play row \(i\) ) about the recommended action \(j\) to player II. If player \(\mathrm{I}\) follows the recommended action \(i\) , he gets a certain expected payoff. This payoff is optimal for player I if no other row \(k\) gives him higher expected payoff, using the same conditional distribution about player II’s random choice \(j\) (because player I does not get any further information when he deviates from \(i\) to \(k\) ). That is, for all other rows \(k\) we require that

\begin{tcolorbox}
只要 \(p\) 以正概率选择第 \(i\) 行，这个表达式就是有定义的，因为此时至少存在一个列 \(j\) 使得 \(p_{ij}>0\) ，从而(12.2)式中的分母为正。这个条件概率分布就是参与者 I(当被告知选择第 \(i\) 行时)对参与者 II 被推荐采取的行动 \(j\) 的推断。如果参与者 I 遵循被推荐的行动 \(i\) ，他将获得一定的期望效用(expected payoff)。如果在使用关于参与者 II 随机选择 \(j\) 的相同条件分布下，没有其他行 \(i'\) 能给参与者 I 带来更高的期望效用(因为当参与者 I 从 \(i\) 偏离到 \(i'\) 时，他不会获得更多信息)，那么这个收益对参与者 I 来说就是最优的。也就是说，对于所有其他行 \(i'\) ，我们要求
\end{tcolorbox}

\[
\mathop{\sum }\limits_{{j = 1}}^{n}{a}_{ij}\frac{{z}_{ij}}{\mathop{\sum }\limits_{{l = 1}}^{n}{z}_{il}} \geq  \mathop{\sum }\limits_{{j = 1}}^{n}{a}_{kj}\frac{{z}_{ij}}{\mathop{\sum }\limits_{{l = 1}}^{n}{z}_{il}}. \tag{12.3}
\]

Multiplication by \(\mathop{\sum }\limits_{{l = 1}}^{n}{z}_{il}\) on both sides shows that (12.3) is equivalent to

\begin{tcolorbox}
两边同时乘以 \(p_{i\cdot}\) 表明(12.3)式等价于
\end{tcolorbox}

\[
\mathop{\sum }\limits_{{j = 1}}^{n}\left( {{a}_{ij} - {a}_{kj}}\right) {z}_{ij} \geq  0\;\left( {1 \leq  i,k \leq  m}\right) . \tag{12.4}
\]

The \({m}^{2}\) inequalities for all pairs \(i,k\) in (12.4) hold also trivially if \(i = k\) , and also if \(\mathop{\sum }\limits_{{l = 1}}^{n}{z}_{il} = 0\) where row \(i\) has probability zero under the distribution \(z\) , because this implies \({z}_{ij} = 0\) for all \(j\) for this \(i\) . If \(\mathop{\sum }\limits_{{l = 1}}^{n}{z}_{il} > 0\) , then (12.4) implies (12.3) and states that player I has no incentive to deviate from the recommended row \(i\) to row \(k\) . The inequalities (12.4) are therefore called the incentive constraints for player I. The corresponding incentive constraints for player II are

\begin{tcolorbox}
对于(12.4)式中所有的策略对 \((i,i')\) ，当 \(p_{i\cdot}=0\) 时， \(n\) 个不等式显然成立；当 \(p_{i'\cdot}=0\) (即在分布 \(p\) 下第 \(i'\) 行的概率为零)时，这些不等式也显然成立，因为这意味着对于这个 \(i'\) 的所有 \(j\) 都有 \(p_{i'j}=0\) 。如果 \(p_{i\cdot}>0\) ，那么(12.4)式蕴含(12.3)式，并且表明参与者 I 没有动机从被推荐的第 \(i\) 行偏离到第 \(i'\) 行。因此，不等式(12.4)被称为参与者 I 的激励约束(incentive constraints)。参与者 II 对应的激励约束为
\end{tcolorbox}

\[
\mathop{\sum }\limits_{{i = 1}}^{m}\left( {{b}_{ij} - {b}_{il}}\right) {z}_{ij} \geq  0\;\left( {1 \leq  j,l \leq  n}\right)  \tag{12.5}
\]

and say that player II, when recommended to play column \(j\) , has no incentive to deviate to column \(l\) .

\begin{tcolorbox}
并且表明，当参与者 II 被推荐选择第 \(j\) 列时，他没有动机偏离到第 \(j'\) 列。
\end{tcolorbox}

Definition 12.1 (Correlated equilibrium, two players). For an \(m \times  n\) bimatrix game (A, B), a correlated equilibrium is an \(m \times  n\) matrix \(z\) of joint probabilities for strategy pairs(i, j)that fulfills(12.1)and the incentive constraints(12.4)for player \(\mathrm{I}\) and(12.5) for player II. The matrix \(z\) defines a correlation device that selects the strategy pair (i, j)with probability \({z}_{ij}\) , and then privately (without telling the other player) tells \(i\) to player \(\mathrm{I}\) and \(j\) to player \(\mathrm{{II}}\) .

\begin{tcolorbox}
定义 12.1(两个参与者的相关均衡)。对于一个 \(m\times n\) 双矩阵博弈 \((A,B)\) ，相关均衡是一个 \(m\times n\) 矩阵 \(p\) ，它是策略对 \((i,j)\) 的联合概率矩阵，满足(12.1)式以及参与者 I 的激励约束(12.4)式和参与者 II 的激励约束(12.5)式。矩阵 \(p\) 定义了一个相关机制，该机制以概率 \(p_{ij}\) 选择策略对 \((i,j)\) ，然后私下(不告知另一个参与者)告知参与者 I 选择 \(i\) ，告知参与者 II 选择 \(j\) 。
\end{tcolorbox}

The incentive constraints (12.4) for player \(\mathrm{I}\) are about each row \(i\) of the matrix \(z\) . Up to a normalization factor (as it appears in (12.3)), this row of \(z\) is player I’s "posterior" about the behavior of player II after player I gets the recommendation to play \(i\) . This can be seen as an inferred "mixed strategy" of player II, against which row \(i\) is a best response; in general, this inferred mixed strategy depends on \(i\) . In the same way, the incentive constraints (12.5) for player II are about each column \(j\) of \(z\) , which describes (up to normalization) an inferred mixed strategy of player I, against which column \(j\) is a best response. If \(z\) is obtained from a Nash equilibrium via \({z}_{ij} = {x}_{i}{y}_{j}\) , then these inferred mixed strategies of the opponent are always the same, namely the actual mixed strategies.

\begin{tcolorbox}
玩家 \(\mathrm{I}\) 的激励约束(12.4)与矩阵 \(z\) 的每一行 \(i\) 有关。直至一个归一化因子(如(12.3)中所示)， \(z\) 的这一行是玩家 I 在得到选择 \(i\) 的建议后，对玩家 II 行为的“后验”。这可以看作是玩家 II 的一个推断“混合策略”，而 \(i\) 行是针对该策略的最优反应；一般来说，这个推断的混合策略依赖于 \(i\) 。同样地，玩家 II 的激励约束(12.5)与 \(z\) 的每一列 \(j\) 有关，这一列(直至归一化)描述了玩家 I 的一个推断混合策略，而 \(j\) 列是针对该策略的最优反应。如果 \(z\) 是通过 \({z}_{ij} = {x}_{i}{y}_{j}\) 从纳什均衡(Nash equilibrium)中得到的，那么对手的这些推断混合策略总是相同的，即实际的混合策略。
\end{tcolorbox}

Proposition 12.2. Let(A, B)be an \(m \times  n\) bimatrix game, let \(x\) be a mixed strategy of player \(I\) and \(y\) be a mixed strategy of player \({II}\) , and \(z \in  {\mathbb{R}}^{m \times  n}\) be defined by \({z}_{ij} = {x}_{i}{y}_{j}\) for all \(i\) and \(j\) . Then \(z\) is a correlated equilibrium of(A, B)if and only if(x, y)is a Nash equilibrium of(A, B).

\begin{tcolorbox}
命题 12.2。设 (A, B) 是一个 \(m \times  n\) 双矩阵博弈(bimatrix game)，设 \(x\) 是玩家 \(I\) 的一个混合策略(mixed strategy)， \(y\) 是玩家 \({II}\) 的一个混合策略，并且对于所有的 \(i\) 和 \(j\) ，由 \({z}_{ij} = {x}_{i}{y}_{j}\) 定义 \(z \in  {\mathbb{R}}^{m \times  n}\) 。那么 \(z\) 是 (A, B) 的一个相关均衡(correlated equilibrium)，当且仅当 (x, y) 是 (A, B) 的一个纳什均衡。
\end{tcolorbox}

Proof. Clearly, \(\mathop{\sum }\limits_{{i = 1}}^{m}{x}_{i} = 1\) and \(\mathop{\sum }\limits_{{j = 1}}^{n}{y}_{j} = 1\) imply that \(z\) is a probability distribution on pairs(i, j)as stated in (12.1). Because \(\mathop{\sum }\limits_{{l = 1}}^{n}{y}_{l} = 1\) and \({z}_{il} = {x}_{i}{y}_{l}\) , we have \(\mathop{\sum }\limits_{{l = 1}}^{n}{z}_{il} = {x}_{i}\) for \(1 \leq  i \leq  m\) .

\begin{tcolorbox}
证明。显然， \(\mathop{\sum }\limits_{{i = 1}}^{m}{x}_{i} = 1\) 和 \(\mathop{\sum }\limits_{{j = 1}}^{n}{y}_{j} = 1\) 意味着 \(z\) 是如(12.1)中所述的关于对 (i, j) 的一个概率分布。因为 \(\mathop{\sum }\limits_{{l = 1}}^{n}{y}_{l} = 1\) 和 \({z}_{il} = {x}_{i}{y}_{l}\) ，我们有对于 \(1 \leq  i \leq  m\) 的 \(\mathop{\sum }\limits_{{l = 1}}^{n}{z}_{il} = {x}_{i}\) 。
\end{tcolorbox}

Suppose \(z\) is a correlated equilibrium. If \({x}_{i} > 0\) , then (12.4) implies (12.3), which states that \(\mathop{\sum }\limits_{{j = 1}}{a}_{ij}{y}_{j} \geq  \mathop{\sum }\limits_{{j = 1}}{a}_{kj}{y}_{j}\) , that is, the expected payoff to player \(\mathrm{I}\) in row \(i\) is at least as large as in any other row \(k\) . According to the best-response condition (6.11), this means that \(x\) is a best response to \(y\) . Similarly,(12.5) implies that \(y\) is a best response to \(x\) . Hence,(x, y)is a Nash equilibrium.

\begin{tcolorbox}
假设 \(z\) 是一个相关均衡。如果 \({x}_{i} > 0\) ，那么(12.4)意味着(12.3)，即 \(\mathop{\sum }\limits_{{j = 1}}{a}_{ij}{y}_{j} \geq  \mathop{\sum }\limits_{{j = 1}}{a}_{kj}{y}_{j}\) ，也就是说，玩家 \(\mathrm{I}\) 在 \(i\) 行的期望效用(expected payoff)至少和在任何其他行 \(k\) 的一样大。根据最优反应条件(best - response condition)(6.11)，这意味着 \(x\) 是对 \(y\) 的最优反应。类似地，(12.5)意味着 \(y\) 是对 \(x\) 的最优反应。因此，(x, y) 是一个纳什均衡。
\end{tcolorbox}

If(x, y)is a Nash equilibrium, these implications hold in reverse: If \({x}_{i} = 0\) , (12.4) holds trivially, and if \({x}_{i} > 0\) , the fact that the expected payoff in row \(i\) is at least the expected payoff in any other row \(k\) implies (12.3) and therefore (12.4). Similarly, \(y\) being a best response to \(x\) is equivalent to (12.5). Therefore, the incentive constraints hold for both players, and \(z\) is a correlated equilibrium.

\begin{tcolorbox}
如果 (x, y) 是一个纳什均衡，这些蕴含关系反过来成立:如果 \({x}_{i} = 0\) ，(12.4)显然成立，并且如果 \({x}_{i} > 0\) ， \(i\) 行的期望效用至少是任何其他行 \(k\) 的期望效用这一事实意味着(12.3)，从而意味着(12.4)。类似地， \(y\) 是对 \(x\) 的最优反应等价于(12.5)。因此，两个玩家的激励约束都成立，并且 \(z\) 是一个相关均衡。
\end{tcolorbox}

The incentive constraints (12.4) and (12.5) involve separately the two payoff matrices \(A\) and \(B\) but affect together the matrix \(z\) . There are \(m\left( {m - 1}\right)\) nontrivial constraints in (12.4) (for \(i = k\) they are trivial) and \(n\left( {n - 1}\right)\) nontrivial constraints in (12.5). For the game in Figure 12.1, they state

\begin{tcolorbox}
激励约束(12.4)和(12.5)分别涉及两个收益矩阵 \(A\) 和 \(B\) ，但共同影响矩阵 \(z\) 。(12.4)中有 \(m\left( {m - 1}\right)\) 个非平凡约束(对于 \(i = k\) 它们是平凡的)，(12.5)中有 \(n\left( {n - 1}\right)\) 个非平凡约束。对于图 12.1 中的博弈，它们表述为
\end{tcolorbox}

\[
- {z}_{11} + {z}_{12} \geq  0,\;{z}_{21} - {z}_{22} \geq  0, \tag{12.6}
\]

\[
- {z}_{11} + {z}_{21} \geq  0,\;{z}_{12} - {z}_{22} \geq  0,
\]

that is, both \({z}_{12}\) and \({z}_{21}\) have to be at least as large as \({z}_{11}\) and \({z}_{22}\) . The set of all joint distributions \(z\) on pure strategy pairs as in (12.1) is a simplex of dimension \({mn} - 1\) , here a tetrahedron of dimension 3 . The additional incentive constraints (12.4) and (12.5) define the set of all correlated equilibria, which is a subset of that simplex and a polytope.

\begin{tcolorbox}
也就是说， \({z}_{12}\) 和 \({z}_{21}\) 都必须至少与 \({z}_{11}\) 和 \({z}_{22}\) 一样大。如 (12.1) 所示，纯策略对的所有联合分布 \(z\) 的集合是一个维度为 \({mn} - 1\) 的单纯形，这里是一个三维的四面体。额外的激励约束 (12.4) 和 (12.5) 定义了所有相关均衡(correlated equilibrium)的集合，它是该单纯形的一个子集，且是一个有界多面体(polytope)。
\end{tcolorbox}

For our \(2 \times  2\) game with incentive constraints (12.6), that polytope is shown in Figure 12.3. It has five vertices, including \(\left( {{z}_{11},{z}_{12},{z}_{21},{z}_{22}}\right)  = \left( {0,0,1,0}\right) ,\left( {0,1,0,0}\right)\) , and \(\left( {\frac{1}{4},\frac{1}{4},\frac{1}{4},\frac{1}{4}}\right)\) . These are the pure Nash equilibria and the mixed Nash equilibrium in Figure 12.1, marked in Figure 12.3 as \({BL},{TR}\) , and \(M\) . The two other vertices of the correlated equilibrium set are \(\left( {0,\frac{1}{3},\frac{1}{3},\frac{1}{3}}\right)\) and \(\left( {\frac{1}{3},\frac{1}{3},\frac{1}{3},0}\right)\) , which are the two middle panels in Figure 12.2 (its left panel \(\left( {0,\frac{1}{2},\frac{1}{2},0}\right)\) is a convex combination of (0,0,1,0)and(0,1,0,0)and not a vertex). Any such vertex is recognized by having at least three of the inequalities in (12.1) and (12.6) holding as equalities, which then have a unique solution.

\begin{tcolorbox}
对于我们具有激励约束 (12.6) 的 \(2 \times  2\) 博弈，该有界多面体如图 12.3 所示。它有五个顶点，包括 \(\left( {{z}_{11},{z}_{12},{z}_{21},{z}_{22}}\right)  = \left( {0,0,1,0}\right) ,\left( {0,1,0,0}\right)\) 和 \(\left( {\frac{1}{4},\frac{1}{4},\frac{1}{4},\frac{1}{4}}\right)\) 。这些是图 12.1 中的纯纳什均衡(Nash equilibrium)和混合纳什均衡，在图 12.3 中标记为 \({BL},{TR}\) 和 \(M\) 。相关均衡集的另外两个顶点是 \(\left( {0,\frac{1}{3},\frac{1}{3},\frac{1}{3}}\right)\) 和 \(\left( {\frac{1}{3},\frac{1}{3},\frac{1}{3},0}\right)\) ，它们是图 12.2 中的两个中间面板(其左面板 \(\left( {0,\frac{1}{2},\frac{1}{2},0}\right)\) 是 (0,0,1,0) 和 (0,1,0,0) 的凸组合(convex combination)，不是顶点)。任何这样的顶点都可以通过 (12.1) 和 (12.6) 中至少有三个不等式取等号来识别，此时有唯一解。
\end{tcolorbox}

Correlated equilibrium has an analogous definition for a game with \(N\) players. Recall from Definition 3.1 that a strategy profile \(s = \left( {{s}_{1},\ldots ,{s}_{N}}\right)\) in an \(N\) -player game is also written as \(\left( {{s}_{i},{s}_{-i}}\right)\) with strategy \({s}_{i}\) of player \(i\) and the partial profile \({s}_{-i}\) of strategies of the remaining players, where the set of these partial profiles is \({S}_{-i} = { \times  }_{j = 1,j \neq  i}^{n}{S}_{j}.\)

\begin{tcolorbox}
对于有 \(N\) 个参与者的博弈，相关均衡(correlated equilibrium)有类似的定义。回顾定义 3.1，在一个 \(N\) 参与者博弈中，一个策略组合 \(s = \left( {{s}_{1},\ldots ,{s}_{N}}\right)\) 也可以写成 \(\left( {{s}_{i},{s}_{-i}}\right)\) ，其中包含参与者 \(i\) 的策略 \({s}_{i}\) 以及其余参与者的部分策略组合 \({s}_{-i}\) ，这些部分策略组合的集合是 \({S}_{-i} = { \times  }_{j = 1,j \neq  i}^{n}{S}_{j}.\)
\end{tcolorbox}

Definition 12.3 (Correlated equilibrium, \(N\) players). Let \(G\) be an \(N\) -player game in strategic form with finite strategy set \({S}_{i}\) for player \(i = 1,\ldots ,N\) and payoff \({u}_{i}\left( s\right)\) for each strategy profile \(s = \left( {{s}_{1},\ldots ,{s}_{N}}\right)\) in \(S = {S}_{1} \times  \cdots  \times  {S}_{N}\) . A correlation device is a probability distribution \(z\) on strategy profiles, that is, a function \(z : S \rightarrow  \mathbb{R}\) so that

\begin{tcolorbox}
定义 12.3( \(N\) 个参与者的相关均衡(correlated equilibrium))。设 \(G\) 是一个策略形式(strategic form)的 \(N\) 参与者博弈，参与者 \(i = 1,\ldots ,N\) 的有限策略集为 \({S}_{i}\) ，对于 \(S = {S}_{1} \times  \cdots  \times  {S}_{N}\) 中的每个策略组合 \(s = \left( {{s}_{1},\ldots ,{s}_{N}}\right)\) ，其收益为 \({u}_{i}\left( s\right)\) 。一个相关机制(correlation device)是策略组合上的一个概率分布 \(z\) ，即一个函数 \(z : S \rightarrow  \mathbb{R}\) ，使得
\end{tcolorbox}

\[
z\left( s\right)  \geq  0\;\left( {s \in  S}\right) ,\;\mathop{\sum }\limits_{{s \in  S}}z\left( s\right)  = 1. \tag{12.7}
\]

\begin{center}
\includegraphics[max width=0.9\textwidth]{images/bo_d2e5cmbef24c73b03tn0_351_329_195_1267_1032_0.jpg}
\end{center}
\hspace*{3em} 

Figure 12.3 Two perspectives of the simplex of joint distributions \(z\) for the game in Figure 12.1, with \({TL},{TR},{BL},{BR}\) (standing for the four pure-strategy pairs) and the mixed Nash equilibrium \(M\) . The left pictures show the uncorrelated strategies defined by \({z}_{ij} = {x}_{i}{y}_{j}\) , the middle pictures the polytope of correlated equilibria defined by the additional incentive constraints (12.6), and the right pictures both sets, whose intersection defines the Nash equilibria \({BL},{TR},M\) , as stated in Proposition 12.2.

\begin{tcolorbox}
图 12.3 展示了图 12.1 中博弈的联合分布 \(z\) 单纯形的两个视角，其中有 \({TL},{TR},{BL},{BR}\) (代表四个纯策略对)和混合纳什均衡 \(M\) 。左图显示了由 \({z}_{ij} = {x}_{i}{y}_{j}\) 定义的非相关策略(uncorrelated strategies)，中图显示了由额外激励约束 (12.6) 定义的相关均衡(correlated equilibrium)的有界多面体，右图同时显示了这两个集合，它们的交集定义了纳什均衡(Nash equilibrium) \({BL},{TR},M\) ，如命题 12.2 所述。
\end{tcolorbox}

The players know \(z\) . The device selects the strategy profile \(s\) with probability \(z\left( s\right)\) and then privately tells each player \(i\) the respective strategy \({s}_{i}\) in \(s\) . If for each player \(i = 1,\ldots ,N\) the incentive constraints

\begin{tcolorbox}
玩家们知道 \(z\) 。相关机制以概率 \(z\left( s\right)\) 选择策略组合 \(s\) ，然后私下告知每个玩家 \(i\) 在 \(s\) 中各自的策略 \({s}_{i}\) 。如果对于每个玩家 \(i = 1,\ldots ,N\) ，激励约束
\end{tcolorbox}

\[
\mathop{\sum }\limits_{{{s}_{-i} \in  {S}_{-i}}}\left( {{u}_{i}\left( {{s}_{i},{s}_{-i}}\right)  - {u}_{i}\left( {{t}_{i},{s}_{-i}}\right) }\right) z\left( {{s}_{i},{s}_{-i}}\right)  \geq  0\;\left( {{s}_{i},{t}_{i} \in  {S}_{i}}\right)  \tag{12.8}
\]

hold, then \(z\) is called a correlated equilibrium of \(G\) .

\begin{tcolorbox}
成立，那么 \(z\) 被称为 \(G\) 的相关均衡。
\end{tcolorbox}

The incentive constraints (12.8) generalize those in (12.4) and (12.5) for two players. In (12.8), player \(i\) is told to play \({s}_{i}\) , and can infer from this information the absolute probabilities \(z\left( {{s}_{i},{s}_{-i}}\right)\) for all possible partial profiles \({s}_{-i}\) of the other players. The sum of these probabilities is positive because otherwise player \(i\) could not have been told \({s}_{i}\) . Division by this sum then gives the conditional probabilities of what the other players have been told, and similar to (12.3) the inequality expresses that player \(i\) has nothing to gain be deviating from \({s}_{i}\) to any other strategy \({t}_{i}\) . This is the equilibrium property.

\begin{tcolorbox}
激励约束(12.8)推广了两个玩家情形下(12.4)和(12.5)中的约束。在(12.8)中，玩家 \(i\) 被要求采用策略 \({s}_{i}\) ，并能从这个信息推断出其他玩家所有可能的部分策略组合 \({s}_{-i}\) 的绝对概率 \(z\left( {{s}_{i},{s}_{-i}}\right)\) 。这些概率之和为正，因为否则玩家 \(i\) 不可能被要求采用 \({s}_{i}\) 。用这个和去除这些概率，就得到了其他玩家被要求采用的策略的条件概率，与(12.3)类似，该不等式表明玩家 \(i\) 从 \({s}_{i}\) 偏离到任何其他策略 \({t}_{i}\) 都无利可图。这就是均衡性质。
\end{tcolorbox}

Proposition 12.2 holds analogously for \(N\) players.

\begin{tcolorbox}
命题12.2对于 \(N\) 个玩家同样成立。
\end{tcolorbox}

Proposition 12.4. Consider an \(N\) -player game \(G\) in strategic form as in Definition 12.3 and a mixed strategy \({\sigma }_{i}\) for each player \(i\) . Suppose \(z\) is a correlation device so that \(z\left( s\right)  = \mathop{\prod }\limits_{{i = 1}}^{N}{\sigma }_{i}\left( {s}_{i}\right)\) for all \(s = \left( {{s}_{1},\ldots ,{s}_{N}}\right)  \in  S\) . Then \(z\) is a correlated equilibrium of \(G\) if and only if \(\left( {{\sigma }_{1},\ldots ,{\sigma }_{N}}\right)\) is a Nash equilibrium of \(G\) .

\begin{tcolorbox}
命题12.4。考虑定义12.3中策略形式的 \(N\) 人博弈 \(G\) ，以及每个玩家 \(i\) 的一个混合策略 \({\sigma }_{i}\) 。假设 \(z\) 是一个相关机制，使得对于所有 \(s = \left( {{s}_{1},\ldots ,{s}_{N}}\right)  \in  S\) 有 \(z\left( s\right)  = \mathop{\prod }\limits_{{i = 1}}^{N}{\sigma }_{i}\left( {s}_{i}\right)\) 。那么 \(z\) 是 \(G\) 的相关均衡，当且仅当 \(\left( {{\sigma }_{1},\ldots ,{\sigma }_{N}}\right)\) 是 \(G\) 的纳什均衡。
\end{tcolorbox}

Proof. The proof of Proposition 12.2 carries over by applying the best-response condition for \(N\) players (see Proposition 6.3).

\begin{tcolorbox}
证明。通过应用 \(N\) 个玩家的最优反应条件(见命题6.3)，命题12.2的证明可以照搬过来。
\end{tcolorbox}

A correlated equilibrium \(z\) of the game \(G\) as in Definition 12.3 is defined by its probabilities \(z\left( s\right)\) for strategy profiles so that the incentive constraints (12.8) hold. These constraints, and the conditions (12.7) that state that \(z \in  \Delta \left( S\right)\) , are linear in the variables \(z\left( s\right)\) for \(s \in  S\) . Hence, the set of correlated equilibria of \(G\) is a polyhedron, and as a subset of the simplex \(\Delta \left( S\right)\) bounded and therefore a polytope. Mathematically, this is a simpler object than the set of Nash equilibria of \(G\) , which is in general not convex. However, so far the only proof that an \(N\) -player game \(G\) has a correlated equilibrium is via the existence of a Nash equilibrium, which relies on Brouwer's fixed-point theorem. An elegant alternative existence proof of correlated equilibria that uses von Neumann's minimax theorem will be given in Section 12.5.

\begin{tcolorbox}
定义12.3中博弈 \(G\) 的相关均衡 \(z\) 由其策略组合的概率 \(z\left( s\right)\) 定义，使得激励约束(12.8)成立。这些约束，以及表明 \(z \in  \Delta \left( S\right)\) 的条件(12.7)，对于 \(s \in  S\) 的变量 \(z\left( s\right)\) 是线性的。因此， \(G\) 的相关均衡集是一个多面体，并且作为单纯形 \(\Delta \left( S\right)\) 的子集是有界的，所以是一个有界多面体。在数学上，这比 \(G\) 的纳什均衡集更简单，后者一般不是凸集。然而，到目前为止， \(N\) 人博弈 \(G\) 存在相关均衡的唯一证明是通过纳什均衡的存在性，这依赖于布劳威尔不动点定理。第12.5节将给出一个使用冯·诺伊曼极小极大定理的、关于相关均衡存在性的优雅替代证明。
\end{tcolorbox}

\section*{12.4 Coarse Correlated Equilibrium}

\begin{tcolorbox}
\section*{12.4 粗略相关均衡}
\end{tcolorbox}

Correlated equilibrium needs a correlation device. For example, a traffic light sends correlated signals to the different road users that are recommendations for actions. The equilibrium condition is that it is optimal to follow these recommendations, assuming that the other players do so.

\begin{tcolorbox}
相关均衡(correlated equilibrium)需要一个相关机制(correlation device)。例如，交通信号灯向不同的道路使用者发送相关信号，这些信号是对行动的建议。均衡条件是，假设其他参与者都遵循这些建议，那么遵循这些建议是最优选择。
\end{tcolorbox}

In this section, we consider a concept that goes one step further. The correlation device not only recommends the action, but each player commits to following the recommended action (for example, by letting the device choose the action, like in a self-driving car). In order to make this a game, the player's choice comes at an earlier stage, namely whether to participate in the usage of this device.

\begin{tcolorbox}
在本节中，我们考虑一个更进一步的概念。相关机制不仅建议行动，而且每个参与者都承诺遵循所建议的行动(例如，让该机制选择行动，就像在自动驾驶汽车中那样)。为了将其构成一个博弈，参与者的选择发生在更早的阶段，即是否参与使用这个机制。
\end{tcolorbox}

Suppose the correlation device is a distribution \(z \in  \Delta \left( S\right)\) on strategy profiles for a game \(G\) as in Definition 12.3. If all players participate, then some profile \(s\) will be chosen by \(z\) and used by the players. If player \(i\) does not participate, he only knows the marginal probability distribution on the partial profiles \({s}_{-i}\) of the other players given by \(z\) . Then player \(i\) can choose his own optimal action \({s}_{i} \in  {S}_{i}\) , but does not know what action the device would have chosen. If against the marginal distribution he cannot gain by participating in the scheme, and this applies to every player \(i\) , then \(z\) is called a coarse correlated equilibrium.

\begin{tcolorbox}
假设相关机制是如定义12.3中所述的博弈 \(G\) 的策略组合上的一个分布 \(z \in  \Delta \left( S\right)\) 。如果所有参与者都参与，那么某个组合 \(s\) 将由 \(z\) 选择并被参与者使用。如果参与者 \(i\) 不参与，他只知道由 \(z\) 给出的其他参与者的部分组合 \({s}_{-i}\) 的边际概率分布。然后参与者 \(i\) 可以选择他自己的最优行动 \({s}_{i} \in  {S}_{i}\) ，但不知道该机制会选择什么行动。如果相对于边际分布，他参与该方案无法获得收益，并且这适用于每个参与者 \(i\) ，那么 \(z\) 就被称为粗相关均衡(coarse correlated equilibrium)。
\end{tcolorbox}

\begin{center}
\includegraphics[max width=0.8\textwidth]{images/bo_d2e5cmbef24c73b03tn0_353_419_643_1093_490_0.jpg}
\end{center}
\hspace*{3em} 

Figure 12.4 Left: \(3 \times  3\) game which is dominance solvable, with unique Nash and correlated equilibrium(T, T). Right: Distribution on strategy pairs that is a coarse correlated equilibrium with higher payoff than in the Nash equilibrium.

\begin{tcolorbox}
图12.4 左:占优可解(dominance solvable)的 \(3 \times  3\) 博弈，具有唯一的纳什均衡(Nash equilibrium)和相关均衡(T, T)。右:策略对的分布，它是一个粗相关均衡，其收益高于纳什均衡。
\end{tcolorbox}

Figure 12.4 shows an example for a symmetric \(3 \times  3\) game. The game is dominance solvable (see Proposition 3.4): Strategy \(B\) is strictly dominated by strategy \(T\) (for both players), and after eliminating \(B\) , strategy \(M\) is strictly dominated in the remaining \(2 \times  2\) game by \(T\) , so that(T, T)is the unique Nash equilibrium of the game. It is also the unique correlated equilibrium, by the following proposition.

\begin{tcolorbox}
图12.4展示了一个对称的 \(3 \times  3\) 博弈的例子。该博弈是占优可解的(见命题3.4):策略 \(B\) 被策略 \(T\) 严格占优(对两个参与者都是如此)，并且在剔除 \(B\) 之后，在剩余的 \(2 \times  2\) 博弈中，策略 \(M\) 被 \(T\) 严格占优，因此(T, T)是该博弈的唯一纳什均衡。根据以下命题，它也是唯一的相关均衡。
\end{tcolorbox}

Proposition 12.5. Consider an \(N\) -player game where the pure strategy \({t}_{i}\) of player \(i\) strictly dominates his pure strategy \({s}_{i}\) . Then \({s}_{i}\) is never recommended to player \(i\) in any correlated equilibrium \(z\) of the game.

\begin{tcolorbox}
命题12.5. 考虑一个 \(N\) 参与者博弈，其中参与者 \(i\) 的纯策略 \({t}_{i}\) 严格占优他的纯策略 \({s}_{i}\) 。那么在该博弈的任何相关均衡 \(z\) 中，策略 \({s}_{i}\) 永远不会被推荐给参与者 \(i\) 。
\end{tcolorbox}

Proof. The strict domination of \({s}_{i}\) by \({t}_{i}\) is stated in (3.4), which is equivalent to \({u}_{i}\left( {{s}_{i},{s}_{-i}}\right)  - {u}_{i}\left( {{t}_{i},{s}_{-i}}\right)\) in (12.8) being negative for all \({s}_{-i} \in  {S}_{-i}\) . Hence, the inequality (12.8) can only hold if \(z\left( {{s}_{i},{s}_{-i}}\right)  = 0\) for all \({s}_{-i}\) . This means that \({s}_{i}\) is never part of any strategy profile \(s\) chosen by \(z\) , and is therefore never told to player \(i\) by \(z\) .

\begin{tcolorbox}
证明。 \({t}_{i}\) 对 \({s}_{i}\) 的严格占优在(3.4)中给出，这等价于(12.8)中的 \({u}_{i}\left( {{s}_{i},{s}_{-i}}\right)  - {u}_{i}\left( {{t}_{i},{s}_{-i}}\right)\) 对所有 \({s}_{-i} \in  {S}_{-i}\) 都是负的。因此，不等式(12.8)只有在对所有 \({s}_{-i}\) 都有 \(z\left( {{s}_{i},{s}_{-i}}\right)  = 0\) 时才成立。这意味着 \({s}_{i}\) 永远不会是 \(z\) 选择的任何策略组合 \(s\) 的一部分，因此 \(z\) 永远不会告诉参与者 \(i\) 选择 \({s}_{i}\) 。
\end{tcolorbox}

Consider now the correlation device shown on the right in Figure 12.4, which selects the strategy profiles(M, B)and(B, M)with probability \(\frac{1}{2}\) each. It defines a coarse correlated equilibrium: The expected payoff to each player is \(\frac{7}{2}\) if they commit to playing the action chosen by the device. Alternatively, a player may choose to opt out from using the device, and then, assuming the other player uses the device, only knows that the other player will choose \(M\) and \(B\) with probability \(\frac{1}{2}\) each. The respective payoffs to the player who opted out when playing \(T,M\) , or \(B\) are \(\frac{5}{2},\frac{5}{2}\) and 1, respectively, which are lower, so it is optimal for the player to commit to using the device. Moreover, each player's payoff is higher than in the unique Nash and correlated equilibrium.

\begin{tcolorbox}
现在考虑图12.4右侧所示的相关机制，它以概率 \(\frac{1}{2}\) 分别选择策略组合(M, B)和(B, M)。它定义了一个粗相关均衡:如果每个参与者都承诺执行该机制选择的行动，那么每个参与者的期望效用(expected payoff)是 \(\frac{7}{2}\) 。或者，参与者可以选择不使用该机制，然后，假设另一个参与者使用该机制，他只知道另一个参与者将以概率 \(\frac{1}{2}\) 分别选择 \(M\) 和 \(B\) 。选择不使用该机制的参与者在选择 \(T,M\) 或 \(B\) 时的相应收益分别是 \(\frac{5}{2},\frac{5}{2}\) 和1，这两个收益都更低，因此参与者承诺使用该机制是最优的。此外，每个参与者的收益都高于唯一的纳什均衡和相关均衡。
\end{tcolorbox}

The following is a formal definition of coarse correlated equilibrium

\begin{tcolorbox}
以下是粗相关均衡的正式定义
\end{tcolorbox}

Definition 12.6. Consider an \(N\) -player game, with the notation of Definition 12.3. Consider a correlation device \(z\) , known to the players, where each player has in advance the choice to commit to the pure strategy told to him by the device from the profile selected using \(z\) , or alternatively to opt out from this and choose any strategy without knowledge of the selection by the device. This is called a coarse correlated equilibrium if every player chooses to commit to the selection by the device and has no incentive to unilaterally deviate from this.

\begin{tcolorbox}
定义12.6。考虑一个 \(N\) 人博弈，采用定义12.3中的符号。考虑一个玩家们都知晓的相关机制 \(z\) ，每个玩家事先都可以选择遵循该机制从使用 \(z\) 选出的策略组中告知他的纯策略，或者选择不遵循该策略，在不知道机制选择的情况下选择任何策略。如果每个玩家都选择遵循该机制的选择，并且没有单方面偏离的动机，那么这被称为一个粗相关均衡(coarse correlated equilibrium)。
\end{tcolorbox}

The following states that coarse correlated equilibrium is more general than correlated equilibrium, and also gives in (12.9) the constraints that characterize coarse correlated equilibria.

\begin{tcolorbox}
以下内容表明粗相关均衡比相关均衡(correlated equilibrium)更具一般性，并且在式(12.9)中给出了刻画粗相关均衡的约束条件。
\end{tcolorbox}

Proposition 12.7. Every correlated equilibrium is a coarse correlated equilibrium.

\begin{tcolorbox}
命题12.7。每个相关均衡都是一个粗相关均衡。
\end{tcolorbox}

Proof. Consider a coarse correlated equilibrium \(z\) . If each player commits to playing the action that is selected by \(z\) , then player \(i\) ’s expected payoff is \(\mathop{\sum }\limits_{{s \in  S}}{u}_{i}\left( s\right) z\left( s\right)\) . If player \(i\) opts out, then player \(i\) only knows that every partial profile \({s}_{-i}\) is played with probability \(\mathop{\sum }\limits_{{{s}_{i} \in  {S}_{i}}}z\left( {{s}_{i},{s}_{-i}}\right)\) if the other players commit to playing the action that is selected by \(z\) . Hence, the constraints for a coarse correlated equilibrium state that for each player \(i\) and each \({t}_{i} \in  {S}_{i}\) , player \(i\) does not get a higher payoff when playing \({t}_{i}\) , that is,

\begin{tcolorbox}
证明。考虑一个粗相关均衡 \(z\) 。如果每个玩家都承诺执行 \(z\) 所选择的行动，那么玩家 \(i\) 的期望效用(expected payoff)是 \(\mathop{\sum }\limits_{{s \in  S}}{u}_{i}\left( s\right) z\left( s\right)\) 。如果玩家 \(i\) 选择不遵循，那么玩家 \(i\) 只知道，如果其他玩家承诺执行 \(z\) 所选择的行动，那么每个部分策略组 \({s}_{-i}\) 以概率 \(\mathop{\sum }\limits_{{{s}_{i} \in  {S}_{i}}}z\left( {{s}_{i},{s}_{-i}}\right)\) 被执行。因此，粗相关均衡的约束条件表明，对于每个玩家 \(i\) 和每个 \({t}_{i} \in  {S}_{i}\) ，玩家 \(i\) 在执行 \({t}_{i}\) 时不会获得更高的收益，即
\end{tcolorbox}

\[
\mathop{\sum }\limits_{{s \in  S}}{u}_{i}\left( s\right) z\left( s\right)  = \mathop{\sum }\limits_{{{s}_{-i} \in  {S}_{-i}}}\mathop{\sum }\limits_{{{s}_{i} \in  {S}_{i}}}{u}_{i}\left( {{s}_{i},{s}_{-i}}\right) z\left( {{s}_{i},{s}_{-i}}\right)  \tag{12.9}
\]

\[
\geq  \mathop{\sum }\limits_{{{s}_{-i} \in  {S}_{-i}}}{u}_{i}\left( {{t}_{i},{s}_{-i}}\right) \mathop{\sum }\limits_{{{s}_{i} \in  {S}_{i}}}z\left( {{s}_{i},{s}_{-i}}\right) .
\]

The incentive constraints (12.8) of a correlated equilibrium state that for all \({s}_{i}\) and

\begin{tcolorbox}
相关均衡的激励约束(incentive constraints)(12.8)表明，对于所有的 \({s}_{i}\) 和
\end{tcolorbox}

\({t}_{i}\) in \({S}_{i}\) ,

\begin{tcolorbox}
\({t}_{i}\) 属于 \({S}_{i}\) ，
\end{tcolorbox}

\[
\mathop{\sum }\limits_{{{s}_{-i} \in  {S}_{-i}}}{u}_{i}\left( {{s}_{i},{s}_{-i}}\right) z\left( {{s}_{i},{s}_{-i}}\right)  \geq  \mathop{\sum }\limits_{{{s}_{-i} \in  {S}_{-i}}}{u}_{i}\left( {{t}_{i},{s}_{-i}}\right) z\left( {{s}_{i},{s}_{-i}}\right)  \tag{12.10}
\]

which imply (12.9) by summation over all \({s}_{i} \in  {S}_{i}\) .

\begin{tcolorbox}
通过对所有 \({s}_{i} \in  {S}_{i}\) 求和，这意味着(12.9)式成立。
\end{tcolorbox}

Coarse correlated equilibria are considered further in Exercises 12.3 and 12.2.

\begin{tcolorbox}
在练习12.3和12.2中进一步探讨了粗相关均衡。
\end{tcolorbox}

\section*{12.5 Existence of a Correlated Equilibrium}

\begin{tcolorbox}
\section*{12.5 相关均衡的存在性}
\end{tcolorbox}

In this section, we prove the existence of a correlated equilibrium in an \(N\) -player game without assuming the existence of a Nash equilibrium (which relies on Brouwer's fixed-point theorem), using only the minimax Theorem 8.4.

\begin{tcolorbox}
在本节中，我们仅使用极小极大定理(minimax theorem)8.4，在不假设纳什均衡(Nash equilibrium)存在(纳什均衡的存在依赖于布劳威尔不动点定理)的情况下，证明一个 \(N\) 人博弈中相关均衡的存在性。
\end{tcolorbox}

As a warm-up, and because we will use it in the existence proof of a correlated equilibrium, we consider the concepts of a Markov chain \(Y\) on a finite set \(T\) and a stationary distribution for \(Y\) . Recall that1is the vector with all components equal to 1 . If \(T = \{ 1,\ldots ,n\}\) , then we write \(x \in  \Delta \left( T\right)\) as \(x \in  {\mathbb{R}}^{n},x \geq  \mathbf{0},{\mathbf{1}}^{\top }x = 1\) .

\begin{tcolorbox}
作为热身，并且因为我们将在相关均衡存在性的证明中使用它，我们考虑有限集 \(T\) 上的马尔可夫链(Markov chain) \(Y\) 和平稳分布(stationary distribution)的概念。回想一下，1 是所有分量都等于1的向量。如果 \(T = \{ 1,\ldots ,n\}\) ，那么我们将 \(x \in  \Delta \left( T\right)\) 写作 \(x \in  {\mathbb{R}}^{n},x \geq  \mathbf{0},{\mathbf{1}}^{\top }x = 1\) 。
\end{tcolorbox}

Definition 12.8. Let \(n \geq  1\) and \(T = \{ 1,\ldots ,n\}\) . A Markov chain on \(T\) is an \(n \times  n\) matrix \(Y\) so that \(Y \geq  0\) and \({Y1} = \mathbf{1}\) . A stationary distribution for a Markov chain \(Y\) on \(T\) is \(x \in  \Delta \left( T\right)\) so that \({x}^{\top } = {x}^{\top }Y\) .

\begin{tcolorbox}
定义12.8。设 \(n \geq  1\) 和 \(T = \{ 1,\ldots ,n\}\) 。 \(T\) 上的马尔可夫链是一个 \(n \times  n\) 矩阵 \(Y\) ，使得 \(Y \geq  0\) 且 \({Y1} = \mathbf{1}\) 。 \(T\) 上的马尔可夫链 \(Y\) 的平稳分布是 \(x \in  \Delta \left( T\right)\) ，使得 \({x}^{\top } = {x}^{\top }Y\) 。
\end{tcolorbox}

With \(T\) seen as a set of states, a Markov chain \(Y\) on \(T\) specifies probabilistic state transitions. From a state chosen as a row of \(Y\) , the next state is reached according to the distribution defined by that row. If the rows are chosen according to a stationary distribution \(x\) on \(T\) , then the condition \({x}^{\top } = {x}^{\top }Y\) means that every state has the same probability before and after the state transitions defined by \(Y\) . An example of a Markov chain \(Y\) and its stationary distribution \(x\) (see also Exercise 12.5) is

\begin{tcolorbox}
将 \(T\) 视为一组状态， \(T\) 上的马尔可夫链 \(Y\) 指定了概率性的状态转移。从 \(Y\) 的某一行所选择的状态开始，下一个状态根据该行所定义的分布来确定。如果根据 \(T\) 上的平稳分布 \(x\) 选择行，那么条件 \({x}^{\top } = {x}^{\top }Y\) 意味着在 \(Y\) 所定义的状态转移前后，每个状态的概率相同。一个马尔可夫链 \(Y\) 及其平稳分布 \(x\) 的例子(另见练习12.5)是
\end{tcolorbox}

\[
Y = \left( \begin{array}{ll} {0.2} & {0.8} \\  {0.5} & {0.5} \end{array}\right) ,\;x = \left( \begin{matrix} \frac{5}{13} \\  \frac{8}{13} \end{matrix}\right) .
\]

Theorem 12.9. Every Markov chain has a stationary distribution.

\begin{tcolorbox}
定理12.9。每个马尔可夫链都有一个平稳分布。
\end{tcolorbox}

Proof. Let \(Y\) be a Markov chain on \(T = \{ 1,\ldots ,n\}\) and \(x \in  \Delta \left( T\right)\) . Because \({x}^{\top }Y\mathbf{1} =\)  \({x}^{\top }\mathbf{1} = 1\) , the map \({x}^{\top } \mapsto  {x}^{\top }Y\) is a continuous function \(\Delta \left( T\right)  \rightarrow  \Delta \left( T\right)\) , which has a fixed point by Brouwer's fixed-point theorem. This fixed point is a stationary distribution for \(Y\) . However, we do not use the fixed-point theorem but instead the minimax Theorem 8.4 for the \(n \times  n\) zero-sum game with payoff matrix \(I - Y\) to the row player, were \(I\) is the \(n \times  n\) identity matrix.

\begin{tcolorbox}
证明。设 \(Y\) 是 \(T = \{ 1,\ldots ,n\}\) 和 \(x \in  \Delta \left( T\right)\) 上的一个马尔可夫链(Markov chain)。因为 \({x}^{\top }Y\mathbf{1} =\)  \({x}^{\top }\mathbf{1} = 1\) ，映射 \({x}^{\top } \mapsto  {x}^{\top }Y\) 是一个连续函数 \(\Delta \left( T\right)  \rightarrow  \Delta \left( T\right)\) ，根据布劳威尔不动点定理，它有一个不动点。这个不动点是 \(Y\) 的平稳分布(stationary distribution)。然而，我们不使用不动点定理，而是使用极小极大定理(minimax theorem)8.4 来处理这个具有支付矩阵 \(I - Y\) 的零和博弈(zero-sum game)，其中 \(I\) 是 \(n \times  n\) 单位矩阵(identity matrix)。
\end{tcolorbox}

We prove that the game with matrix \(I - Y\) has value zero. Let \(s \in  \Delta \left( T\right)\) be a mixed strategy of the minimizing column player in this game. Taking \(s = \mathbf{1}\frac{1}{n}\) shows \(\left( {I - Y}\right) s = \mathbf{1}\frac{1}{n} - Y\mathbf{1}\frac{1}{n} = \mathbf{0}\) , so by playing \(s\) the column player pays at most zero in each row, and therefore the value of the game is at most zero.

\begin{tcolorbox}
我们证明具有矩阵 \(I - Y\) 的博弈值为零。设 \(s \in  \Delta \left( T\right)\) 是该博弈中最小化的列玩家的混合策略(mixed strategy)。取 \(s = \mathbf{1}\frac{1}{n}\) 表明 \(\left( {I - Y}\right) s = \mathbf{1}\frac{1}{n} - Y\mathbf{1}\frac{1}{n} = \mathbf{0}\) ，因此通过采用 \(s\) ，列玩家在每一行的支付最多为零，所以该博弈的值最多为零。
\end{tcolorbox}

If the value of the game \(I - Y\) is negative, then the column player has a mixed strategy \(s\) so that \(\left( {I - Y}\right) s < \mathbf{0}\) , that is, \(s < {Ys}\) . But this cannot be the case: Because \(Y\mathbf{1} = \mathbf{1}\) , every row \({Y}_{a}\) of \(Y\) for \(a \in  T\) is a probability distribution on \(T\) , and \({Y}_{a}s\) is a weighted average of the components \({s}_{b}\) of \(s\) , where \(s < {Ys}\) would mean \({s}_{a} < {Y}_{a}s\) for all \(a \in  T\) . This certainly fails for the largest component \({s}_{a}\) of \(s\) (as proved in (6.15), for example). Hence, the value of the game is indeed zero.

\begin{tcolorbox}
如果博弈 \(I - Y\) 的值为负，那么列玩家有一个混合策略 \(s\) 使得 \(\left( {I - Y}\right) s < \mathbf{0}\) ，即 \(s < {Ys}\) 。但这不可能成立:因为 \(Y\mathbf{1} = \mathbf{1}\) ，对于 \(a \in  T\) ， \(Y\) 的每一行 \({Y}_{a}\) 是 \(T\) 上的一个概率分布，并且 \({Y}_{a}s\) 是 \(s\) 的分量 \({s}_{b}\) 的加权平均值，其中 \(s < {Ys}\) 意味着对于所有 \(a \in  T\) 都有 \({s}_{a} < {Y}_{a}s\) 。对于 \(s\) 的最大分量 \({s}_{a}\) 来说，这肯定不成立(例如，如 (6.15) 中所证明的)。因此，该博弈的值确实为零。
\end{tcolorbox}

Therefore, the row player has a max-min strategy \(x \in  \Delta \left( T\right)\) so that \({x}^{\top }\left( {I - Y}\right)  \geq  {\mathbf{0}}^{\top }\) . If the row vector \({x}^{\top }\left( {I - Y}\right)\) has some positive component, then \({x}^{\top }\left( {I - Y}\right) \mathbf{1} > 0\) , which contradicts \(\left( {I - Y}\right) \mathbf{1} = \mathbf{0}\) . Hence, \({x}^{\top }\left( {I - Y}\right)  = {\mathbf{0}}^{\top }\) , which is equivalent to \({x}^{\top } = {x}^{\top }Y\) , that is, \(x\) is a stationary distribution for \(Y\) .

\begin{tcolorbox}
因此，行玩家有一个极大极小策略(max-min strategy) \(x \in  \Delta \left( T\right)\) 使得 \({x}^{\top }\left( {I - Y}\right)  \geq  {\mathbf{0}}^{\top }\) 。如果行向量 \({x}^{\top }\left( {I - Y}\right)\) 有某个正分量，那么 \({x}^{\top }\left( {I - Y}\right) \mathbf{1} > 0\) ，这与 \(\left( {I - Y}\right) \mathbf{1} = \mathbf{0}\) 矛盾。因此， \({x}^{\top }\left( {I - Y}\right)  = {\mathbf{0}}^{\top }\) ，这等价于 \({x}^{\top } = {x}^{\top }Y\) ，即 \(x\) 是 \(Y\) 的平稳分布。
\end{tcolorbox}

The next lemma provides a key step in the proof that a finite game has a correlated equilibrium. It states that if each row of a nonnegative square matrix has a sum of at most 1, then increasing the diagonal element in that row so that the sum equals 1 as in (12.13) creates a Markov chain whose stationary distribution fulfills the equation (12.12) that we need subsequently (this works because the diagonal elements in this equation cancel out).

\begin{tcolorbox}
下一个引理为证明有限博弈存在相关均衡(correlated equilibrium)提供了关键步骤。它表明，如果一个非负方阵的每一行的和至多为 1，那么像 (12.13) 那样增加该行的对角元素使得和等于 1，就会得到一个马尔可夫链，其平稳分布满足我们后续所需的方程 (12.12)(这是可行的，因为该方程中的对角元素会相互抵消)。
\end{tcolorbox}

Lemma 12.10. Let \(n \geq  1\) and \(T = \{ 1,\ldots ,n\}\) . Let \(y\left( {a,b}\right)  \geq  0\) for \(a,b \in  T\) so that

\begin{tcolorbox}
引理 12.10。设 \(n \geq  1\) 和 \(T = \{ 1,\ldots ,n\}\) 。设 \(y\left( {a,b}\right)  \geq  0\) 对于 \(a,b \in  T\) 使得
\end{tcolorbox}

\[
\forall a \in  T : \;\mathop{\sum }\limits_{{b \in  T}}y\left( {a,b}\right)  \leq  1. \tag{12.11}
\]

Then there exists \(x \in  \Delta \left( T\right)\) so that

\begin{tcolorbox}
那么存在 \(x \in  \Delta \left( T\right)\) 使得
\end{tcolorbox}

\[
\forall c \in  T : \;x\left( c\right) \mathop{\sum }\limits_{{b \in  T}}y\left( {c,b}\right)  = \mathop{\sum }\limits_{{a \in  T}}x\left( a\right) y\left( {a,c}\right) . \tag{12.12}
\]

Define the \(n \times  n\) matrix \(Y\) with entries \({Y}_{ab}\) by

\begin{tcolorbox}
通过以下方式定义具有元素 \({Y}_{ab}\) 的 \(n \times  n\) 矩阵 \(Y\)
\end{tcolorbox}

\[
{Y}_{ab} = \left\{  \begin{array}{ll} y\left( {a,b}\right) & \text{ if }a \neq  b, \\  1 - \mathop{\sum }\limits_{{c \in  T,c \neq  a}}y\left( {a,c}\right) & \text{ if }a = b. \end{array}\right.  \tag{12.13}
\]

Then \(Y\) is a Markov chain on \(T\) , and \(x \in  \Delta \left( T\right)\) fulfills (12.12) if and only if \(x\) is a stationary distribution for \(Y\) .

\begin{tcolorbox}
那么 \(Y\) 是 \(T\) 上的一个马尔可夫链(Markov chain)，并且当且仅当 \(x\) 是 \(Y\) 的平稳分布(stationary distribution)时， \(x \in  \Delta \left( T\right)\) 满足 (12.12)。
\end{tcolorbox}

Proof. Clearly,(12.11) implies \(\mathop{\sum }\limits_{{c \in  T,c \neq  a}}y\left( {a,c}\right)  \leq  1\) and therefore \({Y}_{aa} \geq  0\) in (12.13), and thus \({Y}_{ab} \geq  0\) for all \(a,b \in  T\) . By (12.13),

\begin{tcolorbox}
证明:显然，(12.11) 意味着 \(\mathop{\sum }\limits_{{c \in  T,c \neq  a}}y\left( {a,c}\right)  \leq  1\) ，因此在 (12.13) 中意味着 \({Y}_{aa} \geq  0\) ，从而对于所有的 \(a,b \in  T\) 都有 \({Y}_{ab} \geq  0\) 。根据 (12.13)，
\end{tcolorbox}

\[
\forall a \in  T : \;\mathop{\sum }\limits_{{b \in  T}}{Y}_{ab} = 1, \tag{12.14}
\]

that is, \(Y \geq  0\) and \(Y\mathbf{1} = \mathbf{1}\) , so \(Y\) is a Markov chain on \(T\) .

\begin{tcolorbox}
即 \(Y \geq  0\) 且 \(Y\mathbf{1} = \mathbf{1}\) ，所以 \(Y\) 是 \(T\) 上的一个马尔可夫链(Markov chain)。
\end{tcolorbox}

Equation (12.12) holds irrespective of the values \(y\left( {c,c}\right)\) for \(c \in  T\) , because the term \(x\left( c\right) y\left( {c,c}\right)\) on both sides of the equation cancels out. We replace \(y\left( {c,c}\right)\) by \({Y}_{cc}\) as defined in (12.13), which does not affect the validity of (12.12). Then the numbers \(y\left( {a,b}\right)\) are just the entries \({Y}_{ab}\) of \(Y\) , and \(\mathop{\sum }\limits_{{b \in  T}}y\left( {c,b}\right)  = 1\) by (12.14), so that (12.12) is equivalent to \({x}^{\top } = {x}^{\top }Y\) , that is, \(x\) is a stationary distribution for \(Y\) . By Theorem 12.9, such a distribution exists.

\begin{tcolorbox}
方程 (12.12) 成立与否与 \(c \in  T\) 对应的 \(y\left( {c,c}\right)\) 的值无关，因为方程两边的项 \(x\left( c\right) y\left( {c,c}\right)\) 相互抵消。我们用 (12.13) 中定义的 \({Y}_{cc}\) 替换 \(y\left( {c,c}\right)\) ，这并不影响 (12.12) 的有效性。那么数字 \(y\left( {a,b}\right)\) 恰好是 \(Y\) 的元素 \({Y}_{ab}\) ，并且根据 (12.14) 有 \(\mathop{\sum }\limits_{{b \in  T}}y\left( {c,b}\right)  = 1\) ，所以 (12.12) 等价于 \({x}^{\top } = {x}^{\top }Y\) ，即 \(x\) 是 \(Y\) 的平稳分布(stationary distribution)。根据定理 12.9，这样的分布存在。
\end{tcolorbox}

Consider an \(N\) -player game \(G\) in strategic form as in Definition 12.3. We show the existence of a correlated equilibrium for \(G\) with help of an auxiliary zero-sum game derived from \(G\) that we call the deviation game. The deviation game has a row player COR (the "correlator") and a column player DEV (the "deviator"). The set of pure strategies of COR is the set \(S\) of pure strategy profiles of \(G\) . The set \(D\) of pure strategies of DEV is defined as

\begin{tcolorbox}
考虑定义 12.3 中策略形式(strategic form)的 \(N\) 人博弈 \(G\) 。我们借助从 \(G\) 导出的一个辅助零和博弈(zero - sum game)来证明 \(G\) 存在相关均衡(correlated equilibrium)，我们将这个辅助博弈称为偏离博弈。偏离博弈有一个行玩家 COR(“相关者”)和一个列玩家 DEV(“偏离者”)。COR 的纯策略集是 \(G\) 的纯策略组合集 \(S\) 。DEV 的纯策略集 \(D\) 定义为
\end{tcolorbox}

\[
D = \mathop{\bigcup }\limits_{{i = 1}}^{N}\left( {\{ i\}  \times  {S}_{i} \times  {S}_{i}}\right)  \tag{12.15}
\]

where a pure strategy \(\left( {i,a,b}\right)  \in  D\) means that DEV selects a player \(i\) and two pure strategies \(a\) and \(b\) of player \(i\) .

\begin{tcolorbox}
其中纯策略 \(\left( {i,a,b}\right)  \in  D\) 表示 DEV 选择一个玩家 \(i\) 以及玩家 \(i\) 的两个纯策略 \(a\) 和 \(b\) 。
\end{tcolorbox}

The payoffs to the row player COR in the deviation game are given in an \(\left| S\right|  \times  \left| D\right|\) matrix \(A\) with the entries

\begin{tcolorbox}
偏离博弈中行玩家 COR 的收益由一个 \(\left| S\right|  \times  \left| D\right|\) 矩阵 \(A\) 给出，其元素为
\end{tcolorbox}

\[
A\left( {s,\left( {i,a,b}\right) }\right)  = \left\{  \begin{array}{ll} {u}_{i}\left( {a,{s}_{-i}}\right)  - {u}_{i}\left( {b,{s}_{-i}}\right) & \text{ if }{s}_{i} = a, \\  0 & \text{ if }{s}_{i} \neq  a. \end{array}\right.  \tag{12.16}
\]

The interpretation is that a mixed strategy of COR is considered as a correlation device, which is altered by the strategy(i, a, b)chosen by DEV as follows: In the strategy profile \(s\) chosen by \(\mathrm{{COR}}\) , if \({s}_{i} = a\) (that is, the strategy \({s}_{i}\) in the profile \(s\) recommended to player \(i\) matches \(a\) ), then player \(i\) deviates from this recommendation \(a\) and instead plays \(b\) . In the deviation game, the payoff gain for player \(i\) when he deviates this way is paid by COR to DEV (if player \(i\) is worse off by deviating, then COR receives the difference as his payoff).

\begin{tcolorbox}
其解释是，COR 的混合策略(mixed strategy)被视为一种相关机制(correlation device)，它会被 DEV 选择的策略 (i, a, b) 按如下方式改变:在 \(\mathrm{{COR}}\) 选择的策略组合 \(s\) 中，如果 \({s}_{i} = a\) (即，在组合 \(s\) 中推荐给玩家 \(i\) 的策略 \({s}_{i}\) 与 \(a\) 匹配)，那么玩家 \(i\) 偏离这个推荐策略 \(a\) 而选择 \(b\) 。在偏离博弈中，玩家 \(i\) 以这种方式偏离时的收益增益由 COR 支付给 DEV(如果玩家 \(i\) 偏离后情况变差，那么 COR 将获得差值作为其收益)。
\end{tcolorbox}

The value of the deviation game is at most zero by letting DEV play any pure strategy of the form(i, a, a)(that is, no player in \(G\) ever deviates). In that case, according to (12.16), the payoff in row \(s\) of \(A\) is zero for every \(s \in  S\) .

\begin{tcolorbox}
通过让 DEV 采用任何形式为 (i, a, a) 的纯策略(即， \(G\) 中没有玩家偏离)，偏离博弈的值至多为零。在这种情况下，根据 (12.16)，对于每个 \(s \in  S\) ， \(A\) 的第 \(s\) 行的收益为零。
\end{tcolorbox}

The value of the deviation game equals zero if and only if the row player COR has a mixed strategy \(z \in  \Delta \left( S\right)\) so that

\begin{tcolorbox}
当且仅当行玩家 COR 有一个混合策略(mixed strategy) \(z \in  \Delta \left( S\right)\) 使得时，偏离博弈的值等于零
\end{tcolorbox}

\[
{z}^{\top }A \geq  {0}^{\top }\text{ . } \tag{12.17}
\]

Then every column(i, a, b)of the game corresponds to exactly one of the incentive constraints (12.8) with \(\left( {i,{s}_{i},{t}_{i}}\right)  = \left( {i,a,b}\right)\) . Hence, such an optimal mixed strategy \(z\) of COR defines a correlated equilibrium of \(G\) . Therefore, \(G\) has a correlated equilibrium if and only if the deviation game has value zero. This is the case, as stated in the following theorem. It has a nice intuition that we explain afterwards.

\begin{tcolorbox}
那么，该博弈的每一列 (i, a, b) 都恰好对应一个带有 \(\left( {i,{s}_{i},{t}_{i}}\right)  = \left( {i,a,b}\right)\) 的激励约束 (12.8)。因此，COR 的这种最优混合策略 \(z\) 定义了 \(G\) 的一个相关均衡(correlated equilibrium)。所以，当且仅当偏离博弈(deviation game)的值为零时， \(G\) 存在一个相关均衡。如下述定理所述，情况确实如此。我们随后会解释其中一个很好理解的直觉。
\end{tcolorbox}

Theorem 12.11. Consider an N-player game G in strategic form as in Definition 12.3. Then the deviation game between COR and DEV with the matrix \(A\) of payoffs in (12.16) to the row player COR has value zero. Any max-min strategy of COR is a correlated equilibrium of \(G\) , which shows that \(G\) has at last one correlated equilibrium.

\begin{tcolorbox}
定理 12.11。考虑定义 12.3 中策略形式(strategic form)的 N 人博弈 G。那么，COR 和 DEV 之间的偏离博弈，其支付矩阵 \(A\) 如 (12.16) 所示，行玩家 COR 的博弈值为零。COR 的任何极大极小策略(max - min strategy)都是 \(G\) 的一个相关均衡，这表明 \(G\) 至少有一个相关均衡。
\end{tcolorbox}

Proof. The value of the deviation game is at most zero. If its value is less than zero, then the column player DEV has a mixed strategy \(y \in  \Delta \left( D\right)\) so that \({Ay} < \mathbf{0}\) , which implies that \({z}^{\top }{Ay} < 0\) for every response \(z \in  \Delta \left( S\right)\) of COR. We show that this is not the case, by finding, as a function of \(y\) , a mixed strategy \(x\) of COR so that \({x}^{\top }{Ay} = 0\) . The distribution \(x\) on \(S\) will be a product distribution given by a profile of \(N\) mixed strategies \({x}_{i} \in  \Delta \left( {S}_{i}\right)\) for each player \(i = 1,\ldots ,N\) , that is, \(x\left( s\right)  = \mathop{\prod }\limits_{{i = 1}}^{n}{x}_{i}\left( {s}_{i}\right)\) . We write \(x\left( s\right)  = {x}_{i}\left( {s}_{i}\right) {x}_{-i}\left( {s}_{-i}\right)\) with \({x}_{-i}\left( {s}_{-i}\right)  = \mathop{\prod }\limits_{{k = 1,k \neq  i}}^{n}{x}_{k}\left( {s}_{k}\right)\) . It turns out that \(x\) can be chosen solely as a function of \(y\) , independently of the players’ payoffs in the original game \(G\) . For such a product distribution \(x\) , we have

\begin{tcolorbox}
证明。偏离博弈的值至多为零。如果其值小于零，那么列玩家 DEV 有一个混合策略 \(y \in  \Delta \left( D\right)\) ，使得 \({Ay} < \mathbf{0}\) ，这意味着对于 COR 的每一个响应 \(z \in  \Delta \left( S\right)\) 都有 \({z}^{\top }{Ay} < 0\) 。我们通过找到一个作为 \(y\) 的函数的 COR 的混合策略 \(x\) ，使得 \({x}^{\top }{Ay} = 0\) ，来证明情况并非如此。 \(S\) 上的分布 \(x\) 将是一个由每个玩家 \(i = 1,\ldots ,N\) 的 \(N\) 混合策略 \({x}_{i} \in  \Delta \left( {S}_{i}\right)\) 的组合所给出的乘积分布，即 \(x\left( s\right)  = \mathop{\prod }\limits_{{i = 1}}^{n}{x}_{i}\left( {s}_{i}\right)\) 。我们写成 \(x\left( s\right)  = {x}_{i}\left( {s}_{i}\right) {x}_{-i}\left( {s}_{-i}\right)\) ，其中 \({x}_{-i}\left( {s}_{-i}\right)  = \mathop{\prod }\limits_{{k = 1,k \neq  i}}^{n}{x}_{k}\left( {s}_{k}\right)\) 。结果表明， \(x\) 可以仅作为 \(y\) 的函数来选择，而与原始博弈 \(G\) 中玩家的支付无关。对于这样的乘积分布 \(x\) ，我们有
\end{tcolorbox}

\[
{x}^{\top }{Ay} = \mathop{\sum }\limits_{{\left( {i,a,b}\right)  \in  D}}\mathop{\sum }\limits_{{s \in  S}}x\left( s\right) y\left( {i,a,b}\right) A\left( {s,\left( {i,a,b}\right) }\right)
\]

\[
= \mathop{\sum }\limits_{{\left( {i,a,b}\right)  \in  D}}\mathop{\sum }\limits_{{{s}_{-i} \in  {S}_{-i}}}x\left( {a,{s}_{-i}}\right) y\left( {i,a,b}\right) \left( {{u}_{i}\left( {a,{s}_{-i}}\right)  - {u}_{i}\left( {b,{s}_{-i}}\right) }\right)
\]

\[
= \mathop{\sum }\limits_{{i = 1}}^{N}\mathop{\sum }\limits_{{a \in  {S}_{i}}}\mathop{\sum }\limits_{{b \in  {S}_{i}}}\mathop{\sum }\limits_{{{s}_{-i} \in  {S}_{-i}}}{x}_{i}\left( a\right) {x}_{-i}\left( {s}_{-i}\right) y\left( {i,a,b}\right) \left( {{u}_{i}\left( {a,{s}_{-i}}\right)  - {u}_{i}\left( {b,{s}_{-i}}\right) }\right)
\]

\[
= \mathop{\sum }\limits_{{i = 1}}^{N}\mathop{\sum }\limits_{{{s}_{-i} \in  {S}_{-i}}}{x}_{-i}\left( {s}_{-i}\right) \mathop{\sum }\limits_{{a \in  {S}_{i}}}\mathop{\sum }\limits_{{b \in  {S}_{i}}}{x}_{i}\left( a\right) y\left( {i,a,b}\right) \left( {{u}_{i}\left( {a,{s}_{-i}}\right)  - {u}_{i}\left( {b,{s}_{-i}}\right) }\right) . \tag{12.18}
\]

For player \(i\) and \({s}_{-i} \in  {S}_{-i}\) , we arrange the last double sum according to the payoffs \({u}_{i}\left( {c,{s}_{-i}}\right)\) for \(c \in  {S}_{i}\) :

\begin{tcolorbox}
对于玩家 \(i\) 和 \({s}_{-i} \in  {S}_{-i}\) ，我们根据 \(c \in  {S}_{i}\) 的支付 \({u}_{i}\left( {c,{s}_{-i}}\right)\) 来排列最后一个双重求和:
\end{tcolorbox}

\[
\mathop{\sum }\limits_{{a \in  {S}_{i}}}\mathop{\sum }\limits_{{b \in  {S}_{i}}}{x}_{i}\left( a\right) y\left( {i,a,b}\right) \left( {{u}_{i}\left( {a,{s}_{-i}}\right)  - {u}_{i}\left( {b,{s}_{-i}}\right) }\right)
\]

\[
= \mathop{\sum }\limits_{{a \in  {S}_{i}}}\mathop{\sum }\limits_{{b \in  {S}_{i}}}{x}_{i}\left( a\right) y\left( {i,a,b}\right) {u}_{i}\left( {a,{s}_{-i}}\right)  - \mathop{\sum }\limits_{{a \in  {S}_{i}}}\mathop{\sum }\limits_{{b \in  {S}_{i}}}{x}_{i}\left( a\right) y\left( {i,a,b}\right) {u}_{i}\left( {b,{s}_{-i}}\right)
\]

\[
= \mathop{\sum }\limits_{{c \in  {S}_{i}}}\mathop{\sum }\limits_{{b \in  {S}_{i}}}{x}_{i}\left( c\right) y\left( {i,c,b}\right) {u}_{i}\left( {c,{s}_{-i}}\right)  - \mathop{\sum }\limits_{{a \in  {S}_{i}}}\mathop{\sum }\limits_{{c \in  {S}_{i}}}{x}_{i}\left( a\right) y\left( {i,a,c}\right) {u}_{i}\left( {c,{s}_{-i}}\right)
\]

\[
= \mathop{\sum }\limits_{{c \in  {S}_{i}}}{u}_{i}\left( {c,{s}_{-i}}\right) \left( {{x}_{i}\left( c\right) \mathop{\sum }\limits_{{b \in  {S}_{i}}}y\left( {i,c,b}\right)  - \mathop{\sum }\limits_{{a \in  {S}_{i}}}{x}_{i}\left( a\right) y\left( {i,a,c}\right) }\right) . \tag{12.19}
\]

We now invoke Lemma 12.10 with \(T = {S}_{i}\) and \(y\left( {a,b}\right)  = y\left( {i,a,b}\right)\) : The numbers \(y\left( {a,b}\right)\) are the mixed strategy probabilities \(y\left( {i,a,b}\right)\) of DEV, which implies (12.11). Hence, there are probabilities \(x\left( c\right)  = {x}_{i}\left( c\right)  \in  \Delta \left( {S}_{i}\right)\) so that (12.12) holds, that is, the difference of the sums in (12.19) is zero for all \(c \in  {S}_{i}\) . Therefore, for any payoffs \({u}_{i}\left( {c,{s}_{-i}}\right)\) , the entire sum in (12.19) is zero, and therefore \({x}^{\top }{Ay} = 0\) in (12.18).

\begin{tcolorbox}
我们现在对 \(T = {S}_{i}\) 和 \(y\left( {a,b}\right)  = y\left( {i,a,b}\right)\) 应用引理 12.10:数字 \(y\left( {a,b}\right)\) 是 DEV 的混合策略概率 \(y\left( {i,a,b}\right)\) ，这意味着 (12.11)。因此，存在概率 \(x\left( c\right)  = {x}_{i}\left( c\right)  \in  \Delta \left( {S}_{i}\right)\) 使得 (12.12) 成立，即 (12.19) 中求和的差值对于所有 \(c \in  {S}_{i}\) 都为零。因此，对于任何支付 \({u}_{i}\left( {c,{s}_{-i}}\right)\) ，(12.19) 中的整个求和都为零，因此 (12.18) 中的 \({x}^{\top }{Ay} = 0\) 也为零。
\end{tcolorbox}

This shows that the column player DEV does not have a mixed strategy \(y\) with \({Ay} < \mathbf{0}\) . Therefore, as claimed, the deviation game has value zero. Any max-min strategy \(z \in  \Delta \left( S\right)\) of the row player COR fulfills \({z}^{\top }A \geq  {\mathbf{0}}^{\top }\) and defines a correlated equilibrium of the original game \(G\) .

\begin{tcolorbox}
这表明列玩家DEV没有满足 \({Ay} < \mathbf{0}\) 的混合策略 \(y\) 。因此，如所声称的，偏离博弈的值为零。行玩家COR的任何极大极小策略 \(z \in  \Delta \left( S\right)\) 都满足 \({z}^{\top }A \geq  {\mathbf{0}}^{\top }\) ，并定义了原始博弈 \(G\) 的一个相关均衡。
\end{tcolorbox}

The proof of Theorem 12.11 has the following intuition. Similarly to the proof of Theorem 12.9, it uses the minimax Theorem 8.4 by showing that the column player cannot simultaneously make all payoff rows negative. Consider any mixed strategy \(y\) of the column player DEV in the deviation game, and some player \(i\) in the original game \(G\) . The pure strategies \(a,b \in  {S}_{i}\) of player \(i\) are seen as states of a Markov chain, and the probabilities \(y\left( {i,a,b}\right)\) define transition probabilities from state \(a\) to state \(b\) . For each \(a\) , the sum \({p}_{a}\) of these probabilities \(y\left( {i,a,b}\right)\) over all \(b \neq  a\) is usually less than 1 . Hence, the complementary probability \(1 - {p}_{a}\) can be considered as the probability that player \(i\) does not deviate from his recommendation to play \(a\) . This interpretation agrees with the entries of the matrix \(A\) in (12.16) and with the definition of the deviation game (which in effect says that if DEV does not choose \(\left( {i,{s}_{i},{t}_{i}}\right)\) for the strategy \({s}_{i}\) in \(s\) chosen by COR, then player \(i\) does not deviate from \({s}_{i}\) ). With \({Y}_{aa} = 1 - {p}_{a}\) and \({Y}_{ab} = y\left( {i,a,b}\right)\) , this defines the Markov chain \(Y\) in Lemma 12.10. Let \(x\) be a stationary distribution for \(Y\) , which defines a mixed strategy \({x}_{i}\) for player \(i\) . Crucially, this stationary distribution neutralizes the Markov chain. In other words, the "deviation plan" \(y\) by DEV is countered, independently for each player \(i\) , by a mixed strategy \({x}_{i}\) so that there is no observable difference, in terms of probabilities, between all players playing according to the profile \(x = \left( {{x}_{1},\ldots ,{x}_{N}}\right)\) before or after applying the deviations imposed by \(y\) . Hence, the gain for deviating is zero, expressed by the equation \({x}^{\top }{Ay} = 0\) , and hence it is not possible that \({Ay} < \mathbf{0}\) .

\begin{tcolorbox}
定理12.11的证明有如下直观解释。与定理12.9的证明类似，它运用了极小极大定理(定理8.4)，通过表明列玩家无法同时使所有收益行都为负。考虑偏离博弈中列玩家DEV的任何混合策略 \(y\) ，以及原始博弈 \(G\) 中的某个玩家 \(i\) 。玩家 \(i\) 的纯策略 \(a,b \in  {S}_{i}\) 被视为马尔可夫链(Markov chain)的状态，概率 \(y\left( {i,a,b}\right)\) 定义了从状态 \(a\) 到状态 \(b\) 的转移概率。对于每个 \(a\) ，这些概率 \(y\left( {i,a,b}\right)\) 在所有 \(b \neq  a\) 上的和 \({p}_{a}\) 通常小于1。因此，互补概率 \(1 - {p}_{a}\) 可以被视为玩家 \(i\) 不偏离其被建议采用的策略 \(a\) 的概率。这种解释与(12.16)中矩阵 \(A\) 的元素以及偏离博弈的定义一致(实际上，这意味着如果DEV没有为COR在 \(s\) 中选择的策略 \({s}_{i}\) 选择 \(\left( {i,{s}_{i},{t}_{i}}\right)\) ，那么玩家 \(i\) 不会偏离 \({s}_{i}\) )。结合 \({Y}_{aa} = 1 - {p}_{a}\) 和 \({Y}_{ab} = y\left( {i,a,b}\right)\) ，这就定义了引理12.10中的马尔可夫链 \(Y\) 。设 \(x\) 是 \(Y\) 的平稳分布(stationary distribution)，它为玩家 \(i\) 定义了一个混合策略 \({x}_{i}\) 。至关重要的是，这个平稳分布使马尔可夫链达到平衡。换句话说，DEV的“偏离计划” \(y\) 会被每个玩家 \(i\) 各自的混合策略 \({x}_{i}\) 所抵消，以至于在应用 \(y\) 所施加的偏离前后，从概率角度来看，所有玩家按照策略组合 \(x = \left( {{x}_{1},\ldots ,{x}_{N}}\right)\) 进行博弈没有可观测到的差异。因此，偏离的收益为零，由方程 \({x}^{\top }{Ay} = 0\) 表示，所以不可能出现 \({Ay} < \mathbf{0}\) 的情况。
\end{tcolorbox}

The "neutralizing behavior" \(x\) is chosen as a function of \(y\) , independently of the payoffs in \(G\) . Moreover, \(x\) does not represent an optimal behavior of COR in the deviation game (otherwise, it would be a Nash equilibrium of \(G\) by Proposition 12.4, with an existence proof by the minimax theorem, which would be too good to be true). In fact, for the optimal strategy \(y\) of COR that never asks for any deviation, the Markov chain on each state set \({S}_{i}\) is trivial and has every \({x}_{i} \in  \Delta \left( {S}_{i}\right)\) as a stationary distribution.

\begin{tcolorbox}
“中和行为” \(x\) 是作为 \(y\) 的函数来选择的，与 \(G\) 中的收益无关。此外， \(x\) 并不代表COR在偏离博弈中的最优行为(否则，根据命题12.4，它将是 \(G\) 的一个纳什均衡(Nash equilibrium)，并且可以通过极小极大定理证明其存在性，这就好得令人难以置信了)。事实上，对于COR从不要求任何偏离的最优策略 \(y\) ，每个状态集 \({S}_{i}\) 上的马尔可夫链是平凡的，并且每个 \({x}_{i} \in  \Delta \left( {S}_{i}\right)\) 都是一个平稳分布。
\end{tcolorbox}

\section*{12.6 Further Reading}

\begin{tcolorbox}
\section*{12.6 进一步阅读}
\end{tcolorbox}

Correlated equilibrium was introduced by Aumann (1974). We have described its canonical form where the correlation device directly generates strategy profiles with the strategies in them as private recommendations to the players, rather than using a device that sends arbitrary messages that are interpreted by the players. The canonical form suffices because only the players' actions matter (Aumann, 1987). As argued by Aumann (1987), correlated equilibrium is the expression of "Bayesian rationality" with a "common prior" about the state of the world, where this prior is in effect the correlation device. In contrast, playing Nash equilibrium requires much stronger assumptions about commonly accepted norms of behavior; see Aumann and Brandenburger (1995) and Aumann and Dreze (2008).

\begin{tcolorbox}
相关均衡(correlated equilibrium)由奥曼(Aumann，1974)提出。我们已经描述了它的标准形式，在这种形式中，相关机制(correlation device)直接生成策略组合，并将其中的策略作为私人建议提供给参与者，而不是使用发送任意消息、由参与者自行解读的机制。标准形式就足够了，因为只有参与者的行动才是重要的(奥曼，1987)。正如奥曼(1987)所论证的，相关均衡是“贝叶斯理性”(Bayesian rationality)在对世界状态有“共同先验”情况下的表达，而这个先验实际上就是相关机制。相比之下，达成纳什均衡(Nash equilibrium)需要对普遍接受的行为规范做出更强的假设；参见奥曼和布兰登伯格(Aumann and Brandenburger，1995)以及奥曼和德雷兹(Aumann and Dreze，2008)的研究。
\end{tcolorbox}

Games similar to the modified Chicken game in Figure 12.1 have been described by Aumann (1974, p. 72) and Myerson (1991, p. 250). Nau, Gomez Canovas, and Hansen (2004) prove that Nash equilibria lie on the relative boundary of the polytope of correlated equilibria, and show a picture similar to Figure 12.3.

\begin{tcolorbox}
与图12.1中修改后的胆小鬼博弈(Chicken game)类似的博弈，已由奥曼(1974，第72页)和迈尔森(Myerson，1991，第250页)进行过描述。瑙(Nau)、戈麦斯·卡诺瓦斯(Gomez Canovas)和汉森(Hansen，2004)证明，纳什均衡位于相关均衡有界多面体(polytope)的相对边界上，并展示了一幅类似于图12.3的图。
\end{tcolorbox}

The term "coarse correlated equilibrium" is due to Young (2004). The game in Figure 12.4 is taken from Moulin and Vial (1978, p. 205), who introduced the concept and called it the "simple extension" of a correlated equilibrium.

\begin{tcolorbox}
“粗相关均衡”(coarse correlated equilibrium)这一术语由扬(Young，2004)提出。图12.4中的博弈取自穆林和维亚尔(Moulin and Vial，1978，第205页)，他们引入了这一概念，并将其称为相关均衡的“简单扩展”。
\end{tcolorbox}

The existence proof of Theorem 12.11 is due to Hart and Schmeidler (1989) and, very similarly and independently, Nau and McCardle (1990). Lemma 12.10 with its construction of the Markov chain \(Y\) from \(y\) by increasing the diagonal elements is better than scaling the rows of \(y\) , which would require a case distinction for all-zero rows, and an, easily overlooked, scaling of the stationary distribution \(x\) . A polynomial-time algorithm for computing a correlated equilibrium based on this existence proof, even for a large class of compactly specified games where the strategic form is exponentially large, is described in Papadimitriou and Roughgarden (2008). The computational complexity of Nash and correlated equilibria is discussed by Gilboa and Zemel (1989). In general, correlated equilibria can be found by linear programming and are computationally more tractable than Nash equilibria and fixed-point problems.

\begin{tcolorbox}
定理12.11的存在性证明归功于哈特和施梅德勒(Hart and Schmeidler，1989)，以及非常相似且独立完成证明的瑙和麦卡德尔(Nau and McCardle，1990)。引理12.10通过增加对角元素从 \(Y\) 构造马尔可夫链(Markov chain) \(y\) ，这比缩放 \(y\) 的行要好，因为缩放行需要对全零行进行情况区分，并且容易忽略对平稳分布(stationary distribution) \(x\) 的缩放。帕帕迪米特里欧和拉夫加登(Papadimitriou and Roughgarden，2008)描述了一种基于此存在性证明来计算相关均衡的多项式时间算法，即使对于一类策略形式呈指数级大的紧凑指定博弈也是如此。吉尔博阿和泽梅尔(Gilboa and Zemel，1989)讨论了纳什均衡和相关均衡的计算复杂性。一般来说，相关均衡可以通过线性规划(linear programming)找到，并且在计算上比纳什均衡和不动点问题更容易处理。
\end{tcolorbox}

Correlated equilibria are important because they appear as limiting behavior of players who learn their strategies in repeated plays of the game. The limit distribution on strategy profiles is thereby given by how often they are played in these repeated interactions. For zero-sum games, fictitious play, where players use their best response against the current history, converges to a Nash equilibrium (Robinson,1951). For non-zero-sum games, however, this fails already for \(3 \times  3\) games, even if the game has a unique equilibrium, such as the game in Exercise 12.2, as shown by Shapley (1963).

\begin{tcolorbox}
相关均衡很重要，因为它们表现为参与者在重复博弈中学习策略时的极限行为。策略组合上的极限分布由此由它们在这些重复交互中被采用的频率决定。对于零和博弈(zero-sum game)，虚构博弈(fictitious play)(即参与者根据当前历史使用最优反应(best response))会收敛到一个纳什均衡(罗宾逊(Robinson)，1951)。然而，对于非零和博弈，即使博弈有唯一的均衡，如练习12.2中的博弈，这种收敛也会失败，正如沙普利(Shapley，1963)所证明的。
\end{tcolorbox}

In learning based on regret-matching, a player checks for any two of his strategies whether or not he should have played one instead of the other, and adapts his behavior accordingly. Hart and Mas-Colell (2000) introduced this learning procedure and showed that it approaches the set of correlated equilibria. It has its origin in the proof of Theorem 12.11, with the idea of applying ficticious play to the auxiliary zero-sum game (Hart and Mas-Colell, 2000, p. 1142). For a survey of learning in games see Young (2004).

\begin{tcolorbox}
在基于后悔匹配(regret-matching)的学习中，参与者会检查他的任意两个策略，判断自己是否应该选择其中一个而不是另一个，并相应地调整自己的行为。哈特和马斯 - 科莱尔(Hart and Mas-Colell，2000)引入了这种学习过程，并证明它会趋近于相关均衡集。它源于定理12.11的证明，其思路是将虚构博弈应用于辅助零和博弈(哈特和马斯 - 科莱尔，2000，第1142页)。关于博弈学习的综述可参见扬(Young，2004)的研究。
\end{tcolorbox}

\section*{12.7 Exercises for Chapter 12}

\begin{tcolorbox}
\section*{12.7 第12章练习}
\end{tcolorbox}

Exercise 12.1. Find the Nash and correlated equilibria of the following zero-sum game (with payoffs to the row player):

\begin{tcolorbox}
练习12.1. 找出以下零和博弈(行参与者的收益)的纳什均衡和相关均衡:
\end{tcolorbox}

\begin{center}
\adjustbox{max width=\textwidth}{
\begin{tabular}{|c|c|c|}
\hline
II I & \(l\) & \(\gamma\) \\
\cline{1-3}
\(T\) & 4 & 2 \\
\cline{1-3}
\(B\) & 0 & 3 \\
\cline{1-3}
\hline
\end{tabular}
}
\end{center}

\begin{tcolorbox}

\begin{center}
\adjustbox{max width=\textwidth}{
\begin{tabular}{|c|c|c|}
\hline
二 一 & \(l\) & \(\gamma\) \\
\cline{1-3}
\(T\) & 4 & 2 \\
\cline{1-3}
\(B\) & 0 & 3 \\
\cline{1-3}
\hline
\end{tabular}
}
\end{center}

\end{tcolorbox}

Exercise 12.2. Consider the following variant of the Rock-Paper-Scissors game where two equal strategies tie with payoff 0 for both players, but two unequal strategies give payoff 1 to the winner and -2 to the loser.

\begin{tcolorbox}
练习12.2。考虑以下石头剪刀布游戏的变体:当两个玩家选择相同策略时，双方收益均为0；当两个玩家选择不同策略时，获胜者收益为1，失败者收益为 -2。
\end{tcolorbox}

\begin{center}
\includegraphics[max width=0.3\textwidth]{images/bo_d2e5cmbef24c73b03tn0_361_722_372_486_489_0.jpg}
\end{center}
\hspace*{3em} 

(a) Find a Nash equilibrium of this game, and its payoff to the two players. Show that when this Nash equilibrium is considered as a correlated equilibrium, all incentive constraints (12.4) and (12.5) are tight (that is, hold as equalities). This fact may suffice to indicate that the Nash and correlated equilibrium is unique; for a different proof see von Stengel and Zamir (2010, Remark 13).

\begin{tcolorbox}
(a) 找出该游戏的一个纳什均衡(Nash equilibrium)，并计算两个玩家在此均衡下的收益。证明当把这个纳什均衡视为相关均衡(correlated equilibrium)时，所有激励约束(12.4)和(12.5)都是紧的(即取等号)。这一事实可能足以表明该纳什均衡和相关均衡是唯一的；另一种证明方法见冯·施滕格尔和扎米尔(2010，注13)。
\end{tcolorbox}

(b) Find a coarse correlated equilibrium of this game that has higher expected payoff to the players than the Nash equilibrium in (a).

\begin{tcolorbox}
(b) 找出该游戏的一个粗相关均衡(coarse correlated equilibrium)，使得玩家在此均衡下的期望效用(expected payoff)高于 (a) 中的纳什均衡。
\end{tcolorbox}

Exercise 12.3. For the game in Figure 12.4, the coarse correlated equilibrium \(z\) on the right has only two probabilities \(a = {z}_{32}\) for the strategy pair(B, M)and \(b = {z}_{23}\) for the strategy pair(M, B)that are positive; all other components \({z}_{ij}\) are zero. Find all possible values of \(a\) and \(b = 1 - a\) so that the correlation devices \(z\) that are defined in this way define a coarse correlated equilibrium of the game.

\begin{tcolorbox}
练习12.3。对于图12.4中的游戏，右侧的粗相关均衡 \(z\) 仅对策略组合 (B, M) 赋予正概率 \(a = {z}_{32}\) ，对策略组合 (M, B) 赋予正概率 \(b = {z}_{23}\) ，其他分量 \({z}_{ij}\) 均为零。找出所有可能的 \(a\) 和 \(b = 1 - a\) 的值，使得以这种方式定义的相关机制 \(z\) 构成该游戏的一个粗相关均衡。
\end{tcolorbox}

Exercise 12.4. Show that in an \(N\) -player game where every player has two strategies, every coarse correlated equilibrium is a correlated equilibrium.

\begin{tcolorbox}
练习12.4。证明在一个有 \(N\) 个玩家且每个玩家有两种策略的游戏中，每个粗相关均衡都是相关均衡。
\end{tcolorbox}

Exercise 12.5. Show that if \(Y\) is a Markov chain on two states, then the off-diagonal elements of \(Y\) , in their columns, are proportional to the probabilities of a stationary distribution \(x\) of \(Y\) . That is, if

\begin{tcolorbox}
练习12.5。证明如果 \(Y\) 是一个两状态的马尔可夫链(Markov chain)，那么 \(Y\) 的非对角元素在其所在列中与 \(Y\) 的平稳分布(stationary distribution) \(x\) 的概率成比例。即，如果
\end{tcolorbox}

\[
Y = \left( \begin{array}{ll} a & b \\  c & d \end{array}\right)
\]

then \(\left( \begin{array}{ll} c & b \end{array}\right)  = \left( {c + b}\right) {x}^{\top }\) . Exactly when is \(x\) unique? Argue carefully.

\begin{tcolorbox}
那么 \(\left( \begin{array}{ll} c & b \end{array}\right)  = \left( {c + b}\right) {x}^{\top }\) 。 \(x\) 何时是唯一的？请给出详细论证。
\end{tcolorbox}

\section*{References}

\begin{tcolorbox}
\section*{参考文献}
\end{tcolorbox}

The page numbers in this book where a reference is cited, for example on page 33, are given after the reference with \(\uparrow  {33}\) .

\begin{tcolorbox}
本书中引用参考文献的页码，例如第33页，会在参考文献后用 \(\uparrow  {33}\) 标注。
\end{tcolorbox}

Adler, I. (2013). The equivalence of linear programs and zero-sum games. International Journal of Game Theory 42(1), 165-177. \(\uparrow  {221}\)

\begin{tcolorbox}
阿德勒，I.(2013)。线性规划与零和博弈(zero - sum game)的等价性。《国际博弈论杂志》42(1)，165 - 177。 \(\uparrow  {221}\)
\end{tcolorbox}

Albert, M. H., R. J. Nowakowski, and D. Wolfe (2007). Lessons in Play: An Introduction to Combinatorial Game Theory. AK Peters, Wellesley, MA. ↑30, 31, 38

\begin{tcolorbox}
阿尔伯特，M. H.，R. J. 诺瓦科夫斯基，和 D. 沃尔夫(2007)。《游戏教程:组合博弈论导论》。AK 彼得斯出版社，马萨诸塞州韦尔斯利。↑30, 31, 38
\end{tcolorbox}

Allais, M. (1953). Le comportement de l'homme rationnel devant le risque: critique des postulats et axiomes de l'école américaine [with English summary]. Econometrica 21(4), 503-546. ↑ 133

\begin{tcolorbox}
阿莱，M.(1953)。《理性人在风险面前的行为:对美国学派假设和公理的批判》(附英文摘要)。《计量经济学》21(4)，503 - 546。↑133
\end{tcolorbox}

Aumann, R. J. (1974). Subjectivity and correlation in randomized strategies. Journal of Mathematical Economics 1(1), 67-96. ↑ 334, 348

\begin{tcolorbox}
奥曼，R. J.(1974)。随机策略中的主观性与相关性。《数学经济学杂志》1(1)，67 - 96。↑334, 348
\end{tcolorbox}

Aumann, R. J. (1987). Correlated equilibrium as an expression of Bayesian rationality. Econometrica 55(1), 1-18. \(\uparrow  {348}\)

\begin{tcolorbox}
奥曼，R. J.(1987)。相关均衡作为贝叶斯理性(Bayesian rationality)的一种表达。《计量经济学》55(1)，1 - 18。 \(\uparrow  {348}\)
\end{tcolorbox}

Aumann, R. and A. Brandenburger (1995). Epistemic conditions for Nash equilibrium. Econometrica 63(5), 1161-1180. ↑ 132, 348

\begin{tcolorbox}
奥曼，R. 和 A. 布兰登伯格(1995)。纳什均衡的认知条件。《计量经济学》63(5)，1161 - 1180。↑132, 348
\end{tcolorbox}

Aumann, R. J. and J. H. Dreze (2008). Rational expectations in games. American Economic Review 98(1), 72-86. ↑ 132, 348

\begin{tcolorbox}
奥曼，R. J. 和 J. H. 德雷兹(2008)。博弈中的理性预期。《美国经济评论》98(1)，72 - 86。↑132, 348
\end{tcolorbox}

Avenhaus, R. and M. J. Canty (1996). Compliance Quantified: An Introduction to Data Verification. Cambridge University Press, Cambridge, UK. ↑166

\begin{tcolorbox}
阿文豪斯(Avenhaus)，R.和M. J. 坎蒂(Canty)(1996)。合规量化:数据验证导论。英国剑桥:剑桥大学出版社。↑166
\end{tcolorbox}

Avis, D., G. D. Rosenberg, R. Savani, and B. von Stengel (2010). Enumeration of Nash equilibria for two-player games. Economic Theory 42(1), 9-37. \(\uparrow  {259},{260}\)

\begin{tcolorbox}
阿维斯(Avis)，D.、G. D. 罗森伯格(Rosenberg)、R. 萨瓦尼(Savani)和B. 冯·斯滕格尔(von Stengel)(2010)。双人博弈的纳什均衡枚举。《经济理论》42(1)，9 - 37。 \(\uparrow  {259},{260}\)
\end{tcolorbox}

Berlekamp, E. R., J. H. Conway, and R. K. Guy (2001-2004). Winning Ways for Your Mathematical Plays, 2nd ed., Volumes 1-4. AK Peters, Wellesley, MA. (First edition 1982). ↑x, 31

\begin{tcolorbox}
伯利坎普(Berlekamp)，E. R.、J. H. 康威(Conway)和R. K. 盖伊(Guy)(2001 - 2004)。《数学博弈的制胜之道》，第2版，第1 - 4卷。美国马萨诸塞州韦尔斯利:AK彼得斯出版社。(第1版于1982年出版)。↑x，31
\end{tcolorbox}

Bernoulli, D. (1738). Specimen theoriae novae de mensura sortis. Commentarii Academiae Scientiarum Imperiales Petropolitanae 5, 175-192. Translated as: Exposition of a new theory on the measurement of risk, Econometrica 22(1), 23-36, 1954. \(\uparrow  {133}\)

\begin{tcolorbox}
伯努利(Bernoulli)，D.(1738)。《关于风险度量的新理论阐述》。《圣彼得堡帝国科学院评论》5，175 - 192。1954年翻译为:《关于风险度量的新理论阐述》，《计量经济学》22(1)，23 - 36。 \(\uparrow  {133}\)
\end{tcolorbox}

Bewersdorff, J. (2005). Luck, Logic, and White Lies: The Mathematics of Games. AK Peters, Wellesley, MA. ↑31

\begin{tcolorbox}
贝沃斯多夫(Bewersdorff)，J.(2005)。《运气、逻辑与善意的谎言:博弈数学》。美国马萨诸塞州韦尔斯利:AK彼得斯出版社。↑31
\end{tcolorbox}

Biggs, N., E. K. Lloyd, and R. J. Wilson (1976). Graph Theory 1736-1936. Clarendon Press, Oxford. \(\uparrow  {203},{353}\)

\begin{tcolorbox}
比格斯(Biggs)，N.、E. K. 劳埃德(Lloyd)和R. J. 威尔逊(Wilson)(1976)。《图论(1736 - 1936)》。英国牛津:克拉伦登出版社。 \(\uparrow  {203},{353}\)
\end{tcolorbox}

Binmore, K. (1992). Fun and Games. D. C. Heath, Lexington, MA. ↑327, 330

\begin{tcolorbox}
宾默尔(Binmore)，K.(1992)。《趣味与博弈》。美国马萨诸塞州列克星敦:D. C. 希思出版社。↑327，330
\end{tcolorbox}

Binmore, K. (2007). Playing for Real. Oxford University Press, New York. ↑4, 330

\begin{tcolorbox}
宾默尔(Binmore)，K.(2007)。《真实博弈》。美国纽约:牛津大学出版社。↑4，330
\end{tcolorbox}

Bouton, C. L. (1901). Nim, a game with a complete mathematical theory. Annals of Mathematics 3(1/4), 35-39. ↑5, 31

\begin{tcolorbox}
布顿(Bouton)，C. L.(1901)。《尼姆游戏(Nim):一个具有完整数学理论的博弈》。《数学年刊》3(1/4)，35 - 39。↑5，31
\end{tcolorbox}

Braess, D. (1968). Über ein Paradoxon aus der Verkehrsplanung. Unternehmensforschung 12(1), 258-268. Translated in Braess, Nagurney, and Wakolbinger (2005). \(\uparrow  {44}\) , 49

\begin{tcolorbox}
布雷斯(Braess)，D.(1968)。《交通规划中的一个悖论》。《企业研究》12(1)，258 - 268。转载于布雷斯、纳古尔内和瓦科尔宾格(2005)。 \(\uparrow  {44}\) ，49
\end{tcolorbox}

Braess, D., A. Nagurney, and T. Wakolbinger (2005). On a paradox of traffic planning. Transportation Science 39(4), 446-450. ↑ 352

\begin{tcolorbox}
布雷斯(Braess)，D.、A. 纳古尔内(Nagurney)和T. 瓦科尔宾格(Wakolbinger)(2005)。《关于交通规划的一个悖论》。《运输科学》39(4)，446 - 450。↑ 352
\end{tcolorbox}

Brouwer, L. E. J. (1911). Über Abbildung von Mannigfaltigkeiten. Mathematische Annalen 71(1), 97-115. ↑203

\begin{tcolorbox}
布劳威尔(Brouwer)，L. E. J.(1911)。《关于流形的映射》。《数学年刊》71(1)，97 - 115。↑203
\end{tcolorbox}

Brouwer, L. E. J. (1976). Collected Works 2: Geometry, Analysis, Topology and Mechanics. North-Holland, Amsterdam. Edited by H. Freudenthal. \(\uparrow  {203}\)

\begin{tcolorbox}
布劳威尔(Brouwer)，L. E. J.(1976)。《文集2:几何、分析、拓扑与力学》。荷兰阿姆斯特丹:北荷兰出版社。由H. 弗罗伊登塔尔编辑。 \(\uparrow  {203}\)
\end{tcolorbox}

Bryant, V. (1990). Yet Another Introduction to Analysis. Cambridge University Press, Cambridge, UK. † 204

\begin{tcolorbox}
布赖恩特(Bryant)，V.(1990)。《再谈分析导论》。英国剑桥:剑桥大学出版社。† 204
\end{tcolorbox}

Chen, X., X. Deng, and S.-H. Teng (2009). Settling the complexity of computing two-player Nash equilibria. Journal of the ACM 56(3), Article 14. \(\uparrow  {259}\)

\begin{tcolorbox}
陈(Chen)，X.、邓(Deng)，X.和滕(Teng)，S. - H.(2009)。《确定计算双人纳什均衡的复杂性》。《美国计算机协会期刊》56(3)，文章编号14。 \(\uparrow  {259}\)
\end{tcolorbox}

Cheng, S.-F., D. M. Reeves, Y. Vorobeychik, and M. P. Wellman (2004). Notes on equilibria in symmetric games. In: Sixth Workshop on Game Theoretic and Decision Theoretic Agents at the 3rd Conference on Autonomous Agents and Multi-Agent Systems (AAMAS '04). New York. ↑74

\begin{tcolorbox}
程(Cheng)，S. - F.、D. M. 里夫斯(Reeves)、Y. 沃罗贝奇克(Vorobeychik)和M. P. 韦尔曼(Wellman)(2004)。《对称博弈(symmetric game)中的均衡注记》。载于:第三届自治代理与多代理系统会议(AAMAS '04)第六届博弈论与决策论代理研讨会。美国纽约。↑74
\end{tcolorbox}

Chvátal, V. (1983). Linear Programming. W. H. Freeman, New York. ↑221, 260

\begin{tcolorbox}
赫瓦塔尔(Chvátal, V.)(1983 年)。《线性规划》(Linear Programming)。纽约:W. H. Freeman 出版社。↑221, 260
\end{tcolorbox}

Cohen, D. I. A. (1967). On the Sperner lemma. Journal of Combinatorial Theory 2(4), 585-587. \(\uparrow  {203}\)

\begin{tcolorbox}
科恩(Cohen, D. I. A.)(1967 年)。关于斯佩纳引理。《组合理论杂志》(Journal of Combinatorial Theory)2(4)，585 - 587。 \(\uparrow  {203}\)
\end{tcolorbox}

Conway, J. H. (1977). All games bright and beautiful. American Mathematical Monthly 84(6), 417-434. \(\uparrow  {31}\)

\begin{tcolorbox}
康威(Conway, J. H.)(1977 年)。所有美好而光明的游戏。《美国数学月刊》(American Mathematical Monthly)84(6)，417 - 434。 \(\uparrow  {31}\)
\end{tcolorbox}

Conway, J. H. (2001). On Numbers and Games, 2nd ed. AK Peters, Natick, MA. (First edition 1976). \(\uparrow  {24},{31}\)

\begin{tcolorbox}
康威(Conway, J. H.)(2001 年)。《论数字与游戏》(On Numbers and Games)，第 2 版。马萨诸塞州纳蒂克:AK Peters 出版社。(第 1 版于 1976 年出版)。 \(\uparrow  {24},{31}\)
\end{tcolorbox}

Cottle, R. W., J.-S. Pang, and R. E. Stone (1992). The Linear Complementarity Problem. Academic Press, San Diego, CA. \(\uparrow  {260}\)

\begin{tcolorbox}
科特尔(Cottle, R. W.)、庞(J.-S. Pang)和斯通(R. E. Stone)(1992 年)。《线性互补问题》(The Linear Complementarity Problem)。加利福尼亚州圣地亚哥:学术出版社(Academic Press)。 \(\uparrow  {260}\)
\end{tcolorbox}

Cournot, A. A. (1838). Recherches sur les principes mathématiques de la théorie des richesses. Hachette, Paris. Translated as: Researches into the Mathematical Principles of Wealth, A. M. Kelly, New York, 1960. \(\uparrow  {66},{74}\)

\begin{tcolorbox}
古诺(Cournot, A. A.)(1838 年)。《财富理论的数学原理研究》(Recherches sur les principes mathématiques de la théorie des richesses)。巴黎:阿歇特出版社(Hachette)。英译本:《财富的数学原理研究》(Researches into the Mathematical Principles of Wealth)，纽约:A. M. Kelly 出版社，1960 年。 \(\uparrow  {66},{74}\)
\end{tcolorbox}

Dantzig, G. B. (1963). Linear Programming and Extensions. Princeton University Press, Princeton, NJ. \(\uparrow  {220},{260}\)

\begin{tcolorbox}
丹齐格(Dantzig, G. B.)(1963 年)。《线性规划及其扩展》(Linear Programming and Extensions)。新泽西州普林斯顿:普林斯顿大学出版社(Princeton University Press)。 \(\uparrow  {220},{260}\)
\end{tcolorbox}

Dixit, A. K. and B. J. Nalebuff (1991). Thinking Strategically: The Competitive Edge in Business, Politics, and Everyday Life. W. W. Norton, New York. ↑74

\begin{tcolorbox}
迪克西特(Dixit, A. K.)和奈尔伯夫(B. J. Nalebuff)(1991 年)。《策略思维:商业、政治与日常生活中的竞争优势》(Thinking Strategically: The Competitive Edge in Business, Politics, and Everyday Life)。纽约:W. W. 诺顿出版社(W. W. Norton)。↑74
\end{tcolorbox}

Euler, L. (1741). Solutio problematis ad geometriam situs pertinentis. Commentarii Academiae Scientiarum Imperiales Petropolitanae 8, 128-140 + Plate VIII. Translated as: The solution of a problem relating to the geometry of position (Biggs, Lloyd, and Wilson, 1976, 3-8). ↑203

\begin{tcolorbox}
欧拉(Euler, L.)(1741 年)。《与位置几何相关问题的解》(Solutio problematis ad geometriam situs pertinentis)。《圣彼得堡帝国科学院评论》(Commentarii Academiae Scientiarum Imperiales Petropolitanae)8，128 - 140 + 图版 VIII。英译本:《与位置几何相关问题的解》(The solution of a problem relating to the geometry of position)(比格斯、劳埃德和威尔逊，1976 年，3 - 8)。↑203
\end{tcolorbox}

Fishburn, P. C. (1970). Utility Theory for Decision Making. Wiley, New York. ↑133

\begin{tcolorbox}
菲什伯恩(Fishburn, P. C.)(1970 年)。《决策的效用理论》(Utility Theory for Decision Making)。纽约:威利出版社(Wiley)。↑133
\end{tcolorbox}

Fraenkel, A. S. (1996). Scenic trails ascending from sea-level Nim to alpine Chess. In: Games of No Chance, edited by R. J. Nowakowski, volume 29 of MSRI Publications, 13-42. Cambridge University Press, Cambridge, UK. ↑31

\begin{tcolorbox}
弗伦克尔(Fraenkel, A. S.)(1996 年)。从海平面的尼姆游戏(Nim)到高山象棋的风景小径。载于:R. J. 诺瓦科夫斯基编辑的《无机会游戏》(Games of No Chance)，数学科学研究所出版物(MSRI Publications)第 29 卷，13 - 42。英国剑桥:剑桥大学出版社(Cambridge University Press)。↑31
\end{tcolorbox}

Freudenthal, H. (1942). Simplizialzerlegungen von beschränkter Flachheit. Annals of Mathematics 43(3), 580-582. \(\uparrow  {170},{196},{204}\)

\begin{tcolorbox}
弗罗伊登塔尔(Freudenthal, H.)(1942 年)。有限平坦性的单纯分解。《数学年刊》(Annals of Mathematics)43(3)，580 - 582。 \(\uparrow  {170},{196},{204}\)
\end{tcolorbox}

Friedman, D. (1998). Monty Hall's three doors: Construction and deconstruction of a choice anomaly. The American Economic Review 88(4), 933-946. \(\uparrow  {296}\)

\begin{tcolorbox}
弗里德曼(Friedman, D.)(1998 年)。蒙提·霍尔的三扇门:选择异常的构建与解构。《美国经济评论》(The American Economic Review)88(4)，933 - 946。 \(\uparrow  {296}\)
\end{tcolorbox}

Gale, D. (1960). The Theory of Linear Economic Models. McGraw-Hill, New York. ↑220

\begin{tcolorbox}
盖尔(Gale, D.)(1960 年)。《线性经济模型理论》(The Theory of Linear Economic Models)。纽约:麦格劳 - 希尔出版社(McGraw - Hill)。↑220
\end{tcolorbox}

Gale, D. (1974). A curious Nim-type game. The American Mathematical Monthly 81(8), 876-879. \(\uparrow  {31}\)

\begin{tcolorbox}
盖尔(Gale, D.)(1974 年)。一个奇特的尼姆类型游戏(Nim - type game)。《美国数学月刊》(The American Mathematical Monthly)81(8)，876 - 879。 \(\uparrow  {31}\)
\end{tcolorbox}

Gale, D. and A. Neyman (1982). Nim-type games. International Journal of Game Theory 11(1), 17-20. \(\uparrow  {31}\)

\begin{tcolorbox}
盖尔(Gale, D.)和内曼(A. Neyman)(1982 年)。尼姆类型游戏(Nim - type games)。《国际博弈论杂志》(International Journal of Game Theory)11(1)，17 - 20。 \(\uparrow  {31}\)
\end{tcolorbox}

Gibbons, R. (1992). Game Theory for Applied Economists. Princeton University Press, Princeton, NJ. \(\uparrow\) ix, 74, 98, 295

\begin{tcolorbox}
吉本斯，R.(1992)。《应用经济学家的博弈论》。普林斯顿大学出版社，新泽西州普林斯顿。 \(\uparrow\) 第ix页，74页，98页，295页
\end{tcolorbox}

Gilboa, I. and E. Zemel (1989). Nash and correlated equilibria: Some complexity considerations. Games and Economic Behavior 1(1), 80-93. ↑349

\begin{tcolorbox}
吉尔博阿，I. 和 E. 泽梅尔(1989)。纳什均衡与相关均衡(Nash and correlated equilibria):一些复杂性考量。《博弈与经济行为》1(1)，80 - 93。↑349
\end{tcolorbox}

Grundy, P. M. (1939). Mathematics and games. Eureka 2, \(6 - 8. \uparrow  1,{16},{31}\)

\begin{tcolorbox}
格伦迪，P. M.(1939)。数学与博弈。《尤里卡》2， \(6 - 8. \uparrow  1,{16},{31}\)
\end{tcolorbox}

Hart, S. and A. Mas-Colell (2000). A simple adaptive procedure leading to correlated equilibrium. Econometrica 68(5), 1127-1150. ↑ 349

\begin{tcolorbox}
哈特，S. 和 A. 马斯 - 克莱尔(2000)。一种通向相关均衡(correlated equilibrium)的简单自适应程序。《计量经济学》68(5)，1127 - 1150。↑ 349
\end{tcolorbox}

Hart, S. and D. Schmeidler (1989). Existence of correlated equilibria. Mathematics of Operations Research 14(1), 18-25. Tx, 333, 349

\begin{tcolorbox}
哈特，S. 和 D. 施梅德勒(1989)。相关均衡(correlated equilibrium)的存在性。《运筹学数学》14(1)，18 - 25。Tx，333，349
\end{tcolorbox}

Herstein, I. N. and J. Milnor (1953). An axiomatic approach to measurable utility. Econometrica 21(2), 291-297. ↑ 133

\begin{tcolorbox}
赫斯坦，I. N. 和 J. 米尔诺(1953)。可测效用的公理化方法。《计量经济学》21(2)，291 - 297。↑ 133
\end{tcolorbox}

Hong, S.-K., I.-J. Song, and J. Wu (2007). Fengshui theory in urban landscape planning. Urban Ecosystems 10(3), 221-237. \(\uparrow  {44},{49}\)

\begin{tcolorbox}
洪，S.-K.，宋，I.-J.，和吴，J.(2007)。城市景观规划中的风水理论。《城市生态系统》10(3)，221 - 237。 \(\uparrow  {44},{49}\)
\end{tcolorbox}

Huizinga, J. (1949). Homo Ludens: A Study of the Play-Element in Culture. Routledge, London. ↑330

\begin{tcolorbox}
赫伊津哈，J.(1949)。《游戏的人:对文化中游戏元素的研究》。劳特利奇出版社，伦敦。↑330
\end{tcolorbox}

Knaster, B., C. Kuratowski, and S. Mazurkiewicz (1929). Ein Beweis des Fixpunktsatzes für n-dimensionale Simplexe. Fundamenta Mathematicae 14(1), 132-137. ↑ 169, 184, 187, 203

\begin{tcolorbox}
克纳斯特，B.，C. 库拉托夫斯基，和 S. 马祖尔克维奇(1929)。n维单纯形(n - dimensional Simplex)不动点定理的一个证明。《数学基础》14(1)，132 - 137。↑ 169，184，187，203
\end{tcolorbox}

Kuhn, H. W. (1950). Simplified two-person poker. In: Contributions to the Theory of Games, Vol. I, edited by H. W. Kuhn and A. W. Tucker, volume 24 of Annals of Mathematics Studies, 97-103. Princeton University Press, Princeton, NJ. \(\uparrow  {297}\)

\begin{tcolorbox}
库恩，H. W.(1950)。简化的两人扑克(Poker)游戏。载于:《博弈论论文集，第一卷》，由 H. W. 库恩和 A. W. 塔克编辑，《数学研究年刊》第24卷，97 - 103。普林斯顿大学出版社，新泽西州普林斯顿。 \(\uparrow  {297}\)
\end{tcolorbox}

Kuhn, H. W. (1953). Extensive games and the problem of information. In: Contributions to the Theory of Games, Vol. II, edited by H. W. Kuhn and A. W. Tucker, volume 28 of Annals of Mathematics Studies, 193-216. Princeton University Press, Princeton, NJ. Reprinted in Kuhn (1997). \(\uparrow  {263},{283},{285},{295}\)

\begin{tcolorbox}
库恩，H. W.(1953)。扩展型博弈与信息问题。载于:《博弈论论文集，第二卷》，由 H. W. 库恩和 A. W. 塔克编辑，《数学研究年刊》第28卷，193 - 216。普林斯顿大学出版社，新泽西州普林斯顿。重印于库恩(1997)。 \(\uparrow  {263},{283},{285},{295}\)
\end{tcolorbox}

Kuhn, H. W. (1960). Some combinatorial lemmas in topology. IBM Journal of Research and Development 4(5), 518-524. \(\uparrow  {204}\)

\begin{tcolorbox}
库恩，H. W.(1960)。拓扑学中的一些组合引理。《IBM研究与开发期刊》4(5)，518 - 524。 \(\uparrow  {204}\)
\end{tcolorbox}

Kuhn, H. W. (1969). Approximate search for fixed points. In: Computing Methods in Optimization Problems 2, edited by L. A. Zadeh, L. W. Neustadt, and A. V. Balakrishnan, 199-211. Academic Press, New York. ↑203

\begin{tcolorbox}
库恩，H. W.(1969)。不动点的近似搜索。载于:《优化问题的计算方法2》，由 L. A. 扎德、L. W. 诺伊施塔特和 A. V. 巴拉科里什南编辑，199 - 211。学术出版社，纽约。↑203
\end{tcolorbox}

Kuhn, H. W. (ed.) (1997). Classics in Game Theory. Princeton University Press, Princeton, NJ. \(\uparrow  {354},{355},{356}\)

\begin{tcolorbox}
库恩，H. W.(编)(1997)。《博弈论经典著作》。普林斯顿大学出版社，新泽西州普林斯顿。 \(\uparrow  {354},{355},{356}\)
\end{tcolorbox}

Lemke, C. E. (1965). Bimatrix equilibrium points and mathematical programming. Management Science 11(7), 681-689. ↑ 260

\begin{tcolorbox}
莱姆基，C. E.(1965)。双矩阵博弈(bimatrix game)的均衡点与数学规划。《管理科学》11(7)，681 - 689。↑ 260
\end{tcolorbox}

Lemke, C. E. and J. T. Howson, Jr (1964). Equilibrium points of bimatrix games. Journal of the Society for Industrial and Applied Mathematics 12(2), 413-423. \(\uparrow  \mathrm{x},{223},{230},{232},{259}\)

\begin{tcolorbox}
莱姆基，C. E. 和小 J. T. 豪森(1964)。双矩阵博弈(bimatrix game)的均衡点。《工业与应用数学学会期刊》12(2)，413 - 423。 \(\uparrow  \mathrm{x},{223},{230},{232},{259}\)
\end{tcolorbox}

Lewis, M. (2016). The Undoing Project: A Friendship That Changed Our Minds. W. W. Norton, New York. \(\uparrow  {133}\)

\begin{tcolorbox}
刘易斯，M.(2016)。《解构计划:一段改变我们思维的友谊》。W. W. 诺顿出版社，纽约。 \(\uparrow  {133}\)
\end{tcolorbox}

Liu, L. A. (1996). The invariance of best reply correspondences in two-player games. Working paper. City University of Hong Kong, Faculty of Business, Department of Economics and Finance. \(\uparrow  {259}\)

\begin{tcolorbox}
刘，L. A.(1996)。双人博弈中最佳反应对应关系的不变性。工作论文。香港城市大学商学院经济与金融系。 \(\uparrow  {259}\)
\end{tcolorbox}

Loomis, L. H. (1946). On a theorem of von Neumann. Proceedings of the National Academy of Sciences of the United States of America 32(8), 213-215. ↑ 220

\begin{tcolorbox}
卢米斯，L. H.(1946)。关于冯·诺伊曼的一个定理。《美国国家科学院院刊》32(8)，213 - 215。↑ 220
\end{tcolorbox}

Luce, R. D. and H. Raiffa (1957). Games and Decisions: Introduction and Critical Survey. Wiley, New York. \(\uparrow  {133}\)

\begin{tcolorbox}
卢斯，R. D. 和 H. 雷法(1957)。《博弈与决策:导论与批判性综述》。威利出版社，纽约。 \(\uparrow  {133}\)
\end{tcolorbox}

Mangasarian, O. L. (1964). Equilibrium points of bimatrix games. Journal of the Society for Industrial and Applied Mathematics 12(4), 778-780. \(\uparrow  {260}\)

\begin{tcolorbox}
曼加萨里安，O. L.(1964)。双矩阵博弈(bimatrix game)的均衡点。《工业与应用数学学会杂志》12(4)，778 - 780。 \(\uparrow  {260}\)
\end{tcolorbox}

Maschler, M., E. Solan, and S. Zamir (2013). Game Theory. Cambridge University Press, Cambridge, UK. ↑ 166, 330

\begin{tcolorbox}
马施勒，M.、E. 索兰和 S. 扎米尔(2013)。《博弈论》。剑桥大学出版社，英国剑桥。↑ 166, 330
\end{tcolorbox}

Matoušek, J. and B. Gärtner (2007). Understanding and Using Linear Programming. Springer, Berlin. \(\uparrow  {221}\)

\begin{tcolorbox}
马托塞克，J. 和 B. 加特纳(2007)。《理解与应用线性规划(linear programming)》。施普林格出版社，柏林。 \(\uparrow  {221}\)
\end{tcolorbox}

Monderer, D. and L. S. Shapley (1996). Potential games. Games and Economic Behavior 14(1), 124-143. \(\uparrow  {50},{74}\)

\begin{tcolorbox}
蒙德勒，D. 和 L. S. 沙普利(1996)。势博弈(potential game)。《博弈与经济行为》14(1)，124 - 143。 \(\uparrow  {50},{74}\)
\end{tcolorbox}

Morgenstern, O. (1973). Ingolf Ståhl: Bargaining theory (book review). The Swedish Journal of Economics 75(4), 410-413. \(\uparrow  {356}\)

\begin{tcolorbox}
摩根斯坦，O.(1973)。英戈尔夫·斯塔尔:《讨价还价理论》(书评)。《瑞典经济学杂志》75(4)，410 - 413。 \(\uparrow  {356}\)
\end{tcolorbox}

Moscati, I. (2016). Retrospectives: How economists came to accept expected utility theory: The case of Samuelson and Savage. Journal of Economic Perspectives 30(2), 219-236. \(\uparrow  {133}\)

\begin{tcolorbox}
莫斯卡蒂，I.(2016)。回顾:经济学家如何接受期望效用理论:以萨缪尔森和萨维奇为例。《经济展望杂志》30(2)，219 - 236。 \(\uparrow  {133}\)
\end{tcolorbox}

Moulin, H. and J.-P. Vial (1978). Strategically zero-sum games: The class of games whose completely mixed equilibria cannot be improved upon. International Journal of Game Theory 7(3/4), 201-221. \(\uparrow  {259},{349}\)

\begin{tcolorbox}
莫林，H. 和 J.-P. 维亚尔(1978)。策略性零和博弈(zero - sum game):完全混合均衡无法改进的博弈类。《国际博弈论杂志》7(3/4)，201 - 221。 \(\uparrow  {259},{349}\)
\end{tcolorbox}

Myerson, R. B. (1991). Game Theory: Analysis of Conflict. Harvard University Press, Cambridge, MA. \(\uparrow  {133},{348}\)

\begin{tcolorbox}
迈尔森，R. B.(1991)。《博弈论:冲突分析》。哈佛大学出版社，马萨诸塞州剑桥。 \(\uparrow  {133},{348}\)
\end{tcolorbox}

Nash, J. F., Jr. (1950). The bargaining problem. Econometrica 18(2), 155-162. Reprinted in Kuhn (1997). \(\uparrow  {299},{330}\)

\begin{tcolorbox}
小纳什，J. F.(1950)。讨价还价问题。《计量经济学》18(2)，155 - 162。转载于库恩(1997)。 \(\uparrow  {299},{330}\)
\end{tcolorbox}

Nash, J. (1951). Non-cooperative games. Annals of Mathematics 54(2), 286-295. Reprinted in Kuhn (1997). \(\uparrow  \mathrm{x},{53},{59},{135},{136},{149},{152},{166},{169},{205},{214}\)

\begin{tcolorbox}
纳什，J.(1951)。非合作博弈。《数学年刊》54(2)，286 - 295。转载于库恩(1997)。 \(\uparrow  \mathrm{x},{53},{59},{135},{136},{149},{152},{166},{169},{205},{214}\)
\end{tcolorbox}

Nau, R., S. Gomez Canovas, and P. Hansen (2004). On the geometry of Nash equilibria and correlated equilibria. International Journal of Game Theory 32(4), 443-453. \(\uparrow  {348}\)

\begin{tcolorbox}
瑙，R.、S. 戈麦斯·卡诺瓦斯和 P. 汉森(2004)。关于纳什均衡(Nash equilibrium)和相关均衡的几何性质。《国际博弈论杂志》32(4)，443 - 453。 \(\uparrow  {348}\)
\end{tcolorbox}

Nau, R. F. and K. F. McCardle (1990). Coherent behavior in noncooperative games. Journal of Economic Theory 50(2), 424-444. \(\uparrow  {349}\)

\begin{tcolorbox}
瑙，R. F. 和 K. F. 麦卡德尔(1990)。非合作博弈中的一致行为。《经济理论杂志》50(2)，424 - 444。 \(\uparrow  {349}\)
\end{tcolorbox}

Neyman, A. (1997). Correlated equilibrium and potential games. International Journal of Game Theory 26(2), 223-227. \(\uparrow  {74}\)

\begin{tcolorbox}
内曼，A.(1997年)。相关均衡(correlated equilibrium)与势博弈(potential game)。《国际博弈论杂志》26(2)，223 - 227。 \(\uparrow  {74}\)
\end{tcolorbox}

Osborne, M. J. (2004). An Introduction to Game Theory. Oxford University Press, New York. ↑ix, 74, 76, 166, 330

\begin{tcolorbox}
奥斯本，M. J.(2004年)。《博弈论导论》。牛津大学出版社，纽约。↑ix，74，76，166，330
\end{tcolorbox}

Osborne, M. J. and A. Rubinstein (1994). A Course in Game Theory. MIT Press, Cambridge, MA. \(\uparrow  {133},{295},{330}\)

\begin{tcolorbox}
奥斯本，M. J.和A. 鲁宾斯坦(1994年)。《博弈论教程》。麻省理工学院出版社，马萨诸塞州剑桥。 \(\uparrow  {133},{295},{330}\)
\end{tcolorbox}

Owen, G. (1967). Communications to the editor: An elementary proof of the minimax theorem. Management Science 13(9), 765-765. ↑ 220

\begin{tcolorbox}
欧文，G.(1967年)。致编辑的信:极小极大定理(minimax theorem)的一个初等证明。《管理科学》13(9)，765 - 765。↑ 220
\end{tcolorbox}

Papadimitriou, C. H. (1994). On the complexity of the parity argument and other inefficient proofs of existence. Journal of Computer and System Sciences 48(3), 498-532. \(\uparrow  {203},{259}\)

\begin{tcolorbox}
帕帕迪米特里欧，C. H.(1994年)。关于奇偶性论证(parity argument)的复杂性及其他低效存在性证明。《计算机与系统科学杂志》48(3)，498 - 532。 \(\uparrow  {203},{259}\)
\end{tcolorbox}

Papadimitriou, C. H. and T. Roughgarden (2008). Computing correlated equilibria in multi-player games. Journal of the ACM 55(3), Article 14. \(\uparrow  {349}\)

\begin{tcolorbox}
帕帕迪米特里欧，C. H.和T. 拉夫加登(2008年)。多人博弈中相关均衡(correlated equilibrium)的计算。《美国计算机协会杂志》55(3)，第14篇。 \(\uparrow  {349}\)
\end{tcolorbox}

Park, S. (1999). Ninety years of the Brouwer fixed point theorem. Vietnam Journal of Mathematics 27(3), 187-222. \(\uparrow  {203}\)

\begin{tcolorbox}
朴，S.(1999年)。布劳威尔不动点定理的九十年。《越南数学杂志》27(3)，187 - 222。 \(\uparrow  {203}\)
\end{tcolorbox}

Perea, A. (2012). Epistemic Game Theory: Reasoning and Choice. Cambridge University Press. ↑132

\begin{tcolorbox}
佩雷阿，A.(2012年)。《认知博弈论:推理与选择》。剑桥大学出版社。↑132
\end{tcolorbox}

Pigou, A. (1920). The Economics of Welfare. McMillan, London. ↑ 49

\begin{tcolorbox}
庇古，A.(1920年)。《福利经济学》。麦克米伦出版社，伦敦。↑ 49
\end{tcolorbox}

Pratt, J. W. (1964). Risk aversion in the small and in the large. Econometrica 32(1/2), 122-136. ↑133

\begin{tcolorbox}
普拉特，J. W.(1964年)。小范围和大范围的风险规避。《计量经济学》32(1/2)，122 - 136。↑133
\end{tcolorbox}

Robinson, J. (1951). An iterative method of solving a game. Annals of Mathematics 54(2), 296-301. Reprinted in Kuhn (1997). ↑349

\begin{tcolorbox}
罗宾逊，J.(1951年)。一种求解博弈的迭代方法。《数学年刊》54(2)，296 - 301。转载于库恩(1997年)。↑349
\end{tcolorbox}

Rosenthal, R. W. (1973). A class of games possessing pure-strategy Nash equilibria. International Journal of Game Theory 2(1), 65-67. \(\uparrow  {49}\)

\begin{tcolorbox}
罗森塔尔，R. W.(1973年)。一类具有纯策略纳什均衡(Nash equilibrium)的博弈。《国际博弈论杂志》2(1)，65 - 67。 \(\uparrow  {49}\)
\end{tcolorbox}

Roughgarden, T. (2016). Twenty Lectures on Algorithmic Game Theory. Cambridge University Press, Cambridge, UK. ↑ xii, 49

\begin{tcolorbox}
拉夫加登，T.(2016年)。《算法博弈论二十讲》。剑桥大学出版社，英国剑桥。↑ xii，49
\end{tcolorbox}

Roughgarden, T. and É. Tardos (2002). How bad is selfish routing? Journal of the ACM 49(2), 236-259. ↑ 49

\begin{tcolorbox}
拉夫加登，T.和É. 塔尔多斯(2002年)。自私路由有多糟糕？《美国计算机协会杂志》49(2)，236 - 259。↑ 49
\end{tcolorbox}

Rubinstein, A. (1982). Perfect equilibrium in a bargaining model. Econometrica 50(1), 97-109. \(\uparrow  {299},{330}\)

\begin{tcolorbox}
鲁宾斯坦，A.(1982年)。一个讨价还价模型中的完美均衡。《计量经济学》50(1)，97 - 109。 \(\uparrow  {299},{330}\)
\end{tcolorbox}

Savani, R. and B. von Stengel (2006). Hard-to-solve bimatrix games. Econometrica 74(2), 397-429. \(\uparrow  {260}\)

\begin{tcolorbox}
萨瓦尼，R.和B. 冯·施滕格尔(2006年)。难解的双矩阵博弈(bimatrix game)。《计量经济学》74(2)，397 - 429。 \(\uparrow  {260}\)
\end{tcolorbox}

Savani, R. and B. von Stengel (2016). Unit vector games. International Journal of Economic Theory 12,7-27. \(\uparrow  {260}\)

\begin{tcolorbox}
萨瓦尼(Savani)，R. 和 B. 冯·施滕格尔(von Stengel)(2016 年)。单位向量博弈。《国际经济理论杂志》(International Journal of Economic Theory)12 卷，第 7 - 27 页。 \(\uparrow  {260}\)
\end{tcolorbox}

Schrijver, A. (1986). Theory of Linear and Integer Programming. John Wiley \& Sons, Chichester, UK. \(\uparrow  {221}\)

\begin{tcolorbox}
施里弗(Schrijver)，A.(1986 年)。《线性与整数规划理论》(Theory of Linear and Integer Programming)。约翰·威利父子出版公司(John Wiley \& Sons)，英国奇切斯特。 \(\uparrow  {221}\)
\end{tcolorbox}

Schwalbe, U. and P. Walker (2001). Zermelo and the early history of game theory. Games and Economic Behavior 34(1), 123-137. ↑ 98, 357

\begin{tcolorbox}
施瓦尔贝(Schwalbe)，U. 和 P. 沃克(Walker)(2001 年)。策梅洛与博弈论早期历史。《博弈与经济行为》(Games and Economic Behavior)34(1) 卷，第 123 - 137 页。 ↑ 98, 357
\end{tcolorbox}

Selten, R. (1975). Reexamination of the perfectness concept for equilibrium points in extensive games. International Journal of Game Theory 4(1), 25-55. Reprinted in Kuhn (1997). ↑295

\begin{tcolorbox}
泽尔腾(Selten)，R.(1975 年)。对扩展型博弈中均衡点完美性概念的重新审视。《国际博弈论杂志》(International Journal of Game Theory)4(1) 卷，第 25 - 55 页。转载于库恩(Kuhn)(1997 年)。 ↑295
\end{tcolorbox}

Selvin, S. (1975). Letters to the editor: A problem in probability. The American Statistician 29(1), 67. ↑ 296

\begin{tcolorbox}
塞尔文(Selvin)，S.(1975 年)。致编辑的信:一个概率问题。《美国统计学家》(The American Statistician)29(1) 卷，第 67 页。 ↑ 296
\end{tcolorbox}

Shapley, L. S. (1963). Some topics in two-person games. Memorandum RM-3672-PR. The Rand Corporation, Santa Monica, CA. ↑349

\begin{tcolorbox}
沙普利(Shapley)，L. S.(1963 年)。两人博弈的一些主题。备忘录 RM - 3672 - PR。兰德公司(The Rand Corporation)，美国加利福尼亚州圣莫尼卡。 ↑349
\end{tcolorbox}

Shapley, L. S. (1974). A note on the Lemke-Howson algorithm. Mathematical Programming Study 1: Pivoting and Extensions, 175-189. \(\uparrow  {223},{259}\)

\begin{tcolorbox}
沙普利(Shapley)，L. S.(1974 年)。关于莱姆基 - 豪森算法(Lemke - Howson algorithm)的一则笔记。《数学规划研究》(Mathematical Programming Study)1:换基操作(Pivoting)与扩展，第 175 - 189 页。 \(\uparrow  {223},{259}\)
\end{tcolorbox}

Siegel, A. N. (2013). Combinatorial Game Theory, volume 146 of Graduate Studies in Mathematics. American Mathematical Society, Providence, RI. \(\uparrow  {30}\)

\begin{tcolorbox}
西格尔(Siegel)，A. N.(2013 年)。《组合博弈论》(Combinatorial Game Theory)，《数学研究生教材》(Graduate Studies in Mathematics)第 146 卷。美国数学学会(American Mathematical Society)，美国罗得岛州普罗维登斯。 \(\uparrow  {30}\)
\end{tcolorbox}

Sperner, E. (1928). Neuer Beweis für die Invarianz der Dimensionszahl und des Gebietes. Abhandlungen aus dem Mathematischen Seminar der Universität Hamburg 6(1), 265-272. \(\uparrow  {169},{170},{184},{203}\)

\begin{tcolorbox}
施佩纳(Sperner)，E.(1928 年)。维数(Dimension)和区域不变性的新证明。《汉堡大学数学研讨会论文集》(Abhandlungen aus dem Mathematischen Seminar der Universität Hamburg)6(1) 卷，第 265 - 272 页。 \(\uparrow  {169},{170},{184},{203}\)
\end{tcolorbox}

Sprague, R. (1935). Über mathematische Kampfspiele. Tohoku Mathematical Journal (First Series) 41, 438-444. \(\uparrow  1,{16},{31}\)

\begin{tcolorbox}
施普拉格(Sprague)，R.(1935 年)。关于数学博弈。《东北数学杂志》(Tohoku Mathematical Journal)(第一辑)41 卷，第 438 - 444 页。 \(\uparrow  1,{16},{31}\)
\end{tcolorbox}

Ståhl, I. (1972). Bargaining Theory. The Economic Research Institute at the Stockholm School of Economics, Stockholm. Reviewed in Morgenstern (1973). ↑ 330

\begin{tcolorbox}
斯塔尔(Ståhl)，I.(1972 年)。《讨价还价理论》(Bargaining Theory)。斯德哥尔摩经济学院经济研究所(The Economic Research Institute at the Stockholm School of Economics)，瑞典斯德哥尔摩。摩根斯坦(Morgenstern)(1973 年)对此书进行了评论。 ↑ 330
\end{tcolorbox}

Szpiro, G. G. (2010). Numbers Rule: The Vexing Mathematics of Democracy, from Plato to the Present. Princeton University Press, Princeton, NJ. \(\uparrow  {133}\)

\begin{tcolorbox}
斯皮罗(Szpiro)，G. G.(2010 年)。《数字统治:从柏拉图到现在的民主数学难题》(Numbers Rule: The Vexing Mathematics of Democracy, from Plato to the Present)。普林斯顿大学出版社(Princeton University Press)，美国新泽西州普林斯顿。 \(\uparrow  {133}\)
\end{tcolorbox}

Szpiro, G. G. (2020). Risk, Choice and Uncertainty: Three Centuries of Economic Decision-Making. Columbia University Press, New York. ↑ 133

\begin{tcolorbox}
斯皮罗(Szpiro)，G. G.(2020 年)。《风险、选择与不确定性:三个世纪的经济决策》(Risk, Choice and Uncertainty: Three Centuries of Economic Decision - Making)。哥伦比亚大学出版社(Columbia University Press)，美国纽约。 ↑ 133
\end{tcolorbox}

Turocy, T. L. and B. von Stengel (2002). Game theory. In: Encyclopedia of Information Systems, volume 2,403-420. Elsevier Science (USA). \(\uparrow  {295}\)

\begin{tcolorbox}
图罗茨基(Turocy)，T. L. 和 B. 冯·施滕格尔(von Stengel)(2002 年)。博弈论。载于《信息系统百科全书》(Encyclopedia of Information Systems)，第 2 卷，第 403 - 420 页。爱思唯尔科学出版社(Elsevier Science)(美国)。 \(\uparrow  {295}\)
\end{tcolorbox}

Tversky, A. (1969). Intransitivity of preferences. Psychological Review 76(1), 31-48. \(\uparrow  {133}\)

\begin{tcolorbox}
特沃斯基(Tversky)，A.(1969 年)。偏好(Preference)的非传递性。《心理学评论》(Psychological Review)76(1) 卷，第 31 - 48 页。 \(\uparrow  {133}\)
\end{tcolorbox}

van Damme, E. (1987). Stability and Perfection of Nash Equilibria. Springer, Berlin. ↑74, 166

\begin{tcolorbox}
范·达姆(van Damme)，E.(1987 年)。《纳什均衡(Nash Equilibrium)的稳定性与完美性》(Stability and Perfection of Nash Equilibria)。施普林格出版社(Springer)，德国柏林。 ↑74, 166
\end{tcolorbox}

von Neumann, J. (1928). Zur Theorie der Gesellschaftsspiele. Mathematische Annalen 100(1), 295-320. Translated as: On the theory of games of strategy. In: Contributions to the Theory of Games, Vol. IV, eds. A. W. Tucker and R. D. Luce, Annals of Mathematics Studies vol. 40, Princeton University Press, 1959, pp. 13-42. \(\uparrow  {205},{206},{213},{215},{216},{220}\)

\begin{tcolorbox}
冯·诺伊曼，J.(1928)。关于社会博弈理论。《数学年刊》100(1)，295 - 320。翻译为:论策略博弈理论。载于:《博弈论贡献》，第四卷，A. W. 塔克和R. D. 卢斯编，《数学研究年刊》第40卷，普林斯顿大学出版社，1959年，第13 - 42页。 \(\uparrow  {205},{206},{213},{215},{216},{220}\)
\end{tcolorbox}

von Neumann, J. and M. Fréchet (1953). Communication on the Borel notes. Econometrica 21(1), 124-127. ↑220

\begin{tcolorbox}
冯·诺伊曼，J. 和M. 弗雷歇(1953)。关于博雷尔笔记的通信。《计量经济学》21(1)，124 - 127。↑220
\end{tcolorbox}

von Neumann, J. and O. Morgenstern (1947). Theory of Games and Economic Behavior, 2nd ed. Princeton University Press, Princeton, NJ. \(\uparrow  {103},{104},{133},{220}\)

\begin{tcolorbox}
冯·诺伊曼，J. 和O. 摩根斯坦(1947)。《博弈论与经济行为》，第二版。普林斯顿大学出版社，新泽西州普林斯顿。 \(\uparrow  {103},{104},{133},{220}\)
\end{tcolorbox}

von Stackelberg, H. (1934). Marktform und Gleichgewicht. Springer, Berlin. Translated as: Market Structure and Equilibrium, Springer, Heidelberg, 2011. \(\uparrow  {98}\)

\begin{tcolorbox}
冯·施塔克尔贝格，H.(1934)。《市场形式与均衡》。柏林:施普林格出版社。翻译为:《市场结构与均衡》，海德堡:施普林格出版社，2011年。 \(\uparrow  {98}\)
\end{tcolorbox}

von Stengel, B. (1996). Efficient computation of behavior strategies. Games and Economic Behavior 14(2), 220-246. ↑ 295

\begin{tcolorbox}
冯·施滕格尔，B.(1996)。行为策略(behavior strategy)的高效计算。《博弈与经济行为》14(2)，220 - 246。↑ 295
\end{tcolorbox}

von Stengel, B. (1999). New maximal numbers of equilibria in bimatrix games. Discrete and Computational Geometry 21(4), 557-568. \(\uparrow  {260}\)

\begin{tcolorbox}
冯·施滕格尔，B.(1999)。双矩阵博弈(bimatrix game)中均衡(equilibrium)数量的新最大值。《离散与计算几何》21(4)，557 - 568。 \(\uparrow  {260}\)
\end{tcolorbox}

von Stengel, B. (2002). Computing equilibria for two-person games. In: Handbook of Game Theory with Economic Applications, edited by R. J. Aumann and S. Hart, volume 3, 1723-1759. North-Holland, Amsterdam. \(\uparrow  {259},{260}\)

\begin{tcolorbox}
冯·施滕格尔，B.(2002)。两人博弈的均衡计算。载于:《博弈论及其经济应用手册》，R. J. 奥曼和S. 哈特编，第三卷，1723 - 1759页。荷兰北荷兰出版社，阿姆斯特丹。 \(\uparrow  {259},{260}\)
\end{tcolorbox}

von Stengel, B. (2010). Follower payoffs in symmetric duopoly games. Games and Economic Behavior 69(2), 512-516. \(\uparrow  {98}\)

\begin{tcolorbox}
冯·施滕格尔，B.(2010)。对称双头垄断博弈中跟随者(follower)的收益。《博弈与经济行为》69(2)，512 - 516。 \(\uparrow  {98}\)
\end{tcolorbox}

von Stengel, B. (2021). Finding Nash equilibria of two-player games. Preprint arXiv:2102.04580. ↑ 166, 259

\begin{tcolorbox}
冯·施滕格尔，B.(2021)。寻找两人博弈的纳什均衡。预印本arXiv:2102.04580。↑ 166, 259
\end{tcolorbox}

von Stengel, B. and S. Zamir (2010). Leadership games with convex strategy sets. Games and Economic Behavior 69(2), 446-457. ↑ 350

\begin{tcolorbox}
冯·施滕格尔，B. 和S. 扎米尔(2010)。具有凸策略集的领导博弈(leadership game)。《博弈与经济行为》69(2)，446 - 457。↑ 350
\end{tcolorbox}

Wardrop, J. G. (1952). Some theoretical aspects of road traffic research. Proceedings of the Institution of Civil Engineers 1(3), 325-362. \(\uparrow  {50}\)

\begin{tcolorbox}
沃德罗普，J. G.(1952)。道路交通研究的一些理论方面。《土木工程师学会会刊》1(3)，325 - 362。 \(\uparrow  {50}\)
\end{tcolorbox}

Wythoff, W. A. (1907). A modification of the game of Nim. Nieuw Archief voor Wiskunde 7(2), 199-202. \(\uparrow  {31}\)

\begin{tcolorbox}
威索夫，W. A.(1907)。尼姆游戏(Nim)的一种变体。《新数学档案馆》7(2)，199 - 202。 \(\uparrow  {31}\)
\end{tcolorbox}

Young, H. P. (2004). Strategic Learning and Its Limits. Oxford University Press, Oxford. ↑349

\begin{tcolorbox}
扬，H. P.(2004)。《战略学习及其局限》。牛津大学出版社，牛津。↑349
\end{tcolorbox}

Zermelo, E. (1913). Über eine Anwendung der Mengenlehre auf die Theorie des Schachspiels. In: Proc. Fifth Congress Mathematicians, Cambridge 1912, 501-504. Cambridge University Press, Cambridge, UK. Translated in Schwalbe and Walker (2001). ↑98

\begin{tcolorbox}
策梅洛，E.(1913)。关于集合论在国际象棋博弈理论中的应用。载于:《1912年第五届数学家大会会议记录》，剑桥，1912年，501 - 504页。剑桥大学出版社，英国剑桥。译文见施瓦尔贝和沃克(2001)。↑98
\end{tcolorbox}

Ziegler, G. M. (1995). Lectures on Polytopes, volume 152 of Graduate Text in Mathematics. Springer, New York. ↑260

\begin{tcolorbox}
齐格勒，G. M.(1995)。《有界多面体(polytope)讲义》，《数学研究生教材》第152卷。施普林格出版社，纽约。↑260
\end{tcolorbox}

\section*{Index}

\begin{tcolorbox}
\section*{索引}
\end{tcolorbox}

abelian group, 1, 10, 14

\begin{tcolorbox}
阿贝尔群(abelian group), 1, 10, 14
\end{tcolorbox}

affine combination, 174

\begin{tcolorbox}
仿射组合(affine combination), 174
\end{tcolorbox}

Allais paradox, 121, 133

\begin{tcolorbox}
阿莱悖论(Allais paradox), 121, 133
\end{tcolorbox}

almost complementary, 251

\begin{tcolorbox}
几乎互补(almost complementary), 251
\end{tcolorbox}

almost completely labeled, 231, 251, 259

\begin{tcolorbox}
几乎完全标记(almost completely labeled), 231, 251, 259
\end{tcolorbox}

angle, 192, 194

\begin{tcolorbox}
角(angle), 192, 194
\end{tcolorbox}

antisymmetric, 8

\begin{tcolorbox}
反对称(antisymmetric), 8
\end{tcolorbox}

artificial equilibrium, 230, 231, 235, 250,

\begin{tcolorbox}
人工均衡(artificial equilibrium), 230, 231, 235, 250,
\end{tcolorbox}

251

associative, 10

\begin{tcolorbox}
结合律(associative), 10
\end{tcolorbox}

Aumann, Robert, xii, 74, 334

\begin{tcolorbox}
罗伯特·奥曼(Aumann, Robert), xii, 74, 334
\end{tcolorbox}

background material, xi

\begin{tcolorbox}
背景材料(background material), xi
\end{tcolorbox}

converging subsequences, 188

\begin{tcolorbox}
收敛子列(converging subsequences), 188
\end{tcolorbox}

directed graphs and trees, 80

\begin{tcolorbox}
有向图与树图(directed graphs and trees), 80
\end{tcolorbox}

Euclidean norm and distance, 187

\begin{tcolorbox}
欧几里得范数与距离(Euclidean norm and distance), 187
\end{tcolorbox}

expected value, 81

\begin{tcolorbox}
期望值(expected value), 81
\end{tcolorbox}

line segments, 130

\begin{tcolorbox}
线段(line segments), 130
\end{tcolorbox}

matrix multiplication, 141, 143

\begin{tcolorbox}
矩阵乘法，141，143
\end{tcolorbox}

partial orders, 8

\begin{tcolorbox}
偏序，8
\end{tcolorbox}

undirected graphs, 170

\begin{tcolorbox}
无向图，170
\end{tcolorbox}

visualizing matrix multiplication,

\begin{tcolorbox}
矩阵乘法的可视化
\end{tcolorbox}

143

backward induction, 83, 314, 319

\begin{tcolorbox}
逆向归纳法，83，314，319
\end{tcolorbox}

more special than SPE, 284, 293

\begin{tcolorbox}
比子博弈完美均衡更特殊，284，293
\end{tcolorbox}

bargaining axioms, 302, 303

\begin{tcolorbox}
讨价还价公理，302，303
\end{tcolorbox}

bargaining solution, 303, 329

\begin{tcolorbox}
讨价还价解，303，329
\end{tcolorbox}

barycenter, 176-178, 190

\begin{tcolorbox}
重心，176 - 178，190
\end{tcolorbox}

barycentric coordinates, 176

\begin{tcolorbox}
重心坐标，176
\end{tcolorbox}

basic solution, 252

\begin{tcolorbox}
基解，252
\end{tcolorbox}

basic variables, 251

\begin{tcolorbox}
基变量，251
\end{tcolorbox}

Battle of the Sexes, 60

\begin{tcolorbox}
性别战，60
\end{tcolorbox}

behavior strategy, 280, 285

\begin{tcolorbox}
行为策略，280，285
\end{tcolorbox}

best response, 46, 57, 58

\begin{tcolorbox}
最优反应，46，57，58
\end{tcolorbox}

best-response condition, 135, 145, 148

\begin{tcolorbox}
最优反应条件，135，145，148
\end{tcolorbox}

best-response diagram, 227

\begin{tcolorbox}
最优反应示意图(best-response diagram)，227
\end{tcolorbox}

bimatrix game, 142

\begin{tcolorbox}
双矩阵博弈(bimatrix game)，142
\end{tcolorbox}

binary system, 1, 5

\begin{tcolorbox}
二进制(binary system)，1，5
\end{tcolorbox}

bluff, 269

\begin{tcolorbox}
虚张声势(bluff)，269
\end{tcolorbox}

Bolzano-Weierstrass theorem, 188

\begin{tcolorbox}
波尔查诺 - 魏尔施特拉斯定理(Bolzano-Weierstrass theorem)，188
\end{tcolorbox}

Braess paradox, 43, 44, 49

\begin{tcolorbox}
布雷斯悖论(Braess paradox)，43，44，49
\end{tcolorbox}

Brouwer's fixed-point theorem, 151, 153,

\begin{tcolorbox}
布劳威尔不动点定理(Brouwer's fixed-point theorem)，151，153，
\end{tcolorbox}

189,195,203,258

castle, 234

\begin{tcolorbox}
城堡(castle)，234
\end{tcolorbox}

Chicken, 60, 74, 334

\begin{tcolorbox}
胆小鬼博弈(Chicken)，60，74，334
\end{tcolorbox}

Chomp, 31, 33, 35

\begin{tcolorbox}
巧取游戏(Chomp)，31，33，35
\end{tcolorbox}

cold game, 30

\begin{tcolorbox}
冷静博弈(cold game)，30
\end{tcolorbox}

column shift, 239, 240

\begin{tcolorbox}
列移位(column shift)，239，240
\end{tcolorbox}

commitment game, 93

\begin{tcolorbox}
承诺博弈(commitment game)，93
\end{tcolorbox}

common knowledge, 55, 132

\begin{tcolorbox}
共同知识(common knowledge)，55，132
\end{tcolorbox}

commutative, 10

\begin{tcolorbox}
交换律(commutative)，10
\end{tcolorbox}

compact set, 149, 151, 171, 190, 212, 248,

\begin{tcolorbox}
紧集(compact set)，149，151，171，190，212，248，
\end{tcolorbox}

302

complementary, 251

\begin{tcolorbox}
互补性, 251
\end{tcolorbox}

complementary pivoting, 252, 255, 260

\begin{tcolorbox}
互补换基法, 252, 255, 260
\end{tcolorbox}

completely labeled, 172, 173, 181,

\begin{tcolorbox}
完全标记, 172, 173, 181,
\end{tcolorbox}

227-229, 247, 249, 251

concave function, 130

\begin{tcolorbox}
凹函数, 130
\end{tcolorbox}

concave utility function, 131, 312

\begin{tcolorbox}
凹效用函数, 131, 312
\end{tcolorbox}

conditional probability, 269, 337

\begin{tcolorbox}
条件概率, 269, 337
\end{tcolorbox}

Condorcet Paradox, 110, 133

\begin{tcolorbox}
孔多塞悖论, 110, 133
\end{tcolorbox}

congestion game, 45

\begin{tcolorbox}
拥塞博弈, 45
\end{tcolorbox}

congestion network, 40, 43, 44

\begin{tcolorbox}
拥塞网络, 40, 43, 44
\end{tcolorbox}

continuity axiom, 108

\begin{tcolorbox}
连续性公理, 108
\end{tcolorbox}

Continuous Ultimatum game, 316

\begin{tcolorbox}
连续最后通牒博弈, 316
\end{tcolorbox}

contractive map, 195

\begin{tcolorbox}
压缩映射, 195
\end{tcolorbox}

convex combination, 150, 165, 174

\begin{tcolorbox}
凸组合, 150, 165, 174
\end{tcolorbox}

convex hull, 174, 301

\begin{tcolorbox}
凸包, 174, 301
\end{tcolorbox}

convex set, 149, 174

\begin{tcolorbox}
凸集, 149, 174
\end{tcolorbox}

cooperative game theory, 299, 330

\begin{tcolorbox}
合作博弈论, 299, 330
\end{tcolorbox}

coordination game, 61, 236

\begin{tcolorbox}
协调博弈, 61, 236
\end{tcolorbox}

copycat principle, 13

\begin{tcolorbox}
模仿原则, 13
\end{tcolorbox}

correlated equilibrium, 338, 339

\begin{tcolorbox}
相关均衡, 338, 339
\end{tcolorbox}

correlated moves, 282, 290

\begin{tcolorbox}
相关行动, 282, 290
\end{tcolorbox}

correlated strategies, 336

\begin{tcolorbox}
相关策略, 336
\end{tcolorbox}

correlation device, 337, 338

\begin{tcolorbox}
相关机制, 337, 338
\end{tcolorbox}

cost, 41, 206

\begin{tcolorbox}
成本, 41, 206
\end{tcolorbox}

Cournot duopoly, 66, 74

\begin{tcolorbox}
古诺双寡头博弈, 66, 74
\end{tcolorbox}

Cram, 25, 32, 34, 36

\begin{tcolorbox}
塞砖棋, 25, 32, 34, 36
\end{tcolorbox}

decision theory, 110

\begin{tcolorbox}
决策理论, 110
\end{tcolorbox}

Dedekind cut, 31

\begin{tcolorbox}
戴德金分割, 31
\end{tcolorbox}

degenerate game, 162, 164, 233, 238

\begin{tcolorbox}
退化博弈, 162, 164, 233, 238
\end{tcolorbox}

degree of a map (Brouwer), 203

\begin{tcolorbox}
映射度 (布劳威尔), 203
\end{tcolorbox}

degree of graph node, 170

\begin{tcolorbox}
图节点的度, 170
\end{tcolorbox}

descendant, 80, 91

\begin{tcolorbox}
后代节点, 80, 91
\end{tcolorbox}

deviation, 41, 46

\begin{tcolorbox}
偏离, 41, 46
\end{tcolorbox}

dictionary, 251

\begin{tcolorbox}
字典，251
\end{tcolorbox}

difference trick, 156, 209

\begin{tcolorbox}
差分技巧，156，209
\end{tcolorbox}

digraph, 45

\begin{tcolorbox}
有向图，45
\end{tcolorbox}

digraph game, 17, 31, 35

\begin{tcolorbox}
有向图博弈，17，31，35
\end{tcolorbox}

dimension, 174

\begin{tcolorbox}
维度，174
\end{tcolorbox}

dominance solvable, 53, 65, 74, 342

\begin{tcolorbox}
占优可解，53，65，74，342
\end{tcolorbox}

dominated

\begin{tcolorbox}
被占优
\end{tcolorbox}

strategy, 342

\begin{tcolorbox}
策略，342
\end{tcolorbox}

dominated option, 27

\begin{tcolorbox}
被占优选项，27
\end{tcolorbox}

Domineering, 24

\begin{tcolorbox}
纵横棋，24
\end{tcolorbox}

domino, 24, 32, 34, 36

\begin{tcolorbox}
多米诺骨牌，24，32，34，36
\end{tcolorbox}

downward closed set, 313

\begin{tcolorbox}
下闭集，313
\end{tcolorbox}

duality of linear programming, 218

\begin{tcolorbox}
线性规划的对偶性，218
\end{tcolorbox}

duplicate label, 231

\begin{tcolorbox}
重复标签，231
\end{tcolorbox}

edge, 40, 170, 176, 181, 230

\begin{tcolorbox}
边，40，170，176，181，230
\end{tcolorbox}

end node of a path, 170, 232

\begin{tcolorbox}
路径的终止节点，170，232
\end{tcolorbox}

ending condition, 2, 7, 10

\begin{tcolorbox}
终止条件, 2, 7, 10
\end{tcolorbox}

endpoint of an edge, 170

\begin{tcolorbox}
一条边的端点, 170
\end{tcolorbox}

entering variable, 252

\begin{tcolorbox}
入基变量, 252
\end{tcolorbox}

equilibria (plural of equilibrium), 59

\begin{tcolorbox}
均衡 (equilibrium 的复数形式), 59
\end{tcolorbox}

equilibrium, 53, 58, 73

\begin{tcolorbox}
均衡, 53, 58, 73
\end{tcolorbox}

correlated, 338, 339

\begin{tcolorbox}
相关的, 338, 339
\end{tcolorbox}

in a congestion game, 41, 46, 48

\begin{tcolorbox}
在拥塞博弈中, 41, 46, 48
\end{tcolorbox}

mixed, 148, 152

\begin{tcolorbox}
混合的, 148, 152
\end{tcolorbox}

Nash, 59, 338

\begin{tcolorbox}
纳什, 59, 338
\end{tcolorbox}

pure, 53

\begin{tcolorbox}
纯的, 53
\end{tcolorbox}

refinement, 62, 74 selection, 62 equilibrium path, 95 equivalence of impartial games, 11 of partizan games, 25, 31 relation, 11 strategic, 239, 259 strong strategic, 241, 259

\begin{tcolorbox}
精炼, 62, 74 选择, 62 均衡路径, 95 无偏博弈的等价性, 11 有偏博弈的等价性, 25, 31 关系, 11 策略性的, 239, 259 强策略性的, 241, 259
\end{tcolorbox}

Euclidean distance, 187

\begin{tcolorbox}
欧几里得距离, 187
\end{tcolorbox}

Euclidean norm, 187

\begin{tcolorbox}
欧几里得范数, 187
\end{tcolorbox}

expected payoff, 85, 148

\begin{tcolorbox}
期望效用, 85, 148
\end{tcolorbox}

expected value, 81

\begin{tcolorbox}
期望值, 81
\end{tcolorbox}

expected-utility function, 103, 106, 112, 117, 124, 208

\begin{tcolorbox}
期望效用函数, 103, 106, 112, 117, 124, 208
\end{tcolorbox}

exponential size of game tree, 90

\begin{tcolorbox}
博弈树的指数级规模，90
\end{tcolorbox}

extreme point, 176, 204

\begin{tcolorbox}
极点，176，204
\end{tcolorbox}

face, 176, 245

\begin{tcolorbox}
面，176，245
\end{tcolorbox}

facet, 176

\begin{tcolorbox}
极面，176
\end{tcolorbox}

Farkas's Lemma, 220

\begin{tcolorbox}
法卡斯引理(Farkas's Lemma)，220
\end{tcolorbox}

first-mover advantage, 96

\begin{tcolorbox}
先行优势，96
\end{tcolorbox}

flow, 41

\begin{tcolorbox}
流，41
\end{tcolorbox}

atomic (non-splittable), 39, 48 splittable, 40, 48, 50

\begin{tcolorbox}
原子(不可拆分)，39，48 可拆分，40，48，50
\end{tcolorbox}

follower, 96

\begin{tcolorbox}
跟随者，96
\end{tcolorbox}

Freudenthal triangulation, 197

\begin{tcolorbox}
弗赖登塔尔三角剖分，197
\end{tcolorbox}

full-dimensional, 197

\begin{tcolorbox}
满维，197
\end{tcolorbox}

Gambit software, xii

\begin{tcolorbox}
Gambit 软件，xii
\end{tcolorbox}

game position, 7

\begin{tcolorbox}
博弈位置，7
\end{tcolorbox}

game sum, 10

\begin{tcolorbox}
博弈之和，10
\end{tcolorbox}

Game Theory Explorer (GTE), xii

\begin{tcolorbox}
博弈论探索器(GTE)，xii
\end{tcolorbox}

game tree, 3, 79, 270

\begin{tcolorbox}
博弈树，3，79，270
\end{tcolorbox}

generic game, 84

\begin{tcolorbox}
通有博弈 (generic game), 84
\end{tcolorbox}

goalpost method, 159, 208, 226

\begin{tcolorbox}
球门柱法 (goalpost method), 159, 208, 226
\end{tcolorbox}

graph, 170

\begin{tcolorbox}
图 (graph), 170
\end{tcolorbox}

Grundy value, 16, 31

\begin{tcolorbox}
Grundy 值 (Grundy value), 16, 31
\end{tcolorbox}

Hackenbush

\begin{tcolorbox}
哈肯布什游戏 (Hackenbush)
\end{tcolorbox}

Black-White, 30, 37

\begin{tcolorbox}
黑白哈肯布什游戏 (Black-White), 30, 37
\end{tcolorbox}

Gray, 36

\begin{tcolorbox}
灰色哈肯布什游戏 (Gray), 36
\end{tcolorbox}

halfline, 192

\begin{tcolorbox}
射线 (halfline), 192
\end{tcolorbox}

halfspace, 193, 245

\begin{tcolorbox}
半空间 (halfspace), 193, 245
\end{tcolorbox}

heap, see Nim heap history of moves, 3, 81

\begin{tcolorbox}
堆，见尼姆堆 行动的历史 (history of moves), 3, 81
\end{tcolorbox}

Homo Ludens (Huizinga), 330

\begin{tcolorbox}
《游戏的人》(Homo Ludens (Huizinga)), 330
\end{tcolorbox}

hot game, 30

\begin{tcolorbox}
热情博弈 (hot game), 30
\end{tcolorbox}

hotel with seaview, 188, 204

\begin{tcolorbox}
海景酒店 (hotel with seaview), 188, 204
\end{tcolorbox}

hyperplane, 193

\begin{tcolorbox}
超平面 (hyperplane), 193
\end{tcolorbox}

identity matrix, 344

\begin{tcolorbox}
单位矩阵 (identity matrix), 344
\end{tcolorbox}

impartial game, 2

\begin{tcolorbox}
无偏博弈 (impartial game), 2
\end{tcolorbox}

imperfect recall, 275

\begin{tcolorbox}
不完美记忆(imperfect recall)，275
\end{tcolorbox}

incentive constraints, 338

\begin{tcolorbox}
激励约束(incentive constraints)，338
\end{tcolorbox}

independence axiom, 104, 108, 121

\begin{tcolorbox}
独立性公理(independence axiom)，104，108，121
\end{tcolorbox}

index of an equilibrium, 259

\begin{tcolorbox}
均衡的指数(index of an equilibrium)，259
\end{tcolorbox}

indifference, 106, 112, 139

\begin{tcolorbox}
无差异(indifference)，106，112，139
\end{tcolorbox}

information set, 265, 266

\begin{tcolorbox}
信息集(information set)，265，266
\end{tcolorbox}

Inspection game, 137

\begin{tcolorbox}
监管博弈(Inspection game)，137
\end{tcolorbox}

insurance

\begin{tcolorbox}
保险
\end{tcolorbox}

accidental breakdown, 109

\begin{tcolorbox}
意外故障，109
\end{tcolorbox}

car rental, 109, 110

\begin{tcolorbox}
汽车租赁，109，110
\end{tcolorbox}

intransitive preferences, 110, 114, 133

\begin{tcolorbox}
非传递偏好(intransitive preferences)，110，114，133
\end{tcolorbox}

invariant under utility scaling, 303

\begin{tcolorbox}
效用缩放不变性(invariant under utility scaling)，303
\end{tcolorbox}

irrelevant alternatives, 303

\begin{tcolorbox}
无关选项(irrelevant alternatives)，303
\end{tcolorbox}

iterated elimination of dominated

\begin{tcolorbox}
重复剔除劣势策略(iterated elimination of dominated)
\end{tcolorbox}

strategies, 64, 65, 69, 71, 222

\begin{tcolorbox}
策略，64，65，69，71，222
\end{tcolorbox}

iterated-offers bargaining game, 317-330

\begin{tcolorbox}
重复讨价还价博弈(iterated-offers bargaining game)，317 - 330
\end{tcolorbox}

Kayles, 17, 23, 31

\begin{tcolorbox}
凯尔斯博弈，17，23，31
\end{tcolorbox}

Kuhn's theorem, 284

\begin{tcolorbox}
库恩定理，284
\end{tcolorbox}

label, 172, 227

\begin{tcolorbox}
标签，172，227
\end{tcolorbox}

Last Year in Marienbad, 4

\begin{tcolorbox}
《去年在马里昂巴德》，4
\end{tcolorbox}

leader, 96

\begin{tcolorbox}
领导，96
\end{tcolorbox}

leadership game, 96

\begin{tcolorbox}
领导博弈，96
\end{tcolorbox}

leaf, 80

\begin{tcolorbox}
叶节点，80
\end{tcolorbox}

learning, 349

\begin{tcolorbox}
学习，349
\end{tcolorbox}

leaving variable, 252

\begin{tcolorbox}
离基变量，252
\end{tcolorbox}

Left player, 24

\begin{tcolorbox}
左玩家，24
\end{tcolorbox}

Lemke's algorithm, 260

\begin{tcolorbox}
莱姆基算法，260
\end{tcolorbox}

Lemke-Howson

\begin{tcolorbox}
莱姆基 - 豪森
\end{tcolorbox}

algorithm, 232

\begin{tcolorbox}
算法，232
\end{tcolorbox}

graph, 233, 235, 259

\begin{tcolorbox}
图，233，235，259
\end{tcolorbox}

path, 231, 232, 244

\begin{tcolorbox}
路径，231，232，244
\end{tcolorbox}

lexico-minimum ratio test, 258

\begin{tcolorbox}
字典序最小比率检验，258
\end{tcolorbox}

lexico-positive, 257

\begin{tcolorbox}
字典正序(lexico-positive)，257
\end{tcolorbox}

lexicographic order, 105, 114, 123, 124

\begin{tcolorbox}
字典序(lexicographic order)，105，114，123，124
\end{tcolorbox}

LH, see Lemke-Howson

\begin{tcolorbox}
LH，见莱姆克 - 豪森法(Lemke - Howson)
\end{tcolorbox}

line segment, 130

\begin{tcolorbox}
线段(line segment)，130
\end{tcolorbox}

linear complementarity problem, 260

\begin{tcolorbox}
线性互补问题(linear complementarity problem)，260
\end{tcolorbox}

linear programming, 218, 221, 255, 260,

\begin{tcolorbox}
线性规划(linear programming)，218，221，255，260
\end{tcolorbox}

296

linear programming duality, 206, 218

\begin{tcolorbox}
线性规划对偶(linear programming duality)，206，218
\end{tcolorbox}

losing position, 3, 25

\begin{tcolorbox}
失败位置(losing position)，3，25
\end{tcolorbox}

lottery, 106, 111

\begin{tcolorbox}
彩票(lottery)，106，111
\end{tcolorbox}

simple, 104, 117

\begin{tcolorbox}
简单，104，117
\end{tcolorbox}

lower envelope, 209

\begin{tcolorbox}
下包络(lower envelope)，209
\end{tcolorbox}

L-position, 25

\begin{tcolorbox}
左胜位置(L - position)，25
\end{tcolorbox}

majority vote, 110

\begin{tcolorbox}
多数人规则(majority vote)，110
\end{tcolorbox}

Marienbad, 4

\begin{tcolorbox}
马里昂巴德(Marienbad)，4
\end{tcolorbox}

Markov chain, 344

\begin{tcolorbox}
马尔可夫链(Markov chain)，344
\end{tcolorbox}

Maschler, Michael, 334

\begin{tcolorbox}
迈克尔·马施勒(Maschler, Michael)，334
\end{tcolorbox}

Matching Pennies, 72

\begin{tcolorbox}
硬币匹配 (博弈)，72
\end{tcolorbox}

matrix game, 207

\begin{tcolorbox}
矩阵博弈，207
\end{tcolorbox}

matrix multiplication, 141, 143

\begin{tcolorbox}
矩阵乘法，141，143
\end{tcolorbox}

max-min strategy, 138, 210, 211

\begin{tcolorbox}
极大极小策略，138，210，211
\end{tcolorbox}

maximizer, 207

\begin{tcolorbox}
最大值点，207
\end{tcolorbox}

mex rule, 15, 16, 31

\begin{tcolorbox}
mex 规则，15，16，31
\end{tcolorbox}

min-max strategy, 211, 212

\begin{tcolorbox}
极小极大策略，211，212
\end{tcolorbox}

minimal element, 8

\begin{tcolorbox}
极小元，8
\end{tcolorbox}

minimax theorem, 213, 214, 216, 344, 347

\begin{tcolorbox}
极小极大定理，213，214，216，344，347
\end{tcolorbox}

minimizer, 207

\begin{tcolorbox}
最小值点，207
\end{tcolorbox}

minimum ratio test, 255

\begin{tcolorbox}
最小比值检验，255
\end{tcolorbox}

misère Nim, 4, 32

\begin{tcolorbox}
反向尼姆游戏，4，32
\end{tcolorbox}

misère play, 2, 14

\begin{tcolorbox}
反向玩法，2，14
\end{tcolorbox}

missing label, 231

\begin{tcolorbox}
缺失标签，231
\end{tcolorbox}

mixed equilibrium, 148, 152, 228

\begin{tcolorbox}
混合均衡，148，152，228
\end{tcolorbox}

mixed extension, 139, 148, 212

\begin{tcolorbox}
混合扩展，139，148，212
\end{tcolorbox}

mixed strategy, 143

\begin{tcolorbox}
混合策略 (mixed strategy)，143
\end{tcolorbox}

mixed-strategy profile, 148

\begin{tcolorbox}
混合策略组合 (mixed-strategy profile)，148
\end{tcolorbox}

mixed-strategy set, 144

\begin{tcolorbox}
混合策略集 (mixed-strategy set)，144
\end{tcolorbox}

negative game, 11, 13, 25

\begin{tcolorbox}
负博弈 (negative game)，11，13，25
\end{tcolorbox}

neighbor, 170

\begin{tcolorbox}
邻居 (neighbor)，170
\end{tcolorbox}

Nim, 2, 32

\begin{tcolorbox}
尼姆游戏 (Nim)，2，32
\end{tcolorbox}

Nim heap, 2, 15

\begin{tcolorbox}
尼姆堆 (Nim heap)，2，15
\end{tcolorbox}

notation \(n * ,{15}\)

\begin{tcolorbox}
符号 \(n * ,{15}\)
\end{tcolorbox}

Nim sum, 5, 20

\begin{tcolorbox}
尼姆和 (Nim sum)，5，20
\end{tcolorbox}

Nim value, 16, 27

\begin{tcolorbox}
尼姆值 (Nim value)，16，27
\end{tcolorbox}

nimber, 27

\begin{tcolorbox}
尼姆数 (nimber)，27
\end{tcolorbox}

node, 170, 230

\begin{tcolorbox}
节点 (node)，170，230
\end{tcolorbox}

nonbasic variables, 252

\begin{tcolorbox}
非基变量 (nonbasic variables)，252
\end{tcolorbox}

nondegenerate, 254

\begin{tcolorbox}
非退化的 (nondegenerate)，254
\end{tcolorbox}

normal form, 74

\begin{tcolorbox}
标准型 (normal form)，74
\end{tcolorbox}

normal play, 2

\begin{tcolorbox}
正常规则 (normal play)，2
\end{tcolorbox}

normal vector, 193, 245

\begin{tcolorbox}
法向量 (normal vector), 193, 245
\end{tcolorbox}

normative, 112, 114, 120, 133

\begin{tcolorbox}
规范的 (normative), 112, 114, 120, 133
\end{tcolorbox}

Northcott's game, 31, 34

\begin{tcolorbox}
诺斯科特博弈 (Northcott's game), 31, 34
\end{tcolorbox}

N-position, 25

\begin{tcolorbox}
先手胜位置 (N-position), 25
\end{tcolorbox}

number (partizan game), 27, 28

\begin{tcolorbox}
数 (有偏博弈中的) (number (partizan game)), 27, 28
\end{tcolorbox}

optimal strategy, 215

\begin{tcolorbox}
最优策略 (optimal strategy), 215
\end{tcolorbox}

option, 3, 7, 10, 26

\begin{tcolorbox}
选项 (option), 3, 7, 10, 26
\end{tcolorbox}

order

\begin{tcolorbox}
顺序 (order)
\end{tcolorbox}

partial, 8

\begin{tcolorbox}
部分的 (partial), 8
\end{tcolorbox}

total, 8

\begin{tcolorbox}
全部的 (total), 8
\end{tcolorbox}

origin of utility scale, 104, 107, 116

\begin{tcolorbox}
效用尺度原点 (origin of utility scale), 104, 107, 116
\end{tcolorbox}

orthogonal, 191

\begin{tcolorbox}
正交的 (orthogonal), 191
\end{tcolorbox}

outcome class, 3, 25

\begin{tcolorbox}
结果类 (outcome class), 3, 25
\end{tcolorbox}

outcome set for decisions, 111

\begin{tcolorbox}
决策的结果集 (outcome set for decisions), 111
\end{tcolorbox}

parallel edges, 40, 41, 49

\begin{tcolorbox}
平行边 (parallel edges), 40, 41, 49
\end{tcolorbox}

Pareto-optimal, 303

\begin{tcolorbox}
帕累托最优 (Pareto-optimal), 303
\end{tcolorbox}

parity argument, 170, 184, 203, 232, 234

\begin{tcolorbox}
奇偶性论证(parity argument)，170，184，203，232，234
\end{tcolorbox}

partial order, 8, 32, 142

\begin{tcolorbox}
偏序(partial order)，8，32，142
\end{tcolorbox}

partial strategy profile, 339

\begin{tcolorbox}
部分策略组合(partial strategy profile)，339
\end{tcolorbox}

partizan game,2,24,37,38

\begin{tcolorbox}
有偏博弈(partizan game)，2，24，37，38
\end{tcolorbox}

path, 45, 80, 170

\begin{tcolorbox}
路径(path)，45，80，170
\end{tcolorbox}

payoff equivalent, 62

\begin{tcolorbox}
收益等价(payoff equivalent)，62
\end{tcolorbox}

perfect information, 2

\begin{tcolorbox}
完美信息(perfect information)，2
\end{tcolorbox}

perfect recall, 275-279, 285

\begin{tcolorbox}
完美记忆(perfect recall)，275 - 279，285
\end{tcolorbox}

permutation, 198

\begin{tcolorbox}
排列(permutation)，198
\end{tcolorbox}

Pigou network, 40, 42, 44, 49

\begin{tcolorbox}
庇古网络(Pigou network)，40，42，44，49
\end{tcolorbox}

pivoting, 252

\begin{tcolorbox}
换基操作(pivoting)，252
\end{tcolorbox}

play of a game, 77

\begin{tcolorbox}
博弈的进行(play of a game)，77
\end{tcolorbox}

Poker, 268, 297

\begin{tcolorbox}
扑克(Poker)，268，297
\end{tcolorbox}

Poker Nim, 15, 16, 17

\begin{tcolorbox}
扑克尼姆游戏(Poker Nim)，15，16，17
\end{tcolorbox}

polyhedron, 245, 341

\begin{tcolorbox}
多面体(polyhedron)，245，341
\end{tcolorbox}

polytope, 175, 245

\begin{tcolorbox}
有界多面体(polytope)，175，245
\end{tcolorbox}

potential function, 39, 46, 48, 49

\begin{tcolorbox}
势函数(potential function)，39，46，48，49
\end{tcolorbox}

potential game, 50, 74

\begin{tcolorbox}
势博弈(potential game)，50，74
\end{tcolorbox}

PPAD, 259

P-position, 25

\begin{tcolorbox}
后手胜位置(P-position)，25
\end{tcolorbox}

preference

\begin{tcolorbox}
偏好(preference)
\end{tcolorbox}

relation, 111

\begin{tcolorbox}
关系，111
\end{tcolorbox}

strict, 106, 112

\begin{tcolorbox}
严格，106，112
\end{tcolorbox}

Price of Anarchy, 49

\begin{tcolorbox}
无政府代价(Price of Anarchy)，49
\end{tcolorbox}

prime number decomposition theorem, 9

\begin{tcolorbox}
质数分解定理(prime number decomposition theorem)，9
\end{tcolorbox}

Prisoner's Dilemma, 55, 73, 300

\begin{tcolorbox}
囚徒困境(Prisoner's Dilemma)，55，73，300
\end{tcolorbox}

product graph, 231

\begin{tcolorbox}
乘积图(product graph)，231
\end{tcolorbox}

projection, 171, 191

\begin{tcolorbox}
投影(projection)，171，191
\end{tcolorbox}

Quality game, 62

\begin{tcolorbox}
产品质量博弈(Quality game)，62
\end{tcolorbox}

quantity competition, 66

\begin{tcolorbox}
产量竞争(quantity competition)，66
\end{tcolorbox}

Queen-move game, 22, 34

\begin{tcolorbox}
皇后移动博弈(Queen-move game)，22，34
\end{tcolorbox}

random variable, 81

\begin{tcolorbox}
随机变量(random variable)，81
\end{tcolorbox}

rationality, 66, 132

\begin{tcolorbox}
理性(rationality)，66，132
\end{tcolorbox}

sequential, 91

\begin{tcolorbox}
顺序的，91
\end{tcolorbox}

ray, 192

\begin{tcolorbox}
射线，192
\end{tcolorbox}

ray termination, 255

\begin{tcolorbox}
射线终止，255
\end{tcolorbox}

realization equivalence, 284

\begin{tcolorbox}
实现等价性，284
\end{tcolorbox}

recursive, 4, 10

\begin{tcolorbox}
递归的，4，10
\end{tcolorbox}

reduced strategic form, 87, 275

\begin{tcolorbox}
精简策略型，87，275
\end{tcolorbox}

reduced strategy, 78, 87, 274

\begin{tcolorbox}
精简策略，78，87，274
\end{tcolorbox}

reflexive, 8, 11

\begin{tcolorbox}
自反的，8，11
\end{tcolorbox}

regret, 109, 122, 349

\begin{tcolorbox}
后悔，109，122，349
\end{tcolorbox}

rental car insurance, 109

\begin{tcolorbox}
租车保险，109
\end{tcolorbox}

Right player, 24

\begin{tcolorbox}
右玩家，24
\end{tcolorbox}

Rock-Paper-Scissors, 72, 74

\begin{tcolorbox}
石头剪刀布，72，74
\end{tcolorbox}

Rook-move game, 17

\begin{tcolorbox}
战车移动游戏，17
\end{tcolorbox}

R-position, 25

\begin{tcolorbox}
右胜位置，25
\end{tcolorbox}

scalar product, 142, 144

\begin{tcolorbox}
标量积，142，144
\end{tcolorbox}

scale-equivalent, 115, 116, 118

\begin{tcolorbox}
尺度等价的，115，116，118
\end{tcolorbox}

Schelling, Thomas, 74

\begin{tcolorbox}
谢林，托马斯，74
\end{tcolorbox}

seaview hotel, 188, 204

\begin{tcolorbox}
海景酒店，188，204
\end{tcolorbox}

security strategy, 212

\begin{tcolorbox}
安全策略，212
\end{tcolorbox}

Seoul (Korea), 44, 49

\begin{tcolorbox}
首尔 (韩国)，44，49
\end{tcolorbox}

separating hyperplane, 193, 220

\begin{tcolorbox}
分离超平面，193，220
\end{tcolorbox}

sequence form, 295

\begin{tcolorbox}
序列形式，295
\end{tcolorbox}

set inclusion, 8

\begin{tcolorbox}
集合包含，8
\end{tcolorbox}

sg-value, 16

\begin{tcolorbox}
sg 值，16
\end{tcolorbox}

simple lottery, 104, 117

\begin{tcolorbox}
简单彩票，104，117
\end{tcolorbox}

simplex, 175, 339

\begin{tcolorbox}
单纯形，175，339
\end{tcolorbox}

simplex algorithm, 255, 260

\begin{tcolorbox}
单纯形算法，255，260
\end{tcolorbox}

simplicial subdivision, 176

\begin{tcolorbox}
单纯剖分，176
\end{tcolorbox}

slack variables, 251

\begin{tcolorbox}
松弛变量，251
\end{tcolorbox}

Social Choice, 110, 133

\begin{tcolorbox}
社会选择，110，133
\end{tcolorbox}

social optimum, 41, 42, 49

\begin{tcolorbox}
社会最优，41，42，49
\end{tcolorbox}

SPE, see subgame-perfect equilibrium

\begin{tcolorbox}
子博弈完美均衡(SPE)，见子博弈完美均衡
\end{tcolorbox}

Sperner condition,180,184

\begin{tcolorbox}
斯佩纳条件,180,184
\end{tcolorbox}

Sperner labeling, 180

\begin{tcolorbox}
斯佩纳标号, 180
\end{tcolorbox}

Sprague-Grundy value, 16, 31

\begin{tcolorbox}
斯普拉格 - 格鲁迪值, 16, 31
\end{tcolorbox}

Stackelberg game, 96, 98

\begin{tcolorbox}
斯塔克尔伯格博弈, 96, 98
\end{tcolorbox}

Stag Hunt, 73

\begin{tcolorbox}
猎鹿博弈, 73
\end{tcolorbox}

staggered payoffs, 56, 74

\begin{tcolorbox}
交错收益, 56, 74
\end{tcolorbox}

stationary distribution, 344

\begin{tcolorbox}
平稳分布, 344
\end{tcolorbox}

stationary strategy, 325

\begin{tcolorbox}
平稳策略, 325
\end{tcolorbox}

strategic equivalence, 239

\begin{tcolorbox}
策略等价, 239
\end{tcolorbox}

strategic form, 74, 273

\begin{tcolorbox}
策略形式, 74, 273
\end{tcolorbox}

strict order, 8, 32

\begin{tcolorbox}
严格序, 8, 32
\end{tcolorbox}

strict preference, 106, 112

\begin{tcolorbox}
严格偏好, 106, 112
\end{tcolorbox}

strictly concave, 131

\begin{tcolorbox}
严格凹, 131
\end{tcolorbox}

strictly dominated strategy, 62

\begin{tcolorbox}
严格被占优策略, 62
\end{tcolorbox}

strong strategic equivalence, 241

\begin{tcolorbox}
强策略等价, 241
\end{tcolorbox}

sub-facet, 176, 178, 182

\begin{tcolorbox}
子极面, 176, 178, 182
\end{tcolorbox}

sub-segment, 173

\begin{tcolorbox}
子线段，173
\end{tcolorbox}

sub-simplex, 176

\begin{tcolorbox}
子单纯形，176
\end{tcolorbox}

subgame, 91, 291

\begin{tcolorbox}
子博弈，91，291
\end{tcolorbox}

subgame-perfect equilibrium, 91, 292

\begin{tcolorbox}
子博弈完美均衡，91，292
\end{tcolorbox}

subjective probability, 111

\begin{tcolorbox}
主观概率，111
\end{tcolorbox}

subtree, 91, 291

\begin{tcolorbox}
子树，91，291
\end{tcolorbox}

symmetric

\begin{tcolorbox}
对称的
\end{tcolorbox}

bimatrix game, 142

\begin{tcolorbox}
双矩阵博弈，142
\end{tcolorbox}

equilibrium, 61, 73

\begin{tcolorbox}
均衡，61，73
\end{tcolorbox}

game, 55, 60, 66, 72

\begin{tcolorbox}
博弈，55，60，66，72
\end{tcolorbox}

symmetry axiom, 303

\begin{tcolorbox}
对称性公理，303
\end{tcolorbox}

temperature scale, 104, 107, 116

\begin{tcolorbox}
温度刻度，104，107，116
\end{tcolorbox}

terminal node, 80

\begin{tcolorbox}
终端节点，80
\end{tcolorbox}

threat point, 302

\begin{tcolorbox}
威胁点，302
\end{tcolorbox}

tight inequality, 245

\begin{tcolorbox}
紧不等式，245
\end{tcolorbox}

top-down induction, 8

\begin{tcolorbox}
自上而下归纳，8
\end{tcolorbox}

total order, 8

\begin{tcolorbox}
全序(total order), 8
\end{tcolorbox}

total relation, 112

\begin{tcolorbox}
全关系(total relation), 112
\end{tcolorbox}

tradeoff, 114

\begin{tcolorbox}
权衡(tradeoff), 114
\end{tcolorbox}

transformation

\begin{tcolorbox}
变换(transformation)
\end{tcolorbox}

positive-affine, 115

\begin{tcolorbox}
正仿射(positive-affine), 115
\end{tcolorbox}

projective, 249

\begin{tcolorbox}
射影(projective), 249
\end{tcolorbox}

strictly increasing, 115

\begin{tcolorbox}
严格递增(strictly increasing), 115
\end{tcolorbox}

transitive, 8, 11

\begin{tcolorbox}
传递的(transitive), 8, 11
\end{tcolorbox}

transposed matrix, 141

\begin{tcolorbox}
转置矩阵(transposed matrix), 141
\end{tcolorbox}

triangulation, 176

\begin{tcolorbox}
三角剖分(triangulation), 176
\end{tcolorbox}

Trust Dilemma, 61

\begin{tcolorbox}
信任困境(Trust Dilemma), 61
\end{tcolorbox}

Ultimatum game, 314

\begin{tcolorbox}
最后通牒博弈(Ultimatum game), 314
\end{tcolorbox}

unit cube, 196

\begin{tcolorbox}
单位立方体(unit cube), 196
\end{tcolorbox}

unit of utility scale,104,107,116

\begin{tcolorbox}
效用刻度单位(unit of utility scale),104,107,116
\end{tcolorbox}

unit simplex, 151, 169-171, 175, 190

\begin{tcolorbox}
单位单纯形(unit simplex), 151, 169 - 171, 175, 190
\end{tcolorbox}

unit vector, 150, 169, 171

\begin{tcolorbox}
单位向量(unit vector), 150, 169, 171
\end{tcolorbox}

upper envelope, 158, 225, 241, 246

\begin{tcolorbox}
上包络, 158, 225, 241, 246
\end{tcolorbox}

utility function, 103-133

\begin{tcolorbox}
效用函数, 103 - 133
\end{tcolorbox}

cardinal, 104, 107, 115

\begin{tcolorbox}
基数(效用函数), 104, 107, 115
\end{tcolorbox}

ordinal, 106, 115

\begin{tcolorbox}
序数(效用函数), 106, 115
\end{tcolorbox}

utility scale, 104, 116, 137, 240

\begin{tcolorbox}
效用尺度, 104, 116, 137, 240
\end{tcolorbox}

vertex, 176, 204, 245

\begin{tcolorbox}
顶点, 176, 204, 245
\end{tcolorbox}

walk, 80

\begin{tcolorbox}
行走, 80
\end{tcolorbox}

Wardrop equilibrium, 50

\begin{tcolorbox}
用户均衡, 50
\end{tcolorbox}

weakly dominated strategy, 62, 63, 64, 75,

\begin{tcolorbox}
弱被占优策略, 62, 63, 64, 75,
\end{tcolorbox}

219

elimination may lose equilibria, 63 ,

\begin{tcolorbox}
消除(弱被占优策略)可能导致均衡丢失, 63 ,
\end{tcolorbox}

220

Wilson, Robert, 259

\begin{tcolorbox}
罗伯特·威尔逊, 259
\end{tcolorbox}

winning move, 3, 4

\begin{tcolorbox}
必胜步, 3, 4
\end{tcolorbox}

winning position, 3, 25

\begin{tcolorbox}
必胜局势, 3, 25
\end{tcolorbox}

zero-sum game, 207

\begin{tcolorbox}
零和博弈, 207
\end{tcolorbox}

\end{document}